\PassOptionsToPackage{hyperfootnotes=false}{hyperref}
\documentclass[oneside]{amsart}

\usepackage[all]{xy}
\usepackage{tikz-cd}
\usepackage[T1]{fontenc}
\usepackage[spanish,es-noquoting,es-noshorthands]{babel}
\usepackage{xstring}
\usepackage{xparse}
\usepackage{xcolor}
\definecolor{brightmaroon}{rgb}{0.76, 0.13, 0.28}
\usepackage[linktocpage=true,colorlinks=true,hyperindex,hypertexnames=false,citecolor=blue,linkcolor=brightmaroon]{hyperref}
\usepackage[left=1.25in,right=1.25in,top=0.75in,bottom=0.75in]{geometry}
%\usepackage[charter,greekfamily=didot]{mathdesign}
%\usepackage{Baskervaldx}
\usepackage{amssymb}
\usepackage{mathrsfs}
\usepackage{mathpazo}
\linespread{1.05}

\usepackage[nobottomtitles]{titlesec}
\usepackage{marginnote}
\usepackage[shortlabels]{enumitem}
\usepackage{longtable}
\usepackage{multicol}
\usepackage{aurical}
\usepackage{microtype}

\newtheoremstyle{ega-env-style}%
  {}{}{\rmfamily}{}{\bfseries}{.}{ }{\thmnote{(#3)}}%

\newtheoremstyle{ega-thm-env-style}%
  {}{}{\itshape}{}{\bfseries}{. --- }{ }{\thmname{#1}\thmnote{ (#3)}}%

\newtheoremstyle{ega-defn-env-style}%
  {}{}{\rmfamily}{}{\bfseries}{. --- }{ }{\thmname{#1}\thmnote{ (#3)}}%

\theoremstyle{ega-env-style}
\newtheorem*{env}{---}

\theoremstyle{ega-thm-env-style}
\newtheorem*{theorem}{Teorema}
\newtheorem*{proposition}{Proposición}
\newtheorem*{lemma}{Lema}
\newtheorem*{corollary}{Corolario}

\theoremstyle{ega-defn-env-style}
\newtheorem*{definition}{Definición}
\newtheorem*{example}{Ejemplo}
\newtheorem*{examples}{Ejemplos}
\newtheorem*{remark}{Observación}
\newtheorem*{remarks}{Observaciones}
\newtheorem*{scholium}{Escolio}
\newtheorem*{notation}{Notación}

% indent subsections, see https://tex.stackexchange.com/questions/177290/.
% also make section titles bigger.
% also add § to \thesection, https://tex.stackexchange.com/questions/119667/ and https://tex.stackexchange.com/questions/308737/.
\makeatletter
\def\l@subsection{\@tocline{2}{0pt}{2.5pc}{2.2pc}{}}
\def\section{\@startsection{section}{1}%
  \z@{.7\linespacing\@plus\linespacing}{.5\linespacing}%
  {\normalfont\bfseries\Large\scshape\centering}}
\renewcommand{\@seccntformat}[1]{%
  \ifnum\pdfstrcmp{#1}{section}=0\textsection\fi%
  \csname the#1\endcsname.~}
\makeatother

%\allowdisplaybreaks[1]
%\binoppenalty=9999
%\relpenalty=9999

% for Chapter 0, Chapter I, etc.
% credit for ZeroRoman https://tex.stackexchange.com/questions/211414/
% added into scripts/make_book.py
%\newcommand{\ZeroRoman}[1]{\ifcase\value{#1}\relax 0\else\Roman{#1}\fi}
%\renewcommand{\thechapter}{\ZeroRoman{chapter}}

\def\mathcal{\mathscr}
\def\sh{\mathcal}                   % sheaf font
\def\bb{\mathbf}                    % bold font
\def\cat{\mathtt}                   % category font
\def\leq{\leqslant}                 % <=
\def\geq{\geqslant}                 % >=
\def\setmin{-}                      % set minus
\def\rad{\mathfrak{r}}              % radical
\def\nilrad{\mathfrak{N}}           % nilradical
\def\emp{\varnothing}               % empty set
\def\vphi{\phi}                     % for switching \phi and \varphi, change if needed
\def\HH{\mathrm{H}}                 % cohomology H
\def\CHH{\check{\HH}}               % Čech cohomology H
\def\RR{\mathrm{R}}                 % right derived R
\def\LL{\mathrm{L}}                 % left derived L
\def\dual#1{{#1}^\vee}              % dual
\def\kres{k}                        % residue field k
\def\C{\cat{C}}                     % category C
\def\op{^\cat{op}}                  % opposite category
\def\Set{\cat{Set}}                 % category of sets
\def\CHom{\cat{Hom}}                % functor category
\def\supertilde{{\,\widetilde{\,}\,}}   % use \supertilde instead of ^\sim
\def\GL{\bb{GL}}
\def\bbGamma{\boldsymbol\Gamma}
\def\red{\mathrm{red}}
\def\rg{\operatorname{rg}}
\def\gr{\operatorname{gr}}
\def\Gr{\operatorname{Gr}}
\def\Hom{\operatorname{Hom}}
\def\Proj{\operatorname{Proj}}
\def\Tor{\operatorname{Tor}}
\def\Ext{\operatorname{Ext}}
\def\Supp{\operatorname{Supp}}
\def\Ker{\operatorname{Ker}}
\def\Im{\operatorname{Im}}
\def\Coker{\operatorname{Coker}}
\def\Spec{\operatorname{Spec}}
\def\Spm{\operatorname{Spm}}
\def\Spf{\operatorname{Spf}}
\def\grad{\operatorname{grad}}
\def\dimc{\operatorname{dimc}}
\def\codim{\operatorname{codim}}
\def\id{\operatorname{id}}
\def\Der{\operatorname{Der}}
\def\Diff{\operatorname{Diff}}
\def\Hyp{\operatorname{Hyp}}
\def\Tr{\operatorname{Tr}}

\renewcommand{\to}{\mathchoice{\longrightarrow}{\rightarrow}{\rightarrow}{\rightarrow}}
\newcommand{\from}{\mathchoice{\longleftarrow}{\leftarrow}{\leftarrow}{\leftarrow}}
\let\mapstoo\mapsto
\renewcommand{\mapsto}{\mathchoice{\longmapsto}{\mapstoo}{\mapstoo}{\mapstoo}}
\def\isoto{\simeq} % isomorphism
\def\simto{\xrightarrow{\sim}}

\renewcommand{\hat}[1]{\widehat{#1}}
\renewcommand{\tilde}[1]{\widetilde{#1}}


\def\shHom{\sh{H}\textup{\kern-2.2pt{\Fontauri\slshape om}}}   % sheaf Hom
\def\shProj{\sh{P}\textup{\kern-2.2pt{\Fontauri\slshape roj}}} % sheaf Proj
\def\shExt{\sh{E}\textup{\kern-2.2pt{\Fontauri\slshape xt}}}   % sheaf Ext
\def\shTor{\sh{T}\textup{\kern-2.2pt{\Fontauri\slshape or}}}   % sheaf Tor
\def\shGr{\sh{G}\textup{\kern-2.2pt{\Fontauri\slshape r}}}     % sheaf Gr
\def\shDer{\sh{D}\,\textup{\kern-2.2pt{\Fontauri\slshape er}}} % sheaf Der
\def\shDiff{\sh{D}\,\textup{\kern-2.2pt{\Fontauri\slshape if{}f}}\,} % sheaf Diff
\def\shHomcont{\sh{H}\textup{\kern-2.2pt{\Fontauri\slshape om.\,cont}}}   % sheaf Hom.cont
\def\shAut{\sh{A}\textup{\kern-2.2pt{\Fontauri\slshape ut}}}   % sheaf Aut

% if unsure of a translation
%\newcommand{\unsure}[2][]{\hl{#2}\marginpar{#1}}
%\newcommand{\completelyunsure}{\unsure{[\ldots]}}
\def\unsure#1{#1 {\color{red}(?)}}
\def\completelyunsure{{\color{red}(???)}}

% use to mark where original page starts
\newcommand{\oldpage}[2][]{{\marginnote{\normalfont{\textbf{#1}~|~#2}}}\ignorespaces}
\def\sectionbreak{\begin{center}***\end{center}}

% Give every source label its own zero-size PDF destination.  The inherited
% unnumbered theorem environments otherwise leave many labels pointing to the
% preceding subsection or display anchor rather than to their exact statement.
\let\EGAOriginalLabel\label
\renewcommand{\label}[1]{\phantomsection\EGAOriginalLabel{#1}}

% for referencing environments.
% use as \sref{chapter-number.x.y.z}, with optional args
% for volume and indices, e.g. \sref[volume]{chapter-number.x.y.z}[i].
\makeatletter
\newcommand{\EGAUnavailableHyperref}[2]{%
  \begingroup
  \let\@latex@warning\@gobble
  \let\@latex@warning@no@line\@gobble
  \let\G@refundefinedtrue\relax
  \hyperref[#1]{#2}%
  \endgroup
}
\NewDocumentCommand{\sref}{o m o}{%
  \@ifundefined{r@#2}%
    {\IfNoValueTF{#1}%
      {\IfNoValueTF{#3}%
        {\EGAUnavailableHyperref{#2}{\normalfont{(\StrBehind{#2}{.})}}}%
        {\EGAUnavailableHyperref{#2}{\normalfont{(\StrBehind{#2}{.},~{#3})}}}}%
      {\IfNoValueTF{#3}%
        {\EGAUnavailableHyperref{#2}{\normalfont{(\textbf{#1},~\StrBehind{#2}{.})}}}%
        {\EGAUnavailableHyperref{#2}{\normalfont{(\textbf{#1},~\StrBehind{#2}{.},~{#3})}}}}}%
    {\IfNoValueTF{#1}%
      {\IfNoValueTF{#3}%
        {\hyperref[#2]{\normalfont{(\StrBehind{#2}{.})}}}%
        {\hyperref[#2]{\normalfont{(\StrBehind{#2}{.},~{#3})}}}}%
      {\IfNoValueTF{#3}%
        {\hyperref[#2]{\normalfont{(\textbf{#1},~\StrBehind{#2}{.})}}}%
        {\hyperref[#2]{\normalfont{(\textbf{#1},~\StrBehind{#2}{.},~{#3})}}}}}%
}
\makeatother

% for marking changes following errata
% use as \erratum[volume]{correction} to say where the erratum is given
\newcommand{\erratum}[2][]{{#2}\marginpar{\footnotesize\textbf{Err}\textsubscript{#1}}}

% for referencing equations, use as \eref{eq:chapter-number.x.y.z}.
\newcommand{\eref}[1]{\hyperref[#1]{\normalfont{(\StrBehind{#1}{.})}}}

% Generated stable semantic destinations for the EGA IV reader.
% Visible text and numbering are unchanged.
\makeatletter
\newcommand{\sgaStableTarget}[1]{%
  \Hy@raisedlink{\hyper@anchorstart{#1}\hyper@anchorend}}
\expandafter\def\csname sga@target@IV.1.1.1\endcsname{ega4.c04.def.1-1-1}
\expandafter\def\csname sga@target@IV.1.1.2\endcsname{ega4.c04.prop.1-1-2}
\expandafter\def\csname sga@target@IV.1.1.3\endcsname{ega4.c04.prop.1-1-3}
\expandafter\def\csname sga@target@IV.1.1.4\endcsname{ega4.c04.cor.1-1-4}
\expandafter\def\csname sga@target@IV.1.1.5\endcsname{ega4.c04.prop.1-1-5}
\expandafter\def\csname sga@target@IV.1.1.6\endcsname{ega4.c04.prop.1-1-6}
\expandafter\def\csname sga@target@IV.1.10.1\endcsname{ega4.c04.thm.1-10-1}
\expandafter\def\csname sga@target@IV.1.10.2\endcsname{ega4.c04.cor.1-10-2}
\expandafter\def\csname sga@target@IV.1.10.3\endcsname{ega4.c04.prop.1-10-3}
\expandafter\def\csname sga@target@IV.1.10.4\endcsname{ega4.c04.cor.1-10-4}
\expandafter\def\csname sga@target@IV.1.2.1\endcsname{ega4.c04.def.1-2-1}
\expandafter\def\csname sga@target@IV.1.2.2\endcsname{ega4.c04.prop.1-2-2}
\expandafter\def\csname sga@target@IV.1.2.3\endcsname{ega4.c04.cor.1-2-3}
\expandafter\def\csname sga@target@IV.1.2.4\endcsname{ega4.c04.prop.1-2-4}
\expandafter\def\csname sga@target@IV.1.2.4.1\endcsname{ega4.c04.eq.1-2-4-1}
\expandafter\def\csname sga@target@IV.1.2.5\endcsname{ega4.c04.prop.1-2-5}
\expandafter\def\csname sga@target@IV.1.2.6\endcsname{ega4.c04.prop.1-2-6}
\expandafter\def\csname sga@target@IV.1.2.7\endcsname{ega4.c04.prop.1-2-7}
\expandafter\def\csname sga@target@IV.1.2.8\endcsname{ega4.c04.cor.1-2-8}
\expandafter\def\csname sga@target@IV.1.2.9\endcsname{ega4.c04.prop.1-2-9}
\expandafter\def\csname sga@target@IV.1.3.1\endcsname{ega4.c04.statement.1-3-1}
\expandafter\def\csname sga@target@IV.1.3.10\endcsname{ega4.c04.statement.1-3-10}
\expandafter\def\csname sga@target@IV.1.3.11\endcsname{ega4.c04.prop.1-3-11}
\expandafter\def\csname sga@target@IV.1.3.2\endcsname{ega4.c04.def.1-3-2}
\expandafter\def\csname sga@target@IV.1.3.3\endcsname{ega4.c04.prop.1-3-3}
\expandafter\def\csname sga@target@IV.1.3.4\endcsname{ega4.c04.prop.1-3-4}
\expandafter\def\csname sga@target@IV.1.3.5\endcsname{ega4.c04.cor.1-3-5}
\expandafter\def\csname sga@target@IV.1.3.6\endcsname{ega4.c04.prop.1-3-6}
\expandafter\def\csname sga@target@IV.1.3.7\endcsname{ega4.c04.prop.1-3-7}
\expandafter\def\csname sga@target@IV.1.3.8\endcsname{ega4.c04.statement.1-3-8}
\expandafter\def\csname sga@target@IV.1.3.9\endcsname{ega4.c04.prop.1-3-9}
\expandafter\def\csname sga@target@IV.1.4.1\endcsname{ega4.c04.statement.1-4-1}
\expandafter\def\csname sga@target@IV.1.4.2\endcsname{ega4.c04.def.1-4-2}
\expandafter\def\csname sga@target@IV.1.4.3\endcsname{ega4.c04.prop.1-4-3}
\expandafter\def\csname sga@target@IV.1.4.3.1\endcsname{ega4.c04.cor.1-4-3-1}
\expandafter\def\csname sga@target@IV.1.4.4\endcsname{ega4.c04.prop.1-4-4}
\expandafter\def\csname sga@target@IV.1.4.4.1\endcsname{ega4.c04.lem.1-4-4-1}
\expandafter\def\csname sga@target@IV.1.4.5\endcsname{ega4.c04.cor.1-4-5}
\expandafter\def\csname sga@target@IV.1.4.6\endcsname{ega4.c04.prop.1-4-6}
\expandafter\def\csname sga@target@IV.1.4.7\endcsname{ega4.c04.prop.1-4-7}
\expandafter\def\csname sga@target@IV.1.4.7.1\endcsname{ega4.c04.lem.1-4-7-1}
\expandafter\def\csname sga@target@IV.1.5.1\endcsname{ega4.c04.prop.1-5-1}
\expandafter\def\csname sga@target@IV.1.5.2\endcsname{ega4.c04.def.1-5-2}
\expandafter\def\csname sga@target@IV.1.5.3\endcsname{ega4.c04.prop.1-5-3}
\expandafter\def\csname sga@target@IV.1.5.4\endcsname{ega4.c04.prop.1-5-4}
\expandafter\def\csname sga@target@IV.1.5.5\endcsname{ega4.c04.cor.1-5-5}
\expandafter\def\csname sga@target@IV.1.5.6\endcsname{ega4.c04.cor.1-5-6}
\expandafter\def\csname sga@target@IV.1.5.7\endcsname{ega4.c04.prop.1-5-7}
\expandafter\def\csname sga@target@IV.1.6.1\endcsname{ega4.c04.def.1-6-1}
\expandafter\def\csname sga@target@IV.1.6.2\endcsname{ega4.c04.prop.1-6-2}
\expandafter\def\csname sga@target@IV.1.6.3\endcsname{ega4.c04.cor.1-6-3}
\expandafter\def\csname sga@target@IV.1.6.4\endcsname{ega4.c04.rem.1-6-4}
\expandafter\def\csname sga@target@IV.1.6.5\endcsname{ega4.c04.prop.1-6-5}
\expandafter\def\csname sga@target@IV.1.7.1\endcsname{ega4.c04.statement.1-7-1}
\expandafter\def\csname sga@target@IV.1.7.10\endcsname{ega4.c04.statement.1-7-10}
\expandafter\def\csname sga@target@IV.1.7.11\endcsname{ega4.c04.statement.1-7-11}
\expandafter\def\csname sga@target@IV.1.7.12\endcsname{ega4.c04.statement.1-7-12}
\expandafter\def\csname sga@target@IV.1.7.13\endcsname{ega4.c04.statement.1-7-13}
\expandafter\def\csname sga@target@IV.1.7.14\endcsname{ega4.c04.statement.1-7-14}
\expandafter\def\csname sga@target@IV.1.7.15\endcsname{ega4.c04.statement.1-7-15}
\expandafter\def\csname sga@target@IV.1.7.16\endcsname{ega4.c04.statement.1-7-16}
\expandafter\def\csname sga@target@IV.1.7.17\endcsname{ega4.c04.statement.1-7-17}
\expandafter\def\csname sga@target@IV.1.7.18\endcsname{ega4.c04.statement.1-7-18}
\expandafter\def\csname sga@target@IV.1.7.19\endcsname{ega4.c04.statement.1-7-19}
\expandafter\def\csname sga@target@IV.1.7.2\endcsname{ega4.c04.statement.1-7-2}
\expandafter\def\csname sga@target@IV.1.7.20\endcsname{ega4.c04.statement.1-7-20}
\expandafter\def\csname sga@target@IV.1.7.21\endcsname{ega4.c04.statement.1-7-21}
\expandafter\def\csname sga@target@IV.1.7.22\endcsname{ega4.c04.statement.1-7-22}
\expandafter\def\csname sga@target@IV.1.7.3\endcsname{ega4.c04.statement.1-7-3}
\expandafter\def\csname sga@target@IV.1.7.4\endcsname{ega4.c04.statement.1-7-4}
\expandafter\def\csname sga@target@IV.1.7.5\endcsname{ega4.c04.statement.1-7-5}
\expandafter\def\csname sga@target@IV.1.7.6\endcsname{ega4.c04.statement.1-7-6}
\expandafter\def\csname sga@target@IV.1.7.7\endcsname{ega4.c04.statement.1-7-7}
\expandafter\def\csname sga@target@IV.1.7.8\endcsname{ega4.c04.statement.1-7-8}
\expandafter\def\csname sga@target@IV.1.7.9\endcsname{ega4.c04.statement.1-7-9}
\expandafter\def\csname sga@target@IV.1.8.1\endcsname{ega4.c04.statement.1-8-1}
\expandafter\def\csname sga@target@IV.1.8.2\endcsname{ega4.c04.prop.1-8-2}
\expandafter\def\csname sga@target@IV.1.8.3\endcsname{ega4.c04.lem.1-8-3}
\expandafter\def\csname sga@target@IV.1.8.4\endcsname{ega4.c04.thm.1-8-4}
\expandafter\def\csname sga@target@IV.1.8.4.1\endcsname{ega4.c04.lem.1-8-4-1}
\expandafter\def\csname sga@target@IV.1.8.4.2\endcsname{ega4.c04.lem.1-8-4-2}
\expandafter\def\csname sga@target@IV.1.8.5\endcsname{ega4.c04.cor.1-8-5}
\expandafter\def\csname sga@target@IV.1.8.5.1\endcsname{ega4.c04.lem.1-8-5-1}
\expandafter\def\csname sga@target@IV.1.8.6\endcsname{ega4.c04.cor.1-8-6}
\expandafter\def\csname sga@target@IV.1.8.7\endcsname{ega4.c04.prop.1-8-7}
\expandafter\def\csname sga@target@IV.1.8.7.1\endcsname{ega4.c04.lem.1-8-7-1}
\expandafter\def\csname sga@target@IV.1.9.1\endcsname{ega4.c04.lem.1-9-1}
\expandafter\def\csname sga@target@IV.1.9.10\endcsname{ega4.c04.prop.1-9-10}
\expandafter\def\csname sga@target@IV.1.9.11\endcsname{ega4.c04.prop.1-9-11}
\expandafter\def\csname sga@target@IV.1.9.12\endcsname{ega4.c04.cor.1-9-12}
\expandafter\def\csname sga@target@IV.1.9.13\endcsname{ega4.c04.statement.1-9-13}
\expandafter\def\csname sga@target@IV.1.9.14\endcsname{ega4.c04.prop.1-9-14}
\expandafter\def\csname sga@target@IV.1.9.15\endcsname{ega4.c04.prop.1-9-15}
\expandafter\def\csname sga@target@IV.1.9.16\endcsname{ega4.c04.statement.1-9-16}
\expandafter\def\csname sga@target@IV.1.9.2\endcsname{ega4.c04.prop.1-9-2}
\expandafter\def\csname sga@target@IV.1.9.2.1\endcsname{ega4.c04.statement.1-9-2-1}
\expandafter\def\csname sga@target@IV.1.9.2.i\endcsname{ega4.c04.statement.1-9-2-i}
\expandafter\def\csname sga@target@IV.1.9.2.ii\endcsname{ega4.c04.statement.1-9-2-ii}
\expandafter\def\csname sga@target@IV.1.9.3\endcsname{ega4.c04.prop.1-9-3}
\expandafter\def\csname sga@target@IV.1.9.3.1\endcsname{ega4.c04.lem.1-9-3-1}
\expandafter\def\csname sga@target@IV.1.9.3.2\endcsname{ega4.c04.lem.1-9-3-2}
\expandafter\def\csname sga@target@IV.1.9.4\endcsname{ega4.c04.def.1-9-4}
\expandafter\def\csname sga@target@IV.1.9.5\endcsname{ega4.c04.prop.1-9-5}
\expandafter\def\csname sga@target@IV.1.9.6\endcsname{ega4.c04.statement.1-9-6}
\expandafter\def\csname sga@target@IV.1.9.7\endcsname{ega4.c04.thm.1-9-7}
\expandafter\def\csname sga@target@IV.1.9.8\endcsname{ega4.c04.cor.1-9-8}
\expandafter\def\csname sga@target@IV.1.9.9\endcsname{ega4.c04.cor.1-9-9}
\expandafter\def\csname sga@target@IV.10.1.1\endcsname{ega4.c04.statement.10-1-1}
\expandafter\def\csname sga@target@IV.10.1.2\endcsname{ega4.c04.prop.10-1-2}
\expandafter\def\csname sga@target@IV.10.1.3\endcsname{ega4.c04.def.10-1-3}
\expandafter\def\csname sga@target@IV.10.1.4\endcsname{ega4.c04.cor.10-1-4}
\expandafter\def\csname sga@target@IV.10.10.1\endcsname{ega4.c04.statement.10-10-1}
\expandafter\def\csname sga@target@IV.10.10.2\endcsname{ega4.c04.statement.10-10-2}
\expandafter\def\csname sga@target@IV.10.10.3\endcsname{ega4.c04.statement.10-10-3}
\expandafter\def\csname sga@target@IV.10.10.4\endcsname{ega4.c04.statement.10-10-4}
\expandafter\def\csname sga@target@IV.10.10.5\endcsname{ega4.c04.statement.10-10-5}
\expandafter\def\csname sga@target@IV.10.2.1\endcsname{ega4.c04.prop.10-2-1}
\expandafter\def\csname sga@target@IV.10.2.2\endcsname{ega4.c04.def.10-2-2}
\expandafter\def\csname sga@target@IV.10.2.3\endcsname{ega4.c04.cor.10-2-3}
\expandafter\def\csname sga@target@IV.10.2.4\endcsname{ega4.c04.cor.10-2-4}
\expandafter\def\csname sga@target@IV.10.2.5\endcsname{ega4.c04.cor.10-2-5}
\expandafter\def\csname sga@target@IV.10.2.6\endcsname{ega4.c04.cor.10-2-6}
\expandafter\def\csname sga@target@IV.10.2.7\endcsname{ega4.c04.prop.10-2-7}
\expandafter\def\csname sga@target@IV.10.2.8\endcsname{ega4.c04.item.10-2-8}
\expandafter\def\csname sga@target@IV.10.2.8.1\endcsname{ega4.c04.eq.10-2-8-1}
\expandafter\def\csname sga@target@IV.10.3.1\endcsname{ega4.c04.def.10-3-1}
\expandafter\def\csname sga@target@IV.10.3.2\endcsname{ega4.c04.prop.10-3-2}
\expandafter\def\csname sga@target@IV.10.3.3\endcsname{ega4.c04.prop.10-3-3}
\expandafter\def\csname sga@target@IV.10.4.1\endcsname{ega4.c04.def.10-4-1}
\expandafter\def\csname sga@target@IV.10.4.10\endcsname{ega4.c04.statement.10-4-10}
\expandafter\def\csname sga@target@IV.10.4.11\endcsname{ega4.c04.statement.10-4-11}
\expandafter\def\csname sga@target@IV.10.4.11.1\endcsname{ega4.c04.lem.10-4-11-1}
\expandafter\def\csname sga@target@IV.10.4.12\endcsname{ega4.c04.cor.10-4-12}
\expandafter\def\csname sga@target@IV.10.4.12.1\endcsname{ega4.c04.sec.10-4-12-1}
\expandafter\def\csname sga@target@IV.10.4.12.1.errata\endcsname{ega4.c04.eq.10-4-12-1-errata}
\expandafter\def\csname sga@target@IV.10.4.13\endcsname{ega4.c04.cor.10-4-13}
\expandafter\def\csname sga@target@IV.10.4.14\endcsname{ega4.c04.rem.10-4-14}
\expandafter\def\csname sga@target@IV.10.4.2\endcsname{ega4.c04.prop.10-4-2}
\expandafter\def\csname sga@target@IV.10.4.3\endcsname{ega4.c04.statement.10-4-3}
\expandafter\def\csname sga@target@IV.10.4.4\endcsname{ega4.c04.prop.10-4-4}
\expandafter\def\csname sga@target@IV.10.4.5\endcsname{ega4.c04.prop.10-4-5}
\expandafter\def\csname sga@target@IV.10.4.5.1\endcsname{ega4.c04.lem.10-4-5-1}
\expandafter\def\csname sga@target@IV.10.4.6\endcsname{ega4.c04.cor.10-4-6}
\expandafter\def\csname sga@target@IV.10.4.7\endcsname{ega4.c04.cor.10-4-7}
\expandafter\def\csname sga@target@IV.10.4.8\endcsname{ega4.c04.cor.10-4-8}
\expandafter\def\csname sga@target@IV.10.4.9\endcsname{ega4.c04.statement.10-4-9}
\expandafter\def\csname sga@target@IV.10.5.1\endcsname{ega4.c04.prop.10-5-1}
\expandafter\def\csname sga@target@IV.10.5.10\endcsname{ega4.c04.prop.10-5-10}
\expandafter\def\csname sga@target@IV.10.5.2\endcsname{ega4.c04.cor.10-5-2}
\expandafter\def\csname sga@target@IV.10.5.3\endcsname{ega4.c04.cor.10-5-3}
\expandafter\def\csname sga@target@IV.10.5.4\endcsname{ega4.c04.rem.10-5-4}
\expandafter\def\csname sga@target@IV.10.5.5\endcsname{ega4.c04.prop.10-5-5}
\expandafter\def\csname sga@target@IV.10.5.6\endcsname{ega4.c04.rem.10-5-6}
\expandafter\def\csname sga@target@IV.10.5.7\endcsname{ega4.c04.prop.10-5-7}
\expandafter\def\csname sga@target@IV.10.5.8\endcsname{ega4.c04.cor.10-5-8}
\expandafter\def\csname sga@target@IV.10.5.9\endcsname{ega4.c04.cor.10-5-9}
\expandafter\def\csname sga@target@IV.10.6.1\endcsname{ega4.c04.prop.10-6-1}
\expandafter\def\csname sga@target@IV.10.6.1.1\endcsname{ega4.c04.eq.10-6-1-1}
\expandafter\def\csname sga@target@IV.10.6.1.2\endcsname{ega4.c04.eq.10-6-1-2}
\expandafter\def\csname sga@target@IV.10.6.1.3\endcsname{ega4.c04.eq.10-6-1-3}
\expandafter\def\csname sga@target@IV.10.6.1.4\endcsname{ega4.c04.eq.10-6-1-4}
\expandafter\def\csname sga@target@IV.10.6.2\endcsname{ega4.c04.cor.10-6-2}
\expandafter\def\csname sga@target@IV.10.6.3\endcsname{ega4.c04.prop.10-6-3}
\expandafter\def\csname sga@target@IV.10.6.3.1\endcsname{ega4.c04.eq.10-6-3-1}
\expandafter\def\csname sga@target@IV.10.6.3.2\endcsname{ega4.c04.eq.10-6-3-2}
\expandafter\def\csname sga@target@IV.10.6.3.3\endcsname{ega4.c04.eq.10-6-3-3}
\expandafter\def\csname sga@target@IV.10.6.4\endcsname{ega4.c04.cor.10-6-4}
\expandafter\def\csname sga@target@IV.10.6.4.1\endcsname{ega4.c04.eq.10-6-4-1}
\expandafter\def\csname sga@target@IV.10.7.1\endcsname{ega4.c04.statement.10-7-1}
\expandafter\def\csname sga@target@IV.10.7.2\endcsname{ega4.c04.statement.10-7-2}
\expandafter\def\csname sga@target@IV.10.7.3\endcsname{ega4.c04.statement.10-7-3}
\expandafter\def\csname sga@target@IV.10.8.1\endcsname{ega4.c04.def.10-8-1}
\expandafter\def\csname sga@target@IV.10.8.1.1\endcsname{ega4.c04.eq.10-8-1-1}
\expandafter\def\csname sga@target@IV.10.8.1.2\endcsname{ega4.c04.eq.10-8-1-2}
\expandafter\def\csname sga@target@IV.10.8.2\endcsname{ega4.c04.statement.10-8-2}
\expandafter\def\csname sga@target@IV.10.8.3\endcsname{ega4.c04.lem.10-8-3}
\expandafter\def\csname sga@target@IV.10.8.3.1\endcsname{ega4.c04.eq.10-8-3-1}
\expandafter\def\csname sga@target@IV.10.8.4\endcsname{ega4.c04.cor.10-8-4}
\expandafter\def\csname sga@target@IV.10.8.5\endcsname{ega4.c04.prop.10-8-5}
\expandafter\def\csname sga@target@IV.10.8.6\endcsname{ega4.c04.cor.10-8-6}
\expandafter\def\csname sga@target@IV.10.8.7\endcsname{ega4.c04.prop.10-8-7}
\expandafter\def\csname sga@target@IV.10.8.7.1\endcsname{ega4.c04.eq.10-8-7-1}
\expandafter\def\csname sga@target@IV.10.8.8\endcsname{ega4.c04.cor.10-8-8}
\expandafter\def\csname sga@target@IV.10.9.1\endcsname{ega4.c04.statement.10-9-1}
\expandafter\def\csname sga@target@IV.10.9.2\endcsname{ega4.c04.statement.10-9-2}
\expandafter\def\csname sga@target@IV.10.9.3\endcsname{ega4.c04.statement.10-9-3}
\expandafter\def\csname sga@target@IV.10.9.4\endcsname{ega4.c04.prop.10-9-4}
\expandafter\def\csname sga@target@IV.10.9.5\endcsname{ega4.c04.statement.10-9-5}
\expandafter\def\csname sga@target@IV.10.9.6\endcsname{ega4.c04.prop.10-9-6}
\expandafter\def\csname sga@target@IV.11.1.1\endcsname{ega4.c04.thm.11-1-1}
\expandafter\def\csname sga@target@IV.11.1.1.1\endcsname{ega4.c04.sec.11-1-1-1}
\expandafter\def\csname sga@target@IV.11.1.1.2\endcsname{ega4.c04.lem.11-1-1-2}
\expandafter\def\csname sga@target@IV.11.1.2\endcsname{ega4.c04.cor.11-1-2}
\expandafter\def\csname sga@target@IV.11.1.3\endcsname{ega4.c04.cor.11-1-3}
\expandafter\def\csname sga@target@IV.11.1.4\endcsname{ega4.c04.cor.11-1-4}
\expandafter\def\csname sga@target@IV.11.1.5\endcsname{ega4.c04.cor.11-1-5}
\expandafter\def\csname sga@target@IV.11.1.6\endcsname{ega4.c04.cor.11-1-6}
\expandafter\def\csname sga@target@IV.11.10.1\endcsname{ega4.c04.prop.11-10-1}
\expandafter\def\csname sga@target@IV.11.10.1.1\endcsname{ega4.c04.eq.11-10-1-1}
\expandafter\def\csname sga@target@IV.11.10.1.2\endcsname{ega4.c04.eq.11-10-1-2}
\expandafter\def\csname sga@target@IV.11.10.10\endcsname{ega4.c04.prop.11-10-10}
\expandafter\def\csname sga@target@IV.11.10.2\endcsname{ega4.c04.def.11-10-2}
\expandafter\def\csname sga@target@IV.11.10.3\endcsname{ega4.c04.item.11-10-3}
\expandafter\def\csname sga@target@IV.11.10.4\endcsname{ega4.c04.prop.11-10-4}
\expandafter\def\csname sga@target@IV.11.10.5\endcsname{ega4.c04.thm.11-10-5}
\expandafter\def\csname sga@target@IV.11.10.6\endcsname{ega4.c04.prop.11-10-6}
\expandafter\def\csname sga@target@IV.11.10.7\endcsname{ega4.c04.cor.11-10-7}
\expandafter\def\csname sga@target@IV.11.10.8\endcsname{ega4.c04.def.11-10-8}
\expandafter\def\csname sga@target@IV.11.10.9\endcsname{ega4.c04.thm.11-10-9}
\expandafter\def\csname sga@target@IV.11.2.1\endcsname{ega4.c04.statement.11-2-1}
\expandafter\def\csname sga@target@IV.11.2.1.1\endcsname{ega4.c04.eq.11-2-1-1}
\expandafter\def\csname sga@target@IV.11.2.1.2\endcsname{ega4.c04.eq.11-2-1-2}
\expandafter\def\csname sga@target@IV.11.2.1.3\endcsname{ega4.c04.eq.11-2-1-3}
\expandafter\def\csname sga@target@IV.11.2.1.4\endcsname{ega4.c04.eq.11-2-1-4}
\expandafter\def\csname sga@target@IV.11.2.1.5\endcsname{ega4.c04.eq.11-2-1-5}
\expandafter\def\csname sga@target@IV.11.2.10\endcsname{ega4.c04.rem.11-2-10}
\expandafter\def\csname sga@target@IV.11.2.2\endcsname{ega4.c04.statement.11-2-2}
\expandafter\def\csname sga@target@IV.11.2.2.1\endcsname{ega4.c04.eq.11-2-2-1}
\expandafter\def\csname sga@target@IV.11.2.2.2\endcsname{ega4.c04.lem.11-2-2-2}
\expandafter\def\csname sga@target@IV.11.2.3\endcsname{ega4.c04.statement.11-2-3}
\expandafter\def\csname sga@target@IV.11.2.3.1\endcsname{ega4.c04.eq.11-2-3-1}
\expandafter\def\csname sga@target@IV.11.2.3.2\endcsname{ega4.c04.eq.11-2-3-2}
\expandafter\def\csname sga@target@IV.11.2.3.3\endcsname{ega4.c04.eq.11-2-3-3}
\expandafter\def\csname sga@target@IV.11.2.4\endcsname{ega4.c04.lem.11-2-4}
\expandafter\def\csname sga@target@IV.11.2.5\endcsname{ega4.c04.lem.11-2-5}
\expandafter\def\csname sga@target@IV.11.2.6\endcsname{ega4.c04.thm.11-2-6}
\expandafter\def\csname sga@target@IV.11.2.6.1\endcsname{ega4.c04.cor.11-2-6-1}
\expandafter\def\csname sga@target@IV.11.2.7\endcsname{ega4.c04.cor.11-2-7}
\expandafter\def\csname sga@target@IV.11.2.8\endcsname{ega4.c04.prop.11-2-8}
\expandafter\def\csname sga@target@IV.11.2.9\endcsname{ega4.c04.thm.11-2-9}
\expandafter\def\csname sga@target@IV.11.2.9.1\endcsname{ega4.c04.statement.11-2-9-1}
\expandafter\def\csname sga@target@IV.11.2.9.2\endcsname{ega4.c04.lem.11-2-9-2}
\expandafter\def\csname sga@target@IV.11.2.9.3\endcsname{ega4.c04.lem.11-2-9-3}
\expandafter\def\csname sga@target@IV.11.2.9.4\endcsname{ega4.c04.cor.11-2-9-4}
\expandafter\def\csname sga@target@IV.11.2.9.5\endcsname{ega4.c04.lem.11-2-9-5}
\expandafter\def\csname sga@target@IV.11.2.9.6\endcsname{ega4.c04.lem.11-2-9-6}
\expandafter\def\csname sga@target@IV.11.2.9.7\endcsname{ega4.c04.eq.11-2-9-7}
\expandafter\def\csname sga@target@IV.11.2.9.8\endcsname{ega4.c04.eq.11-2-9-8}
\expandafter\def\csname sga@target@IV.11.2.9.9\endcsname{ega4.c04.eq.11-2-9-9}
\expandafter\def\csname sga@target@IV.11.3\endcsname{ega4.c04.subsec.11-3.key-iv-11-3}
\expandafter\def\csname sga@target@IV.11.3.1\endcsname{ega4.c04.thm.11-3-1}
\expandafter\def\csname sga@target@IV.11.3.10\endcsname{ega4.c04.thm.11-3-10}
\expandafter\def\csname sga@target@IV.11.3.10.1\endcsname{ega4.c04.lem.11-3-10-1}
\expandafter\def\csname sga@target@IV.11.3.10.2\endcsname{ega4.c04.lem.11-3-10-2}
\expandafter\def\csname sga@target@IV.11.3.10.3\endcsname{ega4.c04.eq.11-3-10-3}
\expandafter\def\csname sga@target@IV.11.3.10.4\endcsname{ega4.c04.sec.11-3-10-4}
\expandafter\def\csname sga@target@IV.11.3.11\endcsname{ega4.c04.cor.11-3-11}
\expandafter\def\csname sga@target@IV.11.3.12\endcsname{ega4.c04.rem.11-3-12}
\expandafter\def\csname sga@target@IV.11.3.13\endcsname{ega4.c04.prop.11-3-13}
\expandafter\def\csname sga@target@IV.11.3.14\endcsname{ega4.c04.cor.11-3-14}
\expandafter\def\csname sga@target@IV.11.3.15\endcsname{ega4.c04.prop.11-3-15}
\expandafter\def\csname sga@target@IV.11.3.16\endcsname{ega4.c04.prop.11-3-16}
\expandafter\def\csname sga@target@IV.11.3.17\endcsname{ega4.c04.cor.11-3-17}
\expandafter\def\csname sga@target@IV.11.3.2\endcsname{ega4.c04.cor.11-3-2}
\expandafter\def\csname sga@target@IV.11.3.3\endcsname{ega4.c04.cor.11-3-3}
\expandafter\def\csname sga@target@IV.11.3.4\endcsname{ega4.c04.prop.11-3-4}
\expandafter\def\csname sga@target@IV.11.3.5\endcsname{ega4.c04.cor.11-3-5}
\expandafter\def\csname sga@target@IV.11.3.6\endcsname{ega4.c04.prop.11-3-6}
\expandafter\def\csname sga@target@IV.11.3.7\endcsname{ega4.c04.prop.11-3-7}
\expandafter\def\csname sga@target@IV.11.3.8\endcsname{ega4.c04.thm.11-3-8}
\expandafter\def\csname sga@target@IV.11.3.9\endcsname{ega4.c04.prop.11-3-9}
\expandafter\def\csname sga@target@IV.11.3.9.1\endcsname{ega4.c04.lem.11-3-9-1}
\expandafter\def\csname sga@target@IV.11.4.1\endcsname{ega4.c04.thm.11-4-1}
\expandafter\def\csname sga@target@IV.11.4.1.1\endcsname{ega4.c04.lem.11-4-1-1}
\expandafter\def\csname sga@target@IV.11.4.1.2\endcsname{ega4.c04.lem.11-4-1-2}
\expandafter\def\csname sga@target@IV.11.4.10\endcsname{ega4.c04.cor.11-4-10}
\expandafter\def\csname sga@target@IV.11.4.11\endcsname{ega4.c04.prop.11-4-11}
\expandafter\def\csname sga@target@IV.11.4.12\endcsname{ega4.c04.prop.11-4-12}
\expandafter\def\csname sga@target@IV.11.4.13\endcsname{ega4.c04.cor.11-4-13}
\expandafter\def\csname sga@target@IV.11.4.2\endcsname{ega4.c04.rem.11-4-2}
\expandafter\def\csname sga@target@IV.11.4.3\endcsname{ega4.c04.cor.11-4-3}
\expandafter\def\csname sga@target@IV.11.4.4\endcsname{ega4.c04.cor.11-4-4}
\expandafter\def\csname sga@target@IV.11.4.5\endcsname{ega4.c04.prop.11-4-5}
\expandafter\def\csname sga@target@IV.11.4.6\endcsname{ega4.c04.cor.11-4-6}
\expandafter\def\csname sga@target@IV.11.4.7\endcsname{ega4.c04.prop.11-4-7}
\expandafter\def\csname sga@target@IV.11.4.8\endcsname{ega4.c04.rem.11-4-8}
\expandafter\def\csname sga@target@IV.11.4.9\endcsname{ega4.c04.cor.11-4-9}
\expandafter\def\csname sga@target@IV.11.5.1\endcsname{ega4.c04.thm.11-5-1}
\expandafter\def\csname sga@target@IV.11.5.2\endcsname{ega4.c04.cor.11-5-2}
\expandafter\def\csname sga@target@IV.11.5.3\endcsname{ega4.c04.prop.11-5-3}
\expandafter\def\csname sga@target@IV.11.5.4\endcsname{ega4.c04.cor.11-5-4}
\expandafter\def\csname sga@target@IV.11.5.5\endcsname{ega4.c04.prop.11-5-5}
\expandafter\def\csname sga@target@IV.11.5.5.1\endcsname{ega4.c04.lem.11-5-5-1}
\expandafter\def\csname sga@target@IV.11.5.5.2\endcsname{ega4.c04.lem.11-5-5-2}
\expandafter\def\csname sga@target@IV.11.6.1\endcsname{ega4.c04.thm.11-6-1}
\expandafter\def\csname sga@target@IV.11.6.2\endcsname{ega4.c04.cor.11-6-2}
\expandafter\def\csname sga@target@IV.11.7.1\endcsname{ega4.c04.statement.11-7-1}
\expandafter\def\csname sga@target@IV.11.7.2\endcsname{ega4.c04.statement.11-7-2}
\expandafter\def\csname sga@target@IV.11.7.3\endcsname{ega4.c04.statement.11-7-3}
\expandafter\def\csname sga@target@IV.11.7.4\endcsname{ega4.c04.statement.11-7-4}
\expandafter\def\csname sga@target@IV.11.7.5\endcsname{ega4.c04.statement.11-7-5}
\expandafter\def\csname sga@target@IV.11.8.1\endcsname{ega4.c04.thm.11-8-1}
\expandafter\def\csname sga@target@IV.11.8.2\endcsname{ega4.c04.subsec.11-8-2}
\expandafter\def\csname sga@target@IV.11.9.1\endcsname{ega4.c04.statement.11-9-1}
\expandafter\def\csname sga@target@IV.11.9.1.1\endcsname{ega4.c04.statement.11-9-1-1}
\expandafter\def\csname sga@target@IV.11.9.10\endcsname{ega4.c04.thm.11-9-10}
\expandafter\def\csname sga@target@IV.11.9.11\endcsname{ega4.c04.rem.11-9-11}
\expandafter\def\csname sga@target@IV.11.9.12\endcsname{ega4.c04.lem.11-9-12}
\expandafter\def\csname sga@target@IV.11.9.13\endcsname{ega4.c04.prop.11-9-13}
\expandafter\def\csname sga@target@IV.11.9.14\endcsname{ega4.c04.statement.11-9-14}
\expandafter\def\csname sga@target@IV.11.9.15\endcsname{ega4.c04.prop.11-9-15}
\expandafter\def\csname sga@target@IV.11.9.16\endcsname{ega4.c04.thm.11-9-16}
\expandafter\def\csname sga@target@IV.11.9.17\endcsname{ega4.c04.prop.11-9-17}
\expandafter\def\csname sga@target@IV.11.9.17.1\endcsname{ega4.c04.lem.11-9-17-1}
\expandafter\def\csname sga@target@IV.11.9.18\endcsname{ega4.c04.item.11-9-18}
\expandafter\def\csname sga@target@IV.11.9.2\endcsname{ega4.c04.statement.11-9-2}
\expandafter\def\csname sga@target@IV.11.9.3\endcsname{ega4.c04.statement.11-9-3}
\expandafter\def\csname sga@target@IV.11.9.4\endcsname{ega4.c04.statement.11-9-4}
\expandafter\def\csname sga@target@IV.11.9.5\endcsname{ega4.c04.statement.11-9-5}
\expandafter\def\csname sga@target@IV.11.9.6\endcsname{ega4.c04.statement.11-9-6}
\expandafter\def\csname sga@target@IV.11.9.7\endcsname{ega4.c04.prop.11-9-7}
\expandafter\def\csname sga@target@IV.11.9.7.1\endcsname{ega4.c04.lem.11-9-7-1}
\expandafter\def\csname sga@target@IV.11.9.8\endcsname{ega4.c04.subsec.11-9-8}
\expandafter\def\csname sga@target@IV.11.9.9\endcsname{ega4.c04.cor.11-9-9}
\expandafter\def\csname sga@target@IV.12.0.1\endcsname{ega4.c04.statement.12-0-1}
\expandafter\def\csname sga@target@IV.12.0.2\endcsname{ega4.c04.statement.12-0-2}
\expandafter\def\csname sga@target@IV.12.1.1\endcsname{ega4.c04.thm.12-1-1}
\expandafter\def\csname sga@target@IV.12.1.1.1\endcsname{ega4.c04.statement.12-1-1-1}
\expandafter\def\csname sga@target@IV.12.1.1.2\endcsname{ega4.c04.statement.12-1-1-2}
\expandafter\def\csname sga@target@IV.12.1.1.3\endcsname{ega4.c04.statement.12-1-1-3}
\expandafter\def\csname sga@target@IV.12.1.1.4\endcsname{ega4.c04.statement.12-1-1-4}
\expandafter\def\csname sga@target@IV.12.1.1.5\endcsname{ega4.c04.prop.12-1-1-5}
\expandafter\def\csname sga@target@IV.12.1.2\endcsname{ega4.c04.item.12-1-2}
\expandafter\def\csname sga@target@IV.12.1.3\endcsname{ega4.c04.cor.12-1-3}
\expandafter\def\csname sga@target@IV.12.1.4\endcsname{ega4.c04.rem.12-1-4}
\expandafter\def\csname sga@target@IV.12.1.5\endcsname{ega4.c04.cor.12-1-5}
\expandafter\def\csname sga@target@IV.12.1.6\endcsname{ega4.c04.thm.12-1-6}
\expandafter\def\csname sga@target@IV.12.1.7\endcsname{ega4.c04.cor.12-1-7}
\expandafter\def\csname sga@target@IV.12.1.8\endcsname{ega4.c04.item.12-1-8}
\expandafter\def\csname sga@target@IV.12.2.1\endcsname{ega4.c04.thm.12-2-1}
\expandafter\def\csname sga@target@IV.12.2.1.1\endcsname{ega4.c04.statement.12-2-1-1}
\expandafter\def\csname sga@target@IV.12.2.1.2\endcsname{ega4.c04.statement.12-2-1-2}
\expandafter\def\csname sga@target@IV.12.2.1.3\endcsname{ega4.c04.statement.12-2-1-3}
\expandafter\def\csname sga@target@IV.12.2.1.4\endcsname{ega4.c04.lem.12-2-1-4}
\expandafter\def\csname sga@target@IV.12.2.2\endcsname{ega4.c04.cor.12-2-2}
\expandafter\def\csname sga@target@IV.12.2.3\endcsname{ega4.c04.item.12-2-3}
\expandafter\def\csname sga@target@IV.12.2.4\endcsname{ega4.c04.thm.12-2-4}
\expandafter\def\csname sga@target@IV.12.3.1\endcsname{ega4.c04.lem.12-3-1}
\expandafter\def\csname sga@target@IV.12.3.2\endcsname{ega4.c04.prop.12-3-2}
\expandafter\def\csname sga@target@IV.12.3.2.1\endcsname{ega4.c04.eq.12-3-2-1}
\expandafter\def\csname sga@target@IV.12.3.3\endcsname{ega4.c04.prop.12-3-3}
\expandafter\def\csname sga@target@IV.12.3.3.1\endcsname{ega4.c04.lem.12-3-3-1}
\expandafter\def\csname sga@target@IV.12.3.3.2\endcsname{ega4.c04.eq.12-3-3-2}
\expandafter\def\csname sga@target@IV.12.3.3.3\endcsname{ega4.c04.lem.12-3-3-3}
\expandafter\def\csname sga@target@IV.12.3.3.4\endcsname{ega4.c04.eq.12-3-3-4}
\expandafter\def\csname sga@target@IV.12.3.3.5\endcsname{ega4.c04.lem.12-3-3-5}
\expandafter\def\csname sga@target@IV.12.3.4\endcsname{ega4.c04.cor.12-3-4}
\expandafter\def\csname sga@target@IV.13.1.1\endcsname{ega4.c04.lem.13-1-1}
\expandafter\def\csname sga@target@IV.13.1.2\endcsname{ega4.c04.lem.13-1-2}
\expandafter\def\csname sga@target@IV.13.1.3\endcsname{ega4.c04.thm.13-1-3}
\expandafter\def\csname sga@target@IV.13.1.4\endcsname{ega4.c04.cor.13-1-4}
\expandafter\def\csname sga@target@IV.13.1.5\endcsname{ega4.c04.cor.13-1-5}
\expandafter\def\csname sga@target@IV.13.1.6\endcsname{ega4.c04.cor.13-1-6}
\expandafter\def\csname sga@target@IV.13.1.7\endcsname{ega4.c04.prop.13-1-7}
\expandafter\def\csname sga@target@IV.13.2.1\endcsname{ega4.c04.statement.13-2-1}
\expandafter\def\csname sga@target@IV.13.2.10\endcsname{ega4.c04.prop.13-2-10}
\expandafter\def\csname sga@target@IV.13.2.10.1\endcsname{ega4.c04.eq.13-2-10-1}
\expandafter\def\csname sga@target@IV.13.2.10.2\endcsname{ega4.c04.eq.13-2-10-2}
\expandafter\def\csname sga@target@IV.13.2.11\endcsname{ega4.c04.cor.13-2-11}
\expandafter\def\csname sga@target@IV.13.2.12\endcsname{ega4.c04.item.13-2-12}
\expandafter\def\csname sga@target@IV.13.2.2\endcsname{ega4.c04.def.13-2-2}
\expandafter\def\csname sga@target@IV.13.2.3\endcsname{ega4.c04.prop.13-2-3}
\expandafter\def\csname sga@target@IV.13.2.3.1\endcsname{ega4.c04.prop.13-2-3-1}
\expandafter\def\csname sga@target@IV.13.2.4\endcsname{ega4.c04.statement.13-2-4}
\expandafter\def\csname sga@target@IV.13.2.4.1\endcsname{ega4.c04.statement.13-2-4-1}
\expandafter\def\csname sga@target@IV.13.2.4.2\endcsname{ega4.c04.eq.13-2-4-2}
\expandafter\def\csname sga@target@IV.13.2.4.3\endcsname{ega4.c04.eq.13-2-4-3}
\expandafter\def\csname sga@target@IV.13.2.4.4\endcsname{ega4.c04.eq.13-2-4-4}
\expandafter\def\csname sga@target@IV.13.2.4.5\endcsname{ega4.c04.eq.13-2-4-5}
\expandafter\def\csname sga@target@IV.13.2.5\endcsname{ega4.c04.lem.13-2-5}
\expandafter\def\csname sga@target@IV.13.2.6\endcsname{ega4.c04.statement.13-2-6}
\expandafter\def\csname sga@target@IV.13.2.6.1\endcsname{ega4.c04.statement.13-2-6-1}
\expandafter\def\csname sga@target@IV.13.2.7\endcsname{ega4.c04.prop.13-2-7}
\expandafter\def\csname sga@target@IV.13.2.7.1\endcsname{ega4.c04.prop.13-2-7-1}
\expandafter\def\csname sga@target@IV.13.2.8\endcsname{ega4.c04.cor.13-2-8}
\expandafter\def\csname sga@target@IV.13.2.9\endcsname{ega4.c04.cor.13-2-9}
\expandafter\def\csname sga@target@IV.13.2.9.1\endcsname{ega4.c04.eq.13-2-9-1}
\expandafter\def\csname sga@target@IV.13.2.9.2\endcsname{ega4.c04.eq.13-2-9-2}
\expandafter\def\csname sga@target@IV.13.3.1\endcsname{ega4.c04.prop.13-3-1}
\expandafter\def\csname sga@target@IV.13.3.1.1\endcsname{ega4.c04.lem.13-3-1-1}
\expandafter\def\csname sga@target@IV.13.3.2\endcsname{ega4.c04.def.13-3-2}
\expandafter\def\csname sga@target@IV.13.3.3\endcsname{ega4.c04.prop.13-3-3}
\expandafter\def\csname sga@target@IV.13.3.4\endcsname{ega4.c04.cor.13-3-4}
\expandafter\def\csname sga@target@IV.13.3.4.1\endcsname{ega4.c04.cor.13-3-4-1}
\expandafter\def\csname sga@target@IV.13.3.5\endcsname{ega4.c04.cor.13-3-5}
\expandafter\def\csname sga@target@IV.13.3.6\endcsname{ega4.c04.prop.13-3-6}
\expandafter\def\csname sga@target@IV.13.3.6.1\endcsname{ega4.c04.prop.13-3-6-1}
\expandafter\def\csname sga@target@IV.13.3.6.2\endcsname{ega4.c04.eq.13-3-6-2}
\expandafter\def\csname sga@target@IV.13.3.7\endcsname{ega4.c04.prop.13-3-7}
\expandafter\def\csname sga@target@IV.13.3.7.1\endcsname{ega4.c04.prop.13-3-7-1}
\expandafter\def\csname sga@target@IV.13.3.8\endcsname{ega4.c04.prop.13-3-8}
\expandafter\def\csname sga@target@IV.13.3.9\endcsname{ega4.c04.statement.13-3-9}
\expandafter\def\csname sga@target@IV.14.1.1\endcsname{ega4.c04.statement.14-1-1}
\expandafter\def\csname sga@target@IV.14.1.2\endcsname{ega4.c04.prop.14-1-2}
\expandafter\def\csname sga@target@IV.14.1.3\endcsname{ega4.c04.item.14-1-3}
\expandafter\def\csname sga@target@IV.14.1.4\endcsname{ega4.c04.prop.14-1-4}
\expandafter\def\csname sga@target@IV.14.1.5\endcsname{ega4.c04.rem.14-1-5}
\expandafter\def\csname sga@target@IV.14.1.6\endcsname{ega4.c04.statement.14-1-6}
\expandafter\def\csname sga@target@IV.14.2.1\endcsname{ega4.c04.thm.14-2-1}
\expandafter\def\csname sga@target@IV.14.2.1.1\endcsname{ega4.c04.eq.14-2-1-1}
\expandafter\def\csname sga@target@IV.14.2.2\endcsname{ega4.c04.cor.14-2-2}
\expandafter\def\csname sga@target@IV.14.2.3\endcsname{ega4.c04.prop.14-2-3}
\expandafter\def\csname sga@target@IV.14.2.4\endcsname{ega4.c04.cor.14-2-4}
\expandafter\def\csname sga@target@IV.14.2.5\endcsname{ega4.c04.cor.14-2-5}
\expandafter\def\csname sga@target@IV.14.2.6\endcsname{ega4.c04.cor.14-2-6}
\expandafter\def\csname sga@target@IV.14.2.7\endcsname{ega4.c04.rem.14-2-7}
\expandafter\def\csname sga@target@IV.14.3\endcsname{ega4.c04.subsec.14-3.key-iv-14-3}
\expandafter\def\csname sga@target@IV.14.3.1\endcsname{ega4.c04.statement.14-3-1}
\expandafter\def\csname sga@target@IV.14.3.10\endcsname{ega4.c04.lem.14-3-10}
\expandafter\def\csname sga@target@IV.14.3.11\endcsname{ega4.c04.prop.14-3-11}
\expandafter\def\csname sga@target@IV.14.3.12\endcsname{ega4.c04.rem.14-3-12}
\expandafter\def\csname sga@target@IV.14.3.13\endcsname{ega4.c04.prop.14-3-13}
\expandafter\def\csname sga@target@IV.14.3.14\endcsname{ega4.c04.rem.14-3-14}
\expandafter\def\csname sga@target@IV.14.3.2\endcsname{ega4.c04.prop.14-3-2}
\expandafter\def\csname sga@target@IV.14.3.3\endcsname{ega4.c04.def.14-3-3}
\expandafter\def\csname sga@target@IV.14.3.3.1\endcsname{ega4.c04.subsec.14-3-3-1}
\expandafter\def\csname sga@target@IV.14.3.4\endcsname{ega4.c04.prop.14-3-4}
\expandafter\def\csname sga@target@IV.14.3.5\endcsname{ega4.c04.prop.14-3-5}
\expandafter\def\csname sga@target@IV.14.3.6\endcsname{ega4.c04.prop.14-3-6}
\expandafter\def\csname sga@target@IV.14.3.7\endcsname{ega4.c04.cor.14-3-7}
\expandafter\def\csname sga@target@IV.14.3.8\endcsname{ega4.c04.prop.14-3-8}
\expandafter\def\csname sga@target@IV.14.3.9\endcsname{ega4.c04.item.14-3-9}
\expandafter\def\csname sga@target@IV.14.4.1\endcsname{ega4.c04.thm.14-4-1}
\expandafter\def\csname sga@target@IV.14.4.1.1\endcsname{ega4.c04.lem.14-4-1-1}
\expandafter\def\csname sga@target@IV.14.4.1.2\endcsname{ega4.c04.lem.14-4-1-2}
\expandafter\def\csname sga@target@IV.14.4.1.3\endcsname{ega4.c04.lem.14-4-1-3}
\expandafter\def\csname sga@target@IV.14.4.10\endcsname{ega4.c04.item.14-4-10}
\expandafter\def\csname sga@target@IV.14.4.2\endcsname{ega4.c04.cor.14-4-2}
\expandafter\def\csname sga@target@IV.14.4.3\endcsname{ega4.c04.cor.14-4-3}
\expandafter\def\csname sga@target@IV.14.4.4\endcsname{ega4.c04.cor.14-4-4}
\expandafter\def\csname sga@target@IV.14.4.5\endcsname{ega4.c04.rem.14-4-5}
\expandafter\def\csname sga@target@IV.14.4.6\endcsname{ega4.c04.cor.14-4-6}
\expandafter\def\csname sga@target@IV.14.4.7\endcsname{ega4.c04.prop.14-4-7}
\expandafter\def\csname sga@target@IV.14.4.7.1\endcsname{ega4.c04.eq.14-4-7-1}
\expandafter\def\csname sga@target@IV.14.4.7.2\endcsname{ega4.c04.eq.14-4-7-2}
\expandafter\def\csname sga@target@IV.14.4.8\endcsname{ega4.c04.prop.14-4-8}
\expandafter\def\csname sga@target@IV.14.4.8.1\endcsname{ega4.c04.lem.14-4-8-1}
\expandafter\def\csname sga@target@IV.14.4.9\endcsname{ega4.c04.cor.14-4-9}
\expandafter\def\csname sga@target@IV.14.5\endcsname{ega4.c04.subsec.14-5.key-iv-14-5}
\expandafter\def\csname sga@target@IV.14.5.1\endcsname{ega4.c04.lem.14-5-1}
\expandafter\def\csname sga@target@IV.14.5.1.1\endcsname{ega4.c04.eq.14-5-1-1}
\expandafter\def\csname sga@target@IV.14.5.1.2\endcsname{ega4.c04.eq.14-5-1-2}
\expandafter\def\csname sga@target@IV.14.5.1.3\endcsname{ega4.c04.eq.14-5-1-3}
\expandafter\def\csname sga@target@IV.14.5.1.4\endcsname{ega4.c04.eq.14-5-1-4}
\expandafter\def\csname sga@target@IV.14.5.10\endcsname{ega4.c04.prop.14-5-10}
\expandafter\def\csname sga@target@IV.14.5.11\endcsname{ega4.c04.cor.14-5-11}
\expandafter\def\csname sga@target@IV.14.5.11.1\endcsname{ega4.c04.lem.14-5-11-1}
\expandafter\def\csname sga@target@IV.14.5.12\endcsname{ega4.c04.cor.14-5-12}
\expandafter\def\csname sga@target@IV.14.5.2\endcsname{ega4.c04.cor.14-5-2}
\expandafter\def\csname sga@target@IV.14.5.2.1\endcsname{ega4.c04.eq.14-5-2-1}
\expandafter\def\csname sga@target@IV.14.5.3\endcsname{ega4.c04.prop.14-5-3}
\expandafter\def\csname sga@target@IV.14.5.4\endcsname{ega4.c04.cor.14-5-4}
\expandafter\def\csname sga@target@IV.14.5.5\endcsname{ega4.c04.rem.14-5-5}
\expandafter\def\csname sga@target@IV.14.5.6\endcsname{ega4.c04.prop.14-5-6}
\expandafter\def\csname sga@target@IV.14.5.7\endcsname{ega4.c04.rem.14-5-7}
\expandafter\def\csname sga@target@IV.14.5.8\endcsname{ega4.c04.prop.14-5-8}
\expandafter\def\csname sga@target@IV.14.5.9\endcsname{ega4.c04.prop.14-5-9}
\expandafter\def\csname sga@target@IV.15.1.1\endcsname{ega4.c04.prop.15-1-1}
\expandafter\def\csname sga@target@IV.15.1.1-proof\endcsname{ega4.c04.statement.15-1-1-proof}
\expandafter\def\csname sga@target@IV.15.1.1.1\endcsname{ega4.c04.eq.15-1-1-1}
\expandafter\def\csname sga@target@IV.15.1.1.2\endcsname{ega4.c04.eq.15-1-1-2}
\expandafter\def\csname sga@target@IV.15.1.1.3\endcsname{ega4.c04.eq.15-1-1-3}
\expandafter\def\csname sga@target@IV.15.1.1.4\endcsname{ega4.c04.eq.15-1-1-4}
\expandafter\def\csname sga@target@IV.15.1.1.5\endcsname{ega4.c04.eq.15-1-1-5}
\expandafter\def\csname sga@target@IV.15.1.1.6\endcsname{ega4.c04.lem.15-1-1-6}
\expandafter\def\csname sga@target@IV.15.1.2\endcsname{ega4.c04.cor.15-1-2}
\expandafter\def\csname sga@target@IV.15.1.3\endcsname{ega4.c04.cor.15-1-3}
\expandafter\def\csname sga@target@IV.15.1.4\endcsname{ega4.c04.statement.15-1-4}
\expandafter\def\csname sga@target@IV.15.2.1\endcsname{ega4.c04.statement.15-2-1}
\expandafter\def\csname sga@target@IV.15.2.2\endcsname{ega4.c04.thm.15-2-2}
\expandafter\def\csname sga@target@IV.15.2.2-proof\endcsname{ega4.c04.item.15-2-2-proof}
\expandafter\def\csname sga@target@IV.15.2.2.1\endcsname{ega4.c04.lem.15-2-2-1}
\expandafter\def\csname sga@target@IV.15.2.3\endcsname{ega4.c04.cor.15-2-3}
\expandafter\def\csname sga@target@IV.15.2.4\endcsname{ega4.c04.item.15-2-4}
\expandafter\def\csname sga@target@IV.15.3.1\endcsname{ega4.c04.prop.15-3-1}
\expandafter\def\csname sga@target@IV.15.3.1-proof\endcsname{ega4.c04.item.15-3-1-proof}
\expandafter\def\csname sga@target@IV.15.3.2\endcsname{ega4.c04.prop.15-3-2}
\expandafter\def\csname sga@target@IV.15.3.2-proof\endcsname{ega4.c04.item.15-3-2-proof}
\expandafter\def\csname sga@target@IV.15.3.3\endcsname{ega4.c04.cor.15-3-3}
\expandafter\def\csname sga@target@IV.15.3.3-proof\endcsname{ega4.c04.item.15-3-3-proof}
\expandafter\def\csname sga@target@IV.15.3.4\endcsname{ega4.c04.item.15-3-4}
\expandafter\def\csname sga@target@IV.15.4.1\endcsname{ega4.c04.prop.15-4-1}
\expandafter\def\csname sga@target@IV.15.4.2\endcsname{ega4.c04.prop.15-4-2}
\expandafter\def\csname sga@target@IV.15.4.3\endcsname{ega4.c04.prop.15-4-3}
\expandafter\def\csname sga@target@IV.15.4.3-proof\endcsname{ega4.c04.item.15-4-3-proof}
\expandafter\def\csname sga@target@IV.15.5.1\endcsname{ega4.c04.prop.15-5-1}
\expandafter\def\csname sga@target@IV.15.5.10\endcsname{ega4.c04.item.15-5-10}
\expandafter\def\csname sga@target@IV.15.5.2\endcsname{ega4.c04.cor.15-5-2}
\expandafter\def\csname sga@target@IV.15.5.2.1\endcsname{ega4.c04.eq.15-5-2-1}
\expandafter\def\csname sga@target@IV.15.5.2.2\endcsname{ega4.c04.lem.15-5-2-2}
\expandafter\def\csname sga@target@IV.15.5.2.3\endcsname{ega4.c04.eq.15-5-2-3}
\expandafter\def\csname sga@target@IV.15.5.2.4\endcsname{ega4.c04.eq.15-5-2-4}
\expandafter\def\csname sga@target@IV.15.5.3\endcsname{ega4.c04.prop.15-5-3}
\expandafter\def\csname sga@target@IV.15.5.4\endcsname{ega4.c04.cor.15-5-4}
\expandafter\def\csname sga@target@IV.15.5.5\endcsname{ega4.c04.statement.15-5-5}
\expandafter\def\csname sga@target@IV.15.5.6\endcsname{ega4.c04.lem.15-5-6}
\expandafter\def\csname sga@target@IV.15.5.6.1\endcsname{ega4.c04.lem.15-5-6-1}
\expandafter\def\csname sga@target@IV.15.5.7\endcsname{ega4.c04.prop.15-5-7}
\expandafter\def\csname sga@target@IV.15.5.7.1\endcsname{ega4.c04.eq.15-5-7-1}
\expandafter\def\csname sga@target@IV.15.5.8\endcsname{ega4.c04.statement.15-5-8}
\expandafter\def\csname sga@target@IV.15.5.9\endcsname{ega4.c04.prop.15-5-9}
\expandafter\def\csname sga@target@IV.15.6.1\endcsname{ega4.c04.prop.15-6-1}
\expandafter\def\csname sga@target@IV.15.6.1.1\endcsname{ega4.c04.eq.15-6-1-1}
\expandafter\def\csname sga@target@IV.15.6.10\endcsname{ega4.c04.item.15-6-10}
\expandafter\def\csname sga@target@IV.15.6.2\endcsname{ega4.c04.cor.15-6-2}
\expandafter\def\csname sga@target@IV.15.6.3\endcsname{ega4.c04.prop.15-6-3}
\expandafter\def\csname sga@target@IV.15.6.4\endcsname{ega4.c04.prop.15-6-4}
\expandafter\def\csname sga@target@IV.15.6.5\endcsname{ega4.c04.cor.15-6-5}
\expandafter\def\csname sga@target@IV.15.6.6\endcsname{ega4.c04.prop.15-6-6}
\expandafter\def\csname sga@target@IV.15.6.7\endcsname{ega4.c04.cor.15-6-7}
\expandafter\def\csname sga@target@IV.15.6.8\endcsname{ega4.c04.prop.15-6-8}
\expandafter\def\csname sga@target@IV.15.6.9\endcsname{ega4.c04.cor.15-6-9}
\expandafter\def\csname sga@target@IV.15.7.1\endcsname{ega4.c04.def.15-7-1}
\expandafter\def\csname sga@target@IV.15.7.10\endcsname{ega4.c04.cor.15-7-10}
\expandafter\def\csname sga@target@IV.15.7.11\endcsname{ega4.c04.cor.15-7-11}
\expandafter\def\csname sga@target@IV.15.7.11.1\endcsname{ega4.c04.lem.15-7-11-1}
\expandafter\def\csname sga@target@IV.15.7.2\endcsname{ega4.c04.prop.15-7-2}
\expandafter\def\csname sga@target@IV.15.7.3\endcsname{ega4.c04.rem.15-7-3}
\expandafter\def\csname sga@target@IV.15.7.4\endcsname{ega4.c04.cor.15-7-4}
\expandafter\def\csname sga@target@IV.15.7.5\endcsname{ega4.c04.cor.15-7-5}
\expandafter\def\csname sga@target@IV.15.7.6\endcsname{ega4.c04.cor.15-7-6}
\expandafter\def\csname sga@target@IV.15.7.7\endcsname{ega4.c04.item.15-7-7}
\expandafter\def\csname sga@target@IV.15.7.8\endcsname{ega4.c04.statement.15-7-8}
\expandafter\def\csname sga@target@IV.15.7.9\endcsname{ega4.c04.prop.15-7-9}
\expandafter\def\csname sga@target@IV.16.1\endcsname{ega4.c04.subsec.16-1}
\expandafter\def\csname sga@target@IV.16.1.1\endcsname{ega4.c04.statement.16-1-1}
\expandafter\def\csname sga@target@IV.16.1.10\endcsname{ega4.c04.cor.16-1-10}
\expandafter\def\csname sga@target@IV.16.1.11\endcsname{ega4.c04.rem.16-1-11}
\expandafter\def\csname sga@target@IV.16.1.2\endcsname{ega4.c04.def.16-1-2}
\expandafter\def\csname sga@target@IV.16.1.2.1\endcsname{ega4.c04.def.16-1-2-1}
\expandafter\def\csname sga@target@IV.16.1.2.2\endcsname{ega4.c04.eq.16-1-2-2}
\expandafter\def\csname sga@target@IV.16.1.3\endcsname{ega4.c04.statement.16-1-3}
\expandafter\def\csname sga@target@IV.16.1.4\endcsname{ega4.c04.statement.16-1-4}
\expandafter\def\csname sga@target@IV.16.1.5\endcsname{ega4.c04.prop.16-1-5}
\expandafter\def\csname sga@target@IV.16.1.6\endcsname{ega4.c04.cor.16-1-6}
\expandafter\def\csname sga@target@IV.16.1.7\endcsname{ega4.c04.cor.16-1-7}
\expandafter\def\csname sga@target@IV.16.1.8\endcsname{ega4.c04.prop.16-1-8}
\expandafter\def\csname sga@target@IV.16.1.9\endcsname{ega4.c04.prop.16-1-9}
\expandafter\def\csname sga@target@IV.16.10\endcsname{ega4.c04.subsec.16-10}
\expandafter\def\csname sga@target@IV.16.10.1\endcsname{ega4.c04.def.16-10-1}
\expandafter\def\csname sga@target@IV.16.10.2\endcsname{ega4.c04.prop.16-10-2}
\expandafter\def\csname sga@target@IV.16.10.3\endcsname{ega4.c04.prop.16-10-3}
\expandafter\def\csname sga@target@IV.16.10.4\endcsname{ega4.c04.prop.16-10-4}
\expandafter\def\csname sga@target@IV.16.10.5\endcsname{ega4.c04.prop.16-10-5}
\expandafter\def\csname sga@target@IV.16.10.6\endcsname{ega4.c04.rem.16-10-6}
\expandafter\def\csname sga@target@IV.16.11\endcsname{ega4.c04.subsec.16-11}
\expandafter\def\csname sga@target@IV.16.11.1\endcsname{ega4.c04.statement.16-11-1}
\expandafter\def\csname sga@target@IV.16.11.1.1\endcsname{ega4.c04.statement.16-11-1-1}
\expandafter\def\csname sga@target@IV.16.11.1.2\endcsname{ega4.c04.statement.16-11-1-2}
\expandafter\def\csname sga@target@IV.16.11.1.3\endcsname{ega4.c04.eq.16-11-1-3}
\expandafter\def\csname sga@target@IV.16.11.1.4\endcsname{ega4.c04.eq.16-11-1-4}
\expandafter\def\csname sga@target@IV.16.11.1.5\endcsname{ega4.c04.eq.16-11-1-5}
\expandafter\def\csname sga@target@IV.16.11.1.6\endcsname{ega4.c04.eq.16-11-1-6}
\expandafter\def\csname sga@target@IV.16.11.1.7\endcsname{ega4.c04.eq.16-11-1-7}
\expandafter\def\csname sga@target@IV.16.11.2\endcsname{ega4.c04.thm.16-11-2}
\expandafter\def\csname sga@target@IV.16.11.2.1\endcsname{ega4.c04.eq.16-11-2-1}
\expandafter\def\csname sga@target@IV.16.11.2.2\endcsname{ega4.c04.eq.16-11-2-2}
\expandafter\def\csname sga@target@IV.16.11.2.3\endcsname{ega4.c04.eq.16-11-2-3}
\expandafter\def\csname sga@target@IV.16.11.3\endcsname{ega4.c04.statement.16-11-3}
\expandafter\def\csname sga@target@IV.16.12\endcsname{ega4.c04.subsec.16-12}
\expandafter\def\csname sga@target@IV.16.12.1\endcsname{ega4.c04.statement.16-12-1}
\expandafter\def\csname sga@target@IV.16.12.2\endcsname{ega4.c04.thm.16-12-2}
\expandafter\def\csname sga@target@IV.16.12.3\endcsname{ega4.c04.statement.16-12-3}
\expandafter\def\csname sga@target@IV.16.2\endcsname{ega4.c04.subsec.16-2}
\expandafter\def\csname sga@target@IV.16.2.1\endcsname{ega4.c04.statement.16-2-1}
\expandafter\def\csname sga@target@IV.16.2.1.1\endcsname{ega4.c04.statement.16-2-1-1}
\expandafter\def\csname sga@target@IV.16.2.1.2\endcsname{ega4.c04.eq.16-2-1-2}
\expandafter\def\csname sga@target@IV.16.2.1.3\endcsname{ega4.c04.eq.16-2-1-3}
\expandafter\def\csname sga@target@IV.16.2.1.4\endcsname{ega4.c04.eq.16-2-1-4}
\expandafter\def\csname sga@target@IV.16.2.2\endcsname{ega4.c04.prop.16-2-2}
\expandafter\def\csname sga@target@IV.16.2.3\endcsname{ega4.c04.cor.16-2-3}
\expandafter\def\csname sga@target@IV.16.2.3.1\endcsname{ega4.c04.eq.16-2-3-1}
\expandafter\def\csname sga@target@IV.16.2.3.2\endcsname{ega4.c04.eq.16-2-3-2}
\expandafter\def\csname sga@target@IV.16.2.4\endcsname{ega4.c04.cor.16-2-4}
\expandafter\def\csname sga@target@IV.16.2.5\endcsname{ega4.c04.statement.16-2-5}
\expandafter\def\csname sga@target@IV.16.2.6\endcsname{ega4.c04.statement.16-2-6}
\expandafter\def\csname sga@target@IV.16.2.6.1\endcsname{ega4.c04.statement.16-2-6-1}
\expandafter\def\csname sga@target@IV.16.2.6.2\endcsname{ega4.c04.eq.16-2-6-2}
\expandafter\def\csname sga@target@IV.16.2.7\endcsname{ega4.c04.prop.16-2-7}
\expandafter\def\csname sga@target@IV.16.2.7.1\endcsname{ega4.c04.prop.16-2-7-1}
\expandafter\def\csname sga@target@IV.16.3\endcsname{ega4.c04.subsec.16-3}
\expandafter\def\csname sga@target@IV.16.3.1\endcsname{ega4.c04.def.16-3-1}
\expandafter\def\csname sga@target@IV.16.3.1.1\endcsname{ega4.c04.eq.16-3-1-1}
\expandafter\def\csname sga@target@IV.16.3.2\endcsname{ega4.c04.statement.16-3-2}
\expandafter\def\csname sga@target@IV.16.3.3\endcsname{ega4.c04.statement.16-3-3}
\expandafter\def\csname sga@target@IV.16.3.3.1\endcsname{ega4.c04.statement.16-3-3-1}
\expandafter\def\csname sga@target@IV.16.3.4\endcsname{ega4.c04.prop.16-3-4}
\expandafter\def\csname sga@target@IV.16.3.5\endcsname{ega4.c04.statement.16-3-5}
\expandafter\def\csname sga@target@IV.16.3.6\endcsname{ega4.c04.def.16-3-6}
\expandafter\def\csname sga@target@IV.16.3.6.1\endcsname{ega4.c04.eq.16-3-6-1}
\expandafter\def\csname sga@target@IV.16.3.7\endcsname{ega4.c04.statement.16-3-7}
\expandafter\def\csname sga@target@IV.16.3.7.1\endcsname{ega4.c04.eq.16-3-7-1}
\expandafter\def\csname sga@target@IV.16.3.8\endcsname{ega4.c04.prop.16-3-8}
\expandafter\def\csname sga@target@IV.16.3.9\endcsname{ega4.c04.prop.16-3-9}
\expandafter\def\csname sga@target@IV.16.4\endcsname{ega4.c04.subsec.16-4}
\expandafter\def\csname sga@target@IV.16.4.1\endcsname{ega4.c04.statement.16-4-1}
\expandafter\def\csname sga@target@IV.16.4.1.1\endcsname{ega4.c04.statement.16-4-1-1}
\expandafter\def\csname sga@target@IV.16.4.1.2\endcsname{ega4.c04.statement.16-4-1-2}
\expandafter\def\csname sga@target@IV.16.4.1.3\endcsname{ega4.c04.eq.16-4-1-3}
\expandafter\def\csname sga@target@IV.16.4.1.4\endcsname{ega4.c04.eq.16-4-1-4}
\expandafter\def\csname sga@target@IV.16.4.1.5\endcsname{ega4.c04.eq.16-4-1-5}
\expandafter\def\csname sga@target@IV.16.4.10\endcsname{ega4.c04.cor.16-4-10}
\expandafter\def\csname sga@target@IV.16.4.11\endcsname{ega4.c04.prop.16-4-11}
\expandafter\def\csname sga@target@IV.16.4.11.1\endcsname{ega4.c04.prop.16-4-11-1}
\expandafter\def\csname sga@target@IV.16.4.11.2\endcsname{ega4.c04.prop.16-4-11-2}
\expandafter\def\csname sga@target@IV.16.4.11.3\endcsname{ega4.c04.eq.16-4-11-3}
\expandafter\def\csname sga@target@IV.16.4.11.4\endcsname{ega4.c04.eq.16-4-11-4}
\expandafter\def\csname sga@target@IV.16.4.12\endcsname{ega4.c04.cor.16-4-12}
\expandafter\def\csname sga@target@IV.16.4.13\endcsname{ega4.c04.cor.16-4-13}
\expandafter\def\csname sga@target@IV.16.4.14\endcsname{ega4.c04.prop.16-4-14}
\expandafter\def\csname sga@target@IV.16.4.14.1\endcsname{ega4.c04.prop.16-4-14-1}
\expandafter\def\csname sga@target@IV.16.4.15\endcsname{ega4.c04.cor.16-4-15}
\expandafter\def\csname sga@target@IV.16.4.15.1\endcsname{ega4.c04.cor.16-4-15-1}
\expandafter\def\csname sga@target@IV.16.4.15.2\endcsname{ega4.c04.eq.16-4-15-2}
\expandafter\def\csname sga@target@IV.16.4.16\endcsname{ega4.c04.cor.16-4-16}
\expandafter\def\csname sga@target@IV.16.4.16.1\endcsname{ega4.c04.cor.16-4-16-1}
\expandafter\def\csname sga@target@IV.16.4.16.2\endcsname{ega4.c04.eq.16-4-16-2}
\expandafter\def\csname sga@target@IV.16.4.17\endcsname{ega4.c04.cor.16-4-17}
\expandafter\def\csname sga@target@IV.16.4.18\endcsname{ega4.c04.statement.16-4-18}
\expandafter\def\csname sga@target@IV.16.4.18.1\endcsname{ega4.c04.statement.16-4-18-1}
\expandafter\def\csname sga@target@IV.16.4.18.2\endcsname{ega4.c04.statement.16-4-18-2}
\expandafter\def\csname sga@target@IV.16.4.18.3\endcsname{ega4.c04.eq.16-4-18-3}
\expandafter\def\csname sga@target@IV.16.4.19\endcsname{ega4.c04.cor.16-4-19}
\expandafter\def\csname sga@target@IV.16.4.19.1\endcsname{ega4.c04.cor.16-4-19-1}
\expandafter\def\csname sga@target@IV.16.4.2\endcsname{ega4.c04.statement.16-4-2}
\expandafter\def\csname sga@target@IV.16.4.20\endcsname{ega4.c04.prop.16-4-20}
\expandafter\def\csname sga@target@IV.16.4.21\endcsname{ega4.c04.cor.16-4-21}
\expandafter\def\csname sga@target@IV.16.4.21.1\endcsname{ega4.c04.cor.16-4-21-1}
\expandafter\def\csname sga@target@IV.16.4.22\endcsname{ega4.c04.cor.16-4-22}
\expandafter\def\csname sga@target@IV.16.4.23\endcsname{ega4.c04.prop.16-4-23}
\expandafter\def\csname sga@target@IV.16.4.23.1\endcsname{ega4.c04.prop.16-4-23-1}
\expandafter\def\csname sga@target@IV.16.4.24\endcsname{ega4.c04.statement.16-4-24}
\expandafter\def\csname sga@target@IV.16.4.3\endcsname{ega4.c04.statement.16-4-3}
\expandafter\def\csname sga@target@IV.16.4.3.1\endcsname{ega4.c04.statement.16-4-3-1}
\expandafter\def\csname sga@target@IV.16.4.3.2\endcsname{ega4.c04.statement.16-4-3-2}
\expandafter\def\csname sga@target@IV.16.4.3.3\endcsname{ega4.c04.eq.16-4-3-3}
\expandafter\def\csname sga@target@IV.16.4.3.4\endcsname{ega4.c04.eq.16-4-3-4}
\expandafter\def\csname sga@target@IV.16.4.3.5\endcsname{ega4.c04.eq.16-4-3-5}
\expandafter\def\csname sga@target@IV.16.4.3.6\endcsname{ega4.c04.eq.16-4-3-6}
\expandafter\def\csname sga@target@IV.16.4.3.7\endcsname{ega4.c04.eq.16-4-3-7}
\expandafter\def\csname sga@target@IV.16.4.4\endcsname{ega4.c04.statement.16-4-4}
\expandafter\def\csname sga@target@IV.16.4.5\endcsname{ega4.c04.prop.16-4-5}
\expandafter\def\csname sga@target@IV.16.4.6\endcsname{ega4.c04.cor.16-4-6}
\expandafter\def\csname sga@target@IV.16.4.7\endcsname{ega4.c04.statement.16-4-7}
\expandafter\def\csname sga@target@IV.16.4.7.1\endcsname{ega4.c04.statement.16-4-7-1}
\expandafter\def\csname sga@target@IV.16.4.7.2\endcsname{ega4.c04.eq.16-4-7-2}
\expandafter\def\csname sga@target@IV.16.4.7.3\endcsname{ega4.c04.eq.16-4-7-3}
\expandafter\def\csname sga@target@IV.16.4.8\endcsname{ega4.c04.prop.16-4-8}
\expandafter\def\csname sga@target@IV.16.4.8.1\endcsname{ega4.c04.prop.16-4-8-1}
\expandafter\def\csname sga@target@IV.16.4.8.2\endcsname{ega4.c04.eq.16-4-8-2}
\expandafter\def\csname sga@target@IV.16.4.8.3\endcsname{ega4.c04.lem.16-4-8-3}
\expandafter\def\csname sga@target@IV.16.4.8.4\endcsname{ega4.c04.eq.16-4-8-4}
\expandafter\def\csname sga@target@IV.16.4.9\endcsname{ega4.c04.cor.16-4-9}
\expandafter\def\csname sga@target@IV.16.4.9.1\endcsname{ega4.c04.cor.16-4-9-1}
\expandafter\def\csname sga@target@IV.16.5\endcsname{ega4.c04.subsec.16-5}
\expandafter\def\csname sga@target@IV.16.5.1\endcsname{ega4.c04.statement.16-5-1}
\expandafter\def\csname sga@target@IV.16.5.1.1\endcsname{ega4.c04.eq.16-5-1-1}
\expandafter\def\csname sga@target@IV.16.5.1.2\endcsname{ega4.c04.eq.16-5-1-2}
\expandafter\def\csname sga@target@IV.16.5.10\endcsname{ega4.c04.statement.16-5-10}
\expandafter\def\csname sga@target@IV.16.5.10.1\endcsname{ega4.c04.statement.16-5-10-1}
\expandafter\def\csname sga@target@IV.16.5.10.2\endcsname{ega4.c04.statement.16-5-10-2}
\expandafter\def\csname sga@target@IV.16.5.11\endcsname{ega4.c04.prop.16-5-11}
\expandafter\def\csname sga@target@IV.16.5.12\endcsname{ega4.c04.statement.16-5-12}
\expandafter\def\csname sga@target@IV.16.5.12.1\endcsname{ega4.c04.statement.16-5-12-1}
\expandafter\def\csname sga@target@IV.16.5.12.2\endcsname{ega4.c04.eq.16-5-12-2}
\expandafter\def\csname sga@target@IV.16.5.13\endcsname{ega4.c04.statement.16-5-13}
\expandafter\def\csname sga@target@IV.16.5.13.1\endcsname{ega4.c04.statement.16-5-13-1}
\expandafter\def\csname sga@target@IV.16.5.13.2\endcsname{ega4.c04.eq.16-5-13-2}
\expandafter\def\csname sga@target@IV.16.5.13.3\endcsname{ega4.c04.eq.16-5-13-3}
\expandafter\def\csname sga@target@IV.16.5.13.4\endcsname{ega4.c04.eq.16-5-13-4}
\expandafter\def\csname sga@target@IV.16.5.13.5\endcsname{ega4.c04.eq.16-5-13-5}
\expandafter\def\csname sga@target@IV.16.5.14\endcsname{ega4.c04.statement.16-5-14}
\expandafter\def\csname sga@target@IV.16.5.14.1\endcsname{ega4.c04.statement.16-5-14-1}
\expandafter\def\csname sga@target@IV.16.5.14.2\endcsname{ega4.c04.eq.16-5-14-2}
\expandafter\def\csname sga@target@IV.16.5.14.3\endcsname{ega4.c04.eq.16-5-14-3}
\expandafter\def\csname sga@target@IV.16.5.15\endcsname{ega4.c04.statement.16-5-15}
\expandafter\def\csname sga@target@IV.16.5.16\endcsname{ega4.c04.prop.16-5-16}
\expandafter\def\csname sga@target@IV.16.5.17\endcsname{ega4.c04.prop.16-5-17}
\expandafter\def\csname sga@target@IV.16.5.18\endcsname{ega4.c04.cor.16-5-18}
\expandafter\def\csname sga@target@IV.16.5.2\endcsname{ega4.c04.prop.16-5-2}
\expandafter\def\csname sga@target@IV.16.5.3\endcsname{ega4.c04.prop.16-5-3}
\expandafter\def\csname sga@target@IV.16.5.3.1\endcsname{ega4.c04.eq.16-5-3-1}
\expandafter\def\csname sga@target@IV.16.5.4\endcsname{ega4.c04.statement.16-5-4}
\expandafter\def\csname sga@target@IV.16.5.5\endcsname{ega4.c04.cor.16-5-5}
\expandafter\def\csname sga@target@IV.16.5.5.1\endcsname{ega4.c04.cor.16-5-5-1}
\expandafter\def\csname sga@target@IV.16.5.6\endcsname{ega4.c04.cor.16-5-6}
\expandafter\def\csname sga@target@IV.16.5.7\endcsname{ega4.c04.statement.16-5-7}
\expandafter\def\csname sga@target@IV.16.5.7.1\endcsname{ega4.c04.statement.16-5-7-1}
\expandafter\def\csname sga@target@IV.16.5.8\endcsname{ega4.c04.statement.16-5-8}
\expandafter\def\csname sga@target@IV.16.5.8.1\endcsname{ega4.c04.eq.16-5-8-1}
\expandafter\def\csname sga@target@IV.16.5.8.2\endcsname{ega4.c04.eq.16-5-8-2}
\expandafter\def\csname sga@target@IV.16.5.9\endcsname{ega4.c04.statement.16-5-9}
\expandafter\def\csname sga@target@IV.16.5.9.1\endcsname{ega4.c04.statement.16-5-9-1}
\expandafter\def\csname sga@target@IV.16.6\endcsname{ega4.c04.subsec.16-6}
\expandafter\def\csname sga@target@IV.16.6.1\endcsname{ega4.c04.statement.16-6-1}
\expandafter\def\csname sga@target@IV.16.6.1.1\endcsname{ega4.c04.statement.16-6-1-1}
\expandafter\def\csname sga@target@IV.16.6.1.2\endcsname{ega4.c04.statement.16-6-1-2}
\expandafter\def\csname sga@target@IV.16.6.2\endcsname{ega4.c04.thm.16-6-2}
\expandafter\def\csname sga@target@IV.16.6.2.1\endcsname{ega4.c04.eq.16-6-2-1}
\expandafter\def\csname sga@target@IV.16.6.2.2\endcsname{ega4.c04.eq.16-6-2-2}
\expandafter\def\csname sga@target@IV.16.6.2.3\endcsname{ega4.c04.eq.16-6-2-3}
\expandafter\def\csname sga@target@IV.16.6.2.4\endcsname{ega4.c04.eq.16-6-2-4}
\expandafter\def\csname sga@target@IV.16.6.2.5\endcsname{ega4.c04.eq.16-6-2-5}
\expandafter\def\csname sga@target@IV.16.6.3\endcsname{ega4.c04.def.16-6-3}
\expandafter\def\csname sga@target@IV.16.6.4\endcsname{ega4.c04.prop.16-6-4}
\expandafter\def\csname sga@target@IV.16.6.4.1\endcsname{ega4.c04.prop.16-6-4-1}
\expandafter\def\csname sga@target@IV.16.6.5\endcsname{ega4.c04.statement.16-6-5}
\expandafter\def\csname sga@target@IV.16.6.5.1\endcsname{ega4.c04.statement.16-6-5-1}
\expandafter\def\csname sga@target@IV.16.7\endcsname{ega4.c04.subsec.16-7}
\expandafter\def\csname sga@target@IV.16.7.1\endcsname{ega4.c04.statement.16-7-1}
\expandafter\def\csname sga@target@IV.16.7.1.1\endcsname{ega4.c04.eq.16-7-1-1}
\expandafter\def\csname sga@target@IV.16.7.1.2\endcsname{ega4.c04.eq.16-7-1-2}
\expandafter\def\csname sga@target@IV.16.7.10\endcsname{ega4.c04.rem.16-7-10}
\expandafter\def\csname sga@target@IV.16.7.2\endcsname{ega4.c04.statement.16-7-2}
\expandafter\def\csname sga@target@IV.16.7.2.1\endcsname{ega4.c04.statement.16-7-2-1}
\expandafter\def\csname sga@target@IV.16.7.3\endcsname{ega4.c04.prop.16-7-3}
\expandafter\def\csname sga@target@IV.16.7.4\endcsname{ega4.c04.statement.16-7-4}
\expandafter\def\csname sga@target@IV.16.7.4.1\endcsname{ega4.c04.statement.16-7-4-1}
\expandafter\def\csname sga@target@IV.16.7.5\endcsname{ega4.c04.statement.16-7-5}
\expandafter\def\csname sga@target@IV.16.7.5.1\endcsname{ega4.c04.statement.16-7-5-1}
\expandafter\def\csname sga@target@IV.16.7.5.2\endcsname{ega4.c04.statement.16-7-5-2}
\expandafter\def\csname sga@target@IV.16.7.5.3\endcsname{ega4.c04.statement.16-7-5-3}
\expandafter\def\csname sga@target@IV.16.7.5.4\endcsname{ega4.c04.statement.16-7-5-4}
\expandafter\def\csname sga@target@IV.16.7.6\endcsname{ega4.c04.prop.16-7-6}
\expandafter\def\csname sga@target@IV.16.7.7\endcsname{ega4.c04.statement.16-7-7}
\expandafter\def\csname sga@target@IV.16.7.7.1\endcsname{ega4.c04.eq.16-7-7-1}
\expandafter\def\csname sga@target@IV.16.7.7.2\endcsname{ega4.c04.eq.16-7-7-2}
\expandafter\def\csname sga@target@IV.16.7.8\endcsname{ega4.c04.statement.16-7-8}
\expandafter\def\csname sga@target@IV.16.7.8.1\endcsname{ega4.c04.statement.16-7-8-1}
\expandafter\def\csname sga@target@IV.16.7.8.2\endcsname{ega4.c04.eq.16-7-8-2}
\expandafter\def\csname sga@target@IV.16.7.9\endcsname{ega4.c04.statement.16-7-9}
\expandafter\def\csname sga@target@IV.16.7.9.1\endcsname{ega4.c04.statement.16-7-9-1}
\expandafter\def\csname sga@target@IV.16.8\endcsname{ega4.c04.subsec.16-8}
\expandafter\def\csname sga@target@IV.16.8.1\endcsname{ega4.c04.def.16-8-1}
\expandafter\def\csname sga@target@IV.16.8.10\endcsname{ega4.c04.cor.16-8-10}
\expandafter\def\csname sga@target@IV.16.8.11\endcsname{ega4.c04.statement.16-8-11}
\expandafter\def\csname sga@target@IV.16.8.2\endcsname{ega4.c04.statement.16-8-2}
\expandafter\def\csname sga@target@IV.16.8.3\endcsname{ega4.c04.statement.16-8-3}
\expandafter\def\csname sga@target@IV.16.8.3.1\endcsname{ega4.c04.eq.16-8-3-1}
\expandafter\def\csname sga@target@IV.16.8.4\endcsname{ega4.c04.prop.16-8-4}
\expandafter\def\csname sga@target@IV.16.8.4.1\endcsname{ega4.c04.prop.16-8-4-1}
\expandafter\def\csname sga@target@IV.16.8.5\endcsname{ega4.c04.statement.16-8-5}
\expandafter\def\csname sga@target@IV.16.8.5.1\endcsname{ega4.c04.eq.16-8-5-1}
\expandafter\def\csname sga@target@IV.16.8.5.2\endcsname{ega4.c04.eq.16-8-5-2}
\expandafter\def\csname sga@target@IV.16.8.6\endcsname{ega4.c04.prop.16-8-6}
\expandafter\def\csname sga@target@IV.16.8.7\endcsname{ega4.c04.statement.16-8-7}
\expandafter\def\csname sga@target@IV.16.8.7.1\endcsname{ega4.c04.eq.16-8-7-1}
\expandafter\def\csname sga@target@IV.16.8.8\endcsname{ega4.c04.prop.16-8-8}
\expandafter\def\csname sga@target@IV.16.8.8.1\endcsname{ega4.c04.eq.16-8-8-1}
\expandafter\def\csname sga@target@IV.16.8.8.2\endcsname{ega4.c04.eq.16-8-8-2}
\expandafter\def\csname sga@target@IV.16.8.9\endcsname{ega4.c04.prop.16-8-9}
\expandafter\def\csname sga@target@IV.16.8.9.1\endcsname{ega4.c04.lem.16-8-9-1}
\expandafter\def\csname sga@target@IV.16.8.9.2\endcsname{ega4.c04.lem.16-8-9-2}
\expandafter\def\csname sga@target@IV.16.8.9.3\endcsname{ega4.c04.lem.16-8-9-3}
\expandafter\def\csname sga@target@IV.16.8.9.4\endcsname{ega4.c04.eq.16-8-9-4}
\expandafter\def\csname sga@target@IV.16.9\endcsname{ega4.c04.subsec.16-9}
\expandafter\def\csname sga@target@IV.16.9.1\endcsname{ega4.c04.def.16-9-1}
\expandafter\def\csname sga@target@IV.16.9.10\endcsname{ega4.c04.prop.16-9-10}
\expandafter\def\csname sga@target@IV.16.9.11\endcsname{ega4.c04.cor.16-9-11}
\expandafter\def\csname sga@target@IV.16.9.12\endcsname{ega4.c04.statement.16-9-12}
\expandafter\def\csname sga@target@IV.16.9.13\endcsname{ega4.c04.prop.16-9-13}
\expandafter\def\csname sga@target@IV.16.9.13.1\endcsname{ega4.c04.prop.16-9-13-1}
\expandafter\def\csname sga@target@IV.16.9.13.2\endcsname{ega4.c04.lem.16-9-13-2}
\expandafter\def\csname sga@target@IV.16.9.13.3\endcsname{ega4.c04.lem.16-9-13-3}
\expandafter\def\csname sga@target@IV.16.9.13.4\endcsname{ega4.c04.eq.16-9-13-4}
\expandafter\def\csname sga@target@IV.16.9.13.5\endcsname{ega4.c04.cor.16-9-13-5}
\expandafter\def\csname sga@target@IV.16.9.13.6\endcsname{ega4.c04.cor.16-9-13-6}
\expandafter\def\csname sga@target@IV.16.9.13.7\endcsname{ega4.c04.eq.16-9-13-7}
\expandafter\def\csname sga@target@IV.16.9.2\endcsname{ega4.c04.def.16-9-2}
\expandafter\def\csname sga@target@IV.16.9.3\endcsname{ega4.c04.prop.16-9-3}
\expandafter\def\csname sga@target@IV.16.9.4\endcsname{ega4.c04.cor.16-9-4}
\expandafter\def\csname sga@target@IV.16.9.4.1\endcsname{ega4.c04.eq.16-9-4-1}
\expandafter\def\csname sga@target@IV.16.9.5\endcsname{ega4.c04.cor.16-9-5}
\expandafter\def\csname sga@target@IV.16.9.6\endcsname{ega4.c04.statement.16-9-6}
\expandafter\def\csname sga@target@IV.16.9.7\endcsname{ega4.c04.statement.16-9-7}
\expandafter\def\csname sga@target@IV.16.9.8\endcsname{ega4.c04.prop.16-9-8}
\expandafter\def\csname sga@target@IV.16.9.9\endcsname{ega4.c04.cor.16-9-9}
\expandafter\def\csname sga@target@IV.17.1.1\endcsname{ega4.c04.def.17-1-1}
\expandafter\def\csname sga@target@IV.17.1.1.1\endcsname{ega4.c04.def.17-1-1-1}
\expandafter\def\csname sga@target@IV.17.1.2\endcsname{ega4.c04.rem.17-1-2}
\expandafter\def\csname sga@target@IV.17.1.2.1\endcsname{ega4.c04.eq.17-1-2-1}
\expandafter\def\csname sga@target@IV.17.1.3\endcsname{ega4.c04.prop.17-1-3}
\expandafter\def\csname sga@target@IV.17.1.4\endcsname{ega4.c04.prop.17-1-4}
\expandafter\def\csname sga@target@IV.17.1.5\endcsname{ega4.c04.cor.17-1-5}
\expandafter\def\csname sga@target@IV.17.1.6\endcsname{ega4.c04.prop.17-1-6}
\expandafter\def\csname sga@target@IV.17.10.1\endcsname{ega4.c04.def.17-10-1}
\expandafter\def\csname sga@target@IV.17.10.2\endcsname{ega4.c04.prop.17-10-2}
\expandafter\def\csname sga@target@IV.17.10.3\endcsname{ega4.c04.cor.17-10-3}
\expandafter\def\csname sga@target@IV.17.10.3.1\endcsname{ega4.c04.cor.17-10-3-1}
\expandafter\def\csname sga@target@IV.17.10.4\endcsname{ega4.c04.cor.17-10-4}
\expandafter\def\csname sga@target@IV.17.10.4.1\endcsname{ega4.c04.eq.17-10-4-1}
\expandafter\def\csname sga@target@IV.17.11.1\endcsname{ega4.c04.thm.17-11-1}
\expandafter\def\csname sga@target@IV.17.11.1.1\endcsname{ega4.c04.eq.17-11-1-1}
\expandafter\def\csname sga@target@IV.17.11.1.2\endcsname{ega4.c04.eq.17-11-1-2}
\expandafter\def\csname sga@target@IV.17.11.2\endcsname{ega4.c04.cor.17-11-2}
\expandafter\def\csname sga@target@IV.17.11.3\endcsname{ega4.c04.cor.17-11-3}
\expandafter\def\csname sga@target@IV.17.11.4\endcsname{ega4.c04.cor.17-11-4}
\expandafter\def\csname sga@target@IV.17.11.5\endcsname{ega4.c04.prop.17-11-5}
\expandafter\def\csname sga@target@IV.17.12.1\endcsname{ega4.c04.thm.17-12-1}
\expandafter\def\csname sga@target@IV.17.12.2\endcsname{ega4.c04.cor.17-12-2}
\expandafter\def\csname sga@target@IV.17.12.3\endcsname{ega4.c04.cor.17-12-3}
\expandafter\def\csname sga@target@IV.17.12.4\endcsname{ega4.c04.prop.17-12-4}
\expandafter\def\csname sga@target@IV.17.12.5\endcsname{ega4.c04.prop.17-12-5}
\expandafter\def\csname sga@target@IV.17.12.6\endcsname{ega4.c04.cor.17-12-6}
\expandafter\def\csname sga@target@IV.17.12.7\endcsname{ega4.c04.ex.17-12-7}
\expandafter\def\csname sga@target@IV.17.12.8\endcsname{ega4.c04.rem.17-12-8}
\expandafter\def\csname sga@target@IV.17.13\endcsname{ega4.c04.subsec.17-13.key-iv-17-13}
\expandafter\def\csname sga@target@IV.17.13.1\endcsname{ega4.c04.statement.17-13-1}
\expandafter\def\csname sga@target@IV.17.13.1.1\endcsname{ega4.c04.statement.17-13-1-1}
\expandafter\def\csname sga@target@IV.17.13.1.2\endcsname{ega4.c04.statement.17-13-1-2}
\expandafter\def\csname sga@target@IV.17.13.1.3\endcsname{ega4.c04.eq.17-13-1-3}
\expandafter\def\csname sga@target@IV.17.13.10\endcsname{ega4.c04.cor.17-13-10}
\expandafter\def\csname sga@target@IV.17.13.11\endcsname{ega4.c04.statement.17-13-11}
\expandafter\def\csname sga@target@IV.17.13.11.1\endcsname{ega4.c04.statement.17-13-11-1}
\expandafter\def\csname sga@target@IV.17.13.12\endcsname{ega4.c04.cor.17-13-12}
\expandafter\def\csname sga@target@IV.17.13.12.1\endcsname{ega4.c04.cor.17-13-12-1}
\expandafter\def\csname sga@target@IV.17.13.12.2\endcsname{ega4.c04.eq.17-13-12-2}
\expandafter\def\csname sga@target@IV.17.13.2\endcsname{ega4.c04.prop.17-13-2}
\expandafter\def\csname sga@target@IV.17.13.2.1\endcsname{ega4.c04.eq.17-13-2-1}
\expandafter\def\csname sga@target@IV.17.13.2.2\endcsname{ega4.c04.eq.17-13-2-2}
\expandafter\def\csname sga@target@IV.17.13.3\endcsname{ega4.c04.def.17-13-3}
\expandafter\def\csname sga@target@IV.17.13.4\endcsname{ega4.c04.statement.17-13-4}
\expandafter\def\csname sga@target@IV.17.13.5\endcsname{ega4.c04.statement.17-13-5}
\expandafter\def\csname sga@target@IV.17.13.5.1\endcsname{ega4.c04.statement.17-13-5-1}
\expandafter\def\csname sga@target@IV.17.13.5.2\endcsname{ega4.c04.eq.17-13-5-2}
\expandafter\def\csname sga@target@IV.17.13.5.3\endcsname{ega4.c04.eq.17-13-5-3}
\expandafter\def\csname sga@target@IV.17.13.5.4\endcsname{ega4.c04.eq.17-13-5-4}
\expandafter\def\csname sga@target@IV.17.13.6\endcsname{ega4.c04.cor.17-13-6}
\expandafter\def\csname sga@target@IV.17.13.6.1\endcsname{ega4.c04.eq.17-13-6-1}
\expandafter\def\csname sga@target@IV.17.13.6.2\endcsname{ega4.c04.eq.17-13-6-2}
\expandafter\def\csname sga@target@IV.17.13.7\endcsname{ega4.c04.statement.17-13-7}
\expandafter\def\csname sga@target@IV.17.13.7.1\endcsname{ega4.c04.eq.17-13-7-1}
\expandafter\def\csname sga@target@IV.17.13.8\endcsname{ega4.c04.prop.17-13-8}
\expandafter\def\csname sga@target@IV.17.13.8.1\endcsname{ega4.c04.eq.17-13-8-1}
\expandafter\def\csname sga@target@IV.17.13.9\endcsname{ega4.c04.statement.17-13-9}
\expandafter\def\csname sga@target@IV.17.13.9.1\endcsname{ega4.c04.eq.17-13-9-1}
\expandafter\def\csname sga@target@IV.17.13.9.2\endcsname{ega4.c04.eq.17-13-9-2}
\expandafter\def\csname sga@target@IV.17.14.1\endcsname{ega4.c04.prop.17-14-1}
\expandafter\def\csname sga@target@IV.17.14.2\endcsname{ega4.c04.prop.17-14-2}
\expandafter\def\csname sga@target@IV.17.14.3\endcsname{ega4.c04.rem.17-14-3}
\expandafter\def\csname sga@target@IV.17.15.1\endcsname{ega4.c04.prop.17-15-1}
\expandafter\def\csname sga@target@IV.17.15.10\endcsname{ega4.c04.statement.17-15-10}
\expandafter\def\csname sga@target@IV.17.15.11\endcsname{ega4.c04.prop.17-15-11}
\expandafter\def\csname sga@target@IV.17.15.12\endcsname{ega4.c04.prop.17-15-12}
\expandafter\def\csname sga@target@IV.17.15.13\endcsname{ega4.c04.cor.17-15-13}
\expandafter\def\csname sga@target@IV.17.15.14\endcsname{ega4.c04.prop.17-15-14}
\expandafter\def\csname sga@target@IV.17.15.14.1\endcsname{ega4.c04.lem.17-15-14-1}
\expandafter\def\csname sga@target@IV.17.15.14.2\endcsname{ega4.c04.lem.17-15-14-2}
\expandafter\def\csname sga@target@IV.17.15.15\endcsname{ega4.c04.prop.17-15-15}
\expandafter\def\csname sga@target@IV.17.15.2\endcsname{ega4.c04.cor.17-15-2}
\expandafter\def\csname sga@target@IV.17.15.3\endcsname{ega4.c04.prop.17-15-3}
\expandafter\def\csname sga@target@IV.17.15.4\endcsname{ega4.c04.cor.17-15-4}
\expandafter\def\csname sga@target@IV.17.15.5\endcsname{ega4.c04.prop.17-15-5}
\expandafter\def\csname sga@target@IV.17.15.6\endcsname{ega4.c04.cor.17-15-6}
\expandafter\def\csname sga@target@IV.17.15.7\endcsname{ega4.c04.cor.17-15-7}
\expandafter\def\csname sga@target@IV.17.15.8\endcsname{ega4.c04.prop.17-15-8}
\expandafter\def\csname sga@target@IV.17.15.8.1\endcsname{ega4.c04.eq.17-15-8-1}
\expandafter\def\csname sga@target@IV.17.15.8.2\endcsname{ega4.c04.eq.17-15-8-2}
\expandafter\def\csname sga@target@IV.17.15.9\endcsname{ega4.c04.cor.17-15-9}
\expandafter\def\csname sga@target@IV.17.16\endcsname{ega4.c04.subsec.17-16.key-iv-17-16}
\expandafter\def\csname sga@target@IV.17.16.1\endcsname{ega4.c04.prop.17-16-1}
\expandafter\def\csname sga@target@IV.17.16.2\endcsname{ega4.c04.cor.17-16-2}
\expandafter\def\csname sga@target@IV.17.16.3\endcsname{ega4.c04.cor.17-16-3}
\expandafter\def\csname sga@target@IV.17.16.4\endcsname{ega4.c04.prop.17-16-4}
\expandafter\def\csname sga@target@IV.17.16.5\endcsname{ega4.c04.cor.17-16-5}
\expandafter\def\csname sga@target@IV.17.16.6\endcsname{ega4.c04.prop.17-16-6}
\expandafter\def\csname sga@target@IV.17.2.1\endcsname{ega4.c04.prop.17-2-1}
\expandafter\def\csname sga@target@IV.17.2.2\endcsname{ega4.c04.cor.17-2-2}
\expandafter\def\csname sga@target@IV.17.2.3\endcsname{ega4.c04.prop.17-2-3}
\expandafter\def\csname sga@target@IV.17.2.3.1\endcsname{ega4.c04.eq.17-2-3-1}
\expandafter\def\csname sga@target@IV.17.2.4\endcsname{ega4.c04.cor.17-2-4}
\expandafter\def\csname sga@target@IV.17.2.5\endcsname{ega4.c04.prop.17-2-5}
\expandafter\def\csname sga@target@IV.17.2.5.1\endcsname{ega4.c04.prop.17-2-5-1}
\expandafter\def\csname sga@target@IV.17.2.6\endcsname{ega4.c04.prop.17-2-6}
\expandafter\def\csname sga@target@IV.17.3.1\endcsname{ega4.c04.def.17-3-1}
\expandafter\def\csname sga@target@IV.17.3.2\endcsname{ega4.c04.subsec.17-3-2}
\expandafter\def\csname sga@target@IV.17.3.3\endcsname{ega4.c04.prop.17-3-3}
\expandafter\def\csname sga@target@IV.17.3.4\endcsname{ega4.c04.prop.17-3-4}
\expandafter\def\csname sga@target@IV.17.3.5\endcsname{ega4.c04.cor.17-3-5}
\expandafter\def\csname sga@target@IV.17.3.6\endcsname{ega4.c04.prop.17-3-6}
\expandafter\def\csname sga@target@IV.17.3.7\endcsname{ega4.c04.def.17-3-7}
\expandafter\def\csname sga@target@IV.17.3.8\endcsname{ega4.c04.prop.17-3-8}
\expandafter\def\csname sga@target@IV.17.3.9\endcsname{ega4.c04.cor.17-3-9}
\expandafter\def\csname sga@target@IV.17.4.1\endcsname{ega4.c04.thm.17-4-1}
\expandafter\def\csname sga@target@IV.17.4.1.1\endcsname{ega4.c04.eq.17-4-1-1}
\expandafter\def\csname sga@target@IV.17.4.1.2\endcsname{ega4.c04.rem.17-4-1-2}
\expandafter\def\csname sga@target@IV.17.4.10\endcsname{ega4.c04.rem.17-4-10}
\expandafter\def\csname sga@target@IV.17.4.2\endcsname{ega4.c04.cor.17-4-2}
\expandafter\def\csname sga@target@IV.17.4.3\endcsname{ega4.c04.cor.17-4-3}
\expandafter\def\csname sga@target@IV.17.4.4\endcsname{ega4.c04.prop.17-4-4}
\expandafter\def\csname sga@target@IV.17.4.5\endcsname{ega4.c04.statement.17-4-5}
\expandafter\def\csname sga@target@IV.17.4.6\endcsname{ega4.c04.prop.17-4-6}
\expandafter\def\csname sga@target@IV.17.4.7\endcsname{ega4.c04.cor.17-4-7}
\expandafter\def\csname sga@target@IV.17.4.8\endcsname{ega4.c04.cor.17-4-8}
\expandafter\def\csname sga@target@IV.17.4.9\endcsname{ega4.c04.prop.17-4-9}
\expandafter\def\csname sga@target@IV.17.5.1\endcsname{ega4.c04.thm.17-5-1}
\expandafter\def\csname sga@target@IV.17.5.1.1\endcsname{ega4.c04.eq.17-5-1-1}
\expandafter\def\csname sga@target@IV.17.5.1.2\endcsname{ega4.c04.eq.17-5-1-2}
\expandafter\def\csname sga@target@IV.17.5.2\endcsname{ega4.c04.cor.17-5-2}
\expandafter\def\csname sga@target@IV.17.5.3\endcsname{ega4.c04.prop.17-5-3}
\expandafter\def\csname sga@target@IV.17.5.4\endcsname{ega4.c04.rem.17-5-4}
\expandafter\def\csname sga@target@IV.17.5.5\endcsname{ega4.c04.prop.17-5-5}
\expandafter\def\csname sga@target@IV.17.5.6\endcsname{ega4.c04.cor.17-5-6}
\expandafter\def\csname sga@target@IV.17.5.7\endcsname{ega4.c04.prop.17-5-7}
\expandafter\def\csname sga@target@IV.17.5.8\endcsname{ega4.c04.prop.17-5-8}
\expandafter\def\csname sga@target@IV.17.6.1\endcsname{ega4.c04.thm.17-6-1}
\expandafter\def\csname sga@target@IV.17.6.2\endcsname{ega4.c04.cor.17-6-2}
\expandafter\def\csname sga@target@IV.17.6.3\endcsname{ega4.c04.prop.17-6-3}
\expandafter\def\csname sga@target@IV.17.6.4\endcsname{ega4.c04.prop.17-6-4}
\expandafter\def\csname sga@target@IV.17.7\endcsname{ega4.c04.subsec.17-7.key-iv-17-7}
\expandafter\def\csname sga@target@IV.17.7.1\endcsname{ega4.c04.prop.17-7-1}
\expandafter\def\csname sga@target@IV.17.7.10\endcsname{ega4.c04.prop.17-7-10}
\expandafter\def\csname sga@target@IV.17.7.11\endcsname{ega4.c04.prop.17-7-11}
\expandafter\def\csname sga@target@IV.17.7.11.1\endcsname{ega4.c04.lem.17-7-11-1}
\expandafter\def\csname sga@target@IV.17.7.2\endcsname{ega4.c04.cor.17-7-2}
\expandafter\def\csname sga@target@IV.17.7.3\endcsname{ega4.c04.cor.17-7-3}
\expandafter\def\csname sga@target@IV.17.7.4\endcsname{ega4.c04.prop.17-7-4}
\expandafter\def\csname sga@target@IV.17.7.5\endcsname{ega4.c04.lem.17-7-5}
\expandafter\def\csname sga@target@IV.17.7.6\endcsname{ega4.c04.rem.17-7-6}
\expandafter\def\csname sga@target@IV.17.7.7\endcsname{ega4.c04.prop.17-7-7}
\expandafter\def\csname sga@target@IV.17.7.8\endcsname{ega4.c04.prop.17-7-8}
\expandafter\def\csname sga@target@IV.17.7.9\endcsname{ega4.c04.cor.17-7-9}
\expandafter\def\csname sga@target@IV.17.8\endcsname{ega4.c04.subsec.17-8.key-iv-17-8}
\expandafter\def\csname sga@target@IV.17.8.1\endcsname{ega4.c04.prop.17-8-1}
\expandafter\def\csname sga@target@IV.17.8.2\endcsname{ega4.c04.prop.17-8-2}
\expandafter\def\csname sga@target@IV.17.8.3\endcsname{ega4.c04.rem.17-8-3}
\expandafter\def\csname sga@target@IV.17.9\endcsname{ega4.c04.subsec.17-9.key-iv-17-9}
\expandafter\def\csname sga@target@IV.17.9.1\endcsname{ega4.c04.thm.17-9-1}
\expandafter\def\csname sga@target@IV.17.9.2\endcsname{ega4.c04.cor.17-9-2}
\expandafter\def\csname sga@target@IV.17.9.3\endcsname{ega4.c04.cor.17-9-3}
\expandafter\def\csname sga@target@IV.17.9.4\endcsname{ega4.c04.cor.17-9-4}
\expandafter\def\csname sga@target@IV.17.9.5\endcsname{ega4.c04.cor.17-9-5}
\expandafter\def\csname sga@target@IV.17.9.6\endcsname{ega4.c04.prop.17-9-6}
\expandafter\def\csname sga@target@IV.18\endcsname{ega4.c04.sec.18.key-iv-18}
\expandafter\def\csname sga@target@IV.18.1.1\endcsname{ega4.c04.prop.18-1-1}
\expandafter\def\csname sga@target@IV.18.1.1.1\endcsname{ega4.c04.eq.18-1-1-1}
\expandafter\def\csname sga@target@IV.18.1.2\endcsname{ega4.c04.thm.18-1-2}
\expandafter\def\csname sga@target@IV.18.1.3\endcsname{ega4.c04.cor.18-1-3}
\expandafter\def\csname sga@target@IV.18.10.1\endcsname{ega4.c04.thm.18-10-1}
\expandafter\def\csname sga@target@IV.18.10.10\endcsname{ega4.c04.statement.18-10-10}
\expandafter\def\csname sga@target@IV.18.10.11\endcsname{ega4.c04.rem.18-10-11}
\expandafter\def\csname sga@target@IV.18.10.12\endcsname{ega4.c04.cor.18-10-12}
\expandafter\def\csname sga@target@IV.18.10.13\endcsname{ega4.c04.prop.18-10-13}
\expandafter\def\csname sga@target@IV.18.10.14\endcsname{ega4.c04.cor.18-10-14}
\expandafter\def\csname sga@target@IV.18.10.15\endcsname{ega4.c04.statement.18-10-15}
\expandafter\def\csname sga@target@IV.18.10.16\endcsname{ega4.c04.thm.18-10-16}
\expandafter\def\csname sga@target@IV.18.10.17\endcsname{ega4.c04.rem.18-10-17}
\expandafter\def\csname sga@target@IV.18.10.17.1\endcsname{ega4.c04.statement.18-10-17-1}
\expandafter\def\csname sga@target@IV.18.10.17.2\endcsname{ega4.c04.statement.18-10-17-2}
\expandafter\def\csname sga@target@IV.18.10.18\endcsname{ega4.c04.lem.18-10-18}
\expandafter\def\csname sga@target@IV.18.10.19\endcsname{ega4.c04.prop.18-10-19}
\expandafter\def\csname sga@target@IV.18.10.2\endcsname{ega4.c04.rem.18-10-2}
\expandafter\def\csname sga@target@IV.18.10.20\endcsname{ega4.c04.rem.18-10-20}
\expandafter\def\csname sga@target@IV.18.10.21\endcsname{ega4.c04.prop.18-10-21}
\expandafter\def\csname sga@target@IV.18.10.3\endcsname{ega4.c04.cor.18-10-3}
\expandafter\def\csname sga@target@IV.18.10.3.1\endcsname{ega4.c04.rem.18-10-3-1}
\expandafter\def\csname sga@target@IV.18.10.4\endcsname{ega4.c04.cor.18-10-4}
\expandafter\def\csname sga@target@IV.18.10.5\endcsname{ega4.c04.cor.18-10-5}
\expandafter\def\csname sga@target@IV.18.10.6\endcsname{ega4.c04.lem.18-10-6}
\expandafter\def\csname sga@target@IV.18.10.7\endcsname{ega4.c04.prop.18-10-7}
\expandafter\def\csname sga@target@IV.18.10.8\endcsname{ega4.c04.prop.18-10-8}
\expandafter\def\csname sga@target@IV.18.10.9\endcsname{ega4.c04.cor.18-10-9}
\expandafter\def\csname sga@target@IV.18.11.1\endcsname{ega4.c04.lem.18-11-1}
\expandafter\def\csname sga@target@IV.18.11.10\endcsname{ega4.c04.thm.18-11-10}
\expandafter\def\csname sga@target@IV.18.11.11\endcsname{ega4.c04.rem.18-11-11}
\expandafter\def\csname sga@target@IV.18.11.12\endcsname{ega4.c04.cor.18-11-12}
\expandafter\def\csname sga@target@IV.18.11.12.1\endcsname{ega4.c04.lem.18-11-12-1}
\expandafter\def\csname sga@target@IV.18.11.2\endcsname{ega4.c04.lem.18-11-2}
\expandafter\def\csname sga@target@IV.18.11.2.1\endcsname{ega4.c04.eq.18-11-2-1}
\expandafter\def\csname sga@target@IV.18.11.2.2\endcsname{ega4.c04.eq.18-11-2-2}
\expandafter\def\csname sga@target@IV.18.11.3\endcsname{ega4.c04.lem.18-11-3}
\expandafter\def\csname sga@target@IV.18.11.3.1\endcsname{ega4.c04.eq.18-11-3-1}
\expandafter\def\csname sga@target@IV.18.11.3.2\endcsname{ega4.c04.eq.18-11-3-2}
\expandafter\def\csname sga@target@IV.18.11.3.3\endcsname{ega4.c04.lem.18-11-3-3}
\expandafter\def\csname sga@target@IV.18.11.3.4\endcsname{ega4.c04.lem.18-11-3-4}
\expandafter\def\csname sga@target@IV.18.11.3.5\endcsname{ega4.c04.lem.18-11-3-5}
\expandafter\def\csname sga@target@IV.18.11.3.6\endcsname{ega4.c04.eq.18-11-3-6}
\expandafter\def\csname sga@target@IV.18.11.4\endcsname{ega4.c04.lem.18-11-4}
\expandafter\def\csname sga@target@IV.18.11.5\endcsname{ega4.c04.prop.18-11-5}
\expandafter\def\csname sga@target@IV.18.11.5.1\endcsname{ega4.c04.lem.18-11-5-1}
\expandafter\def\csname sga@target@IV.18.11.5.2\endcsname{ega4.c04.lem.18-11-5-2}
\expandafter\def\csname sga@target@IV.18.11.6\endcsname{ega4.c04.rem.18-11-6}
\expandafter\def\csname sga@target@IV.18.11.7\endcsname{ega4.c04.cor.18-11-7}
\expandafter\def\csname sga@target@IV.18.11.8\endcsname{ega4.c04.cor.18-11-8}
\expandafter\def\csname sga@target@IV.18.11.9\endcsname{ega4.c04.cor.18-11-9}
\expandafter\def\csname sga@target@IV.18.12.1\endcsname{ega4.c04.thm.18-12-1}
\expandafter\def\csname sga@target@IV.18.12.10\endcsname{ega4.c04.cor.18-12-10}
\expandafter\def\csname sga@target@IV.18.12.11\endcsname{ega4.c04.cor.18-12-11}
\expandafter\def\csname sga@target@IV.18.12.12\endcsname{ega4.c04.prop.18-12-12}
\expandafter\def\csname sga@target@IV.18.12.13\endcsname{ega4.c04.cor.18-12-13}
\expandafter\def\csname sga@target@IV.18.12.14\endcsname{ega4.c04.rem.18-12-14}
\expandafter\def\csname sga@target@IV.18.12.15\endcsname{ega4.c04.prop.18-12-15}
\expandafter\def\csname sga@target@IV.18.12.16\endcsname{ega4.c04.prop.18-12-16}
\expandafter\def\csname sga@target@IV.18.12.17\endcsname{ega4.c04.cor.18-12-17}
\expandafter\def\csname sga@target@IV.18.12.2\endcsname{ega4.c04.rem.18-12-2}
\expandafter\def\csname sga@target@IV.18.12.3\endcsname{ega4.c04.cor.18-12-3}
\expandafter\def\csname sga@target@IV.18.12.4\endcsname{ega4.c04.cor.18-12-4}
\expandafter\def\csname sga@target@IV.18.12.5\endcsname{ega4.c04.rem.18-12-5}
\expandafter\def\csname sga@target@IV.18.12.6\endcsname{ega4.c04.cor.18-12-6}
\expandafter\def\csname sga@target@IV.18.12.7\endcsname{ega4.c04.prop.18-12-7}
\expandafter\def\csname sga@target@IV.18.12.7.1\endcsname{ega4.c04.lem.18-12-7-1}
\expandafter\def\csname sga@target@IV.18.12.8\endcsname{ega4.c04.prop.18-12-8}
\expandafter\def\csname sga@target@IV.18.12.9\endcsname{ega4.c04.rem.18-12-9}
\expandafter\def\csname sga@target@IV.18.2.1\endcsname{ega4.c04.statement.18-2-1}
\expandafter\def\csname sga@target@IV.18.2.1.1\endcsname{ega4.c04.statement.18-2-1-1}
\expandafter\def\csname sga@target@IV.18.2.1.2\endcsname{ega4.c04.eq.18-2-1-2}
\expandafter\def\csname sga@target@IV.18.2.10\endcsname{ega4.c04.statement.18-2-10}
\expandafter\def\csname sga@target@IV.18.2.2\endcsname{ega4.c04.statement.18-2-2}
\expandafter\def\csname sga@target@IV.18.2.2.1\endcsname{ega4.c04.eq.18-2-2-1}
\expandafter\def\csname sga@target@IV.18.2.2.2\endcsname{ega4.c04.eq.18-2-2-2}
\expandafter\def\csname sga@target@IV.18.2.2.3\endcsname{ega4.c04.eq.18-2-2-3}
\expandafter\def\csname sga@target@IV.18.2.2.4\endcsname{ega4.c04.eq.18-2-2-4}
\expandafter\def\csname sga@target@IV.18.2.3\endcsname{ega4.c04.prop.18-2-3}
\expandafter\def\csname sga@target@IV.18.2.4\endcsname{ega4.c04.cor.18-2-4}
\expandafter\def\csname sga@target@IV.18.2.5\endcsname{ega4.c04.cor.18-2-5}
\expandafter\def\csname sga@target@IV.18.2.6\endcsname{ega4.c04.cor.18-2-6}
\expandafter\def\csname sga@target@IV.18.2.7\endcsname{ega4.c04.item.18-2-7}
\expandafter\def\csname sga@target@IV.18.2.7.1\endcsname{ega4.c04.eq.18-2-7-1}
\expandafter\def\csname sga@target@IV.18.2.7.2\endcsname{ega4.c04.eq.18-2-7-2}
\expandafter\def\csname sga@target@IV.18.2.8\endcsname{ega4.c04.prop.18-2-8}
\expandafter\def\csname sga@target@IV.18.2.9\endcsname{ega4.c04.cor.18-2-9}
\expandafter\def\csname sga@target@IV.18.3.1\endcsname{ega4.c04.prop.18-3-1}
\expandafter\def\csname sga@target@IV.18.3.2\endcsname{ega4.c04.prop.18-3-2}
\expandafter\def\csname sga@target@IV.18.3.2.1\endcsname{ega4.c04.lem.18-3-2-1}
\expandafter\def\csname sga@target@IV.18.3.2.2\endcsname{ega4.c04.eq.18-3-2-2}
\expandafter\def\csname sga@target@IV.18.3.3\endcsname{ega4.c04.statement.18-3-3}
\expandafter\def\csname sga@target@IV.18.3.4\endcsname{ega4.c04.thm.18-3-4}
\expandafter\def\csname sga@target@IV.18.4.1\endcsname{ega4.c04.lem.18-4-1}
\expandafter\def\csname sga@target@IV.18.4.10\endcsname{ega4.c04.prop.18-4-10}
\expandafter\def\csname sga@target@IV.18.4.11\endcsname{ega4.c04.statement.18-4-11}
\expandafter\def\csname sga@target@IV.18.4.12\endcsname{ega4.c04.cor.18-4-12}
\expandafter\def\csname sga@target@IV.18.4.13\endcsname{ega4.c04.cor.18-4-13}
\expandafter\def\csname sga@target@IV.18.4.14\endcsname{ega4.c04.cor.18-4-14}
\expandafter\def\csname sga@target@IV.18.4.2\endcsname{ega4.c04.prop.18-4-2}
\expandafter\def\csname sga@target@IV.18.4.3\endcsname{ega4.c04.cor.18-4-3}
\expandafter\def\csname sga@target@IV.18.4.4\endcsname{ega4.c04.lem.18-4-4}
\expandafter\def\csname sga@target@IV.18.4.5\endcsname{ega4.c04.prop.18-4-5}
\expandafter\def\csname sga@target@IV.18.4.6\endcsname{ega4.c04.thm.18-4-6}
\expandafter\def\csname sga@target@IV.18.4.6.1\endcsname{ega4.c04.lem.18-4-6-1}
\expandafter\def\csname sga@target@IV.18.4.7\endcsname{ega4.c04.cor.18-4-7}
\expandafter\def\csname sga@target@IV.18.4.8\endcsname{ega4.c04.cor.18-4-8}
\expandafter\def\csname sga@target@IV.18.4.9\endcsname{ega4.c04.prop.18-4-9}
\expandafter\def\csname sga@target@IV.18.5.1\endcsname{ega4.c04.statement.18-5-1}
\expandafter\def\csname sga@target@IV.18.5.1.1\endcsname{ega4.c04.eq.18-5-1-1}
\expandafter\def\csname sga@target@IV.18.5.10\endcsname{ega4.c04.prop.18-5-10}
\expandafter\def\csname sga@target@IV.18.5.11\endcsname{ega4.c04.thm.18-5-11}
\expandafter\def\csname sga@target@IV.18.5.12\endcsname{ega4.c04.cor.18-5-12}
\expandafter\def\csname sga@target@IV.18.5.13\endcsname{ega4.c04.rem.18-5-13}
\expandafter\def\csname sga@target@IV.18.5.13.1\endcsname{ega4.c04.lem.18-5-13-1}
\expandafter\def\csname sga@target@IV.18.5.14\endcsname{ega4.c04.prop.18-5-14}
\expandafter\def\csname sga@target@IV.18.5.15\endcsname{ega4.c04.prop.18-5-15}
\expandafter\def\csname sga@target@IV.18.5.16\endcsname{ega4.c04.item.18-5-16}
\expandafter\def\csname sga@target@IV.18.5.17\endcsname{ega4.c04.thm.18-5-17}
\expandafter\def\csname sga@target@IV.18.5.18\endcsname{ega4.c04.rem.18-5-18}
\expandafter\def\csname sga@target@IV.18.5.19\endcsname{ega4.c04.prop.18-5-19}
\expandafter\def\csname sga@target@IV.18.5.2\endcsname{ega4.c04.statement.18-5-2}
\expandafter\def\csname sga@target@IV.18.5.3\endcsname{ega4.c04.lem.18-5-3}
\expandafter\def\csname sga@target@IV.18.5.3.1\endcsname{ega4.c04.lem.18-5-3-1}
\expandafter\def\csname sga@target@IV.18.5.4\endcsname{ega4.c04.prop.18-5-4}
\expandafter\def\csname sga@target@IV.18.5.4.1\endcsname{ega4.c04.eq.18-5-4-1}
\expandafter\def\csname sga@target@IV.18.5.4.2\endcsname{ega4.c04.eq.18-5-4-2}
\expandafter\def\csname sga@target@IV.18.5.4.3\endcsname{ega4.c04.lem.18-5-4-3}
\expandafter\def\csname sga@target@IV.18.5.4.4\endcsname{ega4.c04.rem.18-5-4-4}
\expandafter\def\csname sga@target@IV.18.5.5\endcsname{ega4.c04.def.18-5-5}
\expandafter\def\csname sga@target@IV.18.5.6\endcsname{ega4.c04.prop.18-5-6}
\expandafter\def\csname sga@target@IV.18.5.7\endcsname{ega4.c04.rem.18-5-7}
\expandafter\def\csname sga@target@IV.18.5.8\endcsname{ega4.c04.def.18-5-8}
\expandafter\def\csname sga@target@IV.18.5.9\endcsname{ega4.c04.prop.18-5-9}
\expandafter\def\csname sga@target@IV.18.6.1\endcsname{ega4.c04.statement.18-6-1}
\expandafter\def\csname sga@target@IV.18.6.10\endcsname{ega4.c04.cor.18-6-10}
\expandafter\def\csname sga@target@IV.18.6.11\endcsname{ega4.c04.cor.18-6-11}
\expandafter\def\csname sga@target@IV.18.6.12\endcsname{ega4.c04.prop.18-6-12}
\expandafter\def\csname sga@target@IV.18.6.13\endcsname{ega4.c04.cor.18-6-13}
\expandafter\def\csname sga@target@IV.18.6.14\endcsname{ega4.c04.prop.18-6-14}
\expandafter\def\csname sga@target@IV.18.6.15\endcsname{ega4.c04.cor.18-6-15}
\expandafter\def\csname sga@target@IV.18.6.2\endcsname{ega4.c04.lem.18-6-2}
\expandafter\def\csname sga@target@IV.18.6.3\endcsname{ega4.c04.lem.18-6-3}
\expandafter\def\csname sga@target@IV.18.6.4\endcsname{ega4.c04.statement.18-6-4}
\expandafter\def\csname sga@target@IV.18.6.5\endcsname{ega4.c04.def.18-6-5}
\expandafter\def\csname sga@target@IV.18.6.6\endcsname{ega4.c04.thm.18-6-6}
\expandafter\def\csname sga@target@IV.18.6.7\endcsname{ega4.c04.statement.18-6-7}
\expandafter\def\csname sga@target@IV.18.6.8\endcsname{ega4.c04.prop.18-6-8}
\expandafter\def\csname sga@target@IV.18.6.9\endcsname{ega4.c04.thm.18-6-9}
\expandafter\def\csname sga@target@IV.18.6.9.1\endcsname{ega4.c04.lem.18-6-9-1}
\expandafter\def\csname sga@target@IV.18.6.9.2\endcsname{ega4.c04.lem.18-6-9-2}
\expandafter\def\csname sga@target@IV.18.7.1\endcsname{ega4.c04.statement.18-7-1}
\expandafter\def\csname sga@target@IV.18.7.2\endcsname{ega4.c04.prop.18-7-2}
\expandafter\def\csname sga@target@IV.18.7.3\endcsname{ega4.c04.cor.18-7-3}
\expandafter\def\csname sga@target@IV.18.7.4\endcsname{ega4.c04.cor.18-7-4}
\expandafter\def\csname sga@target@IV.18.7.5\endcsname{ega4.c04.prop.18-7-5}
\expandafter\def\csname sga@target@IV.18.7.5.1\endcsname{ega4.c04.lem.18-7-5-1}
\expandafter\def\csname sga@target@IV.18.7.6\endcsname{ega4.c04.cor.18-7-6}
\expandafter\def\csname sga@target@IV.18.7.7\endcsname{ega4.c04.rem.18-7-7}
\expandafter\def\csname sga@target@IV.18.8.1\endcsname{ega4.c04.prop.18-8-1}
\expandafter\def\csname sga@target@IV.18.8.10\endcsname{ega4.c04.prop.18-8-10}
\expandafter\def\csname sga@target@IV.18.8.11\endcsname{ega4.c04.rem.18-8-11}
\expandafter\def\csname sga@target@IV.18.8.12\endcsname{ega4.c04.prop.18-8-12}
\expandafter\def\csname sga@target@IV.18.8.13\endcsname{ega4.c04.cor.18-8-13}
\expandafter\def\csname sga@target@IV.18.8.14\endcsname{ega4.c04.cor.18-8-14}
\expandafter\def\csname sga@target@IV.18.8.15\endcsname{ega4.c04.prop.18-8-15}
\expandafter\def\csname sga@target@IV.18.8.16\endcsname{ega4.c04.cor.18-8-16}
\expandafter\def\csname sga@target@IV.18.8.17\endcsname{ega4.c04.prop.18-8-17}
\expandafter\def\csname sga@target@IV.18.8.18\endcsname{ega4.c04.prop.18-8-18}
\expandafter\def\csname sga@target@IV.18.8.2\endcsname{ega4.c04.def.18-8-2}
\expandafter\def\csname sga@target@IV.18.8.3\endcsname{ega4.c04.rem.18-8-3}
\expandafter\def\csname sga@target@IV.18.8.4\endcsname{ega4.c04.lem.18-8-4}
\expandafter\def\csname sga@target@IV.18.8.5\endcsname{ega4.c04.lem.18-8-5}
\expandafter\def\csname sga@target@IV.18.8.6\endcsname{ega4.c04.statement.18-8-6}
\expandafter\def\csname sga@target@IV.18.8.6.1\endcsname{ega4.c04.eq.18-8-6-1}
\expandafter\def\csname sga@target@IV.18.8.6.2\endcsname{ega4.c04.eq.18-8-6-2}
\expandafter\def\csname sga@target@IV.18.8.7\endcsname{ega4.c04.def.18-8-7}
\expandafter\def\csname sga@target@IV.18.8.8\endcsname{ega4.c04.prop.18-8-8}
\expandafter\def\csname sga@target@IV.18.8.8.1\endcsname{ega4.c04.statement.18-8-8-1}
\expandafter\def\csname sga@target@IV.18.8.9\endcsname{ega4.c04.statement.18-8-9}
\expandafter\def\csname sga@target@IV.18.9.1\endcsname{ega4.c04.thm.18-9-1}
\expandafter\def\csname sga@target@IV.18.9.10\endcsname{ega4.c04.rem.18-9-10}
\expandafter\def\csname sga@target@IV.18.9.11\endcsname{ega4.c04.cor.18-9-11}
\expandafter\def\csname sga@target@IV.18.9.2\endcsname{ega4.c04.cor.18-9-2}
\expandafter\def\csname sga@target@IV.18.9.3\endcsname{ega4.c04.cor.18-9-3}
\expandafter\def\csname sga@target@IV.18.9.4\endcsname{ega4.c04.cor.18-9-4}
\expandafter\def\csname sga@target@IV.18.9.5\endcsname{ega4.c04.cor.18-9-5}
\expandafter\def\csname sga@target@IV.18.9.5.1\endcsname{ega4.c04.eq.18-9-5-1}
\expandafter\def\csname sga@target@IV.18.9.6\endcsname{ega4.c04.statement.18-9-6}
\expandafter\def\csname sga@target@IV.18.9.7\endcsname{ega4.c04.thm.18-9-7}
\expandafter\def\csname sga@target@IV.18.9.7.1\endcsname{ega4.c04.rem.18-9-7-1}
\expandafter\def\csname sga@target@IV.18.9.7.2\endcsname{ega4.c04.statement.18-9-7-2}
\expandafter\def\csname sga@target@IV.18.9.7.3\endcsname{ega4.c04.lem.18-9-7-3}
\expandafter\def\csname sga@target@IV.18.9.7.4\endcsname{ega4.c04.eq.18-9-7-4}
\expandafter\def\csname sga@target@IV.18.9.7.5\endcsname{ega4.c04.lem.18-9-7-5}
\expandafter\def\csname sga@target@IV.18.9.7.6\endcsname{ega4.c04.rem.18-9-7-6}
\expandafter\def\csname sga@target@IV.18.9.7.7\endcsname{ega4.c04.lem.18-9-7-7}
\expandafter\def\csname sga@target@IV.18.9.7.8\endcsname{ega4.c04.rem.18-9-7-8}
\expandafter\def\csname sga@target@IV.18.9.8\endcsname{ega4.c04.cor.18-9-8}
\expandafter\def\csname sga@target@IV.18.9.9\endcsname{ega4.c04.cor.18-9-9}
\expandafter\def\csname sga@target@IV.19.1.1\endcsname{ega4.c04.prop.19-1-1}
\expandafter\def\csname sga@target@IV.19.1.2\endcsname{ega4.c04.cor.19-1-2}
\expandafter\def\csname sga@target@IV.19.1.3\endcsname{ega4.c04.def.19-1-3}
\expandafter\def\csname sga@target@IV.19.1.4\endcsname{ega4.c04.prop.19-1-4}
\expandafter\def\csname sga@target@IV.19.1.4.1\endcsname{ega4.c04.eq.19-1-4-1}
\expandafter\def\csname sga@target@IV.19.1.5\endcsname{ega4.c04.prop.19-1-5}
\expandafter\def\csname sga@target@IV.19.1.5.1\endcsname{ega4.c04.item.19-1-5-1}
\expandafter\def\csname sga@target@IV.19.1.6\endcsname{ega4.c04.cor.19-1-6}
\expandafter\def\csname sga@target@IV.19.1.6.1\endcsname{ega4.c04.lem.19-1-6-1}
\expandafter\def\csname sga@target@IV.19.1.7\endcsname{ega4.c04.item.19-1-7}
\expandafter\def\csname sga@target@IV.19.1.8\endcsname{ega4.c04.statement.19-1-8}
\expandafter\def\csname sga@target@IV.19.2.1\endcsname{ega4.c04.def.19-2-1}
\expandafter\def\csname sga@target@IV.19.2.2\endcsname{ega4.c04.def.19-2-2}
\expandafter\def\csname sga@target@IV.19.2.3\endcsname{ega4.c04.prop.19-2-3}
\expandafter\def\csname sga@target@IV.19.2.4\endcsname{ega4.c04.prop.19-2-4}
\expandafter\def\csname sga@target@IV.19.2.5\endcsname{ega4.c04.item.19-2-5}
\expandafter\def\csname sga@target@IV.19.2.6\endcsname{ega4.c04.cor.19-2-6}
\expandafter\def\csname sga@target@IV.19.2.7\endcsname{ega4.c04.prop.19-2-7}
\expandafter\def\csname sga@target@IV.19.2.8\endcsname{ega4.c04.prop.19-2-8}
\expandafter\def\csname sga@target@IV.19.2.9\endcsname{ega4.c04.prop.19-2-9}
\expandafter\def\csname sga@target@IV.19.3.1\endcsname{ega4.c04.def.19-3-1}
\expandafter\def\csname sga@target@IV.19.3.10\endcsname{ega4.c04.cor.19-3-10}
\expandafter\def\csname sga@target@IV.19.3.2\endcsname{ega4.c04.prop.19-3-2}
\expandafter\def\csname sga@target@IV.19.3.2.1\endcsname{ega4.c04.lem.19-3-2-1}
\expandafter\def\csname sga@target@IV.19.3.3\endcsname{ega4.c04.cor.19-3-3}
\expandafter\def\csname sga@target@IV.19.3.4\endcsname{ega4.c04.cor.19-3-4}
\expandafter\def\csname sga@target@IV.19.3.5\endcsname{ega4.c04.cor.19-3-5}
\expandafter\def\csname sga@target@IV.19.3.6\endcsname{ega4.c04.def.19-3-6}
\expandafter\def\csname sga@target@IV.19.3.7\endcsname{ega4.c04.prop.19-3-7}
\expandafter\def\csname sga@target@IV.19.3.8\endcsname{ega4.c04.cor.19-3-8}
\expandafter\def\csname sga@target@IV.19.3.9\endcsname{ega4.c04.prop.19-3-9}
\expandafter\def\csname sga@target@IV.19.4.1\endcsname{ega4.c04.statement.19-4-1}
\expandafter\def\csname sga@target@IV.19.4.1.1\endcsname{ega4.c04.subsec.19-4-1-1}
\expandafter\def\csname sga@target@IV.19.4.10\endcsname{ega4.c04.rem.19-4-10}
\expandafter\def\csname sga@target@IV.19.4.11\endcsname{ega4.c04.prop.19-4-11}
\expandafter\def\csname sga@target@IV.19.4.11.1\endcsname{ega4.c04.eq.19-4-11-1}
\expandafter\def\csname sga@target@IV.19.4.2\endcsname{ega4.c04.lem.19-4-2}
\expandafter\def\csname sga@target@IV.19.4.3\endcsname{ega4.c04.cor.19-4-3}
\expandafter\def\csname sga@target@IV.19.4.4\endcsname{ega4.c04.prop.19-4-4}
\expandafter\def\csname sga@target@IV.19.4.5\endcsname{ega4.c04.cor.19-4-5}
\expandafter\def\csname sga@target@IV.19.4.6\endcsname{ega4.c04.prop.19-4-6}
\expandafter\def\csname sga@target@IV.19.4.7\endcsname{ega4.c04.cor.19-4-7}
\expandafter\def\csname sga@target@IV.19.4.8\endcsname{ega4.c04.prop.19-4-8}
\expandafter\def\csname sga@target@IV.19.4.9\endcsname{ega4.c04.cor.19-4-9}
\expandafter\def\csname sga@target@IV.19.5.1\endcsname{ega4.c04.thm.19-5-1}
\expandafter\def\csname sga@target@IV.19.5.1.1\endcsname{ega4.c04.eq.19-5-1-1}
\expandafter\def\csname sga@target@IV.19.5.2\endcsname{ega4.c04.cor.19-5-2}
\expandafter\def\csname sga@target@IV.19.5.3\endcsname{ega4.c04.prop.19-5-3}
\expandafter\def\csname sga@target@IV.19.5.4\endcsname{ega4.c04.statement.19-5-4}
\expandafter\def\csname sga@target@IV.19.5.5\endcsname{ega4.c04.thm.19-5-5}
\expandafter\def\csname sga@target@IV.19.5.5.1\endcsname{ega4.c04.statement.19-5-5-1}
\expandafter\def\csname sga@target@IV.19.5.5.2\endcsname{ega4.c04.statement.19-5-5-2}
\expandafter\def\csname sga@target@IV.19.5.5.3\endcsname{ega4.c04.statement.19-5-5-3}
\expandafter\def\csname sga@target@IV.19.5.5.4\endcsname{ega4.c04.statement.19-5-5-4}
\expandafter\def\csname sga@target@IV.19.5.5.5\endcsname{ega4.c04.statement.19-5-5-5}
\expandafter\def\csname sga@target@IV.19.5.6\endcsname{ega4.c04.statement.19-5-6}
\expandafter\def\csname sga@target@IV.19.6.1\endcsname{ega4.c04.statement.19-6-1}
\expandafter\def\csname sga@target@IV.19.6.1.1\endcsname{ega4.c04.eq.19-6-1-1}
\expandafter\def\csname sga@target@IV.19.6.2\endcsname{ega4.c04.statement.19-6-2}
\expandafter\def\csname sga@target@IV.19.6.2.1\endcsname{ega4.c04.statement.19-6-2-1}
\expandafter\def\csname sga@target@IV.19.6.3\endcsname{ega4.c04.prop.19-6-3}
\expandafter\def\csname sga@target@IV.19.6.4\endcsname{ega4.c04.cor.19-6-4}
\expandafter\def\csname sga@target@IV.19.7\endcsname{ega4.c04.subsec.19-7.key-iv-19-7}
\expandafter\def\csname sga@target@IV.19.7.1\endcsname{ega4.c04.thm.19-7-1}
\expandafter\def\csname sga@target@IV.19.7.1.1\endcsname{ega4.c04.statement.19-7-1-1}
\expandafter\def\csname sga@target@IV.19.7.1.10\endcsname{ega4.c04.statement.19-7-1-10}
\expandafter\def\csname sga@target@IV.19.7.1.11\endcsname{ega4.c04.statement.19-7-1-11}
\expandafter\def\csname sga@target@IV.19.7.1.12\endcsname{ega4.c04.statement.19-7-1-12}
\expandafter\def\csname sga@target@IV.19.7.1.13\endcsname{ega4.c04.statement.19-7-1-13}
\expandafter\def\csname sga@target@IV.19.7.1.14\endcsname{ega4.c04.statement.19-7-1-14}
\expandafter\def\csname sga@target@IV.19.7.1.15\endcsname{ega4.c04.statement.19-7-1-15}
\expandafter\def\csname sga@target@IV.19.7.1.16\endcsname{ega4.c04.statement.19-7-1-16}
\expandafter\def\csname sga@target@IV.19.7.1.17\endcsname{ega4.c04.statement.19-7-1-17}
\expandafter\def\csname sga@target@IV.19.7.1.18\endcsname{ega4.c04.statement.19-7-1-18}
\expandafter\def\csname sga@target@IV.19.7.1.19\endcsname{ega4.c04.statement.19-7-1-19}
\expandafter\def\csname sga@target@IV.19.7.1.2\endcsname{ega4.c04.statement.19-7-1-2}
\expandafter\def\csname sga@target@IV.19.7.1.20\endcsname{ega4.c04.statement.19-7-1-20}
\expandafter\def\csname sga@target@IV.19.7.1.21\endcsname{ega4.c04.statement.19-7-1-21}
\expandafter\def\csname sga@target@IV.19.7.1.22\endcsname{ega4.c04.statement.19-7-1-22}
\expandafter\def\csname sga@target@IV.19.7.1.23\endcsname{ega4.c04.statement.19-7-1-23}
\expandafter\def\csname sga@target@IV.19.7.1.3\endcsname{ega4.c04.statement.19-7-1-3}
\expandafter\def\csname sga@target@IV.19.7.1.5\endcsname{ega4.c04.statement.19-7-1-5}
\expandafter\def\csname sga@target@IV.19.7.1.6\endcsname{ega4.c04.statement.19-7-1-6}
\expandafter\def\csname sga@target@IV.19.7.1.7\endcsname{ega4.c04.lem.19-7-1-7}
\expandafter\def\csname sga@target@IV.19.7.1.8\endcsname{ega4.c04.statement.19-7-1-8}
\expandafter\def\csname sga@target@IV.19.7.1.9\endcsname{ega4.c04.statement.19-7-1-9}
\expandafter\def\csname sga@target@IV.19.7.2\endcsname{ega4.c04.statement.19-7-2}
\expandafter\def\csname sga@target@IV.19.7.3\endcsname{ega4.c04.cor.19-7-3}
\expandafter\def\csname sga@target@IV.19.7.3.1\endcsname{ega4.c04.eq.19-7-3-1}
\expandafter\def\csname sga@target@IV.19.7.4\endcsname{ega4.c04.cor.19-7-4}
\expandafter\def\csname sga@target@IV.19.7.4.1\endcsname{ega4.c04.eq.19-7-4-1}
\expandafter\def\csname sga@target@IV.19.7.5\endcsname{ega4.c04.cor.19-7-5}
\expandafter\def\csname sga@target@IV.19.8.1\endcsname{ega4.c04.prop.19-8-1}
\expandafter\def\csname sga@target@IV.19.8.2\endcsname{ega4.c04.prop.19-8-2}
\expandafter\def\csname sga@target@IV.19.8.3\endcsname{ega4.c04.rem.19-8-3}
\expandafter\def\csname sga@target@IV.19.9.1\endcsname{ega4.c04.statement.19-9-1}
\expandafter\def\csname sga@target@IV.19.9.1.1\endcsname{ega4.c04.eq.19-9-1-1}
\expandafter\def\csname sga@target@IV.19.9.1.2\endcsname{ega4.c04.eq.19-9-1-2}
\expandafter\def\csname sga@target@IV.19.9.1.3\endcsname{ega4.c04.eq.19-9-1-3}
\expandafter\def\csname sga@target@IV.19.9.2\endcsname{ega4.c04.lem.19-9-2}
\expandafter\def\csname sga@target@IV.19.9.3\endcsname{ega4.c04.prop.19-9-3}
\expandafter\def\csname sga@target@IV.19.9.4\endcsname{ega4.c04.cor.19-9-4}
\expandafter\def\csname sga@target@IV.19.9.5\endcsname{ega4.c04.prop.19-9-5}
\expandafter\def\csname sga@target@IV.19.9.6\endcsname{ega4.c04.prop.19-9-6}
\expandafter\def\csname sga@target@IV.19.9.7\endcsname{ega4.c04.cor.19-9-7}
\expandafter\def\csname sga@target@IV.19.9.8\endcsname{ega4.c04.prop.19-9-8}
\expandafter\def\csname sga@target@IV.19.9.9\endcsname{ega4.c04.rem.19-9-9}
\expandafter\def\csname sga@target@IV.2.1.1\endcsname{ega4.c04.statement.2-1-1}
\expandafter\def\csname sga@target@IV.2.1.10\endcsname{ega4.c04.cor.2-1-10}
\expandafter\def\csname sga@target@IV.2.1.11\endcsname{ega4.c04.prop.2-1-11}
\expandafter\def\csname sga@target@IV.2.1.12\endcsname{ega4.c04.prop.2-1-12}
\expandafter\def\csname sga@target@IV.2.1.13\endcsname{ega4.c04.prop.2-1-13}
\expandafter\def\csname sga@target@IV.2.1.14\endcsname{ega4.c04.prop.2-1-14}
\expandafter\def\csname sga@target@IV.2.1.2\endcsname{ega4.c04.prop.2-1-2}
\expandafter\def\csname sga@target@IV.2.1.3\endcsname{ega4.c04.prop.2-1-3}
\expandafter\def\csname sga@target@IV.2.1.4\endcsname{ega4.c04.prop.2-1-4}
\expandafter\def\csname sga@target@IV.2.1.5\endcsname{ega4.c04.prop.2-1-5}
\expandafter\def\csname sga@target@IV.2.1.6\endcsname{ega4.c04.cor.2-1-6}
\expandafter\def\csname sga@target@IV.2.1.7\endcsname{ega4.c04.cor.2-1-7}
\expandafter\def\csname sga@target@IV.2.1.8\endcsname{ega4.c04.prop.2-1-8}
\expandafter\def\csname sga@target@IV.2.1.9\endcsname{ega4.c04.cor.2-1-9}
\expandafter\def\csname sga@target@IV.2.2.1\endcsname{ega4.c04.prop.2-2-1}
\expandafter\def\csname sga@target@IV.2.2.10\endcsname{ega4.c04.prop.2-2-10}
\expandafter\def\csname sga@target@IV.2.2.10.1\endcsname{ega4.c04.eq.2-2-10-1}
\expandafter\def\csname sga@target@IV.2.2.11\endcsname{ega4.c04.cor.2-2-11}
\expandafter\def\csname sga@target@IV.2.2.12\endcsname{ega4.c04.cor.2-2-12}
\expandafter\def\csname sga@target@IV.2.2.13\endcsname{ega4.c04.cor.2-2-13}
\expandafter\def\csname sga@target@IV.2.2.14\endcsname{ega4.c04.prop.2-2-14}
\expandafter\def\csname sga@target@IV.2.2.15\endcsname{ega4.c04.prop.2-2-15}
\expandafter\def\csname sga@target@IV.2.2.16\endcsname{ega4.c04.cor.2-2-16}
\expandafter\def\csname sga@target@IV.2.2.17\endcsname{ega4.c04.prop.2-2-17}
\expandafter\def\csname sga@target@IV.2.2.2\endcsname{ega4.c04.prop.2-2-2}
\expandafter\def\csname sga@target@IV.2.2.2.1\endcsname{ega4.c04.eq.2-2-2-1}
\expandafter\def\csname sga@target@IV.2.2.3\endcsname{ega4.c04.cor.2-2-3}
\expandafter\def\csname sga@target@IV.2.2.4\endcsname{ega4.c04.def.2-2-4}
\expandafter\def\csname sga@target@IV.2.2.5\endcsname{ega4.c04.cor.2-2-5}
\expandafter\def\csname sga@target@IV.2.2.6\endcsname{ega4.c04.cor.2-2-6}
\expandafter\def\csname sga@target@IV.2.2.7\endcsname{ega4.c04.prop.2-2-7}
\expandafter\def\csname sga@target@IV.2.2.8\endcsname{ega4.c04.cor.2-2-8}
\expandafter\def\csname sga@target@IV.2.2.8.1\endcsname{ega4.c04.eq.2-2-8-1}
\expandafter\def\csname sga@target@IV.2.2.9\endcsname{ega4.c04.rem.2-2-9}
\expandafter\def\csname sga@target@IV.2.3.1\endcsname{ega4.c04.lem.2-3-1}
\expandafter\def\csname sga@target@IV.2.3.1.1\endcsname{ega4.c04.eq.2-3-1-1}
\expandafter\def\csname sga@target@IV.2.3.10\endcsname{ega4.c04.thm.2-3-10}
\expandafter\def\csname sga@target@IV.2.3.11\endcsname{ega4.c04.cor.2-3-11}
\expandafter\def\csname sga@target@IV.2.3.12\endcsname{ega4.c04.cor.2-3-12}
\expandafter\def\csname sga@target@IV.2.3.13\endcsname{ega4.c04.cor.2-3-13}
\expandafter\def\csname sga@target@IV.2.3.14\endcsname{ega4.c04.cor.2-3-14}
\expandafter\def\csname sga@target@IV.2.3.15\endcsname{ega4.c04.rem.2-3-15}
\expandafter\def\csname sga@target@IV.2.3.2\endcsname{ega4.c04.prop.2-3-2}
\expandafter\def\csname sga@target@IV.2.3.3\endcsname{ega4.c04.statement.2-3-3}
\expandafter\def\csname sga@target@IV.2.3.4\endcsname{ega4.c04.prop.2-3-4}
\expandafter\def\csname sga@target@IV.2.3.5\endcsname{ega4.c04.cor.2-3-5}
\expandafter\def\csname sga@target@IV.2.3.6\endcsname{ega4.c04.prop.2-3-6}
\expandafter\def\csname sga@target@IV.2.3.7\endcsname{ega4.c04.prop.2-3-7}
\expandafter\def\csname sga@target@IV.2.3.8\endcsname{ega4.c04.cor.2-3-8}
\expandafter\def\csname sga@target@IV.2.3.9\endcsname{ega4.c04.rem.2-3-9}
\expandafter\def\csname sga@target@IV.2.4.1\endcsname{ega4.c04.statement.2-4-1}
\expandafter\def\csname sga@target@IV.2.4.10\endcsname{ega4.c04.cor.2-4-10}
\expandafter\def\csname sga@target@IV.2.4.11\endcsname{ega4.c04.rem.2-4-11}
\expandafter\def\csname sga@target@IV.2.4.2\endcsname{ega4.c04.def.2-4-2}
\expandafter\def\csname sga@target@IV.2.4.3\endcsname{ega4.c04.prop.2-4-3}
\expandafter\def\csname sga@target@IV.2.4.4\endcsname{ega4.c04.prop.2-4-4}
\expandafter\def\csname sga@target@IV.2.4.5\endcsname{ega4.c04.prop.2-4-5}
\expandafter\def\csname sga@target@IV.2.4.6\endcsname{ega4.c04.thm.2-4-6}
\expandafter\def\csname sga@target@IV.2.4.7\endcsname{ega4.c04.cor.2-4-7}
\expandafter\def\csname sga@target@IV.2.4.8\endcsname{ega4.c04.rem.2-4-8}
\expandafter\def\csname sga@target@IV.2.4.9\endcsname{ega4.c04.prop.2-4-9}
\expandafter\def\csname sga@target@IV.2.5.1\endcsname{ega4.c04.prop.2-5-1}
\expandafter\def\csname sga@target@IV.2.5.2\endcsname{ega4.c04.prop.2-5-2}
\expandafter\def\csname sga@target@IV.2.5.3\endcsname{ega4.c04.lem.2-5-3}
\expandafter\def\csname sga@target@IV.2.5.4\endcsname{ega4.c04.rem.2-5-4}
\expandafter\def\csname sga@target@IV.2.5.4.1\endcsname{ega4.c04.lem.2-5-4-1}
\expandafter\def\csname sga@target@IV.2.5.5\endcsname{ega4.c04.prop.2-5-5}
\expandafter\def\csname sga@target@IV.2.5.5.1\endcsname{ega4.c04.statement.2-5-5-1}
\expandafter\def\csname sga@target@IV.2.5.5.2\endcsname{ega4.c04.lem.2-5-5-2}
\expandafter\def\csname sga@target@IV.2.5.5.3\endcsname{ega4.c04.statement.2-5-5-3}
\expandafter\def\csname sga@target@IV.2.5.5.4\endcsname{ega4.c04.rem.2-5-5-4}
\expandafter\def\csname sga@target@IV.2.5.6\endcsname{ega4.c04.lem.2-5-6}
\expandafter\def\csname sga@target@IV.2.5.7\endcsname{ega4.c04.cor.2-5-7}
\expandafter\def\csname sga@target@IV.2.5.8\endcsname{ega4.c04.prop.2-5-8}
\expandafter\def\csname sga@target@IV.2.5.8.1\endcsname{ega4.c04.statement.2-5-8-1}
\expandafter\def\csname sga@target@IV.2.5.8.2\endcsname{ega4.c04.statement.2-5-8-2}
\expandafter\def\csname sga@target@IV.2.5.8.3\endcsname{ega4.c04.statement.2-5-8-3}
\expandafter\def\csname sga@target@IV.2.5.8.4\endcsname{ega4.c04.statement.2-5-8-4}
\expandafter\def\csname sga@target@IV.2.5.8.5\endcsname{ega4.c04.statement.2-5-8-5}
\expandafter\def\csname sga@target@IV.2.5.8.6\endcsname{ega4.c04.statement.2-5-8-6}
\expandafter\def\csname sga@target@IV.2.5.8.7\endcsname{ega4.c04.statement.2-5-8-7}
\expandafter\def\csname sga@target@IV.2.5.8.8\endcsname{ega4.c04.statement.2-5-8-8}
\expandafter\def\csname sga@target@IV.2.5.8.9\endcsname{ega4.c04.statement.2-5-8-9}
\expandafter\def\csname sga@target@IV.2.6.1\endcsname{ega4.c04.prop.2-6-1}
\expandafter\def\csname sga@target@IV.2.6.2\endcsname{ega4.c04.prop.2-6-2}
\expandafter\def\csname sga@target@IV.2.6.3\endcsname{ega4.c04.rem.2-6-3}
\expandafter\def\csname sga@target@IV.2.6.4\endcsname{ega4.c04.cor.2-6-4}
\expandafter\def\csname sga@target@IV.2.7.1\endcsname{ega4.c04.prop.2-7-1}
\expandafter\def\csname sga@target@IV.2.7.1.1\endcsname{ega4.c04.lem.2-7-1-1}
\expandafter\def\csname sga@target@IV.2.7.1.viii\endcsname{ega4.c04.item.2-7-1-viii}
\expandafter\def\csname sga@target@IV.2.7.2\endcsname{ega4.c04.cor.2-7-2}
\expandafter\def\csname sga@target@IV.2.7.3\endcsname{ega4.c04.sec.2-7-3}
\expandafter\def\csname sga@target@IV.2.8.1\endcsname{ega4.c04.prop.2-8-1}
\expandafter\def\csname sga@target@IV.2.8.1.1\endcsname{ega4.c04.lem.2-8-1-1}
\expandafter\def\csname sga@target@IV.2.8.2\endcsname{ega4.c04.cor.2-8-2}
\expandafter\def\csname sga@target@IV.2.8.3\endcsname{ega4.c04.statement.2-8-3}
\expandafter\def\csname sga@target@IV.2.8.4\endcsname{ega4.c04.prop.2-8-4}
\expandafter\def\csname sga@target@IV.2.8.4.1\endcsname{ega4.c04.statement.2-8-4-1}
\expandafter\def\csname sga@target@IV.2.8.5\endcsname{ega4.c04.prop.2-8-5}
\expandafter\def\csname sga@target@IV.2.8.6\endcsname{ega4.c04.cor.2-8-6}
\expandafter\def\csname sga@target@IV.2.8.6.1\endcsname{ega4.c04.statement.2-8-6-1}
\expandafter\def\csname sga@target@IV.20\endcsname{ega4.c04.sec.20}
\expandafter\def\csname sga@target@IV.20.0\endcsname{ega4.c04.subsec.20-0}
\expandafter\def\csname sga@target@IV.20.1\endcsname{ega4.c04.subsec.20-1}
\expandafter\def\csname sga@target@IV.20.1.1\endcsname{ega4.c04.statement.20-1-1}
\expandafter\def\csname sga@target@IV.20.1.1.1\endcsname{ega4.c04.eq.20-1-1-1}
\expandafter\def\csname sga@target@IV.20.1.10\endcsname{ega4.c04.statement.20-1-10}
\expandafter\def\csname sga@target@IV.20.1.11\endcsname{ega4.c04.statement.20-1-11}
\expandafter\def\csname sga@target@IV.20.1.11.1\endcsname{ega4.c04.eq.20-1-11-1}
\expandafter\def\csname sga@target@IV.20.1.11.2\endcsname{ega4.c04.eq.20-1-11-2}
\expandafter\def\csname sga@target@IV.20.1.12\endcsname{ega4.c04.prop.20-1-12}
\expandafter\def\csname sga@target@IV.20.1.13\endcsname{ega4.c04.rem.20-1-13}
\expandafter\def\csname sga@target@IV.20.1.2\endcsname{ega4.c04.statement.20-1-2}
\expandafter\def\csname sga@target@IV.20.1.2.1\endcsname{ega4.c04.statement.20-1-2-1}
\expandafter\def\csname sga@target@IV.20.1.2.2\endcsname{ega4.c04.statement.20-1-2-2}
\expandafter\def\csname sga@target@IV.20.1.3\endcsname{ega4.c04.statement.20-1-3}
\expandafter\def\csname sga@target@IV.20.1.3.1\endcsname{ega4.c04.statement.20-1-3-1}
\expandafter\def\csname sga@target@IV.20.1.4\endcsname{ega4.c04.statement.20-1-4}
\expandafter\def\csname sga@target@IV.20.1.4.1\endcsname{ega4.c04.statement.20-1-4-1}
\expandafter\def\csname sga@target@IV.20.1.5\endcsname{ega4.c04.statement.20-1-5}
\expandafter\def\csname sga@target@IV.20.1.5.1\endcsname{ega4.c04.statement.20-1-5-1}
\expandafter\def\csname sga@target@IV.20.1.6\endcsname{ega4.c04.prop.20-1-6}
\expandafter\def\csname sga@target@IV.20.1.7\endcsname{ega4.c04.statement.20-1-7}
\expandafter\def\csname sga@target@IV.20.1.8\endcsname{ega4.c04.statement.20-1-8}
\expandafter\def\csname sga@target@IV.20.1.9\endcsname{ega4.c04.statement.20-1-9}
\expandafter\def\csname sga@target@IV.20.2\endcsname{ega4.c04.subsec.20-2}
\expandafter\def\csname sga@target@IV.20.2.1\endcsname{ega4.c04.statement.20-2-1}
\expandafter\def\csname sga@target@IV.20.2.10\endcsname{ega4.c04.statement.20-2-10}
\expandafter\def\csname sga@target@IV.20.2.10.1\endcsname{ega4.c04.eq.20-2-10-1}
\expandafter\def\csname sga@target@IV.20.2.10.2\endcsname{ega4.c04.eq.20-2-10-2}
\expandafter\def\csname sga@target@IV.20.2.11\endcsname{ega4.c04.prop.20-2-11}
\expandafter\def\csname sga@target@IV.20.2.11.1\endcsname{ega4.c04.eq.20-2-11-1}
\expandafter\def\csname sga@target@IV.20.2.12\endcsname{ega4.c04.cor.20-2-12}
\expandafter\def\csname sga@target@IV.20.2.13\endcsname{ega4.c04.statement.20-2-13}
\expandafter\def\csname sga@target@IV.20.2.14\endcsname{ega4.c04.statement.20-2-14}
\expandafter\def\csname sga@target@IV.20.2.15\endcsname{ega4.c04.prop.20-2-15}
\expandafter\def\csname sga@target@IV.20.2.16\endcsname{ega4.c04.rem.20-2-16}
\expandafter\def\csname sga@target@IV.20.2.2\endcsname{ega4.c04.statement.20-2-2}
\expandafter\def\csname sga@target@IV.20.2.3\endcsname{ega4.c04.statement.20-2-3}
\expandafter\def\csname sga@target@IV.20.2.3.1\endcsname{ega4.c04.statement.20-2-3-1}
\expandafter\def\csname sga@target@IV.20.2.4\endcsname{ega4.c04.prop.20-2-4}
\expandafter\def\csname sga@target@IV.20.2.5\endcsname{ega4.c04.cor.20-2-5}
\expandafter\def\csname sga@target@IV.20.2.6\endcsname{ega4.c04.cor.20-2-6}
\expandafter\def\csname sga@target@IV.20.2.7\endcsname{ega4.c04.rem.20-2-7}
\expandafter\def\csname sga@target@IV.20.2.8\endcsname{ega4.c04.statement.20-2-8}
\expandafter\def\csname sga@target@IV.20.2.9\endcsname{ega4.c04.lem.20-2-9}
\expandafter\def\csname sga@target@IV.20.3\endcsname{ega4.c04.subsec.20-3}
\expandafter\def\csname sga@target@IV.20.3.1\endcsname{ega4.c04.statement.20-3-1}
\expandafter\def\csname sga@target@IV.20.3.10\endcsname{ega4.c04.lem.20-3-10}
\expandafter\def\csname sga@target@IV.20.3.11\endcsname{ega4.c04.prop.20-3-11}
\expandafter\def\csname sga@target@IV.20.3.11.1\endcsname{ega4.c04.eq.20-3-11-1}
\expandafter\def\csname sga@target@IV.20.3.12\endcsname{ega4.c04.cor.20-3-12}
\expandafter\def\csname sga@target@IV.20.3.13\endcsname{ega4.c04.rem.20-3-13}
\expandafter\def\csname sga@target@IV.20.3.2\endcsname{ega4.c04.statement.20-3-2}
\expandafter\def\csname sga@target@IV.20.3.2.1\endcsname{ega4.c04.eq.20-3-2-1}
\expandafter\def\csname sga@target@IV.20.3.3\endcsname{ega4.c04.statement.20-3-3}
\expandafter\def\csname sga@target@IV.20.3.4\endcsname{ega4.c04.statement.20-3-4}
\expandafter\def\csname sga@target@IV.20.3.5\endcsname{ega4.c04.prop.20-3-5}
\expandafter\def\csname sga@target@IV.20.3.6\endcsname{ega4.c04.statement.20-3-6}
\expandafter\def\csname sga@target@IV.20.3.6.1\endcsname{ega4.c04.eq.20-3-6-1}
\expandafter\def\csname sga@target@IV.20.3.7\endcsname{ega4.c04.prop.20-3-7}
\expandafter\def\csname sga@target@IV.20.3.8\endcsname{ega4.c04.prop.20-3-8}
\expandafter\def\csname sga@target@IV.20.3.8.1\endcsname{ega4.c04.eq.20-3-8-1}
\expandafter\def\csname sga@target@IV.20.3.8.2\endcsname{ega4.c04.lem.20-3-8-2}
\expandafter\def\csname sga@target@IV.20.3.8.3\endcsname{ega4.c04.lem.20-3-8-3}
\expandafter\def\csname sga@target@IV.20.3.8.4\endcsname{ega4.c04.lem.20-3-8-4}
\expandafter\def\csname sga@target@IV.20.3.8.5\endcsname{ega4.c04.sec.20-3-8-5}
\expandafter\def\csname sga@target@IV.20.3.8.6\endcsname{ega4.c04.sec.20-3-8-6}
\expandafter\def\csname sga@target@IV.20.3.8.7\endcsname{ega4.c04.sec.20-3-8-7}
\expandafter\def\csname sga@target@IV.20.3.8.8\endcsname{ega4.c04.lem.20-3-8-8}
\expandafter\def\csname sga@target@IV.20.3.9\endcsname{ega4.c04.cor.20-3-9}
\expandafter\def\csname sga@target@IV.20.4\endcsname{ega4.c04.subsec.20-4}
\expandafter\def\csname sga@target@IV.20.4.1\endcsname{ega4.c04.statement.20-4-1}
\expandafter\def\csname sga@target@IV.20.4.10\endcsname{ega4.c04.cor.20-4-10}
\expandafter\def\csname sga@target@IV.20.4.11\endcsname{ega4.c04.prop.20-4-11}
\expandafter\def\csname sga@target@IV.20.4.12\endcsname{ega4.c04.cor.20-4-12}
\expandafter\def\csname sga@target@IV.20.4.2\endcsname{ega4.c04.statement.20-4-2}
\expandafter\def\csname sga@target@IV.20.4.3\endcsname{ega4.c04.statement.20-4-3}
\expandafter\def\csname sga@target@IV.20.4.4\endcsname{ega4.c04.prop.20-4-4}
\expandafter\def\csname sga@target@IV.20.4.5\endcsname{ega4.c04.prop.20-4-5}
\expandafter\def\csname sga@target@IV.20.4.6\endcsname{ega4.c04.prop.20-4-6}
\expandafter\def\csname sga@target@IV.20.4.7\endcsname{ega4.c04.cor.20-4-7}
\expandafter\def\csname sga@target@IV.20.4.8\endcsname{ega4.c04.rem.20-4-8}
\expandafter\def\csname sga@target@IV.20.4.9\endcsname{ega4.c04.cor.20-4-9}
\expandafter\def\csname sga@target@IV.20.5\endcsname{ega4.c04.subsec.20-5}
\expandafter\def\csname sga@target@IV.20.5.1\endcsname{ega4.c04.statement.20-5-1}
\expandafter\def\csname sga@target@IV.20.5.2\endcsname{ega4.c04.statement.20-5-2}
\expandafter\def\csname sga@target@IV.20.5.2.1\endcsname{ega4.c04.statement.20-5-2-1}
\expandafter\def\csname sga@target@IV.20.5.3\endcsname{ega4.c04.prop.20-5-3}
\expandafter\def\csname sga@target@IV.20.5.4\endcsname{ega4.c04.statement.20-5-4}
\expandafter\def\csname sga@target@IV.20.5.5\endcsname{ega4.c04.statement.20-5-5}
\expandafter\def\csname sga@target@IV.20.5.6\endcsname{ega4.c04.statement.20-5-6}
\expandafter\def\csname sga@target@IV.20.5.7\endcsname{ega4.c04.statement.20-5-7}
\expandafter\def\csname sga@target@IV.20.5.8\endcsname{ega4.c04.statement.20-5-8}
\expandafter\def\csname sga@target@IV.20.6\endcsname{ega4.c04.subsec.20-6}
\expandafter\def\csname sga@target@IV.20.6.1\endcsname{ega4.c04.statement.20-6-1}
\expandafter\def\csname sga@target@IV.20.6.1.1\endcsname{ega4.c04.eq.20-6-1-1}
\expandafter\def\csname sga@target@IV.20.6.2\endcsname{ega4.c04.statement.20-6-2}
\expandafter\def\csname sga@target@IV.20.6.2.1\endcsname{ega4.c04.statement.20-6-2-1}
\expandafter\def\csname sga@target@IV.20.6.2.2\endcsname{ega4.c04.eq.20-6-2-2}
\expandafter\def\csname sga@target@IV.20.6.3\endcsname{ega4.c04.prop.20-6-3}
\expandafter\def\csname sga@target@IV.20.6.4\endcsname{ega4.c04.prop.20-6-4}
\expandafter\def\csname sga@target@IV.20.6.5\endcsname{ega4.c04.statement.20-6-5}
\expandafter\def\csname sga@target@IV.20.6.5.1\endcsname{ega4.c04.eq.20-6-5-1}
\expandafter\def\csname sga@target@IV.20.6.6\endcsname{ega4.c04.prop.20-6-6}
\expandafter\def\csname sga@target@IV.20.6.7\endcsname{ega4.c04.statement.20-6-7}
\expandafter\def\csname sga@target@IV.20.6.8\endcsname{ega4.c04.prop.20-6-8}
\expandafter\def\csname sga@target@IV.20.6.8.1\endcsname{ega4.c04.eq.20-6-8-1}
\expandafter\def\csname sga@target@IV.20.6.9\endcsname{ega4.c04.statement.20-6-9}
\expandafter\def\csname sga@target@IV.20.6.9.1\endcsname{ega4.c04.statement.20-6-9-1}
\expandafter\def\csname sga@target@IV.21.1\endcsname{ega4.c04.subsec.21-1.key-iv-21-1}
\expandafter\def\csname sga@target@IV.21.1.1\endcsname{ega4.c04.statement.21-1-1}
\expandafter\def\csname sga@target@IV.21.1.10\endcsname{ega4.c04.prop.21-1-10}
\expandafter\def\csname sga@target@IV.21.1.2\endcsname{ega4.c04.statement.21-1-2}
\expandafter\def\csname sga@target@IV.21.1.3\endcsname{ega4.c04.statement.21-1-3}
\expandafter\def\csname sga@target@IV.21.1.3.1\endcsname{ega4.c04.statement.21-1-3-1}
\expandafter\def\csname sga@target@IV.21.1.3.2\endcsname{ega4.c04.statement.21-1-3-2}
\expandafter\def\csname sga@target@IV.21.1.3.3\endcsname{ega4.c04.statement.21-1-3-3}
\expandafter\def\csname sga@target@IV.21.1.3.4\endcsname{ega4.c04.statement.21-1-3-4}
\expandafter\def\csname sga@target@IV.21.1.4\endcsname{ega4.c04.statement.21-1-4}
\expandafter\def\csname sga@target@IV.21.1.4.1\endcsname{ega4.c04.eq.21-1-4-1}
\expandafter\def\csname sga@target@IV.21.1.4.2\endcsname{ega4.c04.eq.21-1-4-2}
\expandafter\def\csname sga@target@IV.21.1.4.3\endcsname{ega4.c04.eq.21-1-4-3}
\expandafter\def\csname sga@target@IV.21.1.5\endcsname{ega4.c04.statement.21-1-5}
\expandafter\def\csname sga@target@IV.21.1.5.1\endcsname{ega4.c04.statement.21-1-5-1}
\expandafter\def\csname sga@target@IV.21.1.5.2\endcsname{ega4.c04.eq.21-1-5-2}
\expandafter\def\csname sga@target@IV.21.1.6\endcsname{ega4.c04.def.21-1-6}
\expandafter\def\csname sga@target@IV.21.1.6.1\endcsname{ega4.c04.eq.21-1-6-1}
\expandafter\def\csname sga@target@IV.21.1.6.2\endcsname{ega4.c04.eq.21-1-6-2}
\expandafter\def\csname sga@target@IV.21.1.6.3\endcsname{ega4.c04.eq.21-1-6-3}
\expandafter\def\csname sga@target@IV.21.1.6.4\endcsname{ega4.c04.eq.21-1-6-4}
\expandafter\def\csname sga@target@IV.21.1.7\endcsname{ega4.c04.statement.21-1-7}
\expandafter\def\csname sga@target@IV.21.1.7.1\endcsname{ega4.c04.statement.21-1-7-1}
\expandafter\def\csname sga@target@IV.21.1.8\endcsname{ega4.c04.prop.21-1-8}
\expandafter\def\csname sga@target@IV.21.1.9\endcsname{ega4.c04.cor.21-1-9}
\expandafter\def\csname sga@target@IV.21.10.1\endcsname{ega4.c04.statement.21-10-1}
\expandafter\def\csname sga@target@IV.21.10.1.1\endcsname{ega4.c04.eq.21-10-1-1}
\expandafter\def\csname sga@target@IV.21.10.1.2\endcsname{ega4.c04.eq.21-10-1-2}
\expandafter\def\csname sga@target@IV.21.10.10\endcsname{ega4.c04.statement.21-10-10}
\expandafter\def\csname sga@target@IV.21.10.11\endcsname{ega4.c04.statement.21-10-11}
\expandafter\def\csname sga@target@IV.21.10.11.1\endcsname{ega4.c04.eq.21-10-11-1}
\expandafter\def\csname sga@target@IV.21.10.12\endcsname{ega4.c04.rem.21-10-12}
\expandafter\def\csname sga@target@IV.21.10.13\endcsname{ega4.c04.lem.21-10-13}
\expandafter\def\csname sga@target@IV.21.10.13.1\endcsname{ega4.c04.prop.21-10-13-1}
\expandafter\def\csname sga@target@IV.21.10.13bis\endcsname{ega4.c04.prop.21-10-13bis}
\expandafter\def\csname sga@target@IV.21.10.14\endcsname{ega4.c04.statement.21-10-14}
\expandafter\def\csname sga@target@IV.21.10.14.1\endcsname{ega4.c04.eq.21-10-14-1}
\expandafter\def\csname sga@target@IV.21.10.14.2\endcsname{ega4.c04.eq.21-10-14-2}
\expandafter\def\csname sga@target@IV.21.10.15\endcsname{ega4.c04.prop.21-10-15}
\expandafter\def\csname sga@target@IV.21.10.16\endcsname{ega4.c04.prop.21-10-16}
\expandafter\def\csname sga@target@IV.21.10.16.1\endcsname{ega4.c04.prop.21-10-16-1}
\expandafter\def\csname sga@target@IV.21.10.16.2\endcsname{ega4.c04.eq.21-10-16-2}
\expandafter\def\csname sga@target@IV.21.10.17\endcsname{ega4.c04.prop.21-10-17}
\expandafter\def\csname sga@target@IV.21.10.17.1\endcsname{ega4.c04.prop.21-10-17-1}
\expandafter\def\csname sga@target@IV.21.10.17.10\endcsname{ega4.c04.eq.21-10-17-10}
\expandafter\def\csname sga@target@IV.21.10.17.2\endcsname{ega4.c04.eq.21-10-17-2}
\expandafter\def\csname sga@target@IV.21.10.17.3\endcsname{ega4.c04.lem.21-10-17-3}
\expandafter\def\csname sga@target@IV.21.10.17.4\endcsname{ega4.c04.lem.21-10-17-4}
\expandafter\def\csname sga@target@IV.21.10.17.5\endcsname{ega4.c04.eq.21-10-17-5}
\expandafter\def\csname sga@target@IV.21.10.17.6\endcsname{ega4.c04.eq.21-10-17-6}
\expandafter\def\csname sga@target@IV.21.10.17.7\endcsname{ega4.c04.lem.21-10-17-7}
\expandafter\def\csname sga@target@IV.21.10.17.8\endcsname{ega4.c04.lem.21-10-17-8}
\expandafter\def\csname sga@target@IV.21.10.17.9\endcsname{ega4.c04.eq.21-10-17-9}
\expandafter\def\csname sga@target@IV.21.10.18\endcsname{ega4.c04.prop.21-10-18}
\expandafter\def\csname sga@target@IV.21.10.18.1\endcsname{ega4.c04.prop.21-10-18-1}
\expandafter\def\csname sga@target@IV.21.10.19\endcsname{ega4.c04.cor.21-10-19}
\expandafter\def\csname sga@target@IV.21.10.19.1\endcsname{ega4.c04.cor.21-10-19-1}
\expandafter\def\csname sga@target@IV.21.10.2\endcsname{ega4.c04.statement.21-10-2}
\expandafter\def\csname sga@target@IV.21.10.3\endcsname{ega4.c04.statement.21-10-3}
\expandafter\def\csname sga@target@IV.21.10.3.1\endcsname{ega4.c04.statement.21-10-3-1}
\expandafter\def\csname sga@target@IV.21.10.4\endcsname{ega4.c04.prop.21-10-4}
\expandafter\def\csname sga@target@IV.21.10.4.1\endcsname{ega4.c04.prop.21-10-4-1}
\expandafter\def\csname sga@target@IV.21.10.5\endcsname{ega4.c04.statement.21-10-5}
\expandafter\def\csname sga@target@IV.21.10.6\endcsname{ega4.c04.prop.21-10-6}
\expandafter\def\csname sga@target@IV.21.10.7\endcsname{ega4.c04.rem.21-10-7}
\expandafter\def\csname sga@target@IV.21.10.8\endcsname{ega4.c04.prop.21-10-8}
\expandafter\def\csname sga@target@IV.21.10.9\endcsname{ega4.c04.statement.21-10-9}
\expandafter\def\csname sga@target@IV.21.11\endcsname{ega4.c04.subsec.21-11.key-iv-21-11}
\expandafter\def\csname sga@target@IV.21.11.1\endcsname{ega4.c04.thm.21-11-1}
\expandafter\def\csname sga@target@IV.21.11.1.1\endcsname{ega4.c04.eq.21-11-1-1}
\expandafter\def\csname sga@target@IV.21.11.2\endcsname{ega4.c04.lem.21-11-2}
\expandafter\def\csname sga@target@IV.21.12\endcsname{ega4.c04.subsec.21-12.key-iv-21-12}
\expandafter\def\csname sga@target@IV.21.12.1\endcsname{ega4.c04.statement.21-12-1}
\expandafter\def\csname sga@target@IV.21.12.1.1\endcsname{ega4.c04.eq.21-12-1-1}
\expandafter\def\csname sga@target@IV.21.12.10\endcsname{ega4.c04.prop.21-12-10}
\expandafter\def\csname sga@target@IV.21.12.12\endcsname{ega4.c04.thm.21-12-12}
\expandafter\def\csname sga@target@IV.21.12.13\endcsname{ega4.c04.cor.21-12-13}
\expandafter\def\csname sga@target@IV.21.12.14\endcsname{ega4.c04.item.21-12-14}
\expandafter\def\csname sga@target@IV.21.12.15\endcsname{ega4.c04.prop.21-12-15}
\expandafter\def\csname sga@target@IV.21.12.16\endcsname{ega4.c04.cor.21-12-16}
\expandafter\def\csname sga@target@IV.21.12.17\endcsname{ega4.c04.item.21-12-17}
\expandafter\def\csname sga@target@IV.21.12.2\endcsname{ega4.c04.prop.21-12-2}
\expandafter\def\csname sga@target@IV.21.12.2.1\endcsname{ega4.c04.prop.21-12-2-1}
\expandafter\def\csname sga@target@IV.21.12.3\endcsname{ega4.c04.cor.21-12-3}
\expandafter\def\csname sga@target@IV.21.12.4\endcsname{ega4.c04.eq.21-12-4}
\expandafter\def\csname sga@target@IV.21.12.5\endcsname{ega4.c04.statement.21-12-5}
\expandafter\def\csname sga@target@IV.21.12.5.1\endcsname{ega4.c04.eq.21-12-5-1}
\expandafter\def\csname sga@target@IV.21.12.5.2\endcsname{ega4.c04.eq.21-12-5-2}
\expandafter\def\csname sga@target@IV.21.12.5.3\endcsname{ega4.c04.eq.21-12-5-3}
\expandafter\def\csname sga@target@IV.21.12.6\endcsname{ega4.c04.thm.21-12-6}
\expandafter\def\csname sga@target@IV.21.12.7\endcsname{ega4.c04.cor.21-12-7}
\expandafter\def\csname sga@target@IV.21.12.8\endcsname{ega4.c04.statement.21-12-8}
\expandafter\def\csname sga@target@IV.21.12.9\endcsname{ega4.c04.statement.21-12-9}
\expandafter\def\csname sga@target@IV.21.13\endcsname{ega4.c04.subsec.21-13.key-iv-21-13}
\expandafter\def\csname sga@target@IV.21.13.1\endcsname{ega4.c04.def.21-13-1}
\expandafter\def\csname sga@target@IV.21.13.10\endcsname{ega4.c04.prop.21-13-10}
\expandafter\def\csname sga@target@IV.21.13.11\endcsname{ega4.c04.cor.21-13-11}
\expandafter\def\csname sga@target@IV.21.13.12\endcsname{ega4.c04.prop.21-13-12}
\expandafter\def\csname sga@target@IV.21.13.12.1\endcsname{ega4.c04.statement.21-13-12-1}
\expandafter\def\csname sga@target@IV.21.13.13\endcsname{ega4.c04.rem.21-13-13}
\expandafter\def\csname sga@target@IV.21.13.14\endcsname{ega4.c04.prop.21-13-14}
\expandafter\def\csname sga@target@IV.21.13.15\endcsname{ega4.c04.cor.21-13-15}
\expandafter\def\csname sga@target@IV.21.13.16\endcsname{ega4.c04.rem.21-13-16}
\expandafter\def\csname sga@target@IV.21.13.2\endcsname{ega4.c04.lem.21-13-2}
\expandafter\def\csname sga@target@IV.21.13.2.1\endcsname{ega4.c04.eq.21-13-2-1}
\expandafter\def\csname sga@target@IV.21.13.3\endcsname{ega4.c04.cor.21-13-3}
\expandafter\def\csname sga@target@IV.21.13.4\endcsname{ega4.c04.lem.21-13-4}
\expandafter\def\csname sga@target@IV.21.13.5\endcsname{ega4.c04.prop.21-13-5}
\expandafter\def\csname sga@target@IV.21.13.6\endcsname{ega4.c04.cor.21-13-6}
\expandafter\def\csname sga@target@IV.21.13.7\endcsname{ega4.c04.def.21-13-7}
\expandafter\def\csname sga@target@IV.21.13.8\endcsname{ega4.c04.prop.21-13-8}
\expandafter\def\csname sga@target@IV.21.13.9\endcsname{ega4.c04.statement.21-13-9}
\expandafter\def\csname sga@target@IV.21.13.9.1\endcsname{ega4.c04.eq.21-13-9-1}
\expandafter\def\csname sga@target@IV.21.14\endcsname{ega4.c04.subsec.21-14.key-iv-21-14}
\expandafter\def\csname sga@target@IV.21.14.1\endcsname{ega4.c04.thm.21-14-1}
\expandafter\def\csname sga@target@IV.21.14.2\endcsname{ega4.c04.cor.21-14-2}
\expandafter\def\csname sga@target@IV.21.14.3\endcsname{ega4.c04.prop.21-14-3}
\expandafter\def\csname sga@target@IV.21.14.4\endcsname{ega4.c04.statement.21-14-4}
\expandafter\def\csname sga@target@IV.21.14.4bis\endcsname{ega4.c04.statement.21-14-4bis}
\expandafter\def\csname sga@target@IV.21.15\endcsname{ega4.c04.subsec.21-15.key-iv-21-15}
\expandafter\def\csname sga@target@IV.21.15.1\endcsname{ega4.c04.statement.21-15-1}
\expandafter\def\csname sga@target@IV.21.15.2\endcsname{ega4.c04.def.21-15-2}
\expandafter\def\csname sga@target@IV.21.15.3\endcsname{ega4.c04.statement.21-15-3}
\expandafter\def\csname sga@target@IV.21.15.3.1\endcsname{ega4.c04.eq.21-15-3-1}
\expandafter\def\csname sga@target@IV.21.15.3.2\endcsname{ega4.c04.eq.21-15-3-2}
\expandafter\def\csname sga@target@IV.21.15.3bis\endcsname{ega4.c04.statement.21-15-3bis}
\expandafter\def\csname sga@target@IV.21.15.4\endcsname{ega4.c04.prop.21-15-4}
\expandafter\def\csname sga@target@IV.21.15.5\endcsname{ega4.c04.prop.21-15-5}
\expandafter\def\csname sga@target@IV.21.15.6\endcsname{ega4.c04.prop.21-15-6}
\expandafter\def\csname sga@target@IV.21.15.7\endcsname{ega4.c04.statement.21-15-7}
\expandafter\def\csname sga@target@IV.21.15.8\endcsname{ega4.c04.statement.21-15-8}
\expandafter\def\csname sga@target@IV.21.15.9\endcsname{ega4.c04.statement.21-15-9}
\expandafter\def\csname sga@target@IV.21.15.9.1\endcsname{ega4.c04.eq.21-15-9-1}
\expandafter\def\csname sga@target@IV.21.2.1\endcsname{ega4.c04.statement.21-2-1}
\expandafter\def\csname sga@target@IV.21.2.10\endcsname{ega4.c04.statement.21-2-10}
\expandafter\def\csname sga@target@IV.21.2.10.1\endcsname{ega4.c04.eq.21-2-10-1}
\expandafter\def\csname sga@target@IV.21.2.11\endcsname{ega4.c04.prop.21-2-11}
\expandafter\def\csname sga@target@IV.21.2.11.1\endcsname{ega4.c04.prop.21-2-11-1}
\expandafter\def\csname sga@target@IV.21.2.12\endcsname{ega4.c04.statement.21-2-12}
\expandafter\def\csname sga@target@IV.21.2.2\endcsname{ega4.c04.prop.21-2-2}
\expandafter\def\csname sga@target@IV.21.2.3\endcsname{ega4.c04.cor.21-2-3}
\expandafter\def\csname sga@target@IV.21.2.4\endcsname{ega4.c04.statement.21-2-4}
\expandafter\def\csname sga@target@IV.21.2.5\endcsname{ega4.c04.statement.21-2-5}
\expandafter\def\csname sga@target@IV.21.2.5.1\endcsname{ega4.c04.statement.21-2-5-1}
\expandafter\def\csname sga@target@IV.21.2.5.2\endcsname{ega4.c04.eq.21-2-5-2}
\expandafter\def\csname sga@target@IV.21.2.5.3\endcsname{ega4.c04.eq.21-2-5-3}
\expandafter\def\csname sga@target@IV.21.2.5.4\endcsname{ega4.c04.eq.21-2-5-4}
\expandafter\def\csname sga@target@IV.21.2.6\endcsname{ega4.c04.prop.21-2-6}
\expandafter\def\csname sga@target@IV.21.2.6.1\endcsname{ega4.c04.eq.21-2-6-1}
\expandafter\def\csname sga@target@IV.21.2.7\endcsname{ega4.c04.statement.21-2-7}
\expandafter\def\csname sga@target@IV.21.2.7.1\endcsname{ega4.c04.statement.21-2-7-1}
\expandafter\def\csname sga@target@IV.21.2.7.2\endcsname{ega4.c04.eq.21-2-7-2}
\expandafter\def\csname sga@target@IV.21.2.8\endcsname{ega4.c04.statement.21-2-8}
\expandafter\def\csname sga@target@IV.21.2.8.1\endcsname{ega4.c04.statement.21-2-8-1}
\expandafter\def\csname sga@target@IV.21.2.8.2\endcsname{ega4.c04.statement.21-2-8-2}
\expandafter\def\csname sga@target@IV.21.2.8.3\endcsname{ega4.c04.statement.21-2-8-3}
\expandafter\def\csname sga@target@IV.21.2.9\endcsname{ega4.c04.statement.21-2-9}
\expandafter\def\csname sga@target@IV.21.2.9.1\endcsname{ega4.c04.statement.21-2-9-1}
\expandafter\def\csname sga@target@IV.21.2.9.2\endcsname{ega4.c04.eq.21-2-9-2}
\expandafter\def\csname sga@target@IV.21.2.9.3\endcsname{ega4.c04.eq.21-2-9-3}
\expandafter\def\csname sga@target@IV.21.2.9.4\endcsname{ega4.c04.eq.21-2-9-4}
\expandafter\def\csname sga@target@IV.21.3.1\endcsname{ega4.c04.statement.21-3-1}
\expandafter\def\csname sga@target@IV.21.3.2\endcsname{ega4.c04.statement.21-3-2}
\expandafter\def\csname sga@target@IV.21.3.2.1\endcsname{ega4.c04.statement.21-3-2-1}
\expandafter\def\csname sga@target@IV.21.3.2.2\endcsname{ega4.c04.eq.21-3-2-2}
\expandafter\def\csname sga@target@IV.21.3.2.3\endcsname{ega4.c04.eq.21-3-2-3}
\expandafter\def\csname sga@target@IV.21.3.2.4\endcsname{ega4.c04.eq.21-3-2-4}
\expandafter\def\csname sga@target@IV.21.3.3\endcsname{ega4.c04.prop.21-3-3}
\expandafter\def\csname sga@target@IV.21.3.3.1\endcsname{ega4.c04.eq.21-3-3-1}
\expandafter\def\csname sga@target@IV.21.3.4\endcsname{ega4.c04.prop.21-3-4}
\expandafter\def\csname sga@target@IV.21.3.5\endcsname{ega4.c04.cor.21-3-5}
\expandafter\def\csname sga@target@IV.21.3.6\endcsname{ega4.c04.rem.21-3-6}
\expandafter\def\csname sga@target@IV.21.3.6.1\endcsname{ega4.c04.eq.21-3-6-1}
\expandafter\def\csname sga@target@IV.21.4\endcsname{ega4.c04.subsec.21-4.key-iv-21-4}
\expandafter\def\csname sga@target@IV.21.4.1\endcsname{ega4.c04.statement.21-4-1}
\expandafter\def\csname sga@target@IV.21.4.10\endcsname{ega4.c04.cor.21-4-10}
\expandafter\def\csname sga@target@IV.21.4.11\endcsname{ega4.c04.cor.21-4-11}
\expandafter\def\csname sga@target@IV.21.4.12\endcsname{ega4.c04.rem.21-4-12}
\expandafter\def\csname sga@target@IV.21.4.13\endcsname{ega4.c04.cor.21-4-13}
\expandafter\def\csname sga@target@IV.21.4.2\endcsname{ega4.c04.def.21-4-2}
\expandafter\def\csname sga@target@IV.21.4.2.1\endcsname{ega4.c04.eq.21-4-2-1}
\expandafter\def\csname sga@target@IV.21.4.3\endcsname{ega4.c04.statement.21-4-3}
\expandafter\def\csname sga@target@IV.21.4.4\endcsname{ega4.c04.statement.21-4-4}
\expandafter\def\csname sga@target@IV.21.4.5\endcsname{ega4.c04.prop.21-4-5}
\expandafter\def\csname sga@target@IV.21.4.6\endcsname{ega4.c04.cor.21-4-6}
\expandafter\def\csname sga@target@IV.21.4.6.1\endcsname{ega4.c04.eq.21-4-6-1}
\expandafter\def\csname sga@target@IV.21.4.7\endcsname{ega4.c04.statement.21-4-7}
\expandafter\def\csname sga@target@IV.21.4.8\endcsname{ega4.c04.prop.21-4-8}
\expandafter\def\csname sga@target@IV.21.4.9\endcsname{ega4.c04.prop.21-4-9}
\expandafter\def\csname sga@target@IV.21.5\endcsname{ega4.c04.subsec.21-5.key-iv-21-5}
\expandafter\def\csname sga@target@IV.21.5.1\endcsname{ega4.c04.statement.21-5-1}
\expandafter\def\csname sga@target@IV.21.5.2\endcsname{ega4.c04.lem.21-5-2}
\expandafter\def\csname sga@target@IV.21.5.3\endcsname{ega4.c04.statement.21-5-3}
\expandafter\def\csname sga@target@IV.21.5.3.1\endcsname{ega4.c04.eq.21-5-3-1}
\expandafter\def\csname sga@target@IV.21.5.3.2\endcsname{ega4.c04.eq.21-5-3-2}
\expandafter\def\csname sga@target@IV.21.5.4\endcsname{ega4.c04.statement.21-5-4}
\expandafter\def\csname sga@target@IV.21.5.5\endcsname{ega4.c04.def.21-5-5}
\expandafter\def\csname sga@target@IV.21.5.5.1\endcsname{ega4.c04.eq.21-5-5-1}
\expandafter\def\csname sga@target@IV.21.5.5.2\endcsname{ega4.c04.eq.21-5-5-2}
\expandafter\def\csname sga@target@IV.21.5.5.3\endcsname{ega4.c04.eq.21-5-5-3}
\expandafter\def\csname sga@target@IV.21.5.6\endcsname{ega4.c04.prop.21-5-6}
\expandafter\def\csname sga@target@IV.21.5.6.1\endcsname{ega4.c04.prop.21-5-6-1}
\expandafter\def\csname sga@target@IV.21.5.7\endcsname{ega4.c04.prop.21-5-7}
\expandafter\def\csname sga@target@IV.21.5.7.1\endcsname{ega4.c04.prop.21-5-7-1}
\expandafter\def\csname sga@target@IV.21.5.7.2\endcsname{ega4.c04.lem.21-5-7-2}
\expandafter\def\csname sga@target@IV.21.5.7.3\endcsname{ega4.c04.lem.21-5-7-3}
\expandafter\def\csname sga@target@IV.21.5.7.4\endcsname{ega4.c04.eq.21-5-7-4}
\expandafter\def\csname sga@target@IV.21.5.8\endcsname{ega4.c04.prop.21-5-8}
\expandafter\def\csname sga@target@IV.21.5.8.1\endcsname{ega4.c04.eq.21-5-8-1}
\expandafter\def\csname sga@target@IV.21.6\endcsname{ega4.c04.subsec.21-6.key-iv-21-6}
\expandafter\def\csname sga@target@IV.21.6.1\endcsname{ega4.c04.statement.21-6-1}
\expandafter\def\csname sga@target@IV.21.6.10\endcsname{ega4.c04.cor.21-6-10}
\expandafter\def\csname sga@target@IV.21.6.10.1\endcsname{ega4.c04.eq.21-6-10-1}
\expandafter\def\csname sga@target@IV.21.6.11\endcsname{ega4.c04.cor.21-6-11}
\expandafter\def\csname sga@target@IV.21.6.12\endcsname{ega4.c04.prop.21-6-12}
\expandafter\def\csname sga@target@IV.21.6.12.1\endcsname{ega4.c04.eq.21-6-12-1}
\expandafter\def\csname sga@target@IV.21.6.12.2\endcsname{ega4.c04.eq.21-6-12-2}
\expandafter\def\csname sga@target@IV.21.6.13\endcsname{ega4.c04.cor.21-6-13}
\expandafter\def\csname sga@target@IV.21.6.14\endcsname{ega4.c04.cor.21-6-14}
\expandafter\def\csname sga@target@IV.21.6.2\endcsname{ega4.c04.statement.21-6-2}
\expandafter\def\csname sga@target@IV.21.6.3\endcsname{ega4.c04.statement.21-6-3}
\expandafter\def\csname sga@target@IV.21.6.4\endcsname{ega4.c04.statement.21-6-4}
\expandafter\def\csname sga@target@IV.21.6.4.1\endcsname{ega4.c04.statement.21-6-4-1}
\expandafter\def\csname sga@target@IV.21.6.4.2\endcsname{ega4.c04.eq.21-6-4-2}
\expandafter\def\csname sga@target@IV.21.6.5\endcsname{ega4.c04.statement.21-6-5}
\expandafter\def\csname sga@target@IV.21.6.5.1\endcsname{ega4.c04.eq.21-6-5-1}
\expandafter\def\csname sga@target@IV.21.6.6\endcsname{ega4.c04.prop.21-6-6}
\expandafter\def\csname sga@target@IV.21.6.6.2\endcsname{ega4.c04.statement.21-6-6-2}
\expandafter\def\csname sga@target@IV.21.6.7\endcsname{ega4.c04.statement.21-6-7}
\expandafter\def\csname sga@target@IV.21.6.7.1\endcsname{ega4.c04.eq.21-6-7-1}
\expandafter\def\csname sga@target@IV.21.6.7.2\endcsname{ega4.c04.eq.21-6-7-2}
\expandafter\def\csname sga@target@IV.21.6.7.3\endcsname{ega4.c04.eq.21-6-7-3}
\expandafter\def\csname sga@target@IV.21.6.7.4\endcsname{ega4.c04.eq.21-6-7-4}
\expandafter\def\csname sga@target@IV.21.6.7.5\endcsname{ega4.c04.eq.21-6-7-5}
\expandafter\def\csname sga@target@IV.21.6.8\endcsname{ega4.c04.statement.21-6-8}
\expandafter\def\csname sga@target@IV.21.6.9\endcsname{ega4.c04.thm.21-6-9}
\expandafter\def\csname sga@target@IV.21.6.9.1\endcsname{ega4.c04.eq.21-6-9-1}
\expandafter\def\csname sga@target@IV.21.6.9.2\endcsname{ega4.c04.statement.21-6-9-2}
\expandafter\def\csname sga@target@IV.21.7\endcsname{ega4.c04.subsec.21-7.key-iv-21-7}
\expandafter\def\csname sga@target@IV.21.7.1\endcsname{ega4.c04.statement.21-7-1}
\expandafter\def\csname sga@target@IV.21.7.2\endcsname{ega4.c04.prop.21-7-2}
\expandafter\def\csname sga@target@IV.21.7.2.1\endcsname{ega4.c04.eq.21-7-2-1}
\expandafter\def\csname sga@target@IV.21.7.3\endcsname{ega4.c04.cor.21-7-3}
\expandafter\def\csname sga@target@IV.21.7.3.1\endcsname{ega4.c04.statement.21-7-3-1}
\expandafter\def\csname sga@target@IV.21.7.4\endcsname{ega4.c04.prop.21-7-4}
\expandafter\def\csname sga@target@IV.21.7.5\endcsname{ega4.c04.rem.21-7-5}
\expandafter\def\csname sga@target@IV.21.7.6\endcsname{ega4.c04.cor.21-7-6}
\expandafter\def\csname sga@target@IV.21.7.7\endcsname{ega4.c04.rem.21-7-7}
\expandafter\def\csname sga@target@IV.21.8.1\endcsname{ega4.c04.lem.21-8-1}
\expandafter\def\csname sga@target@IV.21.8.2\endcsname{ega4.c04.cor.21-8-2}
\expandafter\def\csname sga@target@IV.21.8.3\endcsname{ega4.c04.prop.21-8-3}
\expandafter\def\csname sga@target@IV.21.8.3.1\endcsname{ega4.c04.eq.21-8-3-1}
\expandafter\def\csname sga@target@IV.21.8.4\endcsname{ega4.c04.cor.21-8-4}
\expandafter\def\csname sga@target@IV.21.8.5\endcsname{ega4.c04.prop.21-8-5}
\expandafter\def\csname sga@target@IV.21.8.5.1\endcsname{ega4.c04.eq.21-8-5-1}
\expandafter\def\csname sga@target@IV.21.8.5.2\endcsname{ega4.c04.eq.21-8-5-2}
\expandafter\def\csname sga@target@IV.21.8.5.3\endcsname{ega4.c04.eq.21-8-5-3}
\expandafter\def\csname sga@target@IV.21.8.6\endcsname{ega4.c04.rem.21-8-6}
\expandafter\def\csname sga@target@IV.21.9\endcsname{ega4.c04.subsec.21-9.key-iv-21-9}
\expandafter\def\csname sga@target@IV.21.9.1\endcsname{ega4.c04.statement.21-9-1}
\expandafter\def\csname sga@target@IV.21.9.1.1\endcsname{ega4.c04.eq.21-9-1-1}
\expandafter\def\csname sga@target@IV.21.9.1.2\endcsname{ega4.c04.eq.21-9-1-2}
\expandafter\def\csname sga@target@IV.21.9.10\endcsname{ega4.c04.statement.21-9-10}
\expandafter\def\csname sga@target@IV.21.9.11\endcsname{ega4.c04.prop.21-9-11}
\expandafter\def\csname sga@target@IV.21.9.12\endcsname{ega4.c04.cor.21-9-12}
\expandafter\def\csname sga@target@IV.21.9.13\endcsname{ega4.c04.rem.21-9-13}
\expandafter\def\csname sga@target@IV.21.9.2\endcsname{ega4.c04.lem.21-9-2}
\expandafter\def\csname sga@target@IV.21.9.3\endcsname{ega4.c04.rem.21-9-3}
\expandafter\def\csname sga@target@IV.21.9.4\endcsname{ega4.c04.prop.21-9-4}
\expandafter\def\csname sga@target@IV.21.9.4.1\endcsname{ega4.c04.prop.21-9-4-1}
\expandafter\def\csname sga@target@IV.21.9.5\endcsname{ega4.c04.cor.21-9-5}
\expandafter\def\csname sga@target@IV.21.9.6\endcsname{ega4.c04.cor.21-9-6}
\expandafter\def\csname sga@target@IV.21.9.7\endcsname{ega4.c04.cor.21-9-7}
\expandafter\def\csname sga@target@IV.21.9.8\endcsname{ega4.c04.statement.21-9-8}
\expandafter\def\csname sga@target@IV.21.9.9\endcsname{ega4.c04.prop.21-9-9}
\expandafter\def\csname sga@target@IV.3.1.1\endcsname{ega4.c04.def.3-1-1}
\expandafter\def\csname sga@target@IV.3.1.1.1\endcsname{ega4.c04.eq.3-1-1-1}
\expandafter\def\csname sga@target@IV.3.1.10\endcsname{ega4.c04.prop.3-1-10}
\expandafter\def\csname sga@target@IV.3.1.10.1\endcsname{ega4.c04.lem.3-1-10-1}
\expandafter\def\csname sga@target@IV.3.1.10.2\endcsname{ega4.c04.rem.3-1-10-2}
\expandafter\def\csname sga@target@IV.3.1.10.3\endcsname{ega4.c04.eq.3-1-10-3}
\expandafter\def\csname sga@target@IV.3.1.10.4\endcsname{ega4.c04.eq.3-1-10-4}
\expandafter\def\csname sga@target@IV.3.1.10.5\endcsname{ega4.c04.rem.3-1-10-5}
\expandafter\def\csname sga@target@IV.3.1.11\endcsname{ega4.c04.cor.3-1-11}
\expandafter\def\csname sga@target@IV.3.1.12\endcsname{ega4.c04.rem.3-1-12}
\expandafter\def\csname sga@target@IV.3.1.13\endcsname{ega4.c04.prop.3-1-13}
\expandafter\def\csname sga@target@IV.3.1.2\endcsname{ega4.c04.prop.3-1-2}
\expandafter\def\csname sga@target@IV.3.1.3\endcsname{ega4.c04.prop.3-1-3}
\expandafter\def\csname sga@target@IV.3.1.4\endcsname{ega4.c04.cor.3-1-4}
\expandafter\def\csname sga@target@IV.3.1.5\endcsname{ega4.c04.cor.3-1-5}
\expandafter\def\csname sga@target@IV.3.1.6\endcsname{ega4.c04.prop.3-1-6}
\expandafter\def\csname sga@target@IV.3.1.7\endcsname{ega4.c04.prop.3-1-7}
\expandafter\def\csname sga@target@IV.3.1.8\endcsname{ega4.c04.prop.3-1-8}
\expandafter\def\csname sga@target@IV.3.1.9\endcsname{ega4.c04.cor.3-1-9}
\expandafter\def\csname sga@target@IV.3.2.1\endcsname{ega4.c04.prop.3-2-1}
\expandafter\def\csname sga@target@IV.3.2.2\endcsname{ega4.c04.def.3-2-2}
\expandafter\def\csname sga@target@IV.3.2.3\endcsname{ega4.c04.prop.3-2-3}
\expandafter\def\csname sga@target@IV.3.2.4\endcsname{ega4.c04.def.3-2-4}
\expandafter\def\csname sga@target@IV.3.2.5\endcsname{ega4.c04.def.3-2-5}
\expandafter\def\csname sga@target@IV.3.2.6\endcsname{ega4.c04.prop.3-2-6}
\expandafter\def\csname sga@target@IV.3.2.7\endcsname{ega4.c04.cor.3-2-7}
\expandafter\def\csname sga@target@IV.3.2.8\endcsname{ega4.c04.cor.3-2-8}
\expandafter\def\csname sga@target@IV.3.2.9\endcsname{ega4.c04.rem.3-2-9}
\expandafter\def\csname sga@target@IV.3.3.1\endcsname{ega4.c04.prop.3-3-1}
\expandafter\def\csname sga@target@IV.3.3.1.1\endcsname{ega4.c04.prop.3-3-1-1}
\expandafter\def\csname sga@target@IV.3.3.2\endcsname{ega4.c04.cor.3-3-2}
\expandafter\def\csname sga@target@IV.3.3.3\endcsname{ega4.c04.cor.3-3-3}
\expandafter\def\csname sga@target@IV.3.3.4\endcsname{ega4.c04.prop.3-3-4}
\expandafter\def\csname sga@target@IV.3.3.5\endcsname{ega4.c04.cor.3-3-5}
\expandafter\def\csname sga@target@IV.3.3.6\endcsname{ega4.c04.prop.3-3-6}
\expandafter\def\csname sga@target@IV.3.3.6.1\endcsname{ega4.c04.eq.3-3-6-1}
\expandafter\def\csname sga@target@IV.3.3.6.2\endcsname{ega4.c04.eq.3-3-6-2}
\expandafter\def\csname sga@target@IV.3.3.6.3\endcsname{ega4.c04.eq.3-3-6-3}
\expandafter\def\csname sga@target@IV.3.3.6.4\endcsname{ega4.c04.eq.3-3-6-4}
\expandafter\def\csname sga@target@IV.3.3.7\endcsname{ega4.c04.cor.3-3-7}
\expandafter\def\csname sga@target@IV.3.3.7.1\endcsname{ega4.c04.cor.3-3-7-1}
\expandafter\def\csname sga@target@IV.3.3.8\endcsname{ega4.c04.sec.3-3-8}
\expandafter\def\csname sga@target@IV.3.3.9\endcsname{ega4.c04.cor.3-3-9}
\expandafter\def\csname sga@target@IV.3.3.9.1\endcsname{ega4.c04.cor.3-3-9-1}
\expandafter\def\csname sga@target@IV.3.4\endcsname{ega4.c04.subsec.3-4.key-iv-3-4}
\expandafter\def\csname sga@target@IV.3.4.1\endcsname{ega4.c04.prop.3-4-1}
\expandafter\def\csname sga@target@IV.3.4.1.1\endcsname{ega4.c04.eq.3-4-1-1}
\expandafter\def\csname sga@target@IV.3.4.1.2\endcsname{ega4.c04.lem.3-4-1-2}
\expandafter\def\csname sga@target@IV.3.4.1.3\endcsname{ega4.c04.lem.3-4-1-3}
\expandafter\def\csname sga@target@IV.3.4.1.4\endcsname{ega4.c04.lem.3-4-1-4}
\expandafter\def\csname sga@target@IV.3.4.1.5\endcsname{ega4.c04.eq.3-4-1-5}
\expandafter\def\csname sga@target@IV.3.4.2\endcsname{ega4.c04.cor.3-4-2}
\expandafter\def\csname sga@target@IV.3.4.3\endcsname{ega4.c04.prop.3-4-3}
\expandafter\def\csname sga@target@IV.3.4.4\endcsname{ega4.c04.prop.3-4-4}
\expandafter\def\csname sga@target@IV.3.4.5\endcsname{ega4.c04.prop.3-4-5}
\expandafter\def\csname sga@target@IV.3.4.6\endcsname{ega4.c04.prop.3-4-6}
\expandafter\def\csname sga@target@IV.3.4.6.1\endcsname{ega4.c04.lem.3-4-6-1}
\expandafter\def\csname sga@target@IV.3.4.7\endcsname{ega4.c04.cor.3-4-7}
\expandafter\def\csname sga@target@IV.3.4.8\endcsname{ega4.c04.rem.3-4-8}
\expandafter\def\csname sga@target@IV.3.4.9\endcsname{ega4.c04.prop.3-4-9}
\expandafter\def\csname sga@target@IV.4.1.1\endcsname{ega4.c04.def.4-1-1}
\expandafter\def\csname sga@target@IV.4.1.1.1\endcsname{ega4.c04.eq.4-1-1-1}
\expandafter\def\csname sga@target@IV.4.1.1.2\endcsname{ega4.c04.eq.4-1-1-2}
\expandafter\def\csname sga@target@IV.4.1.1.3\endcsname{ega4.c04.eq.4-1-1-3}
\expandafter\def\csname sga@target@IV.4.1.2\endcsname{ega4.c04.thm.4-1-2}
\expandafter\def\csname sga@target@IV.4.1.2.1\endcsname{ega4.c04.cor.4-1-2-1}
\expandafter\def\csname sga@target@IV.4.1.3\endcsname{ega4.c04.rem.4-1-3}
\expandafter\def\csname sga@target@IV.4.1.4\endcsname{ega4.c04.cor.4-1-4}
\expandafter\def\csname sga@target@IV.4.1.5\endcsname{ega4.c04.cor.4-1-5}
\expandafter\def\csname sga@target@IV.4.2.1\endcsname{ega4.c04.prop.4-2-1}
\expandafter\def\csname sga@target@IV.4.2.1.1\endcsname{ega4.c04.eq.4-2-1-1}
\expandafter\def\csname sga@target@IV.4.2.1.2\endcsname{ega4.c04.eq.4-2-1-2}
\expandafter\def\csname sga@target@IV.4.2.1.3\endcsname{ega4.c04.eq.4-2-1-3}
\expandafter\def\csname sga@target@IV.4.2.1.4\endcsname{ega4.c04.rem.4-2-1-4}
\expandafter\def\csname sga@target@IV.4.2.1.5\endcsname{ega4.c04.eq.4-2-1-5}
\expandafter\def\csname sga@target@IV.4.2.10\endcsname{ega4.c04.cor.4-2-10}
\expandafter\def\csname sga@target@IV.4.2.2\endcsname{ega4.c04.cor.4-2-2}
\expandafter\def\csname sga@target@IV.4.2.3\endcsname{ega4.c04.cor.4-2-3}
\expandafter\def\csname sga@target@IV.4.2.4\endcsname{ega4.c04.prop.4-2-4}
\expandafter\def\csname sga@target@IV.4.2.4.i\endcsname{ega4.c04.sec.4-2-4-i}
\expandafter\def\csname sga@target@IV.4.2.4.ii\endcsname{ega4.c04.sec.4-2-4-ii}
\expandafter\def\csname sga@target@IV.4.2.4.iii\endcsname{ega4.c04.sec.4-2-4-iii}
\expandafter\def\csname sga@target@IV.4.2.5\endcsname{ega4.c04.prop.4-2-5}
\expandafter\def\csname sga@target@IV.4.2.6\endcsname{ega4.c04.cor.4-2-6}
\expandafter\def\csname sga@target@IV.4.2.7\endcsname{ega4.c04.prop.4-2-7}
\expandafter\def\csname sga@target@IV.4.2.7.i\endcsname{ega4.c04.sec.4-2-7-i}
\expandafter\def\csname sga@target@IV.4.2.7.ii\endcsname{ega4.c04.sec.4-2-7-ii}
\expandafter\def\csname sga@target@IV.4.2.7.iii\endcsname{ega4.c04.sec.4-2-7-iii}
\expandafter\def\csname sga@target@IV.4.2.7.iv\endcsname{ega4.c04.sec.4-2-7-iv}
\expandafter\def\csname sga@target@IV.4.2.8\endcsname{ega4.c04.cor.4-2-8}
\expandafter\def\csname sga@target@IV.4.2.9\endcsname{ega4.c04.prop.4-2-9}
\expandafter\def\csname sga@target@IV.4.3.1\endcsname{ega4.c04.statement.4-3-1}
\expandafter\def\csname sga@target@IV.4.3.2\endcsname{ega4.c04.prop.4-3-2}
\expandafter\def\csname sga@target@IV.4.3.3\endcsname{ega4.c04.cor.4-3-3}
\expandafter\def\csname sga@target@IV.4.3.4\endcsname{ega4.c04.cor.4-3-4}
\expandafter\def\csname sga@target@IV.4.3.5\endcsname{ega4.c04.prop.4-3-5}
\expandafter\def\csname sga@target@IV.4.3.6\endcsname{ega4.c04.cor.4-3-6}
\expandafter\def\csname sga@target@IV.4.3.7\endcsname{ega4.c04.cor.4-3-7}
\expandafter\def\csname sga@target@IV.4.3.8\endcsname{ega4.c04.cor.4-3-8}
\expandafter\def\csname sga@target@IV.4.3.9\endcsname{ega4.c04.prop.4-3-9}
\expandafter\def\csname sga@target@IV.4.4.1\endcsname{ega4.c04.statement.4-4-1}
\expandafter\def\csname sga@target@IV.4.4.2\endcsname{ega4.c04.lem.4-4-2}
\expandafter\def\csname sga@target@IV.4.4.2.i\endcsname{ega4.c04.sec.4-4-2-i}
\expandafter\def\csname sga@target@IV.4.4.2.ii\endcsname{ega4.c04.sec.4-4-2-ii}
\expandafter\def\csname sga@target@IV.4.4.3\endcsname{ega4.c04.rem.4-4-3}
\expandafter\def\csname sga@target@IV.4.4.4\endcsname{ega4.c04.thm.4-4-4}
\expandafter\def\csname sga@target@IV.4.4.5\endcsname{ega4.c04.cor.4-4-5}
\expandafter\def\csname sga@target@IV.4.4.6\endcsname{ega4.c04.cor.4-4-6}
\expandafter\def\csname sga@target@IV.4.5\endcsname{ega4.c04.subsec.4-5.key-iv-4-5}
\expandafter\def\csname sga@target@IV.4.5.1\endcsname{ega4.c04.prop.4-5-1}
\expandafter\def\csname sga@target@IV.4.5.10\endcsname{ega4.c04.cor.4-5-10}
\expandafter\def\csname sga@target@IV.4.5.11\endcsname{ega4.c04.cor.4-5-11}
\expandafter\def\csname sga@target@IV.4.5.12\endcsname{ega4.c04.sec.4-5-12}
\expandafter\def\csname sga@target@IV.4.5.12.i\endcsname{ega4.c04.sec.4-5-12-i}
\expandafter\def\csname sga@target@IV.4.5.12.ii\endcsname{ega4.c04.sec.4-5-12-ii}
\expandafter\def\csname sga@target@IV.4.5.13\endcsname{ega4.c04.prop.4-5-13}
\expandafter\def\csname sga@target@IV.4.5.13.1\endcsname{ega4.c04.cor.4-5-13-1}
\expandafter\def\csname sga@target@IV.4.5.13.1.i\endcsname{ega4.c04.sec.4-5-13-1-i}
\expandafter\def\csname sga@target@IV.4.5.13.1.ii\endcsname{ega4.c04.sec.4-5-13-1-ii}
\expandafter\def\csname sga@target@IV.4.5.14\endcsname{ega4.c04.cor.4-5-14}
\expandafter\def\csname sga@target@IV.4.5.15\endcsname{ega4.c04.prop.4-5-15}
\expandafter\def\csname sga@target@IV.4.5.15.i\endcsname{ega4.c04.sec.4-5-15-i}
\expandafter\def\csname sga@target@IV.4.5.15.ii\endcsname{ega4.c04.sec.4-5-15-ii}
\expandafter\def\csname sga@target@IV.4.5.16\endcsname{ega4.c04.cor.4-5-16}
\expandafter\def\csname sga@target@IV.4.5.16.i\endcsname{ega4.c04.sec.4-5-16-i}
\expandafter\def\csname sga@target@IV.4.5.16.ii\endcsname{ega4.c04.sec.4-5-16-ii}
\expandafter\def\csname sga@target@IV.4.5.17\endcsname{ega4.c04.cor.4-5-17}
\expandafter\def\csname sga@target@IV.4.5.18\endcsname{ega4.c04.rem.4-5-18}
\expandafter\def\csname sga@target@IV.4.5.19\endcsname{ega4.c04.prop.4-5-19}
\expandafter\def\csname sga@target@IV.4.5.19.1\endcsname{ega4.c04.lem.4-5-19-1}
\expandafter\def\csname sga@target@IV.4.5.19.2\endcsname{ega4.c04.eq.4-5-19-2}
\expandafter\def\csname sga@target@IV.4.5.19.3\endcsname{ega4.c04.cor.4-5-19-3}
\expandafter\def\csname sga@target@IV.4.5.2\endcsname{ega4.c04.def.4-5-2}
\expandafter\def\csname sga@target@IV.4.5.2.a\endcsname{ega4.c04.sec.4-5-2-a}
\expandafter\def\csname sga@target@IV.4.5.2.b\endcsname{ega4.c04.sec.4-5-2-b}
\expandafter\def\csname sga@target@IV.4.5.2.c\endcsname{ega4.c04.sec.4-5-2-c}
\expandafter\def\csname sga@target@IV.4.5.20\endcsname{ega4.c04.rem.4-5-20}
\expandafter\def\csname sga@target@IV.4.5.21\endcsname{ega4.c04.prop.4-5-21}
\expandafter\def\csname sga@target@IV.4.5.3\endcsname{ega4.c04.statement.4-5-3}
\expandafter\def\csname sga@target@IV.4.5.4\endcsname{ega4.c04.prop.4-5-4}
\expandafter\def\csname sga@target@IV.4.5.5\endcsname{ega4.c04.def.4-5-5}
\expandafter\def\csname sga@target@IV.4.5.6\endcsname{ega4.c04.prop.4-5-6}
\expandafter\def\csname sga@target@IV.4.5.6.i\endcsname{ega4.c04.sec.4-5-6-i}
\expandafter\def\csname sga@target@IV.4.5.6.ii\endcsname{ega4.c04.sec.4-5-6-ii}
\expandafter\def\csname sga@target@IV.4.5.7\endcsname{ega4.c04.prop.4-5-7}
\expandafter\def\csname sga@target@IV.4.5.8\endcsname{ega4.c04.cor.4-5-8}
\expandafter\def\csname sga@target@IV.4.5.8.i\endcsname{ega4.c04.sec.4-5-8-i}
\expandafter\def\csname sga@target@IV.4.5.8.ii\endcsname{ega4.c04.sec.4-5-8-ii}
\expandafter\def\csname sga@target@IV.4.5.9\endcsname{ega4.c04.prop.4-5-9}
\expandafter\def\csname sga@target@IV.4.5.9.a\endcsname{ega4.c04.sec.4-5-9-a}
\expandafter\def\csname sga@target@IV.4.5.9.b\endcsname{ega4.c04.sec.4-5-9-b}
\expandafter\def\csname sga@target@IV.4.5.9.c\endcsname{ega4.c04.sec.4-5-9-c}
\expandafter\def\csname sga@target@IV.4.6\endcsname{ega4.c04.subsec.4-6.key-iv-4-6}
\expandafter\def\csname sga@target@IV.4.6.1\endcsname{ega4.c04.prop.4-6-1}
\expandafter\def\csname sga@target@IV.4.6.1.a\endcsname{ega4.c04.sec.4-6-1-a}
\expandafter\def\csname sga@target@IV.4.6.1.b\endcsname{ega4.c04.sec.4-6-1-b}
\expandafter\def\csname sga@target@IV.4.6.1.c\endcsname{ega4.c04.sec.4-6-1-c}
\expandafter\def\csname sga@target@IV.4.6.1.d\endcsname{ega4.c04.sec.4-6-1-d}
\expandafter\def\csname sga@target@IV.4.6.1.e\endcsname{ega4.c04.sec.4-6-1-e}
\expandafter\def\csname sga@target@IV.4.6.10\endcsname{ega4.c04.prop.4-6-10}
\expandafter\def\csname sga@target@IV.4.6.11\endcsname{ega4.c04.cor.4-6-11}
\expandafter\def\csname sga@target@IV.4.6.12\endcsname{ega4.c04.prop.4-6-12}
\expandafter\def\csname sga@target@IV.4.6.12.a\endcsname{ega4.c04.sec.4-6-12-a}
\expandafter\def\csname sga@target@IV.4.6.12.b\endcsname{ega4.c04.sec.4-6-12-b}
\expandafter\def\csname sga@target@IV.4.6.12.c\endcsname{ega4.c04.sec.4-6-12-c}
\expandafter\def\csname sga@target@IV.4.6.12.d\endcsname{ega4.c04.sec.4-6-12-d}
\expandafter\def\csname sga@target@IV.4.6.12.e\endcsname{ega4.c04.sec.4-6-12-e}
\expandafter\def\csname sga@target@IV.4.6.13\endcsname{ega4.c04.cor.4-6-13}
\expandafter\def\csname sga@target@IV.4.6.13.f\endcsname{ega4.c04.sec.4-6-13-f}
\expandafter\def\csname sga@target@IV.4.6.14\endcsname{ega4.c04.statement.4-6-14}
\expandafter\def\csname sga@target@IV.4.6.15\endcsname{ega4.c04.prop.4-6-15}
\expandafter\def\csname sga@target@IV.4.6.15.i\endcsname{ega4.c04.sec.4-6-15-i}
\expandafter\def\csname sga@target@IV.4.6.15.ii\endcsname{ega4.c04.sec.4-6-15-ii}
\expandafter\def\csname sga@target@IV.4.6.16\endcsname{ega4.c04.prop.4-6-16}
\expandafter\def\csname sga@target@IV.4.6.16.a\endcsname{ega4.c04.sec.4-6-16-a}
\expandafter\def\csname sga@target@IV.4.6.16.b\endcsname{ega4.c04.sec.4-6-16-b}
\expandafter\def\csname sga@target@IV.4.6.16.c\endcsname{ega4.c04.sec.4-6-16-c}
\expandafter\def\csname sga@target@IV.4.6.16.d\endcsname{ega4.c04.sec.4-6-16-d}
\expandafter\def\csname sga@target@IV.4.6.17\endcsname{ega4.c04.def.4-6-17}
\expandafter\def\csname sga@target@IV.4.6.18\endcsname{ega4.c04.prop.4-6-18}
\expandafter\def\csname sga@target@IV.4.6.18.a\endcsname{ega4.c04.sec.4-6-18-a}
\expandafter\def\csname sga@target@IV.4.6.18.b\endcsname{ega4.c04.sec.4-6-18-b}
\expandafter\def\csname sga@target@IV.4.6.18.c\endcsname{ega4.c04.sec.4-6-18-c}
\expandafter\def\csname sga@target@IV.4.6.18.d\endcsname{ega4.c04.sec.4-6-18-d}
\expandafter\def\csname sga@target@IV.4.6.19\endcsname{ega4.c04.def.4-6-19}
\expandafter\def\csname sga@target@IV.4.6.2\endcsname{ega4.c04.def.4-6-2}
\expandafter\def\csname sga@target@IV.4.6.20\endcsname{ega4.c04.prop.4-6-20}
\expandafter\def\csname sga@target@IV.4.6.21\endcsname{ega4.c04.prop.4-6-21}
\expandafter\def\csname sga@target@IV.4.6.21.i\endcsname{ega4.c04.sec.4-6-21-i}
\expandafter\def\csname sga@target@IV.4.6.21.ii\endcsname{ega4.c04.sec.4-6-21-ii}
\expandafter\def\csname sga@target@IV.4.6.22\endcsname{ega4.c04.def.4-6-22}
\expandafter\def\csname sga@target@IV.4.6.3\endcsname{ega4.c04.cor.4-6-3}
\expandafter\def\csname sga@target@IV.4.6.4\endcsname{ega4.c04.cor.4-6-4}
\expandafter\def\csname sga@target@IV.4.6.5\endcsname{ega4.c04.prop.4-6-5}
\expandafter\def\csname sga@target@IV.4.6.5.i\endcsname{ega4.c04.sec.4-6-5-i}
\expandafter\def\csname sga@target@IV.4.6.5.ii\endcsname{ega4.c04.sec.4-6-5-ii}
\expandafter\def\csname sga@target@IV.4.6.6\endcsname{ega4.c04.prop.4-6-6}
\expandafter\def\csname sga@target@IV.4.6.7\endcsname{ega4.c04.cor.4-6-7}
\expandafter\def\csname sga@target@IV.4.6.8\endcsname{ega4.c04.cor.4-6-8}
\expandafter\def\csname sga@target@IV.4.6.9\endcsname{ega4.c04.def.4-6-9}
\expandafter\def\csname sga@target@IV.4.7.1\endcsname{ega4.c04.lem.4-7-1}
\expandafter\def\csname sga@target@IV.4.7.1.1\endcsname{ega4.c04.eq.4-7-1-1}
\expandafter\def\csname sga@target@IV.4.7.10\endcsname{ega4.c04.prop.4-7-10}
\expandafter\def\csname sga@target@IV.4.7.11\endcsname{ega4.c04.prop.4-7-11}
\expandafter\def\csname sga@target@IV.4.7.12\endcsname{ega4.c04.def.4-7-12}
\expandafter\def\csname sga@target@IV.4.7.13\endcsname{ega4.c04.prop.4-7-13}
\expandafter\def\csname sga@target@IV.4.7.2\endcsname{ega4.c04.cor.4-7-2}
\expandafter\def\csname sga@target@IV.4.7.2.1\endcsname{ega4.c04.eq.4-7-2-1}
\expandafter\def\csname sga@target@IV.4.7.3\endcsname{ega4.c04.prop.4-7-3}
\expandafter\def\csname sga@target@IV.4.7.3.1\endcsname{ega4.c04.eq.4-7-3-1}
\expandafter\def\csname sga@target@IV.4.7.3.i\endcsname{ega4.c04.sec.4-7-3-i}
\expandafter\def\csname sga@target@IV.4.7.3.ii\endcsname{ega4.c04.sec.4-7-3-ii}
\expandafter\def\csname sga@target@IV.4.7.3.iii\endcsname{ega4.c04.sec.4-7-3-iii}
\expandafter\def\csname sga@target@IV.4.7.4\endcsname{ega4.c04.def.4-7-4}
\expandafter\def\csname sga@target@IV.4.7.5\endcsname{ega4.c04.def.4-7-5}
\expandafter\def\csname sga@target@IV.4.7.6\endcsname{ega4.c04.prop.4-7-6}
\expandafter\def\csname sga@target@IV.4.7.6.1\endcsname{ega4.c04.eq.4-7-6-1}
\expandafter\def\csname sga@target@IV.4.7.7\endcsname{ega4.c04.statement.4-7-7}
\expandafter\def\csname sga@target@IV.4.7.8\endcsname{ega4.c04.prop.4-7-8}
\expandafter\def\csname sga@target@IV.4.7.9\endcsname{ega4.c04.cor.4-7-9}
\expandafter\def\csname sga@target@IV.4.8.1\endcsname{ega4.c04.statement.4-8-1}
\expandafter\def\csname sga@target@IV.4.8.10\endcsname{ega4.c04.cor.4-8-10}
\expandafter\def\csname sga@target@IV.4.8.11\endcsname{ega4.c04.cor.4-8-11}
\expandafter\def\csname sga@target@IV.4.8.12\endcsname{ega4.c04.cor.4-8-12}
\expandafter\def\csname sga@target@IV.4.8.13\endcsname{ega4.c04.prop.4-8-13}
\expandafter\def\csname sga@target@IV.4.8.14\endcsname{ega4.c04.prop.4-8-14}
\expandafter\def\csname sga@target@IV.4.8.15\endcsname{ega4.c04.sec.4-8-15}
\expandafter\def\csname sga@target@IV.4.8.15.i\endcsname{ega4.c04.sec.4-8-15-i}
\expandafter\def\csname sga@target@IV.4.8.15.ii\endcsname{ega4.c04.sec.4-8-15-ii}
\expandafter\def\csname sga@target@IV.4.8.2\endcsname{ega4.c04.statement.4-8-2}
\expandafter\def\csname sga@target@IV.4.8.2.1\endcsname{ega4.c04.eq.4-8-2-1}
\expandafter\def\csname sga@target@IV.4.8.2.2\endcsname{ega4.c04.eq.4-8-2-2}
\expandafter\def\csname sga@target@IV.4.8.2.3\endcsname{ega4.c04.eq.4-8-2-3}
\expandafter\def\csname sga@target@IV.4.8.2.4\endcsname{ega4.c04.eq.4-8-2-4}
\expandafter\def\csname sga@target@IV.4.8.2.5\endcsname{ega4.c04.eq.4-8-2-5}
\expandafter\def\csname sga@target@IV.4.8.2.6\endcsname{ega4.c04.eq.4-8-2-6}
\expandafter\def\csname sga@target@IV.4.8.3\endcsname{ega4.c04.prop.4-8-3}
\expandafter\def\csname sga@target@IV.4.8.4\endcsname{ega4.c04.def.4-8-4}
\expandafter\def\csname sga@target@IV.4.8.5\endcsname{ega4.c04.prop.4-8-5}
\expandafter\def\csname sga@target@IV.4.8.6\endcsname{ega4.c04.lem.4-8-6}
\expandafter\def\csname sga@target@IV.4.8.6.a\endcsname{ega4.c04.sec.4-8-6-a}
\expandafter\def\csname sga@target@IV.4.8.6.b\endcsname{ega4.c04.sec.4-8-6-b}
\expandafter\def\csname sga@target@IV.4.8.6.c\endcsname{ega4.c04.sec.4-8-6-c}
\expandafter\def\csname sga@target@IV.4.8.7\endcsname{ega4.c04.cor.4-8-7}
\expandafter\def\csname sga@target@IV.4.8.8\endcsname{ega4.c04.lem.4-8-8}
\expandafter\def\csname sga@target@IV.4.8.8.i\endcsname{ega4.c04.sec.4-8-8-i}
\expandafter\def\csname sga@target@IV.4.8.8.ii\endcsname{ega4.c04.sec.4-8-8-ii}
\expandafter\def\csname sga@target@IV.4.8.8.iii\endcsname{ega4.c04.sec.4-8-8-iii}
\expandafter\def\csname sga@target@IV.4.8.9\endcsname{ega4.c04.prop.4-8-9}
\expandafter\def\csname sga@target@IV.4.9.1\endcsname{ega4.c04.prop.4-9-1}
\expandafter\def\csname sga@target@IV.4.9.1.a\endcsname{ega4.c04.sec.4-9-1-a}
\expandafter\def\csname sga@target@IV.4.9.1.b\endcsname{ega4.c04.sec.4-9-1-b}
\expandafter\def\csname sga@target@IV.4.9.1.bprime\endcsname{ega4.c04.sec.4-9-1-bprime}
\expandafter\def\csname sga@target@IV.4.9.1.c\endcsname{ega4.c04.sec.4-9-1-c}
\expandafter\def\csname sga@target@IV.4.9.2\endcsname{ega4.c04.cor.4-9-2}
\expandafter\def\csname sga@target@IV.4.9.3\endcsname{ega4.c04.rem.4-9-3}
\expandafter\def\csname sga@target@IV.4.9.4\endcsname{ega4.c04.prop.4-9-4}
\expandafter\def\csname sga@target@IV.4.9.5\endcsname{ega4.c04.cor.4-9-5}
\expandafter\def\csname sga@target@IV.4.9.5.a\endcsname{ega4.c04.sec.4-9-5-a}
\expandafter\def\csname sga@target@IV.4.9.5.b\endcsname{ega4.c04.sec.4-9-5-b}
\expandafter\def\csname sga@target@IV.4.9.6\endcsname{ega4.c04.prop.4-9-6}
\expandafter\def\csname sga@target@IV.4.9.6.a\endcsname{ega4.c04.sec.4-9-6-a}
\expandafter\def\csname sga@target@IV.4.9.6.aprime\endcsname{ega4.c04.sec.4-9-6-aprime}
\expandafter\def\csname sga@target@IV.4.9.6.b\endcsname{ega4.c04.sec.4-9-6-b}
\expandafter\def\csname sga@target@IV.4.9.6.bprime\endcsname{ega4.c04.sec.4-9-6-bprime}
\expandafter\def\csname sga@target@IV.4.9.7\endcsname{ega4.c04.cor.4-9-7}
\expandafter\def\csname sga@target@IV.5.1.1\endcsname{ega4.c04.statement.5-1-1}
\expandafter\def\csname sga@target@IV.5.1.1.1\endcsname{ega4.c04.eq.5-1-1-1}
\expandafter\def\csname sga@target@IV.5.1.1.2\endcsname{ega4.c04.eq.5-1-1-2}
\expandafter\def\csname sga@target@IV.5.1.1.3\endcsname{ega4.c04.eq.5-1-1-3}
\expandafter\def\csname sga@target@IV.5.1.1.4\endcsname{ega4.c04.sec.5-1-1-4}
\expandafter\def\csname sga@target@IV.5.1.1.5\endcsname{ega4.c04.sec.5-1-1-5}
\expandafter\def\csname sga@target@IV.5.1.1.6\endcsname{ega4.c04.sec.5-1-1-6}
\expandafter\def\csname sga@target@IV.5.1.10\endcsname{ega4.c04.prop.5-1-10}
\expandafter\def\csname sga@target@IV.5.1.10.1\endcsname{ega4.c04.lem.5-1-10-1}
\expandafter\def\csname sga@target@IV.5.1.10.i\endcsname{ega4.c04.sec.5-1-10-i}
\expandafter\def\csname sga@target@IV.5.1.10.ii\endcsname{ega4.c04.sec.5-1-10-ii}
\expandafter\def\csname sga@target@IV.5.1.11\endcsname{ega4.c04.cor.5-1-11}
\expandafter\def\csname sga@target@IV.5.1.12\endcsname{ega4.c04.statement.5-1-12}
\expandafter\def\csname sga@target@IV.5.1.12.1\endcsname{ega4.c04.eq.5-1-12-1}
\expandafter\def\csname sga@target@IV.5.1.12.2\endcsname{ega4.c04.eq.5-1-12-2}
\expandafter\def\csname sga@target@IV.5.1.12.3\endcsname{ega4.c04.eq.5-1-12-3}
\expandafter\def\csname sga@target@IV.5.1.13\endcsname{ega4.c04.prop.5-1-13}
\expandafter\def\csname sga@target@IV.5.1.13.1\endcsname{ega4.c04.eq.5-1-13-1}
\expandafter\def\csname sga@target@IV.5.1.2\endcsname{ega4.c04.prop.5-1-2}
\expandafter\def\csname sga@target@IV.5.1.2.1\endcsname{ega4.c04.eq.5-1-2-1}
\expandafter\def\csname sga@target@IV.5.1.3\endcsname{ega4.c04.cor.5-1-3}
\expandafter\def\csname sga@target@IV.5.1.3.1\endcsname{ega4.c04.eq.5-1-3-1}
\expandafter\def\csname sga@target@IV.5.1.3.2\endcsname{ega4.c04.eq.5-1-3-2}
\expandafter\def\csname sga@target@IV.5.1.4\endcsname{ega4.c04.prop.5-1-4}
\expandafter\def\csname sga@target@IV.5.1.4.1\endcsname{ega4.c04.eq.5-1-4-1}
\expandafter\def\csname sga@target@IV.5.1.4.2\endcsname{ega4.c04.eq.5-1-4-2}
\expandafter\def\csname sga@target@IV.5.1.5\endcsname{ega4.c04.cor.5-1-5}
\expandafter\def\csname sga@target@IV.5.1.6\endcsname{ega4.c04.statement.5-1-6}
\expandafter\def\csname sga@target@IV.5.1.7\endcsname{ega4.c04.prop.5-1-7}
\expandafter\def\csname sga@target@IV.5.1.8\endcsname{ega4.c04.cor.5-1-8}
\expandafter\def\csname sga@target@IV.5.1.9\endcsname{ega4.c04.prop.5-1-9}
\expandafter\def\csname sga@target@IV.5.1.9.1\endcsname{ega4.c04.statement.5-1-9-1}
\expandafter\def\csname sga@target@IV.5.10.1\endcsname{ega4.c04.statement.5-10-1}
\expandafter\def\csname sga@target@IV.5.10.1.1\endcsname{ega4.c04.eq.5-10-1-1}
\expandafter\def\csname sga@target@IV.5.10.10\endcsname{ega4.c04.prop.5-10-10}
\expandafter\def\csname sga@target@IV.5.10.11\endcsname{ega4.c04.statement.5-10-11}
\expandafter\def\csname sga@target@IV.5.10.11.1\endcsname{ega4.c04.eq.5-10-11-1}
\expandafter\def\csname sga@target@IV.5.10.12\endcsname{ega4.c04.prop.5-10-12}
\expandafter\def\csname sga@target@IV.5.10.13\endcsname{ega4.c04.statement.5-10-13}
\expandafter\def\csname sga@target@IV.5.10.14\endcsname{ega4.c04.prop.5-10-14}
\expandafter\def\csname sga@target@IV.5.10.15\endcsname{ega4.c04.cor.5-10-15}
\expandafter\def\csname sga@target@IV.5.10.16\endcsname{ega4.c04.prop.5-10-16}
\expandafter\def\csname sga@target@IV.5.10.17\endcsname{ega4.c04.prop.5-10-17}
\expandafter\def\csname sga@target@IV.5.10.2\endcsname{ega4.c04.prop.5-10-2}
\expandafter\def\csname sga@target@IV.5.10.3\endcsname{ega4.c04.cor.5-10-3}
\expandafter\def\csname sga@target@IV.5.10.4\endcsname{ega4.c04.cor.5-10-4}
\expandafter\def\csname sga@target@IV.5.10.5\endcsname{ega4.c04.thm.5-10-5}
\expandafter\def\csname sga@target@IV.5.10.6\endcsname{ega4.c04.cor.5-10-6}
\expandafter\def\csname sga@target@IV.5.10.7\endcsname{ega4.c04.thm.5-10-7}
\expandafter\def\csname sga@target@IV.5.10.8\endcsname{ega4.c04.cor.5-10-8}
\expandafter\def\csname sga@target@IV.5.10.8.1\endcsname{ega4.c04.lem.5-10-8-1}
\expandafter\def\csname sga@target@IV.5.10.9\endcsname{ega4.c04.cor.5-10-9}
\expandafter\def\csname sga@target@IV.5.11.1\endcsname{ega4.c04.prop.5-11-1}
\expandafter\def\csname sga@target@IV.5.11.1.1\endcsname{ega4.c04.lem.5-11-1-1}
\expandafter\def\csname sga@target@IV.5.11.2\endcsname{ega4.c04.cor.5-11-2}
\expandafter\def\csname sga@target@IV.5.11.3\endcsname{ega4.c04.rem.5-11-3}
\expandafter\def\csname sga@target@IV.5.11.4\endcsname{ega4.c04.cor.5-11-4}
\expandafter\def\csname sga@target@IV.5.11.5\endcsname{ega4.c04.cor.5-11-5}
\expandafter\def\csname sga@target@IV.5.11.6\endcsname{ega4.c04.cor.5-11-6}
\expandafter\def\csname sga@target@IV.5.11.7\endcsname{ega4.c04.item.5-11-7}
\expandafter\def\csname sga@target@IV.5.12\endcsname{ega4.c04.subsec.5-12.key-iv-5-12}
\expandafter\def\csname sga@target@IV.5.12.1\endcsname{ega4.c04.prop.5-12-1}
\expandafter\def\csname sga@target@IV.5.12.1.1\endcsname{ega4.c04.eq.5-12-1-1}
\expandafter\def\csname sga@target@IV.5.12.10\endcsname{ega4.c04.cor.5-12-10}
\expandafter\def\csname sga@target@IV.5.12.2\endcsname{ega4.c04.prop.5-12-2}
\expandafter\def\csname sga@target@IV.5.12.3\endcsname{ega4.c04.rem.5-12-3}
\expandafter\def\csname sga@target@IV.5.12.4\endcsname{ega4.c04.cor.5-12-4}
\expandafter\def\csname sga@target@IV.5.12.5\endcsname{ega4.c04.prop.5-12-5}
\expandafter\def\csname sga@target@IV.5.12.6\endcsname{ega4.c04.rem.5-12-6}
\expandafter\def\csname sga@target@IV.5.12.7\endcsname{ega4.c04.cor.5-12-7}
\expandafter\def\csname sga@target@IV.5.12.8\endcsname{ega4.c04.prop.5-12-8}
\expandafter\def\csname sga@target@IV.5.12.9\endcsname{ega4.c04.rem.5-12-9}
\expandafter\def\csname sga@target@IV.5.13.1\endcsname{ega4.c04.prop.5-13-1}
\expandafter\def\csname sga@target@IV.5.13.2\endcsname{ega4.c04.cor.5-13-2}
\expandafter\def\csname sga@target@IV.5.13.3\endcsname{ega4.c04.prop.5-13-3}
\expandafter\def\csname sga@target@IV.5.13.4\endcsname{ega4.c04.cor.5-13-4}
\expandafter\def\csname sga@target@IV.5.13.5\endcsname{ega4.c04.statement.5-13-5}
\expandafter\def\csname sga@target@IV.5.13.6\endcsname{ega4.c04.cor.5-13-6}
\expandafter\def\csname sga@target@IV.5.13.7\endcsname{ega4.c04.prop.5-13-7}
\expandafter\def\csname sga@target@IV.5.13.7.1\endcsname{ega4.c04.lem.5-13-7-1}
\expandafter\def\csname sga@target@IV.5.13.8\endcsname{ega4.c04.cor.5-13-8}
\expandafter\def\csname sga@target@IV.5.2.1\endcsname{ega4.c04.prop.5-2-1}
\expandafter\def\csname sga@target@IV.5.2.1.1\endcsname{ega4.c04.eq.5-2-1-1}
\expandafter\def\csname sga@target@IV.5.2.2\endcsname{ega4.c04.cor.5-2-2}
\expandafter\def\csname sga@target@IV.5.2.3\endcsname{ega4.c04.cor.5-2-3}
\expandafter\def\csname sga@target@IV.5.2.3.1\endcsname{ega4.c04.eq.5-2-3-1}
\expandafter\def\csname sga@target@IV.5.2.4\endcsname{ega4.c04.cor.5-2-4}
\expandafter\def\csname sga@target@IV.5.2.5\endcsname{ega4.c04.sec.5-2-5}
\expandafter\def\csname sga@target@IV.5.2.5.i\endcsname{ega4.c04.sec.5-2-5-i}
\expandafter\def\csname sga@target@IV.5.2.5.ii\endcsname{ega4.c04.sec.5-2-5-ii}
\expandafter\def\csname sga@target@IV.5.3.1\endcsname{ega4.c04.prop.5-3-1}
\expandafter\def\csname sga@target@IV.5.3.1.1\endcsname{ega4.c04.eq.5-3-1-1}
\expandafter\def\csname sga@target@IV.5.3.1.2\endcsname{ega4.c04.eq.5-3-1-2}
\expandafter\def\csname sga@target@IV.5.3.1.3\endcsname{ega4.c04.lem.5-3-1-3}
\expandafter\def\csname sga@target@IV.5.4.1\endcsname{ega4.c04.prop.5-4-1}
\expandafter\def\csname sga@target@IV.5.4.1.i\endcsname{ega4.c04.sec.5-4-1-i}
\expandafter\def\csname sga@target@IV.5.4.1.ii\endcsname{ega4.c04.sec.5-4-1-ii}
\expandafter\def\csname sga@target@IV.5.4.2\endcsname{ega4.c04.cor.5-4-2}
\expandafter\def\csname sga@target@IV.5.4.3\endcsname{ega4.c04.sec.5-4-3}
\expandafter\def\csname sga@target@IV.5.4.3.1\endcsname{ega4.c04.eq.5-4-3-1}
\expandafter\def\csname sga@target@IV.5.4.3.i\endcsname{ega4.c04.sec.5-4-3-i}
\expandafter\def\csname sga@target@IV.5.4.3.ii\endcsname{ega4.c04.sec.5-4-3-ii}
\expandafter\def\csname sga@target@IV.5.5\endcsname{ega4.c04.sec.5-5}
\expandafter\def\csname sga@target@IV.5.5.1\endcsname{ega4.c04.statement.5-5-1}
\expandafter\def\csname sga@target@IV.5.5.1.1\endcsname{ega4.c04.eq.5-5-1-1}
\expandafter\def\csname sga@target@IV.5.5.2\endcsname{ega4.c04.prop.5-5-2}
\expandafter\def\csname sga@target@IV.5.5.2.1\endcsname{ega4.c04.eq.5-5-2-1}
\expandafter\def\csname sga@target@IV.5.5.2.2\endcsname{ega4.c04.eq.5-5-2-2}
\expandafter\def\csname sga@target@IV.5.5.3\endcsname{ega4.c04.prop.5-5-3}
\expandafter\def\csname sga@target@IV.5.5.3.1\endcsname{ega4.c04.eq.5-5-3-1}
\expandafter\def\csname sga@target@IV.5.5.4\endcsname{ega4.c04.cor.5-5-4}
\expandafter\def\csname sga@target@IV.5.5.4.1\endcsname{ega4.c04.eq.5-5-4-1}
\expandafter\def\csname sga@target@IV.5.5.5\endcsname{ega4.c04.cor.5-5-5}
\expandafter\def\csname sga@target@IV.5.5.6\endcsname{ega4.c04.cor.5-5-6}
\expandafter\def\csname sga@target@IV.5.5.6.1\endcsname{ega4.c04.eq.5-5-6-1}
\expandafter\def\csname sga@target@IV.5.5.7\endcsname{ega4.c04.lem.5-5-7}
\expandafter\def\csname sga@target@IV.5.5.7.1\endcsname{ega4.c04.eq.5-5-7-1}
\expandafter\def\csname sga@target@IV.5.5.7.2\endcsname{ega4.c04.statement.5-5-7-2}
\expandafter\def\csname sga@target@IV.5.5.7.3\endcsname{ega4.c04.eq.5-5-7-3}
\expandafter\def\csname sga@target@IV.5.5.7.i\endcsname{ega4.c04.sec.5-5-7-i}
\expandafter\def\csname sga@target@IV.5.5.7.ii\endcsname{ega4.c04.sec.5-5-7-ii}
\expandafter\def\csname sga@target@IV.5.5.8\endcsname{ega4.c04.thm.5-5-8}
\expandafter\def\csname sga@target@IV.5.5.8.1\endcsname{ega4.c04.eq.5-5-8-1}
\expandafter\def\csname sga@target@IV.5.5.8.2\endcsname{ega4.c04.eq.5-5-8-2}
\expandafter\def\csname sga@target@IV.5.5.8.3\endcsname{ega4.c04.eq.5-5-8-3}
\expandafter\def\csname sga@target@IV.5.6\endcsname{ega4.c04.sec.5-6}
\expandafter\def\csname sga@target@IV.5.6.1\endcsname{ega4.c04.prop.5-6-1}
\expandafter\def\csname sga@target@IV.5.6.1.1\endcsname{ega4.c04.eq.5-6-1-1}
\expandafter\def\csname sga@target@IV.5.6.1.2\endcsname{ega4.c04.eq.5-6-1-2}
\expandafter\def\csname sga@target@IV.5.6.1.3\endcsname{ega4.c04.statement.5-6-1-3}
\expandafter\def\csname sga@target@IV.5.6.1.4\endcsname{ega4.c04.statement.5-6-1-4}
\expandafter\def\csname sga@target@IV.5.6.1.5\endcsname{ega4.c04.statement.5-6-1-5}
\expandafter\def\csname sga@target@IV.5.6.1.a\endcsname{ega4.c04.sec.5-6-1-a}
\expandafter\def\csname sga@target@IV.5.6.1.b\endcsname{ega4.c04.sec.5-6-1-b}
\expandafter\def\csname sga@target@IV.5.6.1.c\endcsname{ega4.c04.sec.5-6-1-c}
\expandafter\def\csname sga@target@IV.5.6.10\endcsname{ega4.c04.prop.5-6-10}
\expandafter\def\csname sga@target@IV.5.6.11\endcsname{ega4.c04.ex.5-6-11}
\expandafter\def\csname sga@target@IV.5.6.11.1\endcsname{ega4.c04.lem.5-6-11-1}
\expandafter\def\csname sga@target@IV.5.6.11.1.i\endcsname{ega4.c04.sec.5-6-11-1-i}
\expandafter\def\csname sga@target@IV.5.6.11.1.ii\endcsname{ega4.c04.sec.5-6-11-1-ii}
\expandafter\def\csname sga@target@IV.5.6.2\endcsname{ega4.c04.def.5-6-2}
\expandafter\def\csname sga@target@IV.5.6.3\endcsname{ega4.c04.sec.5-6-3}
\expandafter\def\csname sga@target@IV.5.6.3.i\endcsname{ega4.c04.sec.5-6-3-i}
\expandafter\def\csname sga@target@IV.5.6.3.ii\endcsname{ega4.c04.sec.5-6-3-ii}
\expandafter\def\csname sga@target@IV.5.6.3.iii\endcsname{ega4.c04.sec.5-6-3-iii}
\expandafter\def\csname sga@target@IV.5.6.3.iv\endcsname{ega4.c04.sec.5-6-3-iv}
\expandafter\def\csname sga@target@IV.5.6.4\endcsname{ega4.c04.prop.5-6-4}
\expandafter\def\csname sga@target@IV.5.6.5\endcsname{ega4.c04.prop.5-6-5}
\expandafter\def\csname sga@target@IV.5.6.5.1\endcsname{ega4.c04.eq.5-6-5-1}
\expandafter\def\csname sga@target@IV.5.6.5.2\endcsname{ega4.c04.eq.5-6-5-2}
\expandafter\def\csname sga@target@IV.5.6.5.3\endcsname{ega4.c04.eq.5-6-5-3}
\expandafter\def\csname sga@target@IV.5.6.6\endcsname{ega4.c04.cor.5-6-6}
\expandafter\def\csname sga@target@IV.5.6.6.1\endcsname{ega4.c04.eq.5-6-6-1}
\expandafter\def\csname sga@target@IV.5.6.6.2\endcsname{ega4.c04.eq.5-6-6-2}
\expandafter\def\csname sga@target@IV.5.6.7\endcsname{ega4.c04.cor.5-6-7}
\expandafter\def\csname sga@target@IV.5.6.7.1\endcsname{ega4.c04.eq.5-6-7-1}
\expandafter\def\csname sga@target@IV.5.6.8\endcsname{ega4.c04.cor.5-6-8}
\expandafter\def\csname sga@target@IV.5.6.8.1\endcsname{ega4.c04.eq.5-6-8-1}
\expandafter\def\csname sga@target@IV.5.6.8.2\endcsname{ega4.c04.eq.5-6-8-2}
\expandafter\def\csname sga@target@IV.5.6.9\endcsname{ega4.c04.sec.5-6-9}
\expandafter\def\csname sga@target@IV.5.6.9.i\endcsname{ega4.c04.sec.5-6-9-i}
\expandafter\def\csname sga@target@IV.5.6.9.ii\endcsname{ega4.c04.sec.5-6-9-ii}
\expandafter\def\csname sga@target@IV.5.7.1\endcsname{ega4.c04.def.5-7-1}
\expandafter\def\csname sga@target@IV.5.7.1.1\endcsname{ega4.c04.eq.5-7-1-1}
\expandafter\def\csname sga@target@IV.5.7.10\endcsname{ega4.c04.cor.5-7-10}
\expandafter\def\csname sga@target@IV.5.7.11\endcsname{ega4.c04.cor.5-7-11}
\expandafter\def\csname sga@target@IV.5.7.11.i\endcsname{ega4.c04.sec.5-7-11-i}
\expandafter\def\csname sga@target@IV.5.7.11.ii\endcsname{ega4.c04.sec.5-7-11-ii}
\expandafter\def\csname sga@target@IV.5.7.12\endcsname{ega4.c04.statement.5-7-12}
\expandafter\def\csname sga@target@IV.5.7.13\endcsname{ega4.c04.cor.5-7-13}
\expandafter\def\csname sga@target@IV.5.7.13.1\endcsname{ega4.c04.eq.5-7-13-1}
\expandafter\def\csname sga@target@IV.5.7.2\endcsname{ega4.c04.def.5-7-2}
\expandafter\def\csname sga@target@IV.5.7.2.1\endcsname{ega4.c04.eq.5-7-2-1}
\expandafter\def\csname sga@target@IV.5.7.2.2\endcsname{ega4.c04.eq.5-7-2-2}
\expandafter\def\csname sga@target@IV.5.7.3\endcsname{ega4.c04.sec.5-7-3}
\expandafter\def\csname sga@target@IV.5.7.3.i\endcsname{ega4.c04.sec.5-7-3-i}
\expandafter\def\csname sga@target@IV.5.7.3.ii\endcsname{ega4.c04.sec.5-7-3-ii}
\expandafter\def\csname sga@target@IV.5.7.3.iii\endcsname{ega4.c04.sec.5-7-3-iii}
\expandafter\def\csname sga@target@IV.5.7.3.iv\endcsname{ega4.c04.sec.5-7-3-iv}
\expandafter\def\csname sga@target@IV.5.7.3.v\endcsname{ega4.c04.sec.5-7-3-v}
\expandafter\def\csname sga@target@IV.5.7.3.vi\endcsname{ega4.c04.sec.5-7-3-vi}
\expandafter\def\csname sga@target@IV.5.7.4\endcsname{ega4.c04.prop.5-7-4}
\expandafter\def\csname sga@target@IV.5.7.4.1\endcsname{ega4.c04.eq.5-7-4-1}
\expandafter\def\csname sga@target@IV.5.7.4.2\endcsname{ega4.c04.eq.5-7-4-2}
\expandafter\def\csname sga@target@IV.5.7.4.i\endcsname{ega4.c04.sec.5-7-4-i}
\expandafter\def\csname sga@target@IV.5.7.4.ii\endcsname{ega4.c04.sec.5-7-4-ii}
\expandafter\def\csname sga@target@IV.5.7.5\endcsname{ega4.c04.prop.5-7-5}
\expandafter\def\csname sga@target@IV.5.7.6\endcsname{ega4.c04.cor.5-7-6}
\expandafter\def\csname sga@target@IV.5.7.7\endcsname{ega4.c04.cor.5-7-7}
\expandafter\def\csname sga@target@IV.5.7.8\endcsname{ega4.c04.sec.5-7-8}
\expandafter\def\csname sga@target@IV.5.7.9\endcsname{ega4.c04.prop.5-7-9}
\expandafter\def\csname sga@target@IV.5.7.9.1\endcsname{ega4.c04.eq.5-7-9-1}
\expandafter\def\csname sga@target@IV.5.7.9.2\endcsname{ega4.c04.eq.5-7-9-2}
\expandafter\def\csname sga@target@IV.5.8.1\endcsname{ega4.c04.statement.5-8-1}
\expandafter\def\csname sga@target@IV.5.8.2\endcsname{ega4.c04.def.5-8-2}
\expandafter\def\csname sga@target@IV.5.8.3\endcsname{ega4.c04.prop.5-8-3}
\expandafter\def\csname sga@target@IV.5.8.4\endcsname{ega4.c04.statement.5-8-4}
\expandafter\def\csname sga@target@IV.5.8.5\endcsname{ega4.c04.prop.5-8-5}
\expandafter\def\csname sga@target@IV.5.8.6\endcsname{ega4.c04.thm.5-8-6}
\expandafter\def\csname sga@target@IV.5.8.6.1\endcsname{ega4.c04.eq.5-8-6-1}
\expandafter\def\csname sga@target@IV.5.8.6.i\endcsname{ega4.c04.item.5-8-6-i}
\expandafter\def\csname sga@target@IV.5.8.6.ii\endcsname{ega4.c04.item.5-8-6-ii}
\expandafter\def\csname sga@target@IV.5.9.1\endcsname{ega4.c04.statement.5-9-1}
\expandafter\def\csname sga@target@IV.5.9.1.1\endcsname{ega4.c04.eq.5-9-1-1}
\expandafter\def\csname sga@target@IV.5.9.1.2\endcsname{ega4.c04.eq.5-9-1-2}
\expandafter\def\csname sga@target@IV.5.9.1.3\endcsname{ega4.c04.eq.5-9-1-3}
\expandafter\def\csname sga@target@IV.5.9.10\endcsname{ega4.c04.cor.5-9-10}
\expandafter\def\csname sga@target@IV.5.9.10.i\endcsname{ega4.c04.item.5-9-10-i}
\expandafter\def\csname sga@target@IV.5.9.10.ii\endcsname{ega4.c04.item.5-9-10-ii}
\expandafter\def\csname sga@target@IV.5.9.11\endcsname{ega4.c04.prop.5-9-11}
\expandafter\def\csname sga@target@IV.5.9.12\endcsname{ega4.c04.item.5-9-12}
\expandafter\def\csname sga@target@IV.5.9.12.i\endcsname{ega4.c04.item.5-9-12-i}
\expandafter\def\csname sga@target@IV.5.9.12.ii\endcsname{ega4.c04.item.5-9-12-ii}
\expandafter\def\csname sga@target@IV.5.9.12.iii\endcsname{ega4.c04.item.5-9-12-iii}
\expandafter\def\csname sga@target@IV.5.9.2\endcsname{ega4.c04.prop.5-9-2}
\expandafter\def\csname sga@target@IV.5.9.2.i\endcsname{ega4.c04.item.5-9-2-i}
\expandafter\def\csname sga@target@IV.5.9.2.ii\endcsname{ega4.c04.item.5-9-2-ii}
\expandafter\def\csname sga@target@IV.5.9.3\endcsname{ega4.c04.rem.5-9-3}
\expandafter\def\csname sga@target@IV.5.9.3.1\endcsname{ega4.c04.eq.5-9-3-1}
\expandafter\def\csname sga@target@IV.5.9.4\endcsname{ega4.c04.prop.5-9-4}
\expandafter\def\csname sga@target@IV.5.9.4.1\endcsname{ega4.c04.eq.5-9-4-1}
\expandafter\def\csname sga@target@IV.5.9.5\endcsname{ega4.c04.cor.5-9-5}
\expandafter\def\csname sga@target@IV.5.9.6\endcsname{ega4.c04.cor.5-9-6}
\expandafter\def\csname sga@target@IV.5.9.6.1\endcsname{ega4.c04.eq.5-9-6-1}
\expandafter\def\csname sga@target@IV.5.9.7\endcsname{ega4.c04.statement.5-9-7}
\expandafter\def\csname sga@target@IV.5.9.7.1\endcsname{ega4.c04.eq.5-9-7-1}
\expandafter\def\csname sga@target@IV.5.9.8\endcsname{ega4.c04.prop.5-9-8}
\expandafter\def\csname sga@target@IV.5.9.8.1\endcsname{ega4.c04.eq.5-9-8-1}
\expandafter\def\csname sga@target@IV.5.9.8.a\endcsname{ega4.c04.item.5-9-8-a}
\expandafter\def\csname sga@target@IV.5.9.8.a-prime\endcsname{ega4.c04.item.5-9-8-a-prime}
\expandafter\def\csname sga@target@IV.5.9.8.b\endcsname{ega4.c04.item.5-9-8-b}
\expandafter\def\csname sga@target@IV.5.9.8.b-prime\endcsname{ega4.c04.item.5-9-8-b-prime}
\expandafter\def\csname sga@target@IV.5.9.9\endcsname{ega4.c04.def.5-9-9}
\expandafter\def\csname sga@target@IV.6.1.1\endcsname{ega4.c04.prop.6-1-1}
\expandafter\def\csname sga@target@IV.6.1.1.1\endcsname{ega4.c04.eq.6-1-1-1}
\expandafter\def\csname sga@target@IV.6.1.2\endcsname{ega4.c04.cor.6-1-2}
\expandafter\def\csname sga@target@IV.6.1.2.1\endcsname{ega4.c04.eq.6-1-2-1}
\expandafter\def\csname sga@target@IV.6.1.3\endcsname{ega4.c04.cor.6-1-3}
\expandafter\def\csname sga@target@IV.6.1.4\endcsname{ega4.c04.cor.6-1-4}
\expandafter\def\csname sga@target@IV.6.1.4.1\endcsname{ega4.c04.eq.6-1-4-1}
\expandafter\def\csname sga@target@IV.6.1.5\endcsname{ega4.c04.prop.6-1-5}
\expandafter\def\csname sga@target@IV.6.10.1\endcsname{ega4.c04.statement.6-10-1}
\expandafter\def\csname sga@target@IV.6.10.2\endcsname{ega4.c04.prop.6-10-2}
\expandafter\def\csname sga@target@IV.6.10.3\endcsname{ega4.c04.prop.6-10-3}
\expandafter\def\csname sga@target@IV.6.10.3.1\endcsname{ega4.c04.eq.6-10-3-1}
\expandafter\def\csname sga@target@IV.6.10.3.2\endcsname{ega4.c04.eq.6-10-3-2}
\expandafter\def\csname sga@target@IV.6.10.3.3\endcsname{ega4.c04.eq.6-10-3-3}
\expandafter\def\csname sga@target@IV.6.10.4\endcsname{ega4.c04.cor.6-10-4}
\expandafter\def\csname sga@target@IV.6.10.5\endcsname{ega4.c04.prop.6-10-5}
\expandafter\def\csname sga@target@IV.6.10.5.1\endcsname{ega4.c04.eq.6-10-5-1}
\expandafter\def\csname sga@target@IV.6.10.5.2\endcsname{ega4.c04.eq.6-10-5-2}
\expandafter\def\csname sga@target@IV.6.10.5.3\endcsname{ega4.c04.eq.6-10-5-3}
\expandafter\def\csname sga@target@IV.6.10.5.4\endcsname{ega4.c04.eq.6-10-5-4}
\expandafter\def\csname sga@target@IV.6.10.6\endcsname{ega4.c04.prop.6-10-6}
\expandafter\def\csname sga@target@IV.6.10.6.1\endcsname{ega4.c04.eq.6-10-6-1}
\expandafter\def\csname sga@target@IV.6.10.6.2\endcsname{ega4.c04.eq.6-10-6-2}
\expandafter\def\csname sga@target@IV.6.10.6.3\endcsname{ega4.c04.eq.6-10-6-3}
\expandafter\def\csname sga@target@IV.6.11\endcsname{ega4.c04.subsec.6-11.key-iv-6-11}
\expandafter\def\csname sga@target@IV.6.11.1\endcsname{ega4.c04.lem.6-11-1}
\expandafter\def\csname sga@target@IV.6.11.2\endcsname{ega4.c04.prop.6-11-2}
\expandafter\def\csname sga@target@IV.6.11.2.1\endcsname{ega4.c04.eq.6-11-2-1}
\expandafter\def\csname sga@target@IV.6.11.2.2\endcsname{ega4.c04.eq.6-11-2-2}
\expandafter\def\csname sga@target@IV.6.11.3\endcsname{ega4.c04.cor.6-11-3}
\expandafter\def\csname sga@target@IV.6.11.4\endcsname{ega4.c04.statement.6-11-4}
\expandafter\def\csname sga@target@IV.6.11.5\endcsname{ega4.c04.prop.6-11-5}
\expandafter\def\csname sga@target@IV.6.11.5.1\endcsname{ega4.c04.eq.6-11-5-1}
\expandafter\def\csname sga@target@IV.6.11.6\endcsname{ega4.c04.prop.6-11-6}
\expandafter\def\csname sga@target@IV.6.11.6.1\endcsname{ega4.c04.lem.6-11-6-1}
\expandafter\def\csname sga@target@IV.6.11.7\endcsname{ega4.c04.cor.6-11-7}
\expandafter\def\csname sga@target@IV.6.11.8\endcsname{ega4.c04.prop.6-11-8}
\expandafter\def\csname sga@target@IV.6.11.8.1\endcsname{ega4.c04.eq.6-11-8-1}
\expandafter\def\csname sga@target@IV.6.11.9\endcsname{ega4.c04.statement.6-11-9}
\expandafter\def\csname sga@target@IV.6.12\endcsname{ega4.c04.subsec.6-12.key-iv-6-12}
\expandafter\def\csname sga@target@IV.6.12.1\endcsname{ega4.c04.statement.6-12-1}
\expandafter\def\csname sga@target@IV.6.12.10\endcsname{ega4.c04.prop.6-12-10}
\expandafter\def\csname sga@target@IV.6.12.2\endcsname{ega4.c04.prop.6-12-2}
\expandafter\def\csname sga@target@IV.6.12.3\endcsname{ega4.c04.cor.6-12-3}
\expandafter\def\csname sga@target@IV.6.12.4\endcsname{ega4.c04.thm.6-12-4}
\expandafter\def\csname sga@target@IV.6.12.4.1\endcsname{ega4.c04.lem.6-12-4-1}
\expandafter\def\csname sga@target@IV.6.12.5\endcsname{ega4.c04.cor.6-12-5}
\expandafter\def\csname sga@target@IV.6.12.6\endcsname{ega4.c04.cor.6-12-6}
\expandafter\def\csname sga@target@IV.6.12.7\endcsname{ega4.c04.thm.6-12-7}
\expandafter\def\csname sga@target@IV.6.12.8\endcsname{ega4.c04.cor.6-12-8}
\expandafter\def\csname sga@target@IV.6.12.9\endcsname{ega4.c04.prop.6-12-9}
\expandafter\def\csname sga@target@IV.6.13\endcsname{ega4.c04.subsec.6-13.key-iv-6-13}
\expandafter\def\csname sga@target@IV.6.13.1\endcsname{ega4.c04.statement.6-13-1}
\expandafter\def\csname sga@target@IV.6.13.2\endcsname{ega4.c04.prop.6-13-2}
\expandafter\def\csname sga@target@IV.6.13.3\endcsname{ega4.c04.cor.6-13-3}
\expandafter\def\csname sga@target@IV.6.13.4\endcsname{ega4.c04.prop.6-13-4}
\expandafter\def\csname sga@target@IV.6.13.5\endcsname{ega4.c04.cor.6-13-5}
\expandafter\def\csname sga@target@IV.6.13.6\endcsname{ega4.c04.prop.6-13-6}
\expandafter\def\csname sga@target@IV.6.13.7\endcsname{ega4.c04.prop.6-13-7}
\expandafter\def\csname sga@target@IV.6.13.7.1\endcsname{ega4.c04.lem.6-13-7-1}
\expandafter\def\csname sga@target@IV.6.14.1\endcsname{ega4.c04.prop.6-14-1}
\expandafter\def\csname sga@target@IV.6.14.1.1\endcsname{ega4.c04.lem.6-14-1-1}
\expandafter\def\csname sga@target@IV.6.14.1.2\endcsname{ega4.c04.lem.6-14-1-2}
\expandafter\def\csname sga@target@IV.6.14.1.2.env\endcsname{ega4.c04.statement.6-14-1-2-env}
\expandafter\def\csname sga@target@IV.6.14.1.3\endcsname{ega4.c04.eq.6-14-1-3}
\expandafter\def\csname sga@target@IV.6.14.1.4\endcsname{ega4.c04.lem.6-14-1-4}
\expandafter\def\csname sga@target@IV.6.14.2\endcsname{ega4.c04.cor.6-14-2}
\expandafter\def\csname sga@target@IV.6.14.3\endcsname{ega4.c04.cor.6-14-3}
\expandafter\def\csname sga@target@IV.6.14.4\endcsname{ega4.c04.prop.6-14-4}
\expandafter\def\csname sga@target@IV.6.14.4.1\endcsname{ega4.c04.lem.6-14-4-1}
\expandafter\def\csname sga@target@IV.6.14.5\endcsname{ega4.c04.cor.6-14-5}
\expandafter\def\csname sga@target@IV.6.15.1\endcsname{ega4.c04.statement.6-15-1}
\expandafter\def\csname sga@target@IV.6.15.10\endcsname{ega4.c04.prop.6-15-10}
\expandafter\def\csname sga@target@IV.6.15.11\endcsname{ega4.c04.item.6-15-11}
\expandafter\def\csname sga@target@IV.6.15.2\endcsname{ega4.c04.statement.6-15-2}
\expandafter\def\csname sga@target@IV.6.15.3\endcsname{ega4.c04.statement.6-15-3}
\expandafter\def\csname sga@target@IV.6.15.3.1\endcsname{ega4.c04.lem.6-15-3-1}
\expandafter\def\csname sga@target@IV.6.15.4\endcsname{ega4.c04.statement.6-15-4}
\expandafter\def\csname sga@target@IV.6.15.4.1\endcsname{ega4.c04.lem.6-15-4-1}
\expandafter\def\csname sga@target@IV.6.15.5\endcsname{ega4.c04.prop.6-15-5}
\expandafter\def\csname sga@target@IV.6.15.6\endcsname{ega4.c04.prop.6-15-6}
\expandafter\def\csname sga@target@IV.6.15.7\endcsname{ega4.c04.prop.6-15-7}
\expandafter\def\csname sga@target@IV.6.15.8\endcsname{ega4.c04.lem.6-15-8}
\expandafter\def\csname sga@target@IV.6.15.9\endcsname{ega4.c04.cor.6-15-9}
\expandafter\def\csname sga@target@IV.6.2.1\endcsname{ega4.c04.prop.6-2-1}
\expandafter\def\csname sga@target@IV.6.2.1.1\endcsname{ega4.c04.eq.6-2-1-1}
\expandafter\def\csname sga@target@IV.6.2.1.2\endcsname{ega4.c04.eq.6-2-1-2}
\expandafter\def\csname sga@target@IV.6.2.2\endcsname{ega4.c04.cor.6-2-2}
\expandafter\def\csname sga@target@IV.6.2.3\endcsname{ega4.c04.prop.6-2-3}
\expandafter\def\csname sga@target@IV.6.2.3.1\endcsname{ega4.c04.eq.6-2-3-1}
\expandafter\def\csname sga@target@IV.6.2.4\endcsname{ega4.c04.prop.6-2-4}
\expandafter\def\csname sga@target@IV.6.2.4.1\endcsname{ega4.c04.eq.6-2-4-1}
\expandafter\def\csname sga@target@IV.6.3.1\endcsname{ega4.c04.prop.6-3-1}
\expandafter\def\csname sga@target@IV.6.3.1.1\endcsname{ega4.c04.eq.6-3-1-1}
\expandafter\def\csname sga@target@IV.6.3.2\endcsname{ega4.c04.cor.6-3-2}
\expandafter\def\csname sga@target@IV.6.3.2.1\endcsname{ega4.c04.eq.6-3-2-1}
\expandafter\def\csname sga@target@IV.6.3.3\endcsname{ega4.c04.cor.6-3-3}
\expandafter\def\csname sga@target@IV.6.3.4\endcsname{ega4.c04.cor.6-3-4}
\expandafter\def\csname sga@target@IV.6.3.4.1\endcsname{ega4.c04.eq.6-3-4-1}
\expandafter\def\csname sga@target@IV.6.3.4.2\endcsname{ega4.c04.eq.6-3-4-2}
\expandafter\def\csname sga@target@IV.6.3.5\endcsname{ega4.c04.cor.6-3-5}
\expandafter\def\csname sga@target@IV.6.3.6\endcsname{ega4.c04.cor.6-3-6}
\expandafter\def\csname sga@target@IV.6.3.7\endcsname{ega4.c04.cor.6-3-7}
\expandafter\def\csname sga@target@IV.6.3.8\endcsname{ega4.c04.prop.6-3-8}
\expandafter\def\csname sga@target@IV.6.4.1\endcsname{ega4.c04.prop.6-4-1}
\expandafter\def\csname sga@target@IV.6.4.2\endcsname{ega4.c04.cor.6-4-2}
\expandafter\def\csname sga@target@IV.6.4.3\endcsname{ega4.c04.rem.6-4-3}
\expandafter\def\csname sga@target@IV.6.5.1\endcsname{ega4.c04.prop.6-5-1}
\expandafter\def\csname sga@target@IV.6.5.2\endcsname{ega4.c04.cor.6-5-2}
\expandafter\def\csname sga@target@IV.6.5.3\endcsname{ega4.c04.prop.6-5-3}
\expandafter\def\csname sga@target@IV.6.5.4\endcsname{ega4.c04.cor.6-5-4}
\expandafter\def\csname sga@target@IV.6.5.5\endcsname{ega4.c04.item.6-5-5}
\expandafter\def\csname sga@target@IV.6.6.1\endcsname{ega4.c04.prop.6-6-1}
\expandafter\def\csname sga@target@IV.6.6.2\endcsname{ega4.c04.rem.6-6-2}
\expandafter\def\csname sga@target@IV.6.7.1\endcsname{ega4.c04.prop.6-7-1}
\expandafter\def\csname sga@target@IV.6.7.1.1\endcsname{ega4.c04.lem.6-7-1-1}
\expandafter\def\csname sga@target@IV.6.7.2\endcsname{ega4.c04.cor.6-7-2}
\expandafter\def\csname sga@target@IV.6.7.3\endcsname{ega4.c04.cor.6-7-3}
\expandafter\def\csname sga@target@IV.6.7.4\endcsname{ega4.c04.prop.6-7-4}
\expandafter\def\csname sga@target@IV.6.7.4.1\endcsname{ega4.c04.lem.6-7-4-1}
\expandafter\def\csname sga@target@IV.6.7.5\endcsname{ega4.c04.item.6-7-5}
\expandafter\def\csname sga@target@IV.6.7.6\endcsname{ega4.c04.def.6-7-6}
\expandafter\def\csname sga@target@IV.6.7.7\endcsname{ega4.c04.prop.6-7-7}
\expandafter\def\csname sga@target@IV.6.7.8\endcsname{ega4.c04.cor.6-7-8}
\expandafter\def\csname sga@target@IV.6.8.1\endcsname{ega4.c04.def.6-8-1}
\expandafter\def\csname sga@target@IV.6.8.2\endcsname{ega4.c04.prop.6-8-2}
\expandafter\def\csname sga@target@IV.6.8.3\endcsname{ega4.c04.prop.6-8-3}
\expandafter\def\csname sga@target@IV.6.8.4\endcsname{ega4.c04.item.6-8-4}
\expandafter\def\csname sga@target@IV.6.8.5\endcsname{ega4.c04.prop.6-8-5}
\expandafter\def\csname sga@target@IV.6.8.6\endcsname{ega4.c04.thm.6-8-6}
\expandafter\def\csname sga@target@IV.6.8.7\endcsname{ega4.c04.cor.6-8-7}
\expandafter\def\csname sga@target@IV.6.8.8\endcsname{ega4.c04.rem.6-8-8}
\expandafter\def\csname sga@target@IV.6.9.1\endcsname{ega4.c04.thm.6-9-1}
\expandafter\def\csname sga@target@IV.6.9.2\endcsname{ega4.c04.lem.6-9-2}
\expandafter\def\csname sga@target@IV.6.9.3\endcsname{ega4.c04.cor.6-9-3}
\expandafter\def\csname sga@target@IV.7.1.1\endcsname{ega4.c04.def.7-1-1}
\expandafter\def\csname sga@target@IV.7.1.10\endcsname{ega4.c04.prop.7-1-10}
\expandafter\def\csname sga@target@IV.7.1.11\endcsname{ega4.c04.prop.7-1-11}
\expandafter\def\csname sga@target@IV.7.1.12\endcsname{ega4.c04.item.7-1-12}
\expandafter\def\csname sga@target@IV.7.1.13\endcsname{ega4.c04.cor.7-1-13}
\expandafter\def\csname sga@target@IV.7.1.2\endcsname{ega4.c04.prop.7-1-2}
\expandafter\def\csname sga@target@IV.7.1.3\endcsname{ega4.c04.prop.7-1-3}
\expandafter\def\csname sga@target@IV.7.1.3.1\endcsname{ega4.c04.eq.7-1-3-1}
\expandafter\def\csname sga@target@IV.7.1.3.2\endcsname{ega4.c04.eq.7-1-3-2}
\expandafter\def\csname sga@target@IV.7.1.3.3\endcsname{ega4.c04.eq.7-1-3-3}
\expandafter\def\csname sga@target@IV.7.1.3.4\endcsname{ega4.c04.eq.7-1-3-4}
\expandafter\def\csname sga@target@IV.7.1.3.5\endcsname{ega4.c04.eq.7-1-3-5}
\expandafter\def\csname sga@target@IV.7.1.3.6\endcsname{ega4.c04.eq.7-1-3-6}
\expandafter\def\csname sga@target@IV.7.1.3.7\endcsname{ega4.c04.eq.7-1-3-7}
\expandafter\def\csname sga@target@IV.7.1.4\endcsname{ega4.c04.cor.7-1-4}
\expandafter\def\csname sga@target@IV.7.1.5\endcsname{ega4.c04.cor.7-1-5}
\expandafter\def\csname sga@target@IV.7.1.6\endcsname{ega4.c04.rem.7-1-6}
\expandafter\def\csname sga@target@IV.7.1.7\endcsname{ega4.c04.lem.7-1-7}
\expandafter\def\csname sga@target@IV.7.1.8\endcsname{ega4.c04.thm.7-1-8}
\expandafter\def\csname sga@target@IV.7.1.9\endcsname{ega4.c04.def.7-1-9}
\expandafter\def\csname sga@target@IV.7.2.1\endcsname{ega4.c04.statement.7-2-1}
\expandafter\def\csname sga@target@IV.7.2.1.1\endcsname{ega4.c04.statement.7-2-1-1}
\expandafter\def\csname sga@target@IV.7.2.10\endcsname{ega4.c04.statement.7-2-10}
\expandafter\def\csname sga@target@IV.7.2.2\endcsname{ega4.c04.prop.7-2-2}
\expandafter\def\csname sga@target@IV.7.2.3\endcsname{ega4.c04.prop.7-2-3}
\expandafter\def\csname sga@target@IV.7.2.4\endcsname{ega4.c04.prop.7-2-4}
\expandafter\def\csname sga@target@IV.7.2.5\endcsname{ega4.c04.thm.7-2-5}
\expandafter\def\csname sga@target@IV.7.2.5.1\endcsname{ega4.c04.eq.7-2-5-1}
\expandafter\def\csname sga@target@IV.7.2.6\endcsname{ega4.c04.def.7-2-6}
\expandafter\def\csname sga@target@IV.7.2.7\endcsname{ega4.c04.cor.7-2-7}
\expandafter\def\csname sga@target@IV.7.2.8\endcsname{ega4.c04.cor.7-2-8}
\expandafter\def\csname sga@target@IV.7.2.9\endcsname{ega4.c04.cor.7-2-9}
\expandafter\def\csname sga@target@IV.7.3.1\endcsname{ega4.c04.statement.7-3-1}
\expandafter\def\csname sga@target@IV.7.3.10\endcsname{ega4.c04.statement.7-3-10}
\expandafter\def\csname sga@target@IV.7.3.11\endcsname{ega4.c04.rem.7-3-11}
\expandafter\def\csname sga@target@IV.7.3.12\endcsname{ega4.c04.prop.7-3-12}
\expandafter\def\csname sga@target@IV.7.3.12.1\endcsname{ega4.c04.eq.7-3-12-1}
\expandafter\def\csname sga@target@IV.7.3.13\endcsname{ega4.c04.statement.7-3-13}
\expandafter\def\csname sga@target@IV.7.3.14\endcsname{ega4.c04.prop.7-3-14}
\expandafter\def\csname sga@target@IV.7.3.15\endcsname{ega4.c04.prop.7-3-15}
\expandafter\def\csname sga@target@IV.7.3.16\endcsname{ega4.c04.prop.7-3-16}
\expandafter\def\csname sga@target@IV.7.3.16.1\endcsname{ega4.c04.eq.7-3-16-1}
\expandafter\def\csname sga@target@IV.7.3.16.2\endcsname{ega4.c04.lem.7-3-16-2}
\expandafter\def\csname sga@target@IV.7.3.17\endcsname{ega4.c04.prop.7-3-17}
\expandafter\def\csname sga@target@IV.7.3.18\endcsname{ega4.c04.prop.7-3-18}
\expandafter\def\csname sga@target@IV.7.3.19\endcsname{ega4.c04.statement.7-3-19}
\expandafter\def\csname sga@target@IV.7.3.2\endcsname{ega4.c04.lem.7-3-2}
\expandafter\def\csname sga@target@IV.7.3.3\endcsname{ega4.c04.cor.7-3-3}
\expandafter\def\csname sga@target@IV.7.3.4\endcsname{ega4.c04.statement.7-3-4}
\expandafter\def\csname sga@target@IV.7.3.5\endcsname{ega4.c04.item.7-3-5}
\expandafter\def\csname sga@target@IV.7.3.6\endcsname{ega4.c04.statement.7-3-6}
\expandafter\def\csname sga@target@IV.7.3.7\endcsname{ega4.c04.lem.7-3-7}
\expandafter\def\csname sga@target@IV.7.3.8\endcsname{ega4.c04.statement.7-3-8}
\expandafter\def\csname sga@target@IV.7.3.9\endcsname{ega4.c04.item.7-3-9}
\expandafter\def\csname sga@target@IV.7.4.1\endcsname{ega4.c04.statement.7-4-1}
\expandafter\def\csname sga@target@IV.7.4.2\endcsname{ega4.c04.lem.7-4-2}
\expandafter\def\csname sga@target@IV.7.4.3\endcsname{ega4.c04.cor.7-4-3}
\expandafter\def\csname sga@target@IV.7.4.4\endcsname{ega4.c04.thm.7-4-4}
\expandafter\def\csname sga@target@IV.7.4.4.1\endcsname{ega4.c04.lem.7-4-4-1}
\expandafter\def\csname sga@target@IV.7.4.4.2\endcsname{ega4.c04.lem.7-4-4-2}
\expandafter\def\csname sga@target@IV.7.4.5\endcsname{ega4.c04.cor.7-4-5}
\expandafter\def\csname sga@target@IV.7.4.6\endcsname{ega4.c04.prop.7-4-6}
\expandafter\def\csname sga@target@IV.7.4.7\endcsname{ega4.c04.cor.7-4-7}
\expandafter\def\csname sga@target@IV.7.4.8\endcsname{ega4.c04.rem.7-4-8}
\expandafter\def\csname sga@target@IV.7.5.0\endcsname{ega4.c04.statement.7-5-0}
\expandafter\def\csname sga@target@IV.7.5.1\endcsname{ega4.c04.thm.7-5-1}
\expandafter\def\csname sga@target@IV.7.5.1.1\endcsname{ega4.c04.lem.7-5-1-1}
\expandafter\def\csname sga@target@IV.7.5.2\endcsname{ega4.c04.cor.7-5-2}
\expandafter\def\csname sga@target@IV.7.5.3\endcsname{ega4.c04.item.7-5-3}
\expandafter\def\csname sga@target@IV.7.5.4\endcsname{ega4.c04.item.7-5-4}
\expandafter\def\csname sga@target@IV.7.5.5\endcsname{ega4.c04.lem.7-5-5}
\expandafter\def\csname sga@target@IV.7.5.6\endcsname{ega4.c04.prop.7-5-6}
\expandafter\def\csname sga@target@IV.7.5.7\endcsname{ega4.c04.cor.7-5-7}
\expandafter\def\csname sga@target@IV.7.5.7.1\endcsname{ega4.c04.lem.7-5-7-1}
\expandafter\def\csname sga@target@IV.7.5.8\endcsname{ega4.c04.rem.7-5-8}
\expandafter\def\csname sga@target@IV.7.6.1\endcsname{ega4.c04.prop.7-6-1}
\expandafter\def\csname sga@target@IV.7.6.2\endcsname{ega4.c04.cor.7-6-2}
\expandafter\def\csname sga@target@IV.7.6.3\endcsname{ega4.c04.cor.7-6-3}
\expandafter\def\csname sga@target@IV.7.6.4\endcsname{ega4.c04.thm.7-6-4}
\expandafter\def\csname sga@target@IV.7.6.4.1\endcsname{ega4.c04.lem.7-6-4-1}
\expandafter\def\csname sga@target@IV.7.6.5\endcsname{ega4.c04.cor.7-6-5}
\expandafter\def\csname sga@target@IV.7.6.6\endcsname{ega4.c04.cor.7-6-6}
\expandafter\def\csname sga@target@IV.7.6.7\endcsname{ega4.c04.item.7-6-7}
\expandafter\def\csname sga@target@IV.7.7.1\endcsname{ega4.c04.statement.7-7-1}
\expandafter\def\csname sga@target@IV.7.7.2\endcsname{ega4.c04.thm.7-7-2}
\expandafter\def\csname sga@target@IV.7.7.3\endcsname{ega4.c04.cor.7-7-3}
\expandafter\def\csname sga@target@IV.7.7.4\endcsname{ega4.c04.cor.7-7-4}
\expandafter\def\csname sga@target@IV.7.8.1\endcsname{ega4.c04.statement.7-8-1}
\expandafter\def\csname sga@target@IV.7.8.2\endcsname{ega4.c04.def.7-8-2}
\expandafter\def\csname sga@target@IV.7.8.3\endcsname{ega4.c04.statement.7-8-3}
\expandafter\def\csname sga@target@IV.7.8.3.1\endcsname{ega4.c04.statement.7-8-3-1}
\expandafter\def\csname sga@target@IV.7.8.4\endcsname{ega4.c04.item.7-8-4}
\expandafter\def\csname sga@target@IV.7.8.4.1\endcsname{ega4.c04.statement.7-8-4-1}
\expandafter\def\csname sga@target@IV.7.8.5\endcsname{ega4.c04.statement.7-8-5}
\expandafter\def\csname sga@target@IV.7.8.6\endcsname{ega4.c04.prop.7-8-6}
\expandafter\def\csname sga@target@IV.7.9.1\endcsname{ega4.c04.statement.7-9-1}
\expandafter\def\csname sga@target@IV.7.9.10\endcsname{ega4.c04.item.7-9-10}
\expandafter\def\csname sga@target@IV.7.9.2\endcsname{ega4.c04.statement.7-9-2}
\expandafter\def\csname sga@target@IV.7.9.3\endcsname{ega4.c04.prop.7-9-3}
\expandafter\def\csname sga@target@IV.7.9.3.1\endcsname{ega4.c04.lem.7-9-3-1}
\expandafter\def\csname sga@target@IV.7.9.4\endcsname{ega4.c04.cor.7-9-4}
\expandafter\def\csname sga@target@IV.7.9.5\endcsname{ega4.c04.prop.7-9-5}
\expandafter\def\csname sga@target@IV.7.9.6\endcsname{ega4.c04.rem.7-9-6}
\expandafter\def\csname sga@target@IV.7.9.7\endcsname{ega4.c04.statement.7-9-7}
\expandafter\def\csname sga@target@IV.7.9.8\endcsname{ega4.c04.prop.7-9-8}
\expandafter\def\csname sga@target@IV.7.9.8.1\endcsname{ega4.c04.statement.7-9-8-1}
\expandafter\def\csname sga@target@IV.7.9.9\endcsname{ega4.c04.cor.7-9-9}
\expandafter\def\csname sga@target@IV.8.1.1\endcsname{ega4.c04.statement.8-1-1}
\expandafter\def\csname sga@target@IV.8.1.2\endcsname{ega4.c04.statement.8-1-2}
\expandafter\def\csname sga@target@IV.8.1.3\endcsname{ega4.c04.statement.8-1-3}
\expandafter\def\csname sga@target@IV.8.10.1\endcsname{ega4.c04.prop.8-10-1}
\expandafter\def\csname sga@target@IV.8.10.2\endcsname{ega4.c04.cor.8-10-2}
\expandafter\def\csname sga@target@IV.8.10.3\endcsname{ega4.c04.prop.8-10-3}
\expandafter\def\csname sga@target@IV.8.10.3.1\endcsname{ega4.c04.rem.8-10-3-1}
\expandafter\def\csname sga@target@IV.8.10.4\endcsname{ega4.c04.prop.8-10-4}
\expandafter\def\csname sga@target@IV.8.10.5\endcsname{ega4.c04.thm.8-10-5}
\expandafter\def\csname sga@target@IV.8.10.5.1\endcsname{ega4.c04.prop.8-10-5-1}
\expandafter\def\csname sga@target@IV.8.10.5.2\endcsname{ega4.c04.lem.8-10-5-2}
\expandafter\def\csname sga@target@IV.8.11\endcsname{ega4.c04.subsec.8-11.key-iv-8-11}
\expandafter\def\csname sga@target@IV.8.11.1\endcsname{ega4.c04.thm.8-11-1}
\expandafter\def\csname sga@target@IV.8.11.2\endcsname{ega4.c04.thm.8-11-2}
\expandafter\def\csname sga@target@IV.8.11.3\endcsname{ega4.c04.rem.8-11-3}
\expandafter\def\csname sga@target@IV.8.11.4\endcsname{ega4.c04.statement.8-11-4}
\expandafter\def\csname sga@target@IV.8.11.5\endcsname{ega4.c04.prop.8-11-5}
\expandafter\def\csname sga@target@IV.8.11.5.1\endcsname{ega4.c04.rem.8-11-5-1}
\expandafter\def\csname sga@target@IV.8.11.6\endcsname{ega4.c04.prop.8-11-6}
\expandafter\def\csname sga@target@IV.8.12\endcsname{ega4.c04.subsec.8-12.key-iv-8-12}
\expandafter\def\csname sga@target@IV.8.12.1\endcsname{ega4.c04.lem.8-12-1}
\expandafter\def\csname sga@target@IV.8.12.10\endcsname{ega4.c04.cor.8-12-10}
\expandafter\def\csname sga@target@IV.8.12.10.1\endcsname{ega4.c04.lem.8-12-10-1}
\expandafter\def\csname sga@target@IV.8.12.11\endcsname{ega4.c04.cor.8-12-11}
\expandafter\def\csname sga@target@IV.8.12.12\endcsname{ega4.c04.cor.8-12-12}
\expandafter\def\csname sga@target@IV.8.12.2\endcsname{ega4.c04.lem.8-12-2}
\expandafter\def\csname sga@target@IV.8.12.3\endcsname{ega4.c04.prop.8-12-3}
\expandafter\def\csname sga@target@IV.8.12.4\endcsname{ega4.c04.cor.8-12-4}
\expandafter\def\csname sga@target@IV.8.12.5\endcsname{ega4.c04.statement.8-12-5}
\expandafter\def\csname sga@target@IV.8.12.6\endcsname{ega4.c04.thm.8-12-6}
\expandafter\def\csname sga@target@IV.8.12.6.1\endcsname{ega4.c04.eq.8-12-6-1}
\expandafter\def\csname sga@target@IV.8.12.7\endcsname{ega4.c04.rem.8-12-7}
\expandafter\def\csname sga@target@IV.8.12.8\endcsname{ega4.c04.cor.8-12-8}
\expandafter\def\csname sga@target@IV.8.12.9\endcsname{ega4.c04.cor.8-12-9}
\expandafter\def\csname sga@target@IV.8.13.1\endcsname{ega4.c04.prop.8-13-1}
\expandafter\def\csname sga@target@IV.8.13.1.1\endcsname{ega4.c04.eq.8-13-1-1}
\expandafter\def\csname sga@target@IV.8.13.2\endcsname{ega4.c04.cor.8-13-2}
\expandafter\def\csname sga@target@IV.8.13.2.1\endcsname{ega4.c04.eq.8-13-2-1}
\expandafter\def\csname sga@target@IV.8.13.3\endcsname{ega4.c04.statement.8-13-3}
\expandafter\def\csname sga@target@IV.8.13.4\endcsname{ega4.c04.statement.8-13-4}
\expandafter\def\csname sga@target@IV.8.13.4.1\endcsname{ega4.c04.eq.8-13-4-1}
\expandafter\def\csname sga@target@IV.8.13.4.2\endcsname{ega4.c04.eq.8-13-4-2}
\expandafter\def\csname sga@target@IV.8.13.5\endcsname{ega4.c04.prop.8-13-5}
\expandafter\def\csname sga@target@IV.8.13.6\endcsname{ega4.c04.rem.8-13-6}
\expandafter\def\csname sga@target@IV.8.14.1\endcsname{ega4.c04.statement.8-14-1}
\expandafter\def\csname sga@target@IV.8.14.2\endcsname{ega4.c04.prop.8-14-2}
\expandafter\def\csname sga@target@IV.8.14.2.1\endcsname{ega4.c04.eq.8-14-2-1}
\expandafter\def\csname sga@target@IV.8.14.2.2\endcsname{ega4.c04.cor.8-14-2-2}
\expandafter\def\csname sga@target@IV.8.14.2.3\endcsname{ega4.c04.eq.8-14-2-3}
\expandafter\def\csname sga@target@IV.8.14.2.4\endcsname{ega4.c04.eq.8-14-2-4}
\expandafter\def\csname sga@target@IV.8.14.3\endcsname{ega4.c04.rem.8-14-3}
\expandafter\def\csname sga@target@IV.8.2.1\endcsname{ega4.c04.statement.8-2-1}
\expandafter\def\csname sga@target@IV.8.2.1.1\endcsname{ega4.c04.eq.8-2-1-1}
\expandafter\def\csname sga@target@IV.8.2.10\endcsname{ega4.c04.cor.8-2-10}
\expandafter\def\csname sga@target@IV.8.2.11\endcsname{ega4.c04.cor.8-2-11}
\expandafter\def\csname sga@target@IV.8.2.12\endcsname{ega4.c04.cor.8-2-12}
\expandafter\def\csname sga@target@IV.8.2.12.1\endcsname{ega4.c04.eq.8-2-12-1}
\expandafter\def\csname sga@target@IV.8.2.13\endcsname{ega4.c04.prop.8-2-13}
\expandafter\def\csname sga@target@IV.8.2.14\endcsname{ega4.c04.rem.8-2-14}
\expandafter\def\csname sga@target@IV.8.2.2\endcsname{ega4.c04.statement.8-2-2}
\expandafter\def\csname sga@target@IV.8.2.3\endcsname{ega4.c04.prop.8-2-3}
\expandafter\def\csname sga@target@IV.8.2.3.1\endcsname{ega4.c04.eq.8-2-3-1}
\expandafter\def\csname sga@target@IV.8.2.4\endcsname{ega4.c04.lem.8-2-4}
\expandafter\def\csname sga@target@IV.8.2.5\endcsname{ega4.c04.prop.8-2-5}
\expandafter\def\csname sga@target@IV.8.2.6\endcsname{ega4.c04.lem.8-2-6}
\expandafter\def\csname sga@target@IV.8.2.7\endcsname{ega4.c04.rem.8-2-7}
\expandafter\def\csname sga@target@IV.8.2.8\endcsname{ega4.c04.statement.8-2-8}
\expandafter\def\csname sga@target@IV.8.2.9\endcsname{ega4.c04.prop.8-2-9}
\expandafter\def\csname sga@target@IV.8.3.1\endcsname{ega4.c04.statement.8-3-1}
\expandafter\def\csname sga@target@IV.8.3.10\endcsname{ega4.c04.statement.8-3-10}
\expandafter\def\csname sga@target@IV.8.3.10.1\endcsname{ega4.c04.eq.8-3-10-1}
\expandafter\def\csname sga@target@IV.8.3.11\endcsname{ega4.c04.thm.8-3-11}
\expandafter\def\csname sga@target@IV.8.3.12\endcsname{ega4.c04.cor.8-3-12}
\expandafter\def\csname sga@target@IV.8.3.13\endcsname{ega4.c04.prop.8-3-13}
\expandafter\def\csname sga@target@IV.8.3.2\endcsname{ega4.c04.thm.8-3-2}
\expandafter\def\csname sga@target@IV.8.3.2.1\endcsname{ega4.c04.eq.8-3-2-1}
\expandafter\def\csname sga@target@IV.8.3.3\endcsname{ega4.c04.cor.8-3-3}
\expandafter\def\csname sga@target@IV.8.3.4\endcsname{ega4.c04.cor.8-3-4}
\expandafter\def\csname sga@target@IV.8.3.5\endcsname{ega4.c04.cor.8-3-5}
\expandafter\def\csname sga@target@IV.8.3.6\endcsname{ega4.c04.cor.8-3-6}
\expandafter\def\csname sga@target@IV.8.3.7\endcsname{ega4.c04.cor.8-3-7}
\expandafter\def\csname sga@target@IV.8.3.7.1\endcsname{ega4.c04.eq.8-3-7-1}
\expandafter\def\csname sga@target@IV.8.3.8\endcsname{ega4.c04.prop.8-3-8}
\expandafter\def\csname sga@target@IV.8.3.9\endcsname{ega4.c04.statement.8-3-9}
\expandafter\def\csname sga@target@IV.8.3.9.1\endcsname{ega4.c04.eq.8-3-9-1}
\expandafter\def\csname sga@target@IV.8.3.9.2\endcsname{ega4.c04.eq.8-3-9-2}
\expandafter\def\csname sga@target@IV.8.3.9.3\endcsname{ega4.c04.eq.8-3-9-3}
\expandafter\def\csname sga@target@IV.8.3.9.4\endcsname{ega4.c04.eq.8-3-9-4}
\expandafter\def\csname sga@target@IV.8.3.9.5\endcsname{ega4.c04.eq.8-3-9-5}
\expandafter\def\csname sga@target@IV.8.4.1\endcsname{ega4.c04.prop.8-4-1}
\expandafter\def\csname sga@target@IV.8.4.2\endcsname{ega4.c04.prop.8-4-2}
\expandafter\def\csname sga@target@IV.8.4.2.1\endcsname{ega4.c04.lem.8-4-2-1}
\expandafter\def\csname sga@target@IV.8.4.3\endcsname{ega4.c04.cor.8-4-3}
\expandafter\def\csname sga@target@IV.8.4.4\endcsname{ega4.c04.prop.8-4-4}
\expandafter\def\csname sga@target@IV.8.4.5\endcsname{ega4.c04.cor.8-4-5}
\expandafter\def\csname sga@target@IV.8.4.6\endcsname{ega4.c04.rem.8-4-6}
\expandafter\def\csname sga@target@IV.8.5.1\endcsname{ega4.c04.statement.8-5-1}
\expandafter\def\csname sga@target@IV.8.5.1.1\endcsname{ega4.c04.eq.8-5-1-1}
\expandafter\def\csname sga@target@IV.8.5.1.2\endcsname{ega4.c04.eq.8-5-1-2}
\expandafter\def\csname sga@target@IV.8.5.1.3\endcsname{ega4.c04.eq.8-5-1-3}
\expandafter\def\csname sga@target@IV.8.5.1.4\endcsname{ega4.c04.eq.8-5-1-4}
\expandafter\def\csname sga@target@IV.8.5.10\endcsname{ega4.c04.statement.8-5-10}
\expandafter\def\csname sga@target@IV.8.5.10.1\endcsname{ega4.c04.eq.8-5-10-1}
\expandafter\def\csname sga@target@IV.8.5.10.2\endcsname{ega4.c04.eq.8-5-10-2}
\expandafter\def\csname sga@target@IV.8.5.11\endcsname{ega4.c04.prop.8-5-11}
\expandafter\def\csname sga@target@IV.8.5.2\endcsname{ega4.c04.thm.8-5-2}
\expandafter\def\csname sga@target@IV.8.5.2.1\endcsname{ega4.c04.lem.8-5-2-1}
\expandafter\def\csname sga@target@IV.8.5.2.2\endcsname{ega4.c04.eq.8-5-2-2}
\expandafter\def\csname sga@target@IV.8.5.2.3\endcsname{ega4.c04.eq.8-5-2-3}
\expandafter\def\csname sga@target@IV.8.5.2.4\endcsname{ega4.c04.cor.8-5-2-4}
\expandafter\def\csname sga@target@IV.8.5.2.5\endcsname{ega4.c04.cor.8-5-2-5}
\expandafter\def\csname sga@target@IV.8.5.3\endcsname{ega4.c04.statement.8-5-3}
\expandafter\def\csname sga@target@IV.8.5.4\endcsname{ega4.c04.cor.8-5-4}
\expandafter\def\csname sga@target@IV.8.5.5\endcsname{ega4.c04.prop.8-5-5}
\expandafter\def\csname sga@target@IV.8.5.6\endcsname{ega4.c04.prop.8-5-6}
\expandafter\def\csname sga@target@IV.8.5.7\endcsname{ega4.c04.cor.8-5-7}
\expandafter\def\csname sga@target@IV.8.5.8\endcsname{ega4.c04.cor.8-5-8}
\expandafter\def\csname sga@target@IV.8.5.9\endcsname{ega4.c04.lem.8-5-9}
\expandafter\def\csname sga@target@IV.8.6.1\endcsname{ega4.c04.statement.8-6-1}
\expandafter\def\csname sga@target@IV.8.6.2\endcsname{ega4.c04.statement.8-6-2}
\expandafter\def\csname sga@target@IV.8.6.2.1\endcsname{ega4.c04.eq.8-6-2-1}
\expandafter\def\csname sga@target@IV.8.6.2.2\endcsname{ega4.c04.eq.8-6-2-2}
\expandafter\def\csname sga@target@IV.8.6.2.3\endcsname{ega4.c04.eq.8-6-2-3}
\expandafter\def\csname sga@target@IV.8.6.3\endcsname{ega4.c04.prop.8-6-3}
\expandafter\def\csname sga@target@IV.8.6.3.1\endcsname{ega4.c04.cor.8-6-3-1}
\expandafter\def\csname sga@target@IV.8.6.4\endcsname{ega4.c04.cor.8-6-4}
\expandafter\def\csname sga@target@IV.8.7.1\endcsname{ega4.c04.prop.8-7-1}
\expandafter\def\csname sga@target@IV.8.7.2\endcsname{ega4.c04.prop.8-7-2}
\expandafter\def\csname sga@target@IV.8.7.3\endcsname{ega4.c04.cor.8-7-3}
\expandafter\def\csname sga@target@IV.8.7.4\endcsname{ega4.c04.rem.8-7-4}
\expandafter\def\csname sga@target@IV.8.8.1\endcsname{ega4.c04.statement.8-8-1}
\expandafter\def\csname sga@target@IV.8.8.1.1\endcsname{ega4.c04.eq.8-8-1-1}
\expandafter\def\csname sga@target@IV.8.8.2\endcsname{ega4.c04.thm.8-8-2}
\expandafter\def\csname sga@target@IV.8.8.2.1\endcsname{ega4.c04.lem.8-8-2-1}
\expandafter\def\csname sga@target@IV.8.8.2.2\endcsname{ega4.c04.lem.8-8-2-2}
\expandafter\def\csname sga@target@IV.8.8.2.3\endcsname{ega4.c04.lem.8-8-2-3}
\expandafter\def\csname sga@target@IV.8.8.2.4\endcsname{ega4.c04.cor.8-8-2-4}
\expandafter\def\csname sga@target@IV.8.8.2.5\endcsname{ega4.c04.cor.8-8-2-5}
\expandafter\def\csname sga@target@IV.8.8.3\endcsname{ega4.c04.statement.8-8-3}
\expandafter\def\csname sga@target@IV.8.9.1\endcsname{ega4.c04.prop.8-9-1}
\expandafter\def\csname sga@target@IV.8.9.2\endcsname{ega4.c04.statement.8-9-2}
\expandafter\def\csname sga@target@IV.8.9.3\endcsname{ega4.c04.prop.8-9-3}
\expandafter\def\csname sga@target@IV.8.9.4\endcsname{ega4.c04.prop.8-9-4}
\expandafter\def\csname sga@target@IV.8.9.4.1\endcsname{ega4.c04.lem.8-9-4-1}
\expandafter\def\csname sga@target@IV.8.9.5\endcsname{ega4.c04.cor.8-9-5}
\expandafter\def\csname sga@target@IV.9.1.1\endcsname{ega4.c04.prop.9-1-1}
\expandafter\def\csname sga@target@IV.9.1.2\endcsname{ega4.c04.cor.9-1-2}
\expandafter\def\csname sga@target@IV.9.1.3\endcsname{ega4.c04.rem.9-1-3}
\expandafter\def\csname sga@target@IV.9.1.4\endcsname{ega4.c04.prop.9-1-4}
\expandafter\def\csname sga@target@IV.9.1.5\endcsname{ega4.c04.cor.9-1-5}
\expandafter\def\csname sga@target@IV.9.1.6\endcsname{ega4.c04.prop.9-1-6}
\expandafter\def\csname sga@target@IV.9.2\endcsname{ega4.c04.subsec.9-2.key-iv-9-2}
\expandafter\def\csname sga@target@IV.9.2.1\endcsname{ega4.c04.def.9-2-1}
\expandafter\def\csname sga@target@IV.9.2.2\endcsname{ega4.c04.item.9-2-2}
\expandafter\def\csname sga@target@IV.9.2.3\endcsname{ega4.c04.prop.9-2-3}
\expandafter\def\csname sga@target@IV.9.2.4\endcsname{ega4.c04.cor.9-2-4}
\expandafter\def\csname sga@target@IV.9.2.5\endcsname{ega4.c04.rem.9-2-5}
\expandafter\def\csname sga@target@IV.9.2.6\endcsname{ega4.c04.prop.9-2-6}
\expandafter\def\csname sga@target@IV.9.2.6.1\endcsname{ega4.c04.prop.9-2-6-1}
\expandafter\def\csname sga@target@IV.9.2.6.2\endcsname{ega4.c04.cor.9-2-6-2}
\expandafter\def\csname sga@target@IV.9.3.1\endcsname{ega4.c04.prop.9-3-1}
\expandafter\def\csname sga@target@IV.9.3.2\endcsname{ega4.c04.prop.9-3-2}
\expandafter\def\csname sga@target@IV.9.3.3\endcsname{ega4.c04.prop.9-3-3}
\expandafter\def\csname sga@target@IV.9.3.4\endcsname{ega4.c04.rem.9-3-4}
\expandafter\def\csname sga@target@IV.9.3.5\endcsname{ega4.c04.prop.9-3-5}
\expandafter\def\csname sga@target@IV.9.3.6\endcsname{ega4.c04.statement.9-3-6}
\expandafter\def\csname sga@target@IV.9.4.1\endcsname{ega4.c04.statement.9-4-1}
\expandafter\def\csname sga@target@IV.9.4.2\endcsname{ega4.c04.prop.9-4-2}
\expandafter\def\csname sga@target@IV.9.4.3\endcsname{ega4.c04.cor.9-4-3}
\expandafter\def\csname sga@target@IV.9.4.3.1\endcsname{ega4.c04.eq.9-4-3-1}
\expandafter\def\csname sga@target@IV.9.4.4\endcsname{ega4.c04.prop.9-4-4}
\expandafter\def\csname sga@target@IV.9.4.5\endcsname{ega4.c04.cor.9-4-5}
\expandafter\def\csname sga@target@IV.9.4.6\endcsname{ega4.c04.cor.9-4-6}
\expandafter\def\csname sga@target@IV.9.4.7\endcsname{ega4.c04.prop.9-4-7}
\expandafter\def\csname sga@target@IV.9.4.7.1\endcsname{ega4.c04.lem.9-4-7-1}
\expandafter\def\csname sga@target@IV.9.4.7.2\endcsname{ega4.c04.rem.9-4-7-2}
\expandafter\def\csname sga@target@IV.9.4.8\endcsname{ega4.c04.prop.9-4-8}
\expandafter\def\csname sga@target@IV.9.4.9\endcsname{ega4.c04.rem.9-4-9}
\expandafter\def\csname sga@target@IV.9.5.1\endcsname{ega4.c04.prop.9-5-1}
\expandafter\def\csname sga@target@IV.9.5.2\endcsname{ega4.c04.cor.9-5-2}
\expandafter\def\csname sga@target@IV.9.5.3\endcsname{ega4.c04.prop.9-5-3}
\expandafter\def\csname sga@target@IV.9.5.3.case2\endcsname{ega4.c04.item.9-5-3-case2}
\expandafter\def\csname sga@target@IV.9.5.4\endcsname{ega4.c04.cor.9-5-4}
\expandafter\def\csname sga@target@IV.9.5.5\endcsname{ega4.c04.prop.9-5-5}
\expandafter\def\csname sga@target@IV.9.5.6\endcsname{ega4.c04.cor.9-5-6}
\expandafter\def\csname sga@target@IV.9.5.7\endcsname{ega4.c04.rem.9-5-7}
\expandafter\def\csname sga@target@IV.9.6.1\endcsname{ega4.c04.prop.9-6-1}
\expandafter\def\csname sga@target@IV.9.6.2\endcsname{ega4.c04.prop.9-6-2}
\expandafter\def\csname sga@target@IV.9.6.3\endcsname{ega4.c04.prop.9-6-3}
\expandafter\def\csname sga@target@IV.9.6.3.1\endcsname{ega4.c04.eq.9-6-3-1}
\expandafter\def\csname sga@target@IV.9.6.4\endcsname{ega4.c04.cor.9-6-4}
\expandafter\def\csname sga@target@IV.9.6.5\endcsname{ega4.c04.cor.9-6-5}
\expandafter\def\csname sga@target@IV.9.6.6\endcsname{ega4.c04.sec.9-6-6}
\expandafter\def\csname sga@target@IV.9.6.7\endcsname{ega4.c04.sec.9-6-7}
\expandafter\def\csname sga@target@IV.9.7\endcsname{ega4.c04.subsec.9-7.key-iv-9-7}
\expandafter\def\csname sga@target@IV.9.7.1\endcsname{ega4.c04.lem.9-7-1}
\expandafter\def\csname sga@target@IV.9.7.10\endcsname{ega4.c04.cor.9-7-10}
\expandafter\def\csname sga@target@IV.9.7.11\endcsname{ega4.c04.rem.9-7-11}
\expandafter\def\csname sga@target@IV.9.7.12\endcsname{ega4.c04.prop.9-7-12}
\expandafter\def\csname sga@target@IV.9.7.2\endcsname{ega4.c04.lem.9-7-2}
\expandafter\def\csname sga@target@IV.9.7.3\endcsname{ega4.c04.subsec.9-7-3}
\expandafter\def\csname sga@target@IV.9.7.4\endcsname{ega4.c04.lem.9-7-4}
\expandafter\def\csname sga@target@IV.9.7.5\endcsname{ega4.c04.lem.9-7-5}
\expandafter\def\csname sga@target@IV.9.7.5.1\endcsname{ega4.c04.lem.9-7-5-1}
\expandafter\def\csname sga@target@IV.9.7.6\endcsname{ega4.c04.prop.9-7-6}
\expandafter\def\csname sga@target@IV.9.7.7\endcsname{ega4.c04.thm.9-7-7}
\expandafter\def\csname sga@target@IV.9.7.8\endcsname{ega4.c04.prop.9-7-8}
\expandafter\def\csname sga@target@IV.9.7.9\endcsname{ega4.c04.cor.9-7-9}
\expandafter\def\csname sga@target@IV.9.8\endcsname{ega4.c04.subsec.9-8.key-iv-9-8}
\expandafter\def\csname sga@target@IV.9.8.1\endcsname{ega4.c04.prop.9-8-1}
\expandafter\def\csname sga@target@IV.9.8.1-proof\endcsname{ega4.c04.subsec.9-8-1-proof}
\expandafter\def\csname sga@target@IV.9.8.2\endcsname{ega4.c04.subsec.9-8-2}
\expandafter\def\csname sga@target@IV.9.8.3\endcsname{ega4.c04.prop.9-8-3}
\expandafter\def\csname sga@target@IV.9.8.3-proof\endcsname{ega4.c04.subsec.9-8-3-proof}
\expandafter\def\csname sga@target@IV.9.8.4\endcsname{ega4.c04.cor.9-8-4}
\expandafter\def\csname sga@target@IV.9.8.4-proof\endcsname{ega4.c04.item.9-8-4-proof}
\expandafter\def\csname sga@target@IV.9.8.5\endcsname{ega4.c04.prop.9-8-5}
\expandafter\def\csname sga@target@IV.9.8.5-proof\endcsname{ega4.c04.item.9-8-5-proof}
\expandafter\def\csname sga@target@IV.9.8.6\endcsname{ega4.c04.prop.9-8-6}
\expandafter\def\csname sga@target@IV.9.8.6-proof\endcsname{ega4.c04.item.9-8-6-proof}
\expandafter\def\csname sga@target@IV.9.8.7\endcsname{ega4.c04.cor.9-8-7}
\expandafter\def\csname sga@target@IV.9.8.7-proof\endcsname{ega4.c04.item.9-8-7-proof}
\expandafter\def\csname sga@target@IV.9.8.8\endcsname{ega4.c04.prop.9-8-8}
\expandafter\def\csname sga@target@IV.9.8.8-proof\endcsname{ega4.c04.item.9-8-8-proof}
\expandafter\def\csname sga@target@IV.9.8.9\endcsname{ega4.c04.rem.9-8-9}
\expandafter\def\csname sga@target@IV.9.8.9.1\endcsname{ega4.c04.item.9-8-9-1}
\expandafter\def\csname sga@target@IV.9.9\endcsname{ega4.c04.subsec.9-9.key-iv-9-9}
\expandafter\def\csname sga@target@IV.9.9.1\endcsname{ega4.c04.prop.9-9-1}
\expandafter\def\csname sga@target@IV.9.9.1-proof\endcsname{ega4.c04.item.9-9-1-proof}
\expandafter\def\csname sga@target@IV.9.9.1.1\endcsname{ega4.c04.eq.9-9-1-1}
\expandafter\def\csname sga@target@IV.9.9.2\endcsname{ega4.c04.prop.9-9-2}
\expandafter\def\csname sga@target@IV.9.9.2-proof\endcsname{ega4.c04.item.9-9-2-proof}
\expandafter\def\csname sga@target@IV.9.9.2.1\endcsname{ega4.c04.eq.9-9-2-1}
\expandafter\def\csname sga@target@IV.9.9.2.2\endcsname{ega4.c04.eq.9-9-2-2}
\expandafter\def\csname sga@target@IV.9.9.3\endcsname{ega4.c04.cor.9-9-3}
\expandafter\def\csname sga@target@IV.9.9.3-proof\endcsname{ega4.c04.statement.9-9-3-proof}
\expandafter\def\csname sga@target@IV.9.9.4\endcsname{ega4.c04.prop.9-9-4}
\expandafter\def\csname sga@target@IV.9.9.4-proof\endcsname{ega4.c04.item.9-9-4-proof}
\expandafter\def\csname sga@target@IV.9.9.5\endcsname{ega4.c04.cor.9-9-5}
\expandafter\def\csname sga@target@IV.9.9.5-proof\endcsname{ega4.c04.item.9-9-5-proof}
\expandafter\def\csname sga@target@IV.9.9.6\endcsname{ega4.c04.prop.9-9-6}
\expandafter\def\csname sga@target@IV.9.9.6-proof\endcsname{ega4.c04.item.9-9-6-proof}
\expandafter\def\csname sga@target@section:IV.1\endcsname{ega4.c04.sec.1}
\expandafter\def\csname sga@target@section:IV.10\endcsname{ega4.c04.sec.10}
\expandafter\def\csname sga@target@section:IV.11\endcsname{ega4.c04.sec.11}
\expandafter\def\csname sga@target@section:IV.12\endcsname{ega4.c04.sec.12}
\expandafter\def\csname sga@target@section:IV.13\endcsname{ega4.c04.sec.13}
\expandafter\def\csname sga@target@section:IV.14\endcsname{ega4.c04.sec.14}
\expandafter\def\csname sga@target@section:IV.15\endcsname{ega4.c04.sec.15}
\expandafter\def\csname sga@target@section:IV.16\endcsname{ega4.c04.sec.16}
\expandafter\def\csname sga@target@section:IV.17\endcsname{ega4.c04.sec.17}
\expandafter\def\csname sga@target@section:IV.18\endcsname{ega4.c04.sec.18}
\expandafter\def\csname sga@target@section:IV.19\endcsname{ega4.c04.sec.19}
\expandafter\def\csname sga@target@section:IV.2\endcsname{ega4.c04.sec.2}
\expandafter\def\csname sga@target@section:IV.21\endcsname{ega4.c04.sec.21}
\expandafter\def\csname sga@target@section:IV.3\endcsname{ega4.c04.sec.3}
\expandafter\def\csname sga@target@section:IV.4\endcsname{ega4.c04.sec.4}
\expandafter\def\csname sga@target@section:IV.5\endcsname{ega4.c04.sec.5}
\expandafter\def\csname sga@target@section:IV.6\endcsname{ega4.c04.sec.6}
\expandafter\def\csname sga@target@section:IV.7\endcsname{ega4.c04.sec.7}
\expandafter\def\csname sga@target@section:IV.8\endcsname{ega4.c04.sec.8}
\expandafter\def\csname sga@target@section:IV.9\endcsname{ega4.c04.sec.9}
\expandafter\def\csname sga@target@subsection:IV.1.1\endcsname{ega4.c04.subsec.1-1}
\expandafter\def\csname sga@target@subsection:IV.1.10\endcsname{ega4.c04.subsec.1-10}
\expandafter\def\csname sga@target@subsection:IV.1.2\endcsname{ega4.c04.subsec.1-2}
\expandafter\def\csname sga@target@subsection:IV.1.3\endcsname{ega4.c04.subsec.1-3}
\expandafter\def\csname sga@target@subsection:IV.1.4\endcsname{ega4.c04.subsec.1-4}
\expandafter\def\csname sga@target@subsection:IV.1.5\endcsname{ega4.c04.subsec.1-5}
\expandafter\def\csname sga@target@subsection:IV.1.6\endcsname{ega4.c04.subsec.1-6}
\expandafter\def\csname sga@target@subsection:IV.1.7\endcsname{ega4.c04.subsec.1-7}
\expandafter\def\csname sga@target@subsection:IV.1.8\endcsname{ega4.c04.subsec.1-8}
\expandafter\def\csname sga@target@subsection:IV.1.9\endcsname{ega4.c04.subsec.1-9}
\expandafter\def\csname sga@target@subsection:IV.10.1\endcsname{ega4.c04.subsec.10-1}
\expandafter\def\csname sga@target@subsection:IV.10.10\endcsname{ega4.c04.subsec.10-10}
\expandafter\def\csname sga@target@subsection:IV.10.2\endcsname{ega4.c04.subsec.10-2}
\expandafter\def\csname sga@target@subsection:IV.10.3\endcsname{ega4.c04.subsec.10-3}
\expandafter\def\csname sga@target@subsection:IV.10.4\endcsname{ega4.c04.subsec.10-4}
\expandafter\def\csname sga@target@subsection:IV.10.5\endcsname{ega4.c04.subsec.10-5}
\expandafter\def\csname sga@target@subsection:IV.10.6\endcsname{ega4.c04.subsec.10-6}
\expandafter\def\csname sga@target@subsection:IV.10.7\endcsname{ega4.c04.subsec.10-7}
\expandafter\def\csname sga@target@subsection:IV.10.8\endcsname{ega4.c04.subsec.10-8}
\expandafter\def\csname sga@target@subsection:IV.10.9\endcsname{ega4.c04.subsec.10-9}
\expandafter\def\csname sga@target@subsection:IV.11.1\endcsname{ega4.c04.subsec.11-1}
\expandafter\def\csname sga@target@subsection:IV.11.10\endcsname{ega4.c04.subsec.11-10}
\expandafter\def\csname sga@target@subsection:IV.11.2\endcsname{ega4.c04.subsec.11-2}
\expandafter\def\csname sga@target@subsection:IV.11.3\endcsname{ega4.c04.subsec.11-3}
\expandafter\def\csname sga@target@subsection:IV.11.4\endcsname{ega4.c04.subsec.11-4}
\expandafter\def\csname sga@target@subsection:IV.11.5\endcsname{ega4.c04.subsec.11-5}
\expandafter\def\csname sga@target@subsection:IV.11.6\endcsname{ega4.c04.subsec.11-6}
\expandafter\def\csname sga@target@subsection:IV.11.7\endcsname{ega4.c04.subsec.11-7}
\expandafter\def\csname sga@target@subsection:IV.11.8\endcsname{ega4.c04.subsec.11-8}
\expandafter\def\csname sga@target@subsection:IV.11.9\endcsname{ega4.c04.subsec.11-9}
\expandafter\def\csname sga@target@subsection:IV.12.0\endcsname{ega4.c04.subsec.12-0}
\expandafter\def\csname sga@target@subsection:IV.12.1\endcsname{ega4.c04.subsec.12-1}
\expandafter\def\csname sga@target@subsection:IV.12.2\endcsname{ega4.c04.subsec.12-2}
\expandafter\def\csname sga@target@subsection:IV.12.3\endcsname{ega4.c04.subsec.12-3}
\expandafter\def\csname sga@target@subsection:IV.13.1\endcsname{ega4.c04.subsec.13-1}
\expandafter\def\csname sga@target@subsection:IV.13.2\endcsname{ega4.c04.subsec.13-2}
\expandafter\def\csname sga@target@subsection:IV.13.3\endcsname{ega4.c04.subsec.13-3}
\expandafter\def\csname sga@target@subsection:IV.14.1\endcsname{ega4.c04.subsec.14-1}
\expandafter\def\csname sga@target@subsection:IV.14.2\endcsname{ega4.c04.subsec.14-2}
\expandafter\def\csname sga@target@subsection:IV.14.3\endcsname{ega4.c04.subsec.14-3}
\expandafter\def\csname sga@target@subsection:IV.14.4\endcsname{ega4.c04.subsec.14-4}
\expandafter\def\csname sga@target@subsection:IV.14.5\endcsname{ega4.c04.subsec.14-5}
\expandafter\def\csname sga@target@subsection:IV.15.1\endcsname{ega4.c04.subsec.15-1}
\expandafter\def\csname sga@target@subsection:IV.15.2\endcsname{ega4.c04.subsec.15-2}
\expandafter\def\csname sga@target@subsection:IV.15.3\endcsname{ega4.c04.subsec.15-3}
\expandafter\def\csname sga@target@subsection:IV.15.4\endcsname{ega4.c04.subsec.15-4}
\expandafter\def\csname sga@target@subsection:IV.15.5\endcsname{ega4.c04.subsec.15-5}
\expandafter\def\csname sga@target@subsection:IV.15.6\endcsname{ega4.c04.subsec.15-6}
\expandafter\def\csname sga@target@subsection:IV.15.7\endcsname{ega4.c04.subsec.15-7}
\expandafter\def\csname sga@target@subsection:IV.17.1\endcsname{ega4.c04.subsec.17-1}
\expandafter\def\csname sga@target@subsection:IV.17.10\endcsname{ega4.c04.subsec.17-10}
\expandafter\def\csname sga@target@subsection:IV.17.11\endcsname{ega4.c04.subsec.17-11}
\expandafter\def\csname sga@target@subsection:IV.17.12\endcsname{ega4.c04.subsec.17-12}
\expandafter\def\csname sga@target@subsection:IV.17.13\endcsname{ega4.c04.subsec.17-13}
\expandafter\def\csname sga@target@subsection:IV.17.14\endcsname{ega4.c04.subsec.17-14}
\expandafter\def\csname sga@target@subsection:IV.17.15\endcsname{ega4.c04.subsec.17-15}
\expandafter\def\csname sga@target@subsection:IV.17.16\endcsname{ega4.c04.subsec.17-16}
\expandafter\def\csname sga@target@subsection:IV.17.2\endcsname{ega4.c04.subsec.17-2}
\expandafter\def\csname sga@target@subsection:IV.17.3\endcsname{ega4.c04.subsec.17-3}
\expandafter\def\csname sga@target@subsection:IV.17.4\endcsname{ega4.c04.subsec.17-4}
\expandafter\def\csname sga@target@subsection:IV.17.5\endcsname{ega4.c04.subsec.17-5}
\expandafter\def\csname sga@target@subsection:IV.17.6\endcsname{ega4.c04.subsec.17-6}
\expandafter\def\csname sga@target@subsection:IV.17.7\endcsname{ega4.c04.subsec.17-7}
\expandafter\def\csname sga@target@subsection:IV.17.8\endcsname{ega4.c04.subsec.17-8}
\expandafter\def\csname sga@target@subsection:IV.17.9\endcsname{ega4.c04.subsec.17-9}
\expandafter\def\csname sga@target@subsection:IV.18.1\endcsname{ega4.c04.subsec.18-1}
\expandafter\def\csname sga@target@subsection:IV.18.10\endcsname{ega4.c04.subsec.18-10}
\expandafter\def\csname sga@target@subsection:IV.18.11\endcsname{ega4.c04.subsec.18-11}
\expandafter\def\csname sga@target@subsection:IV.18.12\endcsname{ega4.c04.subsec.18-12}
\expandafter\def\csname sga@target@subsection:IV.18.2\endcsname{ega4.c04.subsec.18-2}
\expandafter\def\csname sga@target@subsection:IV.18.3\endcsname{ega4.c04.subsec.18-3}
\expandafter\def\csname sga@target@subsection:IV.18.4\endcsname{ega4.c04.subsec.18-4}
\expandafter\def\csname sga@target@subsection:IV.18.5\endcsname{ega4.c04.subsec.18-5}
\expandafter\def\csname sga@target@subsection:IV.18.6\endcsname{ega4.c04.subsec.18-6}
\expandafter\def\csname sga@target@subsection:IV.18.7\endcsname{ega4.c04.subsec.18-7}
\expandafter\def\csname sga@target@subsection:IV.18.8\endcsname{ega4.c04.subsec.18-8}
\expandafter\def\csname sga@target@subsection:IV.18.9\endcsname{ega4.c04.subsec.18-9}
\expandafter\def\csname sga@target@subsection:IV.19.1\endcsname{ega4.c04.subsec.19-1}
\expandafter\def\csname sga@target@subsection:IV.19.2\endcsname{ega4.c04.subsec.19-2}
\expandafter\def\csname sga@target@subsection:IV.19.3\endcsname{ega4.c04.subsec.19-3}
\expandafter\def\csname sga@target@subsection:IV.19.4\endcsname{ega4.c04.subsec.19-4}
\expandafter\def\csname sga@target@subsection:IV.19.5\endcsname{ega4.c04.subsec.19-5}
\expandafter\def\csname sga@target@subsection:IV.19.6\endcsname{ega4.c04.subsec.19-6}
\expandafter\def\csname sga@target@subsection:IV.19.7\endcsname{ega4.c04.subsec.19-7}
\expandafter\def\csname sga@target@subsection:IV.19.8\endcsname{ega4.c04.subsec.19-8}
\expandafter\def\csname sga@target@subsection:IV.19.9\endcsname{ega4.c04.subsec.19-9}
\expandafter\def\csname sga@target@subsection:IV.2.1\endcsname{ega4.c04.subsec.2-1}
\expandafter\def\csname sga@target@subsection:IV.2.2\endcsname{ega4.c04.subsec.2-2}
\expandafter\def\csname sga@target@subsection:IV.2.3\endcsname{ega4.c04.subsec.2-3}
\expandafter\def\csname sga@target@subsection:IV.2.4\endcsname{ega4.c04.subsec.2-4}
\expandafter\def\csname sga@target@subsection:IV.2.5\endcsname{ega4.c04.subsec.2-5}
\expandafter\def\csname sga@target@subsection:IV.2.6\endcsname{ega4.c04.subsec.2-6}
\expandafter\def\csname sga@target@subsection:IV.2.7\endcsname{ega4.c04.subsec.2-7}
\expandafter\def\csname sga@target@subsection:IV.2.8\endcsname{ega4.c04.subsec.2-8}
\expandafter\def\csname sga@target@subsection:IV.21.1\endcsname{ega4.c04.subsec.21-1}
\expandafter\def\csname sga@target@subsection:IV.21.10\endcsname{ega4.c04.subsec.21-10}
\expandafter\def\csname sga@target@subsection:IV.21.11\endcsname{ega4.c04.subsec.21-11}
\expandafter\def\csname sga@target@subsection:IV.21.12\endcsname{ega4.c04.subsec.21-12}
\expandafter\def\csname sga@target@subsection:IV.21.13\endcsname{ega4.c04.subsec.21-13}
\expandafter\def\csname sga@target@subsection:IV.21.14\endcsname{ega4.c04.subsec.21-14}
\expandafter\def\csname sga@target@subsection:IV.21.15\endcsname{ega4.c04.subsec.21-15}
\expandafter\def\csname sga@target@subsection:IV.21.2\endcsname{ega4.c04.subsec.21-2}
\expandafter\def\csname sga@target@subsection:IV.21.3\endcsname{ega4.c04.subsec.21-3}
\expandafter\def\csname sga@target@subsection:IV.21.4\endcsname{ega4.c04.subsec.21-4}
\expandafter\def\csname sga@target@subsection:IV.21.5\endcsname{ega4.c04.subsec.21-5}
\expandafter\def\csname sga@target@subsection:IV.21.6\endcsname{ega4.c04.subsec.21-6}
\expandafter\def\csname sga@target@subsection:IV.21.7\endcsname{ega4.c04.subsec.21-7}
\expandafter\def\csname sga@target@subsection:IV.21.8\endcsname{ega4.c04.subsec.21-8}
\expandafter\def\csname sga@target@subsection:IV.21.9\endcsname{ega4.c04.subsec.21-9}
\expandafter\def\csname sga@target@subsection:IV.3.1\endcsname{ega4.c04.subsec.3-1}
\expandafter\def\csname sga@target@subsection:IV.3.2\endcsname{ega4.c04.subsec.3-2}
\expandafter\def\csname sga@target@subsection:IV.3.3\endcsname{ega4.c04.subsec.3-3}
\expandafter\def\csname sga@target@subsection:IV.3.4\endcsname{ega4.c04.subsec.3-4}
\expandafter\def\csname sga@target@subsection:IV.4.1\endcsname{ega4.c04.subsec.4-1}
\expandafter\def\csname sga@target@subsection:IV.4.2\endcsname{ega4.c04.subsec.4-2}
\expandafter\def\csname sga@target@subsection:IV.4.3\endcsname{ega4.c04.subsec.4-3}
\expandafter\def\csname sga@target@subsection:IV.4.4\endcsname{ega4.c04.subsec.4-4}
\expandafter\def\csname sga@target@subsection:IV.4.5\endcsname{ega4.c04.subsec.4-5}
\expandafter\def\csname sga@target@subsection:IV.4.6\endcsname{ega4.c04.subsec.4-6}
\expandafter\def\csname sga@target@subsection:IV.4.7\endcsname{ega4.c04.subsec.4-7}
\expandafter\def\csname sga@target@subsection:IV.4.8\endcsname{ega4.c04.subsec.4-8}
\expandafter\def\csname sga@target@subsection:IV.4.9\endcsname{ega4.c04.subsec.4-9}
\expandafter\def\csname sga@target@subsection:IV.5.1\endcsname{ega4.c04.subsec.5-1}
\expandafter\def\csname sga@target@subsection:IV.5.10\endcsname{ega4.c04.subsec.5-10}
\expandafter\def\csname sga@target@subsection:IV.5.11\endcsname{ega4.c04.subsec.5-11}
\expandafter\def\csname sga@target@subsection:IV.5.12\endcsname{ega4.c04.subsec.5-12}
\expandafter\def\csname sga@target@subsection:IV.5.13\endcsname{ega4.c04.subsec.5-13}
\expandafter\def\csname sga@target@subsection:IV.5.2\endcsname{ega4.c04.subsec.5-2}
\expandafter\def\csname sga@target@subsection:IV.5.3\endcsname{ega4.c04.subsec.5-3}
\expandafter\def\csname sga@target@subsection:IV.5.4\endcsname{ega4.c04.subsec.5-4}
\expandafter\def\csname sga@target@subsection:IV.5.5\endcsname{ega4.c04.subsec.5-5}
\expandafter\def\csname sga@target@subsection:IV.5.6\endcsname{ega4.c04.subsec.5-6}
\expandafter\def\csname sga@target@subsection:IV.5.7\endcsname{ega4.c04.subsec.5-7}
\expandafter\def\csname sga@target@subsection:IV.5.8\endcsname{ega4.c04.subsec.5-8}
\expandafter\def\csname sga@target@subsection:IV.5.9\endcsname{ega4.c04.subsec.5-9}
\expandafter\def\csname sga@target@subsection:IV.6.1\endcsname{ega4.c04.subsec.6-1}
\expandafter\def\csname sga@target@subsection:IV.6.10\endcsname{ega4.c04.subsec.6-10}
\expandafter\def\csname sga@target@subsection:IV.6.11\endcsname{ega4.c04.subsec.6-11}
\expandafter\def\csname sga@target@subsection:IV.6.12\endcsname{ega4.c04.subsec.6-12}
\expandafter\def\csname sga@target@subsection:IV.6.13\endcsname{ega4.c04.subsec.6-13}
\expandafter\def\csname sga@target@subsection:IV.6.14\endcsname{ega4.c04.subsec.6-14}
\expandafter\def\csname sga@target@subsection:IV.6.15\endcsname{ega4.c04.subsec.6-15}
\expandafter\def\csname sga@target@subsection:IV.6.2\endcsname{ega4.c04.subsec.6-2}
\expandafter\def\csname sga@target@subsection:IV.6.3\endcsname{ega4.c04.subsec.6-3}
\expandafter\def\csname sga@target@subsection:IV.6.4\endcsname{ega4.c04.subsec.6-4}
\expandafter\def\csname sga@target@subsection:IV.6.5\endcsname{ega4.c04.subsec.6-5}
\expandafter\def\csname sga@target@subsection:IV.6.6\endcsname{ega4.c04.subsec.6-6}
\expandafter\def\csname sga@target@subsection:IV.6.7\endcsname{ega4.c04.subsec.6-7}
\expandafter\def\csname sga@target@subsection:IV.6.8\endcsname{ega4.c04.subsec.6-8}
\expandafter\def\csname sga@target@subsection:IV.6.9\endcsname{ega4.c04.subsec.6-9}
\expandafter\def\csname sga@target@subsection:IV.7.1\endcsname{ega4.c04.subsec.7-1}
\expandafter\def\csname sga@target@subsection:IV.7.2\endcsname{ega4.c04.subsec.7-2}
\expandafter\def\csname sga@target@subsection:IV.7.3\endcsname{ega4.c04.subsec.7-3}
\expandafter\def\csname sga@target@subsection:IV.7.4\endcsname{ega4.c04.subsec.7-4}
\expandafter\def\csname sga@target@subsection:IV.7.5\endcsname{ega4.c04.subsec.7-5}
\expandafter\def\csname sga@target@subsection:IV.7.6\endcsname{ega4.c04.subsec.7-6}
\expandafter\def\csname sga@target@subsection:IV.7.7\endcsname{ega4.c04.subsec.7-7}
\expandafter\def\csname sga@target@subsection:IV.7.8\endcsname{ega4.c04.subsec.7-8}
\expandafter\def\csname sga@target@subsection:IV.7.9\endcsname{ega4.c04.subsec.7-9}
\expandafter\def\csname sga@target@subsection:IV.8.1\endcsname{ega4.c04.subsec.8-1}
\expandafter\def\csname sga@target@subsection:IV.8.10\endcsname{ega4.c04.subsec.8-10}
\expandafter\def\csname sga@target@subsection:IV.8.11\endcsname{ega4.c04.subsec.8-11}
\expandafter\def\csname sga@target@subsection:IV.8.12\endcsname{ega4.c04.subsec.8-12}
\expandafter\def\csname sga@target@subsection:IV.8.13\endcsname{ega4.c04.subsec.8-13}
\expandafter\def\csname sga@target@subsection:IV.8.14\endcsname{ega4.c04.subsec.8-14}
\expandafter\def\csname sga@target@subsection:IV.8.2\endcsname{ega4.c04.subsec.8-2}
\expandafter\def\csname sga@target@subsection:IV.8.3\endcsname{ega4.c04.subsec.8-3}
\expandafter\def\csname sga@target@subsection:IV.8.4\endcsname{ega4.c04.subsec.8-4}
\expandafter\def\csname sga@target@subsection:IV.8.5\endcsname{ega4.c04.subsec.8-5}
\expandafter\def\csname sga@target@subsection:IV.8.6\endcsname{ega4.c04.subsec.8-6}
\expandafter\def\csname sga@target@subsection:IV.8.7\endcsname{ega4.c04.subsec.8-7}
\expandafter\def\csname sga@target@subsection:IV.8.8\endcsname{ega4.c04.subsec.8-8}
\expandafter\def\csname sga@target@subsection:IV.8.9\endcsname{ega4.c04.subsec.8-9}
\expandafter\def\csname sga@target@subsection:IV.9.1\endcsname{ega4.c04.subsec.9-1}
\expandafter\def\csname sga@target@subsection:IV.9.2\endcsname{ega4.c04.subsec.9-2}
\expandafter\def\csname sga@target@subsection:IV.9.3\endcsname{ega4.c04.subsec.9-3}
\expandafter\def\csname sga@target@subsection:IV.9.4\endcsname{ega4.c04.subsec.9-4}
\expandafter\def\csname sga@target@subsection:IV.9.5\endcsname{ega4.c04.subsec.9-5}
\expandafter\def\csname sga@target@subsection:IV.9.6\endcsname{ega4.c04.subsec.9-6}
\expandafter\def\csname sga@target@subsection:IV.9.7\endcsname{ega4.c04.subsec.9-7}
\expandafter\def\csname sga@target@subsection:IV.9.8\endcsname{ega4.c04.subsec.9-8}
\expandafter\def\csname sga@target@subsection:IV.9.9\endcsname{ega4.c04.subsec.9-9}
\expandafter\def\csname sga@bib@I-4\endcsname{ega4.reader.bib.i-4}
\expandafter\def\csname sga@bib@IV-32\endcsname{ega4.reader.bib.iv-32}
\expandafter\def\csname sga@bib@IV-34\endcsname{ega4.reader.bib.iv-34}
\expandafter\def\csname sga@bib@IV-37\endcsname{ega4.reader.bib.iv-37}
\expandafter\def\csname sga@bib@IV-41\endcsname{ega4.reader.bib.iv-41}
\expandafter\def\csname sga@bib@IV-42\endcsname{ega4.reader.bib.iv-42}
\expandafter\def\csname sga@bib@IV-43\endcsname{ega4.reader.bib.iv-43}
\expandafter\def\csname sga@bib@IV-44\endcsname{ega4.reader.bib.iv-44}
\expandafter\def\csname sga@bib@SGA1\endcsname{ega4.reader.bib.sga1}
\let\sgaOriginalLabel\label
\renewcommand{\label}[1]{%
  \sgaOriginalLabel{#1}%
  \ifcsname sga@target@#1\endcsname
    \expandafter\sgaStableTarget\expandafter{\csname sga@target@#1\endcsname}%
  \fi
}
\let\sgaOriginalBibitem\bibitem
\newcommand{\sgaBibTarget}[1]{%
  \ifcsname sga@bib@#1\endcsname
    \expandafter\sgaStableTarget\expandafter{\csname sga@bib@#1\endcsname}%
  \fi
}
\newcommand{\sgaBibitemPlain}[1]{%
  \sgaOriginalBibitem{#1}\sgaBibTarget{#1}}
\newcommand{\sgaBibitemOptional}[2][]{%
  \sgaOriginalBibitem[#1]{#2}\sgaBibTarget{#2}}
\renewcommand{\bibitem}{%
  \@ifnextchar[\sgaBibitemOptional\sgaBibitemPlain}
\makeatother

% Generated stable-link dispatch for the EGA IV reader.
% Stable destinations are emitted by EGA4_reference_v2_registry.tex.
\makeatletter
\expandafter\def\csname sga@target@IV.1.1.1\endcsname{ega4.c04.def.1-1-1}
\expandafter\def\csname sga@target@IV.1.1.2\endcsname{ega4.c04.prop.1-1-2}
\expandafter\def\csname sga@target@IV.1.1.3\endcsname{ega4.c04.prop.1-1-3}
\expandafter\def\csname sga@target@IV.1.1.4\endcsname{ega4.c04.cor.1-1-4}
\expandafter\def\csname sga@target@IV.1.1.5\endcsname{ega4.c04.prop.1-1-5}
\expandafter\def\csname sga@target@IV.1.1.6\endcsname{ega4.c04.prop.1-1-6}
\expandafter\def\csname sga@target@IV.1.10.1\endcsname{ega4.c04.thm.1-10-1}
\expandafter\def\csname sga@target@IV.1.10.2\endcsname{ega4.c04.cor.1-10-2}
\expandafter\def\csname sga@target@IV.1.10.3\endcsname{ega4.c04.prop.1-10-3}
\expandafter\def\csname sga@target@IV.1.10.4\endcsname{ega4.c04.cor.1-10-4}
\expandafter\def\csname sga@target@IV.1.2.1\endcsname{ega4.c04.def.1-2-1}
\expandafter\def\csname sga@target@IV.1.2.2\endcsname{ega4.c04.prop.1-2-2}
\expandafter\def\csname sga@target@IV.1.2.3\endcsname{ega4.c04.cor.1-2-3}
\expandafter\def\csname sga@target@IV.1.2.4\endcsname{ega4.c04.prop.1-2-4}
\expandafter\def\csname sga@target@IV.1.2.4.1\endcsname{ega4.c04.eq.1-2-4-1}
\expandafter\def\csname sga@target@IV.1.2.5\endcsname{ega4.c04.prop.1-2-5}
\expandafter\def\csname sga@target@IV.1.2.6\endcsname{ega4.c04.prop.1-2-6}
\expandafter\def\csname sga@target@IV.1.2.7\endcsname{ega4.c04.prop.1-2-7}
\expandafter\def\csname sga@target@IV.1.2.8\endcsname{ega4.c04.cor.1-2-8}
\expandafter\def\csname sga@target@IV.1.2.9\endcsname{ega4.c04.prop.1-2-9}
\expandafter\def\csname sga@target@IV.1.3.1\endcsname{ega4.c04.statement.1-3-1}
\expandafter\def\csname sga@target@IV.1.3.10\endcsname{ega4.c04.statement.1-3-10}
\expandafter\def\csname sga@target@IV.1.3.11\endcsname{ega4.c04.prop.1-3-11}
\expandafter\def\csname sga@target@IV.1.3.2\endcsname{ega4.c04.def.1-3-2}
\expandafter\def\csname sga@target@IV.1.3.3\endcsname{ega4.c04.prop.1-3-3}
\expandafter\def\csname sga@target@IV.1.3.4\endcsname{ega4.c04.prop.1-3-4}
\expandafter\def\csname sga@target@IV.1.3.5\endcsname{ega4.c04.cor.1-3-5}
\expandafter\def\csname sga@target@IV.1.3.6\endcsname{ega4.c04.prop.1-3-6}
\expandafter\def\csname sga@target@IV.1.3.7\endcsname{ega4.c04.prop.1-3-7}
\expandafter\def\csname sga@target@IV.1.3.8\endcsname{ega4.c04.statement.1-3-8}
\expandafter\def\csname sga@target@IV.1.3.9\endcsname{ega4.c04.prop.1-3-9}
\expandafter\def\csname sga@target@IV.1.4.1\endcsname{ega4.c04.statement.1-4-1}
\expandafter\def\csname sga@target@IV.1.4.2\endcsname{ega4.c04.def.1-4-2}
\expandafter\def\csname sga@target@IV.1.4.3\endcsname{ega4.c04.prop.1-4-3}
\expandafter\def\csname sga@target@IV.1.4.3.1\endcsname{ega4.c04.cor.1-4-3-1}
\expandafter\def\csname sga@target@IV.1.4.4\endcsname{ega4.c04.prop.1-4-4}
\expandafter\def\csname sga@target@IV.1.4.4.1\endcsname{ega4.c04.lem.1-4-4-1}
\expandafter\def\csname sga@target@IV.1.4.5\endcsname{ega4.c04.cor.1-4-5}
\expandafter\def\csname sga@target@IV.1.4.6\endcsname{ega4.c04.prop.1-4-6}
\expandafter\def\csname sga@target@IV.1.4.7\endcsname{ega4.c04.prop.1-4-7}
\expandafter\def\csname sga@target@IV.1.4.7.1\endcsname{ega4.c04.lem.1-4-7-1}
\expandafter\def\csname sga@target@IV.1.5.1\endcsname{ega4.c04.prop.1-5-1}
\expandafter\def\csname sga@target@IV.1.5.2\endcsname{ega4.c04.def.1-5-2}
\expandafter\def\csname sga@target@IV.1.5.3\endcsname{ega4.c04.prop.1-5-3}
\expandafter\def\csname sga@target@IV.1.5.4\endcsname{ega4.c04.prop.1-5-4}
\expandafter\def\csname sga@target@IV.1.5.5\endcsname{ega4.c04.cor.1-5-5}
\expandafter\def\csname sga@target@IV.1.5.6\endcsname{ega4.c04.cor.1-5-6}
\expandafter\def\csname sga@target@IV.1.5.7\endcsname{ega4.c04.prop.1-5-7}
\expandafter\def\csname sga@target@IV.1.6.1\endcsname{ega4.c04.def.1-6-1}
\expandafter\def\csname sga@target@IV.1.6.2\endcsname{ega4.c04.prop.1-6-2}
\expandafter\def\csname sga@target@IV.1.6.3\endcsname{ega4.c04.cor.1-6-3}
\expandafter\def\csname sga@target@IV.1.6.4\endcsname{ega4.c04.rem.1-6-4}
\expandafter\def\csname sga@target@IV.1.6.5\endcsname{ega4.c04.prop.1-6-5}
\expandafter\def\csname sga@target@IV.1.7.1\endcsname{ega4.c04.statement.1-7-1}
\expandafter\def\csname sga@target@IV.1.7.10\endcsname{ega4.c04.statement.1-7-10}
\expandafter\def\csname sga@target@IV.1.7.11\endcsname{ega4.c04.statement.1-7-11}
\expandafter\def\csname sga@target@IV.1.7.12\endcsname{ega4.c04.statement.1-7-12}
\expandafter\def\csname sga@target@IV.1.7.13\endcsname{ega4.c04.statement.1-7-13}
\expandafter\def\csname sga@target@IV.1.7.14\endcsname{ega4.c04.statement.1-7-14}
\expandafter\def\csname sga@target@IV.1.7.15\endcsname{ega4.c04.statement.1-7-15}
\expandafter\def\csname sga@target@IV.1.7.16\endcsname{ega4.c04.statement.1-7-16}
\expandafter\def\csname sga@target@IV.1.7.17\endcsname{ega4.c04.statement.1-7-17}
\expandafter\def\csname sga@target@IV.1.7.18\endcsname{ega4.c04.statement.1-7-18}
\expandafter\def\csname sga@target@IV.1.7.19\endcsname{ega4.c04.statement.1-7-19}
\expandafter\def\csname sga@target@IV.1.7.2\endcsname{ega4.c04.statement.1-7-2}
\expandafter\def\csname sga@target@IV.1.7.20\endcsname{ega4.c04.statement.1-7-20}
\expandafter\def\csname sga@target@IV.1.7.21\endcsname{ega4.c04.statement.1-7-21}
\expandafter\def\csname sga@target@IV.1.7.22\endcsname{ega4.c04.statement.1-7-22}
\expandafter\def\csname sga@target@IV.1.7.3\endcsname{ega4.c04.statement.1-7-3}
\expandafter\def\csname sga@target@IV.1.7.4\endcsname{ega4.c04.statement.1-7-4}
\expandafter\def\csname sga@target@IV.1.7.5\endcsname{ega4.c04.statement.1-7-5}
\expandafter\def\csname sga@target@IV.1.7.6\endcsname{ega4.c04.statement.1-7-6}
\expandafter\def\csname sga@target@IV.1.7.7\endcsname{ega4.c04.statement.1-7-7}
\expandafter\def\csname sga@target@IV.1.7.8\endcsname{ega4.c04.statement.1-7-8}
\expandafter\def\csname sga@target@IV.1.7.9\endcsname{ega4.c04.statement.1-7-9}
\expandafter\def\csname sga@target@IV.1.8.1\endcsname{ega4.c04.statement.1-8-1}
\expandafter\def\csname sga@target@IV.1.8.2\endcsname{ega4.c04.prop.1-8-2}
\expandafter\def\csname sga@target@IV.1.8.3\endcsname{ega4.c04.lem.1-8-3}
\expandafter\def\csname sga@target@IV.1.8.4\endcsname{ega4.c04.thm.1-8-4}
\expandafter\def\csname sga@target@IV.1.8.4.1\endcsname{ega4.c04.lem.1-8-4-1}
\expandafter\def\csname sga@target@IV.1.8.4.2\endcsname{ega4.c04.lem.1-8-4-2}
\expandafter\def\csname sga@target@IV.1.8.5\endcsname{ega4.c04.cor.1-8-5}
\expandafter\def\csname sga@target@IV.1.8.5.1\endcsname{ega4.c04.lem.1-8-5-1}
\expandafter\def\csname sga@target@IV.1.8.6\endcsname{ega4.c04.cor.1-8-6}
\expandafter\def\csname sga@target@IV.1.8.7\endcsname{ega4.c04.prop.1-8-7}
\expandafter\def\csname sga@target@IV.1.8.7.1\endcsname{ega4.c04.lem.1-8-7-1}
\expandafter\def\csname sga@target@IV.1.9.1\endcsname{ega4.c04.lem.1-9-1}
\expandafter\def\csname sga@target@IV.1.9.10\endcsname{ega4.c04.prop.1-9-10}
\expandafter\def\csname sga@target@IV.1.9.11\endcsname{ega4.c04.prop.1-9-11}
\expandafter\def\csname sga@target@IV.1.9.12\endcsname{ega4.c04.cor.1-9-12}
\expandafter\def\csname sga@target@IV.1.9.13\endcsname{ega4.c04.statement.1-9-13}
\expandafter\def\csname sga@target@IV.1.9.14\endcsname{ega4.c04.prop.1-9-14}
\expandafter\def\csname sga@target@IV.1.9.15\endcsname{ega4.c04.prop.1-9-15}
\expandafter\def\csname sga@target@IV.1.9.16\endcsname{ega4.c04.statement.1-9-16}
\expandafter\def\csname sga@target@IV.1.9.2\endcsname{ega4.c04.prop.1-9-2}
\expandafter\def\csname sga@target@IV.1.9.2.1\endcsname{ega4.c04.statement.1-9-2-1}
\expandafter\def\csname sga@target@IV.1.9.2.i\endcsname{ega4.c04.statement.1-9-2-i}
\expandafter\def\csname sga@target@IV.1.9.2.ii\endcsname{ega4.c04.statement.1-9-2-ii}
\expandafter\def\csname sga@target@IV.1.9.3\endcsname{ega4.c04.prop.1-9-3}
\expandafter\def\csname sga@target@IV.1.9.3.1\endcsname{ega4.c04.lem.1-9-3-1}
\expandafter\def\csname sga@target@IV.1.9.3.2\endcsname{ega4.c04.lem.1-9-3-2}
\expandafter\def\csname sga@target@IV.1.9.4\endcsname{ega4.c04.def.1-9-4}
\expandafter\def\csname sga@target@IV.1.9.5\endcsname{ega4.c04.prop.1-9-5}
\expandafter\def\csname sga@target@IV.1.9.6\endcsname{ega4.c04.statement.1-9-6}
\expandafter\def\csname sga@target@IV.1.9.7\endcsname{ega4.c04.thm.1-9-7}
\expandafter\def\csname sga@target@IV.1.9.8\endcsname{ega4.c04.cor.1-9-8}
\expandafter\def\csname sga@target@IV.1.9.9\endcsname{ega4.c04.cor.1-9-9}
\expandafter\def\csname sga@target@IV.10.1.1\endcsname{ega4.c04.statement.10-1-1}
\expandafter\def\csname sga@target@IV.10.1.2\endcsname{ega4.c04.prop.10-1-2}
\expandafter\def\csname sga@target@IV.10.1.3\endcsname{ega4.c04.def.10-1-3}
\expandafter\def\csname sga@target@IV.10.1.4\endcsname{ega4.c04.cor.10-1-4}
\expandafter\def\csname sga@target@IV.10.10.1\endcsname{ega4.c04.statement.10-10-1}
\expandafter\def\csname sga@target@IV.10.10.2\endcsname{ega4.c04.statement.10-10-2}
\expandafter\def\csname sga@target@IV.10.10.3\endcsname{ega4.c04.statement.10-10-3}
\expandafter\def\csname sga@target@IV.10.10.4\endcsname{ega4.c04.statement.10-10-4}
\expandafter\def\csname sga@target@IV.10.10.5\endcsname{ega4.c04.statement.10-10-5}
\expandafter\def\csname sga@target@IV.10.2.1\endcsname{ega4.c04.prop.10-2-1}
\expandafter\def\csname sga@target@IV.10.2.2\endcsname{ega4.c04.def.10-2-2}
\expandafter\def\csname sga@target@IV.10.2.3\endcsname{ega4.c04.cor.10-2-3}
\expandafter\def\csname sga@target@IV.10.2.4\endcsname{ega4.c04.cor.10-2-4}
\expandafter\def\csname sga@target@IV.10.2.5\endcsname{ega4.c04.cor.10-2-5}
\expandafter\def\csname sga@target@IV.10.2.6\endcsname{ega4.c04.cor.10-2-6}
\expandafter\def\csname sga@target@IV.10.2.7\endcsname{ega4.c04.prop.10-2-7}
\expandafter\def\csname sga@target@IV.10.2.8\endcsname{ega4.c04.item.10-2-8}
\expandafter\def\csname sga@target@IV.10.2.8.1\endcsname{ega4.c04.eq.10-2-8-1}
\expandafter\def\csname sga@target@IV.10.3.1\endcsname{ega4.c04.def.10-3-1}
\expandafter\def\csname sga@target@IV.10.3.2\endcsname{ega4.c04.prop.10-3-2}
\expandafter\def\csname sga@target@IV.10.3.3\endcsname{ega4.c04.prop.10-3-3}
\expandafter\def\csname sga@target@IV.10.4.1\endcsname{ega4.c04.def.10-4-1}
\expandafter\def\csname sga@target@IV.10.4.10\endcsname{ega4.c04.statement.10-4-10}
\expandafter\def\csname sga@target@IV.10.4.11\endcsname{ega4.c04.statement.10-4-11}
\expandafter\def\csname sga@target@IV.10.4.11.1\endcsname{ega4.c04.lem.10-4-11-1}
\expandafter\def\csname sga@target@IV.10.4.12\endcsname{ega4.c04.cor.10-4-12}
\expandafter\def\csname sga@target@IV.10.4.12.1\endcsname{ega4.c04.sec.10-4-12-1}
\expandafter\def\csname sga@target@IV.10.4.12.1.errata\endcsname{ega4.c04.eq.10-4-12-1-errata}
\expandafter\def\csname sga@target@IV.10.4.13\endcsname{ega4.c04.cor.10-4-13}
\expandafter\def\csname sga@target@IV.10.4.14\endcsname{ega4.c04.rem.10-4-14}
\expandafter\def\csname sga@target@IV.10.4.2\endcsname{ega4.c04.prop.10-4-2}
\expandafter\def\csname sga@target@IV.10.4.3\endcsname{ega4.c04.statement.10-4-3}
\expandafter\def\csname sga@target@IV.10.4.4\endcsname{ega4.c04.prop.10-4-4}
\expandafter\def\csname sga@target@IV.10.4.5\endcsname{ega4.c04.prop.10-4-5}
\expandafter\def\csname sga@target@IV.10.4.5.1\endcsname{ega4.c04.lem.10-4-5-1}
\expandafter\def\csname sga@target@IV.10.4.6\endcsname{ega4.c04.cor.10-4-6}
\expandafter\def\csname sga@target@IV.10.4.7\endcsname{ega4.c04.cor.10-4-7}
\expandafter\def\csname sga@target@IV.10.4.8\endcsname{ega4.c04.cor.10-4-8}
\expandafter\def\csname sga@target@IV.10.4.9\endcsname{ega4.c04.statement.10-4-9}
\expandafter\def\csname sga@target@IV.10.5.1\endcsname{ega4.c04.prop.10-5-1}
\expandafter\def\csname sga@target@IV.10.5.10\endcsname{ega4.c04.prop.10-5-10}
\expandafter\def\csname sga@target@IV.10.5.2\endcsname{ega4.c04.cor.10-5-2}
\expandafter\def\csname sga@target@IV.10.5.3\endcsname{ega4.c04.cor.10-5-3}
\expandafter\def\csname sga@target@IV.10.5.4\endcsname{ega4.c04.rem.10-5-4}
\expandafter\def\csname sga@target@IV.10.5.5\endcsname{ega4.c04.prop.10-5-5}
\expandafter\def\csname sga@target@IV.10.5.6\endcsname{ega4.c04.rem.10-5-6}
\expandafter\def\csname sga@target@IV.10.5.7\endcsname{ega4.c04.prop.10-5-7}
\expandafter\def\csname sga@target@IV.10.5.8\endcsname{ega4.c04.cor.10-5-8}
\expandafter\def\csname sga@target@IV.10.5.9\endcsname{ega4.c04.cor.10-5-9}
\expandafter\def\csname sga@target@IV.10.6.1\endcsname{ega4.c04.prop.10-6-1}
\expandafter\def\csname sga@target@IV.10.6.1.1\endcsname{ega4.c04.eq.10-6-1-1}
\expandafter\def\csname sga@target@IV.10.6.1.2\endcsname{ega4.c04.eq.10-6-1-2}
\expandafter\def\csname sga@target@IV.10.6.1.3\endcsname{ega4.c04.eq.10-6-1-3}
\expandafter\def\csname sga@target@IV.10.6.1.4\endcsname{ega4.c04.eq.10-6-1-4}
\expandafter\def\csname sga@target@IV.10.6.2\endcsname{ega4.c04.cor.10-6-2}
\expandafter\def\csname sga@target@IV.10.6.3\endcsname{ega4.c04.prop.10-6-3}
\expandafter\def\csname sga@target@IV.10.6.3.1\endcsname{ega4.c04.eq.10-6-3-1}
\expandafter\def\csname sga@target@IV.10.6.3.2\endcsname{ega4.c04.eq.10-6-3-2}
\expandafter\def\csname sga@target@IV.10.6.3.3\endcsname{ega4.c04.eq.10-6-3-3}
\expandafter\def\csname sga@target@IV.10.6.4\endcsname{ega4.c04.cor.10-6-4}
\expandafter\def\csname sga@target@IV.10.6.4.1\endcsname{ega4.c04.eq.10-6-4-1}
\expandafter\def\csname sga@target@IV.10.7.1\endcsname{ega4.c04.statement.10-7-1}
\expandafter\def\csname sga@target@IV.10.7.2\endcsname{ega4.c04.statement.10-7-2}
\expandafter\def\csname sga@target@IV.10.7.3\endcsname{ega4.c04.statement.10-7-3}
\expandafter\def\csname sga@target@IV.10.8.1\endcsname{ega4.c04.def.10-8-1}
\expandafter\def\csname sga@target@IV.10.8.1.1\endcsname{ega4.c04.eq.10-8-1-1}
\expandafter\def\csname sga@target@IV.10.8.1.2\endcsname{ega4.c04.eq.10-8-1-2}
\expandafter\def\csname sga@target@IV.10.8.2\endcsname{ega4.c04.statement.10-8-2}
\expandafter\def\csname sga@target@IV.10.8.3\endcsname{ega4.c04.lem.10-8-3}
\expandafter\def\csname sga@target@IV.10.8.3.1\endcsname{ega4.c04.eq.10-8-3-1}
\expandafter\def\csname sga@target@IV.10.8.4\endcsname{ega4.c04.cor.10-8-4}
\expandafter\def\csname sga@target@IV.10.8.5\endcsname{ega4.c04.prop.10-8-5}
\expandafter\def\csname sga@target@IV.10.8.6\endcsname{ega4.c04.cor.10-8-6}
\expandafter\def\csname sga@target@IV.10.8.7\endcsname{ega4.c04.prop.10-8-7}
\expandafter\def\csname sga@target@IV.10.8.7.1\endcsname{ega4.c04.eq.10-8-7-1}
\expandafter\def\csname sga@target@IV.10.8.8\endcsname{ega4.c04.cor.10-8-8}
\expandafter\def\csname sga@target@IV.10.9.1\endcsname{ega4.c04.statement.10-9-1}
\expandafter\def\csname sga@target@IV.10.9.2\endcsname{ega4.c04.statement.10-9-2}
\expandafter\def\csname sga@target@IV.10.9.3\endcsname{ega4.c04.statement.10-9-3}
\expandafter\def\csname sga@target@IV.10.9.4\endcsname{ega4.c04.prop.10-9-4}
\expandafter\def\csname sga@target@IV.10.9.5\endcsname{ega4.c04.statement.10-9-5}
\expandafter\def\csname sga@target@IV.10.9.6\endcsname{ega4.c04.prop.10-9-6}
\expandafter\def\csname sga@target@IV.11.1.1\endcsname{ega4.c04.thm.11-1-1}
\expandafter\def\csname sga@target@IV.11.1.1.1\endcsname{ega4.c04.sec.11-1-1-1}
\expandafter\def\csname sga@target@IV.11.1.1.2\endcsname{ega4.c04.lem.11-1-1-2}
\expandafter\def\csname sga@target@IV.11.1.2\endcsname{ega4.c04.cor.11-1-2}
\expandafter\def\csname sga@target@IV.11.1.3\endcsname{ega4.c04.cor.11-1-3}
\expandafter\def\csname sga@target@IV.11.1.4\endcsname{ega4.c04.cor.11-1-4}
\expandafter\def\csname sga@target@IV.11.1.5\endcsname{ega4.c04.cor.11-1-5}
\expandafter\def\csname sga@target@IV.11.1.6\endcsname{ega4.c04.cor.11-1-6}
\expandafter\def\csname sga@target@IV.11.10.1\endcsname{ega4.c04.prop.11-10-1}
\expandafter\def\csname sga@target@IV.11.10.1.1\endcsname{ega4.c04.eq.11-10-1-1}
\expandafter\def\csname sga@target@IV.11.10.1.2\endcsname{ega4.c04.eq.11-10-1-2}
\expandafter\def\csname sga@target@IV.11.10.10\endcsname{ega4.c04.prop.11-10-10}
\expandafter\def\csname sga@target@IV.11.10.2\endcsname{ega4.c04.def.11-10-2}
\expandafter\def\csname sga@target@IV.11.10.3\endcsname{ega4.c04.item.11-10-3}
\expandafter\def\csname sga@target@IV.11.10.4\endcsname{ega4.c04.prop.11-10-4}
\expandafter\def\csname sga@target@IV.11.10.5\endcsname{ega4.c04.thm.11-10-5}
\expandafter\def\csname sga@target@IV.11.10.6\endcsname{ega4.c04.prop.11-10-6}
\expandafter\def\csname sga@target@IV.11.10.7\endcsname{ega4.c04.cor.11-10-7}
\expandafter\def\csname sga@target@IV.11.10.8\endcsname{ega4.c04.def.11-10-8}
\expandafter\def\csname sga@target@IV.11.10.9\endcsname{ega4.c04.thm.11-10-9}
\expandafter\def\csname sga@target@IV.11.2.1\endcsname{ega4.c04.statement.11-2-1}
\expandafter\def\csname sga@target@IV.11.2.1.1\endcsname{ega4.c04.eq.11-2-1-1}
\expandafter\def\csname sga@target@IV.11.2.1.2\endcsname{ega4.c04.eq.11-2-1-2}
\expandafter\def\csname sga@target@IV.11.2.1.3\endcsname{ega4.c04.eq.11-2-1-3}
\expandafter\def\csname sga@target@IV.11.2.1.4\endcsname{ega4.c04.eq.11-2-1-4}
\expandafter\def\csname sga@target@IV.11.2.1.5\endcsname{ega4.c04.eq.11-2-1-5}
\expandafter\def\csname sga@target@IV.11.2.10\endcsname{ega4.c04.rem.11-2-10}
\expandafter\def\csname sga@target@IV.11.2.2\endcsname{ega4.c04.statement.11-2-2}
\expandafter\def\csname sga@target@IV.11.2.2.1\endcsname{ega4.c04.eq.11-2-2-1}
\expandafter\def\csname sga@target@IV.11.2.2.2\endcsname{ega4.c04.lem.11-2-2-2}
\expandafter\def\csname sga@target@IV.11.2.3\endcsname{ega4.c04.statement.11-2-3}
\expandafter\def\csname sga@target@IV.11.2.3.1\endcsname{ega4.c04.eq.11-2-3-1}
\expandafter\def\csname sga@target@IV.11.2.3.2\endcsname{ega4.c04.eq.11-2-3-2}
\expandafter\def\csname sga@target@IV.11.2.3.3\endcsname{ega4.c04.eq.11-2-3-3}
\expandafter\def\csname sga@target@IV.11.2.4\endcsname{ega4.c04.lem.11-2-4}
\expandafter\def\csname sga@target@IV.11.2.5\endcsname{ega4.c04.lem.11-2-5}
\expandafter\def\csname sga@target@IV.11.2.6\endcsname{ega4.c04.thm.11-2-6}
\expandafter\def\csname sga@target@IV.11.2.6.1\endcsname{ega4.c04.cor.11-2-6-1}
\expandafter\def\csname sga@target@IV.11.2.7\endcsname{ega4.c04.cor.11-2-7}
\expandafter\def\csname sga@target@IV.11.2.8\endcsname{ega4.c04.prop.11-2-8}
\expandafter\def\csname sga@target@IV.11.2.9\endcsname{ega4.c04.thm.11-2-9}
\expandafter\def\csname sga@target@IV.11.2.9.1\endcsname{ega4.c04.statement.11-2-9-1}
\expandafter\def\csname sga@target@IV.11.2.9.2\endcsname{ega4.c04.lem.11-2-9-2}
\expandafter\def\csname sga@target@IV.11.2.9.3\endcsname{ega4.c04.lem.11-2-9-3}
\expandafter\def\csname sga@target@IV.11.2.9.4\endcsname{ega4.c04.cor.11-2-9-4}
\expandafter\def\csname sga@target@IV.11.2.9.5\endcsname{ega4.c04.lem.11-2-9-5}
\expandafter\def\csname sga@target@IV.11.2.9.6\endcsname{ega4.c04.lem.11-2-9-6}
\expandafter\def\csname sga@target@IV.11.2.9.7\endcsname{ega4.c04.eq.11-2-9-7}
\expandafter\def\csname sga@target@IV.11.2.9.8\endcsname{ega4.c04.eq.11-2-9-8}
\expandafter\def\csname sga@target@IV.11.2.9.9\endcsname{ega4.c04.eq.11-2-9-9}
\expandafter\def\csname sga@target@IV.11.3\endcsname{ega4.c04.subsec.11-3.key-iv-11-3}
\expandafter\def\csname sga@target@IV.11.3.1\endcsname{ega4.c04.thm.11-3-1}
\expandafter\def\csname sga@target@IV.11.3.10\endcsname{ega4.c04.thm.11-3-10}
\expandafter\def\csname sga@target@IV.11.3.10.1\endcsname{ega4.c04.lem.11-3-10-1}
\expandafter\def\csname sga@target@IV.11.3.10.2\endcsname{ega4.c04.lem.11-3-10-2}
\expandafter\def\csname sga@target@IV.11.3.10.3\endcsname{ega4.c04.eq.11-3-10-3}
\expandafter\def\csname sga@target@IV.11.3.10.4\endcsname{ega4.c04.sec.11-3-10-4}
\expandafter\def\csname sga@target@IV.11.3.11\endcsname{ega4.c04.cor.11-3-11}
\expandafter\def\csname sga@target@IV.11.3.12\endcsname{ega4.c04.rem.11-3-12}
\expandafter\def\csname sga@target@IV.11.3.13\endcsname{ega4.c04.prop.11-3-13}
\expandafter\def\csname sga@target@IV.11.3.14\endcsname{ega4.c04.cor.11-3-14}
\expandafter\def\csname sga@target@IV.11.3.15\endcsname{ega4.c04.prop.11-3-15}
\expandafter\def\csname sga@target@IV.11.3.16\endcsname{ega4.c04.prop.11-3-16}
\expandafter\def\csname sga@target@IV.11.3.17\endcsname{ega4.c04.cor.11-3-17}
\expandafter\def\csname sga@target@IV.11.3.2\endcsname{ega4.c04.cor.11-3-2}
\expandafter\def\csname sga@target@IV.11.3.3\endcsname{ega4.c04.cor.11-3-3}
\expandafter\def\csname sga@target@IV.11.3.4\endcsname{ega4.c04.prop.11-3-4}
\expandafter\def\csname sga@target@IV.11.3.5\endcsname{ega4.c04.cor.11-3-5}
\expandafter\def\csname sga@target@IV.11.3.6\endcsname{ega4.c04.prop.11-3-6}
\expandafter\def\csname sga@target@IV.11.3.7\endcsname{ega4.c04.prop.11-3-7}
\expandafter\def\csname sga@target@IV.11.3.8\endcsname{ega4.c04.thm.11-3-8}
\expandafter\def\csname sga@target@IV.11.3.9\endcsname{ega4.c04.prop.11-3-9}
\expandafter\def\csname sga@target@IV.11.3.9.1\endcsname{ega4.c04.lem.11-3-9-1}
\expandafter\def\csname sga@target@IV.11.4.1\endcsname{ega4.c04.thm.11-4-1}
\expandafter\def\csname sga@target@IV.11.4.1.1\endcsname{ega4.c04.lem.11-4-1-1}
\expandafter\def\csname sga@target@IV.11.4.1.2\endcsname{ega4.c04.lem.11-4-1-2}
\expandafter\def\csname sga@target@IV.11.4.10\endcsname{ega4.c04.cor.11-4-10}
\expandafter\def\csname sga@target@IV.11.4.11\endcsname{ega4.c04.prop.11-4-11}
\expandafter\def\csname sga@target@IV.11.4.12\endcsname{ega4.c04.prop.11-4-12}
\expandafter\def\csname sga@target@IV.11.4.13\endcsname{ega4.c04.cor.11-4-13}
\expandafter\def\csname sga@target@IV.11.4.2\endcsname{ega4.c04.rem.11-4-2}
\expandafter\def\csname sga@target@IV.11.4.3\endcsname{ega4.c04.cor.11-4-3}
\expandafter\def\csname sga@target@IV.11.4.4\endcsname{ega4.c04.cor.11-4-4}
\expandafter\def\csname sga@target@IV.11.4.5\endcsname{ega4.c04.prop.11-4-5}
\expandafter\def\csname sga@target@IV.11.4.6\endcsname{ega4.c04.cor.11-4-6}
\expandafter\def\csname sga@target@IV.11.4.7\endcsname{ega4.c04.prop.11-4-7}
\expandafter\def\csname sga@target@IV.11.4.8\endcsname{ega4.c04.rem.11-4-8}
\expandafter\def\csname sga@target@IV.11.4.9\endcsname{ega4.c04.cor.11-4-9}
\expandafter\def\csname sga@target@IV.11.5.1\endcsname{ega4.c04.thm.11-5-1}
\expandafter\def\csname sga@target@IV.11.5.2\endcsname{ega4.c04.cor.11-5-2}
\expandafter\def\csname sga@target@IV.11.5.3\endcsname{ega4.c04.prop.11-5-3}
\expandafter\def\csname sga@target@IV.11.5.4\endcsname{ega4.c04.cor.11-5-4}
\expandafter\def\csname sga@target@IV.11.5.5\endcsname{ega4.c04.prop.11-5-5}
\expandafter\def\csname sga@target@IV.11.5.5.1\endcsname{ega4.c04.lem.11-5-5-1}
\expandafter\def\csname sga@target@IV.11.5.5.2\endcsname{ega4.c04.lem.11-5-5-2}
\expandafter\def\csname sga@target@IV.11.6.1\endcsname{ega4.c04.thm.11-6-1}
\expandafter\def\csname sga@target@IV.11.6.2\endcsname{ega4.c04.cor.11-6-2}
\expandafter\def\csname sga@target@IV.11.7.1\endcsname{ega4.c04.statement.11-7-1}
\expandafter\def\csname sga@target@IV.11.7.2\endcsname{ega4.c04.statement.11-7-2}
\expandafter\def\csname sga@target@IV.11.7.3\endcsname{ega4.c04.statement.11-7-3}
\expandafter\def\csname sga@target@IV.11.7.4\endcsname{ega4.c04.statement.11-7-4}
\expandafter\def\csname sga@target@IV.11.7.5\endcsname{ega4.c04.statement.11-7-5}
\expandafter\def\csname sga@target@IV.11.8.1\endcsname{ega4.c04.thm.11-8-1}
\expandafter\def\csname sga@target@IV.11.8.2\endcsname{ega4.c04.subsec.11-8-2}
\expandafter\def\csname sga@target@IV.11.9.1\endcsname{ega4.c04.statement.11-9-1}
\expandafter\def\csname sga@target@IV.11.9.1.1\endcsname{ega4.c04.statement.11-9-1-1}
\expandafter\def\csname sga@target@IV.11.9.10\endcsname{ega4.c04.thm.11-9-10}
\expandafter\def\csname sga@target@IV.11.9.11\endcsname{ega4.c04.rem.11-9-11}
\expandafter\def\csname sga@target@IV.11.9.12\endcsname{ega4.c04.lem.11-9-12}
\expandafter\def\csname sga@target@IV.11.9.13\endcsname{ega4.c04.prop.11-9-13}
\expandafter\def\csname sga@target@IV.11.9.14\endcsname{ega4.c04.statement.11-9-14}
\expandafter\def\csname sga@target@IV.11.9.15\endcsname{ega4.c04.prop.11-9-15}
\expandafter\def\csname sga@target@IV.11.9.16\endcsname{ega4.c04.thm.11-9-16}
\expandafter\def\csname sga@target@IV.11.9.17\endcsname{ega4.c04.prop.11-9-17}
\expandafter\def\csname sga@target@IV.11.9.17.1\endcsname{ega4.c04.lem.11-9-17-1}
\expandafter\def\csname sga@target@IV.11.9.18\endcsname{ega4.c04.item.11-9-18}
\expandafter\def\csname sga@target@IV.11.9.2\endcsname{ega4.c04.statement.11-9-2}
\expandafter\def\csname sga@target@IV.11.9.3\endcsname{ega4.c04.statement.11-9-3}
\expandafter\def\csname sga@target@IV.11.9.4\endcsname{ega4.c04.statement.11-9-4}
\expandafter\def\csname sga@target@IV.11.9.5\endcsname{ega4.c04.statement.11-9-5}
\expandafter\def\csname sga@target@IV.11.9.6\endcsname{ega4.c04.statement.11-9-6}
\expandafter\def\csname sga@target@IV.11.9.7\endcsname{ega4.c04.prop.11-9-7}
\expandafter\def\csname sga@target@IV.11.9.7.1\endcsname{ega4.c04.lem.11-9-7-1}
\expandafter\def\csname sga@target@IV.11.9.8\endcsname{ega4.c04.subsec.11-9-8}
\expandafter\def\csname sga@target@IV.11.9.9\endcsname{ega4.c04.cor.11-9-9}
\expandafter\def\csname sga@target@IV.12.0.1\endcsname{ega4.c04.statement.12-0-1}
\expandafter\def\csname sga@target@IV.12.0.2\endcsname{ega4.c04.statement.12-0-2}
\expandafter\def\csname sga@target@IV.12.1.1\endcsname{ega4.c04.thm.12-1-1}
\expandafter\def\csname sga@target@IV.12.1.1.1\endcsname{ega4.c04.statement.12-1-1-1}
\expandafter\def\csname sga@target@IV.12.1.1.2\endcsname{ega4.c04.statement.12-1-1-2}
\expandafter\def\csname sga@target@IV.12.1.1.3\endcsname{ega4.c04.statement.12-1-1-3}
\expandafter\def\csname sga@target@IV.12.1.1.4\endcsname{ega4.c04.statement.12-1-1-4}
\expandafter\def\csname sga@target@IV.12.1.1.5\endcsname{ega4.c04.prop.12-1-1-5}
\expandafter\def\csname sga@target@IV.12.1.2\endcsname{ega4.c04.item.12-1-2}
\expandafter\def\csname sga@target@IV.12.1.3\endcsname{ega4.c04.cor.12-1-3}
\expandafter\def\csname sga@target@IV.12.1.4\endcsname{ega4.c04.rem.12-1-4}
\expandafter\def\csname sga@target@IV.12.1.5\endcsname{ega4.c04.cor.12-1-5}
\expandafter\def\csname sga@target@IV.12.1.6\endcsname{ega4.c04.thm.12-1-6}
\expandafter\def\csname sga@target@IV.12.1.7\endcsname{ega4.c04.cor.12-1-7}
\expandafter\def\csname sga@target@IV.12.1.8\endcsname{ega4.c04.item.12-1-8}
\expandafter\def\csname sga@target@IV.12.2.1\endcsname{ega4.c04.thm.12-2-1}
\expandafter\def\csname sga@target@IV.12.2.1.1\endcsname{ega4.c04.statement.12-2-1-1}
\expandafter\def\csname sga@target@IV.12.2.1.2\endcsname{ega4.c04.statement.12-2-1-2}
\expandafter\def\csname sga@target@IV.12.2.1.3\endcsname{ega4.c04.statement.12-2-1-3}
\expandafter\def\csname sga@target@IV.12.2.1.4\endcsname{ega4.c04.lem.12-2-1-4}
\expandafter\def\csname sga@target@IV.12.2.2\endcsname{ega4.c04.cor.12-2-2}
\expandafter\def\csname sga@target@IV.12.2.3\endcsname{ega4.c04.item.12-2-3}
\expandafter\def\csname sga@target@IV.12.2.4\endcsname{ega4.c04.thm.12-2-4}
\expandafter\def\csname sga@target@IV.12.3.1\endcsname{ega4.c04.lem.12-3-1}
\expandafter\def\csname sga@target@IV.12.3.2\endcsname{ega4.c04.prop.12-3-2}
\expandafter\def\csname sga@target@IV.12.3.2.1\endcsname{ega4.c04.eq.12-3-2-1}
\expandafter\def\csname sga@target@IV.12.3.3\endcsname{ega4.c04.prop.12-3-3}
\expandafter\def\csname sga@target@IV.12.3.3.1\endcsname{ega4.c04.lem.12-3-3-1}
\expandafter\def\csname sga@target@IV.12.3.3.2\endcsname{ega4.c04.eq.12-3-3-2}
\expandafter\def\csname sga@target@IV.12.3.3.3\endcsname{ega4.c04.lem.12-3-3-3}
\expandafter\def\csname sga@target@IV.12.3.3.4\endcsname{ega4.c04.eq.12-3-3-4}
\expandafter\def\csname sga@target@IV.12.3.3.5\endcsname{ega4.c04.lem.12-3-3-5}
\expandafter\def\csname sga@target@IV.12.3.4\endcsname{ega4.c04.cor.12-3-4}
\expandafter\def\csname sga@target@IV.13.1.1\endcsname{ega4.c04.lem.13-1-1}
\expandafter\def\csname sga@target@IV.13.1.2\endcsname{ega4.c04.lem.13-1-2}
\expandafter\def\csname sga@target@IV.13.1.3\endcsname{ega4.c04.thm.13-1-3}
\expandafter\def\csname sga@target@IV.13.1.4\endcsname{ega4.c04.cor.13-1-4}
\expandafter\def\csname sga@target@IV.13.1.5\endcsname{ega4.c04.cor.13-1-5}
\expandafter\def\csname sga@target@IV.13.1.6\endcsname{ega4.c04.cor.13-1-6}
\expandafter\def\csname sga@target@IV.13.1.7\endcsname{ega4.c04.prop.13-1-7}
\expandafter\def\csname sga@target@IV.13.2.1\endcsname{ega4.c04.statement.13-2-1}
\expandafter\def\csname sga@target@IV.13.2.10\endcsname{ega4.c04.prop.13-2-10}
\expandafter\def\csname sga@target@IV.13.2.10.1\endcsname{ega4.c04.eq.13-2-10-1}
\expandafter\def\csname sga@target@IV.13.2.10.2\endcsname{ega4.c04.eq.13-2-10-2}
\expandafter\def\csname sga@target@IV.13.2.11\endcsname{ega4.c04.cor.13-2-11}
\expandafter\def\csname sga@target@IV.13.2.12\endcsname{ega4.c04.item.13-2-12}
\expandafter\def\csname sga@target@IV.13.2.2\endcsname{ega4.c04.def.13-2-2}
\expandafter\def\csname sga@target@IV.13.2.3\endcsname{ega4.c04.prop.13-2-3}
\expandafter\def\csname sga@target@IV.13.2.3.1\endcsname{ega4.c04.prop.13-2-3-1}
\expandafter\def\csname sga@target@IV.13.2.4\endcsname{ega4.c04.statement.13-2-4}
\expandafter\def\csname sga@target@IV.13.2.4.1\endcsname{ega4.c04.statement.13-2-4-1}
\expandafter\def\csname sga@target@IV.13.2.4.2\endcsname{ega4.c04.eq.13-2-4-2}
\expandafter\def\csname sga@target@IV.13.2.4.3\endcsname{ega4.c04.eq.13-2-4-3}
\expandafter\def\csname sga@target@IV.13.2.4.4\endcsname{ega4.c04.eq.13-2-4-4}
\expandafter\def\csname sga@target@IV.13.2.4.5\endcsname{ega4.c04.eq.13-2-4-5}
\expandafter\def\csname sga@target@IV.13.2.5\endcsname{ega4.c04.lem.13-2-5}
\expandafter\def\csname sga@target@IV.13.2.6\endcsname{ega4.c04.statement.13-2-6}
\expandafter\def\csname sga@target@IV.13.2.6.1\endcsname{ega4.c04.statement.13-2-6-1}
\expandafter\def\csname sga@target@IV.13.2.7\endcsname{ega4.c04.prop.13-2-7}
\expandafter\def\csname sga@target@IV.13.2.7.1\endcsname{ega4.c04.prop.13-2-7-1}
\expandafter\def\csname sga@target@IV.13.2.8\endcsname{ega4.c04.cor.13-2-8}
\expandafter\def\csname sga@target@IV.13.2.9\endcsname{ega4.c04.cor.13-2-9}
\expandafter\def\csname sga@target@IV.13.2.9.1\endcsname{ega4.c04.eq.13-2-9-1}
\expandafter\def\csname sga@target@IV.13.2.9.2\endcsname{ega4.c04.eq.13-2-9-2}
\expandafter\def\csname sga@target@IV.13.3.1\endcsname{ega4.c04.prop.13-3-1}
\expandafter\def\csname sga@target@IV.13.3.1.1\endcsname{ega4.c04.lem.13-3-1-1}
\expandafter\def\csname sga@target@IV.13.3.2\endcsname{ega4.c04.def.13-3-2}
\expandafter\def\csname sga@target@IV.13.3.3\endcsname{ega4.c04.prop.13-3-3}
\expandafter\def\csname sga@target@IV.13.3.4\endcsname{ega4.c04.cor.13-3-4}
\expandafter\def\csname sga@target@IV.13.3.4.1\endcsname{ega4.c04.cor.13-3-4-1}
\expandafter\def\csname sga@target@IV.13.3.5\endcsname{ega4.c04.cor.13-3-5}
\expandafter\def\csname sga@target@IV.13.3.6\endcsname{ega4.c04.prop.13-3-6}
\expandafter\def\csname sga@target@IV.13.3.6.1\endcsname{ega4.c04.prop.13-3-6-1}
\expandafter\def\csname sga@target@IV.13.3.6.2\endcsname{ega4.c04.eq.13-3-6-2}
\expandafter\def\csname sga@target@IV.13.3.7\endcsname{ega4.c04.prop.13-3-7}
\expandafter\def\csname sga@target@IV.13.3.7.1\endcsname{ega4.c04.prop.13-3-7-1}
\expandafter\def\csname sga@target@IV.13.3.8\endcsname{ega4.c04.prop.13-3-8}
\expandafter\def\csname sga@target@IV.13.3.9\endcsname{ega4.c04.statement.13-3-9}
\expandafter\def\csname sga@target@IV.14.1.1\endcsname{ega4.c04.statement.14-1-1}
\expandafter\def\csname sga@target@IV.14.1.2\endcsname{ega4.c04.prop.14-1-2}
\expandafter\def\csname sga@target@IV.14.1.3\endcsname{ega4.c04.item.14-1-3}
\expandafter\def\csname sga@target@IV.14.1.4\endcsname{ega4.c04.prop.14-1-4}
\expandafter\def\csname sga@target@IV.14.1.5\endcsname{ega4.c04.rem.14-1-5}
\expandafter\def\csname sga@target@IV.14.1.6\endcsname{ega4.c04.statement.14-1-6}
\expandafter\def\csname sga@target@IV.14.2.1\endcsname{ega4.c04.thm.14-2-1}
\expandafter\def\csname sga@target@IV.14.2.1.1\endcsname{ega4.c04.eq.14-2-1-1}
\expandafter\def\csname sga@target@IV.14.2.2\endcsname{ega4.c04.cor.14-2-2}
\expandafter\def\csname sga@target@IV.14.2.3\endcsname{ega4.c04.prop.14-2-3}
\expandafter\def\csname sga@target@IV.14.2.4\endcsname{ega4.c04.cor.14-2-4}
\expandafter\def\csname sga@target@IV.14.2.5\endcsname{ega4.c04.cor.14-2-5}
\expandafter\def\csname sga@target@IV.14.2.6\endcsname{ega4.c04.cor.14-2-6}
\expandafter\def\csname sga@target@IV.14.2.7\endcsname{ega4.c04.rem.14-2-7}
\expandafter\def\csname sga@target@IV.14.3\endcsname{ega4.c04.subsec.14-3.key-iv-14-3}
\expandafter\def\csname sga@target@IV.14.3.1\endcsname{ega4.c04.statement.14-3-1}
\expandafter\def\csname sga@target@IV.14.3.10\endcsname{ega4.c04.lem.14-3-10}
\expandafter\def\csname sga@target@IV.14.3.11\endcsname{ega4.c04.prop.14-3-11}
\expandafter\def\csname sga@target@IV.14.3.12\endcsname{ega4.c04.rem.14-3-12}
\expandafter\def\csname sga@target@IV.14.3.13\endcsname{ega4.c04.prop.14-3-13}
\expandafter\def\csname sga@target@IV.14.3.14\endcsname{ega4.c04.rem.14-3-14}
\expandafter\def\csname sga@target@IV.14.3.2\endcsname{ega4.c04.prop.14-3-2}
\expandafter\def\csname sga@target@IV.14.3.3\endcsname{ega4.c04.def.14-3-3}
\expandafter\def\csname sga@target@IV.14.3.3.1\endcsname{ega4.c04.subsec.14-3-3-1}
\expandafter\def\csname sga@target@IV.14.3.4\endcsname{ega4.c04.prop.14-3-4}
\expandafter\def\csname sga@target@IV.14.3.5\endcsname{ega4.c04.prop.14-3-5}
\expandafter\def\csname sga@target@IV.14.3.6\endcsname{ega4.c04.prop.14-3-6}
\expandafter\def\csname sga@target@IV.14.3.7\endcsname{ega4.c04.cor.14-3-7}
\expandafter\def\csname sga@target@IV.14.3.8\endcsname{ega4.c04.prop.14-3-8}
\expandafter\def\csname sga@target@IV.14.3.9\endcsname{ega4.c04.item.14-3-9}
\expandafter\def\csname sga@target@IV.14.4.1\endcsname{ega4.c04.thm.14-4-1}
\expandafter\def\csname sga@target@IV.14.4.1.1\endcsname{ega4.c04.lem.14-4-1-1}
\expandafter\def\csname sga@target@IV.14.4.1.2\endcsname{ega4.c04.lem.14-4-1-2}
\expandafter\def\csname sga@target@IV.14.4.1.3\endcsname{ega4.c04.lem.14-4-1-3}
\expandafter\def\csname sga@target@IV.14.4.10\endcsname{ega4.c04.item.14-4-10}
\expandafter\def\csname sga@target@IV.14.4.2\endcsname{ega4.c04.cor.14-4-2}
\expandafter\def\csname sga@target@IV.14.4.3\endcsname{ega4.c04.cor.14-4-3}
\expandafter\def\csname sga@target@IV.14.4.4\endcsname{ega4.c04.cor.14-4-4}
\expandafter\def\csname sga@target@IV.14.4.5\endcsname{ega4.c04.rem.14-4-5}
\expandafter\def\csname sga@target@IV.14.4.6\endcsname{ega4.c04.cor.14-4-6}
\expandafter\def\csname sga@target@IV.14.4.7\endcsname{ega4.c04.prop.14-4-7}
\expandafter\def\csname sga@target@IV.14.4.7.1\endcsname{ega4.c04.eq.14-4-7-1}
\expandafter\def\csname sga@target@IV.14.4.7.2\endcsname{ega4.c04.eq.14-4-7-2}
\expandafter\def\csname sga@target@IV.14.4.8\endcsname{ega4.c04.prop.14-4-8}
\expandafter\def\csname sga@target@IV.14.4.8.1\endcsname{ega4.c04.lem.14-4-8-1}
\expandafter\def\csname sga@target@IV.14.4.9\endcsname{ega4.c04.cor.14-4-9}
\expandafter\def\csname sga@target@IV.14.5\endcsname{ega4.c04.subsec.14-5.key-iv-14-5}
\expandafter\def\csname sga@target@IV.14.5.1\endcsname{ega4.c04.lem.14-5-1}
\expandafter\def\csname sga@target@IV.14.5.1.1\endcsname{ega4.c04.eq.14-5-1-1}
\expandafter\def\csname sga@target@IV.14.5.1.2\endcsname{ega4.c04.eq.14-5-1-2}
\expandafter\def\csname sga@target@IV.14.5.1.3\endcsname{ega4.c04.eq.14-5-1-3}
\expandafter\def\csname sga@target@IV.14.5.1.4\endcsname{ega4.c04.eq.14-5-1-4}
\expandafter\def\csname sga@target@IV.14.5.10\endcsname{ega4.c04.prop.14-5-10}
\expandafter\def\csname sga@target@IV.14.5.11\endcsname{ega4.c04.cor.14-5-11}
\expandafter\def\csname sga@target@IV.14.5.11.1\endcsname{ega4.c04.lem.14-5-11-1}
\expandafter\def\csname sga@target@IV.14.5.12\endcsname{ega4.c04.cor.14-5-12}
\expandafter\def\csname sga@target@IV.14.5.2\endcsname{ega4.c04.cor.14-5-2}
\expandafter\def\csname sga@target@IV.14.5.2.1\endcsname{ega4.c04.eq.14-5-2-1}
\expandafter\def\csname sga@target@IV.14.5.3\endcsname{ega4.c04.prop.14-5-3}
\expandafter\def\csname sga@target@IV.14.5.4\endcsname{ega4.c04.cor.14-5-4}
\expandafter\def\csname sga@target@IV.14.5.5\endcsname{ega4.c04.rem.14-5-5}
\expandafter\def\csname sga@target@IV.14.5.6\endcsname{ega4.c04.prop.14-5-6}
\expandafter\def\csname sga@target@IV.14.5.7\endcsname{ega4.c04.rem.14-5-7}
\expandafter\def\csname sga@target@IV.14.5.8\endcsname{ega4.c04.prop.14-5-8}
\expandafter\def\csname sga@target@IV.14.5.9\endcsname{ega4.c04.prop.14-5-9}
\expandafter\def\csname sga@target@IV.15.1.1\endcsname{ega4.c04.prop.15-1-1}
\expandafter\def\csname sga@target@IV.15.1.1-proof\endcsname{ega4.c04.statement.15-1-1-proof}
\expandafter\def\csname sga@target@IV.15.1.1.1\endcsname{ega4.c04.eq.15-1-1-1}
\expandafter\def\csname sga@target@IV.15.1.1.2\endcsname{ega4.c04.eq.15-1-1-2}
\expandafter\def\csname sga@target@IV.15.1.1.3\endcsname{ega4.c04.eq.15-1-1-3}
\expandafter\def\csname sga@target@IV.15.1.1.4\endcsname{ega4.c04.eq.15-1-1-4}
\expandafter\def\csname sga@target@IV.15.1.1.5\endcsname{ega4.c04.eq.15-1-1-5}
\expandafter\def\csname sga@target@IV.15.1.1.6\endcsname{ega4.c04.lem.15-1-1-6}
\expandafter\def\csname sga@target@IV.15.1.2\endcsname{ega4.c04.cor.15-1-2}
\expandafter\def\csname sga@target@IV.15.1.3\endcsname{ega4.c04.cor.15-1-3}
\expandafter\def\csname sga@target@IV.15.1.4\endcsname{ega4.c04.statement.15-1-4}
\expandafter\def\csname sga@target@IV.15.2.1\endcsname{ega4.c04.statement.15-2-1}
\expandafter\def\csname sga@target@IV.15.2.2\endcsname{ega4.c04.thm.15-2-2}
\expandafter\def\csname sga@target@IV.15.2.2-proof\endcsname{ega4.c04.item.15-2-2-proof}
\expandafter\def\csname sga@target@IV.15.2.2.1\endcsname{ega4.c04.lem.15-2-2-1}
\expandafter\def\csname sga@target@IV.15.2.3\endcsname{ega4.c04.cor.15-2-3}
\expandafter\def\csname sga@target@IV.15.2.4\endcsname{ega4.c04.item.15-2-4}
\expandafter\def\csname sga@target@IV.15.3.1\endcsname{ega4.c04.prop.15-3-1}
\expandafter\def\csname sga@target@IV.15.3.1-proof\endcsname{ega4.c04.item.15-3-1-proof}
\expandafter\def\csname sga@target@IV.15.3.2\endcsname{ega4.c04.prop.15-3-2}
\expandafter\def\csname sga@target@IV.15.3.2-proof\endcsname{ega4.c04.item.15-3-2-proof}
\expandafter\def\csname sga@target@IV.15.3.3\endcsname{ega4.c04.cor.15-3-3}
\expandafter\def\csname sga@target@IV.15.3.3-proof\endcsname{ega4.c04.item.15-3-3-proof}
\expandafter\def\csname sga@target@IV.15.3.4\endcsname{ega4.c04.item.15-3-4}
\expandafter\def\csname sga@target@IV.15.4.1\endcsname{ega4.c04.prop.15-4-1}
\expandafter\def\csname sga@target@IV.15.4.2\endcsname{ega4.c04.prop.15-4-2}
\expandafter\def\csname sga@target@IV.15.4.3\endcsname{ega4.c04.prop.15-4-3}
\expandafter\def\csname sga@target@IV.15.4.3-proof\endcsname{ega4.c04.item.15-4-3-proof}
\expandafter\def\csname sga@target@IV.15.5.1\endcsname{ega4.c04.prop.15-5-1}
\expandafter\def\csname sga@target@IV.15.5.10\endcsname{ega4.c04.item.15-5-10}
\expandafter\def\csname sga@target@IV.15.5.2\endcsname{ega4.c04.cor.15-5-2}
\expandafter\def\csname sga@target@IV.15.5.2.1\endcsname{ega4.c04.eq.15-5-2-1}
\expandafter\def\csname sga@target@IV.15.5.2.2\endcsname{ega4.c04.lem.15-5-2-2}
\expandafter\def\csname sga@target@IV.15.5.2.3\endcsname{ega4.c04.eq.15-5-2-3}
\expandafter\def\csname sga@target@IV.15.5.2.4\endcsname{ega4.c04.eq.15-5-2-4}
\expandafter\def\csname sga@target@IV.15.5.3\endcsname{ega4.c04.prop.15-5-3}
\expandafter\def\csname sga@target@IV.15.5.4\endcsname{ega4.c04.cor.15-5-4}
\expandafter\def\csname sga@target@IV.15.5.5\endcsname{ega4.c04.statement.15-5-5}
\expandafter\def\csname sga@target@IV.15.5.6\endcsname{ega4.c04.lem.15-5-6}
\expandafter\def\csname sga@target@IV.15.5.6.1\endcsname{ega4.c04.lem.15-5-6-1}
\expandafter\def\csname sga@target@IV.15.5.7\endcsname{ega4.c04.prop.15-5-7}
\expandafter\def\csname sga@target@IV.15.5.7.1\endcsname{ega4.c04.eq.15-5-7-1}
\expandafter\def\csname sga@target@IV.15.5.8\endcsname{ega4.c04.statement.15-5-8}
\expandafter\def\csname sga@target@IV.15.5.9\endcsname{ega4.c04.prop.15-5-9}
\expandafter\def\csname sga@target@IV.15.6.1\endcsname{ega4.c04.prop.15-6-1}
\expandafter\def\csname sga@target@IV.15.6.1.1\endcsname{ega4.c04.eq.15-6-1-1}
\expandafter\def\csname sga@target@IV.15.6.10\endcsname{ega4.c04.item.15-6-10}
\expandafter\def\csname sga@target@IV.15.6.2\endcsname{ega4.c04.cor.15-6-2}
\expandafter\def\csname sga@target@IV.15.6.3\endcsname{ega4.c04.prop.15-6-3}
\expandafter\def\csname sga@target@IV.15.6.4\endcsname{ega4.c04.prop.15-6-4}
\expandafter\def\csname sga@target@IV.15.6.5\endcsname{ega4.c04.cor.15-6-5}
\expandafter\def\csname sga@target@IV.15.6.6\endcsname{ega4.c04.prop.15-6-6}
\expandafter\def\csname sga@target@IV.15.6.7\endcsname{ega4.c04.cor.15-6-7}
\expandafter\def\csname sga@target@IV.15.6.8\endcsname{ega4.c04.prop.15-6-8}
\expandafter\def\csname sga@target@IV.15.6.9\endcsname{ega4.c04.cor.15-6-9}
\expandafter\def\csname sga@target@IV.15.7.1\endcsname{ega4.c04.def.15-7-1}
\expandafter\def\csname sga@target@IV.15.7.10\endcsname{ega4.c04.cor.15-7-10}
\expandafter\def\csname sga@target@IV.15.7.11\endcsname{ega4.c04.cor.15-7-11}
\expandafter\def\csname sga@target@IV.15.7.11.1\endcsname{ega4.c04.lem.15-7-11-1}
\expandafter\def\csname sga@target@IV.15.7.2\endcsname{ega4.c04.prop.15-7-2}
\expandafter\def\csname sga@target@IV.15.7.3\endcsname{ega4.c04.rem.15-7-3}
\expandafter\def\csname sga@target@IV.15.7.4\endcsname{ega4.c04.cor.15-7-4}
\expandafter\def\csname sga@target@IV.15.7.5\endcsname{ega4.c04.cor.15-7-5}
\expandafter\def\csname sga@target@IV.15.7.6\endcsname{ega4.c04.cor.15-7-6}
\expandafter\def\csname sga@target@IV.15.7.7\endcsname{ega4.c04.item.15-7-7}
\expandafter\def\csname sga@target@IV.15.7.8\endcsname{ega4.c04.statement.15-7-8}
\expandafter\def\csname sga@target@IV.15.7.9\endcsname{ega4.c04.prop.15-7-9}
\expandafter\def\csname sga@target@IV.16.1\endcsname{ega4.c04.subsec.16-1}
\expandafter\def\csname sga@target@IV.16.1.1\endcsname{ega4.c04.statement.16-1-1}
\expandafter\def\csname sga@target@IV.16.1.10\endcsname{ega4.c04.cor.16-1-10}
\expandafter\def\csname sga@target@IV.16.1.11\endcsname{ega4.c04.rem.16-1-11}
\expandafter\def\csname sga@target@IV.16.1.2\endcsname{ega4.c04.def.16-1-2}
\expandafter\def\csname sga@target@IV.16.1.2.1\endcsname{ega4.c04.def.16-1-2-1}
\expandafter\def\csname sga@target@IV.16.1.2.2\endcsname{ega4.c04.eq.16-1-2-2}
\expandafter\def\csname sga@target@IV.16.1.3\endcsname{ega4.c04.statement.16-1-3}
\expandafter\def\csname sga@target@IV.16.1.4\endcsname{ega4.c04.statement.16-1-4}
\expandafter\def\csname sga@target@IV.16.1.5\endcsname{ega4.c04.prop.16-1-5}
\expandafter\def\csname sga@target@IV.16.1.6\endcsname{ega4.c04.cor.16-1-6}
\expandafter\def\csname sga@target@IV.16.1.7\endcsname{ega4.c04.cor.16-1-7}
\expandafter\def\csname sga@target@IV.16.1.8\endcsname{ega4.c04.prop.16-1-8}
\expandafter\def\csname sga@target@IV.16.1.9\endcsname{ega4.c04.prop.16-1-9}
\expandafter\def\csname sga@target@IV.16.10\endcsname{ega4.c04.subsec.16-10}
\expandafter\def\csname sga@target@IV.16.10.1\endcsname{ega4.c04.def.16-10-1}
\expandafter\def\csname sga@target@IV.16.10.2\endcsname{ega4.c04.prop.16-10-2}
\expandafter\def\csname sga@target@IV.16.10.3\endcsname{ega4.c04.prop.16-10-3}
\expandafter\def\csname sga@target@IV.16.10.4\endcsname{ega4.c04.prop.16-10-4}
\expandafter\def\csname sga@target@IV.16.10.5\endcsname{ega4.c04.prop.16-10-5}
\expandafter\def\csname sga@target@IV.16.10.6\endcsname{ega4.c04.rem.16-10-6}
\expandafter\def\csname sga@target@IV.16.11\endcsname{ega4.c04.subsec.16-11}
\expandafter\def\csname sga@target@IV.16.11.1\endcsname{ega4.c04.statement.16-11-1}
\expandafter\def\csname sga@target@IV.16.11.1.1\endcsname{ega4.c04.statement.16-11-1-1}
\expandafter\def\csname sga@target@IV.16.11.1.2\endcsname{ega4.c04.statement.16-11-1-2}
\expandafter\def\csname sga@target@IV.16.11.1.3\endcsname{ega4.c04.eq.16-11-1-3}
\expandafter\def\csname sga@target@IV.16.11.1.4\endcsname{ega4.c04.eq.16-11-1-4}
\expandafter\def\csname sga@target@IV.16.11.1.5\endcsname{ega4.c04.eq.16-11-1-5}
\expandafter\def\csname sga@target@IV.16.11.1.6\endcsname{ega4.c04.eq.16-11-1-6}
\expandafter\def\csname sga@target@IV.16.11.1.7\endcsname{ega4.c04.eq.16-11-1-7}
\expandafter\def\csname sga@target@IV.16.11.2\endcsname{ega4.c04.thm.16-11-2}
\expandafter\def\csname sga@target@IV.16.11.2.1\endcsname{ega4.c04.eq.16-11-2-1}
\expandafter\def\csname sga@target@IV.16.11.2.2\endcsname{ega4.c04.eq.16-11-2-2}
\expandafter\def\csname sga@target@IV.16.11.2.3\endcsname{ega4.c04.eq.16-11-2-3}
\expandafter\def\csname sga@target@IV.16.11.3\endcsname{ega4.c04.statement.16-11-3}
\expandafter\def\csname sga@target@IV.16.12\endcsname{ega4.c04.subsec.16-12}
\expandafter\def\csname sga@target@IV.16.12.1\endcsname{ega4.c04.statement.16-12-1}
\expandafter\def\csname sga@target@IV.16.12.2\endcsname{ega4.c04.thm.16-12-2}
\expandafter\def\csname sga@target@IV.16.12.3\endcsname{ega4.c04.statement.16-12-3}
\expandafter\def\csname sga@target@IV.16.2\endcsname{ega4.c04.subsec.16-2}
\expandafter\def\csname sga@target@IV.16.2.1\endcsname{ega4.c04.statement.16-2-1}
\expandafter\def\csname sga@target@IV.16.2.1.1\endcsname{ega4.c04.statement.16-2-1-1}
\expandafter\def\csname sga@target@IV.16.2.1.2\endcsname{ega4.c04.eq.16-2-1-2}
\expandafter\def\csname sga@target@IV.16.2.1.3\endcsname{ega4.c04.eq.16-2-1-3}
\expandafter\def\csname sga@target@IV.16.2.1.4\endcsname{ega4.c04.eq.16-2-1-4}
\expandafter\def\csname sga@target@IV.16.2.2\endcsname{ega4.c04.prop.16-2-2}
\expandafter\def\csname sga@target@IV.16.2.3\endcsname{ega4.c04.cor.16-2-3}
\expandafter\def\csname sga@target@IV.16.2.3.1\endcsname{ega4.c04.eq.16-2-3-1}
\expandafter\def\csname sga@target@IV.16.2.3.2\endcsname{ega4.c04.eq.16-2-3-2}
\expandafter\def\csname sga@target@IV.16.2.4\endcsname{ega4.c04.cor.16-2-4}
\expandafter\def\csname sga@target@IV.16.2.5\endcsname{ega4.c04.statement.16-2-5}
\expandafter\def\csname sga@target@IV.16.2.6\endcsname{ega4.c04.statement.16-2-6}
\expandafter\def\csname sga@target@IV.16.2.6.1\endcsname{ega4.c04.statement.16-2-6-1}
\expandafter\def\csname sga@target@IV.16.2.6.2\endcsname{ega4.c04.eq.16-2-6-2}
\expandafter\def\csname sga@target@IV.16.2.7\endcsname{ega4.c04.prop.16-2-7}
\expandafter\def\csname sga@target@IV.16.2.7.1\endcsname{ega4.c04.prop.16-2-7-1}
\expandafter\def\csname sga@target@IV.16.3\endcsname{ega4.c04.subsec.16-3}
\expandafter\def\csname sga@target@IV.16.3.1\endcsname{ega4.c04.def.16-3-1}
\expandafter\def\csname sga@target@IV.16.3.1.1\endcsname{ega4.c04.eq.16-3-1-1}
\expandafter\def\csname sga@target@IV.16.3.2\endcsname{ega4.c04.statement.16-3-2}
\expandafter\def\csname sga@target@IV.16.3.3\endcsname{ega4.c04.statement.16-3-3}
\expandafter\def\csname sga@target@IV.16.3.3.1\endcsname{ega4.c04.statement.16-3-3-1}
\expandafter\def\csname sga@target@IV.16.3.4\endcsname{ega4.c04.prop.16-3-4}
\expandafter\def\csname sga@target@IV.16.3.5\endcsname{ega4.c04.statement.16-3-5}
\expandafter\def\csname sga@target@IV.16.3.6\endcsname{ega4.c04.def.16-3-6}
\expandafter\def\csname sga@target@IV.16.3.6.1\endcsname{ega4.c04.eq.16-3-6-1}
\expandafter\def\csname sga@target@IV.16.3.7\endcsname{ega4.c04.statement.16-3-7}
\expandafter\def\csname sga@target@IV.16.3.7.1\endcsname{ega4.c04.eq.16-3-7-1}
\expandafter\def\csname sga@target@IV.16.3.8\endcsname{ega4.c04.prop.16-3-8}
\expandafter\def\csname sga@target@IV.16.3.9\endcsname{ega4.c04.prop.16-3-9}
\expandafter\def\csname sga@target@IV.16.4\endcsname{ega4.c04.subsec.16-4}
\expandafter\def\csname sga@target@IV.16.4.1\endcsname{ega4.c04.statement.16-4-1}
\expandafter\def\csname sga@target@IV.16.4.1.1\endcsname{ega4.c04.statement.16-4-1-1}
\expandafter\def\csname sga@target@IV.16.4.1.2\endcsname{ega4.c04.statement.16-4-1-2}
\expandafter\def\csname sga@target@IV.16.4.1.3\endcsname{ega4.c04.eq.16-4-1-3}
\expandafter\def\csname sga@target@IV.16.4.1.4\endcsname{ega4.c04.eq.16-4-1-4}
\expandafter\def\csname sga@target@IV.16.4.1.5\endcsname{ega4.c04.eq.16-4-1-5}
\expandafter\def\csname sga@target@IV.16.4.10\endcsname{ega4.c04.cor.16-4-10}
\expandafter\def\csname sga@target@IV.16.4.11\endcsname{ega4.c04.prop.16-4-11}
\expandafter\def\csname sga@target@IV.16.4.11.1\endcsname{ega4.c04.prop.16-4-11-1}
\expandafter\def\csname sga@target@IV.16.4.11.2\endcsname{ega4.c04.prop.16-4-11-2}
\expandafter\def\csname sga@target@IV.16.4.11.3\endcsname{ega4.c04.eq.16-4-11-3}
\expandafter\def\csname sga@target@IV.16.4.11.4\endcsname{ega4.c04.eq.16-4-11-4}
\expandafter\def\csname sga@target@IV.16.4.12\endcsname{ega4.c04.cor.16-4-12}
\expandafter\def\csname sga@target@IV.16.4.13\endcsname{ega4.c04.cor.16-4-13}
\expandafter\def\csname sga@target@IV.16.4.14\endcsname{ega4.c04.prop.16-4-14}
\expandafter\def\csname sga@target@IV.16.4.14.1\endcsname{ega4.c04.prop.16-4-14-1}
\expandafter\def\csname sga@target@IV.16.4.15\endcsname{ega4.c04.cor.16-4-15}
\expandafter\def\csname sga@target@IV.16.4.15.1\endcsname{ega4.c04.cor.16-4-15-1}
\expandafter\def\csname sga@target@IV.16.4.15.2\endcsname{ega4.c04.eq.16-4-15-2}
\expandafter\def\csname sga@target@IV.16.4.16\endcsname{ega4.c04.cor.16-4-16}
\expandafter\def\csname sga@target@IV.16.4.16.1\endcsname{ega4.c04.cor.16-4-16-1}
\expandafter\def\csname sga@target@IV.16.4.16.2\endcsname{ega4.c04.eq.16-4-16-2}
\expandafter\def\csname sga@target@IV.16.4.17\endcsname{ega4.c04.cor.16-4-17}
\expandafter\def\csname sga@target@IV.16.4.18\endcsname{ega4.c04.statement.16-4-18}
\expandafter\def\csname sga@target@IV.16.4.18.1\endcsname{ega4.c04.statement.16-4-18-1}
\expandafter\def\csname sga@target@IV.16.4.18.2\endcsname{ega4.c04.statement.16-4-18-2}
\expandafter\def\csname sga@target@IV.16.4.18.3\endcsname{ega4.c04.eq.16-4-18-3}
\expandafter\def\csname sga@target@IV.16.4.19\endcsname{ega4.c04.cor.16-4-19}
\expandafter\def\csname sga@target@IV.16.4.19.1\endcsname{ega4.c04.cor.16-4-19-1}
\expandafter\def\csname sga@target@IV.16.4.2\endcsname{ega4.c04.statement.16-4-2}
\expandafter\def\csname sga@target@IV.16.4.20\endcsname{ega4.c04.prop.16-4-20}
\expandafter\def\csname sga@target@IV.16.4.21\endcsname{ega4.c04.cor.16-4-21}
\expandafter\def\csname sga@target@IV.16.4.21.1\endcsname{ega4.c04.cor.16-4-21-1}
\expandafter\def\csname sga@target@IV.16.4.22\endcsname{ega4.c04.cor.16-4-22}
\expandafter\def\csname sga@target@IV.16.4.23\endcsname{ega4.c04.prop.16-4-23}
\expandafter\def\csname sga@target@IV.16.4.23.1\endcsname{ega4.c04.prop.16-4-23-1}
\expandafter\def\csname sga@target@IV.16.4.24\endcsname{ega4.c04.statement.16-4-24}
\expandafter\def\csname sga@target@IV.16.4.3\endcsname{ega4.c04.statement.16-4-3}
\expandafter\def\csname sga@target@IV.16.4.3.1\endcsname{ega4.c04.statement.16-4-3-1}
\expandafter\def\csname sga@target@IV.16.4.3.2\endcsname{ega4.c04.statement.16-4-3-2}
\expandafter\def\csname sga@target@IV.16.4.3.3\endcsname{ega4.c04.eq.16-4-3-3}
\expandafter\def\csname sga@target@IV.16.4.3.4\endcsname{ega4.c04.eq.16-4-3-4}
\expandafter\def\csname sga@target@IV.16.4.3.5\endcsname{ega4.c04.eq.16-4-3-5}
\expandafter\def\csname sga@target@IV.16.4.3.6\endcsname{ega4.c04.eq.16-4-3-6}
\expandafter\def\csname sga@target@IV.16.4.3.7\endcsname{ega4.c04.eq.16-4-3-7}
\expandafter\def\csname sga@target@IV.16.4.4\endcsname{ega4.c04.statement.16-4-4}
\expandafter\def\csname sga@target@IV.16.4.5\endcsname{ega4.c04.prop.16-4-5}
\expandafter\def\csname sga@target@IV.16.4.6\endcsname{ega4.c04.cor.16-4-6}
\expandafter\def\csname sga@target@IV.16.4.7\endcsname{ega4.c04.statement.16-4-7}
\expandafter\def\csname sga@target@IV.16.4.7.1\endcsname{ega4.c04.statement.16-4-7-1}
\expandafter\def\csname sga@target@IV.16.4.7.2\endcsname{ega4.c04.eq.16-4-7-2}
\expandafter\def\csname sga@target@IV.16.4.7.3\endcsname{ega4.c04.eq.16-4-7-3}
\expandafter\def\csname sga@target@IV.16.4.8\endcsname{ega4.c04.prop.16-4-8}
\expandafter\def\csname sga@target@IV.16.4.8.1\endcsname{ega4.c04.prop.16-4-8-1}
\expandafter\def\csname sga@target@IV.16.4.8.2\endcsname{ega4.c04.eq.16-4-8-2}
\expandafter\def\csname sga@target@IV.16.4.8.3\endcsname{ega4.c04.lem.16-4-8-3}
\expandafter\def\csname sga@target@IV.16.4.8.4\endcsname{ega4.c04.eq.16-4-8-4}
\expandafter\def\csname sga@target@IV.16.4.9\endcsname{ega4.c04.cor.16-4-9}
\expandafter\def\csname sga@target@IV.16.4.9.1\endcsname{ega4.c04.cor.16-4-9-1}
\expandafter\def\csname sga@target@IV.16.5\endcsname{ega4.c04.subsec.16-5}
\expandafter\def\csname sga@target@IV.16.5.1\endcsname{ega4.c04.statement.16-5-1}
\expandafter\def\csname sga@target@IV.16.5.1.1\endcsname{ega4.c04.eq.16-5-1-1}
\expandafter\def\csname sga@target@IV.16.5.1.2\endcsname{ega4.c04.eq.16-5-1-2}
\expandafter\def\csname sga@target@IV.16.5.10\endcsname{ega4.c04.statement.16-5-10}
\expandafter\def\csname sga@target@IV.16.5.10.1\endcsname{ega4.c04.statement.16-5-10-1}
\expandafter\def\csname sga@target@IV.16.5.10.2\endcsname{ega4.c04.statement.16-5-10-2}
\expandafter\def\csname sga@target@IV.16.5.11\endcsname{ega4.c04.prop.16-5-11}
\expandafter\def\csname sga@target@IV.16.5.12\endcsname{ega4.c04.statement.16-5-12}
\expandafter\def\csname sga@target@IV.16.5.12.1\endcsname{ega4.c04.statement.16-5-12-1}
\expandafter\def\csname sga@target@IV.16.5.12.2\endcsname{ega4.c04.eq.16-5-12-2}
\expandafter\def\csname sga@target@IV.16.5.13\endcsname{ega4.c04.statement.16-5-13}
\expandafter\def\csname sga@target@IV.16.5.13.1\endcsname{ega4.c04.statement.16-5-13-1}
\expandafter\def\csname sga@target@IV.16.5.13.2\endcsname{ega4.c04.eq.16-5-13-2}
\expandafter\def\csname sga@target@IV.16.5.13.3\endcsname{ega4.c04.eq.16-5-13-3}
\expandafter\def\csname sga@target@IV.16.5.13.4\endcsname{ega4.c04.eq.16-5-13-4}
\expandafter\def\csname sga@target@IV.16.5.13.5\endcsname{ega4.c04.eq.16-5-13-5}
\expandafter\def\csname sga@target@IV.16.5.14\endcsname{ega4.c04.statement.16-5-14}
\expandafter\def\csname sga@target@IV.16.5.14.1\endcsname{ega4.c04.statement.16-5-14-1}
\expandafter\def\csname sga@target@IV.16.5.14.2\endcsname{ega4.c04.eq.16-5-14-2}
\expandafter\def\csname sga@target@IV.16.5.14.3\endcsname{ega4.c04.eq.16-5-14-3}
\expandafter\def\csname sga@target@IV.16.5.15\endcsname{ega4.c04.statement.16-5-15}
\expandafter\def\csname sga@target@IV.16.5.16\endcsname{ega4.c04.prop.16-5-16}
\expandafter\def\csname sga@target@IV.16.5.17\endcsname{ega4.c04.prop.16-5-17}
\expandafter\def\csname sga@target@IV.16.5.18\endcsname{ega4.c04.cor.16-5-18}
\expandafter\def\csname sga@target@IV.16.5.2\endcsname{ega4.c04.prop.16-5-2}
\expandafter\def\csname sga@target@IV.16.5.3\endcsname{ega4.c04.prop.16-5-3}
\expandafter\def\csname sga@target@IV.16.5.3.1\endcsname{ega4.c04.eq.16-5-3-1}
\expandafter\def\csname sga@target@IV.16.5.4\endcsname{ega4.c04.statement.16-5-4}
\expandafter\def\csname sga@target@IV.16.5.5\endcsname{ega4.c04.cor.16-5-5}
\expandafter\def\csname sga@target@IV.16.5.5.1\endcsname{ega4.c04.cor.16-5-5-1}
\expandafter\def\csname sga@target@IV.16.5.6\endcsname{ega4.c04.cor.16-5-6}
\expandafter\def\csname sga@target@IV.16.5.7\endcsname{ega4.c04.statement.16-5-7}
\expandafter\def\csname sga@target@IV.16.5.7.1\endcsname{ega4.c04.statement.16-5-7-1}
\expandafter\def\csname sga@target@IV.16.5.8\endcsname{ega4.c04.statement.16-5-8}
\expandafter\def\csname sga@target@IV.16.5.8.1\endcsname{ega4.c04.eq.16-5-8-1}
\expandafter\def\csname sga@target@IV.16.5.8.2\endcsname{ega4.c04.eq.16-5-8-2}
\expandafter\def\csname sga@target@IV.16.5.9\endcsname{ega4.c04.statement.16-5-9}
\expandafter\def\csname sga@target@IV.16.5.9.1\endcsname{ega4.c04.statement.16-5-9-1}
\expandafter\def\csname sga@target@IV.16.6\endcsname{ega4.c04.subsec.16-6}
\expandafter\def\csname sga@target@IV.16.6.1\endcsname{ega4.c04.statement.16-6-1}
\expandafter\def\csname sga@target@IV.16.6.1.1\endcsname{ega4.c04.statement.16-6-1-1}
\expandafter\def\csname sga@target@IV.16.6.1.2\endcsname{ega4.c04.statement.16-6-1-2}
\expandafter\def\csname sga@target@IV.16.6.2\endcsname{ega4.c04.thm.16-6-2}
\expandafter\def\csname sga@target@IV.16.6.2.1\endcsname{ega4.c04.eq.16-6-2-1}
\expandafter\def\csname sga@target@IV.16.6.2.2\endcsname{ega4.c04.eq.16-6-2-2}
\expandafter\def\csname sga@target@IV.16.6.2.3\endcsname{ega4.c04.eq.16-6-2-3}
\expandafter\def\csname sga@target@IV.16.6.2.4\endcsname{ega4.c04.eq.16-6-2-4}
\expandafter\def\csname sga@target@IV.16.6.2.5\endcsname{ega4.c04.eq.16-6-2-5}
\expandafter\def\csname sga@target@IV.16.6.3\endcsname{ega4.c04.def.16-6-3}
\expandafter\def\csname sga@target@IV.16.6.4\endcsname{ega4.c04.prop.16-6-4}
\expandafter\def\csname sga@target@IV.16.6.4.1\endcsname{ega4.c04.prop.16-6-4-1}
\expandafter\def\csname sga@target@IV.16.6.5\endcsname{ega4.c04.statement.16-6-5}
\expandafter\def\csname sga@target@IV.16.6.5.1\endcsname{ega4.c04.statement.16-6-5-1}
\expandafter\def\csname sga@target@IV.16.7\endcsname{ega4.c04.subsec.16-7}
\expandafter\def\csname sga@target@IV.16.7.1\endcsname{ega4.c04.statement.16-7-1}
\expandafter\def\csname sga@target@IV.16.7.1.1\endcsname{ega4.c04.eq.16-7-1-1}
\expandafter\def\csname sga@target@IV.16.7.1.2\endcsname{ega4.c04.eq.16-7-1-2}
\expandafter\def\csname sga@target@IV.16.7.10\endcsname{ega4.c04.rem.16-7-10}
\expandafter\def\csname sga@target@IV.16.7.2\endcsname{ega4.c04.statement.16-7-2}
\expandafter\def\csname sga@target@IV.16.7.2.1\endcsname{ega4.c04.statement.16-7-2-1}
\expandafter\def\csname sga@target@IV.16.7.3\endcsname{ega4.c04.prop.16-7-3}
\expandafter\def\csname sga@target@IV.16.7.4\endcsname{ega4.c04.statement.16-7-4}
\expandafter\def\csname sga@target@IV.16.7.4.1\endcsname{ega4.c04.statement.16-7-4-1}
\expandafter\def\csname sga@target@IV.16.7.5\endcsname{ega4.c04.statement.16-7-5}
\expandafter\def\csname sga@target@IV.16.7.5.1\endcsname{ega4.c04.statement.16-7-5-1}
\expandafter\def\csname sga@target@IV.16.7.5.2\endcsname{ega4.c04.statement.16-7-5-2}
\expandafter\def\csname sga@target@IV.16.7.5.3\endcsname{ega4.c04.statement.16-7-5-3}
\expandafter\def\csname sga@target@IV.16.7.5.4\endcsname{ega4.c04.statement.16-7-5-4}
\expandafter\def\csname sga@target@IV.16.7.6\endcsname{ega4.c04.prop.16-7-6}
\expandafter\def\csname sga@target@IV.16.7.7\endcsname{ega4.c04.statement.16-7-7}
\expandafter\def\csname sga@target@IV.16.7.7.1\endcsname{ega4.c04.eq.16-7-7-1}
\expandafter\def\csname sga@target@IV.16.7.7.2\endcsname{ega4.c04.eq.16-7-7-2}
\expandafter\def\csname sga@target@IV.16.7.8\endcsname{ega4.c04.statement.16-7-8}
\expandafter\def\csname sga@target@IV.16.7.8.1\endcsname{ega4.c04.statement.16-7-8-1}
\expandafter\def\csname sga@target@IV.16.7.8.2\endcsname{ega4.c04.eq.16-7-8-2}
\expandafter\def\csname sga@target@IV.16.7.9\endcsname{ega4.c04.statement.16-7-9}
\expandafter\def\csname sga@target@IV.16.7.9.1\endcsname{ega4.c04.statement.16-7-9-1}
\expandafter\def\csname sga@target@IV.16.8\endcsname{ega4.c04.subsec.16-8}
\expandafter\def\csname sga@target@IV.16.8.1\endcsname{ega4.c04.def.16-8-1}
\expandafter\def\csname sga@target@IV.16.8.10\endcsname{ega4.c04.cor.16-8-10}
\expandafter\def\csname sga@target@IV.16.8.11\endcsname{ega4.c04.statement.16-8-11}
\expandafter\def\csname sga@target@IV.16.8.2\endcsname{ega4.c04.statement.16-8-2}
\expandafter\def\csname sga@target@IV.16.8.3\endcsname{ega4.c04.statement.16-8-3}
\expandafter\def\csname sga@target@IV.16.8.3.1\endcsname{ega4.c04.eq.16-8-3-1}
\expandafter\def\csname sga@target@IV.16.8.4\endcsname{ega4.c04.prop.16-8-4}
\expandafter\def\csname sga@target@IV.16.8.4.1\endcsname{ega4.c04.prop.16-8-4-1}
\expandafter\def\csname sga@target@IV.16.8.5\endcsname{ega4.c04.statement.16-8-5}
\expandafter\def\csname sga@target@IV.16.8.5.1\endcsname{ega4.c04.eq.16-8-5-1}
\expandafter\def\csname sga@target@IV.16.8.5.2\endcsname{ega4.c04.eq.16-8-5-2}
\expandafter\def\csname sga@target@IV.16.8.6\endcsname{ega4.c04.prop.16-8-6}
\expandafter\def\csname sga@target@IV.16.8.7\endcsname{ega4.c04.statement.16-8-7}
\expandafter\def\csname sga@target@IV.16.8.7.1\endcsname{ega4.c04.eq.16-8-7-1}
\expandafter\def\csname sga@target@IV.16.8.8\endcsname{ega4.c04.prop.16-8-8}
\expandafter\def\csname sga@target@IV.16.8.8.1\endcsname{ega4.c04.eq.16-8-8-1}
\expandafter\def\csname sga@target@IV.16.8.8.2\endcsname{ega4.c04.eq.16-8-8-2}
\expandafter\def\csname sga@target@IV.16.8.9\endcsname{ega4.c04.prop.16-8-9}
\expandafter\def\csname sga@target@IV.16.8.9.1\endcsname{ega4.c04.lem.16-8-9-1}
\expandafter\def\csname sga@target@IV.16.8.9.2\endcsname{ega4.c04.lem.16-8-9-2}
\expandafter\def\csname sga@target@IV.16.8.9.3\endcsname{ega4.c04.lem.16-8-9-3}
\expandafter\def\csname sga@target@IV.16.8.9.4\endcsname{ega4.c04.eq.16-8-9-4}
\expandafter\def\csname sga@target@IV.16.9\endcsname{ega4.c04.subsec.16-9}
\expandafter\def\csname sga@target@IV.16.9.1\endcsname{ega4.c04.def.16-9-1}
\expandafter\def\csname sga@target@IV.16.9.10\endcsname{ega4.c04.prop.16-9-10}
\expandafter\def\csname sga@target@IV.16.9.11\endcsname{ega4.c04.cor.16-9-11}
\expandafter\def\csname sga@target@IV.16.9.12\endcsname{ega4.c04.statement.16-9-12}
\expandafter\def\csname sga@target@IV.16.9.13\endcsname{ega4.c04.prop.16-9-13}
\expandafter\def\csname sga@target@IV.16.9.13.1\endcsname{ega4.c04.prop.16-9-13-1}
\expandafter\def\csname sga@target@IV.16.9.13.2\endcsname{ega4.c04.lem.16-9-13-2}
\expandafter\def\csname sga@target@IV.16.9.13.3\endcsname{ega4.c04.lem.16-9-13-3}
\expandafter\def\csname sga@target@IV.16.9.13.4\endcsname{ega4.c04.eq.16-9-13-4}
\expandafter\def\csname sga@target@IV.16.9.13.5\endcsname{ega4.c04.cor.16-9-13-5}
\expandafter\def\csname sga@target@IV.16.9.13.6\endcsname{ega4.c04.cor.16-9-13-6}
\expandafter\def\csname sga@target@IV.16.9.13.7\endcsname{ega4.c04.eq.16-9-13-7}
\expandafter\def\csname sga@target@IV.16.9.2\endcsname{ega4.c04.def.16-9-2}
\expandafter\def\csname sga@target@IV.16.9.3\endcsname{ega4.c04.prop.16-9-3}
\expandafter\def\csname sga@target@IV.16.9.4\endcsname{ega4.c04.cor.16-9-4}
\expandafter\def\csname sga@target@IV.16.9.4.1\endcsname{ega4.c04.eq.16-9-4-1}
\expandafter\def\csname sga@target@IV.16.9.5\endcsname{ega4.c04.cor.16-9-5}
\expandafter\def\csname sga@target@IV.16.9.6\endcsname{ega4.c04.statement.16-9-6}
\expandafter\def\csname sga@target@IV.16.9.7\endcsname{ega4.c04.statement.16-9-7}
\expandafter\def\csname sga@target@IV.16.9.8\endcsname{ega4.c04.prop.16-9-8}
\expandafter\def\csname sga@target@IV.16.9.9\endcsname{ega4.c04.cor.16-9-9}
\expandafter\def\csname sga@target@IV.17.1.1\endcsname{ega4.c04.def.17-1-1}
\expandafter\def\csname sga@target@IV.17.1.1.1\endcsname{ega4.c04.def.17-1-1-1}
\expandafter\def\csname sga@target@IV.17.1.2\endcsname{ega4.c04.rem.17-1-2}
\expandafter\def\csname sga@target@IV.17.1.2.1\endcsname{ega4.c04.eq.17-1-2-1}
\expandafter\def\csname sga@target@IV.17.1.3\endcsname{ega4.c04.prop.17-1-3}
\expandafter\def\csname sga@target@IV.17.1.4\endcsname{ega4.c04.prop.17-1-4}
\expandafter\def\csname sga@target@IV.17.1.5\endcsname{ega4.c04.cor.17-1-5}
\expandafter\def\csname sga@target@IV.17.1.6\endcsname{ega4.c04.prop.17-1-6}
\expandafter\def\csname sga@target@IV.17.10.1\endcsname{ega4.c04.def.17-10-1}
\expandafter\def\csname sga@target@IV.17.10.2\endcsname{ega4.c04.prop.17-10-2}
\expandafter\def\csname sga@target@IV.17.10.3\endcsname{ega4.c04.cor.17-10-3}
\expandafter\def\csname sga@target@IV.17.10.3.1\endcsname{ega4.c04.cor.17-10-3-1}
\expandafter\def\csname sga@target@IV.17.10.4\endcsname{ega4.c04.cor.17-10-4}
\expandafter\def\csname sga@target@IV.17.10.4.1\endcsname{ega4.c04.eq.17-10-4-1}
\expandafter\def\csname sga@target@IV.17.11.1\endcsname{ega4.c04.thm.17-11-1}
\expandafter\def\csname sga@target@IV.17.11.1.1\endcsname{ega4.c04.eq.17-11-1-1}
\expandafter\def\csname sga@target@IV.17.11.1.2\endcsname{ega4.c04.eq.17-11-1-2}
\expandafter\def\csname sga@target@IV.17.11.2\endcsname{ega4.c04.cor.17-11-2}
\expandafter\def\csname sga@target@IV.17.11.3\endcsname{ega4.c04.cor.17-11-3}
\expandafter\def\csname sga@target@IV.17.11.4\endcsname{ega4.c04.cor.17-11-4}
\expandafter\def\csname sga@target@IV.17.11.5\endcsname{ega4.c04.prop.17-11-5}
\expandafter\def\csname sga@target@IV.17.12.1\endcsname{ega4.c04.thm.17-12-1}
\expandafter\def\csname sga@target@IV.17.12.2\endcsname{ega4.c04.cor.17-12-2}
\expandafter\def\csname sga@target@IV.17.12.3\endcsname{ega4.c04.cor.17-12-3}
\expandafter\def\csname sga@target@IV.17.12.4\endcsname{ega4.c04.prop.17-12-4}
\expandafter\def\csname sga@target@IV.17.12.5\endcsname{ega4.c04.prop.17-12-5}
\expandafter\def\csname sga@target@IV.17.12.6\endcsname{ega4.c04.cor.17-12-6}
\expandafter\def\csname sga@target@IV.17.12.7\endcsname{ega4.c04.ex.17-12-7}
\expandafter\def\csname sga@target@IV.17.12.8\endcsname{ega4.c04.rem.17-12-8}
\expandafter\def\csname sga@target@IV.17.13\endcsname{ega4.c04.subsec.17-13.key-iv-17-13}
\expandafter\def\csname sga@target@IV.17.13.1\endcsname{ega4.c04.statement.17-13-1}
\expandafter\def\csname sga@target@IV.17.13.1.1\endcsname{ega4.c04.statement.17-13-1-1}
\expandafter\def\csname sga@target@IV.17.13.1.2\endcsname{ega4.c04.statement.17-13-1-2}
\expandafter\def\csname sga@target@IV.17.13.1.3\endcsname{ega4.c04.eq.17-13-1-3}
\expandafter\def\csname sga@target@IV.17.13.10\endcsname{ega4.c04.cor.17-13-10}
\expandafter\def\csname sga@target@IV.17.13.11\endcsname{ega4.c04.statement.17-13-11}
\expandafter\def\csname sga@target@IV.17.13.11.1\endcsname{ega4.c04.statement.17-13-11-1}
\expandafter\def\csname sga@target@IV.17.13.12\endcsname{ega4.c04.cor.17-13-12}
\expandafter\def\csname sga@target@IV.17.13.12.1\endcsname{ega4.c04.cor.17-13-12-1}
\expandafter\def\csname sga@target@IV.17.13.12.2\endcsname{ega4.c04.eq.17-13-12-2}
\expandafter\def\csname sga@target@IV.17.13.2\endcsname{ega4.c04.prop.17-13-2}
\expandafter\def\csname sga@target@IV.17.13.2.1\endcsname{ega4.c04.eq.17-13-2-1}
\expandafter\def\csname sga@target@IV.17.13.2.2\endcsname{ega4.c04.eq.17-13-2-2}
\expandafter\def\csname sga@target@IV.17.13.3\endcsname{ega4.c04.def.17-13-3}
\expandafter\def\csname sga@target@IV.17.13.4\endcsname{ega4.c04.statement.17-13-4}
\expandafter\def\csname sga@target@IV.17.13.5\endcsname{ega4.c04.statement.17-13-5}
\expandafter\def\csname sga@target@IV.17.13.5.1\endcsname{ega4.c04.statement.17-13-5-1}
\expandafter\def\csname sga@target@IV.17.13.5.2\endcsname{ega4.c04.eq.17-13-5-2}
\expandafter\def\csname sga@target@IV.17.13.5.3\endcsname{ega4.c04.eq.17-13-5-3}
\expandafter\def\csname sga@target@IV.17.13.5.4\endcsname{ega4.c04.eq.17-13-5-4}
\expandafter\def\csname sga@target@IV.17.13.6\endcsname{ega4.c04.cor.17-13-6}
\expandafter\def\csname sga@target@IV.17.13.6.1\endcsname{ega4.c04.eq.17-13-6-1}
\expandafter\def\csname sga@target@IV.17.13.6.2\endcsname{ega4.c04.eq.17-13-6-2}
\expandafter\def\csname sga@target@IV.17.13.7\endcsname{ega4.c04.statement.17-13-7}
\expandafter\def\csname sga@target@IV.17.13.7.1\endcsname{ega4.c04.eq.17-13-7-1}
\expandafter\def\csname sga@target@IV.17.13.8\endcsname{ega4.c04.prop.17-13-8}
\expandafter\def\csname sga@target@IV.17.13.8.1\endcsname{ega4.c04.eq.17-13-8-1}
\expandafter\def\csname sga@target@IV.17.13.9\endcsname{ega4.c04.statement.17-13-9}
\expandafter\def\csname sga@target@IV.17.13.9.1\endcsname{ega4.c04.eq.17-13-9-1}
\expandafter\def\csname sga@target@IV.17.13.9.2\endcsname{ega4.c04.eq.17-13-9-2}
\expandafter\def\csname sga@target@IV.17.14.1\endcsname{ega4.c04.prop.17-14-1}
\expandafter\def\csname sga@target@IV.17.14.2\endcsname{ega4.c04.prop.17-14-2}
\expandafter\def\csname sga@target@IV.17.14.3\endcsname{ega4.c04.rem.17-14-3}
\expandafter\def\csname sga@target@IV.17.15.1\endcsname{ega4.c04.prop.17-15-1}
\expandafter\def\csname sga@target@IV.17.15.10\endcsname{ega4.c04.statement.17-15-10}
\expandafter\def\csname sga@target@IV.17.15.11\endcsname{ega4.c04.prop.17-15-11}
\expandafter\def\csname sga@target@IV.17.15.12\endcsname{ega4.c04.prop.17-15-12}
\expandafter\def\csname sga@target@IV.17.15.13\endcsname{ega4.c04.cor.17-15-13}
\expandafter\def\csname sga@target@IV.17.15.14\endcsname{ega4.c04.prop.17-15-14}
\expandafter\def\csname sga@target@IV.17.15.14.1\endcsname{ega4.c04.lem.17-15-14-1}
\expandafter\def\csname sga@target@IV.17.15.14.2\endcsname{ega4.c04.lem.17-15-14-2}
\expandafter\def\csname sga@target@IV.17.15.15\endcsname{ega4.c04.prop.17-15-15}
\expandafter\def\csname sga@target@IV.17.15.2\endcsname{ega4.c04.cor.17-15-2}
\expandafter\def\csname sga@target@IV.17.15.3\endcsname{ega4.c04.prop.17-15-3}
\expandafter\def\csname sga@target@IV.17.15.4\endcsname{ega4.c04.cor.17-15-4}
\expandafter\def\csname sga@target@IV.17.15.5\endcsname{ega4.c04.prop.17-15-5}
\expandafter\def\csname sga@target@IV.17.15.6\endcsname{ega4.c04.cor.17-15-6}
\expandafter\def\csname sga@target@IV.17.15.7\endcsname{ega4.c04.cor.17-15-7}
\expandafter\def\csname sga@target@IV.17.15.8\endcsname{ega4.c04.prop.17-15-8}
\expandafter\def\csname sga@target@IV.17.15.8.1\endcsname{ega4.c04.eq.17-15-8-1}
\expandafter\def\csname sga@target@IV.17.15.8.2\endcsname{ega4.c04.eq.17-15-8-2}
\expandafter\def\csname sga@target@IV.17.15.9\endcsname{ega4.c04.cor.17-15-9}
\expandafter\def\csname sga@target@IV.17.16\endcsname{ega4.c04.subsec.17-16.key-iv-17-16}
\expandafter\def\csname sga@target@IV.17.16.1\endcsname{ega4.c04.prop.17-16-1}
\expandafter\def\csname sga@target@IV.17.16.2\endcsname{ega4.c04.cor.17-16-2}
\expandafter\def\csname sga@target@IV.17.16.3\endcsname{ega4.c04.cor.17-16-3}
\expandafter\def\csname sga@target@IV.17.16.4\endcsname{ega4.c04.prop.17-16-4}
\expandafter\def\csname sga@target@IV.17.16.5\endcsname{ega4.c04.cor.17-16-5}
\expandafter\def\csname sga@target@IV.17.16.6\endcsname{ega4.c04.prop.17-16-6}
\expandafter\def\csname sga@target@IV.17.2.1\endcsname{ega4.c04.prop.17-2-1}
\expandafter\def\csname sga@target@IV.17.2.2\endcsname{ega4.c04.cor.17-2-2}
\expandafter\def\csname sga@target@IV.17.2.3\endcsname{ega4.c04.prop.17-2-3}
\expandafter\def\csname sga@target@IV.17.2.3.1\endcsname{ega4.c04.eq.17-2-3-1}
\expandafter\def\csname sga@target@IV.17.2.4\endcsname{ega4.c04.cor.17-2-4}
\expandafter\def\csname sga@target@IV.17.2.5\endcsname{ega4.c04.prop.17-2-5}
\expandafter\def\csname sga@target@IV.17.2.5.1\endcsname{ega4.c04.prop.17-2-5-1}
\expandafter\def\csname sga@target@IV.17.2.6\endcsname{ega4.c04.prop.17-2-6}
\expandafter\def\csname sga@target@IV.17.3.1\endcsname{ega4.c04.def.17-3-1}
\expandafter\def\csname sga@target@IV.17.3.2\endcsname{ega4.c04.subsec.17-3-2}
\expandafter\def\csname sga@target@IV.17.3.3\endcsname{ega4.c04.prop.17-3-3}
\expandafter\def\csname sga@target@IV.17.3.4\endcsname{ega4.c04.prop.17-3-4}
\expandafter\def\csname sga@target@IV.17.3.5\endcsname{ega4.c04.cor.17-3-5}
\expandafter\def\csname sga@target@IV.17.3.6\endcsname{ega4.c04.prop.17-3-6}
\expandafter\def\csname sga@target@IV.17.3.7\endcsname{ega4.c04.def.17-3-7}
\expandafter\def\csname sga@target@IV.17.3.8\endcsname{ega4.c04.prop.17-3-8}
\expandafter\def\csname sga@target@IV.17.3.9\endcsname{ega4.c04.cor.17-3-9}
\expandafter\def\csname sga@target@IV.17.4.1\endcsname{ega4.c04.thm.17-4-1}
\expandafter\def\csname sga@target@IV.17.4.1.1\endcsname{ega4.c04.eq.17-4-1-1}
\expandafter\def\csname sga@target@IV.17.4.1.2\endcsname{ega4.c04.rem.17-4-1-2}
\expandafter\def\csname sga@target@IV.17.4.10\endcsname{ega4.c04.rem.17-4-10}
\expandafter\def\csname sga@target@IV.17.4.2\endcsname{ega4.c04.cor.17-4-2}
\expandafter\def\csname sga@target@IV.17.4.3\endcsname{ega4.c04.cor.17-4-3}
\expandafter\def\csname sga@target@IV.17.4.4\endcsname{ega4.c04.prop.17-4-4}
\expandafter\def\csname sga@target@IV.17.4.5\endcsname{ega4.c04.statement.17-4-5}
\expandafter\def\csname sga@target@IV.17.4.6\endcsname{ega4.c04.prop.17-4-6}
\expandafter\def\csname sga@target@IV.17.4.7\endcsname{ega4.c04.cor.17-4-7}
\expandafter\def\csname sga@target@IV.17.4.8\endcsname{ega4.c04.cor.17-4-8}
\expandafter\def\csname sga@target@IV.17.4.9\endcsname{ega4.c04.prop.17-4-9}
\expandafter\def\csname sga@target@IV.17.5.1\endcsname{ega4.c04.thm.17-5-1}
\expandafter\def\csname sga@target@IV.17.5.1.1\endcsname{ega4.c04.eq.17-5-1-1}
\expandafter\def\csname sga@target@IV.17.5.1.2\endcsname{ega4.c04.eq.17-5-1-2}
\expandafter\def\csname sga@target@IV.17.5.2\endcsname{ega4.c04.cor.17-5-2}
\expandafter\def\csname sga@target@IV.17.5.3\endcsname{ega4.c04.prop.17-5-3}
\expandafter\def\csname sga@target@IV.17.5.4\endcsname{ega4.c04.rem.17-5-4}
\expandafter\def\csname sga@target@IV.17.5.5\endcsname{ega4.c04.prop.17-5-5}
\expandafter\def\csname sga@target@IV.17.5.6\endcsname{ega4.c04.cor.17-5-6}
\expandafter\def\csname sga@target@IV.17.5.7\endcsname{ega4.c04.prop.17-5-7}
\expandafter\def\csname sga@target@IV.17.5.8\endcsname{ega4.c04.prop.17-5-8}
\expandafter\def\csname sga@target@IV.17.6.1\endcsname{ega4.c04.thm.17-6-1}
\expandafter\def\csname sga@target@IV.17.6.2\endcsname{ega4.c04.cor.17-6-2}
\expandafter\def\csname sga@target@IV.17.6.3\endcsname{ega4.c04.prop.17-6-3}
\expandafter\def\csname sga@target@IV.17.6.4\endcsname{ega4.c04.prop.17-6-4}
\expandafter\def\csname sga@target@IV.17.7\endcsname{ega4.c04.subsec.17-7.key-iv-17-7}
\expandafter\def\csname sga@target@IV.17.7.1\endcsname{ega4.c04.prop.17-7-1}
\expandafter\def\csname sga@target@IV.17.7.10\endcsname{ega4.c04.prop.17-7-10}
\expandafter\def\csname sga@target@IV.17.7.11\endcsname{ega4.c04.prop.17-7-11}
\expandafter\def\csname sga@target@IV.17.7.11.1\endcsname{ega4.c04.lem.17-7-11-1}
\expandafter\def\csname sga@target@IV.17.7.2\endcsname{ega4.c04.cor.17-7-2}
\expandafter\def\csname sga@target@IV.17.7.3\endcsname{ega4.c04.cor.17-7-3}
\expandafter\def\csname sga@target@IV.17.7.4\endcsname{ega4.c04.prop.17-7-4}
\expandafter\def\csname sga@target@IV.17.7.5\endcsname{ega4.c04.lem.17-7-5}
\expandafter\def\csname sga@target@IV.17.7.6\endcsname{ega4.c04.rem.17-7-6}
\expandafter\def\csname sga@target@IV.17.7.7\endcsname{ega4.c04.prop.17-7-7}
\expandafter\def\csname sga@target@IV.17.7.8\endcsname{ega4.c04.prop.17-7-8}
\expandafter\def\csname sga@target@IV.17.7.9\endcsname{ega4.c04.cor.17-7-9}
\expandafter\def\csname sga@target@IV.17.8\endcsname{ega4.c04.subsec.17-8.key-iv-17-8}
\expandafter\def\csname sga@target@IV.17.8.1\endcsname{ega4.c04.prop.17-8-1}
\expandafter\def\csname sga@target@IV.17.8.2\endcsname{ega4.c04.prop.17-8-2}
\expandafter\def\csname sga@target@IV.17.8.3\endcsname{ega4.c04.rem.17-8-3}
\expandafter\def\csname sga@target@IV.17.9\endcsname{ega4.c04.subsec.17-9.key-iv-17-9}
\expandafter\def\csname sga@target@IV.17.9.1\endcsname{ega4.c04.thm.17-9-1}
\expandafter\def\csname sga@target@IV.17.9.2\endcsname{ega4.c04.cor.17-9-2}
\expandafter\def\csname sga@target@IV.17.9.3\endcsname{ega4.c04.cor.17-9-3}
\expandafter\def\csname sga@target@IV.17.9.4\endcsname{ega4.c04.cor.17-9-4}
\expandafter\def\csname sga@target@IV.17.9.5\endcsname{ega4.c04.cor.17-9-5}
\expandafter\def\csname sga@target@IV.17.9.6\endcsname{ega4.c04.prop.17-9-6}
\expandafter\def\csname sga@target@IV.18\endcsname{ega4.c04.sec.18.key-iv-18}
\expandafter\def\csname sga@target@IV.18.1.1\endcsname{ega4.c04.prop.18-1-1}
\expandafter\def\csname sga@target@IV.18.1.1.1\endcsname{ega4.c04.eq.18-1-1-1}
\expandafter\def\csname sga@target@IV.18.1.2\endcsname{ega4.c04.thm.18-1-2}
\expandafter\def\csname sga@target@IV.18.1.3\endcsname{ega4.c04.cor.18-1-3}
\expandafter\def\csname sga@target@IV.18.10.1\endcsname{ega4.c04.thm.18-10-1}
\expandafter\def\csname sga@target@IV.18.10.10\endcsname{ega4.c04.statement.18-10-10}
\expandafter\def\csname sga@target@IV.18.10.11\endcsname{ega4.c04.rem.18-10-11}
\expandafter\def\csname sga@target@IV.18.10.12\endcsname{ega4.c04.cor.18-10-12}
\expandafter\def\csname sga@target@IV.18.10.13\endcsname{ega4.c04.prop.18-10-13}
\expandafter\def\csname sga@target@IV.18.10.14\endcsname{ega4.c04.cor.18-10-14}
\expandafter\def\csname sga@target@IV.18.10.15\endcsname{ega4.c04.statement.18-10-15}
\expandafter\def\csname sga@target@IV.18.10.16\endcsname{ega4.c04.thm.18-10-16}
\expandafter\def\csname sga@target@IV.18.10.17\endcsname{ega4.c04.rem.18-10-17}
\expandafter\def\csname sga@target@IV.18.10.17.1\endcsname{ega4.c04.statement.18-10-17-1}
\expandafter\def\csname sga@target@IV.18.10.17.2\endcsname{ega4.c04.statement.18-10-17-2}
\expandafter\def\csname sga@target@IV.18.10.18\endcsname{ega4.c04.lem.18-10-18}
\expandafter\def\csname sga@target@IV.18.10.19\endcsname{ega4.c04.prop.18-10-19}
\expandafter\def\csname sga@target@IV.18.10.2\endcsname{ega4.c04.rem.18-10-2}
\expandafter\def\csname sga@target@IV.18.10.20\endcsname{ega4.c04.rem.18-10-20}
\expandafter\def\csname sga@target@IV.18.10.21\endcsname{ega4.c04.prop.18-10-21}
\expandafter\def\csname sga@target@IV.18.10.3\endcsname{ega4.c04.cor.18-10-3}
\expandafter\def\csname sga@target@IV.18.10.3.1\endcsname{ega4.c04.rem.18-10-3-1}
\expandafter\def\csname sga@target@IV.18.10.4\endcsname{ega4.c04.cor.18-10-4}
\expandafter\def\csname sga@target@IV.18.10.5\endcsname{ega4.c04.cor.18-10-5}
\expandafter\def\csname sga@target@IV.18.10.6\endcsname{ega4.c04.lem.18-10-6}
\expandafter\def\csname sga@target@IV.18.10.7\endcsname{ega4.c04.prop.18-10-7}
\expandafter\def\csname sga@target@IV.18.10.8\endcsname{ega4.c04.prop.18-10-8}
\expandafter\def\csname sga@target@IV.18.10.9\endcsname{ega4.c04.cor.18-10-9}
\expandafter\def\csname sga@target@IV.18.11.1\endcsname{ega4.c04.lem.18-11-1}
\expandafter\def\csname sga@target@IV.18.11.10\endcsname{ega4.c04.thm.18-11-10}
\expandafter\def\csname sga@target@IV.18.11.11\endcsname{ega4.c04.rem.18-11-11}
\expandafter\def\csname sga@target@IV.18.11.12\endcsname{ega4.c04.cor.18-11-12}
\expandafter\def\csname sga@target@IV.18.11.12.1\endcsname{ega4.c04.lem.18-11-12-1}
\expandafter\def\csname sga@target@IV.18.11.2\endcsname{ega4.c04.lem.18-11-2}
\expandafter\def\csname sga@target@IV.18.11.2.1\endcsname{ega4.c04.eq.18-11-2-1}
\expandafter\def\csname sga@target@IV.18.11.2.2\endcsname{ega4.c04.eq.18-11-2-2}
\expandafter\def\csname sga@target@IV.18.11.3\endcsname{ega4.c04.lem.18-11-3}
\expandafter\def\csname sga@target@IV.18.11.3.1\endcsname{ega4.c04.eq.18-11-3-1}
\expandafter\def\csname sga@target@IV.18.11.3.2\endcsname{ega4.c04.eq.18-11-3-2}
\expandafter\def\csname sga@target@IV.18.11.3.3\endcsname{ega4.c04.lem.18-11-3-3}
\expandafter\def\csname sga@target@IV.18.11.3.4\endcsname{ega4.c04.lem.18-11-3-4}
\expandafter\def\csname sga@target@IV.18.11.3.5\endcsname{ega4.c04.lem.18-11-3-5}
\expandafter\def\csname sga@target@IV.18.11.3.6\endcsname{ega4.c04.eq.18-11-3-6}
\expandafter\def\csname sga@target@IV.18.11.4\endcsname{ega4.c04.lem.18-11-4}
\expandafter\def\csname sga@target@IV.18.11.5\endcsname{ega4.c04.prop.18-11-5}
\expandafter\def\csname sga@target@IV.18.11.5.1\endcsname{ega4.c04.lem.18-11-5-1}
\expandafter\def\csname sga@target@IV.18.11.5.2\endcsname{ega4.c04.lem.18-11-5-2}
\expandafter\def\csname sga@target@IV.18.11.6\endcsname{ega4.c04.rem.18-11-6}
\expandafter\def\csname sga@target@IV.18.11.7\endcsname{ega4.c04.cor.18-11-7}
\expandafter\def\csname sga@target@IV.18.11.8\endcsname{ega4.c04.cor.18-11-8}
\expandafter\def\csname sga@target@IV.18.11.9\endcsname{ega4.c04.cor.18-11-9}
\expandafter\def\csname sga@target@IV.18.12.1\endcsname{ega4.c04.thm.18-12-1}
\expandafter\def\csname sga@target@IV.18.12.10\endcsname{ega4.c04.cor.18-12-10}
\expandafter\def\csname sga@target@IV.18.12.11\endcsname{ega4.c04.cor.18-12-11}
\expandafter\def\csname sga@target@IV.18.12.12\endcsname{ega4.c04.prop.18-12-12}
\expandafter\def\csname sga@target@IV.18.12.13\endcsname{ega4.c04.cor.18-12-13}
\expandafter\def\csname sga@target@IV.18.12.14\endcsname{ega4.c04.rem.18-12-14}
\expandafter\def\csname sga@target@IV.18.12.15\endcsname{ega4.c04.prop.18-12-15}
\expandafter\def\csname sga@target@IV.18.12.16\endcsname{ega4.c04.prop.18-12-16}
\expandafter\def\csname sga@target@IV.18.12.17\endcsname{ega4.c04.cor.18-12-17}
\expandafter\def\csname sga@target@IV.18.12.2\endcsname{ega4.c04.rem.18-12-2}
\expandafter\def\csname sga@target@IV.18.12.3\endcsname{ega4.c04.cor.18-12-3}
\expandafter\def\csname sga@target@IV.18.12.4\endcsname{ega4.c04.cor.18-12-4}
\expandafter\def\csname sga@target@IV.18.12.5\endcsname{ega4.c04.rem.18-12-5}
\expandafter\def\csname sga@target@IV.18.12.6\endcsname{ega4.c04.cor.18-12-6}
\expandafter\def\csname sga@target@IV.18.12.7\endcsname{ega4.c04.prop.18-12-7}
\expandafter\def\csname sga@target@IV.18.12.7.1\endcsname{ega4.c04.lem.18-12-7-1}
\expandafter\def\csname sga@target@IV.18.12.8\endcsname{ega4.c04.prop.18-12-8}
\expandafter\def\csname sga@target@IV.18.12.9\endcsname{ega4.c04.rem.18-12-9}
\expandafter\def\csname sga@target@IV.18.2.1\endcsname{ega4.c04.statement.18-2-1}
\expandafter\def\csname sga@target@IV.18.2.1.1\endcsname{ega4.c04.statement.18-2-1-1}
\expandafter\def\csname sga@target@IV.18.2.1.2\endcsname{ega4.c04.eq.18-2-1-2}
\expandafter\def\csname sga@target@IV.18.2.10\endcsname{ega4.c04.statement.18-2-10}
\expandafter\def\csname sga@target@IV.18.2.2\endcsname{ega4.c04.statement.18-2-2}
\expandafter\def\csname sga@target@IV.18.2.2.1\endcsname{ega4.c04.eq.18-2-2-1}
\expandafter\def\csname sga@target@IV.18.2.2.2\endcsname{ega4.c04.eq.18-2-2-2}
\expandafter\def\csname sga@target@IV.18.2.2.3\endcsname{ega4.c04.eq.18-2-2-3}
\expandafter\def\csname sga@target@IV.18.2.2.4\endcsname{ega4.c04.eq.18-2-2-4}
\expandafter\def\csname sga@target@IV.18.2.3\endcsname{ega4.c04.prop.18-2-3}
\expandafter\def\csname sga@target@IV.18.2.4\endcsname{ega4.c04.cor.18-2-4}
\expandafter\def\csname sga@target@IV.18.2.5\endcsname{ega4.c04.cor.18-2-5}
\expandafter\def\csname sga@target@IV.18.2.6\endcsname{ega4.c04.cor.18-2-6}
\expandafter\def\csname sga@target@IV.18.2.7\endcsname{ega4.c04.item.18-2-7}
\expandafter\def\csname sga@target@IV.18.2.7.1\endcsname{ega4.c04.eq.18-2-7-1}
\expandafter\def\csname sga@target@IV.18.2.7.2\endcsname{ega4.c04.eq.18-2-7-2}
\expandafter\def\csname sga@target@IV.18.2.8\endcsname{ega4.c04.prop.18-2-8}
\expandafter\def\csname sga@target@IV.18.2.9\endcsname{ega4.c04.cor.18-2-9}
\expandafter\def\csname sga@target@IV.18.3.1\endcsname{ega4.c04.prop.18-3-1}
\expandafter\def\csname sga@target@IV.18.3.2\endcsname{ega4.c04.prop.18-3-2}
\expandafter\def\csname sga@target@IV.18.3.2.1\endcsname{ega4.c04.lem.18-3-2-1}
\expandafter\def\csname sga@target@IV.18.3.2.2\endcsname{ega4.c04.eq.18-3-2-2}
\expandafter\def\csname sga@target@IV.18.3.3\endcsname{ega4.c04.statement.18-3-3}
\expandafter\def\csname sga@target@IV.18.3.4\endcsname{ega4.c04.thm.18-3-4}
\expandafter\def\csname sga@target@IV.18.4.1\endcsname{ega4.c04.lem.18-4-1}
\expandafter\def\csname sga@target@IV.18.4.10\endcsname{ega4.c04.prop.18-4-10}
\expandafter\def\csname sga@target@IV.18.4.11\endcsname{ega4.c04.statement.18-4-11}
\expandafter\def\csname sga@target@IV.18.4.12\endcsname{ega4.c04.cor.18-4-12}
\expandafter\def\csname sga@target@IV.18.4.13\endcsname{ega4.c04.cor.18-4-13}
\expandafter\def\csname sga@target@IV.18.4.14\endcsname{ega4.c04.cor.18-4-14}
\expandafter\def\csname sga@target@IV.18.4.2\endcsname{ega4.c04.prop.18-4-2}
\expandafter\def\csname sga@target@IV.18.4.3\endcsname{ega4.c04.cor.18-4-3}
\expandafter\def\csname sga@target@IV.18.4.4\endcsname{ega4.c04.lem.18-4-4}
\expandafter\def\csname sga@target@IV.18.4.5\endcsname{ega4.c04.prop.18-4-5}
\expandafter\def\csname sga@target@IV.18.4.6\endcsname{ega4.c04.thm.18-4-6}
\expandafter\def\csname sga@target@IV.18.4.6.1\endcsname{ega4.c04.lem.18-4-6-1}
\expandafter\def\csname sga@target@IV.18.4.7\endcsname{ega4.c04.cor.18-4-7}
\expandafter\def\csname sga@target@IV.18.4.8\endcsname{ega4.c04.cor.18-4-8}
\expandafter\def\csname sga@target@IV.18.4.9\endcsname{ega4.c04.prop.18-4-9}
\expandafter\def\csname sga@target@IV.18.5.1\endcsname{ega4.c04.statement.18-5-1}
\expandafter\def\csname sga@target@IV.18.5.1.1\endcsname{ega4.c04.eq.18-5-1-1}
\expandafter\def\csname sga@target@IV.18.5.10\endcsname{ega4.c04.prop.18-5-10}
\expandafter\def\csname sga@target@IV.18.5.11\endcsname{ega4.c04.thm.18-5-11}
\expandafter\def\csname sga@target@IV.18.5.12\endcsname{ega4.c04.cor.18-5-12}
\expandafter\def\csname sga@target@IV.18.5.13\endcsname{ega4.c04.rem.18-5-13}
\expandafter\def\csname sga@target@IV.18.5.13.1\endcsname{ega4.c04.lem.18-5-13-1}
\expandafter\def\csname sga@target@IV.18.5.14\endcsname{ega4.c04.prop.18-5-14}
\expandafter\def\csname sga@target@IV.18.5.15\endcsname{ega4.c04.prop.18-5-15}
\expandafter\def\csname sga@target@IV.18.5.16\endcsname{ega4.c04.item.18-5-16}
\expandafter\def\csname sga@target@IV.18.5.17\endcsname{ega4.c04.thm.18-5-17}
\expandafter\def\csname sga@target@IV.18.5.18\endcsname{ega4.c04.rem.18-5-18}
\expandafter\def\csname sga@target@IV.18.5.19\endcsname{ega4.c04.prop.18-5-19}
\expandafter\def\csname sga@target@IV.18.5.2\endcsname{ega4.c04.statement.18-5-2}
\expandafter\def\csname sga@target@IV.18.5.3\endcsname{ega4.c04.lem.18-5-3}
\expandafter\def\csname sga@target@IV.18.5.3.1\endcsname{ega4.c04.lem.18-5-3-1}
\expandafter\def\csname sga@target@IV.18.5.4\endcsname{ega4.c04.prop.18-5-4}
\expandafter\def\csname sga@target@IV.18.5.4.1\endcsname{ega4.c04.eq.18-5-4-1}
\expandafter\def\csname sga@target@IV.18.5.4.2\endcsname{ega4.c04.eq.18-5-4-2}
\expandafter\def\csname sga@target@IV.18.5.4.3\endcsname{ega4.c04.lem.18-5-4-3}
\expandafter\def\csname sga@target@IV.18.5.4.4\endcsname{ega4.c04.rem.18-5-4-4}
\expandafter\def\csname sga@target@IV.18.5.5\endcsname{ega4.c04.def.18-5-5}
\expandafter\def\csname sga@target@IV.18.5.6\endcsname{ega4.c04.prop.18-5-6}
\expandafter\def\csname sga@target@IV.18.5.7\endcsname{ega4.c04.rem.18-5-7}
\expandafter\def\csname sga@target@IV.18.5.8\endcsname{ega4.c04.def.18-5-8}
\expandafter\def\csname sga@target@IV.18.5.9\endcsname{ega4.c04.prop.18-5-9}
\expandafter\def\csname sga@target@IV.18.6.1\endcsname{ega4.c04.statement.18-6-1}
\expandafter\def\csname sga@target@IV.18.6.10\endcsname{ega4.c04.cor.18-6-10}
\expandafter\def\csname sga@target@IV.18.6.11\endcsname{ega4.c04.cor.18-6-11}
\expandafter\def\csname sga@target@IV.18.6.12\endcsname{ega4.c04.prop.18-6-12}
\expandafter\def\csname sga@target@IV.18.6.13\endcsname{ega4.c04.cor.18-6-13}
\expandafter\def\csname sga@target@IV.18.6.14\endcsname{ega4.c04.prop.18-6-14}
\expandafter\def\csname sga@target@IV.18.6.15\endcsname{ega4.c04.cor.18-6-15}
\expandafter\def\csname sga@target@IV.18.6.2\endcsname{ega4.c04.lem.18-6-2}
\expandafter\def\csname sga@target@IV.18.6.3\endcsname{ega4.c04.lem.18-6-3}
\expandafter\def\csname sga@target@IV.18.6.4\endcsname{ega4.c04.statement.18-6-4}
\expandafter\def\csname sga@target@IV.18.6.5\endcsname{ega4.c04.def.18-6-5}
\expandafter\def\csname sga@target@IV.18.6.6\endcsname{ega4.c04.thm.18-6-6}
\expandafter\def\csname sga@target@IV.18.6.7\endcsname{ega4.c04.statement.18-6-7}
\expandafter\def\csname sga@target@IV.18.6.8\endcsname{ega4.c04.prop.18-6-8}
\expandafter\def\csname sga@target@IV.18.6.9\endcsname{ega4.c04.thm.18-6-9}
\expandafter\def\csname sga@target@IV.18.6.9.1\endcsname{ega4.c04.lem.18-6-9-1}
\expandafter\def\csname sga@target@IV.18.6.9.2\endcsname{ega4.c04.lem.18-6-9-2}
\expandafter\def\csname sga@target@IV.18.7.1\endcsname{ega4.c04.statement.18-7-1}
\expandafter\def\csname sga@target@IV.18.7.2\endcsname{ega4.c04.prop.18-7-2}
\expandafter\def\csname sga@target@IV.18.7.3\endcsname{ega4.c04.cor.18-7-3}
\expandafter\def\csname sga@target@IV.18.7.4\endcsname{ega4.c04.cor.18-7-4}
\expandafter\def\csname sga@target@IV.18.7.5\endcsname{ega4.c04.prop.18-7-5}
\expandafter\def\csname sga@target@IV.18.7.5.1\endcsname{ega4.c04.lem.18-7-5-1}
\expandafter\def\csname sga@target@IV.18.7.6\endcsname{ega4.c04.cor.18-7-6}
\expandafter\def\csname sga@target@IV.18.7.7\endcsname{ega4.c04.rem.18-7-7}
\expandafter\def\csname sga@target@IV.18.8.1\endcsname{ega4.c04.prop.18-8-1}
\expandafter\def\csname sga@target@IV.18.8.10\endcsname{ega4.c04.prop.18-8-10}
\expandafter\def\csname sga@target@IV.18.8.11\endcsname{ega4.c04.rem.18-8-11}
\expandafter\def\csname sga@target@IV.18.8.12\endcsname{ega4.c04.prop.18-8-12}
\expandafter\def\csname sga@target@IV.18.8.13\endcsname{ega4.c04.cor.18-8-13}
\expandafter\def\csname sga@target@IV.18.8.14\endcsname{ega4.c04.cor.18-8-14}
\expandafter\def\csname sga@target@IV.18.8.15\endcsname{ega4.c04.prop.18-8-15}
\expandafter\def\csname sga@target@IV.18.8.16\endcsname{ega4.c04.cor.18-8-16}
\expandafter\def\csname sga@target@IV.18.8.17\endcsname{ega4.c04.prop.18-8-17}
\expandafter\def\csname sga@target@IV.18.8.18\endcsname{ega4.c04.prop.18-8-18}
\expandafter\def\csname sga@target@IV.18.8.2\endcsname{ega4.c04.def.18-8-2}
\expandafter\def\csname sga@target@IV.18.8.3\endcsname{ega4.c04.rem.18-8-3}
\expandafter\def\csname sga@target@IV.18.8.4\endcsname{ega4.c04.lem.18-8-4}
\expandafter\def\csname sga@target@IV.18.8.5\endcsname{ega4.c04.lem.18-8-5}
\expandafter\def\csname sga@target@IV.18.8.6\endcsname{ega4.c04.statement.18-8-6}
\expandafter\def\csname sga@target@IV.18.8.6.1\endcsname{ega4.c04.eq.18-8-6-1}
\expandafter\def\csname sga@target@IV.18.8.6.2\endcsname{ega4.c04.eq.18-8-6-2}
\expandafter\def\csname sga@target@IV.18.8.7\endcsname{ega4.c04.def.18-8-7}
\expandafter\def\csname sga@target@IV.18.8.8\endcsname{ega4.c04.prop.18-8-8}
\expandafter\def\csname sga@target@IV.18.8.8.1\endcsname{ega4.c04.statement.18-8-8-1}
\expandafter\def\csname sga@target@IV.18.8.9\endcsname{ega4.c04.statement.18-8-9}
\expandafter\def\csname sga@target@IV.18.9.1\endcsname{ega4.c04.thm.18-9-1}
\expandafter\def\csname sga@target@IV.18.9.10\endcsname{ega4.c04.rem.18-9-10}
\expandafter\def\csname sga@target@IV.18.9.11\endcsname{ega4.c04.cor.18-9-11}
\expandafter\def\csname sga@target@IV.18.9.2\endcsname{ega4.c04.cor.18-9-2}
\expandafter\def\csname sga@target@IV.18.9.3\endcsname{ega4.c04.cor.18-9-3}
\expandafter\def\csname sga@target@IV.18.9.4\endcsname{ega4.c04.cor.18-9-4}
\expandafter\def\csname sga@target@IV.18.9.5\endcsname{ega4.c04.cor.18-9-5}
\expandafter\def\csname sga@target@IV.18.9.5.1\endcsname{ega4.c04.eq.18-9-5-1}
\expandafter\def\csname sga@target@IV.18.9.6\endcsname{ega4.c04.statement.18-9-6}
\expandafter\def\csname sga@target@IV.18.9.7\endcsname{ega4.c04.thm.18-9-7}
\expandafter\def\csname sga@target@IV.18.9.7.1\endcsname{ega4.c04.rem.18-9-7-1}
\expandafter\def\csname sga@target@IV.18.9.7.2\endcsname{ega4.c04.statement.18-9-7-2}
\expandafter\def\csname sga@target@IV.18.9.7.3\endcsname{ega4.c04.lem.18-9-7-3}
\expandafter\def\csname sga@target@IV.18.9.7.4\endcsname{ega4.c04.eq.18-9-7-4}
\expandafter\def\csname sga@target@IV.18.9.7.5\endcsname{ega4.c04.lem.18-9-7-5}
\expandafter\def\csname sga@target@IV.18.9.7.6\endcsname{ega4.c04.rem.18-9-7-6}
\expandafter\def\csname sga@target@IV.18.9.7.7\endcsname{ega4.c04.lem.18-9-7-7}
\expandafter\def\csname sga@target@IV.18.9.7.8\endcsname{ega4.c04.rem.18-9-7-8}
\expandafter\def\csname sga@target@IV.18.9.8\endcsname{ega4.c04.cor.18-9-8}
\expandafter\def\csname sga@target@IV.18.9.9\endcsname{ega4.c04.cor.18-9-9}
\expandafter\def\csname sga@target@IV.19.1.1\endcsname{ega4.c04.prop.19-1-1}
\expandafter\def\csname sga@target@IV.19.1.2\endcsname{ega4.c04.cor.19-1-2}
\expandafter\def\csname sga@target@IV.19.1.3\endcsname{ega4.c04.def.19-1-3}
\expandafter\def\csname sga@target@IV.19.1.4\endcsname{ega4.c04.prop.19-1-4}
\expandafter\def\csname sga@target@IV.19.1.4.1\endcsname{ega4.c04.eq.19-1-4-1}
\expandafter\def\csname sga@target@IV.19.1.5\endcsname{ega4.c04.prop.19-1-5}
\expandafter\def\csname sga@target@IV.19.1.5.1\endcsname{ega4.c04.item.19-1-5-1}
\expandafter\def\csname sga@target@IV.19.1.6\endcsname{ega4.c04.cor.19-1-6}
\expandafter\def\csname sga@target@IV.19.1.6.1\endcsname{ega4.c04.lem.19-1-6-1}
\expandafter\def\csname sga@target@IV.19.1.7\endcsname{ega4.c04.item.19-1-7}
\expandafter\def\csname sga@target@IV.19.1.8\endcsname{ega4.c04.statement.19-1-8}
\expandafter\def\csname sga@target@IV.19.2.1\endcsname{ega4.c04.def.19-2-1}
\expandafter\def\csname sga@target@IV.19.2.2\endcsname{ega4.c04.def.19-2-2}
\expandafter\def\csname sga@target@IV.19.2.3\endcsname{ega4.c04.prop.19-2-3}
\expandafter\def\csname sga@target@IV.19.2.4\endcsname{ega4.c04.prop.19-2-4}
\expandafter\def\csname sga@target@IV.19.2.5\endcsname{ega4.c04.item.19-2-5}
\expandafter\def\csname sga@target@IV.19.2.6\endcsname{ega4.c04.cor.19-2-6}
\expandafter\def\csname sga@target@IV.19.2.7\endcsname{ega4.c04.prop.19-2-7}
\expandafter\def\csname sga@target@IV.19.2.8\endcsname{ega4.c04.prop.19-2-8}
\expandafter\def\csname sga@target@IV.19.2.9\endcsname{ega4.c04.prop.19-2-9}
\expandafter\def\csname sga@target@IV.19.3.1\endcsname{ega4.c04.def.19-3-1}
\expandafter\def\csname sga@target@IV.19.3.10\endcsname{ega4.c04.cor.19-3-10}
\expandafter\def\csname sga@target@IV.19.3.2\endcsname{ega4.c04.prop.19-3-2}
\expandafter\def\csname sga@target@IV.19.3.2.1\endcsname{ega4.c04.lem.19-3-2-1}
\expandafter\def\csname sga@target@IV.19.3.3\endcsname{ega4.c04.cor.19-3-3}
\expandafter\def\csname sga@target@IV.19.3.4\endcsname{ega4.c04.cor.19-3-4}
\expandafter\def\csname sga@target@IV.19.3.5\endcsname{ega4.c04.cor.19-3-5}
\expandafter\def\csname sga@target@IV.19.3.6\endcsname{ega4.c04.def.19-3-6}
\expandafter\def\csname sga@target@IV.19.3.7\endcsname{ega4.c04.prop.19-3-7}
\expandafter\def\csname sga@target@IV.19.3.8\endcsname{ega4.c04.cor.19-3-8}
\expandafter\def\csname sga@target@IV.19.3.9\endcsname{ega4.c04.prop.19-3-9}
\expandafter\def\csname sga@target@IV.19.4.1\endcsname{ega4.c04.statement.19-4-1}
\expandafter\def\csname sga@target@IV.19.4.1.1\endcsname{ega4.c04.subsec.19-4-1-1}
\expandafter\def\csname sga@target@IV.19.4.10\endcsname{ega4.c04.rem.19-4-10}
\expandafter\def\csname sga@target@IV.19.4.11\endcsname{ega4.c04.prop.19-4-11}
\expandafter\def\csname sga@target@IV.19.4.11.1\endcsname{ega4.c04.eq.19-4-11-1}
\expandafter\def\csname sga@target@IV.19.4.2\endcsname{ega4.c04.lem.19-4-2}
\expandafter\def\csname sga@target@IV.19.4.3\endcsname{ega4.c04.cor.19-4-3}
\expandafter\def\csname sga@target@IV.19.4.4\endcsname{ega4.c04.prop.19-4-4}
\expandafter\def\csname sga@target@IV.19.4.5\endcsname{ega4.c04.cor.19-4-5}
\expandafter\def\csname sga@target@IV.19.4.6\endcsname{ega4.c04.prop.19-4-6}
\expandafter\def\csname sga@target@IV.19.4.7\endcsname{ega4.c04.cor.19-4-7}
\expandafter\def\csname sga@target@IV.19.4.8\endcsname{ega4.c04.prop.19-4-8}
\expandafter\def\csname sga@target@IV.19.4.9\endcsname{ega4.c04.cor.19-4-9}
\expandafter\def\csname sga@target@IV.19.5.1\endcsname{ega4.c04.thm.19-5-1}
\expandafter\def\csname sga@target@IV.19.5.1.1\endcsname{ega4.c04.eq.19-5-1-1}
\expandafter\def\csname sga@target@IV.19.5.2\endcsname{ega4.c04.cor.19-5-2}
\expandafter\def\csname sga@target@IV.19.5.3\endcsname{ega4.c04.prop.19-5-3}
\expandafter\def\csname sga@target@IV.19.5.4\endcsname{ega4.c04.statement.19-5-4}
\expandafter\def\csname sga@target@IV.19.5.5\endcsname{ega4.c04.thm.19-5-5}
\expandafter\def\csname sga@target@IV.19.5.5.1\endcsname{ega4.c04.statement.19-5-5-1}
\expandafter\def\csname sga@target@IV.19.5.5.2\endcsname{ega4.c04.statement.19-5-5-2}
\expandafter\def\csname sga@target@IV.19.5.5.3\endcsname{ega4.c04.statement.19-5-5-3}
\expandafter\def\csname sga@target@IV.19.5.5.4\endcsname{ega4.c04.statement.19-5-5-4}
\expandafter\def\csname sga@target@IV.19.5.5.5\endcsname{ega4.c04.statement.19-5-5-5}
\expandafter\def\csname sga@target@IV.19.5.6\endcsname{ega4.c04.statement.19-5-6}
\expandafter\def\csname sga@target@IV.19.6.1\endcsname{ega4.c04.statement.19-6-1}
\expandafter\def\csname sga@target@IV.19.6.1.1\endcsname{ega4.c04.eq.19-6-1-1}
\expandafter\def\csname sga@target@IV.19.6.2\endcsname{ega4.c04.statement.19-6-2}
\expandafter\def\csname sga@target@IV.19.6.2.1\endcsname{ega4.c04.statement.19-6-2-1}
\expandafter\def\csname sga@target@IV.19.6.3\endcsname{ega4.c04.prop.19-6-3}
\expandafter\def\csname sga@target@IV.19.6.4\endcsname{ega4.c04.cor.19-6-4}
\expandafter\def\csname sga@target@IV.19.7\endcsname{ega4.c04.subsec.19-7.key-iv-19-7}
\expandafter\def\csname sga@target@IV.19.7.1\endcsname{ega4.c04.thm.19-7-1}
\expandafter\def\csname sga@target@IV.19.7.1.1\endcsname{ega4.c04.statement.19-7-1-1}
\expandafter\def\csname sga@target@IV.19.7.1.10\endcsname{ega4.c04.statement.19-7-1-10}
\expandafter\def\csname sga@target@IV.19.7.1.11\endcsname{ega4.c04.statement.19-7-1-11}
\expandafter\def\csname sga@target@IV.19.7.1.12\endcsname{ega4.c04.statement.19-7-1-12}
\expandafter\def\csname sga@target@IV.19.7.1.13\endcsname{ega4.c04.statement.19-7-1-13}
\expandafter\def\csname sga@target@IV.19.7.1.14\endcsname{ega4.c04.statement.19-7-1-14}
\expandafter\def\csname sga@target@IV.19.7.1.15\endcsname{ega4.c04.statement.19-7-1-15}
\expandafter\def\csname sga@target@IV.19.7.1.16\endcsname{ega4.c04.statement.19-7-1-16}
\expandafter\def\csname sga@target@IV.19.7.1.17\endcsname{ega4.c04.statement.19-7-1-17}
\expandafter\def\csname sga@target@IV.19.7.1.18\endcsname{ega4.c04.statement.19-7-1-18}
\expandafter\def\csname sga@target@IV.19.7.1.19\endcsname{ega4.c04.statement.19-7-1-19}
\expandafter\def\csname sga@target@IV.19.7.1.2\endcsname{ega4.c04.statement.19-7-1-2}
\expandafter\def\csname sga@target@IV.19.7.1.20\endcsname{ega4.c04.statement.19-7-1-20}
\expandafter\def\csname sga@target@IV.19.7.1.21\endcsname{ega4.c04.statement.19-7-1-21}
\expandafter\def\csname sga@target@IV.19.7.1.22\endcsname{ega4.c04.statement.19-7-1-22}
\expandafter\def\csname sga@target@IV.19.7.1.23\endcsname{ega4.c04.statement.19-7-1-23}
\expandafter\def\csname sga@target@IV.19.7.1.3\endcsname{ega4.c04.statement.19-7-1-3}
\expandafter\def\csname sga@target@IV.19.7.1.5\endcsname{ega4.c04.statement.19-7-1-5}
\expandafter\def\csname sga@target@IV.19.7.1.6\endcsname{ega4.c04.statement.19-7-1-6}
\expandafter\def\csname sga@target@IV.19.7.1.7\endcsname{ega4.c04.lem.19-7-1-7}
\expandafter\def\csname sga@target@IV.19.7.1.8\endcsname{ega4.c04.statement.19-7-1-8}
\expandafter\def\csname sga@target@IV.19.7.1.9\endcsname{ega4.c04.statement.19-7-1-9}
\expandafter\def\csname sga@target@IV.19.7.2\endcsname{ega4.c04.statement.19-7-2}
\expandafter\def\csname sga@target@IV.19.7.3\endcsname{ega4.c04.cor.19-7-3}
\expandafter\def\csname sga@target@IV.19.7.3.1\endcsname{ega4.c04.eq.19-7-3-1}
\expandafter\def\csname sga@target@IV.19.7.4\endcsname{ega4.c04.cor.19-7-4}
\expandafter\def\csname sga@target@IV.19.7.4.1\endcsname{ega4.c04.eq.19-7-4-1}
\expandafter\def\csname sga@target@IV.19.7.5\endcsname{ega4.c04.cor.19-7-5}
\expandafter\def\csname sga@target@IV.19.8.1\endcsname{ega4.c04.prop.19-8-1}
\expandafter\def\csname sga@target@IV.19.8.2\endcsname{ega4.c04.prop.19-8-2}
\expandafter\def\csname sga@target@IV.19.8.3\endcsname{ega4.c04.rem.19-8-3}
\expandafter\def\csname sga@target@IV.19.9.1\endcsname{ega4.c04.statement.19-9-1}
\expandafter\def\csname sga@target@IV.19.9.1.1\endcsname{ega4.c04.eq.19-9-1-1}
\expandafter\def\csname sga@target@IV.19.9.1.2\endcsname{ega4.c04.eq.19-9-1-2}
\expandafter\def\csname sga@target@IV.19.9.1.3\endcsname{ega4.c04.eq.19-9-1-3}
\expandafter\def\csname sga@target@IV.19.9.2\endcsname{ega4.c04.lem.19-9-2}
\expandafter\def\csname sga@target@IV.19.9.3\endcsname{ega4.c04.prop.19-9-3}
\expandafter\def\csname sga@target@IV.19.9.4\endcsname{ega4.c04.cor.19-9-4}
\expandafter\def\csname sga@target@IV.19.9.5\endcsname{ega4.c04.prop.19-9-5}
\expandafter\def\csname sga@target@IV.19.9.6\endcsname{ega4.c04.prop.19-9-6}
\expandafter\def\csname sga@target@IV.19.9.7\endcsname{ega4.c04.cor.19-9-7}
\expandafter\def\csname sga@target@IV.19.9.8\endcsname{ega4.c04.prop.19-9-8}
\expandafter\def\csname sga@target@IV.19.9.9\endcsname{ega4.c04.rem.19-9-9}
\expandafter\def\csname sga@target@IV.2.1.1\endcsname{ega4.c04.statement.2-1-1}
\expandafter\def\csname sga@target@IV.2.1.10\endcsname{ega4.c04.cor.2-1-10}
\expandafter\def\csname sga@target@IV.2.1.11\endcsname{ega4.c04.prop.2-1-11}
\expandafter\def\csname sga@target@IV.2.1.12\endcsname{ega4.c04.prop.2-1-12}
\expandafter\def\csname sga@target@IV.2.1.13\endcsname{ega4.c04.prop.2-1-13}
\expandafter\def\csname sga@target@IV.2.1.14\endcsname{ega4.c04.prop.2-1-14}
\expandafter\def\csname sga@target@IV.2.1.2\endcsname{ega4.c04.prop.2-1-2}
\expandafter\def\csname sga@target@IV.2.1.3\endcsname{ega4.c04.prop.2-1-3}
\expandafter\def\csname sga@target@IV.2.1.4\endcsname{ega4.c04.prop.2-1-4}
\expandafter\def\csname sga@target@IV.2.1.5\endcsname{ega4.c04.prop.2-1-5}
\expandafter\def\csname sga@target@IV.2.1.6\endcsname{ega4.c04.cor.2-1-6}
\expandafter\def\csname sga@target@IV.2.1.7\endcsname{ega4.c04.cor.2-1-7}
\expandafter\def\csname sga@target@IV.2.1.8\endcsname{ega4.c04.prop.2-1-8}
\expandafter\def\csname sga@target@IV.2.1.9\endcsname{ega4.c04.cor.2-1-9}
\expandafter\def\csname sga@target@IV.2.2.1\endcsname{ega4.c04.prop.2-2-1}
\expandafter\def\csname sga@target@IV.2.2.10\endcsname{ega4.c04.prop.2-2-10}
\expandafter\def\csname sga@target@IV.2.2.10.1\endcsname{ega4.c04.eq.2-2-10-1}
\expandafter\def\csname sga@target@IV.2.2.11\endcsname{ega4.c04.cor.2-2-11}
\expandafter\def\csname sga@target@IV.2.2.12\endcsname{ega4.c04.cor.2-2-12}
\expandafter\def\csname sga@target@IV.2.2.13\endcsname{ega4.c04.cor.2-2-13}
\expandafter\def\csname sga@target@IV.2.2.14\endcsname{ega4.c04.prop.2-2-14}
\expandafter\def\csname sga@target@IV.2.2.15\endcsname{ega4.c04.prop.2-2-15}
\expandafter\def\csname sga@target@IV.2.2.16\endcsname{ega4.c04.cor.2-2-16}
\expandafter\def\csname sga@target@IV.2.2.17\endcsname{ega4.c04.prop.2-2-17}
\expandafter\def\csname sga@target@IV.2.2.2\endcsname{ega4.c04.prop.2-2-2}
\expandafter\def\csname sga@target@IV.2.2.2.1\endcsname{ega4.c04.eq.2-2-2-1}
\expandafter\def\csname sga@target@IV.2.2.3\endcsname{ega4.c04.cor.2-2-3}
\expandafter\def\csname sga@target@IV.2.2.4\endcsname{ega4.c04.def.2-2-4}
\expandafter\def\csname sga@target@IV.2.2.5\endcsname{ega4.c04.cor.2-2-5}
\expandafter\def\csname sga@target@IV.2.2.6\endcsname{ega4.c04.cor.2-2-6}
\expandafter\def\csname sga@target@IV.2.2.7\endcsname{ega4.c04.prop.2-2-7}
\expandafter\def\csname sga@target@IV.2.2.8\endcsname{ega4.c04.cor.2-2-8}
\expandafter\def\csname sga@target@IV.2.2.8.1\endcsname{ega4.c04.eq.2-2-8-1}
\expandafter\def\csname sga@target@IV.2.2.9\endcsname{ega4.c04.rem.2-2-9}
\expandafter\def\csname sga@target@IV.2.3.1\endcsname{ega4.c04.lem.2-3-1}
\expandafter\def\csname sga@target@IV.2.3.1.1\endcsname{ega4.c04.eq.2-3-1-1}
\expandafter\def\csname sga@target@IV.2.3.10\endcsname{ega4.c04.thm.2-3-10}
\expandafter\def\csname sga@target@IV.2.3.11\endcsname{ega4.c04.cor.2-3-11}
\expandafter\def\csname sga@target@IV.2.3.12\endcsname{ega4.c04.cor.2-3-12}
\expandafter\def\csname sga@target@IV.2.3.13\endcsname{ega4.c04.cor.2-3-13}
\expandafter\def\csname sga@target@IV.2.3.14\endcsname{ega4.c04.cor.2-3-14}
\expandafter\def\csname sga@target@IV.2.3.15\endcsname{ega4.c04.rem.2-3-15}
\expandafter\def\csname sga@target@IV.2.3.2\endcsname{ega4.c04.prop.2-3-2}
\expandafter\def\csname sga@target@IV.2.3.3\endcsname{ega4.c04.statement.2-3-3}
\expandafter\def\csname sga@target@IV.2.3.4\endcsname{ega4.c04.prop.2-3-4}
\expandafter\def\csname sga@target@IV.2.3.5\endcsname{ega4.c04.cor.2-3-5}
\expandafter\def\csname sga@target@IV.2.3.6\endcsname{ega4.c04.prop.2-3-6}
\expandafter\def\csname sga@target@IV.2.3.7\endcsname{ega4.c04.prop.2-3-7}
\expandafter\def\csname sga@target@IV.2.3.8\endcsname{ega4.c04.cor.2-3-8}
\expandafter\def\csname sga@target@IV.2.3.9\endcsname{ega4.c04.rem.2-3-9}
\expandafter\def\csname sga@target@IV.2.4.1\endcsname{ega4.c04.statement.2-4-1}
\expandafter\def\csname sga@target@IV.2.4.10\endcsname{ega4.c04.cor.2-4-10}
\expandafter\def\csname sga@target@IV.2.4.11\endcsname{ega4.c04.rem.2-4-11}
\expandafter\def\csname sga@target@IV.2.4.2\endcsname{ega4.c04.def.2-4-2}
\expandafter\def\csname sga@target@IV.2.4.3\endcsname{ega4.c04.prop.2-4-3}
\expandafter\def\csname sga@target@IV.2.4.4\endcsname{ega4.c04.prop.2-4-4}
\expandafter\def\csname sga@target@IV.2.4.5\endcsname{ega4.c04.prop.2-4-5}
\expandafter\def\csname sga@target@IV.2.4.6\endcsname{ega4.c04.thm.2-4-6}
\expandafter\def\csname sga@target@IV.2.4.7\endcsname{ega4.c04.cor.2-4-7}
\expandafter\def\csname sga@target@IV.2.4.8\endcsname{ega4.c04.rem.2-4-8}
\expandafter\def\csname sga@target@IV.2.4.9\endcsname{ega4.c04.prop.2-4-9}
\expandafter\def\csname sga@target@IV.2.5.1\endcsname{ega4.c04.prop.2-5-1}
\expandafter\def\csname sga@target@IV.2.5.2\endcsname{ega4.c04.prop.2-5-2}
\expandafter\def\csname sga@target@IV.2.5.3\endcsname{ega4.c04.lem.2-5-3}
\expandafter\def\csname sga@target@IV.2.5.4\endcsname{ega4.c04.rem.2-5-4}
\expandafter\def\csname sga@target@IV.2.5.4.1\endcsname{ega4.c04.lem.2-5-4-1}
\expandafter\def\csname sga@target@IV.2.5.5\endcsname{ega4.c04.prop.2-5-5}
\expandafter\def\csname sga@target@IV.2.5.5.1\endcsname{ega4.c04.statement.2-5-5-1}
\expandafter\def\csname sga@target@IV.2.5.5.2\endcsname{ega4.c04.lem.2-5-5-2}
\expandafter\def\csname sga@target@IV.2.5.5.3\endcsname{ega4.c04.statement.2-5-5-3}
\expandafter\def\csname sga@target@IV.2.5.5.4\endcsname{ega4.c04.rem.2-5-5-4}
\expandafter\def\csname sga@target@IV.2.5.6\endcsname{ega4.c04.lem.2-5-6}
\expandafter\def\csname sga@target@IV.2.5.7\endcsname{ega4.c04.cor.2-5-7}
\expandafter\def\csname sga@target@IV.2.5.8\endcsname{ega4.c04.prop.2-5-8}
\expandafter\def\csname sga@target@IV.2.5.8.1\endcsname{ega4.c04.statement.2-5-8-1}
\expandafter\def\csname sga@target@IV.2.5.8.2\endcsname{ega4.c04.statement.2-5-8-2}
\expandafter\def\csname sga@target@IV.2.5.8.3\endcsname{ega4.c04.statement.2-5-8-3}
\expandafter\def\csname sga@target@IV.2.5.8.4\endcsname{ega4.c04.statement.2-5-8-4}
\expandafter\def\csname sga@target@IV.2.5.8.5\endcsname{ega4.c04.statement.2-5-8-5}
\expandafter\def\csname sga@target@IV.2.5.8.6\endcsname{ega4.c04.statement.2-5-8-6}
\expandafter\def\csname sga@target@IV.2.5.8.7\endcsname{ega4.c04.statement.2-5-8-7}
\expandafter\def\csname sga@target@IV.2.5.8.8\endcsname{ega4.c04.statement.2-5-8-8}
\expandafter\def\csname sga@target@IV.2.5.8.9\endcsname{ega4.c04.statement.2-5-8-9}
\expandafter\def\csname sga@target@IV.2.6.1\endcsname{ega4.c04.prop.2-6-1}
\expandafter\def\csname sga@target@IV.2.6.2\endcsname{ega4.c04.prop.2-6-2}
\expandafter\def\csname sga@target@IV.2.6.3\endcsname{ega4.c04.rem.2-6-3}
\expandafter\def\csname sga@target@IV.2.6.4\endcsname{ega4.c04.cor.2-6-4}
\expandafter\def\csname sga@target@IV.2.7.1\endcsname{ega4.c04.prop.2-7-1}
\expandafter\def\csname sga@target@IV.2.7.1.1\endcsname{ega4.c04.lem.2-7-1-1}
\expandafter\def\csname sga@target@IV.2.7.1.viii\endcsname{ega4.c04.item.2-7-1-viii}
\expandafter\def\csname sga@target@IV.2.7.2\endcsname{ega4.c04.cor.2-7-2}
\expandafter\def\csname sga@target@IV.2.7.3\endcsname{ega4.c04.sec.2-7-3}
\expandafter\def\csname sga@target@IV.2.8.1\endcsname{ega4.c04.prop.2-8-1}
\expandafter\def\csname sga@target@IV.2.8.1.1\endcsname{ega4.c04.lem.2-8-1-1}
\expandafter\def\csname sga@target@IV.2.8.2\endcsname{ega4.c04.cor.2-8-2}
\expandafter\def\csname sga@target@IV.2.8.3\endcsname{ega4.c04.statement.2-8-3}
\expandafter\def\csname sga@target@IV.2.8.4\endcsname{ega4.c04.prop.2-8-4}
\expandafter\def\csname sga@target@IV.2.8.4.1\endcsname{ega4.c04.statement.2-8-4-1}
\expandafter\def\csname sga@target@IV.2.8.5\endcsname{ega4.c04.prop.2-8-5}
\expandafter\def\csname sga@target@IV.2.8.6\endcsname{ega4.c04.cor.2-8-6}
\expandafter\def\csname sga@target@IV.2.8.6.1\endcsname{ega4.c04.statement.2-8-6-1}
\expandafter\def\csname sga@target@IV.20\endcsname{ega4.c04.sec.20}
\expandafter\def\csname sga@target@IV.20.0\endcsname{ega4.c04.subsec.20-0}
\expandafter\def\csname sga@target@IV.20.1\endcsname{ega4.c04.subsec.20-1}
\expandafter\def\csname sga@target@IV.20.1.1\endcsname{ega4.c04.statement.20-1-1}
\expandafter\def\csname sga@target@IV.20.1.1.1\endcsname{ega4.c04.eq.20-1-1-1}
\expandafter\def\csname sga@target@IV.20.1.10\endcsname{ega4.c04.statement.20-1-10}
\expandafter\def\csname sga@target@IV.20.1.11\endcsname{ega4.c04.statement.20-1-11}
\expandafter\def\csname sga@target@IV.20.1.11.1\endcsname{ega4.c04.eq.20-1-11-1}
\expandafter\def\csname sga@target@IV.20.1.11.2\endcsname{ega4.c04.eq.20-1-11-2}
\expandafter\def\csname sga@target@IV.20.1.12\endcsname{ega4.c04.prop.20-1-12}
\expandafter\def\csname sga@target@IV.20.1.13\endcsname{ega4.c04.rem.20-1-13}
\expandafter\def\csname sga@target@IV.20.1.2\endcsname{ega4.c04.statement.20-1-2}
\expandafter\def\csname sga@target@IV.20.1.2.1\endcsname{ega4.c04.statement.20-1-2-1}
\expandafter\def\csname sga@target@IV.20.1.2.2\endcsname{ega4.c04.statement.20-1-2-2}
\expandafter\def\csname sga@target@IV.20.1.3\endcsname{ega4.c04.statement.20-1-3}
\expandafter\def\csname sga@target@IV.20.1.3.1\endcsname{ega4.c04.statement.20-1-3-1}
\expandafter\def\csname sga@target@IV.20.1.4\endcsname{ega4.c04.statement.20-1-4}
\expandafter\def\csname sga@target@IV.20.1.4.1\endcsname{ega4.c04.statement.20-1-4-1}
\expandafter\def\csname sga@target@IV.20.1.5\endcsname{ega4.c04.statement.20-1-5}
\expandafter\def\csname sga@target@IV.20.1.5.1\endcsname{ega4.c04.statement.20-1-5-1}
\expandafter\def\csname sga@target@IV.20.1.6\endcsname{ega4.c04.prop.20-1-6}
\expandafter\def\csname sga@target@IV.20.1.7\endcsname{ega4.c04.statement.20-1-7}
\expandafter\def\csname sga@target@IV.20.1.8\endcsname{ega4.c04.statement.20-1-8}
\expandafter\def\csname sga@target@IV.20.1.9\endcsname{ega4.c04.statement.20-1-9}
\expandafter\def\csname sga@target@IV.20.2\endcsname{ega4.c04.subsec.20-2}
\expandafter\def\csname sga@target@IV.20.2.1\endcsname{ega4.c04.statement.20-2-1}
\expandafter\def\csname sga@target@IV.20.2.10\endcsname{ega4.c04.statement.20-2-10}
\expandafter\def\csname sga@target@IV.20.2.10.1\endcsname{ega4.c04.eq.20-2-10-1}
\expandafter\def\csname sga@target@IV.20.2.10.2\endcsname{ega4.c04.eq.20-2-10-2}
\expandafter\def\csname sga@target@IV.20.2.11\endcsname{ega4.c04.prop.20-2-11}
\expandafter\def\csname sga@target@IV.20.2.11.1\endcsname{ega4.c04.eq.20-2-11-1}
\expandafter\def\csname sga@target@IV.20.2.12\endcsname{ega4.c04.cor.20-2-12}
\expandafter\def\csname sga@target@IV.20.2.13\endcsname{ega4.c04.statement.20-2-13}
\expandafter\def\csname sga@target@IV.20.2.14\endcsname{ega4.c04.statement.20-2-14}
\expandafter\def\csname sga@target@IV.20.2.15\endcsname{ega4.c04.prop.20-2-15}
\expandafter\def\csname sga@target@IV.20.2.16\endcsname{ega4.c04.rem.20-2-16}
\expandafter\def\csname sga@target@IV.20.2.2\endcsname{ega4.c04.statement.20-2-2}
\expandafter\def\csname sga@target@IV.20.2.3\endcsname{ega4.c04.statement.20-2-3}
\expandafter\def\csname sga@target@IV.20.2.3.1\endcsname{ega4.c04.statement.20-2-3-1}
\expandafter\def\csname sga@target@IV.20.2.4\endcsname{ega4.c04.prop.20-2-4}
\expandafter\def\csname sga@target@IV.20.2.5\endcsname{ega4.c04.cor.20-2-5}
\expandafter\def\csname sga@target@IV.20.2.6\endcsname{ega4.c04.cor.20-2-6}
\expandafter\def\csname sga@target@IV.20.2.7\endcsname{ega4.c04.rem.20-2-7}
\expandafter\def\csname sga@target@IV.20.2.8\endcsname{ega4.c04.statement.20-2-8}
\expandafter\def\csname sga@target@IV.20.2.9\endcsname{ega4.c04.lem.20-2-9}
\expandafter\def\csname sga@target@IV.20.3\endcsname{ega4.c04.subsec.20-3}
\expandafter\def\csname sga@target@IV.20.3.1\endcsname{ega4.c04.statement.20-3-1}
\expandafter\def\csname sga@target@IV.20.3.10\endcsname{ega4.c04.lem.20-3-10}
\expandafter\def\csname sga@target@IV.20.3.11\endcsname{ega4.c04.prop.20-3-11}
\expandafter\def\csname sga@target@IV.20.3.11.1\endcsname{ega4.c04.eq.20-3-11-1}
\expandafter\def\csname sga@target@IV.20.3.12\endcsname{ega4.c04.cor.20-3-12}
\expandafter\def\csname sga@target@IV.20.3.13\endcsname{ega4.c04.rem.20-3-13}
\expandafter\def\csname sga@target@IV.20.3.2\endcsname{ega4.c04.statement.20-3-2}
\expandafter\def\csname sga@target@IV.20.3.2.1\endcsname{ega4.c04.eq.20-3-2-1}
\expandafter\def\csname sga@target@IV.20.3.3\endcsname{ega4.c04.statement.20-3-3}
\expandafter\def\csname sga@target@IV.20.3.4\endcsname{ega4.c04.statement.20-3-4}
\expandafter\def\csname sga@target@IV.20.3.5\endcsname{ega4.c04.prop.20-3-5}
\expandafter\def\csname sga@target@IV.20.3.6\endcsname{ega4.c04.statement.20-3-6}
\expandafter\def\csname sga@target@IV.20.3.6.1\endcsname{ega4.c04.eq.20-3-6-1}
\expandafter\def\csname sga@target@IV.20.3.7\endcsname{ega4.c04.prop.20-3-7}
\expandafter\def\csname sga@target@IV.20.3.8\endcsname{ega4.c04.prop.20-3-8}
\expandafter\def\csname sga@target@IV.20.3.8.1\endcsname{ega4.c04.eq.20-3-8-1}
\expandafter\def\csname sga@target@IV.20.3.8.2\endcsname{ega4.c04.lem.20-3-8-2}
\expandafter\def\csname sga@target@IV.20.3.8.3\endcsname{ega4.c04.lem.20-3-8-3}
\expandafter\def\csname sga@target@IV.20.3.8.4\endcsname{ega4.c04.lem.20-3-8-4}
\expandafter\def\csname sga@target@IV.20.3.8.5\endcsname{ega4.c04.sec.20-3-8-5}
\expandafter\def\csname sga@target@IV.20.3.8.6\endcsname{ega4.c04.sec.20-3-8-6}
\expandafter\def\csname sga@target@IV.20.3.8.7\endcsname{ega4.c04.sec.20-3-8-7}
\expandafter\def\csname sga@target@IV.20.3.8.8\endcsname{ega4.c04.lem.20-3-8-8}
\expandafter\def\csname sga@target@IV.20.3.9\endcsname{ega4.c04.cor.20-3-9}
\expandafter\def\csname sga@target@IV.20.4\endcsname{ega4.c04.subsec.20-4}
\expandafter\def\csname sga@target@IV.20.4.1\endcsname{ega4.c04.statement.20-4-1}
\expandafter\def\csname sga@target@IV.20.4.10\endcsname{ega4.c04.cor.20-4-10}
\expandafter\def\csname sga@target@IV.20.4.11\endcsname{ega4.c04.prop.20-4-11}
\expandafter\def\csname sga@target@IV.20.4.12\endcsname{ega4.c04.cor.20-4-12}
\expandafter\def\csname sga@target@IV.20.4.2\endcsname{ega4.c04.statement.20-4-2}
\expandafter\def\csname sga@target@IV.20.4.3\endcsname{ega4.c04.statement.20-4-3}
\expandafter\def\csname sga@target@IV.20.4.4\endcsname{ega4.c04.prop.20-4-4}
\expandafter\def\csname sga@target@IV.20.4.5\endcsname{ega4.c04.prop.20-4-5}
\expandafter\def\csname sga@target@IV.20.4.6\endcsname{ega4.c04.prop.20-4-6}
\expandafter\def\csname sga@target@IV.20.4.7\endcsname{ega4.c04.cor.20-4-7}
\expandafter\def\csname sga@target@IV.20.4.8\endcsname{ega4.c04.rem.20-4-8}
\expandafter\def\csname sga@target@IV.20.4.9\endcsname{ega4.c04.cor.20-4-9}
\expandafter\def\csname sga@target@IV.20.5\endcsname{ega4.c04.subsec.20-5}
\expandafter\def\csname sga@target@IV.20.5.1\endcsname{ega4.c04.statement.20-5-1}
\expandafter\def\csname sga@target@IV.20.5.2\endcsname{ega4.c04.statement.20-5-2}
\expandafter\def\csname sga@target@IV.20.5.2.1\endcsname{ega4.c04.statement.20-5-2-1}
\expandafter\def\csname sga@target@IV.20.5.3\endcsname{ega4.c04.prop.20-5-3}
\expandafter\def\csname sga@target@IV.20.5.4\endcsname{ega4.c04.statement.20-5-4}
\expandafter\def\csname sga@target@IV.20.5.5\endcsname{ega4.c04.statement.20-5-5}
\expandafter\def\csname sga@target@IV.20.5.6\endcsname{ega4.c04.statement.20-5-6}
\expandafter\def\csname sga@target@IV.20.5.7\endcsname{ega4.c04.statement.20-5-7}
\expandafter\def\csname sga@target@IV.20.5.8\endcsname{ega4.c04.statement.20-5-8}
\expandafter\def\csname sga@target@IV.20.6\endcsname{ega4.c04.subsec.20-6}
\expandafter\def\csname sga@target@IV.20.6.1\endcsname{ega4.c04.statement.20-6-1}
\expandafter\def\csname sga@target@IV.20.6.1.1\endcsname{ega4.c04.eq.20-6-1-1}
\expandafter\def\csname sga@target@IV.20.6.2\endcsname{ega4.c04.statement.20-6-2}
\expandafter\def\csname sga@target@IV.20.6.2.1\endcsname{ega4.c04.statement.20-6-2-1}
\expandafter\def\csname sga@target@IV.20.6.2.2\endcsname{ega4.c04.eq.20-6-2-2}
\expandafter\def\csname sga@target@IV.20.6.3\endcsname{ega4.c04.prop.20-6-3}
\expandafter\def\csname sga@target@IV.20.6.4\endcsname{ega4.c04.prop.20-6-4}
\expandafter\def\csname sga@target@IV.20.6.5\endcsname{ega4.c04.statement.20-6-5}
\expandafter\def\csname sga@target@IV.20.6.5.1\endcsname{ega4.c04.eq.20-6-5-1}
\expandafter\def\csname sga@target@IV.20.6.6\endcsname{ega4.c04.prop.20-6-6}
\expandafter\def\csname sga@target@IV.20.6.7\endcsname{ega4.c04.statement.20-6-7}
\expandafter\def\csname sga@target@IV.20.6.8\endcsname{ega4.c04.prop.20-6-8}
\expandafter\def\csname sga@target@IV.20.6.8.1\endcsname{ega4.c04.eq.20-6-8-1}
\expandafter\def\csname sga@target@IV.20.6.9\endcsname{ega4.c04.statement.20-6-9}
\expandafter\def\csname sga@target@IV.20.6.9.1\endcsname{ega4.c04.statement.20-6-9-1}
\expandafter\def\csname sga@target@IV.21.1\endcsname{ega4.c04.subsec.21-1.key-iv-21-1}
\expandafter\def\csname sga@target@IV.21.1.1\endcsname{ega4.c04.statement.21-1-1}
\expandafter\def\csname sga@target@IV.21.1.10\endcsname{ega4.c04.prop.21-1-10}
\expandafter\def\csname sga@target@IV.21.1.2\endcsname{ega4.c04.statement.21-1-2}
\expandafter\def\csname sga@target@IV.21.1.3\endcsname{ega4.c04.statement.21-1-3}
\expandafter\def\csname sga@target@IV.21.1.3.1\endcsname{ega4.c04.statement.21-1-3-1}
\expandafter\def\csname sga@target@IV.21.1.3.2\endcsname{ega4.c04.statement.21-1-3-2}
\expandafter\def\csname sga@target@IV.21.1.3.3\endcsname{ega4.c04.statement.21-1-3-3}
\expandafter\def\csname sga@target@IV.21.1.3.4\endcsname{ega4.c04.statement.21-1-3-4}
\expandafter\def\csname sga@target@IV.21.1.4\endcsname{ega4.c04.statement.21-1-4}
\expandafter\def\csname sga@target@IV.21.1.4.1\endcsname{ega4.c04.eq.21-1-4-1}
\expandafter\def\csname sga@target@IV.21.1.4.2\endcsname{ega4.c04.eq.21-1-4-2}
\expandafter\def\csname sga@target@IV.21.1.4.3\endcsname{ega4.c04.eq.21-1-4-3}
\expandafter\def\csname sga@target@IV.21.1.5\endcsname{ega4.c04.statement.21-1-5}
\expandafter\def\csname sga@target@IV.21.1.5.1\endcsname{ega4.c04.statement.21-1-5-1}
\expandafter\def\csname sga@target@IV.21.1.5.2\endcsname{ega4.c04.eq.21-1-5-2}
\expandafter\def\csname sga@target@IV.21.1.6\endcsname{ega4.c04.def.21-1-6}
\expandafter\def\csname sga@target@IV.21.1.6.1\endcsname{ega4.c04.eq.21-1-6-1}
\expandafter\def\csname sga@target@IV.21.1.6.2\endcsname{ega4.c04.eq.21-1-6-2}
\expandafter\def\csname sga@target@IV.21.1.6.3\endcsname{ega4.c04.eq.21-1-6-3}
\expandafter\def\csname sga@target@IV.21.1.6.4\endcsname{ega4.c04.eq.21-1-6-4}
\expandafter\def\csname sga@target@IV.21.1.7\endcsname{ega4.c04.statement.21-1-7}
\expandafter\def\csname sga@target@IV.21.1.7.1\endcsname{ega4.c04.statement.21-1-7-1}
\expandafter\def\csname sga@target@IV.21.1.8\endcsname{ega4.c04.prop.21-1-8}
\expandafter\def\csname sga@target@IV.21.1.9\endcsname{ega4.c04.cor.21-1-9}
\expandafter\def\csname sga@target@IV.21.10.1\endcsname{ega4.c04.statement.21-10-1}
\expandafter\def\csname sga@target@IV.21.10.1.1\endcsname{ega4.c04.eq.21-10-1-1}
\expandafter\def\csname sga@target@IV.21.10.1.2\endcsname{ega4.c04.eq.21-10-1-2}
\expandafter\def\csname sga@target@IV.21.10.10\endcsname{ega4.c04.statement.21-10-10}
\expandafter\def\csname sga@target@IV.21.10.11\endcsname{ega4.c04.statement.21-10-11}
\expandafter\def\csname sga@target@IV.21.10.11.1\endcsname{ega4.c04.eq.21-10-11-1}
\expandafter\def\csname sga@target@IV.21.10.12\endcsname{ega4.c04.rem.21-10-12}
\expandafter\def\csname sga@target@IV.21.10.13\endcsname{ega4.c04.lem.21-10-13}
\expandafter\def\csname sga@target@IV.21.10.13.1\endcsname{ega4.c04.prop.21-10-13-1}
\expandafter\def\csname sga@target@IV.21.10.13bis\endcsname{ega4.c04.prop.21-10-13bis}
\expandafter\def\csname sga@target@IV.21.10.14\endcsname{ega4.c04.statement.21-10-14}
\expandafter\def\csname sga@target@IV.21.10.14.1\endcsname{ega4.c04.eq.21-10-14-1}
\expandafter\def\csname sga@target@IV.21.10.14.2\endcsname{ega4.c04.eq.21-10-14-2}
\expandafter\def\csname sga@target@IV.21.10.15\endcsname{ega4.c04.prop.21-10-15}
\expandafter\def\csname sga@target@IV.21.10.16\endcsname{ega4.c04.prop.21-10-16}
\expandafter\def\csname sga@target@IV.21.10.16.1\endcsname{ega4.c04.prop.21-10-16-1}
\expandafter\def\csname sga@target@IV.21.10.16.2\endcsname{ega4.c04.eq.21-10-16-2}
\expandafter\def\csname sga@target@IV.21.10.17\endcsname{ega4.c04.prop.21-10-17}
\expandafter\def\csname sga@target@IV.21.10.17.1\endcsname{ega4.c04.prop.21-10-17-1}
\expandafter\def\csname sga@target@IV.21.10.17.10\endcsname{ega4.c04.eq.21-10-17-10}
\expandafter\def\csname sga@target@IV.21.10.17.2\endcsname{ega4.c04.eq.21-10-17-2}
\expandafter\def\csname sga@target@IV.21.10.17.3\endcsname{ega4.c04.lem.21-10-17-3}
\expandafter\def\csname sga@target@IV.21.10.17.4\endcsname{ega4.c04.lem.21-10-17-4}
\expandafter\def\csname sga@target@IV.21.10.17.5\endcsname{ega4.c04.eq.21-10-17-5}
\expandafter\def\csname sga@target@IV.21.10.17.6\endcsname{ega4.c04.eq.21-10-17-6}
\expandafter\def\csname sga@target@IV.21.10.17.7\endcsname{ega4.c04.lem.21-10-17-7}
\expandafter\def\csname sga@target@IV.21.10.17.8\endcsname{ega4.c04.lem.21-10-17-8}
\expandafter\def\csname sga@target@IV.21.10.17.9\endcsname{ega4.c04.eq.21-10-17-9}
\expandafter\def\csname sga@target@IV.21.10.18\endcsname{ega4.c04.prop.21-10-18}
\expandafter\def\csname sga@target@IV.21.10.18.1\endcsname{ega4.c04.prop.21-10-18-1}
\expandafter\def\csname sga@target@IV.21.10.19\endcsname{ega4.c04.cor.21-10-19}
\expandafter\def\csname sga@target@IV.21.10.19.1\endcsname{ega4.c04.cor.21-10-19-1}
\expandafter\def\csname sga@target@IV.21.10.2\endcsname{ega4.c04.statement.21-10-2}
\expandafter\def\csname sga@target@IV.21.10.3\endcsname{ega4.c04.statement.21-10-3}
\expandafter\def\csname sga@target@IV.21.10.3.1\endcsname{ega4.c04.statement.21-10-3-1}
\expandafter\def\csname sga@target@IV.21.10.4\endcsname{ega4.c04.prop.21-10-4}
\expandafter\def\csname sga@target@IV.21.10.4.1\endcsname{ega4.c04.prop.21-10-4-1}
\expandafter\def\csname sga@target@IV.21.10.5\endcsname{ega4.c04.statement.21-10-5}
\expandafter\def\csname sga@target@IV.21.10.6\endcsname{ega4.c04.prop.21-10-6}
\expandafter\def\csname sga@target@IV.21.10.7\endcsname{ega4.c04.rem.21-10-7}
\expandafter\def\csname sga@target@IV.21.10.8\endcsname{ega4.c04.prop.21-10-8}
\expandafter\def\csname sga@target@IV.21.10.9\endcsname{ega4.c04.statement.21-10-9}
\expandafter\def\csname sga@target@IV.21.11\endcsname{ega4.c04.subsec.21-11.key-iv-21-11}
\expandafter\def\csname sga@target@IV.21.11.1\endcsname{ega4.c04.thm.21-11-1}
\expandafter\def\csname sga@target@IV.21.11.1.1\endcsname{ega4.c04.eq.21-11-1-1}
\expandafter\def\csname sga@target@IV.21.11.2\endcsname{ega4.c04.lem.21-11-2}
\expandafter\def\csname sga@target@IV.21.12\endcsname{ega4.c04.subsec.21-12.key-iv-21-12}
\expandafter\def\csname sga@target@IV.21.12.1\endcsname{ega4.c04.statement.21-12-1}
\expandafter\def\csname sga@target@IV.21.12.1.1\endcsname{ega4.c04.eq.21-12-1-1}
\expandafter\def\csname sga@target@IV.21.12.10\endcsname{ega4.c04.prop.21-12-10}
\expandafter\def\csname sga@target@IV.21.12.12\endcsname{ega4.c04.thm.21-12-12}
\expandafter\def\csname sga@target@IV.21.12.13\endcsname{ega4.c04.cor.21-12-13}
\expandafter\def\csname sga@target@IV.21.12.14\endcsname{ega4.c04.item.21-12-14}
\expandafter\def\csname sga@target@IV.21.12.15\endcsname{ega4.c04.prop.21-12-15}
\expandafter\def\csname sga@target@IV.21.12.16\endcsname{ega4.c04.cor.21-12-16}
\expandafter\def\csname sga@target@IV.21.12.17\endcsname{ega4.c04.item.21-12-17}
\expandafter\def\csname sga@target@IV.21.12.2\endcsname{ega4.c04.prop.21-12-2}
\expandafter\def\csname sga@target@IV.21.12.2.1\endcsname{ega4.c04.prop.21-12-2-1}
\expandafter\def\csname sga@target@IV.21.12.3\endcsname{ega4.c04.cor.21-12-3}
\expandafter\def\csname sga@target@IV.21.12.4\endcsname{ega4.c04.eq.21-12-4}
\expandafter\def\csname sga@target@IV.21.12.5\endcsname{ega4.c04.statement.21-12-5}
\expandafter\def\csname sga@target@IV.21.12.5.1\endcsname{ega4.c04.eq.21-12-5-1}
\expandafter\def\csname sga@target@IV.21.12.5.2\endcsname{ega4.c04.eq.21-12-5-2}
\expandafter\def\csname sga@target@IV.21.12.5.3\endcsname{ega4.c04.eq.21-12-5-3}
\expandafter\def\csname sga@target@IV.21.12.6\endcsname{ega4.c04.thm.21-12-6}
\expandafter\def\csname sga@target@IV.21.12.7\endcsname{ega4.c04.cor.21-12-7}
\expandafter\def\csname sga@target@IV.21.12.8\endcsname{ega4.c04.statement.21-12-8}
\expandafter\def\csname sga@target@IV.21.12.9\endcsname{ega4.c04.statement.21-12-9}
\expandafter\def\csname sga@target@IV.21.13\endcsname{ega4.c04.subsec.21-13.key-iv-21-13}
\expandafter\def\csname sga@target@IV.21.13.1\endcsname{ega4.c04.def.21-13-1}
\expandafter\def\csname sga@target@IV.21.13.10\endcsname{ega4.c04.prop.21-13-10}
\expandafter\def\csname sga@target@IV.21.13.11\endcsname{ega4.c04.cor.21-13-11}
\expandafter\def\csname sga@target@IV.21.13.12\endcsname{ega4.c04.prop.21-13-12}
\expandafter\def\csname sga@target@IV.21.13.12.1\endcsname{ega4.c04.statement.21-13-12-1}
\expandafter\def\csname sga@target@IV.21.13.13\endcsname{ega4.c04.rem.21-13-13}
\expandafter\def\csname sga@target@IV.21.13.14\endcsname{ega4.c04.prop.21-13-14}
\expandafter\def\csname sga@target@IV.21.13.15\endcsname{ega4.c04.cor.21-13-15}
\expandafter\def\csname sga@target@IV.21.13.16\endcsname{ega4.c04.rem.21-13-16}
\expandafter\def\csname sga@target@IV.21.13.2\endcsname{ega4.c04.lem.21-13-2}
\expandafter\def\csname sga@target@IV.21.13.2.1\endcsname{ega4.c04.eq.21-13-2-1}
\expandafter\def\csname sga@target@IV.21.13.3\endcsname{ega4.c04.cor.21-13-3}
\expandafter\def\csname sga@target@IV.21.13.4\endcsname{ega4.c04.lem.21-13-4}
\expandafter\def\csname sga@target@IV.21.13.5\endcsname{ega4.c04.prop.21-13-5}
\expandafter\def\csname sga@target@IV.21.13.6\endcsname{ega4.c04.cor.21-13-6}
\expandafter\def\csname sga@target@IV.21.13.7\endcsname{ega4.c04.def.21-13-7}
\expandafter\def\csname sga@target@IV.21.13.8\endcsname{ega4.c04.prop.21-13-8}
\expandafter\def\csname sga@target@IV.21.13.9\endcsname{ega4.c04.statement.21-13-9}
\expandafter\def\csname sga@target@IV.21.13.9.1\endcsname{ega4.c04.eq.21-13-9-1}
\expandafter\def\csname sga@target@IV.21.14\endcsname{ega4.c04.subsec.21-14.key-iv-21-14}
\expandafter\def\csname sga@target@IV.21.14.1\endcsname{ega4.c04.thm.21-14-1}
\expandafter\def\csname sga@target@IV.21.14.2\endcsname{ega4.c04.cor.21-14-2}
\expandafter\def\csname sga@target@IV.21.14.3\endcsname{ega4.c04.prop.21-14-3}
\expandafter\def\csname sga@target@IV.21.14.4\endcsname{ega4.c04.statement.21-14-4}
\expandafter\def\csname sga@target@IV.21.14.4bis\endcsname{ega4.c04.statement.21-14-4bis}
\expandafter\def\csname sga@target@IV.21.15\endcsname{ega4.c04.subsec.21-15.key-iv-21-15}
\expandafter\def\csname sga@target@IV.21.15.1\endcsname{ega4.c04.statement.21-15-1}
\expandafter\def\csname sga@target@IV.21.15.2\endcsname{ega4.c04.def.21-15-2}
\expandafter\def\csname sga@target@IV.21.15.3\endcsname{ega4.c04.statement.21-15-3}
\expandafter\def\csname sga@target@IV.21.15.3.1\endcsname{ega4.c04.eq.21-15-3-1}
\expandafter\def\csname sga@target@IV.21.15.3.2\endcsname{ega4.c04.eq.21-15-3-2}
\expandafter\def\csname sga@target@IV.21.15.3bis\endcsname{ega4.c04.statement.21-15-3bis}
\expandafter\def\csname sga@target@IV.21.15.4\endcsname{ega4.c04.prop.21-15-4}
\expandafter\def\csname sga@target@IV.21.15.5\endcsname{ega4.c04.prop.21-15-5}
\expandafter\def\csname sga@target@IV.21.15.6\endcsname{ega4.c04.prop.21-15-6}
\expandafter\def\csname sga@target@IV.21.15.7\endcsname{ega4.c04.statement.21-15-7}
\expandafter\def\csname sga@target@IV.21.15.8\endcsname{ega4.c04.statement.21-15-8}
\expandafter\def\csname sga@target@IV.21.15.9\endcsname{ega4.c04.statement.21-15-9}
\expandafter\def\csname sga@target@IV.21.15.9.1\endcsname{ega4.c04.eq.21-15-9-1}
\expandafter\def\csname sga@target@IV.21.2.1\endcsname{ega4.c04.statement.21-2-1}
\expandafter\def\csname sga@target@IV.21.2.10\endcsname{ega4.c04.statement.21-2-10}
\expandafter\def\csname sga@target@IV.21.2.10.1\endcsname{ega4.c04.eq.21-2-10-1}
\expandafter\def\csname sga@target@IV.21.2.11\endcsname{ega4.c04.prop.21-2-11}
\expandafter\def\csname sga@target@IV.21.2.11.1\endcsname{ega4.c04.prop.21-2-11-1}
\expandafter\def\csname sga@target@IV.21.2.12\endcsname{ega4.c04.statement.21-2-12}
\expandafter\def\csname sga@target@IV.21.2.2\endcsname{ega4.c04.prop.21-2-2}
\expandafter\def\csname sga@target@IV.21.2.3\endcsname{ega4.c04.cor.21-2-3}
\expandafter\def\csname sga@target@IV.21.2.4\endcsname{ega4.c04.statement.21-2-4}
\expandafter\def\csname sga@target@IV.21.2.5\endcsname{ega4.c04.statement.21-2-5}
\expandafter\def\csname sga@target@IV.21.2.5.1\endcsname{ega4.c04.statement.21-2-5-1}
\expandafter\def\csname sga@target@IV.21.2.5.2\endcsname{ega4.c04.eq.21-2-5-2}
\expandafter\def\csname sga@target@IV.21.2.5.3\endcsname{ega4.c04.eq.21-2-5-3}
\expandafter\def\csname sga@target@IV.21.2.5.4\endcsname{ega4.c04.eq.21-2-5-4}
\expandafter\def\csname sga@target@IV.21.2.6\endcsname{ega4.c04.prop.21-2-6}
\expandafter\def\csname sga@target@IV.21.2.6.1\endcsname{ega4.c04.eq.21-2-6-1}
\expandafter\def\csname sga@target@IV.21.2.7\endcsname{ega4.c04.statement.21-2-7}
\expandafter\def\csname sga@target@IV.21.2.7.1\endcsname{ega4.c04.statement.21-2-7-1}
\expandafter\def\csname sga@target@IV.21.2.7.2\endcsname{ega4.c04.eq.21-2-7-2}
\expandafter\def\csname sga@target@IV.21.2.8\endcsname{ega4.c04.statement.21-2-8}
\expandafter\def\csname sga@target@IV.21.2.8.1\endcsname{ega4.c04.statement.21-2-8-1}
\expandafter\def\csname sga@target@IV.21.2.8.2\endcsname{ega4.c04.statement.21-2-8-2}
\expandafter\def\csname sga@target@IV.21.2.8.3\endcsname{ega4.c04.statement.21-2-8-3}
\expandafter\def\csname sga@target@IV.21.2.9\endcsname{ega4.c04.statement.21-2-9}
\expandafter\def\csname sga@target@IV.21.2.9.1\endcsname{ega4.c04.statement.21-2-9-1}
\expandafter\def\csname sga@target@IV.21.2.9.2\endcsname{ega4.c04.eq.21-2-9-2}
\expandafter\def\csname sga@target@IV.21.2.9.3\endcsname{ega4.c04.eq.21-2-9-3}
\expandafter\def\csname sga@target@IV.21.2.9.4\endcsname{ega4.c04.eq.21-2-9-4}
\expandafter\def\csname sga@target@IV.21.3.1\endcsname{ega4.c04.statement.21-3-1}
\expandafter\def\csname sga@target@IV.21.3.2\endcsname{ega4.c04.statement.21-3-2}
\expandafter\def\csname sga@target@IV.21.3.2.1\endcsname{ega4.c04.statement.21-3-2-1}
\expandafter\def\csname sga@target@IV.21.3.2.2\endcsname{ega4.c04.eq.21-3-2-2}
\expandafter\def\csname sga@target@IV.21.3.2.3\endcsname{ega4.c04.eq.21-3-2-3}
\expandafter\def\csname sga@target@IV.21.3.2.4\endcsname{ega4.c04.eq.21-3-2-4}
\expandafter\def\csname sga@target@IV.21.3.3\endcsname{ega4.c04.prop.21-3-3}
\expandafter\def\csname sga@target@IV.21.3.3.1\endcsname{ega4.c04.eq.21-3-3-1}
\expandafter\def\csname sga@target@IV.21.3.4\endcsname{ega4.c04.prop.21-3-4}
\expandafter\def\csname sga@target@IV.21.3.5\endcsname{ega4.c04.cor.21-3-5}
\expandafter\def\csname sga@target@IV.21.3.6\endcsname{ega4.c04.rem.21-3-6}
\expandafter\def\csname sga@target@IV.21.3.6.1\endcsname{ega4.c04.eq.21-3-6-1}
\expandafter\def\csname sga@target@IV.21.4\endcsname{ega4.c04.subsec.21-4.key-iv-21-4}
\expandafter\def\csname sga@target@IV.21.4.1\endcsname{ega4.c04.statement.21-4-1}
\expandafter\def\csname sga@target@IV.21.4.10\endcsname{ega4.c04.cor.21-4-10}
\expandafter\def\csname sga@target@IV.21.4.11\endcsname{ega4.c04.cor.21-4-11}
\expandafter\def\csname sga@target@IV.21.4.12\endcsname{ega4.c04.rem.21-4-12}
\expandafter\def\csname sga@target@IV.21.4.13\endcsname{ega4.c04.cor.21-4-13}
\expandafter\def\csname sga@target@IV.21.4.2\endcsname{ega4.c04.def.21-4-2}
\expandafter\def\csname sga@target@IV.21.4.2.1\endcsname{ega4.c04.eq.21-4-2-1}
\expandafter\def\csname sga@target@IV.21.4.3\endcsname{ega4.c04.statement.21-4-3}
\expandafter\def\csname sga@target@IV.21.4.4\endcsname{ega4.c04.statement.21-4-4}
\expandafter\def\csname sga@target@IV.21.4.5\endcsname{ega4.c04.prop.21-4-5}
\expandafter\def\csname sga@target@IV.21.4.6\endcsname{ega4.c04.cor.21-4-6}
\expandafter\def\csname sga@target@IV.21.4.6.1\endcsname{ega4.c04.eq.21-4-6-1}
\expandafter\def\csname sga@target@IV.21.4.7\endcsname{ega4.c04.statement.21-4-7}
\expandafter\def\csname sga@target@IV.21.4.8\endcsname{ega4.c04.prop.21-4-8}
\expandafter\def\csname sga@target@IV.21.4.9\endcsname{ega4.c04.prop.21-4-9}
\expandafter\def\csname sga@target@IV.21.5\endcsname{ega4.c04.subsec.21-5.key-iv-21-5}
\expandafter\def\csname sga@target@IV.21.5.1\endcsname{ega4.c04.statement.21-5-1}
\expandafter\def\csname sga@target@IV.21.5.2\endcsname{ega4.c04.lem.21-5-2}
\expandafter\def\csname sga@target@IV.21.5.3\endcsname{ega4.c04.statement.21-5-3}
\expandafter\def\csname sga@target@IV.21.5.3.1\endcsname{ega4.c04.eq.21-5-3-1}
\expandafter\def\csname sga@target@IV.21.5.3.2\endcsname{ega4.c04.eq.21-5-3-2}
\expandafter\def\csname sga@target@IV.21.5.4\endcsname{ega4.c04.statement.21-5-4}
\expandafter\def\csname sga@target@IV.21.5.5\endcsname{ega4.c04.def.21-5-5}
\expandafter\def\csname sga@target@IV.21.5.5.1\endcsname{ega4.c04.eq.21-5-5-1}
\expandafter\def\csname sga@target@IV.21.5.5.2\endcsname{ega4.c04.eq.21-5-5-2}
\expandafter\def\csname sga@target@IV.21.5.5.3\endcsname{ega4.c04.eq.21-5-5-3}
\expandafter\def\csname sga@target@IV.21.5.6\endcsname{ega4.c04.prop.21-5-6}
\expandafter\def\csname sga@target@IV.21.5.6.1\endcsname{ega4.c04.prop.21-5-6-1}
\expandafter\def\csname sga@target@IV.21.5.7\endcsname{ega4.c04.prop.21-5-7}
\expandafter\def\csname sga@target@IV.21.5.7.1\endcsname{ega4.c04.prop.21-5-7-1}
\expandafter\def\csname sga@target@IV.21.5.7.2\endcsname{ega4.c04.lem.21-5-7-2}
\expandafter\def\csname sga@target@IV.21.5.7.3\endcsname{ega4.c04.lem.21-5-7-3}
\expandafter\def\csname sga@target@IV.21.5.7.4\endcsname{ega4.c04.eq.21-5-7-4}
\expandafter\def\csname sga@target@IV.21.5.8\endcsname{ega4.c04.prop.21-5-8}
\expandafter\def\csname sga@target@IV.21.5.8.1\endcsname{ega4.c04.eq.21-5-8-1}
\expandafter\def\csname sga@target@IV.21.6\endcsname{ega4.c04.subsec.21-6.key-iv-21-6}
\expandafter\def\csname sga@target@IV.21.6.1\endcsname{ega4.c04.statement.21-6-1}
\expandafter\def\csname sga@target@IV.21.6.10\endcsname{ega4.c04.cor.21-6-10}
\expandafter\def\csname sga@target@IV.21.6.10.1\endcsname{ega4.c04.eq.21-6-10-1}
\expandafter\def\csname sga@target@IV.21.6.11\endcsname{ega4.c04.cor.21-6-11}
\expandafter\def\csname sga@target@IV.21.6.12\endcsname{ega4.c04.prop.21-6-12}
\expandafter\def\csname sga@target@IV.21.6.12.1\endcsname{ega4.c04.eq.21-6-12-1}
\expandafter\def\csname sga@target@IV.21.6.12.2\endcsname{ega4.c04.eq.21-6-12-2}
\expandafter\def\csname sga@target@IV.21.6.13\endcsname{ega4.c04.cor.21-6-13}
\expandafter\def\csname sga@target@IV.21.6.14\endcsname{ega4.c04.cor.21-6-14}
\expandafter\def\csname sga@target@IV.21.6.2\endcsname{ega4.c04.statement.21-6-2}
\expandafter\def\csname sga@target@IV.21.6.3\endcsname{ega4.c04.statement.21-6-3}
\expandafter\def\csname sga@target@IV.21.6.4\endcsname{ega4.c04.statement.21-6-4}
\expandafter\def\csname sga@target@IV.21.6.4.1\endcsname{ega4.c04.statement.21-6-4-1}
\expandafter\def\csname sga@target@IV.21.6.4.2\endcsname{ega4.c04.eq.21-6-4-2}
\expandafter\def\csname sga@target@IV.21.6.5\endcsname{ega4.c04.statement.21-6-5}
\expandafter\def\csname sga@target@IV.21.6.5.1\endcsname{ega4.c04.eq.21-6-5-1}
\expandafter\def\csname sga@target@IV.21.6.6\endcsname{ega4.c04.prop.21-6-6}
\expandafter\def\csname sga@target@IV.21.6.6.2\endcsname{ega4.c04.statement.21-6-6-2}
\expandafter\def\csname sga@target@IV.21.6.7\endcsname{ega4.c04.statement.21-6-7}
\expandafter\def\csname sga@target@IV.21.6.7.1\endcsname{ega4.c04.eq.21-6-7-1}
\expandafter\def\csname sga@target@IV.21.6.7.2\endcsname{ega4.c04.eq.21-6-7-2}
\expandafter\def\csname sga@target@IV.21.6.7.3\endcsname{ega4.c04.eq.21-6-7-3}
\expandafter\def\csname sga@target@IV.21.6.7.4\endcsname{ega4.c04.eq.21-6-7-4}
\expandafter\def\csname sga@target@IV.21.6.7.5\endcsname{ega4.c04.eq.21-6-7-5}
\expandafter\def\csname sga@target@IV.21.6.8\endcsname{ega4.c04.statement.21-6-8}
\expandafter\def\csname sga@target@IV.21.6.9\endcsname{ega4.c04.thm.21-6-9}
\expandafter\def\csname sga@target@IV.21.6.9.1\endcsname{ega4.c04.eq.21-6-9-1}
\expandafter\def\csname sga@target@IV.21.6.9.2\endcsname{ega4.c04.statement.21-6-9-2}
\expandafter\def\csname sga@target@IV.21.7\endcsname{ega4.c04.subsec.21-7.key-iv-21-7}
\expandafter\def\csname sga@target@IV.21.7.1\endcsname{ega4.c04.statement.21-7-1}
\expandafter\def\csname sga@target@IV.21.7.2\endcsname{ega4.c04.prop.21-7-2}
\expandafter\def\csname sga@target@IV.21.7.2.1\endcsname{ega4.c04.eq.21-7-2-1}
\expandafter\def\csname sga@target@IV.21.7.3\endcsname{ega4.c04.cor.21-7-3}
\expandafter\def\csname sga@target@IV.21.7.3.1\endcsname{ega4.c04.statement.21-7-3-1}
\expandafter\def\csname sga@target@IV.21.7.4\endcsname{ega4.c04.prop.21-7-4}
\expandafter\def\csname sga@target@IV.21.7.5\endcsname{ega4.c04.rem.21-7-5}
\expandafter\def\csname sga@target@IV.21.7.6\endcsname{ega4.c04.cor.21-7-6}
\expandafter\def\csname sga@target@IV.21.7.7\endcsname{ega4.c04.rem.21-7-7}
\expandafter\def\csname sga@target@IV.21.8.1\endcsname{ega4.c04.lem.21-8-1}
\expandafter\def\csname sga@target@IV.21.8.2\endcsname{ega4.c04.cor.21-8-2}
\expandafter\def\csname sga@target@IV.21.8.3\endcsname{ega4.c04.prop.21-8-3}
\expandafter\def\csname sga@target@IV.21.8.3.1\endcsname{ega4.c04.eq.21-8-3-1}
\expandafter\def\csname sga@target@IV.21.8.4\endcsname{ega4.c04.cor.21-8-4}
\expandafter\def\csname sga@target@IV.21.8.5\endcsname{ega4.c04.prop.21-8-5}
\expandafter\def\csname sga@target@IV.21.8.5.1\endcsname{ega4.c04.eq.21-8-5-1}
\expandafter\def\csname sga@target@IV.21.8.5.2\endcsname{ega4.c04.eq.21-8-5-2}
\expandafter\def\csname sga@target@IV.21.8.5.3\endcsname{ega4.c04.eq.21-8-5-3}
\expandafter\def\csname sga@target@IV.21.8.6\endcsname{ega4.c04.rem.21-8-6}
\expandafter\def\csname sga@target@IV.21.9\endcsname{ega4.c04.subsec.21-9.key-iv-21-9}
\expandafter\def\csname sga@target@IV.21.9.1\endcsname{ega4.c04.statement.21-9-1}
\expandafter\def\csname sga@target@IV.21.9.1.1\endcsname{ega4.c04.eq.21-9-1-1}
\expandafter\def\csname sga@target@IV.21.9.1.2\endcsname{ega4.c04.eq.21-9-1-2}
\expandafter\def\csname sga@target@IV.21.9.10\endcsname{ega4.c04.statement.21-9-10}
\expandafter\def\csname sga@target@IV.21.9.11\endcsname{ega4.c04.prop.21-9-11}
\expandafter\def\csname sga@target@IV.21.9.12\endcsname{ega4.c04.cor.21-9-12}
\expandafter\def\csname sga@target@IV.21.9.13\endcsname{ega4.c04.rem.21-9-13}
\expandafter\def\csname sga@target@IV.21.9.2\endcsname{ega4.c04.lem.21-9-2}
\expandafter\def\csname sga@target@IV.21.9.3\endcsname{ega4.c04.rem.21-9-3}
\expandafter\def\csname sga@target@IV.21.9.4\endcsname{ega4.c04.prop.21-9-4}
\expandafter\def\csname sga@target@IV.21.9.4.1\endcsname{ega4.c04.prop.21-9-4-1}
\expandafter\def\csname sga@target@IV.21.9.5\endcsname{ega4.c04.cor.21-9-5}
\expandafter\def\csname sga@target@IV.21.9.6\endcsname{ega4.c04.cor.21-9-6}
\expandafter\def\csname sga@target@IV.21.9.7\endcsname{ega4.c04.cor.21-9-7}
\expandafter\def\csname sga@target@IV.21.9.8\endcsname{ega4.c04.statement.21-9-8}
\expandafter\def\csname sga@target@IV.21.9.9\endcsname{ega4.c04.prop.21-9-9}
\expandafter\def\csname sga@target@IV.3.1.1\endcsname{ega4.c04.def.3-1-1}
\expandafter\def\csname sga@target@IV.3.1.1.1\endcsname{ega4.c04.eq.3-1-1-1}
\expandafter\def\csname sga@target@IV.3.1.10\endcsname{ega4.c04.prop.3-1-10}
\expandafter\def\csname sga@target@IV.3.1.10.1\endcsname{ega4.c04.lem.3-1-10-1}
\expandafter\def\csname sga@target@IV.3.1.10.2\endcsname{ega4.c04.rem.3-1-10-2}
\expandafter\def\csname sga@target@IV.3.1.10.3\endcsname{ega4.c04.eq.3-1-10-3}
\expandafter\def\csname sga@target@IV.3.1.10.4\endcsname{ega4.c04.eq.3-1-10-4}
\expandafter\def\csname sga@target@IV.3.1.10.5\endcsname{ega4.c04.rem.3-1-10-5}
\expandafter\def\csname sga@target@IV.3.1.11\endcsname{ega4.c04.cor.3-1-11}
\expandafter\def\csname sga@target@IV.3.1.12\endcsname{ega4.c04.rem.3-1-12}
\expandafter\def\csname sga@target@IV.3.1.13\endcsname{ega4.c04.prop.3-1-13}
\expandafter\def\csname sga@target@IV.3.1.2\endcsname{ega4.c04.prop.3-1-2}
\expandafter\def\csname sga@target@IV.3.1.3\endcsname{ega4.c04.prop.3-1-3}
\expandafter\def\csname sga@target@IV.3.1.4\endcsname{ega4.c04.cor.3-1-4}
\expandafter\def\csname sga@target@IV.3.1.5\endcsname{ega4.c04.cor.3-1-5}
\expandafter\def\csname sga@target@IV.3.1.6\endcsname{ega4.c04.prop.3-1-6}
\expandafter\def\csname sga@target@IV.3.1.7\endcsname{ega4.c04.prop.3-1-7}
\expandafter\def\csname sga@target@IV.3.1.8\endcsname{ega4.c04.prop.3-1-8}
\expandafter\def\csname sga@target@IV.3.1.9\endcsname{ega4.c04.cor.3-1-9}
\expandafter\def\csname sga@target@IV.3.2.1\endcsname{ega4.c04.prop.3-2-1}
\expandafter\def\csname sga@target@IV.3.2.2\endcsname{ega4.c04.def.3-2-2}
\expandafter\def\csname sga@target@IV.3.2.3\endcsname{ega4.c04.prop.3-2-3}
\expandafter\def\csname sga@target@IV.3.2.4\endcsname{ega4.c04.def.3-2-4}
\expandafter\def\csname sga@target@IV.3.2.5\endcsname{ega4.c04.def.3-2-5}
\expandafter\def\csname sga@target@IV.3.2.6\endcsname{ega4.c04.prop.3-2-6}
\expandafter\def\csname sga@target@IV.3.2.7\endcsname{ega4.c04.cor.3-2-7}
\expandafter\def\csname sga@target@IV.3.2.8\endcsname{ega4.c04.cor.3-2-8}
\expandafter\def\csname sga@target@IV.3.2.9\endcsname{ega4.c04.rem.3-2-9}
\expandafter\def\csname sga@target@IV.3.3.1\endcsname{ega4.c04.prop.3-3-1}
\expandafter\def\csname sga@target@IV.3.3.1.1\endcsname{ega4.c04.prop.3-3-1-1}
\expandafter\def\csname sga@target@IV.3.3.2\endcsname{ega4.c04.cor.3-3-2}
\expandafter\def\csname sga@target@IV.3.3.3\endcsname{ega4.c04.cor.3-3-3}
\expandafter\def\csname sga@target@IV.3.3.4\endcsname{ega4.c04.prop.3-3-4}
\expandafter\def\csname sga@target@IV.3.3.5\endcsname{ega4.c04.cor.3-3-5}
\expandafter\def\csname sga@target@IV.3.3.6\endcsname{ega4.c04.prop.3-3-6}
\expandafter\def\csname sga@target@IV.3.3.6.1\endcsname{ega4.c04.eq.3-3-6-1}
\expandafter\def\csname sga@target@IV.3.3.6.2\endcsname{ega4.c04.eq.3-3-6-2}
\expandafter\def\csname sga@target@IV.3.3.6.3\endcsname{ega4.c04.eq.3-3-6-3}
\expandafter\def\csname sga@target@IV.3.3.6.4\endcsname{ega4.c04.eq.3-3-6-4}
\expandafter\def\csname sga@target@IV.3.3.7\endcsname{ega4.c04.cor.3-3-7}
\expandafter\def\csname sga@target@IV.3.3.7.1\endcsname{ega4.c04.cor.3-3-7-1}
\expandafter\def\csname sga@target@IV.3.3.8\endcsname{ega4.c04.sec.3-3-8}
\expandafter\def\csname sga@target@IV.3.3.9\endcsname{ega4.c04.cor.3-3-9}
\expandafter\def\csname sga@target@IV.3.3.9.1\endcsname{ega4.c04.cor.3-3-9-1}
\expandafter\def\csname sga@target@IV.3.4\endcsname{ega4.c04.subsec.3-4.key-iv-3-4}
\expandafter\def\csname sga@target@IV.3.4.1\endcsname{ega4.c04.prop.3-4-1}
\expandafter\def\csname sga@target@IV.3.4.1.1\endcsname{ega4.c04.eq.3-4-1-1}
\expandafter\def\csname sga@target@IV.3.4.1.2\endcsname{ega4.c04.lem.3-4-1-2}
\expandafter\def\csname sga@target@IV.3.4.1.3\endcsname{ega4.c04.lem.3-4-1-3}
\expandafter\def\csname sga@target@IV.3.4.1.4\endcsname{ega4.c04.lem.3-4-1-4}
\expandafter\def\csname sga@target@IV.3.4.1.5\endcsname{ega4.c04.eq.3-4-1-5}
\expandafter\def\csname sga@target@IV.3.4.2\endcsname{ega4.c04.cor.3-4-2}
\expandafter\def\csname sga@target@IV.3.4.3\endcsname{ega4.c04.prop.3-4-3}
\expandafter\def\csname sga@target@IV.3.4.4\endcsname{ega4.c04.prop.3-4-4}
\expandafter\def\csname sga@target@IV.3.4.5\endcsname{ega4.c04.prop.3-4-5}
\expandafter\def\csname sga@target@IV.3.4.6\endcsname{ega4.c04.prop.3-4-6}
\expandafter\def\csname sga@target@IV.3.4.6.1\endcsname{ega4.c04.lem.3-4-6-1}
\expandafter\def\csname sga@target@IV.3.4.7\endcsname{ega4.c04.cor.3-4-7}
\expandafter\def\csname sga@target@IV.3.4.8\endcsname{ega4.c04.rem.3-4-8}
\expandafter\def\csname sga@target@IV.3.4.9\endcsname{ega4.c04.prop.3-4-9}
\expandafter\def\csname sga@target@IV.4.1.1\endcsname{ega4.c04.def.4-1-1}
\expandafter\def\csname sga@target@IV.4.1.1.1\endcsname{ega4.c04.eq.4-1-1-1}
\expandafter\def\csname sga@target@IV.4.1.1.2\endcsname{ega4.c04.eq.4-1-1-2}
\expandafter\def\csname sga@target@IV.4.1.1.3\endcsname{ega4.c04.eq.4-1-1-3}
\expandafter\def\csname sga@target@IV.4.1.2\endcsname{ega4.c04.thm.4-1-2}
\expandafter\def\csname sga@target@IV.4.1.2.1\endcsname{ega4.c04.cor.4-1-2-1}
\expandafter\def\csname sga@target@IV.4.1.3\endcsname{ega4.c04.rem.4-1-3}
\expandafter\def\csname sga@target@IV.4.1.4\endcsname{ega4.c04.cor.4-1-4}
\expandafter\def\csname sga@target@IV.4.1.5\endcsname{ega4.c04.cor.4-1-5}
\expandafter\def\csname sga@target@IV.4.2.1\endcsname{ega4.c04.prop.4-2-1}
\expandafter\def\csname sga@target@IV.4.2.1.1\endcsname{ega4.c04.eq.4-2-1-1}
\expandafter\def\csname sga@target@IV.4.2.1.2\endcsname{ega4.c04.eq.4-2-1-2}
\expandafter\def\csname sga@target@IV.4.2.1.3\endcsname{ega4.c04.eq.4-2-1-3}
\expandafter\def\csname sga@target@IV.4.2.1.4\endcsname{ega4.c04.rem.4-2-1-4}
\expandafter\def\csname sga@target@IV.4.2.1.5\endcsname{ega4.c04.eq.4-2-1-5}
\expandafter\def\csname sga@target@IV.4.2.10\endcsname{ega4.c04.cor.4-2-10}
\expandafter\def\csname sga@target@IV.4.2.2\endcsname{ega4.c04.cor.4-2-2}
\expandafter\def\csname sga@target@IV.4.2.3\endcsname{ega4.c04.cor.4-2-3}
\expandafter\def\csname sga@target@IV.4.2.4\endcsname{ega4.c04.prop.4-2-4}
\expandafter\def\csname sga@target@IV.4.2.4.i\endcsname{ega4.c04.sec.4-2-4-i}
\expandafter\def\csname sga@target@IV.4.2.4.ii\endcsname{ega4.c04.sec.4-2-4-ii}
\expandafter\def\csname sga@target@IV.4.2.4.iii\endcsname{ega4.c04.sec.4-2-4-iii}
\expandafter\def\csname sga@target@IV.4.2.5\endcsname{ega4.c04.prop.4-2-5}
\expandafter\def\csname sga@target@IV.4.2.6\endcsname{ega4.c04.cor.4-2-6}
\expandafter\def\csname sga@target@IV.4.2.7\endcsname{ega4.c04.prop.4-2-7}
\expandafter\def\csname sga@target@IV.4.2.7.i\endcsname{ega4.c04.sec.4-2-7-i}
\expandafter\def\csname sga@target@IV.4.2.7.ii\endcsname{ega4.c04.sec.4-2-7-ii}
\expandafter\def\csname sga@target@IV.4.2.7.iii\endcsname{ega4.c04.sec.4-2-7-iii}
\expandafter\def\csname sga@target@IV.4.2.7.iv\endcsname{ega4.c04.sec.4-2-7-iv}
\expandafter\def\csname sga@target@IV.4.2.8\endcsname{ega4.c04.cor.4-2-8}
\expandafter\def\csname sga@target@IV.4.2.9\endcsname{ega4.c04.prop.4-2-9}
\expandafter\def\csname sga@target@IV.4.3.1\endcsname{ega4.c04.statement.4-3-1}
\expandafter\def\csname sga@target@IV.4.3.2\endcsname{ega4.c04.prop.4-3-2}
\expandafter\def\csname sga@target@IV.4.3.3\endcsname{ega4.c04.cor.4-3-3}
\expandafter\def\csname sga@target@IV.4.3.4\endcsname{ega4.c04.cor.4-3-4}
\expandafter\def\csname sga@target@IV.4.3.5\endcsname{ega4.c04.prop.4-3-5}
\expandafter\def\csname sga@target@IV.4.3.6\endcsname{ega4.c04.cor.4-3-6}
\expandafter\def\csname sga@target@IV.4.3.7\endcsname{ega4.c04.cor.4-3-7}
\expandafter\def\csname sga@target@IV.4.3.8\endcsname{ega4.c04.cor.4-3-8}
\expandafter\def\csname sga@target@IV.4.3.9\endcsname{ega4.c04.prop.4-3-9}
\expandafter\def\csname sga@target@IV.4.4.1\endcsname{ega4.c04.statement.4-4-1}
\expandafter\def\csname sga@target@IV.4.4.2\endcsname{ega4.c04.lem.4-4-2}
\expandafter\def\csname sga@target@IV.4.4.2.i\endcsname{ega4.c04.sec.4-4-2-i}
\expandafter\def\csname sga@target@IV.4.4.2.ii\endcsname{ega4.c04.sec.4-4-2-ii}
\expandafter\def\csname sga@target@IV.4.4.3\endcsname{ega4.c04.rem.4-4-3}
\expandafter\def\csname sga@target@IV.4.4.4\endcsname{ega4.c04.thm.4-4-4}
\expandafter\def\csname sga@target@IV.4.4.5\endcsname{ega4.c04.cor.4-4-5}
\expandafter\def\csname sga@target@IV.4.4.6\endcsname{ega4.c04.cor.4-4-6}
\expandafter\def\csname sga@target@IV.4.5\endcsname{ega4.c04.subsec.4-5.key-iv-4-5}
\expandafter\def\csname sga@target@IV.4.5.1\endcsname{ega4.c04.prop.4-5-1}
\expandafter\def\csname sga@target@IV.4.5.10\endcsname{ega4.c04.cor.4-5-10}
\expandafter\def\csname sga@target@IV.4.5.11\endcsname{ega4.c04.cor.4-5-11}
\expandafter\def\csname sga@target@IV.4.5.12\endcsname{ega4.c04.sec.4-5-12}
\expandafter\def\csname sga@target@IV.4.5.12.i\endcsname{ega4.c04.sec.4-5-12-i}
\expandafter\def\csname sga@target@IV.4.5.12.ii\endcsname{ega4.c04.sec.4-5-12-ii}
\expandafter\def\csname sga@target@IV.4.5.13\endcsname{ega4.c04.prop.4-5-13}
\expandafter\def\csname sga@target@IV.4.5.13.1\endcsname{ega4.c04.cor.4-5-13-1}
\expandafter\def\csname sga@target@IV.4.5.13.1.i\endcsname{ega4.c04.sec.4-5-13-1-i}
\expandafter\def\csname sga@target@IV.4.5.13.1.ii\endcsname{ega4.c04.sec.4-5-13-1-ii}
\expandafter\def\csname sga@target@IV.4.5.14\endcsname{ega4.c04.cor.4-5-14}
\expandafter\def\csname sga@target@IV.4.5.15\endcsname{ega4.c04.prop.4-5-15}
\expandafter\def\csname sga@target@IV.4.5.15.i\endcsname{ega4.c04.sec.4-5-15-i}
\expandafter\def\csname sga@target@IV.4.5.15.ii\endcsname{ega4.c04.sec.4-5-15-ii}
\expandafter\def\csname sga@target@IV.4.5.16\endcsname{ega4.c04.cor.4-5-16}
\expandafter\def\csname sga@target@IV.4.5.16.i\endcsname{ega4.c04.sec.4-5-16-i}
\expandafter\def\csname sga@target@IV.4.5.16.ii\endcsname{ega4.c04.sec.4-5-16-ii}
\expandafter\def\csname sga@target@IV.4.5.17\endcsname{ega4.c04.cor.4-5-17}
\expandafter\def\csname sga@target@IV.4.5.18\endcsname{ega4.c04.rem.4-5-18}
\expandafter\def\csname sga@target@IV.4.5.19\endcsname{ega4.c04.prop.4-5-19}
\expandafter\def\csname sga@target@IV.4.5.19.1\endcsname{ega4.c04.lem.4-5-19-1}
\expandafter\def\csname sga@target@IV.4.5.19.2\endcsname{ega4.c04.eq.4-5-19-2}
\expandafter\def\csname sga@target@IV.4.5.19.3\endcsname{ega4.c04.cor.4-5-19-3}
\expandafter\def\csname sga@target@IV.4.5.2\endcsname{ega4.c04.def.4-5-2}
\expandafter\def\csname sga@target@IV.4.5.2.a\endcsname{ega4.c04.sec.4-5-2-a}
\expandafter\def\csname sga@target@IV.4.5.2.b\endcsname{ega4.c04.sec.4-5-2-b}
\expandafter\def\csname sga@target@IV.4.5.2.c\endcsname{ega4.c04.sec.4-5-2-c}
\expandafter\def\csname sga@target@IV.4.5.20\endcsname{ega4.c04.rem.4-5-20}
\expandafter\def\csname sga@target@IV.4.5.21\endcsname{ega4.c04.prop.4-5-21}
\expandafter\def\csname sga@target@IV.4.5.3\endcsname{ega4.c04.statement.4-5-3}
\expandafter\def\csname sga@target@IV.4.5.4\endcsname{ega4.c04.prop.4-5-4}
\expandafter\def\csname sga@target@IV.4.5.5\endcsname{ega4.c04.def.4-5-5}
\expandafter\def\csname sga@target@IV.4.5.6\endcsname{ega4.c04.prop.4-5-6}
\expandafter\def\csname sga@target@IV.4.5.6.i\endcsname{ega4.c04.sec.4-5-6-i}
\expandafter\def\csname sga@target@IV.4.5.6.ii\endcsname{ega4.c04.sec.4-5-6-ii}
\expandafter\def\csname sga@target@IV.4.5.7\endcsname{ega4.c04.prop.4-5-7}
\expandafter\def\csname sga@target@IV.4.5.8\endcsname{ega4.c04.cor.4-5-8}
\expandafter\def\csname sga@target@IV.4.5.8.i\endcsname{ega4.c04.sec.4-5-8-i}
\expandafter\def\csname sga@target@IV.4.5.8.ii\endcsname{ega4.c04.sec.4-5-8-ii}
\expandafter\def\csname sga@target@IV.4.5.9\endcsname{ega4.c04.prop.4-5-9}
\expandafter\def\csname sga@target@IV.4.5.9.a\endcsname{ega4.c04.sec.4-5-9-a}
\expandafter\def\csname sga@target@IV.4.5.9.b\endcsname{ega4.c04.sec.4-5-9-b}
\expandafter\def\csname sga@target@IV.4.5.9.c\endcsname{ega4.c04.sec.4-5-9-c}
\expandafter\def\csname sga@target@IV.4.6\endcsname{ega4.c04.subsec.4-6.key-iv-4-6}
\expandafter\def\csname sga@target@IV.4.6.1\endcsname{ega4.c04.prop.4-6-1}
\expandafter\def\csname sga@target@IV.4.6.1.a\endcsname{ega4.c04.sec.4-6-1-a}
\expandafter\def\csname sga@target@IV.4.6.1.b\endcsname{ega4.c04.sec.4-6-1-b}
\expandafter\def\csname sga@target@IV.4.6.1.c\endcsname{ega4.c04.sec.4-6-1-c}
\expandafter\def\csname sga@target@IV.4.6.1.d\endcsname{ega4.c04.sec.4-6-1-d}
\expandafter\def\csname sga@target@IV.4.6.1.e\endcsname{ega4.c04.sec.4-6-1-e}
\expandafter\def\csname sga@target@IV.4.6.10\endcsname{ega4.c04.prop.4-6-10}
\expandafter\def\csname sga@target@IV.4.6.11\endcsname{ega4.c04.cor.4-6-11}
\expandafter\def\csname sga@target@IV.4.6.12\endcsname{ega4.c04.prop.4-6-12}
\expandafter\def\csname sga@target@IV.4.6.12.a\endcsname{ega4.c04.sec.4-6-12-a}
\expandafter\def\csname sga@target@IV.4.6.12.b\endcsname{ega4.c04.sec.4-6-12-b}
\expandafter\def\csname sga@target@IV.4.6.12.c\endcsname{ega4.c04.sec.4-6-12-c}
\expandafter\def\csname sga@target@IV.4.6.12.d\endcsname{ega4.c04.sec.4-6-12-d}
\expandafter\def\csname sga@target@IV.4.6.12.e\endcsname{ega4.c04.sec.4-6-12-e}
\expandafter\def\csname sga@target@IV.4.6.13\endcsname{ega4.c04.cor.4-6-13}
\expandafter\def\csname sga@target@IV.4.6.13.f\endcsname{ega4.c04.sec.4-6-13-f}
\expandafter\def\csname sga@target@IV.4.6.14\endcsname{ega4.c04.statement.4-6-14}
\expandafter\def\csname sga@target@IV.4.6.15\endcsname{ega4.c04.prop.4-6-15}
\expandafter\def\csname sga@target@IV.4.6.15.i\endcsname{ega4.c04.sec.4-6-15-i}
\expandafter\def\csname sga@target@IV.4.6.15.ii\endcsname{ega4.c04.sec.4-6-15-ii}
\expandafter\def\csname sga@target@IV.4.6.16\endcsname{ega4.c04.prop.4-6-16}
\expandafter\def\csname sga@target@IV.4.6.16.a\endcsname{ega4.c04.sec.4-6-16-a}
\expandafter\def\csname sga@target@IV.4.6.16.b\endcsname{ega4.c04.sec.4-6-16-b}
\expandafter\def\csname sga@target@IV.4.6.16.c\endcsname{ega4.c04.sec.4-6-16-c}
\expandafter\def\csname sga@target@IV.4.6.16.d\endcsname{ega4.c04.sec.4-6-16-d}
\expandafter\def\csname sga@target@IV.4.6.17\endcsname{ega4.c04.def.4-6-17}
\expandafter\def\csname sga@target@IV.4.6.18\endcsname{ega4.c04.prop.4-6-18}
\expandafter\def\csname sga@target@IV.4.6.18.a\endcsname{ega4.c04.sec.4-6-18-a}
\expandafter\def\csname sga@target@IV.4.6.18.b\endcsname{ega4.c04.sec.4-6-18-b}
\expandafter\def\csname sga@target@IV.4.6.18.c\endcsname{ega4.c04.sec.4-6-18-c}
\expandafter\def\csname sga@target@IV.4.6.18.d\endcsname{ega4.c04.sec.4-6-18-d}
\expandafter\def\csname sga@target@IV.4.6.19\endcsname{ega4.c04.def.4-6-19}
\expandafter\def\csname sga@target@IV.4.6.2\endcsname{ega4.c04.def.4-6-2}
\expandafter\def\csname sga@target@IV.4.6.20\endcsname{ega4.c04.prop.4-6-20}
\expandafter\def\csname sga@target@IV.4.6.21\endcsname{ega4.c04.prop.4-6-21}
\expandafter\def\csname sga@target@IV.4.6.21.i\endcsname{ega4.c04.sec.4-6-21-i}
\expandafter\def\csname sga@target@IV.4.6.21.ii\endcsname{ega4.c04.sec.4-6-21-ii}
\expandafter\def\csname sga@target@IV.4.6.22\endcsname{ega4.c04.def.4-6-22}
\expandafter\def\csname sga@target@IV.4.6.3\endcsname{ega4.c04.cor.4-6-3}
\expandafter\def\csname sga@target@IV.4.6.4\endcsname{ega4.c04.cor.4-6-4}
\expandafter\def\csname sga@target@IV.4.6.5\endcsname{ega4.c04.prop.4-6-5}
\expandafter\def\csname sga@target@IV.4.6.5.i\endcsname{ega4.c04.sec.4-6-5-i}
\expandafter\def\csname sga@target@IV.4.6.5.ii\endcsname{ega4.c04.sec.4-6-5-ii}
\expandafter\def\csname sga@target@IV.4.6.6\endcsname{ega4.c04.prop.4-6-6}
\expandafter\def\csname sga@target@IV.4.6.7\endcsname{ega4.c04.cor.4-6-7}
\expandafter\def\csname sga@target@IV.4.6.8\endcsname{ega4.c04.cor.4-6-8}
\expandafter\def\csname sga@target@IV.4.6.9\endcsname{ega4.c04.def.4-6-9}
\expandafter\def\csname sga@target@IV.4.7.1\endcsname{ega4.c04.lem.4-7-1}
\expandafter\def\csname sga@target@IV.4.7.1.1\endcsname{ega4.c04.eq.4-7-1-1}
\expandafter\def\csname sga@target@IV.4.7.10\endcsname{ega4.c04.prop.4-7-10}
\expandafter\def\csname sga@target@IV.4.7.11\endcsname{ega4.c04.prop.4-7-11}
\expandafter\def\csname sga@target@IV.4.7.12\endcsname{ega4.c04.def.4-7-12}
\expandafter\def\csname sga@target@IV.4.7.13\endcsname{ega4.c04.prop.4-7-13}
\expandafter\def\csname sga@target@IV.4.7.2\endcsname{ega4.c04.cor.4-7-2}
\expandafter\def\csname sga@target@IV.4.7.2.1\endcsname{ega4.c04.eq.4-7-2-1}
\expandafter\def\csname sga@target@IV.4.7.3\endcsname{ega4.c04.prop.4-7-3}
\expandafter\def\csname sga@target@IV.4.7.3.1\endcsname{ega4.c04.eq.4-7-3-1}
\expandafter\def\csname sga@target@IV.4.7.3.i\endcsname{ega4.c04.sec.4-7-3-i}
\expandafter\def\csname sga@target@IV.4.7.3.ii\endcsname{ega4.c04.sec.4-7-3-ii}
\expandafter\def\csname sga@target@IV.4.7.3.iii\endcsname{ega4.c04.sec.4-7-3-iii}
\expandafter\def\csname sga@target@IV.4.7.4\endcsname{ega4.c04.def.4-7-4}
\expandafter\def\csname sga@target@IV.4.7.5\endcsname{ega4.c04.def.4-7-5}
\expandafter\def\csname sga@target@IV.4.7.6\endcsname{ega4.c04.prop.4-7-6}
\expandafter\def\csname sga@target@IV.4.7.6.1\endcsname{ega4.c04.eq.4-7-6-1}
\expandafter\def\csname sga@target@IV.4.7.7\endcsname{ega4.c04.statement.4-7-7}
\expandafter\def\csname sga@target@IV.4.7.8\endcsname{ega4.c04.prop.4-7-8}
\expandafter\def\csname sga@target@IV.4.7.9\endcsname{ega4.c04.cor.4-7-9}
\expandafter\def\csname sga@target@IV.4.8.1\endcsname{ega4.c04.statement.4-8-1}
\expandafter\def\csname sga@target@IV.4.8.10\endcsname{ega4.c04.cor.4-8-10}
\expandafter\def\csname sga@target@IV.4.8.11\endcsname{ega4.c04.cor.4-8-11}
\expandafter\def\csname sga@target@IV.4.8.12\endcsname{ega4.c04.cor.4-8-12}
\expandafter\def\csname sga@target@IV.4.8.13\endcsname{ega4.c04.prop.4-8-13}
\expandafter\def\csname sga@target@IV.4.8.14\endcsname{ega4.c04.prop.4-8-14}
\expandafter\def\csname sga@target@IV.4.8.15\endcsname{ega4.c04.sec.4-8-15}
\expandafter\def\csname sga@target@IV.4.8.15.i\endcsname{ega4.c04.sec.4-8-15-i}
\expandafter\def\csname sga@target@IV.4.8.15.ii\endcsname{ega4.c04.sec.4-8-15-ii}
\expandafter\def\csname sga@target@IV.4.8.2\endcsname{ega4.c04.statement.4-8-2}
\expandafter\def\csname sga@target@IV.4.8.2.1\endcsname{ega4.c04.eq.4-8-2-1}
\expandafter\def\csname sga@target@IV.4.8.2.2\endcsname{ega4.c04.eq.4-8-2-2}
\expandafter\def\csname sga@target@IV.4.8.2.3\endcsname{ega4.c04.eq.4-8-2-3}
\expandafter\def\csname sga@target@IV.4.8.2.4\endcsname{ega4.c04.eq.4-8-2-4}
\expandafter\def\csname sga@target@IV.4.8.2.5\endcsname{ega4.c04.eq.4-8-2-5}
\expandafter\def\csname sga@target@IV.4.8.2.6\endcsname{ega4.c04.eq.4-8-2-6}
\expandafter\def\csname sga@target@IV.4.8.3\endcsname{ega4.c04.prop.4-8-3}
\expandafter\def\csname sga@target@IV.4.8.4\endcsname{ega4.c04.def.4-8-4}
\expandafter\def\csname sga@target@IV.4.8.5\endcsname{ega4.c04.prop.4-8-5}
\expandafter\def\csname sga@target@IV.4.8.6\endcsname{ega4.c04.lem.4-8-6}
\expandafter\def\csname sga@target@IV.4.8.6.a\endcsname{ega4.c04.sec.4-8-6-a}
\expandafter\def\csname sga@target@IV.4.8.6.b\endcsname{ega4.c04.sec.4-8-6-b}
\expandafter\def\csname sga@target@IV.4.8.6.c\endcsname{ega4.c04.sec.4-8-6-c}
\expandafter\def\csname sga@target@IV.4.8.7\endcsname{ega4.c04.cor.4-8-7}
\expandafter\def\csname sga@target@IV.4.8.8\endcsname{ega4.c04.lem.4-8-8}
\expandafter\def\csname sga@target@IV.4.8.8.i\endcsname{ega4.c04.sec.4-8-8-i}
\expandafter\def\csname sga@target@IV.4.8.8.ii\endcsname{ega4.c04.sec.4-8-8-ii}
\expandafter\def\csname sga@target@IV.4.8.8.iii\endcsname{ega4.c04.sec.4-8-8-iii}
\expandafter\def\csname sga@target@IV.4.8.9\endcsname{ega4.c04.prop.4-8-9}
\expandafter\def\csname sga@target@IV.4.9.1\endcsname{ega4.c04.prop.4-9-1}
\expandafter\def\csname sga@target@IV.4.9.1.a\endcsname{ega4.c04.sec.4-9-1-a}
\expandafter\def\csname sga@target@IV.4.9.1.b\endcsname{ega4.c04.sec.4-9-1-b}
\expandafter\def\csname sga@target@IV.4.9.1.bprime\endcsname{ega4.c04.sec.4-9-1-bprime}
\expandafter\def\csname sga@target@IV.4.9.1.c\endcsname{ega4.c04.sec.4-9-1-c}
\expandafter\def\csname sga@target@IV.4.9.2\endcsname{ega4.c04.cor.4-9-2}
\expandafter\def\csname sga@target@IV.4.9.3\endcsname{ega4.c04.rem.4-9-3}
\expandafter\def\csname sga@target@IV.4.9.4\endcsname{ega4.c04.prop.4-9-4}
\expandafter\def\csname sga@target@IV.4.9.5\endcsname{ega4.c04.cor.4-9-5}
\expandafter\def\csname sga@target@IV.4.9.5.a\endcsname{ega4.c04.sec.4-9-5-a}
\expandafter\def\csname sga@target@IV.4.9.5.b\endcsname{ega4.c04.sec.4-9-5-b}
\expandafter\def\csname sga@target@IV.4.9.6\endcsname{ega4.c04.prop.4-9-6}
\expandafter\def\csname sga@target@IV.4.9.6.a\endcsname{ega4.c04.sec.4-9-6-a}
\expandafter\def\csname sga@target@IV.4.9.6.aprime\endcsname{ega4.c04.sec.4-9-6-aprime}
\expandafter\def\csname sga@target@IV.4.9.6.b\endcsname{ega4.c04.sec.4-9-6-b}
\expandafter\def\csname sga@target@IV.4.9.6.bprime\endcsname{ega4.c04.sec.4-9-6-bprime}
\expandafter\def\csname sga@target@IV.4.9.7\endcsname{ega4.c04.cor.4-9-7}
\expandafter\def\csname sga@target@IV.5.1.1\endcsname{ega4.c04.statement.5-1-1}
\expandafter\def\csname sga@target@IV.5.1.1.1\endcsname{ega4.c04.eq.5-1-1-1}
\expandafter\def\csname sga@target@IV.5.1.1.2\endcsname{ega4.c04.eq.5-1-1-2}
\expandafter\def\csname sga@target@IV.5.1.1.3\endcsname{ega4.c04.eq.5-1-1-3}
\expandafter\def\csname sga@target@IV.5.1.1.4\endcsname{ega4.c04.sec.5-1-1-4}
\expandafter\def\csname sga@target@IV.5.1.1.5\endcsname{ega4.c04.sec.5-1-1-5}
\expandafter\def\csname sga@target@IV.5.1.1.6\endcsname{ega4.c04.sec.5-1-1-6}
\expandafter\def\csname sga@target@IV.5.1.10\endcsname{ega4.c04.prop.5-1-10}
\expandafter\def\csname sga@target@IV.5.1.10.1\endcsname{ega4.c04.lem.5-1-10-1}
\expandafter\def\csname sga@target@IV.5.1.10.i\endcsname{ega4.c04.sec.5-1-10-i}
\expandafter\def\csname sga@target@IV.5.1.10.ii\endcsname{ega4.c04.sec.5-1-10-ii}
\expandafter\def\csname sga@target@IV.5.1.11\endcsname{ega4.c04.cor.5-1-11}
\expandafter\def\csname sga@target@IV.5.1.12\endcsname{ega4.c04.statement.5-1-12}
\expandafter\def\csname sga@target@IV.5.1.12.1\endcsname{ega4.c04.eq.5-1-12-1}
\expandafter\def\csname sga@target@IV.5.1.12.2\endcsname{ega4.c04.eq.5-1-12-2}
\expandafter\def\csname sga@target@IV.5.1.12.3\endcsname{ega4.c04.eq.5-1-12-3}
\expandafter\def\csname sga@target@IV.5.1.13\endcsname{ega4.c04.prop.5-1-13}
\expandafter\def\csname sga@target@IV.5.1.13.1\endcsname{ega4.c04.eq.5-1-13-1}
\expandafter\def\csname sga@target@IV.5.1.2\endcsname{ega4.c04.prop.5-1-2}
\expandafter\def\csname sga@target@IV.5.1.2.1\endcsname{ega4.c04.eq.5-1-2-1}
\expandafter\def\csname sga@target@IV.5.1.3\endcsname{ega4.c04.cor.5-1-3}
\expandafter\def\csname sga@target@IV.5.1.3.1\endcsname{ega4.c04.eq.5-1-3-1}
\expandafter\def\csname sga@target@IV.5.1.3.2\endcsname{ega4.c04.eq.5-1-3-2}
\expandafter\def\csname sga@target@IV.5.1.4\endcsname{ega4.c04.prop.5-1-4}
\expandafter\def\csname sga@target@IV.5.1.4.1\endcsname{ega4.c04.eq.5-1-4-1}
\expandafter\def\csname sga@target@IV.5.1.4.2\endcsname{ega4.c04.eq.5-1-4-2}
\expandafter\def\csname sga@target@IV.5.1.5\endcsname{ega4.c04.cor.5-1-5}
\expandafter\def\csname sga@target@IV.5.1.6\endcsname{ega4.c04.statement.5-1-6}
\expandafter\def\csname sga@target@IV.5.1.7\endcsname{ega4.c04.prop.5-1-7}
\expandafter\def\csname sga@target@IV.5.1.8\endcsname{ega4.c04.cor.5-1-8}
\expandafter\def\csname sga@target@IV.5.1.9\endcsname{ega4.c04.prop.5-1-9}
\expandafter\def\csname sga@target@IV.5.1.9.1\endcsname{ega4.c04.statement.5-1-9-1}
\expandafter\def\csname sga@target@IV.5.10.1\endcsname{ega4.c04.statement.5-10-1}
\expandafter\def\csname sga@target@IV.5.10.1.1\endcsname{ega4.c04.eq.5-10-1-1}
\expandafter\def\csname sga@target@IV.5.10.10\endcsname{ega4.c04.prop.5-10-10}
\expandafter\def\csname sga@target@IV.5.10.11\endcsname{ega4.c04.statement.5-10-11}
\expandafter\def\csname sga@target@IV.5.10.11.1\endcsname{ega4.c04.eq.5-10-11-1}
\expandafter\def\csname sga@target@IV.5.10.12\endcsname{ega4.c04.prop.5-10-12}
\expandafter\def\csname sga@target@IV.5.10.13\endcsname{ega4.c04.statement.5-10-13}
\expandafter\def\csname sga@target@IV.5.10.14\endcsname{ega4.c04.prop.5-10-14}
\expandafter\def\csname sga@target@IV.5.10.15\endcsname{ega4.c04.cor.5-10-15}
\expandafter\def\csname sga@target@IV.5.10.16\endcsname{ega4.c04.prop.5-10-16}
\expandafter\def\csname sga@target@IV.5.10.17\endcsname{ega4.c04.prop.5-10-17}
\expandafter\def\csname sga@target@IV.5.10.2\endcsname{ega4.c04.prop.5-10-2}
\expandafter\def\csname sga@target@IV.5.10.3\endcsname{ega4.c04.cor.5-10-3}
\expandafter\def\csname sga@target@IV.5.10.4\endcsname{ega4.c04.cor.5-10-4}
\expandafter\def\csname sga@target@IV.5.10.5\endcsname{ega4.c04.thm.5-10-5}
\expandafter\def\csname sga@target@IV.5.10.6\endcsname{ega4.c04.cor.5-10-6}
\expandafter\def\csname sga@target@IV.5.10.7\endcsname{ega4.c04.thm.5-10-7}
\expandafter\def\csname sga@target@IV.5.10.8\endcsname{ega4.c04.cor.5-10-8}
\expandafter\def\csname sga@target@IV.5.10.8.1\endcsname{ega4.c04.lem.5-10-8-1}
\expandafter\def\csname sga@target@IV.5.10.9\endcsname{ega4.c04.cor.5-10-9}
\expandafter\def\csname sga@target@IV.5.11.1\endcsname{ega4.c04.prop.5-11-1}
\expandafter\def\csname sga@target@IV.5.11.1.1\endcsname{ega4.c04.lem.5-11-1-1}
\expandafter\def\csname sga@target@IV.5.11.2\endcsname{ega4.c04.cor.5-11-2}
\expandafter\def\csname sga@target@IV.5.11.3\endcsname{ega4.c04.rem.5-11-3}
\expandafter\def\csname sga@target@IV.5.11.4\endcsname{ega4.c04.cor.5-11-4}
\expandafter\def\csname sga@target@IV.5.11.5\endcsname{ega4.c04.cor.5-11-5}
\expandafter\def\csname sga@target@IV.5.11.6\endcsname{ega4.c04.cor.5-11-6}
\expandafter\def\csname sga@target@IV.5.11.7\endcsname{ega4.c04.item.5-11-7}
\expandafter\def\csname sga@target@IV.5.12\endcsname{ega4.c04.subsec.5-12.key-iv-5-12}
\expandafter\def\csname sga@target@IV.5.12.1\endcsname{ega4.c04.prop.5-12-1}
\expandafter\def\csname sga@target@IV.5.12.1.1\endcsname{ega4.c04.eq.5-12-1-1}
\expandafter\def\csname sga@target@IV.5.12.10\endcsname{ega4.c04.cor.5-12-10}
\expandafter\def\csname sga@target@IV.5.12.2\endcsname{ega4.c04.prop.5-12-2}
\expandafter\def\csname sga@target@IV.5.12.3\endcsname{ega4.c04.rem.5-12-3}
\expandafter\def\csname sga@target@IV.5.12.4\endcsname{ega4.c04.cor.5-12-4}
\expandafter\def\csname sga@target@IV.5.12.5\endcsname{ega4.c04.prop.5-12-5}
\expandafter\def\csname sga@target@IV.5.12.6\endcsname{ega4.c04.rem.5-12-6}
\expandafter\def\csname sga@target@IV.5.12.7\endcsname{ega4.c04.cor.5-12-7}
\expandafter\def\csname sga@target@IV.5.12.8\endcsname{ega4.c04.prop.5-12-8}
\expandafter\def\csname sga@target@IV.5.12.9\endcsname{ega4.c04.rem.5-12-9}
\expandafter\def\csname sga@target@IV.5.13.1\endcsname{ega4.c04.prop.5-13-1}
\expandafter\def\csname sga@target@IV.5.13.2\endcsname{ega4.c04.cor.5-13-2}
\expandafter\def\csname sga@target@IV.5.13.3\endcsname{ega4.c04.prop.5-13-3}
\expandafter\def\csname sga@target@IV.5.13.4\endcsname{ega4.c04.cor.5-13-4}
\expandafter\def\csname sga@target@IV.5.13.5\endcsname{ega4.c04.statement.5-13-5}
\expandafter\def\csname sga@target@IV.5.13.6\endcsname{ega4.c04.cor.5-13-6}
\expandafter\def\csname sga@target@IV.5.13.7\endcsname{ega4.c04.prop.5-13-7}
\expandafter\def\csname sga@target@IV.5.13.7.1\endcsname{ega4.c04.lem.5-13-7-1}
\expandafter\def\csname sga@target@IV.5.13.8\endcsname{ega4.c04.cor.5-13-8}
\expandafter\def\csname sga@target@IV.5.2.1\endcsname{ega4.c04.prop.5-2-1}
\expandafter\def\csname sga@target@IV.5.2.1.1\endcsname{ega4.c04.eq.5-2-1-1}
\expandafter\def\csname sga@target@IV.5.2.2\endcsname{ega4.c04.cor.5-2-2}
\expandafter\def\csname sga@target@IV.5.2.3\endcsname{ega4.c04.cor.5-2-3}
\expandafter\def\csname sga@target@IV.5.2.3.1\endcsname{ega4.c04.eq.5-2-3-1}
\expandafter\def\csname sga@target@IV.5.2.4\endcsname{ega4.c04.cor.5-2-4}
\expandafter\def\csname sga@target@IV.5.2.5\endcsname{ega4.c04.sec.5-2-5}
\expandafter\def\csname sga@target@IV.5.2.5.i\endcsname{ega4.c04.sec.5-2-5-i}
\expandafter\def\csname sga@target@IV.5.2.5.ii\endcsname{ega4.c04.sec.5-2-5-ii}
\expandafter\def\csname sga@target@IV.5.3.1\endcsname{ega4.c04.prop.5-3-1}
\expandafter\def\csname sga@target@IV.5.3.1.1\endcsname{ega4.c04.eq.5-3-1-1}
\expandafter\def\csname sga@target@IV.5.3.1.2\endcsname{ega4.c04.eq.5-3-1-2}
\expandafter\def\csname sga@target@IV.5.3.1.3\endcsname{ega4.c04.lem.5-3-1-3}
\expandafter\def\csname sga@target@IV.5.4.1\endcsname{ega4.c04.prop.5-4-1}
\expandafter\def\csname sga@target@IV.5.4.1.i\endcsname{ega4.c04.sec.5-4-1-i}
\expandafter\def\csname sga@target@IV.5.4.1.ii\endcsname{ega4.c04.sec.5-4-1-ii}
\expandafter\def\csname sga@target@IV.5.4.2\endcsname{ega4.c04.cor.5-4-2}
\expandafter\def\csname sga@target@IV.5.4.3\endcsname{ega4.c04.sec.5-4-3}
\expandafter\def\csname sga@target@IV.5.4.3.1\endcsname{ega4.c04.eq.5-4-3-1}
\expandafter\def\csname sga@target@IV.5.4.3.i\endcsname{ega4.c04.sec.5-4-3-i}
\expandafter\def\csname sga@target@IV.5.4.3.ii\endcsname{ega4.c04.sec.5-4-3-ii}
\expandafter\def\csname sga@target@IV.5.5\endcsname{ega4.c04.sec.5-5}
\expandafter\def\csname sga@target@IV.5.5.1\endcsname{ega4.c04.statement.5-5-1}
\expandafter\def\csname sga@target@IV.5.5.1.1\endcsname{ega4.c04.eq.5-5-1-1}
\expandafter\def\csname sga@target@IV.5.5.2\endcsname{ega4.c04.prop.5-5-2}
\expandafter\def\csname sga@target@IV.5.5.2.1\endcsname{ega4.c04.eq.5-5-2-1}
\expandafter\def\csname sga@target@IV.5.5.2.2\endcsname{ega4.c04.eq.5-5-2-2}
\expandafter\def\csname sga@target@IV.5.5.3\endcsname{ega4.c04.prop.5-5-3}
\expandafter\def\csname sga@target@IV.5.5.3.1\endcsname{ega4.c04.eq.5-5-3-1}
\expandafter\def\csname sga@target@IV.5.5.4\endcsname{ega4.c04.cor.5-5-4}
\expandafter\def\csname sga@target@IV.5.5.4.1\endcsname{ega4.c04.eq.5-5-4-1}
\expandafter\def\csname sga@target@IV.5.5.5\endcsname{ega4.c04.cor.5-5-5}
\expandafter\def\csname sga@target@IV.5.5.6\endcsname{ega4.c04.cor.5-5-6}
\expandafter\def\csname sga@target@IV.5.5.6.1\endcsname{ega4.c04.eq.5-5-6-1}
\expandafter\def\csname sga@target@IV.5.5.7\endcsname{ega4.c04.lem.5-5-7}
\expandafter\def\csname sga@target@IV.5.5.7.1\endcsname{ega4.c04.eq.5-5-7-1}
\expandafter\def\csname sga@target@IV.5.5.7.2\endcsname{ega4.c04.statement.5-5-7-2}
\expandafter\def\csname sga@target@IV.5.5.7.3\endcsname{ega4.c04.eq.5-5-7-3}
\expandafter\def\csname sga@target@IV.5.5.7.i\endcsname{ega4.c04.sec.5-5-7-i}
\expandafter\def\csname sga@target@IV.5.5.7.ii\endcsname{ega4.c04.sec.5-5-7-ii}
\expandafter\def\csname sga@target@IV.5.5.8\endcsname{ega4.c04.thm.5-5-8}
\expandafter\def\csname sga@target@IV.5.5.8.1\endcsname{ega4.c04.eq.5-5-8-1}
\expandafter\def\csname sga@target@IV.5.5.8.2\endcsname{ega4.c04.eq.5-5-8-2}
\expandafter\def\csname sga@target@IV.5.5.8.3\endcsname{ega4.c04.eq.5-5-8-3}
\expandafter\def\csname sga@target@IV.5.6\endcsname{ega4.c04.sec.5-6}
\expandafter\def\csname sga@target@IV.5.6.1\endcsname{ega4.c04.prop.5-6-1}
\expandafter\def\csname sga@target@IV.5.6.1.1\endcsname{ega4.c04.eq.5-6-1-1}
\expandafter\def\csname sga@target@IV.5.6.1.2\endcsname{ega4.c04.eq.5-6-1-2}
\expandafter\def\csname sga@target@IV.5.6.1.3\endcsname{ega4.c04.statement.5-6-1-3}
\expandafter\def\csname sga@target@IV.5.6.1.4\endcsname{ega4.c04.statement.5-6-1-4}
\expandafter\def\csname sga@target@IV.5.6.1.5\endcsname{ega4.c04.statement.5-6-1-5}
\expandafter\def\csname sga@target@IV.5.6.1.a\endcsname{ega4.c04.sec.5-6-1-a}
\expandafter\def\csname sga@target@IV.5.6.1.b\endcsname{ega4.c04.sec.5-6-1-b}
\expandafter\def\csname sga@target@IV.5.6.1.c\endcsname{ega4.c04.sec.5-6-1-c}
\expandafter\def\csname sga@target@IV.5.6.10\endcsname{ega4.c04.prop.5-6-10}
\expandafter\def\csname sga@target@IV.5.6.11\endcsname{ega4.c04.ex.5-6-11}
\expandafter\def\csname sga@target@IV.5.6.11.1\endcsname{ega4.c04.lem.5-6-11-1}
\expandafter\def\csname sga@target@IV.5.6.11.1.i\endcsname{ega4.c04.sec.5-6-11-1-i}
\expandafter\def\csname sga@target@IV.5.6.11.1.ii\endcsname{ega4.c04.sec.5-6-11-1-ii}
\expandafter\def\csname sga@target@IV.5.6.2\endcsname{ega4.c04.def.5-6-2}
\expandafter\def\csname sga@target@IV.5.6.3\endcsname{ega4.c04.sec.5-6-3}
\expandafter\def\csname sga@target@IV.5.6.3.i\endcsname{ega4.c04.sec.5-6-3-i}
\expandafter\def\csname sga@target@IV.5.6.3.ii\endcsname{ega4.c04.sec.5-6-3-ii}
\expandafter\def\csname sga@target@IV.5.6.3.iii\endcsname{ega4.c04.sec.5-6-3-iii}
\expandafter\def\csname sga@target@IV.5.6.3.iv\endcsname{ega4.c04.sec.5-6-3-iv}
\expandafter\def\csname sga@target@IV.5.6.4\endcsname{ega4.c04.prop.5-6-4}
\expandafter\def\csname sga@target@IV.5.6.5\endcsname{ega4.c04.prop.5-6-5}
\expandafter\def\csname sga@target@IV.5.6.5.1\endcsname{ega4.c04.eq.5-6-5-1}
\expandafter\def\csname sga@target@IV.5.6.5.2\endcsname{ega4.c04.eq.5-6-5-2}
\expandafter\def\csname sga@target@IV.5.6.5.3\endcsname{ega4.c04.eq.5-6-5-3}
\expandafter\def\csname sga@target@IV.5.6.6\endcsname{ega4.c04.cor.5-6-6}
\expandafter\def\csname sga@target@IV.5.6.6.1\endcsname{ega4.c04.eq.5-6-6-1}
\expandafter\def\csname sga@target@IV.5.6.6.2\endcsname{ega4.c04.eq.5-6-6-2}
\expandafter\def\csname sga@target@IV.5.6.7\endcsname{ega4.c04.cor.5-6-7}
\expandafter\def\csname sga@target@IV.5.6.7.1\endcsname{ega4.c04.eq.5-6-7-1}
\expandafter\def\csname sga@target@IV.5.6.8\endcsname{ega4.c04.cor.5-6-8}
\expandafter\def\csname sga@target@IV.5.6.8.1\endcsname{ega4.c04.eq.5-6-8-1}
\expandafter\def\csname sga@target@IV.5.6.8.2\endcsname{ega4.c04.eq.5-6-8-2}
\expandafter\def\csname sga@target@IV.5.6.9\endcsname{ega4.c04.sec.5-6-9}
\expandafter\def\csname sga@target@IV.5.6.9.i\endcsname{ega4.c04.sec.5-6-9-i}
\expandafter\def\csname sga@target@IV.5.6.9.ii\endcsname{ega4.c04.sec.5-6-9-ii}
\expandafter\def\csname sga@target@IV.5.7.1\endcsname{ega4.c04.def.5-7-1}
\expandafter\def\csname sga@target@IV.5.7.1.1\endcsname{ega4.c04.eq.5-7-1-1}
\expandafter\def\csname sga@target@IV.5.7.10\endcsname{ega4.c04.cor.5-7-10}
\expandafter\def\csname sga@target@IV.5.7.11\endcsname{ega4.c04.cor.5-7-11}
\expandafter\def\csname sga@target@IV.5.7.11.i\endcsname{ega4.c04.sec.5-7-11-i}
\expandafter\def\csname sga@target@IV.5.7.11.ii\endcsname{ega4.c04.sec.5-7-11-ii}
\expandafter\def\csname sga@target@IV.5.7.12\endcsname{ega4.c04.statement.5-7-12}
\expandafter\def\csname sga@target@IV.5.7.13\endcsname{ega4.c04.cor.5-7-13}
\expandafter\def\csname sga@target@IV.5.7.13.1\endcsname{ega4.c04.eq.5-7-13-1}
\expandafter\def\csname sga@target@IV.5.7.2\endcsname{ega4.c04.def.5-7-2}
\expandafter\def\csname sga@target@IV.5.7.2.1\endcsname{ega4.c04.eq.5-7-2-1}
\expandafter\def\csname sga@target@IV.5.7.2.2\endcsname{ega4.c04.eq.5-7-2-2}
\expandafter\def\csname sga@target@IV.5.7.3\endcsname{ega4.c04.sec.5-7-3}
\expandafter\def\csname sga@target@IV.5.7.3.i\endcsname{ega4.c04.sec.5-7-3-i}
\expandafter\def\csname sga@target@IV.5.7.3.ii\endcsname{ega4.c04.sec.5-7-3-ii}
\expandafter\def\csname sga@target@IV.5.7.3.iii\endcsname{ega4.c04.sec.5-7-3-iii}
\expandafter\def\csname sga@target@IV.5.7.3.iv\endcsname{ega4.c04.sec.5-7-3-iv}
\expandafter\def\csname sga@target@IV.5.7.3.v\endcsname{ega4.c04.sec.5-7-3-v}
\expandafter\def\csname sga@target@IV.5.7.3.vi\endcsname{ega4.c04.sec.5-7-3-vi}
\expandafter\def\csname sga@target@IV.5.7.4\endcsname{ega4.c04.prop.5-7-4}
\expandafter\def\csname sga@target@IV.5.7.4.1\endcsname{ega4.c04.eq.5-7-4-1}
\expandafter\def\csname sga@target@IV.5.7.4.2\endcsname{ega4.c04.eq.5-7-4-2}
\expandafter\def\csname sga@target@IV.5.7.4.i\endcsname{ega4.c04.sec.5-7-4-i}
\expandafter\def\csname sga@target@IV.5.7.4.ii\endcsname{ega4.c04.sec.5-7-4-ii}
\expandafter\def\csname sga@target@IV.5.7.5\endcsname{ega4.c04.prop.5-7-5}
\expandafter\def\csname sga@target@IV.5.7.6\endcsname{ega4.c04.cor.5-7-6}
\expandafter\def\csname sga@target@IV.5.7.7\endcsname{ega4.c04.cor.5-7-7}
\expandafter\def\csname sga@target@IV.5.7.8\endcsname{ega4.c04.sec.5-7-8}
\expandafter\def\csname sga@target@IV.5.7.9\endcsname{ega4.c04.prop.5-7-9}
\expandafter\def\csname sga@target@IV.5.7.9.1\endcsname{ega4.c04.eq.5-7-9-1}
\expandafter\def\csname sga@target@IV.5.7.9.2\endcsname{ega4.c04.eq.5-7-9-2}
\expandafter\def\csname sga@target@IV.5.8.1\endcsname{ega4.c04.statement.5-8-1}
\expandafter\def\csname sga@target@IV.5.8.2\endcsname{ega4.c04.def.5-8-2}
\expandafter\def\csname sga@target@IV.5.8.3\endcsname{ega4.c04.prop.5-8-3}
\expandafter\def\csname sga@target@IV.5.8.4\endcsname{ega4.c04.statement.5-8-4}
\expandafter\def\csname sga@target@IV.5.8.5\endcsname{ega4.c04.prop.5-8-5}
\expandafter\def\csname sga@target@IV.5.8.6\endcsname{ega4.c04.thm.5-8-6}
\expandafter\def\csname sga@target@IV.5.8.6.1\endcsname{ega4.c04.eq.5-8-6-1}
\expandafter\def\csname sga@target@IV.5.8.6.i\endcsname{ega4.c04.item.5-8-6-i}
\expandafter\def\csname sga@target@IV.5.8.6.ii\endcsname{ega4.c04.item.5-8-6-ii}
\expandafter\def\csname sga@target@IV.5.9.1\endcsname{ega4.c04.statement.5-9-1}
\expandafter\def\csname sga@target@IV.5.9.1.1\endcsname{ega4.c04.eq.5-9-1-1}
\expandafter\def\csname sga@target@IV.5.9.1.2\endcsname{ega4.c04.eq.5-9-1-2}
\expandafter\def\csname sga@target@IV.5.9.1.3\endcsname{ega4.c04.eq.5-9-1-3}
\expandafter\def\csname sga@target@IV.5.9.10\endcsname{ega4.c04.cor.5-9-10}
\expandafter\def\csname sga@target@IV.5.9.10.i\endcsname{ega4.c04.item.5-9-10-i}
\expandafter\def\csname sga@target@IV.5.9.10.ii\endcsname{ega4.c04.item.5-9-10-ii}
\expandafter\def\csname sga@target@IV.5.9.11\endcsname{ega4.c04.prop.5-9-11}
\expandafter\def\csname sga@target@IV.5.9.12\endcsname{ega4.c04.item.5-9-12}
\expandafter\def\csname sga@target@IV.5.9.12.i\endcsname{ega4.c04.item.5-9-12-i}
\expandafter\def\csname sga@target@IV.5.9.12.ii\endcsname{ega4.c04.item.5-9-12-ii}
\expandafter\def\csname sga@target@IV.5.9.12.iii\endcsname{ega4.c04.item.5-9-12-iii}
\expandafter\def\csname sga@target@IV.5.9.2\endcsname{ega4.c04.prop.5-9-2}
\expandafter\def\csname sga@target@IV.5.9.2.i\endcsname{ega4.c04.item.5-9-2-i}
\expandafter\def\csname sga@target@IV.5.9.2.ii\endcsname{ega4.c04.item.5-9-2-ii}
\expandafter\def\csname sga@target@IV.5.9.3\endcsname{ega4.c04.rem.5-9-3}
\expandafter\def\csname sga@target@IV.5.9.3.1\endcsname{ega4.c04.eq.5-9-3-1}
\expandafter\def\csname sga@target@IV.5.9.4\endcsname{ega4.c04.prop.5-9-4}
\expandafter\def\csname sga@target@IV.5.9.4.1\endcsname{ega4.c04.eq.5-9-4-1}
\expandafter\def\csname sga@target@IV.5.9.5\endcsname{ega4.c04.cor.5-9-5}
\expandafter\def\csname sga@target@IV.5.9.6\endcsname{ega4.c04.cor.5-9-6}
\expandafter\def\csname sga@target@IV.5.9.6.1\endcsname{ega4.c04.eq.5-9-6-1}
\expandafter\def\csname sga@target@IV.5.9.7\endcsname{ega4.c04.statement.5-9-7}
\expandafter\def\csname sga@target@IV.5.9.7.1\endcsname{ega4.c04.eq.5-9-7-1}
\expandafter\def\csname sga@target@IV.5.9.8\endcsname{ega4.c04.prop.5-9-8}
\expandafter\def\csname sga@target@IV.5.9.8.1\endcsname{ega4.c04.eq.5-9-8-1}
\expandafter\def\csname sga@target@IV.5.9.8.a\endcsname{ega4.c04.item.5-9-8-a}
\expandafter\def\csname sga@target@IV.5.9.8.a-prime\endcsname{ega4.c04.item.5-9-8-a-prime}
\expandafter\def\csname sga@target@IV.5.9.8.b\endcsname{ega4.c04.item.5-9-8-b}
\expandafter\def\csname sga@target@IV.5.9.8.b-prime\endcsname{ega4.c04.item.5-9-8-b-prime}
\expandafter\def\csname sga@target@IV.5.9.9\endcsname{ega4.c04.def.5-9-9}
\expandafter\def\csname sga@target@IV.6.1.1\endcsname{ega4.c04.prop.6-1-1}
\expandafter\def\csname sga@target@IV.6.1.1.1\endcsname{ega4.c04.eq.6-1-1-1}
\expandafter\def\csname sga@target@IV.6.1.2\endcsname{ega4.c04.cor.6-1-2}
\expandafter\def\csname sga@target@IV.6.1.2.1\endcsname{ega4.c04.eq.6-1-2-1}
\expandafter\def\csname sga@target@IV.6.1.3\endcsname{ega4.c04.cor.6-1-3}
\expandafter\def\csname sga@target@IV.6.1.4\endcsname{ega4.c04.cor.6-1-4}
\expandafter\def\csname sga@target@IV.6.1.4.1\endcsname{ega4.c04.eq.6-1-4-1}
\expandafter\def\csname sga@target@IV.6.1.5\endcsname{ega4.c04.prop.6-1-5}
\expandafter\def\csname sga@target@IV.6.10.1\endcsname{ega4.c04.statement.6-10-1}
\expandafter\def\csname sga@target@IV.6.10.2\endcsname{ega4.c04.prop.6-10-2}
\expandafter\def\csname sga@target@IV.6.10.3\endcsname{ega4.c04.prop.6-10-3}
\expandafter\def\csname sga@target@IV.6.10.3.1\endcsname{ega4.c04.eq.6-10-3-1}
\expandafter\def\csname sga@target@IV.6.10.3.2\endcsname{ega4.c04.eq.6-10-3-2}
\expandafter\def\csname sga@target@IV.6.10.3.3\endcsname{ega4.c04.eq.6-10-3-3}
\expandafter\def\csname sga@target@IV.6.10.4\endcsname{ega4.c04.cor.6-10-4}
\expandafter\def\csname sga@target@IV.6.10.5\endcsname{ega4.c04.prop.6-10-5}
\expandafter\def\csname sga@target@IV.6.10.5.1\endcsname{ega4.c04.eq.6-10-5-1}
\expandafter\def\csname sga@target@IV.6.10.5.2\endcsname{ega4.c04.eq.6-10-5-2}
\expandafter\def\csname sga@target@IV.6.10.5.3\endcsname{ega4.c04.eq.6-10-5-3}
\expandafter\def\csname sga@target@IV.6.10.5.4\endcsname{ega4.c04.eq.6-10-5-4}
\expandafter\def\csname sga@target@IV.6.10.6\endcsname{ega4.c04.prop.6-10-6}
\expandafter\def\csname sga@target@IV.6.10.6.1\endcsname{ega4.c04.eq.6-10-6-1}
\expandafter\def\csname sga@target@IV.6.10.6.2\endcsname{ega4.c04.eq.6-10-6-2}
\expandafter\def\csname sga@target@IV.6.10.6.3\endcsname{ega4.c04.eq.6-10-6-3}
\expandafter\def\csname sga@target@IV.6.11\endcsname{ega4.c04.subsec.6-11.key-iv-6-11}
\expandafter\def\csname sga@target@IV.6.11.1\endcsname{ega4.c04.lem.6-11-1}
\expandafter\def\csname sga@target@IV.6.11.2\endcsname{ega4.c04.prop.6-11-2}
\expandafter\def\csname sga@target@IV.6.11.2.1\endcsname{ega4.c04.eq.6-11-2-1}
\expandafter\def\csname sga@target@IV.6.11.2.2\endcsname{ega4.c04.eq.6-11-2-2}
\expandafter\def\csname sga@target@IV.6.11.3\endcsname{ega4.c04.cor.6-11-3}
\expandafter\def\csname sga@target@IV.6.11.4\endcsname{ega4.c04.statement.6-11-4}
\expandafter\def\csname sga@target@IV.6.11.5\endcsname{ega4.c04.prop.6-11-5}
\expandafter\def\csname sga@target@IV.6.11.5.1\endcsname{ega4.c04.eq.6-11-5-1}
\expandafter\def\csname sga@target@IV.6.11.6\endcsname{ega4.c04.prop.6-11-6}
\expandafter\def\csname sga@target@IV.6.11.6.1\endcsname{ega4.c04.lem.6-11-6-1}
\expandafter\def\csname sga@target@IV.6.11.7\endcsname{ega4.c04.cor.6-11-7}
\expandafter\def\csname sga@target@IV.6.11.8\endcsname{ega4.c04.prop.6-11-8}
\expandafter\def\csname sga@target@IV.6.11.8.1\endcsname{ega4.c04.eq.6-11-8-1}
\expandafter\def\csname sga@target@IV.6.11.9\endcsname{ega4.c04.statement.6-11-9}
\expandafter\def\csname sga@target@IV.6.12\endcsname{ega4.c04.subsec.6-12.key-iv-6-12}
\expandafter\def\csname sga@target@IV.6.12.1\endcsname{ega4.c04.statement.6-12-1}
\expandafter\def\csname sga@target@IV.6.12.10\endcsname{ega4.c04.prop.6-12-10}
\expandafter\def\csname sga@target@IV.6.12.2\endcsname{ega4.c04.prop.6-12-2}
\expandafter\def\csname sga@target@IV.6.12.3\endcsname{ega4.c04.cor.6-12-3}
\expandafter\def\csname sga@target@IV.6.12.4\endcsname{ega4.c04.thm.6-12-4}
\expandafter\def\csname sga@target@IV.6.12.4.1\endcsname{ega4.c04.lem.6-12-4-1}
\expandafter\def\csname sga@target@IV.6.12.5\endcsname{ega4.c04.cor.6-12-5}
\expandafter\def\csname sga@target@IV.6.12.6\endcsname{ega4.c04.cor.6-12-6}
\expandafter\def\csname sga@target@IV.6.12.7\endcsname{ega4.c04.thm.6-12-7}
\expandafter\def\csname sga@target@IV.6.12.8\endcsname{ega4.c04.cor.6-12-8}
\expandafter\def\csname sga@target@IV.6.12.9\endcsname{ega4.c04.prop.6-12-9}
\expandafter\def\csname sga@target@IV.6.13\endcsname{ega4.c04.subsec.6-13.key-iv-6-13}
\expandafter\def\csname sga@target@IV.6.13.1\endcsname{ega4.c04.statement.6-13-1}
\expandafter\def\csname sga@target@IV.6.13.2\endcsname{ega4.c04.prop.6-13-2}
\expandafter\def\csname sga@target@IV.6.13.3\endcsname{ega4.c04.cor.6-13-3}
\expandafter\def\csname sga@target@IV.6.13.4\endcsname{ega4.c04.prop.6-13-4}
\expandafter\def\csname sga@target@IV.6.13.5\endcsname{ega4.c04.cor.6-13-5}
\expandafter\def\csname sga@target@IV.6.13.6\endcsname{ega4.c04.prop.6-13-6}
\expandafter\def\csname sga@target@IV.6.13.7\endcsname{ega4.c04.prop.6-13-7}
\expandafter\def\csname sga@target@IV.6.13.7.1\endcsname{ega4.c04.lem.6-13-7-1}
\expandafter\def\csname sga@target@IV.6.14.1\endcsname{ega4.c04.prop.6-14-1}
\expandafter\def\csname sga@target@IV.6.14.1.1\endcsname{ega4.c04.lem.6-14-1-1}
\expandafter\def\csname sga@target@IV.6.14.1.2\endcsname{ega4.c04.lem.6-14-1-2}
\expandafter\def\csname sga@target@IV.6.14.1.2.env\endcsname{ega4.c04.statement.6-14-1-2-env}
\expandafter\def\csname sga@target@IV.6.14.1.3\endcsname{ega4.c04.eq.6-14-1-3}
\expandafter\def\csname sga@target@IV.6.14.1.4\endcsname{ega4.c04.lem.6-14-1-4}
\expandafter\def\csname sga@target@IV.6.14.2\endcsname{ega4.c04.cor.6-14-2}
\expandafter\def\csname sga@target@IV.6.14.3\endcsname{ega4.c04.cor.6-14-3}
\expandafter\def\csname sga@target@IV.6.14.4\endcsname{ega4.c04.prop.6-14-4}
\expandafter\def\csname sga@target@IV.6.14.4.1\endcsname{ega4.c04.lem.6-14-4-1}
\expandafter\def\csname sga@target@IV.6.14.5\endcsname{ega4.c04.cor.6-14-5}
\expandafter\def\csname sga@target@IV.6.15.1\endcsname{ega4.c04.statement.6-15-1}
\expandafter\def\csname sga@target@IV.6.15.10\endcsname{ega4.c04.prop.6-15-10}
\expandafter\def\csname sga@target@IV.6.15.11\endcsname{ega4.c04.item.6-15-11}
\expandafter\def\csname sga@target@IV.6.15.2\endcsname{ega4.c04.statement.6-15-2}
\expandafter\def\csname sga@target@IV.6.15.3\endcsname{ega4.c04.statement.6-15-3}
\expandafter\def\csname sga@target@IV.6.15.3.1\endcsname{ega4.c04.lem.6-15-3-1}
\expandafter\def\csname sga@target@IV.6.15.4\endcsname{ega4.c04.statement.6-15-4}
\expandafter\def\csname sga@target@IV.6.15.4.1\endcsname{ega4.c04.lem.6-15-4-1}
\expandafter\def\csname sga@target@IV.6.15.5\endcsname{ega4.c04.prop.6-15-5}
\expandafter\def\csname sga@target@IV.6.15.6\endcsname{ega4.c04.prop.6-15-6}
\expandafter\def\csname sga@target@IV.6.15.7\endcsname{ega4.c04.prop.6-15-7}
\expandafter\def\csname sga@target@IV.6.15.8\endcsname{ega4.c04.lem.6-15-8}
\expandafter\def\csname sga@target@IV.6.15.9\endcsname{ega4.c04.cor.6-15-9}
\expandafter\def\csname sga@target@IV.6.2.1\endcsname{ega4.c04.prop.6-2-1}
\expandafter\def\csname sga@target@IV.6.2.1.1\endcsname{ega4.c04.eq.6-2-1-1}
\expandafter\def\csname sga@target@IV.6.2.1.2\endcsname{ega4.c04.eq.6-2-1-2}
\expandafter\def\csname sga@target@IV.6.2.2\endcsname{ega4.c04.cor.6-2-2}
\expandafter\def\csname sga@target@IV.6.2.3\endcsname{ega4.c04.prop.6-2-3}
\expandafter\def\csname sga@target@IV.6.2.3.1\endcsname{ega4.c04.eq.6-2-3-1}
\expandafter\def\csname sga@target@IV.6.2.4\endcsname{ega4.c04.prop.6-2-4}
\expandafter\def\csname sga@target@IV.6.2.4.1\endcsname{ega4.c04.eq.6-2-4-1}
\expandafter\def\csname sga@target@IV.6.3.1\endcsname{ega4.c04.prop.6-3-1}
\expandafter\def\csname sga@target@IV.6.3.1.1\endcsname{ega4.c04.eq.6-3-1-1}
\expandafter\def\csname sga@target@IV.6.3.2\endcsname{ega4.c04.cor.6-3-2}
\expandafter\def\csname sga@target@IV.6.3.2.1\endcsname{ega4.c04.eq.6-3-2-1}
\expandafter\def\csname sga@target@IV.6.3.3\endcsname{ega4.c04.cor.6-3-3}
\expandafter\def\csname sga@target@IV.6.3.4\endcsname{ega4.c04.cor.6-3-4}
\expandafter\def\csname sga@target@IV.6.3.4.1\endcsname{ega4.c04.eq.6-3-4-1}
\expandafter\def\csname sga@target@IV.6.3.4.2\endcsname{ega4.c04.eq.6-3-4-2}
\expandafter\def\csname sga@target@IV.6.3.5\endcsname{ega4.c04.cor.6-3-5}
\expandafter\def\csname sga@target@IV.6.3.6\endcsname{ega4.c04.cor.6-3-6}
\expandafter\def\csname sga@target@IV.6.3.7\endcsname{ega4.c04.cor.6-3-7}
\expandafter\def\csname sga@target@IV.6.3.8\endcsname{ega4.c04.prop.6-3-8}
\expandafter\def\csname sga@target@IV.6.4.1\endcsname{ega4.c04.prop.6-4-1}
\expandafter\def\csname sga@target@IV.6.4.2\endcsname{ega4.c04.cor.6-4-2}
\expandafter\def\csname sga@target@IV.6.4.3\endcsname{ega4.c04.rem.6-4-3}
\expandafter\def\csname sga@target@IV.6.5.1\endcsname{ega4.c04.prop.6-5-1}
\expandafter\def\csname sga@target@IV.6.5.2\endcsname{ega4.c04.cor.6-5-2}
\expandafter\def\csname sga@target@IV.6.5.3\endcsname{ega4.c04.prop.6-5-3}
\expandafter\def\csname sga@target@IV.6.5.4\endcsname{ega4.c04.cor.6-5-4}
\expandafter\def\csname sga@target@IV.6.5.5\endcsname{ega4.c04.item.6-5-5}
\expandafter\def\csname sga@target@IV.6.6.1\endcsname{ega4.c04.prop.6-6-1}
\expandafter\def\csname sga@target@IV.6.6.2\endcsname{ega4.c04.rem.6-6-2}
\expandafter\def\csname sga@target@IV.6.7.1\endcsname{ega4.c04.prop.6-7-1}
\expandafter\def\csname sga@target@IV.6.7.1.1\endcsname{ega4.c04.lem.6-7-1-1}
\expandafter\def\csname sga@target@IV.6.7.2\endcsname{ega4.c04.cor.6-7-2}
\expandafter\def\csname sga@target@IV.6.7.3\endcsname{ega4.c04.cor.6-7-3}
\expandafter\def\csname sga@target@IV.6.7.4\endcsname{ega4.c04.prop.6-7-4}
\expandafter\def\csname sga@target@IV.6.7.4.1\endcsname{ega4.c04.lem.6-7-4-1}
\expandafter\def\csname sga@target@IV.6.7.5\endcsname{ega4.c04.item.6-7-5}
\expandafter\def\csname sga@target@IV.6.7.6\endcsname{ega4.c04.def.6-7-6}
\expandafter\def\csname sga@target@IV.6.7.7\endcsname{ega4.c04.prop.6-7-7}
\expandafter\def\csname sga@target@IV.6.7.8\endcsname{ega4.c04.cor.6-7-8}
\expandafter\def\csname sga@target@IV.6.8.1\endcsname{ega4.c04.def.6-8-1}
\expandafter\def\csname sga@target@IV.6.8.2\endcsname{ega4.c04.prop.6-8-2}
\expandafter\def\csname sga@target@IV.6.8.3\endcsname{ega4.c04.prop.6-8-3}
\expandafter\def\csname sga@target@IV.6.8.4\endcsname{ega4.c04.item.6-8-4}
\expandafter\def\csname sga@target@IV.6.8.5\endcsname{ega4.c04.prop.6-8-5}
\expandafter\def\csname sga@target@IV.6.8.6\endcsname{ega4.c04.thm.6-8-6}
\expandafter\def\csname sga@target@IV.6.8.7\endcsname{ega4.c04.cor.6-8-7}
\expandafter\def\csname sga@target@IV.6.8.8\endcsname{ega4.c04.rem.6-8-8}
\expandafter\def\csname sga@target@IV.6.9.1\endcsname{ega4.c04.thm.6-9-1}
\expandafter\def\csname sga@target@IV.6.9.2\endcsname{ega4.c04.lem.6-9-2}
\expandafter\def\csname sga@target@IV.6.9.3\endcsname{ega4.c04.cor.6-9-3}
\expandafter\def\csname sga@target@IV.7.1.1\endcsname{ega4.c04.def.7-1-1}
\expandafter\def\csname sga@target@IV.7.1.10\endcsname{ega4.c04.prop.7-1-10}
\expandafter\def\csname sga@target@IV.7.1.11\endcsname{ega4.c04.prop.7-1-11}
\expandafter\def\csname sga@target@IV.7.1.12\endcsname{ega4.c04.item.7-1-12}
\expandafter\def\csname sga@target@IV.7.1.13\endcsname{ega4.c04.cor.7-1-13}
\expandafter\def\csname sga@target@IV.7.1.2\endcsname{ega4.c04.prop.7-1-2}
\expandafter\def\csname sga@target@IV.7.1.3\endcsname{ega4.c04.prop.7-1-3}
\expandafter\def\csname sga@target@IV.7.1.3.1\endcsname{ega4.c04.eq.7-1-3-1}
\expandafter\def\csname sga@target@IV.7.1.3.2\endcsname{ega4.c04.eq.7-1-3-2}
\expandafter\def\csname sga@target@IV.7.1.3.3\endcsname{ega4.c04.eq.7-1-3-3}
\expandafter\def\csname sga@target@IV.7.1.3.4\endcsname{ega4.c04.eq.7-1-3-4}
\expandafter\def\csname sga@target@IV.7.1.3.5\endcsname{ega4.c04.eq.7-1-3-5}
\expandafter\def\csname sga@target@IV.7.1.3.6\endcsname{ega4.c04.eq.7-1-3-6}
\expandafter\def\csname sga@target@IV.7.1.3.7\endcsname{ega4.c04.eq.7-1-3-7}
\expandafter\def\csname sga@target@IV.7.1.4\endcsname{ega4.c04.cor.7-1-4}
\expandafter\def\csname sga@target@IV.7.1.5\endcsname{ega4.c04.cor.7-1-5}
\expandafter\def\csname sga@target@IV.7.1.6\endcsname{ega4.c04.rem.7-1-6}
\expandafter\def\csname sga@target@IV.7.1.7\endcsname{ega4.c04.lem.7-1-7}
\expandafter\def\csname sga@target@IV.7.1.8\endcsname{ega4.c04.thm.7-1-8}
\expandafter\def\csname sga@target@IV.7.1.9\endcsname{ega4.c04.def.7-1-9}
\expandafter\def\csname sga@target@IV.7.2.1\endcsname{ega4.c04.statement.7-2-1}
\expandafter\def\csname sga@target@IV.7.2.1.1\endcsname{ega4.c04.statement.7-2-1-1}
\expandafter\def\csname sga@target@IV.7.2.10\endcsname{ega4.c04.statement.7-2-10}
\expandafter\def\csname sga@target@IV.7.2.2\endcsname{ega4.c04.prop.7-2-2}
\expandafter\def\csname sga@target@IV.7.2.3\endcsname{ega4.c04.prop.7-2-3}
\expandafter\def\csname sga@target@IV.7.2.4\endcsname{ega4.c04.prop.7-2-4}
\expandafter\def\csname sga@target@IV.7.2.5\endcsname{ega4.c04.thm.7-2-5}
\expandafter\def\csname sga@target@IV.7.2.5.1\endcsname{ega4.c04.eq.7-2-5-1}
\expandafter\def\csname sga@target@IV.7.2.6\endcsname{ega4.c04.def.7-2-6}
\expandafter\def\csname sga@target@IV.7.2.7\endcsname{ega4.c04.cor.7-2-7}
\expandafter\def\csname sga@target@IV.7.2.8\endcsname{ega4.c04.cor.7-2-8}
\expandafter\def\csname sga@target@IV.7.2.9\endcsname{ega4.c04.cor.7-2-9}
\expandafter\def\csname sga@target@IV.7.3.1\endcsname{ega4.c04.statement.7-3-1}
\expandafter\def\csname sga@target@IV.7.3.10\endcsname{ega4.c04.statement.7-3-10}
\expandafter\def\csname sga@target@IV.7.3.11\endcsname{ega4.c04.rem.7-3-11}
\expandafter\def\csname sga@target@IV.7.3.12\endcsname{ega4.c04.prop.7-3-12}
\expandafter\def\csname sga@target@IV.7.3.12.1\endcsname{ega4.c04.eq.7-3-12-1}
\expandafter\def\csname sga@target@IV.7.3.13\endcsname{ega4.c04.statement.7-3-13}
\expandafter\def\csname sga@target@IV.7.3.14\endcsname{ega4.c04.prop.7-3-14}
\expandafter\def\csname sga@target@IV.7.3.15\endcsname{ega4.c04.prop.7-3-15}
\expandafter\def\csname sga@target@IV.7.3.16\endcsname{ega4.c04.prop.7-3-16}
\expandafter\def\csname sga@target@IV.7.3.16.1\endcsname{ega4.c04.eq.7-3-16-1}
\expandafter\def\csname sga@target@IV.7.3.16.2\endcsname{ega4.c04.lem.7-3-16-2}
\expandafter\def\csname sga@target@IV.7.3.17\endcsname{ega4.c04.prop.7-3-17}
\expandafter\def\csname sga@target@IV.7.3.18\endcsname{ega4.c04.prop.7-3-18}
\expandafter\def\csname sga@target@IV.7.3.19\endcsname{ega4.c04.statement.7-3-19}
\expandafter\def\csname sga@target@IV.7.3.2\endcsname{ega4.c04.lem.7-3-2}
\expandafter\def\csname sga@target@IV.7.3.3\endcsname{ega4.c04.cor.7-3-3}
\expandafter\def\csname sga@target@IV.7.3.4\endcsname{ega4.c04.statement.7-3-4}
\expandafter\def\csname sga@target@IV.7.3.5\endcsname{ega4.c04.item.7-3-5}
\expandafter\def\csname sga@target@IV.7.3.6\endcsname{ega4.c04.statement.7-3-6}
\expandafter\def\csname sga@target@IV.7.3.7\endcsname{ega4.c04.lem.7-3-7}
\expandafter\def\csname sga@target@IV.7.3.8\endcsname{ega4.c04.statement.7-3-8}
\expandafter\def\csname sga@target@IV.7.3.9\endcsname{ega4.c04.item.7-3-9}
\expandafter\def\csname sga@target@IV.7.4.1\endcsname{ega4.c04.statement.7-4-1}
\expandafter\def\csname sga@target@IV.7.4.2\endcsname{ega4.c04.lem.7-4-2}
\expandafter\def\csname sga@target@IV.7.4.3\endcsname{ega4.c04.cor.7-4-3}
\expandafter\def\csname sga@target@IV.7.4.4\endcsname{ega4.c04.thm.7-4-4}
\expandafter\def\csname sga@target@IV.7.4.4.1\endcsname{ega4.c04.lem.7-4-4-1}
\expandafter\def\csname sga@target@IV.7.4.4.2\endcsname{ega4.c04.lem.7-4-4-2}
\expandafter\def\csname sga@target@IV.7.4.5\endcsname{ega4.c04.cor.7-4-5}
\expandafter\def\csname sga@target@IV.7.4.6\endcsname{ega4.c04.prop.7-4-6}
\expandafter\def\csname sga@target@IV.7.4.7\endcsname{ega4.c04.cor.7-4-7}
\expandafter\def\csname sga@target@IV.7.4.8\endcsname{ega4.c04.rem.7-4-8}
\expandafter\def\csname sga@target@IV.7.5.0\endcsname{ega4.c04.statement.7-5-0}
\expandafter\def\csname sga@target@IV.7.5.1\endcsname{ega4.c04.thm.7-5-1}
\expandafter\def\csname sga@target@IV.7.5.1.1\endcsname{ega4.c04.lem.7-5-1-1}
\expandafter\def\csname sga@target@IV.7.5.2\endcsname{ega4.c04.cor.7-5-2}
\expandafter\def\csname sga@target@IV.7.5.3\endcsname{ega4.c04.item.7-5-3}
\expandafter\def\csname sga@target@IV.7.5.4\endcsname{ega4.c04.item.7-5-4}
\expandafter\def\csname sga@target@IV.7.5.5\endcsname{ega4.c04.lem.7-5-5}
\expandafter\def\csname sga@target@IV.7.5.6\endcsname{ega4.c04.prop.7-5-6}
\expandafter\def\csname sga@target@IV.7.5.7\endcsname{ega4.c04.cor.7-5-7}
\expandafter\def\csname sga@target@IV.7.5.7.1\endcsname{ega4.c04.lem.7-5-7-1}
\expandafter\def\csname sga@target@IV.7.5.8\endcsname{ega4.c04.rem.7-5-8}
\expandafter\def\csname sga@target@IV.7.6.1\endcsname{ega4.c04.prop.7-6-1}
\expandafter\def\csname sga@target@IV.7.6.2\endcsname{ega4.c04.cor.7-6-2}
\expandafter\def\csname sga@target@IV.7.6.3\endcsname{ega4.c04.cor.7-6-3}
\expandafter\def\csname sga@target@IV.7.6.4\endcsname{ega4.c04.thm.7-6-4}
\expandafter\def\csname sga@target@IV.7.6.4.1\endcsname{ega4.c04.lem.7-6-4-1}
\expandafter\def\csname sga@target@IV.7.6.5\endcsname{ega4.c04.cor.7-6-5}
\expandafter\def\csname sga@target@IV.7.6.6\endcsname{ega4.c04.cor.7-6-6}
\expandafter\def\csname sga@target@IV.7.6.7\endcsname{ega4.c04.item.7-6-7}
\expandafter\def\csname sga@target@IV.7.7.1\endcsname{ega4.c04.statement.7-7-1}
\expandafter\def\csname sga@target@IV.7.7.2\endcsname{ega4.c04.thm.7-7-2}
\expandafter\def\csname sga@target@IV.7.7.3\endcsname{ega4.c04.cor.7-7-3}
\expandafter\def\csname sga@target@IV.7.7.4\endcsname{ega4.c04.cor.7-7-4}
\expandafter\def\csname sga@target@IV.7.8.1\endcsname{ega4.c04.statement.7-8-1}
\expandafter\def\csname sga@target@IV.7.8.2\endcsname{ega4.c04.def.7-8-2}
\expandafter\def\csname sga@target@IV.7.8.3\endcsname{ega4.c04.statement.7-8-3}
\expandafter\def\csname sga@target@IV.7.8.3.1\endcsname{ega4.c04.statement.7-8-3-1}
\expandafter\def\csname sga@target@IV.7.8.4\endcsname{ega4.c04.item.7-8-4}
\expandafter\def\csname sga@target@IV.7.8.4.1\endcsname{ega4.c04.statement.7-8-4-1}
\expandafter\def\csname sga@target@IV.7.8.5\endcsname{ega4.c04.statement.7-8-5}
\expandafter\def\csname sga@target@IV.7.8.6\endcsname{ega4.c04.prop.7-8-6}
\expandafter\def\csname sga@target@IV.7.9.1\endcsname{ega4.c04.statement.7-9-1}
\expandafter\def\csname sga@target@IV.7.9.10\endcsname{ega4.c04.item.7-9-10}
\expandafter\def\csname sga@target@IV.7.9.2\endcsname{ega4.c04.statement.7-9-2}
\expandafter\def\csname sga@target@IV.7.9.3\endcsname{ega4.c04.prop.7-9-3}
\expandafter\def\csname sga@target@IV.7.9.3.1\endcsname{ega4.c04.lem.7-9-3-1}
\expandafter\def\csname sga@target@IV.7.9.4\endcsname{ega4.c04.cor.7-9-4}
\expandafter\def\csname sga@target@IV.7.9.5\endcsname{ega4.c04.prop.7-9-5}
\expandafter\def\csname sga@target@IV.7.9.6\endcsname{ega4.c04.rem.7-9-6}
\expandafter\def\csname sga@target@IV.7.9.7\endcsname{ega4.c04.statement.7-9-7}
\expandafter\def\csname sga@target@IV.7.9.8\endcsname{ega4.c04.prop.7-9-8}
\expandafter\def\csname sga@target@IV.7.9.8.1\endcsname{ega4.c04.statement.7-9-8-1}
\expandafter\def\csname sga@target@IV.7.9.9\endcsname{ega4.c04.cor.7-9-9}
\expandafter\def\csname sga@target@IV.8.1.1\endcsname{ega4.c04.statement.8-1-1}
\expandafter\def\csname sga@target@IV.8.1.2\endcsname{ega4.c04.statement.8-1-2}
\expandafter\def\csname sga@target@IV.8.1.3\endcsname{ega4.c04.statement.8-1-3}
\expandafter\def\csname sga@target@IV.8.10.1\endcsname{ega4.c04.prop.8-10-1}
\expandafter\def\csname sga@target@IV.8.10.2\endcsname{ega4.c04.cor.8-10-2}
\expandafter\def\csname sga@target@IV.8.10.3\endcsname{ega4.c04.prop.8-10-3}
\expandafter\def\csname sga@target@IV.8.10.3.1\endcsname{ega4.c04.rem.8-10-3-1}
\expandafter\def\csname sga@target@IV.8.10.4\endcsname{ega4.c04.prop.8-10-4}
\expandafter\def\csname sga@target@IV.8.10.5\endcsname{ega4.c04.thm.8-10-5}
\expandafter\def\csname sga@target@IV.8.10.5.1\endcsname{ega4.c04.prop.8-10-5-1}
\expandafter\def\csname sga@target@IV.8.10.5.2\endcsname{ega4.c04.lem.8-10-5-2}
\expandafter\def\csname sga@target@IV.8.11\endcsname{ega4.c04.subsec.8-11.key-iv-8-11}
\expandafter\def\csname sga@target@IV.8.11.1\endcsname{ega4.c04.thm.8-11-1}
\expandafter\def\csname sga@target@IV.8.11.2\endcsname{ega4.c04.thm.8-11-2}
\expandafter\def\csname sga@target@IV.8.11.3\endcsname{ega4.c04.rem.8-11-3}
\expandafter\def\csname sga@target@IV.8.11.4\endcsname{ega4.c04.statement.8-11-4}
\expandafter\def\csname sga@target@IV.8.11.5\endcsname{ega4.c04.prop.8-11-5}
\expandafter\def\csname sga@target@IV.8.11.5.1\endcsname{ega4.c04.rem.8-11-5-1}
\expandafter\def\csname sga@target@IV.8.11.6\endcsname{ega4.c04.prop.8-11-6}
\expandafter\def\csname sga@target@IV.8.12\endcsname{ega4.c04.subsec.8-12.key-iv-8-12}
\expandafter\def\csname sga@target@IV.8.12.1\endcsname{ega4.c04.lem.8-12-1}
\expandafter\def\csname sga@target@IV.8.12.10\endcsname{ega4.c04.cor.8-12-10}
\expandafter\def\csname sga@target@IV.8.12.10.1\endcsname{ega4.c04.lem.8-12-10-1}
\expandafter\def\csname sga@target@IV.8.12.11\endcsname{ega4.c04.cor.8-12-11}
\expandafter\def\csname sga@target@IV.8.12.12\endcsname{ega4.c04.cor.8-12-12}
\expandafter\def\csname sga@target@IV.8.12.2\endcsname{ega4.c04.lem.8-12-2}
\expandafter\def\csname sga@target@IV.8.12.3\endcsname{ega4.c04.prop.8-12-3}
\expandafter\def\csname sga@target@IV.8.12.4\endcsname{ega4.c04.cor.8-12-4}
\expandafter\def\csname sga@target@IV.8.12.5\endcsname{ega4.c04.statement.8-12-5}
\expandafter\def\csname sga@target@IV.8.12.6\endcsname{ega4.c04.thm.8-12-6}
\expandafter\def\csname sga@target@IV.8.12.6.1\endcsname{ega4.c04.eq.8-12-6-1}
\expandafter\def\csname sga@target@IV.8.12.7\endcsname{ega4.c04.rem.8-12-7}
\expandafter\def\csname sga@target@IV.8.12.8\endcsname{ega4.c04.cor.8-12-8}
\expandafter\def\csname sga@target@IV.8.12.9\endcsname{ega4.c04.cor.8-12-9}
\expandafter\def\csname sga@target@IV.8.13.1\endcsname{ega4.c04.prop.8-13-1}
\expandafter\def\csname sga@target@IV.8.13.1.1\endcsname{ega4.c04.eq.8-13-1-1}
\expandafter\def\csname sga@target@IV.8.13.2\endcsname{ega4.c04.cor.8-13-2}
\expandafter\def\csname sga@target@IV.8.13.2.1\endcsname{ega4.c04.eq.8-13-2-1}
\expandafter\def\csname sga@target@IV.8.13.3\endcsname{ega4.c04.statement.8-13-3}
\expandafter\def\csname sga@target@IV.8.13.4\endcsname{ega4.c04.statement.8-13-4}
\expandafter\def\csname sga@target@IV.8.13.4.1\endcsname{ega4.c04.eq.8-13-4-1}
\expandafter\def\csname sga@target@IV.8.13.4.2\endcsname{ega4.c04.eq.8-13-4-2}
\expandafter\def\csname sga@target@IV.8.13.5\endcsname{ega4.c04.prop.8-13-5}
\expandafter\def\csname sga@target@IV.8.13.6\endcsname{ega4.c04.rem.8-13-6}
\expandafter\def\csname sga@target@IV.8.14.1\endcsname{ega4.c04.statement.8-14-1}
\expandafter\def\csname sga@target@IV.8.14.2\endcsname{ega4.c04.prop.8-14-2}
\expandafter\def\csname sga@target@IV.8.14.2.1\endcsname{ega4.c04.eq.8-14-2-1}
\expandafter\def\csname sga@target@IV.8.14.2.2\endcsname{ega4.c04.cor.8-14-2-2}
\expandafter\def\csname sga@target@IV.8.14.2.3\endcsname{ega4.c04.eq.8-14-2-3}
\expandafter\def\csname sga@target@IV.8.14.2.4\endcsname{ega4.c04.eq.8-14-2-4}
\expandafter\def\csname sga@target@IV.8.14.3\endcsname{ega4.c04.rem.8-14-3}
\expandafter\def\csname sga@target@IV.8.2.1\endcsname{ega4.c04.statement.8-2-1}
\expandafter\def\csname sga@target@IV.8.2.1.1\endcsname{ega4.c04.eq.8-2-1-1}
\expandafter\def\csname sga@target@IV.8.2.10\endcsname{ega4.c04.cor.8-2-10}
\expandafter\def\csname sga@target@IV.8.2.11\endcsname{ega4.c04.cor.8-2-11}
\expandafter\def\csname sga@target@IV.8.2.12\endcsname{ega4.c04.cor.8-2-12}
\expandafter\def\csname sga@target@IV.8.2.12.1\endcsname{ega4.c04.eq.8-2-12-1}
\expandafter\def\csname sga@target@IV.8.2.13\endcsname{ega4.c04.prop.8-2-13}
\expandafter\def\csname sga@target@IV.8.2.14\endcsname{ega4.c04.rem.8-2-14}
\expandafter\def\csname sga@target@IV.8.2.2\endcsname{ega4.c04.statement.8-2-2}
\expandafter\def\csname sga@target@IV.8.2.3\endcsname{ega4.c04.prop.8-2-3}
\expandafter\def\csname sga@target@IV.8.2.3.1\endcsname{ega4.c04.eq.8-2-3-1}
\expandafter\def\csname sga@target@IV.8.2.4\endcsname{ega4.c04.lem.8-2-4}
\expandafter\def\csname sga@target@IV.8.2.5\endcsname{ega4.c04.prop.8-2-5}
\expandafter\def\csname sga@target@IV.8.2.6\endcsname{ega4.c04.lem.8-2-6}
\expandafter\def\csname sga@target@IV.8.2.7\endcsname{ega4.c04.rem.8-2-7}
\expandafter\def\csname sga@target@IV.8.2.8\endcsname{ega4.c04.statement.8-2-8}
\expandafter\def\csname sga@target@IV.8.2.9\endcsname{ega4.c04.prop.8-2-9}
\expandafter\def\csname sga@target@IV.8.3.1\endcsname{ega4.c04.statement.8-3-1}
\expandafter\def\csname sga@target@IV.8.3.10\endcsname{ega4.c04.statement.8-3-10}
\expandafter\def\csname sga@target@IV.8.3.10.1\endcsname{ega4.c04.eq.8-3-10-1}
\expandafter\def\csname sga@target@IV.8.3.11\endcsname{ega4.c04.thm.8-3-11}
\expandafter\def\csname sga@target@IV.8.3.12\endcsname{ega4.c04.cor.8-3-12}
\expandafter\def\csname sga@target@IV.8.3.13\endcsname{ega4.c04.prop.8-3-13}
\expandafter\def\csname sga@target@IV.8.3.2\endcsname{ega4.c04.thm.8-3-2}
\expandafter\def\csname sga@target@IV.8.3.2.1\endcsname{ega4.c04.eq.8-3-2-1}
\expandafter\def\csname sga@target@IV.8.3.3\endcsname{ega4.c04.cor.8-3-3}
\expandafter\def\csname sga@target@IV.8.3.4\endcsname{ega4.c04.cor.8-3-4}
\expandafter\def\csname sga@target@IV.8.3.5\endcsname{ega4.c04.cor.8-3-5}
\expandafter\def\csname sga@target@IV.8.3.6\endcsname{ega4.c04.cor.8-3-6}
\expandafter\def\csname sga@target@IV.8.3.7\endcsname{ega4.c04.cor.8-3-7}
\expandafter\def\csname sga@target@IV.8.3.7.1\endcsname{ega4.c04.eq.8-3-7-1}
\expandafter\def\csname sga@target@IV.8.3.8\endcsname{ega4.c04.prop.8-3-8}
\expandafter\def\csname sga@target@IV.8.3.9\endcsname{ega4.c04.statement.8-3-9}
\expandafter\def\csname sga@target@IV.8.3.9.1\endcsname{ega4.c04.eq.8-3-9-1}
\expandafter\def\csname sga@target@IV.8.3.9.2\endcsname{ega4.c04.eq.8-3-9-2}
\expandafter\def\csname sga@target@IV.8.3.9.3\endcsname{ega4.c04.eq.8-3-9-3}
\expandafter\def\csname sga@target@IV.8.3.9.4\endcsname{ega4.c04.eq.8-3-9-4}
\expandafter\def\csname sga@target@IV.8.3.9.5\endcsname{ega4.c04.eq.8-3-9-5}
\expandafter\def\csname sga@target@IV.8.4.1\endcsname{ega4.c04.prop.8-4-1}
\expandafter\def\csname sga@target@IV.8.4.2\endcsname{ega4.c04.prop.8-4-2}
\expandafter\def\csname sga@target@IV.8.4.2.1\endcsname{ega4.c04.lem.8-4-2-1}
\expandafter\def\csname sga@target@IV.8.4.3\endcsname{ega4.c04.cor.8-4-3}
\expandafter\def\csname sga@target@IV.8.4.4\endcsname{ega4.c04.prop.8-4-4}
\expandafter\def\csname sga@target@IV.8.4.5\endcsname{ega4.c04.cor.8-4-5}
\expandafter\def\csname sga@target@IV.8.4.6\endcsname{ega4.c04.rem.8-4-6}
\expandafter\def\csname sga@target@IV.8.5.1\endcsname{ega4.c04.statement.8-5-1}
\expandafter\def\csname sga@target@IV.8.5.1.1\endcsname{ega4.c04.eq.8-5-1-1}
\expandafter\def\csname sga@target@IV.8.5.1.2\endcsname{ega4.c04.eq.8-5-1-2}
\expandafter\def\csname sga@target@IV.8.5.1.3\endcsname{ega4.c04.eq.8-5-1-3}
\expandafter\def\csname sga@target@IV.8.5.1.4\endcsname{ega4.c04.eq.8-5-1-4}
\expandafter\def\csname sga@target@IV.8.5.10\endcsname{ega4.c04.statement.8-5-10}
\expandafter\def\csname sga@target@IV.8.5.10.1\endcsname{ega4.c04.eq.8-5-10-1}
\expandafter\def\csname sga@target@IV.8.5.10.2\endcsname{ega4.c04.eq.8-5-10-2}
\expandafter\def\csname sga@target@IV.8.5.11\endcsname{ega4.c04.prop.8-5-11}
\expandafter\def\csname sga@target@IV.8.5.2\endcsname{ega4.c04.thm.8-5-2}
\expandafter\def\csname sga@target@IV.8.5.2.1\endcsname{ega4.c04.lem.8-5-2-1}
\expandafter\def\csname sga@target@IV.8.5.2.2\endcsname{ega4.c04.eq.8-5-2-2}
\expandafter\def\csname sga@target@IV.8.5.2.3\endcsname{ega4.c04.eq.8-5-2-3}
\expandafter\def\csname sga@target@IV.8.5.2.4\endcsname{ega4.c04.cor.8-5-2-4}
\expandafter\def\csname sga@target@IV.8.5.2.5\endcsname{ega4.c04.cor.8-5-2-5}
\expandafter\def\csname sga@target@IV.8.5.3\endcsname{ega4.c04.statement.8-5-3}
\expandafter\def\csname sga@target@IV.8.5.4\endcsname{ega4.c04.cor.8-5-4}
\expandafter\def\csname sga@target@IV.8.5.5\endcsname{ega4.c04.prop.8-5-5}
\expandafter\def\csname sga@target@IV.8.5.6\endcsname{ega4.c04.prop.8-5-6}
\expandafter\def\csname sga@target@IV.8.5.7\endcsname{ega4.c04.cor.8-5-7}
\expandafter\def\csname sga@target@IV.8.5.8\endcsname{ega4.c04.cor.8-5-8}
\expandafter\def\csname sga@target@IV.8.5.9\endcsname{ega4.c04.lem.8-5-9}
\expandafter\def\csname sga@target@IV.8.6.1\endcsname{ega4.c04.statement.8-6-1}
\expandafter\def\csname sga@target@IV.8.6.2\endcsname{ega4.c04.statement.8-6-2}
\expandafter\def\csname sga@target@IV.8.6.2.1\endcsname{ega4.c04.eq.8-6-2-1}
\expandafter\def\csname sga@target@IV.8.6.2.2\endcsname{ega4.c04.eq.8-6-2-2}
\expandafter\def\csname sga@target@IV.8.6.2.3\endcsname{ega4.c04.eq.8-6-2-3}
\expandafter\def\csname sga@target@IV.8.6.3\endcsname{ega4.c04.prop.8-6-3}
\expandafter\def\csname sga@target@IV.8.6.3.1\endcsname{ega4.c04.cor.8-6-3-1}
\expandafter\def\csname sga@target@IV.8.6.4\endcsname{ega4.c04.cor.8-6-4}
\expandafter\def\csname sga@target@IV.8.7.1\endcsname{ega4.c04.prop.8-7-1}
\expandafter\def\csname sga@target@IV.8.7.2\endcsname{ega4.c04.prop.8-7-2}
\expandafter\def\csname sga@target@IV.8.7.3\endcsname{ega4.c04.cor.8-7-3}
\expandafter\def\csname sga@target@IV.8.7.4\endcsname{ega4.c04.rem.8-7-4}
\expandafter\def\csname sga@target@IV.8.8.1\endcsname{ega4.c04.statement.8-8-1}
\expandafter\def\csname sga@target@IV.8.8.1.1\endcsname{ega4.c04.eq.8-8-1-1}
\expandafter\def\csname sga@target@IV.8.8.2\endcsname{ega4.c04.thm.8-8-2}
\expandafter\def\csname sga@target@IV.8.8.2.1\endcsname{ega4.c04.lem.8-8-2-1}
\expandafter\def\csname sga@target@IV.8.8.2.2\endcsname{ega4.c04.lem.8-8-2-2}
\expandafter\def\csname sga@target@IV.8.8.2.3\endcsname{ega4.c04.lem.8-8-2-3}
\expandafter\def\csname sga@target@IV.8.8.2.4\endcsname{ega4.c04.cor.8-8-2-4}
\expandafter\def\csname sga@target@IV.8.8.2.5\endcsname{ega4.c04.cor.8-8-2-5}
\expandafter\def\csname sga@target@IV.8.8.3\endcsname{ega4.c04.statement.8-8-3}
\expandafter\def\csname sga@target@IV.8.9.1\endcsname{ega4.c04.prop.8-9-1}
\expandafter\def\csname sga@target@IV.8.9.2\endcsname{ega4.c04.statement.8-9-2}
\expandafter\def\csname sga@target@IV.8.9.3\endcsname{ega4.c04.prop.8-9-3}
\expandafter\def\csname sga@target@IV.8.9.4\endcsname{ega4.c04.prop.8-9-4}
\expandafter\def\csname sga@target@IV.8.9.4.1\endcsname{ega4.c04.lem.8-9-4-1}
\expandafter\def\csname sga@target@IV.8.9.5\endcsname{ega4.c04.cor.8-9-5}
\expandafter\def\csname sga@target@IV.9.1.1\endcsname{ega4.c04.prop.9-1-1}
\expandafter\def\csname sga@target@IV.9.1.2\endcsname{ega4.c04.cor.9-1-2}
\expandafter\def\csname sga@target@IV.9.1.3\endcsname{ega4.c04.rem.9-1-3}
\expandafter\def\csname sga@target@IV.9.1.4\endcsname{ega4.c04.prop.9-1-4}
\expandafter\def\csname sga@target@IV.9.1.5\endcsname{ega4.c04.cor.9-1-5}
\expandafter\def\csname sga@target@IV.9.1.6\endcsname{ega4.c04.prop.9-1-6}
\expandafter\def\csname sga@target@IV.9.2\endcsname{ega4.c04.subsec.9-2.key-iv-9-2}
\expandafter\def\csname sga@target@IV.9.2.1\endcsname{ega4.c04.def.9-2-1}
\expandafter\def\csname sga@target@IV.9.2.2\endcsname{ega4.c04.item.9-2-2}
\expandafter\def\csname sga@target@IV.9.2.3\endcsname{ega4.c04.prop.9-2-3}
\expandafter\def\csname sga@target@IV.9.2.4\endcsname{ega4.c04.cor.9-2-4}
\expandafter\def\csname sga@target@IV.9.2.5\endcsname{ega4.c04.rem.9-2-5}
\expandafter\def\csname sga@target@IV.9.2.6\endcsname{ega4.c04.prop.9-2-6}
\expandafter\def\csname sga@target@IV.9.2.6.1\endcsname{ega4.c04.prop.9-2-6-1}
\expandafter\def\csname sga@target@IV.9.2.6.2\endcsname{ega4.c04.cor.9-2-6-2}
\expandafter\def\csname sga@target@IV.9.3.1\endcsname{ega4.c04.prop.9-3-1}
\expandafter\def\csname sga@target@IV.9.3.2\endcsname{ega4.c04.prop.9-3-2}
\expandafter\def\csname sga@target@IV.9.3.3\endcsname{ega4.c04.prop.9-3-3}
\expandafter\def\csname sga@target@IV.9.3.4\endcsname{ega4.c04.rem.9-3-4}
\expandafter\def\csname sga@target@IV.9.3.5\endcsname{ega4.c04.prop.9-3-5}
\expandafter\def\csname sga@target@IV.9.3.6\endcsname{ega4.c04.statement.9-3-6}
\expandafter\def\csname sga@target@IV.9.4.1\endcsname{ega4.c04.statement.9-4-1}
\expandafter\def\csname sga@target@IV.9.4.2\endcsname{ega4.c04.prop.9-4-2}
\expandafter\def\csname sga@target@IV.9.4.3\endcsname{ega4.c04.cor.9-4-3}
\expandafter\def\csname sga@target@IV.9.4.3.1\endcsname{ega4.c04.eq.9-4-3-1}
\expandafter\def\csname sga@target@IV.9.4.4\endcsname{ega4.c04.prop.9-4-4}
\expandafter\def\csname sga@target@IV.9.4.5\endcsname{ega4.c04.cor.9-4-5}
\expandafter\def\csname sga@target@IV.9.4.6\endcsname{ega4.c04.cor.9-4-6}
\expandafter\def\csname sga@target@IV.9.4.7\endcsname{ega4.c04.prop.9-4-7}
\expandafter\def\csname sga@target@IV.9.4.7.1\endcsname{ega4.c04.lem.9-4-7-1}
\expandafter\def\csname sga@target@IV.9.4.7.2\endcsname{ega4.c04.rem.9-4-7-2}
\expandafter\def\csname sga@target@IV.9.4.8\endcsname{ega4.c04.prop.9-4-8}
\expandafter\def\csname sga@target@IV.9.4.9\endcsname{ega4.c04.rem.9-4-9}
\expandafter\def\csname sga@target@IV.9.5.1\endcsname{ega4.c04.prop.9-5-1}
\expandafter\def\csname sga@target@IV.9.5.2\endcsname{ega4.c04.cor.9-5-2}
\expandafter\def\csname sga@target@IV.9.5.3\endcsname{ega4.c04.prop.9-5-3}
\expandafter\def\csname sga@target@IV.9.5.3.case2\endcsname{ega4.c04.item.9-5-3-case2}
\expandafter\def\csname sga@target@IV.9.5.4\endcsname{ega4.c04.cor.9-5-4}
\expandafter\def\csname sga@target@IV.9.5.5\endcsname{ega4.c04.prop.9-5-5}
\expandafter\def\csname sga@target@IV.9.5.6\endcsname{ega4.c04.cor.9-5-6}
\expandafter\def\csname sga@target@IV.9.5.7\endcsname{ega4.c04.rem.9-5-7}
\expandafter\def\csname sga@target@IV.9.6.1\endcsname{ega4.c04.prop.9-6-1}
\expandafter\def\csname sga@target@IV.9.6.2\endcsname{ega4.c04.prop.9-6-2}
\expandafter\def\csname sga@target@IV.9.6.3\endcsname{ega4.c04.prop.9-6-3}
\expandafter\def\csname sga@target@IV.9.6.3.1\endcsname{ega4.c04.eq.9-6-3-1}
\expandafter\def\csname sga@target@IV.9.6.4\endcsname{ega4.c04.cor.9-6-4}
\expandafter\def\csname sga@target@IV.9.6.5\endcsname{ega4.c04.cor.9-6-5}
\expandafter\def\csname sga@target@IV.9.6.6\endcsname{ega4.c04.sec.9-6-6}
\expandafter\def\csname sga@target@IV.9.6.7\endcsname{ega4.c04.sec.9-6-7}
\expandafter\def\csname sga@target@IV.9.7\endcsname{ega4.c04.subsec.9-7.key-iv-9-7}
\expandafter\def\csname sga@target@IV.9.7.1\endcsname{ega4.c04.lem.9-7-1}
\expandafter\def\csname sga@target@IV.9.7.10\endcsname{ega4.c04.cor.9-7-10}
\expandafter\def\csname sga@target@IV.9.7.11\endcsname{ega4.c04.rem.9-7-11}
\expandafter\def\csname sga@target@IV.9.7.12\endcsname{ega4.c04.prop.9-7-12}
\expandafter\def\csname sga@target@IV.9.7.2\endcsname{ega4.c04.lem.9-7-2}
\expandafter\def\csname sga@target@IV.9.7.3\endcsname{ega4.c04.subsec.9-7-3}
\expandafter\def\csname sga@target@IV.9.7.4\endcsname{ega4.c04.lem.9-7-4}
\expandafter\def\csname sga@target@IV.9.7.5\endcsname{ega4.c04.lem.9-7-5}
\expandafter\def\csname sga@target@IV.9.7.5.1\endcsname{ega4.c04.lem.9-7-5-1}
\expandafter\def\csname sga@target@IV.9.7.6\endcsname{ega4.c04.prop.9-7-6}
\expandafter\def\csname sga@target@IV.9.7.7\endcsname{ega4.c04.thm.9-7-7}
\expandafter\def\csname sga@target@IV.9.7.8\endcsname{ega4.c04.prop.9-7-8}
\expandafter\def\csname sga@target@IV.9.7.9\endcsname{ega4.c04.cor.9-7-9}
\expandafter\def\csname sga@target@IV.9.8\endcsname{ega4.c04.subsec.9-8.key-iv-9-8}
\expandafter\def\csname sga@target@IV.9.8.1\endcsname{ega4.c04.prop.9-8-1}
\expandafter\def\csname sga@target@IV.9.8.1-proof\endcsname{ega4.c04.subsec.9-8-1-proof}
\expandafter\def\csname sga@target@IV.9.8.2\endcsname{ega4.c04.subsec.9-8-2}
\expandafter\def\csname sga@target@IV.9.8.3\endcsname{ega4.c04.prop.9-8-3}
\expandafter\def\csname sga@target@IV.9.8.3-proof\endcsname{ega4.c04.subsec.9-8-3-proof}
\expandafter\def\csname sga@target@IV.9.8.4\endcsname{ega4.c04.cor.9-8-4}
\expandafter\def\csname sga@target@IV.9.8.4-proof\endcsname{ega4.c04.item.9-8-4-proof}
\expandafter\def\csname sga@target@IV.9.8.5\endcsname{ega4.c04.prop.9-8-5}
\expandafter\def\csname sga@target@IV.9.8.5-proof\endcsname{ega4.c04.item.9-8-5-proof}
\expandafter\def\csname sga@target@IV.9.8.6\endcsname{ega4.c04.prop.9-8-6}
\expandafter\def\csname sga@target@IV.9.8.6-proof\endcsname{ega4.c04.item.9-8-6-proof}
\expandafter\def\csname sga@target@IV.9.8.7\endcsname{ega4.c04.cor.9-8-7}
\expandafter\def\csname sga@target@IV.9.8.7-proof\endcsname{ega4.c04.item.9-8-7-proof}
\expandafter\def\csname sga@target@IV.9.8.8\endcsname{ega4.c04.prop.9-8-8}
\expandafter\def\csname sga@target@IV.9.8.8-proof\endcsname{ega4.c04.item.9-8-8-proof}
\expandafter\def\csname sga@target@IV.9.8.9\endcsname{ega4.c04.rem.9-8-9}
\expandafter\def\csname sga@target@IV.9.8.9.1\endcsname{ega4.c04.item.9-8-9-1}
\expandafter\def\csname sga@target@IV.9.9\endcsname{ega4.c04.subsec.9-9.key-iv-9-9}
\expandafter\def\csname sga@target@IV.9.9.1\endcsname{ega4.c04.prop.9-9-1}
\expandafter\def\csname sga@target@IV.9.9.1-proof\endcsname{ega4.c04.item.9-9-1-proof}
\expandafter\def\csname sga@target@IV.9.9.1.1\endcsname{ega4.c04.eq.9-9-1-1}
\expandafter\def\csname sga@target@IV.9.9.2\endcsname{ega4.c04.prop.9-9-2}
\expandafter\def\csname sga@target@IV.9.9.2-proof\endcsname{ega4.c04.item.9-9-2-proof}
\expandafter\def\csname sga@target@IV.9.9.2.1\endcsname{ega4.c04.eq.9-9-2-1}
\expandafter\def\csname sga@target@IV.9.9.2.2\endcsname{ega4.c04.eq.9-9-2-2}
\expandafter\def\csname sga@target@IV.9.9.3\endcsname{ega4.c04.cor.9-9-3}
\expandafter\def\csname sga@target@IV.9.9.3-proof\endcsname{ega4.c04.statement.9-9-3-proof}
\expandafter\def\csname sga@target@IV.9.9.4\endcsname{ega4.c04.prop.9-9-4}
\expandafter\def\csname sga@target@IV.9.9.4-proof\endcsname{ega4.c04.item.9-9-4-proof}
\expandafter\def\csname sga@target@IV.9.9.5\endcsname{ega4.c04.cor.9-9-5}
\expandafter\def\csname sga@target@IV.9.9.5-proof\endcsname{ega4.c04.item.9-9-5-proof}
\expandafter\def\csname sga@target@IV.9.9.6\endcsname{ega4.c04.prop.9-9-6}
\expandafter\def\csname sga@target@IV.9.9.6-proof\endcsname{ega4.c04.item.9-9-6-proof}
\expandafter\def\csname sga@target@section:IV.1\endcsname{ega4.c04.sec.1}
\expandafter\def\csname sga@target@section:IV.10\endcsname{ega4.c04.sec.10}
\expandafter\def\csname sga@target@section:IV.11\endcsname{ega4.c04.sec.11}
\expandafter\def\csname sga@target@section:IV.12\endcsname{ega4.c04.sec.12}
\expandafter\def\csname sga@target@section:IV.13\endcsname{ega4.c04.sec.13}
\expandafter\def\csname sga@target@section:IV.14\endcsname{ega4.c04.sec.14}
\expandafter\def\csname sga@target@section:IV.15\endcsname{ega4.c04.sec.15}
\expandafter\def\csname sga@target@section:IV.16\endcsname{ega4.c04.sec.16}
\expandafter\def\csname sga@target@section:IV.17\endcsname{ega4.c04.sec.17}
\expandafter\def\csname sga@target@section:IV.18\endcsname{ega4.c04.sec.18}
\expandafter\def\csname sga@target@section:IV.19\endcsname{ega4.c04.sec.19}
\expandafter\def\csname sga@target@section:IV.2\endcsname{ega4.c04.sec.2}
\expandafter\def\csname sga@target@section:IV.21\endcsname{ega4.c04.sec.21}
\expandafter\def\csname sga@target@section:IV.3\endcsname{ega4.c04.sec.3}
\expandafter\def\csname sga@target@section:IV.4\endcsname{ega4.c04.sec.4}
\expandafter\def\csname sga@target@section:IV.5\endcsname{ega4.c04.sec.5}
\expandafter\def\csname sga@target@section:IV.6\endcsname{ega4.c04.sec.6}
\expandafter\def\csname sga@target@section:IV.7\endcsname{ega4.c04.sec.7}
\expandafter\def\csname sga@target@section:IV.8\endcsname{ega4.c04.sec.8}
\expandafter\def\csname sga@target@section:IV.9\endcsname{ega4.c04.sec.9}
\expandafter\def\csname sga@target@subsection:IV.1.1\endcsname{ega4.c04.subsec.1-1}
\expandafter\def\csname sga@target@subsection:IV.1.10\endcsname{ega4.c04.subsec.1-10}
\expandafter\def\csname sga@target@subsection:IV.1.2\endcsname{ega4.c04.subsec.1-2}
\expandafter\def\csname sga@target@subsection:IV.1.3\endcsname{ega4.c04.subsec.1-3}
\expandafter\def\csname sga@target@subsection:IV.1.4\endcsname{ega4.c04.subsec.1-4}
\expandafter\def\csname sga@target@subsection:IV.1.5\endcsname{ega4.c04.subsec.1-5}
\expandafter\def\csname sga@target@subsection:IV.1.6\endcsname{ega4.c04.subsec.1-6}
\expandafter\def\csname sga@target@subsection:IV.1.7\endcsname{ega4.c04.subsec.1-7}
\expandafter\def\csname sga@target@subsection:IV.1.8\endcsname{ega4.c04.subsec.1-8}
\expandafter\def\csname sga@target@subsection:IV.1.9\endcsname{ega4.c04.subsec.1-9}
\expandafter\def\csname sga@target@subsection:IV.10.1\endcsname{ega4.c04.subsec.10-1}
\expandafter\def\csname sga@target@subsection:IV.10.10\endcsname{ega4.c04.subsec.10-10}
\expandafter\def\csname sga@target@subsection:IV.10.2\endcsname{ega4.c04.subsec.10-2}
\expandafter\def\csname sga@target@subsection:IV.10.3\endcsname{ega4.c04.subsec.10-3}
\expandafter\def\csname sga@target@subsection:IV.10.4\endcsname{ega4.c04.subsec.10-4}
\expandafter\def\csname sga@target@subsection:IV.10.5\endcsname{ega4.c04.subsec.10-5}
\expandafter\def\csname sga@target@subsection:IV.10.6\endcsname{ega4.c04.subsec.10-6}
\expandafter\def\csname sga@target@subsection:IV.10.7\endcsname{ega4.c04.subsec.10-7}
\expandafter\def\csname sga@target@subsection:IV.10.8\endcsname{ega4.c04.subsec.10-8}
\expandafter\def\csname sga@target@subsection:IV.10.9\endcsname{ega4.c04.subsec.10-9}
\expandafter\def\csname sga@target@subsection:IV.11.1\endcsname{ega4.c04.subsec.11-1}
\expandafter\def\csname sga@target@subsection:IV.11.10\endcsname{ega4.c04.subsec.11-10}
\expandafter\def\csname sga@target@subsection:IV.11.2\endcsname{ega4.c04.subsec.11-2}
\expandafter\def\csname sga@target@subsection:IV.11.3\endcsname{ega4.c04.subsec.11-3}
\expandafter\def\csname sga@target@subsection:IV.11.4\endcsname{ega4.c04.subsec.11-4}
\expandafter\def\csname sga@target@subsection:IV.11.5\endcsname{ega4.c04.subsec.11-5}
\expandafter\def\csname sga@target@subsection:IV.11.6\endcsname{ega4.c04.subsec.11-6}
\expandafter\def\csname sga@target@subsection:IV.11.7\endcsname{ega4.c04.subsec.11-7}
\expandafter\def\csname sga@target@subsection:IV.11.8\endcsname{ega4.c04.subsec.11-8}
\expandafter\def\csname sga@target@subsection:IV.11.9\endcsname{ega4.c04.subsec.11-9}
\expandafter\def\csname sga@target@subsection:IV.12.0\endcsname{ega4.c04.subsec.12-0}
\expandafter\def\csname sga@target@subsection:IV.12.1\endcsname{ega4.c04.subsec.12-1}
\expandafter\def\csname sga@target@subsection:IV.12.2\endcsname{ega4.c04.subsec.12-2}
\expandafter\def\csname sga@target@subsection:IV.12.3\endcsname{ega4.c04.subsec.12-3}
\expandafter\def\csname sga@target@subsection:IV.13.1\endcsname{ega4.c04.subsec.13-1}
\expandafter\def\csname sga@target@subsection:IV.13.2\endcsname{ega4.c04.subsec.13-2}
\expandafter\def\csname sga@target@subsection:IV.13.3\endcsname{ega4.c04.subsec.13-3}
\expandafter\def\csname sga@target@subsection:IV.14.1\endcsname{ega4.c04.subsec.14-1}
\expandafter\def\csname sga@target@subsection:IV.14.2\endcsname{ega4.c04.subsec.14-2}
\expandafter\def\csname sga@target@subsection:IV.14.3\endcsname{ega4.c04.subsec.14-3}
\expandafter\def\csname sga@target@subsection:IV.14.4\endcsname{ega4.c04.subsec.14-4}
\expandafter\def\csname sga@target@subsection:IV.14.5\endcsname{ega4.c04.subsec.14-5}
\expandafter\def\csname sga@target@subsection:IV.15.1\endcsname{ega4.c04.subsec.15-1}
\expandafter\def\csname sga@target@subsection:IV.15.2\endcsname{ega4.c04.subsec.15-2}
\expandafter\def\csname sga@target@subsection:IV.15.3\endcsname{ega4.c04.subsec.15-3}
\expandafter\def\csname sga@target@subsection:IV.15.4\endcsname{ega4.c04.subsec.15-4}
\expandafter\def\csname sga@target@subsection:IV.15.5\endcsname{ega4.c04.subsec.15-5}
\expandafter\def\csname sga@target@subsection:IV.15.6\endcsname{ega4.c04.subsec.15-6}
\expandafter\def\csname sga@target@subsection:IV.15.7\endcsname{ega4.c04.subsec.15-7}
\expandafter\def\csname sga@target@subsection:IV.17.1\endcsname{ega4.c04.subsec.17-1}
\expandafter\def\csname sga@target@subsection:IV.17.10\endcsname{ega4.c04.subsec.17-10}
\expandafter\def\csname sga@target@subsection:IV.17.11\endcsname{ega4.c04.subsec.17-11}
\expandafter\def\csname sga@target@subsection:IV.17.12\endcsname{ega4.c04.subsec.17-12}
\expandafter\def\csname sga@target@subsection:IV.17.13\endcsname{ega4.c04.subsec.17-13}
\expandafter\def\csname sga@target@subsection:IV.17.14\endcsname{ega4.c04.subsec.17-14}
\expandafter\def\csname sga@target@subsection:IV.17.15\endcsname{ega4.c04.subsec.17-15}
\expandafter\def\csname sga@target@subsection:IV.17.16\endcsname{ega4.c04.subsec.17-16}
\expandafter\def\csname sga@target@subsection:IV.17.2\endcsname{ega4.c04.subsec.17-2}
\expandafter\def\csname sga@target@subsection:IV.17.3\endcsname{ega4.c04.subsec.17-3}
\expandafter\def\csname sga@target@subsection:IV.17.4\endcsname{ega4.c04.subsec.17-4}
\expandafter\def\csname sga@target@subsection:IV.17.5\endcsname{ega4.c04.subsec.17-5}
\expandafter\def\csname sga@target@subsection:IV.17.6\endcsname{ega4.c04.subsec.17-6}
\expandafter\def\csname sga@target@subsection:IV.17.7\endcsname{ega4.c04.subsec.17-7}
\expandafter\def\csname sga@target@subsection:IV.17.8\endcsname{ega4.c04.subsec.17-8}
\expandafter\def\csname sga@target@subsection:IV.17.9\endcsname{ega4.c04.subsec.17-9}
\expandafter\def\csname sga@target@subsection:IV.18.1\endcsname{ega4.c04.subsec.18-1}
\expandafter\def\csname sga@target@subsection:IV.18.10\endcsname{ega4.c04.subsec.18-10}
\expandafter\def\csname sga@target@subsection:IV.18.11\endcsname{ega4.c04.subsec.18-11}
\expandafter\def\csname sga@target@subsection:IV.18.12\endcsname{ega4.c04.subsec.18-12}
\expandafter\def\csname sga@target@subsection:IV.18.2\endcsname{ega4.c04.subsec.18-2}
\expandafter\def\csname sga@target@subsection:IV.18.3\endcsname{ega4.c04.subsec.18-3}
\expandafter\def\csname sga@target@subsection:IV.18.4\endcsname{ega4.c04.subsec.18-4}
\expandafter\def\csname sga@target@subsection:IV.18.5\endcsname{ega4.c04.subsec.18-5}
\expandafter\def\csname sga@target@subsection:IV.18.6\endcsname{ega4.c04.subsec.18-6}
\expandafter\def\csname sga@target@subsection:IV.18.7\endcsname{ega4.c04.subsec.18-7}
\expandafter\def\csname sga@target@subsection:IV.18.8\endcsname{ega4.c04.subsec.18-8}
\expandafter\def\csname sga@target@subsection:IV.18.9\endcsname{ega4.c04.subsec.18-9}
\expandafter\def\csname sga@target@subsection:IV.19.1\endcsname{ega4.c04.subsec.19-1}
\expandafter\def\csname sga@target@subsection:IV.19.2\endcsname{ega4.c04.subsec.19-2}
\expandafter\def\csname sga@target@subsection:IV.19.3\endcsname{ega4.c04.subsec.19-3}
\expandafter\def\csname sga@target@subsection:IV.19.4\endcsname{ega4.c04.subsec.19-4}
\expandafter\def\csname sga@target@subsection:IV.19.5\endcsname{ega4.c04.subsec.19-5}
\expandafter\def\csname sga@target@subsection:IV.19.6\endcsname{ega4.c04.subsec.19-6}
\expandafter\def\csname sga@target@subsection:IV.19.7\endcsname{ega4.c04.subsec.19-7}
\expandafter\def\csname sga@target@subsection:IV.19.8\endcsname{ega4.c04.subsec.19-8}
\expandafter\def\csname sga@target@subsection:IV.19.9\endcsname{ega4.c04.subsec.19-9}
\expandafter\def\csname sga@target@subsection:IV.2.1\endcsname{ega4.c04.subsec.2-1}
\expandafter\def\csname sga@target@subsection:IV.2.2\endcsname{ega4.c04.subsec.2-2}
\expandafter\def\csname sga@target@subsection:IV.2.3\endcsname{ega4.c04.subsec.2-3}
\expandafter\def\csname sga@target@subsection:IV.2.4\endcsname{ega4.c04.subsec.2-4}
\expandafter\def\csname sga@target@subsection:IV.2.5\endcsname{ega4.c04.subsec.2-5}
\expandafter\def\csname sga@target@subsection:IV.2.6\endcsname{ega4.c04.subsec.2-6}
\expandafter\def\csname sga@target@subsection:IV.2.7\endcsname{ega4.c04.subsec.2-7}
\expandafter\def\csname sga@target@subsection:IV.2.8\endcsname{ega4.c04.subsec.2-8}
\expandafter\def\csname sga@target@subsection:IV.21.1\endcsname{ega4.c04.subsec.21-1}
\expandafter\def\csname sga@target@subsection:IV.21.10\endcsname{ega4.c04.subsec.21-10}
\expandafter\def\csname sga@target@subsection:IV.21.11\endcsname{ega4.c04.subsec.21-11}
\expandafter\def\csname sga@target@subsection:IV.21.12\endcsname{ega4.c04.subsec.21-12}
\expandafter\def\csname sga@target@subsection:IV.21.13\endcsname{ega4.c04.subsec.21-13}
\expandafter\def\csname sga@target@subsection:IV.21.14\endcsname{ega4.c04.subsec.21-14}
\expandafter\def\csname sga@target@subsection:IV.21.15\endcsname{ega4.c04.subsec.21-15}
\expandafter\def\csname sga@target@subsection:IV.21.2\endcsname{ega4.c04.subsec.21-2}
\expandafter\def\csname sga@target@subsection:IV.21.3\endcsname{ega4.c04.subsec.21-3}
\expandafter\def\csname sga@target@subsection:IV.21.4\endcsname{ega4.c04.subsec.21-4}
\expandafter\def\csname sga@target@subsection:IV.21.5\endcsname{ega4.c04.subsec.21-5}
\expandafter\def\csname sga@target@subsection:IV.21.6\endcsname{ega4.c04.subsec.21-6}
\expandafter\def\csname sga@target@subsection:IV.21.7\endcsname{ega4.c04.subsec.21-7}
\expandafter\def\csname sga@target@subsection:IV.21.8\endcsname{ega4.c04.subsec.21-8}
\expandafter\def\csname sga@target@subsection:IV.21.9\endcsname{ega4.c04.subsec.21-9}
\expandafter\def\csname sga@target@subsection:IV.3.1\endcsname{ega4.c04.subsec.3-1}
\expandafter\def\csname sga@target@subsection:IV.3.2\endcsname{ega4.c04.subsec.3-2}
\expandafter\def\csname sga@target@subsection:IV.3.3\endcsname{ega4.c04.subsec.3-3}
\expandafter\def\csname sga@target@subsection:IV.3.4\endcsname{ega4.c04.subsec.3-4}
\expandafter\def\csname sga@target@subsection:IV.4.1\endcsname{ega4.c04.subsec.4-1}
\expandafter\def\csname sga@target@subsection:IV.4.2\endcsname{ega4.c04.subsec.4-2}
\expandafter\def\csname sga@target@subsection:IV.4.3\endcsname{ega4.c04.subsec.4-3}
\expandafter\def\csname sga@target@subsection:IV.4.4\endcsname{ega4.c04.subsec.4-4}
\expandafter\def\csname sga@target@subsection:IV.4.5\endcsname{ega4.c04.subsec.4-5}
\expandafter\def\csname sga@target@subsection:IV.4.6\endcsname{ega4.c04.subsec.4-6}
\expandafter\def\csname sga@target@subsection:IV.4.7\endcsname{ega4.c04.subsec.4-7}
\expandafter\def\csname sga@target@subsection:IV.4.8\endcsname{ega4.c04.subsec.4-8}
\expandafter\def\csname sga@target@subsection:IV.4.9\endcsname{ega4.c04.subsec.4-9}
\expandafter\def\csname sga@target@subsection:IV.5.1\endcsname{ega4.c04.subsec.5-1}
\expandafter\def\csname sga@target@subsection:IV.5.10\endcsname{ega4.c04.subsec.5-10}
\expandafter\def\csname sga@target@subsection:IV.5.11\endcsname{ega4.c04.subsec.5-11}
\expandafter\def\csname sga@target@subsection:IV.5.12\endcsname{ega4.c04.subsec.5-12}
\expandafter\def\csname sga@target@subsection:IV.5.13\endcsname{ega4.c04.subsec.5-13}
\expandafter\def\csname sga@target@subsection:IV.5.2\endcsname{ega4.c04.subsec.5-2}
\expandafter\def\csname sga@target@subsection:IV.5.3\endcsname{ega4.c04.subsec.5-3}
\expandafter\def\csname sga@target@subsection:IV.5.4\endcsname{ega4.c04.subsec.5-4}
\expandafter\def\csname sga@target@subsection:IV.5.5\endcsname{ega4.c04.subsec.5-5}
\expandafter\def\csname sga@target@subsection:IV.5.6\endcsname{ega4.c04.subsec.5-6}
\expandafter\def\csname sga@target@subsection:IV.5.7\endcsname{ega4.c04.subsec.5-7}
\expandafter\def\csname sga@target@subsection:IV.5.8\endcsname{ega4.c04.subsec.5-8}
\expandafter\def\csname sga@target@subsection:IV.5.9\endcsname{ega4.c04.subsec.5-9}
\expandafter\def\csname sga@target@subsection:IV.6.1\endcsname{ega4.c04.subsec.6-1}
\expandafter\def\csname sga@target@subsection:IV.6.10\endcsname{ega4.c04.subsec.6-10}
\expandafter\def\csname sga@target@subsection:IV.6.11\endcsname{ega4.c04.subsec.6-11}
\expandafter\def\csname sga@target@subsection:IV.6.12\endcsname{ega4.c04.subsec.6-12}
\expandafter\def\csname sga@target@subsection:IV.6.13\endcsname{ega4.c04.subsec.6-13}
\expandafter\def\csname sga@target@subsection:IV.6.14\endcsname{ega4.c04.subsec.6-14}
\expandafter\def\csname sga@target@subsection:IV.6.15\endcsname{ega4.c04.subsec.6-15}
\expandafter\def\csname sga@target@subsection:IV.6.2\endcsname{ega4.c04.subsec.6-2}
\expandafter\def\csname sga@target@subsection:IV.6.3\endcsname{ega4.c04.subsec.6-3}
\expandafter\def\csname sga@target@subsection:IV.6.4\endcsname{ega4.c04.subsec.6-4}
\expandafter\def\csname sga@target@subsection:IV.6.5\endcsname{ega4.c04.subsec.6-5}
\expandafter\def\csname sga@target@subsection:IV.6.6\endcsname{ega4.c04.subsec.6-6}
\expandafter\def\csname sga@target@subsection:IV.6.7\endcsname{ega4.c04.subsec.6-7}
\expandafter\def\csname sga@target@subsection:IV.6.8\endcsname{ega4.c04.subsec.6-8}
\expandafter\def\csname sga@target@subsection:IV.6.9\endcsname{ega4.c04.subsec.6-9}
\expandafter\def\csname sga@target@subsection:IV.7.1\endcsname{ega4.c04.subsec.7-1}
\expandafter\def\csname sga@target@subsection:IV.7.2\endcsname{ega4.c04.subsec.7-2}
\expandafter\def\csname sga@target@subsection:IV.7.3\endcsname{ega4.c04.subsec.7-3}
\expandafter\def\csname sga@target@subsection:IV.7.4\endcsname{ega4.c04.subsec.7-4}
\expandafter\def\csname sga@target@subsection:IV.7.5\endcsname{ega4.c04.subsec.7-5}
\expandafter\def\csname sga@target@subsection:IV.7.6\endcsname{ega4.c04.subsec.7-6}
\expandafter\def\csname sga@target@subsection:IV.7.7\endcsname{ega4.c04.subsec.7-7}
\expandafter\def\csname sga@target@subsection:IV.7.8\endcsname{ega4.c04.subsec.7-8}
\expandafter\def\csname sga@target@subsection:IV.7.9\endcsname{ega4.c04.subsec.7-9}
\expandafter\def\csname sga@target@subsection:IV.8.1\endcsname{ega4.c04.subsec.8-1}
\expandafter\def\csname sga@target@subsection:IV.8.10\endcsname{ega4.c04.subsec.8-10}
\expandafter\def\csname sga@target@subsection:IV.8.11\endcsname{ega4.c04.subsec.8-11}
\expandafter\def\csname sga@target@subsection:IV.8.12\endcsname{ega4.c04.subsec.8-12}
\expandafter\def\csname sga@target@subsection:IV.8.13\endcsname{ega4.c04.subsec.8-13}
\expandafter\def\csname sga@target@subsection:IV.8.14\endcsname{ega4.c04.subsec.8-14}
\expandafter\def\csname sga@target@subsection:IV.8.2\endcsname{ega4.c04.subsec.8-2}
\expandafter\def\csname sga@target@subsection:IV.8.3\endcsname{ega4.c04.subsec.8-3}
\expandafter\def\csname sga@target@subsection:IV.8.4\endcsname{ega4.c04.subsec.8-4}
\expandafter\def\csname sga@target@subsection:IV.8.5\endcsname{ega4.c04.subsec.8-5}
\expandafter\def\csname sga@target@subsection:IV.8.6\endcsname{ega4.c04.subsec.8-6}
\expandafter\def\csname sga@target@subsection:IV.8.7\endcsname{ega4.c04.subsec.8-7}
\expandafter\def\csname sga@target@subsection:IV.8.8\endcsname{ega4.c04.subsec.8-8}
\expandafter\def\csname sga@target@subsection:IV.8.9\endcsname{ega4.c04.subsec.8-9}
\expandafter\def\csname sga@target@subsection:IV.9.1\endcsname{ega4.c04.subsec.9-1}
\expandafter\def\csname sga@target@subsection:IV.9.2\endcsname{ega4.c04.subsec.9-2}
\expandafter\def\csname sga@target@subsection:IV.9.3\endcsname{ega4.c04.subsec.9-3}
\expandafter\def\csname sga@target@subsection:IV.9.4\endcsname{ega4.c04.subsec.9-4}
\expandafter\def\csname sga@target@subsection:IV.9.5\endcsname{ega4.c04.subsec.9-5}
\expandafter\def\csname sga@target@subsection:IV.9.6\endcsname{ega4.c04.subsec.9-6}
\expandafter\def\csname sga@target@subsection:IV.9.7\endcsname{ega4.c04.subsec.9-7}
\expandafter\def\csname sga@target@subsection:IV.9.8\endcsname{ega4.c04.subsec.9-8}
\expandafter\def\csname sga@target@subsection:IV.9.9\endcsname{ega4.c04.subsec.9-9}
\let\sgaOriginalRef\ref
\let\sgaOriginalPageref\pageref
\renewcommand{\ref}[1]{%
  \ifcsname sga@target@#1\endcsname
    \hyperlink{\csname sga@target@#1\endcsname}{\sgaOriginalRef*{#1}}%
  \else
    \sgaOriginalRef{#1}%
  \fi
}
\renewcommand{\pageref}[1]{%
  \ifcsname sga@target@#1\endcsname
    \hyperlink{\csname sga@target@#1\endcsname}{\sgaOriginalPageref*{#1}}%
  \else
    \sgaOriginalPageref{#1}%
  \fi
}
\renewcommand{\eqref}[1]{%
  \ifcsname sga@target@#1\endcsname
    \hyperlink{\csname sga@target@#1\endcsname}{\textup{\tagform@{\sgaOriginalRef*{#1}}}}%
  \else
    \textup{\tagform@{\sgaOriginalRef{#1}}}%
  \fi
}
% Stable semantic routing for the reader's prose reference helpers.
\newcommand{\sgaStableOrOriginalLink}[2]{%
  \ifcsname sga@target@#1\endcsname
    \hyperlink{\csname sga@target@#1\endcsname}{#2}%
  \else
    \@ifundefined{r@#1}{#2}{\hyperref[#1]{#2}}%
  \fi
}
\RenewDocumentCommand{\sref}{o m o}{%
  \IfNoValueTF{#1}%
    {\IfNoValueTF{#3}%
      {\sgaStableOrOriginalLink{#2}{\normalfont{(\StrBehind{#2}{.})}}}%
      {\sgaStableOrOriginalLink{#2}{\normalfont{(\StrBehind{#2}{.},~{#3})}}}}%
    {\IfNoValueTF{#3}%
      {\sgaStableOrOriginalLink{#2}{\normalfont{(\textbf{#1},~\StrBehind{#2}{.})}}}%
      {\sgaStableOrOriginalLink{#2}{\normalfont{(\textbf{#1},~\StrBehind{#2}{.},~{#3})}}}}%
}
\renewcommand{\eref}[1]{%
  \sgaStableOrOriginalLink{#1}{\normalfont{(\StrBehind{#1}{.})}}}
\makeatother

\makeatletter
\expandafter\def\csname EGAExternal@0.1.3.2\endcsname{https://github.com/KokunoYumeto/ega-fr/releases/download/ega-fr-2026-10-23-r16/EGA_French_Current_Cumulative_Reader.pdf\#page=11}
\expandafter\def\csname EGAExternal@0.14.1.10\endcsname{https://github.com/KokunoYumeto/ega-fr/releases/download/ega-fr-2026-10-23-r16/EGA_French_Current_Cumulative_Reader.pdf\#page=493}
\expandafter\def\csname EGAExternal@0.14.1.2.1\endcsname{https://github.com/KokunoYumeto/ega-fr/releases/download/ega-fr-2026-10-23-r16/EGA_French_Current_Cumulative_Reader.pdf\#page=492}
\expandafter\def\csname EGAExternal@0.14.1.4\endcsname{https://github.com/KokunoYumeto/ega-fr/releases/download/ega-fr-2026-10-23-r16/EGA_French_Current_Cumulative_Reader.pdf\#page=492}
\expandafter\def\csname EGAExternal@0.14.1.4.1\endcsname{https://github.com/KokunoYumeto/ega-fr/releases/download/ega-fr-2026-10-23-r16/EGA_French_Current_Cumulative_Reader.pdf\#page=493}
\expandafter\def\csname EGAExternal@0.14.1.5\endcsname{https://github.com/KokunoYumeto/ega-fr/releases/download/ega-fr-2026-10-23-r16/EGA_French_Current_Cumulative_Reader.pdf\#page=493}
\expandafter\def\csname EGAExternal@0.14.1.6\endcsname{https://github.com/KokunoYumeto/ega-fr/releases/download/ega-fr-2026-10-23-r16/EGA_French_Current_Cumulative_Reader.pdf\#page=493}
\expandafter\def\csname EGAExternal@0.14.1.7\endcsname{https://github.com/KokunoYumeto/ega-fr/releases/download/ega-fr-2026-10-23-r16/EGA_French_Current_Cumulative_Reader.pdf\#page=493}
\expandafter\def\csname EGAExternal@0.14.2.1\endcsname{https://github.com/KokunoYumeto/ega-fr/releases/download/ega-fr-2026-10-23-r16/EGA_French_Current_Cumulative_Reader.pdf\#page=494}
\expandafter\def\csname EGAExternal@0.14.2.1.1\endcsname{https://github.com/KokunoYumeto/ega-fr/releases/download/ega-fr-2026-10-23-r16/EGA_French_Current_Cumulative_Reader.pdf\#page=494}
\expandafter\def\csname EGAExternal@0.14.2.2\endcsname{https://github.com/KokunoYumeto/ega-fr/releases/download/ega-fr-2026-10-23-r16/EGA_French_Current_Cumulative_Reader.pdf\#page=494}
\expandafter\def\csname EGAExternal@0.14.2.2.2\endcsname{https://github.com/KokunoYumeto/ega-fr/releases/download/ega-fr-2026-10-23-r16/EGA_French_Current_Cumulative_Reader.pdf\#page=494}
\expandafter\def\csname EGAExternal@0.14.3.2\endcsname{https://github.com/KokunoYumeto/ega-fr/releases/download/ega-fr-2026-10-23-r16/EGA_French_Current_Cumulative_Reader.pdf\#page=495}
\expandafter\def\csname EGAExternal@0.14.3.3\endcsname{https://github.com/KokunoYumeto/ega-fr/releases/download/ega-fr-2026-10-23-r16/EGA_French_Current_Cumulative_Reader.pdf\#page=495}
\expandafter\def\csname EGAExternal@0.14.3.3.2\endcsname{https://github.com/KokunoYumeto/ega-fr/releases/download/ega-fr-2026-10-23-r16/EGA_French_Current_Cumulative_Reader.pdf\#page=496}
\expandafter\def\csname EGAExternal@0.14.3.5\endcsname{https://github.com/KokunoYumeto/ega-fr/releases/download/ega-fr-2026-10-23-r16/EGA_French_Current_Cumulative_Reader.pdf\#page=496}
\expandafter\def\csname EGAExternal@0.14.3.5.1\endcsname{https://github.com/KokunoYumeto/ega-fr/releases/download/ega-fr-2026-10-23-r16/EGA_French_Current_Cumulative_Reader.pdf\#page=496}
\expandafter\def\csname EGAExternal@0.16.1.2.1\endcsname{https://github.com/KokunoYumeto/ega-fr/releases/download/ega-fr-2026-10-23-r16/EGA_French_Current_Cumulative_Reader.pdf\#page=505}
\expandafter\def\csname EGAExternal@0.16.1.5\endcsname{https://github.com/KokunoYumeto/ega-fr/releases/download/ega-fr-2026-10-23-r16/EGA_French_Current_Cumulative_Reader.pdf\#page=506}
\expandafter\def\csname EGAExternal@0.16.3.3\endcsname{https://github.com/KokunoYumeto/ega-fr/releases/download/ega-fr-2026-10-23-r16/EGA_French_Current_Cumulative_Reader.pdf\#page=509}
\expandafter\def\csname EGAExternal@0.16.3.9\endcsname{https://github.com/KokunoYumeto/ega-fr/releases/download/ega-fr-2026-10-23-r16/EGA_French_Current_Cumulative_Reader.pdf\#page=510}
\expandafter\def\csname EGAExternal@0.16.4.10\endcsname{https://github.com/KokunoYumeto/ega-fr/releases/download/ega-fr-2026-10-23-r16/EGA_French_Current_Cumulative_Reader.pdf\#page=514}
\expandafter\def\csname EGAExternal@0.16.4.9\endcsname{https://github.com/KokunoYumeto/ega-fr/releases/download/ega-fr-2026-10-23-r16/EGA_French_Current_Cumulative_Reader.pdf\#page=514}
\expandafter\def\csname EGAExternal@0.16.5.12\endcsname{https://github.com/KokunoYumeto/ega-fr/releases/download/ega-fr-2026-10-23-r16/EGA_French_Current_Cumulative_Reader.pdf\#page=516}
\expandafter\def\csname EGAExternal@0.16.5.4\endcsname{https://github.com/KokunoYumeto/ega-fr/releases/download/ega-fr-2026-10-23-r16/EGA_French_Current_Cumulative_Reader.pdf\#page=515}
\expandafter\def\csname EGAExternal@0.17.1.1\endcsname{https://github.com/KokunoYumeto/ega-fr/releases/download/ega-fr-2026-10-23-r16/EGA_French_Current_Cumulative_Reader.pdf\#page=517}
\expandafter\def\csname EGAExternal@0.17.2.2\endcsname{https://github.com/KokunoYumeto/ega-fr/releases/download/ega-fr-2026-10-23-r16/EGA_French_Current_Cumulative_Reader.pdf\#page=519}
\expandafter\def\csname EGAExternal@0.17.2.4\endcsname{https://github.com/KokunoYumeto/ega-fr/releases/download/ega-fr-2026-10-23-r16/EGA_French_Current_Cumulative_Reader.pdf\#page=519}
\expandafter\def\csname EGAExternal@0.17.3.4\endcsname{https://github.com/KokunoYumeto/ega-fr/releases/download/ega-fr-2026-10-23-r16/EGA_French_Current_Cumulative_Reader.pdf\#page=523}
\expandafter\def\csname EGAExternal@0.17.3.7\endcsname{https://github.com/KokunoYumeto/ega-fr/releases/download/ega-fr-2026-10-23-r16/EGA_French_Current_Cumulative_Reader.pdf\#page=524}
\expandafter\def\csname EGAExternal@0.19.1.12\endcsname{https://github.com/KokunoYumeto/ega-fr/releases/download/ega-fr-2026-10-23-r16/EGA_French_Current_Cumulative_Reader.pdf\#page=539}
\expandafter\def\csname EGAExternal@0.2.1.3\endcsname{https://github.com/KokunoYumeto/ega-fr/releases/download/ega-fr-2026-10-23-r16/EGA_French_Current_Cumulative_Reader.pdf\#page=16}
\expandafter\def\csname EGAExternal@0.2.1.6\endcsname{https://github.com/KokunoYumeto/ega-fr/releases/download/ega-fr-2026-10-23-r16/EGA_French_Current_Cumulative_Reader.pdf\#page=16}
\expandafter\def\csname EGAExternal@0.2.1.7\endcsname{https://github.com/KokunoYumeto/ega-fr/releases/download/ega-fr-2026-10-23-r16/EGA_French_Current_Cumulative_Reader.pdf\#page=16}
\expandafter\def\csname EGAExternal@0.2.1.8\endcsname{https://github.com/KokunoYumeto/ega-fr/releases/download/ega-fr-2026-10-23-r16/EGA_French_Current_Cumulative_Reader.pdf\#page=17}
\expandafter\def\csname EGAExternal@0.2.2.2\endcsname{https://github.com/KokunoYumeto/ega-fr/releases/download/ega-fr-2026-10-23-r16/EGA_French_Current_Cumulative_Reader.pdf\#page=17}
\expandafter\def\csname EGAExternal@0.23.1.5\endcsname{https://github.com/KokunoYumeto/ega-fr/releases/download/ega-fr-2026-10-23-r16/EGA_French_Current_Cumulative_Reader.pdf\#page=641}
\expandafter\def\csname EGAExternal@0.23.1.6\endcsname{https://github.com/KokunoYumeto/ega-fr/releases/download/ega-fr-2026-10-23-r16/EGA_French_Current_Cumulative_Reader.pdf\#page=641}
\expandafter\def\csname EGAExternal@0.23.2.5\endcsname{https://github.com/KokunoYumeto/ega-fr/releases/download/ega-fr-2026-10-23-r16/EGA_French_Current_Cumulative_Reader.pdf\#page=643}
\expandafter\def\csname EGAExternal@0.3.3.1\endcsname{https://github.com/KokunoYumeto/ega-fr/releases/download/ega-fr-2026-10-23-r16/EGA_French_Current_Cumulative_Reader.pdf\#page=21}
\expandafter\def\csname EGAExternal@0.3.7.1\endcsname{https://github.com/KokunoYumeto/ega-fr/releases/download/ega-fr-2026-10-23-r16/EGA_French_Current_Cumulative_Reader.pdf\#page=26}
\expandafter\def\csname EGAExternal@0.3.7.2\endcsname{https://github.com/KokunoYumeto/ega-fr/releases/download/ega-fr-2026-10-23-r16/EGA_French_Current_Cumulative_Reader.pdf\#page=26}
\expandafter\def\csname EGAExternal@0.4.1.7\endcsname{https://github.com/KokunoYumeto/ega-fr/releases/download/ega-fr-2026-10-23-r16/EGA_French_Current_Cumulative_Reader.pdf\#page=27}
\expandafter\def\csname EGAExternal@0.4.3.1\endcsname{https://github.com/KokunoYumeto/ega-fr/releases/download/ega-fr-2026-10-23-r16/EGA_French_Current_Cumulative_Reader.pdf\#page=30}
\expandafter\def\csname EGAExternal@0.4.3.3\endcsname{https://github.com/KokunoYumeto/ega-fr/releases/download/ega-fr-2026-10-23-r16/EGA_French_Current_Cumulative_Reader.pdf\#page=31}
\expandafter\def\csname EGAExternal@0.4.4.3.2\endcsname{https://github.com/KokunoYumeto/ega-fr/releases/download/ega-fr-2026-10-23-r16/EGA_French_Current_Cumulative_Reader.pdf\#page=32}
\expandafter\def\csname EGAExternal@0.4.4.3.3\endcsname{https://github.com/KokunoYumeto/ega-fr/releases/download/ega-fr-2026-10-23-r16/EGA_French_Current_Cumulative_Reader.pdf\#page=32}
\expandafter\def\csname EGAExternal@0.5.1.1\endcsname{https://github.com/KokunoYumeto/ega-fr/releases/download/ega-fr-2026-10-23-r16/EGA_French_Current_Cumulative_Reader.pdf\#page=33}
\expandafter\def\csname EGAExternal@0.5.2.2\endcsname{https://github.com/KokunoYumeto/ega-fr/releases/download/ega-fr-2026-10-23-r16/EGA_French_Current_Cumulative_Reader.pdf\#page=34}
\expandafter\def\csname EGAExternal@0.5.2.5\endcsname{https://github.com/KokunoYumeto/ega-fr/releases/download/ega-fr-2026-10-23-r16/EGA_French_Current_Cumulative_Reader.pdf\#page=34}
\expandafter\def\csname EGAExternal@0.5.3.8\endcsname{https://github.com/KokunoYumeto/ega-fr/releases/download/ega-fr-2026-10-23-r16/EGA_French_Current_Cumulative_Reader.pdf\#page=35}
\expandafter\def\csname EGAExternal@0.5.5.1\endcsname{https://github.com/KokunoYumeto/ega-fr/releases/download/ega-fr-2026-10-23-r16/EGA_French_Current_Cumulative_Reader.pdf\#page=39}
\expandafter\def\csname EGAExternal@0.5.5.5\endcsname{https://github.com/KokunoYumeto/ega-fr/releases/download/ega-fr-2026-10-23-r16/EGA_French_Current_Cumulative_Reader.pdf\#page=39}
\expandafter\def\csname EGAExternal@0.6.2.3\endcsname{https://github.com/KokunoYumeto/ega-fr/releases/download/ega-fr-2026-10-23-r16/EGA_French_Current_Cumulative_Reader.pdf\#page=41}
\expandafter\def\csname EGAExternal@0.6.3.4\endcsname{https://github.com/KokunoYumeto/ega-fr/releases/download/ega-fr-2026-10-23-r16/EGA_French_Current_Cumulative_Reader.pdf\#page=41}
\expandafter\def\csname EGAExternal@0.6.5.1\endcsname{https://github.com/KokunoYumeto/ega-fr/releases/download/ega-fr-2026-10-23-r16/EGA_French_Current_Cumulative_Reader.pdf\#page=43}
\expandafter\def\csname EGAExternal@0.6.6.2\endcsname{https://github.com/KokunoYumeto/ega-fr/releases/download/ega-fr-2026-10-23-r16/EGA_French_Current_Cumulative_Reader.pdf\#page=43}
\expandafter\def\csname EGAExternal@0.6.7.2\endcsname{https://github.com/KokunoYumeto/ega-fr/releases/download/ega-fr-2026-10-23-r16/EGA_French_Current_Cumulative_Reader.pdf\#page=43}
\expandafter\def\csname EGAExternal@0.7.3.5\endcsname{https://github.com/KokunoYumeto/ega-fr/releases/download/ega-fr-2026-10-23-r16/EGA_French_Current_Cumulative_Reader.pdf\#page=48}
\expandafter\def\csname EGAExternal@0.8.2.5\endcsname{https://github.com/KokunoYumeto/ega-fr/releases/download/ega-fr-2026-10-23-r16/EGA_French_Current_Cumulative_Reader.pdf\#page=304}
\expandafter\def\csname EGAExternal@0.9.1.1\endcsname{https://github.com/KokunoYumeto/ega-fr/releases/download/ega-fr-2026-10-23-r16/EGA_French_Current_Cumulative_Reader.pdf\#page=305}
\expandafter\def\csname EGAExternal@0.9.1.3\endcsname{https://github.com/KokunoYumeto/ega-fr/releases/download/ega-fr-2026-10-23-r16/EGA_French_Current_Cumulative_Reader.pdf\#page=306}
\expandafter\def\csname EGAExternal@0.9.1.5\endcsname{https://github.com/KokunoYumeto/ega-fr/releases/download/ega-fr-2026-10-23-r16/EGA_French_Current_Cumulative_Reader.pdf\#page=305}
\expandafter\def\csname EGAExternal@0.9.2.3\endcsname{https://github.com/KokunoYumeto/ega-fr/releases/download/ega-fr-2026-10-23-r16/EGA_French_Current_Cumulative_Reader.pdf\#page=307}
\expandafter\def\csname EGAExternal@0.9.2.5\endcsname{https://github.com/KokunoYumeto/ega-fr/releases/download/ega-fr-2026-10-23-r16/EGA_French_Current_Cumulative_Reader.pdf\#page=307}
\expandafter\def\csname EGAExternal@I.1.1.4\endcsname{https://github.com/KokunoYumeto/ega-fr/releases/download/ega-fr-2026-10-23-r16/EGA_French_Current_Cumulative_Reader.pdf\#page=58}
\expandafter\def\csname EGAExternal@I.1.1.5\endcsname{https://github.com/KokunoYumeto/ega-fr/releases/download/ega-fr-2026-10-23-r16/EGA_French_Current_Cumulative_Reader.pdf\#page=58}
\expandafter\def\csname EGAExternal@I.1.2.2\endcsname{https://github.com/KokunoYumeto/ega-fr/releases/download/ega-fr-2026-10-23-r16/EGA_French_Current_Cumulative_Reader.pdf\#page=60}
\expandafter\def\csname EGAExternal@I.1.2.3\endcsname{https://github.com/KokunoYumeto/ega-fr/releases/download/ega-fr-2026-10-23-r16/EGA_French_Current_Cumulative_Reader.pdf\#page=60}
\expandafter\def\csname EGAExternal@I.1.2.7\endcsname{https://github.com/KokunoYumeto/ega-fr/releases/download/ega-fr-2026-10-23-r16/EGA_French_Current_Cumulative_Reader.pdf\#page=60}
\expandafter\def\csname EGAExternal@I.1.3.13\endcsname{https://github.com/KokunoYumeto/ega-fr/releases/download/ega-fr-2026-10-23-r16/EGA_French_Current_Cumulative_Reader.pdf\#page=63}
\expandafter\def\csname EGAExternal@I.1.3.4\endcsname{https://github.com/KokunoYumeto/ega-fr/releases/download/ega-fr-2026-10-23-r16/EGA_French_Current_Cumulative_Reader.pdf\#page=61}
\expandafter\def\csname EGAExternal@I.1.3.9\endcsname{https://github.com/KokunoYumeto/ega-fr/releases/download/ega-fr-2026-10-23-r16/EGA_French_Current_Cumulative_Reader.pdf\#page=62}
\expandafter\def\csname EGAExternal@I.1.4.3\endcsname{https://github.com/KokunoYumeto/ega-fr/releases/download/ega-fr-2026-10-23-r16/EGA_French_Current_Cumulative_Reader.pdf\#page=65}
\expandafter\def\csname EGAExternal@I.1.5.4\endcsname{https://github.com/KokunoYumeto/ega-fr/releases/download/ega-fr-2026-10-23-r16/EGA_French_Current_Cumulative_Reader.pdf\#page=66}
\expandafter\def\csname EGAExternal@I.1.6.1\endcsname{https://github.com/KokunoYumeto/ega-fr/releases/download/ega-fr-2026-10-23-r16/EGA_French_Current_Cumulative_Reader.pdf\#page=67}
\expandafter\def\csname EGAExternal@I.1.6.2\endcsname{https://github.com/KokunoYumeto/ega-fr/releases/download/ega-fr-2026-10-23-r16/EGA_French_Current_Cumulative_Reader.pdf\#page=67}
\expandafter\def\csname EGAExternal@I.1.6.3\endcsname{https://github.com/KokunoYumeto/ega-fr/releases/download/ega-fr-2026-10-23-r16/EGA_French_Current_Cumulative_Reader.pdf\#page=67}
\expandafter\def\csname EGAExternal@I.1.6.5\endcsname{https://github.com/KokunoYumeto/ega-fr/releases/download/ega-fr-2026-10-23-r16/EGA_French_Current_Cumulative_Reader.pdf\#page=67}
\expandafter\def\csname EGAExternal@I.1.7.3\endcsname{https://github.com/KokunoYumeto/ega-fr/releases/download/ega-fr-2026-10-23-r16/EGA_French_Current_Cumulative_Reader.pdf\#page=69}
\expandafter\def\csname EGAExternal@I.1.8.2\endcsname{https://github.com/KokunoYumeto/ega-fr/releases/download/ega-fr-2026-10-23-r16/EGA_French_Current_Cumulative_Reader.pdf\#page=297}
\expandafter\def\csname EGAExternal@I.1.8.4\endcsname{https://github.com/KokunoYumeto/ega-fr/releases/download/ega-fr-2026-10-23-r16/EGA_French_Current_Cumulative_Reader.pdf\#page=298}
\expandafter\def\csname EGAExternal@I.1.8.5\endcsname{https://github.com/KokunoYumeto/ega-fr/releases/download/ega-fr-2026-10-23-r16/EGA_French_Current_Cumulative_Reader.pdf\#page=298}
\expandafter\def\csname EGAExternal@I.10.4.8\endcsname{https://github.com/KokunoYumeto/ega-fr/releases/download/ega-fr-2026-10-23-r16/EGA_French_Current_Cumulative_Reader.pdf\#page=129}
\expandafter\def\csname EGAExternal@I.2.1.3\endcsname{https://github.com/KokunoYumeto/ega-fr/releases/download/ega-fr-2026-10-23-r16/EGA_French_Current_Cumulative_Reader.pdf\#page=70}
\expandafter\def\csname EGAExternal@I.2.1.4\endcsname{https://github.com/KokunoYumeto/ega-fr/releases/download/ega-fr-2026-10-23-r16/EGA_French_Current_Cumulative_Reader.pdf\#page=70}
\expandafter\def\csname EGAExternal@I.2.2.4\endcsname{https://github.com/KokunoYumeto/ega-fr/releases/download/ega-fr-2026-10-23-r16/EGA_French_Current_Cumulative_Reader.pdf\#page=70}
\expandafter\def\csname EGAExternal@I.2.4.2\endcsname{https://github.com/KokunoYumeto/ega-fr/releases/download/ega-fr-2026-10-23-r16/EGA_French_Current_Cumulative_Reader.pdf\#page=72}
\expandafter\def\csname EGAExternal@I.2.4.4\endcsname{https://github.com/KokunoYumeto/ega-fr/releases/download/ega-fr-2026-10-23-r16/EGA_French_Current_Cumulative_Reader.pdf\#page=72}
\expandafter\def\csname EGAExternal@I.2.5.5\endcsname{https://github.com/KokunoYumeto/ega-fr/releases/download/ega-fr-2026-10-23-r16/EGA_French_Current_Cumulative_Reader.pdf\#page=74}
\expandafter\def\csname EGAExternal@I.3.3.11\endcsname{https://github.com/KokunoYumeto/ega-fr/releases/download/ega-fr-2026-10-23-r16/EGA_French_Current_Cumulative_Reader.pdf\#page=78}
\expandafter\def\csname EGAExternal@I.3.3.14\endcsname{https://github.com/KokunoYumeto/ega-fr/releases/download/ega-fr-2026-10-23-r16/EGA_French_Current_Cumulative_Reader.pdf\#page=79}
\expandafter\def\csname EGAExternal@I.3.3.15\endcsname{https://github.com/KokunoYumeto/ega-fr/releases/download/ega-fr-2026-10-23-r16/EGA_French_Current_Cumulative_Reader.pdf\#page=79}
\expandafter\def\csname EGAExternal@I.3.4.3.1\endcsname{https://github.com/KokunoYumeto/ega-fr/releases/download/ega-fr-2026-10-23-r16/EGA_French_Current_Cumulative_Reader.pdf\#page=79}
\expandafter\def\csname EGAExternal@I.3.4.4\endcsname{https://github.com/KokunoYumeto/ega-fr/releases/download/ega-fr-2026-10-23-r16/EGA_French_Current_Cumulative_Reader.pdf\#page=80}
\expandafter\def\csname EGAExternal@I.3.4.5\endcsname{https://github.com/KokunoYumeto/ega-fr/releases/download/ega-fr-2026-10-23-r16/EGA_French_Current_Cumulative_Reader.pdf\#page=80}
\expandafter\def\csname EGAExternal@I.3.4.7\endcsname{https://github.com/KokunoYumeto/ega-fr/releases/download/ega-fr-2026-10-23-r16/EGA_French_Current_Cumulative_Reader.pdf\#page=80}
\expandafter\def\csname EGAExternal@I.3.4.8\endcsname{https://github.com/KokunoYumeto/ega-fr/releases/download/ega-fr-2026-10-23-r16/EGA_French_Current_Cumulative_Reader.pdf\#page=80}
\expandafter\def\csname EGAExternal@I.3.4.9\endcsname{https://github.com/KokunoYumeto/ega-fr/releases/download/ega-fr-2026-10-23-r16/EGA_French_Current_Cumulative_Reader.pdf\#page=81}
\expandafter\def\csname EGAExternal@I.3.5.2\endcsname{https://github.com/KokunoYumeto/ega-fr/releases/download/ega-fr-2026-10-23-r16/EGA_French_Current_Cumulative_Reader.pdf\#page=81}
\expandafter\def\csname EGAExternal@I.3.5.7\endcsname{https://github.com/KokunoYumeto/ega-fr/releases/download/ega-fr-2026-10-23-r16/EGA_French_Current_Cumulative_Reader.pdf\#page=82}
\expandafter\def\csname EGAExternal@I.3.5.8\endcsname{https://github.com/KokunoYumeto/ega-fr/releases/download/ega-fr-2026-10-23-r16/EGA_French_Current_Cumulative_Reader.pdf\#page=82}
\expandafter\def\csname EGAExternal@I.3.6.1\endcsname{https://github.com/KokunoYumeto/ega-fr/releases/download/ega-fr-2026-10-23-r16/EGA_French_Current_Cumulative_Reader.pdf\#page=83}
\expandafter\def\csname EGAExternal@I.3.6.4\endcsname{https://github.com/KokunoYumeto/ega-fr/releases/download/ega-fr-2026-10-23-r16/EGA_French_Current_Cumulative_Reader.pdf\#page=83}
\expandafter\def\csname EGAExternal@I.3.6.5\endcsname{https://github.com/KokunoYumeto/ega-fr/releases/download/ega-fr-2026-10-23-r16/EGA_French_Current_Cumulative_Reader.pdf\#page=83}
\expandafter\def\csname EGAExternal@I.3.7\endcsname{https://github.com/KokunoYumeto/ega-fr/releases/download/ega-fr-2026-10-23-r16/EGA_French_Current_Cumulative_Reader.pdf\#page=84}
\expandafter\def\csname EGAExternal@I.4.1.1\endcsname{https://github.com/KokunoYumeto/ega-fr/releases/download/ega-fr-2026-10-23-r16/EGA_French_Current_Cumulative_Reader.pdf\#page=84}
\expandafter\def\csname EGAExternal@I.4.1.10\endcsname{https://github.com/KokunoYumeto/ega-fr/releases/download/ega-fr-2026-10-23-r16/EGA_French_Current_Cumulative_Reader.pdf\#page=84}
\expandafter\def\csname EGAExternal@I.4.1.9\endcsname{https://github.com/KokunoYumeto/ega-fr/releases/download/ega-fr-2026-10-23-r16/EGA_French_Current_Cumulative_Reader.pdf\#page=84}
\expandafter\def\csname EGAExternal@I.4.2.3\endcsname{https://github.com/KokunoYumeto/ega-fr/releases/download/ega-fr-2026-10-23-r16/EGA_French_Current_Cumulative_Reader.pdf\#page=87}
\expandafter\def\csname EGAExternal@I.4.2.4\endcsname{https://github.com/KokunoYumeto/ega-fr/releases/download/ega-fr-2026-10-23-r16/EGA_French_Current_Cumulative_Reader.pdf\#page=87}
\expandafter\def\csname EGAExternal@I.4.3.2\endcsname{https://github.com/KokunoYumeto/ega-fr/releases/download/ega-fr-2026-10-23-r16/EGA_French_Current_Cumulative_Reader.pdf\#page=88}
\expandafter\def\csname EGAExternal@I.4.4.1\endcsname{https://github.com/KokunoYumeto/ega-fr/releases/download/ega-fr-2026-10-23-r16/EGA_French_Current_Cumulative_Reader.pdf\#page=88}
\expandafter\def\csname EGAExternal@I.5.1.4\endcsname{https://github.com/KokunoYumeto/ega-fr/releases/download/ega-fr-2026-10-23-r16/EGA_French_Current_Cumulative_Reader.pdf\#page=90}
\expandafter\def\csname EGAExternal@I.5.1.8\endcsname{https://github.com/KokunoYumeto/ega-fr/releases/download/ega-fr-2026-10-23-r16/EGA_French_Current_Cumulative_Reader.pdf\#page=90}
\expandafter\def\csname EGAExternal@I.5.2.1\endcsname{https://github.com/KokunoYumeto/ega-fr/releases/download/ega-fr-2026-10-23-r16/EGA_French_Current_Cumulative_Reader.pdf\#page=92}
\expandafter\def\csname EGAExternal@I.5.2.2\endcsname{https://github.com/KokunoYumeto/ega-fr/releases/download/ega-fr-2026-10-23-r16/EGA_French_Current_Cumulative_Reader.pdf\#page=92}
\expandafter\def\csname EGAExternal@I.5.2.3\endcsname{https://github.com/KokunoYumeto/ega-fr/releases/download/ega-fr-2026-10-23-r16/EGA_French_Current_Cumulative_Reader.pdf\#page=92}
\expandafter\def\csname EGAExternal@I.5.3.4\endcsname{https://github.com/KokunoYumeto/ega-fr/releases/download/ega-fr-2026-10-23-r16/EGA_French_Current_Cumulative_Reader.pdf\#page=93}
\expandafter\def\csname EGAExternal@I.5.3.8\endcsname{https://github.com/KokunoYumeto/ega-fr/releases/download/ega-fr-2026-10-23-r16/EGA_French_Current_Cumulative_Reader.pdf\#page=93}
\expandafter\def\csname EGAExternal@I.5.4\endcsname{https://github.com/KokunoYumeto/ega-fr/releases/download/ega-fr-2026-10-23-r16/EGA_French_Current_Cumulative_Reader.pdf\#page=95}
\expandafter\def\csname EGAExternal@I.5.4.1\endcsname{https://github.com/KokunoYumeto/ega-fr/releases/download/ega-fr-2026-10-23-r16/EGA_French_Current_Cumulative_Reader.pdf\#page=95}
\expandafter\def\csname EGAExternal@I.5.5.1\endcsname{https://github.com/KokunoYumeto/ega-fr/releases/download/ega-fr-2026-10-23-r16/EGA_French_Current_Cumulative_Reader.pdf\#page=96}
\expandafter\def\csname EGAExternal@I.5.5.11\endcsname{https://github.com/KokunoYumeto/ega-fr/releases/download/ega-fr-2026-10-23-r16/EGA_French_Current_Cumulative_Reader.pdf\#page=96}
\expandafter\def\csname EGAExternal@I.6.1.10\endcsname{https://github.com/KokunoYumeto/ega-fr/releases/download/ega-fr-2026-10-23-r16/EGA_French_Current_Cumulative_Reader.pdf\#page=99}
\expandafter\def\csname EGAExternal@I.6.2\endcsname{https://github.com/KokunoYumeto/ega-fr/releases/download/ega-fr-2026-10-23-r16/EGA_French_Current_Cumulative_Reader.pdf\#page=101}
\expandafter\def\csname EGAExternal@I.6.2.2\endcsname{https://github.com/KokunoYumeto/ega-fr/releases/download/ega-fr-2026-10-23-r16/EGA_French_Current_Cumulative_Reader.pdf\#page=101}
\expandafter\def\csname EGAExternal@I.6.3.5\endcsname{https://github.com/KokunoYumeto/ega-fr/releases/download/ega-fr-2026-10-23-r16/EGA_French_Current_Cumulative_Reader.pdf\#page=101}
\expandafter\def\csname EGAExternal@I.6.4.12\endcsname{https://github.com/KokunoYumeto/ega-fr/releases/download/ega-fr-2026-10-23-r16/EGA_French_Current_Cumulative_Reader.pdf\#page=104}
\expandafter\def\csname EGAExternal@I.6.4.2\endcsname{https://github.com/KokunoYumeto/ega-fr/releases/download/ega-fr-2026-10-23-r16/EGA_French_Current_Cumulative_Reader.pdf\#page=103}
\expandafter\def\csname EGAExternal@I.6.4.4\endcsname{https://github.com/KokunoYumeto/ega-fr/releases/download/ega-fr-2026-10-23-r16/EGA_French_Current_Cumulative_Reader.pdf\#page=103}
\expandafter\def\csname EGAExternal@I.6.4.8\endcsname{https://github.com/KokunoYumeto/ega-fr/releases/download/ega-fr-2026-10-23-r16/EGA_French_Current_Cumulative_Reader.pdf\#page=104}
\expandafter\def\csname EGAExternal@I.6.5.1\endcsname{https://github.com/KokunoYumeto/ega-fr/releases/download/ega-fr-2026-10-23-r16/EGA_French_Current_Cumulative_Reader.pdf\#page=105}
\expandafter\def\csname EGAExternal@I.6.5.4\endcsname{https://github.com/KokunoYumeto/ega-fr/releases/download/ega-fr-2026-10-23-r16/EGA_French_Current_Cumulative_Reader.pdf\#page=105}
\expandafter\def\csname EGAExternal@I.7.1.2\endcsname{https://github.com/KokunoYumeto/ega-fr/releases/download/ega-fr-2026-10-23-r16/EGA_French_Current_Cumulative_Reader.pdf\#page=109}
\expandafter\def\csname EGAExternal@I.7.2.2\endcsname{https://github.com/KokunoYumeto/ega-fr/releases/download/ega-fr-2026-10-23-r16/EGA_French_Current_Cumulative_Reader.pdf\#page=111}
\expandafter\def\csname EGAExternal@I.7.4.6\endcsname{https://github.com/KokunoYumeto/ega-fr/releases/download/ega-fr-2026-10-23-r16/EGA_French_Current_Cumulative_Reader.pdf\#page=114}
\expandafter\def\csname EGAExternal@I.8.1.4\endcsname{https://github.com/KokunoYumeto/ega-fr/releases/download/ega-fr-2026-10-23-r16/EGA_French_Current_Cumulative_Reader.pdf\#page=115}
\expandafter\def\csname EGAExternal@I.8.2.2\endcsname{https://github.com/KokunoYumeto/ega-fr/releases/download/ega-fr-2026-10-23-r16/EGA_French_Current_Cumulative_Reader.pdf\#page=116}
\expandafter\def\csname EGAExternal@I.8.2.8\endcsname{https://github.com/KokunoYumeto/ega-fr/releases/download/ega-fr-2026-10-23-r16/EGA_French_Current_Cumulative_Reader.pdf\#page=116}
\expandafter\def\csname EGAExternal@I.9.1.12\endcsname{https://github.com/KokunoYumeto/ega-fr/releases/download/ega-fr-2026-10-23-r16/EGA_French_Current_Cumulative_Reader.pdf\#page=119}
\expandafter\def\csname EGAExternal@I.9.1.13\endcsname{https://github.com/KokunoYumeto/ega-fr/releases/download/ega-fr-2026-10-23-r16/EGA_French_Current_Cumulative_Reader.pdf\#page=119}
\expandafter\def\csname EGAExternal@I.9.1.13.1\endcsname{https://github.com/KokunoYumeto/ega-fr/releases/download/ega-fr-2026-10-23-r16/EGA_French_Current_Cumulative_Reader.pdf\#page=119}
\expandafter\def\csname EGAExternal@I.9.2.2\endcsname{https://github.com/KokunoYumeto/ega-fr/releases/download/ega-fr-2026-10-23-r16/EGA_French_Current_Cumulative_Reader.pdf\#page=120}
\expandafter\def\csname EGAExternal@I.9.3.5\endcsname{https://github.com/KokunoYumeto/ega-fr/releases/download/ega-fr-2026-10-23-r16/EGA_French_Current_Cumulative_Reader.pdf\#page=120}
\expandafter\def\csname EGAExternal@I.9.4.7\endcsname{https://github.com/KokunoYumeto/ega-fr/releases/download/ega-fr-2026-10-23-r16/EGA_French_Current_Cumulative_Reader.pdf\#page=121}
\expandafter\def\csname EGAExternal@I.9.4.9\endcsname{https://github.com/KokunoYumeto/ega-fr/releases/download/ega-fr-2026-10-23-r16/EGA_French_Current_Cumulative_Reader.pdf\#page=121}
\expandafter\def\csname EGAExternal@I.9.5.1\endcsname{https://github.com/KokunoYumeto/ega-fr/releases/download/ega-fr-2026-10-23-r16/EGA_French_Current_Cumulative_Reader.pdf\#page=123}
\expandafter\def\csname EGAExternal@I.9.5.10\endcsname{https://github.com/KokunoYumeto/ega-fr/releases/download/ega-fr-2026-10-23-r16/EGA_French_Current_Cumulative_Reader.pdf\#page=123}
\expandafter\def\csname EGAExternal@I.9.5.2\endcsname{https://github.com/KokunoYumeto/ega-fr/releases/download/ega-fr-2026-10-23-r16/EGA_French_Current_Cumulative_Reader.pdf\#page=123}
\expandafter\def\csname EGAExternal@I.9.5.3\endcsname{https://github.com/KokunoYumeto/ega-fr/releases/download/ega-fr-2026-10-23-r16/EGA_French_Current_Cumulative_Reader.pdf\#page=123}
\expandafter\def\csname EGAExternal@I.9.6.6\endcsname{https://github.com/KokunoYumeto/ega-fr/releases/download/ega-fr-2026-10-23-r16/EGA_French_Current_Cumulative_Reader.pdf\#page=125}
\expandafter\def\csname EGAExternal@II.1.1.2\endcsname{https://github.com/KokunoYumeto/ega-fr/releases/download/ega-fr-2026-10-23-r16/EGA_French_Current_Cumulative_Reader.pdf\#page=150}
\expandafter\def\csname EGAExternal@II.1.2.7\endcsname{https://github.com/KokunoYumeto/ega-fr/releases/download/ega-fr-2026-10-23-r16/EGA_French_Current_Cumulative_Reader.pdf\#page=151}
\expandafter\def\csname EGAExternal@II.1.3.1\endcsname{https://github.com/KokunoYumeto/ega-fr/releases/download/ega-fr-2026-10-23-r16/EGA_French_Current_Cumulative_Reader.pdf\#page=152}
\expandafter\def\csname EGAExternal@II.1.3.6\endcsname{https://github.com/KokunoYumeto/ega-fr/releases/download/ega-fr-2026-10-23-r16/EGA_French_Current_Cumulative_Reader.pdf\#page=152}
\expandafter\def\csname EGAExternal@II.1.4.10\endcsname{https://github.com/KokunoYumeto/ega-fr/releases/download/ega-fr-2026-10-23-r16/EGA_French_Current_Cumulative_Reader.pdf\#page=154}
\expandafter\def\csname EGAExternal@II.1.5.5\endcsname{https://github.com/KokunoYumeto/ega-fr/releases/download/ega-fr-2026-10-23-r16/EGA_French_Current_Cumulative_Reader.pdf\#page=154}
\expandafter\def\csname EGAExternal@II.1.6.2\endcsname{https://github.com/KokunoYumeto/ega-fr/releases/download/ega-fr-2026-10-23-r16/EGA_French_Current_Cumulative_Reader.pdf\#page=156}
\expandafter\def\csname EGAExternal@II.1.7.9\endcsname{https://github.com/KokunoYumeto/ega-fr/releases/download/ega-fr-2026-10-23-r16/EGA_French_Current_Cumulative_Reader.pdf\#page=156}
\expandafter\def\csname EGAExternal@II.2.8.10\endcsname{https://github.com/KokunoYumeto/ega-fr/releases/download/ega-fr-2026-10-23-r16/EGA_French_Current_Cumulative_Reader.pdf\#page=178}
\expandafter\def\csname EGAExternal@II.3.1.7\endcsname{https://github.com/KokunoYumeto/ega-fr/releases/download/ega-fr-2026-10-23-r16/EGA_French_Current_Cumulative_Reader.pdf\#page=182}
\expandafter\def\csname EGAExternal@II.3.5.3\endcsname{https://github.com/KokunoYumeto/ega-fr/releases/download/ega-fr-2026-10-23-r16/EGA_French_Current_Cumulative_Reader.pdf\#page=190}
\expandafter\def\csname EGAExternal@II.3.8.4\endcsname{https://github.com/KokunoYumeto/ega-fr/releases/download/ega-fr-2026-10-23-r16/EGA_French_Current_Cumulative_Reader.pdf\#page=194}
\expandafter\def\csname EGAExternal@II.4.1.1\endcsname{https://github.com/KokunoYumeto/ega-fr/releases/download/ega-fr-2026-10-23-r16/EGA_French_Current_Cumulative_Reader.pdf\#page=196}
\expandafter\def\csname EGAExternal@II.4.1.2\endcsname{https://github.com/KokunoYumeto/ega-fr/releases/download/ega-fr-2026-10-23-r16/EGA_French_Current_Cumulative_Reader.pdf\#page=196}
\expandafter\def\csname EGAExternal@II.4.1.3\endcsname{https://github.com/KokunoYumeto/ega-fr/releases/download/ega-fr-2026-10-23-r16/EGA_French_Current_Cumulative_Reader.pdf\#page=196}
\expandafter\def\csname EGAExternal@II.4.1.3.2\endcsname{https://github.com/KokunoYumeto/ega-fr/releases/download/ega-fr-2026-10-23-r16/EGA_French_Current_Cumulative_Reader.pdf\#page=196}
\expandafter\def\csname EGAExternal@II.4.4.1\endcsname{https://github.com/KokunoYumeto/ega-fr/releases/download/ega-fr-2026-10-23-r16/EGA_French_Current_Cumulative_Reader.pdf\#page=201}
\expandafter\def\csname EGAExternal@II.4.4.10\endcsname{https://github.com/KokunoYumeto/ega-fr/releases/download/ega-fr-2026-10-23-r16/EGA_French_Current_Cumulative_Reader.pdf\#page=201}
\expandafter\def\csname EGAExternal@II.4.4.2\endcsname{https://github.com/KokunoYumeto/ega-fr/releases/download/ega-fr-2026-10-23-r16/EGA_French_Current_Cumulative_Reader.pdf\#page=201}
\expandafter\def\csname EGAExternal@II.4.4.4\endcsname{https://github.com/KokunoYumeto/ega-fr/releases/download/ega-fr-2026-10-23-r16/EGA_French_Current_Cumulative_Reader.pdf\#page=201}
\expandafter\def\csname EGAExternal@II.4.4.5\endcsname{https://github.com/KokunoYumeto/ega-fr/releases/download/ega-fr-2026-10-23-r16/EGA_French_Current_Cumulative_Reader.pdf\#page=201}
\expandafter\def\csname EGAExternal@II.4.5.3\endcsname{https://github.com/KokunoYumeto/ega-fr/releases/download/ega-fr-2026-10-23-r16/EGA_French_Current_Cumulative_Reader.pdf\#page=204}
\expandafter\def\csname EGAExternal@II.4.6.11\endcsname{https://github.com/KokunoYumeto/ega-fr/releases/download/ega-fr-2026-10-23-r16/EGA_French_Current_Cumulative_Reader.pdf\#page=209}
\expandafter\def\csname EGAExternal@II.4.6.13\endcsname{https://github.com/KokunoYumeto/ega-fr/releases/download/ega-fr-2026-10-23-r16/EGA_French_Current_Cumulative_Reader.pdf\#page=209}
\expandafter\def\csname EGAExternal@II.4.6.17\endcsname{https://github.com/KokunoYumeto/ega-fr/releases/download/ega-fr-2026-10-23-r16/EGA_French_Current_Cumulative_Reader.pdf\#page=209}
\expandafter\def\csname EGAExternal@II.5.1.1\endcsname{https://github.com/KokunoYumeto/ega-fr/releases/download/ega-fr-2026-10-23-r16/EGA_French_Current_Cumulative_Reader.pdf\#page=212}
\expandafter\def\csname EGAExternal@II.5.1.10\endcsname{https://github.com/KokunoYumeto/ega-fr/releases/download/ega-fr-2026-10-23-r16/EGA_French_Current_Cumulative_Reader.pdf\#page=212}
\expandafter\def\csname EGAExternal@II.5.1.2\endcsname{https://github.com/KokunoYumeto/ega-fr/releases/download/ega-fr-2026-10-23-r16/EGA_French_Current_Cumulative_Reader.pdf\#page=212}
\expandafter\def\csname EGAExternal@II.5.1.9\endcsname{https://github.com/KokunoYumeto/ega-fr/releases/download/ega-fr-2026-10-23-r16/EGA_French_Current_Cumulative_Reader.pdf\#page=212}
\expandafter\def\csname EGAExternal@II.5.3.2\endcsname{https://github.com/KokunoYumeto/ega-fr/releases/download/ega-fr-2026-10-23-r16/EGA_French_Current_Cumulative_Reader.pdf\#page=215}
\expandafter\def\csname EGAExternal@II.5.3.3\endcsname{https://github.com/KokunoYumeto/ega-fr/releases/download/ega-fr-2026-10-23-r16/EGA_French_Current_Cumulative_Reader.pdf\#page=215}
\expandafter\def\csname EGAExternal@II.5.3.4\endcsname{https://github.com/KokunoYumeto/ega-fr/releases/download/ega-fr-2026-10-23-r16/EGA_French_Current_Cumulative_Reader.pdf\#page=215}
\expandafter\def\csname EGAExternal@II.5.4.1\endcsname{https://github.com/KokunoYumeto/ega-fr/releases/download/ega-fr-2026-10-23-r16/EGA_French_Current_Cumulative_Reader.pdf\#page=216}
\expandafter\def\csname EGAExternal@II.5.4.10\endcsname{https://github.com/KokunoYumeto/ega-fr/releases/download/ega-fr-2026-10-23-r16/EGA_French_Current_Cumulative_Reader.pdf\#page=216}
\expandafter\def\csname EGAExternal@II.5.4.2\endcsname{https://github.com/KokunoYumeto/ega-fr/releases/download/ega-fr-2026-10-23-r16/EGA_French_Current_Cumulative_Reader.pdf\#page=216}
\expandafter\def\csname EGAExternal@II.5.4.3\endcsname{https://github.com/KokunoYumeto/ega-fr/releases/download/ega-fr-2026-10-23-r16/EGA_French_Current_Cumulative_Reader.pdf\#page=216}
\expandafter\def\csname EGAExternal@II.5.4.5\endcsname{https://github.com/KokunoYumeto/ega-fr/releases/download/ega-fr-2026-10-23-r16/EGA_French_Current_Cumulative_Reader.pdf\#page=216}
\expandafter\def\csname EGAExternal@II.5.4.6\endcsname{https://github.com/KokunoYumeto/ega-fr/releases/download/ega-fr-2026-10-23-r16/EGA_French_Current_Cumulative_Reader.pdf\#page=216}
\expandafter\def\csname EGAExternal@II.5.5.3\endcsname{https://github.com/KokunoYumeto/ega-fr/releases/download/ega-fr-2026-10-23-r16/EGA_French_Current_Cumulative_Reader.pdf\#page=218}
\expandafter\def\csname EGAExternal@II.5.5.5\endcsname{https://github.com/KokunoYumeto/ega-fr/releases/download/ega-fr-2026-10-23-r16/EGA_French_Current_Cumulative_Reader.pdf\#page=218}
\expandafter\def\csname EGAExternal@II.5.6.1\endcsname{https://github.com/KokunoYumeto/ega-fr/releases/download/ega-fr-2026-10-23-r16/EGA_French_Current_Cumulative_Reader.pdf\#page=220}
\expandafter\def\csname EGAExternal@II.5.6.3\endcsname{https://github.com/KokunoYumeto/ega-fr/releases/download/ega-fr-2026-10-23-r16/EGA_French_Current_Cumulative_Reader.pdf\#page=221}
\expandafter\def\csname EGAExternal@II.6.1.10\endcsname{https://github.com/KokunoYumeto/ega-fr/releases/download/ega-fr-2026-10-23-r16/EGA_French_Current_Cumulative_Reader.pdf\#page=222}
\expandafter\def\csname EGAExternal@II.6.1.2\endcsname{https://github.com/KokunoYumeto/ega-fr/releases/download/ega-fr-2026-10-23-r16/EGA_French_Current_Cumulative_Reader.pdf\#page=222}
\expandafter\def\csname EGAExternal@II.6.1.3\endcsname{https://github.com/KokunoYumeto/ega-fr/releases/download/ega-fr-2026-10-23-r16/EGA_French_Current_Cumulative_Reader.pdf\#page=222}
\expandafter\def\csname EGAExternal@II.6.1.5\endcsname{https://github.com/KokunoYumeto/ega-fr/releases/download/ega-fr-2026-10-23-r16/EGA_French_Current_Cumulative_Reader.pdf\#page=222}
\expandafter\def\csname EGAExternal@II.6.2.2\endcsname{https://github.com/KokunoYumeto/ega-fr/releases/download/ega-fr-2026-10-23-r16/EGA_French_Current_Cumulative_Reader.pdf\#page=225}
\expandafter\def\csname EGAExternal@II.6.2.3\endcsname{https://github.com/KokunoYumeto/ega-fr/releases/download/ega-fr-2026-10-23-r16/EGA_French_Current_Cumulative_Reader.pdf\#page=225}
\expandafter\def\csname EGAExternal@II.6.2.4\endcsname{https://github.com/KokunoYumeto/ega-fr/releases/download/ega-fr-2026-10-23-r16/EGA_French_Current_Cumulative_Reader.pdf\#page=225}
\expandafter\def\csname EGAExternal@II.6.2.5\endcsname{https://github.com/KokunoYumeto/ega-fr/releases/download/ega-fr-2026-10-23-r16/EGA_French_Current_Cumulative_Reader.pdf\#page=225}
\expandafter\def\csname EGAExternal@II.6.2.6\endcsname{https://github.com/KokunoYumeto/ega-fr/releases/download/ega-fr-2026-10-23-r16/EGA_French_Current_Cumulative_Reader.pdf\#page=225}
\expandafter\def\csname EGAExternal@II.6.3.4\endcsname{https://github.com/KokunoYumeto/ega-fr/releases/download/ega-fr-2026-10-23-r16/EGA_French_Current_Cumulative_Reader.pdf\#page=227}
\expandafter\def\csname EGAExternal@II.6.3.8\endcsname{https://github.com/KokunoYumeto/ega-fr/releases/download/ega-fr-2026-10-23-r16/EGA_French_Current_Cumulative_Reader.pdf\#page=228}
\expandafter\def\csname EGAExternal@II.6.3.9\endcsname{https://github.com/KokunoYumeto/ega-fr/releases/download/ega-fr-2026-10-23-r16/EGA_French_Current_Cumulative_Reader.pdf\#page=228}
\expandafter\def\csname EGAExternal@II.6.7.1\endcsname{https://github.com/KokunoYumeto/ega-fr/releases/download/ega-fr-2026-10-23-r16/EGA_French_Current_Cumulative_Reader.pdf\#page=239}
\expandafter\def\csname EGAExternal@II.7.1.1\endcsname{https://github.com/KokunoYumeto/ega-fr/releases/download/ega-fr-2026-10-23-r16/EGA_French_Current_Cumulative_Reader.pdf\#page=241}
\expandafter\def\csname EGAExternal@II.7.1.2\endcsname{https://github.com/KokunoYumeto/ega-fr/releases/download/ega-fr-2026-10-23-r16/EGA_French_Current_Cumulative_Reader.pdf\#page=241}
\expandafter\def\csname EGAExternal@II.7.1.4\endcsname{https://github.com/KokunoYumeto/ega-fr/releases/download/ega-fr-2026-10-23-r16/EGA_French_Current_Cumulative_Reader.pdf\#page=241}
\expandafter\def\csname EGAExternal@II.7.1.6\endcsname{https://github.com/KokunoYumeto/ega-fr/releases/download/ega-fr-2026-10-23-r16/EGA_French_Current_Cumulative_Reader.pdf\#page=241}
\expandafter\def\csname EGAExternal@II.7.1.7\endcsname{https://github.com/KokunoYumeto/ega-fr/releases/download/ega-fr-2026-10-23-r16/EGA_French_Current_Cumulative_Reader.pdf\#page=241}
\expandafter\def\csname EGAExternal@II.7.1.9\endcsname{https://github.com/KokunoYumeto/ega-fr/releases/download/ega-fr-2026-10-23-r16/EGA_French_Current_Cumulative_Reader.pdf\#page=241}
\expandafter\def\csname EGAExternal@II.7.2.1\endcsname{https://github.com/KokunoYumeto/ega-fr/releases/download/ega-fr-2026-10-23-r16/EGA_French_Current_Cumulative_Reader.pdf\#page=243}
\expandafter\def\csname EGAExternal@II.7.2.3\endcsname{https://github.com/KokunoYumeto/ega-fr/releases/download/ega-fr-2026-10-23-r16/EGA_French_Current_Cumulative_Reader.pdf\#page=243}
\expandafter\def\csname EGAExternal@II.7.2.4\endcsname{https://github.com/KokunoYumeto/ega-fr/releases/download/ega-fr-2026-10-23-r16/EGA_French_Current_Cumulative_Reader.pdf\#page=243}
\expandafter\def\csname EGAExternal@II.7.3.10\endcsname{https://github.com/KokunoYumeto/ega-fr/releases/download/ega-fr-2026-10-23-r16/EGA_French_Current_Cumulative_Reader.pdf\#page=244}
\expandafter\def\csname EGAExternal@II.7.3.3\endcsname{https://github.com/KokunoYumeto/ega-fr/releases/download/ega-fr-2026-10-23-r16/EGA_French_Current_Cumulative_Reader.pdf\#page=244}
\expandafter\def\csname EGAExternal@II.7.3.8\endcsname{https://github.com/KokunoYumeto/ega-fr/releases/download/ega-fr-2026-10-23-r16/EGA_French_Current_Cumulative_Reader.pdf\#page=244}
\expandafter\def\csname EGAExternal@II.7.3.9\endcsname{https://github.com/KokunoYumeto/ega-fr/releases/download/ega-fr-2026-10-23-r16/EGA_French_Current_Cumulative_Reader.pdf\#page=244}
\expandafter\def\csname EGAExternal@II.8.1\endcsname{https://github.com/KokunoYumeto/ega-fr/releases/download/ega-fr-2026-10-23-r16/EGA_French_Current_Cumulative_Reader.pdf\#page=250}
\expandafter\def\csname EGAExternal@II.8.1.3\endcsname{https://github.com/KokunoYumeto/ega-fr/releases/download/ega-fr-2026-10-23-r16/EGA_French_Current_Cumulative_Reader.pdf\#page=250}
\expandafter\def\csname EGAExternal@II.8.3.9\endcsname{https://github.com/KokunoYumeto/ega-fr/releases/download/ega-fr-2026-10-23-r16/EGA_French_Current_Cumulative_Reader.pdf\#page=259}
\expandafter\def\csname EGAExternal@II.8.4.4\endcsname{https://github.com/KokunoYumeto/ega-fr/releases/download/ega-fr-2026-10-23-r16/EGA_French_Current_Cumulative_Reader.pdf\#page=260}
\expandafter\def\csname EGAExternal@III.1.4.12\endcsname{https://github.com/KokunoYumeto/ega-fr/releases/download/ega-fr-2026-10-23-r16/EGA_French_Current_Cumulative_Reader.pdf\#page=362}
\expandafter\def\csname EGAExternal@III.2.6.2\endcsname{https://github.com/KokunoYumeto/ega-fr/releases/download/ega-fr-2026-10-23-r16/EGA_French_Current_Cumulative_Reader.pdf\#page=377}
\expandafter\def\csname EGAExternal@III.3.2.1\endcsname{https://github.com/KokunoYumeto/ega-fr/releases/download/ega-fr-2026-10-23-r16/EGA_French_Current_Cumulative_Reader.pdf\#page=381}
\expandafter\def\csname EGAExternal@III.4.3\endcsname{https://github.com/KokunoYumeto/ega-fr/releases/download/ega-fr-2026-10-23-r16/EGA_French_Current_Cumulative_Reader.pdf\#page=391}
\expandafter\def\csname EGAExternal@III.4.3.10\endcsname{https://github.com/KokunoYumeto/ega-fr/releases/download/ega-fr-2026-10-23-r16/EGA_French_Current_Cumulative_Reader.pdf\#page=391}
\expandafter\def\csname EGAExternal@III.4.3.3\endcsname{https://github.com/KokunoYumeto/ega-fr/releases/download/ega-fr-2026-10-23-r16/EGA_French_Current_Cumulative_Reader.pdf\#page=391}
\expandafter\def\csname EGAExternal@III.4.3.4\endcsname{https://github.com/KokunoYumeto/ega-fr/releases/download/ega-fr-2026-10-23-r16/EGA_French_Current_Cumulative_Reader.pdf\#page=391}
\expandafter\def\csname EGAExternal@III.4.3.5\endcsname{https://github.com/KokunoYumeto/ega-fr/releases/download/ega-fr-2026-10-23-r16/EGA_French_Current_Cumulative_Reader.pdf\#page=391}
\expandafter\def\csname EGAExternal@III.4.3.7\endcsname{https://github.com/KokunoYumeto/ega-fr/releases/download/ega-fr-2026-10-23-r16/EGA_French_Current_Cumulative_Reader.pdf\#page=391}
\expandafter\def\csname EGAExternal@III.4.3.9\endcsname{https://github.com/KokunoYumeto/ega-fr/releases/download/ega-fr-2026-10-23-r16/EGA_French_Current_Cumulative_Reader.pdf\#page=391}
\expandafter\def\csname EGAExternal@III.4.4.10\endcsname{https://github.com/KokunoYumeto/ega-fr/releases/download/ega-fr-2026-10-23-r16/EGA_French_Current_Cumulative_Reader.pdf\#page=394}
\expandafter\def\csname EGAExternal@III.4.4.2\endcsname{https://github.com/KokunoYumeto/ega-fr/releases/download/ega-fr-2026-10-23-r16/EGA_French_Current_Cumulative_Reader.pdf\#page=394}
\expandafter\def\csname EGAExternal@III.4.4.3\endcsname{https://github.com/KokunoYumeto/ega-fr/releases/download/ega-fr-2026-10-23-r16/EGA_French_Current_Cumulative_Reader.pdf\#page=394}
\expandafter\def\csname EGAExternal@III.4.7.1\endcsname{https://github.com/KokunoYumeto/ega-fr/releases/download/ega-fr-2026-10-23-r16/EGA_French_Current_Cumulative_Reader.pdf\#page=401}
\expandafter\def\csname EGAExternal@III.5.5.2\endcsname{https://github.com/KokunoYumeto/ega-fr/releases/download/ega-fr-2026-10-23-r16/EGA_French_Current_Cumulative_Reader.pdf\#page=412}
\expandafter\def\csname EGAExternal@III.6.2.2\endcsname{https://github.com/KokunoYumeto/ega-fr/releases/download/ega-fr-2026-10-23-r16/EGA_French_Current_Cumulative_Reader.pdf\#page=378}
\expandafter\def\csname EGAExternal@III.6.2.3\endcsname{https://github.com/KokunoYumeto/ega-fr/releases/download/ega-fr-2026-10-23-r16/EGA_French_Current_Cumulative_Reader.pdf\#page=378}
\expandafter\def\csname EGAExternal@III.6.9.9\endcsname{https://github.com/KokunoYumeto/ega-fr/releases/download/ega-fr-2026-10-23-r16/EGA_French_Current_Cumulative_Reader.pdf\#page=446}
\expandafter\def\csname EGAExternal@III.7.4.7.2\endcsname{https://github.com/KokunoYumeto/ega-fr/releases/download/ega-fr-2026-10-23-r16/EGA_French_Current_Cumulative_Reader.pdf\#page=462}
\expandafter\def\csname EGAExternal@III.7.8.6\endcsname{https://github.com/KokunoYumeto/ega-fr/releases/download/ega-fr-2026-10-23-r16/EGA_French_Current_Cumulative_Reader.pdf\#page=476}
\expandafter\def\csname EGAExternal@III.7.8.7\endcsname{https://github.com/KokunoYumeto/ega-fr/releases/download/ega-fr-2026-10-23-r16/EGA_French_Current_Cumulative_Reader.pdf\#page=477}
\expandafter\def\csname EGAExternal@IV.1.1.2\endcsname{https://github.com/KokunoYumeto/ega-fr/releases/download/ega-fr-2026-10-23-r16/EGA_French_Current_Cumulative_Reader.pdf\#page=648}
\expandafter\def\csname EGAExternal@IV.1.10.1\endcsname{https://github.com/KokunoYumeto/ega-fr/releases/download/ega-fr-2026-10-23-r16/EGA_French_Current_Cumulative_Reader.pdf\#page=664}
\expandafter\def\csname EGAExternal@IV.1.10.2\endcsname{https://github.com/KokunoYumeto/ega-fr/releases/download/ega-fr-2026-10-23-r16/EGA_French_Current_Cumulative_Reader.pdf\#page=664}
\expandafter\def\csname EGAExternal@IV.1.10.3\endcsname{https://github.com/KokunoYumeto/ega-fr/releases/download/ega-fr-2026-10-23-r16/EGA_French_Current_Cumulative_Reader.pdf\#page=664}
\expandafter\def\csname EGAExternal@IV.1.10.4\endcsname{https://github.com/KokunoYumeto/ega-fr/releases/download/ega-fr-2026-10-23-r16/EGA_French_Current_Cumulative_Reader.pdf\#page=664}
\expandafter\def\csname EGAExternal@IV.1.2.2\endcsname{https://github.com/KokunoYumeto/ega-fr/releases/download/ega-fr-2026-10-23-r16/EGA_French_Current_Cumulative_Reader.pdf\#page=649}
\expandafter\def\csname EGAExternal@IV.1.2.3\endcsname{https://github.com/KokunoYumeto/ega-fr/releases/download/ega-fr-2026-10-23-r16/EGA_French_Current_Cumulative_Reader.pdf\#page=649}
\expandafter\def\csname EGAExternal@IV.1.2.4\endcsname{https://github.com/KokunoYumeto/ega-fr/releases/download/ega-fr-2026-10-23-r16/EGA_French_Current_Cumulative_Reader.pdf\#page=649}
\expandafter\def\csname EGAExternal@IV.1.2.6\endcsname{https://github.com/KokunoYumeto/ega-fr/releases/download/ega-fr-2026-10-23-r16/EGA_French_Current_Cumulative_Reader.pdf\#page=649}
\expandafter\def\csname EGAExternal@IV.1.2.7\endcsname{https://github.com/KokunoYumeto/ega-fr/releases/download/ega-fr-2026-10-23-r16/EGA_French_Current_Cumulative_Reader.pdf\#page=649}
\expandafter\def\csname EGAExternal@IV.1.2.8\endcsname{https://github.com/KokunoYumeto/ega-fr/releases/download/ega-fr-2026-10-23-r16/EGA_French_Current_Cumulative_Reader.pdf\#page=649}
\expandafter\def\csname EGAExternal@IV.1.3.4\endcsname{https://github.com/KokunoYumeto/ega-fr/releases/download/ega-fr-2026-10-23-r16/EGA_French_Current_Cumulative_Reader.pdf\#page=650}
\expandafter\def\csname EGAExternal@IV.1.3.8\endcsname{https://github.com/KokunoYumeto/ega-fr/releases/download/ega-fr-2026-10-23-r16/EGA_French_Current_Cumulative_Reader.pdf\#page=650}
\expandafter\def\csname EGAExternal@IV.1.4.3\endcsname{https://github.com/KokunoYumeto/ega-fr/releases/download/ega-fr-2026-10-23-r16/EGA_French_Current_Cumulative_Reader.pdf\#page=652}
\expandafter\def\csname EGAExternal@IV.1.4.5\endcsname{https://github.com/KokunoYumeto/ega-fr/releases/download/ega-fr-2026-10-23-r16/EGA_French_Current_Cumulative_Reader.pdf\#page=652}
\expandafter\def\csname EGAExternal@IV.1.4.6\endcsname{https://github.com/KokunoYumeto/ega-fr/releases/download/ega-fr-2026-10-23-r16/EGA_French_Current_Cumulative_Reader.pdf\#page=653}
\expandafter\def\csname EGAExternal@IV.1.4.7\endcsname{https://github.com/KokunoYumeto/ega-fr/releases/download/ega-fr-2026-10-23-r16/EGA_French_Current_Cumulative_Reader.pdf\#page=652}
\expandafter\def\csname EGAExternal@IV.1.5.4\endcsname{https://github.com/KokunoYumeto/ega-fr/releases/download/ega-fr-2026-10-23-r16/EGA_French_Current_Cumulative_Reader.pdf\#page=654}
\expandafter\def\csname EGAExternal@IV.1.6.1\endcsname{https://github.com/KokunoYumeto/ega-fr/releases/download/ega-fr-2026-10-23-r16/EGA_French_Current_Cumulative_Reader.pdf\#page=655}
\expandafter\def\csname EGAExternal@IV.1.6.2\endcsname{https://github.com/KokunoYumeto/ega-fr/releases/download/ega-fr-2026-10-23-r16/EGA_French_Current_Cumulative_Reader.pdf\#page=655}
\expandafter\def\csname EGAExternal@IV.1.7.19\endcsname{https://github.com/KokunoYumeto/ega-fr/releases/download/ega-fr-2026-10-23-r16/EGA_French_Current_Cumulative_Reader.pdf\#page=656}
\expandafter\def\csname EGAExternal@IV.1.7.4\endcsname{https://github.com/KokunoYumeto/ega-fr/releases/download/ega-fr-2026-10-23-r16/EGA_French_Current_Cumulative_Reader.pdf\#page=656}
\expandafter\def\csname EGAExternal@IV.1.7.5\endcsname{https://github.com/KokunoYumeto/ega-fr/releases/download/ega-fr-2026-10-23-r16/EGA_French_Current_Cumulative_Reader.pdf\#page=656}
\expandafter\def\csname EGAExternal@IV.1.7.9\endcsname{https://github.com/KokunoYumeto/ega-fr/releases/download/ega-fr-2026-10-23-r16/EGA_French_Current_Cumulative_Reader.pdf\#page=656}
\expandafter\def\csname EGAExternal@IV.1.8.1\endcsname{https://github.com/KokunoYumeto/ega-fr/releases/download/ega-fr-2026-10-23-r16/EGA_French_Current_Cumulative_Reader.pdf\#page=657}
\expandafter\def\csname EGAExternal@IV.1.8.2\endcsname{https://github.com/KokunoYumeto/ega-fr/releases/download/ega-fr-2026-10-23-r16/EGA_French_Current_Cumulative_Reader.pdf\#page=657}
\expandafter\def\csname EGAExternal@IV.1.8.4\endcsname{https://github.com/KokunoYumeto/ega-fr/releases/download/ega-fr-2026-10-23-r16/EGA_French_Current_Cumulative_Reader.pdf\#page=657}
\expandafter\def\csname EGAExternal@IV.1.8.4.1\endcsname{https://github.com/KokunoYumeto/ega-fr/releases/download/ega-fr-2026-10-23-r16/EGA_French_Current_Cumulative_Reader.pdf\#page=657}
\expandafter\def\csname EGAExternal@IV.1.8.4.2\endcsname{https://github.com/KokunoYumeto/ega-fr/releases/download/ega-fr-2026-10-23-r16/EGA_French_Current_Cumulative_Reader.pdf\#page=658}
\expandafter\def\csname EGAExternal@IV.1.8.6\endcsname{https://github.com/KokunoYumeto/ega-fr/releases/download/ega-fr-2026-10-23-r16/EGA_French_Current_Cumulative_Reader.pdf\#page=658}
\expandafter\def\csname EGAExternal@IV.1.8.7\endcsname{https://github.com/KokunoYumeto/ega-fr/releases/download/ega-fr-2026-10-23-r16/EGA_French_Current_Cumulative_Reader.pdf\#page=657}
\expandafter\def\csname EGAExternal@IV.1.8.7.1\endcsname{https://github.com/KokunoYumeto/ega-fr/releases/download/ega-fr-2026-10-23-r16/EGA_French_Current_Cumulative_Reader.pdf\#page=659}
\expandafter\def\csname EGAExternal@IV.1.9.10\endcsname{https://github.com/KokunoYumeto/ega-fr/releases/download/ega-fr-2026-10-23-r16/EGA_French_Current_Cumulative_Reader.pdf\#page=659}
\expandafter\def\csname EGAExternal@IV.1.9.11\endcsname{https://github.com/KokunoYumeto/ega-fr/releases/download/ega-fr-2026-10-23-r16/EGA_French_Current_Cumulative_Reader.pdf\#page=659}
\expandafter\def\csname EGAExternal@IV.1.9.2.1\endcsname{https://github.com/KokunoYumeto/ega-fr/releases/download/ega-fr-2026-10-23-r16/EGA_French_Current_Cumulative_Reader.pdf\#page=659}
\expandafter\def\csname EGAExternal@IV.1.9.4\endcsname{https://github.com/KokunoYumeto/ega-fr/releases/download/ega-fr-2026-10-23-r16/EGA_French_Current_Cumulative_Reader.pdf\#page=659}
\expandafter\def\csname EGAExternal@IV.1.9.5\endcsname{https://github.com/KokunoYumeto/ega-fr/releases/download/ega-fr-2026-10-23-r16/EGA_French_Current_Cumulative_Reader.pdf\#page=659}
\expandafter\def\csname EGAExternal@IV.1.9.6\endcsname{https://github.com/KokunoYumeto/ega-fr/releases/download/ega-fr-2026-10-23-r16/EGA_French_Current_Cumulative_Reader.pdf\#page=659}
\expandafter\def\csname EGAExternal@IV.1.9.9\endcsname{https://github.com/KokunoYumeto/ega-fr/releases/download/ega-fr-2026-10-23-r16/EGA_French_Current_Cumulative_Reader.pdf\#page=659}
\expandafter\def\csname EGAExternal@IV.11.8\endcsname{https://github.com/KokunoYumeto/ega-fr/releases/download/ega-fr-2026-10-23-r16/EGA_French_Current_Cumulative_Reader.pdf\#page=941}
\expandafter\def\csname EGAExternal@IV.15.2\endcsname{https://github.com/KokunoYumeto/ega-fr/releases/download/ega-fr-2026-10-23-r16/EGA_French_Current_Cumulative_Reader.pdf\#page=986}
\expandafter\def\csname EGAExternal@IV.16.4.11\endcsname{https://github.com/KokunoYumeto/ega-fr/releases/download/ega-fr-2026-10-23-r16/EGA_French_Current_Cumulative_Reader.pdf\#page=1019}
\expandafter\def\csname EGAExternal@IV.17.3.7\endcsname{https://github.com/KokunoYumeto/ega-fr/releases/download/ega-fr-2026-10-23-r16/EGA_French_Current_Cumulative_Reader.pdf\#page=1048}
\expandafter\def\csname EGAExternal@IV.17.4.2\endcsname{https://github.com/KokunoYumeto/ega-fr/releases/download/ega-fr-2026-10-23-r16/EGA_French_Current_Cumulative_Reader.pdf\#page=1049}
\expandafter\def\csname EGAExternal@IV.17.6.3\endcsname{https://github.com/KokunoYumeto/ega-fr/releases/download/ega-fr-2026-10-23-r16/EGA_French_Current_Cumulative_Reader.pdf\#page=1053}
\expandafter\def\csname EGAExternal@IV.17.9.6\endcsname{https://github.com/KokunoYumeto/ega-fr/releases/download/ega-fr-2026-10-23-r16/EGA_French_Current_Cumulative_Reader.pdf\#page=1060}
\expandafter\def\csname EGAExternal@IV.2.1.10\endcsname{https://github.com/KokunoYumeto/ega-fr/releases/download/ega-fr-2026-10-23-r16/EGA_French_Current_Cumulative_Reader.pdf\#page=676}
\expandafter\def\csname EGAExternal@IV.2.1.13\endcsname{https://github.com/KokunoYumeto/ega-fr/releases/download/ega-fr-2026-10-23-r16/EGA_French_Current_Cumulative_Reader.pdf\#page=676}
\expandafter\def\csname EGAExternal@IV.2.1.2\endcsname{https://github.com/KokunoYumeto/ega-fr/releases/download/ega-fr-2026-10-23-r16/EGA_French_Current_Cumulative_Reader.pdf\#page=676}
\expandafter\def\csname EGAExternal@IV.2.1.4\endcsname{https://github.com/KokunoYumeto/ega-fr/releases/download/ega-fr-2026-10-23-r16/EGA_French_Current_Cumulative_Reader.pdf\#page=676}
\expandafter\def\csname EGAExternal@IV.2.1.6\endcsname{https://github.com/KokunoYumeto/ega-fr/releases/download/ega-fr-2026-10-23-r16/EGA_French_Current_Cumulative_Reader.pdf\#page=676}
\expandafter\def\csname EGAExternal@IV.2.1.8\endcsname{https://github.com/KokunoYumeto/ega-fr/releases/download/ega-fr-2026-10-23-r16/EGA_French_Current_Cumulative_Reader.pdf\#page=676}
\expandafter\def\csname EGAExternal@IV.2.2.11\endcsname{https://github.com/KokunoYumeto/ega-fr/releases/download/ega-fr-2026-10-23-r16/EGA_French_Current_Cumulative_Reader.pdf\#page=680}
\expandafter\def\csname EGAExternal@IV.2.2.13\endcsname{https://github.com/KokunoYumeto/ega-fr/releases/download/ega-fr-2026-10-23-r16/EGA_French_Current_Cumulative_Reader.pdf\#page=680}
\expandafter\def\csname EGAExternal@IV.2.2.7\endcsname{https://github.com/KokunoYumeto/ega-fr/releases/download/ega-fr-2026-10-23-r16/EGA_French_Current_Cumulative_Reader.pdf\#page=679}
\expandafter\def\csname EGAExternal@IV.2.2.8\endcsname{https://github.com/KokunoYumeto/ega-fr/releases/download/ega-fr-2026-10-23-r16/EGA_French_Current_Cumulative_Reader.pdf\#page=679}
\expandafter\def\csname EGAExternal@IV.2.3.1\endcsname{https://github.com/KokunoYumeto/ega-fr/releases/download/ega-fr-2026-10-23-r16/EGA_French_Current_Cumulative_Reader.pdf\#page=682}
\expandafter\def\csname EGAExternal@IV.2.3.10\endcsname{https://github.com/KokunoYumeto/ega-fr/releases/download/ega-fr-2026-10-23-r16/EGA_French_Current_Cumulative_Reader.pdf\#page=682}
\expandafter\def\csname EGAExternal@IV.2.3.11\endcsname{https://github.com/KokunoYumeto/ega-fr/releases/download/ega-fr-2026-10-23-r16/EGA_French_Current_Cumulative_Reader.pdf\#page=682}
\expandafter\def\csname EGAExternal@IV.2.3.12\endcsname{https://github.com/KokunoYumeto/ega-fr/releases/download/ega-fr-2026-10-23-r16/EGA_French_Current_Cumulative_Reader.pdf\#page=682}
\expandafter\def\csname EGAExternal@IV.2.3.14\endcsname{https://github.com/KokunoYumeto/ega-fr/releases/download/ega-fr-2026-10-23-r16/EGA_French_Current_Cumulative_Reader.pdf\#page=682}
\expandafter\def\csname EGAExternal@IV.2.3.3\endcsname{https://github.com/KokunoYumeto/ega-fr/releases/download/ega-fr-2026-10-23-r16/EGA_French_Current_Cumulative_Reader.pdf\#page=682}
\expandafter\def\csname EGAExternal@IV.2.3.4\endcsname{https://github.com/KokunoYumeto/ega-fr/releases/download/ega-fr-2026-10-23-r16/EGA_French_Current_Cumulative_Reader.pdf\#page=682}
\expandafter\def\csname EGAExternal@IV.2.3.5\endcsname{https://github.com/KokunoYumeto/ega-fr/releases/download/ega-fr-2026-10-23-r16/EGA_French_Current_Cumulative_Reader.pdf\#page=682}
\expandafter\def\csname EGAExternal@IV.2.3.8\endcsname{https://github.com/KokunoYumeto/ega-fr/releases/download/ega-fr-2026-10-23-r16/EGA_French_Current_Cumulative_Reader.pdf\#page=682}
\expandafter\def\csname EGAExternal@IV.2.4.10\endcsname{https://github.com/KokunoYumeto/ega-fr/releases/download/ega-fr-2026-10-23-r16/EGA_French_Current_Cumulative_Reader.pdf\#page=685}
\expandafter\def\csname EGAExternal@IV.2.4.2\endcsname{https://github.com/KokunoYumeto/ega-fr/releases/download/ega-fr-2026-10-23-r16/EGA_French_Current_Cumulative_Reader.pdf\#page=685}
\expandafter\def\csname EGAExternal@IV.2.4.3\endcsname{https://github.com/KokunoYumeto/ega-fr/releases/download/ega-fr-2026-10-23-r16/EGA_French_Current_Cumulative_Reader.pdf\#page=685}
\expandafter\def\csname EGAExternal@IV.2.4.4\endcsname{https://github.com/KokunoYumeto/ega-fr/releases/download/ega-fr-2026-10-23-r16/EGA_French_Current_Cumulative_Reader.pdf\#page=685}
\expandafter\def\csname EGAExternal@IV.2.4.5\endcsname{https://github.com/KokunoYumeto/ega-fr/releases/download/ega-fr-2026-10-23-r16/EGA_French_Current_Cumulative_Reader.pdf\#page=686}
\expandafter\def\csname EGAExternal@IV.2.4.6\endcsname{https://github.com/KokunoYumeto/ega-fr/releases/download/ega-fr-2026-10-23-r16/EGA_French_Current_Cumulative_Reader.pdf\#page=685}
\expandafter\def\csname EGAExternal@IV.2.7.1\endcsname{https://github.com/KokunoYumeto/ega-fr/releases/download/ega-fr-2026-10-23-r16/EGA_French_Current_Cumulative_Reader.pdf\#page=694}
\expandafter\def\csname EGAExternal@IV.2.7.2\endcsname{https://github.com/KokunoYumeto/ega-fr/releases/download/ega-fr-2026-10-23-r16/EGA_French_Current_Cumulative_Reader.pdf\#page=694}
\expandafter\def\csname EGAExternal@IV.3.1.3\endcsname{https://github.com/KokunoYumeto/ega-fr/releases/download/ega-fr-2026-10-23-r16/EGA_French_Current_Cumulative_Reader.pdf\#page=699}
\expandafter\def\csname EGAExternal@IV.3.1.4\endcsname{https://github.com/KokunoYumeto/ega-fr/releases/download/ega-fr-2026-10-23-r16/EGA_French_Current_Cumulative_Reader.pdf\#page=699}
\expandafter\def\csname EGAExternal@IV.3.1.6\endcsname{https://github.com/KokunoYumeto/ega-fr/releases/download/ega-fr-2026-10-23-r16/EGA_French_Current_Cumulative_Reader.pdf\#page=699}
\expandafter\def\csname EGAExternal@IV.3.1.7\endcsname{https://github.com/KokunoYumeto/ega-fr/releases/download/ega-fr-2026-10-23-r16/EGA_French_Current_Cumulative_Reader.pdf\#page=700}
\expandafter\def\csname EGAExternal@IV.3.2.1\endcsname{https://github.com/KokunoYumeto/ega-fr/releases/download/ega-fr-2026-10-23-r16/EGA_French_Current_Cumulative_Reader.pdf\#page=701}
\expandafter\def\csname EGAExternal@IV.3.2.2\endcsname{https://github.com/KokunoYumeto/ega-fr/releases/download/ega-fr-2026-10-23-r16/EGA_French_Current_Cumulative_Reader.pdf\#page=701}
\expandafter\def\csname EGAExternal@IV.3.2.5\endcsname{https://github.com/KokunoYumeto/ega-fr/releases/download/ega-fr-2026-10-23-r16/EGA_French_Current_Cumulative_Reader.pdf\#page=701}
\expandafter\def\csname EGAExternal@IV.3.2.6\endcsname{https://github.com/KokunoYumeto/ega-fr/releases/download/ega-fr-2026-10-23-r16/EGA_French_Current_Cumulative_Reader.pdf\#page=701}
\expandafter\def\csname EGAExternal@IV.3.3.4\endcsname{https://github.com/KokunoYumeto/ega-fr/releases/download/ega-fr-2026-10-23-r16/EGA_French_Current_Cumulative_Reader.pdf\#page=703}
\expandafter\def\csname EGAExternal@IV.3.4\endcsname{https://github.com/KokunoYumeto/ega-fr/releases/download/ega-fr-2026-10-23-r16/EGA_French_Current_Cumulative_Reader.pdf\#page=705}
\expandafter\def\csname EGAExternal@IV.3.4.1.1\endcsname{https://github.com/KokunoYumeto/ega-fr/releases/download/ega-fr-2026-10-23-r16/EGA_French_Current_Cumulative_Reader.pdf\#page=705}
\expandafter\def\csname EGAExternal@IV.3.4.1.4\endcsname{https://github.com/KokunoYumeto/ega-fr/releases/download/ega-fr-2026-10-23-r16/EGA_French_Current_Cumulative_Reader.pdf\#page=705}
\expandafter\def\csname EGAExternal@IV.3.4.3\endcsname{https://github.com/KokunoYumeto/ega-fr/releases/download/ega-fr-2026-10-23-r16/EGA_French_Current_Cumulative_Reader.pdf\#page=705}
\expandafter\def\csname EGAExternal@IV.3.4.4\endcsname{https://github.com/KokunoYumeto/ega-fr/releases/download/ega-fr-2026-10-23-r16/EGA_French_Current_Cumulative_Reader.pdf\#page=705}
\expandafter\def\csname EGAExternal@IV.3.4.5\endcsname{https://github.com/KokunoYumeto/ega-fr/releases/download/ega-fr-2026-10-23-r16/EGA_French_Current_Cumulative_Reader.pdf\#page=705}
\expandafter\def\csname EGAExternal@IV.3.4.6\endcsname{https://github.com/KokunoYumeto/ega-fr/releases/download/ega-fr-2026-10-23-r16/EGA_French_Current_Cumulative_Reader.pdf\#page=705}
\expandafter\def\csname EGAExternal@IV.3.4.6.1\endcsname{https://github.com/KokunoYumeto/ega-fr/releases/download/ega-fr-2026-10-23-r16/EGA_French_Current_Cumulative_Reader.pdf\#page=705}
\expandafter\def\csname EGAExternal@IV.4.1.1\endcsname{https://github.com/KokunoYumeto/ega-fr/releases/download/ega-fr-2026-10-23-r16/EGA_French_Current_Cumulative_Reader.pdf\#page=710}
\expandafter\def\csname EGAExternal@IV.4.1.1.3\endcsname{https://github.com/KokunoYumeto/ega-fr/releases/download/ega-fr-2026-10-23-r16/EGA_French_Current_Cumulative_Reader.pdf\#page=710}
\expandafter\def\csname EGAExternal@IV.4.1.2\endcsname{https://github.com/KokunoYumeto/ega-fr/releases/download/ega-fr-2026-10-23-r16/EGA_French_Current_Cumulative_Reader.pdf\#page=710}
\expandafter\def\csname EGAExternal@IV.4.1.2.1\endcsname{https://github.com/KokunoYumeto/ega-fr/releases/download/ega-fr-2026-10-23-r16/EGA_French_Current_Cumulative_Reader.pdf\#page=710}
\expandafter\def\csname EGAExternal@IV.4.1.4\endcsname{https://github.com/KokunoYumeto/ega-fr/releases/download/ega-fr-2026-10-23-r16/EGA_French_Current_Cumulative_Reader.pdf\#page=710}
\expandafter\def\csname EGAExternal@IV.4.2.2\endcsname{https://github.com/KokunoYumeto/ega-fr/releases/download/ega-fr-2026-10-23-r16/EGA_French_Current_Cumulative_Reader.pdf\#page=712}
\expandafter\def\csname EGAExternal@IV.4.2.6\endcsname{https://github.com/KokunoYumeto/ega-fr/releases/download/ega-fr-2026-10-23-r16/EGA_French_Current_Cumulative_Reader.pdf\#page=712}
\expandafter\def\csname EGAExternal@IV.4.2.7\endcsname{https://github.com/KokunoYumeto/ega-fr/releases/download/ega-fr-2026-10-23-r16/EGA_French_Current_Cumulative_Reader.pdf\#page=712}
\expandafter\def\csname EGAExternal@IV.4.2.8\endcsname{https://github.com/KokunoYumeto/ega-fr/releases/download/ega-fr-2026-10-23-r16/EGA_French_Current_Cumulative_Reader.pdf\#page=712}
\expandafter\def\csname EGAExternal@IV.4.3.2\endcsname{https://github.com/KokunoYumeto/ega-fr/releases/download/ega-fr-2026-10-23-r16/EGA_French_Current_Cumulative_Reader.pdf\#page=714}
\expandafter\def\csname EGAExternal@IV.4.4.1\endcsname{https://github.com/KokunoYumeto/ega-fr/releases/download/ega-fr-2026-10-23-r16/EGA_French_Current_Cumulative_Reader.pdf\#page=715}
\expandafter\def\csname EGAExternal@IV.4.5\endcsname{https://github.com/KokunoYumeto/ega-fr/releases/download/ega-fr-2026-10-23-r16/EGA_French_Current_Cumulative_Reader.pdf\#page=716}
\expandafter\def\csname EGAExternal@IV.4.5.1\endcsname{https://github.com/KokunoYumeto/ega-fr/releases/download/ega-fr-2026-10-23-r16/EGA_French_Current_Cumulative_Reader.pdf\#page=716}
\expandafter\def\csname EGAExternal@IV.4.5.11\endcsname{https://github.com/KokunoYumeto/ega-fr/releases/download/ega-fr-2026-10-23-r16/EGA_French_Current_Cumulative_Reader.pdf\#page=717}
\expandafter\def\csname EGAExternal@IV.4.5.13\endcsname{https://github.com/KokunoYumeto/ega-fr/releases/download/ega-fr-2026-10-23-r16/EGA_French_Current_Cumulative_Reader.pdf\#page=716}
\expandafter\def\csname EGAExternal@IV.4.5.13.1\endcsname{https://github.com/KokunoYumeto/ega-fr/releases/download/ega-fr-2026-10-23-r16/EGA_French_Current_Cumulative_Reader.pdf\#page=716}
\expandafter\def\csname EGAExternal@IV.4.5.2\endcsname{https://github.com/KokunoYumeto/ega-fr/releases/download/ega-fr-2026-10-23-r16/EGA_French_Current_Cumulative_Reader.pdf\#page=716}
\expandafter\def\csname EGAExternal@IV.4.5.6\endcsname{https://github.com/KokunoYumeto/ega-fr/releases/download/ega-fr-2026-10-23-r16/EGA_French_Current_Cumulative_Reader.pdf\#page=716}
\expandafter\def\csname EGAExternal@IV.4.5.8\endcsname{https://github.com/KokunoYumeto/ega-fr/releases/download/ega-fr-2026-10-23-r16/EGA_French_Current_Cumulative_Reader.pdf\#page=716}
\expandafter\def\csname EGAExternal@IV.4.5.9\endcsname{https://github.com/KokunoYumeto/ega-fr/releases/download/ega-fr-2026-10-23-r16/EGA_French_Current_Cumulative_Reader.pdf\#page=716}
\expandafter\def\csname EGAExternal@IV.4.6\endcsname{https://github.com/KokunoYumeto/ega-fr/releases/download/ega-fr-2026-10-23-r16/EGA_French_Current_Cumulative_Reader.pdf\#page=721}
\expandafter\def\csname EGAExternal@IV.4.6.1\endcsname{https://github.com/KokunoYumeto/ega-fr/releases/download/ega-fr-2026-10-23-r16/EGA_French_Current_Cumulative_Reader.pdf\#page=721}
\expandafter\def\csname EGAExternal@IV.4.6.10\endcsname{https://github.com/KokunoYumeto/ega-fr/releases/download/ega-fr-2026-10-23-r16/EGA_French_Current_Cumulative_Reader.pdf\#page=721}
\expandafter\def\csname EGAExternal@IV.4.6.17\endcsname{https://github.com/KokunoYumeto/ega-fr/releases/download/ega-fr-2026-10-23-r16/EGA_French_Current_Cumulative_Reader.pdf\#page=721}
\expandafter\def\csname EGAExternal@IV.4.6.2\endcsname{https://github.com/KokunoYumeto/ega-fr/releases/download/ega-fr-2026-10-23-r16/EGA_French_Current_Cumulative_Reader.pdf\#page=721}
\expandafter\def\csname EGAExternal@IV.4.6.22\endcsname{https://github.com/KokunoYumeto/ega-fr/releases/download/ega-fr-2026-10-23-r16/EGA_French_Current_Cumulative_Reader.pdf\#page=721}
\expandafter\def\csname EGAExternal@IV.4.6.3\endcsname{https://github.com/KokunoYumeto/ega-fr/releases/download/ega-fr-2026-10-23-r16/EGA_French_Current_Cumulative_Reader.pdf\#page=721}
\expandafter\def\csname EGAExternal@IV.4.6.5\endcsname{https://github.com/KokunoYumeto/ega-fr/releases/download/ega-fr-2026-10-23-r16/EGA_French_Current_Cumulative_Reader.pdf\#page=721}
\expandafter\def\csname EGAExternal@IV.4.6.6\endcsname{https://github.com/KokunoYumeto/ega-fr/releases/download/ega-fr-2026-10-23-r16/EGA_French_Current_Cumulative_Reader.pdf\#page=721}
\expandafter\def\csname EGAExternal@IV.4.6.8\endcsname{https://github.com/KokunoYumeto/ega-fr/releases/download/ega-fr-2026-10-23-r16/EGA_French_Current_Cumulative_Reader.pdf\#page=721}
\expandafter\def\csname EGAExternal@IV.4.6.9\endcsname{https://github.com/KokunoYumeto/ega-fr/releases/download/ega-fr-2026-10-23-r16/EGA_French_Current_Cumulative_Reader.pdf\#page=721}
\expandafter\def\csname EGAExternal@IV.4.7\endcsname{https://github.com/KokunoYumeto/ega-fr/releases/download/ega-fr-2026-10-23-r16/EGA_French_Current_Cumulative_Reader.pdf\#page=726}
\expandafter\def\csname EGAExternal@IV.4.7.10\endcsname{https://github.com/KokunoYumeto/ega-fr/releases/download/ega-fr-2026-10-23-r16/EGA_French_Current_Cumulative_Reader.pdf\#page=728}
\expandafter\def\csname EGAExternal@IV.4.7.11\endcsname{https://github.com/KokunoYumeto/ega-fr/releases/download/ega-fr-2026-10-23-r16/EGA_French_Current_Cumulative_Reader.pdf\#page=728}
\expandafter\def\csname EGAExternal@IV.4.7.12\endcsname{https://github.com/KokunoYumeto/ega-fr/releases/download/ega-fr-2026-10-23-r16/EGA_French_Current_Cumulative_Reader.pdf\#page=728}
\expandafter\def\csname EGAExternal@IV.4.7.13\endcsname{https://github.com/KokunoYumeto/ega-fr/releases/download/ega-fr-2026-10-23-r16/EGA_French_Current_Cumulative_Reader.pdf\#page=728}
\expandafter\def\csname EGAExternal@IV.4.7.3\endcsname{https://github.com/KokunoYumeto/ega-fr/releases/download/ega-fr-2026-10-23-r16/EGA_French_Current_Cumulative_Reader.pdf\#page=726}
\expandafter\def\csname EGAExternal@IV.4.7.4\endcsname{https://github.com/KokunoYumeto/ega-fr/releases/download/ega-fr-2026-10-23-r16/EGA_French_Current_Cumulative_Reader.pdf\#page=726}
\expandafter\def\csname EGAExternal@IV.4.7.5\endcsname{https://github.com/KokunoYumeto/ega-fr/releases/download/ega-fr-2026-10-23-r16/EGA_French_Current_Cumulative_Reader.pdf\#page=726}
\expandafter\def\csname EGAExternal@IV.4.7.8\endcsname{https://github.com/KokunoYumeto/ega-fr/releases/download/ega-fr-2026-10-23-r16/EGA_French_Current_Cumulative_Reader.pdf\#page=728}
\expandafter\def\csname EGAExternal@IV.4.7.9\endcsname{https://github.com/KokunoYumeto/ega-fr/releases/download/ega-fr-2026-10-23-r16/EGA_French_Current_Cumulative_Reader.pdf\#page=728}
\expandafter\def\csname EGAExternal@IV.4.8\endcsname{https://github.com/KokunoYumeto/ega-fr/releases/download/ega-fr-2026-10-23-r16/EGA_French_Current_Cumulative_Reader.pdf\#page=729}
\expandafter\def\csname EGAExternal@IV.5.1.1.5\endcsname{https://github.com/KokunoYumeto/ega-fr/releases/download/ega-fr-2026-10-23-r16/EGA_French_Current_Cumulative_Reader.pdf\#page=734}
\expandafter\def\csname EGAExternal@IV.5.1.10\endcsname{https://github.com/KokunoYumeto/ega-fr/releases/download/ega-fr-2026-10-23-r16/EGA_French_Current_Cumulative_Reader.pdf\#page=735}
\expandafter\def\csname EGAExternal@IV.5.1.10.1\endcsname{https://github.com/KokunoYumeto/ega-fr/releases/download/ega-fr-2026-10-23-r16/EGA_French_Current_Cumulative_Reader.pdf\#page=735}
\expandafter\def\csname EGAExternal@IV.5.1.11\endcsname{https://github.com/KokunoYumeto/ega-fr/releases/download/ega-fr-2026-10-23-r16/EGA_French_Current_Cumulative_Reader.pdf\#page=735}
\expandafter\def\csname EGAExternal@IV.5.1.12\endcsname{https://github.com/KokunoYumeto/ega-fr/releases/download/ega-fr-2026-10-23-r16/EGA_French_Current_Cumulative_Reader.pdf\#page=735}
\expandafter\def\csname EGAExternal@IV.5.1.12.2\endcsname{https://github.com/KokunoYumeto/ega-fr/releases/download/ega-fr-2026-10-23-r16/EGA_French_Current_Cumulative_Reader.pdf\#page=736}
\expandafter\def\csname EGAExternal@IV.5.1.2\endcsname{https://github.com/KokunoYumeto/ega-fr/releases/download/ega-fr-2026-10-23-r16/EGA_French_Current_Cumulative_Reader.pdf\#page=734}
\expandafter\def\csname EGAExternal@IV.5.1.2.1\endcsname{https://github.com/KokunoYumeto/ega-fr/releases/download/ega-fr-2026-10-23-r16/EGA_French_Current_Cumulative_Reader.pdf\#page=734}
\expandafter\def\csname EGAExternal@IV.5.1.4\endcsname{https://github.com/KokunoYumeto/ega-fr/releases/download/ega-fr-2026-10-23-r16/EGA_French_Current_Cumulative_Reader.pdf\#page=734}
\expandafter\def\csname EGAExternal@IV.5.1.5\endcsname{https://github.com/KokunoYumeto/ega-fr/releases/download/ega-fr-2026-10-23-r16/EGA_French_Current_Cumulative_Reader.pdf\#page=734}
\expandafter\def\csname EGAExternal@IV.5.1.8\endcsname{https://github.com/KokunoYumeto/ega-fr/releases/download/ega-fr-2026-10-23-r16/EGA_French_Current_Cumulative_Reader.pdf\#page=734}
\expandafter\def\csname EGAExternal@IV.5.10.2\endcsname{https://github.com/KokunoYumeto/ega-fr/releases/download/ega-fr-2026-10-23-r16/EGA_French_Current_Cumulative_Reader.pdf\#page=753}
\expandafter\def\csname EGAExternal@IV.5.10.7\endcsname{https://github.com/KokunoYumeto/ega-fr/releases/download/ega-fr-2026-10-23-r16/EGA_French_Current_Cumulative_Reader.pdf\#page=753}
\expandafter\def\csname EGAExternal@IV.5.10.8.1\endcsname{https://github.com/KokunoYumeto/ega-fr/releases/download/ega-fr-2026-10-23-r16/EGA_French_Current_Cumulative_Reader.pdf\#page=753}
\expandafter\def\csname EGAExternal@IV.5.11.1.1\endcsname{https://github.com/KokunoYumeto/ega-fr/releases/download/ega-fr-2026-10-23-r16/EGA_French_Current_Cumulative_Reader.pdf\#page=758}
\expandafter\def\csname EGAExternal@IV.5.12\endcsname{https://github.com/KokunoYumeto/ega-fr/releases/download/ega-fr-2026-10-23-r16/EGA_French_Current_Cumulative_Reader.pdf\#page=761}
\expandafter\def\csname EGAExternal@IV.5.12.2\endcsname{https://github.com/KokunoYumeto/ega-fr/releases/download/ega-fr-2026-10-23-r16/EGA_French_Current_Cumulative_Reader.pdf\#page=761}
\expandafter\def\csname EGAExternal@IV.5.12.3\endcsname{https://github.com/KokunoYumeto/ega-fr/releases/download/ega-fr-2026-10-23-r16/EGA_French_Current_Cumulative_Reader.pdf\#page=761}
\expandafter\def\csname EGAExternal@IV.5.12.4\endcsname{https://github.com/KokunoYumeto/ega-fr/releases/download/ega-fr-2026-10-23-r16/EGA_French_Current_Cumulative_Reader.pdf\#page=761}
\expandafter\def\csname EGAExternal@IV.5.12.5\endcsname{https://github.com/KokunoYumeto/ega-fr/releases/download/ega-fr-2026-10-23-r16/EGA_French_Current_Cumulative_Reader.pdf\#page=761}
\expandafter\def\csname EGAExternal@IV.5.12.6\endcsname{https://github.com/KokunoYumeto/ega-fr/releases/download/ega-fr-2026-10-23-r16/EGA_French_Current_Cumulative_Reader.pdf\#page=761}
\expandafter\def\csname EGAExternal@IV.5.13.1\endcsname{https://github.com/KokunoYumeto/ega-fr/releases/download/ega-fr-2026-10-23-r16/EGA_French_Current_Cumulative_Reader.pdf\#page=765}
\expandafter\def\csname EGAExternal@IV.5.13.2\endcsname{https://github.com/KokunoYumeto/ega-fr/releases/download/ega-fr-2026-10-23-r16/EGA_French_Current_Cumulative_Reader.pdf\#page=765}
\expandafter\def\csname EGAExternal@IV.5.13.3\endcsname{https://github.com/KokunoYumeto/ega-fr/releases/download/ega-fr-2026-10-23-r16/EGA_French_Current_Cumulative_Reader.pdf\#page=765}
\expandafter\def\csname EGAExternal@IV.5.13.7.1\endcsname{https://github.com/KokunoYumeto/ega-fr/releases/download/ega-fr-2026-10-23-r16/EGA_French_Current_Cumulative_Reader.pdf\#page=765}
\expandafter\def\csname EGAExternal@IV.5.2.1\endcsname{https://github.com/KokunoYumeto/ega-fr/releases/download/ega-fr-2026-10-23-r16/EGA_French_Current_Cumulative_Reader.pdf\#page=736}
\expandafter\def\csname EGAExternal@IV.5.2.2\endcsname{https://github.com/KokunoYumeto/ega-fr/releases/download/ega-fr-2026-10-23-r16/EGA_French_Current_Cumulative_Reader.pdf\#page=736}
\expandafter\def\csname EGAExternal@IV.5.2.3\endcsname{https://github.com/KokunoYumeto/ega-fr/releases/download/ega-fr-2026-10-23-r16/EGA_French_Current_Cumulative_Reader.pdf\#page=737}
\expandafter\def\csname EGAExternal@IV.5.2.5\endcsname{https://github.com/KokunoYumeto/ega-fr/releases/download/ega-fr-2026-10-23-r16/EGA_French_Current_Cumulative_Reader.pdf\#page=737}
\expandafter\def\csname EGAExternal@IV.5.4.2\endcsname{https://github.com/KokunoYumeto/ega-fr/releases/download/ega-fr-2026-10-23-r16/EGA_French_Current_Cumulative_Reader.pdf\#page=738}
\expandafter\def\csname EGAExternal@IV.5.5\endcsname{https://github.com/KokunoYumeto/ega-fr/releases/download/ega-fr-2026-10-23-r16/EGA_French_Current_Cumulative_Reader.pdf\#page=739}
\expandafter\def\csname EGAExternal@IV.5.5.3\endcsname{https://github.com/KokunoYumeto/ega-fr/releases/download/ega-fr-2026-10-23-r16/EGA_French_Current_Cumulative_Reader.pdf\#page=687}
\expandafter\def\csname EGAExternal@IV.5.6\endcsname{https://github.com/KokunoYumeto/ega-fr/releases/download/ega-fr-2026-10-23-r16/EGA_French_Current_Cumulative_Reader.pdf\#page=741}
\expandafter\def\csname EGAExternal@IV.5.6.1\endcsname{https://github.com/KokunoYumeto/ega-fr/releases/download/ega-fr-2026-10-23-r16/EGA_French_Current_Cumulative_Reader.pdf\#page=741}
\expandafter\def\csname EGAExternal@IV.5.6.10\endcsname{https://github.com/KokunoYumeto/ega-fr/releases/download/ega-fr-2026-10-23-r16/EGA_French_Current_Cumulative_Reader.pdf\#page=744}
\expandafter\def\csname EGAExternal@IV.5.6.2\endcsname{https://github.com/KokunoYumeto/ega-fr/releases/download/ega-fr-2026-10-23-r16/EGA_French_Current_Cumulative_Reader.pdf\#page=742}
\expandafter\def\csname EGAExternal@IV.5.6.3\endcsname{https://github.com/KokunoYumeto/ega-fr/releases/download/ega-fr-2026-10-23-r16/EGA_French_Current_Cumulative_Reader.pdf\#page=742}
\expandafter\def\csname EGAExternal@IV.5.6.4\endcsname{https://github.com/KokunoYumeto/ega-fr/releases/download/ega-fr-2026-10-23-r16/EGA_French_Current_Cumulative_Reader.pdf\#page=742}
\expandafter\def\csname EGAExternal@IV.5.6.5\endcsname{https://github.com/KokunoYumeto/ega-fr/releases/download/ega-fr-2026-10-23-r16/EGA_French_Current_Cumulative_Reader.pdf\#page=743}
\expandafter\def\csname EGAExternal@IV.5.6.5.2\endcsname{https://github.com/KokunoYumeto/ega-fr/releases/download/ega-fr-2026-10-23-r16/EGA_French_Current_Cumulative_Reader.pdf\#page=743}
\expandafter\def\csname EGAExternal@IV.5.6.5.3\endcsname{https://github.com/KokunoYumeto/ega-fr/releases/download/ega-fr-2026-10-23-r16/EGA_French_Current_Cumulative_Reader.pdf\#page=743}
\expandafter\def\csname EGAExternal@IV.5.6.7\endcsname{https://github.com/KokunoYumeto/ega-fr/releases/download/ega-fr-2026-10-23-r16/EGA_French_Current_Cumulative_Reader.pdf\#page=743}
\expandafter\def\csname EGAExternal@IV.5.7.1\endcsname{https://github.com/KokunoYumeto/ega-fr/releases/download/ega-fr-2026-10-23-r16/EGA_French_Current_Cumulative_Reader.pdf\#page=745}
\expandafter\def\csname EGAExternal@IV.5.7.2\endcsname{https://github.com/KokunoYumeto/ega-fr/releases/download/ega-fr-2026-10-23-r16/EGA_French_Current_Cumulative_Reader.pdf\#page=745}
\expandafter\def\csname EGAExternal@IV.5.7.3\endcsname{https://github.com/KokunoYumeto/ega-fr/releases/download/ega-fr-2026-10-23-r16/EGA_French_Current_Cumulative_Reader.pdf\#page=745}
\expandafter\def\csname EGAExternal@IV.5.7.4\endcsname{https://github.com/KokunoYumeto/ega-fr/releases/download/ega-fr-2026-10-23-r16/EGA_French_Current_Cumulative_Reader.pdf\#page=745}
\expandafter\def\csname EGAExternal@IV.5.8.6\endcsname{https://github.com/KokunoYumeto/ega-fr/releases/download/ega-fr-2026-10-23-r16/EGA_French_Current_Cumulative_Reader.pdf\#page=687}
\expandafter\def\csname EGAExternal@IV.6.1.2\endcsname{https://github.com/KokunoYumeto/ega-fr/releases/download/ega-fr-2026-10-23-r16/EGA_French_Current_Cumulative_Reader.pdf\#page=767}
\expandafter\def\csname EGAExternal@IV.6.1.4\endcsname{https://github.com/KokunoYumeto/ega-fr/releases/download/ega-fr-2026-10-23-r16/EGA_French_Current_Cumulative_Reader.pdf\#page=768}
\expandafter\def\csname EGAExternal@IV.6.11.1\endcsname{https://github.com/KokunoYumeto/ega-fr/releases/download/ega-fr-2026-10-23-r16/EGA_French_Current_Cumulative_Reader.pdf\#page=784}
\expandafter\def\csname EGAExternal@IV.6.11.2\endcsname{https://github.com/KokunoYumeto/ega-fr/releases/download/ega-fr-2026-10-23-r16/EGA_French_Current_Cumulative_Reader.pdf\#page=784}
\expandafter\def\csname EGAExternal@IV.6.11.2.1\endcsname{https://github.com/KokunoYumeto/ega-fr/releases/download/ega-fr-2026-10-23-r16/EGA_French_Current_Cumulative_Reader.pdf\#page=785}
\expandafter\def\csname EGAExternal@IV.6.11.5\endcsname{https://github.com/KokunoYumeto/ega-fr/releases/download/ega-fr-2026-10-23-r16/EGA_French_Current_Cumulative_Reader.pdf\#page=785}
\expandafter\def\csname EGAExternal@IV.6.12.4.1\endcsname{https://github.com/KokunoYumeto/ega-fr/releases/download/ega-fr-2026-10-23-r16/EGA_French_Current_Cumulative_Reader.pdf\#page=788}
\expandafter\def\csname EGAExternal@IV.6.14.4\endcsname{https://github.com/KokunoYumeto/ega-fr/releases/download/ega-fr-2026-10-23-r16/EGA_French_Current_Cumulative_Reader.pdf\#page=795}
\expandafter\def\csname EGAExternal@IV.6.15.1\endcsname{https://github.com/KokunoYumeto/ega-fr/releases/download/ega-fr-2026-10-23-r16/EGA_French_Current_Cumulative_Reader.pdf\#page=796}
\expandafter\def\csname EGAExternal@IV.6.15.11\endcsname{https://github.com/KokunoYumeto/ega-fr/releases/download/ega-fr-2026-10-23-r16/EGA_French_Current_Cumulative_Reader.pdf\#page=796}
\expandafter\def\csname EGAExternal@IV.6.15.2\endcsname{https://github.com/KokunoYumeto/ega-fr/releases/download/ega-fr-2026-10-23-r16/EGA_French_Current_Cumulative_Reader.pdf\#page=796}
\expandafter\def\csname EGAExternal@IV.6.15.5\endcsname{https://github.com/KokunoYumeto/ega-fr/releases/download/ega-fr-2026-10-23-r16/EGA_French_Current_Cumulative_Reader.pdf\#page=796}
\expandafter\def\csname EGAExternal@IV.6.15.8\endcsname{https://github.com/KokunoYumeto/ega-fr/releases/download/ega-fr-2026-10-23-r16/EGA_French_Current_Cumulative_Reader.pdf\#page=796}
\expandafter\def\csname EGAExternal@IV.6.15.9\endcsname{https://github.com/KokunoYumeto/ega-fr/releases/download/ega-fr-2026-10-23-r16/EGA_French_Current_Cumulative_Reader.pdf\#page=796}
\expandafter\def\csname EGAExternal@IV.6.2.1\endcsname{https://github.com/KokunoYumeto/ega-fr/releases/download/ega-fr-2026-10-23-r16/EGA_French_Current_Cumulative_Reader.pdf\#page=769}
\expandafter\def\csname EGAExternal@IV.6.2.3\endcsname{https://github.com/KokunoYumeto/ega-fr/releases/download/ega-fr-2026-10-23-r16/EGA_French_Current_Cumulative_Reader.pdf\#page=769}
\expandafter\def\csname EGAExternal@IV.6.3.3\endcsname{https://github.com/KokunoYumeto/ega-fr/releases/download/ega-fr-2026-10-23-r16/EGA_French_Current_Cumulative_Reader.pdf\#page=770}
\expandafter\def\csname EGAExternal@IV.6.4.1\endcsname{https://github.com/KokunoYumeto/ega-fr/releases/download/ega-fr-2026-10-23-r16/EGA_French_Current_Cumulative_Reader.pdf\#page=771}
\expandafter\def\csname EGAExternal@IV.6.5.4\endcsname{https://github.com/KokunoYumeto/ega-fr/releases/download/ega-fr-2026-10-23-r16/EGA_French_Current_Cumulative_Reader.pdf\#page=773}
\expandafter\def\csname EGAExternal@IV.6.5.5\endcsname{https://github.com/KokunoYumeto/ega-fr/releases/download/ega-fr-2026-10-23-r16/EGA_French_Current_Cumulative_Reader.pdf\#page=773}
\expandafter\def\csname EGAExternal@IV.6.7.1\endcsname{https://github.com/KokunoYumeto/ega-fr/releases/download/ega-fr-2026-10-23-r16/EGA_French_Current_Cumulative_Reader.pdf\#page=775}
\expandafter\def\csname EGAExternal@IV.6.7.6\endcsname{https://github.com/KokunoYumeto/ega-fr/releases/download/ega-fr-2026-10-23-r16/EGA_French_Current_Cumulative_Reader.pdf\#page=775}
\expandafter\def\csname EGAExternal@IV.6.7.8\endcsname{https://github.com/KokunoYumeto/ega-fr/releases/download/ega-fr-2026-10-23-r16/EGA_French_Current_Cumulative_Reader.pdf\#page=775}
\expandafter\def\csname EGAExternal@IV.6.8.1\endcsname{https://github.com/KokunoYumeto/ega-fr/releases/download/ega-fr-2026-10-23-r16/EGA_French_Current_Cumulative_Reader.pdf\#page=778}
\expandafter\def\csname EGAExternal@IV.6.8.7\endcsname{https://github.com/KokunoYumeto/ega-fr/releases/download/ega-fr-2026-10-23-r16/EGA_French_Current_Cumulative_Reader.pdf\#page=778}
\expandafter\def\csname EGAExternal@IV.6.9.1\endcsname{https://github.com/KokunoYumeto/ega-fr/releases/download/ega-fr-2026-10-23-r16/EGA_French_Current_Cumulative_Reader.pdf\#page=780}
\expandafter\def\csname EGAExternal@IV.6.9.2\endcsname{https://github.com/KokunoYumeto/ega-fr/releases/download/ega-fr-2026-10-23-r16/EGA_French_Current_Cumulative_Reader.pdf\#page=780}
\expandafter\def\csname EGAExternal@IV.6.9.3\endcsname{https://github.com/KokunoYumeto/ega-fr/releases/download/ega-fr-2026-10-23-r16/EGA_French_Current_Cumulative_Reader.pdf\#page=780}
\expandafter\def\csname EGAExternal@IV.7.1.13\endcsname{https://github.com/KokunoYumeto/ega-fr/releases/download/ega-fr-2026-10-23-r16/EGA_French_Current_Cumulative_Reader.pdf\#page=801}
\expandafter\def\csname EGAExternal@IV.7.2.9\endcsname{https://github.com/KokunoYumeto/ega-fr/releases/download/ega-fr-2026-10-23-r16/EGA_French_Current_Cumulative_Reader.pdf\#page=806}
\expandafter\def\csname EGAExternal@IV.7.8.3\endcsname{https://github.com/KokunoYumeto/ega-fr/releases/download/ega-fr-2026-10-23-r16/EGA_French_Current_Cumulative_Reader.pdf\#page=824}
\expandafter\def\csname EGAExternal@IV.7.8.5\endcsname{https://github.com/KokunoYumeto/ega-fr/releases/download/ega-fr-2026-10-23-r16/EGA_French_Current_Cumulative_Reader.pdf\#page=823}
\expandafter\def\csname EGAExternal@T.1.8\endcsname{https://doi.org/10.2748/tmj/1178244839}
\expandafter\def\csname EGAExternal@subsection:I.3.4\endcsname{https://github.com/KokunoYumeto/ega-fr/releases/download/ega-fr-2026-10-23-r16/EGA_French_Current_Cumulative_Reader.pdf\#page=79}
\renewcommand{\sgaStableOrOriginalLink}[2]{%
  \ifcsname EGAExternal@#1\endcsname
    \href{\csname EGAExternal@#1\endcsname}{#2}%
  \else
    \@ifundefined{r@#1}{#2}{\hyperref[#1]{#2}}%
  \fi}
\newcommand{\EGAReadLink}[2]{\sgaStableOrOriginalLink{#1}{#2}}
\makeatother

% Freeze the pdfTeX trailer identifier for byte-identical replay builds.
\pdftrailerid{0123456789ABCDEF0123456789ABCDEF}
\setlength{\emergencystretch}{3em}

\hypersetup{
  pdftitle={Elementos de geometría algebraica IV, tercera parte -- edición española},
  pdfauthor={Alexandre Grothendieck y Jean Dieudonné},
  pdfsubject={Estudio local de los esquemas y de los morfismos de esquemas: secciones 8 a 15}
}

\setlength{\emergencystretch}{6em}
\begingroup\fontencoding{T1}\fontfamily{ppl}\selectfont\endgroup
\DeclareFontShape{T1}{ppl}{m}{scit}{<->ssub*ppl/m/sc}{}
\begin{document}
% BEGIN RIGHTS AND AI TRANSLATION NOTICE
\begingroup
\small
\noindent\textbf{Derechos, disponibilidad y autoría}

Los derechos sobre las obras originales y sobre las contribuciones protegidas a estas traducciones y transcripciones derivadas corresponden a sus respectivos titulares. Todo aquello sobre lo que el editor tenga derechos en estas ediciones queda dedicado al dominio público bajo CC0 1.0 (https://creativecommons.org/publicdomain/zero/1.0/deed.es). Las menciones de derechos de terceros y las autorizaciones ya concedidas siguen siendo aplicables.

Estas ediciones se ofrecen gratuitamente, con fines no comerciales y exclusivamente educativos, bajo la presunción de que los respectivos titulares de derechos no se opondrían a que estas traducciones y transcripciones estuvieran disponibles gratuitamente junto a los originales franceses en NUMDAM (www.numdam.org/). Esta presunción no significa que se haya obtenido autorización; el acceso gratuito a los originales no autoriza por sí solo su traducción, adaptación o republicación.

Si usted es titular de derechos sobre algún material presentado aquí y se opone a su disponibilidad, puede abrir una incidencia en GitHub (github.com/KokunoYumeto/ega-fga-es/issues). El responsable de la publicación retirará el material afectado a petición del titular de los derechos.

La traducción al español está producida por IA. Las correcciones, la edición y la composición también están producidas por IA. El responsable de la publicación no tradujo el texto. Se conservan la autoría y las atribuciones de las obras originales y de las contribuciones heredadas; no se atribuyen a la IA. Los modelos documentados se identifican en cada edición.
\par
\endgroup
\clearpage
% END RIGHTS AND AI TRANSLATION NOTICE

\begin{titlepage}
\thispagestyle{empty}
\begin{center}
{\Large INSTITUTO DE ALTOS ESTUDIOS CIENTÍFICOS\par}
\vspace{4.5em}
{\LARGE ELEMENTOS DE GEOMETRÍA ALGEBRAICA\par}
\vspace{1.5em}
{\large por A. GROTHENDIECK\par}
\vspace{0.5em}
{\large\itshape Redactados con la colaboración de J. DIEUDONNÉ\par}
\vspace{3em}
{\Large IV\par}
\vspace{1.25em}
{\Large\bfseries ESTUDIO LOCAL DE LOS ESQUEMAS\\
Y DE LOS MORFISMOS DE ESQUEMAS\par}
\vspace{0.75em}
{\large\itshape (Tercera parte)\par}
\vfill
{\Large 1966\par}
\vspace{1.5em}
{\Large PUBLICATIONS MATHÉMATIQUES, N.º 28\par}
\vspace{0.75em}
{\large LE BOIS-MARIE --- BURES-SUR-YVETTE (S.-et-O.)\par}
\end{center}
\end{titlepage}

\clearpage
\thispagestyle{empty}
\vspace*{0.27\textheight}
\begin{center}
\fbox{\begin{minipage}{0.64\textwidth}
\centering
\textbf{DEPÓSITO LEGAL}\\
Primera edición \quad Primer trimestre de 1966\\[0.5em]
\textbf{TODOS LOS DERECHOS}\\
reservados en todos los países\\[0.4em]
\copyright\ 1966, \textit{Institut des Hautes Études Scientifiques}
\end{minipage}}
\end{center}

\vfill

\clearpage


\tableofcontents

\setcounter{section}{7}
\section{Límites proyectivos de preesquemas}
\label{section:IV.8}

\subsection{Introducción}
\label{subsection:IV.8.1}

\oldpage[IV-3]{5}

\begin{env}[8.1.1]
\label{IV.8.1.1}
En este párrafo estudiaremos sistemáticamente la siguiente situación.
Sea $\mathrm{I}$ un conjunto creciente preordenado filtrado, $(\mathrm{A}_\lambda,\varphi_{\mu\lambda})$ un sistema inductivo de anillos cuyo conjunto de índices es $\mathrm{I}$, y sea $\mathrm{A}=\varinjlim\mathrm{A}_\lambda$ su límite inductivo.
Para cada $\alpha\in\mathrm{I}$ y cada $\mathrm{A}_\alpha$-preesquema $\mathrm{X}_\alpha$, consideremos los $\mathrm{A}_\lambda$-preesquemas $\mathrm{X}_\lambda=\mathrm{X}_\alpha\otimes_{\mathrm{A}_\alpha}\mathrm{A}_\lambda$ para $\lambda\geqslant\alpha$, y el $\mathrm{A}$-preesquema $\mathrm{X}=\mathrm{X}_\alpha\otimes_{\mathrm{A}_\alpha}\mathrm{A}$; está claro que los preesquemas $\mathrm{X}_\lambda$ (para $\lambda\geqslant\alpha$) forman un \emph{sistema proyectivo}, y veremos en \sref{IV.8.2.5} que $\mathrm{X}$ es un \emph{límite proyectivo} de este sistema en la categoría de preesquemas.
Nos proponemos encontrar condiciones sobre $\mathrm{X}_\alpha$ o sobre los $\mathrm{A}_\lambda$ que permitan demostrar propiedades del tipo siguiente:
para que $\mathrm{X}$ tenga una propiedad $\mathbf{P}$ (por ejemplo, ser propio sobre $\mathrm{S}=\operatorname{Spec}(\mathrm{A})$, o irreducible, o conexo, etc.), es necesario y suficiente que exista un índice $\mu\geqslant\alpha$ tal que, para todo $\lambda\geqslant\mu$, $\mathrm{X}_\lambda$ tenga (con respecto a $\mathrm{S}_\lambda=\operatorname{Spec}(\mathrm{A}_\lambda)$, según el caso) la misma propiedad $\mathbf{P}$.
Obtendremos enunciados análogos para propiedades de $\sh{O}_\mathrm{X}$-Módulos, de $\mathrm{A}$-morfismos de preesquemas $\mathrm{A}$, etc.
También mostraremos en \sref{IV.8.9.1} que el dato de un preesquema $\mathrm{A}$ de presentación finita \sref{IV.1.6.1} es esencialmente equivalente al dato de un preesquema $\mathrm{A}_\lambda$ de presentación finita $\mathrm{X}_\lambda$ para algún $\lambda$ suficientemente grande, siendo entonces $\mathrm{X}$ \emph{isomorfo a} $\mathrm{X}_\lambda\otimes_{\mathrm{A}_\lambda}\mathrm{A}$.
También existen enunciados análogos para $\sh{O}_\mathrm{X}$-Módulos de presentación finita, sus homomorfismos, los $\mathrm{A}$-morfismos de preesquemas $\mathrm{A}$ de presentación finita, etc.
\end{env}

\begin{env}[8.1.2]
\label{IV.8.1.2}
La utilidad de estos resultados aparecerá, por ejemplo, en las preguntas siguientes:
\begin{enumerate}[label=\emph{\alph*)}]
  \item
  Sea $\mathrm{Y}$ un preesquema, $y$ un punto de $\mathrm{Y}$, y $(\mathrm{U}_\lambda)$ el sistema proyectivo de vecindades afines abiertas filtradas decrecientes de $y$ en $\mathrm{Y}$; si $\mathrm{A}_\lambda$ es el anillo de $\mathrm{U}_\lambda$, los $\mathrm{A}_\lambda$ forman un sistema inductivo creciente filtrado, cuyo límite inductivo $\mathrm{A}$ es el anillo local $\sh{O}_y$; además, si denotamos por $\mathfrak{p}_\lambda$ el ideal primo $\mathfrak{j}_y$ en el anillo $\mathrm{A}_\lambda$, el sistema inductivo $(\mathrm{A}_\lambda)$ es cofinal en todo sistema inductivo $(\mathrm{A}_\alpha)_f$, donde $f$ recorre $\mathrm{A}_\alpha-\mathfrak{p}_\alpha$
\oldpage[IV-3]{6}
(para un $\alpha$ fijo), pues los $\mathrm{D}(f)$ forman un sistema fundamental de vecindades de $y$ en $\mathrm{U}_\alpha$ y, por tanto, en $\mathrm{Y}$.
Por consiguiente, los resultados de este párrafo implican que la geometría algebraica de los preesquemas $\sh{O}_y$ de presentación finita (y la teoría de los Módulos de presentación finita sobre estos preesquemas) es esencialmente equivalente a la geometría algebraica de los preesquemas de presentación finita en vecindades abiertas «suficientemente pequeñas» de $y$.
Así, el enunciado \sref{IV.8.10.5},~(xii) implica que, si un morfismo de presentación finita $f:\mathrm{X}\to\mathrm{Y}$ es tal que $\mathrm{X}\times_\mathrm{Y}\operatorname{Spec}(\sh{O}_y)$ es propio sobre $\operatorname{Spec}(\sh{O}_y)$, es necesario y suficiente que exista una vecindad abierta $\mathrm{U}$ de $y$ en $\mathrm{Y}$ tal que $f^{-1}(\mathrm{U})$ sea propio sobre $\mathrm{U}$.

Un caso particular importante, clásico en cierto sentido, es aquel en que $\mathrm{Y}$ es íntegro y $y$ es su punto genérico, de modo que $\sh{O}_y$ no es otra cosa que \emph{el cuerpo} $\mathrm{R}(\mathrm{Y})=\mathrm{K}$ \emph{de funciones racionales sobre} $\mathrm{Y}$.
Los resultados de este párrafo equivalen entonces a interpretar la geometría algebraica sobre $\mathrm{K}$ en términos de geometría algebraica sobre subconjuntos abiertos no vacíos «suficientemente pequeños» de $\mathrm{Y}$, es decir, intuitivamente, en términos de «familias» de objetos geométricos indexadas por los puntos de tal abierto.
Además, este punto de vista se ha utilizado habitualmente desde hace mucho tiempo, no solo en geometría algebraica sobre cuerpos algebraicamente cerrados, sino también en el estudio aritmético de variedades definidas sobre un \emph{cuerpo de números} $\mathrm{K}$ (una extensión finita de $\mathbf{Q}$), considerando este último como el cuerpo de fracciones de su anillo de enteros $\mathrm{A}$ (la «teoría de la \emph{reducción módulo} $\mathfrak{p}$» ($\mathfrak{p}$ un ideal primo de $\mathrm{A}$); cf.\ \sref[I]{I.3.7}).
Los resultados de las secciones \textsection\textsection\,8 y~9 proporcionan así, entre otras cosas, algunos fundamentos del lenguaje de la «reducción módulo $\mathfrak{p}$» en aritmética.

Observamos que, en el ejemplo considerado aquí, los morfismos $\mathrm{S}_\mu\to\mathrm{S}_\lambda$ (para $\lambda\leqslant\mu$) son las \emph{inmersiones abiertas} canónicas $\mathrm{U}_\mu\hookrightarrow\mathrm{U}_\lambda$ y, \emph{a fortiori}, son morfismos \emph{planos} (aunque en general no fielmente planos), lo que explica el interés de los enunciados que recurren a tal restricción.

  \item
  Supongamos que los $\mathrm{A}_\lambda$ son \emph{cuerpos}, de modo que $\mathrm{A}=\varinjlim\mathrm{A}_\lambda$ también es un cuerpo.
Este caso aparece generalmente cuando se parte de datos geométricos sobre un cuerpo arbitrario $\mathrm{K}$, que se considera como extensión de un cuerpo $k$ (por ejemplo, el subcuerpo primo de $\mathrm{K}$).
Resulta entonces útil considerar $\mathrm{K}$ como límite inductivo de sus subextensiones que son \emph{de tipo finito sobre $k$}, lo que permite reducir muchas cuestiones al caso en que $\mathrm{K}$ es una extensión de tipo finito de $k$.
Usando también el método esbozado en \emph{a)}, se puede reducir en general al caso de un anillo de base $\mathrm{A}$ que sea una \emph{álgebra íntegra de tipo finito sobre $k$}.

Observamos que, en este ejemplo, los morfismos $\mathrm{S}_\mu\to\mathrm{S}_\lambda$ son \emph{fielmente planos}.

  \item
  Supongamos que nos interesan propiedades, locales sobre $\mathrm{Y}$, de preesquemas de presentación finita sobre un preesquema arbitrario $\mathrm{Y}$, que en adelante podemos suponer un esquema afín de anillo $\mathrm{A}$.
Resulta entonces útil considerar $\mathrm{A}$ como límite inductivo de sus subanillos que son \emph{$\mathbf{Z}$-álgebras de tipo finito}, lo que permite reducir muchas cuestiones al caso en que $\mathrm{Y}$ es el espectro de una de tales álgebras.
Esta es la explicación del «punto de vista kroneckeriano», según el cual la geometría algebraica se reduce a la
\oldpage[IV-3]{7}
geometría algebraica de preesquemas de tipo finito sobre $\mathbf{Z}$ (que a veces se califica de «geometría algebraica absoluta»).
Este ejemplo muestra, en particular, que en la mayoría de las cuestiones «relativas» a un preesquema de base $\mathrm{Y}$ se puede reducir al caso en que $\mathrm{Y}$ es \emph{noetheriano}.

Observamos que, en este ejemplo, a diferencia de los anteriores, los morfismos $\mathrm{S}_\mu\to\mathrm{S}_\lambda$ no tienen en general ninguna propiedad particular de regularidad.
\end{enumerate}

En lo que sigue, al aplicar los resultados que siguen a cualquiera de las tres situaciones particulares que acabamos de describir, no entraremos en detalles sobre el procedimiento que permite realizar estas aplicaciones, y nos limitaremos a remitirnos a lo anterior.
\end{env}

\begin{env}[8.1.3]
\label{IV.8.1.3}
En el ejemplo \emph{a)} de \sref{IV.8.1.2} vimos que, si $\mathrm{Y}$ es un preesquema íntegro con punto genérico $y$, $f:\mathrm{X}\to\mathrm{Y}$ un morfismo de presentación finita, y la fibra genérica $f^{-1}(y)$ es propia sobre $\boldsymbol{k}(y)$, entonces existe una vecindad abierta $\mathrm{U}$ de $y$ en $\mathrm{Y}$ tal que $f^{-1}(\mathrm{U})$ es propia sobre $\mathrm{U}$; se deduce \emph{a fortiori} que, para todo $s\in\mathrm{U}$, $f^{-1}(s)$ es propia sobre $\boldsymbol{k}(s)$.
Resulta que se necesita un recíproco que asegure que, si $f^{-1}(s)$ es propia sobre $\boldsymbol{k}(s)$ para «suficientemente muchos» puntos $s\in\mathrm{U}$, entonces $f^{-1}(y)$ es propia sobre $\boldsymbol{k}(y)$.
Por ejemplo, supongamos que $\mathrm{X}$ y $\mathrm{Y}$ son preesquemas algebraicos sobre un cuerpo algebraicamente cerrado $k$ (para fijar ideas, se puede tomar como $k$ el cuerpo $\mathbf{C}$ de los números complejos); a veces se necesita saber que, si para todo $s\in\mathrm{Y}$, \emph{racional sobre $k$}, la fibra $f^{-1}(s)$ es propia sobre $\boldsymbol{k}(s)$, entonces $f^{-1}(y)$ es propia sobre $\boldsymbol{k}(y)$ y, por consiguiente, $f^{-1}(\mathrm{U})$ es propia sobre $\mathrm{U}$ para una vecindad $\mathrm{U}$ de $y$\footnote{Observamos que tal enunciado es, en último término, de naturaleza puramente geométrica, pues solo recurre a puntos racionales sobre $k$, y no a puntos genéricos; por ejemplo, cuando $k=\mathbf{C}$, este enunciado tiene un significado topológico evidente para el analista, interpretando «propio» en el sentido topológico del término para los espacios analíticos subyacentes formados por los puntos de $\mathrm{X}$ y $\mathrm{Y}$ racionales sobre $\mathbf{C}$.}.
Ahora bien, este enunciado se deduce fácilmente de lo siguiente: el conjunto $\mathrm{E}$ de puntos $s\in\mathrm{Y}$ tales que $f^{-1}(s)$ es propio sobre $\boldsymbol{k}(s)$ es \emph{constructible} (y, por tanto, coincide con todo $\mathrm{Y}$ si contiene los puntos cerrados de $\mathrm{Y}$, por el teorema de los ceros de Hilbert \sref[I]{I.10.4.8}); esto equivale una vez más a decir que, si $f^{-1}(y)$ \emph{no es propio} sobre $\boldsymbol{k}(y)$, existe una vecindad abierta $\mathrm{U}$ de $y$ tal que $f^{-1}(s)$ no es propio sobre $\boldsymbol{k}(s)$ para todo $s\in\mathrm{U}$ (cf.\ \sref{IV.9.6.1},~(iv))).
Este ejemplo ilustra el interés de desarrollar sistemáticamente \emph{criterios de constructibilidad}: esto es lo que se hará en la \textsection\,9.
\end{env}

\subsection{Límites proyectivos de preesquemas}
\label{subsection:IV.8.2}

\begin{env}[8.2.1]
\label{IV.8.2.1}
Sea $\mathrm{S}_0$ un espacio anillado, $\mathrm{L}$ un conjunto creciente preordenado filtrado, $(\sh{A}_\lambda,\varphi_{\mu\lambda})$ un sistema inductivo de $\sh{O}_{\mathrm{S}_0}$-Álgebras (no necesariamente conmutativas), cuyo conjunto de índices es $\mathrm{L}$.
Sabemos que, considerado como sistema inductivo de $\sh{O}_{\mathrm{S}_0}$-Módulos, $(\sh{A}_\lambda,\varphi_{\mu\lambda})$ admite un límite inductivo $\sh{A}$; denotemos por $\varphi_\lambda:\sh{A}_\lambda\to\sh{A}$ el homomorfismo canónico de $\sh{O}_{\mathrm{S}_0}$-Módulos.
Sea $m_\lambda:\sh{A}_\lambda\otimes\sh{A}_\lambda\to\sh{A}_\lambda$ el homomorfismo de $\sh{O}_{\mathrm{S}_0}$-Módulos que define la multiplicación en el $\sh{O}_{\mathrm{S}_0}$-Álgebra $\sh{A}_\lambda$; la hipótesis sobre los $\varphi_{\mu\lambda}$ implica que los $m_\lambda$ forman un sistema inductivo de homomorfismos,
\oldpage[IV-3]{8}
y, puesto que el funtor $\varinjlim$ conmuta con los productos tensoriales, $m=\varinjlim m_\lambda$ es un homomorfismo $\sh{A}\otimes\sh{A}\to\sh{A}$ de $\sh{O}_{\mathrm{S}_0}$-Módulos; pasando al límite en los diagramas conmutativos que expresan la asociatividad de $m_\lambda$ y la existencia de la sección identidad en $\sh{A}_\lambda$, se ve que $m$ dota a $\sh{A}$ de una estructura de $\sh{O}_{\mathrm{S}_0}$-Álgebra y que $\varphi_\lambda$ es un homomorfismo de $\sh{O}_{\mathrm{S}_0}$-Álgebras para todo $\lambda\in\mathrm{L}$.
Además, $\sh{A}$ es el límite inductivo del sistema $(\sh{A}_\lambda,\varphi_{\mu\lambda})$ \emph{en la categoría de $\sh{O}_{\mathrm{S}_0}$-Álgebras}; en otras palabras, para todo $\sh{O}_{\mathrm{S}_0}$-Álgebra $\sh{B}$, la aplicación canónica
\begin{equation}
  \operatorname{Hom}_{\sh{O}_{\mathrm{S}_0}\text{-}\mathrm{Alg.}}(\sh{A},\sh{B})\to\varprojlim\operatorname{Hom}_{\sh{O}_{\mathrm{S}_0}\text{-}\mathrm{Alg.}}(\sh{A}_\lambda,\sh{B})
  \tag{8.2.1.1}
\label{IV.8.2.1.1}
\end{equation}
que asocia a cada homomorfismo $f:\sh{A}\to\sh{B}$ de $\sh{O}_{\mathrm{S}_0}$-Álgebras la familia $(f\circ\varphi_\lambda)$ es \emph{biyectiva}.
En efecto, ya sabemos que es inyectiva e identifica $\operatorname{Hom}_{\sh{O}_{\mathrm{S}_0}\text{-}\mathrm{Alg.}}(\sh{A},\sh{B})$ con un subconjunto de $\varprojlim\operatorname{Hom}_{\sh{O}_{\mathrm{S}_0}\text{-}\mathrm{Mod.}}(\sh{A}_\lambda,\sh{B})$; todo se reduce a ver que, si $(f_\lambda)$ es un sistema inductivo de homomorfismos de $\sh{O}_{\mathrm{S}_0}$-Álgebras, $f_\lambda:\sh{A}_\lambda\to\sh{B}$, su límite inductivo $f:\sh{A}\to\sh{B}$, que por definición es un homomorfismo de $\sh{O}_{\mathrm{S}_0}$-\emph{Módulos}, es también un homomorfismo de $\sh{O}_{\mathrm{S}_0}$-\emph{Álgebras}; pero esto resulta de pasar al límite inductivo en el diagrama conmutativo de $\sh{O}_{\mathrm{S}_0}$-Módulos
\[
\xymatrix{
  \sh{A}_\lambda\otimes\sh{A}_\lambda \ar[r]^{f_\lambda\otimes f_\lambda} \ar[d]_{m_\lambda} & \sh{B}\otimes\sh{B} \ar[d] \\
  \sh{A}_\lambda \ar[r]_{f_\lambda} & \sh{B}
}
\]
y del hecho de que el funtor $\varinjlim$ conmuta con los productos tensoriales.
Observamos finalmente que, si los $\sh{A}_\lambda$ son $\sh{O}_{\mathrm{S}_0}$-Álgebras conmutativas, lo mismo vale para $\sh{A}$.
\end{env}

\begin{env}[8.2.2]
\label{IV.8.2.2}
Supongamos ahora que $\mathrm{S}_0$ es un \emph{preesquema} y que las $\sh{A}_\lambda$ son $\sh{O}_{\mathrm{S}_0}$-Álgebras (conmutativas) \emph{cuasi-coherentes}; sabemos entonces que $\sh{A}=\varinjlim\sh{A}_\lambda$ es un $\sh{O}_{\mathrm{S}_0}$-Álgebra cuasi-coherente \sref[I]{I.4.1.1}.
Denotemos por $\mathrm{S}_\lambda$ (resp. $\mathrm{S}$) el espectro del $\sh{O}_{\mathrm{S}_0}$-Álgebra $\sh{A}_\lambda$ (resp. $\sh{A}$) \sref[II]{II.1.3.1}, y sean $u_{\lambda\mu}:\mathrm{S}_\mu\to\mathrm{S}_\lambda$ (para $\lambda\leqslant\mu$) y $u_\lambda:\mathrm{S}\to\mathrm{S}_\lambda$ los $\mathrm{S}_0$-morfismos correspondientes a los homomorfismos $\varphi_{\mu\lambda}$ y $\varphi_\lambda$, respectivamente \sref[II]{II.1.2.7}; está claro que $(\mathrm{S}_\lambda,u_{\lambda\mu})$ es un \emph{sistema proyectivo} en la categoría de preesquemas $\mathrm{S}_0$.
Observamos que los $u_{\lambda\mu}$ y $u_\lambda$ son morfismos \emph{afines} \sref[II]{II.1.6.2}, por tanto \emph{cuasi-compactos} y \emph{separados}.
\end{env}

\begin{proposition}[8.2.3]
\label{IV.8.2.3}
Con las notaciones de \sref{IV.8.2.2}, los morfismos $u_\lambda:\mathrm{S}\to\mathrm{S}_\lambda$ hacen de $\mathrm{S}$ un límite proyectivo del sistema proyectivo $(\mathrm{S}_\lambda,u_{\lambda\mu})$ en la categoría de preesquemas.
Además, si $h:\mathrm{S}_0\to\mathrm{T}$ es un morfismo que convierte todo preesquema $\mathrm{S}_0$ en un preesquema $\mathrm{T}$, $\mathrm{S}$ es también un límite proyectivo del sistema $(\mathrm{S}_\lambda,u_{\lambda\mu})$ en la categoría de preesquemas $\mathrm{T}$.
\end{proposition}

\begin{proof}
Demostremos primero la segunda afirmación del enunciado en el caso $\mathrm{T}=\mathrm{S}_0$.
\oldpage[IV-3]{9}
Todo se reduce a demostrar que, si $\mathrm{X}$ es un preesquema $\mathrm{S}_0$ cualquiera, la aplicación canónica
\begin{equation}
  \operatorname{Hom}_{\mathrm{S}_0}(\mathrm{X},\mathrm{S})\to\varprojlim\operatorname{Hom}_{\mathrm{S}_0}(\mathrm{X},\mathrm{S}_\lambda)
  \tag{8.2.3.1}
\label{IV.8.2.3.1}
\end{equation}
que asocia a cada morfismo $\mathrm{S}_0$ $v:\mathrm{X}\to\mathrm{S}$ la familia $(u_\lambda\circ v)$ es \emph{biyectiva}.
Ahora bien, si $g:\mathrm{X}\to\mathrm{S}_0$ es el morfismo estructural y ponemos $\sh{B}=g_*(\sh{O}_\mathrm{X})$, que es un $\sh{O}_{\mathrm{S}_0}$-Álgebra, la aplicación \sref{IV.8.2.3.1} se identifica canónicamente con \sref{IV.8.2.1.1} \sref[II]{II.1.2.7}, y la conclusión se sigue de lo visto en \sref{IV.8.2.1}.

Las demás afirmaciones de \sref{IV.8.2.3} son consecuencias del siguiente lema general:
\end{proof}

\begin{lemma}[8.2.4]
\label{IV.8.2.4}
Sea $\mathbf{C}$ una categoría, $\mathrm{T}$ un objeto de $\mathbf{C}$ y $\mathbf{C}_\mathrm{T}$ la subcategoría de objetos de $\mathbf{C}$ sobre $\mathrm{T}$.
Sea $(\mathrm{S}_\lambda,u_{\lambda\mu})$ un sistema proyectivo en $\mathbf{C}_\mathrm{T}$; entonces todo límite proyectivo de este sistema en $\mathbf{C}_\mathrm{T}$ es también un límite proyectivo en $\mathbf{C}$, y recíprocamente.
\end{lemma}

\begin{proof}
Sea $f_\lambda:\mathrm{S}_\lambda\to\mathrm{T}$ el morfismo estructural.
Supongamos que $\mathrm{S}$ es un límite proyectivo de $(\mathrm{S}_\lambda,u_{\lambda\mu})$ en $\mathbf{C}$, y denotemos por $u_\lambda:\mathrm{S}\to\mathrm{S}_\lambda$ los morfismos canónicos correspondientes.
Consideremos un sistema proyectivo de morfismos $\mathrm{T}$ (en $\mathbf{C}$) $w_\lambda:\mathrm{Y}\to\mathrm{S}_\lambda$, donde $\mathrm{Y}\in\mathbf{C}_\mathrm{T}$.
Por hipótesis existe un único morfismo (en $\mathbf{C}$) $w:\mathrm{Y}\to\mathrm{S}$ tal que $w_\lambda=u_\lambda\circ w$ para todo $\lambda$.
La hipótesis de que los $u_\lambda$ son morfismos $\mathrm{T}$ implica que los morfismos $f_\lambda\circ u_\lambda:\mathrm{S}\to\mathrm{T}$ son todos iguales, y este compuesto $f$ dota a $\mathrm{S}$ de una estructura de objeto $\mathrm{T}$.
Si $g:\mathrm{Y}\to\mathrm{T}$ es el morfismo estructural de $\mathrm{Y}$, se tiene $f\circ w=f_\lambda\circ u_\lambda\circ w=f_\lambda\circ w_\lambda=g$ para todo $\lambda$, lo que demuestra que $w$ es un morfismo $\mathrm{T}$.
Recíprocamente, supongamos (con las mismas notaciones) que $\mathrm{S}$ es un límite proyectivo de $(\mathrm{S}_\lambda,u_{\lambda\mu})$ en $\mathbf{C}_\mathrm{T}$, y consideremos ahora un sistema proyectivo de morfismos (en $\mathbf{C}$) $w_\lambda:\mathrm{Y}\to\mathrm{S}_\lambda$.
Los compuestos $f_\lambda\circ w_\lambda:\mathrm{Y}\to\mathrm{T}$ son entonces todos iguales: en efecto, para dos índices cualesquiera $\lambda$, $\mu$, existe un índice $\nu$ tal que $\lambda\leqslant\nu$ y $\mu\leqslant\nu$, de donde $f_\nu=f_\lambda\circ u_{\lambda\nu}=f_\mu\circ u_{\mu\nu}$; como $w_\lambda=u_{\lambda\nu}\circ w_\nu$ y $w_\mu=u_{\mu\nu}\circ w_\nu$, se tiene $f_\lambda\circ w_\lambda=f_\lambda\circ u_{\lambda\nu}\circ w_\nu=f_\nu\circ w_\nu$ y $f_\mu\circ w_\mu=f_\mu\circ u_{\mu\nu}\circ w_\nu=f_\nu\circ w_\nu$.
Si $g:\mathrm{Y}\to\mathrm{T}$ es el único morfismo así definido, $g$ dota a $\mathrm{Y}$ de una estructura de objeto $\mathrm{T}$ y los $w_\lambda$ son morfismos $\mathrm{T}$; por tanto tienen un límite proyectivo $w:\mathrm{Y}\to\mathrm{S}$ que es un morfismo $\mathrm{T}$ y, \emph{a fortiori}, un morfismo de $\mathbf{C}$; además, la primera parte del argumento muestra que todo límite proyectivo $w'$ (en $\mathbf{C}$) del sistema proyectivo $(w_\lambda)$ es también un morfismo $\mathrm{T}$ y, por consiguiente, igual a $w$, lo que concluye la demostración del lema.
\end{proof}

\begin{proposition}[8.2.5]
\label{IV.8.2.5}
Bajo las condiciones de \sref{IV.8.2.2}, sea $\alpha$ un elemento de $\mathrm{L}$ y $\mathrm{X}_\alpha$ un preesquema $\mathrm{S}_\alpha$.
Para todo $\lambda\geqslant\alpha$, pongamos $\mathrm{X}_\lambda=\mathrm{X}_\alpha\times_{\mathrm{S}_\alpha}\mathrm{S}_\lambda$ y, para $\alpha\leqslant\lambda\leqslant\mu$, pongamos $v_{\lambda\mu}=1_{\mathrm{X}_\alpha}\times u_{\lambda\mu}$, de modo que $(\mathrm{X}_\lambda,v_{\lambda\mu})$ sea un sistema proyectivo de preesquemas $\mathrm{X}_\alpha$, cuyo conjunto de índices está formado por los $\lambda\geqslant\alpha$ de $\mathrm{L}$.
Pongamos asimismo $\mathrm{X}=\mathrm{X}_\alpha\times_{\mathrm{S}_\alpha}\mathrm{S}$ y $v_\lambda=1_{\mathrm{X}_\alpha}\times u_\lambda$.
Entonces los morfismos $\mathrm{X}_\alpha$ $v_\lambda:\mathrm{X}\to\mathrm{X}_\lambda$ hacen de $\mathrm{X}$ un límite proyectivo del sistema proyectivo $(\mathrm{X}_\lambda,v_{\lambda\mu})$ en la categoría de preesquemas $\mathrm{X}_\alpha$, o en la categoría de todos los preesquemas.
\end{proposition}

\begin{proof}
Esto también se deducirá del siguiente lema general:
\end{proof}

\begin{lemma}[8.2.6]
\label{IV.8.2.6}
Sea $\mathbf{C}$ una categoría en la que existen productos fibrados, $q:\mathrm{T}'\to\mathrm{T}$ un morfismo de $\mathbf{C}$, y $\mathbf{C}_\mathrm{T}$ (resp. $\mathbf{C}_{\mathrm{T}'}$) la categoría de objetos de $\mathbf{C}$ sobre
\oldpage[IV-3]{10}
$\mathrm{T}$ (resp. $\mathrm{T}'$).
Sea $(\mathrm{S}_\lambda,u_{\lambda\mu})$ un sistema proyectivo (no necesariamente filtrado) en $\mathbf{C}_\mathrm{T}$, y pongamos $\mathrm{S}'_\lambda=\mathrm{S}_\lambda\times_\mathrm{T}\mathrm{T}'$, $u'_{\lambda\mu}=u_{\lambda\mu}\times 1_{\mathrm{T}'}$, de modo que $(\mathrm{S}'_\lambda,u'_{\lambda\mu})$ sea un sistema proyectivo en $\mathbf{C}_{\mathrm{T}'}$.
Entonces, si $(\mathrm{S},u_\lambda)$ es un límite proyectivo de $(\mathrm{S}_\lambda,u_{\lambda\mu})$ en $\mathbf{C}_\mathrm{T}$, $(\mathrm{S}\times_\mathrm{T}\mathrm{T}',u_\lambda\times 1_{\mathrm{T}'})$ es un límite proyectivo de $(\mathrm{S}'_\lambda,u'_{\lambda\mu})$ en $\mathbf{C}_{\mathrm{T}'}$.
\end{lemma}

\begin{proof}
Por hipótesis, para todo $\lambda$ se tiene un diagrama conmutativo
\[
\xymatrix{
  \mathrm{S}' \ar[r]^{u'_\lambda} \ar[d]_p & \mathrm{S}'_\lambda \ar[r]^{h'_\lambda} \ar[d]_{p_\lambda} & \mathrm{T}' \ar[d]^q \\
  \mathrm{S} \ar[r]_{u_\lambda} & \mathrm{S}_\lambda \ar[r]_{h_\lambda} & \mathrm{T}
}
\]
donde hemos puesto $\mathrm{S}'=\mathrm{S}\times_\mathrm{T}\mathrm{T}'$, $u'_\lambda=u_\lambda\times 1_{\mathrm{T}'}$, $h'_\lambda=h_\lambda\times 1_{\mathrm{T}'}$.
Sea $\mathrm{Y}$ un objeto $\mathrm{T}'$, $g':\mathrm{Y}\to\mathrm{T}'$, y consideremos un sistema proyectivo de morfismos $\mathrm{T}'$ $w'_\lambda:\mathrm{Y}\to\mathrm{S}'_\lambda$.
Entonces $\mathrm{Y}$ es un objeto $\mathrm{T}$ para el morfismo $g=q\circ g'$, y los $w_\lambda=p_\lambda\circ w'_\lambda$ son morfismos $\mathrm{T}$, pues $h_\lambda\circ w_\lambda=h_\lambda\circ p_\lambda\circ w'_\lambda=q\circ h'_\lambda\circ w'_\lambda=q\circ g'$ por hipótesis.
Además, se verifica inmediatamente que $(w_\lambda)$ es un sistema proyectivo.
Por hipótesis existe entonces un único morfismo $\mathrm{T}$ $w:\mathrm{Y}\to\mathrm{S}$ tal que $u_\lambda\circ w=w_\lambda$ para todo $\lambda$.
Por la definición del producto fibrado, existe un único morfismo $\mathrm{T}'$ $w':\mathrm{Y}\to\mathrm{S}'$ tal que $p\circ w'=w$.
Entonces $u_\lambda\circ p\circ w'=u_\lambda\circ w=w_\lambda=p_\lambda\circ w'_\lambda$, lo que también puede escribirse $p_\lambda\circ u'_\lambda\circ w'=p_\lambda\circ w'_\lambda$; por otra parte, como $w'$ es un morfismo $\mathrm{T}'$, $h'_\lambda\circ u'_\lambda\circ w'=g'=h'_\lambda\circ w'_\lambda$. La definición de $\mathrm{S}'_\lambda$ como producto fibrado $\mathrm{S}_\lambda\times_\mathrm{T}\mathrm{T}'$ da por tanto $u'_\lambda\circ w'=w'_\lambda$, y es inmediato que $w'$ es el único morfismo $\mathrm{T}'$ que satisface estas relaciones, de donde se sigue el lema.
\end{proof}

\begin{remark}[8.2.7]
\label{IV.8.2.7}
Dado un espacio anillado $\mathrm{S}$, los límites inductivos respecto de un conjunto preordenado $\mathrm{L}$ (no necesariamente filtrado) existen en la categoría de $\sh{O}_\mathrm{S}$-Álgebras conmutativas, pues el límite inductivo filtrado existe por \sref{IV.8.2.1} y, por otra parte, para dos homomorfismos cualesquiera de $\sh{O}_\mathrm{S}$-Álgebras $\sh{A}\to\sh{B}$, $\sh{A}\to\sh{C}$, el producto tensorial $\sh{B}\otimes_{\sh{A}}\sh{C}$ es la «suma amalgamada» correspondiente en esta categoría.

Cuando $\mathrm{S}$ es un preesquema, sabemos que el producto tensorial $\sh{B}\otimes_{\sh{A}}\sh{C}$ es un $\sh{O}_\mathrm{S}$-Álgebra cuasi-coherente cuando esto vale para $\sh{A}$, $\sh{B}$ y $\sh{C}$ \sref[I]{I.1.3.13}; de ello se concluye que, en la categoría de $\sh{O}_\mathrm{S}$-Álgebras cuasi-coherentes, existen siempre límites inductivos para cualquier conjunto de índices preordenado.
Esto permite generalizar la definición de límite proyectivo de preesquemas y las proposiciones \sref{IV.8.2.3} y \sref{IV.8.2.5} al caso en que el conjunto preordenado $\mathrm{L}$ no sea necesariamente filtrado.
\end{remark}

\begin{env}[8.2.8]
\label{IV.8.2.8}
Con las notaciones de \sref{IV.8.2.2}, pongamos $u_{\lambda\mu}=(\psi_{\lambda\mu},\theta_{\lambda\mu})$ y $u_\lambda=(\psi_\lambda,\theta_\lambda)$; $(\mathrm{S}_\lambda,\psi_{\lambda\mu})$ es así un sistema proyectivo de espacios topológicos y $(\psi_\lambda)$ un sistema proyectivo de aplicaciones continuas $\mathrm{S}\to\mathrm{S}_\lambda$ entre los espacios subyacentes, respectivamente, de los preesquemas $\mathrm{S}$ y $\mathrm{S}_\lambda$.
\end{env}

\begin{proposition}[8.2.9]
\label{IV.8.2.9}
Con las notaciones de \sref{IV.8.2.8}, el límite proyectivo del sistema proyectivo $(\psi_\lambda)$ de aplicaciones continuas es un homeomorfismo del espacio subyacente de $\mathrm{S}$ sobre el límite proyectivo del sistema proyectivo $(\mathrm{S}_\lambda,\psi_{\lambda\mu})$ de espacios topológicos.
\end{proposition}

\begin{proof}
\oldpage[IV-3]{11}
Sea $\mathrm{T}$ el espacio límite topológico del sistema $(\mathrm{S}_\lambda,\psi_{\lambda\mu})$ y pongamos $\psi=\varprojlim\psi_\lambda:\mathrm{S}\to\mathrm{T}$.
Podemos restringirnos al caso en que $\mathrm{S}_0=\mathrm{S}_\alpha$ para algún $\alpha\in\mathrm{L}$ y $\lambda\geqslant\alpha$.

\medskip
\noindent
(i) Mostremos primero que la topología de $\mathrm{S}$ es la imagen inversa por $\psi$ de la topología de $\mathrm{T}$; en otras palabras, si $\pi_\lambda:\mathrm{T}\to\mathrm{S}_\lambda$ es la aplicación canónica, se trata de demostrar que todo abierto de $\mathrm{S}$ es unión de abiertos de la forma $\psi^{-1}(\pi_\lambda^{-1}(\mathrm{U}_\lambda))$, donde $\mathrm{U}_\lambda$ es abierto en $\mathrm{S}_\lambda$.
Ahora bien, por definición todo abierto de $\mathrm{S}$ es unión de abiertos obtenidos del modo siguiente: se considera un abierto afín $\mathrm{U}_0$ de $\mathrm{S}_0$, de anillo $\mathrm{A}_0$, de modo que $\psi_\alpha^{-1}(\mathrm{U}_0)$ sea el abierto afín de $\mathrm{S}$ de anillo $\mathrm{A}=\Gamma(\mathrm{U}_0,\sh{A})$; después se toma un elemento $f\in\mathrm{A}$ y se considera en $\operatorname{Spec}(\mathrm{A})$, identificado con $\psi_\alpha^{-1}(\mathrm{U}_0)$, el abierto $\mathrm{D}(f)$.
Estos son los abiertos $\mathrm{D}(f)$ que forman una base de la topología de $\mathrm{S}$ \sref[II]{II.1.3.1}.
Ahora, si ponemos $\mathrm{A}_\lambda=\Gamma(\mathrm{U}_0,\sh{A}_\lambda)$, se tiene $\mathrm{A}=\varinjlim\mathrm{A}_\lambda$ \sref[I]{I.1.3.9}, por lo que existe un índice $\lambda$ tal que $f$ es la imagen canónica de un elemento $f_\lambda\in\mathrm{A}_\lambda$; entonces $\mathrm{D}(f)=\psi_\lambda^{-1}(\mathrm{D}(f_\lambda))$ \sref[I]{I.1.2.2}, y como $\psi_\lambda=\pi_\lambda\circ\psi$, queda demostrada nuestra afirmación.

\medskip
\noindent
(ii) Mostremos ahora que $\psi$ es \emph{biyectiva}, lo que concluirá la demostración.
Como $\mathrm{S}$ es un espacio de Kolmogoroff, de (i) se deduce ya que $\psi$ es inyectiva, y queda por demostrar que $\psi$ es sobreyectiva.
Evidentemente podemos sustituir $\mathrm{T}$ por un abierto $\pi_\alpha^{-1}(\mathrm{U}_0)$, donde $\mathrm{U}_0$ es un abierto afín de $\mathrm{S}_\alpha=\mathrm{S}_0$; por tanto podemos restringirnos al caso en que los $\mathrm{S}_\lambda$ y $\mathrm{S}$ sean afines, es decir, $\sh{A}_\lambda$ sea el haz de álgebras asociado a un $\mathrm{A}_0$-álgebra $\mathrm{A}_\lambda$, y $\sh{A}$ el haz de álgebras asociado a $\mathrm{A}=\varinjlim\mathrm{A}_\lambda$; seguiremos denotando por $\varphi_{\mu\lambda}:\mathrm{A}_\lambda\to\mathrm{A}_\mu$ y $\varphi_\lambda:\mathrm{A}_\lambda\to\mathrm{A}$ los homomorfismos canónicos.
Por definición, un elemento de $\mathrm{T}$ es una familia $(\mathfrak{p}_\lambda)_{\lambda\in\mathrm{L}}$, donde $\mathfrak{p}_\lambda$ es un ideal primo de $\mathrm{A}_\lambda$ y donde $\mathfrak{p}_\lambda=\varphi_{\mu\lambda}^{-1}(\mathfrak{p}_\mu)$ para $\lambda\leqslant\mu$.
Sabemos (\sref{IV.5.13.3} y \sref{IV.5.13.1}) que existe un ideal primo $\mathfrak{p}$ de $\mathrm{A}$ tal que $\mathfrak{p}_\lambda=\varphi_\lambda^{-1}(\mathfrak{p})$ para todo $\lambda\in\mathrm{L}$, lo que concluye la demostración.
\end{proof}

En particular, hemos demostrado así el siguiente

\begin{corollary}[8.2.10]
\label{IV.8.2.10}
Sea $(\mathrm{A}_\lambda)_{\lambda\in\mathrm{L}}$ un sistema inductivo filtrado de anillos, y sean $\mathrm{A}=\varinjlim\mathrm{A}_\lambda$, $\varphi_\lambda:\mathrm{A}_\lambda\to\mathrm{A}$ los homomorfismos canónicos.
La aplicación canónica $\mathfrak{p}\mapsto(\varphi_\lambda^{-1}(\mathfrak{p}))$ es un homeomorfismo de $\operatorname{Spec}(\mathrm{A})$ sobre el espacio topológico $\varprojlim\operatorname{Spec}(\mathrm{A}_\lambda)$.
\end{corollary}

\begin{corollary}[8.2.11]
\label{IV.8.2.11}
Con las notaciones de \sref{IV.8.2.8}, para todo abierto cuasi-compacto $\mathrm{U}$ de $\mathrm{S}$ existen un índice $\lambda$ y un abierto cuasi-compacto $\mathrm{U}_\lambda$ de $\mathrm{S}_\lambda$ tales que $\mathrm{U}=\psi_\lambda^{-1}(\mathrm{U}_\lambda)$.
\end{corollary}

\begin{proof}
Esto se deduce del hecho de que, por definición de la topología del límite proyectivo, los $\psi_\lambda^{-1}(\mathrm{U}_\lambda)$ ($\mathrm{U}_\lambda$ abierto cuasi-compacto de $\mathrm{S}_\lambda$) forman una base de la topología de $\mathrm{S}$, y del hecho de que el conjunto de índices $\mathrm{L}$ es filtrado.
\end{proof}

\begin{corollary}[8.2.12]
\label{IV.8.2.12}
Con las notaciones de \sref{IV.8.2.8}, el límite inductivo del sistema inductivo de homomorfismos $\theta_\lambda^\sharp:\psi_\lambda^*(\sh{O}_{\mathrm{S}_\lambda})\to\sh{O}_\mathrm{S}$ de haces de anillos sobre $\mathrm{S}$ es un isomorfismo
\begin{equation}
  \varinjlim_\lambda\psi_\lambda^*(\sh{O}_{\mathrm{S}_\lambda})\xrightarrow{\sim}\sh{O}_\mathrm{S}.
  \label{IV.8.2.12.1}
\tag{8.2.12.1}
\end{equation}
\end{corollary}

\begin{proof}
Podemos suponer evidentemente que los $\mathrm{S}_\lambda$ son afines; con las notaciones de la demostración de \sref{IV.8.2.9}, todo se reduce a comprobar que el límite inductivo del sistema inductivo de aplicaciones canónicas $(\mathrm{A}_\lambda)_{\mathfrak{p}_\lambda}\to\mathrm{A}_\mathfrak{p}$ es un isomorfismo, lo cual no es otra cosa que \sref{IV.5.13.3},~(ii).
\end{proof}

\begin{proposition}[8.2.13]
\label{IV.8.2.13}
Supongamos que los morfismos $u_{\lambda\mu}:\mathrm{S}_\mu\to\mathrm{S}_\lambda$ son inmersiones abiertas, de modo que $\mathrm{S}_\mu$ se identifica con el subesquema inducido en un abierto de $\mathrm{S}_\lambda$ para $\lambda\leqslant\mu$.
Entonces, para todo $\alpha\in\mathrm{L}$, el espacio subyacente del preesquema $\mathrm{S}$ se identifica con el subespacio de $\mathrm{S}_\alpha$ dado por la intersección de los $\mathrm{S}_\lambda$ para $\lambda\geqslant\alpha$, y el haz estructural $\sh{O}_\mathrm{S}$ se identifica con el haz inducido \textnormal{(}$\mathrm{G}$, II, $1.5$\textnormal{)} por $\sh{O}_{\mathrm{S}_\alpha}$ sobre esta intersección; más generalmente, para todo $\sh{O}_{\mathrm{S}_\alpha}$-Módulo $\sh{F}_\alpha$, $u_\alpha^*(\sh{F}_\alpha)$ se identifica con el $\sh{O}_\mathrm{S}$-Módulo inducido por $\sh{F}_\alpha$ sobre $\mathrm{S}$.
\end{proposition}

\begin{proof}
\oldpage[IV-3]{12}
La primera afirmación se deduce de \sref{IV.8.2.9}, teniendo en cuenta la definición de límite proyectivo de espacios topológicos; además, todos los $\psi_\lambda^*(\sh{O}_{\mathrm{S}_\lambda})$ son iguales al haz inducido por $\sh{O}_{\mathrm{S}_\alpha}$ sobre $\mathrm{S}$ por definición \sref[$0_{\mathrm{I}}$]{0.3.7.1}, y, con las notaciones de la demostración de \sref{IV.8.2.9}, los homomorfismos $(\mathrm{A}_\lambda)_{\mathfrak{p}_\lambda}\to(\mathrm{A}_\mu)_{\mathfrak{p}_\mu}$ son biyectivos para un sistema $(\mathfrak{p}_\lambda)$ de ideales primos correspondiente a un único punto de $\mathrm{S}$; la afirmación relativa a $\sh{O}_\mathrm{S}$ se deduce entonces de \sref{IV.8.2.12}.
La última afirmación se deduce entonces de la definición de $u_\alpha^*(\sh{F}_\alpha)$ \sref[$0_{\mathrm{I}}$]{0.4.3.1}.
\end{proof}

\begin{remark}[8.2.14]
\label{IV.8.2.14}
Los resultados de \sref{IV.8.2.9} y \sref{IV.8.2.12} muestran que $\mathrm{S}$ es un límite proyectivo del sistema proyectivo $(\mathrm{S}_\lambda,u_{\lambda\mu})$ en la categoría de \emph{todos} los espacios anillados (o de todos los espacios localmente anillados).
En efecto, sea $\mathrm{Y}$ un espacio anillado y consideremos un sistema proyectivo de morfismos de espacios anillados $w_\lambda:\mathrm{Y}\to\mathrm{S}_\lambda$.
Si se escribe $w_\lambda=(\rho_\lambda,\omega_\lambda)$, los $\rho_\lambda$ forman un sistema proyectivo de aplicaciones continuas y, por \sref{IV.8.2.9}, su límite proyectivo se identifica con una aplicación continua $\rho:\mathrm{Y}\to\mathrm{S}$ tal que $\rho_\lambda=\psi_\lambda\circ\rho$. Por otra parte, los homomorfismos
\[
 \omega_\lambda^\sharp:\rho_\lambda^*(\sh{O}_{\mathrm{S}_\lambda})\longrightarrow\sh{O}_\mathrm{Y}
\]
forman un sistema inductivo de homomorfismos de haces de anillos. Puesto que
\[
 \rho_\lambda^*(\sh{O}_{\mathrm{S}_\lambda})
 =\rho^*\bigl(\psi_\lambda^*(\sh{O}_{\mathrm{S}_\lambda})\bigr)
\]
y el funtor $\rho^*$ es exacto, el límite inductivo de los $\rho_\lambda^*(\sh{O}_{\mathrm{S}_\lambda})$ es $\rho^*(\sh{O}_\mathrm{S})$ por \sref{IV.8.2.12}. Existe por tanto un único homomorfismo
\[
 \omega^\sharp:\rho^*(\sh{O}_\mathrm{S})\longrightarrow\sh{O}_\mathrm{Y}
\]
tal que $\omega_\lambda^\sharp=\omega^\sharp\circ\rho^*(\theta_\lambda^\sharp)$ para todo $\lambda$, lo que demuestra la afirmación.
\end{remark}

\subsection{Subconjuntos constructibles de un límite proyectivo de preesquemas}
\label{subsection:IV.8.3}

\begin{env}[8.3.1]
\label{IV.8.3.1}
En todo lo que sigue en este párrafo suponemos satisfechas las condiciones de \sref{IV.8.2.2} y conservamos las notaciones.
\end{env}

\begin{theorem}[8.3.2]
\label{IV.8.3.2}
Para todo $\lambda$, sean $\mathrm{E}_\lambda$, $\mathrm{F}_\lambda$ dos subconjuntos de $\mathrm{S}_\lambda$.
Pongamos
\begin{equation}
  \mathrm{E}=\bigcap_\lambda u_\lambda^{-1}(\mathrm{E}_\lambda),\qquad \mathrm{F}=\bigcup_\lambda u_\lambda^{-1}(\mathrm{F}_\lambda).
  \tag{8.3.2.1}
\label{IV.8.3.2.1}
\end{equation}
Supongamos satisfechas las condiciones siguientes:
\begin{enumerate}[label=(\roman*)]
   \item Para todo $\lambda$, $\mathrm{E}_\lambda$ es pro-constructible y $\mathrm{F}_\lambda$ es ind-constructible \sref{IV.1.9.4}.
   \item Para $\lambda\leqslant\mu$, se tiene
  \[
    \mathrm{E}_\mu\subset u_{\lambda\mu}^{-1}(\mathrm{E}_\lambda),\qquad \mathrm{F}_\mu\supset u_{\lambda\mu}^{-1}(\mathrm{F}_\lambda).
  \]
   \item Existe $\alpha\in\mathrm{L}$ tal que $\mathrm{S}_\alpha$ es cuasi-compacto (lo que implica que $\mathrm{S}_\lambda$ es cuasi-compacto para $\lambda\geqslant\alpha$).
\end{enumerate}
Entonces las propiedades siguientes son equivalentes:
\begin{enumerate}[label=\emph{\alph*)}]
   \item $\mathrm{E}\subset\mathrm{F}$.
   \item Existe $\lambda\geqslant\alpha$ tal que $u_\lambda^{-1}(\mathrm{E}_\lambda)\subset u_\lambda^{-1}(\mathrm{F}_\lambda)$ (y entonces también se tiene $u_\mu^{-1}(\mathrm{E}_\mu)\subset u_\mu^{-1}(\mathrm{F}_\mu)$ para
\oldpage[IV-3]{13}
$\mu\geqslant\lambda$).
   \item Existe $\lambda\geqslant\alpha$ tal que $\mathrm{E}_\lambda\subset\mathrm{F}_\lambda$ (y entonces también se tiene $\mathrm{E}_\mu\subset\mathrm{F}_\mu$ para $\mu\geqslant\lambda$).
\end{enumerate}
Las observaciones entre paréntesis de \emph{b)} y \emph{c)} se deducen de (ii).
Pongamos
\[
  \mathrm{G}_\lambda=\mathrm{E}_\lambda\cap(\mathrm{S}_\lambda-\mathrm{F}_\lambda),\qquad \mathrm{G}=\mathrm{E}\cap(\mathrm{S}-\mathrm{F}).
\]
Entonces $\mathrm{G}_\lambda$ es un subconjunto pro-constructible de $\mathrm{S}_\lambda$ \sref{IV.1.9.5},~(ii), y, en virtud de \sref{IV.8.3.2.1} y de la condición (ii), se tiene
\[
  \mathrm{G}_\mu\subset u_{\lambda\mu}^{-1}(\mathrm{G}_\lambda)\qquad\text{para }\lambda\leqslant\mu
\]
\[
  \mathrm{G}=\bigcap_\lambda u_\lambda^{-1}(\mathrm{G}_\lambda).
\]
Quedamos así reducidos a demostrar el caso particular de \sref{IV.8.3.2} correspondiente a $\mathrm{F}_\lambda=\varnothing$ para todo $\lambda$:
\end{theorem}

\begin{corollary}[8.3.3]
\label{IV.8.3.3}
Para todo $\lambda$, sea $\mathrm{E}_\lambda$ un subconjunto pro-constructible de $\mathrm{S}_\lambda$ tal que, para $\lambda\leqslant\mu$, $\mathrm{E}_\mu\subset u_{\lambda\mu}^{-1}(\mathrm{E}_\lambda)$.
Supongamos que existe $\alpha\in\mathrm{L}$ tal que $\mathrm{S}_\alpha$ es cuasi-compacto.
Entonces las condiciones siguientes son equivalentes:
\begin{enumerate}[label=\emph{\alph*)}]
   \item $\mathrm{E}=\bigcap_\lambda u_\lambda^{-1}(\mathrm{E}_\lambda)=\varnothing$.
   \item Existe $\lambda$ tal que $u_\lambda^{-1}(\mathrm{E}_\lambda)=\varnothing$ (y entonces $u_\mu^{-1}(\mathrm{E}_\mu)=\varnothing$ para $\mu\geqslant\lambda$).
   \item Existe $\lambda$ tal que $\mathrm{E}_\lambda=\varnothing$ (y entonces $\mathrm{E}_\mu=\varnothing$ para $\mu\geqslant\lambda$).
\end{enumerate}
\end{corollary}

\begin{proof}
Es claro que \emph{c)} implica \emph{a)}.
Demostremos que \emph{a)} implica \emph{b)}: como $\mathrm{S}_\lambda$ es cuasi-compacto, lo mismo vale para $\mathrm{S}$ \sref{IV.8.2.2}; como $\mathrm{E}_\lambda$ es pro-constructible, lo mismo vale para $u_\lambda^{-1}(\mathrm{E}_\lambda)$ \sref{IV.1.9.5},~(vi); por tanto, la familia filtrada decreciente de subconjuntos pro-constructibles $u_\lambda^{-1}(\mathrm{E}_\lambda)$ tiene intersección vacía, y por \sref{IV.1.9.9} uno de ellos es vacío.

Finalmente, demostremos que \emph{b)} implica \emph{c)}.
Como $\mathrm{S}_\alpha$ es cuasi-compacto y $\mathrm{L}$ es filtrado, podemos sustituir $\mathrm{S}_\alpha$ por un afín abierto; supongamos entonces \sref{IV.8.2.2} que $\mathrm{S}$ y los $\mathrm{S}_\lambda$ para $\lambda\geqslant\alpha$ son afines; se tiene entonces \sref{IV.1.9.2.1}, para $\lambda\geqslant\alpha$,
\[
  u_\lambda(\mathrm{S})=\bigcap_{\mu\geqslant\lambda}u_{\lambda\mu}(\mathrm{S}_\mu)
\]
de donde
\[
  \mathrm{E}_\lambda\cap u_\lambda(\mathrm{S})=\bigcap_{\mu\geqslant\lambda}(\mathrm{E}_\lambda\cap u_{\lambda\mu}(\mathrm{S}_\mu)).
\]
Ahora bien, como $u_\lambda$ y los $u_{\lambda\mu}$ son cuasi-compactos, $u_\lambda(\mathrm{S})$ y $u_{\lambda\mu}(\mathrm{S}_\mu)$ son pro-constructibles en $\mathrm{S}_\lambda$ \sref{IV.1.9.5},~(viii); por tanto, los conjuntos $\mathrm{E}_\lambda\cap u_{\lambda\mu}(\mathrm{S}_\mu)$ para $\mu\geqslant\lambda$ forman una familia filtrada decreciente de subconjuntos pro-constructibles de $\mathrm{S}_\lambda$.
Como $\mathrm{S}_\lambda$ es cuasi-compacto, la hipótesis \emph{b)} implica que la intersección de esta familia es vacía y, por \sref{IV.1.9.9}, uno de los conjuntos $\mathrm{E}_\lambda\cap u_{\lambda\mu}(\mathrm{S}_\mu)$ es vacío; por tanto, $\mathrm{E}_\mu\subset u_{\lambda\mu}^{-1}(\mathrm{E}_\lambda)$ es vacío.
\end{proof}

\begin{corollary}[8.3.4]
\label{IV.8.3.4}
Para todo $\lambda$, sea $\mathrm{F}_\lambda$ un subconjunto ind-constructible de $\mathrm{S}_\lambda$ tal que, para $\lambda\leqslant\mu$, se tenga $u_{\lambda\mu}^{-1}(\mathrm{F}_\lambda)\subset\mathrm{F}_\mu$.
Supongamos que existe $\alpha\in\mathrm{L}$ tal que $\mathrm{S}_\alpha$ es cuasi-compacto.
Entonces las condiciones siguientes son equivalentes:
\begin{enumerate}[label=\emph{\alph*)}]
   \item El conjunto $\mathrm{F}=\bigcup_\lambda u_\lambda^{-1}(\mathrm{F}_\lambda)$ es igual a $\mathrm{S}$.
   \item Existe $\lambda$ tal que $u_\lambda^{-1}(\mathrm{F}_\lambda)=\mathrm{S}$ (y entonces $u_\mu^{-1}(\mathrm{F}_\mu)=\mathrm{S}$ para $\mu\geqslant\lambda$).
\oldpage[IV-3]{14}
   \item Existe $\lambda$ tal que $\mathrm{F}_\lambda=\mathrm{S}_\lambda$ (y entonces $\mathrm{F}_\mu=\mathrm{S}_\mu$ para $\mu\geqslant\lambda$).
\end{enumerate}
Esto se deduce inmediatamente de \sref{IV.8.3.3} pasando a los complementarios.
\end{corollary}

\begin{corollary}[8.3.5]
\label{IV.8.3.5}
Para todo $\lambda$, sean $\mathrm{E}_\lambda$, $\mathrm{F}_\lambda$ dos subconjuntos constructibles de $\mathrm{S}_\lambda$ tales que, para $\mu\geqslant\lambda$, se tenga $\mathrm{E}_\mu\subset u_{\lambda\mu}^{-1}(\mathrm{E}_\lambda)$ y $\mathrm{F}_\mu\supset u_{\lambda\mu}^{-1}(\mathrm{F}_\lambda)$.
Supongamos que existe $\alpha$ tal que $\mathrm{S}_\alpha$ es cuasi-compacto.
Entonces, para que $\mathrm{E}\subset\mathrm{F}$ \textnormal{(}resp. $\mathrm{E}=\mathrm{F}$\textnormal{)}, es necesario y suficiente que exista $\lambda$ tal que $\mathrm{E}_\lambda\subset\mathrm{F}_\lambda$ \textnormal{(}resp. $\mathrm{E}_\lambda=\mathrm{F}_\lambda$\textnormal{)}, en cuyo caso también se tiene $\mathrm{E}_\mu\subset\mathrm{F}_\mu$ \textnormal{(}resp. $\mathrm{E}_\mu=\mathrm{F}_\mu$\textnormal{)} para $\mu\geqslant\lambda$.

Este es un caso particular de \sref{IV.8.3.2}.
\end{corollary}

\begin{corollary}[8.3.6]
\label{IV.8.3.6}
Supongamos que existe $\alpha$ tal que $\mathrm{S}_\alpha$ es cuasi-compacto.
Para que $\mathrm{S}=\varnothing$, es necesario y suficiente que exista $\lambda$ tal que $\mathrm{S}_\lambda=\varnothing$.
\end{corollary}

\begin{corollary}[8.3.7]
\label{IV.8.3.7}
Para todo $\lambda$, se tiene
\begin{equation}
  u_\lambda(\mathrm{S})=\bigcap_{\mu\geqslant\lambda}u_{\lambda\mu}(\mathrm{S}_\mu).
  \tag{8.3.7.1}
\label{IV.8.3.7.1}
\end{equation}
\end{corollary}

\begin{proof}
Es claro que el primer miembro de \sref{IV.8.3.7.1} está contenido en el segundo.
Sea $s$ un punto de $\mathrm{S}_\lambda$ y pongamos $\mathrm{X}_\lambda=\operatorname{Spec}(\boldsymbol{k}(s))$; consideremos el sistema proyectivo $(\mathrm{X}_\mu,v_{\mu\nu})$ donde $\mathrm{X}_\mu=\mathrm{X}_\lambda\times_{\mathrm{S}_\lambda}\mathrm{S}_\mu$ y $v_{\mu\nu}=1\times u_{\mu\nu}$ para $\lambda\leqslant\mu\leqslant\nu$; su límite proyectivo es $\mathrm{X}=\mathrm{X}_\lambda\times_{\mathrm{S}_\lambda}\mathrm{S}$, y $v_\lambda=1\times u_\lambda$ es la aplicación canónica $\mathrm{X}\to\mathrm{X}_\lambda$ \sref{IV.8.2.5}.
Si $s\in\bigcap_{\mu\geqslant\lambda}u_{\lambda\mu}(\mathrm{S}_\mu)$, esto implica que $\mathrm{X}_\mu\neq\varnothing$ para todo $\mu\geqslant\lambda$ \sref[I]{I.3.4.8}; se sigue entonces de \sref{IV.8.3.6} que $\mathrm{X}\neq\varnothing$, y por tanto $s\in u_\lambda(\mathrm{S})$ por \sref[I]{I.3.4.8}.
\end{proof}

\begin{proposition}[8.3.8]
\label{IV.8.3.8}
\begin{enumerate}[label=(\roman*)]
   \item Para que el morfismo $u_\lambda:\mathrm{S}\to\mathrm{S}_\lambda$ sea dominante (resp. sobreyectivo), es necesario y suficiente que, para $\mu\geqslant\lambda$, el morfismo $u_{\lambda\mu}:\mathrm{S}_\mu\to\mathrm{S}_\lambda$ sea dominante (resp. sobreyectivo).
   \item Si, para algún índice $\lambda$, los morfismos $u_{\lambda\mu}:\mathrm{S}_\mu\to\mathrm{S}_\lambda$ son planos (resp. fielmente planos) para todo $\mu\geqslant\lambda$, entonces el morfismo $u_\lambda:\mathrm{S}\to\mathrm{S}_\lambda$ es plano (resp. fielmente plano).
   \item Supongamos que los morfismos $u_\mu:\mathrm{S}\to\mathrm{S}_\mu$ son sobreyectivos para $\mu\geqslant\lambda$.
Para que $u_\lambda$ sea un morfismo abierto (resp. universalmente abierto), es necesario y suficiente que, para todo $\mu\geqslant\lambda$, $u_{\lambda\mu}$ sea un morfismo abierto (resp. universalmente abierto).
\end{enumerate}
\end{proposition}

\begin{proof}
\begin{enumerate}[label=(\roman*)]
   \item Como $u_\lambda(\mathrm{S})\subset u_{\lambda\mu}(\mathrm{S}_\mu)$ para todo $\mu\geqslant\lambda$, la necesidad de las condiciones es trivial, y se deduce inmediatamente de \sref{IV.8.3.7.1} que, si los $u_{\lambda\mu}$ son sobreyectivos, también lo es $u_\lambda$.
Supongamos ahora que los $u_{\lambda\mu}$ son dominantes para $\mu\geqslant\lambda$, y consideremos en $\mathrm{S}_\lambda$ un subconjunto abierto, cuasi-compacto y no vacío $\mathrm{U}_\lambda$; entonces los $\mathrm{U}_\mu=u_{\lambda\mu}^{-1}(\mathrm{U}_\lambda)$ para $\mu\geqslant\lambda$ forman un sistema proyectivo cuyo límite proyectivo es $\mathrm{U}=u_\lambda^{-1}(\mathrm{U}_\lambda)$ \sref{IV.8.2.5}.
Por hipótesis los $\mathrm{U}_\mu$ son todos no vacíos, y lo mismo vale para $\mathrm{U}$ por \sref{IV.8.3.6}, lo que demuestra que $u_\lambda$ es dominante.

   \item En virtud de (i), basta considerar el caso en que los $u_{\lambda\mu}$ son planos; entonces se puede restringir al caso en que $\mathrm{S}_\lambda$ es afín, y por tanto también lo son los $\mathrm{S}_\mu$ para $\mu\geqslant\lambda$ y $\mathrm{S}$, y la afirmación se deduce de \sref{IV.2.1.2} y \sref[0]{0.6.2.3}.

   \item En virtud de \sref{IV.8.2.5} y de \sref[I]{I.3.5.2}, basta tratar el caso de morfismos abiertos.
Como $u_\lambda=u_{\lambda\mu}\circ u_\mu$ y $u_\mu$ es sobreyectivo, sabemos que, si $u_\lambda$ es abierto, también lo es $u_{\lambda\mu}$ para $\mu\geqslant\lambda$ (Bourbaki, \emph{Top.\ g\'en.}, chap.~I\textsuperscript{er}, 3\textsuperscript{e}~ed., \textsection\,5, n\textsuperscript{o}\,1, prop.~1).
\oldpage[IV-3]{15}
Recíprocamente, para demostrar que $u_\lambda$ es abierto cuando todos los $u_{\lambda\mu}$ son abiertos para $\mu\geqslant\lambda$, basta comprobar que, para todo abierto cuasi-compacto $\mathrm{V}$ de $\mathrm{S}$, $u_\lambda(\mathrm{V})$ es abierto en $\mathrm{S}_\lambda$; entonces existe $\mu\geqslant\lambda$ tal que $\mathrm{V}=u_\mu^{-1}(\mathrm{V}_\mu)$, donde $\mathrm{V}_\mu$ es abierto en $\mathrm{S}_\mu$ \sref{IV.8.2.11}; como $u_\mu$ es sobreyectivo, se tiene $\mathrm{V}_\mu=u_\mu(\mathrm{V})$ y $u_\lambda(\mathrm{V})=u_{\lambda\mu}(u_\mu(\mathrm{V}))$ es, por tanto, abierto por hipótesis.
\end{enumerate}

Observamos que puede ocurrir que todos los $u_{\lambda\mu}$ sean abiertos sin que $u_\lambda$ sea abierto, cuando los $u_\mu$ no son sobreyectivos.
En efecto, se obtiene un ejemplo considerando un anillo íntegro $\mathrm{A}$ que no sea un cuerpo, y su cuerpo de fracciones $\mathrm{K}$, que es el límite inductivo de los $\mathrm{A}_f$, donde $f$ recorre $\mathrm{A}-\{0\}$; tomando $\mathrm{S}_1=\operatorname{Spec}(\mathrm{A})$, $\mathrm{S}_f=\operatorname{Spec}(\mathrm{A}_f)$ y $\mathrm{S}=\varprojlim\mathrm{S}_f=\operatorname{Spec}(\mathrm{K})$, el morfismo $\mathrm{S}\to\mathrm{S}_1$ no es abierto, aunque los morfismos $\mathrm{S}_f\to\mathrm{S}_1$ son abiertos.
\end{proof}

\begin{env}[8.3.9]
\label{IV.8.3.9}
Para todo preesquema $\mathrm{X}$, denotaremos como de costumbre por $\mathfrak{P}(\mathrm{X})$ el conjunto de todos los subconjuntos del conjunto subyacente de $\mathrm{X}$, y por $\mathfrak{C}(\mathrm{X})$ (resp. $\mathfrak{D}_c(\mathrm{X})$, $\mathfrak{F}_c(\mathrm{X})$, $\mathfrak{LF}_c(\mathrm{X})$) el conjunto de subconjuntos constructibles (resp. constructibles y abiertos, constructibles y cerrados, y constructibles y localmente cerrados) de $\mathrm{X}$.
Es claro que $(\mathfrak{P}(\mathrm{S}_\lambda),u_{\lambda\mu}^{-1})$ es un \emph{sistema inductivo de conjuntos} y que las aplicaciones $u_\lambda^{-1}:\mathfrak{P}(\mathrm{S}_\lambda)\to\mathfrak{P}(\mathrm{S})$ forman un sistema inductivo de aplicaciones; por tanto, pasando al límite inductivo, se obtiene una aplicación canónica
\begin{equation}
  u_\mathfrak{P}:\varinjlim\mathfrak{P}(\mathrm{S}_\lambda)\to\mathfrak{P}(\mathrm{S}).
  \tag{8.3.9.1}
\label{IV.8.3.9.1}
\end{equation}

Además, se deduce de \sref{IV.1.8.2} que $u_{\lambda\mu}^{-1}$ envía $\mathfrak{C}(\mathrm{S}_\lambda)$ (resp. $\mathfrak{D}_c(\mathrm{S}_\lambda)$, $\mathfrak{F}_c(\mathrm{S}_\lambda)$, $\mathfrak{LF}_c(\mathrm{S}_\lambda)$) en $\mathfrak{C}(\mathrm{S}_\mu)$ (resp. $\mathfrak{D}_c(\mathrm{S}_\mu)$, $\mathfrak{F}_c(\mathrm{S}_\mu)$, $\mathfrak{LF}_c(\mathrm{S}_\mu)$), y que $u_\lambda^{-1}$ envía $\mathfrak{C}(\mathrm{S}_\lambda)$ (resp. $\mathfrak{D}_c(\mathrm{S}_\lambda)$, $\mathfrak{F}_c(\mathrm{S}_\lambda)$, $\mathfrak{LF}_c(\mathrm{S}_\lambda)$) en $\mathfrak{C}(\mathrm{S})$ (resp. $\mathfrak{D}_c(\mathrm{S})$, $\mathfrak{F}_c(\mathrm{S})$, $\mathfrak{LF}_c(\mathrm{S})$).
Tenemos, pues, por restricción de \sref{IV.8.3.9.1}, las aplicaciones canónicas
\begin{align}
  &\varinjlim\mathfrak{C}(\mathrm{S}_\lambda)\to\mathfrak{C}(\mathrm{S}) \tag{8.3.9.2}\label{IV.8.3.9.2}\\
  &\varinjlim\mathfrak{D}_c(\mathrm{S}_\lambda)\to\mathfrak{D}_c(\mathrm{S}) \tag{8.3.9.3}\label{IV.8.3.9.3}\\
  &\varinjlim\mathfrak{F}_c(\mathrm{S}_\lambda)\to\mathfrak{F}_c(\mathrm{S}) \tag{8.3.9.4}\label{IV.8.3.9.4}\\
  &\varinjlim\mathfrak{LF}_c(\mathrm{S}_\lambda)\to\mathfrak{LF}_c(\mathrm{S}). \tag{8.3.9.5}
\label{IV.8.3.9.5}
\end{align}
\end{env}

\begin{env}[8.3.10]
\label{IV.8.3.10}
Sea $g_\alpha:\mathrm{X}_\alpha\to\mathrm{S}_\alpha$ un morfismo; con las notaciones de \sref{IV.8.2.5} se tiene, por una parte, una aplicación canónica $v_\mathfrak{P}:\varinjlim\mathfrak{P}(\mathrm{X}_\lambda)\to\mathfrak{P}(\mathrm{X})$; por otra parte, se tienen morfismos de proyección $g_\lambda:\mathrm{X}_\lambda\to\mathrm{S}_\lambda$ para todo $\lambda\geqslant\alpha$ y un morfismo de proyección $g:\mathrm{X}\to\mathrm{S}$.
Es claro que los $g_\lambda^{-1}:\mathfrak{P}(\mathrm{S}_\lambda)\to\mathfrak{P}(\mathrm{X}_\lambda)$ forman un sistema inductivo de aplicaciones, y que los diagramas
\[
\xymatrix{
  \mathfrak{P}(\mathrm{S}_\lambda) \ar[r]^{u_\lambda^{-1}} \ar[d]_{g_\lambda^{-1}} & \mathfrak{P}(\mathrm{S}) \ar[d]^{g^{-1}} \\
  \mathfrak{P}(\mathrm{X}_\lambda) \ar[r]_{v_\lambda^{-1}} & \mathfrak{P}(\mathrm{X})
}
\]
\oldpage[IV-3]{16}
son conmutativos; pasando al límite inductivo se deduce de ello un diagrama conmutativo
\begin{equation}
\xymatrix{
  \varinjlim\mathfrak{P}(\mathrm{S}_\lambda) \ar[r]^{u_\mathfrak{P}} \ar[d]_{\varinjlim g_\lambda^{-1}} & \mathfrak{P}(\mathrm{S}) \ar[d]^{g^{-1}} \\
  \varinjlim\mathfrak{P}(\mathrm{X}_\lambda) \ar[r]_{v_\mathfrak{P}} & \mathfrak{P}(\mathrm{X})
}
\tag{8.3.10.1}
\label{IV.8.3.10.1}
\end{equation}
y se deduce de \sref{IV.1.8.2} que se tienen diagramas análogos sustituyendo $\mathfrak{P}$ por $\mathfrak{C}$, $\mathfrak{D}_c$, $\mathfrak{F}_c$ o $\mathfrak{LF}_c$.
\end{env}

Se sigue de \sref{IV.8.3.5} que, bajo la hipótesis de que para algún $\alpha\in\mathrm{L}$, $\mathrm{S}_\alpha$ es \emph{cuasi-compacto}, la aplicación canónica \sref{IV.8.3.9.2} es \emph{inyectiva}.
Furthermore:

\begin{theorem}[8.3.11]
\label{IV.8.3.11}
Supongamos que existe $\alpha\in\mathrm{L}$ tal que $\mathrm{S}_\alpha$ es cuasi-compacto y cuasi-separado.
Entonces las aplicaciones canónicas \sref{IV.8.3.9.2}, \sref{IV.8.3.9.3}, \sref{IV.8.3.9.4} y \sref{IV.8.3.9.5} son \emph{biyectivas}.
\end{theorem}

\begin{proof}
En virtud de la observación precedente, queda por demostrar que estas aplicaciones son \emph{sobreyectivas}; como todo subconjunto constructible de $\mathrm{S}$ es unión finita de conjuntos de la forma $\mathrm{U}\cap\complement\mathrm{V}$, donde $\mathrm{U}$ y $\mathrm{V}$ son abiertos y constructibles, bastará demostrar que \sref{IV.8.3.9.3} es sobreyectiva para que lo sea también \sref{IV.8.3.9.2} (y también \sref{IV.8.3.9.4}, pasando a los complementarios).
Ahora bien, como los morfismos $\mathrm{S}_\lambda\to\mathrm{S}_\alpha$ y $\mathrm{S}\to\mathrm{S}_\alpha$ son afines y, por tanto, separados, los $\mathrm{S}_\lambda$ para $\lambda\geqslant\alpha$ y $\mathrm{S}$ son cuasi-compactos y cuasi-separados \sref{IV.1.2.2}, y se sabe que los subconjuntos abiertos constructibles de tal preesquema no son otros que los abiertos cuasi-compactos \sref{IV.1.8.1}.
La conclusión se sigue entonces de \sref{IV.8.2.11}, salvo para la aplicación \sref{IV.8.3.9.5}.
Para demostrar que esta última aplicación es sobreyectiva, consideremos un subconjunto localmente cerrado y constructible $\mathrm{Z}$ de $\mathrm{S}$; por tanto, $\mathrm{Z}$ es cuasi-compacto ($\mathbf{0}_{\mathrm{III}}$,~9.1.4).
Como, por hipótesis, todo punto $x\in\mathrm{Z}$ admite una vecindad abierta cuasi-compacta $\mathrm{V}_x$ en $\mathrm{S}$ tal que $\mathrm{Z}\cap\mathrm{V}_x$ es cerrado en $\mathrm{V}_x$, podemos cubrir $\mathrm{Z}$ con un número finito de los $\mathrm{V}_x$; en otras palabras, existe un abierto cuasi-compacto $\mathrm{U}$ que contiene a $\mathrm{Z}$ tal que $\mathrm{Z}$ es cerrado en $\mathrm{U}$; como $\mathrm{Z}$ es constructible en $\mathrm{S}$, también es constructible en $\mathrm{U}$ ($\mathbf{0}_{\mathrm{III}}$,~9.1.8).
Sabemos por \sref{IV.8.2.11} que existen $\lambda$ y un abierto cuasi-compacto $\mathrm{U}_\lambda$ de $\mathrm{S}_\lambda$ tales que $\mathrm{U}=u_\lambda^{-1}(\mathrm{U}_\lambda)$.
Aplicando a $\mathrm{U}$ (que es el límite proyectivo de los $\mathrm{U}_\mu=u_{\lambda\mu}^{-1}(\mathrm{U}_\lambda)$ para $\mu\geqslant\lambda$) el hecho de que la aplicación \sref{IV.8.3.9.4} es sobreyectiva, se ve que existen $\mu\geqslant\lambda$ y un subconjunto cerrado constructible $\mathrm{Z}_\mu$ en $\mathrm{U}_\mu$ tales que $\mathrm{Z}=u_\mu^{-1}(\mathrm{Z}_\mu)$.
Pero, como la inmersión canónica $\mathrm{U}_\mu\to\mathrm{S}_\mu$ es cuasi-compacta por hipótesis \sref{IV.1.2.7}, es de presentación finita \sref{IV.1.6.2},~(i), y $\mathrm{Z}_\mu$ es también un subconjunto constructible \emph{de} $\mathrm{S}_\mu$ en virtud de \sref{IV.1.8.4} y \sref{IV.1.8.1}; como $\mathrm{Z}_\mu$ es evidentemente localmente cerrado en $\mathrm{S}_\mu$, queda demostrada la afirmación.
\end{proof}

\begin{corollary}[8.3.12]
\label{IV.8.3.12}
Supongamos que existe $\alpha$ tal que $\mathrm{S}_\alpha$ es cuasi-compacto, y sea, para todo $\lambda$, $\mathrm{Z}_\lambda$ un subconjunto constructible de $\mathrm{S}_\lambda$ tal que $\mathrm{Z}_\mu=u_{\lambda\mu}^{-1}(\mathrm{Z}_\lambda)$ para $\mu\geqslant\lambda$.
Si $\mathrm{Z}=u_\lambda^{-1}(\mathrm{Z}_\lambda)$,
\oldpage[IV-3]{17}
entonces, para que $\mathrm{Z}$ sea abierto (resp. cerrado, resp. localmente cerrado) en $\mathrm{S}$, es necesario y suficiente que exista $\lambda\geqslant\alpha$ tal que $\mathrm{Z}_\lambda$ lo sea en $\mathrm{S}_\lambda$.
\end{corollary}

\begin{proof}
Cubramos $\mathrm{S}_\alpha$ con un número finito de abiertos afines $\mathrm{U}_\alpha^{(i)}$; entonces los $\mathrm{U}^{(i)}=u_\alpha^{-1}(\mathrm{U}_\alpha^{(i)})$ forman un recubrimiento afín abierto de $\mathrm{S}$, y para que $\mathrm{Z}$ sea abierto (resp. cerrado, resp. localmente cerrado) en $\mathrm{S}$ es necesario y suficiente que cada $\mathrm{Z}\cap\mathrm{U}^{(i)}$ lo sea en $\mathrm{U}^{(i)}$.
Como $\mathrm{L}$ es filtrado y cada $\mathrm{Z}\cap\mathrm{U}^{(i)}$ es constructible en $\mathrm{U}^{(i)}$ ($\mathbf{0}_{\mathrm{III}}$,~9.1.8), podemos restringirnos a demostrar el corolario cuando $\mathrm{S}_\alpha$ es afín, y por tanto cuasi-compacto y cuasi-separado; pero entonces se sigue de \sref{IV.8.3.11}.
\end{proof}

\begin{proposition}[8.3.13]
\label{IV.8.3.13}
Supongamos que los morfismos $u_{\lambda\mu}:\mathrm{S}_\mu\to\mathrm{S}_\lambda$ son planos para $\lambda\leqslant\mu$, y que existe $\alpha$ tal que $\mathrm{S}_\alpha$ es cuasi-compacto.
Para todo $\lambda$, sean $\mathrm{Z}'_\lambda$, $\mathrm{Z}''_\lambda$ dos subconjuntos pro-constructibles de $\mathrm{S}_\lambda$ tales que $\mathrm{Z}'_\mu=u_{\lambda\mu}^{-1}(\mathrm{Z}'_\lambda)$ y $\mathrm{Z}''_\mu=u_{\lambda\mu}^{-1}(\mathrm{Z}''_\lambda)$ para $\mu\geqslant\lambda$; supongamos además que $\overline{\mathrm{Z}'_\alpha}$ es constructible en $\mathrm{S}_\alpha$.
Pongamos $\mathrm{Z}'=u_\lambda^{-1}(\mathrm{Z}'_\lambda)$, $\mathrm{Z}''=u_\lambda^{-1}(\mathrm{Z}''_\lambda)$; para que $\mathrm{Z}''\subset\overline{\mathrm{Z}'}$, es necesario y suficiente que exista $\lambda\geqslant\alpha$ tal que $\mathrm{Z}''_\lambda\subset\overline{\mathrm{Z}'_\lambda}$.
\end{proposition}

\begin{proof}
En efecto, se sabe que $u_\lambda:\mathrm{S}\to\mathrm{S}_\lambda$ es también un morfismo plano para todo $\lambda$ \sref{IV.8.3.8}; como $\mathrm{Z}'_\lambda$ es pro-constructible, la clausura de $\mathrm{Z}'_\mu$ para $\mu\geqslant\lambda$ (resp. de $\mathrm{Z}'$) en $\mathrm{S}_\mu$ (resp. $\mathrm{S}$) es igual a $u_{\lambda\mu}^{-1}(\overline{\mathrm{Z}'_\lambda})$ (resp. $u_\lambda^{-1}(\overline{\mathrm{Z}'_\lambda})$) \sref{IV.2.3.10}.
Como $u_{\lambda\mu}^{-1}(\overline{\mathrm{Z}'_\lambda})$ y $u_\lambda^{-1}(\overline{\mathrm{Z}'_\lambda})$ son constructibles \sref{IV.1.8.2}, la conclusión se sigue de \sref{IV.8.3.2}.
\end{proof}

\subsection{Criterios de irreducibilidad y conexión para los límites proyectivos de preesquemas}
\label{subsection:IV.8.4}

\begin{proposition}[8.4.1]
\label{IV.8.4.1}
Supongamos que existe un índice $\alpha$ tal que $\mathrm{S}_\alpha$ es cuasi-compacto.
\begin{enumerate}[label=(\roman*)]
  \item Si $\mathrm{S}$ no es irreducible y si, además, el espacio subyacente de $\mathrm{S}$ es noetheriano y $\mathrm{S}_\alpha$ es cuasi-separado, existe $\lambda\geqslant\alpha$ tal que, para $\mu\geqslant\lambda$, $\mathrm{S}_\mu$ no es irreducible.
  \item Si $\mathrm{S}$ no es conexo, existe $\lambda\geqslant\alpha$ tal que, para $\mu\geqslant\lambda$, $\mathrm{S}_\mu$ no es conexo.
\end{enumerate}
\end{proposition}

\begin{proof}
Supongamos que $\mathrm{S}$ es la unión de dos subconjuntos cerrados distintos $\mathrm{S}'$, $\mathrm{S}''$ de $\mathrm{S}$ (resp.\ de dos subconjuntos cerrados no vacíos y disjuntos).
En el caso (i), $\mathrm{S}'$ y $\mathrm{S}''$ son constructibles, puesto que el espacio $\mathrm{S}$ es noetheriano.
En virtud de \sref{IV.8.3.11}, existen $\lambda\geqslant\alpha$ y dos subconjuntos cerrados constructibles $\mathrm{S}'_\lambda$, $\mathrm{S}''_\lambda$ de $\mathrm{S}_\lambda$ tales que $\mathrm{S}'=u_\lambda^{-1}(\mathrm{S}'_\lambda)$, $\mathrm{S}''=u_\lambda^{-1}(\mathrm{S}''_\lambda)$; como $\mathrm{S}=\mathrm{S}'\cup\mathrm{S}''$, se deduce de \sref{IV.8.3.11} que se puede suponer $\mathrm{S}_\lambda=\mathrm{S}'_\lambda\cup\mathrm{S}''_\lambda$; como $\mathrm{S}'_\lambda$ y $\mathrm{S}''_\lambda$ son distintos de $\mathrm{S}_\lambda$, esto prueba que $\mathrm{S}_\lambda$ no es irreducible.

En el caso (ii), $\mathrm{S}'$ y $\mathrm{S}''$ son abiertos y cuasi-compactos; por tanto, en virtud de \sref{IV.8.2.11}, existen $\lambda\geqslant\alpha$ y dos subconjuntos abiertos cuasi-compactos $\mathrm{S}'_\lambda$, $\mathrm{S}''_\lambda$ de $\mathrm{S}_\lambda$ tales que $\mathrm{S}'=u_\lambda^{-1}(\mathrm{S}'_\lambda)$, $\mathrm{S}''=u_\lambda^{-1}(\mathrm{S}''_\lambda)$.
Además, como $\mathrm{S}'$ y $\mathrm{S}''$ son abiertos y cerrados en $\mathrm{S}$, ambos son proconstructibles e indconstructibles \sref{IV.1.9.6}, y por tanto constructibles \sref{IV.1.9.11}; se deduce entonces de \sref{IV.8.3.5} que se puede suponer $\lambda$ elegido de modo que $\mathrm{S}_\lambda=\mathrm{S}'_\lambda\cup\mathrm{S}''_\lambda$ y $\mathrm{S}'_\lambda\cap\mathrm{S}''_\lambda=\varnothing$, lo que muestra que $\mathrm{S}_\lambda$ no es conexo.
\end{proof}

\begin{proposition}[8.4.2]
\label{IV.8.4.2}
Supongamos que el espacio subyacente de $\mathrm{S}$ es noetheriano y que se satisface una de las dos condiciones siguientes:
\begin{enumerate}[label=\emph{\alph*)}]
  \item Para $\lambda\leqslant\mu$, $u_{\lambda\mu}:\mathrm{S}_\mu\to\mathrm{S}_\lambda$ es dominante, y existe $\alpha$ tal que $\mathrm{S}_\alpha$ es cuasi-compacto.
\oldpage[IV-3]{18}
  \item Existe $\alpha$ tal que el espacio subyacente de $\mathrm{S}_\alpha$ es noetheriano y, para $\mu\geqslant\lambda$, $u_{\lambda\mu}$ es un homeomorfismo de $\mathrm{S}_\mu$ sobre un subespacio de $\mathrm{S}_\lambda$.
\end{enumerate}
En estas condiciones, existe $\lambda$ tal que, para todo $\mu\geqslant\lambda$:
\begin{enumerate}[label=(\roman*)]
  \item Para toda componente irreducible $\mathrm{Y}_i$ de $\mathrm{S}$ ($1\leqslant i\leqslant m$), $\overline{u_\mu(\mathrm{Y}_i)}$ es una componente irreducible de $\mathrm{S}_\mu$, y la aplicación $\mathrm{Y}_i\mapsto\overline{u_\mu(\mathrm{Y}_i)}$ es una biyección del conjunto de las componentes irreducibles de $\mathrm{S}$ sobre el conjunto de las componentes irreducibles de $\mathrm{S}_\mu$.
  \item Para toda componente conexa $\mathrm{C}_j$ de $\mathrm{S}$ ($1\leqslant j\leqslant n$), $\overline{u_\mu(\mathrm{C}_j)}$ es una componente conexa de $\mathrm{S}_\mu$, y la aplicación $\mathrm{C}_j\mapsto\overline{u_\mu(\mathrm{C}_j)}$ es una biyección del conjunto de las componentes conexas de $\mathrm{S}$ sobre el conjunto de las componentes conexas de $\mathrm{S}_\mu$.
\end{enumerate}
\end{proposition}

\begin{proof}
Demostremos primero el

\begin{lemma}[8.4.2.1]
\label{IV.8.4.2.1}
Bajo la condición \emph{a)} o \emph{b)} de \sref{IV.8.4.2}, existe $\lambda$ tal que, para $\mu\geqslant\lambda$, $u_\mu:\mathrm{S}\to\mathrm{S}_\mu$ es dominante.
\end{lemma}

En el caso \emph{a)}, esto ya se ha demostrado sin suponer noetheriano el espacio $\mathrm{S}$ \sref{IV.8.3.8},~(i).
En el caso \emph{b)}, pongamos $\mathrm{Z}_\alpha=\overline{u_\alpha(\mathrm{S})}$; como subconjunto cerrado del espacio noetheriano $\mathrm{S}_\alpha$, $\mathrm{Z}_\alpha$ es constructible y, como $u_\alpha^{-1}(\mathrm{Z}_\alpha)=\mathrm{S}$, se deduce de \sref{IV.8.3.5} que existe $\lambda\geqslant\alpha$ tal que, para $\mu\geqslant\lambda$, $\mathrm{Z}_\mu=u_{\alpha\mu}^{-1}(\mathrm{Z}_\alpha)=\mathrm{S}_\mu$.
Pero, como $u_{\alpha\mu}$ es un homeomorfismo de $\mathrm{S}_\mu$ sobre un subespacio de $\mathrm{S}_\alpha$ y la aplicación compuesta $\mathrm{S}\to\mathrm{Z}_\mu\to\mathrm{Z}_\alpha$ es dominante, lo mismo vale para $\mathrm{S}\to\mathrm{Z}_\mu=\mathrm{S}_\mu$.

Demostrado este lema, se puede suponer que, para todo $\lambda$, $u_\lambda$ es un morfismo dominante.
\begin{enumerate}[label=(\roman*)]
  \item Cada $\mathrm{S}_\lambda$ es la unión de los $\overline{u_\lambda(\mathrm{Y}_i)}$, que son irreducibles.
Por otro lado, si $\mathrm{U}_i$ es el abierto de $\mathrm{S}$ complementario de la unión de los $\mathrm{Y}_j$ de índice $j\neq i$ ($1\leqslant i\leqslant m$), los $\mathrm{U}_i$ son dos a dos disjuntos y $\mathrm{Y}_i=\overline{\mathrm{U}_i}$ \sref[0]{0.2.1.6}.
Como el espacio subyacente de $\mathrm{S}$ es noetheriano, los $\mathrm{U}_i$ son cuasi-compactos; por tanto, existen un índice $\lambda$ y abiertos $\mathrm{U}_{i\lambda}$ de $\mathrm{S}_\lambda$ tales que $\mathrm{U}_i=u_\lambda^{-1}(\mathrm{U}_{i\lambda})$ para $1\leqslant i\leqslant m$ \sref{IV.8.2.11}.
De esto se deduce que, si se pone $\mathrm{U}_{i\mu}=u_{\lambda\mu}^{-1}(\mathrm{U}_{i\lambda})$ para $\lambda\leqslant\mu$, los $\mathrm{U}_{i\mu}$ son dos a dos disjuntos, pues los $\mathrm{U}_i=u_\mu^{-1}(\mathrm{U}_{i\mu})$ lo son y $u_\mu$ es dominante.
En consecuencia, ninguna de las clausuras $\overline{\mathrm{U}_{i\mu}}$ está contenida en otra, y $u_\mu(\mathrm{U}_i)$ es denso en $\mathrm{U}_{i\mu}$ puesto que $u_\mu$ es dominante; de ahí $\overline{\mathrm{U}_{i\mu}}=\overline{u_\mu(\mathrm{Y}_i)}$, lo que prueba que los $\overline{\mathrm{U}_{i\mu}}$ son las componentes irreducibles de $\mathrm{S}_\mu$ \sref[0]{0.2.1.7} y concluye la demostración.

  \item Como el espacio subyacente $\mathrm{S}$ es noetheriano, los $\mathrm{C}_j$ son abiertos y cerrados en $\mathrm{S}$ y cuasi-compactos; el mismo razonamiento que en (i) muestra que existen $\lambda$ y abiertos $\mathrm{V}_{j\lambda}$ de $\mathrm{S}_\lambda$ tales que $\mathrm{C}_j=u_\lambda^{-1}(\mathrm{V}_{j\lambda})$ para $1\leqslant j\leqslant n$.
Se ve asimismo, como en (i), que si se pone $\mathrm{V}_{j\mu}=u_{\lambda\mu}^{-1}(\mathrm{V}_{j\lambda})$ para $\mu\geqslant\lambda$, los $\mathrm{V}_{j\mu}$ son dos a dos disjuntos y $u_\mu(\mathrm{C}_j)$ es denso en $\mathrm{V}_{j\mu}$; esto implica que $\mathrm{V}_{j\mu}$ es conexo.
Además, se deduce de \sref{IV.8.3.4} que, para $\mu$ suficientemente grande, la unión de los $\mathrm{V}_{j\mu}$ es $\mathrm{S}_\mu$, ya que todo abierto de un preesquema es indconstructible \sref{IV.1.9.6}.
Los $\mathrm{V}_{j\mu}$ son, pues, las componentes conexas de $\mathrm{S}_\mu$, y esto concluye la demostración.
\end{enumerate}
\end{proof}

Observemos que, si los morfismos $u_{\lambda\mu}$ son \emph{inmersiones}, satisfacen en particular la condición \emph{b)} de \sref{IV.8.4.2}.

\oldpage[IV-3]{19}
\begin{corollary}[8.4.3]
\label{IV.8.4.3}
Supongamos que se satisface una de las condiciones \emph{a)}, \emph{b)} de \sref{IV.8.4.2}, siendo noetheriano el espacio subyacente de $\mathrm{S}$; entonces, para que $\mathrm{S}$ sea irreducible, es necesario y suficiente que exista $\lambda$ tal que $\mathrm{S}_\mu$ lo sea para todo $\mu\geqslant\lambda$.
\end{corollary}

\begin{proposition}[8.4.4]
\label{IV.8.4.4}
Supongamos que existe $\alpha$ tal que $\mathrm{S}_\alpha$ es cuasi-compacto y que, para $\lambda\leqslant\mu$, $u_{\lambda\mu}:\mathrm{S}_\mu\to\mathrm{S}_\lambda$ es dominante.
Entonces, para que $\mathrm{S}$ sea conexo, es necesario y suficiente que exista $\lambda$ tal que $\mathrm{S}_\mu$ lo sea para todo $\mu\geqslant\lambda$.
\end{proposition}

\begin{proof}
La condición es suficiente en virtud de \sref{IV.8.4.1}; por otro lado, hemos visto en \sref{IV.8.3.8},~(i), que $u_\lambda:\mathrm{S}\to\mathrm{S}_\lambda$ es dominante para $\lambda$ suficientemente grande; por tanto, si $\mathrm{S}$ es conexo, lo mismo ocurre con $\mathrm{S}_\lambda$, ya que $u_\lambda(\mathrm{S})$ es denso en $\mathrm{S}_\lambda$ y es conexo.
\end{proof}

\begin{corollary}[8.4.5]
\label{IV.8.4.5}
Sea $k$ un cuerpo y $\mathrm{X}$ un preesquema $k$ cuasi-compacto.
Para que $\mathrm{X}$ sea geométricamente conexo \sref{IV.4.5.2}, es necesario y suficiente que, para toda extensión finita y separable $\mathrm{K}$ de $k$, $\mathrm{X}\otimes_k\mathrm{K}$ sea conexo.
\end{corollary}

\begin{proof}
La condición es trivialmente necesaria.
Para ver que es suficiente, debemos demostrar que si $\Omega$ es una clausura algebraica de $k$, $\mathrm{X}\otimes_k\Omega$ es conexo \sref{IV.4.5.1}.
Ahora bien, $\Omega$ es un límite inductivo filtrante de las subextensiones finitas $k'$ de $k$, y para $k\subset k'\subset k''\subset\Omega$, el morfismo $\mathrm{X}\otimes_k k''\to\mathrm{X}\otimes_k k'$ es sobreyectivo.
Así pues, en virtud de \sref{IV.8.4.4}, basta demostrar que $\mathrm{X}\otimes_k k'$ es conexo para toda extensión finita $k'$ de $k$.
Pero si $\mathrm{K}$ es la mayor extensión separable contenida en $k'$, el morfismo $\mathrm{X}\otimes_k k'\to\mathrm{X}\otimes_k\mathrm{K}$ es finito, sobreyectivo y radicial, y por tanto \sref{IV.2.4.5} un homeomorfismo; como $\mathrm{X}\otimes_k\mathrm{K}$ es conexo por hipótesis, también lo es $\mathrm{X}\otimes_k k'$.
\end{proof}

\begin{remark}[8.4.6]
\label{IV.8.4.6}
\begin{enumerate}[label=(\roman*)]
  \item La demostración de \sref{IV.8.4.2} muestra que la conclusión de esta proposición sigue siendo válida si se supone que el espacio subyacente de $\mathrm{S}$ es noetheriano, que existe $\alpha$ tal que $\mathrm{S}_\alpha$ es cuasi-compacto y, finalmente, que los $u_\lambda:\mathrm{S}\to\mathrm{S}_\lambda$ son dominantes.
  \item Por otro lado, la conclusión de \sref{IV.8.4.2} puede fallar cuando los $u_\lambda$ no son dominantes para $\lambda$ suficientemente grande, incluso cuando los $\mathrm{S}_\lambda$ y $\mathrm{S}$ son noetherianos, como muestra el ejemplo siguiente.
Tómese $\mathbf{N}$ como conjunto de índices, todos los $\mathrm{S}_n$ iguales a $\operatorname{Spec}(\mathrm{A}\times\mathrm{K})=\operatorname{Spec}(\mathrm{A})\amalg\operatorname{Spec}(\mathrm{K})$, donde $\mathrm{K}$ es un cuerpo, $\mathrm{A}$ cualquier $\mathrm{K}$-álgebra, y todos los morfismos $u_{n,n+1}$ iguales al mismo morfismo correspondiente al homomorfismo $(x,y)\mapsto(j(y),y)$ de $\mathrm{A}\times\mathrm{K}$ en sí misma, donde $j:\mathrm{K}\to\mathrm{A}$ es el homomorfismo canónico.
Se verifica fácilmente que el límite inductivo de este sistema de anillos es $\mathrm{K}$, siendo el homomorfismo canónico $\mathrm{A}\times\mathrm{K}\to\mathrm{K}$ la segunda proyección.
Así se ve que $\mathrm{S}=\operatorname{Spec}(\mathrm{K})$ es irreducible aunque ninguno de los $\mathrm{S}_n$ sea conexo.
\end{enumerate}
\end{remark}

\subsection{Módulos de presentación finita sobre un límite proyectivo de preesquemas}
\label{subsection:IV.8.5}

\begin{env}[8.5.1]
\label{IV.8.5.1}
Conservamos siempre las notaciones de \sref{IV.8.2.2}; además nos restringiremos al caso en que $\mathrm{S}_0$ sea uno de los $\mathrm{S}_\lambda$, al que siempre se puede reducir.

\emph{Siempre que, en esta sección, consideremos una familia $(\sh{F}_\lambda)$, donde, para cada $\lambda\in\mathrm{L}$, $\sh{F}_\lambda$ es un Módulo sobre $\sh{O}_{\mathrm{S}_\lambda}$, se entenderá que esta familia satisface la condición}
\begin{equation}
  \sh{F}_\mu=u_{\lambda\mu}^*(\sh{F}_\lambda) \qquad\text{para }\lambda\leqslant\mu.
  \tag{8.5.1.1}
\label{IV.8.5.1.1}
\end{equation}
\oldpage[IV-3]{20}
Pondremos entonces
\begin{equation}
  \sh{F}=u_\lambda^*(\sh{F}_\lambda)
  \label{IV.8.5.1.2}
\tag{8.5.1.2}
\end{equation}
que es un Módulo sobre $\sh{O}_\mathrm{S}$ que no depende del índice $\lambda\in\mathrm{L}$, en virtud de la hipótesis \sref{IV.8.5.1.1}.

Sean $(\sh{F}_\lambda)$ y $(\sh{G}_\lambda)$ dos familias de ese tipo de Módulos sobre $\sh{O}_{\mathrm{S}_\lambda}$.
Es claro que las aplicaciones $f_\lambda\mapsto u_{\lambda\mu}^*(f_\lambda)$ de $\operatorname{Hom}_{\sh{O}_{\mathrm{S}_\lambda}}(\sh{F}_\lambda,\sh{G}_\lambda)$ en $\operatorname{Hom}_{\sh{O}_{\mathrm{S}_\mu}}(\sh{F}_\mu,\sh{G}_\mu)$ definen un \emph{sistema inductivo} de grupos abelianos $(\operatorname{Hom}_{\sh{O}_{\mathrm{S}_\lambda}}(\sh{F}_\lambda,\sh{G}_\lambda),u_{\lambda\mu}^*)$, y que las aplicaciones $f_\lambda\mapsto u_\lambda^*(f_\lambda)$ forman un sistema inductivo de homomorfismos de grupos abelianos
\[
  u_\lambda^*:\operatorname{Hom}_{\sh{O}_{\mathrm{S}_\lambda}}(\sh{F}_\lambda,\sh{G}_\lambda)\to\operatorname{Hom}_\mathrm{S}(\sh{F},\sh{G});
\]
de donde, pasando al límite inductivo, se obtiene un homomorfismo canónico de grupos abelianos
\begin{equation}
  u_{\sh{F},\sh{G}}^*:\varinjlim\operatorname{Hom}_{\sh{O}_{\mathrm{S}_\lambda}}(\sh{F}_\lambda,\sh{G}_\lambda)\to\operatorname{Hom}_\mathrm{S}(\sh{F},\sh{G}).
  \tag{8.5.1.3}
\label{IV.8.5.1.3}
\end{equation}

Obsérvese que cuando $\sh{F}_\lambda=\sh{O}_{\mathrm{S}_\lambda}$ se satisface la condición \sref{IV.8.5.1.1}, y se tiene $\sh{F}=\sh{O}_\mathrm{S}$; por tanto, el homomorfismo \sref{IV.8.5.1.3} proporciona \sref[0]{0.5.1.1} un homomorfismo canónico de grupos abelianos
\begin{equation}
  u_{\sh{G}}:\varinjlim\Gamma(\mathrm{S}_\lambda,\sh{G}_\lambda)\to\Gamma(\mathrm{S},\sh{G}).
  \tag{8.5.1.4}
\label{IV.8.5.1.4}
\end{equation}
\end{env}

\begin{theorem}[8.5.2]
\label{IV.8.5.2}
\begin{enumerate}[label=(\roman*)]
  \item Supóngase que $\mathrm{S}_0$ es cuasi-compacto (resp.\ cuasi-compacto y cuasi-separado) y que, para algún $\lambda\in\mathrm{L}$, $\sh{F}_\lambda$ es cuasi-coherente y de tipo finito (resp.\ de presentación finita) y $\sh{G}_\lambda$ cuasi-coherente.
Entonces el homomorfismo $u_{\sh{F},\sh{G}}^*$ es inyectivo (resp.\ biyectivo).
  \item Supóngase que $\mathrm{S}_0$ es cuasi-compacto y cuasi-separado.
Para todo Módulo sobre $\sh{O}_\mathrm{S}$ cuasi-coherente y de presentación finita $\sh{F}$, existen $\lambda\in\mathrm{L}$ y un Módulo sobre $\sh{O}_{\mathrm{S}_\lambda}$ cuasi-coherente y de presentación finita $\sh{F}_\lambda$ tales que $\sh{F}$ es isomorfo a $u_\lambda^*(\sh{F}_\lambda)$.
\end{enumerate}
\end{theorem}

\begin{proof}
\begin{enumerate}[label=(\roman*)]
  \item Podemos evidentemente restringirnos al caso $\mathrm{S}_0=\mathrm{S}_\lambda$, ya que los morfismos $u_{0\lambda}:\mathrm{S}_\lambda\to\mathrm{S}_0$ son afines, y por tanto cuasi-compactos y separados.
Consideremos primero el caso en que $\mathrm{S}_0=\operatorname{Spec}(\mathrm{A}_0)$ es \emph{afín}.
Entonces la afirmación (i) es equivalente al

\begin{lemma}[8.5.2.1]
\label{IV.8.5.2.1}
Sea $\mathrm{A}_0$ un anillo, $(\mathrm{A}_\lambda)_{\lambda\in\mathrm{L}}$ un sistema inductivo de $\mathrm{A}_0$-álgebras, $\mathrm{A}=\varinjlim\mathrm{A}_\lambda$; sean $\mathrm{M}_0$, $\mathrm{N}_0$ dos $\mathrm{A}_0$-módulos, y pongamos $\mathrm{M}_\lambda=\mathrm{M}_0\otimes_{\mathrm{A}_0}\mathrm{A}_\lambda$, $\mathrm{N}_\lambda=\mathrm{N}_0\otimes_{\mathrm{A}_0}\mathrm{A}_\lambda$, $\mathrm{M}=\mathrm{M}_0\otimes_{\mathrm{A}_0}\mathrm{A}=\varinjlim\mathrm{M}_\lambda$, $\mathrm{N}=\mathrm{N}_0\otimes_{\mathrm{A}_0}\mathrm{A}=\varinjlim\mathrm{N}_\lambda$.
Si $\mathrm{M}_0$ es de tipo finito (resp.\ de presentación finita), el homomorfismo canónico
\begin{equation}
  \varinjlim\operatorname{Hom}_{\mathrm{A}_\lambda}(\mathrm{M}_\lambda,\mathrm{N}_\lambda)\to\operatorname{Hom}_\mathrm{A}(\mathrm{M},\mathrm{N})
  \tag{8.5.2.2}
\label{IV.8.5.2.2}
\end{equation}
es inyectivo (resp.\ biyectivo).
\end{lemma}

Sabemos, en efecto (Bourbaki, \emph{Alg.}, cap.~II, 3\textsuperscript{a}~ed., \textsection\,5, n\textsuperscript{o}\,1), que existen isomorfismos canónicos functoriales
\[
  \operatorname{Hom}_{\mathrm{A}_\lambda}(\mathrm{M}_\lambda,\mathrm{N}_\lambda)\xrightarrow{\sim}\operatorname{Hom}_{\mathrm{A}_0}(\mathrm{M}_0,\mathrm{N}_\lambda),\qquad\operatorname{Hom}_\mathrm{A}(\mathrm{M},\mathrm{N})\xrightarrow{\sim}\operatorname{Hom}_{\mathrm{A}_0}(\mathrm{M}_0,\mathrm{N})
\]
de modo que el homomorfismo \sref{IV.8.5.2.2} no es otro, salvo isomorfismos canónicos, que el homomorfismo canónico
\begin{equation}
  \varinjlim\operatorname{Hom}_{\mathrm{A}_0}(\mathrm{M}_0,\mathrm{N}_\lambda)\to\operatorname{Hom}_{\mathrm{A}_0}(\mathrm{M}_0,\varinjlim\mathrm{N}_\lambda),
  \tag{8.5.2.3}
\label{IV.8.5.2.3}
\end{equation}
\oldpage[IV-3]{21}
que asocia a cada sistema inductivo de homomorfismos de $\mathrm{A}_0$-módulos $\theta_\lambda:\mathrm{M}_0\to\mathrm{N}_\lambda$ su límite inductivo.

Ahora bien, si $\mathrm{M}_0$ es de tipo finito (resp.\ de presentación finita), existe una sucesión exacta $\mathrm{A}_0^m\to\mathrm{M}_0\to 0$ (resp.\ $\mathrm{A}_0^n\to\mathrm{A}_0^m\to\mathrm{M}_0\to 0$); como es claro que \sref{IV.8.5.2.3} es biyectivo cuando $\mathrm{M}_0$ es de la forma $\mathrm{A}_0^q$, basta utilizar la exactitud por la izquierda del functor $\mathrm{M}_0\mapsto\operatorname{Hom}_{\mathrm{A}_0}(\mathrm{M}_0,\mathrm{P})$ y la exactitud del functor $\varinjlim$ (en la categoría de grupos abelianos) para concluir.

\medskip
Pasemos al caso en que $\mathrm{S}_0$ es cuasi-compacto, y sea $(\mathrm{U}_i)$ un recubrimiento finito de $\mathrm{S}_0$ por abiertos afines; para cada $\lambda$, los $\mathrm{U}_{i\lambda}=u_{0\lambda}^{-1}(\mathrm{U}_i)$ forman un recubrimiento abierto afín de $\mathrm{S}_\lambda$, y los $\mathrm{V}_i=u_0^{-1}(\mathrm{U}_i)$ un recubrimiento abierto afín de $\mathrm{S}$.
Para demostrar que $u_{\sh{F},\sh{G}}^*$ es inyectivo, basta demostrar que si $f_\lambda:\sh{F}_\lambda\to\sh{G}_\lambda$ satisface $f=u_\lambda^*(f_\lambda)=0$, entonces existe $\mu\geqslant\lambda$ tal que $u_{\lambda\mu}^*(f_\lambda)=0$.
En virtud del lema \sref{IV.8.5.2.1}, para cada $i$ existe un índice $\lambda_i$ tal que $f_\mu|_{\mathrm{U}_{i\mu}}=0$ para $\mu\geqslant\lambda_i$.
Por tanto basta tomar $\mu$ mayor que todos los $\lambda_i$.

Supongamos además que $\mathrm{S}_0$ es cuasi-separado y que $\sh{F}_0$ es de presentación finita, y sea $f:\sh{F}\to\sh{G}$ un homomorfismo de Módulos sobre $\sh{O}_\mathrm{S}$.
En virtud del lema \sref{IV.8.5.2.1}, para cada $i$ existen un índice $\lambda_i$ y un homomorfismo $f_{\lambda_i}^{(i)}:\sh{F}_{\lambda_i}|\mathrm{U}_{i,\lambda_i}\to\sh{G}_{\lambda_i}|\mathrm{U}_{i,\lambda_i}$ tales que $u_{\lambda_i}^*(f_{\lambda_i}^{(i)})=f|\mathrm{V}_i$.
Como $\mathrm{L}$ es filtrado, podemos suponer además que todos los $\lambda_i$ son iguales a un único $\lambda$.
Ahora $\mathrm{S}_\lambda$ es cuasi-separado \sref{IV.1.2.3}, y $\sh{F}_\lambda$ es un Módulo sobre $\sh{O}_{\mathrm{S}_\lambda}$ cuasi-coherente y de presentación finita \sref[0]{0.5.2.5}; puesto que, para cada par de índices $i,j$, $\mathrm{U}_{ij\lambda}=\mathrm{U}_{i\lambda}\cap\mathrm{U}_{j\lambda}$ es cuasi-compacto y, por definición,
\[
  u_\lambda^*(f_\lambda^{(i)}|\mathrm{U}_{ij\lambda})=u_\lambda^*(f_\lambda^{(j)}|\mathrm{U}_{ij\lambda})=f|(\mathrm{V}_i\cap\mathrm{V}_j),
\]
se deduce de lo demostrado anteriormente que existe un índice $\lambda_{ij}$ tal que
\[
  u_{\lambda\mu}^*(f_\lambda^{(i)}|\mathrm{U}_{ij\lambda})=u_{\lambda\mu}^*(f_\lambda^{(j)}|\mathrm{U}_{ij\lambda})\quad\text{para }\mu\geqslant\lambda_{ij}.
\]
Tomando $\mu$ mayor que todos los $\lambda_{ij}$, se ve que $u_{\lambda\mu}^*(f_\lambda^{(i)})$ y $u_{\lambda\mu}^*(f_\lambda^{(j)})$ coinciden en $\mathrm{U}_{i\mu}\cap\mathrm{U}_{j\mu}$ para cada par $(i,j)$, y por tanto definen un homomorfismo $f_\mu:\sh{F}_\mu\to\sh{G}_\mu$ tal que $f=u_\mu^*(f_\mu)$.

\end{enumerate}
\end{proof}

Antes de pasar a la demostración de (ii), observemos los corolarios siguientes de (i):

\begin{corollary}[8.5.2.4]
\label{IV.8.5.2.4}
Supongamos que $\mathrm{S}_0$ es cuasi-compacto, $\sh{F}_\lambda$ cuasi-coherente de tipo finito y $\sh{G}_\lambda$ cuasi-coherente de presentación finita.
Sea $f_\lambda:\sh{F}_\lambda\to\sh{G}_\lambda$ un homomorfismo.
Para que $f=u_\lambda^*(f_\lambda):\sh{F}\to\sh{G}$ sea un isomorfismo, es necesario y suficiente que exista $\mu\geqslant\lambda$ tal que $f_\mu=u_{\lambda\mu}^*(f_\lambda):\sh{F}_\mu\to\sh{G}_\mu$ lo sea.
\end{corollary}

\begin{proof}
Siempre podemos suponer $\mathrm{S}_\lambda=\mathrm{S}_0$; como la cuestión es local en $\mathrm{S}_0$, podemos además (puesto que $\mathrm{S}_0$ es cuasi-compacto y $\mathrm{L}$ es filtrado) reducirnos al caso en que $\mathrm{S}_0$ es afín, y por tanto cuasi-separado.
La condición es trivialmente suficiente; queda demostrar que es necesaria: por hipótesis existe un homomorfismo de $\sh{O}_\mathrm{S}$-Módulos $g:\sh{G}\to\sh{F}$ tal que $g\circ f=1_\sh{F}$ y $f\circ g=1_\sh{G}$.
Como $\sh{G}$ es de presentación finita, existen un índice $\nu\geqslant\lambda$ y un homomorfismo $g_\nu:\sh{G}_\nu\to\sh{F}_\nu$ tales que $g=u_\nu^*(g_\nu)$ en virtud de \sref{IV.8.5.2},~(i); por tanto se tiene
\[
  u_\nu^*(g_\nu\circ f_\nu)=1_\sh{F}
\]
\oldpage[IV-3]{22}
y $u_\nu^*(f_\nu\circ g_\nu)=1_\sh{G}$; como $\sh{F}_\nu$ y $\sh{G}_\nu$ son de tipo finito, se deduce de \sref{IV.8.5.2},~(i), que existe $\mu\geqslant\nu$ tal que $g_\mu\circ f_\mu=1_{\sh{F}_\mu}$ y $f_\mu\circ g_\mu=1_{\sh{G}_\mu}$, de donde se obtiene el corolario.
\end{proof}

\begin{corollary}[8.5.2.5]
\label{IV.8.5.2.5}
Supongamos que $\mathrm{S}_0$ es cuasi-compacto y cuasi-separado.
Supongamos que $\sh{F}_\lambda$ y $\sh{G}_\lambda$ son Módulos sobre $\sh{O}_{\mathrm{S}_\lambda}$ cuasi-coherentes y de presentación finita.
Para que $\sh{F}$ y $\sh{G}$ sean isomorfos, es necesario y suficiente que exista $\mu\geqslant\lambda$ tal que $\sh{F}_\mu$ y $\sh{G}_\mu$ sean isomorfos.
Además, para todo isomorfismo $f:\sh{F}\xrightarrow{\sim}\sh{G}$, existen $\nu\geqslant\mu$ y un isomorfismo $f_\nu:\sh{F}_\nu\xrightarrow{\sim}\sh{G}_\nu$ tales que $f=u_\nu^*(f_\nu)$.
\end{corollary}

\begin{proof}
Esto se deduce de \sref{IV.8.5.2.4} y de \sref{IV.8.5.2},~(i), ya que todo homomorfismo $f:\sh{F}\to\sh{G}$ es de la forma $u_\mu^*(f_\mu)$ para algún $\mu\geqslant\lambda$ y algún homomorfismo $f_\mu:\sh{F}_\mu\to\sh{G}_\mu$.
\end{proof}

\begin{proof}[Demostración de \sref{IV.8.5.2},~(ii)]
Consideremos de nuevo primero el caso en que $\mathrm{S}_0=\operatorname{Spec}(\mathrm{A}_0)$ es afín.
Entonces la afirmación es equivalente al lema \sref{IV.5.13.7.1}.

En el caso general, partiendo de un recubrimiento abierto afín finito $(\mathrm{U}_i)$ de $\mathrm{S}_0$, se deduce de \sref{IV.5.13.7.1} que para cada $i$ existen un índice $\lambda(i)$ y un Módulo sobre $\sh{O}_{\mathrm{U}_{i,\lambda(i)}}$ cuasi-coherente y de presentación finita $\sh{F}_i$ tales que $u_{\lambda(i)}^*(\sh{F}_i)=\sh{F}|_{\mathrm{V}_i}$ (con las notaciones de (i)).
Además, como $\mathrm{L}$ es filtrado, podemos suponer que todos los $\lambda_i$ son iguales a un único $\lambda$.
Como $\mathrm{U}_{ij\lambda}=\mathrm{U}_{i\lambda}\cap\mathrm{U}_{j\lambda}$ es cuasi-compacto y cuasi-separado \sref{IV.1.2.7}, se deduce de \sref{IV.8.5.2.5} que para cada par $(i,j)$ existen un índice $\lambda_{ij}\geqslant\lambda$ y un isomorfismo
\[
  \theta_{ij}:u_{\lambda,\lambda_{ij}}^*(\sh{F}_i|\mathrm{U}_{ij\lambda})\xrightarrow{\sim}u_{\lambda,\lambda_{ij}}^*(\sh{F}_j|\mathrm{U}_{ij\lambda})
\]
tal que $u_{\lambda_{ij}}^*(\theta_{ij})$ es el automorfismo identidad de $\sh{F}|(\mathrm{V}_i\cap\mathrm{V}_j)$; podemos suponer además que todos los $\lambda_{ij}$ son iguales a $\lambda$.
Cambiando las notaciones, podemos suponer por tanto que para cada par $(i,j)$ existe un isomorfismo $\theta_{ij}:\sh{F}_i|_{\mathrm{U}_{ij\lambda}}\xrightarrow{\sim}\sh{F}_j|_{\mathrm{U}_{ij\lambda}}$ tal que $u_\lambda^*(\theta_{ij})$ es el automorfismo identidad de $\sh{F}|_{(\mathrm{V}_i\cap\mathrm{V}_j)}$.
Finalmente, para cualesquiera tres índices $i$, $j$, $k$, si ponemos $\mathrm{U}_{ijk,\lambda}=\mathrm{U}_{i\lambda}\cap\mathrm{U}_{j\lambda}\cap\mathrm{U}_{k\lambda}$, $\mathrm{U}_{ijk,\lambda}$ es cuasi-compacto, y si $\theta'_{ij}$, $\theta'_{jk}$ y $\theta'_{ik}$ denotan las restricciones de $\theta_{ij}$, $\theta_{jk}$ y $\theta_{ik}$ a $\mathrm{U}_{ijk,\lambda}$, se tiene $u_\lambda^*(\theta'_{ij}\circ\theta'_{jk})=u_\lambda^*(\theta'_{ik})$.
Por tanto, en virtud de (i), existe un índice $\mu\geqslant\lambda$ tal que $u_{\lambda\mu}^*(\theta'_{ik})=u_{\lambda\mu}^*(\theta'_{ij}\circ\theta'_{jk})$; en consecuencia, los isomorfismos $u_{\lambda\mu}^*(\theta'_{ij}):u_{\lambda\mu}^*(\sh{F}_i)|_{\mathrm{U}_{ij\mu}}\xrightarrow{\sim}u_{\lambda\mu}^*(\sh{F}_j)|_{\mathrm{U}_{ij\mu}}$ satisfacen la condición de pegado y definen en $\mathrm{S}_\mu$ un Módulo sobre $\sh{O}_{\mathrm{S}_\mu}$ cuasi-coherente y de presentación finita $\sh{F}_\mu$ \sref[0]{0.3.3.1} tal que $\sh{F}$ es isomorfo a $u_\mu^*(\sh{F}_\mu)$.
\end{proof}

\begin{env}[8.5.3]
\label{IV.8.5.3}
El resultado de \sref{IV.8.5.2} también puede expresarse diciendo que, si $\mathrm{S}_0$ es cuasi-compacto y cuasi-separado, la categoría de Módulos sobre $\sh{O}_\mathrm{S}$ cuasi-coherentes y de presentación finita queda determinada, salvo equivalencia, por los datos de las categorías de Módulos sobre $\sh{O}_{\mathrm{S}_\lambda}$ cuasi-coherentes y de presentación finita, los funtores $u_{\lambda\mu}^*$ entre estas categorías y los isomorfismos de transición $u_{\mu\nu}^*\circ u_{\lambda\mu}^*\xrightarrow{\sim}u_{\lambda\nu}^*$.
De manera figurada, puede decirse que el dato de un Módulo sobre $\sh{O}_\mathrm{S}$ cuasi-coherente y de presentación finita $\sh{F}$ es ``funcionalmente'' equivalente al dato de un Módulo sobre $\sh{O}_{\mathrm{S}_\lambda}$ cuasi-coherente y de presentación finita $\sh{F}'_\lambda$ para algún $\lambda$ suficientemente grande y que, si otro Módulo sobre $\sh{O}_{\mathrm{S}_\mu}$ cuasi-coherente y de presentación finita $\sh{F}''_\mu$ tiene también a $\sh{F}$ como imagen inversa, entonces $\sh{F}'_\lambda$ y $\sh{F}''_\mu$ tienen la misma imagen inversa en algún $\mathrm{S}_\nu$ adecuado ($\nu\geqslant\lambda$, $\nu\geqslant\mu$).

A partir de este punto interpretaremos desde este punto de vista diversas nociones relacionadas con los Módulos sobre $\sh{O}_\mathrm{S}$ cuasi-coherentes.
\end{env}

\oldpage[IV-3]{23}
\begin{corollary}[8.5.4]
\label{IV.8.5.4}
Supongamos que $\mathrm{S}_0$ es cuasi-compacto y cuasi-separado; entonces, para todo Módulo sobre $\sh{O}_{\mathrm{S}_\lambda}$ cuasi-coherente $\sh{G}_\lambda$, el homomorfismo canónico \sref{IV.8.5.1.4} es biyectivo.
\end{corollary}

\begin{proof}
En efecto, basta aplicar \sref{IV.8.5.2},~(i), al caso $\sh{F}_\lambda=\sh{O}_{\mathrm{S}_\lambda}$, que es de presentación finita.
\end{proof}

\begin{proposition}[8.5.5]
\label{IV.8.5.5}
Supongamos que $\mathrm{S}_0$ es cuasi-compacto y que $\sh{F}_\lambda$ es un Módulo sobre $\sh{O}_{\mathrm{S}_\lambda}$ cuasi-coherente y de presentación finita.
Para que $\sh{F}$ sea localmente libre (resp.\ localmente libre de rango $n$), es necesario y suficiente que exista $\mu\geqslant\lambda$ tal que $\sh{F}_\mu$ lo sea.
\end{proposition}

\begin{proof}
La condición es trivialmente suficiente; demostremos que es necesaria.
Si $\sh{F}$ es localmente libre (resp.\ localmente libre de rango $n$), existe un recubrimiento abierto afín finito $(\mathrm{V}_i)$ de $\mathrm{S}$ tal que $\sh{F}|_{\mathrm{V}_i}$ es isomorfo a $\sh{O}_\mathrm{S}^{n_i}|_{\mathrm{V}_i}$ (resp.\ $\sh{O}_\mathrm{S}^n|_{\mathrm{V}_i}$) para cada $i$.
En virtud de \sref{IV.8.2.11}, existen $\nu\geqslant\lambda$ y, para cada $i$, un abierto cuasi-compacto $\mathrm{U}_{i\nu}$ de $\mathrm{S}_\nu$ tal que $\mathrm{V}_i=u_\nu^{-1}(\mathrm{U}_{i\nu})$.
Como $\mathrm{S}_\nu$ es cuasi-compacto, cada $\mathrm{U}_{i\nu}$ es unión finita de afines abiertos; así nos reducimos al caso en que $\mathrm{S}_0$ es afín y $\sh{F}=\sh{O}_\mathrm{S}^n$; entonces sabemos que existe $\mu\geqslant\lambda$ tal que $\sh{F}_\mu$ es isomorfo a $\sh{O}_{\mathrm{S}_\mu}^n$ \sref{IV.8.5.2.5}.
\end{proof}

\begin{proposition}[8.5.6]
\label{IV.8.5.6}
Supongamos que $\mathrm{S}_0$ es cuasi-compacto y consideremos una sucesión
\[
  \sh{F}_\lambda\to\sh{G}_\lambda\to\sh{H}_\lambda\to 0
\]
de homomorfismos de Módulos sobre $\sh{O}_{\mathrm{S}_\lambda}$ cuasi-coherentes, donde $\sh{F}_\lambda$ y $\sh{G}_\lambda$ son de tipo finito y $\sh{H}_\lambda$ de presentación finita.
Para que la sucesión correspondiente $\sh{F}\to\sh{G}\to\sh{H}\to 0$ sea exacta, es necesario y suficiente que exista $\mu\geqslant\lambda$ tal que la sucesión $\sh{F}_\mu\to\sh{G}_\mu\to\sh{H}_\mu\to 0$ lo sea (en cuyo caso lo mismo vale para $\nu\geqslant\mu$).
\end{proposition}

\begin{proof}
El hecho de que la condición sea suficiente y la última afirmación se deben a que el functor $u_\lambda^*$ (resp.\ $u_{\mu\nu}^*$) es exacto por la derecha.
Para demostrar que la condición es necesaria, observemos que de la hipótesis y de \sref{IV.8.5.2},~(i), se deduce que existe $\nu\geqslant\lambda$ tal que la composición $\sh{F}_\nu\to\sh{G}_\nu\to\sh{H}_\nu$ es nula.
Si ponemos $\sh{H}'_\nu=\operatorname{Coker}(\sh{F}_\nu\to\sh{G}_\nu)$, tenemos por tanto un homomorfismo $f_\nu:\sh{H}'_\nu\to\sh{H}_\nu$; por hipótesis, $u_\nu^*(f_\nu)$ es un isomorfismo, y se deduce de \sref{IV.8.5.2.4} que existe $\mu\geqslant\nu$ tal que $u_{\nu\mu}^*(f_\nu)$ es un isomorfismo, lo que concluye la demostración.
\end{proof}

\begin{corollary}[8.5.7]
\label{IV.8.5.7}
Supongamos que $\mathrm{S}_0$ es cuasi-compacto, $\sh{F}_\lambda$ cuasi-coherente, $\sh{G}_\lambda$ cuasi-coherente de tipo finito, y sea $f_\lambda:\sh{F}_\lambda\to\sh{G}_\lambda$ un homomorfismo.
Para que $f=u_\lambda^*(f_\lambda):\sh{F}\to\sh{G}$ sea sobreyectivo, es necesario y suficiente que exista $\mu\geqslant\lambda$ tal que $f_\mu=u_{\lambda\mu}^*(f_\lambda):\sh{F}_\mu\to\sh{G}_\mu$ lo sea.
\end{corollary}

\begin{proof}
Este es el caso particular de \sref{IV.8.5.6} aplicado a la sucesión $0\to\sh{H}_\lambda\to 0\to 0$, donde $\sh{H}_\lambda=\operatorname{Coker}(f_\lambda)$, que es cuasi-coherente y de tipo finito (teniendo en cuenta que $\sh{H}=\operatorname{Coker}(f)$ y $\sh{H}_\mu=\operatorname{Coker}(f_\mu)$, en virtud de la exactitud por la derecha de los funtores $u_\lambda^*$ y $u_{\lambda\mu}^*$).
\end{proof}

\begin{corollary}[8.5.8]
\label{IV.8.5.8}
Supongamos que $\mathrm{S}_0$ es cuasi-compacto y que los morfismos $u_{\lambda\mu}:\mathrm{S}_\mu\to\mathrm{S}_\lambda$ son planos.
Entonces:
\begin{enumerate}[label=(\roman*)]
  \item Sea $\sh{F}_\lambda\xrightarrow{f_\lambda}\sh{G}_\lambda\xrightarrow{g_\lambda}\sh{H}_\lambda$ una sucesión de homomorfismos de Módulos sobre $\sh{O}_{\mathrm{S}_\lambda}$ cuasi-coherentes, tales que $\operatorname{Im} f_\lambda$ y $\operatorname{Ker} g_\lambda$ son de tipo finito.
Para que la sucesión correspondiente $\sh{F}\xrightarrow{f}\sh{G}\xrightarrow{g}\sh{H}$ sea
\oldpage[IV-3]{24}
exacta, es necesario y suficiente que exista $\mu\geqslant\lambda$ tal que la sucesión $\sh{F}_\mu\xrightarrow{f_\mu}\sh{G}_\mu\xrightarrow{g_\mu}\sh{H}_\mu$ lo sea (en cuyo caso lo mismo vale para $\nu\geqslant\mu$).
  \item Sea $f_\lambda:\sh{F}_\lambda\to\sh{G}_\lambda$ un homomorfismo de Módulos sobre $\sh{O}_{\mathrm{S}_\lambda}$ cuasi-coherentes, y supongamos que $\operatorname{Ker} f_\lambda$ es de tipo finito.
Para que $f=u_\lambda^*(f_\lambda):\sh{F}\to\sh{G}$ sea inyectivo, es necesario y suficiente que exista $\mu\geqslant\lambda$ tal que $f_\mu=u_{\lambda\mu}^*(f_\lambda):\sh{F}_\mu\to\sh{G}_\mu$ lo sea.
\end{enumerate}
\end{corollary}

\begin{proof}
(i) Teniendo en cuenta \sref{IV.8.3.8},~(ii), observemos que, por planitud, $\operatorname{Im} f$ y $\operatorname{Ker} g$ (resp. $\operatorname{Im} f_\mu$ y $\operatorname{Ker} g_\mu$ para $\mu\geqslant\lambda$) son las imágenes inversas de $\operatorname{Im} f_\lambda$ y $\operatorname{Ker} g_\lambda$ \sref[0\textsubscript{I}]{0.6.7.2}. Supongamos que la sucesión $\sh{F}\xrightarrow{f}\sh{G}\xrightarrow{g}\sh{H}$ es exacta. Como $\operatorname{Im} f_\lambda$ es de tipo finito, se deduce de \sref{IV.8.5.2},~(i), que existe $\mu\geqslant\lambda$ tal que la composición $\sh{F}_\mu\xrightarrow{f_\mu}\sh{G}_\mu\xrightarrow{g_\mu}\sh{H}_\mu$ es nula. Cambiando la notación, podemos suponer ya que $g_\lambda\circ f_\lambda=0$. Como $\operatorname{Ker} g_\lambda$ es de tipo finito, se deduce entonces de \sref{IV.8.5.7} que existe $\mu\geqslant\lambda$ tal que el homomorfismo $\sh{F}_\mu\to\operatorname{Ker} g_\mu$ es sobreyectivo, lo que demuestra (i).

(ii) Este es el caso particular de (i) aplicado a la sucesión $0\to\sh{F}_\lambda\to\sh{G}_\lambda$.
\end{proof}

\begin{lemma}[8.5.9]
\label{IV.8.5.9}
Supongamos que $\mathrm{S}_0$ es cuasi-compacto y que $\sh{F}_\lambda$ es cuasi-coherente y de tipo finito. Sean $\sh{G}'_\lambda$ y $\sh{G}''_\lambda$ dos cocientes cuasi-coherentes de $\sh{F}_\lambda$, suponiendo además que $\sh{G}'_\lambda$ es de presentación finita. Para que $\sh{G}''$ sea un cociente de $\sh{G}'$, es necesario y suficiente que exista $\mu\geqslant\lambda$ tal que $\sh{G}''_\mu$ sea un cociente de $\sh{G}'_\mu$.
\end{lemma}

\begin{proof}
Por hipótesis existen dos homomorfismos sobreyectivos $p'_\lambda:\sh{F}_\lambda\to\sh{G}'_\lambda$ y $p''_\lambda:\sh{F}_\lambda\to\sh{G}''_\lambda$. Por la exactitud por la derecha de $u_\lambda^*$ y $u_{\lambda\mu}^*$, los homomorfismos $p'=u_\lambda^*(p'_\lambda)$, $p''=u_\lambda^*(p''_\lambda)$, $p'_\mu=u_{\lambda\mu}^*(p'_\lambda)$ y $p''_\mu=u_{\lambda\mu}^*(p''_\lambda)$ son también sobreyectivos. Además, si existe un homomorfismo $f:\sh{G}'\to\sh{G}''$ (resp. $f_\mu:\sh{G}'_\mu\to\sh{G}''_\mu$) tal que $p''=f\circ p'$ (resp. $p''_\mu=f_\mu\circ p'_\mu$), dicho homomorfismo es necesariamente único; esto muestra que la cuestión es local en $\mathrm{S}_\lambda$ y, en consecuencia, como $\mathrm{S}_0$ es cuasi-compacto y $\mathbf{L}$ es filtrado, podemos suponer que $\mathrm{S}_\lambda$ es afín y por tanto cuasi-separado. La condición del enunciado es claramente suficiente. Recíprocamente, como $\sh{G}'_\lambda$ es de presentación finita y $\mathrm{S}_\lambda$ es cuasi-compacto y cuasi-separado, se deduce de \sref{IV.8.5.2},~(i), que si existe $f:\sh{G}'\to\sh{G}''$ tal que $p''=f\circ p'$, entonces existen $\mu\geqslant\lambda$ y $f_\mu:\sh{G}'_\mu\to\sh{G}''_\mu$ tales que $f=u_\mu^*(f_\mu)$ y $p''_\mu=f_\mu\circ p'_\mu$, lo que demuestra el lema.
\end{proof}

\begin{env}[8.5.10]
\label{IV.8.5.10}
Para el resto de esta sección, para todo Módulo cuasi-coherente $\sh{F}$ sobre un preesquema, denotemos por $\mathfrak{Q}(\sh{F})$ el conjunto de los Módulos cociente de $\sh{F}$ que son de presentación finita. Si $\sh{F}_\lambda$ es cuasi-coherente y $\sh{G}_\lambda\in\mathfrak{Q}(\sh{F}_\lambda)$, se deduce de la exactitud por la derecha de $u_{\lambda\mu}^*$ y $u_\lambda^*$ que $\sh{G}_\mu\in\mathfrak{Q}(\sh{F}_\mu)$ para $\mu\geqslant\lambda$ y $\sh{G}\in\mathfrak{Q}(\sh{F})$. Es claro que $(\mathfrak{Q}(\sh{F}_\lambda),u_{\lambda\mu}^*)$ es un sistema inductivo de conjuntos y que los $u_\lambda^*:\mathfrak{Q}(\sh{F}_\lambda)\to\mathfrak{Q}(\sh{F})$ forman un sistema inductivo de aplicaciones; por tanto, pasando al límite inductivo, existe una aplicación canónica
\[
u_{\mathfrak{Q}}:\varinjlim\mathfrak{Q}(\sh{F}_\lambda)\longrightarrow\mathfrak{Q}(\sh{F}).
\tag{8.5.10.1}\label{IV.8.5.10.1}
\]
Además, si $(\sh{F}'_\lambda)$ es una segunda familia de Módulos sobre $\sh{O}_{\mathrm{S}_\lambda}$ cuasi-coherentes y, para cada $\lambda$, $\sh{F}'_\lambda$ es un cociente de $\sh{F}_\lambda$, entonces $\sh{F}'$ es un cociente de $\sh{F}$, y se tiene un diagrama conmutativo

\oldpage[IV-3]{25}

\[
\begin{tikzcd}
\varinjlim\mathfrak{Q}(\sh{F}_\lambda) \ar[r,"u_{\mathfrak{Q}}"] \ar[d] & \mathfrak{Q}(\sh{F}) \ar[d] \\
\varinjlim\mathfrak{Q}(\sh{F}'_\lambda) \ar[r,"u_{\mathfrak{Q}}"'] & \mathfrak{Q}(\sh{F}')
\end{tikzcd}
\tag{8.5.10.2}\label{IV.8.5.10.2}
\]
\end{env}

\begin{proposition}[8.5.11]
\label{IV.8.5.11}
Supongamos que $\mathrm{S}_0$ es cuasi-compacto (resp. cuasi-compacto y cuasi-separado). Supongamos que $\sh{F}_\lambda$ es cuasi-coherente y de tipo finito (resp. de presentación finita) para cada $\lambda$; entonces la aplicación canónica \eqref{IV.8.5.10.1} es inyectiva (resp. biyectiva).
\end{proposition}

\begin{proof}
La primera afirmación se deduce con mayor precisión del lema \sref{IV.8.5.9}. Para la segunda, consideremos un Módulo cociente $\sh{O}_{\mathrm{S}}$-cuasi-coherente $\sh{G}$ de $\sh{F}$ que sea de presentación finita. Se deduce de \sref{IV.8.5.2},~(ii), que existen $\lambda'\in\mathbf{L}$ y un Módulo sobre $\sh{O}_{\mathrm{S}_{\lambda'}}$ cuasi-coherente y de presentación finita $\sh{G}_{\lambda'}$ tales que $\sh{G}=u_{\lambda'}^*(\sh{G}_{\lambda'})$; como $\mathbf{L}$ es filtrado, podemos reemplazar $\lambda$ y $\lambda'$ por un índice que domine a ambos y suponer que $\lambda'=\lambda$. Consideremos entonces el homomorfismo canónico $p:\sh{F}\to\sh{G}$. Se deduce de \sref{IV.8.5.2},~(i), que existen $\mu\geqslant\lambda$ y un homomorfismo $p_\mu:\sh{F}_\mu\to\sh{G}_\mu$ tales que $p=u_\mu^*(p_\mu)$. En virtud de \sref{IV.8.5.7}, podemos suponer además que $\mu$ es suficientemente grande para que $p_\mu$ sea sobreyectivo, lo que concluye la demostración.
\end{proof}

\subsection{Subpreesquemas de presentación finita de un límite proyectivo de preesquemas}
\label{subsection:IV.8.6}

\begin{env}[8.6.1]
\label{IV.8.6.1}
Dado un preesquema $Y$, denotamos en esta sección por $\mathfrak{Spr}(Y)$ el conjunto ordenado
\sref[I]{I.4.1.10} de los subpreesquemas de $Y$ que son \emph{de presentación finita}
\sref{IV.1.6.1}, y por $\mathfrak{Spro}(Y)$ (resp.\ $\mathfrak{Sprf}(Y)$) el subconjunto de $\mathfrak{Spr}(Y)$ formado por los subpreesquemas inducidos en abiertos (resp.\ en subpreesquemas cerrados) de $Y$, de presentación finita sobre $Y$; equivale a decir que un subpreesquema $Z$ de $Y$ pertenece a $\mathfrak{Spro}(Y)$ (resp.\ $\mathfrak{Sprf}(Y)$) o bien que está inducido en un abierto y que el espacio subyacente es \emph{retrocompacto} en $Y$ (resp.\ que es cerrado y que el ideal $\sh{J}$ de $\sh{O}_Y$ que lo define es \emph{de tipo finito}, lo que también significa que $j_*(\sh{O}_Z)\in\mathfrak{Q}(\sh{O}_Y)$, donde $j:Z\to Y$ es la inyección canónica) \sref{IV.1.6.1} y \sref{IV.1.4.5}.
\end{env}

\begin{env}[8.6.2]
\label{IV.8.6.2}
Conservamos siempre las notaciones de \sref{IV.8.2.2} y suponemos que $S_0$ es uno de los $S_\lambda$.
Dado un subpreesquema $Y_\lambda$ de $S_\lambda$, entonces $Y_\mu=u_{\lambda\mu}^{-1}(Y_\lambda)$ (resp.\ $Y=u_\lambda^{-1}(Y_\lambda)$) es un subpreesquema de $S_\mu$ para $\mu\geqslant\lambda$ (resp.\ de $S$). Es inducido en un abierto (resp.\ es cerrado) si $Y_\lambda$ lo es \sref[I]{I.4.3.2}, y es de presentación finita sobre $S_\mu$ (resp.\ $S$) si $Y_\lambda$ es de presentación finita sobre $S_\lambda$ \sref{IV.1.6.2},~(iii). Se sigue que $(\mathfrak{Spr}(S_\lambda),u_{\lambda\mu}^{-1})$ (resp.\ $(\mathfrak{Spro}(S_\lambda),u_{\lambda\mu}^{-1})$, $(\mathfrak{Sprf}(S_\lambda),u_{\lambda\mu}^{-1})$) es un \emph{sistema inductivo} de conjuntos y que las aplicaciones $u_\lambda^{-1}:\mathfrak{Spr}(S_\lambda)\to\mathfrak{Spr}(S)$ (resp.\ $\mathfrak{Spro}(S_\lambda)\to\mathfrak{Spro}(S)$, $\mathfrak{Sprf}(S_\lambda)\to\mathfrak{Sprf}(S)$) forman un sistema inductivo de aplicaciones; de donde, pasando al límite inductivo, se obtienen aplicaciones canónicas
\end{env}

\begin{align}
  \varinjlim \mathfrak{Spr}(S_\lambda) &\to \mathfrak{Spr}(S) \tag{8.6.2.1}\label{IV.8.6.2.1}\\
  \varinjlim \mathfrak{Spro}(S_\lambda) &\to \mathfrak{Spro}(S) \tag{8.6.2.2}\label{IV.8.6.2.2}\\
  \varinjlim \mathfrak{Sprf}(S_\lambda) &\to \mathfrak{Sprf}(S). \tag{8.6.2.3}\label{IV.8.6.2.3}
\end{align}

\oldpage[IV-3]{26}

Recordemos de \sref[I]{I.4.1.10} que $\mathfrak{Spr}(X)$, para todo preesquema $X$, es un conjunto ordenado por la relación «$Z$ es un subpreesquema del subpreesquema $Y$», que escribimos $Z\leqslant Y$.
Las aplicaciones $u_{\lambda\mu}^{-1}$ y $u_\lambda^{-1}$ son \emph{crecientes} para las relaciones de orden correspondientes en $\mathfrak{Spr}(S_\lambda)$, $\mathfrak{Spr}(S_\mu)$ y $\mathfrak{Spr}(S)$.
Además, escribiendo $\eta\leqslant\eta'$ para dos elementos de este conjunto siempre que existan $\lambda$ y dos elementos $Y_\lambda$, $Y'_\lambda$ de $\mathfrak{Spr}(S_\lambda)$, cuyas imágenes canónicas sean $\eta$ y $\eta'$, y tales que $Y_\lambda\leqslant Y'_\lambda$, se verifica fácilmente que esto no depende de la elección de representantes y que, en efecto, se obtiene una \emph{relación de orden}.
El hecho de que las $u_\lambda^{-1}$ sean crecientes implica inmediatamente que la aplicación canónica \eqref{IV.8.6.2.1} es \emph{creciente}; evidentemente lo mismo vale para \eqref{IV.8.6.2.2} y \eqref{IV.8.6.2.3}.

\begin{proposition}[8.6.3]
\label{IV.8.6.3}
Supongamos que $S_0$ es cuasi-compacto (resp.\ cuasi-compacto y cuasi-separado).
Entonces las aplicaciones \eqref{IV.8.6.2.1}, \eqref{IV.8.6.2.2} y \eqref{IV.8.6.2.3} son inyectivas (resp.\ biyectivas).
\end{proposition}

\begin{proof}
Teniendo en cuenta las observaciones de \sref{IV.8.6.1}, las afirmaciones relativas a \eqref{IV.8.6.2.3} se deducen de \sref{IV.8.5.11} aplicado a $\sh{F}_\lambda=\sh{O}_{S_\lambda}$; asimismo, las afirmaciones relativas a \eqref{IV.8.6.2.2} son casos particulares de \sref{IV.8.3.5} y \sref{IV.8.3.11}, teniendo en cuenta que los $S_\lambda$ y $S$ son cuasi-compactos.
Queda considerar la aplicación \eqref{IV.8.6.2.1}.
Demostremos primero que es sobreyectiva cuando $S_0$ es cuasi-compacto y cuasi-separado.
Sea $Z$ un subpreesquema de $S$, de presentación finita sobre $S$; como $S$ es cuasi-compacto, también lo es $Z$, y por tanto existe un abierto cuasi-compacto $U$ de $S$ tal que $Z$ es un subpreesquema de $U$, cerrado en $U$, de presentación finita sobre $U$.
Existen entonces un índice $\lambda$ y un abierto cuasi-compacto $U_\lambda$ de $S_\lambda$ tales que $U=u_\lambda^{-1}(U_\lambda)$ \sref{IV.8.2.11}; como $S_\lambda$ es cuasi-separado, $U_\lambda$ también lo es \sref[I]{I.1.2.7}, y por tanto podemos restringirnos al caso $U=S$; pero en este caso nos reducimos al hecho de que \eqref{IV.8.6.2.3} es sobreyectiva.

Finalmente, para ver que \eqref{IV.8.6.2.1} es inyectiva cuando $S_0$ es cuasi-compacto, bastará demostrar el resultado más preciso siguiente:
\end{proof}

\begin{corollary}[8.6.3.1]
\label{IV.8.6.3.1}
Supongamos que $S_0$ es cuasi-compacto y sean $Z'_\lambda$, $Z''_\lambda$ dos subpreesquemas de $S_\lambda$, de presentación finita.
Para que $Z'=u_\lambda^{-1}(Z'_\lambda)$ esté dominado por $Z''=u_\lambda^{-1}(Z''_\lambda)$ \sref[I]{I.4.1.10}, es necesario y suficiente que exista $\mu\geqslant\lambda$ tal que $Z'_\mu=u_{\lambda\mu}^{-1}(Z'_\lambda)$ esté dominado por $Z''_\mu=u_{\lambda\mu}^{-1}(Z''_\lambda)$.
\end{corollary}

\begin{proof}
Es trivial que la condición es suficiente.
Para ver que es necesaria, observemos primero que los conjuntos subyacentes $Z'_\lambda$ y $Z''_\lambda$ son localmente constructibles en $S_\lambda$ por hipótesis \sref[I]{I.1.8.4}; por tanto, la hipótesis implica que existe un $\nu\geqslant\lambda$ tal que $Z'_\nu\subset Z''_\nu$ \sref{IV.8.3.5}; reemplazando $\lambda$ por $\nu$, podemos suponer ya que $Z'_\lambda\subset Z''_\lambda$.
Además, por hipótesis \sref[I]{I.1.6.1}, los subespacios $Z'_\lambda$ y $Z''_\lambda$ de $S_\lambda$ son cuasi-compactos; para cada punto $x\in Z'_\lambda$, existe un entorno abierto cuasi-compacto $V(x)$ en $S_\lambda$ tal que $V(x)\cap Z'_\lambda$ y $V(x)\cap Z''_\lambda$ son cerrados en $V(x)$.
Recubriendo $S_\lambda$ por un número finito de tales entornos, vemos que existe un abierto cuasi-compacto $U_\lambda$ de $S_\lambda$ que contiene a $Z'_\lambda$ y tal que $Z'_\lambda$ y $Z''_\lambda\cap U_\lambda$ son cerrados en $U_\lambda$.
Si denotamos por $Y_\lambda$ el subpreesquema de $S_\lambda$ inducido por $Z''_\lambda$ en $U_\lambda\cap Z''_\lambda$, es claro que, con las notaciones usuales, $Y_\mu$ (resp.\ $Y$) es inducido por $Z''_\mu$ en $U_\mu\cap Z''_\mu$ para $\mu\geqslant\lambda$ (resp.\ por $Z''$ en $U\cap Z''$); además, $Z'$ está dominado por $Y$ \sref[I]{I.4.4.1}, y como basta demostrar que $Z'_\mu$ está dominado por $Y_\mu$ para $\mu$
\oldpage[IV-3]{27}
 suficientemente grande, vemos finalmente que (reemplazando $S_\lambda$ por $U_\lambda$) nos reducimos al caso en que $Z'_\lambda$ y $Z''_\lambda$ son subpreesquemas cerrados de $S_\lambda$.
Pero esto ya se ha demostrado, pues \eqref{IV.8.6.2.3} es creciente e inyectiva.
\end{proof}

\begin{corollary}[8.6.4]
\label{IV.8.6.4}
Supongamos que $S_0$ es cuasi-compacto y sea $Z_\lambda$ un subpreesquema de $S_\lambda$, de presentación finita sobre $S_\lambda$.
Para que $Z=u_\lambda^{-1}(Z_\lambda)$ sea el subpreesquema inducido en un abierto de $S$ (resp.\ un subpreesquema cerrado de $S$), es necesario y suficiente que exista $\mu\geqslant\lambda$ tal que $Z_\mu=u_{\lambda\mu}^{-1}(Z_\lambda)$ sea el subpreesquema inducido en un abierto de $S_\mu$ (resp.\ un subpreesquema cerrado de $S_\mu$).
\end{corollary}

\begin{proof}
Sea $(U_\lambda^{(i)})$ un recubrimiento abierto afín finito de $S_\lambda$, y pongamos $U^{(i)}_\mu=u_{\lambda\mu}^{-1}(U_\lambda^{(i)})$ para $\mu\geqslant\lambda$ y $U^{(i)}=u_\lambda^{-1}(U_\lambda^{(i)})$.
Si $Z$ es abierto (resp.\ cerrado) en $S$, $Z\cap U^{(i)}$ es abierto (resp.\ cerrado) en $U^{(i)}$; recíprocamente, si cada $Z_\mu\cap U^{(i)}_\mu$ es abierto (resp.\ cerrado) en $U^{(i)}_\mu$, entonces $Z_\mu$ es abierto (resp.\ cerrado) en $S_\mu$.
Como $L$ es filtrado, basta demostrar el corolario cuando $S_\lambda$ es afín, y por tanto cuasi-separado; el resultado se deduce del hecho de que las aplicaciones \eqref{IV.8.6.2.1}, \eqref{IV.8.6.2.2} y \eqref{IV.8.6.2.3} son biyectivas.
\end{proof}

\subsection{Criterios para que un límite proyectivo de preesquemas sea un preesquema reducido (resp.\ íntegro)}
\label{subsection:IV.8.7}

Conservamos las hipótesis y notaciones de \sref{IV.8.2.2} y suponemos siempre que $S_0$ es uno de los $S_\lambda$.

\begin{proposition}[8.7.1]
\label{IV.8.7.1}
Supongamos que $S$ no es reducido.
Entonces existe $\lambda_0$ tal que, para $\lambda\geqslant\lambda_0$, $S_\lambda$ no es reducido.
\end{proposition}

\begin{proof}
Como la cuestión es local en $S_0$, podemos suponer que $S_0=\operatorname{Spec}(A_0)$ es afín, de donde $S=\operatorname{Spec}(A)$, donde $A=\varinjlim A_\lambda$ es el límite inductivo de un sistema inductivo de $A_0$-álgebras $(A_\lambda)$.
Sabemos entonces por \sref{IV.5.13.2} que el nilradical de $A$ es el límite inductivo de los nilradicales de los $A_\lambda$; si no es nulo, uno de los $A_\lambda$ contiene por tanto un elemento nilpotente $a_\lambda\neq 0$ cuya imagen en $A$ es nilpotente y no nula, y en consecuencia la imagen de $a_\lambda$ en $A_\mu$ es nilpotente y no nula para $\mu\geqslant\lambda$.
\end{proof}

\begin{proposition}[8.7.2]
\label{IV.8.7.2}
Supongamos que se satisface una de las hipótesis siguientes:
\begin{enumerate}[label=\emph{\alph*)}]
\item $S_0$ es cuasi-compacto, el nilradical $\sh{N}_0$ de $\sh{O}_{S_0}$ es un ideal de tipo finito (lo que se satisface, por ejemplo, cuando $S_0$ es noetheriano), y los morfismos $u_{\lambda\mu}:S_\mu\to S_\lambda$ son inmersiones abiertas.
\item Los morfismos $u_{\lambda\mu}:S_\mu\to S_\lambda$ son fielmente planos.
\end{enumerate}
Bajo estas condiciones, para que $S$ sea reducido, es necesario y suficiente que exista $\lambda_0$ tal que $S_\lambda$ sea reducido para $\lambda\geqslant\lambda_0$.

Además, en el caso b), si $S$ es reducido, todos los $S_\lambda$ son reducidos.
\end{proposition}

\begin{proof}
La última afirmación se deduce de que el morfismo $S\to S_\lambda$ es entonces fielmente plano para todo $\lambda$ \sref{IV.8.3.8}, y de \sref{IV.2.1.13}.
Por otro lado, \sref{IV.8.7.1} demuestra que la condición del enunciado es suficiente (sin hipótesis sobre $S_0$ ni sobre los $u_{\lambda\mu}$).
Queda, pues, ver que la condición es necesaria bajo la hipótesis $a)$; entonces, por \sref{IV.8.2.13}, el espacio subyacente de $S$ se identifica con la intersección de los espacios subyacentes de los $S_\lambda$ (identificándose los $S_\lambda$ con subpreesquemas inducidos en abiertos de $S_0$), y el haz estructural $\sh{O}_S$ se identifica con
\oldpage[IV-3]{28}
el haz inducido en $S$ por cada uno de los $\sh{O}_{S_\lambda}$; en particular, para todo $s\in S$, el anillo local $\sh{O}_s$ es el mismo para $S$ y para todos los $S_\lambda$.
Si $\sh{N}_\lambda$ es el nilradical de $\sh{O}_{S_\lambda}$, el nilradical $\sh{N}$ de $\sh{O}_S$ tiene, por tanto, en cada punto de $S$ la misma fibra $\sh{N}_s$ (el nilradical de $\sh{O}_s$) que $u_\lambda^*(\sh{N}_\lambda)$ (inducido en $S$ por $\sh{N}_\lambda$).
La hipótesis de que $S$ es reducido implica entonces $u_\lambda^*(\sh{N}_\lambda)=0$; como se supone que $\sh{N}_0$ es de tipo finito, lo mismo vale para los $\sh{N}_\lambda$, y la conclusión se deduce de \sref{IV.8.5.7}.
\end{proof}

\begin{corollary}[8.7.3]
\label{IV.8.7.3}
Supongamos que se satisface una de las hipótesis siguientes:
\begin{enumerate}[label=\emph{\alph*)}]
\item $S_0$ es un preesquema noetheriano y los morfismos $u_{\lambda\mu}:S_\mu\to S_\lambda$ son inmersiones abiertas.
\item Los morfismos $u_{\lambda\mu}:S_\mu\to S_\lambda$ son fielmente planos.
\end{enumerate}
Entonces, para que $S$ sea íntegro, es necesario y suficiente que exista $\lambda_0$ tal que $S_\lambda$ sea íntegro para $\lambda\geqslant\lambda_0$.
\end{corollary}

\begin{proof}
Decir que un preesquema es íntegro significa que es a la vez reducido e irreducible; el corolario se deduce, pues, de \sref{IV.8.7.2} y de \sref{IV.8.4.3}.
\end{proof}

\begin{remark}[8.7.4]
\label{IV.8.7.4}
Si no se formula ninguna hipótesis sobre los $u_{\lambda\mu}$, puede ocurrir que $S$ sea íntegro aunque ninguno de los $S_\lambda$ sea reducido ni conexo, como muestra el ejemplo \sref{IV.8.4.4}, tomando el anillo $A$ no reducido.
\end{remark}

\subsection{Preesquemas de presentación finita sobre un límite proyectivo de preesquemas}
\label{subsection:IV.8.8}

\begin{env}[8.8.1]
\label{IV.8.8.1}
Conservando siempre las notaciones e hipótesis de \sref{IV.8.2.2}, supondremos dados en esta sección dos $S_\alpha$-preesquemas $X_\alpha$, $Y_\alpha$, lo que define \sref{IV.8.2.5} dos sistemas proyectivos de preesquemas $(X_\lambda,v_{\lambda\mu})$ e $(Y_\lambda,w_{\lambda\mu})$, poniendo $X_\lambda=X_\alpha\times_{S_\alpha}S_\lambda$, $v_{\lambda\mu}=1_{X_\alpha}\times u_{\lambda\mu}$, $w_{\lambda\mu}=1_{Y_\alpha}\times u_{\lambda\mu}$ (para $\alpha\leqslant\lambda\leqslant\mu$), cuyos límites proyectivos son respectivamente $X=X_\alpha\times_{S_\alpha}S$, $Y=Y_\alpha\times_{S_\alpha}S$, con los morfismos canónicos $v_\lambda:X\to X_\lambda$ y $w_\lambda:Y\to Y_\lambda$.
Para $\alpha\leqslant\lambda\leqslant\mu$, existe una aplicación canónica $e_{\mu\lambda}:\operatorname{Hom}_{S_\lambda}(X_\lambda,Y_\lambda)\to\operatorname{Hom}_{S_\mu}(X_\mu,Y_\mu)$, que a todo $S_\lambda$-morfismo $f_\lambda:X_\lambda\to Y_\lambda$ asocia $f_\mu=f_\lambda\times 1_{S_\mu}:X_\lambda\times_{S_\lambda}S_\mu\to Y_\lambda\times_{S_\lambda}S_\mu$, y es claro que $(\operatorname{Hom}_{S_\lambda}(X_\lambda,Y_\lambda),e_{\mu\lambda})$ es un \emph{sistema inductivo} de conjuntos.
Análogamente, existe una aplicación canónica $e_\lambda:\operatorname{Hom}_{S_\lambda}(X_\lambda,Y_\lambda)\to\operatorname{Hom}_S(X,Y)$ que a $f_\lambda$ asocia $f=f_\lambda\times 1_S:X_\lambda\times_{S_\lambda}S\to Y_\lambda\times_{S_\lambda}S$, y $(e_\lambda)$ es un sistema inductivo de aplicaciones; de ahí, pasando al límite inductivo, una aplicación canónica, functorial en $S_\alpha$, $X_\alpha$ e $Y_\alpha$:
\[
  e:\varinjlim\operatorname{Hom}_{S_\lambda}(X_\lambda,Y_\lambda)\to\operatorname{Hom}_S(X,Y).
  \tag{8.8.1.1}\label{IV.8.8.1.1}
\]
\end{env}

\begin{theorem}[8.8.2]
\label{IV.8.8.2}
\begin{enumerate}[label=(\roman*)]
\item Supongamos que $X_\alpha$ es cuasi-compacto (resp.\ cuasi-compacto y cuasi-separado), y que $Y_\alpha$ es localmente de tipo finito (resp.\ localmente de presentación finita) sobre $S_\alpha$.
Entonces la aplicación \eqref{IV.8.8.1.1} es inyectiva (resp.\ biyectiva).
\item Supongamos que $S_0$ es cuasi-compacto y cuasi-separado.
Para todo preesquema $X$ de presentación finita sobre $S$, existen un $\lambda\in L$, un preesquema $X_\lambda$ de presentación finita sobre $S_\lambda$ y un $S$-isomorfismo $X\xrightarrow{\sim}X_\lambda\times_{S_\lambda}S$.
\end{enumerate}
\end{theorem}

\begin{proof}
\oldpage[IV-3]{29}
(i) Consideremos primero el caso en que $S_0=\operatorname{Spec}(A_0)$, $X_\alpha=\operatorname{Spec}(B_\alpha)$ e $Y_\alpha=\operatorname{Spec}(C_\alpha)$ son afines; entonces los $S_\lambda=\operatorname{Spec}(A_\lambda)$ y $S=\operatorname{Spec}(A)$ son también afines, con $A=\varinjlim A_\lambda$, y las afirmaciones de (i) son equivalentes al

\begin{lemma}[8.8.2.1]
\label{IV.8.8.2.1}
Sea $A_0$ un anillo, $(A_\lambda)_{\lambda\in L}$ un sistema inductivo de $A_0$-álgebras, $A=\varinjlim A_\lambda$; sea $B_\alpha$ una $A_\alpha$-álgebra y $C_\alpha$ una $A_\alpha$-álgebra de tipo finito (resp.\ de presentación finita).
Entonces el homomorfismo canónico
\[
  \varinjlim_\lambda\operatorname{Hom}_{A_\lambda\text{-alg.}}(C_\alpha\otimes_{A_\alpha}A_\lambda,B_\alpha\otimes_{A_\alpha}A_\lambda)\to\operatorname{Hom}_{A\text{-alg.}}(C_\alpha\otimes_{A_\alpha}A,B_\alpha\otimes_{A_\alpha}A)
  \tag{8.8.2.2}\label{IV.8.8.2.2}
\]
es inyectivo (resp.\ biyectivo).
\end{lemma}

\begin{proof}
Sabemos que existen isomorfismos canónicos functoriales
\[
  \operatorname{Hom}_{A_\lambda\text{-alg.}}(C_\alpha\otimes_{A_\alpha}A_\lambda,B_\alpha\otimes_{A_\alpha}A_\lambda)\xrightarrow{\sim}\operatorname{Hom}_{A_\alpha\text{-alg.}}(C_\alpha,B_\alpha\otimes_{A_\alpha}A_\lambda)
\]
\[
  \operatorname{Hom}_{A\text{-alg.}}(C_\alpha\otimes_{A_\alpha}A,B_\alpha\otimes_{A_\alpha}A)\xrightarrow{\sim}\operatorname{Hom}_{A_\alpha\text{-alg.}}(C_\alpha,B_\alpha\otimes_{A_\alpha}A)
\]
en virtud de la propiedad universal del producto tensorial de dos álgebras.
Basta, por tanto, demostrar el
\end{proof}

\begin{lemma}[8.8.2.3]
\label{IV.8.8.2.3}
Sea $E$ un anillo, $G$ una $E$-álgebra de tipo finito (resp.\ de presentación finita), $(F_\lambda)$ un sistema inductivo de $E$-álgebras.
Entonces el homomorfismo canónico
\[
  \varinjlim\operatorname{Hom}_{E\text{-alg.}}(G,F_\lambda)\to\operatorname{Hom}_{E\text{-alg.}}(G,\varinjlim F_\lambda)
\]
que a todo sistema inductivo de homomorfismos $\theta_\lambda:G\to F_\lambda$ de $E$-álgebras asocia su límite inductivo, es inyectivo (resp.\ biyectivo).
\end{lemma}

\begin{proof}
Supongamos primero que la $E$-álgebra $G$ es de tipo finito, y sea $(t_i)_{1\leqslant i\leqslant n}$ un sistema de generadores de esta $E$-álgebra; demostremos que, si $(\theta'_\lambda)$, $(\theta''_\lambda)$ son dos sistemas inductivos de homomorfismos $G\to F_\lambda$, con $\varinjlim\theta'_\lambda=\varinjlim\theta''_\lambda$, existe $\mu\geqslant\lambda$ tal que $\theta'_\mu=\theta''_\mu$.
En efecto, si $\varphi_{\mu\lambda}:F_\lambda\to F_\mu$ y $\varphi_\lambda:F_\lambda\to F=\varinjlim F_\lambda$ son los homomorfismos del sistema inductivo $(F_\lambda)$, por hipótesis, para cada índice $i$, existe un índice $\lambda_i$ tal que $\varphi_{\lambda_i}(\theta'_{\lambda_i}(t_i))=\varphi_{\lambda_i}(\theta''_{\lambda_i}(t_i))$, y se puede suponer que todos los $\lambda_i$ son iguales a un único $\lambda$; se deduce del mismo modo la existencia de $\mu\geqslant\lambda$ tal que $\varphi_{\mu\lambda}(\theta'_\lambda(t_i))=\varphi_{\mu\lambda}(\theta''_\lambda(t_i))$ para $1\leqslant i\leqslant n$, es decir, $\theta'_\mu(t_i)=\theta''_\mu(t_i)$ para $1\leqslant i\leqslant n$, de donde $\theta'_\mu=\theta''_\mu$.

Supongamos, en segundo lugar, que $G$ es de presentación finita, de modo que se tiene
\[
  G=E[T_1,\ldots,T_n]/\mathfrak{J},
\]
donde $\mathfrak{J}$ es un ideal de tipo finito, siendo $t_i$ la clase de $T_i$ módulo $\mathfrak{J}$.
Sea $(P_j)_{1\leqslant j\leqslant m}$ un sistema de generadores de $\mathfrak{J}$.
Sea dado un homomorfismo de $E$-álgebras $\theta:G\to F$; pongamos $b_i=\theta(t_i)$; por definición, se tiene entonces $P_j(b_1,b_2,\ldots,b_n)=\theta(P_j(t_1,\ldots,t_n))=0$ para $1\leqslant j\leqslant m$.
Ahora bien, existen $\lambda$ y elementos $x_1,\ldots,x_n$ de $F_\lambda$ tales que $b_i=\varphi_\lambda(x_i)$ para $1\leqslant i\leqslant n$; se tiene, pues, $\varphi_\lambda(P_j(x_1,\ldots,x_n))=P_j(b_1,\ldots,b_n)=0$, y en consecuencia existe $\mu\geqslant\lambda$ tal que $\varphi_{\mu\lambda}(P_j(x_1,\ldots,x_n))=P_j(\varphi_{\mu\lambda}(x_1),\ldots,\varphi_{\mu\lambda}(x_n))=0$ para $1\leqslant j\leqslant m$; se concluye que existe un homomorfismo de $E$-álgebras $\theta_\mu:G\to F_\mu$ tal que $\theta_\mu(t_i)=\varphi_{\mu\lambda}(x_i)$
\oldpage[IV-3]{30}
para $1\leqslant i\leqslant n$; se deduce, por tanto, para todo $\nu\geqslant\mu$, un homomorfismo de $E$-álgebras $\theta_\nu=\varphi_{\nu\mu}\circ\theta_\mu$ de $G$ en $F_\nu$, y es claro que $\theta$ es el límite inductivo de este sistema inductivo de homomorfismos, lo que concluye la demostración del lema.
\end{proof}

Pasemos ahora al caso en que $X_\alpha$ es cuasi-compacto y $Y_\alpha$ localmente de tipo finito sobre $S_\alpha$.
Pongamos $Z_\alpha=X_\alpha\times_{S_\alpha}Y_\alpha$ e introduzcamos el sistema proyectivo correspondiente a $Z_\lambda=X_\lambda\times_{S_\lambda}Y_\lambda$ y su límite $Z=X\times_S Y$; las biyecciones canónicas \sref[I]{I.3.3.14} dan diagramas conmutativos
\[
\xymatrix{
  \operatorname{Hom}_{S_\lambda}(X_\lambda,Y_\lambda) \ar[r]^-{e_{\mu\lambda}} \ar[d]^{\wr} & \operatorname{Hom}_{S_\mu}(X_\mu,Y_\mu) \ar[r]^-{e_\mu} \ar[d]^{\wr} & \operatorname{Hom}_S(X,Y) \ar[d]^{\wr}\\
  \operatorname{Hom}_{X_\lambda}(X_\lambda,Z_\lambda) \ar[r]_-{e_{\mu\lambda}} & \operatorname{Hom}_{X_\mu}(X_\mu,Z_\mu) \ar[r]_-{e_\mu} & \operatorname{Hom}_X(X,Z)
}
\]
y, en consecuencia, nos reducimos a demostrar que \eqref{IV.8.8.1.1} es inyectiva en el caso particular en que $S_\alpha=X_\alpha$ (teniendo en cuenta \sref[I]{I.1.3.4}).
Además, como $X_\alpha$ es cuasi-compacto, y por tanto unión finita de abiertos afines, podemos suponer $X_\alpha$ afín (puesto que $L$ es filtrado). Sean, pues, dos $X_\alpha$-morfismos $f'_\alpha:X_\alpha\to Y_\alpha$, $f''_\alpha:X_\alpha\to Y_\alpha$ tales que $f'$ y $f''$ son $X$-morfismos iguales de $X$ en $Y$; debemos demostrar que $f'_\mu=f''_\mu$ para $\mu\geqslant\alpha$ suficientemente grande. Como $X_\alpha$ es cuasi-compacto, también lo es $f'_\alpha(X_\alpha)\cup f''_\alpha(X_\alpha)$ y, como $Y_\alpha$ es localmente de tipo finito sobre $X_\alpha$, esta unión está contenida en una unión finita de abiertos afines $V_{i\alpha}$ de $Y_\alpha$, de tipo finito sobre $X_\alpha$. Pongamos $U'_{i\alpha}=f_\alpha^{'-1}(V_{i\alpha})$, $U''_{i\alpha}=f_\alpha^{''-1}(V_{i\alpha})$, $U_{i\alpha}=U'_{i\alpha}\cap U''_{i\alpha}$ y $U_\alpha=\bigcup_iU_{i\alpha}$. La hipótesis $f'=f''$ implica $v_\alpha^{-1}(U'_{i\alpha})=v_\alpha^{-1}(U''_{i\alpha})$, pues estos dos conjuntos son respectivamente iguales a $f^{'-1}(w_\alpha^{-1}(V_{i\alpha}))$ y $f^{''-1}(w_\alpha^{-1}(V_{i\alpha}))$. Como los $V_{i\alpha}$ recubren $Y_\alpha$, se tiene $v_\alpha^{-1}(U_\alpha)=f^{'-1}(Y)=X$. Puesto que $X_\alpha$ es cuasi-compacto y todo abierto de $X_\lambda$ es indconstructible, se deduce de \sref{IV.8.3.4} que existe un índice $\mu\geqslant\alpha$ tal que los $U_{i\mu}=v_{\alpha\mu}^{-1}(U_{i\alpha})$ forman un recubrimiento de $X_\mu$. Reemplazando $\alpha$ por $\mu$, podemos suponer, por tanto, que los $U_{i\alpha}$ recubren $X_\alpha$; esto implica que, para todo $x\in X_\alpha$, existe un entorno abierto afín $W(x)$ contenido en uno de los $U_{i\alpha}$, es decir, tal que las restricciones de $f'_\alpha$ y $f''_\alpha$ a $W(x)$ envían $W(x)$ al mismo abierto afín $V_{i\alpha}$. Como $X_\alpha$ es cuasi-compacto, está recubierto por un número finito de abiertos afines $W(x_k)$; en virtud del lema~\sref{IV.8.8.2.1} y de que $L$ es filtrado, existe $\lambda\geqslant\alpha$ tal que las restricciones de $f'_\lambda$ y $f''_\lambda$ a cada uno de los abiertos $v_{\alpha\lambda}^{-1}(W(x_k))$ son iguales, de donde $f'_\lambda=f''_\lambda$.

Supongamos ahora que $X_\alpha$ es cuasi-compacto y cuasi-separado y que $Y_\alpha$ es localmente de presentación finita sobre $S_\alpha$, y demostremos que \eqref{IV.8.8.1.1} es sobreyectiva.
Sea, pues, un $X$-morfismo $f:X\to Y$.
Como $X$ es cuasi-compacto, también lo es
\oldpage[IV-3]{31}
$f(X)$ y, en consecuencia, existe un abierto cuasi-compacto $Y'$ de $Y$ que contiene a $f(X)$; existen entonces $\lambda\geqslant\alpha$ y un abierto cuasi-compacto $Y'_\lambda$ de $Y_\lambda$ tales que $Y'=w_\lambda^{-1}(Y'_\lambda)$ \sref{IV.8.2.11}.
Reemplazando, si es necesario, $\alpha$ por $\lambda$ e $Y_\lambda$ por $Y'_\lambda$, podemos, pues, restringirnos al caso en que $Y_\alpha$ es \emph{cuasi-compacto}, y por tanto también $Y$ y los $Y_\lambda$.
Como $Y$ es cuasi-compacto, es unión finita de abiertos afines $V_i$, y en consecuencia $X$ es la unión de los abiertos $f^{-1}(V_i)$. Puesto que todo punto de $X$ posee un entorno abierto cuasi-compacto contenido en uno de los $f^{-1}(V_i)$ y $X$ es cuasi-compacto, podemos, teniendo en cuenta \sref{IV.8.2.11} y reemplazando $\alpha$, si es necesario, por un índice $\lambda\geqslant\alpha$, suponer que $Y$ es unión finita de abiertos $V_i=w_\alpha^{-1}(V_{i\alpha})$, donde los $V_{i\alpha}$ son abiertos afines de $Y_\alpha$; en consecuencia, $X$ es la unión de los abiertos $f^{-1}(V_i)$. Como todo punto de $X$ posee un entorno abierto cuasi-compacto contenido en uno de los $f^{-1}(V_i)$ y $X$ es cuasi-compacto, podemos recubrir $X$ por un número finito de tales entornos y, repitiendo algunos de los $V_i$ si es necesario, suponer que existe un recubrimiento $(U_i)$ de $X$ por abiertos cuasi-compactos, con el mismo conjunto de índices que $(V_i)$ y tal que $f(U_i)\subset V_i$ para todo $i$. Además, utilizando \sref{IV.8.2.11} (y reemplazando $\alpha$, si es necesario, por un índice $\lambda\geqslant\alpha$), podemos suponer que $U_i=v_\alpha^{-1}(U_{i\alpha})$, donde los $U_{i\alpha}$ son abiertos cuasi-compactos de $X_\alpha$; utilizando asimismo \sref{IV.8.3.4} como antes, podemos suponer que $(U_{i\alpha})$ es un \emph{recubrimiento} de $X_\alpha$.

Bastará entonces demostrar que existen $\lambda\geqslant\alpha$ y, para cada $i$, un morfismo $f_{i\lambda}:U_{i\lambda}\to V_{i\lambda}$ (donde $U_{i\lambda}=v_{\alpha\lambda}^{-1}(U_{i\alpha})$ y $V_{i\lambda}=w_{\alpha\lambda}^{-1}(V_{i\alpha})$) tal que el morfismo correspondiente $f_i=e_\lambda(f_{i\lambda}):U_i\to V_i$ sea igual a la \emph{restricción} de $f$ a $U_i$.
En efecto, si esto es así, como $X_\lambda$ es cuasi-separado \sref[I]{I.1.2.3}, los $U_{i\lambda}\cap U_{j\lambda}$ son \emph{cuasi-compactos}, y el resultado de unicidad demostrado anteriormente (aplicable puesto que $Y_\lambda$ es localmente de tipo finito sobre $S_\lambda$) prueba que existe $\mu\geqslant\lambda$ tal que $f_{i\mu}$ y $f_{j\mu}$ coinciden en $U_{i\mu}\cap U_{j\mu}$ para todo par $(i,j)$ y definen, en consecuencia, un $S_\mu$-morfismo $f_\mu:X_\mu\to Y_\mu$ tal que $e_\mu(f_\mu)=f$.

Nos reducimos así al caso en que $Y_\alpha$ es \emph{afín} y, como además podemos suponer que las imágenes de los $V_{i\alpha}$ en $S_\alpha$ están contenidas en abiertos afines, también podemos suponer que $S_\alpha$ es afín; sean, pues, $S_\alpha=\operatorname{Spec}(A_\alpha)$, $Y_\alpha=\operatorname{Spec}(C_\alpha)$, con $C_\alpha$ una $A_\alpha$-álgebra de presentación finita, y sean $S=\operatorname{Spec}(A)$, $Y=\operatorname{Spec}(C)$, donde $A=\varinjlim A_\lambda$ y $C=\varinjlim(C_\alpha\otimes_{A_\alpha}A_\lambda)=C_\alpha\otimes_{A_\alpha}A$.
Se tiene entonces
\[
  \operatorname{Hom}_S(X,Y)=\operatorname{Hom}_{A\text{-alg.}}(C,\Gamma(X,\sh{O}_X))=\operatorname{Hom}_{A_\alpha\text{-alg.}}(C_\alpha,\Gamma(X,\sh{O}_X))
\]
\sref[I]{I.2.2.4} y, análogamente,
\[
  \operatorname{Hom}_{S_\lambda}(X_\lambda,Y_\lambda)=\operatorname{Hom}_{A_\lambda\text{-alg.}}(C_\alpha\otimes_{A_\alpha}A_\lambda,\Gamma(X_\lambda,\sh{O}_{X_\lambda}))=\operatorname{Hom}_{A_\alpha\text{-alg.}}(C_\alpha,\Gamma(X_\lambda,\sh{O}_{X_\lambda})).
\]
Pero, como $X_\alpha$ es cuasi-compacto y cuasi-separado, se sabe por \sref{IV.8.5.4} que $\varinjlim\Gamma(X_\lambda,\sh{O}_{X_\lambda})=\Gamma(X,\sh{O}_X)$; puesto que $C_\alpha$ es una $A_\alpha$-álgebra de presentación finita, la biyectividad de \eqref{IV.8.8.1.1} se deduce entonces de \sref{IV.8.8.2.3}.

Antes de pasar a la demostración de (ii), observemos los siguientes corolarios de (i):
\end{proof}

\begin{corollary}[8.8.2.4]
\label{IV.8.8.2.4}
Supongamos que $S_0$ es cuasi-compacto, que $X_\alpha$ es de presentación finita sobre $S_\alpha$, y que $Y_\alpha$ es de tipo finito sobre $S_\alpha$ y cuasi-separado sobre $S_\alpha$ (lo que ocurre, por ejemplo, si $Y_\alpha$ es también
\oldpage[IV-3]{32}
de presentación finita sobre $S_\alpha$).
Sea $f_\alpha:X_\alpha\to Y_\alpha$ un $S_\alpha$-morfismo.
Para que $f:X\to Y$ sea un isomorfismo, es necesario y suficiente que exista $\lambda\geqslant\alpha$ tal que $f_\lambda:X_\lambda\to Y_\lambda$ sea un isomorfismo.
\end{corollary}

\begin{proof}
La condición es evidentemente suficiente. Para demostrar que es necesaria, observemos primero que, como la cuestión es local en $S_0$ (puesto que $L$ es filtrado), siempre podemos suponer $S_0$ afín y, por tanto, cuasi-separado. Por hipótesis existe un $S$-morfismo $g:Y\to X$ tal que $g\circ f=1_X$ y $f\circ g=1_Y$. Como $X_\alpha$ es de presentación finita sobre $S_\alpha$, e $Y_\alpha$ es cuasi-compacto y cuasi-separado (porque $S_\alpha$ es cuasi-compacto y cuasi-separado), se deduce de \sref{IV.8.8.2},~(i) que existen $\mu\geqslant\alpha$ y un $S_\mu$-morfismo $g_\mu:Y_\mu\to X_\mu$ tales que $g=g_\mu\times 1_S$. Se deduce además de \sref{IV.8.8.2},~(i), junto con las relaciones $g\circ f=1_X$ y $f\circ g=1_Y$, que existe $\nu\geqslant\mu$ tal que $g_\nu\circ f_\nu=1_{X_\nu}$ y $f_\nu\circ g_\nu=1_{Y_\nu}$, puesto que $X_\alpha$ e $Y_\alpha$ son de tipo finito sobre $S_\alpha$ y cuasi-compactos. Así, $f_\nu:X_\nu\to Y_\nu$ es un isomorfismo, lo que demuestra el corolario.
\end{proof}

\begin{corollary}[8.8.2.5]
\label{IV.8.8.2.5}
Supongamos que $S_0$ es cuasi-compacto y cuasi-separado, y que $X_\alpha$ e $Y_\alpha$ son de presentación finita sobre $S_\alpha$.
Para que $X$ e $Y$ sean $S$-isomorfos, es necesario y suficiente que exista $\lambda\geqslant\alpha$ tal que $X_\lambda$ e $Y_\lambda$ sean $S_\lambda$-isomorfos. Además, para todo $S$-isomorfismo $f:X\xrightarrow{\sim}Y$, existen $\mu\geqslant\lambda$ y un $S_\mu$-isomorfismo $f_\mu:X_\mu\xrightarrow{\sim}Y_\mu$ tales que $f=f_\mu\times 1_S$.
\end{corollary}

\begin{proof}
La condición es evidentemente suficiente. Recíprocamente, si existe un $S$-isomorfismo $f:X\to Y$, se deduce de \sref{IV.8.8.2},~(i) que $f$ es de la forma $f_\mu\times 1_S$ para algún $\mu\geqslant\alpha$ y algún $S_\mu$-morfismo $f_\mu:X_\mu\to Y_\mu$. Como $f$ es un isomorfismo, se deduce de \sref{IV.8.8.2.4} que existe $\nu\geqslant\mu$ tal que $f_\nu:X_\nu\to Y_\nu$ es un isomorfismo.
\end{proof}

\begin{proof}[Demostración de \sref{IV.8.8.2}, (ii)]
Consideremos primero el caso en que $S_0=\operatorname{Spec}(A_0)$ y $X=\operatorname{Spec}(B)$ son afines; entonces la afirmación~(ii) es equivalente al lema~\sref{IV.1.8.4.2}.

Para demostrar (ii) en el caso general, observemos que $S$ es cuasi-compacto y cuasi-separado. Como $X$ es de presentación finita sobre $S$, y $S$ es afín sobre $S_0$, existe un recubrimiento abierto finito $(U_i)$ de $S_0$ y, si $(W_i)$ denota el recubrimiento abierto afín de $S$ formado por las imágenes inversas de los $U_i$ por $S\to S_0$, un recubrimiento finito $(X_r)$ de $X$ por abiertos afines tal que la imagen de cada $X_r$ por $X\to S$ está contenida en algún $W_{i(r)}$. El anillo $A(X_r)$ es entonces un álgebra de presentación finita sobre $A(W_{i(r)})$ \sref{IV.1.4.6}.

Por el lema~\sref{IV.1.8.4.2} y el hecho de que $L$ es filtrado, existen un índice $\lambda\in L$ y, para todo $r$, un esquema afín $Z_{r\lambda}$ de presentación finita sobre la imagen inversa $W_{i(r),\lambda}$ de $U_{i(r)}$ en $S_\lambda$, junto con un $S$-isomorfismo
\[
 g_r:Z_{r\lambda}\times_{S_\lambda}S\xrightarrow{\sim}X_r.
\]
Sea $Z_{rs}$ la imagen inversa por $g_r$ del subpreesquema inducido en el abierto $X_r\cap X_s$ de $X_r$, que es cuasi-compacto porque $X$ es cuasi-separado, y sea
\[
 g'_{rs}:Z_{rs}\xrightarrow{\sim}X_r\cap X_s
\]
la restricción de $g_r$. Por \sref{IV.8.2.11}, existen $\mu\geqslant\lambda$ y, para todo par $(r,s)$, un abierto cuasi-compacto $Z_{rs\mu}$ de $Z_{r\mu}=v_{\lambda\mu}^{-1}(Z_{r\lambda})$ tal que $Z_{rs}$ es la imagen inversa de $Z_{rs\mu}$. Además, como $S_\mu$ es cuasi-separado y $W_{i(r),\mu}$ es un abierto cuasi-compacto de $S_\mu$, cada $Z_{rs\mu}$ es de presentación finita sobre $S_\mu$ \sref[I]{I.1.6.2}.

Para todo par $(r,s)$, consideremos el isomorfismo
\[
 h_{sr}=(g'_{sr})^{-1}\circ g'_{rs}:Z_{rs}\xrightarrow{\sim}Z_{sr}.
\]
Por \sref{IV.8.8.2.4}, existen $\nu\geqslant\mu$ y, para todo par $(r,s)$, un isomorfismo
\[
 h_{sr\nu}:Z_{rs\nu}\xrightarrow{\sim}Z_{sr\nu}
\]
tal que $h_{sr}=h_{sr\nu}\times 1_S$. Finalmente, para toda terna $(r,s,t)$, denotemos por $h'_{sr}$ la restricción de $h_{sr}$ a

\oldpage[IV-3]{33}

$Z_{rs}\cap Z_{rt}$. Se deduce inmediatamente de las definiciones que $h'_{sr}$ es un isomorfismo de $Z_{rs}\cap Z_{rt}$ sobre $Z_{sr}\cap Z_{st}$ y que
\[
 h'_{ts}\circ h'_{sr}=h'_{tr}.
\]
Por \sref{IV.8.8.2},~(i), existe entonces $\rho\geqslant\nu$ tal que, para toda terna $(r,s,t)$,
\[
 h'_{ts\rho}\circ h'_{sr\rho}=h'_{tr\rho}.
\]
En consecuencia, los isomorfismos $h_{sr\rho}$ satisfacen la condición de pegado \sref[0\textsubscript{I}]{0.4.1.7} y definen, por tanto, un preesquema $X_\rho$ tal que $X$ es isomorfo a $X_\rho\times_{S_\rho}S$. Además, los $Z_{r\rho}$ son de presentación finita sobre $S_\rho$ y, al identificarlos con abiertos de $X_\rho$, sus intersecciones $Z_{r\rho}\cap Z_{s\rho}$, que son isomorfas a los $Z_{rs\rho}$, son cuasi-compactas. Así, $X_\rho$ es de presentación finita sobre $S_\rho$ \sref[I]{I.1.6.3}.
\end{proof}

\begin{env}[8.8.3]
\label{IV.8.8.3}
El contenido esencial de \sref{IV.8.8.2} se expresa también diciendo que, si $S_0$ es cuasi-compacto y cuasi-separado, la categoría de los $S$-preesquemas de presentación finita queda determinada, salvo equivalencia, por los datos de las categorías de los $S_\lambda$-preesquemas de presentación finita, los funtores $\rho_{\mu\lambda}:X_\lambda\mapsto X_\lambda\times_{S_\lambda}S_\mu$ (para $\lambda\leqslant\mu$) entre estas categorías y los isomorfismos de transitividad $\rho_{\nu\mu}\circ\rho_{\mu\lambda}\xrightarrow{\sim}\rho_{\nu\lambda}$.
De manera figurada, puede decirse que el dato de un $S$-preesquema de presentación finita $X$ es ``funtorialmente'' equivalente al dato de un $S_\lambda$-preesquema de presentación finita $X'_\lambda$; si un $S_\mu$-preesquema de presentación finita $X''_\mu$ es tal que $X$ es $S$-isomorfo a $X''_\mu\times_{S_\mu}S$, existe $\nu$ tal que $\lambda\leqslant\nu$, $\mu\leqslant\nu$, y $X'_\lambda\times_{S_\lambda}S_\nu$ y $X''_\mu\times_{S_\mu}S_\nu$ son $S_\nu$-isomorfos.
El hecho de que $L$ sea filtrado implica además que, si se da una \emph{familia} finita $(X^{(i)})_{i\in I}$ de $S$-preesquemas de presentación finita y una familia finita $(f^{(j)})_{j\in J}$ de $S$-morfismos entre estos preesquemas (siendo, por tanto, $f^{(j)}$ de la forma $X^{(\sigma(j))}\to X^{(\tau(j))}$, donde $\sigma$ y $\tau$ son dos aplicaciones de $J$ en $I$), entonces existen un índice $\lambda\in L$, una familia $(X^{(i)}_\lambda)_{i\in I}$ de $S_\lambda$-preesquemas de presentación finita y una familia $(f^{(j)}_\lambda)_{j\in J}$ de $S_\lambda$-morfismos tales que $X^{(i)}$ se identifica con $X^{(i)}_\lambda\times_{S_\lambda}S$ y $f^{(j)}$ con $f^{(j)}_\lambda\times 1_S$ para todos $i$ y $j$; además, si se tiene una relación $f^{(j)}=f^{(k)}$, se puede suponer $\lambda$ elegido de modo que $f^{(j)}_\lambda=f^{(k)}_\lambda$.

Consideremos, en particular, tres $S$-preesquemas de presentación finita $X$, $Y$, $Z$ y dos $S$-morfismos $f:X\to Z$, $g:Y\to Z$, de modo que el producto $X\times_Z Y$ relativo a estos morfismos es también un $S$-preesquema de presentación finita \sref[I]{I.1.6.2}.
Se deduce entonces de lo anterior y de \sref[I]{I.3.3.11} que se puede escribir $X\times_Z Y=(X_\lambda\times_{Z_\lambda}Y_\lambda)\times_{S_\lambda}S$ para un $\lambda$ adecuado, siendo $X_\lambda$, $Y_\lambda$, $Z_\lambda$ $S_\lambda$-preesquemas de presentación finita; puede decirse, por tanto, que la determinación de los $S$-preesquemas de presentación finita mediante los datos de los $S_\lambda$-preesquemas de presentación finita es ``\emph{compatible con los productos fibrados}''.
Por otro lado, hemos visto \sref{IV.8.6.3} que, si $g$ es una inmersión, se puede suponer lo mismo para $g_\lambda$; identificando $Y$ (resp.\ $Y_\lambda$) con un subpreesquema de $Z$ (resp.\ $Z_\lambda$), se ve que se puede escribir, para un $\lambda$ adecuado, $f^{-1}(Y)=f_\lambda^{-1}(Y_\lambda)\times_{S_\lambda}S$ \sref[I]{I.4.4.1};
existe, pues, también ``\emph{compatibilidad con la formación de imágenes inversas de subpreesquemas}''.
Más particularmente, si $f_1$, $f_2$ son dos $S$-morfismos de $X$ en $Y$, se llama \emph{núcleo} de este par de morfismos a la imagen inversa $N$ del subpreesquema diagonal de $Y\times_S Y$ por el $S$-morfismo $(f_1,f_2)_S$; se deduce de lo anterior que se tendrá, para un $\lambda$ adecuado, $N=N_\lambda\times_{S_\lambda}S$, donde $N_\lambda$ es el núcleo del par de $S_\lambda$-morfismos $(f_{1\lambda},f_{2\lambda})$ de $X_\lambda$ en $Y_\lambda$.
Estas observaciones se extienden a los productos finitos y a los ``núcleos'' de sistemas finitos arbitrarios de $S$-morfismos de un $S$-preesquema de presentación finita en
\oldpage[IV-3]{34}
otro; se concluye que, de manera general, la formación de $S$-preesquemas de presentación finita mediante los datos de $S_\lambda$-preesquemas de presentación finita es ``\emph{compatible con los límites proyectivos finitos}'', siendo tal límite, por definición, el núcleo de un sistema finito de morfismos de un $S$-preesquema en un producto de $S$-preesquemas \sref[T]{T.1.8}.

En lo que sigue nos permitiremos efectuar libremente las traslaciones implicadas por las propiedades anteriores (así como por \sref{IV.8.3.11}, \sref{IV.8.5.2} y \sref{IV.8.6.3}), sin esforzarnos siempre por justificarlas como se ha hecho anteriormente.
Por ejemplo, el dato de un ``preesquema en grupos'' $G$ de presentación finita sobre $S$ equivale al dato de un preesquema en grupos $G_\lambda$ de presentación finita sobre algún $S_\lambda$ para un $\lambda$ suficientemente grande; en efecto, escribir la condición de asociatividad para el $S$-morfismo ``ley de composición'' $G\times_S G\to G$ del preesquema en grupos equivale a escribir que el núcleo de dos $S$-morfismos de la forma $G\times_S G\times_S G\to G$ es igual a $G\times_S G\times_S G$ \sref[II]{II.8.3.9}, y las demás condiciones que intervienen en la definición de un preesquema en grupos se interpretan de manera análoga.
\end{env}

\subsection{Primeras aplicaciones a la eliminación de las hipótesis noetherianas}
\label{subsection:IV.8.9}

\begin{proposition}[8.9.1]
\label{IV.8.9.1}
Sea $A$ un anillo y sea $X$ un $A$-preesquema.
\begin{enumerate}[label=(\roman*)]
\item Las condiciones siguientes son equivalentes:
\begin{enumerate}[label=\emph{\alph*)}]
\item $X$ es de presentación finita sobre $A$.
\item Existen un anillo noetheriano $A_0$, un preesquema $X_0$ de tipo finito sobre $A_0$, un homomorfismo de anillos $A_0\to A$ y un $A$-isomorfismo $X_0\otimes_{A_0}A\xrightarrow{\sim}X$.
\item Existen un subanillo $A_0$ de $A$, que es una $\mathbf{Z}$-álgebra de tipo finito, un preesquema $X_0$ de tipo finito sobre $A_0$ y un $A$-isomorfismo $X_0\otimes_{A_0}A\xrightarrow{\sim}X$.
\end{enumerate}
\item Si $\sh{F}$ es un $\sh{O}_X$-Módulo cuasicoherente de presentación finita, se puede suponer elegido el subanillo $A_0$ de modo que exista un $\sh{O}_{X_0}$-Módulo coherente $\sh{F}_0$ tal que $\sh{F}$ sea isomorfo a $\sh{F}_0\otimes_{A_0}A$; $\operatorname{Supp}(\sh{F})$ es construible y cerrado en $X$, y existe un subpreesquema cerrado $Z$ de $X$, de espacio subyacente $\operatorname{Supp}(\sh{F})$, tal que la inmersión canónica $Z\to X$ sea de presentación finita.
\item Si $Y$ es otro $A$-preesquema de presentación finita y $f:X\to Y$ un $A$-morfismo, se puede suponer elegido el subanillo $A_0$ de $A$ de modo que existan un preesquema $Y_0$ de tipo finito sobre $A_0$, un $A$-isomorfismo $Y_0\otimes_{A_0}A\xrightarrow{\sim}Y$ y un $A_0$-morfismo $f_0:X_0\to Y_0$ (necesariamente de tipo finito) tal que $f$ se identifique con $f_0\otimes 1_A$.
\end{enumerate}
\end{proposition}

\begin{proof}
\begin{enumerate}[label=(\roman*)]
\item Como $A$ es el límite inductivo de sus subanillos de tipo finito sobre $\mathbf{Z}$, \emph{a)} entraña \emph{c)} en virtud de \sref{IV.8.8.2},~(ii); \emph{c)} entraña \emph{b)}, pues una $\mathbf{Z}$-álgebra de tipo finito es un anillo noetheriano; por último, si $A_0$ es noetheriano, todo $A_0$-preesquema de tipo finito es de presentación finita \sref[I]{I.1.6.1}, de modo que \emph{b)} entraña \emph{a)} en virtud de \sref[I]{I.1.6.2},~(iii).
\end{enumerate}
La aserción~(ii) se sigue inmediatamente de \sref{IV.8.5.2},~(ii), \sref{IV.8.3.11} y \sref[I]{I.1.8.2}, y la aserción~(iii), de \sref{IV.8.8.2},~(i).
\end{proof}

\begin{env}[8.9.2]
\label{IV.8.9.2}
La proposición \sref{IV.8.9.1} y los resultados de \sref{IV.8.5.2} y \sref{IV.8.8.2} permiten extender en numerosos casos a los morfismos de presentación finita $X\to Y$ los resultados demostrados mediante técnicas noetherianas suponiendo $Y$ \emph{localmente noetheriano}.
\oldpage[IV-3]{35}
Veremos numerosos ejemplos en lo que sigue y nos limitaremos aquí a dar algunos resultados de este tipo.
\end{env}

\begin{proposition}[8.9.3]
\label{IV.8.9.3}
Sea $A$ un anillo y sea $M$ un $A$-módulo de presentación finita.
Entonces todo $A$-endomorfismo sobreyectivo de $M$ es biyectivo.
\end{proposition}

\begin{proof}
En efecto, consideremos $A$ como límite inductivo de sus sub-$\mathbf{Z}$-álgebras de tipo finito.
De \sref{IV.8.5.2},~(ii) se sigue que existen una de estas subálgebras $A_0$ y un $A_0$-módulo $M_0$ de presentación finita tales que $M$ es $A$-isomorfo a $M_0\otimes_{A_0}A$; además, si $f:M\to M$ es un $A$-endomorfismo sobreyectivo, se puede suponer \sref{IV.8.5.2},~(i) que existe un $A_0$-endomorfismo $f_0:M_0\to M_0$ tal que $f=f_0\otimes 1_A$; por último, en virtud de \sref{IV.8.5.7}, se puede suponer que $f_0$ es \emph{sobreyectivo}.
Pero, como $A_0$ es noetheriano y $M_0$ un $A_0$-módulo de tipo finito, $M_0$ es un módulo noetheriano; por tanto (Bourbaki, \emph{Alg.}, cap.~VIII, \textsection\,2, n\textsuperscript{o}\,2, lema~3), $f_0$ es biyectivo, y lo mismo ocurre con $f$ por cambio de base.
\end{proof}

\begin{proposition}[8.9.4]
\label{IV.8.9.4}
(``teorema de planitud genérica''). --- Sea $Y$ un preesquema íntegro, sea $u:X\to Y$ un morfismo de tipo finito y localmente de presentación finita, y sea $\sh{F}$ un $\sh{O}_X$-Módulo cuasicoherente de presentación finita.
Entonces existe un abierto no vacío $U$ de $Y$ tal que $\sh{F}|u^{-1}(U)$ es plano sobre $U$.
\end{proposition}

\begin{proof}
El razonamiento del comienzo de \sref{IV.6.9.1} conduce a demostrar el

\begin{lemma}[8.9.4.1]
\label{IV.8.9.4.1}
Sea $A$ un anillo íntegro, sea $B$ una $A$-álgebra de presentación finita y sea $M$ un $B$-módulo de presentación finita.
Entonces existe $f\neq 0$ en $A$ tal que $M_f$ es un $A_f$-módulo libre.
\end{lemma}

\begin{proof}
En efecto, por \sref{IV.8.9.1} existen una $\mathbf{Z}$-subálgebra $A_0$ de $A$ de tipo finito, una $A_0$-álgebra $B_0$ de tipo finito y un $B_0$-módulo $M_0$ de tipo finito tales que $B$ es $A$-isomorfo a $B_0\otimes_{A_0}A$ y $M$ es $B$-isomorfo a $M_0\otimes_{A_0}A$; por \sref{IV.6.9.2}, existe $f_0\neq 0$ en $A_0$ tal que $(M_0)_{f_0}$ es un $(A_0)_{f_0}$-módulo libre.
Como $M_{f_0}=(M_0)_{f_0}\otimes_{A_0}A$ y $A_{f_0}=(A_0)_{f_0}\otimes_{A_0}A$, $M_{f_0}$ es un $A_{f_0}$-módulo libre.
\end{proof}
\end{proof}

\begin{corollary}[8.9.5]
\label{IV.8.9.5}
Sea $S$ un preesquema cuasicompacto y cuasiseparado, sea $u:X\to S$ un morfismo de presentación finita y sea $\sh{F}$ un $\sh{O}_X$-Módulo cuasicoherente de presentación finita.
Entonces existe una partición $(S_i)_{1\leqslant i\leqslant n}$ de $S$ en un número finito de subconjuntos localmente cerrados de $S$ tal que, para $1\leqslant i\leqslant n$, existe un subpreesquema $S'_i$ de $S$, de espacio subyacente $S_i$, que es de presentación finita sobre $S$, y tal que, poniendo $X_i=X\times_S S'_i$, el $\sh{O}_{X_i}$-Módulo $\sh{F}_i=\sh{F}\otimes_{\sh{O}_X}\sh{O}_{X_i}$ es plano sobre $S'_i$.
\end{corollary}

\begin{proof}
Consideremos un recubrimiento finito $(U_j)_{1\leqslant j\leqslant m}$ de $S$ por abiertos afines y definamos, por inducción, $T_1=U_1$ y
\[
T_j=U_j\setminus\bigcup_{h<j}(U_j\cap U_h)\qquad(j\geqslant 2).
\]
Cada $T_j$ es cerrado en el abierto afín $U_j$, y los $T_j$ son disjuntos dos a dos.
Además, los $U_j\cap U_h$ son cuasicompactos porque $S$ es cuasiseparado y, por tanto, como $S$ es también cuasicompacto, son construibles en $S$ \sref{IV.1.8.1}; en consecuencia, cada $T_j$ es asimismo construible.
Bastará, pues, demostrar el corolario para un subpreesquema adecuado $T'_j$ de $S$, de espacio subyacente $T_j$ y de presentación finita sobre $S$, y para el morfismo y el Módulo deducidos respectivamente de $u$ y de $\sh{F}$ por el cambio de base $T'_j\to S$.
Para ello, observemos que, si se dota a $U_j$ de la estructura de preesquema inducida por la de $S$, la inmersión abierta $U_j\to S$ es cuasicompacta porque $S$ es cuasiseparado \sref{IV.1.2.7}, y por tanto es de
\oldpage[IV-3]{36}
presentación finita \sref{IV.1.6.2},~(i).
Basta, pues, con que $T'_j$ sea de presentación finita sobre $U_j$; dicho de otro modo, podemos reducirnos ya al caso en que $U_j=S$ y $T_j=T$ es cerrado y construible en $S$.

Escribamos $S=\operatorname{Spec}(A)$ y consideremos $A$ como límite inductivo de sus $\mathbf{Z}$-subálgebras de tipo finito, de modo que $S=\varprojlim S_\lambda$, donde los $S_\lambda$ son afines y noetherianos.
Por \sref{IV.8.3.11}, existen $\lambda$ y un subconjunto cerrado (necesariamente construible) $T_\lambda$ de $S_\lambda$ tales que $T=u_\lambda^{-1}(T_\lambda)$, donde $u_\lambda:S\to S_\lambda$ es el morfismo canónico.
Dotemos a $T_\lambda$ de una estructura de subpreesquema y tomemos $T'=T_\lambda\times_{S_\lambda}S$.
Como la inmersión canónica $T_\lambda\to S_\lambda$ es de presentación finita \sref{IV.1.6.1}, también lo es la inmersión $T'\to S$.
Puesto que $T'$ es afín, vemos finalmente que basta tratar el caso en que $T'=S$ es afín.

Entonces, con la notación de \sref{IV.8.9.1}, existen $\lambda$, un morfismo de tipo finito $u_\lambda:X_\lambda\to S_\lambda$ y un $\sh{O}_{X_\lambda}$-Módulo coherente $\sh{F}_\lambda$ tales que $X$ es isomorfo a $X_\lambda\times_{S_\lambda}S$ y $\sh{F}$ a $\sh{F}_\lambda\otimes_{\sh{O}_{X_\lambda}}\sh{O}_X$.
Aplicando \sref{IV.6.9.3} a $S_\lambda$, $X_\lambda$ y $\sh{F}_\lambda$, obtenemos subpreesquemas $S'_{\lambda i}$ de $S_\lambda$ cuyos conjuntos subyacentes forman una partición de $S_\lambda$ y tales que, poniendo $X_{\lambda i}=X_\lambda\times_{S_\lambda}S'_{\lambda i}$, el Módulo
\[
\sh{F}_{\lambda i}=\sh{F}_\lambda\otimes_{\sh{O}_{X_\lambda}}\sh{O}_{X_{\lambda i}}
\]
es plano sobre $S'_{\lambda i}$.
Los subpreesquemas $S'_i=S'_{\lambda i}\times_{S_\lambda}S$ de $S$ poseen entonces las propiedades requeridas en virtud de \sref{IV.2.1.4}.
\end{proof}

\subsection{Propiedades de permanencia de los morfismos por paso al límite proyectivo}
\label{subsection:IV.8.10}

En este parágrafo conservamos las hipótesis generales y las notaciones de \sref{IV.8.8.1}.

\begin{proposition}[8.10.1]
\label{IV.8.10.1}
Si existe $\lambda$ tal que, para $\mu\geqslant\lambda$, $f_\mu$ es un morfismo abierto (resp.\ universalmente abierto), entonces $f$ es un morfismo abierto (resp.\ universalmente abierto).
\end{proposition}

\begin{proof}
Como $X=X_\lambda\times_{Y_\lambda}Y$, la aserción relativa a los morfismos universalmente abiertos es un caso particular de \sref{IV.2.4.3},~(iii).
Supongamos, pues, que $f_\mu$ es abierto para $\mu\geqslant\lambda$; basta ver que, para todo abierto \emph{cuasicompacto} $U$ de $X$, $f(U)$ es abierto en $Y$.
Ahora bien, existe $\mu\geqslant\lambda$ tal que $U=v_\mu^{-1}(U_\mu)$, donde $U_\mu$ es un abierto cuasicompacto de $X_\mu$ \sref{IV.2.3.11}; se tiene entonces $f(v_\mu^{-1}(U_\mu))=w_\mu^{-1}(f_\mu(U_\mu))$ \sref[I]{I.3.4.8}, y por tanto $f(U)$ es abierto en $Y$.
\end{proof}

\begin{corollary}[8.10.2]
\label{IV.8.10.2}
Sea $f:X\to Y$ un morfismo.
Para que $f$ sea universalmente abierto, basta que, para todo entero $n>0$, poniendo $Y_n=Y\otimes_\mathbf{Z}\mathbf{Z}[T_1,\ldots,T_n](=\mathbf{V}_Y^n)$, $X_n=X\times_Y Y_n$ y $f_n=f\times 1_{Y_n}:X_n\to Y_n$, la proyección $f_n$ sea un morfismo abierto.
\end{corollary}

\begin{proof}
Para demostrar que $f$ es universalmente abierto, basta demostrar que sigue siéndolo sobre todo abierto afín $U$ de $Y$; por hipótesis, si $U_n=U\otimes_\mathbf{Z}\mathbf{Z}[T_1,\ldots,T_n]$ es la imagen inversa de $U$ en $Y_n$, el morfismo $f_n^{-1}(U_n)\to U_n$, restricción de $f_n$, es abierto, y se ve que podemos reducirnos al caso en que $Y=\operatorname{Spec}(A)$ es afín.
Además, basta evidentemente demostrar que, para todo morfismo $Y'\to Y$, donde $Y'=\operatorname{Spec}(A')$ es también \emph{afín}, $f'=f_{|Y'}$ es abierto.
Supongamos primero que $A'$ es una $A$-álgebra de tipo finito, y por tanto un \emph{cociente} de un álgebra de polinomios $A[T_1,\ldots,T_n]$; entonces $Y'$ es un subpreesquema cerrado de $Y_n$ y $f'$ la restricción
\oldpage[IV-3]{37}
de $f_n$ a $f_n^{-1}(Y')$; pero, para todo abierto $V$ de $X_n$, se tiene $f_n(V\cap f_n^{-1}(Y'))=f_n(V)\cap Y'$, y como por hipótesis $f_n(V)$ es abierto en $Y_n$, esto demuestra que $f'$ es también un morfismo abierto.
Cuando $A'$ es arbitraria, puede considerarse como límite inductivo de sus sub-$A$-álgebras de tipo finito $A'_\lambda$, y el hecho de que $f'$ sea un morfismo abierto se sigue de lo anterior y de \sref{IV.8.10.1}.
\end{proof}

\begin{proposition}[8.10.3]
\label{IV.8.10.3}
Supongamos que existe $\alpha$ tal que: $1^\circ$ $S_\alpha$ es cuasicompacto; $2^\circ$ los morfismos $X_\alpha\to S_\alpha$, $Y_\alpha\to S_\alpha$ son cuasicompactos y el morfismo $Y_\alpha\to S_\alpha$ es cuasiseparado; $3^\circ$ para $\alpha\leqslant\lambda\leqslant\mu$, los morfismos $u_{\lambda\mu}:S_\mu\to S_\lambda$ son planos; $4^\circ$ $\overline{f_\alpha(X_\alpha)}$ es construible en $Y_\alpha$.
Entonces, para que $f$ sea dominante, es necesario y suficiente que exista $\lambda\geqslant\alpha$ tal que $f_\lambda$ sea dominante.
\end{proposition}

\begin{proof}
Las hipótesis implican que $Y_\alpha$ es cuasicompacto y que el morfismo $f_\alpha$ es cuasicompacto \sref{IV.1.2.4}; por tanto, $\overline{f_\alpha(X_\alpha)}=Z_\alpha$ es proconstruible \sref{IV.1.9.5},~(vii) en $Y_\alpha$.
Si se pone $Z_\lambda=w_{\alpha\lambda}^{-1}(Z_\alpha)$ para $\lambda\geqslant\alpha$ y $Z=w_\alpha^{-1}(Z_\alpha)$, se tiene $Z_\lambda=\overline{f_\lambda(X_\lambda)}$ y $Z=\overline{f(X)}$ \sref[I]{I.3.4.8}, y $Z_\lambda$ es proconstruible en $Y_\lambda$ \sref{IV.1.9.5},~(vii).
Basta entonces aplicar \sref{IV.8.3.13}, sustituyendo $S_\lambda$, $Z'_\lambda$ y $Z''_\lambda$ por $Y_\lambda$, $Y_\lambda$ y $Z_\lambda$, respectivamente.
\end{proof}

\begin{proposition}[8.10.4]
\label{IV.8.10.4}
Supongamos que existe $\alpha$ tal que $Y_\alpha$ es cuasicompacto y $f_\alpha$ es de tipo finito y cuasiseparado.
Para que el morfismo $f$ sea separado, es necesario y suficiente que exista $\lambda\geqslant\alpha$ tal que $f_\lambda$ sea separado.
\end{proposition}

\begin{proof}
Como la cuestión es local sobre $Y_\alpha$ (pues $Y_\alpha$ es cuasicompacto y $L$ filtrante), podemos reducirnos al caso en que $Y_\alpha$ es afín, y por tanto cuasiseparado; la hipótesis implica entonces que $X_\alpha$ (y por consiguiente los $X_\lambda$ y $X$) es cuasicompacto y cuasiseparado.
Pongamos $X'_\lambda=X_\lambda\times_{Y_\lambda}X_\lambda$ para $\lambda\geqslant\alpha$ y $X'=X\times_Y X$; se tiene $X'_\lambda=X'_\alpha\times_{Y_\alpha}Y_\lambda$ y $X'=X'_\alpha\times_{Y_\alpha}Y$. La primera proyección $X'_\alpha\to X_\alpha$ es cuasicompacta y cuasiseparada por hipótesis \sref{IV.1.2.3},~(iii), y por tanto $X'_\lambda$ es cuasicompacto y cuasiseparado.
Observemos ahora que, si denotamos por $\Delta_\lambda$ (resp.\ $\Delta$) la diagonal de $X_\lambda\times_{Y_\lambda}X_\lambda$ (resp.\ $X\times_Y X$), se sigue de \sref[I]{I.5.3.4} y \sref[I]{I.3.4.8} que $\Delta_\mu$ (resp.\ $\Delta$) es la imagen inversa de $\Delta_\lambda$ por la aplicación $v'_{\lambda\mu}:X'_\mu\to X'_\lambda$ (resp.\ $v'_\lambda:X'\to X'_\lambda$).
Por otra parte, $\Delta_\lambda$ es \emph{construible} en $X'_\lambda$: en efecto, como $f_\lambda$ es cuasiseparado, la inmersión diagonal $X_\lambda\to X'_\lambda$ es cuasicompacta, y localmente de presentación finita porque $f_\lambda$ es de tipo finito \sref[I]{I.1.4.3} y \sref[I]{I.1.5.4},~(iii); se sigue entonces de \sref{IV.1.8.4.1} que $\Delta_\lambda$ es localmente construible, y por tanto construible, ya que $X'_\lambda$ es cuasicompacto.
Podemos aplicar ahora \sref{IV.8.3.12}, sustituyendo $S_\lambda$ y $Z_\lambda$ por $X'_\lambda$ y $\Delta_\lambda$, respectivamente.
\end{proof}

\begin{theorem}[8.10.5]
\label{IV.8.10.5}
Supongamos $S_0$ cuasicompacto, $X_\alpha$ e $Y_\alpha$ de presentación finita sobre $S_\alpha$, y sea $f_\alpha:X_\alpha\to Y_\alpha$ un $S_\alpha$-morfismo.
Consideremos, para un morfismo, la propiedad de ser:
\begin{enumerate}[label=(\roman*)]
\item un isomorfismo;
\item[\rm(i\,\emph{bis})] un monomorfismo;
\item una inmersión;
\item una inmersión abierta;
\item una inmersión cerrada;
\item separado;
\item sobreyectivo;
\oldpage[IV-3]{38}
\item radicial;
\item afín;
\item cuasi-afín;
\item finito;
\item cuasifinito;
\item propio.
\end{enumerate}
Entonces, si $\mathbf{P}$ denota una de las propiedades anteriores, para que $f$ posea la propiedad $\mathbf{P}$, es necesario y suficiente que exista $\lambda\geqslant\alpha$ tal que $f_\lambda$ posea la propiedad $\mathbf{P}$ (en cuyo caso $f_\mu$ la posee también para $\mu\geqslant\lambda$).

Si se supone además que $S_0$ es cuasiseparado, la misma conclusión vale cuando $\mathbf{P}$ es la propiedad de ser:
\begin{enumerate}[label=(\roman*)]
\setcounter{enumi}{12}
\item proyectivo;
\item cuasiproyectivo.
\end{enumerate}
\end{theorem}

\begin{proof}
El caso en que $\mathbf{P}$ es una de las propiedades~(i) o~(v) se ha incluido en el enunciado solo a título recordatorio; en estos casos, el teorema se sigue de lo demostrado respectivamente en \sref{IV.8.8.2.4} y \sref{IV.8.10.4}.
Además, teniendo en cuenta \sref[I]{I.5.4.1} y \sref[I]{I.5.3.4}, (v) se sigue también de (iv).
El caso (i\,\emph{bis}) se sigue inmediatamente de (i), utilizando \sref[I]{I.5.3.8} y observando (como ya hemos utilizado en \sref{IV.8.10.4}) que la diagonal $\Delta_f$ se deduce de $\Delta_{f_\lambda}$ por el cambio de base $S\to S_\lambda$.

Observaremos, por otra parte, que (vi), (vii) y (xi) son en realidad condiciones sobre las \emph{fibras} $f^{-1}(y)$ de los morfismos considerados, teniendo en cuenta la transitividad de las fibras por cambio de base \sref[I]{I.3.6.4}: la condición~(vi) significa, en efecto, que todas las fibras deben ser no vacías; la condición~(vii), que deben ser radiciales \sref[I]{I.3.5.8}; y la condición~(xi), que deben ser finitas \sref[III]{III.6.2.2}, \sref[III]{III.6.2.3} y \sref[I]{I.6.4.4}, teniendo en cuenta que $f$ y los $f_\lambda$ son morfismos de tipo finito por \sref{IV.1.5.4},~(v).
El teorema en estos tres casos será todavía consecuencia de un resultado general sobre este tipo de propiedades, que se establecerá en \sref{IV.9.3.3}; aplazamos, pues, hasta entonces la demostración del teorema en el caso~(xi) (el lector podrá comprobar, naturalmente, que salvo en los n\textsuperscript{os}~{\hyperlink{ega4.c04.subsec.8-11}{8.11}} y {\hyperlink{ega4.c04.subsec.8-12}{8.12}}, no utilizaremos el teorema en estos casos antes de \sref{IV.9.3.3}, y que \sref{IV.8.11} y \sref{IV.8.12} no se utilizarán antes de \sref{IV.9.3.3}).

Para los casos que quedan por demostrar, podemos limitarnos a probar que la condición del enunciado es necesaria, pues todas las propiedades $\mathbf{P}$ consideradas son invariantes por cambio de base (véanse los caps.~I\textsuperscript{er} y II para las referencias relativas a cada una de estas propiedades).
Podemos suponer además que $S_\alpha=S_0$ y que $Y_\alpha=S_\alpha$, y por tanto $Y_\lambda=S_\lambda$ para todo $\lambda\geqslant\alpha$.
Por último, las propiedades~(i) a~(xii) son locales sobre $S_0$; así, como $S_0$ es unión finita de abiertos afines y $L$ es filtrante, podemos limitarnos a demostrarlas para $S_0=\operatorname{Spec}(A_0)$ \emph{afín} (y por tanto cuasiseparado).
Denotaremos por $A_\lambda$ (resp.\ $A$) el anillo de $S_\lambda$ (resp.\ $S$).

\emph{Casos} (ii), (iii), (iv): Supongamos que $f$ es una inmersión (resp.\ una inmersión abierta, una inmersión cerrada), y sea $X'$ el subpreesquema (resp.\ inducido sobre un abierto, resp.\ cerrado) de $S$ asociado a $f$, que es por tanto un $S$-preesquema de presentación finita.
\oldpage[IV-3]{39}
En virtud de \sref{IV.8.6.3}, existen entonces $\lambda\geqslant\alpha$ y un subpreesquema $X'_\lambda$ de $S_\lambda$, de presentación finita sobre $S_\lambda$, tales que $X'$ es isomorfo a $X'_\lambda\times_{S_\lambda}S$.
Para todo $\mu\geqslant\lambda$, $X'_\mu=X'_\lambda\times_{S_\lambda}S_\mu$ es, pues, un subpreesquema (resp.\ inducido sobre un abierto, resp.\ cerrado) de $S_\mu$, y se sigue entonces de \sref{IV.8.8.2.4} y \sref{IV.8.8.2.5} que existen $\mu\geqslant\lambda$ y un isomorfismo $g_\mu:X_\mu\to X'_\mu$ tales que $g=g_\mu\times 1_S$ es el isomorfismo $X\to X'$ asociado a $f$; de donde la conclusión.

\emph{Casos} (vi) \emph{y} (vii): Sabemos por \sref{IV.1.8.4.1} que $Z_\alpha=f_\alpha(X_\alpha)$ es construible en $S_\alpha$; si ponemos $Z_\lambda=u_{\alpha\lambda}^{-1}(Z_\alpha)$ para $\lambda\geqslant\alpha$ y $Z=u_\alpha^{-1}(Z_\alpha)$, se tiene $Z_\lambda=f_\lambda(X_\lambda)$ y $Z=f(X)$ \sref[I]{I.3.4.8}.
Como $Z_\alpha$ es proconstruible en $S_\alpha$, y en virtud de \sref{IV.8.3.11} la aplicación canónica $\varinjlim\mathfrak{C}(S_\lambda)\to\mathfrak{C}(S)$ es inyectiva, la relación $f(X)=S$ implica la existencia de $\lambda\geqslant\alpha$ tal que $f_\lambda(X_\lambda)=S_\lambda$, lo que demuestra el teorema en el caso~(vi).
Para demostrarlo en el caso~(vii), observemos que el morfismo estructural $X_\alpha\times_{S_\alpha}X_\alpha\to S_\alpha$ es de presentación finita, pues la primera proyección $X_\alpha\times_{S_\alpha}X_\alpha\to X_\alpha$ lo es \sref{IV.1.6.2}. Basta, pues, en virtud de \sref{IV.1.8.7.1}, aplicar el caso~(vi) al morfismo diagonal
\[
  \Delta_{f_\alpha}:X_\alpha\longrightarrow X_\alpha\times_{S_\alpha}X_\alpha,
\]
observando que $\Delta_{f_\lambda}=\Delta_{f_\alpha}\times1_{S_\lambda}$ y $\Delta_f=\Delta_{f_\alpha}\times1_S$ \sref[I]{I.5.3.4} y \sref[I]{I.3.3.11}.

\emph{Casos} (viii) \emph{y} (ix): Como $S=\operatorname{Spec}(A)$ es afín, decir que $f:X\to S$ es afín (resp.\ cuasi-afín) significa que existen un entero $r$ y una inmersión cerrada (resp.\ una inmersión) $j:X\to\mathbf{V}^r_S=\operatorname{Spec}(A[T_1,\ldots,T_r])$ de $S$-preesquemas, pues $f$ es de tipo finito y $S$ cuasicompacto \sref[II]{II.5.1.9}.
Como $\mathbf{V}^r_S=\mathbf{V}^r_{S_\lambda}\times_{S_\lambda}S$, y $\mathbf{V}^r_{S_\lambda}$ es un $S_0$-preesquema de presentación finita, se sigue de \sref{IV.8.8.2},~(i), aplicado al $S$-morfismo $j$, que existen $\lambda$ y un $S_\lambda$-morfismo $j_\lambda:X_\lambda\to\mathbf{V}^r_{S_\lambda}$ tales que $j=j_\lambda\times 1_S$; aplicando (iv) (resp.\ (ii)) a $j$, se deduce que existe $\mu\geqslant\lambda$ tal que $j_\mu$ es una inmersión cerrada (resp.\ una inmersión); por tanto, $f_\mu$ es afín (resp.\ cuasi-afín).

\emph{Caso} (x): Por hipótesis, $X=\operatorname{Spec}(B)$, donde $B$ es una $A$-álgebra que, como $A$-módulo, es de presentación finita \sref{IV.1.4.7}; se sigue, pues, de \sref{IV.8.5.2},~(ii) que existen $\lambda$ y un $A_\lambda$-módulo de presentación finita $B_\lambda$ tales que el $A$-módulo $B$ es isomorfo a $B_\lambda\otimes_{A_\lambda}A$.
La estructura de $A$-álgebra de $B$ está definida por un $A$-homomorfismo $m:B\otimes_A B\to B$; como $B\otimes_A B=(B_\lambda\otimes_{A_\lambda}B_\lambda)\otimes_{A_\lambda} A$, existen, por \sref{IV.8.5.2},~(i), $\mu\geqslant\lambda$ y un $A_\mu$-homomorfismo $m_\mu:B_\mu\otimes_{A_\mu}B_\mu\to B_\mu$ tales que $m=m_\mu\otimes 1$.
Considerando los diagramas habituales que expresan la asociatividad y la conmutatividad de $m$, se ve, aplicando de nuevo \sref{IV.8.5.2},~(i), que podemos suponer $\nu$ suficientemente grande para que $m_\nu$ defina sobre $B_\nu$ una multiplicación asociativa y conmutativa.
De igual modo, podemos suponer $\nu$ suficientemente grande para que $B_\nu$, así definida, admita un elemento unidad.
Si $X'_\nu=\operatorname{Spec}(B_\nu)$, está claro entonces que $X$ es $S$-isomorfo a $X'_\nu\times_{S_\nu}S$; por tanto, en virtud de (i), existe $\rho\geqslant\nu$ tal que $X_\rho$ y $X'_\rho$ son $S_\rho$-isomorfos, lo que termina la demostración en este caso.

Para demostrar el teorema en el caso~(xii), demostraremos primero la proposición siguiente:
\end{proof}

\begin{proposition}[8.10.5.1]
\label{IV.8.10.5.1}
(Lema de Chow para morfismos de presentación finita). --- Sea $\mathrm{A}$ un anillo, sean $X$ e $Y$ dos $\mathrm{A}$-preesquemas de presentación finita y sea $f:X\to Y$ un $\mathrm{A}$-morfismo separado. Entonces existen dos $\mathrm{A}$-preesquemas $X'$ y $P$ de presentación finita, y $\mathrm{A}$-morfismos $p:P\to Y$, $j:X'\to P$ y $g:X'\to X$, tales que el diagrama
\oldpage[IV-3]{40}
\[
\begin{tikzcd}[column sep=large,row sep=large]
X' \arrow[r,"j"] \arrow[d,"g"'] & P \arrow[d,"p"] \\
X \arrow[r,"f"'] & Y
\end{tikzcd}
\]
es conmutativo y tales que: $1^\circ$ $p$ es proyectivo; $2^\circ$ $g$ es proyectivo y sobreyectivo; $3^\circ$ $j$ es una inmersión abierta.
\end{proposition}

\begin{proof}
Elijamos $\mathrm{A}_0\subset\mathrm{A}$, $X_0$, $Y_0$ y $f_0$ como en \sref{IV.8.9.1}, con $Y_0$ noetheriano y $f_0$ de tipo finito; por \sref{IV.8.10.4}, podemos suponer además que $f_0$ es separado. El lema de Chow \sref[II]{II.5.6.1} proporciona entonces morfismos de tipo finito $p_0:P_0\to Y_0$, $g_0:X'_0\to X_0$ y $j_0:X'_0\to P_0$ tales que
\[
\begin{tikzcd}[column sep=large,row sep=large]
X'_0 \arrow[r,"j_0"] \arrow[d,"g_0"'] & P_0 \arrow[d,"p_0"] \\
X_0 \arrow[r,"f_0"'] & Y_0
\end{tikzcd}
\]
es conmutativo, $p_0$ es proyectivo, $g_0$ es proyectivo y sobreyectivo, y $j_0$ es una inmersión abierta. Las aserciones se siguen de la invariancia de estas propiedades por cambio de base \sref[II]{II.5.5.5},~(iii), \sref[I]{I.3.5.2} y \sref{IV.4.3.2}.
\end{proof}

\begin{proof}[Continuación de la demostración de \sref{IV.8.10.5}]
\emph{Caso (xii).} Apliquemos \sref{IV.8.10.5.1} a $f_0:X_0\to S_0$. Obtenemos un diagrama conmutativo
\[
\begin{tikzcd}[column sep=large,row sep=large]
X'_0 \arrow[r,"j_0"] \arrow[d,"g_0"'] & P_0 \arrow[d,"p_0"] \\
X_0 \arrow[r,"f_0"'] & S_0
\end{tikzcd}
\]
donde $p_0$ es proyectivo, $g_0$ es proyectivo y sobreyectivo, y $j_0$ es una inmersión abierta. Para cada $\lambda$ se obtiene el diagrama análogo, en el que $p_\lambda=p_0\times1_{S_\lambda}$, $g_\lambda=g_0\times1_{S_\lambda}$ y $j_\lambda=j_0\times1_{S_\lambda}$ poseen las mismas propiedades respectivas, y lo mismo ocurre con los morfismos límite inverso $p=p_0\times1_S$, $g=g_0\times1_S$ y $j=j_0\times1_S$. Como $g$ es propio \sref[II]{II.5.5.3}, también lo es $p\circ j=f\circ g$ \sref[II]{II.5.4.2}; y como $p$ es separado, $j$ es propio, y por tanto una inmersión cerrada. Apliquemos el caso~(iv) a $j$, observando que $X'_0$ y $P_0$ son $S_0$-preesquemas de presentación finita por \sref{IV.8.10.5.1} y \sref{IV.1.6.2}. Existe $\lambda$ tal que $j_\lambda$ es una inmersión cerrada, y por tanto propia \sref[II]{II.5.4.2}. Así,
\[
  f_\lambda\circ g_\lambda=p_\lambda\circ j_\lambda
\]
es propio \sref[II]{II.5.5.3} y \sref[II]{II.5.4.2}. Como $g_\lambda$ es sobreyectivo y, por la hipótesis sobre $f$ y el caso~(v), se puede suponer que $f_\lambda$ es separado, \sref[II]{II.5.4.3} demuestra que $f_\lambda$ es propio.

\emph{Casos (xiii) y (xiv).} Por el caso~(xii) y \sref[II]{II.5.5.3}, que se aplica porque los $S_\lambda$ son cuasicompactos y cuasiseparados por \sref{IV.1.7.19}, basta considerar el caso en que $f$ es cuasiproyectivo. Supongamos que existe un $\sh{O}_X$-Módulo invertible $\sh{L}$ que es amplio para $f$. Como $S_0$ es cuasicompacto y cuasiseparado, también lo es $X_0$ \sref{IV.1.2.3}; existen, pues, $\lambda$ y un $\sh{O}_{X_\lambda}$-Módulo cuasicoherente de presentación finita $\sh{L}_\lambda$ tales que $\sh{L}=v_\lambda^*(\sh{L}_\lambda)$ \sref{IV.8.5.2},~(ii). Por \sref{IV.8.5.5}, podemos suponer además que $\sh{L}_\lambda$ es invertible. La aserción se sigue ahora del lema más preciso que figura a continuación.

\begin{lemma}[8.10.5.2]
\label{IV.8.10.5.2}
\oldpage[IV-3]{41}
Supongamos que $S_0$ es cuasicompacto, y sea $\sh{L}_\lambda$ un $\sh{O}_{X_\lambda}$-Módulo invertible. Para que $\sh{L}$ sea amplio para $f$ (resp.\ muy amplio para $f$), es necesario y suficiente que exista $\mu\geqslant\lambda$ tal que $\sh{L}_\mu$ sea amplio para $f_\mu$ (resp.\ muy amplio para $f_\mu$).
\end{lemma}

\begin{proof}
La condición es manifiestamente suficiente \sref[II]{II.4.4.10} y \sref[II]{II.4.6.13}. Recíprocamente, como los $S_\lambda$ son cuasicompactos y los $f_\lambda$ son de tipo finito, tras sustituir $\sh{L}$ por una potencia tensorial adecuada, podemos reducirnos al caso en que $\sh{L}$ es muy amplio \sref[II]{II.4.6.11}. La cuestión es local sobre $S_0$ por \sref[II]{II.4.4.5} y por el hecho de que $L$ es filtrante, de modo que podemos suponer que $S_0$, y por tanto $S$, es afín. Por \sref[II]{II.4.4.1},~(ii) y \sref[II]{II.4.1.2}, existe una $S$-inmersión
\[
  j:X\longrightarrow\mathbf{P}^r_S=P
\]
tal que $\sh{L}\simeq j^*(\sh{O}_P(1))$. Por \sref{IV.8.8.2},~(i), el caso~(ii) y \sref[II]{II.4.1.3}, existen $\mu\geqslant\lambda$ y una inmersión $j_\mu:X_\mu\to\mathbf{P}^r_{S_\mu}=P_\mu$ tales que $j=j_\mu\times1_S$. Entonces \sref[II]{II.4.1.3.2} y \sref{IV.8.5.2.5} proporcionan $\nu\geqslant\mu$ tal que
\[
  \sh{L}_\nu\simeq j_\nu^*(\sh{O}_{P_\nu}(1)),
\]
lo que demuestra que $\sh{L}_\nu$ es muy amplio para $f_\nu$ \sref[II]{II.4.4.2}.
\end{proof}
\end{proof}

\subsection{Aplicación a los morfismos cuasifinitos}
\label{subsection:IV.8.11}
\label{IV.8.11}

Nos proponemos demostrar en este parágrafo los dos teoremas siguientes:

\begin{theorem}[8.11.1]
\label{IV.8.11.1}
Sea $f:X\to Y$ un morfismo propio, localmente de presentación finita y cuasifinito.
Entonces el morfismo $f$ es finito.
\end{theorem}

\begin{theorem}[8.11.2]
\label{IV.8.11.2}
Sea $f:X\to Y$ un morfismo localmente de presentación finita, cuasifinito y separado.
Entonces el morfismo $f$ es cuasi-afín y, con mayor razón, cuasiproyectivo.
\end{theorem}

\begin{remark}[8.11.3]
\label{IV.8.11.3}
Veremos más adelante que, para la demostración de \sref{IV.8.11.1} y \sref{IV.8.11.2}, podemos reducirnos al caso en que $Y$ es localmente noetheriano; obsérvese que en este caso se obtiene así otra demostración del teorema de Chevalley \sref[III]{III.4.4.2}.
\end{remark}

\begin{env}[8.11.4]
\label{IV.8.11.4}
Las hipótesis y las conclusiones de \sref{IV.8.11.1} y \sref{IV.8.11.2} son todas locales sobre $Y$.
En efecto, esto se sigue de \sref{IV.1.6.1}, \sref{IV.1.2.6}, \sref[II]{II.5.1.1}, \sref[II]{II.5.4.1} y \sref[II]{II.6.2.2}.
Podemos suponer, por tanto, $Y=\operatorname{Spec}(A)$ afín.
Sabemos que existen entonces un subanillo $A_0$ de $A$, que es una $\bb{Z}$-álgebra de tipo finito, y un morfismo de tipo finito $f_0:X_0\to\operatorname{Spec}(A_0)$ tales que $X$ se identifica con $X_0\otimes_{A_0}A$ y $f$ con $f_0\times 1$ \sref{IV.8.9.1}.
Además, $A$ puede considerarse como límite inductivo de sus subanillos que contienen $A_0$ y que son $\bb{Z}$-álgebras de tipo finito; utilizando el método de \sref{IV.8.1.2},~$c$)) y \sref{IV.8.10.5},~(v),~(xi) y (xii)), se ve que basta demostrar los teoremas \sref{IV.8.11.1} y \sref{IV.8.11.2} para $f_0$.
Supongamos en adelante $Y$ noetheriano; utilizando ahora el método de \sref{IV.8.1.2},~$a$)) y \sref{IV.8.10.5},~(v),~(ix),~(x),~(xi) y (xii)), se puede sustituir $Y$ por $\operatorname{Spec}(\sh{O}_y)$, donde $y$ es un
\oldpage[IV-3]{42}
punto arbitrario de $Y$, de modo que finalmente podemos suponer que $Y=\operatorname{Spec}(A)$, donde $A$ es un \emph{anillo local noetheriano}.
Sea $\mathfrak{m}$ el ideal maximal de $A$, y sea $\hat{A}$ el completado de $A$ para la topología $\mathfrak{m}$-ádica. El anillo $\hat{A}$ es local noetheriano y $\operatorname{Spec}(\hat{A})\to\operatorname{Spec}(A)$ es fielmente plano y cuasicompacto \sref[0]{0.7.3.5}; aplicando \sref{IV.2.7.1},~(i), (vii), (xiv), (xv) y (xvi), podemos reducirnos, por tanto, al caso en que $A$ es \emph{completo}.
Se sigue entonces de \sref[II]{II.6.2.6} que $X=X'\sqcup X''$, donde $X'$ es un $Y$-esquema finito y $X''$ un $Y$-esquema cuasifinito tal que $X''\cap f^{-1}(y)=\varnothing$.
\end{env}

Situémonos primero bajo las hipótesis de \sref{IV.8.11.1}; como $f$ es propio, $X''$, que es cerrado en $X$, es propio sobre $Y$ \sref[II]{II.5.4.10}, de modo que $f(X'')$ es cerrado en $Y$; pero $y$ no pertenece a $f(X'')$ y, además, $y$ es adherente a todo punto de $Y$, por lo que $f(X'')=\varnothing$; en consecuencia, $X''$ es vacío y $f$ es \emph{finito}.

Situémonos ahora bajo las hipótesis de \sref{IV.8.11.2} y, restringiéndonos aún más (como permite lo anterior) al caso en que $Y=\operatorname{Spec}(A)$ es afín y noetheriano de dimensión finita, razonemos además por inducción sobre la dimensión de $Y$.
Reduciéndonos como antes al caso en que $A$ es además local y completo, se tiene $\dim(\sh{O}_y)=\dim(A)=\dim(Y)$ y, para todo $z\neq y$, $\dim(\sh{O}_z)<\dim(\sh{O}_y)$, de modo que $\dim(Y-\{y\})<\dim(Y)$.
Ahora bien, por hipótesis se tiene $f(X'')\subset Y-\{y\}$ y la restricción de $f$ a $X''$ es evidentemente un morfismo cuasifinito y separado; aplicando la hipótesis de inducción a $Y-\{y\}$ y a $X''$, se ve que $X''$ es cuasi-afín sobre $Y-\{y\}$; pero, como el abierto $Y-\{y\}$ es cuasi-afín sobre $Y$ porque $Y$ es noetheriano, $X''$ es también cuasi-afín sobre $Y$ \sref[II]{II.5.1.10},~(ii)); y como además $X'$ es finito (y \emph{a fortiori} afín) sobre $Y$, $X$ es cuasi-afín sobre $Y$ \sref[II]{II.4.6.17} y \sref[II]{II.5.1.2},~$c'$)).

\begin{proposition}[8.11.5]
\label{IV.8.11.5}
Sea $f:X\to Y$ un morfismo de presentación finita.
Las propiedades siguientes son equivalentes:
\begin{enumerate}[label=\emph{\alph*)}]
  \item $f$ es una inmersión cerrada.
  \item $f$ es un monomorfismo propio.
  \item $f$ es propio y, para todo $y\in Y$, $f^{-1}(y)$ es radicial y geométricamente reducido sobre $k(y)$ (es decir, vacío o $k(y)$-isomorfo a $\operatorname{Spec}(k(y))$).
\end{enumerate}
\end{proposition}

Está claro que \emph{a)} entraña \emph{b)}.
Para ver que \emph{b)} entraña \emph{c)}, observemos \sref[I]{I.5.3.8} que, para todo $y\in Y$, el morfismo $f^{-1}(y)\to\operatorname{Spec}(k(y))$ deducido de $f$ por extensión de la base es un monomorfismo, y por tanto inyectivo; en consecuencia, $f^{-1}(y)$ es vacío o se reduce a un punto, y en todo caso es afín.
Además, si $A$ es el anillo de $f^{-1}(y)$, el homomorfismo canónico $A\otimes_k A\to A$ es biyectivo \sref[I]{I.5.3.8}.
Esto implica evidentemente que $A=k$, pues de lo contrario existiría un elemento $a\in A$ que no perteneciera a $k$, y las imágenes de $a\otimes 1$ y de $1\otimes a$ en $A$ serían ambas iguales a $a$, mientras que $a\otimes 1\neq1\otimes a$, ya que $1$ y $a$ forman un sistema libre sobre $k$.

Queda por demostrar que \emph{c)} entraña \emph{a)}.
Se sigue primero de \sref{IV.8.11.1} que $f$ es un morfismo \emph{finito}, de modo que $X=\operatorname{Spec}(\sh{A})$, donde $\sh{A}$ es una $\sh{O}_Y$-Álgebra finita.
Basta, pues, demostrar que el homomorfismo canónico $\sh{O}_Y\to\sh{A}$ es \emph{sobreyectivo} \sref[II]{II.1.4.10}, o, equivalentemente, que para todo $y\in Y$ el homomorfismo $\sh{O}_y\to\sh{A}_y$ es sobreyectivo.
Pero, por hipótesis,
\oldpage[IV-3]{43}
$f^{-1}(y)=\operatorname{Spec}(\sh{A}_y/\mathfrak{m}_y\sh{A}_y)$ \sref[II]{II.1.5.5} es tal que el homomorfismo correspondiente $k(y)=\sh{O}_y/\mathfrak{m}_y\to\sh{A}_y/\mathfrak{m}_y\sh{A}_y$ es biyectivo; como $\sh{A}_y$ es un $\sh{O}_y$-módulo de tipo finito, el lema de Nakayama demuestra que $\sh{O}_y\to\sh{A}_y$ es sobreyectivo, lo que termina la demostración.

\begin{remark}[8.11.5.1]
\label{IV.8.11.5.1}
Obsérvese que el razonamiento anterior demuestra que, si $f$ es un \emph{monomorfismo}, entonces, para todo $y\in Y$, $f^{-1}(y)$ es vacío o $k(y)$-isomorfo a $\operatorname{Spec}(k(y))$.
\end{remark}

\begin{proposition}[8.11.6]
\label{IV.8.11.6}
Si un morfismo $f:X\to Y$ de presentación finita es un homeomorfismo universal, es finito, sobreyectivo y radicial (siendo cierto el recíproco por \sref{IV.2.4.5},~(i))).
\end{proposition}

\begin{proof}
En efecto, como $f$ es de tipo finito, universalmente cerrado y separado en virtud de \sref{IV.2.4.4}, es propio por definición \sref[II]{II.5.4.1}.
Como es evidentemente cuasifinito \sref[II]{II.6.2.3}, es finito por \sref{IV.8.11.1}.
Sabemos además que es radicial \sref{IV.2.4.4}, y evidentemente sobreyectivo.
\end{proof}

\subsection{Nueva demostración y generalización del ``Teorema principal'' de Zariski}
\label{subsection:IV.8.12}
\label{IV.8.12}

\begin{lemma}[8.12.1]
\label{IV.8.12.1}
Sea $f:X\to Y$ un morfismo cuasicompacto y cuasiseparado, sea $\sh{C}$ una $\sh{O}_Y$-Álgebra cuasicoherente y sea $Z=\operatorname{Spec}(\sh{C})$, que es un $Y$-preesquema afín sobre $Y$.
Sea $g:X\to Z$ un $Y$-morfismo y sea $\varphi=\sh{A}(g):\sh{C}\to f_*(\sh{O}_X)$ el $\sh{O}_Y$-homomorfismo correspondiente de $\sh{O}_Y$-Álgebras \sref[II]{II.1.2.7}.
Supongamos que $g$ es una inmersión.
Entonces, para que la imagen cerrada de $X$ por $g$ \sref[I]{I.9.5.3} sea igual a $Z$, es necesario y suficiente que $\varphi$ sea inyectivo; $g$ es entonces una inmersión abierta.
\end{lemma}

\begin{proof}
La hipótesis implica que $f_*(\sh{O}_X)$ es una $\sh{O}_Y$-Álgebra cuasicoherente \sref{IV.1.7.5}; además, como el morfismo canónico $h:Z\to Y$ es afín, y por tanto cuasicompacto y cuasiseparado \sref{IV.1.2.2} y \sref{IV.1.1.2}, $g_*(\sh{O}_X)$ es una $\sh{O}_Z$-Álgebra cuasicoherente \sref{IV.1.7.5}.
Dicho esto, afirmar que la imagen cerrada de $X$ por $g$ es igual a $Z$ significa \sref[I]{I.9.5.2} que el homomorfismo canónico $\theta:\sh{O}_Z\to g_*(\sh{O}_X)$ es inyectivo.
Pero se tiene $h_*(\sh{O}_Z)=\sh{C}$ por definición de $Z$ \sref[II]{II.1.3.1}, y $h_*(g_*(\sh{O}_X))=f_*(\sh{O}_X)$.
Como $Z$ es afín sobre $Y$, equivale a decir que el homomorfismo $\theta:\sh{O}_Z\to g_*(\sh{O}_X)$ es inyectivo o que el homomorfismo correspondiente $\varphi=h_*(\theta):\sh{C}\to f_*(\sh{O}_X)$ es inyectivo \sref[I]{I.1.3.9}.
El hecho de que $g$ sea entonces una inmersión abierta se sigue de \sref[I]{I.9.5.10} y de la hipótesis.
\end{proof}

\begin{lemma}[8.12.2]
\label{IV.8.12.2}
Sea $Y$ un preesquema cuasicompacto y cuasiseparado, sea $f:X\to Y$ un morfismo cuasiseparado de tipo finito, sea $\sh{C}$ una $\sh{O}_Y$-Álgebra cuasicoherente y sea $Z=\operatorname{Spec}(\sh{C})$.
Sea $g:X\to Z$ un $Y$-morfismo y sea $\varphi:\sh{C}\to f_*(\sh{O}_X)$ el $\sh{O}_Y$-homomorfismo correspondiente de $\sh{O}_Y$-Álgebras.
Sea $(\sh{C}_\lambda)$ la familia creciente filtrante de sub-$\sh{O}_Y$-Álgebras cuasicoherentes de tipo finito de $\sh{C}$ (cuya unión es $\sh{C}$ \sref[I]{I.9.6.6} y \sref{IV.1.7.9}); pongamos $Z_\lambda=\operatorname{Spec}(\sh{C}_\lambda)$ y sea $g_\lambda$ el morfismo compuesto $X\xrightarrow{g}Z\to Z_\lambda$.
Entonces las condiciones siguientes son equivalentes:
\begin{enumerate}[label=\emph{\alph*)}]
  \item $g$ es una inmersión.
  \item Existe $\lambda$ tal que $g_\lambda$ es una inmersión.
\end{enumerate}
\oldpage[IV-3]{44}
Además, cuando $g_\lambda$ es una inmersión, lo mismo ocurre con $g_\mu$ para $\mu\geqslant\lambda$.

Basta aplicar \sref[II]{II.3.8.4}, sustituyendo $\sh{L}$ por $\sh{O}_X$ y $\sh{S}$ por $\sh{C}[\mathrm{T}]$, y teniendo en cuenta \sref[II]{II.3.1.7}.
\end{lemma}

\begin{proposition}[8.12.3]
\label{IV.8.12.3}
Sea $Y$ un preesquema cuasicompacto y cuasiseparado, y sea $f:X\to Y$ un morfismo separado de tipo finito.
Sea $\sh{B}=f_*(\sh{O}_X)$, que es una $\sh{O}_Y$-Álgebra cuasicoherente \sref[I]{I.9.2.2}; sea $\sh{C}$ la clausura entera de $\sh{O}_Y$ en $\sh{B}$, que es una $\sh{O}_Y$-Álgebra cuasicoherente \sref[II]{II.6.3.4}; pongamos $Z=\operatorname{Spec}(\sh{C})$, y sea $g:X\to Z$ el $Y$-morfismo correspondiente a la inyección $\sh{C}\to\sh{B}=f_*(\sh{O}_X)$; para todo $\lambda$, sea $g_\lambda:X\to Z_\lambda=\operatorname{Spec}(\sh{C}_\lambda)$ el $Y$-morfismo correspondiente a la inyección $\sh{C}_\lambda\to\sh{B}=f_*(\sh{O}_X)$.
Entonces las condiciones siguientes son equivalentes:
\begin{enumerate}[label=\emph{\alph*)}]
  \item[\emph{a)}] Existe una factorización de $f$ de la forma
\[
X\xrightarrow{f'}Y'\xrightarrow{u}Y
\]
donde $f'$ es una inmersión y $u$ un morfismo finito.
  \item[\emph{a$'$)}] Existe una factorización $X\xrightarrow{f'}Y'\xrightarrow{u}Y$ de $f$, donde $f'$ es una inmersión abierta y $u$ un morfismo finito.
  \item[\emph{b)}] El morfismo $g:X\to Z$ es una inmersión.
  \item[\emph{c)}] Existe $\lambda$ tal que $g_\lambda:X\to Z_\lambda$ es una inmersión.
\end{enumerate}
Además, cuando $g_\lambda$ es una inmersión, $g$ es una inmersión abierta, $g(X)$ es denso en $Z$, y existe $\lambda$ tal que, para $\mu\geqslant\lambda$, $g_\mu$ es una inmersión abierta.
\end{proposition}

Como el homomorfismo $\varphi:\sh{C}\to f_*(\sh{O}_X)$ es inyectivo, se sigue de \sref{IV.8.12.1} que, si $g$ es una inmersión, es una inmersión abierta y $g(X)$ es denso en $Z$; lo mismo ocurre con $g_\lambda$.
El hecho de que \emph{a)} entrañe \emph{a$'$)} se sigue también de \sref{IV.8.12.1}, aplicado sustituyendo $Z$ por $Y'$ y $g$ por $f'$ (pues $Y'$ es finito sobre $Y$, y por tanto afín sobre $Y$): en efecto, si $Y''$ es la imagen cerrada de $X$ por $f'$, $Y''$ es finito sobre $Y$ y $f'$ se factoriza como $X\xrightarrow{f''}Y''\xrightarrow{j}Y'$, donde $j$ es la inyección canónica y $f''$ es una inmersión \sref[I]{I.4.1.10}; se sigue, pues, de \sref{IV.8.12.1} que $f''$ es una inmersión abierta.

La equivalencia de \emph{b)} y \emph{c)} se sigue de \sref{IV.8.12.2}, al igual que el hecho de que $g_\lambda$ sea entonces una inmersión para $\lambda$ suficientemente grande.
Está claro que \emph{c)} entraña \emph{a)}, puesto que $Z_\lambda$ es finito sobre $Y$ \sref[II]{II.6.3.4} y \sref[II]{II.6.1.2}.
Demostremos finalmente que \emph{a)} entraña \emph{c)}.
Hemos visto más arriba que $Y'$ es la imagen cerrada de $X$ por $f'$, y se sigue entonces de \sref{IV.8.12.1} que, poniendo $\sh{B}'=u_*(\sh{O}_{Y'})$, el homomorfismo $\varphi':\sh{B}'\to\sh{B}=f_*(\sh{O}_X)$ es \emph{inyectivo}.
Pero, como por hipótesis $\sh{B}'$ es una $\sh{O}_Y$-álgebra finita, se identifica, por definición de $\sh{C}$, con una de las $\sh{O}_Y$-subálgebras $\sh{C}_\lambda$, lo que demuestra \emph{c)}.

Diremos que un morfismo $f:X\to Y$, donde $Y$ es cuasicompacto y cuasiseparado, es \emph{pseudofinito} si es de tipo finito y satisface la condición \emph{a)} de \sref{IV.8.12.3} (en cuyo caso es necesariamente separado).

\begin{corollary}[8.12.4]
\label{IV.8.12.4}
Sea $Y$ un preesquema cuasicompacto y cuasiseparado, y sea $f:X\to Y$ un morfismo.
\oldpage[IV-3]{45}
\begin{enumerate}[label=(\roman*)]
  \item Supongamos $f$ pseudofinito.
Entonces, para todo morfismo $Y'\to Y$, donde $Y'$ es cuasicompacto y cuasiseparado, $f_{(Y')}:X'=X_{(Y')}\to Y'$ es pseudofinito.
  \item Sea $(U_\lambda)$ un recubrimiento de $Y$ por abiertos cuasicompactos.
Para que $f$ sea pseudofinito, es necesario y suficiente que, para todo $\lambda$, la restricción $f_\lambda:f^{-1}(U_\lambda)\to U_\lambda$ de $f$ sea un morfismo pseudofinito.
  \item Supongamos además que $Y$ es noetheriano y $f$ de tipo finito.
Entonces, para que $f$ sea pseudofinito, es necesario y suficiente que, para todo $y\in Y$, el morfismo $f_y=f\times 1:X_y=X\times_Y\operatorname{Spec}(\sh{O}_y)\to\operatorname{Spec}(\sh{O}_y)$ sea pseudofinito.
\end{enumerate}
\end{corollary}

\begin{proof}
(i) Está claro que $f_{(Y')}$ es de tipo finito \sref{IV.1.5.4}. Además, una factorización $X\xrightarrow{g}Z\xrightarrow{u}Y$, con $g$ una inmersión y $u$ finito, induce una factorización $X'\xrightarrow{g'}Z'\xrightarrow{u'}Y'$ de $f_{(Y')}$, donde $Z'=Z_{(Y')}$, $g'=g_{(Y')}$ y $u'=u_{(Y')}$. El morfismo $g'$ es una inmersión \sref[I]{I.4.3.2} y $u'$ es finito \sref[II]{II.6.1.5}; por tanto, $f_{(Y')}$ es pseudofinito.

(ii) La condición es necesaria en virtud de (i), pues los $U_\lambda$ son cuasiseparados al serlo $Y$.
Para ver que es suficiente, observemos (con la notación de \sref{IV.8.12.3}) que, si se pone $X_\lambda=f^{-1}(U_\lambda)$, entonces $\sh{B}|U_\lambda=(f_\lambda)_*(\sh{O}_{X_\lambda})$, y $\sh{C}|U_\lambda$ es la clausura entera de $\sh{O}_{U_\lambda}$ en $\sh{B}|U_\lambda$. En consecuencia, si $h:Z\to Y$ es el morfismo canónico, $Z_\lambda=\operatorname{Spec}(\sh{C}|U_\lambda)$ se identifica con $h^{-1}(U_\lambda)$. Para que $g:X\to Z$ sea una inmersión, es necesario y suficiente que, para todo $\lambda$, su restricción $g_\lambda:f^{-1}(U_\lambda)\to h^{-1}(U_\lambda)$ sea una inmersión \sref[I]{I.4.2.4}. La conclusión se sigue ahora de \sref{IV.8.12.3}.

(iii) Basta, en virtud de (ii), demostrar que, para todo $y\in Y$, existe un entorno $U$ de $y$ tal que la restricción $f^{-1}(U)\to U$ de $f$ sea pseudofinita.
Denotemos por $(U_\lambda)$ el sistema proyectivo filtrante de entornos abiertos afines de $y$, y apliquemos el método de \sref{IV.8.1.2},~$a$)).
Como $Y$ es noetheriano, las restricciones $f_\lambda:f^{-1}(U_\lambda)\to U_\lambda$ de $f$ son de presentación finita, y lo mismo ocurre con $f_y$.
Por hipótesis, $f_y$ se factoriza como $X_y\xrightarrow{g_y}Z_y\xrightarrow{u_y}\operatorname{Spec}(\sh{O}_y)$, donde $u_y$ es finito y $g_y$ una inmersión.
Como $\sh{O}_y$ es noetheriano, también lo es $Z_y$; y como $Z_y$ es de presentación finita sobre $\operatorname{Spec}(\sh{O}_y)$, existen $\lambda$ y un morfismo de presentación finita $u_\lambda:Z_\lambda\to U_\lambda$ tales que $Z_y$ se identifica con $Z_\lambda\times_{U_\lambda}\operatorname{Spec}(\sh{O}_y)$ y $u_y$ con $u_\lambda\times 1$ \sref{IV.8.8.2},~(ii)). Además, existe un morfismo $g_\lambda:X_\lambda\to Z_\lambda$ tal que $g_y=g_\lambda\times 1$ y $f_\lambda=u_\lambda\circ g_\lambda$ \sref{IV.8.8.2},~(i)). Se puede tomar además $\lambda$ de modo que $g_\lambda$ sea una inmersión y $u_\lambda$ un morfismo finito \sref{IV.8.10.5},~(ii) y (x)), lo que demuestra que $f_\lambda$ es pseudofinito.
\end{proof}

\begin{env}[8.12.5]
\label{IV.8.12.5}
Podemos dar ahora del ``Teorema principal'' de Zariski \sref[III]{III.4.4.3}) una demostración que no utiliza los resultados cohomológicos ``globales'' del cap.~III, sino que recurre a las propiedades más finas de los anillos locales noetherianos; generalizaremos además el enunciado del teorema suprimiendo las hipótesis noetherianas:
\end{env}

\begin{theorem}[8.12.6]
\label{IV.8.12.6}
(``Teorema principal'' de Zariski). --- Sea $Y$ un preesquema cuasicompacto y cuasiseparado.
Si un morfismo $f:X\to Y$ es cuasifinito, separado y de presentación finita, existe una factorización de $f$
\begin{equation}
\label{IV.8.12.6.1}
X\xrightarrow{f'}Y'\xrightarrow{u}Y \tag{8.12.6.1}
\end{equation}
donde $f'$ es una inmersión abierta y $u$ un morfismo finito.
\end{theorem}

\oldpage[IV-3]{46}

\begin{proof}
En virtud de \sref{IV.8.12.4},~(ii)) y del carácter local (sobre $Y$) de las nociones de morfismo cuasifinito, separado y de presentación finita, podemos reducirnos al caso en que $Y=\operatorname{Spec}(A)$ es afín.
Aplicando \sref{IV.8.9.1}), podemos suponer que existen un subanillo $A_0$ de $A$, que es una $\bb{Z}$-álgebra de tipo finito, y un $A$-isomorfismo $X_0\otimes_{A_0}A\xrightarrow{\sim}X$, identificándose $f$ mediante este isomorfismo con $f_0\times 1$, donde $f_0:X_0\to\operatorname{Spec}(A_0)$ es un morfismo \emph{de tipo finito}; además, por \sref{IV.8.10.5},~(v) y (xi)), podemos suponer que $f_0$ es separado y cuasifinito; si se demuestra que $f_0$ es pseudofinito, lo mismo ocurrirá con $f$ por \sref{IV.8.12.4},~(i)).
Como $A_0$ es entonces noetheriano y las nociones de morfismo de tipo finito, separado y cuasifinito se conservan por cambio de base, se sigue de \sref{IV.8.12.4},~(iii)) que incluso podemos suponer que $A$ es un anillo local, \emph{esencialmente de tipo finito sobre} $\bb{Z}$ \sref{IV.1.3.8}).
Pongamos $n=\dim(A)$ y procedamos por inducción sobre $n$; para $n=0$, el teorema es evidente, pues $A$ es un cuerpo y el morfismo $f$ ya es finito \sref[II]{II.6.2.2}).
Pongamos $B=\Gamma(X,\sh{O}_X)$; denotemos por $C$ la clausura entera de $A$ en $B$, pongamos $Z=\operatorname{Spec}(C)$ y sea $g:X\to Z$ el $Y$-morfismo correspondiente al $A$-homomorfismo canónico inyectivo $C\to B$; en virtud de \sref{IV.8.12.3}), basta demostrar que $g$ es una inmersión abierta.
Sea $a$ el punto cerrado de $Y$, y sea $U=Y-\{a\}$; $U$ es noetheriano y todos sus anillos locales son esencialmente de tipo finito sobre $\bb{Z}$ y de dimensión $<n$; teniendo en cuenta la hipótesis de inducción, se ve que la restricción $f^{-1}(U)\to U$ de $f$ es un morfismo pseudofinito.
Se concluye de \sref{IV.8.12.3}) que, si $h:Z\to Y$ es el morfismo estructural, la restricción $f^{-1}(U)\to h^{-1}(U)$ de $g$ es una inmersión abierta.
Pongamos $A'=\hat{A}$, $Y'=\operatorname{Spec}(A')$, $X'=X_{(Y')}$ y $f'=f_{(Y')}:X'\to Y'$.
Como el morfismo canónico $u:Y'\to Y$ es plano, se sigue de \sref{IV.2.3.1}) que $B'=\Gamma(X',\sh{O}_{X'})$ se identifica con la $A'$-álgebra $B\otimes_A A'$.
Por otra parte, como $A$ es un anillo local excelente \sref{IV.7.8.3}), el morfismo $u:Y'\to Y$ es regular y, en consecuencia, \sref{IV.6.14.4}) \emph{la clausura entera de $A'$ en $B'$ es igual a} $C'=C\otimes_A A'$.
Se ve, por tanto, que $Z'=\operatorname{Spec}(C')$ es igual a $Z_{(Y')}$ y que el morfismo $g':X'\to Z'$ procedente de la inyección canónica $C'\to B'$ es igual a $g_{(Y')}$.
Como $u:Y'\to Y$ es fielmente plano y cuasicompacto, para demostrar que $g$ es una inmersión abierta basta demostrar que $g'$ es una inmersión abierta \sref{IV.2.7.1},~(x)).
Ahora bien, $u^{-1}(a)$ consta del punto cerrado $a'$ de $Y'$, y en consecuencia $U'=Y'-\{a'\}=u^{-1}(U)$.
Si $h':Z'\to Y'$ es el morfismo canónico, el hecho de que la restricción $f^{-1}(U)\to h^{-1}(U)$ de $g$ sea una inmersión abierta implica que lo mismo ocurre con la restricción $f'^{-1}(U')\to h'^{-1}(U')$ de $g'$.
El morfismo $f'$ es separado y cuasifinito \sref[II]{II.6.2.4}; como $A'$ es \emph{completo}, se deduce de \sref[II]{II.6.2.6} que $X'$ es $Y'$-isomorfo a una suma $X'_1\sqcup X'_2$, donde la restricción $f'|X'_1=f'_1:X'_1\to Y'$ es un morfismo \emph{finito}, y $X'_2\subset f'^{-1}(U')$.
De ello se sigue que $B'$ es el producto directo de las dos $A'$-álgebras $\Gamma(X'_1,\sh{O}_{X'_1})=B'_1$ y $\Gamma(X'_2,\sh{O}_{X'_2})=B'_2$; se concluye inmediatamente que la clausura entera $C'$ de $A'$ en $B'$ es el producto directo de las clausuras enteras $C'_1$, $C'_2$ de $A'$ en $B'_1$, $B'_2$, respectivamente, de donde $Z'=Z'_1\sqcup Z'_2$, donde $Z'_i=\operatorname{Spec}(C'_i)$ ($i=1,2$); y el morfismo canónico $g':X'\to Z'$ es tal que $g'|X'_i$ es el morfismo canónico $g'_i:X'_i\to Z'_i$ ($i=1,2$).
Pero, como $B'_1$ ya es una $A'$-álgebra finita, se tiene $C'_1=B'_1$, y por tanto $g'_1$ es un isomorfismo.
Por otra parte, como $X'_2\subset f'^{-1}(U')$ y es abierto en $f'^{-1}(U')$, se
\oldpage[IV-3]{47}
sabe ya que $g'_2$ es una inmersión abierta.
Se concluye que $g'$ es, en efecto, una inmersión abierta.
\end{proof}

\begin{remark}[8.12.7]
\label{IV.8.12.7}
Cuando, en \sref{IV.8.12.6}), se supone que $Y$ es un \emph{esquema} afín, la demostración por reducción al caso noetheriano muestra que, en la factorización \sref{IV.8.12.6.1}), los morfismos $f'$ y $u$ son también morfismos \emph{de presentación finita} \sref{IV.1.6.2}).
\end{remark}

\begin{corollary}[8.12.8]
\label{IV.8.12.8}
Sea $Y$ un esquema cuasicompacto tal que existe un $\sh{O}_Y$-Módulo amplio \sref[II]{II.4.5.3}), y sea $f:X\to Y$ un morfismo cuasifinito y cuasiproyectivo.
Entonces existe una factorización de $f$ de la forma
\[
X\xrightarrow{f'}Y'\xrightarrow{u}Y
\]
donde $f'$ es una inmersión abierta y $u$ un morfismo finito.
\end{corollary}

\begin{proof}
La hipótesis implica que $X$ se identifica con un sub-$Y$-esquema cuasicompacto de un $Y$-esquema de la forma $Z=\mathbf{P}^r_Y$ \sref[II]{II.5.3.3}).
Existe, pues, un entorno abierto cuasicompacto $U$ de $X$ en $Z$ tal que $X$ es cerrado en $U$; como $Z$ es un esquema, la inyección canónica $U\to Z$ es un morfismo de presentación finita (\sref{IV.1.2.7} y \sref{IV.1.6.2})), y por tanto el morfismo compuesto $g:U\to Z\to Y$ es también de presentación finita (el hecho de que $\mathbf{P}^r_Y$ sea de presentación finita sobre $Y$ se sigue inmediatamente de la definición \sref[II]{II.4.1.1})).
Sea $\sh{J}$ el Ideal cuasicoherente de $\sh{O}_U$ que define el subpreesquema cerrado $X$; como $U$ es un esquema cuasicompacto, $\sh{J}$ es el límite inductivo filtrante de sus sub-Ideales cuasicoherentes de tipo finito $\sh{J}_\lambda$ \sref[I]{I.9.4.9}).
Si $X_\lambda$ es el subpreesquema cerrado de $U$ definido por $\sh{J}_\lambda$, se tiene entonces $X=\bigcap X_\lambda$.
Para todo $y\in Y$, se tiene por tanto $f^{-1}(y)=\bigcap(X_\lambda\cap g^{-1}(y))$ y, como los conjuntos $X_\lambda\cap g^{-1}(y)$ son cerrados en el espacio noetheriano $g^{-1}(y)$, existe para todo $y$ un índice $\lambda(y)$ tal que $f^{-1}(y)=X_{\lambda(y)}\cap g^{-1}(y)$.
Denotemos por $E_\lambda$ el conjunto de los $y\in Y$ tales que la fibra $X_\lambda\cap g^{-1}(y)$ de la restricción de $g$ a $X_\lambda$ es un $k(y)$-preesquema finito.
La hipótesis de que $f$ es cuasifinito implica, en virtud de lo anterior, que $Y=\bigcup_\lambda E_\lambda$.
Ahora bien, cada uno de los $X_\lambda$ es, por definición, \emph{de presentación finita} sobre $Y$; se sigue, pues, de \sref{IV.9.2.3} y \sref{IV.9.2.6})~($^1$) que los $E_\lambda$ son conjuntos \emph{construibles} en el esquema $Y$; como forman una familia creciente filtrante, existe un índice $\lambda$ tal que $E_\lambda=Y$ \sref{IV.1.9.9}), y para este $\lambda$ el morfismo $f_\lambda:X_\lambda\to Y$, restricción de $g$ a $X_\lambda$, es cuasifinito. Como es de presentación finita y separado, se puede aplicar \sref{IV.8.12.6}), y $f_\lambda$ se factoriza como
\[
X_\lambda\xrightarrow{j_\lambda}Y'_\lambda\xrightarrow{u_\lambda}Y
\]
donde $j_\lambda$ es una inmersión abierta y $u_\lambda$ un morfismo finito.
Como $X$ es un subpreesquema cerrado de $X_\lambda$, hemos demostrado así que $f$ posee la propiedad \sref{IV.8.12.3},~$a$)), de donde se sigue el corolario en virtud de la equivalencia de \sref{IV.8.12.3},~$a$)) y \sref{IV.8.12.3},~$a'$)).
\end{proof}

\footnotetext{($^1$) El lector comprobará que los corolarios \sref{IV.8.12.8} a \sref{IV.8.12.11} no se utilizan en la \textsection~9.}

\begin{corollary}[8.12.9]
\label{IV.8.12.9}
\oldpage[IV-3]{48}
Sea $f:X\to Y$ un morfismo localmente cuasifinito $(\mathbf{Err}_{\text{III}},~20)$.
Para todo $x\in X$, existen un entorno abierto $U$ de $x$ en $X$, un entorno abierto $V$ de $y=f(x)$ en $Y$, tales que $f(U)\subset V$, y una factorización
\[
U\xrightarrow{f'}V'\xrightarrow{u}V
\]
de la restricción de $f$ a $U$, donde $f'$ es una inmersión abierta y $u$ un morfismo finito.
\end{corollary}

\begin{proof}
Basta evidentemente tomar para $V$ un entorno afín de $y$ en $Y$, y para $U$ un entorno afín de $x$ en $X$, contenido en $f^{-1}(V)$ y tal que $f|U$ sea cuasifinito.
Como el morfismo $U\to V$, restricción de $f$, es entonces afín (y por tanto cuasiproyectivo), se puede aplicar \sref{IV.8.12.8}).
\end{proof}

\begin{corollary}[8.12.10]
\label{IV.8.12.10}
Sea $Y$ un preesquema íntegro normal, sea $X$ un preesquema íntegro y sea $f:X\to Y$ un morfismo birracional, localmente cuasifinito $(\mathbf{Err}_{\text{III}},~20)$.
Entonces $f$ es un isomorfismo local; para que $f$ sea una inmersión abierta, es necesario y suficiente que $f$ sea además separado.
\end{corollary}

\begin{proof}
La segunda aserción se sigue inmediatamente de la primera y de \sref[I]{I.8.2.8}).
Para demostrar la primera aserción, podemos suponer $X$ e $Y$ afines y $f$ cuasifinito; consideremos la factorización $f=u\circ f'$ de \sref{IV.8.12.8}, que identifica $X$, mediante $f'$, con el subpreesquema inducido sobre un abierto de $Y'$.
Como $X$ es íntegro, podemos, en virtud de \sref[I]{I.5.2.3}), sustituir $Y'$ por el subpreesquema reducido de $Y'$ cuyo espacio subyacente es $\overline{X}$; podemos, pues, suponer también que $Y'$ es \emph{íntegro}.
Además, como $f$ es birracional, también lo es $u$.
La conclusión se sigue entonces del lema siguiente:
\end{proof}

\begin{lemma}[8.12.10.1]
\label{IV.8.12.10.1}
Sea $Y'$ un preesquema íntegro y sea $Y$ un preesquema íntegro y normal; entonces todo morfismo $u:Y'\to Y$ finito y birracional es un isomorfismo.
\end{lemma}

\begin{proof}
Pongamos $\sh{A}=u_*(\sh{O}_{Y'})$; entonces $\sh{A}$ es una $\sh{O}_Y$-álgebra finita y $Y'$ se identifica con $\operatorname{Spec}(\sh{A})$ \sref[II]{II.1.3.6}.
Si $R(Y)$ es el cuerpo de funciones racionales de $Y$, se tiene, por tanto, para todo $y\in Y$, $\sh{O}_{Y,y}\subset\sh{A}_y\subset R(Y)$; pero, como el anillo $\sh{O}_{Y,y}$ es por hipótesis íntegramente cerrado y $R(Y)$ un cuerpo de fracciones, se tiene necesariamente $\sh{A}_y=\sh{O}_{Y,y}$, de donde el lema.
\end{proof}

\begin{corollary}[8.12.11]
\label{IV.8.12.11}
Sea $Y$ un preesquema íntegro, sea $X$ un preesquema íntegro normal y sea $f:X\to Y$ un morfismo dominante, localmente cuasifinito.
Sean $K$ y $L$ (una extensión de $K$) los cuerpos de funciones racionales de $Y$ y $X$, respectivamente; y sea $Y'$ la clausura entera de $Y$ con respecto a $L$ \sref[II]{II.6.3.4}); entonces $f$ se factoriza de manera única como $f:X\xrightarrow{f'}Y'\xrightarrow{u}Y$, donde $f'$ es birracional y corresponde al automorfismo identidad de $L$; $f'$ es entonces un isomorfismo local, y para que $f'$ sea una inmersión abierta, es necesario y suficiente que $f$ sea separado.
\end{corollary}

\begin{proof}
La existencia y la unicidad de la factorización de $f$ se siguen de \sref[II]{II.6.3.9}).
Se sigue además de \sref[II]{II.6.2.4},~(v)), reduciéndose al caso afín, que $f'$ es localmente cuasifinito; por otra parte, se sigue de \sref[I]{I.5.5.1}) que, para que $f'$ sea separado, es necesario y suficiente que lo sea $f$, ya que $u$ es afín, y por tanto separado; las dos últimas aserciones son, pues, consecuencias de \sref{IV.8.12.10}) aplicado a $f'$.
\end{proof}

\oldpage[IV-3]{49}

\subsection{Traducción en términos de pro-objetos}
\label{subsection:IV.8.13}

La proposición siguiente es esencialmente equivalente a \sref{IV.8.8.2},~(i):

\begin{proposition}[8.13.1]
\label{IV.8.13.1}
Sea $S$ un preesquema y sea $(X_\lambda,v_{\lambda\mu})$ un sistema proyectivo filtrante de $S$-preesquemas; supongamos que existe $\alpha$ tal que $v_{\alpha\lambda}$ es un morfismo afín para todo $\lambda\geqslant\alpha$ (lo que implica por \sref[II]{II.1.6.2} que $v_{\lambda\mu}$ es afín para $\alpha\leqslant\lambda\leqslant\mu$), de modo que el límite proyectivo $X=\varprojlim_\lambda X_\lambda$ existe en la categoría de los $S$-preesquemas \sref{IV.8.2.3}.
Sea $Y$ un $S$-preesquema y, para todo $\lambda\geqslant\alpha$, sea $e_\lambda:\operatorname{Hom}_S(X_\lambda,Y)\to\operatorname{Hom}_S(X,Y)$ la aplicación que a cada $S$-morfismo $f_\lambda:X_\lambda\to Y$ hace corresponder $f=f_\lambda\circ v_\lambda$, donde $v_\lambda:X\to X_\lambda$ es el morfismo canónico.
La familia $(e_\lambda)$ es un sistema inductivo de aplicaciones, que define por tanto una aplicación canónica
\begin{equation}
\label{IV.8.13.1.1}
\varinjlim_\lambda\operatorname{Hom}_S(X_\lambda,Y)\to\operatorname{Hom}_S(X,Y). \tag{8.13.1.1}
\end{equation}

Supongamos que $X_\alpha$ es cuasicompacto (resp. cuasicompacto y cuasiseparado), y que el morfismo estructural $Y\to S$ es localmente de tipo finito (resp. localmente de presentación finita).
Entonces la aplicación \sref{IV.8.13.1.1} es inyectiva (resp. biyectiva).
\end{proposition}

\begin{proof}
Para $\lambda\geqslant\alpha$, pongamos $Z_\lambda=Y\times_S X_\lambda$, de modo que $Z_\lambda=Z_\alpha\times_{X_\alpha}X_\lambda$.
Asimismo, pongamos $Z=Y\times_S X=Z_\alpha\times_{X_\alpha}X$; se sigue entonces de \sref{IV.8.2.5} que, si se pone además $w_{\lambda\mu}=1\times v_{\lambda\mu}:Z_\mu\to Z_\lambda$ para $\alpha\leqslant\lambda\leqslant\mu$ y $w_\lambda=1\times v_\lambda:Z\to Z_\lambda$ para $\alpha\leqslant\lambda$, entonces $Z$ es el límite proyectivo del sistema $(Z_\lambda,w_{\lambda\mu})$, y los $w_\lambda$ son los morfismos canónicos correspondientes.
Además, el morfismo $Z_\alpha\to X_\alpha$ es localmente de tipo finito (resp. localmente de presentación finita) por \sref{IV.1.3.4} y \sref{IV.1.4.3}.
Por último, por \sref[I]{I.3.3.14}, se tiene
\[
\operatorname{Hom}_S(X_\lambda,Y)=\operatorname{Hom}_{X_\lambda}(X_\lambda,Z_\lambda) \quad\text{y}\quad \operatorname{Hom}_S(X,Y)=\operatorname{Hom}_X(X,Z)
\]
Basta ahora aplicar \sref{IV.8.8.2},~(i), tomando $X_\lambda=S_\lambda$ y sustituyendo $Y_\lambda$ por $Z_\lambda$.
\end{proof}

\begin{corollary}[8.13.2]
\label{IV.8.13.2}
Con la notación de \sref{IV.8.13.1}, supongamos que $X_\alpha$ es cuasicompacto y cuasiseparado, y que los $v_{\alpha\lambda}$ son afines para $\alpha\leqslant\lambda$; supongamos además que $Y=\varprojlim_\rho Y_\rho$, donde $(Y_\rho,t_{\rho\sigma})$ es un sistema proyectivo filtrante de $S$-preesquemas tal que, para cada $\rho$, el morfismo estructural $Y_\rho\to S$ es localmente de presentación finita.
Entonces se tiene una biyección canónica
\begin{equation}
\label{IV.8.13.2.1}
\operatorname{Hom}_S(X,Y)\xrightarrow{\sim}\varprojlim_\rho\left(\varinjlim_\lambda\operatorname{Hom}_S(X_\lambda,Y_\rho)\right). \tag{8.13.2.1}
\end{equation}
\end{corollary}

\begin{proof}
En efecto, como $Y$ es el límite proyectivo de los $Y_\rho$, la aplicación canónica $\operatorname{Hom}_S(X,Y)\to\varprojlim_\rho\operatorname{Hom}_S(X,Y_\rho)$ es biyectiva; por otra parte, las hipótesis proporcionan, para cada $\rho$, una biyección canónica $\operatorname{Hom}_S(X,Y_\rho)\xrightarrow{\sim}\varinjlim_\lambda\operatorname{Hom}_S(X_\lambda,Y_\rho)$ por \sref{IV.8.13.1}; de donde la conclusión.
\end{proof}

\begin{env}[8.13.3]
\label{IV.8.13.3}
Los resultados anteriores permiten interpretar en la teoría de preesquemas las nociones de ``provariedad'' y ``proesquema'' que aparecen en ciertas aplicaciones (por ejemplo, en la teoría local de cuerpos de clases siguiendo las ideas de Serre~[39], o en la teoría de N\'eron de la reducción de
\oldpage[IV-3]{50}
variedades abelianas~[32]).
Recordemos aquí brevemente la noción de \emph{pro-objeto} de una categoría, remitiendo al cap.~V para desarrollos más amplios (por lo demás, no utilizaremos antes del cap.~V la interpretación que sigue, y el lector puede por tanto omitirla hasta el final de este número).
Dada una categoría $\mathcal{C}$, la categoría $\mathbf{Pro}(\mathcal{C})$ de los \emph{pro-objetos} de $\mathcal{C}$ tiene por objetos los \emph{sistemas proyectivos} (en el universo considerado) $\mathbf{X}=(X_\mu)_{\mu\in M}$ de objetos de $\mathcal{C}$ cuyos conjuntos de índices (que dependen del sistema proyectivo) se suponen conjuntos preordenados filtrantes.
Dados dos pro-objetos de este tipo $\mathbf{X}=(X_\mu)_{\mu\in M}$ y $\mathbf{X}'=(X'_{\mu'})_{\mu'\in M'}$, los \emph{morfismos} de $\mathbf{X}$ en $\mathbf{X}'$ son, por definición, los elementos del conjunto $\varprojlim_{\mu'}(\varinjlim_\mu\operatorname{Hom}(X_\mu,X'_{\mu'}))$; es inmediato que estos conjuntos pueden servir como conjuntos de morfismos. La composición de sistemas de morfismos $u_{\mu'}^\mu:X_\mu\to X'_{\mu'}$ y $u_{\mu''}^{\prime\mu'}:X'_{\mu'}\to X''_{\mu''}$, inductivos respecto del índice superior y proyectivos respecto del índice inferior, se efectúa ``argumento por argumento'', es decir, considerando el sistema $u_{\mu''}^{\prime\mu'}\circ u_{\mu'}^\mu$.
\end{env}

\begin{env}[8.13.4]
\label{IV.8.13.4}
Sea $S$ un preesquema cuasicompacto y cuasiseparado, y denotemos por $\mathcal{C}$ la subcategoría plena de la categoría $(\mathbf{Sch})_{/S}$ de los $S$-preesquemas, formada por los $S$-preesquemas $X$ que poseen la propiedad siguiente: el morfismo estructural $X\to S$ se factoriza como $X\xrightarrow{g}X_0\xrightarrow{f}S$, donde $g:X\to X_0$ es afín y $f:X_0\to S$ es \emph{de presentación finita}; diremos abreviadamente que los preesquemas de $\mathcal{C}$ son \emph{esencialmente afines sobre} $S$.
Consideremos, por otra parte, la subcategoría plena $\mathcal{C}'_0$ de $(\mathbf{Sch})_{/S}$ formada por los $S$-preesquemas \emph{de presentación finita}, y la categoría $\mathbf{Pro}(\mathcal{C}'_0)$ de pro-objetos de $\mathcal{C}'_0$.
Diremos que un objeto $\mathbf{X}=(X_\mu)_{\mu\in M}$ de $\mathbf{Pro}(\mathcal{C}'_0)$ es \emph{esencialmente afín} si existe $\gamma\in M$ tal que, para $\mu\geqslant\gamma$, el morfismo de transición $v_{\gamma\mu}:X_\mu\to X_\gamma$ es afín (lo que implica que, para $\gamma\leqslant\nu\leqslant\mu$, $v_{\nu\mu}:X_\mu\to X_\nu$ es afín).
Obsérvese que un objeto de $\mathbf{Pro}(\mathcal{C}'_0)$ isomorfo a un objeto esencialmente afín \emph{no es} necesariamente él mismo esencialmente afín.
Denotaremos por $\mathcal{C}'$ la subcategoría plena de $\mathbf{Pro}(\mathcal{C}'_0)$ formada por los pro-objetos esencialmente afines.

Dicho esto, se sigue de \sref{IV.8.2.2} y \sref{IV.8.2.3} que, para todo objeto $\mathbf{X}=(X_\mu)_{\mu\in M}$ de $\mathcal{C}'$, existe el $S$-preesquema $X=\varprojlim_\mu X_\mu$; además, como el morfismo $X\to X_\mu$ es afín para todo $\mu$ suficientemente grande \sref{IV.8.2.2}, $X$ es, por definición, \emph{esencialmente afín sobre} $S$.
Pongamos $X=L(\mathbf{X})$; demostremos que se ha definido así un \emph{funtor} canónico
\begin{equation}
\label{IV.8.13.4.1}
L:\mathcal{C}'\to\mathcal{C} \tag{8.13.4.1}
\end{equation}

Para dos objetos $\mathbf{X}=(X_\mu)$ y $\mathbf{X}'=(X'_{\mu'})$ de $\mathcal{C}'$, se tiene, para cada $\mu'$, una aplicación canónica
\[
\varinjlim_\mu\operatorname{Hom}_S(X_\mu,X'_{\mu'})\to\operatorname{Hom}_S(\varprojlim_\mu X_\mu,X'_{\mu'})
\]
definida en \sref{IV.8.13.1.1}; por otra parte, por definición del límite proyectivo, existe una biyección canónica
\[
\varprojlim_{\mu'}\operatorname{Hom}_S(\varprojlim_\mu X_\mu,X'_{\mu'})\xrightarrow{\sim}\operatorname{Hom}_S(\varprojlim_\mu X_\mu,\varprojlim_{\mu'}X'_{\mu'})
\]
\oldpage[IV-3]{51}

de donde se obtiene una aplicación canónica
\begin{equation}
\label{IV.8.13.4.2}
\varprojlim_{\mu'}\varinjlim_\mu\operatorname{Hom}_S(X_\mu,X'_{\mu'})\to\operatorname{Hom}_S(\varprojlim_\mu X_\mu,\varprojlim_{\mu'}X'_{\mu'}) \tag{8.13.4.2}
\end{equation}
evidentemente funtorial en $\mathbf{X}$ y $\mathbf{X}'$, y que completa la definición del funtor $L$.
\end{env}

\begin{proposition}[8.13.5]
\label{IV.8.13.5}
Bajo las hipótesis y con la notación de \sref{IV.8.13.4}, el funtor $L$ es plenamente fiel.
Si además $S$ es un preesquema noetheriano (lo que implica ya que $S$ es cuasicompacto y cuasiseparado \sref{IV.1.2.8}), $L$ es \emph{una equivalencia de categorías}.
\end{proposition}

\begin{proof}
Decir que $L$ es plenamente fiel significa que la aplicación \sref{IV.8.13.4.2} es biyectiva, lo que es un caso particular de \sref{IV.8.13.2}.
En efecto, como los morfismos estructurales $X_\mu\to S$ son de presentación finita, son en particular cuasicompactos y cuasiseparados; por tanto, los $X_\mu$ son cuasicompactos y cuasiseparados.

Para demostrar que, cuando $S$ es noetheriano, $L$ es una equivalencia de categorías, basta, puesto que ya sabemos que $L$ es plenamente fiel, probar que todo $S$-preesquema esencialmente afín $X$ es $S$-\emph{isomorfo} a un objeto de la forma $L(\mathbf{X})$, donde $\mathbf{X}\in\mathcal{C}'$ $(\mathbf{0}_{\text{III}},~8.1.5)$.
Ahora bien, por hipótesis existe una factorización $X\xrightarrow{g}X_0\xrightarrow{f}S$ del morfismo estructural, siendo $f$ de presentación finita y $g$ afín.
Podemos escribir, por tanto, $X=\operatorname{Spec}(\sh{A})$, donde $\sh{A}$ es una $\sh{O}_{X_0}$-Álgebra cuasicoherente \sref[II]{II.1.3.1}.
Ahora bien, como $X_0$ es noetheriano (pues es de tipo finito sobre $S$ y $S$ es noetheriano), $\sh{A}$ es el límite inductivo filtrante de la familia $(\sh{A}_\lambda)$ de sus sub-$\sh{O}_{X_0}$-Álgebras cuasicoherentes de tipo finito \sref[I]{I.9.6.6}.
Pongamos $X_\lambda=\operatorname{Spec}(\sh{A}_\lambda)$; los morfismos $X_\lambda\to X_0$ son de tipo finito, y por tanto de presentación finita puesto que $X_0$ es noetheriano; en consecuencia, lo mismo ocurre con los morfismos compuestos $X_\lambda\to X_0\to S$ \sref{IV.1.6.2}; dicho de otro modo, los $X_\lambda$ pertenecen a $\mathcal{C}'_0$, y como los morfismos $X_\lambda\to X_0$ son afines, $\mathbf{X}=(X_\lambda)$ es un objeto de $\mathcal{C}'$ cuyo límite proyectivo existe y es $S$-isomorfo a $X$ por \sref{IV.8.2.2}.
Esto termina la demostración.
\end{proof}

\begin{remark}[8.13.6]
\label{IV.8.13.6}
Se sigue de \sref{IV.1.6.2} y \sref[II]{II.1.6.2} que, si $X$ e $Y$ son esencialmente afines sobre $S$, lo mismo ocurre con $X\times_S Y$.
Se concluye, por ejemplo, $(\mathbf{0}_{\text{III}},~8.2.5)$ que un $\mathcal{C}$-grupo no es otra cosa que un $(\mathbf{Sch})_{/S}$-grupo que es un preesquema esencialmente afín sobre $S$.
Por otra parte, los productos finitos existen en la categoría $\mathcal{C}'$: en efecto, si $\mathbf{X}=(X_\mu)_{\mu\in M}$, $\mathbf{Y}=(Y_\rho)_{\rho\in R}$ son dos objetos de $\mathcal{C}'$, los productos $X_\mu\times_S Y_\rho$ son $S$-preesquemas de presentación finita y, tomando como morfismos de transición $X_\nu\times_S Y_\sigma\to X_\mu\times_S Y_\rho$ los productos de los morfismos de transición $X_\nu\to X_\mu$ e $Y_\sigma\to Y_\rho$, se ve inmediatamente que $(X_\mu\times_S Y_\rho)$ es un producto de $\mathbf{X}$ y $\mathbf{Y}$ en $\mathbf{Pro}(\mathcal{C}'_0)$; además, por \sref[II]{II.1.6.2}, los morfismos de transición así definidos son afines para $\mu$ y $\rho$ suficientemente grandes, y por tanto el producto $\mathbf{X}\times\mathbf{Y}$ pertenece también a $\mathcal{C}'$.
Se concluye entonces, como antes, que un $\mathcal{C}'$-grupo es un $\mathbf{Pro}(\mathcal{C}'_0)$-grupo que es esencialmente afín.
Se deduce de \sref{IV.8.13.5} que las categorías de los $\mathcal{C}'$-grupos y de los $\mathcal{C}$-grupos son \emph{equivalentes} cuando $S$ es \emph{noetheriano}.
Parece plausible que, cuando $S$ es el espectro de un cuerpo $k$, la categoría de los $\mathcal{C}$-grupos sea \emph{equivalente} a la de los $\boldsymbol{K}$-\emph{grupos}, donde $\boldsymbol{K}$ es la categoría de los $S$-preesquemas \emph{cuasicompactos}; dicho de otro modo, todo preesquema en grupos cuasicompacto sobre $k$ sería esencialmente afín.
Por otra parte, si se denota
\oldpage[IV-3]{52}
por $\mathcal{C}'_0$-$\boldsymbol{gr}$ la categoría de los $\mathcal{C}'_0$-grupos, es muy probable que la categoría de los $\mathcal{C}'$-grupos sea equivalente a la subcategoría plena de $\mathbf{Pro}(\mathcal{C}'_0\text{-}\boldsymbol{gr})$ formada por los ``grupos proalgebraicos esencialmente afines'', es decir, los sistemas proyectivos $\mathbf{G}=(G_\mu)_{\mu\in M}$, donde los $G_\mu$ son grupos algebraicos sobre $k$ y los morfismos de transición $G_\nu\to G_\mu$ son afines para $\mu$ suficientemente grande (lo que también puede expresarse diciendo que $\mathbf{G}$ es una extensión de un grupo algebraico por un grupo proalgebraico afín).
La conjunción de estas dos conjeturas es además equivalente a la afirmación siguiente: todo preesquema en grupos cuasicompacto sobre $k$ es una extensión de un ``grupo algebraico'' (es decir, un preesquema en grupos de tipo finito sobre $k$) por un preesquema en grupos \emph{afín} sobre $k$.
Como los únicos grupos proalgebraicos encontrados hasta ahora en la práctica son de hecho esencialmente afines, sería ventajoso sustituir el estudio de los grupos proalgebraicos generales (introducidos y estudiados por Serre~[40]) por el de los esquemas en grupos cuasicompactos sobre $k$, cuya definición es conceptualmente mucho más sencilla.
\end{remark}

\subsection{Caracterización de un preesquema localmente de presentación finita sobre otro mediante el funtor que representa}
\label{subsection:IV.8.14}

\begin{env}[8.14.1]
\label{IV.8.14.1}
Dado un preesquema $S$, diremos, como en \sref{IV.8.13.4}, que un sistema proyectivo filtrante $(X_\lambda,v_{\lambda\mu})$ de $S$-preesquemas es \emph{esencialmente afín} si existe $\alpha$ tal que $v_{\alpha\lambda}$ es un morfismo afín para $\lambda\geqslant\alpha$.

El enunciado siguiente, que será especialmente útil en el cap.~V, precisa \sref{IV.8.8.2},~(i) proporcionando un recíproco:
\end{env}

\begin{proposition}[8.14.2]
\label{IV.8.14.2}
Sea $S$ un preesquema y sea $f:X\to S$ un morfismo.
Para todo $S$-preesquema $T$, pongamos
\[
h_X(T)=\operatorname{Hom}_S(T,X)
\]
de modo que $h_X$ es un funtor contravariante de la categoría $(\mathbf{Sch})_{/S}$ de los $S$-preesquemas a la categoría $\mathbf{Ens}$ de los conjuntos $(\mathbf{0}_{\text{III}},~8.1.1)$, y $X$ es un objeto que representa este funtor $(\mathbf{0}_{\text{III}},~8.1.8)$.
Las condiciones siguientes son equivalentes:
\begin{enumerate}[label=\emph{\alph*)}]
  \item $f$ es localmente de presentación finita.
  \item Para todo sistema proyectivo filtrante $(Z_\lambda)$ de $S$-preesquemas, esencialmente afín \sref{IV.8.13.4} y formado por preesquemas cuasicompactos y cuasiseparados, la aplicación canónica \sref{IV.8.13.1.1}
\begin{equation}
\label{IV.8.14.2.1}
\varinjlim_\lambda h_X(Z_\lambda)\to h_X(\varprojlim_\lambda Z_\lambda) \tag{8.14.2.1}
\end{equation}
es biyectiva.
  \item Para todo sistema proyectivo filtrante $(Z_\lambda)$ de $S$-preesquemas tal que los $Z_\lambda$ sean esquemas afines, la aplicación \sref{IV.8.14.2.1} es biyectiva.
  \item[\rm c$'$)] Para todo abierto afín $U$ de $S$ y todo sistema proyectivo filtrante $(Z_\lambda)$ de $U$-preesquemas tal que los $Z_\lambda$ sean esquemas afines, la aplicación \sref{IV.8.14.2.1} es biyectiva.
\end{enumerate}
\end{proposition}

\begin{proof}
El hecho de que \emph{a)} entrañe \emph{b)} es precisamente \sref{IV.8.13.1}; es trivial que \emph{b)} entraña \emph{c)} y que \emph{c)} entraña \emph{c$'$)}.
Queda ver que \emph{c$'$)} entraña \emph{a)} y, como la propiedad \emph{a)} es local sobre $S$, podemos reducirnos al caso en que $S$ es \emph{afín}.
\oldpage[IV-3]{53}

Supongamos primero que $X$ es también afín; la aserción que hay que demostrar es entonces equivalente a la siguiente:

\begin{corollary}[8.14.2.2]
\label{IV.8.14.2.2}
Sea $A$ un anillo y sea $B$ una $A$-álgebra.
Para que, para todo sistema inductivo filtrante $(C_\lambda)$ de $A$-álgebras, la aplicación canónica
\begin{equation}
\label{IV.8.14.2.3}
\varinjlim_\lambda\operatorname{Hom}_{A\text{-}\mathrm{alg.}}(B,C_\lambda)\to\operatorname{Hom}_{A\text{-}\mathrm{alg.}}(B,\varinjlim_\lambda C_\lambda) \tag{8.14.2.3}
\end{equation}
sea biyectiva, es necesario y suficiente que $B$ sea una $A$-álgebra de presentación finita.
\end{corollary}

Solo queda demostrar que la condición es necesaria.
Tomemos primero para $(C_\lambda)$ el sistema inductivo filtrante de las sub-$A$-álgebras de tipo finito de $B$, de modo que $\varinjlim_\lambda C_\lambda=B$.
El hecho de que \sref{IV.8.14.2.3} sea biyectiva implica, en particular, que la aplicación identidad $1_B:B\to B$ se factoriza como $B\to C_\lambda\to B$ para un $\lambda$ adecuado, lo que implica que $C_\lambda=B$; por tanto, $B$ es una $A$-álgebra \emph{de tipo finito}.
Escribamos $B=C/\mathfrak{J}$, donde $C=A[T_1,\ldots,T_n]$ y $\mathfrak{J}$ es un ideal de $C$. El ideal $\mathfrak{J}$ es el límite inductivo filtrante de sus ideales de tipo finito $\mathfrak{J}_\lambda\subset\mathfrak{J}$; poniendo $C_\lambda=C/\mathfrak{J}_\lambda$ y utilizando la exactitud de los límites inductivos filtrantes, se ve que $B$ es nuevamente isomorfo al límite inductivo del sistema inductivo filtrante $(C_\lambda)$.
Existen, por tanto, $\lambda$ y un $A$-homomorfismo $u:B\to C_\lambda$ tales que el compuesto $B\xrightarrow{u}C_\lambda\xrightarrow{p_\lambda}B$ (donde $p_\lambda$ es el homomorfismo canónico) es la identidad.
Sea $q_\lambda:C\to C_\lambda$ el homomorfismo canónico, y pongamos
$t_i=p_\lambda(q_\lambda(T_i))$.
Entonces $p_\lambda(u(t_i))=p_\lambda(q_\lambda(T_i))$, es decir,
$u(t_i)-q_\lambda(T_i)\in\mathfrak{J}/\mathfrak{J}_\lambda$.
Existe, pues, $\mu\geqslant\lambda$ tal que estos $n$ elementos pertenecen a
$\mathfrak{J}_\mu/\mathfrak{J}_\lambda$ ($1\leqslant i\leqslant n$).
Si $p_{\mu\lambda}:C_\lambda\to C_\mu$ es el homomorfismo canónico, entonces
$p_{\mu\lambda}(u(t_i))=p_{\mu\lambda}(q_\lambda(T_i))=q_\mu(T_i)$.
Sustituyendo $\lambda$ por $\mu$ y $u$ por $p_{\mu\lambda}\circ u$, podemos por tanto
suponer que $u(t_i)=q_\lambda(T_i)$ para todo $i$; y si
$r=p_\lambda\circ q_\lambda$ es el homomorfismo canónico
$C\to C/\mathfrak{J}=B$, podemos escribir $u(r(T_i))=q_\lambda(T_i)$ para todo $i$,
de donde $q_\lambda=u\circ r$.
Esto implica necesariamente que $\mathfrak{J}=\mathfrak{J}_\lambda$: si $z\in\mathfrak{J}$, entonces $r(z)=0$, y por tanto $q_\lambda(z)=u(r(z))=0$ y $z\in\mathfrak{J}_\lambda$. Así, $B=C_\lambda$.

Pasemos ahora al caso en que $S$ es afín y $X$ arbitrario; lo que hay que demostrar es que todo abierto afín $V$ de $X$ es de presentación finita sobre $S$, y en virtud de lo que acabamos de demostrar, basta probar que, para todo sistema proyectivo filtrante $(Z_\lambda)$ de $S$-preesquemas afines, la aplicación
\begin{equation}
\label{IV.8.14.2.4}
\varinjlim_\lambda\operatorname{Hom}_S(Z_\lambda,V)\to\operatorname{Hom}_S(\varprojlim_\lambda Z_\lambda,V) \tag{8.14.2.4}
\end{equation}
es biyectiva.
Es inmediato que esta aplicación es inyectiva, pues si $(v_\lambda)$, $(v'_\lambda)$ son dos sistemas inductivos de $S$-homomorfismos $v_\lambda:Z_\lambda\to V$, $v'_\lambda:Z_\lambda\to V$ tales que los morfismos correspondientes
\[
Z\xrightarrow{u_\lambda}Z_\lambda\xrightarrow{v_\lambda}V,\qquad
Z\xrightarrow{u_\lambda}Z_\lambda\xrightarrow{v'_\lambda}V
\]
son iguales (siendo $u_\lambda$ el morfismo canónico), entonces los morfismos
\[
Z\xrightarrow{u_\lambda}Z_\lambda\xrightarrow{v_\lambda}V\xrightarrow{j}X,\qquad Z\xrightarrow{u_\lambda}Z_\lambda\xrightarrow{v'_\lambda}V\xrightarrow{j}X
\]
(donde $j$ es la inyección canónica) son iguales, lo que implica por hipótesis $j\circ v_\lambda=j\circ v'_\lambda$ para un $\lambda$ adecuado, y por tanto $v_\lambda=v'_\lambda$.

\oldpage[IV-3]{54}

Queda demostrar que \sref{IV.8.14.2.4} es sobreyectiva.
Sea $v:Z\to V$ un $S$-morfismo.
Por hipótesis, existen un índice $\lambda$ y un $S$-morfismo
$w_\lambda:Z_\lambda\to X$ tales que $j\circ v$ se factoriza como
\[
 Z\xrightarrow{u_\lambda}Z_\lambda\xrightarrow{w_\lambda}X.
\]
Basta demostrar que existe $\mu\geqslant\lambda$ tal que
\[
 Z_\mu\xrightarrow{u_{\lambda\mu}}Z_\lambda\xrightarrow{w_\lambda}X
\]
(donde $u_{\lambda\mu}$ es el morfismo de transición) se factoriza como
\[
 Z_\mu\xrightarrow{v_\mu}V\xrightarrow{j}X.
\]
Para todo $\mu\geqslant\lambda$, pongamos
$U_\mu=u_{\lambda\mu}^{-1}(w_\lambda^{-1}(V))$.
Entonces
\[
 u_\mu^{-1}(U_\mu)=u_\lambda^{-1}(U_\lambda)
 =u_\lambda^{-1}(w_\lambda^{-1}(V))
 =v^{-1}(V)=Z.
\]
Los $Z_\lambda$ son cuasicompactos y los $U_\mu$, por ser abiertos, son
ind-construibles \sref{IV.1.9.6}; por tanto, \sref{IV.8.3.4} proporciona un
$\mu\geqslant\lambda$ para el cual $U_\mu=Z_\mu$.
\end{proof}

\begin{remark}[8.14.3]
\label{IV.8.14.3}
El hecho de que la aplicación \sref{IV.8.14.2.1} sea inyectiva cuando $f$ es localmente de tipo
finito \sref{IV.8.8.2},~(i) conduce naturalmente a preguntarse si este resultado admite un recíproco.
No es así, ni siquiera cuando $S$ y $X$ son afines, pues existen monomorfismos
$X\to S$ que no son de tipo finito \sref[I]{I.2.4.2}, y estos refutan tal
conjetura.
\end{remark}

\section{Propiedades construibles}
\label{section:IV.9}

Sea $S$ un preesquema, sea $f:X\to S$ un morfismo de \emph{presentación finita}~\sref{IV.1.6.1},
y sea $\sh{F}$ un $\sh{O}_X$-Módulo cuasicoherente de \emph{presentación finita}.
Nos proponemos en este parágrafo dar criterios que aseguren, por ejemplo, que el conjunto de los $s\in S$ tales que el $\mathbf{k}(s)$-preesquema $X_s=f^{-1}(s)=X\times_S\operatorname{Spec}(\mathbf{k}(s))$ posee cierta propiedad, o tales que el $\sh{O}_{X_s}$-Módulo $\sh{F}_s=\sh{F}\otimes_{\sh{O}_X}\mathbf{k}(s)$ posee cierta propiedad, es \emph{localmente construible}, o al menos \emph{ind-construible}~\sref{IV.1.9.4}; veremos que este es el caso para la mayoría de las propiedades que aparecen en Geometría Algebraica.
Suponiendo únicamente que $f$ es \emph{localmente de presentación finita}, daremos también en~{\hyperlink{ega4.c04.subsec.9-9}{(9.9)}} criterios para que el conjunto de puntos $x\in X$ en los que la fibra $X_{f(x)}$ (o el $\sh{O}_{X_{f(x)}}$-Módulo $\sh{F}_{f(x)}$) posee cierta propiedad sea \emph{localmente construible}.
Veremos en \textsection12 que estos resultados, junto con la hipótesis adicional de que $f$ sea \emph{plano} (resp.\ \emph{propio y plano}), permiten demostrar que los conjuntos considerados en $X$ (resp.\ en $S$) son incluso \emph{abiertos} en numerosos casos.

\subsection{El principio de extensión finita}
\label{subsection:IV.9.1}

\begin{proposition}[9.1.1]
\label{IV.9.1.1}
\textnormal{(Principio de extensión finita).}
Sea $k$ un cuerpo y sea $\mathcal{C}$ un conjunto de extensiones de $k$.
Supongamos que se satisfacen las condiciones siguientes:
\begin{enumerate}[label=(\roman*)]
  \item Si $K\in\mathcal{C}$ y existe un $k$-homomorfismo $K\to K'$ (donde $K'$ pertenece al universo en el que nos hemos situado), entonces $K'\in\mathcal{C}$.
  \item Si $K\in\mathcal{C}$, existe una subextensión $K'$ de $K$, de tipo finito sobre $k$, tal que $K'\in\mathcal{C}$.
  \item Si $K\in\mathcal{C}$ es el cuerpo de fracciones de una $k$-álgebra íntegra de tipo finito $A$, existe $f\in A-\{0\}$ tal que, para todo ideal maximal $\mathfrak{m}$ de $A$, se tiene $A_f/\mathfrak{m}\in\mathcal{C}$.
\end{enumerate}

\oldpage[IV-3]{55}

Sea entonces $\Omega$ una extensión algebraicamente cerrada de $k$ (en el universo considerado).
Las condiciones siguientes son equivalentes:
\begin{enumerate}[label=\emph{\alph*)}]
  \item $\mathcal{C}$ no es vacío.
  \item Existe $K'\in\mathcal{C}$ que es una extensión finita de $k$.
  \item Se tiene $\Omega\in\mathcal{C}$.
\end{enumerate}
\end{proposition}

La condición~(i) implica evidentemente que \emph{b)} entraña \emph{c)}, y \emph{c)} entraña trivialmente \emph{a)};
demostremos que \emph{a)} entraña \emph{b)}.
En virtud de~(ii) y~(iii), existe una extensión $K'\in\mathcal{C}$ de la forma $A/\mathfrak{m}$, donde $A$ es una $k$-álgebra de tipo finito sobre $k$ y $\mathfrak{m}$ es un ideal maximal de $A$.
Sabemos, por el Nullstellensatz de Hilbert (Bourbaki, \emph{Alg.\ comm.}, cap.~V, \textsection3, n\textsuperscript{o}~1, cor.~2 del teo.~1), que $K'$ es una extensión finita de $k$.

\begin{corollary}[9.1.2]
\label{IV.9.1.2}
Bajo las hipótesis~\emph{(i)}, \emph{(ii)}, \emph{(iii)} de~\sref{IV.9.1.1}, si $k$ es algebraicamente cerrado y $\mathcal{C}$ no es vacío, $\mathcal{C}$ contiene todas las extensiones de $k$ del universo considerado.
\end{corollary}

En efecto, como \emph{a)} entraña \emph{c)}, se tiene $k\in\mathcal{C}$, y la conclusión se sigue de la hipótesis~(i).

\begin{remark}[9.1.3]
\label{IV.9.1.3}
En la práctica se verificará la condición~(ii) de~\sref{IV.9.1.1} observando que $K$ es el límite directo de sus subextensiones de tipo finito $K_\alpha$, y aplicando los resultados de \textsection8, teniendo en cuenta eventualmente que, para $K_\alpha\subset K_\beta$, $K_\beta$ es fielmente plana sobre $K_\alpha$.
Con frecuencia, el conjunto $\mathcal{C}$ está formado por \emph{cuerpos} que pertenecen a un conjunto $\mathcal{C}'$ de $k$-álgebras que satisface la condición siguiente:
\begin{enumerate}[label=(\roman*~bis)]
  \item Si $A\in\mathcal{C}'$ y existe un homomorfismo de $k$-álgebras $A\to A'$ (donde $A'$ pertenece al universo en el que nos hemos situado), entonces $A'\in\mathcal{C}'$.
\end{enumerate}
\end{remark}

En tal caso, la condición~(i) se satisface trivialmente, y se verificará en general la condición~(iii) de~\sref{IV.9.1.1} observando que el cuerpo de fracciones $K$ de $A$ es el límite directo de las álgebras $A_f$ (para la relación $D(f)\supset D(g)$), y aplicando los resultados de \textsection8, teniendo en cuenta eventualmente que los morfismos $D(g)\to D(f)$ son inmersiones abiertas.

Por otra parte, cuando no se satisface (i~bis), la demostración de~(iii) suele ser más delicada y está ligada a criterios de construibilidad que se desarrollarán más adelante.

He aquí algunos ejemplos típicos de aplicación del principio de extensión finita:

\begin{proposition}[9.1.4]
\label{IV.9.1.4}
Sea $k$ un cuerpo, sea $\Omega$ una extensión algebraicamente cerrada de $k$, y sean $X$ e $Y$ dos preesquemas de tipo finito sobre $k$.
Las condiciones siguientes son equivalentes:
\begin{enumerate}[label=\emph{\alph*)}]
  \item Existe un $\Omega$-morfismo $X_{(\Omega)}\to Y_{(\Omega)}$ \textnormal{(}resp.\ un $\Omega$-morfismo que posee una de las propiedades~\emph{(i)} a~\emph{(xiv)} especificada de~\sref{IV.8.10.5}\textnormal{)}.
  \item Existe una extensión finita $k'$ de $k$ y un $k'$-morfismo $X_{(k')}\to Y_{(k')}$ \textnormal{(}resp.\ un $k'$-morfismo que posee la propiedad considerada\textnormal{)}.
  \item Existe una extensión $K$ de $k$ y un $K$-morfismo $X_{(K)}\to Y_{(K)}$ \textnormal{(}resp.\ un $K$-morfismo que posee la propiedad considerada\textnormal{)}.
\end{enumerate}
\end{proposition}

Aplicamos la observación~\sref{IV.9.1.3} tomando como $\mathcal{C}'$ el conjunto de todas las $k$-álgebras $A$ (del universo en el que nos hemos situado) tales que existe un $A$-morfismo $X\otimes_k A\to Y\otimes_k A$ (resp.\ un morfismo que posee la propiedad particular de~\sref{IV.8.10.5} considerada).
La condición~(i~bis) de~\sref{IV.9.1.3} se satisface entonces gracias a que las
\oldpage[IV-3]{56}
propiedades consideradas en~\sref{IV.8.10.5} son todas estables por cambio de base.
La condición~(iii) de~\sref{IV.9.1.1} se satisface así en virtud de~\sref{IV.8.8.2}[i] (resp.\ \sref{IV.8.10.5}), pues $\operatorname{Spec}(k)$ es cuasicompacto y cuasiseparado.
Queda por verificar la condición~(ii) de~\sref{IV.9.1.1}, que se sigue también de~\sref{IV.8.8.2}[i] y~\sref{IV.8.10.5}, considerando $K$ como límite directo de sus subextensiones de tipo finito.
Concluimos, por tanto, mediante~\sref{IV.9.1.1}.

En particular, si existe una extensión $K$ de $k$ y un $K$-isomorfismo $X_{(K)}\xrightarrow{\sim} Y_{(K)}$, decimos que $X$ e $Y$ son \emph{geométricamente isomorfos}.

El corolario siguiente generaliza~(\textbf{II},~6.6.5):

\begin{corollary}[9.1.5]
\label{IV.9.1.5}
Sea $k$ un cuerpo y sea $X$ un $k$-preesquema.
Si existe una extensión $K$ de $k$ tal que $X_{(K)}$ es proyectivo \textnormal{(}resp.\ cuasiproyectivo\textnormal{)} sobre $K$, entonces $X$ es proyectivo \textnormal{(}resp.\ cuasiproyectivo\textnormal{)} sobre $k$.
\end{corollary}

Como el morfismo $\operatorname{Spec}(K)\to\operatorname{Spec}(k)$ es fielmente plano y cuasicompacto, ya se sigue de~\sref{IV.2.7.1}[v] que $X$ es de tipo finito sobre $k$.
La hipótesis significa que existe una inmersión cerrada (resp.\ una inmersión) $X_{(K)}\to\mathbf{P}_K^n=\mathbf{P}_k^n\otimes_k K$ (\textbf{II},~5.5.4,~(ii) y~5.5.3); aplicando~\sref{IV.9.1.4} a la propiedad~(iv) (resp.~(ii)) de~\sref{IV.8.10.5}, deducimos que, para alguna extensión finita $k'$ de $k$, existe una inmersión cerrada (resp.\ una inmersión) $X_{(k')}\to\mathbf{P}_{k'}^n$; en otras palabras, $X_{(k')}$ es proyectivo (resp.\ cuasiproyectivo) sobre $k'$.
Concluimos entonces mediante~(\textbf{II},~6.6.5).

\begin{proposition}[9.1.6]
\label{IV.9.1.6}
Sea $k$ un cuerpo, sea $\Omega$ una extensión algebraicamente cerrada de $k$, sea $X$ un preesquema de tipo finito sobre $k$, y sean $\sh{F}$, $\sh{G}$ dos $\sh{O}_X$-Módulos coherentes.
Supongamos que existe un isomorfismo $\sh{F}\otimes_k\Omega\xrightarrow{\sim}\sh{G}\otimes_k\Omega$.
Entonces existe una extensión finita $k'$ de $k$ y un isomorfismo $\sh{F}\otimes_k k'\xrightarrow{\sim}\sh{G}\otimes_k k'$.
\end{proposition}

El razonamiento es el mismo que en~\sref{IV.9.1.4}, aplicando~\sref{IV.8.5.2}[i] (en la demostración de la propiedad~(iii) de~\sref{IV.9.1.1}, utilizamos aquí el hecho de que los morfismos $D(g)\to D(f)$ (con las notaciones de~\sref{IV.9.1.3}) son inmersiones abiertas y, \emph{a fortiori}, morfismos planos).

\subsection{Propiedades construibles e ind-construibles}
\label{subsection:IV.9.2}
\label{IV.9.2}

\begin{definition}[9.2.1]
\label{IV.9.2.1}
Sea $\mathbf{P}(X,\sh{F},k)$ una relación.
Decimos que $\mathbf{P}$ es una propiedad \emph{construible} (resp.\ \emph{ind-construible}) de los preesquemas algebraicos si se satisfacen las dos condiciones siguientes:
\begin{enumerate}[label=(\roman*)]
  \item Si $k$ es un cuerpo, $X$ un preesquema algebraico sobre $k$, $\sh{F}$ un $\sh{O}_X$-Módulo coherente y $k'$ una extensión de $k$, para que $\mathbf{P}(X,\sh{F},k)$ sea cierta es necesario y suficiente que $\mathbf{P}(X_{(k')},\sh{F}\otimes_k k',k')$ sea cierta.
  \item Sea $S$ un preesquema noetheriano íntegro, de punto genérico $\eta$, sea $u:X\to S$ un morfismo de tipo finito y sea $\sh{F}$ un $\sh{O}_X$-Módulo coherente.
Para todo $s\in S$, pongamos $X_s=u^{-1}(s)=X\times_S\operatorname{Spec}(\mathbf{k}(s))$, $\sh{F}_s=\sh{F}\otimes_{\sh{O}_S}\mathbf{k}(s)$.
Sea $E$ el conjunto de los $s\in S$ tales que $\mathbf{P}(X_s,\sh{F}_s,\mathbf{k}(s))$ es cierta.
Entonces uno de los dos conjuntos $E$, $S-E$ (resp.\ el conjunto $E$) contiene un abierto no vacío (y, por consiguiente, es un entorno de $\eta$) (resp.\ contiene un abierto no vacío si contiene a $\eta$).
\end{enumerate}
\end{definition}

\oldpage[IV-3]{57}

\begin{remarks}[9.2.2]
\label{IV.9.2.2}
\begin{enumerate}[label=(\roman*)]
  \item Se trata, desde luego, de una convención de lenguaje metamatemático y no de una definición matemática propiamente dicha.
Tenemos ``definiciones'' análogas para relaciones entre $k$, uno o varios $k$-preesquemas algebraicos, $k$-morfismos entre estos preesquemas, Módulos coherentes sobre estos preesquemas u homomorfismos entre estos Módulos.
  \item También tendremos que considerar relaciones en las que intervengan \emph{partes construibles} de preesquemas.
Por ejemplo, sea $\mathbf{P}(X,Z,k)$ una relación; diremos (por abuso de lenguaje) que $\mathbf{P}$ es una propiedad construible (resp.\ ind-construible) de la parte construible $Z$ de $X$ si se satisfacen las dos condiciones siguientes:
\begin{enumerate}[label=$\arabic*^{\mathrm{o}}$]
  \item Si $k$ es un cuerpo, $X$ un preesquema algebraico sobre $k$, $Z$ una \emph{parte construible} de $X$ y $k'$ una extensión de $k$, para que $\mathbf{P}(X,Z,k)$ sea cierta es necesario y suficiente que $\mathbf{P}(X_{(k')},p^{-1}(Z),k')$ sea cierta (siendo $p:X_{(k')}\to X$ la proyección canónica).
  \item Sea $S$ un preesquema noetheriano íntegro, de punto genérico $\eta$, sea $u:X\to S$ un morfismo de tipo finito y sea $Z$ una \emph{parte construible de $X$}.
Para todo $s\in S$, pongamos $X_s=u^{-1}(s)$, $Z_s=Z\cap X_s$.
Sea $E$ el conjunto de los $s\in S$ tales que $\mathbf{P}(X_s,Z_s,\mathbf{k}(s))$ es cierta.
Entonces uno de los dos conjuntos $E$, $S-E$ (resp.\ el conjunto $E$) contiene un abierto no vacío (resp.\ contiene un abierto no vacío si contiene a $\eta$).
\end{enumerate}

Observemos que en la condición~$2^\circ$ hay que suponer que $Z$ es una parte construible de $X$, y no solamente que $Z_s$ es una parte construible de $X_s$ para todo $s$; la primera de estas dos propiedades entraña la segunda~\sref{IV.1.8.2}, pero no recíprocamente.

  \item Si $\mathbf{P}$ es una propiedad construible, está claro que lo mismo ocurre con ``no $\mathbf{P}$''.
Si $\mathbf{P}$, $\mathbf{Q}$ son dos propiedades construibles (resp.\ ind-construibles), lo mismo ocurre con las propiedades ``$\mathbf{P}$ o $\mathbf{Q}$'' y ``$\mathbf{P}$ y $\mathbf{Q}$''; en efecto, bajo las hipótesis de~\sref{IV.9.2.1}[ii], si $E$, $E'$ son dos partes de $S$ y $E$ contiene un abierto no vacío, lo mismo ocurre con $E\cup E'$, y si $S-E$ y $S-E'$ contienen sendos abiertos no vacíos, lo mismo ocurre con $S-(E\cup E')=(S-E)\cap(S-E')$.

  \item Sea $\mathbf{P}(X,\sh{F},k)$ una relación que satisface la condición~\sref{IV.9.2.1}[i]; sea $S$ un preesquema y sea $u:X\to S$ un morfismo de tipo finito; con las notaciones de~\sref{IV.9.2.1}[ii], sea $E$ el conjunto de los $s\in S$ tales que $\mathbf{P}(X_s,\sh{F}_s,\mathbf{k}(s))$ es cierta.
Sea, por otra parte, $g:S'\to S$ un morfismo cualquiera, y pongamos $X'=X_{(S')}$, $\sh{F}'=\sh{F}\otimes_{\sh{O}_S}\sh{O}_{S'}$; de la transitividad de las fibras~(\textbf{I},~3.6.4) y de la condición~\sref{IV.9.2.1}[i] se sigue entonces que \emph{el conjunto de los $s'\in S'$ tales que $\mathbf{P}(X'_{s'},\sh{F}'_{s'},\mathbf{k}(s'))$ es cierta es igual a $g^{-1}(E)$}.
Esto se extiende inmediatamente al caso en que aparecen en $\mathbf{P}$ varios preesquemas, Módulos, morfismos de preesquemas u homomorfismos de Módulos, así como a las propiedades del tipo considerado en~(ii).

  \item Como veremos en el resto de este parágrafo y en el resto del cap.~IV, la mayoría de las propiedades $\mathbf{P}$ que satisfacen la condición~\sref{IV.9.2.1}[i] satisfacen también la condición~\sref{IV.9.2.1}[ii].
Como posibles excepciones, señalemos la propiedad de ser \emph{proyectivo}, o \emph{cuasiproyectivo}, o \emph{afín}, o \emph{cuasiafín} sobre el cuerpo base (para un preesquema algebraico); veremos en~\sref{IV.9.6.2} que estas propiedades son \emph{ind-construibles}, pero daremos más adelante un ejemplo en el que $S$ es una parte abierta no vacía de $\operatorname{Spec}(\mathbf{Z})$ (o una parte abierta de una curva elíptica sobre un cuerpo finito) y en el que todas las fibras $X_s$

\oldpage[IV-3]{58}
salvo la fibra genérica $X_\eta$ son proyectivas sobre $\mathbf{k}(s)$ (siendo todos los preesquemas $X_s$ de dimensión~2).
\end{enumerate}
\end{remarks}

\begin{proposition}[9.2.3]
\label{IV.9.2.3}
Sea $\mathbf{P}$ una propiedad construible \textnormal{(}resp.\ ind-construible\textnormal{)} de los preesquemas algebraicos, sea $S$ un preesquema, sea $X$ un preesquema de presentación finita sobre $S$ y sea $\sh{F}$ un $\sh{O}_X$-Módulo cuasicoherente de presentación finita.
Entonces el conjunto $E$ de los $s\in S$ tales que $\mathbf{P}(X_s,\sh{F}_s,\mathbf{k}(s))$ es cierta es localmente construible \textnormal{(}resp.\ ind-construible\textnormal{)}.
Además, si $S$ es irreducible de punto genérico $\eta$, uno de los dos conjuntos $E$, $S-E$ es un entorno de $\eta$ \textnormal{(}resp.\ $E$ es un entorno de $\eta$ si contiene este punto\textnormal{)}.
\end{proposition}

Para demostrar estas afirmaciones, podemos reducirnos al caso en que $S=\operatorname{Spec}(A)$ es afín.
Sabemos entonces que existe un subanillo $A_0$ de $A$ que es una $\mathbf{Z}$-álgebra de tipo finito, un $A_0$-preesquema de tipo finito $X_0$ y un $\sh{O}_{X_0}$-Módulo coherente $\sh{F}_0$ tales que $X$ es isomorfo a $X_0\otimes_{A_0}A$ y $\sh{F}$ a $\sh{F}_0\otimes_{A_0}A$~\sref{IV.8.9.1}.
Sea $p:S\to S_0=\operatorname{Spec}(A_0)$ el morfismo correspondiente a la inyección $A_0\to A$, y sea $E_0$ el conjunto de los $s_0\in S_0$ tales que
\[
  \mathbf{P}((X_0)_{s_0},(\sh{F}_0)_{s_0},\mathbf{k}(s_0))
\]
es cierta; entonces, en virtud de~\sref{IV.9.2.2}[iv], se tiene $E=p^{-1}(E_0)$; por~\sref{IV.1.8.2}, podemos por tanto reducirnos al caso en que $S$ es el espectro de una $\mathbf{Z}$-álgebra de tipo finito y, por tanto, un esquema noetheriano.
Utilicemos el criterio de construibilidad~($\mathbf{0}_{\mathrm{III}}$,~9.2.3) (resp.\ el criterio de ind-construibilidad~\sref{IV.1.9.10}); empleando~\sref{IV.9.2.2}[iv] como antes y sustituyendo $S$ por un subpreesquema cerrado íntegro de $S$, nos reducimos al caso en que $S$ es noetheriano e íntegro y en que hay que demostrar que $E$ no es denso en ninguna parte de $S$ o contiene un subconjunto abierto no vacío de $S$ (resp.\ que $E$ contiene un subconjunto abierto no vacío de $S$ si contiene el punto genérico); esto queda garantizado por la condición~\sref{IV.9.2.1}[ii].

Observemos que hemos utilizado~\sref{IV.9.2.1}[ii] únicamente cuando $S$ es el espectro de un anillo íntegro de tipo finito sobre $\mathbf{Z}$.
Por otra parte, está claro que el enunciado de~\sref{IV.9.2.3} se aplica también cuando aparecen en $\mathbf{P}$ varios preesquemas, Módulos sobre estos preesquemas, morfismos de preesquemas u homomorfismos de Módulos.
Se aplica asimismo cuando aparecen partes (en número finito) de los preesquemas considerados, siempre que se exija que estas partes sean \emph{localmente construibles}.
En efecto, la reducción al caso en que $S$ es afín muestra que podemos reducirnos al caso en que estas partes son \emph{construibles}: se aplica entonces~{\hyperlink{ega4.c04.thm.8-3-11}{(8.3.11)}}, que muestra (con la notación precedente) que una parte construible de $X$ es la imagen inversa de una parte construible de $X_0$ para una elección adecuada de $A_0$.

\begin{corollary}[9.2.4]
\label{IV.9.2.4}
Sea $\mathbf{P}$ una propiedad construible \textnormal{(}resp.\ ind-construible\textnormal{)} de los preesquemas algebraicos, sean $X$, $Y$ dos $S$-preesquemas de presentación finita y sea $f:X\to Y$ un $S$-morfismo.
Para todo $s\in S$, pongamos $X_s=X\times_S\operatorname{Spec}(\mathbf{k}(s))$, $Y_s=Y\times_S\operatorname{Spec}(\mathbf{k}(s))$, $f_s=f\times 1:X_s\to Y_s$.
Entonces el conjunto $E$ de los $s\in S$ tales que, para todo $y\in Y_s$, la propiedad $\mathbf{P}(f_s^{-1}(y),\mathbf{k}(y))$ es cierta, es localmente construible \textnormal{(}resp.\ ind-construible\textnormal{)}.
\end{corollary}

En efecto, sea $Z$ el conjunto de los $y\in Y$ tales que $\mathbf{P}(f^{-1}(y),\mathbf{k}(y))$ es cierta.
Como las fibras $f^{-1}(y)$ y $f_s^{-1}(y)$ son isomorfas, vemos que $E$ es el conjunto de los $s\in S$ tales que $Y_s\subset Z$; si $g:Y\to S$ es el morfismo estructural, se tiene por tanto $E=S-g(Y-Z)$.

\oldpage[IV-3]{59}

Ahora bien, $f$ es de presentación finita~\sref{IV.1.6.2}[v], luego se sigue de~\sref{IV.9.2.3} que $Z$ es localmente construible (resp.\ ind-construible) en $Y$, y por tanto $Y-Z$ es localmente construible (resp.\ pro-construible) en $Y$.
Como $g$ es de presentación finita, $g(Y-Z)$ es localmente construible (resp.\ pro-construible) en $S$, en virtud del teorema de Chevalley~\sref{IV.1.8.4} (resp.\ de~\sref{IV.1.9.5}[vii]); por tanto, $E$ es localmente construible (resp.\ ind-construible) en $S$.

\begin{remark}[9.2.5]
\label{IV.9.2.5}
Observemos que, si $\mathbf{P}$ es una propiedad de los preesquemas algebraicos para la cual es cierta la prop.~\sref{IV.9.2.3}, entonces $\mathbf{P}$ satisface también la condición~\sref{IV.9.2.1}[ii]: esto resulta, en efecto, de que, en un espacio noetheriano irreducible, un conjunto construible o bien no es denso en ninguna parte o bien contiene un abierto no vacío~($\mathbf{0}_{\mathrm{III}}$,~9.2.3).
\end{remark}

\begin{proposition}[9.2.6]
\label{IV.9.2.6}
Sea $\mathbf{P}$ una de las propiedades siguientes de un $k$-preesquema algebraico $X$:
\begin{enumerate}[label=(\roman*)]
  \item $X$ es vacío.
  \item $X$ es finito sobre $k$.
  \item $X$ es radicial sobre $k$.
  \item $\dim(X)$ pertenece a una parte dada $\Phi$ del conjunto $\overline{\mathbf{Z}}=\mathbf{Z}\cup\{-\infty\}$.
\end{enumerate}
Entonces $\mathbf{P}$ es construible.
\end{proposition}

Está claro que (i) y (ii) son casos particulares de (iv), tomando como $\Phi$, respectivamente, el conjunto $\{-\infty\}$ y el conjunto $\{-\infty,0\}$.
Por tanto, solo tenemos que demostrar (iii) y (iv).
En cada uno de estos dos casos, la condición~(i) de~\sref{IV.9.2.1} se satisface en virtud de~\sref{IV.2.7.1}[xv] y~\sref{IV.4.1.4}.
Por otra parte, en el caso (iii), la propiedad $\mathbf{P}$ satisface la conclusión de~\sref{IV.9.2.3} en virtud de~\sref{IV.1.8.7}; queda, pues, por comprobar que lo mismo ocurre en el caso (iv).
Esto se seguirá de la proposición más precisa que damos a continuación:

\begin{proposition}[9.2.6.1]
\label{IV.9.2.6.1}
Si $f:X\to S$ es un morfismo de presentación finita, la función $s\mapsto\dim(f^{-1}(s))$ es localmente construible.
\end{proposition}

La cuestión es local sobre $S$, por lo que podemos suponer que $S=\operatorname{Spec}(A)$ es afín y demostrar que, para todo $n$, el conjunto $E$ de los $s\in S$ tales que $\dim(X_s)=n$ es construible.
El mismo razonamiento que en~\sref{IV.9.2.3} nos reduce al caso en que $A$ es noetheriano e íntegro, y basta entonces demostrar el siguiente

\begin{corollary}[9.2.6.2]
\label{IV.9.2.6.2}
Sea $S$ un preesquema noetheriano íntegro de punto genérico $\eta$, y sea $f:X\to S$ un morfismo de tipo finito.
Entonces existe un entorno de $\eta$ en $S$ tal que la función $s\mapsto\dim(X_s)$ es constante en este entorno.
\end{corollary}

Las imágenes por $f$ de las componentes irreducibles (en número finito) de $X$ que no cortan a $X_\eta$ están contenidas en partes cerradas de $S$ que no contienen a $\eta$ (pues $S$ es íntegro~($\mathbf{0}_{\mathrm{I}}$,~2.1.5)); por tanto, tras sustituir $S$ por un entorno abierto de $\eta$, podemos reducirnos al caso en que todas las componentes irreducibles $X_i$ de $X$ cortan a $X_\eta$.
Denotemos de nuevo por $X_i$ el subpreesquema cerrado reducido de $X$ que tiene a $X_i$ como espacio subyacente; como $\dim(X_s)=\sup_i\dim((X_i)_s)$~\sref{IV.4.1.1}, podemos reducirnos al caso en que $X$ es irreducible.
Existe entonces un recubrimiento finito $(U_j)$ de $X$ por

\oldpage[IV-3]{60}
abiertos afines densos, y los números $\dim((U_j)_\eta)$ son todos iguales a $n=\dim(X_\eta)$~\sref{IV.4.1.1.3}; podemos, por tanto, reducirnos al caso en que $X$ es afín, y en consecuencia también $X_\eta$.
Existe entonces, en virtud de~\sref{IV.4.1.2}, un abierto no vacío $W$ de $X$ tal que $W\supset X_\eta$, y un $\mathbf{k}(\eta)$-morfismo finito y sobreyectivo $h:W_\eta\to\mathbf{V}_{\mathbf{k}(\eta)}^n(=\operatorname{Spec}(\mathbf{k}(\eta)[T_1,\ldots,T_n]))$; aplicando~\sref{IV.8.8.2}[i] y el método de~\sref{IV.8.1.2}[a], deducimos (tras sustituir $S$ por un entorno de $\eta$ si es necesario) que $h=g_\eta$, donde $g:W\to\mathbf{V}_A^n(=\operatorname{Spec}(A[T_1,\ldots,T_n]))$ es un $S$-morfismo, y podemos suponer que este morfismo es finito y sobreyectivo en virtud de~\sref{IV.8.10.5}[vi] y~\sref{IV.8.10.5}[x].
Concluimos, por tanto, que para todo $s\in S$, el morfismo $g_s:W_s\to\mathbf{V}_{\mathbf{k}(s)}^n$ es finito y sobreyectivo, y por consiguiente $\dim(W_s)=n$~\sref{IV.4.1.2}.

\subsection{Propiedades construibles de los morfismos de preesquemas algebraicos}
\label{subsection:IV.9.3}

\begin{proposition}[9.3.1]
\label{IV.9.3.1}
Sea $\mathbf{P}(X,k)$ una propiedad construible \textnormal{(}resp.\ ind-construible\textnormal{)} de los preesquemas algebraicos.
Denotemos por $\mathbf{P}'(f,X,Y,k)$ la relación siguiente: $f:X\to Y$ es un $k$-morfismo de $k$-preesquemas algebraicos tal que, para todo $y\in Y$, la propiedad $\mathbf{P}(f^{-1}(y),\mathbf{k}(y))$ es cierta.
Entonces $\mathbf{P}'$ es una propiedad construible \textnormal{(}resp.\ ind-construible\textnormal{)}.
\end{proposition}

En efecto, como $\mathbf{P}$ satisface la condición~\sref{IV.9.2.1}[i], lo mismo ocurre con $\mathbf{P}'$ en virtud de la transitividad de las fibras~(\textbf{I},~3.6.4); además, el hecho de que $\mathbf{P}'$ satisface la condición~\sref{IV.9.2.1}[ii] se sigue de~\sref{IV.9.2.4}, en razón de la observación~\sref{IV.9.2.5}.

\begin{proposition}[9.3.2]
\label{IV.9.3.2}
Sea $\mathbf{P}$ una de las propiedades siguientes de un $k$-morfismo $f:X\to Y$ de preesquemas algebraicos:
\begin{enumerate}[label=(\roman*)]
  \item $f$ es sobreyectivo.
  \item $f$ es cuasifinito.
  \item $f$ es radicial.
  \item Para todo $y\in Y$, $\dim(f^{-1}(y))$ pertenece a $\Phi$ \textnormal{(}notación de~\sref{IV.9.2.6}\textnormal{)}.
\end{enumerate}
Entonces $\mathbf{P}$ es construible.
\end{proposition}

Esto se sigue inmediatamente de~\sref{IV.9.3.1} y~\sref{IV.9.2.6}, teniendo en cuenta que $f$ es de tipo finito~\sref[I]{I.1.5.4}[v], la caracterización de los morfismos radiciales~\sref[I]{I.3.5.8} y la de los morfismos cuasifinitos~\sref[II]{II.6.2.2}.

\begin{proposition}[9.3.3]
\label{IV.9.3.3}
Supongamos satisfechas las hipótesis de~\sref{IV.8.8.1} y conservemos las mismas notaciones; supongamos además que $S_\alpha$ es cuasicompacto, que $X_\alpha$ e $Y_\alpha$ son de presentación finita sobre $S_\alpha$, y sea $f_\alpha:X_\alpha\to Y_\alpha$ un $S_\alpha$-morfismo.
Sea $\mathbf{P}$ una propiedad ind-construible de los morfismos de preesquemas algebraicos.
Para todo $s\in S$ \textnormal{(}resp.\ $s_\lambda\in S_\lambda$\textnormal{)}, pongamos $X_s=X\times_S\operatorname{Spec}(\mathbf{k}(s))$, $Y_s=Y\times_S\operatorname{Spec}(\mathbf{k}(s))$, $f_s=f\times 1:X_s\to Y_s$ \textnormal{(}resp.\ $X_{\lambda,s_\lambda}=X_\lambda\times_{S_\lambda}\operatorname{Spec}(\mathbf{k}(s_\lambda))$, $Y_{\lambda,s_\lambda}=Y_\lambda\times_{S_\lambda}\operatorname{Spec}(\mathbf{k}(s_\lambda))$, $f_{\lambda,s_\lambda}=f_\lambda\times 1:X_{\lambda,s_\lambda}\to Y_{\lambda,s_\lambda}$\textnormal{)}.
Entonces, para que para todo $s\in S$ la propiedad $\mathbf{P}(f_s)$ sea cierta, es necesario y suficiente que exista $\lambda\geqslant\alpha$ tal que para todo $s_\lambda\in S_\lambda$ se tenga $\mathbf{P}(f_{\lambda,s_\lambda})$.
\end{proposition}

En efecto, sea $E$ (resp.\ $E_\lambda$) el conjunto de los $s\in S$ (resp.\ $s_\lambda\in S_\lambda$) tales que la propiedad $\mathbf{P}(f_s)$ (resp.\ $\mathbf{P}(f_{\lambda,s_\lambda})$) es cierta; se sigue de~\sref{IV.9.2.2}[iv] que $E_\mu=u_{\lambda\mu}^{-1}(E_\lambda)$ para $\lambda\leqslant\mu$, y que $E=u_\lambda^{-1}(E_\lambda)$; además, en virtud de~\sref{IV.9.2.3}, $E$ (resp.\ $E_\lambda$) es ind-construible en $S$ (resp.\ $S_\lambda$); la proposición se sigue entonces de~\sref{IV.8.3.4} aplicado a la parte ind-construible $E$ de $S$.

\oldpage[IV-3]{61}

Este resultado se extiende sin dificultad a las propiedades $\mathbf{P}$ del tipo considerado en (9.2.3, (i) y (ii)).

\begin{remark}[9.3.4]
\label{IV.9.3.4}
La conjunción de~\sref{IV.9.3.3} y~\sref{IV.9.3.2}[ii] demuestra la afirmación~\sref{IV.8.10.5}[xi].
\end{remark}

\begin{proposition}[9.3.5]
\label{IV.9.3.5}
Sea $\mathbf{P}(X,Y,k)$ la propiedad: ``$X$ e $Y$ son dos preesquemas de tipo finito sobre el cuerpo $k$, y existe una extensión $k'$ de $k$ y un $k'$-morfismo $g:X_{(k')}\to Y_{(k')}$ que satisface $\mathbf{Q}(g)$'', donde $\mathbf{Q}$ es una de las propiedades~\emph{(i)} a~\emph{(xiv)} de~\sref{IV.8.10.5}.
Entonces $\mathbf{P}$ es una propiedad ind-construible.
\end{proposition}

La definición de $\mathbf{P}$ muestra, en efecto, que se satisface la condición~\sref{IV.9.2.1}[i], teniendo en cuenta que la propiedad $\mathbf{Q}(g)$ es estable por cambio de base y que dos extensiones de $k$ siempre pueden considerarse como subextensiones de una tercera extensión de $k$.
Para verificar~\sref{IV.9.2.1}[ii], podemos reducirnos al caso en que $S=\operatorname{Spec}(A)$ es afín; si $K=\mathbf{k}(\eta)$ es el cuerpo de fracciones de $A$, existen, por hipótesis y en virtud de~\sref{IV.9.1.4}, una extensión finita $K'$ de $K$ y un $K'$-morfismo $g':(X_\eta)_{(K')}\to(Y_\eta)_{(K')}$ que satisface $\mathbf{Q}(g')$, y $K'$ es evidentemente el cuerpo de fracciones de una $A$-álgebra íntegra finita $A'$.
Si ponemos $S'=\operatorname{Spec}(A')$, se sigue entonces de~\sref{IV.8.10.5} que existe un entorno $U'$ del punto genérico $\eta'$ de $S'$ tal que, si ponemos $X'=X\otimes_A A'$ e $Y'=Y\otimes_A A'$, para todo $s'\in U'$ existe un morfismo $X'_{s'}\to Y'_{s'}$ que posee la propiedad $\mathbf{Q}$.
Pero el morfismo $h:S'\to S$ es finito y, por tanto, cerrado; como $h^{-1}(\eta)=\{\eta'\}$, $h(U')$ contiene un entorno abierto $U$ de $\eta$ en $S$.
Para todo $s\in U$ existe, por tanto, $s'\in U'$ tal que $h(s')=s$, y como $X'_{s'}=X_s\otimes_{\mathbf{k}(s)}\mathbf{k}(s')$ e $Y'_{s'}=Y_s\otimes_{\mathbf{k}(s)}\mathbf{k}(s')$, la propiedad $\mathbf{P}(X_s,Y_s,\mathbf{k}(s))$ es cierta para todo $s\in U$.

\begin{env}[9.3.6]
\label{IV.9.3.6}
\textbf{Ejemplo.}
Tomemos, por ejemplo, como $\mathbf{Q}$ la propiedad de ser un isomorfismo.
Combinando entonces~\sref{IV.9.3.5} y~\sref{IV.9.3.3}, obtenemos la propiedad siguiente: con las notaciones e hipótesis de~\sref{IV.8.8.1}, siendo $S_\alpha$ cuasicompacto y $X_\alpha$ e $Y_\alpha$ de presentación finita sobre $S_\alpha$, para que, para todo $s\in S$, $X_s$ e $Y_s$ sean \emph{geométricamente isomorfos}~\sref{IV.9.1.4}, es necesario y suficiente que exista $\lambda\geqslant\alpha$ tal que, para todo $s_\lambda\in S_\lambda$, $X_{\lambda,s_\lambda}$ e $Y_{\lambda,s_\lambda}$ sean geométricamente isomorfos.
\end{env}

Tenemos un resultado análogo cuando los preesquemas considerados están dotados de ``leyes de composición'' de cierto tipo~($\mathbf{0}_{\mathrm{III}}$,~8.2.1), por ejemplo ``preesquemas en grupos'', ``preesquemas en anillos'', etc.~($\mathbf{0}_{\mathrm{III}}$,~8.2.3).
El enunciado precedente sigue siendo válido cuando por ``isomorfismo'' se entienden isomorfismos de preesquemas que son homomorfismos para las leyes de composición consideradas~($\mathbf{0}_{\mathrm{III}}$,~8.2.2); basta utilizar aquí no solo~\sref{IV.8.10.5}, sino también~\sref{IV.8.8.2}[i], observando que la noción de homomorfismo para una ley de composición se expresa diciendo que ciertos diagramas de morfismos de preesquemas son conmutativos (se sobreentiende que los morfismos de transición $X_\mu\to X_\lambda$ e $Y_\mu\to Y_\lambda$ para $\lambda\leqslant\mu$ son homomorfismos para las leyes de composición consideradas).

También se pueden considerar, en vez de morfismos de preesquemas como en~\sref{IV.9.3.5}, homomorfismos de Módulos, utilizando~\sref{IV.9.1.6} en lugar de~\sref{IV.9.1.5}.

\oldpage[IV-3]{62}

\subsection{Construibilidad de ciertas propiedades de los módulos}
\label{subsection:IV.9.4}

\begin{env}[9.4.1]
\label{IV.9.4.1}
\textbf{Notaciones (9.4.1).}
En este n\textsuperscript{o} y en los siguientes hasta el final de \textsection9, utilizaremos sistemáticamente las notaciones siguientes: dado un morfismo $f:X\to S$, pondremos, para todo $s\in S$, $X_s=f^{-1}(s)=X\times_S\operatorname{Spec}(\mathbf{k}(s))$; para todo $\sh{O}_X$-Módulo cuasicoherente $\sh{F}$, $\sh{F}_s$ denotará el $\sh{O}_{X_s}$-Módulo cuasicoherente $\sh{F}\otimes_{\sh{O}_S}\mathbf{k}(s)$, y para todo homomorfismo $u:\sh{F}\to\sh{G}$ de $\sh{F}$ en un $\sh{O}_X$-Módulo cuasicoherente $\sh{G}$, $u_s:\sh{F}_s\to\sh{G}_s$ será el morfismo $p^*(u)$, designando por $p$ la proyección canónica $X_s\to X$.
Para toda sección $g$ de $\sh{F}$ sobre $X$, denotaremos por $g_s$ la imagen de $g$ por el homomorfismo canónico $\Gamma(X,\sh{F})\to\Gamma(X_s,\sh{F}_s)$.
Para toda parte $Z$ de $X$, denotaremos por $Z_s$ la imagen inversa $p^{-1}(Z)=Z\cap X_s$~(\textbf{I},~3.6.1).
Finalmente, si $Y$ es un segundo $S$-preesquema y $h:X\to Y$ un $S$-morfismo, denotaremos por $h_s$ el morfismo $h\times 1:X_s\to Y_s$.
\end{env}

\begin{proposition}[9.4.2]
\label{IV.9.4.2}
Sea $S$ un preesquema íntegro de punto genérico $\eta$, $f:X\to S$ un morfismo de presentación finita, y sean $\sh{F}$, $\sh{G}$, $\sh{H}$ tres $\sh{O}_X$-Módulos cuasicoherentes de presentación finita.
Sean $u:\sh{F}\to\sh{G}$, $v:\sh{G}\to\sh{H}$ dos homomorfismos de $\sh{O}_X$-Módulos, y supongamos que la sucesión $\sh{F}_\eta\xrightarrow{u_\eta}\sh{G}_\eta\xrightarrow{v_\eta}\sh{H}_\eta$ es exacta.
Entonces existe un entorno $U$ de $\eta$ en $S$ tal que, para todo $s\in U$, la sucesión $\sh{F}_s\xrightarrow{u_s}\sh{G}_s\xrightarrow{v_s}\sh{H}_s$ es exacta.
\end{proposition}

\begin{proof}
Con las notaciones generales de~\sref{IV.9.2.1}, se trata aquí de la relación $\mathbf{P}(X,\sh{F},\sh{G},\sh{H},u,v,k)$: ``$X$ es un preesquema algebraico sobre el cuerpo $k$, $\sh{F}\xrightarrow{u}\sh{G}\xrightarrow{v}\sh{H}$ es una sucesión exacta de $\sh{O}_X$-Módulos cuasicoherentes''.
Como, para toda extensión $k'$ de $k$, la proyección canónica $X_{(k')}\to X$ es un morfismo fielmente plano~\sref{IV.2.2.13}, se satisface la condición~\sref{IV.9.2.1},~(i))~\sref{IV.2.2.7}.
En virtud de~\sref{IV.9.2.3}, podemos restringirnos al caso en que $S$ es íntegro y noetheriano, en cuyo caso $X$ es noetheriano, y $\sh{F}$, $\sh{G}$, $\sh{H}$ son $\sh{O}_X$-Módulos coherentes.
La hipótesis implica que existe un entorno $U$ de $\eta$ en $S$ tal que $\sh{F}|f^{-1}(U)\to\sh{G}|f^{-1}(U)\to\sh{H}|f^{-1}(U)$ es \emph{exacta}, en virtud de~\sref{IV.8.5.8},~(ii)) aplicado según el método general de~\sref{IV.8.1.2},~$a$)), y por tanto podemos suponer ya que $\sh{F}\xrightarrow{u}\sh{G}\xrightarrow{v}\sh{H}$ es exacta; basta evidentemente demostrar que $\operatorname{Ker}(v_s)=(\operatorname{Ker} v)_s$ y $\operatorname{Im}(u_s)=(\operatorname{Im} u)_s$ para $s$ próximo a $\eta$ en $S$.
Teniendo en cuenta la coherencia de $\sh{F}$, $\sh{G}$ y $\sh{H}$ y de
($\mathbf{0}_{\mathrm{I}}$,~5.3.4), nos reducimos al caso particular en que
\[
 0\longrightarrow\sh{F}\longrightarrow\sh{G}\longrightarrow\sh{H}
 \longrightarrow0
\]
es exacta.
Existe entonces un entorno abierto $U$ de $\eta$ en $S$ tal que
$\sh{H}|f^{-1}(U)$ es plano sobre $U$ \sref{IV.6.9.1}; por tanto, \sref{IV.2.1.8}
muestra que para todo $s\in U$ la sucesión
\[
 0\longrightarrow\sh{F}_s\longrightarrow\sh{G}_s\longrightarrow\sh{H}_s
 \longrightarrow0
\]
es exacta.
\end{proof}

\begin{corollary}[9.4.3]
\label{IV.9.4.3}
Sea $S$ un preesquema íntegro, de punto genérico $\eta$, y sea $f:X\to S$ un morfismo de presentación finita.
Sea $\sh{L}^\bullet=(\sh{L}^i)_{i\in\mathbf{Z}}$ un complejo cuyos términos son $\sh{O}_X$-Módulos cuasicoherentes de presentación finita.
Para todo $i$, existe un entorno abierto $U$ de $\eta$ en $S$ tal que los homomorfismos canónicos
\begin{equation}
\label{IV.9.4.3.1}
  (\sh{H}^i(\sh{L}^\bullet))_s\to\sh{H}^i(\sh{L}^\bullet_s) \tag{9.4.3.1}
\end{equation}
son biyectivos para todo $s\in U$.
\end{corollary}

\oldpage[IV-3]{63}

\begin{proof}
Es suficiente considerar un complejo de tres términos en grados $-1$, $0$, $+1$,
\[
 \sh{M}\xrightarrow{u}\sh{N}\xrightarrow{v}\sh{P},
\]
y $i=0$.
El homomorfismo canónico en cuestión es entonces
\[
 (\operatorname{Ker}v/\operatorname{Im}u)_s
 \longrightarrow \operatorname{Ker}(v_s)/\operatorname{Im}(u_s).
\]
Utilizando \sref{IV.8.9.1} y \sref{IV.8.5.2},~(i), podemos reducirnos al caso
en que $S$ (y por tanto $X$) es noetheriano; los Módulos $\sh{M}$, $\sh{N}$ y
$\sh{P}$ son entonces coherentes.
Así, $\operatorname{Im}u$ y $\operatorname{Ker}v$ son coherentes
($\mathbf{0}_{\mathrm{I}}$,~5.3.4), y \sref{IV.9.4.2} proporciona un entorno
$U$ de $\eta$ en el cual
\[
 \operatorname{Ker}(v_s)=(\operatorname{Ker}v)_s,
 \qquad
 \operatorname{Im}(u_s)=(\operatorname{Im}u)_s.
\]
La conclusión se sigue aplicando \sref{IV.9.4.2} a la sucesión exacta
\[
 0\longrightarrow\operatorname{Im}u\longrightarrow\operatorname{Ker}v
 \longrightarrow\operatorname{Ker}v/\operatorname{Im}u\longrightarrow0,
\]
puesto que $\sh{O}_\eta=\boldsymbol{k}(\eta)$ y, por consiguiente, la sucesión
correspondiente sobre $\eta$ es exacta.
\end{proof}

\begin{proposition}[9.4.4]
\label{IV.9.4.4}
Sea $f:X\to S$ un morfismo de presentación finita, y sean $\sh{F}$, $\sh{G}$, $\sh{H}$ tres $\sh{O}_X$-Módulos cuasicoherentes de presentación finita.
Sean $u:\sh{F}\to\sh{G}$, $v:\sh{G}\to\sh{H}$ dos homomorfismos de $\sh{O}_X$-Módulos.
Entonces el conjunto $E$ de los $s\in S$ tales que la sucesión $\sh{F}_s\xrightarrow{u_s}\sh{G}_s\xrightarrow{v_s}\sh{H}_s$ es exacta es localmente construible.
\end{proposition}

Teniendo en cuenta~\sref{IV.9.2.3}, basta establecer que la propiedad $\mathbf{P}$ considerada en~\sref{IV.9.4.2} es construible.
Ya hemos observado en la demostración de~\sref{IV.9.4.2} que la condición~\sref{IV.9.2.1},~(i)) se satisface para esta propiedad, y queda verificar la condición~\sref{IV.9.2.1},~(ii)).
Supongamos, pues, que $S$ es íntegro noetheriano, de punto genérico $\eta$, y demostremos que $E$ o $S-E$ es un entorno de $\eta$.
Si $\eta\in E$, nuestra afirmación se sigue de~\sref{IV.9.4.2}, y por tanto podemos restringirnos al caso en que $\eta\notin E$, es decir, la sucesión $\sh{F}_\eta\xrightarrow{u_\eta}\sh{G}_\eta\xrightarrow{v_\eta}\sh{H}_\eta$ no es exacta.
Distinguimos entonces dos casos:

$1^\circ$ Pongamos $w=v\circ u$, y supongamos primero que $w_\eta=v_\eta\circ u_\eta\neq 0$.
Como $\sh{F}$, $\sh{G}$, $\sh{H}$ son coherentes, también lo es $\sh{N}=\operatorname{Ker}(w)$~($\mathbf{0}_{\mathrm{I}}$,~5.3.4); se sigue entonces de~\sref{IV.9.4.2}, aplicado a la sucesión exacta $0\to\sh{N}\to\sh{F}\xrightarrow{w}\sh{H}$, que existe un entorno $U$ de $\eta$ en $S$ tal que $\operatorname{Ker}(w_s)=\sh{N}_s$; restringiendo $S$, podemos por tanto suponer que esta relación se satisface para todo $s\in S$.
Sea $j$ la inyección canónica $\sh{N}\to\sh{F}$, y pongamos $\sh{M}=\operatorname{Coker}(j)$; la exactitud por la derecha del funtor $\sh{F}\mapsto\sh{F}_s$ implica que $\sh{M}_s=\operatorname{Coker}(j_s)$ para todo $s\in S$.
La hipótesis $w_\eta\neq 0$ significa que $\sh{M}_\eta\neq 0$; como $\sh{M}$ es coherente~($\mathbf{0}_{\mathrm{I}}$,~5.3.4), se sigue de~\sref{IV.1.8.6} que existe un entorno $U$ de $\eta$ en $S$ tal que $\sh{M}_s\neq 0$ para todo $s\in U$, de donde $w_s\neq 0$ para todo $s\in U$, y \emph{a fortiori} $S-E$ es un entorno de $\eta$.

$2^\circ$ Supongamos que $w_\eta=0$; en virtud de~\sref{IV.8.5.2},~(i)), aplicado según el método general de~\sref{IV.8.1.2},~$a$)), existe un entorno $U$ de $\eta$ tal que $w|f^{-1}(U)=0$; sustituyendo $S$ por $U$, podemos suponer ya que $w=0$ en $X$.
Entonces $\sh{F}\xrightarrow{u}\sh{G}\xrightarrow{v}\sh{H}$ es un complejo de tres términos, al que podemos aplicar~\sref{IV.9.4.3}; por hipótesis tenemos $(\sh{H}^0(\sh{L}^\bullet))_\eta=\sh{H}^0(\sh{L}^\bullet_\eta)\neq 0$, y $\sh{H}^0(\sh{L}^\bullet)$ es coherente~($\mathbf{0}_{\mathrm{I}}$,~5.3.4 y~5.3.3), por tanto se sigue de~\sref{IV.1.8.6} que existe un entorno $U$ de $\eta$ tal que $(\sh{H}^0(\sh{L}^\bullet))_s\neq 0$ para todo $s\in U$; pero como podemos suponer que $\sh{H}^0(\sh{L}^\bullet_s)=(\sh{H}^0(\sh{L}^\bullet))_s$ para todo $s\in U$ por~\sref{IV.9.4.3}, vemos de nuevo que $S-E$ es un entorno de $\eta$ en $S$.

\oldpage[IV-3]{64}

\begin{corollary}[9.4.5]
\label{IV.9.4.5}
Sea $f:X\to S$ un morfismo de presentación finita, sean $\sh{F}$, $\sh{G}$ dos $\sh{O}_X$-Módulos cuasicoherentes de presentación finita, y sea $u:\sh{F}\to\sh{G}$ un homomorfismo de $\sh{O}_X$-Módulos.
Entonces el conjunto de los $s\in S$ tales que $u_s$ es inyectivo \textnormal{(}resp.\ sobreyectivo, biyectivo, nulo\textnormal{)} es localmente construible.
\end{corollary}

Basta aplicar~\sref{IV.9.4.4} a las sucesiones $0\to\sh{F}\xrightarrow{u}\sh{G}$, $\sh{F}\xrightarrow{u}\sh{G}\to 0$, $0\to\sh{F}\xrightarrow{u}\sh{G}\to 0$, $\sh{F}\xrightarrow{u}\sh{G}\xrightarrow{1_\sh{G}}\sh{G}$.

\begin{corollary}[9.4.6]
\label{IV.9.4.6}
Sea $f:X\to S$ un morfismo de presentación finita y sea $\sh{F}$ un $\sh{O}_X$-Módulo cuasicoherente de presentación finita.
Sea $h$ una sección de $\sh{F}$ sobre $X$; para todo $s\in S$, sea $h_s$ la sección correspondiente de $\sh{F}_s$ sobre $X_s$ \textnormal{(}para el morfismo de proyección $X_s\to X$\textnormal{)}.
Entonces el conjunto de los $s\in S$ tales que $h_s=0$ es localmente construible.
\end{corollary}

Basta observar que $h$ corresponde a un homomorfismo $u:\sh{O}_X\to\sh{F}$~($\mathbf{0}_{\mathrm{I}}$,~5.1.1) y $h_s$ al homomorfismo $u_s$.

\begin{proposition}[9.4.7]
\label{IV.9.4.7}
Sea $f:X\to S$ un morfismo de presentación finita y sea $\sh{F}$ un $\sh{O}_X$-Módulo cuasicoherente de presentación finita.
El conjunto $E$ \textnormal{(}resp.\ $E'$\textnormal{)} de los $s\in S$ tales que $\sh{F}_s$ es un $\sh{O}_{X_s}$-Módulo localmente libre \textnormal{(}resp.\ localmente libre de rango $n$\textnormal{)} es localmente construible.
\end{proposition}

Si $X$ es un preesquema algebraico sobre un cuerpo $k$, $\sh{F}$ un $\sh{O}_X$-Módulo coherente y $k'$ una extensión de $k$, para que $\sh{F}$ sea localmente libre (resp.\ localmente libre de rango $n$), es necesario y suficiente que lo sea $\sh{F}\otimes_k k'$, puesto que la proyección $X_{(k')}\to X$ es un morfismo fielmente plano~\sref{IV.2.2.7}.
En otras palabras, la condición~\sref{IV.9.2.1},~(i)) se satisface para las propiedades cuya construibilidad queremos demostrar, y queda verificar~\sref{IV.9.2.1},~(ii)); podemos, por tanto, suponer aún que $S$ es afín, noetheriano e íntegro.
Hay entonces cuatro casos que considerar:

$1^\circ$ $\eta\in E$.
En virtud de~\sref{IV.8.5.5}, aplicado según el método general de~\sref{IV.8.1.2},~$a$)), existe un entorno $U$ de $\eta$ en $S$ tal que $\sh{F}|f^{-1}(U)$ es localmente libre; \emph{a fortiori} $\sh{F}_s$ es localmente libre para todo $s\in U$.

$2^\circ$ $\eta\in E'$.
El mismo razonamiento que en $1^\circ$.

$3^\circ$ $\eta\in S-E$.
Como $\sh{F}_\eta$ es un $\sh{O}_{X_\eta}$-Módulo coherente, decir que no es localmente libre significa que no es plano sobre $\sh{O}_{X_\eta}$ (Bourbaki, \emph{Alg.\ comm.}, cap.~II, \textsection5, n\textsuperscript{o}~2, cor.~2 del teorema~1).
El hecho de que $S-E$ sea un entorno de $\eta$ se seguirá entonces del lema más general siguiente:

\begin{lemma}[9.4.7.1]
\label{IV.9.4.7.1}
Sea $S$ un preesquema íntegro noetheriano, de punto genérico $\eta$, sean $X$, $Y$ dos $S$-preesquemas de tipo finito sobre $S$, $g:X\to Y$ un $S$-morfismo y $\sh{F}$ un $\sh{O}_X$-Módulo coherente.
Si $\sh{F}_\eta$ no es $g_\eta$-plano, entonces existe un entorno $U$ de $\eta$ en $S$ tal que, para todo $s\in U$, $\sh{F}_s$ no es $g_s$-plano.
\end{lemma}

Teniendo en cuenta~\sref{IV.2.1.2} y Bourbaki, \emph{Alg.\ comm.}, cap.~I\textsuperscript{er}, \textsection2, n\textsuperscript{o}~3, Remarque~1, la hipótesis significa que existe un abierto no vacío $V$ de $Y_\eta$, y un homomorfismo inyectivo $v:\sh{M}\to\sh{N}$ de $\sh{O}_V$-Módulos \emph{coherentes}, tal que el homomorfismo $1\otimes v:\sh{F}_\eta\otimes_{\sh{O}_V}\sh{M}\to\sh{F}_\eta\otimes_{\sh{O}_V}\sh{N}$ no es inyectivo.
Tenemos $V=Y_\eta\cap W$, donde $W$ es abierto en $Y$~(\textbf{I},~3.6.1), y se sigue de~\sref{IV.8.5.2},~(i) y~(ii)), aplicado según el método de~\sref{IV.8.1.2},~$a$)), que existe un entorno $U_0$ de $\eta$ en $S$, dos $\sh{O}_Z$-Módulos
\oldpage[IV-3]{65}
coherentes $\sh{G}$, $\sh{H}$ (donde $Z=W\cap h^{-1}(U_0)$, siendo $h:Y\to S$ el morfismo estructural) y un $\sh{O}_Z$-homomorfismo $u:\sh{G}\to\sh{H}$ tales que $\sh{M}=\sh{G}_\eta$, $\sh{N}=\sh{H}_\eta$, y $v=u_\eta$; podemos por tanto tomar $U_0$ de modo que $u_s:\sh{G}_s\to\sh{H}_s$ sea inyectivo para todo $s\in U_0$~\sref{IV.9.4.5}.
Pero para todo $s\in U_0$, el homomorfismo $1\otimes u_s:\sh{F}_s\otimes_{\sh{O}_{Z_s}}\sh{G}_s\to\sh{F}_s\otimes_{\sh{O}_{Z_s}}\sh{H}_s$ no es otro que $(1\otimes u)_s$; la hipótesis de que $(1\otimes u)_\eta$ no es inyectivo implica, pues, por~\sref{IV.9.4.5}, la existencia de un abierto no vacío $U\subset U_0$ tal que, para todo $s\in U$, $(1\otimes u)_s$ no es inyectivo, y por tanto $\sh{F}_s$ no es $g_s$-plano para todo $s\in U$.

$4^\circ$ $\eta\in S-E'$.
Está claro que $S-E\subset S-E'$, y si $\eta\in S-E$, $S-E'$ es \emph{a fortiori} un entorno de $\eta$ por $3^\circ$.
Supongamos, pues, que $\eta\in E$, de donde $\sh{F}_\eta$ es localmente libre; estas hipótesis implican que $X_\eta$ no es conexo, y que los rangos del $\sh{O}_{X_\eta}$-Módulo localmente libre $\sh{F}_\eta$ no son los mismos sobre las diversas componentes conexas de $X_\eta$.
Ahora bien, se sigue de \sref{IV.8.4.2}, aplicado según el método de \sref{IV.8.1.2},~$a)$, que podemos suponer (tras sustituir $S$ por un entorno de $\eta$) que $X$ y $X_\eta$ tienen el mismo número de componentes conexas, siendo las componentes conexas de $X_\eta$ las intersecciones de $X_\eta$ con las componentes conexas de $X$.
La conclusión se sigue entonces del razonamiento dado en $2^\circ$, aplicado a cada una de las componentes conexas de $X$ (que son en número finito).

\begin{remark}[9.4.7.2]
\label{IV.9.4.7.2}
El lema~\sref{IV.9.4.7.1} será generalizado más adelante y liberado de las hipótesis noetherianas~\sref{IV.11.2.8}.
\end{remark}

\begin{proposition}[9.4.8]
\label{IV.9.4.8}
Sea $S$ un preesquema localmente noetheriano, $f:X\to S$ un morfismo de tipo finito y $\sh{F}$ un $\sh{O}_X$-Módulo coherente.
Supongamos que, para todo $s\in S$, $X_s$ es un preesquema localmente íntegro.
Entonces el conjunto $E$ \textnormal{(}resp.\ $E'$\textnormal{)} de los $s\in S$ tales que $\sh{F}_s$ es un $\sh{O}_{X_s}$-Módulo de torsión \textnormal{(}resp.\ un $\sh{O}_{X_s}$-Módulo sin torsión\textnormal{)} es localmente construible.
\end{proposition}

\begin{proof}
Podemos suponer evidentemente que $S=\operatorname{Spec}(A)$ es afín y noetheriano y demostrar que $E$ (resp.\ $E'$) es construible utilizando el criterio~($\mathbf{0}_{\mathrm{III}}$,~9.2.3); sustituyendo $S$ por el subpreesquema cerrado reducido de $S$ cuyo espacio subyacente es una parte cerrada irreducible $Y$ de $S$, y observando~(\textbf{I},~3.6.4) que para $s\in Y$, la fibra $(X_{(Y)})_s$ se identifica canónicamente con $X_s$ y el haz $(\sh{F}\otimes_{\sh{O}_S}\sh{O}_Y)_s$ con $\sh{F}_s$, vemos que nos reducimos al caso en que $S$ es íntegro de punto genérico $\eta$, y a demostrar que $E$ o $S-E$ (resp.\ $E'$ o $S-E'$) es un entorno de $\eta$ en $S$.
Además, $X$ es una unión finita de abiertos afines $U_i$ de tipo finito sobre $S$, y cada $(U_i)_s$ es un subpreesquema abierto de $X_s$, por tanto localmente íntegro. Si $(U_i)_\eta$ es vacío (resp.\ no vacío), entonces $(U_i)_s$ es también vacío (resp.\ no vacío) en un entorno de $\eta$ \sref{IV.9.2.6}. Podemos, por tanto, suponer que todos los $(U_i)_\eta$ son no vacíos e íntegros; entonces $\sh{F}_s$ es de torsión (resp.\ sin torsión) si y solo si cada restricción $\sh{F}_s|_{(U_i)_s}$ es de torsión (resp.\ sin torsión).

Nos reducimos así al caso en que $X=\operatorname{Spec}(B)$ es afín y $\sh{F}=\widetilde{M}$, siendo $M$ un $B$-módulo finito. Pongamos $B_s=B\otimes_A\mathbf{k}(s)$ y $M_s=M\otimes_A\mathbf{k}(s)$, y supongamos $B_\eta$ íntegro. Hay cuatro casos.

$1^\circ$ Supongamos $\eta\in E$. Entonces $M_\eta$ es un $B_\eta$-módulo finito de torsión, por lo que existe un $h\neq 0$ en $B_\eta$ tal que $hM_\eta=0$. Por \sref{IV.8.5.2},~(i), aplicado como en \sref{IV.8.1.2},~$a)$,

\oldpage[IV-3]{66}

tras restringir $S$ podemos escribir $h=g_\eta$ para cierto $g\in\Gamma(X,\sh{O}_X)$. Sea $u:\sh{F}\to\sh{F}$ la multiplicación por $g$. Como $u_\eta=0$, \sref{IV.9.4.5} muestra que $u_s=0$ cerca de $\eta$. Por otra parte, la multiplicación por $g_\eta=h$ es inyectiva sobre $\sh{O}_{X_\eta}$ porque $B_\eta$ es íntegro y $h\neq 0$; de nuevo por \sref{IV.9.4.5}, la multiplicación por $g_s$ sobre $\sh{O}_{X_s}$ es inyectiva cerca de $\eta$. Así, $g_s$ es $\sh{O}_{X_s}$-regular y anula $\sh{F}_s$, por lo que $\sh{F}_s$ es de torsión cerca de $\eta$.

$2^\circ$ Supongamos $\eta\in S-E$. Si $T$ es el submódulo de torsión de $M_\eta$, entonces $M_\eta/T\neq 0$ y, siendo un $B_\eta$-módulo finito sin torsión, se sumerge en $B_\eta^n$ para cierto $n$. Existe, por tanto, un homomorfismo no nulo $w:M_\eta\to B_\eta$. Por \sref{IV.8.5.2},~(i), tras restringir $S$ se extiende a un homomorfismo $v:\sh{F}\to\sh{O}_X$. Como $v_\eta\neq 0$, \sref{IV.9.4.5} da $v_s\neq 0$ cerca de $\eta$; como $X_s$ es localmente íntegro, $\sh{F}_s$ no es de torsión para estos $s$.

$3^\circ$ Supongamos $\eta\in E'$. Como $M_\eta$ es un $B_\eta$-módulo finito sin torsión, existe un homomorfismo inyectivo $w:M_\eta\to B_\eta^n$. Utilizando \sref{IV.8.5.2},~(i), y \sref{IV.9.4.5} como en $2^\circ$, tras restringir $S$ obtenemos un homomorfismo $v:\sh{F}\to\sh{O}_X^n$ tal que $v_\eta=w$ y $v_s$ es inyectivo cerca de $\eta$. Así, $\sh{F}_s$ no tiene torsión cerca de $\eta$.

$4^\circ$ Supongamos $\eta\in S-E'$. Sea $T$ el submódulo de torsión de $M_\eta$. Entonces $T\neq 0$ y es finito. Por \sref{IV.8.5.2},~(i) y (ii), tras restringir $S$ existen un $\sh{O}_X$-Módulo coherente $\sh{G}$ y un homomorfismo inyectivo $u:\sh{G}\to\sh{F}$ tales que $\sh{G}_\eta=\widetilde{T}$ y $u_\eta$ es la inclusión canónica. Por $1^\circ$ y \sref{IV.1.8.6}, cerca de $\eta$ el módulo $\sh{G}_s$ es un $\sh{O}_{X_s}$-Módulo de torsión no nulo; por \sref{IV.9.4.5}, $u_s$ es inyectivo cerca de $\eta$. Por tanto, el submódulo de torsión de $\sh{F}_s$ es no nulo cerca de $\eta$.
\end{proof}

\begin{remark}[9.4.9]
\label{IV.9.4.9}
La propiedad ``$X$ es un $k$-preesquema algebraico localmente íntegro'' no satisface la condición \sref{IV.9.2.1},~(i), y por tanto no es seguro que el enunciado de \sref{IV.9.4.8} siga siendo válido cuando no se impone ninguna hipótesis sobre $S$ y se supone únicamente que $f$ es un morfismo de presentación finita y $\sh{F}$ es un $\sh{O}_X$-Módulo de presentación finita. Consideremos, sin embargo, el caso particular siguiente. Sea $S_0$ un preesquema localmente noetheriano, sea $f_0:X_0\to S_0$ un morfismo de tipo finito cuyas fibras $(X_0)_{s_0}$ son localmente íntegras para todo $s_0\in S_0$, y sea $\sh{F}_0$ un $\sh{O}_{X_0}$-Módulo coherente. Sea $g:S\to S_0$ un morfismo cualquiera, pongamos
$X=X_0\times_{S_0}S$, $\sh{F}=\sh{F}_0\otimes_{\sh{O}_{S_0}}\sh{O}_S$, y $f=f_0\times 1_S:X\to S$, y supongamos que, para todo $s\in S$, la fibra $X_s$ es también localmente íntegra. Entonces el conjunto $E$ (resp. $E'$) de los $s\in S$ tales que $\sh{F}_s$ es de torsión (resp. sin torsión)

\oldpage[IV-3]{67}

es también localmente construible.

En efecto, sea $s\in S$ y pongamos $s_0=g(s)$. Teniendo en cuenta \sref[I]{I.1.8.2}, basta demostrar que $\sh{F}_s$ es de torsión (resp. sin torsión) si y solo si $(\sh{F}_0)_{s_0}$ lo es. Ahora bien, $\operatorname{Supp}(\sh{F}_s)$ es la imagen inversa de $\operatorname{Supp}((\sh{F}_0)_{s_0})$ por la proyección $p:X_s\to(X_0)_{s_0}$ \sref[I]{I.9.1.13}. Como $p$ es fielmente plano y cuasicompacto, decir que $\operatorname{Supp}(\sh{F}_s)$ contiene un punto maximal de $X_s$ equivale a decir que $\operatorname{Supp}((\sh{F}_0)_{s_0})$ contiene un punto maximal de $(X_0)_{s_0}$ \sref[I]{I.1.1.5} y \sref{IV.2.3.4}; esto demuestra la afirmación relativa a $E$ \sref[I]{I.7.4.6}. Si el submódulo de torsión $\sh{T}$ de $(\sh{F}_0)_{s_0}$ no es nulo, entonces $\sh{T}\otimes_{\boldsymbol{k}(s_0)}\boldsymbol{k}(s)$, que es de torsión por lo que precede, no es nulo y se identifica con un submódulo de $\sh{F}_s$ \sref{IV.2.2.7}; por tanto, el submódulo de torsión de $\sh{F}_s$ no es nulo. Finalmente, si $(\sh{F}_0)_{s_0}$ no tiene torsión, entonces, tras restringirnos a un abierto afín de $(X_0)_{s_0}$, podemos suponer que es isomorfo a un submódulo de algún $\sh{O}_{(X_0)_{s_0}}^n$; por consiguiente, $\sh{F}_s$ es isomorfo a un submódulo de $\sh{O}_{X_s}^n$ \sref{IV.2.2.7}, lo que demuestra la afirmación relativa a $E'$.
\end{remark}

\subsection{Construibilidad de propiedades topológicas}
\label{subsection:IV.9.5}

\begin{proposition}[9.5.1]
\label{IV.9.5.1}
Sea $f:X\to S$ un morfismo de presentación finita y sea $Z$ una parte localmente construible de $X$.
Entonces el conjunto $E$ de los $s\in S$ tales que $Z_s\neq\varnothing$ es localmente construible.
\end{proposition}

\begin{proof}
En efecto, tenemos $E=f(Z)$, y basta aplicar el teorema de Chevalley \sref{IV.1.8.4}.
\end{proof}

\begin{corollary}[9.5.2]
\label{IV.9.5.2}
Si $Z'$, $Z''$ son dos partes localmente construibles de $X$, el conjunto de los $s\in S$ tales que $Z'_s\subset Z''_s$ (resp.\ $Z'_s=Z''_s$) es localmente construible.
\end{corollary}

\begin{proof}
En efecto, la relación $Z'_s\subset Z''_s$ equivale a $(Z'\cap(\complement Z''))_s=\varnothing$, y $Z'\cap\complement Z''$ es localmente construible.
\end{proof}

\begin{proposition}[9.5.3]
\label{IV.9.5.3}
Sea $f:X\to S$ un morfismo de presentación finita, y sean $Z$, $Z'$ dos partes localmente construibles de $X$ tales que $Z\subset Z'$.
Entonces el conjunto $E$ de los $s\in S$ tales que $Z_s$ es denso en $Z'_s$ es localmente construible en $S$.
\end{proposition}

\begin{proof}
Debemos verificar las dos condiciones de \sref{IV.9.2.2},~(iii).
En cuanto a la primera, consideremos un preesquema $X$ algebraico sobre un cuerpo $k$, dos partes construibles $Z$, $Z'$ de $X$ tales que $Z\subset Z'$, y una extensión $k'$ de $k$.
Entonces la proyección canónica $p:X_{(k')}\to X$ es un morfismo fielmente plano y cuasicompacto, y tenemos $p^{-1}(\overline{Z})=\overline{p^{-1}(Z)}$ y $p^{-1}(\overline{Z'})=\overline{p^{-1}(Z')}$; en virtud de \sref{IV.2.3.10}, como $p$ es sobreyectivo, la relación $\overline{Z}=\overline{Z'}$ equivale a $\overline{p^{-1}(Z)}=\overline{p^{-1}(Z')}$.

Verifiquemos la segunda condición y supongamos, por tanto, que $S$ es afín, noetheriano, íntegro, de punto genérico $\eta$.
Distinguimos dos casos:
\begin{enumerate}[label=\upshape{$\arabic{enumi}^\circ$}]
  \item $\eta\in S-E$, es decir, $Z_\eta$ no es denso en $Z'_\eta$; existe entonces en $X$ un abierto $V$ tal que $V\cap Z_\eta=\varnothing$ y $V\cap Z'_\eta\neq\varnothing$.
  Como $X$ es noetheriano, $V$ es localmente construible, y por tanto también lo es $V\cap Z$; en virtud de \sref{IV.9.5.1}, existe un entorno $U$ de $\eta$ en $S$ tal que, para todo $s\in U$, $(V\cap Z)_s=\varnothing$ y $(V\cap Z')_s\neq\varnothing$; esto implica que $Z_s$ no es denso en $Z'_s$ para $s\in U$, es decir, $U\subset S-E$.

\oldpage[IV-3]{68}

  \item\label{IV.9.5.3.case2} $\eta\in E$, por lo que $Z_\eta$ es denso en $Z'_\eta$.
  Demostremos primero que podemos suponer $Z'$ \emph{cerrado}.
  En efecto, $Z_\eta$ es denso en $\overline{Z'_\eta}$ (adherencia tomada en $X_\eta$); poniendo $V_\eta=X_\eta-\overline{Z_\eta}$, que es abierto en $X_\eta$ y no encuentra $Z'_\eta$, podemos suponer $V_\eta$ de la forma $V\cap X_\eta$, donde $V$ es abierto (y por tanto construible) en $X$, y la hipótesis $V_\eta\cap Z'_\eta=\varnothing$ implica entonces $V_s\cap Z'_s=\varnothing$ para todo $s$ en un entorno de $\eta$ en $S$, en virtud de \sref{IV.9.5.1}.
  Sustituyendo $S$ por un entorno abierto de $\eta$, podemos, por tanto, suponer $V\cap Z'=\varnothing$, de donde $V\cap\overline{Z'}=\varnothing$ (adherencia tomada en $X$), y por consiguiente $(\overline{Z'})_\eta=\overline{Z'_\eta}$, lo que demuestra nuestra afirmación.
  Como la descomposición de $Z'$ en sus componentes irreducibles es una unión finita, y restringiendo aún más $S$ a un entorno de $\eta$, podemos suponer que todas las componentes irreducibles $Z'_i$ de $Z'$ encuentran $X_\eta$, de donde se sigue que $X_\eta$ contiene los puntos genéricos de los $Z'_i$ \sref[0]{0.2.1.8}.
  Decir que $Z_s$ es denso en $Z'_s$ equivale a decir que cada uno de los $(Z\cap Z'_i)_s$ es denso en $(Z'_i)_s$, y nos reducimos así al caso en que $Z'$ es \emph{irreducible}.
  Sustituyendo entonces $X$ por el subpreesquema cuyo espacio subyacente reducido es $Z'$, vemos que podemos suponer $Z'=X$ y que $X$ es \emph{íntegro} y domina $S$.
  Finalmente, recubriendo $X$ por un número finito de abiertos afines $W_j$ y sustituyendo $Z$ por $Z\cap W_j$, podemos suponer que $X=\operatorname{Spec}(A)$, donde $A$ es un anillo noetheriano íntegro.
  Como $X_\eta$ es noetheriano íntegro y $Z_\eta$ es construible y denso en $X_\eta$, $Z_\eta$ contiene un abierto no vacío de $X_\eta$~($\mathbf{0}_{\mathrm{III}}$,~9.2.2), que puede tomarse de la forma $(D(t))_\eta$, donde $t$ es un elemento no nulo de $A$.
  Sustituyendo $S$ por un entorno de $\eta$ en $S$ si es necesario, en virtud de la relación $(D(t))_\eta\subset Z_\eta$, podemos suponer además que $D(t)\subset Z$ \sref{IV.9.5.2}.
  Finalmente, como la homotecia de razón $t_\eta$ en $\sh{O}_{X_\eta}$ es inyectiva, se sigue de \sref{IV.9.4.5} que para $s$ próximo a $\eta$, $t_s$ es $\sh{O}_{X_s}$-regular, de modo que $(X_s)_{t_s}$ es denso en $X_s$, y \emph{a fortiori} lo mismo vale para $Z_s$, que contiene $(X_s)_{t_s}$.
\end{enumerate}
\end{proof}

\begin{corollary}[9.5.4]
\label{IV.9.5.4}
Sea $f:X\to S$ un morfismo de presentación finita y sea $Z$ una parte localmente construible de $X$.
El conjunto $E$ de los $s\in S$ tales que $Z_s$ es cerrado (resp.\ abierto, resp.\ localmente cerrado) en $X_s$ es localmente construible en $S$.
\end{corollary}

\begin{proof}
Decir que $Z_s$ es abierto en $X_s$ significa que $(X-Z)_s=X_s-Z_s$ es cerrado en $X_s$, y como $X-Z$ es localmente construible, podemos limitarnos a considerar el conjunto de los $s\in S$ tales que $Z_s$ es cerrado y el conjunto de los $s\in S$ tales que $Z_s$ es localmente cerrado.

Verifiquemos en cada caso las dos condiciones de \sref{IV.9.2.2},~(ii).
La primera resulta del hecho de que $p:X_{(k')}\to X$ es fielmente plano y cuasicompacto, y de \sref{IV.2.3.12} y \sref{IV.2.3.14}.
Verifiquemos, por tanto, la segunda condición, suponiendo $S$ afín, noetheriano, íntegro, de punto genérico $\eta$.
Pongamos $Z'=\overline{Z}$; el mismo razonamiento que en \sref{IV.9.5.3} muestra que $Z'_\eta$ es igual a la adherencia de $Z_\eta$ en $X_\eta$; en virtud de \sref{IV.9.5.3}, existe, pues, un entorno $U$ de $\eta$ tal que para $s\in U$, $Z_s$ es denso en $Z'_s$, siendo este último cerrado.
Decir que $Z_s$ es cerrado en $X_s$ significa que $Z''_s=\varnothing$, donde $Z''=Z'-Z$; se sigue de \sref{IV.9.5.1} que el conjunto $E$ de los $s\in S$ donde $Z''_s=\varnothing$ es tal que $E$ o $S-E$ contiene un entorno de $\eta$.
Decir que $Z_s$ es localmente cerrado en $X_s$ significa que $Z''_s$ es cerrado en $X_s$; basta, por tanto, aplicar el resultado precedente sustituyendo $Z$ por $Z''$, que es localmente construible en $X$.
\end{proof}

\oldpage[IV-3]{69}

\begin{proposition}[9.5.5]
\label{IV.9.5.5}
Sea $f:X\to S$ un morfismo de presentación finita y sea $Z$ una parte localmente construible de $X$ tal que, para todo $s\in S$, $Z_s$ es localmente cerrado en $X_s$.
Para todo $s\in S$, sea $D(s)\subset\mathbf{Z}\cup\{-\infty\}$ el conjunto de las dimensiones de las componentes irreducibles de $Z_s$.
Entonces la función $s\mapsto D(s)$ es localmente construible en $S$.
\end{proposition}

\begin{proof}
Sea $\Phi$ una parte finita de $\mathbf{Z}\cup\{-\infty\}$; se trata de demostrar que el conjunto de los $s\in S$ tales que $D(s)=\Phi$ es localmente construible; teniendo en cuenta \sref{IV.9.2.3}, debemos verificar una vez más las dos condiciones de \sref{IV.9.2.2},~(iii).

En cuanto a la primera, se trata de ver que si $X$ es un preesquema algebraico sobre un cuerpo $k$, $Z$ una parte localmente cerrada de $X$, $k'$ una extensión de $k$, y $p:X_{(k')}\to X$ la proyección canónica, entonces el conjunto de las dimensiones de las componentes irreducibles de $Z$ es el mismo que el de las dimensiones de las componentes irreducibles de $p^{-1}(Z)$; teniendo en cuenta la existencia de un subpreesquema de $X$ cuyo espacio subyacente es $Z$ \sref[I]{I.5.2.1}, esto se sigue de \sref{IV.4.2.8}.

Para la segunda condición de \sref{IV.9.2.2},~(iii), estamos en el caso en que $S$ es noetheriano, íntegro, de punto genérico $\eta$, y $f:X\to S$ es un morfismo de tipo finito, de modo que $X$ es noetheriano.
El subespacio $Z_\eta$ del espacio noetheriano $X_\eta$ es, por hipótesis, localmente cerrado, por lo que tiene un número finito de componentes irreducibles $Z_{i\eta}$, que son localmente cerradas en $X_\eta$.
Se sigue, para cada índice $i$, que existe una parte localmente cerrada $Z_i$ de $X$ tal que $(Z_i)_\eta=Z_{i\eta}$, por lo que si $Z'=\bigcup_i Z_i$, tenemos $Z'_\eta=Z_\eta$.
Pero como $Z$ y $Z'$ son localmente construibles, podemos, sustituyendo $S$ por un entorno de $\eta$, suponer que $Z'=Z$ \sref{IV.9.5.1}.
Además, para $i\neq j$, $Z_{i\eta}\cap Z_{j\eta}$ es raro y \emph{cerrado} en $Z_{i\eta}$; por tanto, mediante \sref{IV.9.5.3} y \sref{IV.9.5.4}, podemos también suponer, restringiendo $S$, que para $j\neq i$, $(Z_i)_s\cap(Z_j)_s$ es raro y cerrado en $(Z_i)_s$.
Como $Z_i$ es localmente cerrado en $X$, existe un subpreesquema de $X$ cuyo espacio subyacente es $Z_i$ (y que seguimos denotando por $Z_i$), que es de tipo finito sobre $S$ \sref[I]{I.6.3.5}.
Pongamos $U_i=Z_i-\bigcup_{j\neq i}(Z_i\cap Z_j)$, que es abierto en $Z_i$ y tal que $(U_i)_s$ es abierto y denso en todas partes en $(Z_i)_s$; además, por construcción, los $U_i$ no tienen dos a dos ningún punto común.
Como $(Z_i)_s$ es un $\boldsymbol{k}(s)$-preesquema algebraico, el conjunto de las dimensiones de las componentes irreducibles de $Z_s=\bigcup_i(Z_i)_s$ es el mismo que el conjunto de las dimensiones de las componentes irreducibles de la reunión de los $(U_i)_s$ \sref{IV.4.1.1.3}, siendo cada una de estas componentes ya una componente irreducible de uno de los $(U_i)_s$.
Podemos, por tanto, limitarnos al caso en que $Z=U_i$; además, como $Z_\eta$ es entonces irreducible, solo hay una componente irreducible de $Z$ que encuentra $X_\eta$~\sref[0]{0.2.1.8}, y podemos evidentemente, restringiendo $Z$, suponer $Z$ irreducible.
Nos reducimos así finalmente a demostrar el siguiente caso particular de \sref{IV.9.5.5}:
\end{proof}

\begin{corollary}[9.5.6]
\label{IV.9.5.6}
Sea $S$ un preesquema noetheriano, irreducible, de punto genérico $\eta$, sea $X$ un preesquema irreducible y sea $f:X\to S$ un morfismo dominante de tipo finito.
Entonces existe un entorno $U$ de $\eta$ en $S$ tal que, para todo $s\in U$, todas las componentes irreducibles de $X_s$ son de dimensión $n=\dim(X_\eta)$.
\end{corollary}

\begin{proof}
Podemos evidentemente reducirnos al caso en que $S=\operatorname{Spec}(A)$ es afín, siendo por tanto $A$ noetheriano; sustituyendo $f$ por $f_{\mathrm{red}}$, que es de tipo finito \sref[I]{I.5.4},~(vii),

\oldpage[IV-3]{70}

podemos suponer $A$ íntegro y $X$ reducido, por tanto íntegro.
Sabemos por \sref{IV.4.1.2} que existe un abierto $W$ denso en $X_\eta$ y un $\boldsymbol{k}(\eta)$-morfismo \emph{finito} y \emph{sobreyectivo} $h:W\to\mathbf{V}((\boldsymbol{k}(\eta))^n)=\operatorname{Spec}(B')$, donde $B'=\boldsymbol{k}(\eta)[\mathrm{T}_1,\ldots,\mathrm{T}_n]$ (siendo los $\mathrm{T}_i$ indeterminados).
Si $V$ es un abierto de $X$ tal que $V\cap X_\eta=W$, sabemos por \sref{IV.9.5.3} que para $s$ próximo a $\eta$ en $S$, $V_s$ es un abierto denso en $X_s$, y podemos por tanto, por \sref{IV.4.1.1.3}, reducirnos al caso en que $V=X$, $W=X_\eta$.
Pongamos $B=A[\mathrm{T}_1,\ldots,\mathrm{T}_n]$ y $Y=\operatorname{Spec}(B)=\mathbf{V}_A^n$, de modo que $\operatorname{Spec}(B')=Y_\eta$; se sigue de \sref{IV.8.8.2},~(i) y \sref{IV.8.10.5},~(vi) y (x), aplicados según el método de \sref{IV.8.1.2},~\emph{a}), que, sustituyendo $S$ por un entorno de $\eta$, podemos suponer que $h=g_\eta$, donde $g:X\to Y$ es un morfismo finito sobreyectivo; en otras palabras, tenemos $X=\operatorname{Spec}(C)$, donde $C$ es una B-álgebra finita y el homomorfismo $B\to C$ correspondiente a $g$ es \emph{inyectivo}; como $C$ es un anillo íntegro, $C$ es por tanto un B-módulo de tipo finito \emph{sin torsión}, y $C_\eta=C\otimes_A\boldsymbol{k}(\eta)$ es un módulo de tipo finito sin torsión sobre $B_\eta=B'=B\otimes_A\boldsymbol{k}(\eta)=\boldsymbol{k}(\eta)[\mathrm{T}_1,\ldots,\mathrm{T}_n]$ (pues es un anillo de fracciones cuyos denominadores pertenecen a $A-\{0\}\subset B-\{0\}$).
Se sigue de \sref{IV.9.4.8} que existe un entorno $U$ de $\eta$ en $S$ tal que para $s\in S$, $C_s=C\otimes_A\boldsymbol{k}(s)$ es un módulo de tipo finito \emph{sin torsión} sobre $B_s=B\otimes_A\boldsymbol{k}(s)=\boldsymbol{k}(s)[\mathrm{T}_1,\ldots,\mathrm{T}_n]$, y en particular el homomorfismo $g_s:B_s\to C_s$ es inyectivo.
Como ningún elemento de $B_s$ es divisor de cero en $C_s$, para todo ideal primo minimal $\mathfrak{p}_i$ de $C_s$ (cuyos elementos son divisores de cero en $C_s$), tenemos necesariamente $\mathfrak{p}_i\cap B_s=0$, de donde el homomorfismo canónico $B_s\to C_s/\mathfrak{p}_i$ es inyectivo.
Deducimos que para cada componente irreducible $Z_i=\operatorname{Spec}(C_s/\mathfrak{p}_i)$ de $X_s$, la restricción a $Z_i$ de $g$ es un morfismo finito y dominante $Z_i\to Y_s$, por tanto \emph{sobreyectivo} \sref[II]{II.6.1.10}.
Concluimos entonces de \sref{IV.4.1.2} que $\dim Z_i=n$, lo que termina la demostración.
\end{proof}

\begin{remark}[9.5.7]
\label{IV.9.5.7}
Conviene observar que bajo las hipótesis de \sref{IV.9.5.6}, puede ocurrir que $X_s$ no sea irreducible para \emph{ningún} $s\neq\eta$ en un entorno de $\eta$; en otras palabras, la propiedad ``$X$ es un $k$-preesquema algebraico irreducible'' \emph{no es construible}.
Tomemos, por ejemplo, $S=\operatorname{Spec}(k[\mathrm{T}])$, donde $k$ es un cuerpo algebraicamente cerrado y $\mathrm{T}$ un indeterminado; tenemos entonces $\boldsymbol{k}(\eta)=K=k(\mathrm{T})$.
Sea $L$ una extensión separable finita de $K$ de grado $>1$, y sea $X$ la \emph{clausura integral} de $S$ en $L$ \sref[II]{II.6.3.4}; tenemos entonces $X=\operatorname{Spec}(B)$, donde $B$ es la clausura integral de $k[\mathrm{T}]$ en $L$.
Sabemos que $B$ es un anillo de Dedekind, y que todos salvo un número finito de ideales maximales de $k[\mathrm{T}]$ no están ramificados en $B$. Además, el cuerpo residual de todo ideal maximal de $B$ es necesariamente $k$, puesto que es una extensión finita del cuerpo algebraicamente cerrado $k$. Así, para casi todos los ideales maximales $\mathfrak{j}_s$ de $k[\mathrm{T}]$, el anillo $B_s=B/\mathfrak{j}_s B$ es un producto directo de $[L:K]$ copias de $k$; en particular, $X_s$ no es irreducible, aunque $X_\eta=\operatorname{Spec}(L)$ sí lo es.
El mismo ejemplo muestra que la propiedad ``$X$ es un $k$-preesquema algebraico íntegro'' no es construible.
Finalmente, la propiedad ``$X$ es un $k$-preesquema algebraico reducido'' tampoco es construible.
Para ello basta tomar de nuevo $S=\operatorname{Spec}(k[\mathrm{T}])$, donde esta vez $k$ es un cuerpo algebraicamente cerrado de característica $p>0$, y como $X$ la clausura integral de $S$ en $L=K^{1/p}$ (donde $K=k(\mathrm{T})$), de modo que $X=\operatorname{Spec}(B)$ con $B=k[\mathrm{T}^{1/p}]$ (siendo $k$ perfecto); todo ideal maximal de $k[\mathrm{T}]$ es de la forma $(\mathrm{T}-\alpha)$ con $\alpha\in k$; el único ideal de $B$ sobre el ideal $(\mathrm{T}-\alpha)$ es el ideal principal $(\mathrm{T}^{1/p}-\alpha^{1/p})$ y es inmediato
\oldpage[IV-3]{71}
que el anillo cociente $B/(\mathrm{T}-\alpha)B$ tiene elementos nilpotentes; en otras palabras, $X_s$ no es reducido para ningún $s\neq\eta$, mientras que $X_\eta=\operatorname{Spec}(L)$ es íntegro.

Veremos un poco más adelante \sref{IV.9.7} que, por el contrario, se obtienen propiedades construibles cuando se consideran las nociones ``geométricas'' correspondientes a las nociones de preesquema irreducible, reducido o íntegro \sref{IV.4.5} y \sref{IV.4.6}.
\end{remark}

\subsection{Construibilidad de ciertas propiedades de morfismos}
\label{subsection:IV.9.6}

\begin{proposition}[9.6.1]
\label{IV.9.6.1}
Sean $X$, $Y$ dos S-preesquemas de presentación finita y sea $f:X\to Y$ un S-morfismo.
Sea $E$ el conjunto de los $s\in S$ para los cuales $f_s$ posee una de las propiedades siguientes: ser
\begin{enumerate}[label=(\roman*),widest=xii]
  \item sobreyectivo;
  \item dominante;
  \item separado;
  \item propio;
  \item radicial;
  \item finito;
  \item cuasifinito;
  \item una inmersión;
  \item una inmersión cerrada;
  \item una inmersión abierta;
  \item un isomorfismo;
  \item un monomorfismo.
\end{enumerate}
Entonces $E$ es localmente construible en $S$.
\end{proposition}

\begin{proof}
Las afirmaciones relativas a (i), (v) y (vii) se incluyen solo como referencia, pues ya fueron establecidas en \sref{IV.9.3.2}.

(ii): Observemos primero que $f$ es de presentación finita \sref[I]{I.6.2},~(v); por tanto, en virtud del teorema de Chevalley \sref{IV.1.8.4}, $f(X)$ es localmente construible en $Y$.
Además, tenemos $f_s(X_s)=(f(X))_s$ \sref[I]{I.3.6.1}; el conjunto de los $s$ tales que $f_s$ es dominante es el conjunto de los $s$ tales que $(f(X))_s$ es denso en $Y_s$; la conclusión se sigue, pues, de \sref{IV.9.5.3}.

(iii): Como $f:X\to Y$ es de presentación finita, la inmersión diagonal $\Delta_f:X\to X\times_Y X$ es de presentación finita \sref[I]{I.6.2},~(iv) y (v); se sigue de \sref{IV.1.8.4.1} que $\Delta_f(X)=\Delta_X$ es una parte localmente construible de $X\times_Y X$.
Decir que $f_s$ es separado significa (teniendo en cuenta \sref[I]{I.5.3.4}) que $(\Delta_X)_s$ es cerrado en $X_s\times_{Y_s} X_s=(X\times_Y X)_s$; la conclusión se sigue de \sref{IV.9.5.4}.

\oldpage[IV-3]{72}

(iv): Demostremos que la propiedad ``$X$ e $Y$ son preesquemas algebraicos sobre un cuerpo $k$, y $f:X\to Y$ un $k$-morfismo propio'' es construible.
La condición \sref{IV.9.2.1},~(i) se verifica en virtud de \sref{IV.2.7.1},~(vii).
Podemos, por tanto, reducirnos al caso en que $S$ es afín, noetheriano, íntegro, de punto genérico $\eta$, y debemos demostrar que $E$ o $S-E$ es un entorno de $\eta$.
Supongamos primero que $\eta\in E$; se sigue de \sref{IV.8.10.5},~(xii), aplicado según el método de \sref{IV.8.1.2},~\emph{a}), que, sustituyendo $S$ por un entorno de $\eta$, podemos
suponer que el propio $f$ es un morfismo propio; sabemos entonces que lo mismo vale para $f_s$ para todo $s\in S$ \sref[II]{II.5.4.2},~(iii).
Supongamos, por el contrario, que $\eta\in S-E$, y distingamos dos casos:
\begin{enumerate}[label=\upshape{$\arabic{enumi}^\circ$}]
  \item Supongamos que $f_\eta$ no es separado; se sigue entonces de (iii) que $f_s$ no es separado (y \emph{a fortiori} no es propio) en un entorno de $\eta$.

  \item Supongamos $f_\eta$ separado; $Y$ es una unión finita de abiertos afines $V_i$, y para que $f_s$ sea propio, basta que cada una de sus restricciones $f_s^{-1}((V_i)_s)\to(V_i)_s$ lo sea \sref[II]{II.5.4.1}; podemos, por tanto, reducirnos al caso en que $Y$ es \emph{afín}, y por tanto un esquema.
  Decir que $f_\eta$ no es propio significa, por \sref[II]{II.5.6.3}, que existe un morfismo \emph{de tipo finito} $h:Z\to Y_\eta$ tal que el morfismo $(f_\eta)_{(Z)}:X_\eta\times_{Y_\eta} Z\to Z$ no es cerrado.
  Como $Y$ es un esquema, deducimos de \sref{IV.8.8.2},~(i) y (ii) (restringiendo $S$, si es necesario, a un entorno de $\eta$) que existe un morfismo de tipo finito $g:Y'\to Y$ tal que $Z$ es isomorfo a $Y'_\eta$ y que $g_\eta=h$;
  si ponemos $X'=X\times_Y Y'$ y $f'=f_{(Y')}:X'\to Y'$, entonces $(f_\eta)_{(Z)}=f'_\eta$, y por hipótesis existe, por tanto, una parte cerrada $M'$ de $X'_\eta$ tal que $f'_\eta(M')$ no es cerrado en $Y'_\eta$.
  Ahora bien, $M'$ es la traza sobre $X'_\eta$ de una parte cerrada $N'$ de $X'$; como $X'$ es noetheriano y $f'$ de tipo finito, $f'(N')$ es construible en $Y'$ \sref{IV.1.8.4} y, por hipótesis, $(f'(N'))_\eta=f'_\eta(N'_\eta)=f'_\eta(M')$ no es cerrado en $Y'_\eta$.
  Concluimos de \sref{IV.9.5.4} que existe un entorno $U$ de $\eta$ en $S$ tal que para $s\in U$, $(f'(N'))_s=f'_s(N'_s)$ no es cerrado en $Y'_s$; en otras palabras, el morfismo $f'_s$ no es cerrado, y \emph{a fortiori} el morfismo $f_s$ no es propio.
\end{enumerate}

(vi): La propiedad es una conjunción de las propiedades (iv) y (vii) \sref{IV.8.11.1}; por tanto, la proposición se sigue de lo ya demostrado.

(ix): La verificación de la condición \sref{IV.9.2.1},~(i) se sigue de \sref{IV.2.7.1},~(xi).
Podemos, por tanto, reducirnos al caso en que $S$ es afín, noetheriano, íntegro, de punto genérico $\eta$, y basta demostrar que $E$ o $S-E$ es un entorno de $\eta$.
Si $\eta\in E$, podemos (sustituyendo $S$ por un entorno de $\eta$) suponer que $f$ es una inmersión cerrada en virtud de \sref{IV.8.10.5},~(iv), y entonces está claro que $f_s$ es una inmersión cerrada para todo $s\in S$ \sref[I]{I.4.3.2}.
Supongamos, pues, que $\eta\in S-E$, y distingamos dos casos:
\begin{enumerate}[label=\upshape{$\arabic{enumi}^\circ$}]
  \item $f_\eta$ no es un morfismo finito.
  Se sigue entonces de (vi) que en un entorno de $\eta$, $f_s$ no es finito, ni \emph{a fortiori} una inmersión cerrada.

  \item $f_\eta$ es finito; entonces, en virtud de \sref{IV.8.10.5},~(x), podemos suponer (restringiendo $S$ si es necesario) que el propio $f$ es un morfismo \emph{finito}.
  En este caso $\mathscr{A}=\mathscr{A}(X)=f_*(\sh{O}_X)$ es un $\sh{O}_Y$-Módulo coherente \sref[II]{II.6.1.3} y $f=\mathscr{A}(u)$, donde $u:\sh{O}_Y\to\mathscr{A}$ es un homomorfismo de $\sh{O}_Y$-Álgebras \sref[II]{II.1.1.2}; como $f_\eta$ no es una inmersión cerrada por hipótesis, $u_\eta$ no es sobreyectivo \sref[II]{II.1.4.10}, por tanto, por \sref{IV.9.4.5}, existe un entorno de $\eta$ en el cual $u_s$ no es sobreyectivo y, en consecuencia, por \sref[I]{I.4.2.3}, $f_s$ no es una inmersión cerrada.
\end{enumerate}

(viii): La verificación de la condición \sref{IV.9.2.1},~(i) se hace como en (ix), utilizando esta vez \sref{IV.2.7.1},~(x) y el hecho de que toda inmersión de un preesquema noetheriano en otro es cuasicompacta.
Nos reducimos, por tanto, al caso en que $S$ es afín, noetheriano, íntegro, de punto genérico $\eta$, y basta demostrar que $E$ o $S-E$ es un entorno de $\eta$.
Si $\eta\in E$, concluimos como en (ix) con ayuda de \sref{IV.8.10.5},~(ii).
Si $\eta\in S-E$, distinguimos de nuevo dos casos:
\begin{enumerate}[label=\upshape{$\arabic{enumi}^\circ$}]
  \item $f_\eta(X_\eta)$ no es una parte localmente cerrada de $Y_\eta$.
  Como $f(X)$ es construible en $Y$ \sref{IV.1.8.4} y $f_s(X_s)=(f(X))_s$, deducimos de \sref{IV.9.5.4} que para $s$ próximo a $\eta$, $f_s(X_s)$ no es localmente cerrado en $Y_s$, y \emph{a fortiori} $f_s$ no es una inmersión.

  \item\oldpage[IV-3]{73} $f_\eta(X_\eta)$ es localmente cerrado en $Y_\eta$.
  Como $f(X)$ es construible en $Y$ \sref{IV.1.8.4} y lo mismo vale para $\overline{f(X)}$, puesto que $Y$ es noetheriano, se sigue de \sref{IV.8.3.11} que, restringiendo $S$ si es necesario, podemos suponer que $f(X)$ es localmente cerrado en $Y$.
  Existe entonces un abierto $V$ de $Y$ que contiene $f(X)$ y en el cual $f(X)$ es cerrado.
  Como $Y$ es noetheriano, $V$ es de tipo finito sobre $S$; sustituyendo $Y$ por $V$, podemos, por tanto, reducirnos al caso en que $f(X)$ es cerrado en $Y$.
  Pero entonces $f_s(X_s)=(f(X))_s$ es cerrado en $Y_s$ para todo $s\in S$, y decir que $f_s$ es una inmersión equivale a decir que es una inmersión cerrada.
  Nos reducimos, por tanto, al caso~(ix).
\end{enumerate}

(x): Utilizando esta vez \sref{IV.2.7.1},~(ix) y \sref{IV.8.10.5},~(iii), nos reducimos al caso en que $S$ es afín, noetheriano, íntegro, de punto genérico $\eta$, y en que $\eta\in S-E$.
Distinguimos tres casos:
\begin{enumerate}[label=\upshape{$\arabic{enumi}^\circ$}]
  \item $f_\eta(X_\eta)$ no es abierto en $Y_\eta$.
  Como $f(X)$ es construible en $Y$ \sref{IV.1.8.4}, deducimos de \sref{IV.9.5.4} que $f_s(X_s)$ no es abierto en $Y_s$ para $s$ próximo a $\eta$, y \emph{a fortiori} $f_s$ no es una inmersión abierta.

  \item $f_\eta(X_\eta)$ es abierto en $Y_\eta$ pero $f_\eta$ no es una inmersión.
  Se sigue entonces de (viii) que para $s$ próximo a $\eta$, $f_s$ no es una inmersión, ni \emph{a fortiori} una inmersión abierta.

  \item $f_\eta(X_\eta)$ es abierto en $Y_\eta$ y $f_\eta$ es una inmersión.
  Como $f(X)$ es construible en $Y$ \sref{IV.1.8.4}, se sigue de \sref{IV.8.3.11} que, restringiendo $S$ si es necesario, podemos suponer ya que $f(X)$ es abierto en $Y$.
  Como $Y$ es noetheriano, el subpreesquema inducido sobre $f(X)$ es de tipo finito sobre $S$, por lo que podemos reducirnos al caso en que $f$ es \emph{sobreyectivo} sustituyendo $Y$ por $f(X)$.
  Por hipótesis, $f_\eta$ es una inmersión cerrada, por lo que, como en (ix), podemos suponer que $f$ es una \emph{inmersión cerrada}; en otras palabras, que $X$ es un subpreesquema cerrado de $Y$ definido por un ideal coherente $\mathscr{J}$ de $\sh{O}_Y$.
  Por hipótesis, $f_\eta$ no es un isomorfismo, de donde $\mathscr{J}_\eta\neq 0$; concluimos por \sref{IV.9.4.5} que en un entorno de $\eta$, $\mathscr{J}_s\neq 0$, y por consiguiente $f_s$ no es una inmersión abierta.
\end{enumerate}

(xi): La propiedad es una conjunción de las propiedades (i) y (x), y por tanto se sigue de lo demostrado.

(xii): En virtud de \sref[I]{I.5.3.8}, decir que $f_s$ es un monomorfismo significa que $\Delta_{f_s}=(\Delta_f)_s:X_s\to X_s\times_{Y_s} X_s=(X\times_Y X)_s$ es un isomorfismo, y como $X\times_Y X$ es un S-preesquema de presentación finita \sref[I]{I.6.2},~(iv), la conclusión se sigue de (xi).
\end{proof}

\begin{proposition}[9.6.2]
\label{IV.9.6.2}
Sean $X$, $Y$ dos S-preesquemas de presentación finita y sea $f:X\to Y$ un S-morfismo.
\begin{enumerate}[label=\upshape{I)}]
  \item Sea $E$ el conjunto de los $s\in S$ para los cuales $f_s$ posee una de las propiedades siguientes: ser
  \begin{enumerate}[label=(\roman*)]
    \item afín;
    \item cuasiafín;
    \item proyectivo;
    \item cuasiproyectivo.
  \end{enumerate}
  Entonces $E$ es ind-construible en $S$.
\end{enumerate}

\oldpage[IV-3]{74}

\begin{enumerate}[label=\upshape{II)}]
  \item Sea $\sh{L}$ un $\sh{O}_X$-Módulo invertible.
  Entonces el conjunto $E'$ de los $s\in S$ tales que $\sh{L}_s$ es un $\sh{O}_{X_s}$-Módulo amplio (resp.\ muy amplio) relativamente a $f_s$, es ind-construible en $S$.
\end{enumerate}
\end{proposition}

\begin{proof}
I) Verifiquemos las condiciones (i) y (ii) de \sref{IV.9.2.1}.
En cuanto a la condición \sref{IV.9.2.1},~(i), se sigue, para las propiedades (i) y (ii), de \sref{IV.2.7.1},~(xiii) y (xiv); para las propiedades (iii) y (iv), se sigue de \sref{IV.9.1.5}.
Verifiquemos entonces la condición \sref{IV.9.2.1},~(ii), suponiendo, por tanto, $S$ noetheriano, íntegro, de punto genérico $\eta$, y que $f_\eta:X_\eta\to Y_\eta$ posee una de las propiedades (i) a (iv) del enunciado.
Aplicando \sref{IV.8.10.5},~(viii), (ix), (xiii) y (xiv) según el método de \sref{IV.8.1.2},~\emph{a}), vemos primero que existe un entorno abierto $U$ de $\eta$ tal que, si $V$ y $W$ son las imágenes inversas respectivas de $U$ en $X$ e $Y$ por los morfismos estructurales, la restricción $V\to W$ de $f$ posee la propiedad (i), (ii), (iii) o (iv) considerada.
La conclusión se sigue entonces del hecho de que todas las propiedades consideradas son estables por cambio de base.

II) Procedemos del mismo modo.
Para la condición \sref{IV.9.2.1},~(i), se sigue esta vez de \sref{IV.2.7.2}.
Para la condición \sref{IV.9.2.1},~(ii), con las mismas notaciones que en I), se sigue de \sref{IV.8.10.5.2} que para un entorno $U$ de $\eta$ en $S$, la restricción $\sh{L}|V$ es amplia (resp.\ muy amplia) relativamente a la restricción $V\to W$ de $f$.
La conclusión se sigue de nuevo de la estabilidad de las propiedades consideradas por cambio de base \sref[II]{II.4.6.13} y \sref[II]{II.4.4.10}.
\end{proof}

Podemos mejorar \sref{IV.9.6.2},~II) bajo ciertas condiciones:

\begin{proposition}[9.6.3]
\label{IV.9.6.3}
Sean $X$, $Y$ dos S-preesquemas de presentación finita, sea $f:X\to Y$ un S-morfismo propio y sea $\sh{L}$ un $\sh{O}_X$-Módulo invertible.
Entonces el conjunto $E$ (resp.\ $E'$) de los $s\in S$ tales que $\sh{L}_s$ es un $\sh{O}_{X_s}$-Módulo amplio (resp.\ muy amplio) relativamente a $f_s:X_s\to Y_s$ es localmente construible en $S$.
\end{proposition}

\begin{proof}
Sea $k$ un cuerpo, sean $Z$, $T$ dos preesquemas algebraicos sobre $k$, sea $\mathscr{M}$ un $\sh{O}_Z$-Módulo invertible y sea $g:Z\to T$ un $k$-morfismo de tipo finito; entonces, si $\mathbf{P}(Z,T,\mathscr{M},g,k)$ denota la relación ``$\mathscr{M}$ es amplio (resp.\ muy amplio) relativamente a $g$'', ya hemos observado en \sref{IV.9.6.2} que $\mathbf{P}(Z,T,\mathscr{M},g,k)$ satisface la condición \sref{IV.9.2.1},~(i) en virtud de \sref{IV.2.7.2}.
Sabemos, por otra parte, que $E$ y $E'$ son ind-construibles.
Queda, por tanto, ver que si $S$ es noetheriano, íntegro, de punto genérico $\eta$ y si $\eta\in S-E$ (resp.\ $\eta\in S-E'$), entonces $S-E$ (resp.\ $S-E'$) contiene un entorno de $\eta$.
Consideraremos separadamente los casos $E$ y $E'$.

\oldpage[IV-3]{75}

\emph{Caso de $E'$.} ---
Como $f$ es separado y $Y$ es cuasicompacto, existe un entero $h$ tal que
$R^qf_*(\sh{L})=0$ para $q>h$ \sref[III]{III.1.4.12}.
Por otra parte, como $f$ es propio y $Y$ es noetheriano, todos los
$R^qf_*(\sh{L})$ son $\sh{O}_Y$-Módulos coherentes \sref[III]{III.3.2.1}.
Se anulan salvo para un número finito de valores de $q$, por lo que el teorema de planitud genérica
\sref{IV.6.9.1} muestra que, tras restringir $S$ a un entorno de $\eta$, podemos
suponer que $\sh{L}$ y todos los $R^qf_*(\sh{L})$ son planos sobre $S$.
Se sigue entonces de \sref[III]{III.6.9.9} que el homomorfismo canónico
\begin{equation}
\label{IV.9.6.3.1}
 f_*(\sh{L})\otimes_{\sh{O}_S}\boldsymbol{k}(s)
 \longrightarrow (f_s)_*(\sh{L}_s)
 \tag{9.6.3.1}
\end{equation}
es un isomorfismo.

Se sigue de \sref[II]{II.4.4.4} que decir que $\sh{L}_s$ no es muy amplio relativamente a $f_s$ significa que: \emph{o bien} el homomorfismo canónico $(f_s)^*((f_s)_*(\sh{L}_s))\to\sh{L}_s$ no es sobreyectivo; \emph{o bien} el homomorfismo precedente es sobreyectivo y el morfismo canónico $r:X_s\to\mathbf{P}((f_s)_*(\sh{L}_s))$ no es una inmersión.
Teniendo en cuenta el isomorfismo \sref{IV.9.6.3.1}, estas condiciones se escriben respectivamente de la forma: 1$^\circ$ el homomorfismo canónico $(f^*(f_*(\sh{L})))_s\to\sh{L}_s$ no es sobreyectivo; 2$^\circ$ el homomorfismo precedente es sobreyectivo y el morfismo canónico $X_s\to\mathbf{P}((f_*(\sh{L}))_s)$ no es una inmersión.

Supongamos primero que el homomorfismo canónico $(f^*(f_*(\sh{L})))_\eta\to\sh{L}_\eta$ no es sobreyectivo.
Como $f_*(\sh{L})$ es coherente, también lo es $f^*(f_*(\sh{L}))$, y se sigue entonces de \sref{IV.9.4.5} que para todo $s$ en un entorno de $\eta$, el homomorfismo $(f^*(f_*(\sh{L})))_s\to\sh{L}_s$ no es sobreyectivo, lo que demuestra en este caso que $S-E'$ es un entorno de $\eta$.

Supongamos ahora que el homomorfismo canónico $(f^*(f_*(\sh{L})))_\eta\to\sh{L}_\eta$ es sobreyectivo, pero que el morfismo $X_\eta\to\mathbf{P}((f_*(\sh{L}))_\eta)$ no es una inmersión.
Entonces el mismo razonamiento anterior muestra que para $s$ suficientemente próximo a $\eta$, el homomorfismo $(f^*(f_*(\sh{L})))_s\to\sh{L}_s$ es sobreyectivo; por otra parte, en virtud de \sref{IV.9.6.1}, (viii), el morfismo $X_s\to\mathbf{P}((f_*(\sh{L}))_s)$ no es una inmersión para $s$ en un entorno de $\eta$, lo que termina la demostración en el caso de $E'$.

\emph{Caso de $E$.} --- Consideremos primero el caso particular en que $Y=S$:

\begin{corollary}[9.6.4]
\label{IV.9.6.4}
Sea $f:X\to S$ un morfismo propio de presentación finita y sea $\sh{L}$ un $\sh{O}_X$-Módulo invertible.
Entonces el conjunto $E$ de los $s\in S$ tales que $\sh{L}_s$ es amplio (relativamente a $f_s$) es abierto en $S$, y $\sh{L}|f^{-1}(E)$ es amplio relativamente a la restricción $f^{-1}(E)\to E$ de $f$.
\end{corollary}

\begin{proof}
Como la condición \sref{IV.9.2.1},~(i) es verificada por la propiedad $\mathbf{P}$ definida anteriormente, el resultado de \sref{IV.9.2.2},~(iv) y el razonamiento de \sref{IV.9.2.3} muestran que basta reducirse al caso en que $S$ es \emph{noetheriano}; pero entonces el resultado se sigue de \sref[III]{III.4.7.1} y de la estabilidad de la amplitud por cambio de base \sref[II]{II.4.6.13}.
\end{proof}

\begin{corollary}[9.6.5]
\label{IV.9.6.5}
Bajo las hipótesis de \sref{IV.9.6.4}, para que $\sh{L}$ sea amplio relativamente a $f$, es necesario y suficiente que, para todo $s\in S$, $\sh{L}_s$ sea amplio relativamente a $f_s$.
\end{corollary}

\phantomsection\label{IV.9.6.6}%
\textbf{(9.6.6) Fin de la demostración de \sref{IV.9.6.3}.} ---
Volvamos al caso general, siendo $S$ noetheriano, íntegro, de punto genérico $\eta\in S-E$.
Como $f$ es propio, se sigue de \sref{IV.9.6.4} que el conjunto $V$ de los $y\in Y$ tales que $\sh{L}_y$ es un $\sh{O}_{X_y}$-Módulo amplio, relativamente al morfismo $f_y:X_y\to\operatorname{Spec}(\boldsymbol{k}(y))$, es abierto, y por tanto $F=Y-V$ es cerrado en $Y$.
Como $f_s$ es propio y, para todo $y\in Y$ situado sobre $s\in S$, se tiene
$\sh{L}_y=(\sh{L}_s)_y$, se sigue de \sref{IV.9.6.5} que $s$ pertenece a
$S-E$ si y solo si $F_s\neq\varnothing$.
Pero $F$ es cerrado en $Y$ y $F_\eta\neq\varnothing$; por tanto, \sref{IV.9.5.1} muestra
que el conjunto de los $s\in S$ tales que $F_s\neq\varnothing$ es un entorno de
$\eta$ en $S$.
\end{proof}

\begin{remarks}[9.6.7]
\label{IV.9.6.7}
(i) Para cada una de las propiedades $\mathbf{P}$ consideradas en \sref{IV.9.6.1}, es aplicable la proposición \sref{IV.9.3.3}, y estas propiedades (para los morfismos $f_s$) son, por tanto, ``estables'' al pasar del límite proyectivo de un sistema proyectivo esencialmente afín \sref{IV.8.13.4} de preesquemas $X_\lambda$ a uno conveniente de ellos.

(ii) Sea $Z$ una parte localmente construible de $X$ tal que, para todo $s\in S$, $Z\cap X_s$ sea \emph{abierto} en $X_s$, y denotemos por $Z_s$ el subpreesquema de $X_s$ inducido sobre el abierto $Z\cap X_s$.
\oldpage[IV-3]{76}
Entonces, en las proposiciones \sref{IV.9.6.1} y \sref{IV.9.6.2},~(I), podemos sustituir en todas partes $f_s$ por su restricción $f_s|Z_s:Z_s\to Y_s$ sin modificar las conclusiones.
En efecto, la verificación de \sref{IV.9.2.1},~(i) procede como en \sref{IV.9.6.1} y \sref{IV.9.6.2}.
Por otra parte, en la reducción al caso noetheriano, efectuada en \sref{IV.9.2.3}, si $Z=q^{-1}(Z_0)$, donde $q:X\to X_0$ es la proyección canónica y $Z_0$ una parte construible de $X_0$ \sref{IV.8.3.11}, el hecho de que $(Z_0)_{s_0}$ sea abierto en $(X_0)_{s_0}$, para $s_0=p(s)$, se sigue de \sref{IV.2.4.10} y del hecho de que la proyección $X_s\to(X_0)_{s_0}$ es sobreyectiva.
Nos reducimos, por tanto, a verificar \sref{IV.9.2.1},~(ii) bajo las nuevas hipótesis.
Ahora bien, como $Z_\eta$ es abierto en $X_\eta$, existe un abierto $Z'\subset X$ tal que $Z_\eta=Z'\cap X_\eta$; como $X$ es entonces noetheriano, $Z'$ es construible, y $Z$ también es construible por hipótesis; concluimos de \sref{IV.9.5.2} y de ($\mathbf{0}_{\mathrm{III}}$,~9.2.2) que existe un entorno $U$ de $\eta$ en $S$ tal que $Z_s=Z'_s$ para $s\in U$.
Sustituyendo $S$ por $U$, podemos, por tanto, reducirnos al caso en que $Z$ es abierto en $X$, y entonces el resultado ha sido demostrado en \sref{IV.9.6.1} y \sref{IV.9.6.2}.
\end{remarks}

\subsection{Construibilidad de las propiedades de separabilidad, irreducibilidad geométrica y conexidad geométrica}
\label{subsection:IV.9.7}
\label{IV.9.7}

\begin{lemma}[9.7.1]
\label{IV.9.7.1}
Sea $S$ un preesquema irreducible de punto genérico $\eta$, y sea $f:X\to S$ un morfismo de presentación finita.
Si $X_\eta$ tiene $n$ componentes irreducibles (resp.\ $n'$ componentes conexas), existe un entorno abierto $U$ de $\eta$ en $S$ tal que, para todo $s\in U$, $X_s$ tiene al menos $n$ componentes irreducibles (resp.\ $n'$ componentes conexas).
\end{lemma}

\begin{proof}
Podemos reducirnos al caso en que $S$ es afín, de modo que $X$ es cuasicompacto y cuasiseparado.
Sean $V_i$ ($1\leqslant i\leqslant n$) los interiores de las componentes irreducibles de $X_\eta$, que son dos a dos disjuntos y cuasicompactos, y cuya unión es densa en $X_\eta$; en virtud de \sref{IV.8.2.11}, aplicado según el método de \sref{IV.8.1.2},~\emph{a}), existe para cada $i$ un abierto cuasicompacto $W_i$ en $X$ tal que $W_i\cap X_\eta=V_i$; como $X$ es cuasiseparado, las intersecciones $W_i\cap W_j$ son cuasicompactas \sref{IV.1.2.7}, y por tanto, sustituyendo $S$ por un entorno de $\eta$, podemos suponer que $W_i\cap W_j=\varnothing$ para $i\neq j$ en virtud de \sref{IV.8.3.3}.
Además, como los $W_i$ son construibles, existe un entorno $U$ de $\eta$ en $S$ tal que para todo $s\in U$ los $(W_i)_s$ son no vacíos \sref{IV.9.5.1} y \sref{IV.9.2.3} y la unión de los $(W_i)_s$ es densa en $X_s$ \sref{IV.9.5.3} y \sref{IV.9.2.3}.
En estas condiciones, como las componentes irreducibles de cada $(W_i)_s$ (que son en número finito) son también componentes irreducibles de la unión de los $(W_i)_s$ \sref[0]{0.2.1.7}, las adherencias en $X_s$ de estas componentes son las componentes irreducibles de $X_s$ \sref[0]{0.2.1.6} y su número es evidentemente $\geqslant n$.

Sean ahora $C_j$ ($1\leqslant j\leqslant n'$) las componentes conexas de $X_\eta$, que son abiertos cuasicompactos de $X_\eta$, dos a dos disjuntos.
Si sustituimos $V_i$ por $C_j$ en el razonamiento precedente, vemos (utilizando dos veces \sref{IV.8.3.3}) que podemos suponer $X$ una unión de $n'$ abiertos cuasicompactos $M_j$ ($1\leqslant j\leqslant n'$), dos a dos disjuntos, y tales que, en un entorno de $\eta$, los $(M_j)_s$ son no vacíos.
Como la unión de los $(M_j)_s$ es $X_s$, las componentes conexas de los $(M_j)_s$ son las componentes conexas de $X_s$, y por tanto su número es $\geqslant n'$.
\end{proof}

\oldpage[IV-3]{77}

\begin{lemma}[9.7.2]
\label{IV.9.7.2}
Sea $S$ un preesquema irreducible de punto genérico $\eta$, y sea $f:X\to S$ un morfismo de presentación finita.
Si $X_\eta$ no es reducido, existe un entorno abierto $U$ de $\eta$ en $S$ tal que, para todo $s\in U$, $X_s$ no es reducido.
\end{lemma}

\begin{proof}
En efecto, sea $\mathscr{N}$ el nilradical de $\sh{O}_X$; se sigue de \sref{IV.8.2.13}, aplicado según el método de \sref{IV.8.1.2},~\emph{a}), que $\mathscr{N}_\eta$ es el nilradical de $\sh{O}_{X_\eta}$, y la hipótesis es que $\mathscr{N}_\eta\neq 0$.
Concluimos de \sref{IV.9.4.5} y \sref{IV.9.2.3} que existe un entorno $U$ de $\eta$ tal que, para todo $s\in U$, $\mathscr{N}_s$ se identifica con un ideal de $\sh{O}_{X_s}$ y $\mathscr{N}_s\neq 0$; como $\mathscr{N}_s$ está evidentemente contenido en el nilradical de $\sh{O}_{X_s}$, vemos que $X_s$ no es reducido para $s\in U$.
\end{proof}

\textbf{(9.7.3)}
\label{IV.9.7.3}
Dado un polinomio $F\in A[\mathrm{T}_1,\ldots,\mathrm{T}_n]$, donde $A$ es un anillo y los $\mathrm{T}_i$ son indeterminados, para un homomorfismo de anillos $\rho:A\to B$, denotamos por $F_{(\rho)}$ o $F_{(B)}$ el polinomio de $B[\mathrm{T}_1,\ldots,\mathrm{T}_n]$ obtenido sustituyendo cada coeficiente de $F$ por su imagen mediante $\rho$.
Si $k$ es un cuerpo, $F\in k[\mathrm{T}_1,\ldots,\mathrm{T}_n]$ un polinomio no constante y $X=\operatorname{Spec}(k[\mathrm{T}_1,\ldots,\mathrm{T}_n]/(F))$, decir que $X$ es íntegro (o que el ideal $(F)$ es primo) significa que $F$ es \emph{irreducible} (es decir, que en toda factorización $F=F_1 F_2$ en polinomios de $k[\mathrm{T}_1,\ldots,\mathrm{T}_n]$, $F_1$ o $F_2$ es de grado $0$); esto se sigue del hecho de que el anillo de polinomios $k[\mathrm{T}_1,\ldots,\mathrm{T}_n]$ es factorial.
Deducimos inmediatamente de \sref{IV.4.6.2} y \sref{IV.4.5.2}:

\begin{lemma}[9.7.4]
\label{IV.9.7.4}
Sea $k$ un cuerpo, $\Omega$ una extensión algebraicamente cerrada de $k$ y $F$ un polinomio no constante de $k[\mathrm{T}_1,\ldots,\mathrm{T}_n]$.
Las condiciones siguientes son equivalentes:
\begin{enumerate}[label=\emph{\alph*)}]
  \item $X=\operatorname{Spec}(k[\mathrm{T}_1,\ldots,\mathrm{T}_n]/(F))$ es geométricamente íntegro.
  \item $F_{(K)}$ es irreducible para toda extensión $K$ de $k$.
  \item $F_{(\Omega)}$ es irreducible.
\end{enumerate}
Diremos en este caso que $F$ es \emph{geométricamente irreducible}.
\end{lemma}

\begin{lemma}[9.7.5]
\label{IV.9.7.5}
Sea $A$ un anillo íntegro, $K$ el cuerpo de fracciones de $A$ y $F$ un polinomio no constante de $A[\mathrm{T}_1,\ldots,\mathrm{T}_n]$ de grado $d$ tal que $F_{(K)}$ es geométricamente irreducible.
Entonces existe $f\neq 0$ en $A$ tal que para todo $x\in D(f)$, $F_{(\boldsymbol{k}(x))}$ es geométricamente irreducible.
\end{lemma}

\begin{proof}
Pongamos $F=\sum_\alpha c_\alpha \mathrm{T}^\alpha$ como de costumbre, con $c_\alpha=c_{\alpha_1\alpha_2\ldots\alpha_n}$, $\mathrm{T}^\alpha=\mathrm{T}_1^{\alpha_1}\mathrm{T}_2^{\alpha_2}\cdots\mathrm{T}_n^{\alpha_n}$, $|\alpha|=\alpha_1+\alpha_2+\cdots+\alpha_n\leqslant d$ (siendo no nulo al menos uno de los $c_\alpha$ tales que $|\alpha|=d$).
Como los $c_\alpha$ no nulos son invertibles en $K$, podemos suponer, sustituyendo $A$ por $A_g$ ($g\neq 0$ en $A$), que los $c_\alpha$ no nulos son invertibles \emph{en} $A$.
Se sigue entonces que, para todo $x\in\operatorname{Spec}(A)$, $F_{(\boldsymbol{k}(x))}$ es de grado $d$.

Demostraremos el lema mediante un lema preliminar.

\begin{lemma}[9.7.5.1]
\label{IV.9.7.5.1}
Sea $A$ un anillo íntegro, $S=\operatorname{Spec}(A)$, $\eta$ el punto genérico de $S$, $B$ el anillo de polinomios $A[\mathrm{T}_1,\ldots,\mathrm{T}_m]$, $(P_i)_{i\in I}$ una familia finita de elementos de $B$ y $\mathfrak{a}$ el ideal de $B$ generado por los $P_i$.
Para todo $s\in S$, sea $\mathfrak{a}_s$ el ideal de $\boldsymbol{k}(s)[\mathrm{T}_1,\ldots,\mathrm{T}_m]$ generado por los polinomios $(P_i)_{(\boldsymbol{k}(s))}$ ($i\in I$).
Entonces, si $V(\mathfrak{a}_\eta)=\varnothing$, existe un entorno $U$ de $\eta$ en $S$ tal que $V(\mathfrak{a}_s)=\varnothing$ para todo $s\in U$.
\end{lemma}

\begin{proof}
En efecto, sea $Y=\operatorname{Spec}(B)$, y sea $Z$ la parte cerrada $V(\mathfrak{a})$ de $Y$; como $\mathfrak{a}$ es un ideal de tipo finito en $B$, $Z$ es construible en $Y$ \sref[0]{0.9.1.5} y, con las notaciones introducidas en \sref{IV.9.4.1}, tenemos $V(\mathfrak{a}_s)=Z_s$ para todo $s\in S$; como el morfismo estructural $Y\to S$ es de presentación finita, la conclusión del lema se sigue de \sref{IV.9.5.1} y \sref{IV.9.2.3}.
\end{proof}

\oldpage[IV-3]{78}
Sea $(p,q)$ un par de enteros $>0$ tales que $p+q=d$; introducimos indeterminados $\mathrm{T}'_\beta$, $\mathrm{T}''_\gamma$ para todos los sistemas de enteros $\beta$, $\gamma$ tales que $|\beta|\leqslant p$ y $|\gamma|\leqslant q$; para todo sistema de enteros $\alpha$ tal que $|\alpha|\leqslant d$, consideremos el polinomio de $B=A[\mathrm{T}'_\beta, \mathrm{T}''_\gamma \mid |\beta|\leqslant p, |\gamma|\leqslant q]: P_\alpha(\mathrm{T}'_\beta,\mathrm{T}''_\gamma)=\sum_{\beta+\gamma=\alpha}\mathrm{T}'_\beta \mathrm{T}''_\gamma-c_\alpha$.
Sea $\Omega$ una clausura algebraica de $K$; decir que existen dos polinomios $F_1$, $F_2$ de $\Omega[\mathrm{T}_1,\ldots,\mathrm{T}_n]$, de grados respectivos $p$ y $q$, y tales que $F_1 F_2=F_{(\Omega)}$ significa que el sistema de ecuaciones $P_\alpha(t'_\beta,t''_\gamma)=0$ ($|\alpha|\leqslant d$) admite una solución $(t'_\beta,t''_\gamma)$ ($|\beta|\leqslant p$, $|\gamma|\leqslant q$) formada por elementos de $\Omega$.
Sea $\mathfrak{a}$ el ideal de $B$ generado por los $P_\alpha$; la interpretación precedente y el teorema de los ceros de Hilbert muestran que la hipótesis sobre $F_{(K)}$ implica que $V(\mathfrak{a}_\eta)=\varnothing$, designando por $\eta$ el punto genérico de $\operatorname{Spec}(A)$; el lema \sref{IV.9.7.5.1} demuestra entonces que en un entorno de $\eta$, tenemos $V(\mathfrak{a}_s)=\varnothing$, y por consiguiente, para estos valores de $x$, $F_{(\boldsymbol{k}(x))}$ no admite ninguna factorización $F_{(\boldsymbol{k}(x))}=G_1 G_2$, donde $G_1$, $G_2$ son polinomios de grados respectivos $p$ y $q$, cuyos coeficientes pertenecen a una clausura algebraica de $\boldsymbol{k}(x)$.
Basta aplicar este resultado a todos los pares de enteros $(p,q)$ tales que $p>0$, $q>0$ y $p+q=d$, para obtener la conclusión del lema \sref{IV.9.7.5}.
\end{proof}

\begin{proposition}[9.7.6]
\label{IV.9.7.6}
Sea $S$ un preesquema noetheriano íntegro de punto genérico $\eta$, sea $f:X\to S$ un morfismo de tipo finito y sea $\sh{F}$ un $\sh{O}_X$-Módulo coherente.
Si $\sh{F}_\eta$ no tiene ningún ciclo primo inmerso asociado, existe un entorno $U$ de $\eta$ en $S$ tal que, para todo $s\in U$, $\sh{F}_s$ no tiene ningún ciclo primo inmerso asociado.
\end{proposition}

\begin{proof}
Como $X_\eta$ es noetheriano, existe un homomorfismo inyectivo $v:\sh{F}_\eta\to\bigoplus_{i=1}^m\sh{G}_i$, donde cada $\sh{G}_i$ es irredundante y $\operatorname{Ass}(\sh{F}_\eta)$ es el conjunto de los $x_i$, con $\{x_i\}=\operatorname{Ass}(\sh{G}_i)$ \sref{IV.3.2.6}; si $Z_i$ es la adherencia de $\{x_i\}$ en $X_\eta$, entonces $Z_i=\operatorname{Supp}(\sh{G}_i)$ \sref{IV.3.1.4}, y por hipótesis los $Z_i$ son las componentes irreducibles de $Z=\operatorname{Supp}(\sh{F}_\eta)$.
Se sigue de \sref{IV.8.5.2},~(i) y (ii) (restringiendo $S$ si es necesario a un entorno de $\eta$) y de \sref{IV.8.5.8} que existen $\sh{O}_X$-Módulos coherentes $\sh{F}_i$ y un homomorfismo inyectivo $u:\sh{F}\to\bigoplus_{i=1}^m\sh{F}_i$ tales que $(\sh{F}_i)_\eta=\sh{G}_i$ para todo $i$ y $v=u_\eta$.
Si $Y_i=\operatorname{Supp}(\sh{F}_i)$, tenemos $(Y_i)_s=\operatorname{Supp}((\sh{F}_i)_s)$ para todo $s\in S$ \sref[I]{I.9.1.13} y, en particular, $(Y_i)_\eta=Z_i$ para todo $i$.
La hipótesis implica que para $i\neq j$, $Z_i-(Z_i\cap Z_j)$ es abierto y denso en $Z_i$; deducimos, por tanto, de \sref{IV.9.5.3} y \sref{IV.9.5.4} que en un entorno de $\eta$, $(Y_i)_s-((Y_i)_s\cap(Y_j)_s)$ es abierto y denso en $(Y_i)_s$.
Supongamos demostrado que cada uno de los $(\sh{F}_i)_s$ carece de ciclo primo inmerso asociado para $s\in U$.
Entonces los elementos de $\operatorname{Ass}((\sh{F}_i)_s)$ son los puntos maximales de $(Y_i)_s$; ninguno de ellos puede, por tanto, pertenecer a $(Y_j)_s$ para $i\neq j$, y la proposición quedará demostrada. Podemos, pues, suponer que $\sh{F}_\eta$ es \emph{irredundante}; además, podemos reducirnos al caso en que $S=\operatorname{Spec}(A)$ es afín. Entonces $X$ es una unión de un número finito de abiertos afines $V_j$, y si $V_j\cap\operatorname{Supp}(\sh{F}_\eta)=\varnothing$, también tenemos $V_j\cap\operatorname{Supp}(\sh{F}_s)=\varnothing$ para $s$ próximo a $\eta$ \sref{IV.9.5.1}. Podemos, por tanto, reducirnos al caso en que $X=\operatorname{Spec}(B)$ es también afín y $f:X\to S$ es dominante. Así, $A$ es un anillo noetheriano íntegro, subanillo de una $A$-álgebra de tipo finito $B$, y $\sh{F}=\widetilde{M}$ con $M$ un $B$-módulo finito. Por hipótesis, si $K$ es el cuerpo de fracciones de $A$, el $B_{(K)}$-módulo $M_{(K)}$ es irredundante.
Sea $\mathfrak{q}$ el único

\oldpage[IV-3]{79}

elemento de $\operatorname{Ass}(M_{(K)})$, y sea $\mathfrak{p}$ el ideal primo de $B$ que es la imagen inversa de $\mathfrak{q}$ por la aplicación canónica $B\to B_{(K)}$.
Sabemos por \sref{IV.5.11.1.1} que existe una filtración finita $(N_h)_{0\leqslant h\leqslant m}$ de $M_{(K)}$ tal que $N_0=M_{(K)}$, $N_m=0$, y $N_h/N_{h+1}$ es isomorfo a un $B_{(K)}$-submódulo no nulo de $B_{(K)}/\mathfrak{q}$.
Sea $M_h$ la imagen inversa de $N_h$ por la aplicación canónica $M\to M_{(K)}$.
Se sigue de \sref{IV.9.4.4} que para $s$ suficientemente próximo a $\eta$, $(M_h)_{(\boldsymbol{k}(s))}$ se identifica con un sub-$B_{(\boldsymbol{k}(s))}$-módulo de $M_{(\boldsymbol{k}(s))}$, y el cociente $(M_h)_{(\boldsymbol{k}(s))}/(M_{h+1})_{(\boldsymbol{k}(s))}$ con un $B_{(\boldsymbol{k}(s))}$-submódulo no nulo de $B_{(\boldsymbol{k}(s))}/\mathfrak{p}_{(\boldsymbol{k}(s))}=(B/\mathfrak{p})_{(\boldsymbol{k}(s))}$.
Teniendo en cuenta \sref{IV.3.1.7}, con las mismas notaciones, queda demostrar que si $B$ es \emph{íntegro}, entonces $B_{(\boldsymbol{k}(s))}$ no tiene ideales primos asociados inmersos para $s$ próximo a $\eta$.
Ahora bien, sustituyendo $A$ por $A_g$ y $B$ por $B_g$ (donde $g$ es un elemento no nulo de $A$), podemos suponer que $B$ contiene una subálgebra de polinomios $C=A[\mathrm{T}_1,\ldots,\mathrm{T}_n]$ tal que $B$ es un $C$-módulo de tipo finito (Bourbaki, \emph{Alg.\ comm.}, cap.\ V, \S~3, n$^\circ$~1, cor.\ 1 del teorema\ 1).
Como $B_{(K)}$ es un $C_{(K)}$-módulo sin torsión, podemos aplicar una vez más el argumento precedente, sustituyendo $B$, $M$ y $\mathfrak{q}$ por $C$, $B$ y $(0)$ respectivamente; basta demostrar que para $s$ próximo a $\eta$, $C_{(\boldsymbol{k}(s))}$ no tiene ideales primos asociados inmersos.
Pero esto es evidente, puesto que $C_{(\boldsymbol{k}(s))}=\boldsymbol{k}(s)[\mathrm{T}_1,\ldots,\mathrm{T}_n]$ es un anillo íntegro.
\end{proof}

\begin{theorem}[9.7.7]
\label{IV.9.7.7}
Sea $f:X\to S$ un morfismo de presentación finita, y sea $E$ el conjunto de los $s\in S$ para los cuales $X_s$ posee una de las propiedades siguientes: ser
\begin{enumerate}[label=(\roman*)]
  \item geométricamente irreducible;
  \item geométricamente conexo;
  \item geométricamente reducido;
  \item geométricamente íntegro.
\end{enumerate}
Entonces $E$ es localmente construible en $S$.
\end{theorem}

\begin{proof}
Para ver que las propiedades consideradas son construibles, observemos primero que satisfacen trivialmente la condición \sref{IV.9.2.1},~(i), en virtud de \sref{IV.4.5.6},~(i) y \sref{IV.4.6.5},~(i).
Queda, por tanto, verificar \sref{IV.9.2.1},~(ii), de modo que podemos suponer $S=\operatorname{Spec}(A)$ afín, noetheriano, íntegro, de punto genérico $\eta$.
En virtud de \sref{IV.4.6.8}, existe una extensión finita $K'$ de $K=\boldsymbol{k}(\eta)$ tal que las componentes irreducibles (resp.\ componentes conexas) de $(X_\eta)_{(K')}$ son geométricamente irreducibles (resp.\ geométricamente conexas) y $((X_\eta)_{(K')})_\mathrm{red}$ es geométricamente reducido.
Además, existe una base de $K'$ sobre $K$ formada por elementos enteros sobre $A$; el anillo íntegro $A'$ generado por estos elementos es finito sobre $A$ y tiene $K'$ como cuerpo de fracciones.
Pongamos $S'=\operatorname{Spec}(A')$; el morfismo $g:S'\to S$ es finito y dominante, por tanto sobreyectivo \sref[II]{II.6.1.10}. Pongamos $X'=X_{(S')}$, de modo que, si $\eta'$ es el punto genérico de $S'$, entonces $X'_{\eta'}=(X_\eta)_{(K')}$. El conjunto $E'$ de los $s'\in S'$ tales que $X'_{s'}$ posee una de las propiedades (i), (ii), (iii) o (iv) es igual a $g^{-1}(E)$ \sref{IV.9.2.2},~(iv), donde $E$ corresponde, naturalmente, a la misma propiedad. Como $g$ es sobreyectivo, $E=g(E')$ y $S-E=g(S'-E')$. Además, $g$ es cerrado y $g^{-1}(\eta)=\{\eta'\}$ porque $A'$ es íntegro y finito sobre $A$ (Bourbaki, \emph{Alg.\ comm.}, cap.\ V, \S~2, n$^\circ$~1, cor.\ 1 de la prop.\ 1); por tanto, la imagen por $g$ de todo entorno de $\eta'$ es un entorno de $\eta$.
El teorema quedará, pues, demostrado si se prueba que $E'$ o $S'-E'$ es un entorno de $\eta'$.
En otras palabras, podemos suponer en lo sucesivo que las componentes irreducibles (resp.\ componentes conexas) de $X_\eta$ son geométricamente

\oldpage[IV-3]{80}

irreducibles (resp.\ geométricamente conexas) y que $(X_\eta)_\mathrm{red}$ es geométricamente reducido.

Supongamos primero que $\eta\in S-E$ para una de las propiedades (i) a (iv).
Si $X_\eta$ no es geométricamente irreducible (resp.\ geométricamente conexo), no es irreducible (resp.\ conexo) en virtud de la hipótesis precedente, y por tanto lo mismo vale para $X_s$ para $s$ en un entorno de $\eta$ \sref{IV.9.7.1}, y \emph{a fortiori} en este entorno $X_s$ no es geométricamente irreducible (resp.\ geométricamente conexo).
Por otra parte, si $X_\eta$ no es geométricamente reducido, no es reducido (pues de otro modo sería igual a $(X_\eta)_\mathrm{red}$, que es geométricamente reducido por hipótesis); por tanto, en un entorno de $\eta$, $X_s$ no es reducido \sref{IV.9.7.2} y, en consecuencia, no es geométricamente reducido.
Finalmente, si $X_\eta$ no es geométricamente íntegro, o bien no es reducido, en cuyo caso acabamos de ver que $X_s$ no es reducido (ni \emph{a fortiori} íntegro) en un entorno de $\eta$; o bien $X_\eta$ es reducido (por tanto geométricamente reducido por hipótesis), en cuyo caso no es geométricamente irreducible, y acabamos de ver más arriba que lo mismo vale entonces para $X_s$ para $s$ próximo a $\eta$; \emph{a fortiori} $X_s$ no es geométricamente íntegro para estos valores de $s$.

Supongamos ahora que $\eta\in E$ y examinemos por separado cada una de las propiedades consideradas.

$1^\circ$ Supongamos que $X_\eta$ es geométricamente íntegro.
Sea $L$ el cuerpo de funciones racionales sobre $X_\eta$; la hipótesis sobre $X_\eta$ implica que $L$ es una extensión separable de $K$ \sref{IV.4.6.3}, por tanto una extensión separable finita de una extensión pura $K(\mathrm{T}_1,\ldots,\mathrm{T}_n)$ (siendo los $\mathrm{T}_i$ indeterminados); existe entonces un elemento $x\in L$, entero sobre el anillo $K[\mathrm{T}_1,\ldots,\mathrm{T}_n]$, y tal que $L=K(\mathrm{T}_1,\ldots,\mathrm{T}_n)(x)$; sea $G\in K[\mathrm{T}_1,\ldots,\mathrm{T}_n,\mathrm{T}_{n+1}]$ su polinomio minimal.
Existe $g\neq 0$ de $A$ tal que todos los coeficientes $\neq 0$ de $G$ (que pertenecen todos a $K$) pertenecen a $A_g$; sustituyendo $A$ por $A_g$ (lo que equivale a sustituir $S$ por un entorno de $\eta$), podemos, por tanto, suponer que $G$ tiene sus coeficientes en $A$; designando por $F$ el polinomio $G$ considerado como elemento de $A[\mathrm{T}_1,\ldots,\mathrm{T}_n,\mathrm{T}_{n+1}]$, tenemos entonces $G=F_{(\boldsymbol{k}(\eta))}$.
Pongamos $Y=\operatorname{Spec}(A[\mathrm{T}_1,\ldots,\mathrm{T}_{n+1}]/(F))$, de modo que $Y_\eta=\operatorname{Spec}(K[\mathrm{T}_1,\ldots,\mathrm{T}_{n+1}]/(G))$; $Y_\eta$ es un esquema íntegro cuyo cuerpo de funciones racionales es $L$.
Como $X_\eta$ e $Y_\eta$ son noetherianos, existe un abierto no vacío $V\subset Y_\eta$ y una inmersión abierta (necesariamente dominante) $v:V\to X_\eta$ \sref[I]{I.6.5.1},~(ii) y \sref[I]{I.6.5.4},~(ii).
Sea $W$ un abierto de $Y$ tal que $W\cap Y_\eta=V$; aplicando \sref{IV.8.8.2},~(i) y \sref{IV.8.10.5},~(iii) según el método de \sref{IV.8.1.2},~\emph{a}), sustituyendo $S$ por un entorno de $\eta$, podemos suponer que $v=w_\eta$, donde $w:W\to X$ es una inmersión abierta.

Vimos en \sref{IV.4.6.3} que el criterio para que un preesquema algebraico íntegro sea geométricamente íntegro depende únicamente de su cuerpo de funciones racionales. Como $X_\eta$ es geométricamente íntegro por hipótesis, lo mismo vale para $Y_\eta$, y la definición de $G$ implica, por tanto, que este polinomio es geométricamente irreducible \sref{IV.9.7.4}.
Aplicando \sref{IV.9.7.5}, vemos que existe un entorno $U$ de $\eta$ en $S$ tal que $F_{(\boldsymbol{k}(s))}$ es geométricamente irreducible para todo $s\in U$, y por tanto $Y_s$ es geométricamente íntegro para $s\in U$ \sref{IV.9.7.4}; además, podemos también suponer que para $s\in U$, $W_s$ no es vacío \sref{IV.9.5.1}, y es por tanto geométricamente íntegro \sref{IV.4.6.3}; finalmente, podemos también suponer que

\oldpage[IV-3]{81}

para $s\in U$, $w_s:W_s\to X_s$ es una inmersión abierta \sref{IV.9.6.1},~(x), y que su imagen en $X_s$ es densa en todas partes \sref{IV.9.5.3}.
En otras palabras, para $s\in U$, existe en $X_s$ un abierto denso en todas partes $W'_s$ que es geométricamente íntegro; el criterio \sref{IV.4.5.9},~\emph{c}) muestra, por tanto, que $X_s$ es ya geométricamente irreducible para $s\in U$.
Finalmente, como $X_\eta$ es reducido y por tanto no tiene ningún ciclo primo inmerso asociado \sref{IV.3.2.1}, podemos también suponer que para $s\in U$, $X_s$ no tiene ningún ciclo primo inmerso asociado \sref{IV.9.7.6}; en estas condiciones (para $s\in U$ fijado), sea $\Omega$ una extensión algebraicamente cerrada de $\boldsymbol{k}(s)$, y sea $p:(X_s)_{(\Omega)}\to X_s$ la proyección canónica; $p^{-1}(W'_s)$ es un abierto denso en $(X_s)_{(\Omega)}$ \sref{IV.2.3.10} y es íntegro por hipótesis; además, $(X_s)_{(\Omega)}$ no tiene ningún ciclo primo inmerso asociado \sref{IV.4.2.7}, por tanto concluimos de \sref{IV.3.2.1} que $(X_s)_{(\Omega)}$ es reducido; esto termina la demostración de que $X_s$ es geométricamente íntegro \sref{IV.4.6.1}.

$2^\circ$ Supongamos que $X_\eta$ es geométricamente irreducible; como $X_{\mathrm{red}}$ es también de tipo finito sobre $S$ \sref{IV.1.5.4},~(vi), teniendo en cuenta \sref[I]{I.5.1.8}, podemos sustituir $X$ por $X_{\mathrm{red}}$; entonces $X_\eta$ es también íntegro, y como por hipótesis $X_\eta$ es geométricamente reducido, es geométricamente íntegro. Estamos entonces en las condiciones de $1^\circ$, y concluimos (volviendo a las hipótesis iniciales) que $X_s$ es geométricamente irreducible para $s$ próximo a $\eta$.

$3^\circ$ Supongamos $X_\eta$ geométricamente conexo, y sean $Z_i$ ($1\leqslant i\leqslant n$) las componentes irreducibles de $X_\eta$; existe (en virtud de \sref{IV.5.10.8.1} aplicado a $\mathfrak{T}=\{\varnothing\}$) una aplicación \emph{sobreyectiva} $j\mapsto\nu(j)$ de un intervalo $[1,m]$ de $\mathbf{N}$ sobre $[1,n]$ tal que $Z_{\nu(j)}\cap Z_{\nu(j+1)}\neq\varnothing$ para $1\leqslant j\leqslant m$.
Para todo $i$, sea $X_i$ la adherencia de $Z_i$ en $X$, y sea $Y$ la unión de los $X_i$; como $Y_\eta=X_\eta$ por definición, podemos, en virtud de \sref{IV.9.5.1}, suponer que $Y_s=X_s$ para todo $s\in S$, de modo que $X_s$ es la unión de los $(X_i)_s$.
Ahora bien, considerando los subpreesquemas reducidos de $X$ cuyos espacios subyacentes son los $X_i$, vemos por $2^\circ$ que existe un entorno $U$ de $\eta$ en $S$ tal que para $s\in U$ los $(X_i)_s$ son \emph{geométricamente irreducibles} (puesto que podemos suponer los $(X_i)_\eta$ geométricamente irreducibles, como vimos al principio).
Además, podemos también suponer que para $s\in U$, $(X_{\nu(j)})_s\cap(X_{\nu(j+1)})_s\neq\varnothing$ \sref{IV.9.5.1} para $1\leqslant j\leqslant m$; concluimos inmediatamente que $X_s$ es conexo, por tanto, por \sref{IV.4.5.13.1}, geométricamente conexo para $s\in U$.

$4^\circ$ Supongamos $X_\eta$ geométricamente reducido; sean $Z_i$ las componentes irreducibles de $X_\eta$, y sea $W_i$ el interior de $Z_i$ en $X_\eta$; existe para cada $i$ un abierto $V_i$ de $X$ tal que $W_i=V_i\cap X_\eta$ para todo $i$; como los $W_i$ son abiertos, dos a dos disjuntos, y su unión es densa en $X_\eta$, \sref{IV.9.5.1},~\sref{IV.9.5.3} y \sref{IV.9.5.4} nos permiten suponer que para $s$ próximo a $\eta$, los $(V_i)_s$ son abiertos, dos a dos disjuntos, en $X_s$ y su unión es densa en $X_s$.
Además, como los $W_i$ son geométricamente reducidos y se supusieron al principio geométricamente irreducibles, se sigue de $1^\circ$ que para $s$ próximo a $\eta$, los $(V_i)_s$ son geométricamente íntegros y, \emph{a fortiori}, reducidos.
Por otra parte, deducimos de \sref{IV.9.7.6} que para $s$ próximo a $\eta$, $X_s$ no tiene ningún ciclo primo inmerso asociado, puesto que lo mismo vale para el preesquema reducido $X_\eta$ \sref{IV.3.2.1}; concluimos, por tanto, de \sref{IV.3.2.1} que $X_s$ es reducido, y de \sref{IV.4.6.1} que $X_s$ es geométricamente reducido.
\end{proof}

\oldpage[IV-3]{82}

Las partes del enunciado de \sref{IV.9.7.7} relativas a las propiedades (i) y (ii) se generalizan como sigue.

\begin{proposition}[9.7.8]
\label{IV.9.7.8}
Sea $S$ un preesquema irreducible de punto genérico $\eta$, y sea $f:X\to S$ un morfismo de presentación finita. Sea $n$ (resp.\ $n'$) el número geométrico de componentes irreducibles (resp.\ componentes conexas) de $X_\eta$ \sref{IV.4.5.2}. Entonces existe un entorno $U$ de $\eta$ en $S$ tal que, para todo $s\in U$, el número geométrico de componentes irreducibles (resp.\ componentes conexas) de $X_s$ es igual a $n$ (resp.\ $n'$).
\end{proposition}

\begin{proof}
Teniendo en cuenta que el número geométrico de componentes irreducibles (resp.\ componentes conexas) de un preesquema algebraico es invariante por extensión del cuerpo base \sref{IV.4.5.6}, vemos, por el método de \sref{IV.9.2.3}, que podemos reducirnos al caso en que $S$ es afín, noetheriano e íntegro. Además, razonando como al principio de la demostración de \sref{IV.9.7.7}, vemos que podemos suponer que las componentes irreducibles (resp.\ componentes conexas) de $X_\eta$ son geométricamente irreducibles (resp.\ geométricamente conexas). Ya sabemos por \sref{IV.9.7.1} que, para todo $s$ próximo a $\eta$, $X_s$ tiene al menos $n$ componentes irreducibles y $n'$ componentes conexas.

Por otra parte, si $Z_i$ (resp.\ $Z'_j$) son las componentes irreducibles (resp.\ componentes conexas) de $X_\eta$, sea $X_i$ (resp.\ $X'_j$) el subpreesquema cerrado reducido de $X$ cuyo espacio subyacente es la adherencia de $Z_i$ (resp.\ $Z'_j$) en $X$. Se sigue de \sref{IV.9.5.1} que, en un entorno de $\eta$, los $(X_i)_s$ (resp.\ $(X'_j)_s$) recubren $X_s$, y que los $(X'_j)_s$ son dos a dos disjuntos. Como los $(X_i)_s$ (resp.\ $(X'_j)_s$) son geométricamente irreducibles (resp.\ geométricamente conexos) por \sref{IV.9.7.7} para $s$ próximo a $\eta$, se sigue que $X_s$ tiene a lo sumo $n$ componentes irreducibles y a lo sumo $n'$ componentes conexas, lo que demuestra la proposición.
\end{proof}

\begin{corollary}[9.7.9]
\label{IV.9.7.9}
Sea $f:X\to S$ un morfismo de presentación finita. Para todo $s\in S$, sea $n(s)$ (resp.\ $n'(s)$) el número geométrico de componentes irreducibles (resp.\ componentes conexas) de $f^{-1}(s)$ \sref{IV.4.5.2}. Entonces la función $s\mapsto n(s)$ (resp.\ $s\mapsto n'(s)$) es localmente construible sobre $S$.
\end{corollary}

\begin{proof}
Basta demostrar que la propiedad ``$X$ es un preesquema algebraico sobre un cuerpo $k$ y el número geométrico de componentes irreducibles (resp.\ componentes conexas) de $X$ es igual a $n$ (resp.\ $n'$)'' es construible. Se sigue de \sref{IV.4.5.6} que esta propiedad satisface la condición \sref{IV.9.2.1},~(i); nos reducimos, por tanto, al caso en que $S$ es afín, noetheriano e íntegro de punto genérico $\eta$, y la conclusión se sigue entonces de \sref{IV.9.7.8}.
\end{proof}

\begin{corollary}[9.7.10]
\label{IV.9.7.10}
Sea $X$ un preesquema localmente noetheriano tal que, si $X'$ es la normalización de $X_{\mathrm{red}}$, el morfismo canónico $f:X'\to X$ sea finito. Entonces el conjunto de los puntos $x\in X$ tales que $X$ es geométricamente unirrama en $x$ es localmente construible en $X$.
\end{corollary}

\begin{proof}
En efecto, este conjunto es, por definición, el conjunto de los puntos $x\in X$ tales que el número de puntos geométricos de $f^{-1}(x)$ es igual a $1$. Como $f$ es finito, este número es también el número geométrico de componentes irreducibles del espacio discreto $f^{-1}(x)$ (teniendo en cuenta la definición de la normalización \sref[II]{II.6.3.8} y \sref{IV.4.5.11}); la conclusión se sigue, por tanto, de \sref{IV.9.7.9}.
\end{proof}

\begin{remark}[9.7.11]
\label{IV.9.7.11}
Sea $Z$ un subconjunto localmente construible de $X$ tal que, para todo $s\in S$, $Z\cap X_s$ sea abierto en $X_s$, y denotemos por $Z_s$ el subpreesquema de $X_s$ inducido

\oldpage[IV-3]{83}

sobre el subconjunto abierto $Z\cap X_s$. Entonces, en \sref{IV.9.7.7}, se puede sustituir en todas partes $X_s$ por $Z_s$ sin modificar la conclusión; esto se sigue repitiendo el argumento utilizado en \sref{IV.9.6.7},~(ii).
\end{remark}

\begin{proposition}[9.7.12]
\label{IV.9.7.12}
Sea $f:X\to S$ un morfismo de presentación finita, y sea $g:S\to X$ una $S$-sección de $X$ \sref[I]{I.2.5.5}. Para todo $s\in S$, sea $X_s^0$ la componente conexa de $f^{-1}(s)$ que contiene $g(s)$. Entonces la unión $X^0$ de los $X_s^0$, para $s\in S$, es un subconjunto localmente construible de $X$.
\end{proposition}

\begin{proof}
Demostremos primero que podemos reducirnos al caso en que $S$ es afín y noetheriano. Siempre podemos suponer que $S=\operatorname{Spec}(A)$ es afín; con las notaciones de \sref{IV.9.2.3}, se tiene $f=(f_0)_{(S)}$, donde $f_0:X_0\to S_0$ es un morfismo de tipo finito, y podemos además suponer que existe una $S_0$-sección $g_0:S_0\to X_0$ tal que $g=(g_0)_{(S)}$ \sref{IV.8.9.1}. Si $p:S\to S_0$ es el morfismo canónico, entonces, para todo $s_0\in S_0$, la componente conexa $(X_0)^0_{s_0}$ de $f_0^{-1}(s_0)$ que contiene $g_0(s_0)$ es geométricamente conexa \sref{IV.4.5.13}; por consiguiente, si $s_0=p(s)$, se tiene $X_s^0=q^{-1}((X_0)^0_{s_0})$, donde $q:X_s\to (X_0)_{s_0}$ es la proyección canónica \sref{IV.4.5.8} y \sref{IV.4.4.1}. Nuestra afirmación se sigue, por tanto, de \sref{IV.1.8.2}.

Utilizamos ahora el criterio de construibilidad \sref[0\textsubscript{III}]{0.9.2.3}. Sea $x$ un punto de $X$, sea $Z$ el subpreesquema reducido de $X$ cuyo espacio subyacente es $\overline{\{x\}}$, y sea $Y$ el subpreesquema reducido de $S$ cuyo espacio subyacente es $\overline{f(Z)}$. Como la restricción de $f$ a $Z$ se factoriza como $Z\to Y\to S$ \sref[I]{I.5.2.2}, podemos sustituir $S$ por $Y$, es decir, suponer que $f(x)=\eta$, el punto genérico del preesquema íntegro $S$.

Por hipótesis, $X_\eta$ es la suma de dos subpreesquemas $X_\eta^0$ y $X_\eta^1$ inducidos sobre subconjuntos abiertos complementarios de $X_\eta$. Por \sref{IV.9.5.4} y \sref{IV.9.5.1}, tras sustituir $S$, si es necesario, por un entorno de $\eta$, podemos suponer que $X$ es la unión de dos subconjuntos abiertos disjuntos $X^0$ y $X^1$ tales que $(X^0)_\eta=X_\eta^0$ y $(X^1)_\eta=X_\eta^1$. Como $g:S\to X$ es continua e inyectiva, $S$ es la unión de los dos subconjuntos abiertos disjuntos $g^{-1}(X^0)$ y $g^{-1}(X^1)$; pero $S$ es irreducible y, \emph{a fortiori}, conexo, por lo que uno de estos dos subconjuntos abiertos es vacío. Como $g(\eta)\in X_\eta^0$ por definición, se tiene $g^{-1}(X^1)=\varnothing$.

En otras palabras, $g$ es una $S$-sección de $X^0$. Por otra parte, como $(X^0)_\eta$ es geométricamente conexo \sref{IV.4.5.13}, lo mismo vale para $(X^0)_s$ para todo $s$ próximo a $\eta$ por \sref{IV.9.7.7}; y como $g(s)\in (X^0)_s$, se tiene efectivamente $(X^0)_s=X_s^0$.
\end{proof}

\subsection{Descomposición primaria en un entorno de una fibra genérica}
\label{subsection:IV.9.8}
\label{IV.9.8}

\begin{proposition}[9.8.1]
\label{IV.9.8.1}
Sea $f:X\to S$ un morfismo de presentación finita y sea $\sh{F}$ un $\sh{O}_X$-Módulo cuasicoherente de presentación finita.
Entonces el conjunto $E$ de los $s\in S$ tales que $\sh{F}_s$ es un $\sh{O}_{X_s}$-Módulo sin ciclo primo inmerso asociado es localmente construible.
\end{proposition}

\begin{proof}
\label{IV.9.8.1-proof}
Sabemos que, para los Módulos cuasicoherentes sobre preesquemas algebraicos, la propiedad de carecer de ciclo primo inmerso asociado es invariante por cambio del cuerpo base \sref{IV.4.2.7}; podemos, por tanto, restringirnos al caso en que $S$ es afín, noetheriano e íntegro, y a demostrar que en este caso $E$ o $S-E$ es un entorno de $\eta$.
Vimos en \sref{IV.9.7.6} que si $\eta\in E$, $E$ es un entorno de $\eta$; queda, por tanto, considerar el caso en que $\sh{F}_\eta$ tiene ciclos primos inmersos asociados.
Podemos restringirnos al caso en que $X$ es afín, y en que existe un sub-$\sh{O}_{X_\eta}$-Módulo coherente $\sh{H}$ de $\sh{F}_\eta$ cuyo soporte $Z'$ es no vacío y \emph{raro} respecto del soporte $Z$ de $\sh{F}$ \sref{IV.3.1.3}; entonces, en virtud de \sref{IV.8.5.2},~(i) y~(ii),
\oldpage[IV-3]{84}
y \sref{IV.8.5.8}, aplicado según el método de \sref{IV.8.1.2},~$a$), existe un sub-$\sh{O}_X$-Módulo coherente $\sh{G}$ de $\sh{F}$ tal que $\sh{G}_\eta=\sh{H}$; si $Y=\operatorname{Supp}(\sh{F})$ e $Y'=\operatorname{Supp}(\sh{G})$, entonces, por consiguiente, $Y_s=\operatorname{Supp}(\sh{F}_s)$ e $Y'_s=\operatorname{Supp}(\sh{G}_s)$ para todo $s\in S$ \sref[I]{I.9.1.13}, y en particular $Z=Y_\eta$ y $Z'=Y'_\eta$; como $Z-Z'$ es denso en $Z$ y $Z'\neq\varnothing$, existe un entorno $U$ de $\eta$ tal que para $s\in U$, $Y_s-Y'_s$ es denso en $Y_s$ e $Y'_s\neq\varnothing$ \sref{IV.9.5.1} y \sref{IV.9.5.3}; considerando un punto genérico de una componente irreducible de $Y'_s$ y un entorno suficientemente pequeño de este punto en $X_s$, deducimos inmediatamente de \sref{IV.3.1.3} que $\sh{F}_s$ admite ciclos primos inmersos asociados.
\end{proof}

\medskip
\noindent
\emph{\sref{IV.9.8.2}}
\label{IV.9.8.2}
Sea $S$ un preesquema noetheriano, \emph{íntegro}, de punto genérico $\eta$, sea $f:X\to S$ un morfismo de tipo finito y sea $\sh{F}$ un $\sh{O}_X$-Módulo coherente.
Consideremos una descomposición reducida irredundante $(\sh{G}_i)_{i\in I}$ de $\sh{F}_\eta$ \sref{IV.3.2.6}; los $\sh{G}_i$ son, pues, cocientes de $\sh{F}_\eta$ y existe un homomorfismo inyectivo $h:\sh{F}_\eta\to\bigoplus_i\sh{G}_i$; además, $\operatorname{Ass}(\sh{G}_i)$ se reduce a un punto $x_i$.
Sea $Z_i$ la adherencia de $\{x_i\}$ en $X$, de modo que $(Z_i)_\eta=\operatorname{Supp}(\sh{G}_i)$.
Denotemos por $\sh{J}_i$ el Ideal coherente de $\sh{O}_X$ que define el subpreesquema cerrado de $X$ cuyo espacio subyacente es $Z_i$ (subpreesquema que seguimos denotando por $Z_i$).
En virtud de \sref{IV.8.5.2} y \sref{IV.8.5.8}, aplicados según el método de \sref{IV.8.1.2},~$a$), podemos (restringiendo $S$ a un entorno de $\eta$ si es necesario) suponer que existen cocientes $\sh{F}_i$ de $\sh{F}$ ($i\in I$) tales que $\sh{G}_i=(\sh{F}_i)_\eta$ para todo $i$, y un homomorfismo $g:\sh{F}\to\bigoplus_i\sh{F}_i$ tal que $h=g_\eta$.
Además, por \sref[I]{I.9.3.5}, existe un entero $m$ tal que $(\sh{J}_i^m\sh{F}_i)_\eta=(\sh{J}_i)_\eta^m\sh{G}_i=0$ y, restringiendo de nuevo $S$, podemos suponer que $\sh{J}_i^m\sh{F}_i=0$ \sref{IV.8.5.2.5}, de modo que el soporte de $\sh{F}_i$ está contenido en $Z_i$; pero como es cerrado y contiene $x_i$, es necesariamente igual a $Z_i$.

\begin{proposition}[9.8.3]
\label{IV.9.8.3}
Bajo las condiciones de \emph{\sref{IV.9.8.2}}, para todo $s\in S$ y todo $i\in I$, denotemos por $x_{is\alpha}$ ($\alpha\in J_{s,i}$) los puntos maximales de $(Z_i)_s$.
Existe un entorno $U$ de $\eta$ en $S$ tal que, para todo $s\in U$, los $x_{is\alpha}$ (para $i\in I$ y, para cada $i$, $\alpha\in J_{s,i}$) son dos a dos distintos y $\operatorname{Ass}(\sh{F}_s)$ es el conjunto de los $x_{is\alpha}$ (en otras palabras, los ciclos primos asociados a $\sh{F}_s$ son las componentes irreducibles de los $(Z_i)_s$).
Además, se puede tomar $U$ de modo que, para que la adherencia de $\{x_{is\alpha}\}$ en $X_s$ sea un ciclo primo asociado maximal de $\sh{F}_s$, sea necesario y suficiente que $(Z_i)_\eta$ (adherencia de $\{x_i\}$ en $X_\eta$) sea un ciclo primo asociado maximal de $\sh{F}_\eta$.
\end{proposition}

\begin{proof}
\label{IV.9.8.3-proof}
Se sigue de \sref{IV.3.1.3},~$c'$) que para cada $i$, existe un abierto $W_i$ en $X_\eta$ tal que $W_i\cap(Z_i)_\eta$ no es vacío, y un $\sh{O}_{X_\eta}$-Módulo coherente $\sh{H}_i$, de soporte $W_i\cap(Z_i)_\eta$, tal que existe un homomorfismo inyectivo $v_i:\sh{H}_i\to\sh{F}_\eta|W_i$.
Sea $V_i$ un abierto de $X$ tal que $V_i\cap X_\eta=W_i$; aplicando, como en \sref{IV.9.8.2}, los resultados de \sref{IV.8.5.2} y \sref{IV.8.5.8}, podemos (restringiendo $S$) suponer que existe un Módulo coherente $\sh{F}'_i$, de soporte $V_i\cap Z_i$, y un homomorfismo $u_i:\sh{F}'_i\to\sh{F}|V_i$ tales que $\sh{H}_i=(\sh{F}'_i)_\eta$ y $v_i=(u_i)_\eta$.
Mostraremos que existe un entorno $U$ de $\eta$ tal que para $s\in U$, se satisfacen las propiedades siguientes:
\begin{enumerate}[label=(\roman*)]
  \item Los $(\sh{F}_i)_s$ carecen de ciclo primo inmerso asociado.
  \item El homomorfismo $g_s:\sh{F}_s\to\bigoplus_i(\sh{F}_i)_s$ es inyectivo.
  \item $(Z_i)_s\cap(V_i)_s$ es denso en $(Z_i)_s$, $(\sh{F}'_i)_s$ tiene soporte $(Z_i)_s\cap(V_i)_s$, y $(u_i)_s:(\sh{F}'_i)_s\to\sh{F}_s|(V_i)_s$ es inyectivo.
\end{enumerate}
\oldpage[IV-3]{85}

\begin{enumerate}[label=(\roman*)]
  \item[(iv)] Para $i\neq j$, toda componente irreducible de $(Z_i)_s$ es distinta de toda componente irreducible de $(Z_j)_s$.
\end{enumerate}

Ahora bien, (i) ya se ha visto en \sref{IV.9.7.6}; (ii) es un caso particular de \sref{IV.9.4.5}; (iii) se sigue de \sref{IV.9.5.3} y \sref{IV.9.4.5}.
Finalmente, si $i\neq j$, $(Z_i)_\eta\cap(Z_j)_\eta=(Z_i\cap Z_j)_\eta$ es raro en $(Z_i)_\eta$ \emph{o} en $(Z_j)_\eta$; supongamos, por ejemplo, que $(Z_i\cap Z_j)_\eta$ es raro en $(Z_j)_\eta$; se sigue entonces de \sref{IV.9.5.3} y \sref{IV.9.5.4} que para $s$ en un entorno de $\eta$, $(Z_i\cap Z_j)_s=(Z_i)_s\cap(Z_j)_s$ es raro en $(Z_j)_s$, lo que muestra que ninguna componente irreducible de $(Z_j)_s$ puede estar contenida en una componente irreducible de $(Z_i)_s$, ni \emph{a fortiori} ser igual a ella.

En estas condiciones, se sigue de (ii) y \sref{IV.3.1.7} que para $s\in U$, tenemos
\[
  \operatorname{Ass}(\sh{F}_s)\subset\bigcup_i\operatorname{Ass}((\sh{F}_i)_s),
\]
de (i), que $\operatorname{Ass}((\sh{F}_i)_s)$ es el conjunto de los $x_{is\alpha}$ ($\alpha\in J_{s,i}$).
Por otra parte, en virtud de (iii) y del criterio \sref{IV.3.1.3},~$c'$), cada uno de los $x_{is\alpha}$ ($\alpha\in J_{s,i}$, $i\in I$) pertenece a $\operatorname{Ass}(\sh{F}_s)$.
Finalmente, (iv) significa que los $x_{is\alpha}$ son dos a dos distintos.

Queda demostrar la última afirmación del enunciado.
Ahora bien, se sigue de (i) que para $s$ y para $i$ dado, ninguno de los conjuntos $\{x_{is\alpha}\}$ puede estar contenido en otro para $\alpha\in J_{s,i}$.
Por otra parte, si $(Z_j)_\eta\subset(Z_i)_\eta$, podemos suponer $U$ tomado de modo que $(Z_j)_s\subset(Z_i)_s$ para $s\in U$ \sref{IV.9.5.1}, por lo que cada $x_{js\beta}$ pertenece a la adherencia de un $x_{is\alpha}$; por el contrario, si $(Z_j)_\eta\cap(Z_i)_\eta$ es raro en $(Z_j)_\eta$, hemos visto al demostrar (iv) que $(Z_j)_s\cap(Z_i)_s$ es raro en $(Z_j)_s$, por lo que ninguno de los $x_{js\beta}$ es adherente a un $x_{is\alpha}$.
En particular, si $(Z_j)_\eta$ es maximal, esto equivale a decir que $(Z_j)_\eta\cap(Z_i)_\eta$ es raro en $(Z_j)_\eta$ para \emph{todo} $i\neq j$, y concluimos que $(Z_j)_s\cap(Z_i)_s$ es raro en $(Z_j)_s$ para todo $i\neq j$, de modo que todo $x_{js\alpha}$ ($\alpha\in J_{s,j}$) es \emph{maximal}.
\end{proof}

\begin{corollary}[9.8.4]
\label{IV.9.8.4}
Siendo las notaciones e hipótesis las de \emph{\sref{IV.9.8.2}}, existe un entorno $U$ de $\eta$ en $S$ tal que, para todo $s\in U$, cada uno de los $(\sh{F}_i)_s$ carece de ciclo primo inmerso asociado; además, si $(\sh{F}_{is\alpha})_{\alpha\in J_{s,i}}$ es la única descomposición reducida irredundante de $(\sh{F}_i)_s$, entonces, para todo $s\in U$, la familia $(\sh{F}_{is\alpha})_{i\in I,\,\alpha\in J_{s,i}}$ es una descomposición reducida irredundante de $\sh{F}_s$.
\end{corollary}

\begin{proof}
\label{IV.9.8.4-proof}
La primera afirmación se sigue de \sref{IV.9.8.1} y de la definición de $\sh{F}_i$.
Por otra parte, hemos visto en \sref{IV.9.8.3} que el homomorfismo $\sh{F}_s\to\bigoplus_i(\sh{F}_i)_s$ es inyectivo y, por definición, lo mismo vale para cada uno de los homomorfismos $(\sh{F}_i)_s\to\bigoplus_\alpha\sh{F}_{is\alpha}$, por lo que el homomorfismo $\sh{F}_s\to\bigoplus_{i,\alpha}\sh{F}_{is\alpha}$ es inyectivo.
Como podemos suponer por \sref{IV.9.4.5} que para $s\in U$ cada uno de los $(\sh{F}_i)_s$ es un cociente de $\sh{F}_s$, los $\sh{F}_{is\alpha}$ son cocientes de $\sh{F}_s$, y queda verificar por \sref{IV.3.2.5} que los $x_{is\alpha}$ son dos a dos distintos y pertenecen a $\operatorname{Ass}(\sh{F}_s)$, lo que se ha demostrado en \sref{IV.9.8.3}.
\end{proof}

\begin{proposition}[9.8.5]
\label{IV.9.8.5}
Sea $f:X\to S$ un morfismo de presentación finita y sea $\sh{F}$ un $\sh{O}_X$-Módulo cuasicoherente de presentación finita.
Para todo $s\in S$, sea $E(s)$ (resp.~$E'(s)$) el conjunto (subconjunto finito de $\mathbf{Z}\cup\{-\infty\}$) de las dimensiones de los ciclos primos asociados a $\sh{F}_s$ (resp.~de los ciclos primos asociados maximales, es decir, las componentes irreducibles de $\operatorname{Supp}(\sh{F}_s)$).
Entonces las funciones $s\mapsto E(s)$ y $s\mapsto E'(s)$ son localmente construibles en $S$.
\end{proposition}

\begin{proof}
\label{IV.9.8.5-proof}
\oldpage[IV-3]{86}
Se sigue de \sref{IV.4.2.7} y \sref{IV.4.2.8} que la condición \sref{IV.9.2.1},~(i) se satisface para las propiedades cuya construibilidad queremos demostrar.
Nos reducimos así al caso en que $S$ es afín, noetheriano e íntegro, y a demostrar que $E$ y $E'$ son constantes en un entorno del punto genérico $\eta$ de $S$; pero esto se sigue de \sref{IV.9.8.3} y \sref{IV.9.5.5}.
\end{proof}

\begin{proposition}[9.8.6]
\label{IV.9.8.6}
Con las hipótesis y notaciones de \emph{\sref{IV.9.8.2}}, sea $i\in I$ tal que $(Z_i)_\eta$ sea maximal (en otras palabras, una componente irreducible de $\operatorname{Supp}(\sh{F}_\eta)$).
Entonces existe un entorno $U$ de $\eta$ en $S$ tal que para todo $s\in U$ y todo $\alpha\in J_{s,i}$, la longitud geométrica de $\sh{F}_s$ en $x_{is\alpha}$ (relativa a $\mathbf{k}(s)$) \sref{IV.4.7.5} es igual a la longitud geométrica de $\sh{F}_\eta$ en $x_i$ (relativa a $\mathbf{k}(\eta)$).
\end{proposition}

\begin{proof}
\label{IV.9.8.6-proof}
Podemos evidentemente restringirnos al caso en que $S=\operatorname{Spec}(A)$ es afín; demostremos primero que podemos reducirnos al caso en que el subpreesquema $(Z_i)_\eta$ de $X_\eta$, que es reducido, es \emph{geométricamente íntegro}.
En efecto, existe una extensión finita $K'$ de $K=\mathbf{k}(\eta)$ tal que $(((Z_i)_\eta)_{(K')})_{\mathrm{red}}$ es geométricamente reducido y que las componentes irreducibles de $((Z_i)_\eta)_{(K')}$ son geométricamente irreducibles \sref{IV.4.6.8}.
Procedemos como en la demostración de \sref{IV.9.7.7}, considerando una sub-$A$-álgebra $A'$ de $K'$ que tiene $K'$ como cuerpo de fracciones y es finita sobre $A$.
Ponemos $S'=\operatorname{Spec}(A')$ y consideramos el morfismo finito sobreyectivo $g:S'\to S$; sean $X'=X_{(S')}$ y $\sh{F}'=\sh{F}\otimes_{\sh{O}_S}\sh{O}_{S'}$, y sea $\eta'$ el punto genérico de $S'$.
Para todo $s'\in S'$, pongamos $s=g(s')$; si $T$ es una componente irreducible de $\operatorname{Supp}(\sh{F}_s)$, las componentes irreducibles $T'_j$ de $T_{(\mathbf{k}(s'))}$ son componentes irreducibles de $\operatorname{Supp}(\sh{F}'_{s'})$ y dominan $T$ \sref{IV.4.2.7}, y las multiplicidades radiciales de $T$ respecto de $\sh{F}_s$ y de cada $T'_j$ respecto de $\sh{F}'_{s'}$ son las mismas \sref{IV.4.7.9}.
El razonamiento de la primera parte de \sref{IV.9.7.7} muestra entonces que podemos restringirnos a demostrar la proposición para $X'$ y $\sh{F}'$; y en virtud de la elección de $K'$, los subpreesquemas reducidos cuyos espacios subyacentes son las componentes irreducibles de $((Z_i)_\eta)_{(K')}$ son geométricamente íntegros \sref{IV.4.6.1}.

Supongamos ahora que $(Z_i)_\eta$ es geométricamente íntegro; entonces, por \sref{IV.9.7.7}, lo mismo vale para $(Z_i)_s$ para $s$ en un entorno de $\eta$; la definición \sref{IV.4.7.5} muestra que bastará, por tanto, demostrar que la longitud del $\sh{O}_{X_\eta,x_i}$-módulo $(\sh{F}_\eta)_{x_i}$ es igual a la del $\sh{O}_{X_s,x_{is}}$-módulo $(\sh{F}_s)_{x_{is}}$ (donde hemos suprimido el índice $\alpha$, ahora innecesario por hipótesis).
Como la cuestión es evidentemente local sobre $X$, podemos suponer (restringiendo a un entorno de $x_i$) que $\sh{F}=\sh{F}_i$, de modo que $\sh{F}_\eta$ es \emph{irredundante} sobre $X=\operatorname{Spec}(B)$ afín, y escribimos $Z$ en lugar de $Z_i$, y $x$ para el punto genérico de $Z$ (y de $Z_\eta$).
El anillo noetheriano $B$ contiene $A$ como subanillo, y $\sh{F}=\widetilde{M}$, donde $M$ es un $B$-módulo de tipo finito; además, si $K$ es el cuerpo de fracciones de $A$, el $B_{(K)}$-módulo $M_{(K)}$ es $\neq 0$ e irredundante.
Sea $q$ el único elemento de $\operatorname{Ass}(M_{(K)})$ y sea $\mathfrak{p}$ el ideal primo de $B$, imagen inversa de $q$.
Utilizando \sref{IV.5.11.1.1} como en la demostración de \sref{IV.9.7.6}, nos reducimos al caso en que $B$ es \emph{íntegro} y $M$ un submódulo $\neq 0$ de $B$; entonces $\sh{F}_\eta$ es un sub-$\sh{O}_{X_\eta}$-Módulo no nulo de $\sh{O}_{Z_\eta}$, y en virtud de \sref{IV.9.4.5}, para $s$ en un entorno de $\eta$, $\sh{F}_s$ es isomorfo a un sub-$\sh{O}_{X_s}$-Módulo no nulo de $\sh{O}_{Z_s}$; como $Z_s$ es geométricamente íntegro, las longitudes de $(\sh{F}_\eta)_x$ y $(\sh{F}_s)_{x_s}$ son ambas iguales a $1$, lo que termina la demostración.
\end{proof}

\begin{corollary}[9.8.7]
\label{IV.9.8.7}
\oldpage[IV-3]{87}
Sea $f:X\to S$ un morfismo de presentación finita y sea $\sh{F}$ un $\sh{O}_X$-Módulo cuasicoherente de presentación finita.
Para todo $s\in S$, sea $M(s)$ el conjunto de los pares $(d,m)$ tales que existe una componente irreducible de $\operatorname{Supp}(\sh{F}_s)$ de dimensión $d$ y de multiplicidad radicial $m$ para $\sh{F}_s$ \sref{IV.4.7.8}.
Entonces la función $s\mapsto M(s)$ es localmente construible en $S$.
\end{corollary}

\begin{proof}
\label{IV.9.8.7-proof}
Se sigue de \sref{IV.4.2.7} y \sref{IV.4.7.9} que la condición \sref{IV.9.2.1},~(i) se satisface para la propiedad cuya construibilidad queremos demostrar.
Nos reducimos así al caso en que $S$ es afín, noetheriano e íntegro, y a demostrar que $M$ es constante en un entorno del punto genérico de $S$; pero esto se sigue de la proposición, a saber, \sref{IV.9.8.3} y \sref{IV.9.8.6}.
\end{proof}

\begin{proposition}[9.8.8]
\label{IV.9.8.8}
Sea $f:X\to S$ un morfismo de presentación finita y sea $\sh{F}$ un $\sh{O}_X$-Módulo cuasicoherente de presentación finita.
Para todo $s\in S$, sea $l(s)$ la suma de las multiplicidades totales para $\sh{F}_s$ (relativas a $\mathbf{k}(s)$) de los puntos genéricos de las componentes irreducibles de $\operatorname{Supp}(\sh{F}_s)$ \sref{IV.4.7.12}.
Entonces la función $s\mapsto l(s)$ es localmente construible en $S$.
\end{proposition}

\begin{proof}
\label{IV.9.8.8-proof}
Teniendo en cuenta \sref{IV.4.7.12}, la condición \sref{IV.9.2.1},~(i) se satisface para la propiedad cuya construibilidad queremos demostrar, y nos reducimos así al caso en que $S$ es afín, noetheriano e íntegro, y a demostrar que $l$ es constante en un entorno del punto genérico $\eta$.
Utilizando las notaciones de \sref{IV.9.8.2}, esto se sigue de la definición \sref{IV.4.7.12}, del hecho de que el número geométrico de componentes irreducibles de cada $(Z_i)_s$ es constante en un entorno de $\eta$ \sref{IV.9.7.8}, de que la longitud geométrica de $\sh{F}_s$ en $x_{is\alpha}$ es también igual a la de $\sh{F}_\eta$ en $x_i$ para cada $i$ tal que $(Z_i)_\eta$ es maximal \sref{IV.9.8.6}, y finalmente del hecho de que la adherencia de $\{x_{is\alpha}\}$ es un ciclo primo asociado maximal de $\sh{F}_s$ si y solo si $(Z_i)_\eta$ es un ciclo primo asociado maximal de $\sh{F}_\eta$ \sref{IV.9.8.3}.
\end{proof}

\begin{remark}[9.8.9]
\label{IV.9.8.9}
Se pueden precisar de diversas maneras las proposiciones precedentes; limitémonos a un enunciado dado a modo de ejemplo.
Diremos que un subconjunto finito $P$ de un $k$-preesquema algebraico $X$ es \emph{saturado} si, para todo par de puntos $x$, $y$ de $P$, los puntos genéricos de las componentes irreducibles de $\overline{\{x\}}\cap\overline{\{y\}}$ pertenecen también a $P$; para todo subconjunto finito $Q$ de $X$, existe un menor subconjunto finito $P$ de $X$ que contiene $Q$ y es saturado; decimos que $P$ es el \emph{saturado} de $Q$.
Para todo $\sh{O}_X$-Módulo coherente $\sh{F}$, llamaremos \emph{esqueleto primario} de $\sh{F}$ al sistema $(P,Q,\omega,d,m)$ donde $Q=\operatorname{Ass}(\sh{F})$, $P$ es el saturado de $Q$, $\omega$ la relación de orden $\overline{\{x\}}\subset\overline{\{y\}}$ sobre $P$, $d$ la función $x\mapsto\dim\overline{\{x\}}$ sobre $P$, y $m$ la función $x\mapsto\operatorname{length}_{\sh{O}_x}\sh{F}_x$ definida sobre el conjunto de los elementos maximales de $Q$ para la relación $\omega$.
Llamamos \emph{esqueleto virtual} a todo sistema $(P,Q,\omega,d,m)$ donde $P$ es un conjunto, $Q$ un subconjunto de $P$, $\omega$ una relación de orden sobre $P$, $d$ una aplicación de $P$ en $\mathbf{N}$, y $m$ una aplicación del conjunto de los elementos maximales de $Q$ en $\mathbf{N}$; se define de manera evidente la noción de \emph{isomorfismo} de dos esqueletos virtuales.
Finalmente, con las notaciones precedentes, llamamos \emph{tipo primario} de $\sh{F}$ a la clase (para la relación de isomorfismo de esqueletos virtuales) del esqueleto primario de $\sh{F}$.
Se sigue de \sref{IV.4.2.6}, \sref{IV.4.2.7}, \sref{IV.4.2.8}, \sref{IV.4.5.1} y \sref{IV.4.7.9} que si, para una extensión algebraicamente cerrada $k'$ de $k$, se ponen $X'=X\otimes_k k'$ y $\sh{F}'=\sh{F}\otimes_k k'$, el \emph{tipo primario} de $\sh{F}'$ es independiente de la
\oldpage[IV-3]{88}
extensión algebraicamente cerrada $k'$ de $k$ considerada; decimos que este es el \emph{tipo primario geométrico} de $\sh{F}$.
Dadas las definiciones, el enunciado que tenemos en vista es el siguiente:

\medskip
\noindent
\emph{\sref{IV.9.8.9.1}}
\label{IV.9.8.9.1}
Sea $f:X\to S$ un morfismo de presentación finita y sea $\sh{F}$ un $\sh{O}_X$-Módulo cuasicoherente de presentación finita.
Para todo $s\in S$, sea $T(s)$ el tipo primario geométrico de $\sh{F}_s$.
Entonces la función $s\mapsto T(s)$ es localmente construible en $S$.

\medskip
Teniendo en cuenta las observaciones precedentes, nos reducimos como de costumbre a demostrar que si $S$ es afín, noetheriano e íntegro, de punto genérico $\eta$, entonces $T(s)$ es constante en un entorno de $\eta$.
Razonando como al principio de \sref{IV.9.7.7}, podemos suponer que todas las componentes irreducibles de $X_\eta$ que intervienen son geométricamente irreducibles, y entonces la proposición se sigue de \sref{IV.9.5.1}, \sref{IV.9.5.5}, \sref{IV.9.8.3} y \sref{IV.9.8.6}; dejamos los detalles al lector.
Se podría generalizar considerando varios Módulos coherentes y definiendo su ``esqueleto primario simultáneo'', etc.
La conclusión general de este parágrafo es que para todas las propiedades del tipo considerado (y para $S$ irreducible), \emph{las propiedades válidas sobre la ``fibra genérica'' son también válidas sobre todas las fibras vecinas}.
\end{remark}

\subsection{Construibilidad de propiedades locales de las fibras}
\label{subsection:IV.9.9}
\label{IV.9.9}

\begin{proposition}[9.9.1]
\label{IV.9.9.1}
Sea $f:X\to S$ un morfismo localmente de presentación finita, sea $Z$ un subconjunto localmente construible de $X$ tal que para todo $s\in S$, $Z_s$ es cerrado en $X_s$, y sea $\Phi$ un subconjunto finito de $\mathbf{Z}\cup\{\pm\infty\}$.
Entonces los subconjuntos siguientes de $X$ son localmente construibles:
\begin{enumerate}[label=(\roman*)]
  \item El conjunto de los $x\in X$ tales que $\dim_x(Z_{f(x)})$ pertenece a $\Phi$.
  \item El conjunto de los $x\in X$ tales que $\operatorname{codim}_x(Z_{f(x)},X_{f(x)})$ pertenece a $\Phi$.
  \item El conjunto de los $x\in X$ tales que el anillo local $\sh{O}_{X_{f(x)},x}$ es equidimensional.
\end{enumerate}

Observamos que las propiedades (i) y (ii) pueden expresarse también diciendo que las funciones $x\mapsto\dim_x(Z_{f(x)})$ y $x\mapsto\operatorname{codim}_x(Z_{f(x)},X_{f(x)})$ son \emph{localmente construibles} en $X$ ($0_{\mathrm{III}}$, \sref{IV.9.3.1}).
\end{proposition}

\begin{proof}
\label{IV.9.9.1-proof}
Como las cuestiones son locales sobre $X$, podemos restringirnos al caso en que $S=\operatorname{Spec}(A)$ y $X=\operatorname{Spec}(B)$ son afines y en que $f$ es un morfismo de presentación finita; existe entonces un subanillo $A_0$ de $A$ que es una $\mathbf{Z}$-álgebra de tipo finito, un $A_0$-preesquema de tipo finito $X_0$ y un subconjunto construible $Z_0$ de $X_0$ tales que $X=X_0\otimes_{A_0}A$ y $Z=h^{-1}(Z_0)$, donde $h:X\to X_0$ es la proyección canónica (\sref{IV.8.9.1} y \sref{IV.8.3.11}).
Además, para todo $s\in S$, si $s_0$ es la proyección de $s$ en $S_0=\operatorname{Spec}(A_0)$, tenemos
\[
  X_s=(X_0)_{s_0}\otimes_{\mathbf{k}(s_0)}\mathbf{k}(s),
\]
y, si $h_s$ es la proyección $X_s\to(X_0)_{s_0}$, tenemos $Z_s=h_s^{-1}((Z_0)_{s_0})$.
Como el morfismo $h_s$ es fielmente plano y cuasicompacto, la hipótesis de que $Z_s$ es cerrado en $X_s$ implica que $(Z_0)_{s_0}$ es cerrado en $(X_0)_{s_0}$ \sref{IV.2.3.12}.

En estas condiciones, la transitividad de las fibras \sref[I]{I.3.6.4} y la prop.~\sref{IV.4.2.7} implican que el conjunto de las dimensiones de las componentes irreducibles de $Z_s$ que contienen $x$ es el mismo que el conjunto de las dimensiones de las componentes irreducibles de $(Z_0)_{s_0}$ que contienen $x_0=h_s(x)$.
En particular, tenemos $\dim_x(Z_s)=\dim_{x_0}((Z_0)_{s_0})$.
Por otra parte, si $Z_s^{(\beta)}$ son las componentes irreducibles de $Z_s$ que contienen $x$ y $X_s^{(\alpha)}$ las componentes irreducibles de $X_s$
\oldpage[IV-3]{89}
que contienen $x$, tenemos
\[
  \operatorname{codim}_x(Z_s,X_s)
  =\inf_{\beta}\sup_{\alpha}\operatorname{codim}(Z_s^{(\beta)},X_s^{(\alpha)}),
\]
donde $(\alpha,\beta)$ recorre los pares tales que $x\in Z_s^{(\beta)}\subset X_s^{(\alpha)}$ ($0$, \sref{IV.14.2.6}).
Como los preesquemas algebraicos irreducibles son biequidimensionales \sref{IV.5.2.1}, podemos escribir, en virtud de ($0$, \sref{IV.14.3.3.1})
\begin{equation}
\tag{9.9.1.1}
\label{IV.9.9.1.1}
  \operatorname{codim}_x(Z_s,X_s)=\inf_{\beta}\sup_{\alpha}\bigl(\dim(X_s^{(\alpha)})-\dim(Z_s^{(\beta)})\bigr)
\end{equation}
con la misma elección de pares $(\alpha,\beta)$.
Como $h_s$ es fielmente plano y cuasicompacto, para todo par formado por una componente irreducible $(X_0)_{s_0}^{(\lambda)}$ de $(X_0)_{s_0}$ que contiene $x_0$ y una componente irreducible $(Z_0)_{s_0}^{(\mu)}$ de $(Z_0)_{s_0}$ que contiene $x_0$ y está contenida en $(X_0)_{s_0}^{(\lambda)}$, existe un par $(X_s^{(\alpha)},Z_s^{(\beta)})$ del tipo descrito anteriormente tal que $Z_s^{(\beta)}$ domina $(Z_0)_{s_0}^{(\mu)}$ y $X_s^{(\alpha)}$ domina $(X_0)_{s_0}^{(\lambda)}$ \sref{IV.2.3.5}.
La fórmula \sref{IV.9.9.1.1} (y la fórmula análoga aplicada a $(X_0)_{s_0}$) muestra entonces, en virtud de \sref{IV.4.2.7}, que tenemos
\[
  \operatorname{codim}_x(Z_s,X_s)=\operatorname{codim}_{x_0}((Z_0)_{s_0},(X_0)_{s_0}).
\]

Vemos así que si $E$ (resp.~$E_0$) es el conjunto de los $x\in X$ (resp.~$x_0\in X_0$) que satisfacen una de las condiciones (i), (ii), (iii) del enunciado (resp.~la misma condición), entonces $E=h^{-1}(E_0)$; por \sref{IV.1.8.2}, podemos, por tanto, restringirnos al caso en que $A$, y por consiguiente también $B$, es noetheriano.
Teniendo en cuenta ($0_{\mathrm{III}}$, \sref{IV.9.2.3}) y \sref{IV.9.9.1.1}, queda demostrar que para todo $x\in X$ existe un entorno $V$ de $x$ en $\overline{\{x\}}$ tal que, para todo $x'\in V$, los conjuntos de dimensiones de las componentes irreducibles de $X_{f(x')}$ y de $Z_{f(x')}$ que contienen $x'$ son constantes, y asimismo el conjunto de pares
\[
  \bigl(\dim Z_{f(x')}^{(\beta)},\dim X_{f(x')}^{(\alpha)}\bigr),
\]
donde $X_{f(x')}^{(\alpha)}$ es una componente irreducible de $X_{f(x')}$ y $Z_{f(x')}^{(\beta)}$ es una componente irreducible de $Z_{f(x')}$ contenida en ella y que contiene $x'$.
Para ello podemos sustituir $S$ por el subpreesquema cerrado reducido $S'$ cuyo espacio subyacente es $\overline{\{f(x)\}}$, y $X$ por $X'=f^{-1}(S')$, puesto que las fibras de $X$ y $X'$ en los puntos de $S'$ son las mismas.
Así, podemos suponer que $S$ es íntegro y que $\eta=f(x)$ es su punto genérico.

Por hipótesis, $Z_\eta$ es cerrado en $X_\eta$; como $Z$ es construible, se sigue de \sref{IV.8.3.11}, aplicado según el método de \sref{IV.8.1.2},~(a), que tras sustituir $S$ por un entorno abierto de $\eta$ podemos suponer que $Z$ es \emph{cerrado} en $X$.
Sean $X_i$ (resp.~$Z_j$) las componentes irreducibles de $X$ (resp.~$Z$) que contienen $x$.
Por ($0_I$, \sref{IV.2.1.8}), los $X_i\cap X_\eta$ (resp.~$Z_j\cap X_\eta$) son las componentes irreducibles de $X_\eta$ (resp.~$Z_\eta$) que contienen $x$; por \sref{IV.9.5.1}, tras restringir $S$ alrededor de $\eta$, podemos suponer además que los $X_i$ (resp.~$Z_j$) son exactamente las componentes irreducibles de $X$ (resp.~$Z$) que encuentran $\overline{\{x\}}$, y que $(X_i)_s\cap\overline{\{x\}}$ y $(Z_j)_s\cap\overline{\{x\}}$ son no vacíos para todo $s\in S$.
En estas condiciones, se sigue aún de \sref{IV.9.7.1} que podemos suponer, restringiendo $S$, que los pares $(i,j)$ tales que $(Z_j)_s\subset(X_i)_s$ son los mismos para todo $s\in S$.
La conclusión se sigue entonces de \sref{IV.9.5.6}: para todo $s$ suficientemente próximo a $\eta$, todas las componentes irreducibles de $(X_i)_s$ (resp.~$(Z_j)_s$) tienen la misma dimensión, igual a la de $(X_i)_\eta$ (resp.~$(Z_j)_\eta$).
Además, si $(i,j)$ es un par tal que $Z_j\not\subset X_i$, $(X_i)_\eta$ no contiene el
\oldpage[IV-3]{90}
punto genérico de $(Z_j)_\eta$, de donde $\dim((X_i\cap Z_j)_\eta)<\dim((Z_j)_\eta)$; por consiguiente, la dimensión común de las componentes irreducibles de $(X_i)_s\cap(Z_j)_s$ es, para todo $s\in S$, estrictamente menor que la dimensión común de las componentes irreducibles de $(Z_j)_s$, lo que demuestra que ninguna de las componentes irreducibles de $(Z_j)_s$ está contenida en una componente irreducible de $(X_i)_s$ para $s\in S$.
\end{proof}

\begin{proposition}[9.9.2]
\label{IV.9.9.2}
Sea $f:X\to S$ un morfismo localmente de presentación finita, sea $\sh{F}$ un $\sh{O}_X$-Módulo cuasicoherente de presentación finita y sea $\Phi$ un subconjunto finito de $\mathbf{Z}\cup\{\pm\infty\}$.
Los subconjuntos siguientes de $X$ son localmente construibles:
\begin{enumerate}[label=(\roman*)]
  \item El conjunto de los $x\in X$ tales que $\sh{F}_{f(x)}$ es equidimensional en el punto $x$ \sref{IV.5.1.12}.
  \item El conjunto de los $x\in X$ tales que las longitudes geométricas de $\sh{F}_{f(x)}$ relativas a $\mathbf{k}(f(x))$, en los puntos genéricos de las componentes irreducibles de $\operatorname{Supp}(\sh{F}_{f(x)})$ que contienen $x$ \sref{IV.4.7.5}, pertenecen a $\Phi$.
  \item El conjunto de los $x\in X$ tales que las dimensiones de los ciclos primos asociados a $\sh{F}_{f(x)}$ y que contienen $x$ pertenecen a $\Phi$.
  \item El conjunto de los $x\in X$ tales que $\sh{F}_{f(x)}$ es geométricamente reducido en el punto $x$ \sref{IV.4.6.17}.
  \item El conjunto de los $x\in X$ tales que $\sh{F}_{f(x)}$ es geométricamente íntegro en el punto $x$ \sref{IV.4.6.22}.
  \item El conjunto de los $x\in X$ tales que $\dim.\operatorname{proj}((\sh{F}_{f(x)})_x)\in\Phi$.
  \item El conjunto de los $x\in X$ tales que $\operatorname{coprof}((\sh{F}_{f(x)})_x)\in\Phi$.
  \item El conjunto de los $x\in X$ tales que $\sh{F}_{f(x)}$ posee la propiedad ($S_k$) en el punto $x$ \sref{IV.5.7.2}.
  \item El conjunto de los $x\in X$ tales que $\sh{F}_{f(x)}$ es un $\sh{O}_{X_{f(x)}}$-Módulo de Cohen-Macaulay en el punto $x$ \sref{IV.5.7.1}.
\end{enumerate}
\end{proposition}

\begin{proof}
\label{IV.9.9.2-proof}
(i) El soporte $Z$ de $\sh{F}$ es localmente construible y cerrado en $X$ \sref{IV.8.9.1}, y considerando un subpreesquema de $X$ cuyo espacio subyacente es $Z$ y que es de presentación finita sobre $S$ \sref{IV.8.9.1}, vemos que la afirmación (i) es un caso particular de \sref{IV.9.9.1},~(iii).

(ii) Todas las propiedades consideradas son locales sobre $X$, y podemos, por tanto, suponer que $X=\operatorname{Spec}(B)$ y $S=\operatorname{Spec}(A)$ son afines y que $f$ es de presentación finita.
Conservamos las notaciones del comienzo de la demostración de \sref{IV.9.9.1}, y elegimos $A_0$ de modo que exista un $\sh{O}_{X_0}$-Módulo coherente $\sh{F}_0$ con $\sh{F}\simeq\sh{F}_0\otimes_{A_0}A$.
Entonces \sref{IV.4.2.7} y \sref{IV.4.7.9} muestran que el conjunto de longitudes geométricas de $(\sh{F}_0)_{f(x_0)}$ en los puntos genéricos de las componentes irreducibles de su soporte que contienen $x_0$ es el mismo que el conjunto correspondiente para $\sh{F}_{f(x)}$ en $x$.
Así, si $E$ (resp.~$E_0$) denota el conjunto de los $x\in X$ (resp.~$x_0\in X_0$) que satisfacen la condición (ii), entonces $E=h^{-1}(E_0)$, y por \sref{IV.1.8.2} podemos restringirnos al caso en que $A$ es noetheriano.
Como en la demostración de \sref{IV.9.9.1}, queda demostrar que para todo $x\in X$ existe un entorno $V$ de $x$ en $\overline{\{x\}}$ sobre el cual es constante el conjunto de longitudes geométricas de $\sh{F}_{f(x')}$ en los puntos genéricos de las componentes irreducibles de su soporte que contienen $x'$.
Si $S'$ es el subpreesquema cerrado reducido de $S$ cuyo espacio subyacente es $\overline{\{f(x)\}}$, y si $X'=f^{-1}(S')$, las fibras de $X$ y $X'$ en los puntos de $S'$ son las mismas,
\oldpage[IV-3]{91}
por lo que podemos sustituir $S$ por $S'$ y $X$ por $X'$; en otras palabras, suponer que $S$ es íntegro y que $\eta=f(x)$ es su punto genérico.
Si $Z=\operatorname{Supp}(\sh{F})$ y los $Z_i$ son las componentes irreducibles de $Z$ que contienen $x$, el argumento de \sref{IV.9.9.1} nos permite suponer que estas son exactamente las componentes irreducibles de $Z$ que encuentran $\overline{\{x\}}$ y que $(Z_i)_s\cap\overline{\{x\}}\neq\varnothing$ para todo $s\in S$.
La conclusión se sigue de \sref{IV.9.8.3} y \sref{IV.9.8.6}.

(iii) Utilizando \sref{IV.4.2.7}, nos reducimos como en (ii) al caso en que $S$ es noetheriano e íntegro y $\eta=f(x)$ es su punto genérico.
Como en la demostración de \sref{IV.9.9.1}, queda encontrar un entorno $V$ de $x$ en $\overline{\{x\}}$ sobre el cual sea constante el conjunto de dimensiones de los ciclos primos asociados a $\sh{F}_{f(x')}$ y que contienen $x'$.
Sean $Z'_i$ las adherencias en $X$ de los ciclos primos asociados a $\sh{F}_\eta$ que contienen $x$ (cf.~\sref{IV.9.8.2}).
Por \sref{IV.9.8.3} y \sref{IV.9.5.1}, para todo $s$ suficientemente próximo a $\eta$, los ciclos primos asociados a $\sh{F}_s$ que encuentran $\overline{\{x\}}$ son componentes irreducibles de los $(Z'_i)_s$; por \sref{IV.9.5.6}, todas estas componentes tienen dimensión $\dim((Z'_i)_\eta)$, lo que demuestra la afirmación.

(iv) Razonamos como en (iii), utilizando ahora \sref{IV.4.7.11}.
Con las mismas notaciones, los ciclos primos asociados a $\sh{F}_s$ que son componentes irreducibles de $(Z'_i)_s$ están inmersos si y solo si $(Z'_i)_\eta$ es un ciclo primo asociado inmerso de $\sh{F}_\eta$.
Se sigue que, según que $x$ pertenezca a un ciclo primo inmerso asociado de $\sh{F}_\eta$ o a ninguno, los puntos $x'\in\overline{\{x\}}$ que poseen la propiedad análoga para $\sh{F}_{f(x')}$ forman un entorno de $x$ en $\overline{\{x\}}$.
La conclusión se sigue de esta observación, de la caracterización de los Módulos geométricamente reducidos en \sref{IV.4.7.10}, y de (ii).

(v) Teniendo en cuenta \sref{IV.4.7.11}, nos reducimos de nuevo al caso en que $S$ es noetheriano e íntegro y en que $\eta=f(x)$ es su punto genérico.
Sea $Y$ un subpreesquema cerrado de $X$ cuyo espacio subyacente es $\operatorname{Supp}(\sh{F})$.
Como al principio de la demostración de \sref{IV.9.7.7}, existe una $A$-álgebra íntegra finita $A'$ (de modo que, si $S'=\operatorname{Spec}(A')$, el morfismo $S'\to S$ es finito y sobreyectivo) tal que, si $Y'=Y_{(S')}$ y $\eta'$ es el punto genérico de $S'$, las componentes irreducibles de $Y'_{\eta'}$ son geométricamente irreducibles.
Si $X'=X_{(S')}$, la proyección $h:X'\to X$ es finita y sobreyectiva, por tanto cerrada; por consiguiente, para un punto $x'$ situado sobre $x$, se tiene $h(\overline{\{x'\}})=\overline{\{x\}}$, y la imagen por $h$ de un entorno de $x'$ en $\overline{\{x'\}}$ es un entorno de $x$ en $\overline{\{x\}}$.
Escribiendo $\sh{F}'=\sh{F}_{(S')}$ y utilizando \sref{IV.4.7.11}, nos reducimos a demostrar que, según que $\sh{F}'_{\eta'}$ sea o no geométricamente íntegro en $x'$, los puntos $y'\in\overline{\{x'\}}$ donde $\sh{F}'_{f'(y')}$ es o no geométricamente íntegro en $y'$ forman un entorno de $x'$ en $\overline{\{x'\}}$.
Podemos, por tanto, suponer que $X'=X$ y que las componentes irreducibles de $Y_\eta$ son geométricamente irreducibles.
Elijamos subconjuntos cerrados $Z_i$ de $X$ tales que los $Z_i\cap X_\eta$ sean las
  \oldpage[IV-3]{92}
  componentes irreducibles de $Y_\eta$.
Por \sref{IV.9.7.7}, \sref{IV.9.7.8} y \sref{IV.9.5.3}, para todo $s$ próximo a $\eta$ los $(Z_i)_s$ son las componentes irreducibles de $Y_s$ y son geométricamente irreducibles.
En un punto $y\in X_s$, el Módulo $\sh{F}_s$ es geométricamente íntegro exactamente cuando es geométricamente reducido allí y $y$ pertenece a uno solo de los $(Z_i)_s$ \sref{IV.4.6.22}.
La conclusión se sigue de (iv) y de \sref{IV.9.5.1} aplicado a $\overline{\{x\}}\cap Z_i$ para cada $i$.

(vi) Conservando las notaciones de (ii), \sref{IV.6.2.1} da
\[
  \dim.\operatorname{proj}((\sh{F}_s)_x)
  =\dim.\operatorname{proj}(((\sh{F}_0)_{s_0})_{x_0}),
\]
de modo que podemos restringirnos de nuevo al caso en que $A$ es noetheriano. Queda demostrar que para todo $x\in X$ existe un entorno $V$ de $x$ en $\overline{\{x\}}$ tal que
\[
  \dim.\operatorname{proj}((\sh{F}_{f(x')})_{x'})
  =\dim.\operatorname{proj}((\sh{F}_{f(x)})_x)
\]
para todo $x'\in V$. Como antes, podemos suponer que $S$ es íntegro y que $\eta=f(x)$ es su punto genérico, de modo que $(\sh{F}_\eta)_x=\sh{F}_x$. Por planitud genérica \sref{IV.6.9.1}, tras sustituir $S$ por un entorno abierto de $\eta$, podemos suponer que $f$ es plano y que $\sh{F}$ es $f$-plano. Entonces \sref{IV.6.2.3} da
\[
  \dim.\operatorname{proj}((\sh{F}_{f(x')})_{x'})
  =\dim.\operatorname{proj}(\sh{F}_{x'})
\]
para todo $x'\in X$. Por \sref{IV.6.11.1}, tras sustituir $X$ por un entorno abierto de $x$, podemos suponer que
\[
  \dim.\operatorname{proj}(\sh{F}_{x'})
  \leqslant\dim.\operatorname{proj}(\sh{F}_x)
\]
para todo $x'\in X$.

Si $\dim.\operatorname{proj}(\sh{F}_x)=n$, existe un $\sh{O}_x$-módulo finito $M$ tal que $\operatorname{Ext}^n_{\sh{O}_x}(\sh{F}_x,M)\neq 0$ \sref[0]{0.17.2.4}.
Tras sustituir $X$ por un entorno abierto de $x$, existe un $\sh{O}_X$-módulo coherente $\sh{G}$ tal que $M=\sh{G}_x$ \sref[0]{0.5.3.8};
en virtud de (T, \sref{IV.4.2.2}), tenemos $(\shExt^n_{\sh{O}_X}(\sh{F},\sh{G}))_x\neq 0$.
Pero $\shExt^n_{\sh{O}_X}(\sh{F},\sh{G})$ es un $\sh{O}_X$-Módulo coherente ($0_{\mathrm{III}}$, \sref{IV.12.3.3}), por tanto su soporte $Y$ es cerrado ($0_I$, \sref{IV.5.2.2}); como contiene $x$, contiene también $\overline{\{x\}}$, de donde concluimos (aplicando de nuevo (T, \sref{IV.4.2.2})) que $\dim.\operatorname{proj}(\sh{F}_{x'})\geqslant n$ para todo $x'\in\overline{\{x\}}$, lo que termina la demostración de la afirmación en el caso (vi).

(vii) Como $B$ es una $A$-álgebra de tipo finito, $X$ es $S$-isomorfo a un \emph{subesquema cerrado} de un $S$-esquema de la forma $Y=\operatorname{Spec}(A[T_1,\ldots,T_r])$; si $j:X\to Y$ es la inyección canónica y $\sh{G}=j_*(\sh{F})$, entonces $\sh{G}_s=(j_s)_*(\sh{F}_s)$ para todo $s\in S$, y, en virtud de ($0$, \sref{IV.16.4.11}), podemos restringirnos a demostrar la afirmación para $Y$ y $\sh{G}$.
En otras palabras, podemos suponer que $B=A[T_1,\ldots,T_r]$, de modo que los esquemas $X_s=\operatorname{Spec}(\mathbf{k}(s)[T_1,\ldots,T_r])$ son \emph{regulares} ($0$, \sref{IV.17.3.7}).
Sea $W=\operatorname{Supp}(\sh{F})$ y $W_s=\operatorname{Supp}(\sh{F}_s)$ \sref[I]{I.9.1.13}; tenemos, por \sref{IV.6.11.2.1},
\begin{equation}
\tag{9.9.2.1}
\label{IV.9.9.2.1}
  \operatorname{coprof}((\sh{F}_{f(x)})_x)=\dim.\operatorname{proj}((\sh{F}_{f(x)})_x)-\operatorname{codim}_x(W_{f(x)},X_{f(x)}).
\end{equation}

Pero como $W$ es construible \sref{IV.8.9.1} y cada $W_s$ es cerrado, se sigue de (vi) y \sref{IV.9.9.1},~(ii), que las dos funciones del segundo miembro de \sref{IV.9.9.2.1} son construibles; por tanto, también lo es su diferencia, lo que demuestra la afirmación en el caso (vii).

(viii) Sea $U_n$ el conjunto de los $x\in X$ tales que $\operatorname{coprof}((\sh{F}_{f(x)})_x)\leqslant n$, y pongamos $Z_n=X-U_n$; se sigue de (vii) que los $Z_n$ son \emph{construibles}.
\oldpage[IV-3]{93}
Además, como la función $x\mapsto\dim_x(W_{f(x)})$ es construible por \sref{IV.9.9.1},~(i), toma solo un número finito de valores, y los números $\dim(W_{f(x)})$ tienen, por tanto, una cota superior finita $m$ cuando $x$ recorre $X$.
Como $\operatorname{coprof}((\sh{F}_{f(x)})_x)\leqslant\dim((\sh{F}_{f(x)})_x)\leqslant\dim(W_{f(x)})$, tenemos $Z_n=\varnothing$ para $n\geqslant m$.
Finalmente, \sref{IV.6.11.2},~(i), muestra que $(Z_n)_s$ es cerrado en $X_s$ para todo $n$ y todo $s\in S$.
Por \sref{IV.5.7.4}, el conjunto de los $x\in X$ donde $(\sh{F}_{f(x)})_x$ posee la propiedad ($S_k$) es el conjunto de los $x\in X$ que satisfacen todas las relaciones
\begin{equation}
\tag{9.9.2.2}
\label{IV.9.9.2.2}
  \operatorname{codim}_x((Z_n)_{f(x)},W_{f(x)})>n+k
\end{equation}
para todo $n\geqslant 0$; como esta relación se satisface automáticamente para $n\geqslant m$, basta considerar las relaciones \sref{IV.9.9.2.2} para $0\leqslant n<m$.
Pero, en virtud de \sref{IV.9.9.1},~(ii), el conjunto $V_{n,k}$ de los $x$ que satisfacen \sref{IV.9.9.2.2} es construible, y lo mismo vale para la intersección de estos conjuntos para $0\leqslant n<m$.

(ix) La afirmación se sigue inmediatamente de (vii), pues el conjunto de los $x\in X$ donde $(\sh{F}_{f(x)})_x$ es un módulo de Cohen-Macaulay está definido por la relación $\operatorname{coprof}((\sh{F}_{f(x)})_x)=0$.
\end{proof}

\begin{corollary}[9.9.3]
\label{IV.9.9.3}
Sea $f:X\to S$ un morfismo de presentación finita, sea $\sh{F}$ un $\sh{O}_X$-Módulo cuasicoherente de presentación finita y sea $\mathbf{P}$ una de las propiedades (i) a (ix) de \sref{IV.9.9.2}.
Entonces el conjunto de los $s\in S$ tales que la propiedad $\mathbf{P}$ es cierta en todos los puntos $x\in X_s$, es localmente construible en $S$.
\end{corollary}

\begin{proof}
\label{IV.9.9.3-proof}
En efecto, su complemento es la imagen por $f$ del complemento del conjunto $E$ de los puntos donde $\mathbf{P}$ es cierta.
Como $E$ es localmente construible en $X$, también lo es $X-E$, y por tanto $f(X-E)$ es localmente construible en $S$ en virtud del teorema de Chevalley \sref{IV.1.8.4}.
\end{proof}

\begin{proposition}[9.9.4]
\label{IV.9.9.4}
Sea $f:X\to S$ un morfismo localmente de presentación finita.
El conjunto de los $x\in X$ tales que $X_{f(x)}$ posee en el punto $x$ una de las propiedades (determinadas) siguientes:
\begin{enumerate}[label=(\roman*)]
  \item ser geométricamente regular;
  \item poseer la propiedad ($R_k$) geométrica;
  \item ser geométricamente normal;
  \item ser geométricamente reducido (es decir, separable);
  \item ser geométricamente puntualmente íntegro;
\end{enumerate}
es un subconjunto localmente construible de $X$.
\end{proposition}

\begin{proof}
\label{IV.9.9.4-proof}
Para las propiedades (iv) y (v), esto se sigue de \sref{IV.9.9.2},~(iv) y~(v), aplicado a $\sh{F}=\sh{O}_X$.
Para las demás propiedades, teniendo en cuenta \sref{IV.6.7.8}, nos reducimos, como al principio de \sref{IV.9.9.2},~(ii), al caso en que $S$ es noetheriano e íntegro y de punto genérico $\eta=f(x)$.
Además, en virtud del teorema de planitud genérica \sref{IV.6.9.1}, podemos, sustituyendo $S$ por un entorno abierto de $\eta$, suponer que $f$ es un morfismo \emph{plano}.
Decir que $X_{f(y)}$ es geométricamente regular en el punto $y$ significa que $f$ es \emph{regular} en el punto $y$ \sref{IV.6.8.1}, y sabemos que el conjunto $L$ de estos puntos es \emph{abierto} en $X$ \sref{IV.6.8.7}, lo que demuestra la proposición en el caso (i).

Para demostrar el caso (ii), pongamos $F=X-L$, que es \emph{cerrado} en $X$.
Decir que $X_s$
\oldpage[IV-3]{94}
satisface la propiedad geométrica ($R_k$) en el punto $y\in f^{-1}(s)$ significa o bien que $y\in L$, o bien que los puntos genéricos $z_i$ de las componentes irreducibles del conjunto cerrado $F_s$ que contienen $y$ satisfacen la relación $\dim(\sh{O}_{X_s,z_i})\geqslant k+1$ (teniendo en cuenta \sref{IV.4.2.7} y \sref{IV.5.2.3}); en otras palabras, los puntos $y\in F_s$ donde $X_s$ satisface la propiedad ($R_k$) geométrica son aquellos tales que $\operatorname{codim}_y(F_s,X_s)\geqslant k+1$ \sref{IV.5.1.2}.
La conclusión se sigue, pues, de (i) y \sref{IV.9.9.1},~(ii).

Finalmente, en el caso (iii), la conclusión se sigue de (ii), de \sref{IV.9.9.2},~(viii), del hecho de que la propiedad ($S_k$) es estable por extensión del cuerpo base \sref{IV.6.7.1} y, finalmente, del criterio de Serre \sref{IV.5.8.6}.
\end{proof}

\begin{corollary}[9.9.5]
\label{IV.9.9.5}
Sea $f:X\to S$ un morfismo de presentación finita, y denotemos por $\mathbf{P}$ una de las propiedades (i) a (v) de \sref{IV.9.9.4}.
Entonces el conjunto de los $s\in S$ tales que la propiedad $\mathbf{P}$ es cierta en todos los puntos $x\in X_s$, es localmente construible en $S$.
\end{corollary}

\begin{proof}
\label{IV.9.9.5-proof}
La demostración a partir de \sref{IV.9.9.4} es la misma que la de \sref{IV.9.9.3} a partir de \sref{IV.9.9.2}.
\end{proof}

\begin{proposition}[9.9.6]
\label{IV.9.9.6}
Sea $f:X\to S$ un morfismo localmente de presentación finita, y sea $\sh{L}_\bullet$ un complejo formado por $\sh{O}_X$-Módulos cuasicoherentes de presentación finita; para todo $s\in S$, sea $(\sh{L}_\bullet)_s$ el complejo $\sh{L}_\bullet\otimes_{\sh{O}_S}\mathbf{k}(s)$ de $\sh{O}_{X_s}$-Módulos coherentes de tipo finito.
Entonces, para un entero dado $n$, el conjunto de los $x\in X$ tales que $(\sh{H}_n((\sh{L}_\bullet)_{f(x)}))_x=0$ es localmente construible en $X$.
\end{proposition}

\begin{proof}
\label{IV.9.9.6-proof}
Podemos restringirnos al caso en que $\sh{L}_i=0$ salvo para $i=0$, $1$ o $2$, y en que $n=1$.
Además, como la cuestión es local sobre $X$, podemos restringirnos al caso en que $S=\operatorname{Spec}(A)$ y $X=\operatorname{Spec}(B)$ son afines, siendo $B$ una $A$-álgebra de presentación finita.
Existe entonces un subanillo noetheriano $A_0$ de $A$, un $A_0$-preesquema de tipo finito $X_0$ y un complejo $\sh{L}_\bullet^{(0)}$ de $\sh{O}_{X_0}$-Módulos coherentes, nulo salvo en dimensiones $0$, $1$ y $2$ y tal que $X=X_0\otimes_{A_0}A$ y $\sh{L}_\bullet=\sh{L}_\bullet^{(0)}\otimes_{A_0}A$.
Para todo $s\in S$, si $s_0$ es la proyección de $s$ en $S_0=\operatorname{Spec}(A_0)$, entonces
\[
  X_s=(X_0)_{s_0}\otimes_{\mathbf{k}(s_0)}\mathbf{k}(s),
\]
y la proyección $X_s\to(X_0)_{s_0}$ es fielmente plana.
Por consiguiente,
\[
  \sh{H}_n((\sh{L}_\bullet)_s)
  =\sh{H}_n((\sh{L}_\bullet^{(0)})_{s_0})\otimes_{\mathbf{k}(s_0)}\mathbf{k}(s).
\]
Si $E$ (resp.~$E_0$) es el conjunto de los $x\in X$ (resp.~$x_0\in X_0$) para los cuales
\[
  \begin{aligned}
  (\sh{H}_n((\sh{L}_\bullet)_{f(x)}))_x&=0\\
  \bigl(\text{resp. }(\sh{H}_n((\sh{L}_\bullet^{(0)})_{f_0(x_0)}))_{x_0}&=0\bigr),
  \end{aligned}
\]
entonces $E=h^{-1}(E_0)$, donde $h:X\to X_0$ es la proyección canónica.
Por \sref{IV.1.8.2}, podemos, por tanto, restringirnos al caso en que $A$ es noetheriano.
Queda demostrar ($0_{\mathrm{III}}$, \sref{IV.9.2.3}) que si $x\in X$ es tal que $(\sh{H}_n((\sh{L}_\bullet)_{f(x)}))_x=0$ (resp.~$\neq 0$), existe un entorno abierto $V$ de $x$ en $\overline{\{x\}}$ tal que para todo $x'\in V$ se tiene $(\sh{H}_n((\sh{L}_\bullet)_{f(x')}))_{x'}=0$ (resp.~$\neq 0$).
Sustituyendo $S$ por el subpreesquema cerrado reducido cuyo espacio subyacente es $\overline{\{f(x)\}}$, podemos suponer que $S$ es íntegro y que $f(x)=\eta$ es su punto genérico.
Entonces, restringiendo $S$ a un entorno abierto de $\eta$, podemos suponer que para todo $s\in S$,
$(\sh{H}_n(\sh{L}_\bullet))_s=\sh{H}_n((\sh{L}_\bullet)_s)$ \sref{IV.9.4.3}; por consiguiente, si $Z$ es el soporte de $\sh{H}_n(\sh{L}_\bullet)$, entonces el soporte de $\sh{H}_n((\sh{L}_\bullet)_s)$ es $Z_s=Z\cap X_s$ \sref[I]{I.9.1.13.1}.
La hipótesis es que $Z_\eta\cap\overline{\{x\}}=\varnothing$ (resp.~$\neq\varnothing$).
Como $Z\cap\overline{\{x\}}$ es cerrado en el espacio noetheriano $X$, concluimos de \sref{IV.9.5.1} que existe un entorno $U$ de $\eta$ en $S$ tal que, para todo $s\in U$, tenemos $Z_s\cap\overline{\{x\}}=\varnothing$ (resp.~$\neq\varnothing$).
El conjunto $V=f^{-1}(U)\cap\overline{\{x\}}$ responde así a la cuestión.
\end{proof}

\oldpage[IV-3]{95}
\section{Preesquemas de Jacobson}
\label{section:IV.10}

Ya hemos tenido ocasión de observar \sref{IV.5.2.5} que incluso los
preesquemas \emph{excelentes} \sref{IV.7.8.5} no se comportan siempre como las
``variedades'' de la Geometría Algebraica clásica, especialmente en lo relativo
a cuestiones de dimensión; así, si $X$ es el espectro de un anillo de valoración
discreta completo, el conjunto de los puntos cerrados de $X$ (reducido a un solo
elemento) no es denso en todas partes en $X$, mientras que su complementario
(reducido a un solo elemento) es un abierto denso en todas partes de $X$ cuya
dimensión es cero, por tanto $<\dim(X)$. En esta sección examinamos condiciones
generales bajo las cuales no se producen tales fenómenos; como consecuencia,
las relaciones entre dimensión, codimensión, profundidad y coprofundidad en
estos preesquemas se comportan más satisfactoriamente en ciertos aspectos (10.6
y 10.8). Además, en los preesquemas considerados, el hecho de que el conjunto
de los puntos cerrados sea denso en todas partes (e incluso ``muy denso''
\sref{IV.10.1.3}) permite considerar únicamente estos puntos en muchas
demostraciones; recuperamos así el punto de vista clásico de las ``variedades
algebraicas'', que, desde nuestro punto de vista, son los conjuntos de puntos
cerrados de los \emph{preesquemas algebraicos sobre un cuerpo}, y establecemos
el vínculo entre el lenguaje de los esquemas y el de las ``variedades'' o
``espacios algebraicos'' de Serre (10.9 y 10.10).

\subsection{Subconjuntos muy densos de un espacio topológico}
\label{subsection:IV.10.1}

\begin{env}[10.1.1]
\label{IV.10.1.1}
Decimos que un subconjunto de un espacio topológico $X$ es \emph{cuasiconstruible} si es una unión finita de subconjuntos localmente cerrados de $X$.
Decimos que un subconjunto $T$ de $X$ es \emph{localmente cuasiconstruible} si, para todo $x\in X$, existe un entorno abierto $V$ de $x$ tal que $T\cap V$ es cuasiconstruible en $V$.
Está claro que todo subconjunto cuasiconstruible de $X$ es localmente cuasiconstruible; el recíproco es cierto si $X$ es cuasicompacto.
El razonamiento de \sref[$0_{\mathrm{III}}$]{0.9.1.3}, en el que se suprime la palabra ``retrocompacto'', muestra que el conjunto de los subconjuntos cuasiconstruibles (resp.\ localmente cuasiconstruibles) de $X$, que denotaremos por $\mathfrak{Qc}(X)$ (resp.\ $\mathfrak{Lqc}(X)$), es estable por unión finita, intersección finita y paso al complementario.
Si $f:X\to Y$ es una aplicación continua, se deduce inmediatamente de estas definiciones que, para todo subconjunto cuasiconstruible (resp.\ localmente cuasiconstruible) $Z$ de $Y$, $f^{-1}(Z)$ es cuasiconstruible (resp.\ localmente cuasiconstruible) en $X$.

Los subconjuntos construibles (resp.\ localmente construibles) de $X$ son evidentemente cuasiconstruibles (resp.\ localmente cuasiconstruibles); el recíproco es cierto cuando $X$ es noetheriano (resp.\ localmente noetheriano).

En lo que sigue, denotaremos respectivamente por $\mathfrak{O}(X)$, $\mathfrak{F}(X)$, $\mathfrak{Lf}(X)$, $\mathfrak{C}(X)$, $\mathfrak{Lc}(X)$ el conjunto de los subconjuntos de $X$ que son respectivamente abiertos, cerrados, localmente cerrados, construibles y localmente construibles.
\end{env}

\oldpage[IV-3]{96}

\begin{proposition}[10.1.2]
\label{IV.10.1.2}
Sea $X$ un espacio topológico y $X_0$ un subconjunto de $X$.
Las condiciones siguientes son equivalentes:
\begin{enumerate}[label=\emph{\alph*)}]
  \item Para todo subconjunto localmente cerrado $Z\neq\varnothing$ de $X$, se tiene $Z\cap X_0\neq\varnothing$.
  \item[\rm a$'$)] Para todo subconjunto cerrado $Z$ de $X$, se tiene $Z=\overline{Z\cap X_0}$.
  \item Para todo subconjunto localmente cuasiconstruible $Z\neq\varnothing$ de $X$, se tiene $Z\cap X_0\neq\varnothing$.
  \item[\rm b$'$)] Para todo subconjunto localmente cuasiconstruible $Z$ de $X$, se tiene $Z\subset\overline{Z\cap X_0}$ (dicho de otro modo, $Z\cap X_0$ es denso en $Z$).
  \item La aplicación $U\mapsto U\cap X_0$ de $\mathfrak{O}(X)$ en $\mathfrak{O}(X_0)$ es inyectiva (y por tanto biyectiva).
  \item[\rm c$'$)] La aplicación $F\mapsto F\cap X_0$ de $\mathfrak{F}(X)$ en $\mathfrak{F}(X_0)$ es inyectiva (y por tanto biyectiva).
  \item[\rm c$''$)] La aplicación $Z\mapsto Z\cap X_0$ de $\mathfrak{Qc}(X)$ en $\mathfrak{Qc}(X_0)$ es inyectiva (y por tanto biyectiva).
  \item[\rm c$'''$)] La aplicación $Z\mapsto Z\cap X_0$ de $\mathfrak{Lqc}(X)$ en $\mathfrak{Lqc}(X_0)$ es inyectiva.
\end{enumerate}
Además, cuando se satisfacen estas condiciones, la aplicación $Z\mapsto Z\cap X_0$ de $\mathfrak{Lqc}(X)$ en $\mathfrak{Lqc}(X_0)$ es biyectiva.
\end{proposition}

\begin{proof}
Observemos que las afirmaciones de sobreyectividad de $c)$ y $c')$ son triviales; significan que todo subconjunto localmente cerrado de $X_0$ es la traza sobre $X_0$ de un subconjunto localmente cerrado de $X$, y que la aplicación $\mathfrak{Qc}(X)\to\mathfrak{Qc}(X_0)$ definida en $c'')$ también es sobreyectiva.

Demostraremos las implicaciones
\[
  c''')\Rightarrow c'')\Rightarrow c')\Rightarrow c)\Rightarrow b)\Rightarrow b')\Rightarrow a')\Rightarrow a)\Rightarrow c''').
\]

Las tres primeras son triviales.
Para ver que $c)$ implica $b)$, observemos primero que $c)$ significa que $X_0$ es \emph{denso} en $X$.
Sustituyendo respectivamente $X$ y $X_0$ por $U$ y $U\cap X_0$, podemos por tanto limitarnos al caso en que $Z$ es localmente cerrado en $X$, y entonces $Z=V\cap\complement W$, donde $V$ y $W$ son abiertos de $X$; la hipótesis $Z\neq\varnothing$ significa que $V\not\subset W$, o equivalentemente $V\cup W\neq W$.
En virtud de $c)$, se tiene $(V\cup W)\cap X_0\neq W\cap X_0$, luego $V\cap X_0\not\subset W$, y por consiguiente $(V\cap\complement W)\cap X_0\neq\varnothing$.

Para ver que $b)$ implica $b')$, basta aplicar $b)$ a $Z\cap U$, donde $U$ es un entorno abierto de un punto arbitrario de $Z$.
Como $Z\cap\overline{X_0}\subset\overline{Z\cap X_0}$, es trivial que $b')$ implica $a')$.
Para mostrar que $a')$ implica $a)$, observemos que, si $Z$ es localmente cerrado en $X$, podemos escribir $Z=F-F'$, donde $F$ y $F'$ son cerrados de $X$ y $F'\subset F$; de donde $Z\cap X_0=(F\cap X_0)-(F'\cap X_0)$.
Si $Z\cap X_0=\varnothing$, entonces, en virtud de $a')$, tendríamos $F\cap X_0=F'\cap X_0$, luego $F=F'$, es decir, $Z=\varnothing$.

Para ver que $a)$ implica $c''')$, basta mostrar que, si $Z\neq\varnothing$ es localmente cuasiconstruible, entonces $Z\cap X_0\neq\varnothing$: en efecto, la relación $Z'\cap X_0=Z''\cap X_0$ equivale a $((Z'\cup Z'')-(Z'\cap Z''))\cap X_0=\varnothing$.
Dicho de otro modo, basta probar que $a)$ implica $b)$; en efecto, sustituyendo respectivamente $X$ y $X_0$ por $U$ y $U\cap X_0$, nos reducimos al caso en que $Z$ es localmente cerrado en $X$, de donde se sigue la conclusión.

Queda por mostrar que la aplicación $\mathfrak{Lqc}(X)\to\mathfrak{Lqc}(X_0)$ es sobreyectiva.
Sea pues $Z_0$ un subconjunto localmente cuasiconstruible de $X_0$: existe un recubrimiento $(V_\alpha)$ de $X_0$ por abiertos de $X_0$ tal que $Z_0\cap V_\alpha$ es cuasiconstruible en $V_\alpha$ (y
\oldpage[IV-3]{97}
por tanto también en $X_0$).
En virtud de $c)$, para todo $\alpha$ existe un único abierto $U_\alpha$ de $X$ tal que $X_0\cap U_\alpha=V_\alpha$, y en virtud de $c'')$ un único subconjunto cuasiconstruible $Z_\alpha$ de $X$ tal que $Z_0\cap V_\alpha=X_0\cap Z_\alpha$.
Si $\alpha$ y $\beta$ son dos índices cualesquiera, se tiene $Z_\beta\cap U_\alpha\cap X_0=U_\alpha\cap Z_0\cap V_\beta=V_\alpha\cap Z_0\cap V_\beta=Z_\alpha\cap X_0\cap V_\beta\subset Z_\alpha\cap X_0$; como $Z_\beta\cap U_\alpha$ y $Z_\alpha$ son cuasiconstruibles en $X$, se deduce de $c'')$ que $Z_\beta\cap U_\alpha\subset Z_\alpha$.
Si ponemos $Z=\bigcup_\alpha Z_\alpha$, entonces $Z\cap U_\alpha=Z_\alpha$ para todo $\alpha$; además, puesto que los $V_\alpha$ recubren $X_0$, se deduce de $c)$ que los $U_\alpha$ recubren $X$, y vemos por tanto que $Z$ es localmente cuasiconstruible en $X$ y que $Z_0=Z\cap X_0$.
\end{proof}

\begin{definition}[10.1.3]
\label{IV.10.1.3}
Cuando un subconjunto $X_0$ de un espacio topológico satisface las condiciones equivalentes de \sref{IV.10.1.2}, decimos que $X_0$ es \emph{muy denso} en $X$.
\end{definition}

Ya hemos visto en el curso de la demostración de \sref{IV.10.1.2} que entonces $X_0$ es denso en $X$.

\begin{corollary}[10.1.4]
\label{IV.10.1.4}
Si $X_0$ es muy denso en $X$ y $U$ es un abierto de $X$, $U\cap X_0$ es muy denso en $U$.
Recíprocamente, si $(U_\alpha)$ es un recubrimiento abierto de $X$ tal que $U_\alpha\cap X_0$ es muy denso en $U_\alpha$ para todo $\alpha$, entonces $X_0$ es muy denso en $X$.
\end{corollary}

Como todo subconjunto localmente cerrado de $U$ es localmente cerrado en $X$, la primera afirmación se deduce del criterio $a)$ de \sref{IV.10.1.2}; y lo mismo vale para la segunda, pues si $Z\neq\varnothing$ es localmente cerrado en $X$, $Z\cap U_\alpha$ es localmente cerrado en $U_\alpha$ para todo $\alpha$, y $Z\cap U_\alpha\neq\varnothing$ para al menos un $\alpha$.

\subsection{Cuasihomeomorfismos}
\label{subsection:IV.10.2}

\begin{proposition}[10.2.1]
\label{IV.10.2.1}
Sean $X_0$, $X$ dos espacios topológicos y $f:X_0\to X$ una aplicación continua.
Las condiciones siguientes son equivalentes:
\begin{enumerate}[label=\emph{\alph*)}]
  \item La aplicación $U\mapsto f^{-1}(U)$ de $\mathfrak{O}(X)$ en $\mathfrak{O}(X_0)$ es biyectiva.
  \item[\rm a$'$)] La aplicación $F\mapsto f^{-1}(F)$ de $\mathfrak{F}(X)$ en $\mathfrak{F}(X_0)$ es biyectiva.
  \item La topología de $X_0$ es la imagen inversa por $f$ de la de $X$, y $f(X_0)$ es muy denso \sref{IV.10.1.3} en $X$.
  \item El funtor $\sh{F}\mapsto f^*(\sh{F})$ de la categoría de haces de conjuntos (resp.\ haces de grupos abelianos) sobre $X$ en la categoría de haces de conjuntos (resp.\ haces de grupos abelianos) sobre $X_0$ es una equivalencia de categorías.
\end{enumerate}
\end{proposition}

\begin{proof}
Está claro que $a)$ y $a')$ son equivalentes, y que $a)$ implica que la topología de $X_0$ es la imagen inversa por $f$ de la de $X$.
Por otra parte, si $f(X_0)$ no es muy denso en $X$, existen dos abiertos distintos $U_1$, $U_2$ de $X$ tales que $U_1\cap f(X_0)=U_2\cap f(X_0)$, y por tanto $f^{-1}(U_1)=f^{-1}(U_2)$, lo que termina de mostrar que $a)$ implica $b)$.
Recíprocamente, la condición $b)$ implica que las aplicaciones $U\mapsto U\cap f(X_0)$ de $\mathfrak{O}(X)$ en $\mathfrak{O}(f(X_0))$ y $V\mapsto f^{-1}(V)$ de $\mathfrak{O}(f(X_0))$ en $\mathfrak{O}(X_0)$ son biyectivas, y por tanto también lo es su compuesta $U\mapsto f^{-1}(U)$.

Para ver que $a)$ implica $c)$, basta aplicar la definición de $f^*(\sh{F})$ \sref[0]{0.3.7.1} y los axiomas de los haces: $a)$ significa que, para todo abierto $U$ de $X$, la aplicación canónica $\Gamma(U,\sh{F})\to\Gamma(f^{-1}(U),f^*(\sh{F}))$ es una biyección funtorial en $\sh{F}$, de donde $c)$.
Queda por mostrar que $c)$ implica $b)$.

\oldpage[IV-3]{98}
Supongamos primero que $f(X_0)$ no es muy denso en $X$; existen entonces dos cerrados distintos $Y\supset Y'$ de $X$ tales que $Y\cap f(X_0)=Y'\cap f(X_0)$.
Sea $\sh{F}$ (resp.\ $\sh{F}'$) el haz de grupos abelianos sobre $X$ que es la imagen directa por la inyección canónica $Y\to X$ (resp.\ $Y'\to X$) del haz simple asociado al prehaz constante $\mathbf{Z}$ sobre $Y$ (resp.\ $Y'$); la definición de $f^*$ \sref[0]{0.3.7.1} muestra que $f^*(\sh{F})$ es isomorfo a $f^*(\sh{F}')$, pero $\sh{F}$ no es isomorfo a $\sh{F}'$, y por tanto no se satisface la condición $c)$.
(Obsérvese que el funtor $f^*$ ni siquiera es entonces \emph{fiel}, pues si $u$ y $v$ son el automorfismo identidad y el endomorfismo nulo de $\sh{F}/\sh{F}'$, $f^*(u)$ y $f^*(v)$ son iguales.)

Mostremos ahora que $c)$ implica que la topología de $X_0$ es la imagen inversa por $f$ de la de $X$.
Observemos que, si la condición $c)$ se satisface para la categoría de haces de conjuntos, también se satisface para la categoría de haces de grupos abelianos, pues se deduce inmediatamente de las definiciones \sref[$0_{\mathrm{III}}$]{0.8.2.5} que la categoría de los \emph{objetos en grupos abelianos} de la categoría de haces de conjuntos no es otra que la categoría de haces de grupos abelianos.
Bastará pues probar nuestra afirmación suponiendo satisfecha $c)$ para las categorías de haces de grupos abelianos sobre $X$ y $X_0$.
Ahora bien, si $\mathbf{Z}_X$ denota el haz simple sobre $X$ asociado al prehaz constante $\mathbf{Z}$, existe un isomorfismo canónico $\Gamma(X,\sh{F})\xrightarrow{\sim}\operatorname{Hom}(\mathbf{Z}_X,\sh{F})$ funtorial en $\sh{F}$ (por el mismo razonamiento que en \sref[0]{0.5.1.1}); como por definición \sref[0]{0.3.7.1} está claro que $f^*(\mathbf{Z}_X)=\mathbf{Z}_{X_0}$, la hipótesis de que $f^*$ es una equivalencia significa que el homomorfismo canónico $\Gamma(X,\sh{F})\to\Gamma(X_0,f^*(\sh{F}))$ es \emph{biyectivo} para todo haz de grupos abelianos $\sh{F}$ sobre $X$.
Sea ahora $Y_0$ un cerrado de $X_0$; sea $\sh{F}_0$ el haz sobre $X_0$ que es la imagen directa por la inyección canónica $Y_0\to X_0$ del haz simple asociado al prehaz constante $\mathbf{Z}$ sobre $Y_0$.
Como $f^*$ es una equivalencia, $\sh{F}_0$ es isomorfo a un haz de la forma $f^*(\sh{F})$, donde $\sh{F}$ es un haz de grupos abelianos sobre $X$.
Consideremos entonces una sección $s_0$ de $\sh{F}_0$ sobre $X_0$ tal que $(s_0)_{x_0}=1_{x_0}$ si $x_0\in Y_0$, y $(s_0)_{x_0}=0_{x_0}$ si $x_0\notin Y_0$.
Las observaciones precedentes muestran que existe una sección $s$ de $\sh{F}$ sobre $X$ tal que $(s_0)_{x_0}=s_{f(x_0)}$ para todo $x_0\in X_0$ (estando identificadas canónicamente las fibras $(\sh{F}_0)_{x_0}$ y $\sh{F}_{f(x_0)}$ \sref[0]{0.3.7.1}); esto significa que $Y_0$ es la imagen inversa por $f$ del conjunto $Y$ de los $x\in X$ tales que $s_x\neq 0_x$, y $Y$ es cerrado en $X$.
\end{proof}

\begin{definition}[10.2.2]
\label{IV.10.2.2}
Cuando una aplicación continua $f:X_0\to X$ satisface las condiciones equivalentes de \sref{IV.10.2.1}, decimos que $f$ es un \emph{cuasihomeomorfismo} de $X_0$ en $X$.
\end{definition}

En virtud de \sref{IV.10.2.1},~$b)$, decir que un subconjunto $X_0$ de un espacio topológico $X$ es muy denso en $X$ significa que la inyección canónica $X_0\to X$ es un cuasihomeomorfismo.

\begin{corollary}[10.2.3]
\label{IV.10.2.3}
La compuesta de dos cuasihomeomorfismos es un cuasihomeomorfismo.
\end{corollary}

Esto se deduce inmediatamente de \sref{IV.10.2.1},~$a)$.

\begin{corollary}[10.2.4]
\label{IV.10.2.4}
Sea $f:X\to Y$ un cuasihomeomorfismo, $Y'$ un subconjunto localmente cuasiconstruible de $Y$ y $X'=f^{-1}(Y')$; entonces la restricción $f'=f|X':X'\to Y'$ es un cuasihomeomorfismo.
\end{corollary}

\begin{proof}
Está claro, en virtud de \sref{IV.10.2.1},~$b)$, que la topología inducida sobre $X'$ por la de $X$ es la imagen inversa por $f'$ de la topología inducida sobre $Y'$ por la de $Y$.
Por otra parte,
\oldpage[IV-3]{99}
sea $Z\neq\varnothing$ un subconjunto cerrado \emph{en} $Y'$, y sea $z\in Z$; existe un entorno abierto $U$ de $z$ en $Y$ tal que $U\cap Y'$ es la unión de un número finito de cerrados $Y'_j$ de $U$; si $j$ es un índice tal que $z\in Y'_j$, $Z\cap Y'_j$ es por tanto cerrado en $U$.
Como $f(X)\cap U$ es muy denso en $U$ \sref{IV.10.1.4}, $Z\cap Y'_j\cap f(X)$ no es vacío \sref{IV.10.1.2},~$a)$, ni, \emph{a fortiori}, $Z\cap f(X)$; pero, como $Z\subset Y'$, se tiene $Z\cap f(X)=Z\cap f'(X')$; el criterio \sref{IV.10.1.2},~$a)$ muestra que $f'(X')$ es muy denso en $Y'$, y concluimos por \sref{IV.10.2.1},~$b)$.
\end{proof}

\begin{corollary}[10.2.5]
\label{IV.10.2.5}
Sea $f:X\to Y$ una aplicación continua y $(V_\alpha)$ un recubrimiento abierto de $Y$.
Si, para todo $\alpha$, la restricción $f^{-1}(V_\alpha)\to V_\alpha$ de $f$ es un cuasihomeomorfismo, entonces $f$ es un cuasihomeomorfismo.
\end{corollary}

Esto se deduce inmediatamente del criterio \sref{IV.10.2.1},~$b)$ y de \sref{IV.10.1.4}.

\begin{corollary}[10.2.6]
\label{IV.10.2.6}
Sea $f:X\to Y$ un cuasihomeomorfismo, $Y'$ un subconjunto localmente cuasiconstruible de $Y$ y $X'=f^{-1}(Y')$.
Para que $Y'$ sea cuasicompacto (resp.\ noetheriano, resp.\ retrocompacto en $Y$), es necesario y suficiente que $X'$ sea cuasicompacto (resp.\ noetheriano, resp.\ retrocompacto en $X$).
\end{corollary}

\begin{proof}
Probemos las dos primeras afirmaciones; en virtud de \sref{IV.10.2.4}, podemos limitarnos al caso en que $Y'=Y$.
Decir que $X$ es cuasicompacto (resp.\ noetheriano) significa que, para toda familia filtrante creciente $(U_\alpha)$ de $\mathfrak{O}(X)$ que tenga a $X$ como elemento máximo (resp.\ para toda familia filtrante creciente $(U_\alpha)$ de $\mathfrak{O}(X)$), existe $\gamma$ tal que $U_\alpha=U_\gamma$ para $\alpha\geq\gamma$.
Como $U\mapsto f^{-1}(U)$ es una biyección de $\mathfrak{O}(Y)$ sobre $\mathfrak{O}(X)$, nuestra afirmación se deduce inmediatamente de la observación precedente.

Los abiertos cuasicompactos de $X$ son por tanto los conjuntos de la forma $f^{-1}(U)$, donde $U$ es un abierto cuasicompacto de $Y$, en virtud de \sref{IV.10.2.1},~$a)$ y de lo que precede.
Para que $X'$ sea retrocompacto en $X$, es necesario y suficiente que, para todo abierto cuasicompacto $U$ de $Y$, $f^{-1}(U)\cap X'=f^{-1}(U\cap Y')$ sea cuasicompacto \sref[$0_{\mathrm{III}}$]{0.9.1.1}; la primera parte de la demostración muestra que esto equivale a decir que $U\cap Y'$ es cuasicompacto para todo abierto cuasicompacto $U$, es decir, que $Y'$ es retrocompacto en $Y$.
\end{proof}

\begin{proposition}[10.2.7]
\label{IV.10.2.7}
Sea $f:X\to Y$ un cuasihomeomorfismo.
Entonces la aplicación $Z\mapsto f^{-1}(Z)$ de $\mathfrak{P}(Y)$ en $\mathfrak{P}(X)$ define por restricción las biyecciones (cf.\ \sref{IV.10.1.1} para la notación):
\[
  \mathfrak{O}(Y)\xrightarrow{\sim}\mathfrak{O}(X)
\]
\[
  \mathfrak{F}(Y)\xrightarrow{\sim}\mathfrak{F}(X)
\]
\[
  \mathfrak{Lf}(Y)\xrightarrow{\sim}\mathfrak{Lf}(X)
\]
\[
  \mathfrak{Qc}(Y)\xrightarrow{\sim}\mathfrak{Qc}(X)
\]
\[
  \mathfrak{Lqc}(Y)\xrightarrow{\sim}\mathfrak{Lqc}(X)
\]
\[
  \mathfrak{C}(Y)\xrightarrow{\sim}\mathfrak{C}(X)
\]
\[
  \mathfrak{Lc}(Y)\xrightarrow{\sim}\mathfrak{Lc}(X).
\]
\end{proposition}

\begin{proof}
\oldpage[IV-3]{100}
Para las dos primeras, esto no es sino la definición \sref{IV.10.2.2}; como la topología de $X$ es la imagen inversa por $f$ de la de $f(X)$, las cinco aplicaciones restantes, en las que se sustituye $Y$ por $f(X)$, son biyectivas.
Basta pues (por \sref{IV.10.2.1},~$b)$) limitarnos al caso en que $X$ es un subespacio muy denso de $Y$, y ya se ha probado \sref{IV.10.1.2} que las aplicaciones $\mathfrak{Lf}(Y)\to\mathfrak{Lf}(X)$, $\mathfrak{Qc}(Y)\to\mathfrak{Qc}(X)$, $\mathfrak{Lqc}(Y)\to\mathfrak{Lqc}(X)$ son biyectivas.
Como todo subconjunto localmente construible es localmente cuasiconstruible, las aplicaciones $\mathfrak{C}(Y)\to\mathfrak{C}(X)$ y $\mathfrak{Lc}(Y)\to\mathfrak{Lc}(X)$ son inyectivas; además, para todo abierto $U\subset Y$, la restricción $f^{-1}(U)\to U$ de $f$ es un cuasihomeomorfismo \sref{IV.10.2.4}, por lo que, si mostramos que $\mathfrak{C}(Y)\to\mathfrak{C}(X)$ es sobreyectiva, lo mismo será cierto para $\mathfrak{Lc}(Y)\to\mathfrak{Lc}(X)$.
Pero, en virtud de \sref{IV.10.2.6}, todo abierto retrocompacto $Z$ de $X$ es de la forma $f^{-1}(T)$, donde $T$ es un abierto retrocompacto de $Y$; esto prueba evidentemente la sobreyectividad de $\mathfrak{C}(Y)\to\mathfrak{C}(X)$.
\end{proof}

\begin{remarks}[10.2.8]
\label{IV.10.2.8}
(i) Vimos en la demostración de \sref{IV.10.2.1} que, si $f$ es un cuasihomeomorfismo, la aplicación canónica
\begin{equation}
\label{IV.10.2.8.1}
\tag{10.2.8.1}
\Gamma(Y,\sh{F})\to\Gamma(X,f^*(\sh{F}))
\end{equation}
es un isomorfismo funtorial en $\sh{F}$ en la categoría de haces de grupos abelianos sobre $Y$.
Como $f^*$ es exacto en esta categoría \sref[0]{0.3.7.2}, esto implica que los homomorfismos canónicos $(T,3.2.2)$
\[
  \mathrm{H}^i(Y,\sh{F})\xrightarrow{\sim}\mathrm{H}^i(X,f^*(\sh{F}))
\]
son \emph{biyectivos} para todo $i$.

(ii) Si $f:X\to Y$ es una aplicación continua y $\sh{F}$, $\sh{G}$ son dos haces de grupos abelianos sobre $Y$, se tiene $f^*(\sh{F}\otimes_{\mathbf{Z}_Y}\sh{G})=f^*(\sh{F})\otimes_{\mathbf{Z}_X}f^*(\sh{G})$, salvo isomorfismo canónico \sref[0]{0.4.3.3}.
Concluimos que, si $f$ es un cuasihomeomorfismo, $f^*$ también es una \emph{equivalencia} entre la categoría de haces de anillos sobre $Y$ y la categoría de haces de anillos sobre $X$; el dato de una estructura de espacio anillado sobre $Y$ equivale por tanto al de una estructura de espacio anillado sobre $X$.

Dados dos espacios anillados $(X,\sh{O}_X)$, $(Y,\sh{O}_Y)$, diremos que un morfismo de espacios anillados $f=(\psi,\theta):(X,\sh{O}_X)\to(Y,\sh{O}_Y)$ es un \emph{cuasiisomorfismo} si $\psi$ es un cuasihomeomorfismo de $X$ en $Y$ y $\theta^\sharp:\psi^*(\sh{O}_Y)\to\sh{O}_X$ es un isomorfismo de haces de anillos; cuando esto ocurre, el espacio anillado $(X,\sh{O}_X)$ queda enteramente determinado, salvo isomorfismo, por $(Y,\sh{O}_Y)$, el espacio $X$ y el cuasihomeomorfismo $\psi$ (que puede elegirse arbitrariamente).
Cuando $f$ es un cuasiisomorfismo de espacios anillados, el funtor
\[
  \sh{F}\mapsto f^*(\sh{F})
\]
es una \emph{equivalencia} entre la categoría de $\sh{O}_Y$-Módulos y la de $\sh{O}_X$-Módulos, pues aquí $f^*(\sh{F})$ se identifica canónicamente con $\psi^*(\sh{F})$.
Obtenemos, por ejemplo, isomorfismos de bi-$\partial$-funtores
\[
  \operatorname{Ext}^i_{\sh{O}_Y}(\sh{F},\sh{G})\xrightarrow{\sim}\operatorname{Ext}^i_{\sh{O}_X}(f^*(\sh{F}),f^*(\sh{G})).
\]

En general, puede decirse que las construcciones usuales de la teoría de haces y del álgebra homológica, efectuadas sobre el espacio anillado $Y$ o sobre el espacio anillado $X$, son equivalentes.
\end{remarks}

\subsection{Espacios de Jacobson}
\label{subsection:IV.10.3}

\oldpage[IV-3]{101}

\begin{definition}[10.3.1]
\label{IV.10.3.1}
Decimos que un espacio topológico $X$ es un \emph{espacio de Jacobson} si el conjunto $X_0$ de los puntos cerrados de $X$ es muy denso en $X$ (dicho de otro modo, si la inyección canónica $X_0\to X$ es un cuasihomeomorfismo).
\end{definition}

Esto significa por tanto \sref{IV.10.1.2} que todo subconjunto localmente cerrado (o solamente localmente cuasiconstruible) $Z\neq\varnothing$ de $X$ contiene un punto cerrado de $X$, o equivalentemente que \emph{todo subconjunto cerrado de $X$ es la adherencia del conjunto de sus puntos cerrados}.

\begin{proposition}[10.3.2]
\label{IV.10.3.2}
Sea $X$ un espacio de Jacobson y $Z$ un subconjunto localmente cuasiconstruible de $X$; entonces el subespacio $Z$ es un espacio de Jacobson, y para que un punto de $Z$ sea cerrado en $Z$ es necesario y suficiente que sea cerrado en $X$.
\end{proposition}

\begin{proof}
Si $X_0$ es el conjunto de los puntos cerrados de $X$, $Z\cap X_0$ es muy denso en $Z$ en virtud de \sref{IV.10.2.4}, aplicado a la inyección $i:X_0\to X$; como el conjunto $Z_0$ de los puntos cerrados de $Z$ contiene evidentemente a $Z\cap X_0$, es \emph{a fortiori} muy denso en $Z$, luego $Z$ es un espacio de Jacobson.
Mostremos ahora que se tiene $Z_0=Z\cap X_0$; sea $x\in Z$ un punto cerrado en $Z$; sea $\overline{\{x\}}$ su adherencia en $X$. Entonces $\overline{\{x\}}\cap Z=\{x\}$ es un subconjunto localmente cuasiconstruible de $X$, y como su intersección con $X_0$ no es vacía \sref{IV.10.1.2}, se tiene $x\in X_0$.
\end{proof}

\begin{proposition}[10.3.3]
\label{IV.10.3.3}
Sea $X$ un espacio topológico y $(U_\alpha)$ un recubrimiento abierto de $X$.
Para que $X$ sea un espacio de Jacobson, es necesario y suficiente que cada uno de los subespacios $U_\alpha$ sea un espacio de Jacobson.
\end{proposition}

\begin{proof}
La condición es necesaria en virtud de \sref{IV.10.3.2}.
Recíprocamente, mostremos primero que la hipótesis de que los $U_\alpha$ son espacios de Jacobson implica que, para que un punto $x\in U_\alpha$ sea cerrado en $X$, basta que sea cerrado en $U_\alpha$.
Basta en efecto ver que esta condición implica que $x$ también es cerrado en cada uno de los $U_\beta$ que lo contienen; pero $U_\alpha\cap U_\beta$ es abierto en $U_\alpha$, luego $x$ es cerrado en $U_\alpha\cap U_\beta$, y en virtud de \sref{IV.10.3.2}, $x$ también es cerrado en $U_\beta$, lo que termina la demostración.
\end{proof}

\subsection{Preesquemas de Jacobson y anillos de Jacobson}
\label{subsection:IV.10.4}

\begin{definition}[10.4.1]
\label{IV.10.4.1}
Decimos que un preesquema $X$ es un \emph{preesquema de Jacobson} si el espacio topológico subyacente a $X$ es un espacio de Jacobson.
Decimos que un anillo $A$ es un \emph{anillo de Jacobson} si $\operatorname{Spec}(A)$ es un preesquema de Jacobson.
\end{definition}

Todo cerrado de $X=\operatorname{Spec}(A)$ es de la forma $Z=V(\mathfrak{a})$, donde $\mathfrak{a}$ es un ideal igual a su radical, y el conjunto $X_0$ de los puntos cerrados de $X$ es el conjunto de los ideales maximales de $A$; decir que $Z\cap X_0$ es denso en $Z$ significa por tanto que $\mathfrak{a}$ es una \emph{intersección de ideales maximales} \sref[I]{I.1.1.4}; como $\mathfrak{a}$ es una intersección de ideales primos, esto equivale a decir que todo ideal primo de $A$ es una intersección de ideales maximales; en virtud de \sref{IV.10.3.1} y \sref{IV.10.1.2}, la definición usual de los anillos de Jacobson (Bourbaki, \emph{Alg.\ comm.}, cap.~V, \S\,3, n$^\circ$\,4, def.~1) coincide así con la definición \sref{IV.10.4.1}.

\begin{proposition}[10.4.2]
\label{IV.10.4.2}
\oldpage[IV-3]{102}
Sea $X$ un preesquema y $(U_\alpha)$ un recubrimiento de $X$ formado por abiertos afines.
Para que $X$ sea un preesquema de Jacobson, es necesario y suficiente que los anillos $\Gamma(U_\alpha,\sh{O}_X)$ de los $U_\alpha$ sean anillos de Jacobson.
\end{proposition}

Esto se deduce de \sref{IV.10.3.3} y de la definición \sref{IV.10.4.1}.

\begin{env}[10.4.3]
\label{IV.10.4.3}
Un preesquema discreto es un preesquema de Jacobson; un anillo artiniano es por tanto un anillo de Jacobson.
Todo anillo de ideales principales que tenga una infinidad de ideales maximales (por ejemplo $\mathbf{Z}$) es un anillo de Jacobson; un anillo local noetheriano $A$ es un anillo de Jacobson si y solo si su ideal maximal es su único ideal primo, es decir, si $A$ es artiniano.
Todo subpreesquema de un preesquema de Jacobson es un preesquema de Jacobson en virtud de \sref{IV.10.3.2}.
\end{env}

\begin{proposition}[10.4.4]
\label{IV.10.4.4}
Sea $B$ un anillo íntegro.
Las condiciones siguientes son equivalentes:
\begin{enumerate}[label=\emph{\alph*)}]
  \item Existe $f\neq 0$ en $B$ tal que $B_f$ es un cuerpo.
  \item El cuerpo de fracciones de $B$ es una $B$-álgebra de tipo finito.
  \item Existe un cuerpo $K$ que contiene a $B$ y que es una $B$-álgebra de tipo finito.
  \item El punto genérico de $\operatorname{Spec}(B)$ es aislado.
\end{enumerate}
\end{proposition}

\begin{proof}
Está claro que $d)$ equivale a $a)$, pues $d)$ significa que existe $f\neq 0$ en $B$ tal que $D(f)$ se reduce al punto genérico de $\operatorname{Spec}(B)$.
Es trivial que $a)$ implica $b)$ y que $b)$ implica $c)$.
Finalmente, $c)$ implica $a)$, en virtud de Bourbaki, \emph{Alg.\ comm.}, cap.~V, \S\,3, n$^\circ$\,1, cor.~2 del th.~1.
\end{proof}

\begin{proposition}[10.4.5]
\label{IV.10.4.5}
Sea $A$ un anillo.
Las condiciones siguientes son equivalentes:
\begin{enumerate}[label=\emph{\alph*)}]
  \item $A$ es un anillo de Jacobson.
  \item Para todo ideal primo no maximal $\mathfrak{p}$ de $A$ y todo $f\neq 0$ en $B=A/\mathfrak{p}$, $B_f$ no es un cuerpo.
  \item[\rm b$'$)] Toda $A$-álgebra de tipo finito $K$ que sea un cuerpo es una $A$-álgebra finita.
\end{enumerate}
\end{proposition}

\begin{proof}
Sabemos que $a)$ implica $b')$ (Bourbaki, \emph{Alg.\ comm.}, cap.~V, \S\,3, n$^\circ$\,4, cor.~3 del th.~3).
Además, el núcleo del homomorfismo $A\to K$ es entonces un ideal \emph{maximal} $\mathfrak{m}$ de $A$, y $K$ es una extensión finita de $A/\mathfrak{m}$ (\emph{loc.\ cit.}).
Es trivial que $b')$ implica $b)$, pues $A/\mathfrak{p}=B$ no es un cuerpo, toda $B$-álgebra de tipo finito es una $A$-álgebra de tipo finito y $B_f$ es una $B$-álgebra de tipo finito.
Queda por mostrar que $b)$ implica $a)$.
Utilizaremos el lema siguiente:
\end{proof}

\begin{lemma}[10.4.5.1]
\label{IV.10.4.5.1}
Sea $X$ un espacio topológico que tenga la propiedad siguiente: para todo subconjunto localmente cerrado $Z\neq\varnothing$ de $X$, existe un subconjunto $Z'$ de $Z$, localmente cerrado en $X$ (o en $Z$, lo que equivale a lo mismo), y un punto $x\in Z'$, cerrado en $Z'$.
Entonces, para que $X$ sea un espacio de Jacobson, es necesario y suficiente que ningún punto no cerrado $x$ de $X$ sea aislado en $\overline{\{x\}}$.
\end{lemma}

\begin{proof}
Si $X$ es un espacio de Jacobson, un punto no cerrado $x$ de $X$ no puede ser aislado en $\overline{\{x\}}$, pues esto significaría que existe un abierto $U$ de $X$ tal que $U\cap\overline{\{x\}}=\{x\}$; pero $U\cap\overline{\{x\}}$ es localmente cerrado en $X$ y no contendría ningún punto cerrado en $X$, lo que contradice la hipótesis de que $X$ es un espacio de Jacobson.
Recíprocamente, supongamos satisfechas las condiciones del enunciado; entonces el conjunto $X_0$ de los puntos cerrados de $X$ coincide con el conjunto de los $x\in X$ que son aislados en $\overline{\{x\}}$; pero se deduce de \sref{IV.5.1.10.1} que
\oldpage[IV-3]{103}
este conjunto es \emph{muy denso} en $X$, luego $X$ es un espacio de Jacobson por definición.
\end{proof}

Recordemos \sref{IV.5.1.10} que la hipótesis formulada sobre $X$ en \sref{IV.10.4.5.1} siempre se satisface cuando el espacio subyacente a $X$ es un \emph{preesquema}.

Volviendo a la demostración de \sref{IV.10.4.5}, la condición $b)$ significa que, para todo punto no cerrado $x$ de $\operatorname{Spec}(A)$, el punto genérico $x$ de $\operatorname{Spec}(A/\mathfrak{j}_x)=V(\mathfrak{j}_x)=\overline{\{x\}}$ no es aislado en $\overline{\{x\}}$; por tanto, $b)$ implica $a)$ por el lema \sref{IV.10.4.5.1} y \sref{IV.10.4.4}.

\begin{corollary}[10.4.6]
\label{IV.10.4.6}
Toda álgebra de tipo finito $B$ sobre un anillo de Jacobson $A$ es un anillo de Jacobson, y la imagen inversa en $A$ de todo ideal maximal de $B$ es un ideal maximal de $A$.
En particular, toda álgebra de tipo finito sobre un cuerpo o sobre $\mathbf{Z}$ es un anillo de Jacobson.
\end{corollary}

\begin{proof}
Una $B$-álgebra de tipo finito $K$ que sea un cuerpo es también una $A$-álgebra de tipo finito, y por tanto un $A$-módulo de tipo finito y \emph{a fortiori} un $B$-módulo de tipo finito, de donde la primera afirmación; la segunda se probó en el curso de la demostración de \sref{IV.10.4.5}, aplicada a $K=B/\mathfrak{m}$.
\end{proof}

\begin{corollary}[10.4.7]
\label{IV.10.4.7}
Si $X$ es un preesquema de Jacobson y $f:Y\to X$ es un morfismo localmente de tipo finito, entonces $Y$ es un preesquema de Jacobson y la imagen por $f$ de todo punto cerrado de $Y$ es un punto cerrado de $X$.
\end{corollary}

\begin{proof}
Como la cuestión es local sobre $X$ y sobre $Y$, nos reducimos al caso en que $X$ e $Y$ son afines, y el corolario se deduce entonces de \sref{IV.10.4.6}.
\end{proof}

\begin{corollary}[10.4.8]
\label{IV.10.4.8}
Si $k$ es un cuerpo algebraicamente cerrado y $X$ un $k$-preesquema localmente de tipo finito, el conjunto de los puntos racionales de $X$ sobre $k$ es muy denso en $X$.
\end{corollary}

\begin{proof}
En efecto, $X$ es un preesquema de Jacobson \sref{IV.10.4.7} y los puntos cerrados son exactamente los puntos racionales sobre $k$ \sref[I]{I.6.4.2}.
\end{proof}

\begin{env}[10.4.9]
\label{IV.10.4.9}
El hecho de que los preesquemas localmente de tipo finito sobre un cuerpo o sobre $\mathbf{Z}$ sean preesquemas de Jacobson es particularmente importante, por la posibilidad de reducirse a este caso en muchas cuestiones de geometría algebraica \sref{IV.8.1.2},~$c)$.
Daremos dos ejemplos:
\end{env}

\begin{env}[10.4.10]
\label{IV.10.4.10}
\textbf{Aplicaciones (10.4.10) :~I~: Demostración de \sref{IV.6.15.9}.}
Sea $k$ un cuerpo separablemente cerrado y $X$ un $k$-preesquema localmente de tipo finito sobre $k$ y \emph{unirramificado}.
Sabemos que el cierre entero de una $k$-álgebra íntegra de tipo finito $A$ en una extensión finita de su cuerpo de fracciones es una $A$-álgebra finita (Bourbaki, \emph{Alg.\ comm.}, cap.~V, \S\,3, n$^\circ$\,2, th.~2), luego toda $k$-álgebra de tipo finito es un anillo universalmente japonés; concluimos que el conjunto de los puntos $x\in X$ en los que $X$ es geométricamente unirramificado es localmente construible \sref{IV.9.7.10}.
Pero la hipótesis y el lema \sref{IV.6.15.8} implican que este conjunto contiene todos los puntos \emph{cerrados} de $X$.
La conclusión se deduce por tanto de \sref{IV.10.4.6}, \sref{IV.10.3.1} y de la biyectividad de la aplicación canónica $\mathfrak{Lc}(X)\to\mathfrak{Lc}(X_0)$ (donde $X_0$ es el conjunto de los puntos cerrados de $X$) \sref{IV.10.2.7}.
\end{env}

\begin{env}[10.4.11]
\label{IV.10.4.11}
\oldpage[IV-3]{104}
\textbf{Aplicaciones :~II~: Proposición (10.4.11).}
Sea $S$ un preesquema y $X$ un $S$-preesquema de tipo finito.
Todo $S$-endomorfismo de $X$ que sea radicial es sobreyectivo (y por tanto biyectivo).
\end{env}

\begin{proof}
Sea $f:X\to S$ el morfismo estructural, $g:X\to X$ el $S$-endomorfismo considerado y, para todo $s\in S$, sea $g_s$ el morfismo deducido de $g$ por el cambio de base $\operatorname{Spec}(k(s))\to S$, que es un $k(s)$-endomorfismo de la fibra $f^{-1}(s)=X\times_S\operatorname{Spec}(k(s))$.

Para probar que $g$ es sobreyectivo, basta probar que $g_s$ es sobreyectivo para todo $s\in S$; podemos por tanto (en virtud de \sref[I]{I.3.5.7}) reducirnos al caso en que $S=\operatorname{Spec}(k)$ es el espectro de un cuerpo $k$, en cuyo caso $f$ es un morfismo de presentación finita, pues $S$ es noetheriano.
Aplicando \sref{IV.8.9.1} y \sref{IV.8.10.5},~(vii), nos reducimos al caso en que $S=\operatorname{Spec}(A)$, donde $A$ es una sub-$\mathbf{Z}$-álgebra de tipo finito de $k$.
Pero entonces $X$ es un preesquema de Jacobson \sref{IV.10.4.7}, y $g(X)$ es construible en $X$ \sref[I]{I.1.8.5}; por tanto, para probar que $g(X)=X$, basta mostrar que $g(X)$ contiene todos los puntos \emph{cerrados} de $X$ \sref{IV.10.3.1}.

\begin{lemma}[10.4.11.1]
\label{IV.10.4.11.1}
Sea $Y$ un $\mathbf{Z}$-preesquema de tipo finito.
\begin{enumerate}[label=(\roman*)]
  \item Para que un punto $y\in Y$ sea cerrado, es necesario y suficiente que $k(y)$ sea un cuerpo finito.
  \item Para todo número primo $p$ y todo entero $d\geq 1$, el conjunto de los puntos $y\in Y$ tales que $k(y)$ es una extensión de $\mathbf{F}_p$ de grado divisor de $d$ es finito.
\end{enumerate}
\end{lemma}

La afirmación (i) se deduce del hecho de que la imagen de un punto cerrado $y\in Y$ en $\operatorname{Spec}(\mathbf{Z})$ es un punto cerrado \sref{IV.10.4.7}, es decir, un número primo, y de \sref[I]{I.6.4.2}.
Por otra parte, como $Y$ es una unión finita de abiertos afines de tipo finito sobre $\mathbf{Z}$, podemos limitarnos a probar (ii) en el caso en que $Y=\operatorname{Spec}(C)$, donde $C$ es una $\mathbf{Z}$-álgebra de tipo finito.
Ahora bien, los puntos $y\in Y$ tales que el grado de $k(y)$ sobre $\mathbf{F}_p$ divide a $d$ son exactamente los núcleos de los homomorfismos $C\to\mathbf{F}_{p^d}$; pero, si $(t_i)$ es un sistema de generadores de la $\mathbf{Z}$-álgebra $C$, todo homomorfismo de $C$ queda determinado por sus valores sobre los elementos $t_i$, y por tanto solo hay un número finito de homomorfismos de $C$ en un cuerpo finito.

Siendo esto así, como $X$ es un $\mathbf{Z}$-preesquema de tipo finito, sea $T_{p,d}$ el conjunto de los puntos cerrados $z\in X$ tales que $k(z)$ es una extensión de $\mathbf{F}_p$ de grado divisor de $d$; se deduce de \sref{IV.10.4.11.1} que el conjunto $T_{p,d}$ es finito y que el conjunto de los puntos cerrados de $X$ es la unión de los $T_{p,d}$.
Además, si $z\in T_{p,d}$ y $h$ es un endomorfismo cualquiera de $X$, $k(h(z))$ es isomorfo a un subcuerpo de $k(z)$, luego $h(z)\in T_{p,d}$; dicho de otro modo, $T_{p,d}$ es \emph{estable} bajo todo endomorfismo de $X$.
Como $g$ es inyectivo por hipótesis, su restricción a $T_{p,d}$ es \emph{biyectiva}, pues $T_{p,d}$ es finito, lo que termina de probar la proposición.

Veremos más adelante \sref{IV.17.9.6} que, si además se supone que $X$ es un $S$-preesquema de presentación finita y, por otra parte, que $g$ es un monomorfismo, puede afirmarse que $g$ es un \emph{automorfismo} de $X$.
\end{proof}

\subsection{Preesquemas de Jacobson noetherianos}
\label{subsection:IV.10.5}

\begin{proposition}[10.5.1]
\label{IV.10.5.1}
Sea $B$ un anillo noetheriano íntegro.
Las condiciones equivalentes $a)$ a $d)$ de \sref{IV.10.4.4} son entonces también equivalentes a las siguientes:
\begin{enumerate}[label=\emph{\alph*)}]
  \setcounter{enumi}{4}
  \item $\operatorname{Spec}(B)$ es finito.
  \item $B$ es un anillo semilocal de dimensión $\leq 1$.
\end{enumerate}
\end{proposition}

\begin{proof}
\oldpage[IV-3]{105}
En efecto, se deduce del teorema de Artin--Tate \sref[0]{0.16.3.3} que las condiciones $a)$ y $f)$ son equivalentes.
La condición $f)$ implica que $\operatorname{Spec}(B)$ es la unión del conjunto de sus puntos cerrados y de su punto genérico, luego $f)$ implica $e)$ sin suponer que $B$ sea noetheriano.
Finalmente, $e)$ implica $d)$ sin suponer que $B$ sea noetheriano, pues el punto genérico $x$ de $X$ es el complementario de la unión de las adherencias $\overline{\{y\}}$, donde $y$ recorre el conjunto de los puntos $y\neq x$, y como estos puntos son finitos en número, $\{x\}$ es el complementario de un cerrado de $X$.
\end{proof}

\begin{corollary}[10.5.2]
\label{IV.10.5.2}
Para que un anillo noetheriano $A$ sea un anillo de Jacobson, es necesario y suficiente que no exista ningún ideal primo $\mathfrak{p}$ de $A$ tal que $A/\mathfrak{p}$ sea un anillo semilocal de dimensión $1$.
\end{corollary}

\begin{proof}
Esto se deduce inmediatamente de \sref{IV.10.5.1} y de la condición $b)$ de \sref{IV.10.4.5}, pues los ideales primos $\mathfrak{p}$ de $A$ tales que $A/\mathfrak{p}$ es semilocal de dimensión~$0$ son los ideales maximales de $A$.
\end{proof}

\begin{corollary}[10.5.3]
\label{IV.10.5.3}
Sea $X$ un preesquema noetheriano irreducible.
Las condiciones siguientes son equivalentes:
\begin{enumerate}[label=\emph{\alph*)}]
  \item El punto genérico de $X$ es aislado.
  \item $X$ es finito.
  \item $X$ es de dimensión $\leq 1$ y el conjunto de sus puntos cerrados es finito.
\end{enumerate}
\end{corollary}

\begin{proof}
Existe por hipótesis un número finito de abiertos afines irreducibles $U_i$ ($1\leq i\leq n$) que recubren $X$ y cada uno de los cuales contiene el punto genérico de $X$; basta probar la equivalencia de $a)$, $b)$ y $c)$ para cada uno de los $(U_i)_{\mathrm{red}}$ (teniendo en cuenta \sref[0]{0.14.1.7}).
Pero esta equivalencia se deduce de \sref{IV.10.5.1}.
\end{proof}

\begin{remark}[10.5.4]
\label{IV.10.5.4}
Un preesquema noetheriano $X$ que satisfaga las condiciones equivalentes de \sref{IV.10.5.3} no es necesariamente un esquema afín; de hecho, puede incluso ser \emph{no separado}, como se ve sustituyendo, en el ejemplo \sref[I]{I.5.5.11} de la ``recta afín con punto duplicado'', $X_1$ y $X_2$ por el espectro del anillo de valoración discreta $(K[s])_{(s)}$, y $U_{12}$ y $U_{21}$ por el abierto reducido al punto genérico en $X_1$ y $X_2$ respectivamente; el preesquema no separado $X$ así obtenido tiene exactamente $3$ puntos.
\end{remark}

\begin{proposition}[10.5.5]
\label{IV.10.5.5}
Sea $X$ un preesquema noetheriano.
\begin{enumerate}[label=(\roman*)]
  \item El conjunto $X_0$ de los puntos $x\in X$ tales que $\overline{\{x\}}$ es finito es muy denso en $X$.
  \item Para que $X$ sea un preesquema de Jacobson, es necesario y suficiente que no exista ningún subpreesquema de $X$ isomorfo al espectro de un anillo semilocal íntegro de dimensión $1$.
\end{enumerate}
\end{proposition}

\begin{proof}
La condición de que $\overline{\{x\}}$ sea finito equivale aquí, por \sref{IV.10.5.3}, a que $x$ sea aislado en $\overline{\{x\}}$, pues este es el espacio subyacente a un subpreesquema (noetheriano) de $X$; la afirmación (i) se deduce por tanto de \sref{IV.5.1.10.1}.
Análogamente, teniendo en cuenta \sref{IV.10.4.5.1}, para mostrar la afirmación (ii), observemos que, para que un punto no cerrado $x$ de $X$ pertenezca a $X_0$, es necesario y suficiente que el subpreesquema cerrado íntegro $Y$ de $X$ cuyo espacio subyacente es $\overline{\{x\}}$ sea de dimensión $1$ y finito, y por tanto una unión finita de subpreesquemas $U_i$ abiertos (en $Y$) y afines que sean espectros de anillos semilocales íntegros de dimensión $1$.
Recíprocamente, si existe un subpreesquema $Z$ de $X$ que sea el espectro de un anillo semilocal íntegro de dimensión $1$, $Z$ no es un preesquema de Jacobson \sref{IV.10.5.2}, y por tanto tampoco lo es $X$ \sref{IV.10.3.2}.
\end{proof}

\oldpage[IV-3]{106}

\begin{remark}[10.5.6]
\label{IV.10.5.6}
La afirmación (ii) de \sref{IV.10.5.5} sigue siendo válida cuando $X$ es \emph{localmente noetheriano}: en efecto, si $(V_\alpha)$ es un recubrimiento de $X$ por abiertos afines (noetherianos), todo subpreesquema de uno de los $V_\alpha$ es un subpreesquema de $X$; recíprocamente, si un subpreesquema $Z$ de $X$ es isomorfo al espectro de un anillo semilocal íntegro de dimensión $1$, existe un $\alpha$ tal que $V_\alpha\cap Z$ contiene un abierto afín $U$ de $Z$ que no se reduce al punto genérico de $Z$, y que por tanto también es el espectro de un anillo semilocal íntegro de dimensión $1$.
Concluimos por \sref{IV.10.3.3}.
\end{remark}

\begin{proposition}[10.5.7]
\label{IV.10.5.7}
Sea $X$ un preesquema localmente noetheriano y $Y$ un cerrado de $X$ tal que todo cerrado no vacío de $X$ corta a $Y$.
Entonces el subpreesquema inducido sobre el abierto $X-Y$ es un preesquema de Jacobson.
\end{proposition}

\begin{proof}
Apliquemos el criterio \sref{IV.10.5.5},~(ii), y supongamos que existe en $X-Y$ un subpreesquema $Z$ que es el espectro de un anillo semilocal íntegro de dimensión $1$. El punto genérico $z$ de $Z$ es entonces aislado en $Z$ (o en la adherencia $\overline{Z}$ de $Z$ en $X$), y $Z$ es distinto de $\{z\}$.
Sea $y\neq z$ un punto de $Z$; como no pertenece a $Y$, no es un punto cerrado de $X$, y su adherencia $\overline{\{y\}}$ en $X$ corta a $Y$ en un punto $x\neq y$, que es por tanto una especialización de $y$.
La existencia de la cadena $\{x\}\subset\overline{\{y\}}\subset\overline{\{z\}}$ muestra entonces que la dimensión de $\overline{\{z\}}$ sería $\geq 2$, y lo mismo ocurriría con la dimensión de $U\cap\overline{\{z\}}$, donde $U$ es un entorno abierto afín (y por tanto noetheriano) de $x$ en $X$.
Pero esto contradice que el punto genérico $z$ de $U\cap\overline{\{z\}}$ sea aislado en $U\cap\overline{\{z\}}$ \sref{IV.10.5.3}.
\end{proof}

\begin{corollary}[10.5.8]
\label{IV.10.5.8}
Sea $A$ un anillo noetheriano; para todo elemento $f$ del radical $\mathfrak{R}$ de $A$, el anillo $A_f$ es un anillo de Jacobson, y el abierto $\operatorname{Spec}(A)-V(\mathfrak{R})$ es un esquema de Jacobson.
\end{corollary}

\begin{proof}
Si $X=\operatorname{Spec}(A)$ y $Y=V(\mathfrak{J})$, donde $\mathfrak{J}$ es un ideal de $A$, decir que todo cerrado no vacío de $X$ corta a $Y$ equivale \sref[0]{0.2.1.3} a decir que $Y$ contiene todos los \emph{puntos cerrados} de $X$, o equivalentemente que $\mathfrak{J}$ está \emph{contenido en el radical} $\mathfrak{R}$ de $A$.
Si $f\in\mathfrak{R}$, el abierto $D(f)$ no corta a $V(\mathfrak{R})$, y es un espacio de Jacobson en virtud de \sref{IV.10.5.7}.
\end{proof}

\begin{corollary}[10.5.9]
\label{IV.10.5.9}
Sea $A$ un anillo local noetheriano, $\mathfrak{m}$ su ideal maximal, $X=\operatorname{Spec}(A)$ y $Y=X-\{\mathfrak{m}\}$; entonces $Y$ es un esquema de Jacobson, cuyos puntos cerrados son los ideales primos $\mathfrak{p}\in\operatorname{Spec}(A)$ tales que $\dim(A/\mathfrak{p})=1$.
\end{corollary}

\begin{proof}
La primera afirmación es un caso particular de \sref{IV.10.5.8}; por otra parte, los puntos cerrados de $Y$ son los ideales primos $\mathfrak{p}$ de $A$ que son elementos maximales del conjunto de los ideales primos $\neq\mathfrak{m}$, lo que, por definición de la dimensión, significa que $\dim(A/\mathfrak{p})=1$.
\end{proof}

\begin{proposition}[10.5.10]
\label{IV.10.5.10}
Sea $A$ un anillo local noetheriano, reducido, completo y que no sea un cuerpo.
Entonces las intersecciones finitas de los núcleos de los homomorfismos locales de $A$ en anillos de valoración discreta $V$, que hacen de $V$ una $A$-álgebra finita, forman una base de filtro convergente a $0$ para la topología ádica de $A$.
\end{proposition}

\begin{proof}
Basta (Bourbaki, \emph{Alg.\ comm.}, cap.~III, \S\,2, n$^\circ$\,7, prop.~8) probar que la intersección de los núcleos considerados en el enunciado se reduce a $0$.
Supongamos primero que $A$ es íntegro y de dimensión $1$; en virtud del teorema de Nagata \sref[0]{0.23.1.5} y \sref[0]{0.23.1.6}, el cierre entero $A'$ de $A$ es un anillo local íntegro completo,
\oldpage[IV-3]{107}
íntegramente cerrado y de dimensión $1$, y una $A$-álgebra finita; es por tanto un anillo de valoración discreta, y la proposición se deduce inmediatamente en este caso.

Pasemos al caso general; sea $x$ el punto cerrado de $X=\operatorname{Spec}(A)$ y pongamos $Y=X-\{x\}$; sabemos \sref{IV.10.5.9} que $Y$ es un preesquema de Jacobson.
Mostremos que la intersección de los ideales primos $\mathfrak{p}$ de $A$ tales que $\dim(A/\mathfrak{p})=1$ se reduce a $0$; la proposición se seguirá entonces, pues, para cada uno de estos ideales $\mathfrak{p}$, la intersección de los núcleos de los homomorfismos locales de $A/\mathfrak{p}$ en anillos de valoración discreta $V$ que hacen de $V$ una $(A/\mathfrak{p})$-álgebra finita se reduce a $0$.
Pero decir que la intersección de estos ideales primos se reduce a $0$ significa que el conjunto de dichos ideales es denso en $X$, o equivalentemente en $Y$ (puesto que $Y$ es denso en $X$), y esto se deduce inmediatamente de \sref{IV.10.5.9}.
\end{proof}

\subsection{Dimensión en los preesquemas de Jacobson}
\label{subsection:IV.10.6}

Los resultados de esta sección precisan, en ciertos casos, y generalizan resultados del \S\,5.

\begin{proposition}[10.6.1]
\label{IV.10.6.1}
Sea $S$ un preesquema localmente noetheriano que satisfaga además las condiciones: $1^\circ$~$S$ es un preesquema de Jacobson; $2^\circ$~para todo $s\in S$, $\sh{O}_{S,s}$ es universalmente catenario \sref{IV.5.6.2}; $3^\circ$~toda componente irreducible $S'$ de $S$ es equicodimensional (dicho de otro modo, para todo punto cerrado $s$ de $S'$ y todo subpreesquema de $S$ cuyo espacio subyacente sea $S'$, se tiene $\dim(S')=\dim(\sh{O}_{S',s})$).
Se tienen entonces las propiedades siguientes:
\begin{enumerate}[label=(\roman*)]
  \item Para todo morfismo localmente de tipo finito $g:X\to S$, $X$ satisface las condiciones $1^\circ$, $2^\circ$ y $3^\circ$ anteriores.
En particular, si $X$ es equidimensional (por ejemplo, si $X$ es irreducible), $X$ es biequidimensional (dicho de otro modo, $X$ es catenario y para todo punto cerrado $x$ de $X$ se tiene $\dim(\sh{O}_{X,x})=\dim(X)$ \sref[0]{0.14.3.3}).
  \item Sean $X$, $Y$ dos $S$-preesquemas localmente de tipo finito sobre $S$ y $f:X\to Y$ un $S$-morfismo; supongamos que $X$ es irreducible y que $f$ es dominante.
Si $\xi$ (resp.\ $\eta$) es el punto genérico de $X$ (resp.\ $Y$) y $e=\dim(f^{-1}(\eta))=\deg.\mathrm{tr}_{k(\eta)}k(\xi)$, se tiene
\begin{equation}
\label{IV.10.6.1.1}
\tag{10.6.1.1}
\dim(X)=\dim(Y)+e.
\end{equation}
  \item Sean $X$, $Y$ dos $S$-preesquemas localmente de tipo finito sobre $S$, $f:X\to Y$ un $S$-morfismo y $n$ un entero no negativo tal que $\dim(f^{-1}(y))\geq n$ (resp.\ $\dim(f^{-1}(y))\leq n$) para todo $y\in Y$.
Entonces se tiene
\begin{equation}
\label{IV.10.6.1.2}
\tag{10.6.1.2}
\dim(X)\geq\dim(Y)+n
\end{equation}
(resp.
\begin{equation}
\label{IV.10.6.1.3}
\tag{10.6.1.3}
\dim(X)\leq\dim(Y)+n).
\end{equation}
\end{enumerate}
\end{proposition}

\begin{proof}
(i) La propiedad $1^\circ$ para $X$ se deduce de \sref{IV.10.4.7}.
Para todo $x\in X$, $\sh{O}_{X,x}$ es el anillo local, en un ideal primo, de una $\sh{O}_{S,g(x)}$-álgebra de tipo finito, y el homomorfismo $\sh{O}_{S,g(x)}\to\sh{O}_{X,x}$
\oldpage[IV-3]{108}
es local; por tanto \sref{IV.5.6.3},~(iv), $\sh{O}_{X,x}$ es universalmente catenario.
Para mostrar que $X$ satisface la condición $3^\circ$, consideremos varios casos:

\emph{a)} $X$ es un subpreesquema cerrado irreducible de $S$; sea $S'$ una componente irreducible de $S$ que contiene a $X$, $\xi$ el punto genérico de $X$ y $x$ un punto cerrado de $X$; para todo $s\in S'$, $\sh{O}_{S',s}$, cociente de $\sh{O}_{S,s}$, es catenario \sref{IV.5.6.1}, luego las condiciones $2^\circ$ y $3^\circ$ implican que $S'$ es biequidimensional \sref[0]{0.14.3.3}.
En virtud de \sref{IV.5.1.2} y de \sref[0]{0.14.3.3.2}, se tiene por tanto
\[
  \dim(\sh{O}_{X,x})=\dim(\sh{O}_{S',x})-\dim(\sh{O}_{S',\xi})=\dim(S')-\dim(\sh{O}_{S',\xi})
\]
lo que prueba que $\dim(\sh{O}_{X,x})$ no depende del punto cerrado $x$ considerado, de donde la afirmación en este caso \sref{IV.5.1.4}.

\emph{b)} $X$ es irreducible y $g$ dominante.
Entonces, para todo punto cerrado $x$ de $X$, $g(x)$ es cerrado en $S$ por \sref{IV.10.4.7}; como $\sh{O}_{S,g(x)}$ es universalmente catenario, se deduce de \sref{IV.5.6.5.3} que se tiene
\[
  \dim(\sh{O}_{X,x})=\dim(\sh{O}_{S,g(x)})+e
\]
donde $e=\dim(g^{-1}(\zeta))$, siendo $\zeta$ el punto genérico de $S$.
Como $\dim(S)=\dim(\sh{O}_{S,g(x)})$ en virtud de la condición $3^\circ$ para $S$, se tiene $\dim(\sh{O}_{X,x})=\dim(S)+e$ para todo punto cerrado $x\in X$; esto prueba la condición $3^\circ$ para $X$ \sref{IV.5.1.4}, y al mismo tiempo la fórmula
\begin{equation}
\label{IV.10.6.1.4}
\tag{10.6.1.4}
\dim(X)=\dim(S)+e.
\end{equation}

\emph{c) Caso general.}
Considerando un subpreesquema reducido $X'$ de $X$ cuyo espacio subyacente sea una componente irreducible de $X$, nos reducimos al caso en que $X$ es íntegro; utilizando \sref[I]{I.5.2.2} y el caso $a)$ probado anteriormente, podemos sustituir entonces $S$ por el subpreesquema reducido cuyo espacio subyacente es $\overline{g(X)}$; quedamos así reducidos al caso $b)$, lo que termina de probar (i).

(ii) Como el morfismo $f$ es localmente de tipo finito \sref[I]{I.1.3.4}, podemos aplicar los resultados de (i) sustituyendo $S$ por $Y$; además, como $X$ es irreducible y $f$ dominante, también podemos sustituir $S$ por $Y$ en \sref{IV.10.6.1.4}, lo que da \sref{IV.10.6.1.1}.

(iii) La afirmación relativa al caso en que $\dim(f^{-1}(y))\leq n$ para todo $y\in Y$ ya se ha probado bajo hipótesis más generales en \sref{IV.5.6.7}.
Supongamos que $\dim(f^{-1}(y))\geq n$ para todo $y\in Y$, y consideremos un punto genérico $\eta$ de una componente irreducible $Y'$ de $Y$; existe al menos una componente irreducible $Z$ de $f^{-1}(\eta)$ de dimensión $\geq n$; si $\xi$ es el punto genérico de $Z$, $\xi$ también es un punto genérico de una componente irreducible $X'$ de $X$ tal que $f(X')$ es denso en $Y'$ \sref[0]{0.2.1.8}.
Consideremos los subpreesquemas reducidos de $X$, $Y$ que tienen respectivamente por espacios subyacentes $X'$, $Y'$, y la restricción $X'\to Y'$ de $f$ \sref[I]{I.5.2.2}; se deduce entonces de (ii) que $\dim(X')\geq\dim(Y')+n$, y \emph{a fortiori} $\dim(X)\geq\dim(Y')+n$; como esto es cierto para toda componente irreducible $Y'$ de $Y$, concluimos la desigualdad \sref{IV.10.6.1.2}.
\end{proof}

\begin{corollary}[10.6.2]
\label{IV.10.6.2}
Supongamos que $X$ satisface las condiciones $1^\circ$, $2^\circ$ y $3^\circ$ de \sref{IV.10.6.1}.
Entonces, para todo abierto denso $U$ de $X$, se tiene $\dim(U)=\dim(X)$.
\end{corollary}

\begin{proof}
\oldpage[IV-3]{109}
Sabemos en efecto que $\dim(U)\leq\dim(X)$ \sref[0]{0.14.1.4}; además, como $X$ es un preesquema de Jacobson, $U$ contiene un punto cerrado de cada componente irreducible de $X$, luego $\dim(U)=\dim(X)$ en virtud de \sref{IV.10.6.1} y \sref[0]{0.14.1.2.1}.
\end{proof}

\begin{proposition}[10.6.3]
\label{IV.10.6.3}
Supongamos que $X$ satisface las condiciones $1^\circ$, $2^\circ$ y $3^\circ$ de \sref{IV.10.6.1}, y sea $Y$ un cerrado de $X$.
Para todo $x\in Y$ y todo entorno abierto $U$ de $x$ en $X$ que no corte las componentes irreducibles de $Y$ que no contienen a $x$, se tiene
\begin{equation}
\label{IV.10.6.3.1}
\tag{10.6.3.1}
\dim_x(Y)=\dim(U\cap Y)=\dim\overline{\{x\}}+\operatorname{codim}(\overline{\{x\}},Y)=\sup_i\dim(Y_i)
\end{equation}
donde los $Y_i$ ($1\leq i\leq m$) son las componentes irreducibles de $Y$ que contienen a $x$.
\end{proposition}

\begin{proof}
Considerando un subpreesquema cerrado de $X$ cuyo espacio subyacente sea $Y$, podemos limitarnos, en virtud de \sref{IV.10.6.1},~(i), al caso en que $Y=X$.
Por la elección de $U$, $U$ es la unión de los $U\cap X_i$, luego $\dim(U)=\sup_i\dim(U\cap X_i)$; pero, por \sref{IV.10.6.2} aplicado a un subpreesquema cerrado de $X$ cuyo espacio subyacente sea $X_i$, se tiene (teniendo en cuenta \sref{IV.10.6.1},~(i)) $\dim(U\cap X_i)=\dim(X_i)$; esto prueba que el segundo y el cuarto miembros de \sref{IV.10.6.3.1} son iguales.
Como $U\cap X_i$ es biequidimensional \sref{IV.10.6.1},~(i) (pues la inmersión $U\cap X_i\to X_i$ es de tipo finito \sref[I]{I.6.3.5}), se tiene

\begin{equation}
\label{IV.10.6.3.2}
\tag{10.6.3.2}
\dim(X_i)=\dim(U\cap X_i)=\dim(U\cap\overline{\{x\}})+\operatorname{codim}(U\cap\overline{\{x\}},U\cap X_i)
\end{equation}
\sref[0]{0.14.3.5}.
Por \sref{IV.10.6.2}, aplicado a un subpreesquema cerrado de $X$ cuyo espacio subyacente sea $\overline{\{x\}}$, $\dim(U\cap\overline{\{x\}})=\dim(\overline{\{x\}})$; como, por las mismas razones, también se tiene

\begin{equation}
\label{IV.10.6.3.3}
\tag{10.6.3.3}
\dim(X_i)=\dim(\overline{\{x\}})+\operatorname{codim}(\overline{\{x\}},X_i)
\end{equation}
se obtiene
\[
  \dim(U)=\dim(\overline{\{x\}})+\sup_i\operatorname{codim}(\overline{\{x\}},X_i)=\dim(\overline{\{x\}})+\operatorname{codim}(\overline{\{x\}},X)
\]
por definición de la codimensión \sref[0]{0.14.2.1}.
Esto muestra que $\dim(U)$ es \emph{independiente} del entorno abierto $U$ de $x$ que satisface las condiciones del enunciado, y por tanto es igual a $\dim_x(X)$, por \sref[0]{0.14.1.4.1}.
\end{proof}

\begin{corollary}[10.6.4]
\label{IV.10.6.4}
Bajo las hipótesis de \sref{IV.10.6.3}, sea $\sh{F}$ un $\sh{O}_X$-módulo coherente y $Y$ su soporte.
Para todo $x\in Y$, se tiene
\begin{equation}
\label{IV.10.6.4.1}
\tag{10.6.4.1}
\dim(\overline{\{x\}})+\dim(\sh{F}_x)=\dim_x Y.
\end{equation}
\end{corollary}

\begin{proof}
En efecto, esto se deduce de \sref{IV.10.6.3.1} y de la fórmula $\dim(\sh{F}_x)=\operatorname{codim}(\overline{\{x\}},Y)$ \sref{IV.5.1.12.2}.
\end{proof}

\subsection{Ejemplos y contraejemplos}
\label{subsection:IV.10.7}

\begin{env}[10.7.1]
\label{IV.10.7.1}
Sea $S$ un preesquema localmente noetheriano de dimensión $\leq 1$ y supongamos que $S$ es un \emph{preesquema de Jacobson}; cuando $S$ es noetheriano, esto equivale a decir que las componentes irreducibles de $S$ de dimensión $1$ son infinitas, pues todo $x\in S$ que no sea un punto cerrado es el punto genérico de una componente tal \sref{IV.10.4.5} y \sref{IV.10.5.4}.
Entonces $S$ satisface también las condiciones $2^\circ$ y $3^\circ$ de \sref{IV.10.6.1}: en efecto, todo anillo local $\sh{O}_{S,s}$ es de dimensión $0$ o $1$, y por tanto universalmente catenario \sref{IV.7.2.9}; por otra parte, una componente irreducible $S'$ de $S$ o bien se reduce
\oldpage[IV-3]{110}
a un punto o bien es de dimensión $1$, y para todo punto cerrado $s\in S'$, $\sh{O}_{S',s}$ es necesariamente de dimensión $1$.

De estas observaciones y de \sref{IV.10.6.1} se deduce que todo preesquema localmente de tipo finito sobre $S$ satisface también las propiedades $1^\circ$, $2^\circ$ y $3^\circ$; este es en particular el caso de los preesquemas \emph{localmente de tipo finito sobre un cuerpo o sobre $\mathbf{Z}$}.
\end{env}

\begin{env}[10.7.2]
\label{IV.10.7.2}
Sea $A$ un anillo local noetheriano \emph{universalmente catenario} y sea $S$ el complementario del \emph{punto cerrado} $a$ en $X=\operatorname{Spec}(A)$.
Entonces $S$ satisface las condiciones $1^\circ$, $2^\circ$ y $3^\circ$ de \sref{IV.10.6.1}: en efecto, ya se ha visto que $S$ es un preesquema de Jacobson \sref{IV.10.5.9}; como $A$ es universalmente catenario, también lo son todos los anillos locales $A_\mathfrak{p}$ en los ideales primos de $A$ \sref{IV.5.6.3}.
Por otra parte, una componente irreducible $S'$ de $S$ es el complementario de $a$ en una componente irreducible $X'$ de $X$. Para todo punto cerrado $x$ de $S$, la adherencia de $x$ en $X$ es $\{x,a\}$, luego $\dim(\overline{\{x\}})=1$. Como $X'$ es biequidimensional por hipótesis \sref[0]{0.14.3.3}, se deduce que $\dim(\sh{O}_{X',x})=\dim(X')-1$ \sref[0]{0.14.3.2}, lo que prueba que $S'$ es equicodimensional \sref{IV.5.1.4}.
\end{env}

\begin{env}[10.7.3]
\label{IV.10.7.3}
Sea $A$ un anillo de valoración discreta; mostremos que en $B=A[T_1,\ldots,T_n]$, con $n$ indeterminadas, existen dos ideales maximales $\mathfrak{m}$, $\mathfrak{n}$ de \emph{alturas respectivas} $n$ y $n+1$.
Lo hemos visto en \sref{IV.5.2.5},~(i) para $n=1$; probémoslo por inducción sobre $n$.
Como $A[T_1,\ldots,T_{n+1}]$ es un $A[T_1,\ldots,T_n]$-módulo libre, y por tanto fielmente plano, existen ideales maximales $\mathfrak{m}'$ y $\mathfrak{n}'$ de $A[T_1,\ldots,T_{n+1}]$ situados respectivamente sobre $\mathfrak{m}$ y $\mathfrak{n}$ \sref[0]{0.6.5.1}; por \sref{IV.5.5.3}, estos ideales tienen alturas respectivas $n+1$ y $n+2$, lo que prueba la afirmación.
Supongamos $n\geq 2$ en lo que sigue.
Sea $\mathfrak{J}=\mathfrak{m}\cap\mathfrak{n}=\mathfrak{m}\mathfrak{n}$ y sea $R=1+\mathfrak{J}$, que es un subconjunto multiplicativo de $B$.
Si $B'=R^{-1}B$, el ideal $\mathfrak{J}'=R^{-1}\mathfrak{J}$ está contenido en el radical de $B'$ (Bourbaki, \emph{Alg\`ebre commutative}, cap.~III, \S3, n.~5, prop.~12).
Como espacio topológico, $X'=\operatorname{Spec}(B')$ se identifica con un subespacio de $X=\operatorname{Spec}(B)$, y en los puntos $x$ de $X'$ los anillos locales $\sh{O}_{X,x}$ y $\sh{O}_{X',x}$ son los mismos \sref[I]{I.1.6.2}.
Consideremos en $X'$ el cerrado $Y'=V(\mathfrak{J}')$, y pongamos $S=X'\setminus Y'$.
Entonces $S$ es un preesquema de Jacobson \sref{IV.10.5.7}, y es evidentemente irreducible y noetheriano; además, para todo $s\in S$, el anillo local $\sh{O}_{S,s}=\sh{O}_{X,s}$ es universalmente catenario por \sref{IV.5.6.3}, pues $A$ es universalmente catenario \sref{IV.5.6.4}.
Sin embargo, existen dos puntos cerrados $a,b$ de $S$ tales que $\sh{O}_{S,a}$ y $\sh{O}_{S,b}$ \emph{no tienen la misma dimensión}; dicho de otro modo, $S$ no satisface la condición $3^\circ$ de \sref{IV.10.6.1}.
Para verlo, consideremos los dos ideales maximales $\mathfrak{m}'=R^{-1}\mathfrak{m}$ y $\mathfrak{n}'=R^{-1}\mathfrak{n}$ de $B'$, que tienen alturas $n$ y $n+1$, respectivamente.
Se tiene $\mathfrak{J}'=\mathfrak{m}'\cap\mathfrak{n}'$, y por tanto $\mathfrak{J}'$ no está contenido en ningún ideal primo de $B'$ distinto de $\mathfrak{m}'$ y $\mathfrak{n}'$; estos son, en consecuencia, los únicos ideales maximales de $B'$.
Sean $a',b'$ los únicos puntos cerrados de $X'$, correspondientes a $\mathfrak{m}'$ y $\mathfrak{n}'$.
Existe en $B'$ un ideal primo no maximal $\mathfrak{p}'\subset\mathfrak{m}'$ que no está contenido en $\mathfrak{n}'$: basta mostrar que en $B$ existe un ideal primo no maximal contenido en $\mathfrak{m}$ y no en $\mathfrak{n}$; para ello puede, por ejemplo, considerarse la fibra de $\mathfrak{m}$ para el morfismo correspondiente a la inyección $A[T_1]\longrightarrow B=A[T_1,\ldots,T_n]$ y aplicar \sref{IV.6.1.2}.
Considerando una cadena maximal de ideales primos entre $\mathfrak{p}'$ y $\mathfrak{m}'$, y sustituyendo $\mathfrak{p}'$ por el penúltimo ideal de esta cadena, podemos por tanto suponer que el punto $a$ de $X'$ correspondiente a $\mathfrak{p}'$ tiene por adherencia $\{a,a'\}$ en $X'$.
Como $B'_{\mathfrak{m}'}$ es biequidimensional, $\mathfrak{p}'$ tiene entonces altura $n-1$.
Del mismo modo se construye un ideal primo no maximal $\mathfrak{q}'$ de $B'$ de altura $n$ tal que, si $b$ es el punto correspondiente de $X'$, la adherencia de $b$ en $X'$ es $\{b,b'\}$.
Así, $a$ y $b$ pertenecen a $S$ y son cerrados en $S$, y proporcionan por tanto el ejemplo buscado.
\end{env}

\subsection{Profundidad rectificada}
\label{subsection:IV.10.8}

\begin{definition}[10.8.1]
\label{IV.10.8.1}
Sea $X$ un preesquema localmente noetheriano y $\sh{F}$ un $\sh{O}_X$-Módulo coherente.
Para todo $x\in X$, llamamos \emph{profundidad rectificada} de $\sh{F}$ en el punto $x$ y denotamos $\operatorname{prof}^*_x(\sh{F})$ el número (entero $\geq 0$ o $+\infty$) igual a
\begin{equation}
\label{IV.10.8.1.1}
\tag{10.8.1.1}
\operatorname{prof}^*_x(\sh{F})=\operatorname{prof}(\sh{F}_x)+\dim(\overline{\{x\}})
\end{equation}
donde $\overline{\{x\}}$ es la adherencia del punto $x$ en $X$.
Para todo subconjunto $Z$ de $X$, llamamos \emph{profundidad rectificada de $\sh{F}$ a lo largo de $Z$} y denotamos $\operatorname{prof}^*_Z(\sh{F})$ el número
\begin{equation}
\label{IV.10.8.1.2}
\tag{10.8.1.2}
\operatorname{prof}^*_Z(\sh{F})=\inf_{x\in Z}\operatorname{prof}^*_x(\sh{F}).
\end{equation}
\end{definition}

Dicho de otro modo, para todo entero $n$, la relación $\operatorname{prof}^*_Z(\sh{F})\geq n$ equivale a $\operatorname{prof}^*_x(\sh{F})\geq n$ para todo $x\in Z$.
Si $Z=X$, se escribe $\operatorname{prof}^*(\sh{F})$ en lugar de $\operatorname{prof}^*_X(\sh{F})$.

\oldpage[IV-3]{111}

\begin{remarks}[10.8.2]
\label{IV.10.8.2}
(i) En todo punto cerrado $x\in X$, la profundidad rectificada es igual a la profundidad.

(ii) Sea $Y$ un subpreesquema cerrado de $X$, $j:Y\to X$ la inyección canónica y $\sh{G}$ un $\sh{O}_Y$-módulo coherente.
Sabemos \sref{IV.5.7.3},~(vi) que $\operatorname{prof}(\sh{G}_x)=\operatorname{prof}(j_*(\sh{G})_x)$ para todo $x\in Y$; deducimos por tanto que $\operatorname{prof}^*_x(\sh{G})=\operatorname{prof}^*_x(j_*(\sh{G}))$ para todo $x\in Y$.

(iii) La noción de profundidad rectificada solo presenta interés cuando es de carácter \emph{local}, es decir, cuando no cambia al sustituir $X$ por un entorno abierto arbitrario $U$ de $x$.
Esto exige evidentemente que $x$ no sea aislado en $\overline{\{x\}}$ cuando $x$ no es un punto cerrado de $X$, y por consiguiente que $X$ sea un \emph{preesquema de Jacobson} \sref{IV.10.4.5.1}; la mayoría de las veces, también será necesario saber que $\dim(U)=\dim(X)$ para todo abierto denso $U$ de $X$, y habrá que suponer por tanto que $X$ satisface además las condiciones $2^\circ$ y $3^\circ$ de \sref{IV.10.6.1}.
\end{remarks}

\begin{lemma}[10.8.3]
\label{IV.10.8.3}
Sea $X$ un preesquema regular y biequidimensional y $\sh{F}$ un $\sh{O}_X$-Módulo coherente.
Entonces, para todo $x\in X$, se tiene
\begin{equation}
\label{IV.10.8.3.1}
\tag{10.8.3.1}
\operatorname{prof}^*_x(\sh{F})=\dim(X)-\dim.\operatorname{proj}(\sh{F}_x).
\end{equation}
\end{lemma}

\begin{proof}
En efecto, como $X$ es biequidimensional, se tiene \sref[0]{0.14.3.5.1}
\[
  \dim(\overline{\{x\}})=\dim(X)-\operatorname{codim}(\overline{\{x\}},X)=\dim(X)-\dim(\sh{O}_{X,x})
\]
en virtud de \sref{IV.5.1.2}.
Por otra parte, como $X$ es regular, se tiene por \sref[0]{0.17.3.4}
\[
  \operatorname{prof}(\sh{F}_x)=\dim(\sh{O}_{X,x})-\dim.\operatorname{proj}(\sh{F}_x)
\]
de donde el lema.
\end{proof}

\begin{corollary}[10.8.4]
\label{IV.10.8.4}
Bajo las hipótesis de \sref{IV.10.8.3}, la función $x\mapsto\operatorname{prof}^*_x(\sh{F})$ es semicontinua inferiormente.
\end{corollary}

\begin{proof}
Esto se deduce de \sref{IV.10.8.3.1}, pues $x\mapsto\dim.\operatorname{proj}(\sh{F}_x)$ es semicontinua superiormente \sref{IV.6.11.1}.
\end{proof}

\begin{proposition}[10.8.5]
\label{IV.10.8.5}
Sea $S$ un preesquema localmente noetheriano, $X$ un preesquema localmente de tipo finito sobre $S$ y $\sh{F}$ un $\sh{O}_X$-módulo coherente.
Supongamos que $S$ satisface las condiciones: $1^\circ$~$S$ es un preesquema de Jacobson; $2^\circ$~$S$ es regular; $3^\circ$~las componentes irreducibles de $S$ son equicodimensionales.
Entonces la función $x\mapsto\operatorname{prof}^*_x(\sh{F})$ es semicontinua inferiormente; dicho de otro modo, para todo entero $n$, el conjunto $U_n$ de los $x\in X$ tales que $\operatorname{prof}^*_x(\sh{F})\geq n$ es abierto.
\end{proposition}

\begin{proof}
Como los anillos locales $\sh{O}_{S,s}$ de $S$ son regulares, son universalmente catenarios \sref{IV.5.6.4}; dicho de otro modo, $S$ satisface las condiciones $1^\circ$, $2^\circ$ y $3^\circ$ de \sref{IV.10.6.1}, y por tanto lo mismo ocurre con $X$ \sref{IV.10.6.1},~(i).
Como la noción de profundidad rectificada es entonces de carácter local \sref{IV.10.8.2},~(iii), podemos limitarnos al caso en que $S=\operatorname{Spec}(A)$ y $X=\operatorname{Spec}(B)$ son afines, siendo $A$ un anillo regular y $B$ una $A$-álgebra de tipo finito, y por tanto un cociente de un anillo de polinomios $C=A[T_1,\ldots,T_n]$; este último anillo es regular \sref[0]{0.17.3.7}.
Podemos pues suponer que $X$ es un subpreesquema cerrado de un preesquema regular $Y$ que satisface también las condiciones $1^\circ$, $2^\circ$ y $3^\circ$ de \sref{IV.10.6.1}; teniendo en cuenta la observación \sref{IV.10.8.2},~(ii), quedamos así reducidos al caso en que $X$ es además regular y noetheriano.
Pero como los anillos locales de $X$ son entonces íntegros, las componentes irreducibles de $X$ son
\oldpage[IV-3]{112}
abiertas \sref[I]{I.6.1.10}, y por tanto también podemos suponer que $X$ es irreducible.
Entonces, como los anillos locales de $X$ son catenarios \sref[0]{0.16.5.12}, la hipótesis $3^\circ$ de \sref{IV.10.6.1} implica que $X$ es biequidimensional (\sref{IV.5.1.5} y \sref[0]{0.14.3.3}); basta pues aplicar \sref{IV.10.8.4}.
\end{proof}

Obsérvese que, si $S$ es el espectro de un cuerpo o de $\mathbf{Z}$, satisface las condiciones de \sref{IV.10.8.5}.

\begin{corollary}[10.8.6]
\label{IV.10.8.6}
Bajo las hipótesis de \sref{IV.10.8.5}, para todo $x\in X$, el número $\operatorname{prof}^*_x(\sh{F})$ es el único entero $n$ que tiene la propiedad siguiente: existe un entorno abierto $U$ de $x$ tal que, para todo punto $x'\in U\cap\overline{\{x\}}$ cerrado en $U$, se tiene $\operatorname{prof}(\sh{F}_{x'})=n$.
En particular, para que $\operatorname{prof}^*_x(\sh{F})\geq m$ (resp.\ $\operatorname{prof}^*_x(\sh{F})\leq m$), es necesario y suficiente que exista un entorno abierto $V$ de $x$ en $X$ tal que, para todo punto $x'\in V\cap\overline{\{x\}}$ cerrado en $V$, se tenga $\operatorname{prof}(\sh{F}_{x'})\geq m$ (resp.\ $\operatorname{prof}(\sh{F}_{x'})\leq m$).
\end{corollary}

\begin{proof}
En efecto, podemos limitarnos al caso en que $x$ no es un punto cerrado; si $\operatorname{prof}^*_x(\sh{F})=n$, el conjunto de los $y\in X$ tales que $\operatorname{prof}^*_y(\sh{F})\leq n$ es cerrado en virtud de \sref{IV.10.8.5}, y por tanto contiene a $\overline{\{x\}}$; en virtud de la semicontinuidad inferior de $\operatorname{prof}^*_y(\sh{F})$, existe un entorno abierto $U$ de $x$ tal que $\operatorname{prof}^*_y(\sh{F})=n$ para todo $y\in U\cap\overline{\{x\}}$, y por tanto $\operatorname{prof}(\sh{F}_{x'})=n$ si $x'\in U\cap\overline{\{x\}}$ es cerrado (puesto que la noción de profundidad rectificada es local).
\end{proof}

Para los preesquemas que satisfacen las hipótesis de \sref{IV.10.8.5}, la noción de profundidad rectificada puede por tanto definirse mediante los valores de la profundidad en los \emph{puntos cerrados} de $X$ (estos forman un conjunto muy denso en todo cerrado de $X$).

\begin{proposition}[10.8.7]
\label{IV.10.8.7}
Sea $S$ un preesquema localmente noetheriano que satisfaga las condiciones $1^\circ$, $2^\circ$ y $3^\circ$ de \sref{IV.10.6.1}.
Sea $X$ un preesquema localmente de tipo finito sobre $S$, $\sh{F}$ un $\sh{O}_X$-Módulo coherente y $Y=\operatorname{Supp}(\sh{F})$.
Entonces, para todo $x\in Y$, se tiene
\begin{equation}
\label{IV.10.8.7.1}
\tag{10.8.7.1}
\operatorname{prof}^*_x(\sh{F})=\dim_x(Y)-\operatorname{coprof}(\sh{F}_x).
\end{equation}
\end{proposition}

\begin{proof}
En efecto, por definición, $\operatorname{coprof}(\sh{F}_x)=\dim(\sh{F}_x)-\operatorname{prof}(\sh{F}_x)$, y se deduce de \sref{IV.10.6.4} que $\dim_x(Y)=\dim(\overline{\{x\}})+\dim(\sh{F}_x)$; de donde \sref{IV.10.8.7.1} por definición de $\operatorname{prof}^*_x(\sh{F})$.
\end{proof}

\begin{corollary}[10.8.8]
\label{IV.10.8.8}
Supongamos que las hipótesis sobre $f$ y $\sh{F}$ son las de \sref{IV.9.9.1}. Entonces la función $x\mapsto\operatorname{prof}^*_x(\sh{F}_{f(x)})$ es construible.
\end{corollary}

\begin{proof}
Para todo $s\in S$, la fibra $X_s$, al ser localmente de tipo finito sobre $k(s)$, es un preesquema de Jacobson \sref{IV.10.4.7}. Como $\operatorname{Spec}(k(s))$ satisface las condiciones de \sref{IV.10.6.1}, se tiene
\[
  \operatorname{prof}^*_x(\sh{F}_{f(x)})
  =\dim_x(\operatorname{Supp}(\sh{F}_{f(x)}))
   -\operatorname{coprof}((\sh{F}_{f(x)})_x)
\]
por \sref{IV.10.8.7}.
Ahora bien, si $Z=\operatorname{Supp}(\sh{F})$, se tiene $Z_{f(x)}=\operatorname{Supp}(\sh{F}_{f(x)})$ \sref[I]{I.9.1.13} y $Z$ es localmente construible \sref{IV.8.9.1}; por tanto, las funciones $x\mapsto\dim_x(\operatorname{Supp}(\sh{F}_{f(x)}))$ y $x\mapsto\operatorname{coprof}((\sh{F}_{f(x)})_x)$ son localmente construibles (\sref{IV.9.9.1} y \sref{IV.9.9.3}), lo que prueba la proposición.
\end{proof}

\subsection{Espectros maximales y ultrapreesquemas}
\label{subsection:IV.10.9}

Los resultados de esta sección no se utilizarán en lo que sigue.

\begin{env}[10.9.1]
\label{IV.10.9.1}
Sea $X$ un preesquema de Jacobson, y sea $S(X)$ el \emph{espacio anillado} cuyo espacio subyacente es el subespacio de los puntos cerrados de $X$, y cuyo haz de anillos es el haz inducido sobre este subespacio por $\sh{O}_X$, es decir, el haz $\psi^*(\sh{O}_X)$, donde $\psi:S(X)\to X$ denota la inyección canónica.
Como $\psi$ es un
\oldpage[IV-3]{113}
cuasihomeomorfismo, hemos visto \sref{IV.10.2.8},~(ii) que, si $\theta:\sh{O}_X\to\psi_*(\sh{O}_{S(X)})$ es el homomorfismo de haces de anillos tal que $\theta^\sharp:\psi^*(\sh{O}_X)\to\sh{O}_{S(X)}$ es la identidad, entonces $j_X=(\psi,\theta)$ es un \emph{cuasiisomorfismo} de espacios anillados, y $\sh{F}\mapsto j_X^*(\sh{F})$ es una equivalencia entre la categoría de $\sh{O}_X$-Módulos y la de $\sh{O}_{S(X)}$-Módulos.
Está claro que, bajo esta equivalencia, a los $\sh{O}_X$-módulos localmente libres (resp.\ coherentes) corresponden los $\sh{O}_{S(X)}$-módulos localmente libres (resp.\ coherentes). Además, si $\sh{L}$ es un $\sh{O}_X$-módulo localmente libre y $U$ es un abierto de $X$ tal que $\sh{L}|_U\simeq(\sh{O}_X|_U)^n$, entonces
\[
  j_X^*(\sh{L})|_{U\cap S(X)}
  \simeq \bigl(\sh{O}_{S(X)}|_{U\cap S(X)}\bigr)^n.
\]
\end{env}

\begin{env}[10.9.2]
\label{IV.10.9.2}
Sean $X$, $Y$ dos preesquemas de Jacobson y $f=(\rho,\lambda):X\to Y$ un morfismo \emph{localmente de tipo finito}.
Se ha visto \sref{IV.10.4.7} que $\rho(S(X))\subset S(Y)$ y, por restricción de $\rho$ a $S(X)$, se define por tanto una aplicación continua $S(\rho):S(X)\to S(Y)$.
Para todo abierto $V$ de $Y$, se define por composición un homomorfismo de anillos
\[
  \Gamma(V\cap S(Y),\sh{O}_{S(Y)})\to\Gamma(V,\sh{O}_Y)\xrightarrow{\lambda_V}\Gamma(f^{-1}(V),\sh{O}_X)\to\Gamma(f^{-1}(V)\cap S(X),\sh{O}_{S(X)})
\]
donde los dos isomorfismos extremos se han definido en \sref{IV.10.9.1}; está claro que esto define un homomorfismo de haces de anillos $S(\lambda):\sh{O}_{S(Y)}\to\rho_*(\sh{O}_{S(X)})$ (recordando que los abiertos de $S(X)$ (resp.\ $S(Y)$) corresponden biyectivamente a los de $X$ (resp.\ $Y$) \sref{IV.10.2.1}); se obtiene así un morfismo de espacios anillados $S(f)=(S(\rho),S(\lambda)):(S(X),\sh{O}_{S(X)})\to(S(Y),\sh{O}_{S(Y)})$ tal que el diagrama
\[
  \xymatrix{
    S(X)\ar[r]^{S(f)}\ar[d]_{j_X} & S(Y)\ar[d]^{j_Y}\\
    X\ar[r]^{f} & Y
  }
\]
es conmutativo; además, si $Z$ es un tercer preesquema de Jacobson y $g:Y\to Z$ un morfismo localmente de tipo finito, $S(g\circ f)=S(g)\circ S(f)$.
Se ha definido así un \emph{funtor covariante} $S:\mathbf{C}\to\mathbf{C}'$, donde $\mathbf{C}'$ es la categoría de los \emph{espacios anillados en anillos locales}, y $\mathbf{C}$ la categoría cuyos objetos son los \emph{preesquemas de Jacobson} y cuyos morfismos son los morfismos localmente de tipo finito entre preesquemas de Jacobson.
\end{env}

\begin{env}[10.9.3]
\label{IV.10.9.3}
Determinemos ahora la subcategoría $\mathbf{C}''$ de $\mathbf{C}'$ formada por los espacios anillados isomorfos a $S(X)$ y cuyos morfismos provienen de los $S(f)$.
Supongamos primero que $X=\operatorname{Spec}(A)$, donde $A$ es un anillo de Jacobson; entonces $S(X)$ es el conjunto de los ideales \emph{maximales} de $A$, provisto: $1^\circ$~de la topología inducida por la de $X$, de modo que una base de esta topología está formada por los $D^m(h)=D(h)\cap S(X)$, es decir, el conjunto de los ideales maximales $\mathfrak{m}$ de $A$ tales que $h\notin\mathfrak{m}$; $2^\circ$~del haz de anillos $\sh{O}_{S(X)}$ tal que $\Gamma(D^m(h),\sh{O}_{S(X)})=A_h$.
Diremos que este espacio anillado es el \emph{espectro maximal} del anillo de Jacobson $A$ y lo denotaremos $\operatorname{Spm}(A)$.

Obsérvese que, si $j:D(h)\to X$ es la inyección canónica, el espacio anillado inducido sobre $D^m(h)$ por $S(X)$ y $S(D(h))$ son canónicamente isomorfos, y que la inyección canónica $D^m(h)\to S(X)$ de espacios anillados es igual a $S(j)$.

Sea $B$ un segundo anillo de Jacobson, $Y=\operatorname{Spec}(B)$, $\varphi:B\to A$ un homomorfismo de anillos que hace de $A$ una $B$-álgebra de \emph{tipo finito}, $f=({}^a\varphi,\widetilde{\varphi}):X\to Y$ el morfismo correspondiente de preesquemas, y
\[
  S(f):\operatorname{Spm}(A)\to\operatorname{Spm}(B)
\]
el morfismo de espacios anillados correspondiente a $f$.
Está claro que $S(f)=(\psi,\theta)$ es un morfismo de espacios anillados en anillos locales, es decir, $(\mathrm{Err}_{\mathrm{II}},1.8.2)$, para todo $x\in\operatorname{Spm}(A)$, $\theta^\sharp_x$ es un homomorfismo \emph{local}.
Recíprocamente:
\end{env}

\begin{proposition}[10.9.4]
\label{IV.10.9.4}
Sean $A$, $B$ dos anillos de Jacobson.
Si $u=(\psi,\theta):\operatorname{Spm}(A)\to\operatorname{Spm}(B)$ es un morfismo de espacios anillados en anillos locales tal que $\Gamma(\theta):B\to A$ hace de $A$ una $B$-álgebra de tipo finito, existe un único morfismo de preesquemas $f:\operatorname{Spec}(A)\to\operatorname{Spec}(B)$ tal que $u=S(f)$.
\end{proposition}

\begin{proof}
La unicidad de $f$ es evidente, pues si $f=({}^a\varphi,\widetilde{\varphi})$, se debe tener $\varphi=\Gamma(\theta)$; queda por mostrar que $S(f)$ está definido y que $u=S(f)$.
Ahora bien, la primera afirmación se deduce de que se supone que $\varphi$ hace de $A$ una $B$-álgebra de tipo finito, y por tanto $f(\operatorname{Spm}(A))\subset\operatorname{Spm}(B)$; el hecho de que $\theta^\sharp_x$ sea un homomorfismo local para todo $x\in\operatorname{Spm}(A)$ permite repetir el razonamiento de \sref[I]{I.1.7.3}, restringiéndonos a los puntos $x$ de $\operatorname{Spm}(A)$, para mostrar así sucesivamente que $\psi(x)={}^a\varphi(x)$ para todo $x\in\operatorname{Spm}(A)$, y después que $\theta^\sharp_x=\varphi_x:B_{\psi(x)}\to A_x$, lo que termina de probar que $u=S(f)$.
\end{proof}

\oldpage[IV-3]{114}

\begin{env}[10.9.5]
\label{IV.10.9.5}
Consideremos ahora un espacio anillado $(X,\sh{O}_X)$; diremos que un abierto $U$ de $X$ es \emph{ultraafín} si el espacio anillado inducido $(U,\sh{O}_X|U)$ es isomorfo a un espectro maximal $\operatorname{Spm}(A)$, donde $A$ es un anillo de Jacobson.
Diremos que $X$ es un \emph{ultrapreesquema} si todo punto de $X$ admite un entorno abierto ultraafín.
Se muestra, como en \sref[I]{I.2.1.3} y \sref[I]{I.2.1.4}, que los abiertos ultraafines forman una \emph{base} de la topología de $X$ y que $X$ es un espacio de Kolmogórov.
Si $Y$ es un segundo ultrapreesquema, diremos que un morfismo de espacios anillados $f:X\to Y$ es un \emph{morfismo de ultrapreesquemas} si satisface las condiciones siguientes: $1^\circ$~$f$ es un morfismo de espacios anillados en anillos locales; $2^\circ$~para todo $x\in X$, existe un entorno abierto ultraafín $V$ de $f(x)$ en $Y$ y un abierto ultraafín $U$ de $x$ en $X$ tales que $f(U)\subset V$ y que el homomorfismo $\Gamma(V,\sh{O}_Y)\to\Gamma(U,\sh{O}_X)$ correspondiente a $f$ hace de $\Gamma(U,\sh{O}_X)$ una $\Gamma(V,\sh{O}_Y)$-\emph{álgebra de tipo finito}.

Es inmediato que la compuesta de dos morfismos tales está bien definida y es un tercer morfismo tal por la observación final de \sref{IV.10.9.3}.
Está claro que la categoría $C''_0$ así definida es una subcategoría de $C'$ que contiene a $C''$; nos proponemos ahora mostrar que $C''=C''_0$, dicho de otro modo:
\end{env}

\begin{proposition}[10.9.6]
\label{IV.10.9.6}
El funtor $X\mapsto S(X)$ de $\mathbf{C}$ en $\mathbf{C}''_0$ es una equivalencia de categorías.
\end{proposition}

\begin{proof}
$1^\circ$ Mostremos primero que el funtor $X\mapsto S(X)$ es \emph{plenamente fiel}, dicho de otro modo, que para $X$, $Y$ en $\mathbf{C}$, la aplicación canónica
\[
  \operatorname{Hom}_{\mathbf{C}}(X,Y)\to\operatorname{Hom}_{\mathbf{C}''_0}(S(X),S(Y))
\]
es \emph{biyectiva}.
En primer lugar es \emph{inyectiva}: sean $f$, $g$ dos morfismos localmente de tipo finito de $X$ en $Y$, y supongamos que $S(f)=S(g)$; se tiene $f^{-1}(V)\cap S(X)=g^{-1}(V)\cap S(X)$ para todo abierto $V$ de $Y$, luego \sref{IV.10.3.1} y \sref{IV.10.2.7} dan $f^{-1}(V)=g^{-1}(V)$; basta por tanto probar que, para todo abierto afín $V$ de $Y$, $f$ y $g$ coinciden en $f^{-1}(V)=g^{-1}(V)$; dicho de otro modo, nos reducimos al caso en que $Y=\operatorname{Spec}(B)$ es el espectro de un anillo de Jacobson $B$.
Basta \sref[I]{I.2.2.4} mostrar que los homomorfismos de anillos $B\to\Gamma(X,\sh{O}_X)$ correspondientes a $f$ y $g$ son iguales.
Ahora bien, para todo $s\in B$, las imágenes de $s$ por estos homomorfismos son dos secciones de $\sh{O}_X$ sobre $X$ que, por hipótesis, inducen la \emph{misma} sección sobre $S(X)$; sabemos \sref{IV.10.2.8} que estas secciones son idénticas, de donde nuestra afirmación.
En segundo lugar, probemos que todo morfismo $h:S(X)\to S(Y)$ (para la categoría $\mathbf{C}''_0$) es de la forma $S(f)$, donde $f:X\to Y$ es un morfismo localmente de tipo finito.
Por hipótesis, existe un recubrimiento abierto ultraafín $(U'_\lambda)$ (resp.\ $(V'_\mu)$) de $S(X)$ (resp.\ $S(Y)$) tal que, para todo $\lambda$, $h(U'_\lambda)$ está contenido en uno de los $V'_\mu$ y el homomorfismo de anillos correspondiente a la restricción $h_\lambda:U'_\lambda\to V'_\mu$ hace del segundo anillo un álgebra de tipo finito sobre el primero.
Podemos suponer que $U'_\lambda=U_\lambda\cap S(X)$ y $V'_\mu=V_\mu\cap S(Y)$, donde $U_\lambda$ y $V_\mu$ son abiertos afines de $X$ y $Y$, respectivamente, determinados de manera única. Si, para todo $\lambda$, $h_\lambda=S(f_\lambda)$ para un morfismo de tipo finito $f_\lambda:U_\lambda\to V_\mu$, entonces la primera parte del razonamiento, aplicada a las restricciones de $f_\alpha$ y $f_\beta$ a $U_\alpha\cap U_\beta$, muestra que los $f_\lambda$ son restricciones de un único morfismo $f:X\to Y$, y claramente $h=S(f)$. Quedamos así reducidos al caso en que $X$ e $Y$ son afines, y la conclusión se deduce de \sref{IV.10.9.4}.

$2^\circ$ Queda por probar que todo ultrapreesquema $X'$ es de la forma $S(X)$ para un preesquema de Jacobson $X$ (que será necesariamente único salvo isomorfismo, en virtud de $1^\circ$).
Existen abiertos ultraafines $(U'_\alpha)$ de $X'$, cada uno de los cuales es de la forma $S(U_\alpha)$, siendo $U_\alpha$ el espectro de un anillo de Jacobson.
Para cada par de índices $\alpha$, $\beta$, consideremos el único abierto $U_{\alpha\beta}$ de $U_\alpha$ cuya traza sobre $U'_\alpha$ es $U'_\alpha\cap U'_\beta$; en virtud de $1^\circ$, el automorfismo identidad de $U'_\alpha\cap U'_\beta$ es de la forma $S(\theta_{\alpha\beta})$, donde $\theta_{\alpha\beta}:U_{\alpha\beta}\to U_{\beta\alpha}$ es un isomorfismo de preesquemas.
Se verifica fácilmente (en virtud de $1^\circ$) que la familia $(\theta_{\alpha\beta})$ satisface la condición de pegado \sref[0]{0.4.1.7}, y que esta familia define un preesquema $X$; está claro entonces que se tiene $X'=S(X)$, lo que termina la demostración.
\end{proof}

\subsection{Espacios algebraicos de Serre}
\label{subsection:IV.10.10}

\begin{env}[10.10.1]
\label{IV.10.10.1}
El lenguaje introducido por Serre en (FAC) resulta a veces cómodo, especialmente en las cuestiones en que el interés principal recae sobre los puntos racionales sobre el cuerpo base $k$ (algebraicamente cerrado por hipótesis) de las ``variedades algebraicas'' sobre $k$ que se consideran.
Esbozaremos aquí este lenguaje enlazándolo con las consideraciones precedentes, para permitir al lector traducir los enunciados de Serre al lenguaje de los esquemas.
Por lo demás, también es posible desarrollar el lenguaje de Serre para preesquemas sobre un cuerpo no algebraicamente cerrado (e incluso sobre un anillo artiniano); pero esto introduce complicaciones técnicas considerables y, además, sobre un cuerpo base arbitrario desaparecen las ventajas (ante todo psicológicas) del punto de vista de Serre; por ello nos limitaremos al marco fijado por Serre.
La presente sección, al igual que la precedente, no se utilizará en lo que sigue de este Tratado, y nos contentaremos por tanto con indicaciones breves.
\end{env}

\oldpage[IV-3]{115}

\begin{env}[10.10.2]
\label{IV.10.10.2}
Dado un ultrapreesquema $R$, se puede definir naturalmente (como en toda categoría) la noción de $R$-\emph{ultrapreesquema}.
Consideremos en particular un cuerpo algebraicamente cerrado $k$; $\operatorname{Spm}(k)$ es entonces idéntico a $\operatorname{Spec}(k)$; diremos que un $\operatorname{Spm}(k)$-ultrapreesquema $X'$ es un \emph{espacio prealgebraico sobre $k$}: es un espacio $k$-anillado en anillos locales en el que todo punto admite un entorno abierto isomorfo al espectro maximal de una $k$-álgebra de tipo finito; por tanto, esto equivale a decir (por \sref{IV.10.9.6}) que $X'=S(X)$, donde $X$ es un $k$-preesquema localmente de tipo finito.
Si $X$, $Y$, $Z$ son tres $k$-preesquemas localmente de tipo finito, lo mismo ocurre con $X\times_ZY$ \sref[I]{subsection:I.3.4}; por tanto, en virtud de \sref{IV.10.9.6}, existen productos en la categoría de los espacios prealgebraicos sobre $k$, tanto el producto ``sobre $k$'' $X'\times_kY'$ como el ``producto fibrado'' $X'\times_{Z'}Y'$, donde $X'$, $Y'$, $Z'$ son tres espacios prealgebraicos sobre $k$.
Se puede definir, pues, el morfismo diagonal $\Delta_{X'}:X'\to X'\times_k X'$ (que no es otro que $S(\Delta_X)$); se dice que $X'$ es un \emph{espacio algebraico sobre $k$} si $X$ es un esquema, y esto equivale a decir que la imagen de $\Delta_{X'}$ es un subconjunto \emph{cerrado} de $X'\times_k X'$.
\end{env}

\begin{env}[10.10.3]
\label{IV.10.10.3}
Las simplificaciones que provienen de la hipótesis de que $k$ es algebraicamente cerrado consisten ante todo en lo siguiente: para todo $k$-preesquema $X$ localmente de tipo finito, existe una correspondencia biyectiva entre los \emph{puntos cerrados de $X$}, los \emph{puntos de $X$ con valores en $k$} \sref[I]{I.3.4.4} y los \emph{puntos racionales de $X$ sobre $k$} \sref[I]{I.3.4.5}, en virtud de \sref[I]{I.6.4.2}.
Esto muestra en particular \sref[I]{I.3.4.3.1} que, para dos espacios $k$-prealgebraicos $X'$, $Y'$, el conjunto subyacente al producto $X'\times_kY'$ es idéntico al \emph{conjunto producto} $X'\times Y'$ de los conjuntos subyacentes (pero, desde luego, la topología del espacio subyacente a $X'\times_kY'$ no es la \emph{topología producto} de las topologías de los espacios subyacentes a $X'$ e $Y'$ y, en general, es estrictamente más fina que esta última).

Por otra parte, los anillos locales $\sh{O}_x$ en los puntos de un espacio prealgebraico $X'$ sobre $k$ son $k$-álgebras cuyo cuerpo residual es \emph{isomorfo a $k$}; si $A$ y $B$ son dos $k$-álgebras de esta índole, un $k$-homomorfismo $\varphi:A\to B$ es necesariamente \emph{local}: en efecto, si existiera un elemento $x$ del ideal maximal de $A$ tal que $\varphi(x)$ fuese invertible, existiría un $\lambda\in k$ no nulo tal que $\varphi(x-\lambda.1)$ perteneciese al ideal maximal de $B$, lo cual es absurdo puesto que $x-\lambda.1$ es invertible.
De ello concluimos inmediatamente que, si $X'=S(X)$, $Y'=S(Y)$ son dos espacios prealgebraicos sobre $k$, todo morfismo $X'\to Y'$ de \emph{espacios $k$-anillados} es también un morfismo de espacios $k$-anillados \emph{en anillos locales} (\sref[I]{I.1.8.2}; cf.\ $\mathrm{Err}_{\mathrm{II}}$). Además, con la notación precedente, si $A$ y $B$ son $k$-álgebras de tipo finito, el homomorfismo correspondiente hace de $B$ una $A$-álgebra de tipo finito; por tanto, todo morfismo $X'\to Y'$ de espacios $k$-anillados es, por \sref{IV.10.9.6}, de la forma $S(f)$, donde $f:X\to Y$ es un morfismo de $k$-preesquemas.

Por último, para todo abierto $U$ de $X'$, toda sección $s\in\Gamma(U,\sh{O}_{X'})$ y todo $x\in U$, $s(x)$ \sref[0]{0.5.5.1} se identifica con un elemento de $k$; así, $s$ determina una aplicación $\overline{s}:x\mapsto s(x)$ de $U$ en $k$, es decir, una sección sobre $U$ del haz $\sh{A}(X')$ de gérmenes de aplicaciones de $X'$ en $k$. Como $h_U:s\mapsto\overline{s}$ es evidentemente un homomorfismo de anillos $\Gamma(U,\sh{O}_{X'})\to\Gamma(U,\sh{A}(X'))$ y conmuta con la restricción a un abierto $V\subset U$, los $h_U$ definen un homomorfismo de haces de anillos $h:\sh{O}_{X'}\to\sh{A}(X')$.
Si se toma $U$ ultraafín, $\operatorname{Spm}(A)$, donde $A$ es un anillo de Jacobson, decir que $\overline{s}=0$ significa que, para todo ideal maximal $\mathfrak{m}$ de $A$, $s$ pertenece a $\mathfrak{m}$, o equivalentemente que $s$ pertenece al \emph{radical} de $A$; pero, como $A$ es un anillo de Jacobson, su radical es su nilradical. Por consiguiente, $h_U$ es inyectivo si y solo si $A$ es reducido.
\end{env}

\begin{env}[10.10.4]
\label{IV.10.10.4}
Se dice que el espacio $k$-prealgebraico $X'=S(X)$ es \emph{reducido} si $X$ es reducido. Como el conjunto de puntos $x\in X$ en los que $X$ es reducido es abierto \sref[0]{0.5.2.2}, su complementario, si no es vacío, contiene al menos un punto cerrado \sref{IV.5.1.11}; por tanto, equivale a decir que $X'$ es reducido o que cada uno de sus anillos locales $\sh{O}_x$ ($x\in X'$) es reducido. Acabamos de ver en \sref{IV.10.10.3} que, para que el homomorfismo $h:\sh{O}_{X'}\to\sh{A}(X')$ sea \emph{inyectivo}, es necesario y suficiente que $X'$ sea \emph{reducido}.
En (FAC), Serre se limita de hecho a los espacios prealgebraicos \emph{reducidos}, lo que le permite definir $\sh{O}_{X'}$ como un \emph{subhaz} de $\sh{A}(X')$.
Obsérvese que, si $X'$ e $Y'$ son dos espacios $k$-prealgebraicos reducidos, lo mismo ocurre con $X'\times_k Y'$: en efecto, esto equivale a decir que, si $A$ y $B$ son dos $k$-álgebras reducidas de tipo finito, $A\otimes_k B$ es también reducido; pero acabamos de ver que los radicales de $A$ y $B$ son $0$, y como $k$ es algebraicamente cerrado, $A$ y $B$ son álgebras ``separables'' sobre $k$ en el sentido de Bourbaki (\emph{Alg.}, cap.~VIII, \S\,7, n$^\circ$\,5, prop.~5); por tanto, $A\otimes_k B$ carece de radical (\emph{loc.\ cit.}, n$^\circ$\,6, cor.~3 del th.~3), y como es un anillo de Jacobson, es reducido.
Sin embargo, si $Z'$ es un tercer espacio prealgebraico reducido sobre $k$, el ``producto fibrado'' $X'\times_{Z'}Y'$ de dos espacios prealgebraicos reducidos sobre $Z'$ no es reducido en general, lo que implica que la categoría de estos espacios resulta insuficiente en muchas cuestiones (especialmente en la teoría de grupos algebraicos).
Pero, como hemos visto anteriormente, se puede conservar el lenguaje de Serre sin limitarse, como hace este último (quien además solo considera espacios \emph{prealgebraicos} cuasicompactos), al caso de los espacios prealgebraicos reducidos.
\end{env}

\oldpage[IV-3]{116}

\begin{env}[10.10.5]
\label{IV.10.10.5}
Por último, también se pueden considerar ultrapreesquemas sobre un cuerpo arbitrario $k$
conservando un lenguaje próximo al de Serre e introduciendo, como en Weil, una
extensión algebraicamente cerrada fija $K$ de $k$ (elegida suficientemente grande---por ejemplo,
de grado de trascendencia infinito sobre $k$---para disponer de suficientes ``puntos genéricos''
en el sentido de Weil).
A todo preesquema $X$ localmente de tipo finito sobre $k$ se asocia entonces el conjunto
\[
 S_K(X)=S(X\otimes_k K)
\]
de puntos de $X$ con valores en $K$.
Existe una aplicación canónica $j:S_K(X)\to X$, que es un cuasihomeomorfismo cuando
se dota a $S_K(X)$ de la topología imagen inversa de la de $X$ por $j$; al dotar a $S_K(X)$
del haz de $k$-álgebras $j^*(\sh{O}_X)$, se obtiene una subcategoría de la categoría
de espacios $k$-anillados en anillos locales, que se podría llamar categoría de los
\emph{espacios $(k,K)$-prealgebraicos}.
Todavía se pueden definir productos en ella, y los resultados de \sref{IV.10.10.3} y
\sref{IV.10.10.4} se pueden generalizar (sustituyendo aquí $\sh{A}(X')$ por el haz
$\sh{A}_K(X')$ de gérmenes de aplicaciones de $X'$ en $K$).
Sin embargo, este punto de vista atribuye un papel artificial a un cuerpo de extensión
de $k$ elegido arbitrariamente, y lo mencionamos únicamente para descartarlo.
\end{env}

\setcounter{section}{10}
\section{Propiedades topológicas de los morfismos planos de presentación finita. Criterios de planitud}
\label{section:IV.11}

\oldpage[IV-3]{116}

Mientras que en la \S~2 consideramos enunciados relativos a la planitud que no dependen
de ninguna hipótesis de finitud, y la \S~6 estudia la noción de planitud en el
marco de los preesquemas localmente noetherianos (pero sin imponer una hipótesis de finitud
sobre los morfismos), la presente sección está dedicada a la noción de
$f$-planitud cuando el morfismo $f:X\to Y$ es \emph{localmente de presentación
finita}. El interés de la noción de morfismo plano de presentación finita
radica en que parece expresar técnicamente, de la manera más adecuada, la noción
intuitiva de una ``familia de preesquemas algebraicos parametrizada por un esquema
$Y$'', cuyo estudio constituye uno de los objetos principales de la
Geometría Algebraica. Además, aun cuando inicialmente solo interesase el caso
de un esquema base noetheriano, resulta indispensable, por ciertas razones
técnicas (por ejemplo, para algunas aplicaciones de la teoría del
``descenso'', que lleva a introducir esquemas no necesariamente
noetherianos), no limitarse a ese caso al tratar problemas esencialmente
relativos ligados a morfismos localmente de presentación finita. Seguiremos
sistemáticamente este principio, ya respaldado por los resultados de las \S\S~8 y 9,
en todo el resto de este Capítulo e incluso en todo el resto de este Tratado,
aunque ello pueda sacrificar ocasionalmente la sencillez de ciertas demostraciones
que las hipótesis noetherianas permiten a veces aligerar.\footnote{Este principio
se inspira asimismo en la necesidad de dar cabida, como ``espacios de parámetros''
para familias de esquemas algebraicos, a espacios anillados arbitrarios (e incluso
a ``topos'' anillados arbitrarios), para los cuales ya no cabe considerar en general
hipótesis noetherianas. Parece bastante claro que esta nueva ampliación de la
Geometría Algebraica no podrá eludirse por mucho más tiempo, y conviene desarrollar
desde ahora las nociones y técnicas de naturaleza ``relativa'' de la teoría de
esquemas de modo que puedan adaptarse, prácticamente sin cambios, a este marco más
general.} En la presente sección, esto nos lleva a retomar, en el contexto de la
``presentación finita'' (en particular en el n$^\circ$~3), ciertos enunciados sobre
planitud ya obtenidos en el contexto noetheriano. La herramienta técnica esencial
para la reducción al caso noetheriano es el teorema de compatibilidad de la
planitud con los límites proyectivos de preesquemas \sref{IV.11.2.6}, que completa
los resultados generales de la \S~8. Demostramos también de paso en \sref{IV.11.3.1}
un resultado utilizado con frecuencia

\oldpage[IV-3]{117}
más adelante, según el cual el conjunto de puntos de planitud de un morfismo
localmente de presentación finita es abierto.

En los n$^\circ$s~4 a 8 estudiamos la cuestión del ``descenso'' de la
planitud, que consiste en hallar condiciones útiles sobre un morfismo de cambio
de base $Y'\to Y$ (no necesariamente plano) bajo las cuales pueda concluirse,
si $X\times_YY'$ es plano sobre $Y'$, que $X$ es plano sobre $Y$. Estos resultados,
más técnicos que los de los n$^\circ$s~1 a 3, se utilizan con menor frecuencia
en lo que sigue; desempeñan, no obstante, un papel importante en las técnicas de
construcción no proyectivas del capítulo siguiente. El único resultado de los
n$^\circ$s~4 a 8 utilizado posteriormente en el Capítulo IV es el criterio valuativo
de planitud \sref{IV.11.8}, que se aplicará en \sref{IV.15.2}.

Por último, los n$^\circ$s~9 y 10 están dedicados al estudio de una noción que
refina, en la teoría de esquemas, la densidad en sentido topológico: la noción de
una familia de subpreesquemas \emph{esquemáticamente densa} en un preesquema dado,
y en particular al estudio del comportamiento de esta noción bajo cambio de base
(plano o arbitrario). Por el momento, esta noción se utiliza sobre todo en el
estudio de los esquemas en grupos.


\subsection{Lugares de planitud (caso noetheriano)}
\label{subsection:IV.11.1}

\begin{theorem}[11.1.1]
\label{IV.11.1.1}
Sea $\mathrm{Y}$ un preesquema localmente noetheriano, $f:\mathrm{X}\to\mathrm{Y}$ un morfismo localmente de tipo finito y $\sh{F}$ un $\sh{O}_\mathrm{X}$-Módulo coherente. Entonces el conjunto $\mathrm{U}$ de los $x\in\mathrm{X}$ tales que $\sh{F}$ es $f$-plano en el punto $x$ es abierto en $\mathrm{X}$.
\end{theorem}

\begin{proof}
Como la cuestión es evidentemente local sobre $\mathrm{X}$ y $\mathrm{Y}$, podemos suponer $\mathrm{Y}=\Spec(\mathrm{A})$, $\mathrm{X}=\Spec(\mathrm{B})$, con $\mathrm{A}$ noetheriano y $\mathrm{B}$ una A-álgebra de tipo finito. Se tiene entonces $\sh{F}=\widetilde{\mathrm{M}}$, donde $\mathrm{M}$ es un B-módulo de tipo finito. Aplicando el criterio $(\mathbf{0}_{\mathrm{III}},~9.2.6)$, basta por tanto demostrar la afirmación siguiente:

\medskip
\phantomsection\label{IV.11.1.1.1}
\noindent\textbf{(11.1.1.1)}
\emph{Sean $\mathrm{A}$ un anillo noetheriano, $\mathrm{B}$ una $\mathrm{A}$-álgebra de tipo finito, $\mathrm{M}$ un $\mathrm{B}$-módulo de tipo finito, $\mathfrak{q}$ un ideal primo de $\mathrm{B}$ y $\mathfrak{p}$ su imagen inversa en $\mathrm{A}$. Supongamos que $\mathrm{M}_\mathfrak{q}$ es un $\mathrm{A}_\mathfrak{p}$-módulo plano. Existe entonces $g\in\mathrm{B}-\mathfrak{q}$ tal que, para todo ideal primo $\mathfrak{q}'\supset\mathfrak{q}$ de $\mathrm{B}$ con $g\notin\mathfrak{q}'$, $\mathrm{M}_{\mathfrak{q}'}$ es un $\mathrm{A}_{\mathfrak{p}'}$-módulo plano, donde $\mathfrak{p}'$ denota la imagen inversa de $\mathfrak{q}'$ en $\mathrm{A}$ (equivale a lo mismo $(\mathbf{0}_\mathrm{I},~6.3.1)$ decir que $\mathrm{M}_{\mathfrak{q}'}$ es un $\mathrm{A}$-módulo plano).}
\medskip

Para ello, consideremos $\mathrm{B}_{\mathfrak{q}'}$ como una $\mathrm{A}$-álgebra; evidentemente
$\mathfrak{p}\mathrm{B}_{\mathfrak{q}'}\subset
\mathfrak{q}'\mathrm{B}_{\mathfrak{q}'}$. Por
$(\mathbf{0}_{\mathrm{III}},~10.2.2)$, el $\mathrm{A}$-módulo
$\mathrm{M}_{\mathfrak{q}'}$ es plano si y solo si
$\mathrm{M}_{\mathfrak{q}'}/\mathfrak{p}\mathrm{M}_{\mathfrak{q}'}$ es un
$(\mathrm{A}/\mathfrak{p})$-módulo plano y
\[
  \operatorname{Tor}_1^\mathrm{A}
  (\mathrm{M}_{\mathfrak{q}'},\mathrm{A}/\mathfrak{p})=0.
\]
Ahora bien, $\mathrm{M}_{\mathfrak{q}'}=
\mathrm{M}\otimes_{\mathrm{B}}\mathrm{B}_{\mathfrak{q}'}$, y
$\mathrm{B}_{\mathfrak{q}'}$ es plano sobre $\mathrm{B}$; por tanto
\[
  \mathrm{M}_{\mathfrak{q}'}/\mathfrak{p}\mathrm{M}_{\mathfrak{q}'}
    =(\mathrm{M}/\mathfrak{p}\mathrm{M})_{\mathfrak{q}'},
  \qquad
  \operatorname{Tor}_1^\mathrm{A}
  (\mathrm{M}_{\mathfrak{q}'},\mathrm{A}/\mathfrak{p})
    =\bigl(\operatorname{Tor}_1^\mathrm{A}
      (\mathrm{M},\mathrm{A}/\mathfrak{p})\bigr)_{\mathfrak{q}'}.
\]
La segunda igualdad se obtiene definiendo Tor a partir de una resolución
proyectiva de $\mathrm{A}/\mathfrak{p}$. Por la misma razón, puesto que exigimos
$g\notin\mathfrak{q}'$, el anillo $\mathrm{B}_{\mathfrak{q}'}$ es plano sobre
$\mathrm{B}_g$, y por consiguiente
\[
  (\mathrm{M}/\mathfrak{p}\mathrm{M})_{\mathfrak{q}'}
    =\bigl((\mathrm{M}/\mathfrak{p}\mathrm{M})_g
      \bigr)_{\mathfrak{q}'\mathrm{B}_g},
  \qquad
  \bigl(\operatorname{Tor}_1^\mathrm{A}
      (\mathrm{M},\mathrm{A}/\mathfrak{p})\bigr)_{\mathfrak{q}'}
    =\bigl(\bigl(\operatorname{Tor}_1^\mathrm{A}
      (\mathrm{M},\mathrm{A}/\mathfrak{p})\bigr)_g
      \bigr)_{\mathfrak{q}'\mathrm{B}_g},
\]
entendiendo que $\mathrm{M}/\mathfrak{p}\mathrm{M}$ y
$\operatorname{Tor}_1^\mathrm{A}
(\mathrm{M},\mathrm{A}/\mathfrak{p})$ se consideran
como $\mathrm{B}$-módulos. En virtud de $(\mathbf{0}_\mathrm{I},~6.3.2)$, quedamos
así reducidos a demostrar el lema siguiente.

\begin{lemma}[11.1.1.2]
\label{IV.11.1.1.2}
Bajo las condiciones de \emph{{\hyperlink{ega4.c04.sec.11-1-1-1}{(11.1.1.1)}}}, existe $g\in\mathrm{B}-\mathfrak{q}$ tal que:
\begin{enumerate}[label=\emph{\roman*)}]
  \item $(\mathrm{M}/\mathfrak{p}\mathrm{M})_g$ es un $(\mathrm{A}/\mathfrak{p})$-módulo plano;
  \item $(\operatorname{Tor}_1^\mathrm{A}(\mathrm{M},\mathrm{A}/\mathfrak{p}))_g=0$.
\end{enumerate}
\end{lemma}

\oldpage[IV-3]{118}

Por el teorema de planitud genérica $(6.9.1)$, aplicado al anillo íntegro
$\mathrm{A}/\mathfrak{p}$, a la
$(\mathrm{A}/\mathfrak{p})$-álgebra de tipo finito
$\mathrm{B}/\mathfrak{p}\mathrm{B}$ y al
$(\mathrm{B}/\mathfrak{p}\mathrm{B})$-módulo de tipo finito
$\mathrm{M}/\mathfrak{p}\mathrm{M}$, existe
$h\in\mathrm{A}-\mathfrak{p}$ tal que, si $\bar h$ denota su imagen canónica
en $\mathrm{A}/\mathfrak{p}$,
$(\mathrm{M}/\mathfrak{p}\mathrm{M})_{\bar h}$ es un
$(\mathrm{A}/\mathfrak{p})$-módulo plano. Por otra parte,
$\mathrm{M}_{\mathfrak q}$ es plano sobre $\mathrm{A}_{\mathfrak p}$ y,
por tanto, en virtud de $(\mathbf{0}_\mathrm{I},~6.3.1)$, es plano sobre $\mathrm{A}$.
En consecuencia,
\[
  \operatorname{Tor}_1^\mathrm{A}
  (\mathrm{M}_{\mathfrak q},\mathrm{A}/\mathfrak p)=0,
  \qquad\text{o, equivalentemente,}\qquad
  \bigl(\operatorname{Tor}_1^\mathrm{A}
  (\mathrm{M},\mathrm{A}/\mathfrak p)\bigr)_{\mathfrak q}=0.
\]
Como $\mathrm{A}$ y $\mathrm{B}$ son noetherianos,
$\operatorname{Tor}_1^\mathrm{A}
(\mathrm{M},\mathrm{A}/\mathfrak p)$ es un $\mathrm{B}$-módulo de tipo finito.
Así, $(\mathbf{0}_\mathrm{I},~5.2.2)$ proporciona un elemento
$g\in\mathrm{B}-\mathfrak q$ tal que
\[
  \bigl(\operatorname{Tor}_1^\mathrm{A}
  (\mathrm{M},\mathrm{A}/\mathfrak p)\bigr)_g=0.
\]
Además,
$(\mathrm{M}/\mathfrak p\mathrm{M})_{\bar h}
=(\mathrm{M}/\mathfrak p\mathrm{M})_h$, donde a la derecha el módulo se
considera como $\mathrm{A}$-módulo. Como
$(\mathrm{M}/\mathfrak p\mathrm{M})_h$ es también un $\mathrm{B}$-módulo,
$(\mathrm{M}/\mathfrak p\mathrm{M})_{hg}$ sigue siendo plano sobre
$\mathrm{A}/\mathfrak p$: es
$(\mathrm{M}/\mathfrak p\mathrm{M})_{\overline{hg}}$, donde
$\overline{hg}$ es la imagen canónica de $hg$ en
$\mathrm{B}/\mathfrak p\mathrm{B}$, y se aplica
$(\mathbf{0}_\mathrm{I},~6.3.2)$. Finalmente,
\[
  \bigl(\operatorname{Tor}_1^\mathrm{A}
  (\mathrm{M},\mathrm{A}/\mathfrak p)\bigr)_{hg}=0,
\]
y $hg\in\mathrm{B}-\mathfrak q$ porque
$h\in\mathrm{A}-\mathfrak p$. Esto demuestra el lema y el teorema.
\end{proof}

\begin{corollary}[11.1.2]
\label{IV.11.1.2}
Sea $\mathrm{Y}$ un preesquema localmente noetheriano, $f:\mathrm{X}\to\mathrm{Y}$ un morfismo localmente de tipo finito, $\sh{F}$ y $\sh{F}'$ dos $\sh{O}_\mathrm{X}$-Módulos coherentes, y $u:\sh{F}'\to\sh{F}$ un homomorfismo de $\sh{O}_\mathrm{X}$-Módulos. Supongamos que $\sh{F}$ es $f$-plano. Entonces el conjunto $\mathrm{U}$ de los $x\in\mathrm{X}$ tales que, poniendo $y=f(x)$, el homomorfismo $u_x\otimes 1:\sh{F}'_x\otimes_{\sh{O}_y}\boldsymbol{k}(y)\to\sh{F}_x\otimes_{\sh{O}_y}\boldsymbol{k}(y)$ es inyectivo, es abierto en $\mathrm{X}$.
\end{corollary}

\begin{proof}
Sean $\sh{N}$ (resp.\ $\sh{P}$) el núcleo (resp.\ el conúcleo) de $u$.
Apliquemos $(\mathbf{0}_{\mathrm{III}},~10.2.4)$\footnote{Damos aquí una demostración de
$(\mathbf{0}_{\mathrm{III}},~10.2.4)$ que no fue publicada en el
\emph{Alg\`ebre commutative} de N.~Bourbaki. En virtud de
$(\mathbf{0}_\mathrm{I},~6.1.2)$, basta demostrar que b) implica a).
Pongamos $\mathrm{P}=\operatorname{Im}(u)$,
$\mathrm{Q}=\operatorname{Coker}(u)$ y
$\mathrm{R}=\operatorname{Ker}(u)$. El compuesto
$\mathrm{M}\otimes k\to\mathrm{P}\otimes k\to\mathrm{N}\otimes k$ es
inyectivo, mientras que $\mathrm{M}\otimes k\to\mathrm{P}\otimes k$ es
sobreyectivo. Por tanto, $\mathrm{P}\otimes k\to\mathrm{N}\otimes k$ es
inyectivo y $\mathrm{M}\otimes k\to\mathrm{P}\otimes k$ es biyectivo.
La sucesión exacta $0\to\mathrm{P}\to\mathrm{N}\to\mathrm{Q}\to0$ proporciona
la sucesión exacta
\[
0=\operatorname{Tor}_1^\mathrm{A}(\mathrm{N},k)\to\operatorname{Tor}_1^\mathrm{A}(\mathrm{Q},k)\to\mathrm{P}\otimes k\to\mathrm{N}\otimes k\to\mathrm{Q}\otimes k\to 0
\]
y, como $\mathrm{P}\otimes k\to\mathrm{N}\otimes k$ es inyectivo,
$\operatorname{Tor}_1^\mathrm{A}(\mathrm{Q},k)=0$. Puesto que $\mathrm{Q}$ es un
$\mathrm{B}$-módulo de tipo finito,
$(\mathbf{0}_{\mathrm{III}},~10.2.2)$ muestra que es plano sobre
$\mathrm{A}$; entonces $\mathrm{P}$ es también plano sobre $\mathrm{A}$ por
$(\mathbf{0}_\mathrm{I},~6.1.2)$. A partir de la sucesión exacta
$0\to\mathrm{R}\to\mathrm{M}\to\mathrm{P}\to0$, el mismo resultado proporciona una
sucesión exacta
$0\to\mathrm{R}\otimes k\to\mathrm{M}\otimes k\to
\mathrm{P}\otimes k\to0$. Como la última aplicación es biyectiva,
$\mathrm{R}\otimes k=0$. Pero $\mathrm{B}$ es noetheriano, de modo que
$\mathrm{R}$ es un $\mathrm{B}$-módulo de tipo finito; el lema de Nakayama
da por tanto $\mathrm{R}=0$.}
a los anillos locales $\sh{O}_y$ y $\sh{O}_x$, y a los
$\sh{O}_x$-módulos $\sh{F}'_x$ y $\sh{F}_x$. La inyectividad de
$u_x\otimes1$ equivale a las dos condiciones
$\sh{N}_x=0$ y que $\sh{P}$ sea $f$-plano en $x$. Como
$\sh{N}$ y $\sh{P}$ son coherentes $(\mathbf{0}_\mathrm{I},~5.3.4)$,
el lugar donde $\sh{N}_x=0$ es abierto
$(\mathbf{0}_\mathrm{I},~5.2.2)$, y el lugar donde $\sh{P}$ es
$f$-plano es abierto por \sref{IV.11.1.1}. De aquí resulta la afirmación.

En particular:
\end{proof}

\begin{corollary}[11.1.3]
\label{IV.11.1.3}
Sea $\mathrm{Y}$ un preesquema localmente noetheriano,
$f:\mathrm{X}\to\mathrm{Y}$ un morfismo plano localmente de tipo finito, y
$g$ una sección de $\sh{O}_\mathrm{X}$ sobre $\mathrm{X}$. El conjunto de los
$x\in\mathrm{X}$ tales que $g_x\otimes 1$ es un no divisor de cero en
$\sh{O}_x\otimes_{\sh{O}_{f(x)}}\boldsymbol{k}(f(x))$ es abierto en
$\mathrm{X}$.
\end{corollary}

\begin{proof}
Basta aplicar \sref{IV.11.1.2} al endomorfismo de $\sh{O}_\mathrm{X}$ definido por $g$.
\end{proof}

\begin{corollary}[11.1.4]
\label{IV.11.1.4}
Sea $\mathrm{Y}$ un preesquema localmente noetheriano, $f:\mathrm{X}\to\mathrm{Y}$ un morfismo localmente de tipo finito, y $\sh{F}$ un $\sh{O}_\mathrm{X}$-Módulo coherente y $f$-plano. Sea $(g_i)_{1\leqslant i\leqslant n}$ una sucesión de secciones de $\sh{O}_\mathrm{X}$ sobre $\mathrm{X}$. El conjunto $\mathrm{U}$ de los $x\in\mathrm{X}$ tales que la sucesión $((g_i)_x\otimes 1)$ es $(\sh{F}_x\otimes_{\sh{O}_{f(x)}}\boldsymbol{k}(f(x)))$-regular es abierto en $\mathrm{X}$.
\end{corollary}

\begin{proof}
\oldpage[IV-3]{119}
Como $\sh{F}$ es $f$-plano, se deduce de
$(\mathbf{0},~15.1.16)$ que $\mathrm{U}$ es también el conjunto de los puntos
$x\in\mathrm{X}$ para los cuales la sucesión $((g_i)_x)$ es
$\sh{F}_x$-regular y el $\sh{O}_x$-módulo
\[
  \sh{G}_x,\qquad
  \sh{G}=\sh{F}\big/\Bigl(\sum_{i=1}^n g_i\sh{F}\Bigr),
\]
es $f$-plano. Los módulos $\sh{F}$ y $\sh{G}$ son coherentes, de modo que el
corolario se deduce de \sref{IV.11.1.1} y de
$(\mathbf{0},~15.2.4)$.
\end{proof}

\begin{corollary}[11.1.5]
\label{IV.11.1.5}
Sean $\mathrm{X}$, $\mathrm{Y}$ y $\mathrm{Z}$ tres preesquemas localmente noetherianos, $f:\mathrm{X}\to\mathrm{Y}$, $g:\mathrm{Y}\to\mathrm{Z}$ dos morfismos de tipo finito, y $\sh{F}$ un $\sh{O}_\mathrm{X}$-Módulo coherente. Entonces el conjunto $\mathrm{U}$ de los $y\in\mathrm{Y}$ tales que, para toda generalización $y'$ de $y$, $\sh{F}$ es $(g\circ f)$-plano en todos los puntos de $f^{-1}(y')$ (es decir,\ tales que $\sh{F}\otimes_\mathrm{Y}\Spec(\sh{O}_{\mathrm{Y},y})$ es plano respecto del morfismo $\mathrm{X}\times_\mathrm{Y}\Spec(\sh{O}_{\mathrm{Y},y})\to\mathrm{Z}$) es abierto en $\mathrm{Y}$.
\end{corollary}

\begin{proof}
Si $\mathrm{V}$ es el conjunto de los puntos $x\in\mathrm{X}$ en los que $\sh{F}$ es
$(g\circ f)$-plano, entonces $\mathrm{U}$ es el conjunto de los $y\in\mathrm{Y}$ tales
que $f^{-1}(y')\subset\mathrm{V}$ para toda generalización $y'$ de $y$.
Ahora bien, $\mathrm{V}$ es abierto por \sref{IV.11.1.1}, y por tanto localmente
construible en $\mathrm{X}$, y el conjunto $\mathrm{W}$ de los
$y\in\mathrm{Y}$ tales que $f^{-1}(y)\subset\mathrm{V}$ es
\[
  \mathrm{W}=\mathrm{Y}-f(\mathrm{X}-\mathrm{V}).
\]
Es, por consiguiente, localmente construible en $\mathrm{Y}$ por el teorema de
Chevalley $(1.8.4)$. Se deduce de
$(\mathbf{0}_{\mathrm{III}},~9.2.5)$ que los puntos de $\mathrm{U}$ son
los puntos interiores de $\mathrm{W}$, y por tanto que $\mathrm{U}$ es abierto.
\end{proof}

\begin{corollary}[11.1.6]
\label{IV.11.1.6}
Sea $\mathrm{A}$ un anillo noetheriano, $\mathrm{B}$ una $\mathrm{A}$-álgebra de tipo finito, $\mathrm{C}$ una $\mathrm{B}$-álgebra de tipo finito y $\mathrm{M}$ un $\mathrm{C}$-módulo de tipo finito. Entonces el conjunto de los $\mathfrak{q}\in\Spec(\mathrm{B})$ tales que $\mathrm{M}_\mathfrak{q}$ es un $\mathrm{A}$-módulo plano es abierto en $\Spec(\mathrm{B})$.
\end{corollary}

\begin{proof}
Teniendo en cuenta $(2.1.2)$, esto es consecuencia de \sref{IV.11.1.5} aplicado a $\mathrm{Z}=\Spec(\mathrm{A})$, $\mathrm{Y}=\Spec(\mathrm{B})$, $\mathrm{X}=\Spec(\mathrm{C})$, $\sh{F}=\widetilde{\mathrm{M}}$.

Los resultados de este número se formularán sin hipótesis noetherianas en \sref{IV.11.3}.
\end{proof}

\subsection{Planitud de un límite proyectivo de preesquemas}
\label{subsection:IV.11.2}

\oldpage[IV-3]{119}

\begin{env}[11.2.1]
\label{IV.11.2.1}
Sea $\mathrm{A}$ un anillo, sean $\mathrm{M}$ y $\mathrm{N}$ dos A-módulos, y sea $\mathrm{A}'$ una A-álgebra; pongamos $\mathrm{M}'=\mathrm{M}\otimes_\mathrm{A}\mathrm{A}'$, $\mathrm{N}'=\mathrm{N}\otimes_\mathrm{A}\mathrm{A}'$. Recordemos $(\textbf{III},~6.3.8)$ que, para todo $i$, se define un homomorfismo canónico de A-módulos
\begin{equation}
  \psi_i:\operatorname{Tor}_i^\mathrm{A}(\mathrm{M},\mathrm{N})\to\operatorname{Tor}_i^{\mathrm{A}'}(\mathrm{M}',\mathrm{N}')
  \label{IV.11.2.1.1}
\tag{11.2.1.1}
\end{equation}
de la manera siguiente: se considera una resolución por la izquierda de $\mathrm{M}$ mediante A-módulos libres
\begin{equation}
  \ldots\longrightarrow\mathrm{L}_{i+1}\xrightarrow{f_i}\mathrm{L}_i\longrightarrow\ldots\longrightarrow\mathrm{L}_0\xrightarrow{\varepsilon}\mathrm{M}\longrightarrow 0
  \label{IV.11.2.1.2}
\tag{11.2.1.2}
\end{equation}
de la cual se deduce, por producto tensorial con $\mathrm{A}'$, un complejo de $\mathrm{A}'$-módulos
\begin{equation}
  \ldots\longrightarrow\mathrm{L}'_{i+1}\xrightarrow{f'_i}\mathrm{L}'_i\longrightarrow\ldots\longrightarrow\mathrm{L}'_0\xrightarrow{\varepsilon'}\mathrm{M}'\longrightarrow 0
  \label{IV.11.2.1.3}
\tag{11.2.1.3}
\end{equation}
donde hemos puesto $\mathrm{L}'_i=\mathrm{L}_i\otimes_\mathrm{A}\mathrm{A}'$, $f'_i=f_i\otimes 1_{\mathrm{A}'}$, $\varepsilon'=\varepsilon\otimes 1_{\mathrm{A}'}$. Consideremos, por otra parte, una resolución por la izquierda de $\mathrm{M}'$ mediante $\mathrm{A}'$-módulos libres
\begin{equation}
  \ldots\longrightarrow\mathrm{L}''_{i+1}\xrightarrow{f''_i}\mathrm{L}''_i\longrightarrow\ldots\longrightarrow\mathrm{L}''_0\xrightarrow{\varepsilon''}\mathrm{M}'\longrightarrow 0
  \label{IV.11.2.1.4}
\tag{11.2.1.4}
\end{equation}
\end{env}

\oldpage[IV-3]{120}

Como los $\mathrm{L}'_i$ son $\mathrm{A}'$-módulos libres, sabemos $(\mathrm{M},~\mathrm{V},~1.1)$ que existen homomorfismos $\mathrm{A}'$-lineales $u_i$ que forman un diagrama conmutativo
\[
\xymatrix{
  \ldots \ar[r] & \mathrm{L}'_{i+1} \ar[r]^{f'_i} \ar[d]_{u_{i+1}} & \mathrm{L}'_i \ar[r] \ar[d]_{u_i} & \ldots \ar[r] & \mathrm{L}'_0 \ar[r]^{\varepsilon'} \ar[d]_{u_0} & \mathrm{M}' \ar[r] \ar[d]_{1_{\mathrm{M}'}} & 0 \\
  \ldots \ar[r] & \mathrm{L}''_{i+1} \ar[r]_{f''_i} & \mathrm{L}''_i \ar[r] & \ldots \ar[r] & \mathrm{L}''_0 \ar[r]_{\varepsilon''} & \mathrm{M}' \ar[r] & 0
}
\]

Al componer el morfismo de complejos
$u_\bullet:\mathrm{L}'_\bullet\to\mathrm{L}''_\bullet$ con el morfismo canónico
$\mathrm{L}_\bullet\to\mathrm{L}'_\bullet$, se obtiene un morfismo de
complejos de $\mathrm{A}$-módulos
$\mathrm{L}_\bullet\to\mathrm{L}''_\bullet$. Como
\[
  (\mathrm{L}_i\otimes_\mathrm{A}\mathrm{N})
    \otimes_\mathrm{A}\mathrm{A}'
  =\mathrm{L}'_i\otimes_{\mathrm{A}'}\mathrm{N}',
\]
obtenemos un morfismo de complejos
\[
  \mathrm{L}_\bullet\otimes_\mathrm{A}\mathrm{N}
  \longrightarrow
  \mathrm{L}''_\bullet\otimes_{\mathrm{A}'}\mathrm{N}'.
\]
Al pasar a la homología se obtienen los homomorfismos canónicos
\eref{IV.11.2.1.1}. Como los $u_i$ están determinados salvo
homotopía $(\mathrm{M},~\mathrm{V},~1.1)$, estos homomorfismos no dependen
ni de la elección de los $u_i$ ni de la elección de las resoluciones libres
$\mathrm{L}_\bullet$ y $\mathrm{L}''_\bullet$.

Como los $\operatorname{Tor}_i^{\mathrm{A}'}(\mathrm{M}',\mathrm{N}')$ son $\mathrm{A}'$-módulos, de \eref{IV.11.2.1.1} se deducen canónicamente homomorfismos $\mathrm{A}'$-lineales
\begin{equation}
  \varphi_i:\operatorname{Tor}_i^\mathrm{A}(\mathrm{M},\mathrm{N})\otimes_\mathrm{A}\mathrm{A}'\to\operatorname{Tor}_i^{\mathrm{A}'}(\mathrm{M}',\mathrm{N}').
  \label{IV.11.2.1.5}
\tag{11.2.1.5}
\end{equation}

Consideremos ahora dos homomorfismos de anillos
\[
  \rho:\mathrm{A}\to\mathrm{A}^{(1)},\quad \sigma:\mathrm{A}^{(1)}\to\mathrm{A}^{(2)},
\]
y su compuesto $\sigma\circ\rho:\mathrm{A}\to\mathrm{A}^{(2)}$; pongamos $\mathrm{M}^{(j)}=\mathrm{M}\otimes_\mathrm{A}\mathrm{A}^{(j)}$, $\mathrm{N}^{(j)}=\mathrm{N}\otimes_\mathrm{A}\mathrm{A}^{(j)}$ para $j=1,2$. Entonces el homomorfismo canónico compuesto
\[
  \operatorname{Tor}_i^\mathrm{A}(\mathrm{M},\mathrm{N})
  \longrightarrow
  \operatorname{Tor}_i^{\mathrm{A}^{(1)}}
    (\mathrm{M}^{(1)},\mathrm{N}^{(1)})
  \longrightarrow
  \operatorname{Tor}_i^{\mathrm{A}^{(2)}}
    (\mathrm{M}^{(2)},\mathrm{N}^{(2)})
\]
coincide con el homomorfismo canónico deducido de $\sigma\circ\rho$; esto resulta de que, si $\mathrm{L}_\bullet^{(j)}$ es una resolución libre de $\mathrm{M}^{(j)}$, el diagrama
\[
\xymatrix{
  \mathrm{L}_\bullet \ar[r] & \mathrm{L}_\bullet\otimes_\mathrm{A}\mathrm{A}^{(1)} \ar[r] \ar[d] & \mathrm{L}_\bullet^{(1)} \ar[d] \\
  & \mathrm{L}_\bullet\otimes_\mathrm{A}\mathrm{A}^{(2)} \ar[r] &
  \mathrm{L}_\bullet^{(1)}\otimes_{\mathrm{A}^{(1)}}\mathrm{A}^{(2)} \ar[r] & \mathrm{L}_\bullet^{(2)}
}
\]
es conmutativo.

\begin{env}[11.2.2]
\label{IV.11.2.2}
Con la notación de \sref{IV.11.2.1}, sea
$(\mathrm{A}_\alpha,g_{\beta\alpha})$ un sistema inductivo filtrante de
$\mathrm{A}$-álgebras y pongamos
\[
  \mathrm{M}_\alpha=\mathrm{M}\otimes_\mathrm{A}\mathrm{A}_\alpha,
  \qquad
  \mathrm{N}_\alpha=\mathrm{N}\otimes_\mathrm{A}\mathrm{A}_\alpha.
\]
Para cada $i$, los módulos
$\operatorname{Tor}_i^{\mathrm{A}_\alpha}
(\mathrm{M}_\alpha,\mathrm{N}_\alpha)$, con las aplicaciones canónicas proporcionadas
por \sref{IV.11.2.1}, forman un sistema inductivo. Pongamos
\[
  \mathrm{A}'=\varinjlim_\alpha\mathrm{A}_\alpha,\qquad
  \mathrm{M}'=\varinjlim_\alpha\mathrm{M}_\alpha
    =\mathrm{M}\otimes_\mathrm{A}\mathrm{A}',\qquad
  \mathrm{N}'=\varinjlim_\alpha\mathrm{N}_\alpha
    =\mathrm{N}\otimes_\mathrm{A}\mathrm{A}'.
\]
Si $g_\alpha:\mathrm{A}_\alpha\to\mathrm{A}'$ denota el homomorfismo
canónico, obtenemos aplicaciones canónicas
\[
  \psi_{i\alpha}:
  \operatorname{Tor}_i^{\mathrm{A}_\alpha}
    (\mathrm{M}_\alpha,\mathrm{N}_\alpha)
  \longrightarrow
  \operatorname{Tor}_i^{\mathrm{A}'}(\mathrm{M}',\mathrm{N}')
\]
que forman también un sistema inductivo. Completamos el resultado de
$(\mathrm{M},~\mathrm{V},~9.4^*)$ demostrando que
\begin{equation}
  \psi_i=\varinjlim\psi_{i\alpha}:\varinjlim_\alpha\operatorname{Tor}_i^{\mathrm{A}_\alpha}(\mathrm{M}_\alpha,\mathrm{N}_\alpha)\to\operatorname{Tor}_i^{\mathrm{A}'}(\mathrm{M}',\mathrm{N}')
  \label{IV.11.2.2.1}
\tag{11.2.2.1}
\end{equation}
son \emph{isomorfismos de $\mathrm{A}'$-módulos}. Para ello procedemos como en $(\mathrm{M},~\mathrm{V},~9.5^*)$, asociando a cada $\mathrm{M}_\alpha$ su \emph{resolución libre canónica}. Teniendo en cuenta la exactitud del funtor $\varinjlim$, esto equivale a demostrar el
\end{env}

\oldpage[IV-3]{121}

\begin{lemma}[11.2.2.2]
\label{IV.11.2.2.2}
Sea $(\mathrm{A}_\alpha,g_{\beta\alpha})$ un sistema inductivo filtrante
de anillos y $(\mathrm{M}_\alpha,h_{\beta\alpha})$ un sistema inductivo
filtrante de conjuntos. Pongamos
$\mathrm{A}'=\varinjlim\mathrm{A}_\alpha$ y
$\mathrm{M}'=\varinjlim\mathrm{M}_\alpha$, y sean
$g_\alpha:\mathrm{A}_\alpha\to\mathrm{A}'$ y
$h_\alpha:\mathrm{M}_\alpha\to\mathrm{M}'$ las aplicaciones canónicas. Para
cada $\alpha$, sea $\mathrm{F}(\mathrm{M}_\alpha)$ el
$\mathrm{A}_\alpha$-módulo de las combinaciones lineales formales de elementos de
$\mathrm{M}_\alpha$; análogamente, sea $\mathrm{F}(\mathrm{M}')$ el
$\mathrm{A}'$-módulo de las combinaciones lineales formales de elementos de
$\mathrm{M}'$. Si
$h_{\beta\alpha}^0:\mathrm{F}(\mathrm{M}_\alpha)\to
\mathrm{F}(\mathrm{M}_\beta)$ (para $\alpha\leqslant\beta$) y
$h'_\alpha:\mathrm{F}(\mathrm{M}_\alpha)\to\mathrm{F}(\mathrm{M}')$
son los homomorfismos $\mathrm{A}_\alpha$-lineales inducidos respectivamente por
$h_{\beta\alpha}$ y $h_\alpha$, entonces
$(\mathrm{F}(\mathrm{M}_\alpha),h^0_{\beta\alpha})$ es un sistema inductivo
de $\mathrm{A}_\alpha$-módulos y $(h'_\alpha)$ es un sistema inductivo
de homomorfismos. Además,
\[
  h''=\varinjlim h'_\alpha:\varinjlim\mathrm{F}(\mathrm{M}_\alpha)\to\mathrm{F}(\mathrm{M}')=\mathrm{F}(\varinjlim\mathrm{M}_\alpha)
\]
\emph{es un isomorfismo.}
\end{lemma}

\begin{proof}
Para la demostración, véase Bourbaki, \emph{Alg.}, cap.~II, 3\textsuperscript{e}~\'{e}d., \textsection\,6, n\textsuperscript{o}\,6, cor.\ de la prop.~10.
\end{proof}

\begin{env}[11.2.3]
\label{IV.11.2.3}
Retomemos la notación de \sref{IV.11.2.1} y consideremos el caso $i=1$.
Pongamos
$\mathrm{T}=\operatorname{Tor}_1^\mathrm{A}(\mathrm{M},\mathrm{N})$,
$\mathrm{T}'=\mathrm{T}\otimes_\mathrm{A}\mathrm{A}'$, y
$\mathrm{T}''=\operatorname{Tor}_1^{\mathrm{A}'}
(\mathrm{M}',\mathrm{N}')$. Entonces $\psi_1$ es el homomorfismo
\[
  \frac{\operatorname{Ker}(f_0\otimes1_\mathrm{N})}
       {\operatorname{Im}(f_1\otimes1_\mathrm{N})}
  \longrightarrow
  \frac{\operatorname{Ker}(f''_0\otimes1_{\mathrm{N}'})}
       {\operatorname{Im}(f''_1\otimes1_{\mathrm{N}'})}
\]
inducido en los cocientes por la restricción
\[
  \operatorname{Ker}(f_0\otimes1_\mathrm{N})
  \longrightarrow
  \operatorname{Ker}(f''_0\otimes1_{\mathrm{N}'})
\]
de
$u_1\otimes1:\mathrm{L}_1\otimes_\mathrm{A}\mathrm{N}\to
\mathrm{L}''_1\otimes_{\mathrm{A}'}\mathrm{N}'$.

Pongamos $\mathrm{R}=\operatorname{Ker}(\varepsilon)$, $\mathrm{R}''=\operatorname{Ker}(\varepsilon'')$, de modo que se tienen las sucesiones exactas
\begin{equation}
  0\longrightarrow\mathrm{R}\xrightarrow{j}\mathrm{L}_0\xrightarrow{\varepsilon}\mathrm{M}\longrightarrow 0
  \quad\text{y}\quad
  0\longrightarrow\mathrm{R}''\xrightarrow{j''}\mathrm{L}''_0\xrightarrow{\varepsilon''}\mathrm{M}'\longrightarrow 0,
\end{equation}
de las cuales se deducen las sucesiones exactas de homología
\begin{equation}
  0=\operatorname{Tor}_1^\mathrm{A}(\mathrm{L}_0,\mathrm{N})\to\mathrm{T}\xrightarrow{\partial}\mathrm{R}\otimes_\mathrm{A}\mathrm{N}\xrightarrow{j\otimes 1}\mathrm{L}_0\otimes_\mathrm{A}\mathrm{N}\xrightarrow{\varepsilon\otimes 1}\mathrm{M}\otimes_\mathrm{A}\mathrm{N}\to 0
  \label{IV.11.2.3.1}
\tag{11.2.3.1}
\end{equation}
\begin{equation}
  0=\operatorname{Tor}_1^{\mathrm{A}'}(\mathrm{L}''_0,\mathrm{N}')\to\mathrm{T}''\xrightarrow{\partial''}\mathrm{R}''\otimes_{\mathrm{A}'}\mathrm{N}'\xrightarrow{j''\otimes 1}\mathrm{L}''_0\otimes_{\mathrm{A}'}\mathrm{N}'\xrightarrow{\varepsilon''\otimes 1}\mathrm{M}'\otimes_{\mathrm{A}'}\mathrm{N}'\to 0
  \label{IV.11.2.3.2}
\tag{11.2.3.2}
\end{equation}

Por otra parte, existe un homomorfismo de $\mathrm{A}'$-módulos
\[
  v:\mathrm{R}'=\operatorname{Ker}(\varepsilon)\otimes_\mathrm{A}\mathrm{A}'\longrightarrow\operatorname{Ker}(\varepsilon')\xrightarrow{u_0}\operatorname{Ker}(\varepsilon'')=\mathrm{R}''.
\]
\end{env}

\oldpage[IV-3]{122}

Demostremos que el diagrama
\begin{equation}
\xymatrix{
  & \mathrm{T}' \ar[r]^-{\partial\otimes 1} \ar[d]_{\varphi_1} & \mathrm{R}'\otimes_{\mathrm{A}'}\mathrm{N}' \ar[r] \ar[d]_{v\otimes 1} & \mathrm{L}'_0\otimes_{\mathrm{A}'}\mathrm{N}' \ar[r] \ar[d]_{u_0\otimes 1} & \mathrm{M}'\otimes_{\mathrm{A}'}\mathrm{N}' \ar[r] \ar@{=}[d] & 0 \\
  0 \ar[r] & \mathrm{T}'' \ar[r]_{\partial''} & \mathrm{R}''\otimes_{\mathrm{A}'}\mathrm{N}' \ar[r] & \mathrm{L}''_0\otimes_{\mathrm{A}'}\mathrm{N}' \ar[r] & \mathrm{M}'\otimes_{\mathrm{A}'}\mathrm{N}' \ar[r] & 0
}
  \label{IV.11.2.3.3}
\tag{11.2.3.3}
\end{equation}
es \emph{conmutativo}. Para ello se comprueba inmediatamente $(\mathrm{M},~\mathrm{IV},~1)$ que el homomorfismo $\partial:\mathrm{T}\to\mathrm{R}\otimes_\mathrm{A}\mathrm{N}$ se obtiene (en el caso presente), por paso al cociente, del homomorfismo $w:\operatorname{Ker}(f_0\otimes 1_\mathrm{N})\to\mathrm{R}\otimes_\mathrm{A}\mathrm{N}$, restricción del homomorfismo $g_0\otimes 1_\mathrm{N}:\mathrm{L}_1\otimes_\mathrm{A}\mathrm{N}\to\mathrm{R}\otimes_\mathrm{A}\mathrm{N}$, donde $g_0:\mathrm{L}_1\to\mathrm{R}$ es tal que $f_0=j\circ g_0$; y análogamente para $\partial''$. Basta, por tanto, ver que el diagrama
\[
\xymatrix{
  \operatorname{Ker}(f_0\otimes 1_\mathrm{N}) \ar[r]^{w} \ar[d] & \mathrm{R}\otimes_\mathrm{A}\mathrm{N} \ar[r] & \mathrm{R}'\otimes_{\mathrm{A}'}\mathrm{N}'=(\mathrm{R}\otimes_\mathrm{A}\mathrm{N})\otimes_\mathrm{A}\mathrm{A}' \ar[d]_{v\otimes 1} \\
  \operatorname{Ker}(f''_0\otimes 1_{\mathrm{N}'}) \ar[rr] & & \mathrm{R}''\otimes_{\mathrm{A}'}\mathrm{N}'
}
\]
es conmutativo, lo que resulta de la conmutatividad del diagrama
\[
\xymatrix{
  \mathrm{L}_1 \ar[r] \ar[d] & \mathrm{R} \ar[d] \\
  \mathrm{L}''_1 \ar[r] & \mathrm{R}''
}
\]

\begin{lemma}[11.2.4]
\label{IV.11.2.4}
Sea $\mathrm{A}$ un anillo, $\mathrm{C}$ una $\mathrm{A}$-álgebra, $\mathrm{M}$ un $\mathrm{A}$-módulo, $\mathrm{N}$ un $\mathrm{C}$-módulo y $\mathrm{A}'$ una $\mathrm{A}$-álgebra. Pongamos $\mathrm{C}'=\mathrm{C}\otimes_\mathrm{A}\mathrm{A}'$, $\mathrm{M}'=\mathrm{M}\otimes_\mathrm{A}\mathrm{A}'$, $\mathrm{N}'=\mathrm{N}\otimes_\mathrm{A}\mathrm{A}'=\mathrm{N}\otimes_\mathrm{C}\mathrm{C}'$. Supongamos que $\mathrm{M}\otimes_\mathrm{A}\mathrm{N}$ es un $\mathrm{C}$-módulo plano. Entonces el homomorfismo canónico
\[
  \varphi_1:\operatorname{Tor}_1^\mathrm{A}(\mathrm{M},\mathrm{N})\otimes_\mathrm{A}\mathrm{A}'\to\operatorname{Tor}_1^{\mathrm{A}'}(\mathrm{M}',\mathrm{N}')
\]
(cf.\ \eref{IV.11.2.1.5}) \emph{es sobreyectivo}.
\end{lemma}

\begin{proof}
Conservemos la notación de \sref{IV.11.2.3}. La exactitud por la derecha del producto
tensorial muestra que
\[
  \mathrm{L}'_1\xrightarrow{f'_0}\mathrm{L}'_0
  \xrightarrow{\varepsilon'}\mathrm{M}'\longrightarrow0
\]
es exacta. Como $\mathrm{L}'_0$ y $\mathrm{L}'_1$ son
$\mathrm{A}'$-módulos libres, podemos suponer que
$\mathrm{L}''_0=\mathrm{L}'_0$ y
$\mathrm{L}''_1=\mathrm{L}'_1$, que $u_0$ y $u_1$ son las aplicaciones
identidad y que $f''_0=f'_0$. Como
$\operatorname{Ker}(\varepsilon')=\operatorname{Im}(f'_0)$ y
$\mathrm{R}''=\operatorname{Im}(f''_0)=\operatorname{Im}(f'_0)$, el
homomorfismo $v$ es sobreyectivo y, por tanto, también lo es $v\otimes1$. La demostración
quedará completa si la primera fila de
\eref{IV.11.2.3.3} es exacta, pues $v\otimes1$ es
sobreyectivo y $u_0\otimes1$ es inyectivo (Bourbaki,
\emph{Alg.\ comm.}, cap.~I, \S\,1, n.º~4, Cor.~2 de la Prop.~2). Para ello
utilizamos

\oldpage[IV-3]{123}

la hipótesis de que $\mathrm{M}\otimes_\mathrm{A}\mathrm{N}$ es un C-módulo plano. Poniendo $\mathrm{Q}=\operatorname{Ker}(\varepsilon\otimes 1)=\operatorname{Im}(j\otimes 1)$, la sucesión exacta \eref{IV.11.2.3.1} da las dos sucesiones exactas $0\to\mathrm{Q}\to\mathrm{L}_0\otimes_\mathrm{A}\mathrm{N}\to\mathrm{M}\otimes_\mathrm{A}\mathrm{N}\to 0$ y $\mathrm{T}\to\mathrm{R}\otimes_\mathrm{A}\mathrm{N}\to\mathrm{Q}\to 0$, cuyos homomorfismos son homomorfismos de \emph{C-módulos}; utilizando la hipótesis de planitud, se deduce de $(\mathbf{0}_\mathrm{I},~6.1.2)$ la sucesión exacta
\[
  0\to\mathrm{Q}\otimes_\mathrm{C}\mathrm{C}'\to(\mathrm{L}_0\otimes_\mathrm{A}\mathrm{N})\otimes_\mathrm{C}\mathrm{C}'\to(\mathrm{M}\otimes_\mathrm{A}\mathrm{N})\otimes_\mathrm{C}\mathrm{C}'\to 0
\]
y, por otra parte, como el producto tensorial es exacto por la derecha, se tiene la sucesión exacta
\[
  \mathrm{T}\otimes_\mathrm{C}\mathrm{C}'\to(\mathrm{R}\otimes_\mathrm{A}\mathrm{N})\otimes_\mathrm{C}\mathrm{C}'\to\mathrm{Q}\otimes_\mathrm{C}\mathrm{C}'\to 0
\]
de donde se obtiene finalmente la sucesión exacta
\[
  \mathrm{T}\otimes_\mathrm{C}\mathrm{C}'\to(\mathrm{R}\otimes_\mathrm{A}\mathrm{N})\otimes_\mathrm{C}\mathrm{C}'\to(\mathrm{L}_0\otimes_\mathrm{A}\mathrm{N})\otimes_\mathrm{C}\mathrm{C}'\to(\mathrm{M}\otimes_\mathrm{A}\mathrm{N})\otimes_\mathrm{C}\mathrm{C}'\to 0
\]
Pero, por definición, para todo C-módulo $\mathrm{P}$, $\mathrm{P}\otimes_\mathrm{C}\mathrm{C}'=\mathrm{P}\otimes_\mathrm{A}\mathrm{A}'$, de donde resulta la conclusión.
\end{proof}

\begin{lemma}[11.2.5]
\label{IV.11.2.5}
Sea $\mathrm{A}$ un anillo, $\mathfrak{J}$ un ideal de $\mathrm{A}$,
$\mathrm{B}$ una $\mathrm{A}$-álgebra, $\mathrm{M}$ un
$\mathrm{B}$-módulo y $\mathrm{A}\to\mathrm{A}'$ un homomorfismo de
anillos. Pongamos
$\mathfrak{J}'=\mathfrak{J}\mathrm{A}'$,
$\mathrm{B}'=\mathrm{B}\otimes_\mathrm{A}\mathrm{A}'$, y
$\mathrm{M}'=\mathrm{M}\otimes_\mathrm{A}\mathrm{A}'
=\mathrm{M}\otimes_\mathrm{B}\mathrm{B}'$. Sea
$\mathfrak{p}'$ un ideal primo de $\mathrm{B}'$ que contiene
$\mathfrak{J}'\mathrm{B}'$. Supongamos que se cumple una de las hipótesis
siguientes:
\begin{enumerate}[label=\emph{\alph*)}]
  \item $\mathfrak{J}$ es nilpotente.
  \item $\mathrm{A}'$ es noetheriano, $\mathrm{B}'$ es una $\mathrm{A}'$-álgebra de tipo finito y $\mathrm{M}'$ un $\mathrm{B}'$-módulo de tipo finito.
\end{enumerate}
Bajo estas condiciones, supongamos que se verifican las dos propiedades siguientes:
\begin{enumerate}[label=\emph{(\roman*)}]
  \item $\mathrm{M}\otimes_\mathrm{A}(\mathrm{A}/\mathfrak{J})$ es un $(\mathrm{A}/\mathfrak{J})$-módulo plano.
  \item El homomorfismo canónico compuesto
\[
  \operatorname{Tor}_1^\mathrm{A}(\mathrm{M},\mathrm{A}/\mathfrak{J})\xrightarrow{\psi_1}\operatorname{Tor}_1^{\mathrm{A}'}(\mathrm{M}',\mathrm{A}'/\mathfrak{J}')\xrightarrow{\theta}(\operatorname{Tor}_1^{\mathrm{A}'}(\mathrm{M}',\mathrm{A}'/\mathfrak{J}'))_{\mathfrak{p}'}
\]
(donde $\psi_1$ es el homomorfismo \eref{IV.11.2.1.1} y $\theta$ el homomorfismo canónico de un $\mathrm{B}'$-módulo en su localizado en $\mathfrak{p}'$) es nulo.
\end{enumerate}

Entonces $\mathrm{M}'_{\mathfrak{p}'}$ es un $\mathrm{A}'$-módulo plano.
\end{lemma}

\begin{proof}
Bajo la hipótesis b), $\mathrm{M}'_{\mathfrak{p}'}$ es un
$\mathrm{B}'_{\mathfrak{p}'}$-módulo de tipo finito,
$\mathrm{B}'_{\mathfrak{p}'}$ es una
$\mathrm{A}'$-álgebra noetheriana y
$\mathfrak{J}'\mathrm{B}'_{\mathfrak{p}'}$ está contenido en el radical
$\mathfrak{p}'\mathrm{B}'_{\mathfrak{p}'}$ de
$\mathrm{B}'_{\mathfrak{p}'}$. Bajo la hipótesis a),
$\mathfrak{J}'$ es nilpotente. Podemos, pues, aplicar el criterio de planitud
$(\mathbf{0}_{\mathrm{III}},~10.2.1)$ o
$(\mathbf{0}_{\mathrm{III}},~10.2.2)$, según se cumpla a) o b).
En primer lugar,
\[
\begin{split}
  \mathrm{M}'_{\mathfrak{p}'}\otimes_{\mathrm{A}'}
    (\mathrm{A}'/\mathfrak{J}')
  &=(\mathrm{M}'\otimes_{\mathrm{A}'}
    (\mathrm{A}'/\mathfrak{J}'))_{\mathfrak{p}'}\\
  &=\bigl((\mathrm{M}\otimes_\mathrm{A}
    (\mathrm{A}/\mathfrak{J}))
    \otimes_{\mathrm{A}/\mathfrak{J}}
    (\mathrm{A}'/\mathfrak{J}')\bigr)_{\mathfrak{p}'}.
\end{split}
\]
La hipótesis (i) implica, por tanto, en virtud de
$(\mathbf{0}_\mathrm{I},~6.2.1\text{ y }6.3.2)$, que este es un
$(\mathrm{A}'/\mathfrak{J}')$-módulo plano. Queda por demostrar que
\[
  \operatorname{Tor}_1^{\mathrm{A}'}
  (\mathrm{M}'_{\mathfrak{p}'},\mathrm{A}'/\mathfrak{J}')=0.
\]
Por la planitud de $\mathrm{B}'_{\mathfrak{p}'}$ sobre $\mathrm{B}'$, este
módulo es
\[
  \bigl(\operatorname{Tor}_1^{\mathrm{A}'}
  (\mathrm{M}',\mathrm{A}'/\mathfrak{J}')\bigr)_{\mathfrak{p}'}.
\]
La hipótesis (ii) dice que $\theta\circ\varphi_1$ es nulo. Por otra
parte, \sref{IV.11.2.4}, aplicado con
$\mathrm{C}=\mathrm{A}/\mathfrak{J}$ y $\mathrm{N}=\mathrm{C}$,
muestra, teniendo en cuenta (i), que $\varphi_1$ es sobreyectivo
(pues $\mathrm{A}'/\mathfrak{J}'=\mathrm{C}\otimes_\mathrm{A}
\mathrm{A}'$). Por consiguiente, $\theta$ es nulo. Como la imagen por $\theta$
de $\operatorname{Tor}_1^{\mathrm{A}'}
(\mathrm{M}',\mathrm{A}'/\mathfrak{J}')$ genera su localizado como
$\mathrm{B}'_{\mathfrak{p}'}$-módulo, dicho localizado es nulo.
\end{proof}

\begin{theorem}[11.2.6]
\label{IV.11.2.6}
Con la notación de $(8.5.1)$ y $(8.8.1)$, supongamos $\mathrm{S}_\alpha$ cuasicompacto, $\mathrm{X}_\alpha$ y $\mathrm{Y}_\alpha$ de presentación finita sobre $\mathrm{S}_\alpha$; sea $f_\alpha:\mathrm{X}_\alpha\to\mathrm{Y}_\alpha$ un $\mathrm{S}_\alpha$-morfismo y $\sh{F}_\alpha$ un $\sh{O}_{\mathrm{X}_\alpha}$-Módulo cuasicoherente de presentación finita.

\oldpage[IV-3]{124}

\begin{enumerate}[label=\emph{(\roman*)}]
  \item \emph{Sea $x$ un punto de $\mathrm{X}$ y $x_\lambda$ su proyección canónica en $\mathrm{X}_\lambda$. Para que $\sh{F}$ sea $f$-plano en el punto $x$, es necesario y suficiente que exista $\lambda\geqslant\alpha$ tal que $\sh{F}_\lambda$ sea $f_\lambda$-plano en el punto $x_\lambda$.}
  \item \emph{Para que $\sh{F}$ sea $f$-plano, es necesario y suficiente que exista $\lambda\geqslant\alpha$ tal que $\sh{F}_\lambda$ sea $f_\lambda$-plano.}
\end{enumerate}
\end{theorem}

\begin{proof}
Podemos suponer que $\mathrm{S}_\alpha=\mathrm{S}_0$; como $\mathrm{Y}_0$ es de presentación finita sobre $\mathrm{S}_0$, $\mathrm{Y}_0$ es cuasicompacto, y $f_0:\mathrm{X}_0\to\mathrm{Y}_0$ es un morfismo de presentación finita $(1.6.2,~(\mathrm{v}))$, por lo que podemos restringirnos también al caso $\mathrm{S}_0=\mathrm{Y}_0$. Además, por la cuasicompacidad de $\mathrm{S}_0$ y porque el conjunto de índices $\mathrm{L}$ es filtrante, el mismo razonamiento permite restringirse al caso en que $\mathrm{S}_0=\Spec(\mathrm{A}_0)$ es afín. Asimismo, $\mathrm{X}_0$ es también cuasicompacto, y el mismo razonamiento permite suponer $\mathrm{X}_0=\Spec(\mathrm{B}_0)$ afín; se tiene entonces $\sh{F}_0=\widetilde{\mathrm{M}}_0$, donde $\mathrm{M}_0$ es un $\mathrm{B}_0$-módulo de presentación finita, y el enunciado \sref{IV.11.2.6} equivale en este caso al siguiente (por $(\mathbf{0}_\mathrm{I},~6.3.1)$):

\begin{corollary}[11.2.6.1]
\label{IV.11.2.6.1}
Sea $\mathrm{A}_0$ un anillo,
$(\mathrm{A}_\lambda)_{\lambda\in\mathrm{L}}$ un sistema inductivo
filtrante de $\mathrm{A}_0$-álgebras, $\mathrm{B}_0$ una
$\mathrm{A}_0$-álgebra de presentación finita y $\mathrm{M}_0$ un
$\mathrm{B}_0$-módulo de presentación finita. Pongamos
\[
\begin{gathered}
  \mathrm{B}_\lambda
    =\mathrm{B}_0\otimes_{\mathrm{A}_0}\mathrm{A}_\lambda,\qquad
  \mathrm{M}_\lambda
    =\mathrm{M}_0\otimes_{\mathrm{A}_0}\mathrm{A}_\lambda
    =\mathrm{M}_0\otimes_{\mathrm{B}_0}\mathrm{B}_\lambda,\\
  \mathrm{A}=\varinjlim_\lambda\mathrm{A}_\lambda,\qquad
  \mathrm{B}=\varinjlim_\lambda\mathrm{B}_\lambda
    =\mathrm{B}_0\otimes_{\mathrm{A}_0}\mathrm{A},\\
  \mathrm{M}=\varinjlim_\lambda\mathrm{M}_\lambda
    =\mathrm{M}_0\otimes_{\mathrm{A}_0}\mathrm{A}
    =\mathrm{M}_0\otimes_{\mathrm{B}_0}\mathrm{B}.
\end{gathered}
\]
\begin{enumerate}[label=\emph{(\roman*)}]
  \item \emph{Sea $\mathfrak{p}$ un ideal primo de $\mathrm{B}$ y, para todo $\lambda$, sea $\mathfrak{p}_\lambda$ su imagen inversa en $\mathrm{B}_\lambda$. Para que $\mathrm{M}_\mathfrak{p}$ sea un $\mathrm{A}$-módulo plano, es necesario y suficiente que exista $\lambda$ tal que $(\mathrm{M}_\lambda)_{\mathfrak{p}_\lambda}$ sea un $\mathrm{A}_\lambda$-módulo plano.}
  \item \emph{Para que $\mathrm{M}$ sea un $\mathrm{A}$-módulo plano, es necesario y suficiente que exista $\lambda$ tal que $\mathrm{M}_\lambda$ sea un $\mathrm{A}_\lambda$-módulo plano.}
\end{enumerate}
\end{corollary}

Solo hay que demostrar que las condiciones son necesarias $(2.1.4)$. Procederemos en varias etapas.

\medskip
\noindent\textbf{I)} \emph{Reducción al caso en que los
$\mathrm{A}_\lambda$ son noetherianos.} --- Por $(8.9.1)$ existen un
subanillo $\mathrm{A}'_0\subset\mathrm{A}_0$ que es una
$\mathbf{Z}$-álgebra de tipo finito, una $\mathrm{A}'_0$-álgebra de tipo finito
$\mathrm{B}'_0$ y un $\mathrm{B}'_0$-módulo de tipo finito
$\mathrm{M}'_0$, tales que
\[
  \mathrm{B}_0=\mathrm{B}'_0\otimes_{\mathrm{A}'_0}\mathrm{A}_0,
  \qquad
  \mathrm{M}_0=\mathrm{M}'_0\otimes_{\mathrm{A}'_0}\mathrm{A}_0.
\]
Como
$\mathrm{B}_\lambda=\mathrm{B}'_0\otimes_{\mathrm{A}'_0}
\mathrm{A}_\lambda$ y
$\mathrm{M}_\lambda=\mathrm{M}'_0\otimes_{\mathrm{A}'_0}
\mathrm{A}_\lambda$, podemos sustituir
$\mathrm{A}_0,\mathrm{B}_0,\mathrm{M}_0$ por
$\mathrm{A}'_0,\mathrm{B}'_0,\mathrm{M}'_0$, considerando los
$\mathrm{A}_\lambda$ como $\mathrm{A}'_0$-álgebras. Podemos, por tanto,
suponer ya que $\mathrm{A}_0$ es noetheriano.

Sea $\mathrm{H}$ el conjunto de los pares $(\lambda,\mathrm{C}_\lambda)$,
donde $\mathrm{C}_\lambda$ es una $\mathrm{A}_0$-subálgebra de tipo finito
de $\mathrm{A}_\lambda$. Ordenemos $\mathrm{H}$ declarando
$(\lambda,\mathrm{C}_\lambda)\leqslant
(\mu,\mathrm{D}_\mu)$ si $\lambda\leqslant\mu$ y
$\varphi_{\mu\lambda}(\mathrm{C}_\lambda)\subset\mathrm{D}_\mu$.
Así $\mathrm{H}$ es filtrante: dos pares están dominados eligiendo
$\nu\geqslant\lambda,\mu$ y tomando como $\mathrm{E}_\nu$ la
$\mathrm{A}_0$-subálgebra engendrada por
$\varphi_{\nu\lambda}(\mathrm{C}_\lambda)$ y
$\varphi_{\nu\mu}(\mathrm{D}_\mu)$. Para
$\xi=(\lambda,\mathrm{C}_\lambda)$ pongamos
$\mathrm{A}_\xi=\mathrm{C}_\lambda$; las aplicaciones de transición restringidas
evidentes forman un sistema inductivo filtrante de $\mathrm{A}_0$-álgebras.
Pongamos
$\mathrm{B}_\xi=\mathrm{B}_0\otimes_{\mathrm{A}_0}\mathrm{A}_\xi$
y
$\mathrm{M}_\xi=\mathrm{M}_0\otimes_{\mathrm{A}_0}\mathrm{A}_\xi$.
Los anillos $\mathrm{A}_\xi$ son noetherianos, y la fórmula del doble límite
inductivo (Bourbaki, \emph{Alg.}, cap.~II, 3.ª ed., \S\,6, n.º~4,
Prop.~7) da
\[
  \varinjlim_{\mathrm H}\mathrm{A}_\xi=\mathrm{A},\qquad
  \varinjlim_{\mathrm H}\mathrm{B}_\xi=\mathrm{B},\qquad
  \varinjlim_{\mathrm H}\mathrm{M}_\xi=\mathrm{M}.
\]
Supongamos demostrado \sref{IV.11.2.6.1} para el sistema
$(\mathrm{A}_\xi)_{\xi\in\mathrm H}$, y sea $\mathfrak p$ un ideal primo
de $\mathrm B$ tal que $\mathrm M_{\mathfrak p}$ es plano sobre
$\mathrm A$. Existe entonces $\xi\in\mathrm H$ tal que
\oldpage[IV-3]{125}
$(\mathrm{M}_\xi)_{\mathfrak{p}_\xi}$ es plano sobre
$\mathrm{A}_\xi$, donde $\mathfrak p_\xi$ es la imagen inversa de
$\mathfrak p$. Escribamos $\xi=(\lambda,\mathrm{C}_\lambda)$. Como
$\mathfrak{p}_\lambda$ es la imagen inversa de $\mathfrak p$ en
$\mathrm{B}_\lambda$, $\mathfrak p_\xi$ es la imagen inversa de
$\mathfrak p_\lambda$ en $\mathrm B_\xi$, y
\[
  (\mathrm{M}_\lambda)_{\mathfrak p_\lambda}
  =(\mathrm{M}_\xi\otimes_{\mathrm{C}_\lambda}
    \mathrm{A}_\lambda)_{\mathfrak p_\lambda}
  =\bigl((\mathrm{M}_\xi)_{\mathfrak p_\xi}
    \otimes_{\mathrm{C}_\lambda}\mathrm{A}_\lambda
    \bigr)_{\mathfrak p_\lambda}.
\]
Así, $(\mathrm{M}_\lambda)_{\mathfrak p_\lambda}$ es plano sobre
$\mathrm{A}_\lambda$ por
$(\mathbf{0}_\mathrm{I},~6.2.1\text{ y }6.3.2)$. El enunciado (ii) se
trata de la misma manera. En lo sucesivo podemos, pues, suponer que todos los
$\mathrm{A}_\lambda$ son noetherianos, aunque el propio $\mathrm A$ no tenga
por qué serlo.

\medskip
\noindent\textbf{II)} \emph{Reducción del enunciado global (ii) al enunciado puntual (i).} --- Supongamos que $\sh{F}$ es $f$-plano. Para todo $\lambda$, sea $\mathrm{U}_\lambda$ el conjunto de los $x_\lambda\in\mathrm{X}_\lambda$ tales que $\sh{F}_\lambda$ es $f_\lambda$-plano en el punto $x_\lambda$; sabemos que $\mathrm{U}_\lambda$ es abierto en $\mathrm{X}_\lambda$, pues $\mathrm{S}_\lambda$ es noetheriano y $f_\lambda$ es de tipo finito \sref{IV.11.1.1}; sea $\mathrm{V}_\lambda=v_\lambda^{-1}(\mathrm{U}_\lambda)$ su imagen inversa en $\mathrm{X}$. Por hipótesis, para todo $x\in\mathrm{X}$, la proyección de $x$ en $\mathrm{X}_\lambda$ pertenece a $\mathrm{U}_\lambda$ para algún $\lambda$, lo que significa que $\mathrm{X}=\bigcup\mathrm{V}_\lambda$. Además, por $(2.1.4)$, si $\lambda\leqslant\mu$, se tiene $\mathrm{V}_\lambda\subset\mathrm{V}_\mu$; por tanto, como $\mathrm{X}$ es cuasicompacto, existe un índice $\mu$ tal que $\mathrm{X}=\mathrm{V}_\mu$. Como los $\mathrm{X}_\lambda$ son cuasicompactos, se deduce de $(8.3.4)$ que existe $\nu\geqslant\mu$ tal que $v_{\mu\nu}^{-1}(\mathrm{U}_\mu)=\mathrm{X}_\nu$; pero, por $(2.1.4)$, esto implica que $\sh{F}_\nu$ es $f_\nu$-plano.

\medskip
\noindent\textbf{III)} \emph{Fin de la demostración.} --- Queda por demostrar
(i), con $\mathrm S_0$ afín y noetheriano. Sea $y_0$ la
proyección de $x$ en $\mathrm S_0$. Por $(2.1.4)$ y
$(\mathbf I,~3.6.5)$ podemos sustituir $\mathrm S_0$ por
$\Spec(\sh O_{\mathrm S_0,y_0})$; podemos, pues, suponer que
$\mathrm A_0$ es un anillo local noetheriano con ideal maximal
$\mathfrak m_0=\mathfrak j_{y_0}$. Por definición,
$\mathfrak m_0$ es la imagen inversa del ideal primo
$\mathfrak p=\mathfrak j_x$ de $\mathrm B$, y
$\mathrm M_{\mathfrak p}$ es plano sobre $\mathrm A$. Por tanto,
\[
  \operatorname{Tor}_1^\mathrm A
  (\mathrm A/\mathfrak m_0\mathrm A,\mathrm M_{\mathfrak p})=0.
\]
Como los módulos Tor considerados son $\mathrm B$-módulos y
$\mathrm B_{\mathfrak p}$ es plano sobre $\mathrm B$, esto equivale
a
\[
  \bigl(\operatorname{Tor}_1^\mathrm A
  (\mathrm A/\mathfrak m_0\mathrm A,\mathrm M)\bigr)_{\mathfrak p}=0.
\]

El $\mathrm B_0$-módulo
$\operatorname{Tor}_1^{\mathrm A_0}
(\mathrm A_0/\mathfrak m_0,\mathrm M_0)$ es de tipo finito. En efecto, puede
calcularse a partir de una resolución de $\mathrm A_0/\mathfrak m_0$ mediante
$\mathrm A_0$-módulos libres finitos; al tomar el producto tensorial con $\mathrm M_0$ se obtiene un
complejo de $\mathrm B_0$-módulos de tipo finito, cuya homología es también
de tipo finito porque $\mathrm B_0$ es noetheriano. Elijamos generadores
$t_i^0$ $(1\leqslant i\leqslant m)$ de este módulo, y sean $t_i$
sus imágenes canónicas en
$\operatorname{Tor}_1^\mathrm A
(\mathrm A/\mathfrak m_0\mathrm A,\mathrm M)$. La igualdad de localización
precedente proporciona un elemento
$h\in\mathrm B-\mathfrak p$ tal que $ht_i=0$ para todo $i$.

Por \eref{IV.11.2.2.1},
\[
  \operatorname{Tor}_1^\mathrm A
  (\mathrm A/\mathfrak m_0\mathrm A,\mathrm M)
  =\varinjlim_\lambda
  \operatorname{Tor}_1^{\mathrm A_\lambda}
  (\mathrm A_\lambda/\mathfrak m_0\mathrm A_\lambda,
   \mathrm M_\lambda).
\]
Por consiguiente, para algún $\lambda$ existe un elemento
$h_\lambda\in\mathrm B_\lambda$ cuya imagen es $h$ tal que
$h_\lambda t_i^\lambda=0$ para todo $i$, donde $t_i^\lambda$ es la
imagen de $t_i^0$. Si $\mathfrak p_\lambda$ es la imagen inversa de
$\mathfrak p$ en $\mathrm B_\lambda$, entonces
$h_\lambda\notin\mathfrak p_\lambda$. Así, las imágenes de todos los
$t_i^0$ en
\[
  \bigl(\operatorname{Tor}_1^{\mathrm A_\lambda}
  (\mathrm A_\lambda/\mathfrak m_0\mathrm A_\lambda,
   \mathrm M_\lambda)\bigr)_{\mathfrak p_\lambda}
\]
son nulas, y la aplicación canónica de
$\operatorname{Tor}_1^{\mathrm A_0}
(\mathrm A_0/\mathfrak m_0,\mathrm M_0)$ a este localizado es nula.
Se satisfacen, por tanto, las hipótesis de \sref{IV.11.2.5}; la condición
(i) es automática porque $\mathrm A_0/\mathfrak m_0$ es un cuerpo.
Se sigue que $(\mathrm M_\lambda)_{\mathfrak p_\lambda}$ es plano
sobre $\mathrm A_\lambda$, lo que concluye la demostración.

\begin{corollary}[11.2.7]
\label{IV.11.2.7}
Sea $\mathrm{S}=\Spec(\mathrm{A})$ un esquema afín, $f:\mathrm{X}\to\mathrm{S}$ un morfismo, $\sh{F}$ un $\sh{O}_\mathrm{X}$-Módulo cuasicoherente y $x$ un punto de $\mathrm{X}$. Las condiciones siguientes son equivalentes:
\begin{enumerate}[label=\emph{\alph*)}]
  \item $f$ es un morfismo de presentación finita, $\sh{F}$ es un $\sh{O}_\mathrm{X}$-Módulo de presentación finita y $\sh{F}$ es $f$-plano en el punto $x$ (resp.\ $f$-plano).
\oldpage[IV-3]{126}
  \item Existe un esquema afín noetheriano $\mathrm{S}_0=\Spec(\mathrm{A}_0)$, un morfismo de tipo finito $f_0:\mathrm{X}_0\to\mathrm{S}_0$, un $\sh{O}_{\mathrm{X}_0}$-Módulo coherente $\sh{F}_0$ y un morfismo $\mathrm{S}\to\mathrm{S}_0$ tales que el S-preesquema $\mathrm{X}_0\otimes_{\mathrm{S}_0}\mathrm{S}$ es S-isomorfo a $\mathrm{X}$, que $\sh{F}$ es isomorfo a $\sh{F}_0\otimes_{\sh{O}_{\mathrm{S}_0}}\sh{O}_\mathrm{S}$ y que $\sh{F}_0$ es $f_0$-plano en el punto $x_0$, proyección de $x$ en $\mathrm{X}_0$ (resp.\ $f_0$-plano).
  \item Se satisfacen las condiciones de $b)$ y, además, $\mathrm{A}_0$ es una sub-$\mathbf{Z}$-álgebra de tipo finito de $\mathrm{A}$, correspondiendo el morfismo $\mathrm{S}\to\mathrm{S}_0$ a la inyección canónica $\mathrm{A}_0\to\mathrm{A}$.
\end{enumerate}
\end{corollary}

\begin{proof}
Es claro que $c)$ implica $b)$ y que $b)$ implica $a)$ por
$(2.1.4)$. Por otra parte, podemos considerar $\mathrm{A}$ como el límite
inductivo filtrante de sus $\mathbf{Z}$-subálgebras de tipo finito, y por $(8.9.1)$
existen una subálgebra de este tipo $\mathrm{A}_0$ y un morfismo de tipo finito
$f_0:\mathrm{X}_0\to\mathrm{S}_0$ tales que $\mathrm{X}$ es
$\mathrm{S}$-isomorfo a
$\mathrm{X}_0\times_{\mathrm{S}_0}\mathrm{S}$ y $\sh{F}$ es isomorfo
a $\sh{F}_0\otimes_{\sh{O}_{\mathrm{S}_0}}\sh{O}_{\mathrm{S}}$. Podemos,
por tanto, escribir
\[
  \mathrm{X}=\varprojlim_\lambda \mathrm{X}_\lambda,\qquad
  \mathrm{X}_\lambda=\mathrm{X}_0\otimes_{\mathrm{A}_0}\mathrm{A}_\lambda,
\]
donde $\mathrm{A}_\lambda$ recorre las
$\mathbf{Z}$-subálgebras de tipo finito de $\mathrm{A}$ que contienen $\mathrm{A}_0$, y
$\sh{F}=v_\lambda^*(\sh{F}_\lambda)$, donde
$v_\lambda:\mathrm{X}\to\mathrm{X}_\lambda$ es la proyección canónica y
$\sh{F}_\lambda=\sh{F}_0\otimes_{\mathrm{A}_0}\mathrm{A}_\lambda$. Se
deduce de \sref{IV.11.2.6} que existe $\lambda$ tal que
$\sh{F}_\lambda$ es $f_\lambda$-plano en
$x_\lambda=v_\lambda(x)$ (resp.\ $f_\lambda$-plano); el subanillo
$\mathrm{A}_\lambda$ de $\mathrm{A}$ satisface entonces la condición $c)$.
\end{proof}

\begin{proposition}[11.2.8]
\label{IV.11.2.8}
Sean $g:\mathrm{X}\to\mathrm{S}$ y $h:\mathrm{Y}\to\mathrm{S}$ dos morfismos de presentación finita, $f:\mathrm{X}\to\mathrm{Y}$ un S-morfismo y $\sh{F}$ un $\sh{O}_\mathrm{X}$-Módulo cuasicoherente de presentación finita; para todo $s\in\mathrm{S}$, sean $\mathrm{X}_s=g^{-1}(s)=\mathrm{X}\times_\mathrm{S}\Spec(\boldsymbol{k}(s))$, $\mathrm{Y}_s=h^{-1}(s)=\mathrm{Y}\times_\mathrm{S}\Spec(\boldsymbol{k}(s))$, $f_s$ el morfismo $f\times 1:\mathrm{X}_s\to\mathrm{Y}_s$ y $\sh{F}_s=\sh{F}\otimes_{\sh{O}_\mathrm{S}}\boldsymbol{k}(s)$. Entonces el conjunto $\mathrm{E}$ de los $s\in\mathrm{S}$ tales que $\sh{F}_s$ es $f_s$-plano es localmente construible.
\end{proposition}

\begin{proof}
La propiedad cuya constructibilidad queremos demostrar satisface la condición
$(9.2.1,~(\mathrm{i}))$, por $(2.2.13)$ y $(2.5.1)$. En virtud de
$(9.2.3)$, podemos, pues, restringirnos al caso en que $\mathrm{S}$ es
afín, noetheriano e íntegro, con punto genérico $\eta$, y queda por
demostrar que $\mathrm{E}$ o $\mathrm{S}-\mathrm{E}$ es un entorno
de $\eta$ en $\mathrm{S}$. Si
$\eta\in\mathrm{S}-\mathrm{E}$, esto resulta inmediatamente del Lema
$(9.4.7.1)$. Si, por el contrario, $\eta\in\mathrm{E}$, se deduce primero
de \sref{IV.11.2.6}, aplicado mediante el procedimiento de $(8.1.2,~a))$, que existe
un entorno abierto $\mathrm{U}$ de $\eta$ en $\mathrm{S}$ tal que
$\sh{F}|_{g^{-1}(\mathrm{U})}$ es plano sobre $h^{-1}(\mathrm{U})$; \emph{a
fortiori}, por $(2.1.4)$, $\sh{F}_s$ es $f_s$-plano para todo
$s\in\mathrm{U}$, lo que concluye la demostración.
\end{proof}

El teorema siguiente generaliza $(11.2.6.1,~(\mathrm{ii}))$.

\begin{theorem}[11.2.9]
\label{IV.11.2.9}
(Raynaud). --- Sea $\mathrm{A}_0$ un anillo, $(\mathrm{A}_\lambda)_{\lambda\in\mathrm{L}}$ un sistema inductivo filtrante de $\mathrm{A}_0$-álgebras, $\mathrm{B}_0$ una $\mathrm{A}_0$-álgebra de presentación finita, $\mathfrak{J}_0$ un ideal de tipo finito de $\mathrm{B}_0$ y $\mathrm{M}_0$ un $\mathrm{B}_0$-módulo de presentación finita. Pongamos $\mathrm{B}_\lambda=\mathrm{B}_0\otimes_{\mathrm{A}_0}\mathrm{A}_\lambda$, $\mathfrak{J}_\lambda=\mathfrak{J}_0\mathrm{B}_\lambda$, $\mathrm{M}_\lambda=\mathrm{M}_0\otimes_{\mathrm{A}_0}\mathrm{A}_\lambda=\mathrm{M}_0\otimes_{\mathrm{B}_0}\mathrm{B}_\lambda$, $\mathrm{A}=\varinjlim\mathrm{A}_\lambda$, $\mathrm{B}=\varinjlim\mathrm{B}_\lambda=\mathrm{B}_0\otimes_{\mathrm{A}_0}\mathrm{A}$, $\mathfrak{J}=\mathfrak{J}_0\mathrm{B}$ y $\mathrm{M}=\varinjlim\mathrm{M}_\lambda=\mathrm{M}_0\otimes_{\mathrm{A}_0}\mathrm{A}=\mathrm{M}_0\otimes_{\mathrm{B}_0}\mathrm{B}$. Para que $\operatorname{gr}^*_\mathfrak{J}(\mathrm{M})$ sea un $\mathrm{A}$-módulo plano, es necesario y suficiente que exista $\lambda$ tal que $\operatorname{gr}^*_{\mathfrak{J}_\lambda}(\mathrm{M}_\lambda)$ sea un $\mathrm{A}_\lambda$-módulo plano; los homomorfismos canónicos
\begin{equation}
  \begin{cases}
    \operatorname{gr}^*_{\mathfrak{J}_\lambda}(\mathrm{M}_\lambda)\otimes_{\mathrm{A}_\lambda}\mathrm{A} \to \operatorname{gr}^*_\mathfrak{J}(\mathrm{M}) \\
    \operatorname{gr}^*_{\mathfrak{J}_\lambda}(\mathrm{M}_\lambda)\otimes_{\mathrm{A}_\lambda}\mathrm{A}_\mu \to \operatorname{gr}^*_{\mathfrak{J}_\mu}(\mathrm{M}_\mu)
  \end{cases}
  \qquad\text{para }\lambda\leqslant\mu
  \label{IV.11.2.9.1}
\tag{11.2.9.1}
\end{equation}
son entonces biyectivos y $\operatorname{gr}^*_{\mathfrak{J}_\mu}(\mathrm{M}_\mu)$ es un $\mathrm{A}_\mu$-módulo plano para $\lambda\leqslant\mu$.
\end{theorem}

\medskip
\noindent\emph{Demostración.}
Demostremos primero que las condiciones son suficientes, lo que equivale a probar
la biyectividad de $(11.2.9.1)$. Esto se deduce del lema siguiente:

\begin{lemma}[11.2.9.2]
\label{IV.11.2.9.2}
Sea $\mathrm{A}$ un anillo, $\mathrm{B}$ una $\mathrm{A}$-álgebra, $\mathrm{M}$ un $\mathrm{B}$-módulo y $\mathfrak{J}$ un ideal de $\mathrm{B}$. Sea $\mathrm{A}'$ una $\mathrm{A}$-álgebra; pongamos $\mathrm{B}'=\mathrm{B}\otimes_\mathrm{A}\mathrm{A}'$, $\mathrm{M}'=\mathrm{M}\otimes_\mathrm{A}\mathrm{A}'=\mathrm{M}\otimes_\mathrm{B}\mathrm{B}'$ y $\mathfrak{J}'=\mathfrak{J}\mathrm{B}'$. Si $\operatorname{gr}^*_\mathfrak{J}(\mathrm{M})$ es un $\mathrm{A}$-módulo plano, el homomorfismo canónico
\[
  (\operatorname{gr}^*_\mathfrak{J}(\mathrm{M}))\otimes_\mathrm{A}\mathrm{A}'\to\operatorname{gr}^*_{\mathfrak{J}'}(\mathrm{M}')
\]
\emph{es bijetivo.}
\end{lemma}

\oldpage[IV-3]{127}

En efecto, por inducción sobre $k$, la hipótesis de que los $\mathfrak{J}^k\mathrm{M}/\mathfrak{J}^{k+1}\mathrm{M}$ son A-módulos planos implica primero, por $(\mathbf{0}_\mathrm{I},~6.1.2)$, que $\mathrm{M}/\mathfrak{J}^{k+1}\mathrm{M}$ es un A-módulo plano; además, por $(\mathbf{0}_\mathrm{I},~6.1.2)$, la sucesión
\[
  0\to(\mathfrak{J}^{k+1}\mathrm{M})\otimes_\mathrm{A}\mathrm{A}'\to\mathrm{M}\otimes_\mathrm{A}\mathrm{A}'\to(\mathrm{M}/\mathfrak{J}^{k+1}\mathrm{M})\otimes_\mathrm{A}\mathrm{A}'\to 0
\]
es entonces exacta; en otras palabras, $(\mathfrak{J}^{k+1}\mathrm{M})\otimes_\mathrm{A}\mathrm{A}'$ se identifica con su imagen canónica $\mathfrak{J}'^{k+1}\mathrm{M}'$ en $\mathrm{M}'=\mathrm{M}\otimes_\mathrm{A}\mathrm{A}'$. Por otra parte, siempre por $(\mathbf{0}_\mathrm{I},~6.1.2)$, la sucesión
\[
  0\to(\mathfrak{J}^{k+1}\mathrm{M})\otimes_\mathrm{A}\mathrm{A}'\to(\mathfrak{J}^k\mathrm{M})\otimes_\mathrm{A}\mathrm{A}'\to(\mathfrak{J}^k\mathrm{M}/\mathfrak{J}^{k+1}\mathrm{M})\otimes_\mathrm{A}\mathrm{A}'\to 0
\]
es exacta, lo que demuestra el lema.

Para demostrar la necesidad de las condiciones de \sref{IV.11.2.9}, procedemos en varios pasos.

\medskip
\noindent\textbf{I)} \emph{Reducción al caso en que los $\mathrm{A}_\lambda$
son noetherianos.} --- Se procede como en la reducción I) de
\sref{IV.11.2.6.1}, conservando la notación; basta con empezar por sustituir
$\mathrm{A}'_0$ por una $\mathrm{A}'_0$-subálgebra de tipo finito
$\mathrm{A}''_0$ de $\mathrm{A}_0$ tal que, si se pone
$\mathrm{B}''_0=\mathrm{B}'_0\otimes_{\mathrm{A}'_0}\mathrm{A}''_0$, hay
en $\mathrm{B}''_0$ un ideal finitamente generado $\mathfrak{J}''_0$ tal que
$\mathfrak{J}''_0\mathrm{B}_0=\mathfrak{J}_0$. Para ello, consideremos las
$\mathrm{A}'_0$-subálgebras de tipo finito $\mathrm{A}'_\alpha$ de
$\mathrm{A}_0$, que forman una familia filtrante; se tiene
$\mathrm{B}_0=\varinjlim\mathrm{B}'_\alpha$, donde
$\mathrm{B}'_\alpha=\mathrm{B}'_0\otimes_{\mathrm{A}'_0}\mathrm{A}'_\alpha$.
Por tanto, hay un índice $\beta$ tal que un conjunto finito de generadores de
$\mathfrak{J}_0$ está formado por las imágenes en $\mathrm{B}_0$ de elementos de
$\mathrm{B}'_\beta$. Tomemos $\mathrm{A}''_0=\mathrm{A}'_\beta$,
$\mathrm{B}''_0=\mathrm{B}'_\beta$, y como $\mathfrak{J}''_0$ el ideal
generado por estos elementos. Podemos, por tanto, suponer que $\mathrm{A}_0$
(y, por consiguiente, también $\mathrm{B}_0$) es noetheriano. Se definen entonces, como en
\emph{loc.\ cit.}, el conjunto filtrado $\mathrm{H}$ y los $\mathrm{A}_\xi$,
$\mathrm{B}_\xi$, $\mathrm{M}_\xi$ para $\xi\in\mathrm{H}$; pongamos también
$\mathfrak{J}_\xi=\mathfrak{J}''_0\mathrm{B}_\xi$ para todo
$\xi\in\mathrm{H}$. Supongamos demostrado que existe
$\xi=(\lambda,\mathrm{C}_\lambda)$ tal que
$\operatorname{gr}^*_{\mathfrak{J}_\xi}(\mathrm{M}_\xi)$ es un
$\mathrm{C}_\lambda$-módulo plano. Como
$\mathrm{M}_\lambda=\mathrm{M}_\xi\otimes_{\mathrm{C}_\lambda}
\mathrm{A}_\lambda$, se deduce de \sref{IV.11.2.9.2} que
$\operatorname{gr}^*_{\mathfrak{J}_\lambda}(\mathrm{M}_\lambda)$ es un
$\mathrm{A}_\lambda$-módulo plano.

\medskip
\noindent\textbf{II)} \emph{Lemas preliminares.}

\begin{lemma}[11.2.9.3]
\label{IV.11.2.9.3}
Sea $\mathrm{A}$ un anillo, sea
$\mathrm{B}=\bigoplus_{i\geqslant 0}\mathrm{B}_i$ una
$\mathrm{A}$-álgebra graduada en grados no negativos y sea
$\mathrm{M}=\bigoplus_{i\in\mathbf{Z}}\mathrm{M}_i$ un
$\mathrm{B}$-módulo graduado. Supongamos que $\mathrm{B}_0$ es un anillo local y
que cada $\mathrm{M}_i$ es un $\mathrm{B}_0$-módulo de tipo finito.
Para que $\mathrm{M}$ sea un $\mathrm{B}$-módulo de tipo finito, es necesario
y suficiente que, si $\mathfrak{q}$ es la imagen inversa en
$\mathrm{A}$ del ideal maximal $\mathfrak{m}$ de $\mathrm{B}_0$,
$\mathrm{M}\otimes_\mathrm{A}\boldsymbol{k}(\mathfrak{q})$ sea un
$(\mathrm{B}\otimes_\mathrm{A}\boldsymbol{k}(\mathfrak{q}))$-módulo
de tipo finito.
\end{lemma}

\oldpage[IV-3]{128}

\begin{proof}
La condición es evidentemente necesaria; demostremos que es
suficiente. Como $\mathrm{B}_0$ (y, por tanto, $\mathrm{B}$) es una
$\mathrm{A}_{\mathfrak{q}}$-álgebra, podemos sustituir $\mathrm{A}$ por
$\mathrm{A}_{\mathfrak{q}}$; en otras palabras, podemos suponer que
$\mathrm{A}$ es también un anillo local con ideal maximal
$\mathfrak{q}$. Por hipótesis, hay un entero $\mathrm{N}$ tal
que el homomorfismo canónico
\[
 \left(\bigoplus_{-\mathrm{N}\leqslant i\leqslant\mathrm{N}}
   \mathrm{M}_i\otimes_{\mathrm{B}_0}\mathrm{B}\right)
   \otimes_\mathrm{A}\boldsymbol{k}(\mathfrak{q})
 \longrightarrow
 \mathrm{M}\otimes_\mathrm{A}\boldsymbol{k}(\mathfrak{q})
\]
es sobreyectivo. Esto significa también que, para todo entero $n$, el
homomorfismo canónico de $\boldsymbol{k}(\mathfrak{q})$-módulos
\[
 \left(\bigoplus_{-\mathrm{N}\leqslant i\leqslant\mathrm{N}}
   \mathrm{M}_i\otimes_{\mathrm{B}_0}\mathrm{B}_{n-i}\right)
   \otimes_\mathrm{A}\boldsymbol{k}(\mathfrak{q})
 \longrightarrow
 \mathrm{M}_n\otimes_\mathrm{A}\boldsymbol{k}(\mathfrak{q})
\]
es sobreyectivo. Ahora bien, $\mathrm{M}_n$ es un
$\mathrm{B}_0$-módulo de tipo finito, $\mathrm{B}_0$ es local y
$\mathfrak{q}\mathrm{B}_0\subset\mathfrak{m}$; el lema de Nakayama
muestra, por tanto, que cada homomorfismo canónico
\[
 \bigoplus_{-\mathrm{N}\leqslant i\leqslant\mathrm{N}}
   \mathrm{M}_i\otimes_{\mathrm{B}_0}\mathrm{B}_{n-i}
 \longrightarrow \mathrm{M}_n
\]
es sobreyectivo, lo que demuestra la afirmación.
\end{proof}

\begin{corollary}[11.2.9.4]
\label{IV.11.2.9.4}
Bajo las hipótesis de \sref{IV.11.2.9.3}, supongamos además que cada $\mathrm{B}_i$ y cada $\mathrm{M}_i$ es un $\mathrm{B}_0$-módulo de presentación finita, y que $\mathrm{M}$ es un $\mathrm{A}$-módulo plano. Para que $\mathrm{M}$ sea un $\mathrm{B}$-módulo de presentación finita, es necesario y suficiente que $\mathrm{M}\otimes_\mathrm{A}\boldsymbol{k}(\mathfrak{q})$ sea un $(\mathrm{B}\otimes_\mathrm{A}\boldsymbol{k}(\mathfrak{q}))$-módulo de presentación finita.
\end{corollary}

\begin{proof}
Por \sref{IV.11.2.9.3}, si se cumple la condición del enunciado,
$\mathrm{M}$ es un $\mathrm{B}$-módulo de tipo finito. Existe, por tanto, un
$\mathrm{B}$-módulo graduado libre de tipo finito $\mathrm{L}$ (que posee,
pues, una base finita de elementos homogéneos) y un homomorfismo graduado
sobreyectivo de grado $0$,
\[
 u:\mathrm{L}\longrightarrow\mathrm{M}.
\]
Sea $\mathrm{R}=\operatorname{Ker}(u)$, que es un
$\mathrm{B}$-módulo graduado. Hay un número finito de enteros $m_j$
($1\leqslant j\leqslant r$) tales que, para todo entero $i$,
$\mathrm{R}_i$ es el núcleo de un homomorfismo sobreyectivo
\[
 \bigoplus_{1\leqslant j\leqslant r}\mathrm{B}_{i+m_j}
 \longrightarrow \mathrm{M}_i.
\]
Las hipótesis sobre los $\mathrm{M}_i$ y los $\mathrm{B}_i$ implican entonces
que $\mathrm{R}_i$ es un $\mathrm{B}_0$-módulo de tipo finito
(Bourbaki, \emph{Alg.\ comm.}, cap.~I,
\textsection\,2, n\textsuperscript{o}\,8, lema~9). Para demostrar que
$\mathrm{R}$ es un $\mathrm{B}$-módulo de tipo finito, observemos que,
por la hipótesis de planitud y $(\mathbf{0}_\mathrm{I},~6.1.2)$, la
sucesión
\[
 0\longrightarrow
 \mathrm{R}\otimes_\mathrm{A}\boldsymbol{k}(\mathfrak{q})
 \longrightarrow
 \mathrm{L}\otimes_\mathrm{A}\boldsymbol{k}(\mathfrak{q})
 \longrightarrow
 \mathrm{M}\otimes_\mathrm{A}\boldsymbol{k}(\mathfrak{q})
 \longrightarrow 0
\]
es exacta. La hipótesis implica, por tanto (Bourbaki,
\emph{loc.\ cit.}) que
$\mathrm{R}\otimes_\mathrm{A}\boldsymbol{k}(\mathfrak{q})$ es un
$(\mathrm{B}\otimes_\mathrm{A}\boldsymbol{k}(\mathfrak{q}))$-módulo
de tipo finito. Basta, por tanto, aplicar \sref{IV.11.2.9.3} a $\mathrm{R}$.
\end{proof}

\medskip
\noindent\textbf{III)} \emph{Reducción al caso en que los homomorfismos de transición $\varphi_{\mu\lambda}:\mathrm{A}_\lambda\to\mathrm{A}_\mu$ (para $\lambda\leqslant\mu$) son inyectivos.} --- Sea $\mathrm{A}'_\lambda$ la imagen de $\mathrm{A}_\lambda$ por el homomorfismo canónico $\varphi_\lambda:\mathrm{A}_\lambda\to\mathrm{A}$; es claro que los $\mathrm{A}'_\lambda$ forman un sistema inductivo de subanillos \emph{noetherianos} de $\mathrm{A}$ y que los homomorfismos de transición $\mathrm{A}'_\lambda\to\mathrm{A}'_\mu$ (para $\lambda\leqslant\mu$) son inyectivos. Pongamos $\mathrm{B}'_\lambda=\mathrm{B}_0\otimes_{\mathrm{A}_0}\mathrm{A}'_\lambda$, $\mathfrak{J}'_\lambda=\mathfrak{J}_0\mathrm{B}'_\lambda$ y $\mathrm{M}'_\lambda=\mathrm{M}_0\otimes_{\mathrm{A}_0}\mathrm{A}'_\lambda=\mathrm{M}_\lambda\otimes_{\mathrm{A}_\lambda}\mathrm{A}'_\lambda$; se sigue teniendo $\mathrm{B}=\varinjlim\mathrm{B}'_\lambda$ y $\mathrm{M}=\varinjlim\mathrm{M}'_\lambda$. Supongamos demostrado que existe $\lambda$ tal que, para $\mu\geqslant\lambda$, $\operatorname{gr}^*_{\mathfrak{J}'_\mu}(\mathrm{M}'_\mu)$ es un $\mathrm{A}'_\mu$-módulo plano. Sea $a_\lambda$ el núcleo del homomorfismo $\varphi_\lambda$, que es un ideal del anillo noetheriano $\mathrm{A}_\lambda$. De la definición del límite inductivo se deduce que existe un índice $\mu\geqslant\lambda$ tal que $\varphi_{\mu\lambda}(a_\lambda)=0$; en otras palabras, $\varphi_{\mu\lambda}$ se factoriza como $\varphi_{\mu\lambda}:\mathrm{A}_\lambda\to\mathrm{A}'_\lambda\to\mathrm{A}_\mu$, y se puede escribir $\mathrm{B}_\mu=\mathrm{B}'_\lambda\otimes_{\mathrm{A}'_\lambda}\mathrm{A}_\mu$, $\mathfrak{J}_\mu=\mathfrak{J}'_\lambda\mathrm{B}_\mu$ y $\mathrm{M}_\mu=\mathrm{M}'_\lambda\otimes_{\mathrm{A}'_\lambda}\mathrm{A}_\mu$; se deduce entonces de \sref{IV.11.2.9.2} que $\operatorname{gr}^*_{\mathfrak{J}_\mu}(\mathrm{M}_\mu)$ es un $\mathrm{A}_\mu$-módulo plano.

\medskip
\noindent\textbf{IV)} \emph{Reducción al caso en que $\mathrm{B}_0=\mathrm{A}_0[\mathrm{T}_1,\ldots,\mathrm{T}_n]$ y $\operatorname{gr}^*_{\mathfrak{J}_0}(\mathrm{B}_0)=\mathrm{B}_0$.} --- En virtud de la hipótesis, existe un sistema $(t_i)_{1\leqslant i\leqslant n}$ de generadores de la $\mathrm{A}_0$-álgebra de tipo finito $\mathrm{B}_0$, tal que los $t_i$ para $i\leqslant m$ generan el ideal $\mathfrak{J}_0$ de $\mathrm{B}_0$. Pongamos $\mathrm{B}'_0=\mathrm{A}_0[\mathrm{T}_1,\ldots,\mathrm{T}_n]$, un álgebra de polinomios (y, por tanto, noetheriana), y sea $\mathfrak{J}'_0$ el ideal de $\mathrm{B}'_0$ generado por los $\mathrm{T}_i$ con $i\leqslant m$; existe entonces un $\mathrm{A}_0$-homomorfismo sobreyectivo $u_0:\mathrm{B}'_0\to\mathrm{B}_0$ que envía cada $\mathrm{T}_i$ a $t_i$ ($1\leqslant i\leqslant n$) y que, por tanto, satisface $u_0(\mathfrak{J}'_0)=\mathfrak{J}_0$. Este homomorfismo permite considerar $\mathrm{M}_0$ como un $\mathrm{B}'_0$-módulo de tipo finito. Pongamos
\[
  \mathrm{B}'_\lambda=\mathrm{B}'_0\otimes_{\mathrm{A}_0}\mathrm{A}_\lambda
  =\mathrm{A}_\lambda[\mathrm{T}_1,\ldots,\mathrm{T}_n],
  \qquad
  \mathfrak{J}'_\lambda=\mathfrak{J}'_0\mathrm{B}'_\lambda,
\]
\oldpage[IV-3]{129}
y también
\[
  \mathrm{B}'=\mathrm{B}'_0\otimes_{\mathrm{A}_0}\mathrm{A}
  =\mathrm{A}[\mathrm{T}_1,\ldots,\mathrm{T}_n],
  \qquad
  \mathfrak{J}'=\mathfrak{J}'_0\mathrm{B}'.
\]
Así, $\mathfrak{J}'$ es el ideal de $\mathrm{B}'$ generado por los $\mathrm{T}_i$ con $i\leqslant m$. Además, para todo entero $k\geqslant 0$ se tiene claramente
\[
  \mathfrak{J}'^k\mathrm{M}=\mathfrak{J}^k\mathrm{M},
  \qquad
  \mathfrak{J}'_\lambda{}^k\mathrm{M}_\lambda
  =\mathfrak{J}_\lambda^k\mathrm{M}_\lambda
\]
para todo $\lambda$. Por tanto,
\[
  \operatorname{gr}^*_{\mathfrak{J}'}(\mathrm{M})
  =\operatorname{gr}^*_{\mathfrak{J}}(\mathrm{M})
\]
como $\mathrm{A}$-módulos, y
\[
  \operatorname{gr}^*_{\mathfrak{J}'_\lambda}(\mathrm{M}_\lambda)
  =\operatorname{gr}^*_{\mathfrak{J}_\lambda}(\mathrm{M}_\lambda)
\]
como $\mathrm{A}_\lambda$-módulos. Podemos, por tanto, sustituir $\mathrm{B}'$, $\mathrm{B}'_\lambda$, $\mathfrak{J}'$ y $\mathfrak{J}'_\lambda$ por $\mathrm{B}$, $\mathrm{B}_\lambda$, $\mathfrak{J}$ y $\mathfrak{J}_\lambda$, respectivamente, en la demostración. Obsérvese también que, por construcción, $\operatorname{gr}^*_{\mathfrak{J}'}(\mathrm{B}')$ se identifica con $\mathrm{B}'$, y $\operatorname{gr}^*_{\mathfrak{J}'_\lambda}(\mathrm{B}'_\lambda)$ con $\mathrm{B}'_\lambda$.

\medskip
\noindent\textbf{V)} \emph{Demostración de que $\operatorname{gr}^*_\mathfrak{J}(\mathrm{M})$ es un $\operatorname{gr}^*_\mathfrak{J}(\mathrm{B})$-módulo de presentación finita.} --- Supondremos, pues, en lo sucesivo, que $\mathrm{A}_0$ y los $\mathrm{A}_\lambda$ son noetherianos, que los homomorfismos de transición $\mathrm{A}_\lambda\to\mathrm{A}_\mu$ son inyectivos, que $\mathrm{B}_0=\mathrm{A}_0[\mathrm{T}_1,\ldots,\mathrm{T}_n]$ y que $\mathfrak{J}_0$ está generado por los $\mathrm{T}_i$ con $i\leqslant m$. El anillo $\mathrm{C}_0=\operatorname{gr}^*_{\mathfrak{J}_0}(\mathrm{B}_0)$ se identifica, por tanto, con $\mathrm{B}_0$, y $\mathrm{C}=\operatorname{gr}^*_\mathfrak{J}(\mathrm{B})$ se identifica con $\mathrm{B}$; al comienzo de lo que sigue solo utilizaremos el hecho de que $\mathrm{C}\cong\mathrm{C}_0\otimes_{\mathrm{A}_0}\mathrm{A}$.

Necesitaremos la siguiente variante de $(6.9.3)$:

\begin{lemma}[11.2.9.5]
\label{IV.11.2.9.5}
Sea $\mathrm{A}_0$ un anillo noetheriano, $\mathrm{B}_0$ una
$\mathrm{A}_0$-álgebra de tipo finito, $\mathfrak{J}_0$ un ideal de
$\mathrm{B}_0$ y $\mathrm{M}_0$ un $\mathrm{B}_0$-módulo de tipo finito.
Existe una sucesión $(\mathrm{S}_{0i})_{1\leqslant i\leqslant m}$ de
subesquemas de $\mathrm{S}_0=\Spec(\mathrm{A}_0)$ con las propiedades
siguientes:
\begin{enumerate}[label=\arabic{enumi}\textsuperscript{o}]
  \item Los espacios subyacentes de los $\mathrm{S}_{0i}$ son disjuntos dos a dos y forman un recubrimiento de $\mathrm{S}_0$.
  \item Para todo $i$, el conjunto $\mathrm{S}_{0i}$ es abierto en $\bigcup_{j\geqslant i}\mathrm{S}_{0j}$.
  \item Cada esquema $\mathrm{S}_{0i}$ es afín.
  \item Si $\mathrm{A}_{0i}$ es el anillo de $\mathrm{S}_{0i}$ y se pone $\mathrm{B}_{0i}=\mathrm{B}_0\otimes_{\mathrm{A}_0}\mathrm{A}_{0i}$, $\mathfrak{J}_{0i}=\mathfrak{J}_0\mathrm{B}_{0i}$, entonces $\operatorname{gr}^*_{\mathfrak{J}_{0i}}(\mathrm{M}_0\otimes_{\mathrm{A}_0}\mathrm{A}_{0i})$ es un $\mathrm{A}_{0i}$-módulo plano.
\end{enumerate}
\end{lemma}

\begin{proof}
Procedamos como en $(6.9.3)$ por inducción noetheriana, suponiendo conocido el lema tras sustituir $\mathrm{A}_0$ por
\[
  \mathrm{A}'_0=\mathrm{A}_0/\mathfrak{a}_0,
\]
donde $\mathfrak{a}_0$ es un ideal tal que $\mathrm{V}(\mathfrak{a}_0)\neq\mathrm{S}_0$, y sustituyendo $\mathrm{B}_0$, $\mathfrak{J}_0$ y $\mathrm{M}_0$ por
\[
  \mathrm{B}'_0=\mathrm{B}_0\otimes_{\mathrm{A}_0}\mathrm{A}'_0,
  \qquad
  \mathfrak{J}'_0=\mathfrak{J}_0\mathrm{B}'_0,
  \qquad
  \mathrm{M}_0\otimes_{\mathrm{A}_0}\mathrm{A}'_0,
\]
respectivamente. Sea $\mathfrak{N}_0$ el nilradical de $\mathrm{A}_0$. Es
evidente que basta demostrar el lema después de sustituir $\mathrm{A}_0$ por
$\mathrm{A}_0/\mathfrak{N}_0$ y $\mathrm{B}_0$,
$\mathfrak{J}_0$ y $\mathrm{M}_0$ por los objetos correspondientes que acabamos
de describir. En otras palabras, podemos suponer que $\mathrm{A}_0$ es
\emph{reducida}. Ahora
$\operatorname{gr}^*_{\mathfrak{J}_0}(\mathrm{B}_0)$ es una
$\mathrm{A}_0$-álgebra de tipo finito y
$\operatorname{gr}^*_{\mathfrak{J}_0}(\mathrm{M}_0)$ es un módulo de tipo finito sobre
$\operatorname{gr}^*_{\mathfrak{J}_0}(\mathrm{B}_0)$. Por el teorema de
planitud genérica $(6.9.1)$, existe un subconjunto abierto
$\mathrm{D}(t_0)\subset\mathrm{S}_0$ tal que, poniendo
$\mathrm{A}_{01}=(\mathrm{A}_0)_{t_0}$,
\[
  \operatorname{gr}^*_{\mathfrak{J}_0}(\mathrm{M}_0)
  \otimes_{\mathrm{A}_0}\mathrm{A}_{01}
\]
es un $\mathrm{A}_{01}$-módulo plano. Como $\mathrm{A}_{01}$ es un $\mathrm{A}_0$-módulo plano, este módulo se identifica con
\[
  \operatorname{gr}^*_{\mathfrak{J}_{01}}
  (\mathrm{M}_0\otimes_{\mathrm{A}_0}\mathrm{A}_{01}),
\]
donde
\[
  \mathrm{B}_{01}=\mathrm{B}_0\otimes_{\mathrm{A}_0}\mathrm{A}_{01},
  \qquad
  \mathfrak{J}_{01}=\mathfrak{J}_0\mathrm{B}_{01}.
\]
El complemento de $\mathrm{D}(t_0)$ en $\mathrm{S}_0$ es entonces de la forma $\mathrm{V}(\mathfrak{a}_0)$, y la conclusión se obtiene aplicando la hipótesis de inducción noetheriana.
\end{proof}

Apliquemos este lema a la situación actual, conservando las mismas notaciones; pongamos $\mathrm{U}_{0i}=\bigcup_{j\leqslant i}\mathrm{S}_{0j}$, que es un abierto cuasicompacto de $\mathrm{S}_0$. Existe entonces una familia $(f_{ik})_{1\leqslant i\leqslant m,\,1\leqslant k\leqslant n}$ de elementos de $\mathrm{A}_0$ tal que, para todo $i$, $\mathrm{U}_{0i}$ es la unión de los $\mathrm{D}(f_{ik})$ ($1\leqslant k\leqslant n$).

\oldpage[IV-3]{130}

El $\mathrm{C}_0$-módulo $\operatorname{gr}^*_{\mathfrak{J}_0}(\mathrm{M}_0)$ está generado por un número finito de elementos homogéneos de grado $0$, por lo que existe un epimorfismo de $\mathrm{C}_0$-módulos graduados
\[
  u_0:\mathrm{C}_0^r\longrightarrow\operatorname{gr}^*_{\mathfrak{J}_0}(\mathrm{M}_0).
\]
Como $\mathrm{C}=\mathrm{C}_0\otimes_{\mathrm{A}_0}\mathrm{A}$ y el homomorfismo canónico $\operatorname{gr}^*_{\mathfrak{J}_0}(\mathrm{M}_0)\otimes_{\mathrm{A}_0}\mathrm{A}\to\operatorname{gr}^*_\mathfrak{J}(\mathrm{M})$ es sobreyectivo, se deduce de $u_0$ un epimorfismo de $\mathrm{C}$-módulos graduados
\[
  u:\mathrm{C}^r\xrightarrow{u_0\otimes 1_{\mathrm{A}}}\operatorname{gr}^*_{\mathfrak{J}_0}(\mathrm{M}_0)\otimes_{\mathrm{A}_0}\mathrm{A}\longrightarrow\operatorname{gr}^*_\mathfrak{J}(\mathrm{M}),
\]
y queda por demostrar que el $\mathrm{C}$-módulo graduado
$\mathrm{Q}=\operatorname{Ker}(u)$ es \emph{de tipo finito}.
\\[0.4\baselineskip]
\begin{lemma}[11.2.9.6]
\label{IV.11.2.9.6}
Bajo las hipótesis precedentes, sea $\mathfrak{R}_0$ un ideal de $\mathrm{A}_0$; pongamos
\[
  \mathrm{A}'_0=\mathrm{A}_0/\mathfrak{R}_0,\qquad
  \mathrm{B}'_0=\mathrm{B}_0\otimes_{\mathrm{A}_0}\mathrm{A}'_0
  =\mathrm{B}_0/\mathfrak{R}_0\mathrm{B}_0,\qquad
  \mathfrak{J}'_0=\mathfrak{J}_0\mathrm{B}'_0,
\]
y $\mathfrak{R}=\mathfrak{R}_0\mathrm{A}$, $\mathrm{A}'=\mathrm{A}/\mathfrak{R}$, y supongamos además que
\[
  \operatorname{gr}^*_{\mathfrak{J}'_0}
  (\mathrm{M}_0\otimes_{\mathrm{A}_0}\mathrm{A}'_0)
\]
es un $\mathrm{A}'_0$-módulo plano. Entonces existe un número finito de elementos $q_h$ ($1\leqslant h\leqslant s$) de $\mathrm{Q}$ tales que, para todo ideal primo $\mathfrak{p}\supset\mathfrak{R}\mathrm{B}$ de $\mathrm{B}$, las imágenes canónicas de los $q_h$ en $\mathrm{Q}_\mathfrak{p}$ generan el $\mathrm{C}_\mathfrak{p}$-módulo $\mathrm{Q}_\mathfrak{p}$.
\end{lemma}

Supongamos demostrado el lema y observemos que sigue siendo aplicable después de sustituir
$\mathrm{A}_0$ por $(\mathrm{A}_0)_{t_0}$, y $\mathrm{B}_0$,
$\mathfrak{J}_0$, $\mathrm{M}_0$, $\mathrm{A}_\lambda$ y $\mathrm{A}$ por
los objetos localizados correspondientes, para cualquier $t_0\in\mathrm{A}_0$, pues
$(\mathrm{A}_0)_{t_0}$ es plano sobre $\mathrm{A}_0$. Apliquemos el lema
sustituyendo sucesivamente $\mathrm{A}_0$ por cada uno de los anillos
$(\mathrm{A}_0)_{f_{ik}}$, y $\mathfrak{R}_0$ por el ideal
$\mathfrak{R}_{ik}$ que define el subesquema cerrado de $\mathrm{D}(f_{ik})$
inducido por $\mathrm{S}_{0i}$ sobre el subconjunto abierto
$\mathrm{S}_{0i}\cap\mathrm{D}(f_{ik})$ de $\mathrm{S}_{0i}$. Como la
hipótesis de planitud de \sref{IV.11.2.9.6} se satisface en cada uno de estos casos
por \sref{IV.11.2.9.5}, para cada par $(i,k)$ se obtiene una familia finita
$(a_{ikh})_{1\leqslant h\leqslant l}$ de elementos de
$\mathrm{Q}_{f_{ik}}$ cuyas imágenes en $\mathrm{Q}_{\mathfrak{p}_{ik}}$
generan este $\mathrm{C}_{\mathfrak{p}_{ik}}$-módulo para todo ideal primo
$\mathfrak{p}_{ik}$ de $\mathrm{B}_{f_{ik}}$ que contiene
$\mathfrak{R}_{ik}\mathrm{B}_{f_{ik}}$. Para un entero adecuado $d>0$, se
puede escribir
\[
  a_{ikh}=\frac{b_{ikh}}{f_{ik}^{d}}
\]
para todos $i,k,h$, con $b_{ikh}\in\mathrm{Q}$. Como los $\mathrm{S}_{0i}$ recubren $\mathrm{S}_0$, todo ideal primo $\mathfrak{p}$ de $\mathrm{B}$ tiene imagen en algún $\mathrm{S}_{0i}\cap\mathrm{D}(f_{ik})$ y, por tanto, su imagen $\mathfrak{p}_{ik}$ en $\mathrm{B}_{f_{ik}}$ contiene $\mathfrak{R}_{ik}\mathrm{B}_{f_{ik}}$. Las imágenes de los $b_{ikh}$ ($1\leqslant h\leqslant l$) en $\mathrm{Q}_{\mathfrak{p}_{ik}}=\mathrm{Q}_{\mathfrak{p}}$ generan, por consiguiente, este $\mathrm{C}_{\mathfrak{p}}$-módulo. Se deduce que los $b_{ikh}$ ($1\leqslant i\leqslant m$, $1\leqslant k\leqslant n$, $1\leqslant h\leqslant l$) generan el $\mathrm{C}$-módulo $\mathrm{Q}$ (Bourbaki, \emph{Alg.\ comm.}, cap.~II, \S3, n.º~3, teorema~1).

Queda por demostrar el Lema \sref{IV.11.2.9.6}. Pongamos
\[
  \mathrm{C}'=\mathrm{C}\otimes_\mathrm{A}\mathrm{A}'=\mathrm{C}/\mathfrak{R}\mathrm{C},\qquad
  \mathrm{B}'=\mathrm{B}\otimes_\mathrm{A}\mathrm{A}',\qquad
  \mathfrak{J}'=\mathfrak{J}\mathrm{B}',\qquad
  \mathrm{M}'=\mathrm{M}\otimes_\mathrm{A}\mathrm{A}'.
\]
Por hipótesis, $\operatorname{gr}^*_\mathfrak{J}(\mathrm{M})$ es un $\mathrm{A}$-módulo plano; por tanto, $\operatorname{gr}^*_{\mathfrak{J}'}(\mathrm{M}')$ se identifica con $\operatorname{gr}^*_\mathfrak{J}(\mathrm{M})\otimes_\mathrm{A}\mathrm{A}'$ \sref{IV.11.2.9.2}, y se tiene la sucesión exacta $(\mathbf{0}_\mathrm{I},~6.1.2)$
\begin{equation}
  0\longrightarrow\mathrm{Q}\otimes_\mathrm{A}\mathrm{A}'
  \longrightarrow(\mathrm{C}')^r
  \xrightarrow{u\otimes 1_{\mathrm{A}'}}
  \operatorname{gr}^*_{\mathfrak{J}'}(\mathrm{M}')
  \longrightarrow 0.
  \label{IV.11.2.9.7}
\tag{11.2.9.7}
\end{equation}

Pongamos
\[
  \mathrm{C}'_0=\mathrm{C}_0\otimes_{\mathrm{A}_0}\mathrm{A}'_0
  =\mathrm{C}_0/\mathfrak{R}_0\mathrm{C}_0
\]
y consideremos el epimorfismo de $\mathrm{C}'_0$-módulos graduados
\[
  u'_0:(\mathrm{C}'_0)^r
  \xrightarrow{u_0\otimes 1_{\mathrm{A}'_0}}
  \operatorname{gr}^*_{\mathfrak{J}_0}(\mathrm{M}_0)\otimes_{\mathrm{A}_0}\mathrm{A}'_0
  \longrightarrow
  \operatorname{gr}^*_{\mathfrak{J}'_0}
  (\mathrm{M}_0\otimes_{\mathrm{A}_0}\mathrm{A}'_0).
\]

Sea $\mathrm{Q}_0=\operatorname{Ker}(u'_0)$. Utilizando el hecho de que
\[
  \operatorname{gr}^*_{\mathfrak{J}'_0}
  (\mathrm{M}_0\otimes_{\mathrm{A}_0}\mathrm{A}'_0)
\]
es un $\mathrm{A}'_0$-módulo plano, se ve que su producto tensorial sobre $\mathrm{A}'_0$ con $\mathrm{A}'$ se identifica con $\operatorname{gr}^*_{\mathfrak{J}'}(\mathrm{M}')$ \sref{IV.11.2.9.2}, y se tiene la sucesión exacta $(\mathbf{0}_\mathrm{I},~6.1.2)$
\begin{equation}
  0\longrightarrow\mathrm{Q}_0\otimes_{\mathrm{A}'_0}\mathrm{A}'
  \longrightarrow(\mathrm{C}')^r
  \xrightarrow{u\otimes 1_{\mathrm{A}'}}
  \operatorname{gr}^*_{\mathfrak{J}'}(\mathrm{M}')
  \longrightarrow 0.
  \label{IV.11.2.9.8}
\tag{11.2.9.8}
\end{equation}

\oldpage[IV-3]{131}

de donde se obtiene, por comparación con {\hyperlink{ega4.c04.thm.11-2-9}{(11.2.9.7)}}, un isomorfismo de $\mathrm{C}'$-módulos graduados
\[
  \mathrm{Q}\otimes_\mathrm{A}\mathrm{A}'
  \xrightarrow{\sim}
  \mathrm{Q}_0\otimes_{\mathrm{A}'_0}\mathrm{A}'.
\]

Como $\mathrm{C}'_0$ es noetheriano, $\mathrm{Q}_0$ es un $\mathrm{C}'_0$-módulo de tipo finito; por tanto, $\mathrm{Q}_0\otimes_{\mathrm{A}'_0}\mathrm{A}'$ es un $\mathrm{C}'$-módulo graduado de tipo finito y, por consiguiente, también lo es $\mathrm{Q}\otimes_\mathrm{A}\mathrm{A}'$. Sean $q_h$ ($1\leqslant h\leqslant s$) elementos homogéneos de $\mathrm{Q}$ cuyas imágenes en $\mathrm{Q}\otimes_\mathrm{A}\mathrm{A}'$ generan este $\mathrm{C}'$-módulo.

Consideremos ahora un ideal primo arbitrario $\mathfrak{p}\supset\mathfrak{R}\mathrm{B}$ de $\mathrm{B}$. Por planitud,
\[
  \mathrm{C}_\mathfrak{p}
  =\operatorname{gr}^*_{\mathfrak{J}_\mathfrak{p}}(\mathrm{B}_\mathfrak{p}),
\]
y $\operatorname{gr}^0_{\mathfrak{J}_\mathfrak{p}}(\mathrm{B}_\mathfrak{p})=\mathrm{B}_\mathfrak{p}/\mathfrak{J}_\mathfrak{p}$ es el anillo cero o un anillo local; para demostrar el lema, basta evidentemente considerar el segundo caso. Cada $(\mathrm{B}_\mathfrak{p}/\mathfrak{J}_\mathfrak{p})$-módulo $\operatorname{gr}^i_{\mathfrak{J}_\mathfrak{p}}(\mathrm{B}_\mathfrak{p})$ es entonces de presentación finita. Por otra parte, para todo $i$,
\[
  \mathrm{M}/\mathfrak{J}^i\mathrm{M}
  \simeq
  (\mathrm{M}_0/\mathfrak{J}_0^i\mathrm{M}_0)
  \otimes_{\mathrm{A}_0}\mathrm{A},
\]
y es, por tanto, un $\mathrm{B}$-módulo de presentación finita. Por inducción sobre $i$, la hipótesis de que $\operatorname{gr}^*_\mathfrak{J}(\mathrm{M})$ es plano sobre $\mathrm{A}$ implica, por $(\mathbf{0}_\mathrm{I},~6.1.2)$, que $\mathrm{M}/\mathfrak{J}^i\mathrm{M}$ es plano sobre $\mathrm{A}$. Aplicando \sref{IV.11.2.6.1}, con $\mathrm{M}_0$ sustituido por $\mathrm{M}_0/\mathfrak{J}_0^k\mathrm{M}_0$ para $k\leqslant i$, se ve que hay un índice $\lambda_i$ tal que, para $\mu\geqslant\lambda_i$, cada $\mathrm{M}_\mu/\mathfrak{J}_\mu^{k+1}\mathrm{M}_\mu$ es plano sobre $\mathrm{A}_\mu$; por consiguiente, de nuevo por $(\mathbf{0}_\mathrm{I},~6.1.2)$, cada $\operatorname{gr}^k_{\mathfrak{J}_\mu}(\mathrm{M}_\mu)$ es plano sobre $\mathrm{A}_\mu$ para $k\leqslant i$. Se deduce de \sref{IV.11.2.9.2} que cada $(\mathrm{B}/\mathfrak{J})$-módulo $\operatorname{gr}^i_\mathfrak{J}(\mathrm{M})$ es de presentación finita y, por tanto, que $\operatorname{gr}^i_{\mathfrak{J}_\mathfrak{p}}(\mathrm{M}_\mathfrak{p})$ es un $(\mathrm{B}_\mathfrak{p}/\mathfrak{J}_\mathfrak{p})$-módulo de presentación finita.

Por otra parte, las imágenes de los $q_h$ en
\[
  \mathrm{Q}_\mathfrak{p}\otimes_{\mathrm{B}_\mathfrak{p}}\boldsymbol{k}(\mathfrak{p})
  =
  (\mathrm{Q}\otimes_\mathrm{A}\mathrm{A}')_\mathfrak{p}
  \otimes_{\mathrm{A}'}\boldsymbol{k}(\mathfrak{p})
\]
generan este módulo sobre $\mathrm{C}_\mathfrak{p}\otimes_{\mathrm{B}_\mathfrak{p}}\boldsymbol{k}(\mathfrak{p})$, pues
\[
  \mathrm{C}_\mathfrak{p}\otimes_{\mathrm{B}_\mathfrak{p}}\boldsymbol{k}(\mathfrak{p})
  =
  \mathrm{C}'_\mathfrak{p}\otimes_{\mathrm{A}'}\boldsymbol{k}(\mathfrak{p}).
\]
Como se tiene una sucesión exacta
\[
  \mathrm{Q}_\mathfrak{p}\otimes_{\mathrm{B}_\mathfrak{p}}\boldsymbol{k}(\mathfrak{p})
  \longrightarrow
  \mathrm{C}_\mathfrak{p}^{\,r}\otimes_{\mathrm{B}_\mathfrak{p}}\boldsymbol{k}(\mathfrak{p})
  \longrightarrow
  \operatorname{gr}^*_{\mathfrak{J}_\mathfrak{p}}(\mathrm{M}_\mathfrak{p})
  \otimes_{\mathrm{B}_\mathfrak{p}}\boldsymbol{k}(\mathfrak{p})
  \longrightarrow 0,
\]
se deduce que $\operatorname{gr}^*_{\mathfrak{J}_\mathfrak{p}}(\mathrm{M}_\mathfrak{p})\otimes_{\mathrm{B}_\mathfrak{p}}\boldsymbol{k}(\mathfrak{p})$ es un módulo de presentación finita sobre $\mathrm{C}_\mathfrak{p}\otimes_{\mathrm{B}_\mathfrak{p}}\boldsymbol{k}(\mathfrak{p})$. Aplicando \sref{IV.11.2.9.4}, con $\mathrm{A}$, $\mathrm{B}$ y $\mathrm{M}$ sustituidos, respectivamente, por $\mathrm{B}_\mathfrak{p}$, $\mathrm{C}_\mathfrak{p}$ y $\operatorname{gr}^*_{\mathfrak{J}_\mathfrak{p}}(\mathrm{M}_\mathfrak{p})$, concluimos que $\mathrm{Q}_\mathfrak{p}$ es un $\mathrm{C}_\mathfrak{p}$-módulo de tipo finito. El lema de Nakayama (y el hecho de que $\mathrm{C}=\mathrm{B}$) muestra entonces que las imágenes de los $q_h$ en $\mathrm{Q}_\mathfrak{p}$ generan este $\mathrm{C}_\mathfrak{p}$-módulo. Esto completa la demostración de \sref{IV.11.2.9.6} y demuestra también que $\operatorname{gr}^*_\mathfrak{J}(\mathrm{M})$ es un $\mathrm{C}$-módulo de presentación finita.

\medskip
\noindent\textbf{VI)} \emph{Fin de la demostración.} --- Para abreviar, pongamos $\mathrm{C}_\lambda=\mathrm{C}_0\otimes_{\mathrm{A}_0}\mathrm{A}_\lambda$ (de hecho, igual a $\mathrm{B}_\lambda$), $\mathrm{N}_\lambda=\operatorname{gr}^*_{\mathfrak{J}_\lambda}(\mathrm{M}_\lambda)$ y $\mathrm{N}=\operatorname{gr}^*_\mathfrak{J}(\mathrm{M})$. Obsérvese primero que, para todo entero $k$, $\mathfrak{J}^k\mathrm{M}$ se identifica con $\varinjlim_\lambda\mathfrak{J}_\lambda^k\mathrm{M}_\lambda$ por la exactitud del funtor $\varinjlim_\lambda$: la imagen de $\mathfrak{J}_\lambda^k\mathrm{M}_\lambda$ en $\mathrm{M}_\mu$ para $\lambda\leqslant\mu$ (respectivamente, en $\mathrm{M}$) genera $\mathfrak{J}_\mu^k\mathrm{M}_\mu$ (respectivamente, $\mathfrak{J}^k\mathrm{M}$). Utilizando una vez más la exactitud de $\varinjlim$, concluimos que $\mathrm{N}$ se identifica canónicamente con $\varinjlim_\lambda\mathrm{N}_\lambda$. Con esta identificación, demostraremos primero:
\begin{equation}
\label{IV.11.2.9.9}
\tag{11.2.9.9}
\text{\emph{Para $\lambda$ suficientemente grande, el homomorfismo canónico $v_\lambda:\mathrm{N}_\lambda\otimes_{\mathrm{A}_\lambda}\mathrm{A}\to\mathrm{N}$ es biyectivo.}}
\end{equation}
Como los $\mathrm{C}_\lambda$ son noetherianos y $\mathrm{N}_\lambda$ es un $\mathrm{C}_\lambda$-módulo de tipo finito, los $\mathrm{C}$-módulos $\mathrm{N}_\lambda\otimes_{\mathrm{A}_\lambda}\mathrm{A}$ son de presentación finita y forman un sistema inductivo filtrante cuyo límite inductivo se identifica canónicamente con $\mathrm{N}$, pues los límites inductivos filtrantes conmutan con los productos tensoriales. Además, los homomorfismos de transición $v_{\mu\lambda}:\mathrm{N}_\lambda\otimes_{\mathrm{A}_\lambda}\mathrm{A}\to\mathrm{N}_\mu\otimes_{\mathrm{A}_\mu}\mathrm{A}$

\oldpage[IV-3]{132}

para $\lambda\leqslant\mu$, así como los homomorfismos $v_\lambda$, son
\emph{sobreyectivos}. Fijemos $\lambda$ y consideremos
\[
  \operatorname{Ker}(v_{\mu\lambda})=\mathrm{N}'_\mu
\]
como un $\mathrm{C}$-submódulo de
$\mathrm{N}_\lambda\otimes_{\mathrm{A}_\lambda}\mathrm{A}$. Por la
definición del límite inductivo, los $\mathrm{N}'_\mu$ forman una
familia filtrante creciente cuya unión es $\operatorname{Ker}(v_\lambda)$.
Demostramos en V) que $\mathrm{N}$ es un $\mathrm{C}$-módulo de presentación finita;
en consecuencia, $\operatorname{Ker}(v_\lambda)$ es un $\mathrm{C}$-módulo de tipo finito
(Bourbaki, \emph{Alg.\ comm.}, cap.~I, \S~2, no.~8, lema~9).
Hay, pues, un índice $\mu\geqslant\lambda$ tal que
$\mathrm{N}'_\mu=\operatorname{Ker}(v_\lambda)$. Como $v_\mu$ es
sobreyectivo, esta igualdad dice que $v_\mu$ es también inyectivo y,
por tanto, demuestra {\hyperlink{ega4.c04.thm.11-2-9}{(11.2.9.9)}}.

Después de sustituir $\mathrm{A}_0$ y $\mathrm{M}_0$ por
$\mathrm{A}_\lambda$ y $\mathrm{M}_\lambda$ para un $\lambda$ suficientemente grande,
podemos suponer, por tanto, que el homomorfismo canónico
$v_\lambda$ es biyectivo para \emph{todo} $\lambda$. Pongamos
\[
  \mathrm{P}_\lambda
  =\mathrm{N}_0\otimes_{\mathrm{C}_0}\mathrm{C}_\lambda
  =\mathrm{N}_0\otimes_{\mathrm{A}_0}\mathrm{A}_\lambda,
\]
de modo que $\mathrm{N}=\varinjlim\mathrm{P}_\lambda$. Como
$\mathrm{C}_0$ es una $\mathrm{A}_0$-álgebra de presentación finita y
$\mathrm{P}_0$ es un $\mathrm{C}_0$-módulo de presentación finita,
\sref{IV.11.2.6.1}, aplicado con
$\mathrm{B}_0=\mathrm{C}_0$ y $\mathrm{M}_0=\mathrm{P}_0$, muestra que
hay un índice $\lambda$ tal que $\mathrm{P}_\mu$ es un
$\mathrm{A}_\mu$-módulo plano para todo $\mu\geqslant\lambda$. Para tal
$\mu$ tenemos un diagrama conmutativo
\[
\xymatrix{
  \mathrm{P}_\mu \ar[r] \ar[d]^{w_\mu} & \mathrm{N} \ar@{=}[d] \\
  \mathrm{N}_\mu \ar[r] & \mathrm{N}
}
\]
y las definiciones muestran que $w_\mu$ es \emph{sobreyectivo}. Por III), el
homomorfismo $\mathrm{A}_\mu\to\mathrm{A}$ es inyectivo; como
$\mathrm{P}_\mu$ es plano sobre $\mathrm{A}_\mu$, el homomorfismo canónico
\[
  \mathrm{P}_\mu\longrightarrow
  \mathrm{N}
  =\mathrm{P}_\mu\otimes_{\mathrm{A}_\mu}\mathrm{A}
\]
es, por tanto, \emph{inyectivo}. El diagrama anterior muestra entonces que
$w_\mu$ es también inyectivo y, por consiguiente, biyectivo. Así,
$\mathrm{N}_\mu$ es un $\mathrm{A}_\mu$-módulo plano para
$\mu\geqslant\lambda$, lo que completa la demostración de
\sref{IV.11.2.9}.
\end{proof}

\begin{remark}[11.2.10]
\label{IV.11.2.10}
No sabemos si es válido el análogo del teorema de Raynaud para la
generalización de \sref{IV.11.2.6},~(i).
\end{remark}

\subsection{Aplicación a la eliminación de hipótesis noetherianas}
\label{subsection:IV.11.3}
\label{IV.11.3}

\oldpage[IV-3]{132}

\begin{theorem}[11.3.1]
\label{IV.11.3.1}
Sea $f:\mathrm{X}\to\mathrm{Y}$ un morfismo localmente de
presentación finita, y sea $\sh{F}$ un $\sh{O}_\mathrm{X}$-módulo
cuasicoherente de presentación finita. Entonces el conjunto
$\mathrm{U}$ de los puntos $x\in\mathrm{X}$ en los que $\sh{F}$ es
$f$-plano es abierto en $\mathrm{X}$. Si, además,
$\operatorname{Supp}(\sh{F})=\mathrm{X}$, la restricción
$f|_\mathrm{U}$ es un morfismo abierto.
\end{theorem}

\begin{proof}
La segunda afirmación ya se ha demostrado en $(2.4.6)$ y se incluye
aquí únicamente como referencia.

Como la cuestión es local en $\mathrm{X}$ y en $\mathrm{Y}$, podemos
suponer que ambos son afines y que $f$ es de presentación finita. Sea
$x\in\mathrm{X}$ un punto en el que $\sh{F}$ es $f$-plano. Aplicando
\sref{IV.11.2.7}, podemos asumir que
\[
  \mathrm{X}=\mathrm{X}_0\times_{\mathrm{Y}_0}\mathrm{Y},
  \qquad
  f=f_0\times 1,
  \qquad
  \sh{F}
  =\sh{F}_0\otimes_{\sh{O}_{\mathrm{Y}_0}}\sh{O}_\mathrm{Y},
\]
en donde $\mathrm{Y}_0$ es noetheriano,
$f_0:\mathrm{X}_0\to\mathrm{Y}_0$ es de tipo finito, y
$\sh{F}_0$ es un $\sh{O}_{\mathrm{X}_0}$-módulo coherente. Además, si
$x_0$ es la proyección canónica de $x$ en $\mathrm{X}_0$, podemos
suponer que $\sh{F}_0$ es $f_0$-plano en $x_0$. Por
\sref{IV.11.1.1}, el conjunto $\mathrm{U}_0$ de puntos de
$\mathrm{X}_0$ en los que $\sh{F}_0$ es $f_0$-plano es, por tanto, una
vecindad de $x_0$. Se deduce de $(2.1.4)$ que $\sh{F}$ es
$f$-plano en cada punto de la imagen inversa de $\mathrm{U}_0$ en
$\mathrm{X}$; así pues, $\mathrm{U}$ es una vecindad de $x$.
\end{proof}

\oldpage[IV-3]{133}

\begin{corollary}[11.3.2]
\label{IV.11.3.2}
Sean $f:\mathrm{X}\to\mathrm{Y}$ y $g:\mathrm{Y}\to\mathrm{Z}$ dos morfismos de presentación finita, y sea $\sh{F}$ un $\sh{O}_\mathrm{X}$-módulo cuasicoherente de presentación finita. Entonces el conjunto $\mathrm{U}$ de los $y\in\mathrm{Y}$ tales que, para toda generización $y'$ de $y$, $\sh{F}$ es $(g\circ f)$-plano en todos los puntos de $f^{-1}(y')$ (es decir, tales que $\sh{F}\otimes_\mathrm{Y}\Spec(\sh{O}_{\mathrm{Y},y})$ es plano respecto del morfismo $\mathrm{X}\times_\mathrm{Y}\Spec(\sh{O}_{\mathrm{Y},y})\to\mathrm{Z}$) es abierto en $\mathrm{Y}$.
\end{corollary}

\begin{proof}
El mismo razonamiento que en \sref{IV.11.1.5} muestra que, si $\mathrm{V}$ es el conjunto de los $x\in\mathrm{X}$ en los que $\sh{F}$ es $(g\circ f)$-plano, $\mathrm{U}$ es el conjunto de los $y\in\mathrm{Y}$ cuyas generizaciones pertenecen todas a $\mathrm{E}=\mathrm{Y}-f(\mathrm{X}-\mathrm{V})$. Como $g\circ f$ es de presentación finita, $\mathrm{V}$ es abierto en $\mathrm{X}$ en virtud de \sref{IV.11.3.1}, y por consiguiente $\mathrm{V}$ es ind-construible $(1.9.6)$ y $\mathrm{X}-\mathrm{V}$ es proconstruible; ahora bien, como $f$ es cuasicompacto, $f(\mathrm{X}-\mathrm{V})$ es proconstruible en $\mathrm{Y}$ $(1.9.5,~(\mathrm{vii}))$, y por tanto $\mathrm{E}$ es ind-construible en $\mathrm{Y}$. Se deduce entonces de $(1.10.1)$ que $\mathrm{U}$ es el interior de $\mathrm{E}$ y, por consiguiente, es abierto en $\mathrm{Y}$.
\end{proof}

\begin{corollary}[11.3.3]
\label{IV.11.3.3}
Sea $\mathrm{A}$ un anillo, $\mathrm{B}$ una $\mathrm{A}$-álgebra de presentación finita, $\mathrm{C}$ una $\mathrm{B}$-álgebra de presentación finita y $\mathrm{M}$ un $\mathrm{C}$-módulo de presentación finita. Entonces el conjunto de los $\mathfrak{q}\in\Spec(\mathrm{B})$ tales que $\mathrm{M}_\mathfrak{q}$ es un $\mathrm{A}$-módulo plano es abierto en $\Spec(\mathrm{B})$.
\end{corollary}

\begin{proposition}[11.3.4]
\label{IV.11.3.4}
Sea $f:\mathrm{X}\to\mathrm{S}$ un morfismo localmente de presentación finita, $\mathscr{J}$ un ideal cuasicoherente de tipo finito de $\sh{O}_\mathrm{X}$, $\mathrm{Y}$ el subpreesquema cerrado de $\mathrm{X}$ definido por $\mathscr{J}$ y $\sh{F}$ un $\sh{O}_\mathrm{X}$-módulo de presentación finita. Supongamos que $\operatorname{gr}^*_\mathscr{J}(\sh{F})$ es $f$-plano. En estas condiciones:
\begin{enumerate}[label=\emph{(\roman*)}]
  \item $\sh{F}$ es $f$-plano en los puntos de $\mathrm{Y}$.
  \item Si $\operatorname{gr}^*_\mathscr{J}(\sh{O}_\mathrm{X})$ es $f$-plano, $\operatorname{gr}^*_\mathscr{J}(\sh{O}_\mathrm{X})$ es una $(\sh{O}_\mathrm{X}/\mathscr{J})$-álgebra de presentación finita y $\operatorname{gr}^*_\mathscr{J}(\sh{F})$ es un $(\operatorname{gr}^*_\mathscr{J}(\sh{O}_\mathrm{X}))$-módulo de presentación finita.
  \item Supongamos que $\operatorname{gr}^*_\mathscr{J}(\sh{O}_\mathrm{X})$ es $f$-plano (lo cual implica que $\sh{O}_\mathrm{X}/\mathscr{J}$ es $f$-plano). Entonces el conjunto $\mathrm{W}$ de los $y\in\mathrm{Y}$ tales que $(\operatorname{gr}^*_\mathscr{J}(\sh{F}))_y$ es un $(\sh{O}_\mathrm{X}/\mathscr{J})_y$-módulo plano es abierto en $\mathrm{Y}$.
  \item Supongamos que $\operatorname{gr}^*_\mathscr{J}(\sh{O}_\mathrm{X})$ es $f$-plano. Sea $\mathrm{S}'\to\mathrm{S}$ un morfismo cualquiera, sea $\mathrm{X}'=\mathrm{X}\times_\mathrm{S}\mathrm{S}'$ y sea $p:\mathrm{X}'\to\mathrm{X}$ la proyección canónica. Pongamos $\mathrm{Y}'=p^{-1}(\mathrm{Y})=\mathrm{Y}\times_\mathrm{S}\mathrm{S}'$, el subpreesquema cerrado de $\mathrm{X}'$ definido por $\mathscr{J}'=p^*(\mathscr{J})\sh{O}_{\mathrm{X}'}$, y pongamos $\sh{F}'=p^*(\sh{F})=\sh{F}\otimes_{\sh{O}_\mathrm{S}}\sh{O}_{\mathrm{S}'}$. Entonces, si $\mathrm{W}'=p^{-1}(\mathrm{W})$, se tiene $\operatorname{gr}^*_{\mathscr{J}'}(\sh{F}')|_{\mathrm{W}'}\cong p^*(\operatorname{gr}^*_\mathscr{J}(\sh{F})|_\mathrm{W})$, y $\operatorname{gr}^*_{\mathscr{J}'}(\sh{F}')$ es un $(\sh{O}_{\mathrm{X}'}/\mathscr{J}')$-módulo plano en los puntos de $\mathrm{W}'$.
\end{enumerate}
\end{proposition}

\begin{proof}
Como las cuestiones son locales en $\mathrm{X}$ y en $\mathrm{S}$, podemos suponer que
$\mathrm{S}=\Spec(\mathrm{A})$ y $\mathrm{X}=\Spec(\mathrm{B})$ son afines,
en donde $\mathrm{B}$ es una $\mathrm{A}$-álgebra de presentación finita, que
$\sh{F}=\widetilde{\mathrm{M}}$, donde $\mathrm{M}$ es una
$\mathrm{B}$-módulo de presentación finita, y que
$\mathscr{J}=\widetilde{\mathfrak{J}}$, donde $\mathfrak{J}$ es un ideal
de tipo finito de $\mathrm{B}$. En virtud de $(8.9.1)$ y $(8.5.11)$,
existen un subanillo noetheriano $\mathrm{A}_0$ de $\mathrm{A}$, una
$\mathrm{A}_0$-álgebra de tipo finito $\mathrm{B}_0$ de tal manera que
$\mathrm{B}=\mathrm{B}_0\otimes_{\mathrm{A}_0}\mathrm{A}$, un ideal
$\mathfrak{J}_0$ de $\mathrm{B}_0$ tal que
$\mathfrak{J}=\mathfrak{J}_0\mathrm{B}$, y un tipo finito
$\mathrm{B}_0$-módulo $\mathrm{M}_0$ tal que
$\mathrm{M}=\mathrm{M}_0\otimes_{\mathrm{A}_0}\mathrm{A}
=\mathrm{M}_0\otimes_{\mathrm{B}_0}\mathrm{B}$. Además,
$\mathrm{A}$ es el límite inductivo de sus
$\mathrm{A}_0$-subálgebras de tipo finito $\mathrm{A}_\lambda$. Pongamos
$\mathrm{B}_\lambda=\mathrm{B}_0\otimes_{\mathrm{A}_0}\mathrm{A}_\lambda$,
$\mathrm{M}_\lambda=\mathrm{M}_0\otimes_{\mathrm{A}_0}\mathrm{A}_\lambda
=\mathrm{M}_0\otimes_{\mathrm{B}_0}\mathrm{B}_\lambda$, y
$\mathfrak{J}_\lambda=\mathfrak{J}_0\mathrm{B}_\lambda$, de modo que
$\mathrm{B}=\varinjlim\mathrm{B}_\lambda$ y
$\mathrm{M}=\varinjlim\mathrm{M}_\lambda$. Por hipótesis,
$\operatorname{gr}^*_\mathfrak{J}(\mathrm{M})$ es un
$\mathrm{A}$-módulo plano; se deduce, pues, de \sref{IV.11.2.9} que
existe $\lambda$ tal que
$\operatorname{gr}^*_{\mathfrak{J}_\lambda}(\mathrm{M}_\lambda)$ es un
$\mathrm{A}_\lambda$-módulo plano. Para probar que $\sh{F}$ es $f$-plano en
los puntos de $\mathrm{Y}$, podemos por tanto limitarnos al caso en que
$\mathrm{S}$ es noetheriano; pero entonces $(\mathbf{0}_\mathrm{I},~6.1.2)$,
aplicado por inducción sobre $n$, demuestra que los
$\sh{F}/\mathscr{J}^{n+1}\sh{F}$ son $f$-planos, y se concluye mediante
$(\mathbf{0}_{\mathrm{III}},~10.2.6)$.
\oldpage[IV-3]{134}

La demostración de (ii) se reduce análogamente al caso en que $\mathrm{S}$ es noetheriano, utilizando \sref{IV.11.2.9}; la conclusión es entonces evidente, pues en este caso $\operatorname{gr}^*_\mathfrak{J}(\mathrm{B})$ es una $(\mathrm{B}/\mathfrak{J})$-álgebra de tipo finito y $\operatorname{gr}^*_\mathfrak{J}(\mathrm{M})$ es un $\operatorname{gr}^*_\mathfrak{J}(\mathrm{B})$-módulo de tipo finito.

Para demostrar (iii), nos reducimos, como en (i), al caso en que $\mathrm{S}$ y $\mathrm{X}$ son afines. Con las mismas notaciones, podemos suponer, en virtud de \sref{IV.11.2.9}, que $\operatorname{gr}^*_{\mathfrak{J}_\lambda}(\mathrm{B}_\lambda)$ y $\operatorname{gr}^*_{\mathfrak{J}_\lambda}(\mathrm{M}_\lambda)$ son $\mathrm{A}_\lambda$-módulos planos y que
\[
\operatorname{gr}^*_\mathfrak{J}(\mathrm{B})
=\operatorname{gr}^*_{\mathfrak{J}_\lambda}(\mathrm{B}_\lambda)\otimes_{\mathrm{A}_\lambda}\mathrm{A},
\qquad
\operatorname{gr}^*_\mathfrak{J}(\mathrm{M})
=\operatorname{gr}^*_{\mathfrak{J}_\lambda}(\mathrm{M}_\lambda)\otimes_{\mathrm{A}_\lambda}\mathrm{A}.
\]
En consecuencia, $\operatorname{gr}^*_\mathfrak{J}(\mathrm{B})$ es una $(\mathrm{B}/\mathfrak{J})$-álgebra de presentación finita y $\operatorname{gr}^*_\mathfrak{J}(\mathrm{M})$ es un $\operatorname{gr}^*_\mathfrak{J}(\mathrm{B})$-módulo de presentación finita. Si $\mathrm{W}^m$ es el conjunto de los $\mathfrak{p}\in\Spec(\mathrm{B}/\mathfrak{J})$ tales que $\operatorname{gr}^m_\mathfrak{J}(\mathrm{M})_{\mathfrak{p}}$ es un $(\mathrm{B}/\mathfrak{J})_{\mathfrak{p}}$-módulo plano, entonces $\mathrm{W}^m$ es abierto en $\Spec(\mathrm{B}/\mathfrak{J})$ por \sref{IV.11.3.1}, y $\mathrm{W}=\bigcap_m\mathrm{W}^m$; por tanto, $\mathrm{W}$ es estable por generización. La aserción (iii) se deduce ahora de \sref{IV.11.3.2}, aplicado con $\mathrm{X}=\Spec(\operatorname{gr}^*_\mathfrak{J}(\mathrm{B}))$, $\mathrm{Y}=\mathrm{Z}=\Spec(\mathrm{B}/\mathfrak{J})$ y $\sh{F}=(\operatorname{gr}^*_\mathfrak{J}(\mathrm{M}))^\sim$.

Por último, (iv) se deduce inmediatamente de \sref{IV.11.2.9.2}.

Generalizando la definición de $(6.10.1)$, sea $\mathrm{Y}$ un subpreesquema cerrado de un preesquema $\mathrm{X}$, definido por un ideal cuasicoherente $\mathscr{J}$ de $\sh{O}_\mathrm{X}$, y sea $\sh{F}$ un $\sh{O}_\mathrm{X}$-módulo cuasicoherente. Se dice que $\sh{F}$ es \emph{normalmente plano a lo largo de $\mathrm{Y}$ en un punto $y$} si $(\operatorname{gr}^*_\mathscr{J}(\sh{F}))_y$ es un $\sh{O}_{\mathrm{Y},y}$-módulo plano. Se dice que $\sh{F}$ es \emph{normalmente plano a lo largo de $\mathrm{Y}$} si lo es en todo punto de $\mathrm{Y}$.
\end{proof}

\begin{corollary}[11.3.5]
\label{IV.11.3.5}
Bajo las hipótesis generales de \sref{IV.11.3.4}, supongamos que $\operatorname{gr}^*_\mathscr{J}(\sh{O}_\mathrm{X})$ y $\operatorname{gr}^*_\mathscr{J}(\sh{F})$ son $f$-planos. Entonces el conjunto $\mathrm{W}$ de los $y\in\mathrm{Y}$ tales que $\sh{F}$ es normalmente plano a lo largo de $\mathrm{Y}$ en el punto $y$ es abierto en $\mathrm{Y}$ y (con las notaciones de \sref{IV.11.3.4},~(iv)) $\sh{F}'$ es normalmente plano a lo largo de $\mathrm{Y}'$ en todos los puntos de $\mathrm{W}'=p^{-1}(\mathrm{W})$.
\end{corollary}

\begin{proposition}[11.3.6]
\label{IV.11.3.6}
Con las notaciones de $(8.5.1)$ y $(8.8.1)$, supongamos que $\mathrm{S}_\alpha$ es cuasicompacto, que $\mathrm{X}_\alpha$ es de presentación finita sobre $\mathrm{S}_\alpha$ y que $\mathrm{Y}_\alpha$ es un subpreesquema cerrado de $\mathrm{X}_\alpha$ definido por un ideal cuasicoherente $\mathscr{J}_\alpha$ de tipo finito de $\sh{O}_{\mathrm{X}_\alpha}$; supongamos que $\operatorname{gr}^*_{\mathscr{J}_\alpha}(\sh{O}_{\mathrm{X}_\alpha})$ es plano sobre $\mathrm{S}_\alpha$; supongamos, por último, que $\sh{F}_\alpha$ es un $\sh{O}_{\mathrm{X}_\alpha}$-módulo de presentación finita. Para que $\sh{F}$ sea normalmente plano a lo largo de $\mathrm{Y}$, es necesario y suficiente que exista $\lambda$ tal que $\sh{F}_\lambda$ sea normalmente plano a lo largo de $\mathrm{Y}_\lambda$.
\end{proposition}

\begin{proof}
Obsérvese que $\mathrm{Y}$ (resp.\ $\mathrm{Y}_\lambda$) es el subpreesquema cerrado de $\mathrm{X}$ (resp.\ $\mathrm{X}_\lambda$) definido por $\mathscr{J}=\mathscr{J}_\alpha\sh{O}_\mathrm{X}$ (resp.\ $\mathscr{J}_\lambda=\mathscr{J}_\alpha\sh{O}_{\mathrm{X}_\lambda}$). Por la hipótesis y \sref{IV.11.2.9.2}, se tiene
\[
\operatorname{gr}^*_\mathscr{J}(\sh{O}_\mathrm{X})
=\operatorname{gr}^*_{\mathscr{J}_\alpha}(\sh{O}_{\mathrm{X}_\alpha})
 \otimes_{\sh{O}_{\mathrm{S}_\alpha}}\sh{O}_\mathrm{S},
\qquad
\operatorname{gr}^*_{\mathscr{J}_\lambda}(\sh{O}_{\mathrm{X}_\lambda})
=\operatorname{gr}^*_{\mathscr{J}_\alpha}(\sh{O}_{\mathrm{X}_\alpha})
 \otimes_{\sh{O}_{\mathrm{S}_\alpha}}\sh{O}_{\mathrm{S}_\lambda}
\quad(\lambda\geqslant\alpha),
\]
y $\operatorname{gr}^*_\mathscr{J}(\sh{O}_\mathrm{X})$ (resp.\ $\operatorname{gr}^*_{\mathscr{J}_\lambda}(\sh{O}_{\mathrm{X}_\lambda})$) es plano sobre $\mathrm{S}$ (resp.\ $\mathrm{S}_\lambda$). Se deduce que $\mathrm{Y}$ es plano sobre $\mathrm{S}$ (resp.\ $\mathrm{Y}_\lambda$ es plano sobre $\mathrm{S}_\lambda$). Si $\sh{F}_\lambda$ es normalmente plano a lo largo de $\mathrm{Y}_\lambda$, entonces $\left.\operatorname{gr}^*_{\mathscr{J}_\lambda}(\sh{F}_\lambda)\right|_{\mathrm{Y}_\lambda}$ es plano sobre $\mathrm{Y}_\lambda$ y, por tanto, también sobre $\mathrm{S}_\lambda$; y como
\[
\left.\operatorname{gr}^n_{\mathscr{J}_\lambda}(\sh{F}_\lambda)\right|_{\mathrm{X}_\lambda\setminus\mathrm{Y}_\lambda}=0
\quad\text{para }n\geqslant0,
\]
$\operatorname{gr}^*_{\mathscr{J}_\lambda}(\sh{F}_\lambda)$ es plano sobre $\mathrm{S}_\lambda$. Por \sref{IV.11.2.9.2},
\[
\operatorname{gr}^*_\mathscr{J}(\sh{F})
=\operatorname{gr}^*_{\mathscr{J}_\lambda}(\sh{F}_\lambda)
 \otimes_{\sh{O}_{\mathrm{S}_\lambda}}\sh{O}_\mathrm{S},
\]
por lo que $\sh{F}$ es normalmente plano a lo largo de $\mathrm{Y}$.

Recíprocamente, podemos suponer que $\mathrm{S}_\alpha$ y $\mathrm{X}_\alpha$ son afines y adoptar las notaciones de \sref{IV.11.2.9}. Puesto que $\operatorname{gr}^*_\mathfrak{J}(\mathrm{B})$ es un $\mathrm{A}$-módulo plano, lo mismo ocurre con $\operatorname{gr}^*_\mathfrak{J}(\mathrm{M})$; por tanto, en virtud de \sref{IV.11.2.9}, existe $\lambda\geqslant\alpha$ tal que $\operatorname{gr}^*_{\mathfrak{J}_\lambda}(\mathrm{M}_\lambda)$ es un $\mathrm{A}_\lambda$-módulo plano y
\[
\operatorname{gr}^*_\mathfrak{J}(\mathrm{M})
=\operatorname{gr}^*_{\mathfrak{J}_\lambda}(\mathrm{M}_\lambda)
 \otimes_{\mathrm{A}_\lambda}\mathrm{A}.
\]
Además, por \sref{IV.11.3.4},~(ii), $\operatorname{gr}^*_{\mathfrak{J}_\lambda}(\mathrm{M}_\lambda)$ es un $\operatorname{gr}^*_{\mathfrak{J}_\lambda}(\mathrm{B}_\lambda)$-módulo de presentación finita y $\operatorname{gr}^*_{\mathfrak{J}_\lambda}(\mathrm{B}_\lambda)$ es una $(\mathrm{B}_\lambda/\mathfrak{J}_\lambda)$-álgebra de presentación finita. La conclusión se deduce ahora de \sref{IV.11.2.6}.
\end{proof}

\oldpage[IV-3]{135}
\begin{proposition}[11.3.7]
\label{IV.11.3.7}
Sea $f:\mathrm{X}\to\mathrm{Y}$ un morfismo localmente de presentación finita, sean $\sh{F}$ y $\sh{F}'$ dos $\sh{O}_\mathrm{X}$-módulos cuasicoherentes de presentación finita, sea $u:\sh{F}'\to\sh{F}$ un homomorfismo de $\sh{O}_\mathrm{X}$-módulos, sea $x$ un punto de $\mathrm{X}$ y sea $y=f(x)$; supongamos que $\sh{F}$ es $f$-plano en el punto $x$. Las condiciones siguientes son equivalentes:
\begin{enumerate}[label=\emph{\alph*)}]
  \item $(\operatorname{Ker} u)_x=0$ y $\operatorname{Coker} u$ es un $\sh{O}_\mathrm{X}$-módulo $f$-plano en el punto $x$.
  \item El homomorfismo $(u\otimes 1)_x:(\sh{F}'\otimes_{\sh{O}_\mathrm{Y}}\boldsymbol{k}(y))_x\to(\sh{F}\otimes_{\sh{O}_\mathrm{Y}}\boldsymbol{k}(y))_x$ es inyectivo.
\end{enumerate}
Si además $\sh{F}$ es $f$-plano, el conjunto de los puntos $x$ que satisfacen las condiciones equivalentes anteriores es abierto en $\mathrm{X}$.
\end{proposition}

\begin{proof}
La condición $a)$ implica $b)$ en virtud de $(\mathbf{0}_\mathrm{I},~6.1.2)$, sin hipótesis alguna sobre $\sh{F}$. Para demostrar la recíproca, podemos limitarnos al caso en que $\mathrm{X}$ y $\mathrm{Y}$ son afines y, a continuación, en virtud de $(8.9.1)$ y \sref{IV.11.2.7}, suponer que
\[
  \mathrm{X}=\mathrm{X}_0\times_{\mathrm{Y}_0}\mathrm{Y},\qquad
  f=f_0\times 1_\mathrm{Y},\qquad
  \sh{F}=\sh{F}_0\otimes_{\sh{O}_{\mathrm{Y}_0}}\sh{O}_\mathrm{Y},\qquad
  \sh{F}'=\sh{F}'_0\otimes_{\sh{O}_{\mathrm{Y}_0}}\sh{O}_\mathrm{Y},\qquad
  u=u_0\otimes 1_\mathrm{Y},
\]
en donde $\mathrm{Y}_0$ es noetheriano, $f_0:\mathrm{X}_0\to\mathrm{Y}_0$ es un morfismo de tipo finito, $\sh{F}_0$ y $\sh{F}'_0$ son dos $\sh{O}_{\mathrm{X}_0}$-módulos coherentes, y $u_0:\sh{F}'_0\to\sh{F}_0$ es un homomorfismo. Además, si $x_0$ es la proyección de $x$ en $\mathrm{X}_0$, podemos suponer que $\sh{F}_0$ es $f_0$-plano en $x_0$. Pongamos $y_0=f_0(x_0)$. Por la transitividad de las fibras $(\mathbf{I},~3.6.4)$, la proyección $f^{-1}(y)\to f_0^{-1}(y_0)$ es fielmente plana $(2.2.13)$; y como
\[
  (u\otimes 1)_x=((u_0\otimes 1)_{x_0})\otimes 1_{\boldsymbol{k}(y)},
\]
la hipótesis de que $(u\otimes 1)_x$ es inyectivo implica que $(u_0\otimes 1)_{x_0}$ es inyectivo $(2.2.7)$. Aplicando $(\mathbf{0}_{\mathrm{III}},~10.2.4)$ a los anillos locales noetherianos $\sh{O}_{\mathrm{Y}_0,y_0}$ y $\sh{O}_{\mathrm{X}_0,x_0}$ y a los $\sh{O}_{\mathrm{X}_0,x_0}$-módulos $(\sh{F}'_0)_{x_0}$ y $(\sh{F}_0)_{x_0}$, el segundo de los cuales es plano sobre $\sh{O}_{\mathrm{Y}_0,y_0}$, obtenemos
\[
  (\operatorname{Ker}u_0)_{x_0}=0
  \quad\text{y}\quad
  \operatorname{Coker}u_0\ \text{es $f_0$-plano en }x_0.
\]
Se deduce, en primer lugar, que $\operatorname{Coker}u$ es $f$-plano en $x$ $(2.1.4)$. Además, por la exactitud por la derecha del producto tensorial,
\[
  \operatorname{Coker}u=(\operatorname{Coker}u_0)
  \otimes_{\sh{O}_{\mathrm{Y}_0}}\sh{O}_\mathrm{Y}.
\]
Aplicando $(\mathbf{0}_\mathrm{I},~6.1.2)$ a la sucesión exacta
\[
  0\longrightarrow(\sh{F}'_0)_{x_0}
  \longrightarrow(\sh{F}_0)_{x_0}
  \longrightarrow(\operatorname{Coker}u_0)_{x_0}
  \longrightarrow0,
\]
concluimos que $u_x$ es inyectivo, lo que demuestra la condición $a)$.

Por último, se deduce de \sref{IV.11.1.2} que el conjunto $\mathrm{U}_0$ de los puntos $z_0\in\mathrm{X}_0$ tales que $(u_0\otimes 1)_{z_0}$ es inyectivo es abierto. Por planitud, para todo $z\in\mathrm{X}$ situado sobre un punto $z_0\in\mathrm{U}_0$, el homomorfismo $(u\otimes 1)_z$ es inyectivo. El conjunto de tales puntos contiene, por tanto, la imagen inversa $\mathrm{U}$ de $\mathrm{U}_0$ en $\mathrm{X}$ y es, en consecuencia, una vecindad de $x$, lo que termina la demostración.
\end{proof}

\begin{theorem}[11.3.8]
\label{IV.11.3.8}
Sea $f:\mathrm{X}\to\mathrm{Y}$ un morfismo localmente de presentación finita, sea $\sh{F}$ un $\sh{O}_\mathrm{X}$-módulo cuasicoherente de presentación finita y sea $(g_i)_{1\leqslant i\leqslant n}$ una sucesión de secciones de $\sh{O}_\mathrm{X}$ sobre $\mathrm{X}$; pongamos $\sh{F}_i=\sh{F}/\sum_{j=1}^{i}g_j\sh{F}$ para $0\leqslant i\leqslant n$ (con $\sh{F}_0=\sh{F}$). Sea $x$ un punto de $\mathrm{X}$, sea $y=f(x)$ y pongamos $\mathrm{X}_y=f^{-1}(y)=\mathrm{X}\times_\mathrm{Y}\Spec(\boldsymbol{k}(y))$, $\sh{F}_y=\sh{F}\otimes_{\sh{O}_y}\boldsymbol{k}(y)$, que es un $\sh{O}_{\mathrm{X}_y}$-módulo. Supongamos que los $(g_i)_x$ pertenecen al ideal maximal $\mathfrak{m}_x$. Las condiciones siguientes son equivalentes:
\begin{enumerate}[label=\emph{\alph*)}]
  \item La sucesión $((g_i)_x)$ es $\sh{F}_x$-regular y los $\sh{F}_i$ ($0\leqslant i\leqslant n$) son $f$-planos en el punto $x$.
  \item La sucesión $((g_i)_x)$ es $\sh{F}_x$-regular y $\sh{F}_n$ es $f$-plano en el punto $x$.
  \item[\emph{b$'$)}] Existe una vecindad $\mathrm{U}$ de $x$ tal que la sucesión $(g_i|\mathrm{U})$ es $(\sh{F}|\mathrm{U})$-cuasirregular, y $\sh{F}_n$ es $f$-plano en el punto $x$.
\oldpage[IV-3]{136}
  \item $\sh{F}$ es $f$-plano en el punto $x$, y la sucesión $((g_i\otimes 1)_x)$ de elementos de $\sh{O}_{\mathrm{X}_y,x}$, imágenes de los $(g_i)_x$, es $(\sh{F}_y)_x$-regular.
  \item $\sh{F}$ es $f$-plano en el punto $x$, y para todo morfismo $\mathrm{Y}'\to\mathrm{Y}$ y todo punto $x'\in\mathrm{X}'=\mathrm{X}\times_\mathrm{Y}\mathrm{Y}'$ situado sobre $x$, si se pone $\sh{F}'=\sh{F}\otimes_\mathrm{Y}\mathrm{Y}'$, la sucesión $(g_i\otimes 1)_{x'}$ de elementos de $\sh{O}_{\mathrm{X}',x'}$ es $\sh{F}'_{x'}$-regular.
\end{enumerate}
Además, el conjunto de los $x\in\mathrm{X}$ que satisfacen estas condiciones es abierto en el conjunto $\mathrm{Z}$ de los $x$ tales que $g_i(x)=0$ para todo $i$.
\end{theorem}

\begin{proof}
El hecho de que $a)$ implique $d)$ se deduce inmediatamente de $(\mathbf{0},~15.1.15)$, y $c)$ es un caso particular de $d)$; además, $a)$ implica trivialmente $b)$. Demostremos que $b)$ o $c)$ implica $a)$. Como los $\sh{F}_i$ son de presentación finita, el hecho de que $c)$ implique $a)$ se deduce inmediatamente de \sref{IV.11.3.7} por inducción sobre $n$, y esto muestra también que el conjunto de los $x\in\mathrm{X}$ que satisfacen $c)$ es abierto en $\mathrm{Z}$.

Para demostrar que $b)$ implica $a)$, la inducción sobre $n$ nos reduce inmediatamente al caso $n=1$; escribimos $g$ en lugar de $g_1$. Como la cuestión es local en $\mathrm{X}$ y en $\mathrm{Y}$, podemos suponer que $\mathrm{Y}=\Spec(\mathrm{A})$ y $\mathrm{X}=\Spec(\mathrm{B})$ son afines, que $\mathrm{B}$ es una $\mathrm{A}$-álgebra de presentación finita y que $\sh{F}=\widetilde{\mathrm{M}}$, donde $\mathrm{M}$ es un $\mathrm{B}$-módulo de presentación finita. Por $(8.9.1)$, podemos escribir entonces
\[
  \mathrm{B}=\mathrm{B}_0\otimes_{\mathrm{A}_0}\mathrm{A},
  \qquad
  \mathrm{M}=\mathrm{M}_0\otimes_{\mathrm{A}_0}\mathrm{A},
\]
en donde $\mathrm{A}_0$ es un subanillo noetheriano de $\mathrm{A}$, $\mathrm{B}_0$ es una $\mathrm{A}_0$-álgebra de tipo finito y $\mathrm{M}_0$ es un $\mathrm{B}_0$-módulo de tipo finito. Además, $\mathrm{A}$ es el límite inductivo filtrante de sus subanillos $\mathrm{A}_\lambda$ que son $\mathrm{A}_0$-álgebras de tipo finito (y, por tanto, noetherianos). Si ponemos
\[
  \mathrm{B}_\lambda=\mathrm{B}_0\otimes_{\mathrm{A}_0}\mathrm{A}_\lambda,
  \qquad
  \mathrm{M}_\lambda=\mathrm{M}_0\otimes_{\mathrm{A}_0}\mathrm{A}_\lambda,
\]
entonces $\mathrm{B}_\lambda$ es una $\mathrm{A}_\lambda$-álgebra de tipo finito, $\mathrm{M}_\lambda$ es un $\mathrm{B}_\lambda$-módulo de tipo finito, y
\[
  \mathrm{B}=\varinjlim_\lambda\mathrm{B}_\lambda,
  \qquad
  \mathrm{M}=\varinjlim_\lambda\mathrm{M}_\lambda.
\]
Existen, por tanto, un índice $\lambda$ y un elemento $g_\lambda\in\mathrm{B}_\lambda$ tales que $g=g_\lambda\otimes 1$. Volviendo a la notación geométrica y poniendo
\[
  \mathrm{Y}_\lambda=\Spec(\mathrm{A}_\lambda),
  \qquad
  \mathrm{X}_\lambda=\Spec(\mathrm{B}_\lambda),
  \qquad
  \sh{F}_\lambda=\widetilde{\mathrm{M}_\lambda},
\]
bastará demostrar que existe $\mu\geqslant\lambda$ tal que, en el punto $x_\mu\in\mathrm{X}_\mu$ que es la proyección de $x$, $(g_\mu)_{x_\mu}$ es $(\sh{F}_\mu)_{x_\mu}$-regular y $\sh{F}_\mu/g_\mu\sh{F}_\mu$ es plano sobre $\mathrm{Y}_\mu$ en $x_\mu$. En efecto, se deducirá entonces de $(\mathbf{0},~15.1.16)$ que $\sh{F}_\mu$ es plano sobre $\mathrm{Y}_\mu$ en $x_\mu$ y, por consiguiente, que $\sh{F}$ es plano sobre $\mathrm{Y}$ en $x$.

El hecho de que $b)$ implique $a)$ se deducirá, por tanto, de la proposición siguiente.

\begin{proposition}[11.3.9]
\label{IV.11.3.9}
Con las notaciones y las hipótesis generales de $(8.5.1)$ y $(8.8.1)$, supongamos que $f_\alpha:\mathrm{X}_\alpha\to\mathrm{Y}_\alpha$ es un morfismo localmente de presentación finita. Sean $\sh{F}_\alpha$ y $\sh{G}_\alpha$ dos $\sh{O}_{\mathrm{X}_\alpha}$-módulos cuasicoherentes de presentación finita, sea $h_\alpha:\sh{F}_\alpha\to\sh{G}_\alpha$ un homomorfismo de $\sh{O}_{\mathrm{X}_\alpha}$-módulos y sea $\sh{H}_\alpha$ su conúcleo. Por último, sea $x$ un punto de $\mathrm{X}$ y sea $x_\lambda$ su proyección en $\mathrm{X}_\lambda$ para $\lambda\geqslant\alpha$. Para que $h_x$ sea inyectivo y $\sh{H}=\operatorname{Coker}h$ sea $f$-plano en $x$, es necesario y suficiente que exista $\lambda\geqslant\alpha$ tal que $(h_\lambda)_{x_\lambda}$ sea inyectivo y $\sh{H}_\lambda=\operatorname{Coker}h_\lambda$ sea $f_\lambda$-plano en $x_\lambda$. Además, el conjunto de los $x\in\mathrm{X}$ que poseen las propiedades anteriores es abierto en $\mathrm{X}$.
\end{proposition}

\begin{proof}
Recordemos que, por la exactitud por la derecha del producto tensorial,
\[
  \sh{H}_\mu=v^*_{\lambda\mu}(\sh{H}_\lambda)
  \quad\text{para }\lambda\leqslant\mu,
  \qquad
  \sh{H}=v^*_\lambda(\sh{H}_\lambda),
\]
lo que justifica las notaciones.

La suficiencia de la condición se deduce de que, si la sucesión
\[
  0\longrightarrow(\sh{F}_\lambda)_{x_\lambda}
  \longrightarrow(\sh{G}_\lambda)_{x_\lambda}
  \longrightarrow(\sh{H}_\lambda)_{x_\lambda}
  \longrightarrow0
\]
es exacta y $(\sh{H}_\lambda)_{x_\lambda}$ es un $\sh{O}_{f_\lambda(x_\lambda)}$-módulo plano, entonces $\sh{H}_x$ es un $\sh{O}_{f(x)}$-módulo plano por cambio de base $(2.1.4)$ y la sucesión
\[
  0\longrightarrow\sh{F}_x
  \longrightarrow\sh{G}_x
  \longrightarrow\sh{H}_x
  \longrightarrow0
\]
es exacta $(2.1.8)$. Para demostrar que

\oldpage[IV-3]{137}

la condición es necesaria, observemos que $\sh{H}$ es de presentación finita. Como la cuestión es local en $\mathrm{X}$ y en $\mathrm{Y}$, podemos suponer que $\mathrm{X}$ y $\mathrm{Y}$ son afines y, por \sref{IV.11.3.1}, que $\sh{H}$ es $f$-plano. Enunciemos ahora el lema siguiente.

\begin{lemma}[11.3.9.1]
\label{IV.11.3.9.1}
Sea $f:\mathrm{X}\to\mathrm{Y}$ un morfismo localmente de presentación finita, sean $\sh{G}$ y $\sh{H}$ dos $\sh{O}_\mathrm{X}$-módulos cuasicoherentes de presentación finita tales que $\sh{H}$ es $f$-plano, y sea $p:\sh{G}\to\sh{H}$ un homomorfismo sobreyectivo. Entonces $\operatorname{Ker}(p)$ es un $\sh{O}_\mathrm{X}$-módulo de presentación finita.
\end{lemma}

\begin{proof}
Podemos, en efecto, suponer que $\mathrm{X}$ y $\mathrm{Y}$ son afines y que existen un morfismo $\mathrm{Y}\to\mathrm{Y}_0$, donde $\mathrm{Y}_0$ es noetheriano; un morfismo $f_0:\mathrm{X}_0\to\mathrm{Y}_0$ de tipo finito tal que $\mathrm{X}=\mathrm{X}_0\times_{\mathrm{Y}_0}\mathrm{Y}$ y $f=(f_0)_{(\mathrm{Y})}$; dos $\sh{O}_{\mathrm{X}_0}$-módulos coherentes $\sh{G}_0$, $\sh{H}_0$; y un homomorfismo $p_0:\sh{G}_0\to\sh{H}_0$ tales que $\sh{G}$, $\sh{H}$ y $p$ se obtienen de $\sh{G}_0$, $\sh{H}_0$ y $p_0$ por cambio de base $(8.9.1\text{ y }8.5.2)$. Podemos suponer además que $p_0$ es sobreyectivo $(8.5.7)$ y que $\sh{H}_0$ es $f_0$-plano \sref{IV.11.2.7}. Entonces, si $\sh{K}_0=\operatorname{Ker}(p_0)$, $\sh{K}_0$ es un $\sh{O}_{\mathrm{X}_0}$-módulo coherente $(\mathbf{0}_\mathrm{I},~5.3.4)$ y, por $(\mathbf{0}_\mathrm{I},~6.1.2)$, $\sh{K}=\operatorname{Ker}(p)$ se obtiene de $\sh{K}_0$ por cambio de base y es, por tanto, de presentación finita.
\end{proof}

Demostrado este lema, pongamos $\sh{K}=\operatorname{Ker}(\sh{G}\to\sh{H})$, que es, por tanto, de presentación finita. Existe un homomorfismo canónico $q:\sh{F}\to\sh{K}$ y, por hipótesis, $q_x:\sh{F}_x\to\sh{K}_x$ es un isomorfismo. Por consiguiente, según $(\mathbf{0}_\mathrm{I},~5.2.7)$, existe una vecindad $\mathrm{U}$ de $x$ tal que $q|_{\mathrm{U}}$ es un isomorfismo; tras restringir $\mathrm{X}$, podemos suponer que la sucesión
\[
  0\to\sh{F}\xrightarrow{h}\sh{G}\to\sh{H}\to 0
\]
es \emph{exacta}. En estas condiciones, se deduce de \sref{IV.11.2.6} que, para todo $\lambda$ suficientemente grande, $\sh{H}_\lambda$ es $f_\lambda$-plano. Pongamos $\sh{K}_\lambda=\operatorname{Ker}(\sh{G}_\lambda\to\sh{H}_\lambda)$ y $\sh{K}_\mu=v^*_{\lambda\mu}(\sh{K}_\lambda)$ para $\mu\geqslant\lambda$. Se deduce de $(\mathbf{0}_\mathrm{I},~6.1.2)$ que
\[
  \sh{K}_\mu=\operatorname{Ker}(\sh{G}_\mu\to\sh{H}_\mu)
  \quad\text{y}\quad
  \sh{K}=v^*_\lambda(\sh{K}_\lambda)
  =\operatorname{Ker}(\sh{G}\to\sh{H})=\sh{F}.
\]
Por otra parte, para todo $\mu\geqslant\lambda$ existe un homomorfismo canónico $q_\mu:\sh{F}_\mu\to\sh{K}_\mu$ tal que $q_\mu=v^*_{\lambda\mu}(q_\lambda)$ y $q=v^*_\mu(q_\mu)$. Como $q$ es un isomorfismo, se deduce de $(8.5.2.4)$ y \sref{IV.11.3.9.1} que existe $\mu\geqslant\lambda$ para el cual $q_\mu$ es un isomorfismo y, por consiguiente, $h_\mu:\sh{F}_\mu\to\sh{G}_\mu$ es inyectivo.
\end{proof}

Volvamos a la demostración de \sref{IV.11.3.8}.

Como el conjunto de los $x\in\mathrm{X}$ que satisfacen $b)$ es abierto en $\mathrm{X}$, está claro que $b)$ implica $b')$. Demostremos, por último, que $b')$ implica $c)$. En primer lugar, $\sh{F}_n$ es $f$-plano en una vecindad de $x$ \sref{IV.11.3.1}; podemos, por tanto, reducirnos al caso en que $\mathrm{U}=\mathrm{X}$ y $\sh{F}_n$ es $f$-plano. Puesto que, por definición, $\operatorname{gr}^*_\mathscr{J}(\sh{F})$ es isomorfo a $\sh{F}_n\otimes_{\sh{O}_\mathrm{Y}}\sh{O}_\mathrm{Y}[\mathrm{T}_1,\ldots,\mathrm{T}_n]$ $(\mathbf{0},~15.2.2)$, es $f$-plano, y se deduce de \sref{IV.11.3.4},~(i), que el propio $\sh{F}$ es $f$-plano en una vecindad de $x$. Por otra parte, si $(\mathscr{J}_y)_x$ es el ideal de $\sh{O}_{\mathrm{X}_y,x}$ generado por los $(g_i\otimes 1)_x$, se deduce de \sref{IV.11.2.9.2} que, en el diagrama
\[
\xymatrix{
  ((\operatorname{gr}^0_{\mathscr{J}_x}(\sh{F}_x))[\mathrm{T}_1,\ldots,\mathrm{T}_n])\otimes_{\sh{O}_y}\boldsymbol{k}(y) \ar[r] \ar[d] & \operatorname{gr}^*_{\mathscr{J}_x}(\sh{F}_x)\otimes_{\sh{O}_y}\boldsymbol{k}(y) \ar[d] \\
  (\operatorname{gr}^0_{(\mathscr{J}_y)_x}((\sh{F}_y)_x))[\mathrm{T}_1,\ldots,\mathrm{T}_n] \ar[r] & \operatorname{gr}^*_{(\mathscr{J}_y)_x}((\sh{F}_y)_x)
}
\]
\oldpage[IV-3]{138}
las flechas verticales son isomorfismos. Como la primera flecha horizontal es un isomorfismo, también lo es la segunda; por consiguiente, la sucesión $((g_i\otimes 1)_x)$ es $(\sh{F}_y)_x$-cuasirregular y, por tanto, regular, ya que $\mathrm{X}_y$ es localmente de tipo finito sobre $\boldsymbol{k}(y)$.
\end{proof}

\begin{theorem}[11.3.10]
\label{IV.11.3.10}
(Criterio de planitud por las fibras). --- Sea $\mathrm{S}$ un preesquema, sean $g:\mathrm{X}\to\mathrm{S}$ y $h:\mathrm{Y}\to\mathrm{S}$ dos morfismos, sea $f:\mathrm{X}\to\mathrm{Y}$ un $\mathrm{S}$-morfismo, sea $\sh{F}$ un $\sh{O}_\mathrm{X}$-módulo cuasicoherente, sea $x$ un punto de $\mathrm{X}$, sea $y=f(x)$ y sea $s=h(y)=g(x)$. Supongamos que se cumple una de las dos hipótesis siguientes:
\begin{enumerate}[label=\arabic{enumi}\textsuperscript{o}]
  \item $\mathrm{S}$, $\mathrm{X}$ y $\mathrm{Y}$ son localmente noetherianos y $\sh{F}$ es coherente.
  \item $g$ y $h$ son localmente de presentación finita y $\sh{F}$ es de presentación finita.
\end{enumerate}
Entonces, con las notaciones de $(9.4.1)$, si $\sh{F}_x\neq 0$, las condiciones siguientes son equivalentes:
\begin{enumerate}[label=\emph{\alph*)}]
  \item $\sh{F}$ es $g$-plano en el punto $x$ y $\sh{F}_s$ es $f_s$-plano en el punto $x$.
  \item $h$ es plano en el punto $y$ y $\sh{F}$ es $f$-plano en el punto $x$.
\end{enumerate}
Además, bajo la hipótesis 2\textsuperscript{o}, el conjunto de los $x\in\mathrm{X}$ que satisfacen la condición $b)$ es abierto en $\mathrm{X}$.
\end{theorem}

\begin{proof}
La última afirmación resulta de aplicar \sref{IV.11.3.1}, por una parte, a $\sh{O}_\mathrm{Y}$ y al morfismo $h$ y, por otra, a $\sh{F}$ y al morfismo $f$ (que es localmente de presentación finita).

Como $g=h\circ f$, está claro que $b)$ implica $a)$ sin suponer 1\textsuperscript{o} ni 2\textsuperscript{o} $(2.1.6$ y $2.1.4)$. Para demostrar que $a)$ implica $b)$, podemos reducirnos al caso en que $\mathrm{S}$, $\mathrm{X}$ y $\mathrm{Y}$ son afines; bajo la hipótesis 2\textsuperscript{o}, aplicando \sref{IV.11.2.7}, nos reducimos además al caso en que $\mathrm{S}$ es \emph{noetheriano}, es decir, al caso en que se cumple la hipótesis 1\textsuperscript{o}. La afirmación equivale entonces al lema siguiente, que mejora $(\mathbf{0}_{\mathrm{III}},~10.2.5)$:

\begin{lemma}[11.3.10.1]
\label{IV.11.3.10.1}
Sean $\rho:\mathrm{A}\to\mathrm{B}$ y $\sigma:\mathrm{B}\to\mathrm{C}$ homomorfismos de anillos locales noetherianos, sea $k$ el cuerpo residual de $\mathrm{A}$ y sea $\mathrm{M}$ un $\mathrm{C}$-módulo $\neq 0$ de tipo finito. Entonces las condiciones siguientes son equivalentes:
\begin{enumerate}[label=\emph{\alph*)}]
  \item $\mathrm{M}$ es un $\mathrm{A}$-módulo plano y $\mathrm{M}\otimes_\mathrm{A} k$ es un $(\mathrm{B}\otimes_\mathrm{A} k)$-módulo plano.
  \item $\mathrm{B}$ es un $\mathrm{A}$-módulo plano y $\mathrm{M}$ es un $\mathrm{B}$-módulo plano.
\end{enumerate}
\end{lemma}

Estableceremos primero el lema más general siguiente:

\begin{lemma}[11.3.10.2]
\label{IV.11.3.10.2}
Sea $\mathrm{A}$ un anillo conmutativo, sea $\mathrm{B}$ una $\mathrm{A}$-álgebra conmutativa, sea $\mathfrak{J}$ un ideal de $\mathrm{A}$ y sea $\mathrm{M}$ un $\mathrm{B}$-módulo. Consideremos primero las condiciones siguientes:
\begin{enumerate}[label=\emph{(\roman*)}]
  \item $\mathfrak{J}$ es nilpotente.
  \item $\mathrm{B}$ es noetheriano y $\mathrm{M}$ está idealmente separado para la topología $\mathfrak{J}\mathrm{B}$-preádica $(\mathbf{0}_{\mathrm{III}},~10.2.1)$.
  \item $\mathfrak{J}\mathrm{B}$ está contenido en el radical de $\mathrm{B}$.
\end{enumerate}
Consideremos, por otra parte, las cuatro propiedades siguientes:
\begin{enumerate}[label=\emph{\alph*)}]
  \item $\mathrm{M}$ es un $\mathrm{B}$-módulo plano.
  \item $\mathrm{B}$ es un $\mathrm{A}$-módulo plano.
  \item $\mathrm{M}$ es un $\mathrm{A}$-módulo plano y $\mathrm{M}/\mathfrak{J}\mathrm{M}$ es un $(\mathrm{B}/\mathfrak{J}\mathrm{B})$-módulo plano.
  \item $\mathrm{B}/\mathfrak{J}\mathrm{B}$ es un $(\mathrm{A}/\mathfrak{J})$-módulo plano y, para todo ideal maximal $\mathfrak{m}\supset\mathfrak{J}\mathrm{B}$ de $\mathrm{B}$, se tiene $\mathfrak{m}\mathrm{M}\neq\mathrm{M}$.
\end{enumerate}
Entonces:
\begin{enumerate}[label=\arabic{enumi}\textsuperscript{o}]
  \item Si se cumple una de las condiciones (i) o (ii), la conjunción de $a)$ y $b)$ implica $c)$, y $c)$ implica $a)$.
  \item \oldpage[IV-3]{139} Si se cumple la condición (i) o la conjunción de (ii) y (iii), la conjunción de $c)$ y $d)$ implica la conjunción de $a)$ y $b)$.
\end{enumerate}
\end{lemma}

\begin{proof}
1\textsuperscript{o} La primera afirmación es inmediata $(\mathbf{0}_\mathrm{I},~6.2.1)$. Supongamos, pues, que se cumple $c)$ y demostremos $a)$. Consideremos los anillos graduados $\operatorname{gr}^*_\mathfrak{J}(\mathrm{A})$, $\operatorname{gr}^*_\mathfrak{J}(\mathrm{B})$ y el módulo graduado $\operatorname{gr}^*_\mathfrak{J}(\mathrm{M})$ (tanto sobre $\operatorname{gr}^*_\mathfrak{J}(\mathrm{A})$ como sobre $\operatorname{gr}^*_\mathfrak{J}(\mathrm{B})$), relativos a las filtraciones $\mathfrak{J}$-preádicas, así como las aplicaciones canónicas sobreyectivas $(\mathbf{0}_{\mathrm{III}},~10.1.1.2)$
\begin{align*}
  u &: \operatorname{gr}^0(\mathrm{B})\otimes_{\operatorname{gr}^0(\mathrm{A})}\operatorname{gr}^*(\mathrm{A})\to\operatorname{gr}^*(\mathrm{B}) \\
  v &: \operatorname{gr}^0(\mathrm{M})\otimes_{\operatorname{gr}^0(\mathrm{A})}\operatorname{gr}^*(\mathrm{A})\to\operatorname{gr}^*(\mathrm{M}) \\
  w &: \operatorname{gr}^0(\mathrm{M})\otimes_{\operatorname{gr}^0(\mathrm{B})}\operatorname{gr}^*(\mathrm{B})\to\operatorname{gr}^*(\mathrm{M})
\end{align*}
Es claro que se tiene un diagrama conmutativo
\begin{equation}
\xymatrix{
  \operatorname{gr}^0(\mathrm{M})\otimes_{\operatorname{gr}^0(\mathrm{A})}\operatorname{gr}^*(\mathrm{A}) \ar[rr]^{v} \ar[dr]_{1\otimes u} & & \operatorname{gr}^*(\mathrm{M}) \\
  & \operatorname{gr}^0(\mathrm{M})\otimes_{\operatorname{gr}^0(\mathrm{B})}\operatorname{gr}^*(\mathrm{B}) \ar[ur]_{w} &
}
  \label{IV.11.3.10.3}
\tag{11.3.10.3}
\end{equation}

La hipótesis $c)$ implica que $v$ es biyectiva $(\mathbf{0}_{\mathrm{III}},~10.2.1)$; como las otras dos aplicaciones del diagrama son sobreyectivas, también ellas son biyectivas. Puesto que, por la hipótesis $c)$, $\mathrm{M}/\mathfrak{J}\mathrm{M}$ es un $(\mathrm{B}/\mathfrak{J}\mathrm{B})$-módulo plano, se deduce de $(\mathbf{0}_{\mathrm{III}},~10.2.1)$ que $\mathrm{M}$ es un $\mathrm{B}$-módulo plano.

2\textsuperscript{o} Cada una de las condiciones (i), (iii) implica que todo ideal maximal de $\mathrm{B}$ contiene $\mathfrak{J}\mathrm{B}$. Se deduce, por tanto, de 1\textsuperscript{o} y de la conjunción de $c)$ y $d)$ que $\mathrm{M}$ es un $\mathrm{B}$-módulo \emph{fielmente plano} y, en consecuencia, que $\operatorname{gr}^0(\mathrm{M})$ es un $\operatorname{gr}^0(\mathrm{B})$-módulo \emph{fielmente plano} $(\mathbf{0}_{\mathrm{I}},~6.2.1)$. Hemos visto en 1\textsuperscript{o} que la hipótesis $c)$ implica que las tres aplicaciones del diagrama \textup{(11.3.10.3)} son biyectivas; la planitud fiel de $\operatorname{gr}^0(\mathrm{M})$ sobre $\operatorname{gr}^0(\mathrm{B})$ implica, por tanto, que $u$ es biyectiva $(\mathbf{0}_{\mathrm{I}},~6.4.1)$. Por otra parte, las condiciones (ii) y (iii) implican que $\mathrm{B}$, como $\mathrm{A}$-módulo, está idealmente separado para la filtración $\mathfrak{J}$-preádica (Bourbaki, \emph{Alg.\ comm.}, cap.~III, \textsection\,5, n\textsuperscript{o}\,4, prop.~2). Se deduce nuevamente de $(\mathbf{0}_{\mathrm{III}},~10.2.1)$ que, si se cumple la condición (i), o la conjunción de (ii) y (iii), $\mathrm{B}$ es un $\mathrm{A}$-módulo plano.
\end{proof}

\phantomsection\label{IV.11.3.10.4}
\textbf{(11.3.10.4)} Una vez establecido el lema \sref{IV.11.3.10.2}, se deduce \sref{IV.11.3.10.1} tomando como $\mathfrak{J}$ el ideal maximal de $\mathrm{A}$ y observando que se cumplen entonces las condiciones (ii) y (iii) de \sref{IV.11.3.10.2} (Bourbaki, \emph{Alg.\ comm.}, cap.~III, \textsection\,5, n\textsuperscript{o}\,4, prop.~2). Esto termina también la demostración de \sref{IV.11.3.10}.
\end{proof}

\begin{corollary}[11.3.11]
\label{IV.11.3.11}
Sean $g:\mathrm{X}\to\mathrm{S}$ y $h:\mathrm{Y}\to\mathrm{S}$ dos morfismos localmente de presentación finita, y sea $f:\mathrm{X}\to\mathrm{Y}$ un $\mathrm{S}$-morfismo. Las condiciones siguientes son equivalentes:
\begin{enumerate}[label=\emph{\alph*)}]
  \item $g$ es plano y $f_s:\mathrm{X}_s\to\mathrm{Y}_s$ es plano para todo $s\in\mathrm{S}$.
  \item $h$ es plano en todos los puntos de $f(\mathrm{X})$ y $f$ es plano.
\end{enumerate}
\end{corollary}

\begin{proof}
Basta aplicar \sref{IV.11.3.10} con $\sh{F}=\sh{O}_\mathrm{X}$.
\end{proof}

\begin{remark}[11.3.12]
\label{IV.11.3.12}
Sería interesante poder dar demostraciones de \sref{IV.11.3.2} y \sref{IV.11.3.10} que no utilizaran el paso al límite inductivo; para ello bastaría demostrar el criterio siguiente:

Sea $\mathrm{A}$ un anillo, sea $\mathrm{B}$ una $\mathrm{A}$-álgebra de presentación finita, sea $\mathrm{M}$ un $\mathrm{B}$-módulo de presentación finita, sea $\mathfrak{J}$ un \oldpage[IV-3]{140} ideal de $\mathrm{A}$ y sea $\mathfrak{p}$ un ideal primo de $\mathrm{B}$ que contiene $\mathfrak{J}\mathrm{B}$. Entonces las dos condiciones siguientes son equivalentes:
\begin{enumerate}[label=\emph{\alph*)}]
  \item $\mathrm{M}_\mathfrak{p}$ es un $\mathrm{A}$-módulo plano.
  \item $\mathrm{M}_\mathfrak{p}/\mathfrak{J}\mathrm{M}_\mathfrak{p}$ es un $(\mathrm{A}/\mathfrak{J})$-módulo plano y $\operatorname{Tor}_1^\mathrm{A}(\mathrm{A}/\mathfrak{J},\mathrm{M}_\mathfrak{p})=0$.
\end{enumerate}
Cuando $\mathrm{A}$ es \emph{noetheriano}, esto se deduce de $(\mathbf{0}_{\mathrm{III}},~10.2.2)$ aplicado a la $\mathrm{A}$-álgebra noetheriana $\mathrm{B}_\mathfrak{p}$; ¿puede deducirse de aquí el enunciado general mediante el paso al límite inductivo?

Conviene señalar, a este respecto, que tal generalización no es ciertamente posible si se sustituye la condición $b)$ anterior por una cualquiera de las condiciones $c)$, $d)$ de $(\mathbf{0}_{\mathrm{III}},~10.2.1)$, reemplazando $\mathrm{M}$ por $\mathrm{M}_\mathfrak{p}$. Por ejemplo, sea $\mathrm{A}$ un anillo local cuyo ideal maximal $\mathfrak{J}$ es principal y tal que
\[
  \mathfrak{R}=\bigcap_{n\geqslant 0}\mathfrak{J}^{n+1}
\]
no es nulo (por ejemplo, el anillo de gérmenes de funciones reales indefinidamente diferenciables en una vecindad de $0$ en $\mathbf{R}$). Tomemos $\mathrm{B}=\mathrm{A}$, $\mathfrak{p}=\mathfrak{J}$ y $\mathrm{M}=\mathrm{A}/\mathfrak{N}$, donde $\mathfrak{N}$ es un submódulo cíclico no nulo de $\mathfrak{R}$; $\mathrm{M}$ es, por tanto, de presentación finita. Es claro que $\mathrm{M}$ no es un $\mathrm{A}$-módulo plano, pues, al ser de presentación finita, sería entonces libre (Bourbaki, \emph{Alg.\ comm.}, cap.~II, \textsection\,3, n\textsuperscript{o}\,2, cor.~2 a prop.~5), lo cual es absurdo, ya que $\mathfrak{N}\neq 0$. Sin embargo,
\[
  \mathfrak{J}^k\mathrm{M}/\mathfrak{J}^{k+n}\mathrm{M}\simeq \mathrm{A}/\mathfrak{J}^n
\]
para todos los enteros positivos $k$ y $n$; por tanto, $\mathrm{M}$ sí satisface las condiciones $c)$ y $d)$ de $(\mathbf{0}_{\mathrm{III}},~10.2.1)$, ya que $\mathrm{A}/\mathfrak{J}$ es un cuerpo.
\end{remark}

\begin{proposition}[11.3.13]
\label{IV.11.3.13}
Sea $f:\mathrm{X}\to\mathrm{Y}$ un morfismo de preesquemas y sea $x$ un punto de $\mathrm{X}$ tal que $f$ es plano en $x$; pongamos $y=f(x)$.
\begin{enumerate}[label=\emph{(\roman*)}]
  \item Si $\mathrm{X}$ es reducido (resp.\ normal) en el punto $x$, entonces $\mathrm{Y}$ es reducido (resp.\ normal) en el punto $y$.
  \item Supongamos además que $f$ es de presentación finita en el punto $x$. Entonces, si $\mathrm{Y}$ es reducido (resp.\ normal) en el punto $y$ y si el morfismo $f$ es reducido (resp.\ normal) en el punto $x$ $(6.8.1)$, $\mathrm{X}$ es reducido (resp.\ normal) en el punto $x$.
\end{enumerate}
\end{proposition}

\begin{proof}
La primera afirmación se recuerda aquí únicamente por comodidad, pues ya se demostró en $(2.1.13)$. Para demostrar (ii), podemos limitarnos al caso en que $\mathrm{X}=\Spec(\mathrm{B})$ y $\mathrm{Y}=\Spec(\mathrm{A})$, donde $\mathrm{A}$ es un anillo local reducido (resp.\ un anillo local íntegro e íntegramente cerrado) y $\mathrm{B}$ es una $\mathrm{A}$-álgebra de presentación finita. Podemos escribir entonces $(8.9.1)$
\[
  \mathrm{B}=\mathrm{B}_0\otimes_{\mathrm{A}_0}\mathrm{A},
\]
en donde $\mathrm{A}_0$ es una $\mathbf{Z}$-subálgebra de tipo finito de $\mathrm{A}$ y $\mathrm{B}_0$ es una $\mathrm{A}_0$-álgebra de tipo finito. Sea $(\mathrm{C}_\alpha)$ la familia creciente filtrante de las $\mathrm{A}_0$-subálgebras de tipo finito de $\mathrm{A}$, que son, por tanto, $\mathbf{Z}$-álgebras de tipo finito; se tiene $\mathrm{A}=\varinjlim\mathrm{C}_\alpha$. Distinguimos los dos casos:

1\textsuperscript{o} Supongamos que $\mathrm{A}$ es reducido y que $f$ es reducido en $x$. Si $\mathfrak{m}$ es el ideal maximal de $\mathrm{A}$, pongamos $\mathfrak{p}_\alpha=\mathfrak{m}\cap\mathrm{C}_\alpha$ y $\mathrm{A}'_\alpha=(\mathrm{C}_\alpha)_{\mathfrak{p}_\alpha}$, de modo que también se tiene $\mathrm{A}=\varinjlim\mathrm{A}'_\alpha$ $(5.13.3)$. Pongamos $\mathrm{Y}_\alpha=\Spec(\mathrm{A}'_\alpha)$ y $\mathrm{X}_\alpha=\Spec(\mathrm{B}_0\otimes_{\mathrm{A}_0}\mathrm{A}'_\alpha)$, y sean $f_\alpha:\mathrm{X}_\alpha\to\mathrm{Y}_\alpha$, $x_\alpha$ (resp.\ $y_\alpha$) la proyección de $x$ (resp.\ $y$) en $\mathrm{X}_\alpha$ (resp.\ $\mathrm{Y}_\alpha$). Como $f$ es reducido en $x$, existe $\alpha_0$ tal que $f_\alpha$ es reducido en $x_\alpha$ para $\alpha\geqslant\alpha_0$ $((6.7.8)$ y \sref{IV.11.2.6}); como $\mathrm{Y}_\alpha$ y $\mathrm{X}_\alpha$ son noetherianos y $\mathrm{A}'_\alpha$ es reducido (pues $\mathrm{C}_\alpha$ lo es), se deduce de $(3.3.5)$ que $\mathrm{X}_\alpha$ es reducido en $x_\alpha$. Pero
\[
  \mathrm{B}=\varinjlim(\mathrm{B}_0\otimes_{\mathrm{A}_0}\mathrm{A}'_\alpha)
\]
y, en consecuencia, $\sh{O}_{\mathrm{X},x}=\varinjlim\sh{O}_{\mathrm{X}_\alpha,x_\alpha}$ $(5.13.3)$; por tanto, $\sh{O}_{\mathrm{X},x}$ es reducido $(5.13.2)$.

2\textsuperscript{o} Supongamos que $\mathrm{A}$ es íntegramente cerrado y que $f$ es normal en $x$. Como $\mathrm{C}_\alpha$ es universalmente japonés $(7.7.4)$, su clausura integral $\mathrm{C}'_\alpha$ es una $\mathrm{C}_\alpha$-álgebra finita, evidentemente contenida en $\mathrm{A}$. Pongamos $\mathfrak{p}'_\alpha=\mathfrak{m}\cap\mathrm{C}'_\alpha$ y $\mathrm{A}''_\alpha=(\mathrm{C}'_\alpha)_{\mathfrak{p}'_\alpha}$, de modo que $\mathrm{A}''_\alpha$ es un anillo local noetheriano íntegro e íntegramente cerrado, y $\mathrm{A}=\varinjlim\mathrm{A}''_\alpha$ $(5.13.3)$. Pongamos $\mathrm{Y}'_\alpha=\Spec(\mathrm{A}''_\alpha)$ y $\mathrm{X}'_\alpha=\Spec(\mathrm{B}_0\otimes_{\mathrm{A}_0}\mathrm{A}''_\alpha)$, y sean $f'_\alpha:\mathrm{X}'_\alpha\to\mathrm{Y}'_\alpha$, $x'_\alpha$ (resp.\ $y'_\alpha$) la proyección de $x$ (resp.\ $y$) en $\mathrm{X}'_\alpha$ (resp.\ $\mathrm{Y}'_\alpha$). Como $f$ es normal en $x$, existe $\alpha_0$ tal que, para $\alpha\geqslant\alpha_0$, $f'_\alpha$ es normal en $x'_\alpha$ $((6.7.8)$ y \sref{IV.11.2.6}); \oldpage[IV-3]{141} como $\mathrm{X}'_\alpha$ y $\mathrm{Y}'_\alpha$ son noetherianos y $\mathrm{A}''_\alpha$ es íntegramente cerrado, se deduce de $(6.5.4)$ que $\mathrm{X}'_\alpha$ es normal en $x'_\alpha$. Pero los morfismos
\[
  \Spec(\mathrm{A}''_\beta)\longrightarrow\Spec(\mathrm{A}''_\alpha)
  \qquad(\alpha\leqslant\beta)
\]
son dominantes; por tanto, puesto que en virtud de \sref{IV.11.3.1} podemos suponer que los $f'_\alpha$ son planos para $\alpha\geqslant\alpha_0$, se deduce de $(2.3.7)$ que toda componente irreducible de $\mathrm{X}'_\beta$ domina una componente irreducible de $\mathrm{X}'_\alpha$. Concluimos entonces de $(5.13.4)$, aplicado a los esquemas
\[
  \Spec\!\left(\sh{O}_{\mathrm{X}_0,x_0}\otimes_{\mathrm{A}_0}\mathrm{A}''_\alpha\right),
\]
que $\mathrm{X}$ es normal en $x$.
\end{proof}

\begin{corollary}[11.3.14]
\label{IV.11.3.14}
Sea $f:\mathrm{X}\to\mathrm{Y}$ un morfismo localmente de presentación finita, sea $x$ un punto de $\mathrm{X}$ y sea $y=f(x)$. Si $\mathrm{Y}$ es geométricamente unirramificado en $y$ $(6.15.1)$ y si $f$ es normal en $x$ $(6.8.1)$, entonces $\mathrm{X}$ es geométricamente unirramificado en $x$.
\end{corollary}

\begin{proof}
Podemos reducirnos evidentemente al caso en que $\mathrm{Y}=\Spec(\mathrm{A})$, con $\mathrm{A}$ un dominio local íntegro geométricamente unirramificado y $y$ el punto cerrado de $\mathrm{Y}$. Sea $\mathrm{A}'$ la clausura integral de $\mathrm{A}$; pongamos $\mathrm{Y}'=\Spec(\mathrm{A}')$ y sea $y'$ el punto cerrado de $\mathrm{Y}'$. Así, el morfismo $g:\mathrm{Y}'\to\mathrm{Y}$ es radicial en $y$ $(6.15.3)$, integral y birracional. Pongamos $\mathrm{X}'=\mathrm{X}\times_\mathrm{Y}\mathrm{Y}'$, y sean $f':\mathrm{X}'\to\mathrm{Y}'$ y $g':\mathrm{X}'\to\mathrm{X}$ las proyecciones. El morfismo $g'$ es integral y radicial en $x$ $(6.15.3.1)$. Sea, pues, $x'$ el único punto de $g'^{-1}(x)$ situado sobre $y'$. El morfismo $f'$ es de presentación finita y normal en $x'$ $(6.7.8)$; en consecuencia, el anillo local $\sh{O}_{\mathrm{X}',x'}$ es un dominio íntegro e íntegramente cerrado \sref{IV.11.3.13}. Por otra parte, podemos suponer que $f$ es plano \sref{IV.11.3.1}, y entonces $g'$ es birracional $(6.15.4.1)$. Se deduce que $\Spec(\sh{O}_{\mathrm{X},x})$ es irreducible; como $\sh{O}_{\mathrm{X},x}$ es reducido \sref{IV.11.3.13}, es un dominio íntegro y es geométricamente unirramificado por $(6.15.5)$.
\end{proof}

\begin{proposition}[11.3.15]
\label{IV.11.3.15}
Sea $\mathrm{A}$ un anillo, sea $\mathrm{B}$ una $\mathrm{A}$-álgebra de presentación finita y sea $\mathrm{M}$ un $\mathrm{B}$-módulo de presentación finita que es plano como $\mathrm{A}$-módulo. Entonces existe una sucesión finita $(f_i)_{1\leqslant i\leqslant n}$ de elementos de $\mathrm{A}$ tal que el ideal generado por los $f_i$ es $\mathrm{A}$ y, para $0\leqslant i<n$,
\[
  \frac{\mathrm{M}_{f_{i+1}}}
       {\sum_{j=1}^{i}f_j\mathrm{M}_{f_{i+1}}}
\]
es un módulo libre sobre
\[
  \frac{\mathrm{A}_{f_{i+1}}}
       {\sum_{j=1}^{i}f_j\mathrm{A}_{f_{i+1}}}.
\]
\end{proposition}

Equivalentemente, los $\mathrm{D}(f_i)$ forman un recubrimiento abierto de $\mathrm{X}=\Spec(\mathrm{A})$ y, si se pone $\mathrm{Z}_1=\mathrm{D}(f_1)$ y luego, por inducción,
\[
  \mathrm{Z}_{i+1}
  =\mathrm{D}(f_{i+1})
   \cap\mathrm{V}(f_1\mathrm{A}+\cdots+f_i\mathrm{A}),
\]
los $\mathrm{Z}_i$ forman una partición de $\mathrm{X}$ en subconjuntos localmente cerrados, mientras que
\[
  \frac{\mathrm{A}_{f_{i+1}}}
       {\sum_{j=1}^{i}f_j\mathrm{A}_{f_{i+1}}}
\]
es el anillo de un subpreesquema afín de $\mathrm{X}$ cuyo espacio subyacente es $\mathrm{Z}_{i+1}$.

\begin{proof}
En virtud de \sref{IV.11.2.7}, podemos suponer primero que existen un subanillo noetheriano $\mathrm{A}_0$ de $\mathrm{A}$, una $\mathrm{A}_0$-álgebra de tipo finito $\mathrm{B}_0$ y un $\mathrm{B}_0$-módulo de tipo finito $\mathrm{M}_0$, plano sobre $\mathrm{A}_0$, tales que
\[
  \mathrm{B}=\mathrm{B}_0\otimes_{\mathrm{A}_0}\mathrm{A},
  \qquad
  \mathrm{M}=\mathrm{M}_0\otimes_{\mathrm{A}_0}\mathrm{A}.
\]
Es claro que basta demostrar la proposición para $\mathrm{A}_0$, $\mathrm{B}_0$ y $\mathrm{M}_0$: si los elementos $f_i\in\mathrm{A}_0$ satisfacen en ese caso las condiciones del enunciado, también las satisfacen para $\mathrm{A}$, $\mathrm{B}$ y $\mathrm{M}$, pues
\[
  \frac{\mathrm{M}_{f_{i+1}}}
       {\sum_{j=1}^{i}f_j\mathrm{M}_{f_{i+1}}}
  =
  \left(
    \frac{(\mathrm{M}_0)_{f_{i+1}}}
         {\sum_{j=1}^{i}f_j(\mathrm{M}_0)_{f_{i+1}}}
  \right)\otimes_{\mathrm{A}_0}\mathrm{A}
\]
(Bourbaki, \emph{Alg.\ comm.}, cap.~II, \textsection\,2, n\textsuperscript{o}\,7, prop.~18). Podemos, por tanto, limitarnos en lo sucesivo al caso en que $\mathrm{A}$ es noetheriano.

Observemos ahora que, si $\mathrm{C}$ es un anillo noetheriano, $\mathfrak{N}$ su nilradical y $\mathrm{P}$ un $\mathrm{C}$-módulo plano, se deduce de $(\mathbf{0}_{\mathrm{III}},10.1.2)$ que, para que $\mathrm{P}$ sea un $\mathrm{C}$-módulo libre, \oldpage[IV-3]{142} es necesario y suficiente que $\mathrm{P}\otimes_\mathrm{C}(\mathrm{C}/\mathfrak{N})$ sea un $(\mathrm{C}/\mathfrak{N})$-módulo libre. Obsérvese también que, si $\mathrm{C}$ es un anillo noetheriano reducido, existe $g\in\mathrm{C}$ tal que $\mathrm{C}_g$ es un dominio íntegro.

Aplicamos ahora el lema $(6.9.2)$. Puede definirse por inducción una sucesión $(f_i)_{i\geqslant1}$ de elementos de $\mathrm{A}$ como sigue:
\begin{enumerate}
\item[$1^\circ$] se elige $f_1$ de modo que
\[
  \mathrm{A}_1=(\mathrm{A}_{\mathrm{red}})_{f_1}
\]
sea un dominio íntegro y $\mathrm{M}\otimes_\mathrm{A}\mathrm{A}_1$ sea un $\mathrm{A}_1$-módulo libre;
\item[$2^\circ$] si el ideal $\mathfrak{J}_i$ generado por $f_1,\ldots,f_i$ es igual a $\mathrm{A}$, se pone $f_{i+1}=f_i$;
\item[$3^\circ$] si, por el contrario, $\mathfrak{J}_i\neq\mathrm{A}$, se elige $f_{i+1}\notin\mathfrak{J}_i$ de modo que
\[
  \mathrm{A}_{i+1}
  =\bigl((\mathrm{A}/\mathfrak{J}_i)_{\mathrm{red}}\bigr)_{f_{i+1}}
\]
sea un dominio íntegro y $\mathrm{M}\otimes_\mathrm{A}\mathrm{A}_{i+1}$ sea un $\mathrm{A}_{i+1}$-módulo libre.
\end{enumerate}
Como $\mathrm{A}$ es noetheriano, la sucesión creciente $(\mathfrak{J}_i)$ es estacionaria. Existe, por tanto, $n$ tal que $f_1,\ldots,f_n$ generan el ideal unidad de $\mathrm{A}$, y estos elementos satisfacen las condiciones requeridas porque
\[
  \bigl((\mathrm{A}/\mathfrak{J}_i)_{\mathrm{red}}\bigr)_{f_{i+1}}
  =
  \bigl((\mathrm{A}/\mathfrak{J}_i)_{f_{i+1}}\bigr)_{\mathrm{red}}.
\]
\end{proof}

\begin{proposition}[11.3.16]
\label{IV.11.3.16}
Sea $f:\mathrm{X}\to\mathrm{Y}$ un morfismo fielmente plano de presentación finita, y sea $g:\mathrm{Y}\to\mathrm{Z}$ un morfismo tal que $g\circ f:\mathrm{X}\to\mathrm{Z}$ es de tipo finito (resp.\ de presentación finita). Entonces $g$ es de tipo finito (resp.\ de presentación finita).
\end{proposition}

\begin{proof}
Como $f$ es sobreyectivo y $g\circ f$ es cuasicompacto, $g$ es cuasicompacto $(1.1.3)$. Demostremos primero que, si $g\circ f$ es de presentación finita, entonces $g$ es cuasiseparado. Sea $\mathrm{W}\subset\mathrm{Z}$ un abierto afín y sea $(\mathrm{V}_i)$ un recubrimiento abierto afín finito de $g^{-1}(\mathrm{W})$; basta demostrar que los $\mathrm{V}_i\cap\mathrm{V}_j$ son cuasicompactos $(1.2.6,\ 1.2.7)$. Para cada $i$, sea $(\mathrm{U}_{ih})$ un recubrimiento abierto afín finito de $f^{-1}(\mathrm{V}_i)$. Como $g\circ f$ es de presentación finita, todas las intersecciones $\mathrm{U}_{ih}\cap\mathrm{U}_{jk}$ son cuasicompactas. Puesto que $f$ es sobreyectivo, $\mathrm{V}_i\cap\mathrm{V}_j$ es la unión, cuando varían $h$ y $k$, de las imágenes $f(\mathrm{U}_{ih}\cap\mathrm{U}_{jk})$; es, por tanto, cuasicompacto, pues $f$ es continuo.

La cuestión es, pues, local en $\mathrm{Y}$, y podemos suponer que $\mathrm{Y}=\Spec(\mathrm{B})$ y $\mathrm{Z}=\Spec(\mathrm{A})$ son afines, mientras que $\mathrm{X}$ es la unión de un número finito de abiertos afines $\mathrm{X}_i=\Spec(\mathrm{C}_i)$, donde los $\mathrm{C}_i$ son $\mathrm{B}$-álgebras de tipo finito (resp.\ de presentación finita). Sea $\mathrm{X}'$ la unión disjunta de los $\mathrm{X}_i$, y sea $p:\mathrm{X}'\to\mathrm{X}$ el morfismo cuya restricción a cada $\mathrm{X}_i$ es la inmersión abierta canónica. El morfismo $p$ es fielmente plano y de presentación finita $(1.6.5)$; por consiguiente, $g\circ f\circ p$ es de tipo finito (resp.\ de presentación finita), y $f\circ p$ es fielmente plano y de presentación finita. Podemos suponer además, por tanto, que $\mathrm{X}=\Spec(\mathrm{C})$ es afín, y la proposición se reduce al corolario siguiente.
\end{proof}

\begin{corollary}[11.3.17]
\label{IV.11.3.17}
Sea $\mathrm{A}$ un anillo, sea $\mathrm{B}$ una $\mathrm{A}$-álgebra y sea $\mathrm{C}$ una $\mathrm{B}$-álgebra de presentación finita que es fielmente plana como $\mathrm{B}$-módulo. Si $\mathrm{C}$ es una $\mathrm{A}$-álgebra de tipo finito (resp.\ de presentación finita), entonces $\mathrm{B}$ es una $\mathrm{A}$-álgebra de tipo finito (resp.\ de presentación finita).
\end{corollary}

\begin{proof}
Supongamos primero que $\mathrm{C}$ es de tipo finito sobre $\mathrm{A}$. Sea $(\mathrm{B}_\lambda)$ la familia creciente filtrante de las $\mathrm{A}$-subálgebras de tipo finito de $\mathrm{B}$. Por $(8.8.2)$, existe un índice $\alpha$ tal que
\[
  \mathrm{C}=\mathrm{C}_\alpha\otimes_{\mathrm{B}_\alpha}\mathrm{B},
\]
en donde $\mathrm{C}_\alpha$ es una $\mathrm{B}_\alpha$-álgebra de presentación finita. Además, por $(8.10.5)$ y \sref{IV.11.2.6}, podemos suponer que $\mathrm{C}_\alpha$ es fielmente plana como $\mathrm{B}_\alpha$-módulo. Para $\lambda\geqslant\alpha$, pongamos
\[
  \mathrm{C}_\lambda
  =\mathrm{C}_\alpha\otimes_{\mathrm{B}_\alpha}\mathrm{B}_\lambda.
\]
Entonces $\mathrm{C}=\mathrm{C}_\lambda\otimes_{\mathrm{B}_\lambda}\mathrm{B}$ y $\mathrm{C}_\lambda$ es fielmente plana como $\mathrm{B}_\lambda$-módulo. Como $\mathrm{B}_\lambda\to\mathrm{B}$ es inyectivo, también lo es $\mathrm{C}_\lambda\to\mathrm{C}$. Además, $\mathrm{C}=\varinjlim\mathrm{C}_\lambda$, de modo que todo elemento de $\mathrm{C}$ pertenece a algún $\mathrm{C}_\lambda$; como $\mathrm{C}$ es de tipo finito sobre $\mathrm{A}$, existe $\lambda$ para el cual $\mathrm{C}_\lambda\to\mathrm{C}$ \oldpage[IV-3]{143} es biyectivo. La planitud fiel implica entonces que $\mathrm{B}_\lambda\to\mathrm{B}$ es biyectivo. Así, $\mathrm{B}=\mathrm{B}_\lambda$ es de tipo finito sobre $\mathrm{A}$.

Supongamos ahora que $\mathrm{C}$ es de presentación finita sobre $\mathrm{A}$. Por lo que acabamos de demostrar, $\mathrm{B}$ es de tipo finito sobre $\mathrm{A}$; escribamos
\[
  \mathrm{B}
  =\mathrm{A}[\mathrm{T}_1,\ldots,\mathrm{T}_n]/\widetilde{\mathfrak{J}}
\]
para un ideal $\widetilde{\mathfrak{J}}$. Sea $(\widetilde{\mathfrak{J}}_\lambda)$ la familia filtrante de los ideales de tipo finito de $\mathrm{A}[\mathrm{T}_1,\ldots,\mathrm{T}_n]$ contenidos en $\widetilde{\mathfrak{J}}$. Entonces
\[
  \widetilde{\mathfrak{J}}
  =\varinjlim\widetilde{\mathfrak{J}}_\lambda,
  \qquad
  \mathrm{B}
  =\varinjlim\mathrm{B}_\lambda,
  \qquad
  \mathrm{B}_\lambda
  =\mathrm{A}[\mathrm{T}_1,\ldots,\mathrm{T}_n]/
   \widetilde{\mathfrak{J}}_\lambda.
\]
Aplicando $(8.8.2)$, $(8.10.5)$ y \sref{IV.11.2.6} como antes, se tiene
\[
  \mathrm{C}
  =\mathrm{C}_\alpha\otimes_{\mathrm{B}_\alpha}\mathrm{B},
\]
en donde $\mathrm{C}_\alpha$ es una $\mathrm{B}_\alpha$-álgebra de presentación finita y un $\mathrm{B}_\alpha$-módulo fielmente plano. Para $\lambda\geqslant\alpha$, pongamos
\[
  \mathrm{C}_\lambda
  =\mathrm{C}_\alpha\otimes_{\mathrm{B}_\alpha}\mathrm{B}_\lambda.
\]
Por planitud,
\[
  \mathrm{C}_\lambda
  =
  \mathrm{C}_\alpha\Big/
  \Bigl(
    \mathrm{C}_\alpha\otimes_{\mathrm{B}_\alpha}
    (\widetilde{\mathfrak{J}}_\lambda/
     \widetilde{\mathfrak{J}}_\alpha)
  \Bigr),
\]
y, análogamente,
\[
  \mathrm{C}
  =
  \mathrm{C}_\alpha\Big/
  \Bigl(
    \mathrm{C}_\alpha\otimes_{\mathrm{B}_\alpha}
    (\widetilde{\mathfrak{J}}/
     \widetilde{\mathfrak{J}}_\alpha)
  \Bigr).
\]
Como $\mathrm{C}$ y $\mathrm{C}_\alpha$ son ambas $\mathrm{A}$-álgebras de presentación finita, el ideal
\[
  \mathrm{C}_\alpha\otimes_{\mathrm{B}_\alpha}
  (\widetilde{\mathfrak{J}}/
   \widetilde{\mathfrak{J}}_\alpha)
\]
es de tipo finito en $\mathrm{C}_\alpha$ $(1.4.4)$. Es la unión de los ideales de tipo finito
\[
  \mathrm{C}_\alpha\otimes_{\mathrm{B}_\alpha}
  (\widetilde{\mathfrak{J}}_\lambda/
   \widetilde{\mathfrak{J}}_\alpha).
\]
Por consiguiente, existe $\lambda\geqslant\alpha$ tal que
\[
  \mathrm{C}_\alpha\otimes_{\mathrm{B}_\alpha}
  (\widetilde{\mathfrak{J}}/
   \widetilde{\mathfrak{J}}_\alpha)
  =
  \mathrm{C}_\alpha\otimes_{\mathrm{B}_\alpha}
  (\widetilde{\mathfrak{J}}_\lambda/
   \widetilde{\mathfrak{J}}_\alpha).
\]
La planitud fiel da entonces
$\widetilde{\mathfrak{J}}=\widetilde{\mathfrak{J}}_\lambda$, lo que demuestra que $\mathrm{B}$ es de presentación finita sobre $\mathrm{A}$.
\end{proof}

\subsection{Descenso de la planitud por morfismos cualesquiera: caso de un preesquema base artiniano}
\label{subsection:IV.11.4}

\begin{theorem}[11.4.1]
\label{IV.11.4.1}
Sea $\mathrm{A}$ un anillo, sea $\mathfrak{J}$ un ideal nilpotente de $\mathrm{A}$ y sea $u_\lambda:\mathrm{A}\to\mathrm{B}_\lambda$ $(\lambda\in\mathrm{L})$
una familia de homomorfismos de anillos tal que la intersección de los núcleos de los $u_\lambda$ es $0$.
Sea $\mathrm{M}$ un $\mathrm{A}$-módulo tal que, para todo $\lambda\in\mathrm{L}$, $\mathrm{M}\otimes_\mathrm{A}\mathrm{B}_\lambda$ es un $\mathrm{B}_\lambda$-módulo libre y que $\mathrm{M}\otimes_\mathrm{A}(\mathrm{A}/\mathfrak{J})$
es un $(\mathrm{A}/\mathfrak{J})$-módulo libre. Entonces $\mathrm{M}$ es un $\mathrm{A}$-módulo libre. Si, además, el conjunto de índices $\mathrm{L}$ es
finito, en el enunciado anterior puede sustituirse en todas partes ``libre'' por ``plano''.
\end{theorem}

\begin{proof}
En ambos casos bastará demostrar que $\mathrm{M}$ es un $\mathrm{A}$-módulo \emph{plano}: en efecto,
cuando $\mathrm{M}\otimes_\mathrm{A}(\mathrm{A}/\mathfrak{J})$ es un $(\mathrm{A}/\mathfrak{J})$-módulo libre, se deduce que $\mathrm{M}$ es un $\mathrm{A}$-módulo
libre en virtud de $(\mathbf{0}_{\mathrm{III}},~10.1.2)$.

Utilizaremos el lema siguiente, que generaliza $(\mathbf{0}_{\mathrm{III}},~10.2.1)$:

\begin{lemma}[11.4.1.1]
\label{IV.11.4.1.1}
Sea $\mathrm{A}$ un anillo provisto de una filtración finita $(\mathfrak{J}_n)_{0\leq n\leq\mathrm{N}+1}$ con
$\mathfrak{J}_0=\mathrm{A}$, $\mathfrak{J}_{\mathrm{N}+1}=0$. Sea $\mathrm{M}$ un $\mathrm{A}$-módulo provisto de la filtración $(\mathfrak{J}_n\mathrm{M})_{0\leq n\leq\mathrm{N}+1}$, y denotemos
por $\operatorname{gr}(\mathrm{A})$ y $\operatorname{gr}(\mathrm{M})$, respectivamente, el anillo y el módulo graduados. Supongamos que $\mathrm{M}\otimes_\mathrm{A}(\mathrm{A}/\mathfrak{J}_1)$
es un $(\mathrm{A}/\mathfrak{J}_1)$-módulo plano y que el homomorfismo canónico
\[
\label{IV.11.4.1.2}
\operatorname{gr}_0(\mathrm{M})\otimes_{\operatorname{gr}_0(\mathrm{A})}\operatorname{gr}(\mathrm{A}) \to \operatorname{gr}(\mathrm{M})
\tag{11.4.1.2}
\]
es inyectivo. Entonces $\mathrm{M}$ es un $\mathrm{A}$-módulo plano.
\end{lemma}

El homomorfismo canónico {\hyperlink{ega4.c04.lem.11-4-1-2}{(11.4.1.2)}} se define del mismo modo que
el de $(\mathbf{0}_{\mathrm{III}},~10.1.1)$, siendo en grado $n$ el homomorfismo
\[
(\mathrm{M}/\mathfrak{J}_1\mathrm{M})\otimes(\mathfrak{J}_n/\mathfrak{J}_{n+1})=(\mathrm{M}\otimes\mathfrak{J}_n)/(\operatorname{Im}(\mathrm{M}\otimes\mathfrak{J}_{n+1})+\operatorname{Im}(\mathfrak{J}_1\mathrm{M}\otimes\mathfrak{J}_n))\to\mathfrak{J}_n\mathrm{M}/\mathfrak{J}_{n+1}\mathrm{M}
\]
inducido por el homomorfismo canónico $\mathrm{M}\otimes\mathfrak{J}_n\to\mathfrak{J}_n\mathrm{M}$ al pasar a los cocientes.
El lema se demuestra por inducción sobre $\mathrm{N}$, pues no hay nada que demostrar para $\mathrm{N}=0$.
\oldpage[IV-3]{144}
Las condiciones impuestas a $\mathrm{M}$ implican, por la hipótesis de inducción, que $\mathrm{M}\otimes_\mathrm{A}(\mathrm{A}/\mathfrak{J}_\mathrm{N})$
es un $(\mathrm{A}/\mathfrak{J}_\mathrm{N})$-módulo plano. Observemos ahora que $\mathfrak{J}_\mathrm{N}^2\subset\mathfrak{J}_{\mathrm{N}+1}=0$ si $\mathrm{N}\geq 1$, y
$(\mathrm{M}/\mathfrak{J}_1\mathrm{M})\otimes(\mathfrak{J}_\mathrm{N}/\mathfrak{J}_{\mathrm{N}+1})=(\mathrm{M}/\mathfrak{J}_\mathrm{N}\mathrm{M})\otimes(\mathfrak{J}_\mathrm{N}/\mathfrak{J}_{\mathrm{N}+1})=(\mathrm{M}/\mathfrak{J}_\mathrm{N}\mathrm{M})\otimes\mathfrak{J}_\mathrm{N}$; de ahí el homomorfismo canónico
\[
(\mathrm{M}/\mathfrak{J}_\mathrm{N}\mathrm{M})\otimes\mathfrak{J}_\mathrm{N} \to \mathfrak{J}_\mathrm{N}\mathrm{M}=\mathfrak{J}_\mathrm{N}\mathrm{M}/\mathfrak{J}_\mathrm{N}^2\mathrm{M}
\]
es inyectivo. Aplicando $(\mathbf{0}_{\mathrm{III}},~10.2.1)$ a la filtración $\mathfrak{J}_\mathrm{N}$-preádica, concluimos que
$\mathrm{M}$ es un $\mathrm{A}$-módulo plano.

Para aplicar este lema a \sref{IV.11.4.1}, denotemos por $\mathfrak{J}_n$ el ideal de $\mathrm{A}$ que es la intersección
de las imágenes inversas $u_\lambda^{-1}(\mathfrak{J}^n\mathrm{B}_\lambda)$: es inmediato que $\mathfrak{J}_0=\mathrm{A}$, $\mathfrak{J}_m\mathfrak{J}_n\subset\mathfrak{J}_{m+n}$
para $m\geq 0$, $n\geq 0$; además, si $\mathfrak{J}^{\mathrm{N}+1}=0$, se tiene también $\mathfrak{J}_{\mathrm{N}+1}=0$, ya que la intersección de los
núcleos de los $u_\lambda$ es $0$.

Provistamos $\mathrm{A}$ de la filtración $(\mathfrak{J}_n)$ y $\mathrm{M}$ de la filtración $(\mathfrak{J}_n\mathrm{M})$ y, para cada $\lambda$,
provistamos $\mathrm{B}_\lambda$ y $\mathrm{N}_\lambda=\mathrm{M}\otimes_\mathrm{A}\mathrm{B}_\lambda$ de las filtraciones $\mathfrak{J}$-preádicas; consideremos para cada $\lambda$ el
diagrama conmutativo
\[
\xymatrix{
\operatorname{gr}^0(\mathrm{M})\otimes_{\operatorname{gr}^0(\mathrm{A})}\operatorname{gr}(\mathrm{A}) \ar[r]^-{\varphi} \ar[d]_{f_\lambda} & \operatorname{gr}(\mathrm{M}) \ar[d]^{g_\lambda} \\
\operatorname{gr}^0_\mathfrak{J}(\mathrm{N}_\lambda)\otimes_{\operatorname{gr}^0_\mathfrak{J}(\mathrm{B}_\lambda)}\operatorname{gr}_\mathfrak{J}(\mathrm{B}_\lambda) \ar[r]_-{\varphi_\lambda} & \operatorname{gr}_\mathfrak{J}(\mathrm{N}_\lambda)
}
\]
en donde las flechas horizontales son los homomorfismos canónicos {\hyperlink{ega4.c04.lem.11-4-1-2}{(11.4.1.2)}} y las flechas
verticales se deducen de los homomorfismos canónicos $\mathrm{A}\to\mathrm{B}_\lambda$ y $\mathrm{M}\to\mathrm{N}_\lambda$. La hipótesis
de que $\mathrm{N}_\lambda$ es un $\mathrm{B}_\lambda$-módulo plano implica que $\varphi_\lambda$ es inyectivo $(\mathbf{0}_{\mathrm{III}},~10.2.1)$ y, por tanto,
$\operatorname{Ker}(\varphi)\subset\operatorname{Ker}(f_\lambda)$. Pongamos $\mathrm{A}_0=\mathrm{A}/\mathfrak{J}_1$, $\mathrm{M}_0=\mathrm{M}\otimes_\mathrm{A}\mathrm{A}_0$ y observemos que
\[
\operatorname{gr}^0_\mathfrak{J}(\mathrm{N}_\lambda)=\mathrm{N}_\lambda/\mathfrak{J}\mathrm{N}_\lambda=(\mathrm{M}/\mathfrak{J}\mathrm{M})\otimes_\mathrm{A}\mathrm{B}_\lambda=(\mathrm{M}/\mathfrak{J}_1\mathrm{M})\otimes_{\mathrm{A}}(\mathrm{B}_\lambda/\mathfrak{J}\mathrm{B}_\lambda)
\]
y, por tanto,
\[
\operatorname{gr}^0_\mathfrak{J}(\mathrm{N}_\lambda)\otimes_{\operatorname{gr}^0_\mathfrak{J}(\mathrm{B}_\lambda)}\operatorname{gr}_\mathfrak{J}(\mathrm{B}_\lambda)=\mathrm{M}_0\otimes_{\mathrm{A}_0}\operatorname{gr}_\mathfrak{J}(\mathrm{B}_\lambda).
\]
\oldpage[IV-3]{145}
El homomorfismo $f_\lambda$ puede escribirse, por tanto,
\[
1\otimes\operatorname{gr}(u_\lambda) : \mathrm{M}_0\otimes_{\mathrm{A}_0}\operatorname{gr}(\mathrm{A}) \to \mathrm{M}_0\otimes_{\mathrm{A}_0}\operatorname{gr}_\mathfrak{J}(\mathrm{B}_\lambda),
\]
y, como por hipótesis $\mathrm{M}_0$ es un $\mathrm{A}_0$-módulo \emph{plano}, el núcleo de $f_\lambda$ es igual a $\mathrm{M}_0\otimes_{\mathrm{A}_0}\mathrm{R}_\lambda$,
en donde $\mathrm{R}_\lambda$ es el núcleo de $\operatorname{gr}(u_\lambda):\operatorname{gr}(\mathrm{A})\to\operatorname{gr}_\mathfrak{J}(\mathrm{B}_\lambda)$. Basta, por tanto, demostrar que
$\bigcap_{\lambda\in\mathrm{L}}(\mathrm{M}_0\otimes_{\mathrm{A}_0}\mathrm{R}_\lambda)=0$. Ahora bien, por definición de los $\mathfrak{J}_n$, la intersección de los núcleos de los homomorfismos
$\mathfrak{J}_n\mathrm{A}/\mathfrak{J}_{n+1}\mathrm{A}\to\mathfrak{J}^n\mathrm{B}_\lambda/\mathfrak{J}^{n+1}\mathrm{B}_\lambda$, cuando $\lambda$ recorre $\mathrm{L}$, es $0$; en otras
palabras, $\bigcap_{\lambda\in\mathrm{L}}\mathrm{R}_\lambda=0$. Supongamos primero que $\mathrm{M}_0$ es un $\mathrm{A}_0$-módulo libre; tomando
una base de $\mathrm{M}_0$, vemos inmediatamente que $\mathrm{M}_0\otimes_{\mathrm{A}_0}(\bigcap_{\lambda\in\mathrm{L}}\mathrm{R}_\lambda)=\bigcap_{\lambda\in\mathrm{L}}(\mathrm{M}_0\otimes_{\mathrm{A}_0}\mathrm{R}_\lambda)$, de donde se obtiene la
proposición en este caso. Cuando $\mathrm{L}$ es finito, la fórmula anterior sigue siendo cierta bajo
la sola hipótesis de que $\mathrm{M}_0$ sea un $\mathrm{A}_0$-módulo plano $(\mathbf{0}_\mathrm{I},~6.1.3)$, lo que termina la
demostración.
\end{proof}

\begin{remark}[11.4.2]
\label{IV.11.4.2}
La conclusión de \sref{IV.11.4.1} puede ser falsa si $\mathrm{L}$ es infinito
y solo se supone que $\mathrm{M}\otimes_\mathrm{A}(\mathrm{A}/\mathfrak{J})$ es un $(\mathrm{A}/\mathfrak{J})$-módulo plano. Por ejemplo,
sea $\mathrm{V}$ un anillo de valoración discreta, sea $\mathrm{K}$ su cuerpo de fracciones y sea $\mathrm{A}=\mathrm{V}[\mathrm{T}]/(\mathrm{T}^2)$
($\mathrm{T}$ una indeterminada); tomemos como $\mathfrak{J}$ la imagen de $(\mathrm{T})$ en $\mathrm{A}$, de modo que $\mathrm{A}/\mathfrak{J}=\mathrm{V}$, y
tomemos $\mathrm{M}=\mathrm{K}$, que es un $(\mathrm{A}/\mathfrak{J})$-módulo y, por tanto, es igual a $\mathrm{M}\otimes_\mathrm{A}(\mathrm{A}/\mathfrak{J})$; además, $\mathrm{M}$ es
un $(\mathrm{A}/\mathfrak{J})$-módulo plano, pero no un $\mathrm{A}$-módulo plano: la inclusión $\mathfrak{J}\hookrightarrow\mathrm{A}$ induce, al tensorizar con $\mathrm{K}$,
la aplicación nula $\mathfrak{J}\otimes_{\mathrm{A}}\mathrm{K}\simeq\mathrm{K}\longrightarrow\mathrm{K}$,
que no es inyectiva, pues $\mathrm{K}\neq0$. Por otra parte, consideremos el ideal maximal $\mathfrak{m}$ del anillo local noetheriano
$\mathrm{A}$; se tiene $\mathfrak{m}\mathrm{K}=\mathrm{K}$ y, por tanto, $\mathrm{M}\otimes_\mathrm{A}(\mathrm{A}/\mathfrak{m}^n)=0$ para todo entero $n$; los $(\mathrm{A}/\mathfrak{m}^n)$-módulos
$\mathrm{M}\otimes_\mathrm{A}(\mathrm{A}/\mathfrak{m}^n)$ son, pues, planos para todo $n$, y la intersección de los $\mathfrak{m}^n$ es $0$.
\end{remark}

\begin{corollary}[11.4.3]
\label{IV.11.4.3}
Sea $\mathrm{A}$ un anillo semilocal cuyo radical $\mathfrak{J}$ es nilpotente (por
ejemplo, un anillo artiniano), y sea $u_\lambda:\mathrm{A}\to\mathrm{B}_\lambda$ $(\lambda\in\mathrm{L})$ una familia de homomorfismos de anillos tal
que la intersección de los núcleos de los $u_\lambda$ es $0$. Para que un $\mathrm{A}$-módulo $\mathrm{M}$ sea plano, es necesario y
suficiente que, para todo $\lambda\in\mathrm{L}$, $\mathrm{M}\otimes_\mathrm{A}\mathrm{B}_\lambda$ sea un $\mathrm{B}_\lambda$-módulo plano.
\end{corollary}

\begin{proof}
Como $\mathrm{A}/\mathfrak{J}$ es un producto directo de un número finito de cuerpos (Bourbaki, \emph{Alg. comm.},
cap.~II, \S~3, n\textsuperscript{o}~5, prop.~16) y $\mathfrak{J}$ es nilpotente, $\mathrm{A}$ es un producto directo de un
número finito de anillos \emph{locales} $\mathrm{A}_i$ cuyo radical es nilpotente (\emph{loc.\ cit.}, \S~4, n\textsuperscript{o}~3, cor.~1
de la prop.~15), y $\mathrm{M}$ es, por tanto, una suma directa de $\mathrm{A}_i$-módulos $\mathrm{M}_i$, y cada $\mathrm{M}_i$ es
anulado por los $\mathrm{A}_j$ de índice $j\neq i$; para que $\mathrm{M}$ sea un $\mathrm{A}$-módulo plano, es necesario y
suficiente que cada $\mathrm{M}_i$ sea un $\mathrm{A}_i$-módulo plano; además, la intersección de los núcleos
de los homomorfismos $\mathrm{A}_i\longrightarrow\mathrm{A}\xrightarrow{u_\lambda}\mathrm{B}_\lambda$ es $0$, y $\mathrm{M}_i\otimes_{\mathrm{A}_i}\mathrm{B}_\lambda$ es un sumando directo
de $\mathrm{M}\otimes_\mathrm{A}\mathrm{B}_\lambda$. Podemos, por tanto, limitarnos al caso en que $\mathrm{A}$ es también \emph{local}. Entonces $\mathrm{A}/\mathfrak{J}$ es un
cuerpo, de modo que $\mathrm{M}\otimes_\mathrm{A}(\mathrm{A}/\mathfrak{J})$ es un $(\mathrm{A}/\mathfrak{J})$-módulo libre, y basta ver que, para
todo $\lambda$, $\mathrm{M}\otimes_\mathrm{A}\mathrm{B}_\lambda$ es un $\mathrm{B}_\lambda$-módulo \emph{libre}, en virtud de \sref{IV.11.4.1}. Ahora bien, si se pone $\mathfrak{J}_\lambda=\mathfrak{J}\mathrm{B}_\lambda$,
$(\mathrm{M}\otimes_\mathrm{A}\mathrm{B}_\lambda)\otimes_{\mathrm{B}_\lambda}(\mathrm{B}_\lambda/\mathfrak{J}_\lambda)=(\mathrm{M}\otimes_\mathrm{A}(\mathrm{A}/\mathfrak{J}))\otimes_{\mathrm{A}/\mathfrak{J}}(\mathrm{B}_\lambda/\mathfrak{J}_\lambda)$ es un $(\mathrm{B}_\lambda/\mathfrak{J}_\lambda)$-módulo libre, y $\mathfrak{J}_\lambda$ es
nilpotente. La conclusión se deduce, pues, de la hipótesis de que $\mathrm{M}\otimes_\mathrm{A}\mathrm{B}_\lambda$ es un $\mathrm{B}_\lambda$-módulo
plano y de $(\mathbf{0}_{\mathrm{III}},~10.1.2)$.
\end{proof}

\begin{corollary}[11.4.4]
\label{IV.11.4.4}
Sea $\mathrm{A}$ un anillo y sea $\mathrm{M}$ un $\mathrm{A}$-módulo; supongamos que existe
un ideal nilpotente $\mathfrak{J}$ de $\mathrm{A}$ tal que $\mathrm{M}\otimes_\mathrm{A}(\mathrm{A}/\mathfrak{J})$ es un $(\mathrm{A}/\mathfrak{J})$-módulo libre. Entonces el conjunto de los
ideales $\mathfrak{K}$ de $\mathrm{A}$ tales que $\mathrm{M}\otimes_\mathrm{A}(\mathrm{A}/\mathfrak{K})$ es un $(\mathrm{A}/\mathfrak{K})$-módulo libre admite un elemento mínimo $\mathfrak{K}_0$
\oldpage[IV-3]{146}
(que es también el mínimo de los ideales $\mathfrak{K}$ tales que $\mathrm{M}\otimes_\mathrm{A}(\mathrm{A}/\mathfrak{K})$ es un $(\mathrm{A}/\mathfrak{K})$-módulo \emph{plano}).
\end{corollary}

\begin{proof}
Para todo homomorfismo $u:\mathrm{A}\to\mathrm{A}'$ tal que $\mathrm{M}\otimes_\mathrm{A}\mathrm{A}'$ es un $\mathrm{A}'$-módulo libre (resp.\ un
$\mathrm{A}'$-módulo plano), es necesario y suficiente que $u$ se factorice a través de $\mathrm{A}\to\mathrm{A}/\mathfrak{K}_0\to\mathrm{A}'$ (o, equivalentemente, que $\mathfrak{K}_0\mathrm{A}'=0$).

El hecho de que la intersección $\mathfrak{K}_0$ de los ideales $\mathfrak{K}$ para los que $\mathrm{M}\otimes_\mathrm{A}(\mathrm{A}/\mathfrak{K})$ es un
$(\mathrm{A}/\mathfrak{K})$-módulo libre sea ella misma el mínimo de estos ideales se deduce de \sref{IV.11.4.1}, aplicado al
anillo $\mathrm{A}/\mathfrak{K}_0$, al ideal nilpotente $\mathfrak{J}/\mathfrak{K}_0$ y a los homomorfismos $\mathrm{A}/\mathfrak{K}_0\to\mathrm{A}/\mathfrak{K}$,
cuyos núcleos tienen intersección $0$. Si $\mathrm{A}'$ es una $(\mathrm{A}/\mathfrak{K}_0)$-álgebra, entonces
$\mathrm{M}\otimes_\mathrm{A}\mathrm{A}'=(\mathrm{M}\otimes_\mathrm{A}(\mathrm{A}/\mathfrak{K}_0))\otimes_{\mathrm{A}/\mathfrak{K}_0}\mathrm{A}'$ y, por tanto, $\mathrm{M}\otimes_\mathrm{A}\mathrm{A}'$ es un $\mathrm{A}'$-módulo libre. Recíprocamente, si $\mathrm{A}'$ es una $\mathrm{A}$-álgebra tal que $\mathrm{M}\otimes_\mathrm{A}\mathrm{A}'$ es un $\mathrm{A}'$-módulo libre y si $\mathfrak{K}$ es el núcleo
del homomorfismo $\mathrm{A}\to\mathrm{A}'$, se deduce de \sref{IV.11.4.1}, aplicado al anillo $\mathrm{A}/\mathfrak{K}$, al
$(\mathrm{A}/\mathfrak{K})$-módulo $\mathrm{M}\otimes_\mathrm{A}(\mathrm{A}/\mathfrak{K})$, al ideal nilpotente $(\mathfrak{J}+\mathfrak{K})/\mathfrak{K}$ y al homomorfismo \emph{inyectivo}
$\mathrm{A}/\mathfrak{K}\to\mathrm{A}'$, que $\mathrm{M}\otimes_\mathrm{A}(\mathrm{A}/\mathfrak{K})$ es un $(\mathrm{A}/\mathfrak{K})$-módulo libre, y por tanto $\mathfrak{K}\supset\mathfrak{K}_0$. El hecho de que
pueda sustituirse ``libre'' por ``plano'' en lo anterior (conservando, naturalmente,
la hipótesis de que $\mathrm{M}\otimes_\mathrm{A}(\mathrm{A}/\mathfrak{J})$ es un $(\mathrm{A}/\mathfrak{J})$-módulo \emph{libre}) se deduce de $(\mathbf{0}_{\mathrm{III}},~10.1.2)$,
como vimos al comienzo de la demostración de \sref{IV.11.4.1}.
\end{proof}

\begin{proposition}[11.4.5]
\label{IV.11.4.5}
Sea $\mathrm{Y}$ un preesquema irreducible, sea $f:\mathrm{X}\to\mathrm{Y}$ un morfismo de
presentación finita y sea $\sh{F}$ un $\sh{O}_\mathrm{X}$-módulo de presentación finita. Existe un abierto no vacío $\mathrm{U}$
de $\mathrm{Y}$ y un subpreesquema cerrado $\mathrm{Z}$ de $\mathrm{U}$, de presentación finita sobre $\mathrm{U}$, con la propiedad siguiente:
para todo cambio de base $\mathrm{U}'\to\mathrm{U}$, si se pone $\mathrm{X}'=\mathrm{X}\times_\mathrm{Y}\mathrm{U}'$, $f'=f_{(\mathrm{U}')}$ y $\sh{F}'=\sh{F}\otimes_{\sh{O}_\mathrm{Y}}\sh{O}_{\mathrm{U}'}$,
entonces, para que $\sh{F}'$ sea $f'$-plano, es necesario y suficiente que el morfismo $\mathrm{U}'\to\mathrm{U}$ se factorice
a través de $\mathrm{U}'\to\mathrm{Z}\to\mathrm{U}$. Tal esquema $\mathrm{Z}$ tiene el mismo espacio subyacente que $\mathrm{U}$. Supongamos además que $\mathrm{Y}$ es afín
y sea $(\mathrm{W}_i)$ un recubrimiento abierto afín finito de $\mathrm{X}$; entonces puede suponerse elegido $\mathrm{U}$
de modo que, si $\mathrm{U}'=\operatorname{Spec}(\mathrm{A}')$ es un esquema afín y $\mathrm{U}'\to\mathrm{U}$ un morfismo que se factoriza
a través de $\mathrm{U}'\to\mathrm{Z}\to\mathrm{U}$, cada uno de los $\Gamma(\mathrm{W}_i\times_\mathrm{Y}\mathrm{U}',\sh{F}')$ sea un $\mathrm{A}'$-módulo libre.
\end{proposition}

\begin{proof}
Podemos limitarnos evidentemente al caso en que $\mathrm{Y}=\operatorname{Spec}(\mathrm{A})$ es afín. Utilizando
\sref{IV.8.9.1}, existen un subanillo noetheriano $\mathrm{A}_1$ de $\mathrm{A}$, un morfismo de tipo finito
$f_1:\mathrm{X}_1\to\mathrm{Y}_1=\operatorname{Spec}(\mathrm{A}_1)$ y un $\sh{O}_{\mathrm{X}_1}$-módulo coherente $\sh{F}_1$ tales que $f=f_{(\mathrm{Y})}$ y $\sh{F}=\sh{F}_1\otimes_{\sh{O}_{\mathrm{Y}_1}}\sh{O}_\mathrm{Y}$;
podemos suponer además que los $\mathrm{W}_i$ son imágenes inversas de abiertos afines de $\mathrm{X}_1$.
Observemos asimismo que $\mathrm{Y}_1$ es irreducible, pues el morfismo $\mathrm{Y}\to\mathrm{Y}_1$ es dominante (\textbf{I},~$1.2.7$).
Supongamos demostrada la proposición para $\mathrm{Y}_1$, $\mathrm{X}_1$ y $\sh{F}_1$, y sean $\mathrm{U}_1$ y $\mathrm{Z}_1$ el abierto de $\mathrm{Y}_1$
y el subpreesquema cerrado de $\mathrm{U}_1$ con las propiedades requeridas. Pongamos $\mathrm{U}=\mathrm{U}_1\times_{\mathrm{Y}_1}\mathrm{Y}$ y $\mathrm{Z}=\mathrm{Z}_1\times_{\mathrm{Y}_1}\mathrm{Y}$, sus imágenes inversas. Entonces, si $\mathrm{U}'\to\mathrm{U}$ es un cambio de base
tal que $\sh{F}'=\sh{F}\otimes_{\sh{O}_\mathrm{Y}}\sh{O}_{\mathrm{U}'}$ es $f'$-plano, el morfismo $\mathrm{U}'\to\mathrm{U}_1$ se factoriza
a través de $\mathrm{U}'\to\mathrm{Z}_1\to\mathrm{U}_1$; pero, como $\mathrm{U}'\to\mathrm{U}\to\mathrm{U}_1$ se factoriza también a través de $\mathrm{U}'\to\mathrm{U}\to\mathrm{U}_1$, la definición
del producto de preesquemas muestra que $\mathrm{U}'\to\mathrm{U}$ se factoriza a través de $\mathrm{U}'\to\mathrm{Z}\to\mathrm{U}$.

Podemos, por tanto, limitarnos al caso en que $\mathrm{A}$ es \emph{noetheriano}; sea $\mathfrak{N}$ su nilradical, que es
nilpotente, y pongamos $\mathrm{A}_0=\mathrm{A}_{\mathrm{red}}=\mathrm{A}/\mathfrak{N}$, $\mathrm{Y}_0=\operatorname{Spec}(\mathrm{A}_0)=\mathrm{Y}_{\mathrm{red}}$, $\mathrm{X}_0=\mathrm{X}\times_\mathrm{Y}\mathrm{Y}_0$, $\sh{F}_0=\sh{F}\otimes_{\sh{O}_\mathrm{Y}}\sh{O}_{\mathrm{Y}_0}$,
$\mathrm{W}_{i0}=\mathrm{W}_i\times_\mathrm{Y}\mathrm{Y}_0$; si $\mathrm{M}_i=\Gamma(\mathrm{W}_i,\sh{F})$, entonces $\mathrm{M}_{i0}=\Gamma(\mathrm{W}_{i0},\sh{F}_0)$ es igual a $\mathrm{M}_i\otimes_\mathrm{A}\mathrm{A}_0$.
Como $\mathrm{A}_0$ es \emph{íntegro}, podemos, en virtud del teorema de planitud genérica $(6.9.2)$,
sustituyendo si es preciso $\mathrm{Y}$ por un abierto no vacío de $\mathrm{Y}$, suponer que los $\mathrm{M}_{i0}$ son
$\mathrm{A}_0$-módulos libres. En virtud de \sref{IV.11.4.4}, existe, por tanto, para cada $i$ un ideal $\mathfrak{J}_i\subset\mathfrak{N}$ tal
\oldpage[IV-3]{147}
que las $\mathrm{A}$-álgebras $\mathrm{A}'$ para las que los $\mathrm{M}_i\otimes_\mathrm{A}\mathrm{A}'$ son $\mathrm{A}'$-módulos libres (o planos)
son exactamente aquellas para las que $\mathfrak{J}_i\mathrm{A}'=0$. Sea $\mathfrak{J}=\sum_i\mathfrak{J}_i$, que es un ideal contenido
en $\mathfrak{N}$; decir que $\sh{F}\otimes_{\mathrm{A}}\mathrm{A}'$ es $\mathrm{A}'$-plano equivale a decir que todos los $\mathrm{M}_i\otimes_\mathrm{A}\mathrm{A}'$ son
$\mathrm{A}'$-módulos planos, es decir, que $\mathfrak{J}\mathrm{A}'=0$. Se deduce que, tomando $\mathrm{Z}=\mathrm{V}(\mathfrak{J})$, se
resuelve la cuestión, pues, para que un morfismo $\mathrm{U}'\to\mathrm{U}$ tenga la propiedad del enunciado,
es necesario y suficiente que, para todo abierto afín $\mathrm{V}'\subset\mathrm{U}'$, el morfismo $\mathrm{V}'\to\mathrm{U}$
tenga la misma propiedad.
\end{proof}

\begin{corollary}[11.4.6]
\label{IV.11.4.6}
Sea $\mathrm{A}$ un anillo tal que $\operatorname{Spec}(\mathrm{A})$ es irreducible, sea $\mathfrak{p}$ su único
ideal primo minimal, sea $\mathrm{B}$ una $\mathrm{A}$-álgebra de presentación finita y sea $\mathrm{M}$ un $\mathrm{B}$-módulo de presentación finita.
Supongamos que $\mathrm{M}_\mathfrak{p}$ es un $\mathrm{A}_\mathfrak{p}$-módulo plano; entonces existe $t\in\mathrm{A}-\mathfrak{p}$ tal que $\mathrm{M}_t$ es un $\mathrm{A}_t$-módulo libre.
\end{corollary}

\begin{proof}
Aplicando \sref{IV.11.4.5} con $\mathrm{Y}=\operatorname{Spec}(\mathrm{A})$, $\mathrm{X}=\operatorname{Spec}(\mathrm{B})$, $\sh{F}=\widetilde{\mathrm{M}}$, podemos (sustituyendo
si es preciso $\mathrm{A}$ por $\mathrm{A}_h$, donde $h\in\mathrm{A}-\mathfrak{p}$) suponer que existe un ideal de tipo finito $\mathfrak{J}$ de $\mathrm{A}$
tal que las $\mathrm{A}$-álgebras $\mathrm{A}'$ para las que $\mathrm{M}\otimes_\mathrm{A}\mathrm{A}'$ es un $\mathrm{A}'$-módulo libre (o plano) son
exactamente aquellas tales que $\mathfrak{J}\mathrm{A}'=0$. En particular, la hipótesis implica que $\mathrm{M}_\mathfrak{p}$ es un
$\mathrm{A}_\mathfrak{p}$-módulo plano, y por tanto $\mathfrak{J}\mathrm{A}_\mathfrak{p}=0$, o $\mathfrak{J}_\mathfrak{p}=0$. Ahora bien, como $\mathfrak{J}$ es de tipo finito,
existe $t\in\mathrm{A}-\mathfrak{p}$ tal que $\mathfrak{J}_t=0$, o equivalentemente $\mathfrak{J}\mathrm{A}_t=0$, y por tanto $\mathrm{M}_t$ es un $\mathrm{A}_t$-módulo libre.
\end{proof}

\begin{proposition}[11.4.7]
\label{IV.11.4.7}
Sea $\mathrm{A}$ un anillo noetheriano, sea $\mathfrak{J}$ un ideal de $\mathrm{A}$ y sea $\mathrm{M}$ un $\mathrm{A}$-módulo
idealmente separado $(\mathbf{0}_{\mathrm{III}},~10.2.1)$ para la topología $\mathfrak{J}$-preádica. Sea $(\mathfrak{p}_i)_{1\leq i\leq r}$ una familia finita
de ideales primos de $\mathrm{A}$ que contienen $\mathfrak{J}$; para todo entero $n\geq 0$, sea $\mathfrak{p}_i^{(n)}$ la $n$-ésima
potencia simbólica (es decir, el núcleo del homomorfismo $\mathrm{A}\to\mathrm{A}_{\mathfrak{p}_i}/\mathfrak{p}_i^n\mathrm{A}_{\mathfrak{p}_i}$); pongamos $\mathfrak{J}_n=\bigcap_{i=1}^{r}\mathfrak{p}_i^{(n)}$, de modo que $\mathfrak{J}_n\supset\mathfrak{J}^n$,
y supongamos que la topología definida por la filtración $(\mathfrak{J}_n)$ es idéntica a la topología $\mathfrak{J}$-preádica
(en otras palabras, para todo $n$, existe $m$ tal que $\mathfrak{J}^n\supset\mathfrak{J}_m$). Para que $\mathrm{M}$ sea un $\mathrm{A}$-módulo plano,
es necesario y suficiente que $\mathrm{M}\otimes_\mathrm{A}(\mathrm{A}/\mathfrak{J})$ sea un $(\mathrm{A}/\mathfrak{J})$-módulo plano y que, para todo $i$, $\mathrm{M}\otimes_\mathrm{A}\mathrm{A}_{\mathfrak{p}_i}=\mathrm{M}_{\mathfrak{p}_i}$
sea un $\mathrm{A}_{\mathfrak{p}_i}$-módulo plano.
\end{proposition}

\begin{proof}
Como $\mathrm{M}$ está idealmente separado, basta, en virtud de $(\mathbf{0}_{\mathrm{III}},~10.2.1)$, demostrar
que, para todo $n\geq 0$, $\mathrm{M}\otimes_\mathrm{A}(\mathrm{A}/\mathfrak{J}^n)$ es un $(\mathrm{A}/\mathfrak{J}^n)$-módulo plano; como cada $\mathfrak{J}^n$ contiene
un $\mathfrak{J}_m$, esto equivale a demostrar que, para todo $n$, $\mathrm{M}\otimes_\mathrm{A}(\mathrm{A}/\mathfrak{J}_n)$ es un $(\mathrm{A}/\mathfrak{J}_n)$-módulo
plano. Ahora bien, como $\mathfrak{J}_1\supset\mathfrak{J}$, $\mathrm{M}\otimes_\mathrm{A}(\mathrm{A}/\mathfrak{J}_1)$ es un $(\mathrm{A}/\mathfrak{J}_1)$-módulo plano; en el anillo $\mathrm{A}/\mathfrak{J}_n$,
el ideal $\mathfrak{J}_1/\mathfrak{J}_n$ es nilpotente y, por último, la intersección de los núcleos de los homomorfismos
$\mathrm{A}/\mathfrak{J}_n\to\mathrm{A}_{\mathfrak{p}_i}/\mathfrak{p}_i^n\mathrm{A}_{\mathfrak{p}_i}$ es cero en $\mathrm{A}/\mathfrak{J}_n$, por definición de $\mathfrak{J}_n$. Basta, pues, en virtud de \sref{IV.11.4.1},
verificar que $\mathrm{M}\otimes_\mathrm{A}(\mathrm{A}_{\mathfrak{p}_i}/\mathfrak{p}_i^n\mathrm{A}_{\mathfrak{p}_i})$ es un $(\mathrm{A}_{\mathfrak{p}_i}/\mathfrak{p}_i^n\mathrm{A}_{\mathfrak{p}_i})$-módulo plano, lo cual se deduce de
la hipótesis de que $\mathrm{M}\otimes_\mathrm{A}\mathrm{A}_{\mathfrak{p}_i}$ es un $\mathrm{A}_{\mathfrak{p}_i}$-módulo plano.
\end{proof}

\begin{remark}[11.4.8]
\label{IV.11.4.8}
La hipótesis formulada en \sref{IV.11.4.7} sobre la topología definida
por los $\mathfrak{J}_n$ se cumple si, para todo $n$ suficientemente grande, $\operatorname{Ass}(\mathrm{A}/\mathfrak{J}^n)$ \emph{está contenido en el conjunto
de los $\mathfrak{p}_i$}. En efecto, $\mathfrak{J}^n$ es entonces una intersección de ideales primarios correspondientes a los $\mathfrak{p}_i$, cada uno de los cuales
contiene una potencia simbólica de $\mathfrak{p}_i$, de donde se sigue la conclusión. En particular:
\end{remark}

\begin{corollary}[11.4.9]
\label{IV.11.4.9}
Sea $\mathrm{A}$ un anillo noetheriano, sea $\mathfrak{J}$ un ideal nilpotente de $\mathrm{A}$ y sea $\mathrm{M}$ un
$\mathrm{A}$-módulo. Para que $\mathrm{M}$ sea un $\mathrm{A}$-módulo plano (resp.\ libre), es necesario y suficiente que $\mathrm{M}\otimes_\mathrm{A}(\mathrm{A}/\mathfrak{J})$
sea un $(\mathrm{A}/\mathfrak{J})$-módulo plano (resp.\ libre) y que, para todo ideal primo $\mathfrak{p}\in\operatorname{Ass}(\mathrm{A})$, $\mathrm{M}_\mathfrak{p}$ sea un
$\mathrm{A}_\mathfrak{p}$-módulo plano.
\end{corollary}

La afirmación relativa al caso en que $\mathrm{M}\otimes_\mathrm{A}(\mathrm{A}/\mathfrak{J})$ es libre se deduce de nuevo de la afirmación relativa al caso en que $\mathrm{M}\otimes_\mathrm{A}(\mathrm{A}/\mathfrak{J})$ es plano mediante $(\mathbf{0}_{\mathrm{III}},10.1.2)$.

\begin{corollary}[11.4.10]
\label{IV.11.4.10}
Sea $\mathrm{A}$ un anillo noetheriano, sea $\mathfrak{J}$ un ideal de $\mathrm{A}$ y sea $\mathrm{M}$ un $\mathrm{A}$-módulo.
Supongamos que $\mathrm{M}$ está idealmente separado para la topología $\mathfrak{J}$-preádica y que $\operatorname{gr}_\mathfrak{J}(\mathrm{A})$ es un
$(\mathrm{A}/\mathfrak{J})$-módulo plano. Para que $\mathrm{M}$ sea un $\mathrm{A}$-módulo plano, es necesario y suficiente que $\mathrm{M}\otimes_\mathrm{A}(\mathrm{A}/\mathfrak{J})$ sea
un $(\mathrm{A}/\mathfrak{J})$-módulo plano y que, para todo $\mathfrak{p}\in\operatorname{Ass}(\mathrm{A}/\mathfrak{J})$, $\mathrm{M}_\mathfrak{p}$ sea un $\mathrm{A}_\mathfrak{p}$-módulo plano.
\end{corollary}

\begin{proof}
\oldpage[IV-3]{148}
Teniendo en cuenta \sref{IV.11.4.8}, basta demostrar que $\operatorname{Ass}(\mathrm{A}/\mathfrak{J}^n)$ está contenido en
$\operatorname{Ass}(\mathrm{A}/\mathfrak{J})$ para todo $n$. Ahora bien, si $a\in\mathrm{A}$ no pertenece a ningún $\mathfrak{p}\in\operatorname{Ass}(\mathrm{A}/\mathfrak{J})$, la homotecia
de razón $a$ es inyectiva en $\mathrm{A}/\mathfrak{J}$; como cada $\mathfrak{J}^k/\mathfrak{J}^{k+1}$ es un $(\mathrm{A}/\mathfrak{J})$-módulo
plano, $a$ es también un elemento $(\mathfrak{J}^k/\mathfrak{J}^{k+1})$-regular y, por tanto, $a$ es un elemento $(\mathrm{A}/\mathfrak{J}^n)$-regular para todo $n$
(Bourbaki, \emph{Alg.\ comm.}, cap.~III, \S~2, n\textsuperscript{o}~8, cor.~1 del teorema~1) y, por consiguiente, no pertenece
a ningún ideal primo asociado a $\mathrm{A}/\mathfrak{J}^n$, de donde se sigue el corolario (Bourbaki, \emph{Alg.\ comm.}, cap.~II,
\S~1, n\textsuperscript{o}~1, prop.~2).
\end{proof}

\begin{proposition}[11.4.11]
\label{IV.11.4.11}
Sea $\mathrm{A}$ un anillo local artiniano, sea $\mathfrak{m}$ su ideal maximal,
sea $\mathrm{Y}=\operatorname{Spec}(\mathrm{A})$, sea $y$ el único punto de $\mathrm{Y}$, sea $f:\mathrm{X}\to\mathrm{Y}$ un morfismo localmente de tipo finito y sea $\sh{F}$ un
$\sh{O}_\mathrm{X}$-módulo coherente. Sea $(\mathrm{A}'_\alpha)$ una familia de anillos locales y, para cada $\alpha$, sea $u_\alpha:\mathrm{A}\to\mathrm{A}'_\alpha$
un homomorfismo de anillos (necesariamente local). Pongamos $\mathrm{Y}'_\alpha=\operatorname{Spec}(\mathrm{A}'_\alpha)$, $\mathrm{X}'_\alpha=\mathrm{X}\times_\mathrm{Y}\mathrm{Y}'_\alpha$,
$\sh{F}'_\alpha=\sh{F}\otimes_\mathrm{A}\mathrm{A}'_\alpha$. Sea $x$ un punto
de $\mathrm{X}$, y supongamos que se cumplen las siguientes condiciones:
\begin{enumerate}[label=(\roman*)]
  \item La intersección de los núcleos de los $u_\alpha$ es $0$.
  \item La extensión $k(x)$ del cuerpo residual $k$ de $\mathrm{A}$ es primaria $(4.3.1)$ (condición que se cumple automáticamente
  si $k$ es separablemente cerrado).
\end{enumerate}
\begin{enumerate}[label=(\roman*)]
  \setcounter{enumi}{2}
  \item Para todo $\alpha$, existe un punto $x'_\alpha\in\mathrm{X}'_\alpha$ cuyas proyecciones en $\mathrm{X}$ y $\mathrm{Y}'_\alpha$
  son $x$ y el punto cerrado $y'_\alpha$ de $\mathrm{Y}'_\alpha$, y tal que $\sh{F}'_\alpha$ es $\mathrm{Y}'_\alpha$-plano en el punto $x'_\alpha$.
\end{enumerate}

En estas condiciones, $\sh{F}$ es $f$-plano en el punto $x$.
\end{proposition}

\begin{proof}
Como la cuestión es local en $\mathrm{X}$, podemos limitarnos evidentemente al caso en que $f$ es
un morfismo de tipo finito y suponer $\sh{F}_x\neq 0$. Procederemos en varias etapas.

\medskip
\noindent\textbf{I)} \emph{Reducción al caso en que $\mathrm{A}'_\alpha$ es un anillo local de un $\mathrm{Y}$-preesquema de tipo finito y en que el
cuerpo residual $k'_\alpha$ de $\mathrm{A}'_\alpha$ es una extensión finita de $k$.}

\medskip
Como la reducción se efectúa por separado para cada $\mathrm{A}'_\alpha$, podemos omitir el índice $\alpha$ en
esta parte. Sea $\mathfrak{m}'$ el ideal maximal de $\mathrm{A}'$. Consideremos $\mathrm{A}'$ como
límite inductivo de sus $\mathrm{A}$-subálgebras de tipo finito $\mathrm{B}_\lambda$, y pongamos $\mathfrak{n}_\lambda=\mathfrak{m}'\cap\mathrm{B}_\lambda$; $\mathrm{A}'$ es también
límite inductivo de sus subanillos locales $\mathrm{B}'_\lambda=(\mathrm{B}_\lambda)_{\mathfrak{n}_\lambda}$ $(5.13.3)$, y se tiene evidentemente
$\mathfrak{m}'\cap\mathrm{B}'_\lambda=\mathfrak{n}'_\lambda$, ideal maximal de $\mathrm{B}'_\lambda$. Se deduce, pues, de \sref{IV.11.2.6} que, poniendo
$\mathrm{Z}'_\lambda=\operatorname{Spec}(\mathrm{B}'_\lambda)$, $\sh{F}\otimes_\mathrm{A}\mathrm{B}'_\lambda$ es $\mathrm{Z}'_\lambda$-plano en el punto $x'_\lambda$, proyección de $x'$, siendo la proyección
de $x'_\lambda$ en $\mathrm{Z}'_\lambda$ el punto cerrado $z'_\lambda$ de $\mathrm{Z}'_\lambda$.

Podemos, por tanto, suponer que existen un esquema afín $\mathrm{Z}'$ de tipo finito sobre $\mathrm{Y}$ y un
punto $z'$ de $\mathrm{Z}'$ tales que $\mathrm{A}'=\sh{O}_{\mathrm{Z}',z'}$ y que, si se pone $\mathrm{T}'=\mathrm{X}\times_\mathrm{Y}\mathrm{Z}'$, existe un punto
$t'\in\mathrm{T}'$ cuyas proyecciones en $\mathrm{Z}'$ y $\mathrm{X}$ son $z'$ y $x$, y tal que $\sh{F}\otimes_\mathrm{Y}\mathrm{Z}'$ es $\mathrm{Z}'$-plano en el
punto $t'$. Sea $\mathrm{Z}'_1$ (resp.\ $\mathrm{X}_1$) un subpreesquema cerrado de $\mathrm{Z}'$ (resp.\ $\mathrm{X}$) cuyo espacio subyacente sea
la adherencia de $\{z'\}$ (resp.\ $\{x\}$), y pongamos $\mathrm{T}'_1=\mathrm{X}_1\times_\mathrm{Y}\mathrm{Z}'_1$. El conjunto $\mathrm{U}$ de los
puntos de $\mathrm{T}'$ en los que $\sh{F}\otimes_\mathrm{Y}\mathrm{Z}'$ es $\mathrm{Z}'$-plano es abierto en $\mathrm{T}'$ \sref{IV.11.1.1} y contiene $t'$; por tanto,
\oldpage[IV-3]{149}
$\mathrm{V}=\mathrm{U}\cap\mathrm{T}'_1$ es un abierto no vacío de $\mathrm{T}'_1$. Como el anillo $\mathrm{A}$ es artiniano, es un anillo de Jacobson;
por consiguiente, según $(10.4.6)$, $\mathrm{Z}'_1$ y $\mathrm{T}'_1$ son preesquemas de Jacobson. Existe, pues,
un punto $t'_0$ de $\mathrm{V}$ cerrado en $\mathrm{V}$, y su imagen $z'_0$ en $\mathrm{Z}'_1$ es un punto \emph{cerrado} de $\mathrm{Z}'_1$ $(10.4.7)$.
Sea $f_1$ la restricción de $f$ a $\mathrm{X}_1$, y sea $f'_1:\mathrm{T}'_1\to\mathrm{Z}'_1$ la proyección canónica.
$\mathrm{V}\cap f_1'^{-1}(z'_0)$ es un subconjunto abierto no vacío de
$f_1^{-1}(y)\otimes_{k(y)}k(z'_0)$; como este último preesquema es plano sobre $f_1^{-1}(y)$, un punto maximal $t'_1$ de
$\mathrm{V}\cap f_1'^{-1}(z'_0)$ se encuentra necesariamente sobre el único punto maximal $x$ de $\mathrm{X}_1$ $(2.3.4)$.
Por último, $k(z'_0)$ es una extensión \emph{finita} de $k$ (\textbf{I},~$6.4.2$), y el homomorfismo
$\mathrm{A}=\sh{O}_{\mathrm{Y},y}\to\mathrm{A}'=\sh{O}_{\mathrm{Z}',z'}$ se factoriza como
$\sh{O}_{\mathrm{Y},y}\to\sh{O}_{\mathrm{Z}',z'_0}\to\sh{O}_{\mathrm{Z}',z'}$; por lo tanto, su núcleo está contenido en el de
$\sh{O}_{\mathrm{Y},y}\to\sh{O}_{\mathrm{Z}',z'_0}$. Esto completa la reducción anunciada.

\medskip
\noindent\textbf{II)} \emph{Reducción al caso en que solo hay un número finito de $\mathrm{A}'_\alpha$ y estos son $\mathrm{A}$-álgebras finitas.} --- Sea $\mathfrak{m}'_\alpha$
el ideal maximal de $\mathrm{A}'_\alpha$. Como $\mathrm{A}'_\alpha$ es un anillo local noetheriano, la intersección de los $\mathfrak{m}'^n_\alpha$
es $0$ $(\mathbf{0}_\mathrm{I},~7.3.5)$. La intersección de los $u_\alpha^{-1}(\mathfrak{m}'^n_\alpha)$, para todos los índices $n$ y $\alpha$, es, por tanto,
igual a la intersección de los núcleos de los $u_\alpha$ y, en consecuencia, es $0$ por la hipótesis (i). Como $\mathrm{A}$
es artiniano, un número finito de estos ideales tiene ya intersección $0$; denotémoslos por
$u_i^{-1}(\mathfrak{m}'^{n_i}_i)$ $(1\leq i\leq r)$. Como los $\mathrm{A}'_i$ son noetherianos, los $\mathfrak{m}'^h_i/\mathfrak{m}'^{h+1}_i$ son
$(\mathrm{A}'_i/\mathfrak{m}'_i)$-módulos de longitud finita y, como $\mathrm{A}'_i/\mathfrak{m}'_i$ es un $k$-espacio vectorial de dimensión finita,
$\mathrm{A}''_i=\mathrm{A}'_i/\mathfrak{m}'^{n_i}_i$ es una $\mathrm{A}$-álgebra finita y un anillo local artiniano. La reducción anunciada
se deduce, pues, de $(2.1.4)$.

\medskip
\noindent\textbf{III)} \emph{Fin de la demostración.} --- Supongamos ahora que solo hay un número finito de $\mathrm{A}'_i$ $(1\leq i\leq r)$ y que
son $\mathrm{A}$-álgebras finitas. Para todo $i$, el cuerpo residual $k_i$ de $\mathrm{A}'_i$ es una extensión finita de $k$; utilizando
la hipótesis (ii), se concluye que la imagen inversa de $x$ en $\mathrm{X}'_i$ consta del \emph{único} punto $x'_i$ $(4.3.2)$.
Sea $\mathrm{Y}'$ la unión disjunta de los $\mathrm{Y}'_i$, sea $\mathrm{X}'=\mathrm{X}\times_\mathrm{Y}\mathrm{Y}'$ la unión disjunta de los
$\mathrm{X}'_i$ y pongamos $\sh{F}'=\sh{F}\otimes_\mathrm{Y}\mathrm{Y}'$. La hipótesis implica que $\sh{F}'$ es $\mathrm{Y}'$-plano en todo
punto de la imagen inversa de $x$ por la proyección $p:\mathrm{X}'\to\mathrm{X}$. Como $p$ es de tipo finito, existe, por tanto, un
abierto $\mathrm{U}'\supset p^{-1}(x)$ tal que $\sh{F}'$ es $\mathrm{Y}'$-plano en todo punto de $\mathrm{U}'$ \sref{IV.11.1.1}.
Además, el morfismo $\mathrm{Y}'\to\mathrm{Y}$ es finito porque los $\mathrm{A}'_i$ son $\mathrm{A}$-álgebras finitas. Así, $p$ es un
morfismo finito (\textbf{II},~$6.1.5$) y, por tanto, es cerrado (\textbf{II},~$6.1.10$); en consecuencia, existe una vecindad afín $\mathrm{U}$ de
$x$ en $\mathrm{X}$ tal que $p^{-1}(\mathrm{U})\subset\mathrm{U}'$. Sean $\mathrm{B}$ y $\mathrm{A}'$ los anillos de los esquemas
$\mathrm{U}$ y $\mathrm{Y}'$, respectivamente (así, $\mathrm{A}'$ es el producto directo de los $\mathrm{A}'_i$). Sustituyendo $\mathrm{X}$ por
$\mathrm{U}$, se tiene $\sh{F}=\widetilde{\mathrm{M}}$, donde $\mathrm{M}$ es un $\mathrm{B}$-módulo, y por hipótesis
$\mathrm{M}'=\mathrm{M}\otimes_\mathrm{A}\mathrm{A}'$ es un $\mathrm{A}'$-módulo plano $(2.1.2)$. Como el homomorfismo
$\mathrm{A}\to\mathrm{A}'$ es inyectivo por construcción, se puede aplicar \sref{IV.11.4.3}, que demuestra que $\mathrm{M}$ es un
$\mathrm{A}$-módulo plano.
\end{proof}

\begin{proposition}[11.4.12]
\label{IV.11.4.12}
Sea $\mathrm{A}$ un anillo local artiniano con cuerpo residual $k$,
sea $\mathrm{Y}=\operatorname{Spec}(\mathrm{A})$, sea $f:\mathrm{X}\to\mathrm{Y}$ un morfismo localmente de tipo finito y sea $\sh{F}$ un $\sh{O}_\mathrm{X}$-módulo coherente.
Sea $(\mathrm{A}'_\alpha)$ una familia de anillos locales y, para cada $\alpha$, sea $u_\alpha:\mathrm{A}\to\mathrm{A}'_\alpha$ un homomorfismo
de anillos. Pongamos $\mathrm{Y}'_\alpha=\operatorname{Spec}(\mathrm{A}'_\alpha)$, $\mathrm{X}'_\alpha=\mathrm{X}\times_\mathrm{Y}\mathrm{Y}'_\alpha$,
$f'_\alpha=f_{(\mathrm{Y}'_\alpha)}$ y $\sh{F}'_\alpha=\sh{F}\otimes_\mathrm{A}\mathrm{A}'_\alpha$. Sea $x$ un punto
de $\mathrm{X}$, y supongamos que se cumplen las siguientes condiciones:
\begin{enumerate}[label=(\roman*)]
  \item La intersección de los núcleos de los $u_\alpha$ es $0$.
  \item Para todo $\alpha$, $\sh{F}'_\alpha$ es $f'_\alpha$-plano en todo punto $x'_\alpha\in\mathrm{X}'_\alpha$ cuyas proyecciones en $\mathrm{X}$ y $\mathrm{Y}'_\alpha$
  son $x$ y el punto cerrado $y'_\alpha$ de $\mathrm{Y}'_\alpha$.
\end{enumerate}

Entonces $\sh{F}$ es $f$-plano en el punto $x$.
\end{proposition}

\begin{proof}
\oldpage[IV-3]{150}
Por hipótesis, la intersección de los núcleos de los $u_\alpha$ es $0$; como $\mathrm{A}$ es artiniano,
existe ya un número \emph{finito} de estos núcleos cuya intersección es $0$, de modo que podemos
limitarnos al caso en que la familia $(\mathrm{A}'_\alpha)$ es \emph{finita}. Sea $k'$ una clausura algebraica de $k$; se sabe
por $(\mathbf{0}_{\mathrm{III}},~10.3.1)$ que existe un homomorfismo local $\mathrm{A}\to\mathrm{B}$ que hace de $\mathrm{B}$ un $\mathrm{A}$-módulo plano,
tal que $\mathrm{B}$ sea un anillo local artiniano, integral sobre $\mathrm{A}$, y que $\mathrm{B}\otimes_\mathrm{A} k$ es isomorfo a $k'$.
Por planitud, los núcleos de los homomorfismos $\mathrm{B}\to\mathrm{B}'_\alpha=\mathrm{B}\otimes_\mathrm{A}\mathrm{A}'_\alpha$ se deducen de los de
$u_\alpha$ por tensorización con $\mathrm{B}$ y, como son un número finito, su intersección es
$0$ $(\mathbf{0}_\mathrm{I},~6.1.3)$. Consideremos los anillos $\mathrm{B}'_{\alpha\beta}$, localizaciones de $\mathrm{B}'_\alpha$ en sus ideales \emph{maximales};
se sabe (Bourbaki, \emph{Alg.\ comm.}, cap.~II, \S~3, n\textsuperscript{o}~3, cor.~2 del teorema~1) que la intersección
de los núcleos de los homomorfismos $\mathrm{B}'_\alpha\to\mathrm{B}'_{\alpha\beta}$ (para un $\alpha$ dado) es $0$; se concluye, por tanto,
que la intersección de los núcleos de los homomorfismos compuestos $v_{\alpha\beta}:\mathrm{B}\to\mathrm{B}'_\alpha\to\mathrm{B}'_{\alpha\beta}$ (con $\alpha$ y
$\beta$ variables) es $0$. Por otra parte, como $\mathrm{B}$ es integral sobre $\mathrm{A}$, $\mathrm{B}'_\alpha$ es integral sobre $\mathrm{A}'_\alpha$,
de modo que sus ideales maximales se encuentran sobre el ideal maximal de $\mathrm{A}'_\alpha$. Si se pone
$\mathrm{Z}'_{\alpha\beta}=\operatorname{Spec}(\mathrm{B}'_{\alpha\beta})$,
$\mathrm{T}'_{\alpha\beta}=\mathrm{X}'_\alpha\times_{\mathrm{Y}'_\alpha}\mathrm{Z}'_{\alpha\beta}$,
$f'_{\alpha\beta}=f_{(\mathrm{Z}'_{\alpha\beta})}$,
$\sh{G}=\sh{F}\otimes_\mathrm{A}\mathrm{B}$, y
$\sh{G}'_{\alpha\beta}=\sh{G}\otimes_\mathrm{B}\mathrm{B}'_{\alpha\beta}$, se ve
que las hipótesis (i) y (ii) siguen cumpliéndose cuando se sustituyen
$\mathrm{A}$, $\mathrm{A}'_\alpha$, $u_\alpha$, $\sh{F}$ y $x$ por $\mathrm{B}$, $\mathrm{B}'_{\alpha\beta}$, $v_{\alpha\beta}$, $\sh{G}$ y un punto $t$ de $\mathrm{T}=\mathrm{X}\otimes_\mathrm{A}\mathrm{B}$ situado sobre $x$.
Como el cuerpo residual de $\mathrm{B}$ es separablemente cerrado, se deduce de \sref{IV.11.4.11} que
$\sh{G}$ es plano sobre $\mathrm{B}$ en el punto $t$. Pero, como $\mathrm{B}$ es un $\mathrm{A}$-módulo fielmente plano, se deduce de $(2.1.4)$
que $\sh{F}$ es plano sobre $\mathrm{A}$ en el punto $x$, lo que demuestra \sref{IV.11.4.12}.
\end{proof}

\begin{corollary}[11.4.13]
\label{IV.11.4.13}
Sea $\mathrm{A}$ un anillo local artiniano, sea $\mathrm{Y}=\operatorname{Spec}(\mathrm{A})$, sea $f:\mathrm{X}\to\mathrm{Y}$ un
morfismo localmente de tipo finito y sea $\sh{F}$ un $\sh{O}_\mathrm{X}$-módulo coherente. Sea $(\mathrm{Y}'_\alpha)$ una familia de
$\mathrm{Y}$-preesquemas y, para todo $\alpha$, pongamos $\mathrm{X}'_\alpha=\mathrm{X}\times_\mathrm{Y}\mathrm{Y}'_\alpha$, $f'_\alpha=f_{(\mathrm{Y}'_\alpha)}$ y $\sh{F}'_\alpha=\sh{F}\otimes_\mathrm{Y}\mathrm{Y}'_\alpha$. Sea $x$ un punto
de $\mathrm{X}$ y supongamos que se cumplen las siguientes hipótesis:
\begin{enumerate}[label=(\roman*)]
  \item La intersección de los núcleos de los homomorfismos $\mathrm{A}\to\Gamma(\mathrm{Y}'_\alpha,\sh{O}_{\mathrm{Y}'_\alpha})$ correspondientes a los
  morfismos estructurales es $0$.
  \item Para todo $\alpha$, $\sh{F}'_\alpha$ es $f'_\alpha$-plano en todo punto $x'_\alpha\in\mathrm{X}'_\alpha$ cuya proyección en $\mathrm{X}$ es $x$.
\end{enumerate}

Entonces $\sh{F}$ es $f$-plano en el punto $x$.
\end{corollary}

\begin{proof}
En efecto, para todo $y'_\alpha\in\mathrm{Y}'_\alpha$, consideremos el esquema local
$\mathrm{Y}''_{\alpha,y'_\alpha}=\operatorname{Spec}(\sh{O}_{\mathrm{Y}'_\alpha,y'_\alpha})$; en
virtud de $(2.1.4)$, $\sh{F}'_\alpha\otimes_{\mathrm{Y}'_\alpha}\mathrm{Y}''_{\alpha,y'_\alpha}$ es
$\mathrm{Y}''_{\alpha,y'_\alpha}$-plano en los puntos cuyas proyecciones en $\mathrm{X}$ y
$\mathrm{Y}''_{\alpha,y'_\alpha}$ son $x$ y el punto cerrado de $\mathrm{Y}''_{\alpha,y'_\alpha}$.
Además, el núcleo del homomorfismo
$\mathrm{A}\to\Gamma(\mathrm{Y}'_\alpha,\sh{O}_{\mathrm{Y}'_\alpha})$ es la intersección de los núcleos de los homomorfismos
$\mathrm{A}\to\Gamma(\mathrm{Y}''_{\alpha,y'_\alpha},\sh{O}_{\mathrm{Y}''_{\alpha,y'_\alpha}})$:
en efecto, nos reducimos inmediatamente al caso en que $\mathrm{Y}'_\alpha$ es afín, y basta entonces aplicar Bourbaki, \emph{Alg.\ comm.},
cap.~II, \S~3, n\textsuperscript{o}~3, cor.~2 del teorema~1. Sustituyendo la familia $(\mathrm{Y}'_\alpha)$ por la familia de los $\mathrm{Y}''_{\alpha,y'_\alpha}$,
nos reducimos, pues, a \sref{IV.11.4.12}.
\end{proof}

\subsection{Descenso de la planitud por morfismos arbitrarios: el caso general}
\label{subsection:IV.11.5}

\begin{theorem}[11.5.1]
\label{IV.11.5.1}
Sean $\mathrm{Y}$ un preesquema localmente noetheriano, $f:\mathrm{X}\to\mathrm{Y}$ un
morfismo localmente de tipo finito, $\sh{F}$ un $\sh{O}_\mathrm{X}$-módulo coherente, $x$ un punto de $\mathrm{X}$ y $y=f(x)$. Sea $(\mathrm{Y}'_\alpha)$
una familia de $\mathrm{Y}$-preesquemas locales $\mathrm{Y}'_\alpha=\operatorname{Spec}(\mathrm{A}'_\alpha)$ tales que las imágenes de los puntos cerrados $y'_\alpha$
de $\mathrm{Y}'_\alpha$ sean todas iguales a $y$. Para todo $\alpha$, sean $\mathfrak{m}'_\alpha$ el ideal maximal de $\mathrm{A}'_\alpha$ y $u_\alpha:\sh{O}_y\to\mathrm{A}'_\alpha$
\oldpage[IV-3]{151}
el homomorfismo canónico (\textbf{I},~$2.4.4$); supongamos que las intersecciones finitas de los ideales
$u_\alpha^{-1}(\mathfrak{m}'_\alpha{}^{n_\alpha})$
formen un sistema fundamental de entornos de $0$ en $\sh{O}_y$. Pongamos $\mathrm{X}'_\alpha=\mathrm{X}\times_\mathrm{Y}\mathrm{Y}'_\alpha$, $f'_\alpha=f_{(\mathrm{Y}'_\alpha)}$,
$\sh{F}'_\alpha=\sh{F}\otimes_\mathrm{Y}\mathrm{Y}'_\alpha$, y supongamos que se cumple una de las hipótesis siguientes:
\begin{enumerate}[label=(\roman*)]
  \item Para todo $\alpha$, $\sh{F}'_\alpha$ es $f'_\alpha$-plano en todo punto cuya proyección en $\mathrm{X}$ es igual
  a $x$ y cuya proyección en $\mathrm{Y}'_\alpha$ es igual a $y'_\alpha$.
  \item Para todo $\alpha$, existe un punto $x'_\alpha\in\mathrm{X}'_\alpha$ cuya proyección en $\mathrm{X}$ es $x$ y cuya proyección en $\mathrm{Y}'_\alpha$
  es igual a $y'_\alpha$, tal que $\sh{F}'_\alpha$ es $f'_\alpha$-plano en el punto $x'_\alpha$ y $k(x)$ es una extensión primaria de $k(y)$.
\end{enumerate}

En estas condiciones, $\sh{F}$ es $f$-plano en el punto $x$.
\end{theorem}

\begin{proof}
Sea $\mathfrak{m}_y$ el ideal maximal de $\sh{O}_y$; como $\sh{O}_y$ y $\sh{O}_x$ son noetherianos y $\sh{O}_y\to\sh{O}_x$ es un
homomorfismo de anillos locales, basta, en virtud de $(\mathbf{0}_{\mathrm{III}},~10.2.2)$, demostrar que, para
todo $n>0$, $\sh{F}_x/\mathfrak{m}_y^n\sh{F}_x$ es un $(\sh{O}_y/\mathfrak{m}_y^n)$-módulo plano. Denotemos por $(\mathfrak{J}_\lambda)$ la familia de
intersecciones finitas de los $u_\alpha^{-1}(\mathfrak{m}'_\alpha{}^{n_\alpha})$; por hipótesis, existe $\lambda$ tal que
$\mathfrak{J}_\lambda\subset\mathfrak{m}_y^n$, y como
$\sh{F}_x\otimes_{\sh{O}_y}(\sh{O}_y/\mathfrak{m}_y^n)=(\sh{F}_x\otimes_{\sh{O}_y}(\sh{O}_y/\mathfrak{J}_\lambda))\otimes_{\sh{O}_y/\mathfrak{J}_\lambda}(\sh{O}_y/\mathfrak{m}_y^n)$, bastará con probar que $\sh{F}_x/\mathfrak{J}_\lambda\sh{F}_x$ es
un $(\sh{O}_y/\mathfrak{J}_\lambda)$-módulo plano. Ahora bien, para cada $\alpha$ tal que
$\mathfrak{J}_\lambda\subset u_\alpha^{-1}(\mathfrak{m}'_\alpha{}^{n_\alpha})$, el paso a los cocientes da
un homomorfismo $u'_\alpha:\sh{O}_y/\mathfrak{J}_\lambda\to\mathrm{A}'_\alpha/\mathfrak{m}'_\alpha{}^{n_\alpha}$, y la intersección de los núcleos
de los $u'_\alpha$ se reduce a $0$. Teniendo en cuenta (\textbf{I},~$3.6.1$), vemos que nos reducimos a \sref{IV.11.4.11}
en el caso (ii), y a \sref{IV.11.4.12} en el caso (i).
\end{proof}

\begin{corollary}[11.5.2]
\label{IV.11.5.2}
Sean $\mathrm{Y}$ un preesquema localmente noetheriano, $f:\mathrm{X}\to\mathrm{Y}$ un
morfismo localmente de tipo finito, $\sh{F}$ un $\sh{O}_\mathrm{X}$-módulo coherente, $x$ un punto de $\mathrm{X}$ y $y=f(x)$; pongamos
$\mathrm{A}=\sh{O}_y$. Sea $u:\mathrm{A}\to\mathrm{A}'$ un homomorfismo de $\mathrm{A}$ en un anillo de Zariski $\mathrm{A}'$, tal que
la imagen inversa $u^{-1}(\mathfrak{r}')$ del radical $\mathfrak{r}'$ de $\mathrm{A}'$ sea el ideal maximal $\mathfrak{m}$ de $\mathrm{A}$; supongamos además
que el homomorfismo $\hat{u}:\hat{\mathrm{A}}\to\hat{\mathrm{A}}'$ es inyectivo. Pongamos $\mathrm{Y}'=\operatorname{Spec}(\mathrm{A}')$, $\mathrm{X}'=\mathrm{X}\times_\mathrm{Y}\mathrm{Y}'$, $f'=f_{(\mathrm{Y}')}$,
$\sh{F}'=\sh{F}\otimes_\mathrm{Y}\mathrm{Y}'$. Para que $\sh{F}$ sea $f$-plano en el punto $x$, es necesario y suficiente que $\sh{F}'$ sea $f'$-plano en todo
punto cuya proyección en $\mathrm{X}$ es igual a $x$ y cuya proyección en $\mathrm{Y}'$ es igual a un punto cerrado
de $\mathrm{Y}'$.

Si además $\mathrm{A}'$ es un $\mathrm{A}$-álgebra finita, en lo anterior se puede sustituir la hipótesis de que $\hat{u}$
es inyectivo por la hipótesis de que $u$ es inyectivo.
\end{corollary}

\begin{proof}
Como $\mathrm{A}$ (resp.\ $\mathrm{A}'$) se identifica con un subanillo de $\hat{\mathrm{A}}$ (resp.\ $\hat{\mathrm{A}}'$) (Bourbaki,
\emph{Alg.\ comm.}, cap.~III, \S~3, n\textsuperscript{o}~3, prop.~6), vemos primero que $u$ es inyectivo
y que $\hat{u}$ es su extensión por continuidad a $\hat{\mathrm{A}}$.

Sea $(\mathfrak{m}'_\alpha)$ la familia de ideales maximales de $\mathrm{A}'$; como se tiene
\[
\mathfrak{m}'_\alpha{}^n\hat{\mathrm{A}}'=\hat{\mathfrak{m}}'_\alpha{}^n, \qquad \text{y} \qquad \hat{\mathfrak{m}}'_\alpha{}^n\cap\mathrm{A}'=\mathfrak{m}'_\alpha{}^n,
\]
y como los $\mathfrak{m}'_\alpha{}^n$ son abiertos en $\mathrm{A}'$, se tiene
$u^{-1}(\mathfrak{m}'_\alpha{}^n)=\mathrm{A}\cap\hat{u}^{-1}(\hat{\mathfrak{m}}'_\alpha{}^n)$, y será suficiente probar que en
$\hat{\mathrm{A}}$ las intersecciones finitas de los $\hat{u}^{-1}(\hat{\mathfrak{m}}'_\alpha{}^{n_\alpha})$ forman un sistema fundamental de entornos
de $0$. Esto permitirá la aplicación de \sref{IV.11.5.1} a los homomorfismos compuestos
$v_\alpha\circ u:\mathrm{A}\to\mathrm{A}'_{\mathfrak{m}'_\alpha}$, donde
$v_\alpha:\mathrm{A}'\to\mathrm{A}'_{\mathfrak{m}'_\alpha}$ es el homomorfismo canónico. Como $\hat{\mathrm{A}}$ es completo, bastará
probar que la intersección de los $\hat{u}^{-1}(\hat{\mathfrak{m}}'_\alpha{}^{n_\alpha})$ se reduce a $0$ (Bourbaki,
\emph{Alg.\ comm.}, cap.~III, \S~2, n\textsuperscript{o}~7, prop.~8, donde en la demostración la sucesión decreciente puede sustituirse por un conjunto
filtrante arbitrario). Ahora bien, para cada $\alpha$
\oldpage[IV-3]{152}
fijo, la intersección de los ideales $\hat{\mathfrak{m}}'_\alpha{}^n\hat{\mathrm{A}}'_{\hat{\mathfrak{m}}'_\alpha}$ para $n>0$ se reduce a $0$ en el
anillo local noetheriano $\hat{\mathrm{A}}'_{\hat{\mathfrak{m}}'_\alpha}$.
Por otra parte, los $\hat{\mathfrak{m}}'_\alpha$ son los ideales maximales de $\hat{\mathrm{A}}'$; por tanto, el homomorfismo canónico
$\hat{\mathrm{A}}'\to\prod_\alpha\hat{\mathrm{A}}'_{\hat{\mathfrak{m}}'_\alpha}$ es inyector (Bourbaki, \emph{Alg.\ comm.}, cap.~II, \S~3, n\textsuperscript{o}~3, cor.~2 de th.~1), y
y como, por hipótesis, $\hat{u}:\hat{\mathrm{A}}\to\hat{\mathrm{A}}'$ también es inyectivo, esto concluye la demostración en
el caso general. La última afirmación se deduce de que $\hat{\mathrm{A}}$ es un $\mathrm{A}$-módulo fielmente
plano $(\mathbf{0}_\mathrm{I},~7.3.5)$ y $\hat{\mathrm{A}}'=\mathrm{A}'\otimes_\mathrm{A}\hat{\mathrm{A}}$, pues $\mathrm{A}'$ es un $\mathrm{A}$-módulo de tipo finito
(Bourbaki, \emph{Alg.\ comm.}, cap.~III, \S~3, n\textsuperscript{o}~4, th.~3 y cap.~IV, \S~2, n\textsuperscript{o}~5, cor.~3 de prop.~9).
\end{proof}

\begin{proposition}[11.5.3]
\label{IV.11.5.3}
Sean $f:\mathrm{X}\to\mathrm{Y}$ un morfismo localmente de presentación finita,
$\sh{F}$ un $\sh{O}_\mathrm{X}$-módulo cuasicoherente de presentación finita, $x$ un punto de $\mathrm{X}$, $y=f(x)$ y $g=(\psi,\theta):\mathrm{Y}'\to\mathrm{Y}$
un morfismo propio de presentación finita. Supongamos que:
\begin{enumerate}[label=(\roman*)]
  \item El homomorfismo $\theta_y:\sh{O}_y\to(g_*(\sh{O}_{\mathrm{Y}'}))_y$ es inyector.
  \item Para todo $x'\in\mathrm{X}'=\mathrm{X}\times_\mathrm{Y}\mathrm{Y}'$ cuya proyección en $\mathrm{X}$ es $x$, $\sh{F}'=\sh{F}\otimes_\mathrm{Y}\mathrm{Y}'$ es
  plano sobre $\mathrm{Y}'$ en el punto $x'$.
\end{enumerate}

Entonces $\sh{F}$ es $f$-plano en el punto $x$.
\end{proposition}

\begin{proof}
Como la cuestión es local sobre $\mathrm{X}$, podemos suponer que $f$ es de presentación finita. Sea
$p:\mathrm{X}'\to\mathrm{X}$ la proyección canónica. Como
$f'=f_{|\mathrm{Y}'}:\mathrm{X}'\to\mathrm{Y}'$ es de presentación finita
$(1.6.2)$ y $\sh{F}'$ es un $\sh{O}_{\mathrm{X}'}$-módulo de presentación finita $(\mathbf{0}_\mathrm{I},~5.2.5)$, se deduce
de \sref{IV.11.3.1} que el conjunto $\mathrm{U}'$ de puntos de $\mathrm{X}'$ en los que $\sh{F}'$ es $f'$-plano es abierto. Como $\mathrm{U}'$
contiene $p^{-1}(x)$ por hipótesis y $p$ es propio, y por tanto \emph{cerrado}, $\mathrm{U}'$ contiene un conjunto de la forma
$p^{-1}(\mathrm{U})$, donde $\mathrm{U}$ es un entorno de $x$. Sustituyendo $\mathrm{X}$ por $\mathrm{U}$, podemos, por tanto,
suponer que $\sh{F}'$ es $f'$-plano. Por otra parte, teniendo en cuenta (\textbf{I},~$3.6.5$), (\textbf{II},~$5.4.2$)
y $(2.1.4)$, también se puede reemplazar $\mathrm{Y}$ por $\operatorname{Spec}(\sh{O}_y)$, es decir, supongamos que $\mathrm{Y}=\operatorname{Spec}(\mathrm{A})$,
donde $\mathrm{A}$ es un anillo local. En estas condiciones demostraremos que $\sh{F}$ es $f$-plano. En virtud
de $(5.13.3)$, $\mathrm{A}$ es el límite inductivo de subanillos locales noetherianos $\mathrm{A}_\lambda$ tales que
la inyección canónica $\mathrm{A}_\lambda\to\mathrm{A}$ es un homomorfismo \emph{local}. En virtud de $(8.9.1)$, podemos
puede suponer que $\mathrm{X}=\mathrm{X}_\lambda\otimes_{\mathrm{A}_\lambda}\mathrm{A}$, $f=f_\lambda\otimes 1_\mathrm{A}$, $\sh{F}=\sh{F}_\lambda\otimes_{\mathrm{A}_\lambda}\mathrm{A}$,
para un $\lambda$ adecuado, donde $f_\lambda:\mathrm{X}_\lambda\to\mathrm{Y}_\lambda=\operatorname{Spec}(\mathrm{A}_\lambda)$ es un morfismo de tipo finito y
$\sh{F}_\lambda$ es un $\sh{O}_{\mathrm{X}_\lambda}$-módulo coherente.
Del mismo modo, podemos suponer que $\mathrm{Y}'=\mathrm{Y}'_\lambda\otimes_{\mathrm{A}_\lambda}\mathrm{A}$ y $g=g_\lambda\otimes 1_\mathrm{A}$, donde $g_\lambda:\mathrm{Y}'_\lambda\to\mathrm{Y}_\lambda$ es un
morfismo de presentación finita; además, en virtud de $(8.10.5,~(\text{xii}))$, podemos suponer que $g_\lambda$ es
\emph{propio}. Como, por hipótesis, el homomorfismo $\mathrm{A}\to\Gamma(\mathrm{Y}',\sh{O}_{\mathrm{Y}'})$ es inyectivo y $\mathrm{A}_\lambda$
es un subanillo de $\mathrm{A}$, el homomorfismo $\mathrm{A}_\lambda\to\Gamma(\mathrm{Y}',\sh{O}_{\mathrm{Y}'})$ también es inyectivo. Finalmente,
en virtud de \sref{IV.11.2.6}, podemos tomar $\lambda$ de manera que $\sh{F}'_\lambda=\sh{F}_\lambda\otimes_{\mathrm{Y}_\lambda}\mathrm{Y}'_\lambda$ sea $f'_\lambda$-plano,
pues $\sh{F}'=\sh{F}'_\lambda\otimes_{\mathrm{Y}'_\lambda}\mathrm{Y}'$. Estas observaciones muestran que, en adelante, podemos suponer que
el anillo $\mathrm{A}$ es \emph{noetheriano}, manteniéndose las demás hipótesis de \sref{IV.11.5.3}. Pongamos entonces
$\mathrm{B}=\hat{\mathrm{A}}$ y $\mathrm{Z}=\operatorname{Spec}(\mathrm{B})$; como $\mathrm{B}$ es un $\mathrm{A}$-módulo fielmente plano $(\mathbf{0}_\mathrm{I},~7.3.5)$, equivale
a decir que $\sh{F}$ es $f$-plano o que $\sh{F}\otimes_\mathrm{Y}\mathrm{Z}$ es plano sobre $\mathrm{Z}$ $(2.1.4)$; análogamente, si se
pone $\mathrm{Z}'=\mathrm{Y}'\times_\mathrm{Y}\mathrm{Z}$, el morfismo $\mathrm{Z}'\to\mathrm{Y}'$ es fielmente plano $(2.2.13)$, por lo que
equivale a decir que $\sh{F}'$ es $f'$-plano o que $\sh{F}'\otimes_{\mathrm{Y}'}\mathrm{Z}'$ es plano sobre $\mathrm{Z}'$; finalmente, $h=g_{(\mathfrak{m})}:\mathrm{Z}'\to\mathrm{Z}$
es propio (\textbf{II},~$5.4.2$) y de tipo finito $(1.5.2)$, $\hat{\mathrm{A}}$ es noetheriano y, si $z$ es su
\oldpage[IV-3]{153}
único punto cerrado, el homomorfismo $\sh{O}_z\to(h_*(\sh{O}_{\mathrm{Z}'}))_z$ es \emph{inyectivo}, pues se deduce de $(2.3.1)$ que
$h_*(\sh{O}_{\mathrm{Z}'})=g_*(\sh{O}_{\mathrm{Y}'})\otimes_\mathrm{Y}\mathrm{Z}$, y nuestra afirmación se deriva de la hipótesis (i) y de la definición de
módulos planos $(\mathbf{0}_\mathrm{I},~6.1.1)$.

Por tanto, podemos suponer ahora que el anillo local noetheriano $\mathrm{A}$ es \emph{completo}; la demostración
se terminará si se demuestra que la intersección de los núcleos de los homomorfismos
$\mathrm{A}=\sh{O}_y\to\sh{O}_{y'}$ (donde $y'$ recorre $g^{-1}(y)$) se reduce a $0$. En efecto, los $\sh{O}_{y'}$ son anillos locales
noetherianos; por tanto, para cada $y'$ la intersección de los $\mathfrak{m}_{y'}^n$ $(n\geq 0)$ se reduce a $0$;
si $\mathfrak{a}_{n,y'}$ es la imagen inversa de $\mathfrak{m}_{y'}^n$ en $\mathrm{A}$, las intersecciones de finitamente muchos de los $\mathfrak{a}_{n,y'}$ son
entornos de $0$ en $\mathrm{A}$ y la intersección de todos los $\mathfrak{a}_{n,y'}$ se reduce a $0$; las intersecciones de
un número finito de los $\mathfrak{a}_{n,y'}$ formarán, por tanto, un sistema fundamental de entornos de $0$ en $\mathrm{A}$
(Bourbaki, \emph{Alg.\ comm.}, cap.~III, \S~2, n\textsuperscript{o}~7, prop.~8, donde uno puede en la prueba
reemplazar la sucesión decreciente por un conjunto filtrante arbitrario); podremos entonces
aplicar \sref{IV.11.5.1}. Sea ahora $s\in\mathrm{A}$ un elemento perteneciente al núcleo de cada uno de
los homomorfismos $\mathrm{A}\to\sh{O}_{y'}$; la imagen $s'$ de $s$ en $\Gamma(\mathrm{Y}',\sh{O}_{\mathrm{Y}'})$ es, por lo tanto, una sección
de $\sh{O}_{\mathrm{Y}'}$ sobre $\mathrm{Y}'$ tal que $s'_{y'}=0$ para todo $y'\in g^{-1}(y)$; existe entonces un entorno $\mathrm{V}'$
de $g^{-1}(y)$ en $\mathrm{Y}'$ tal que $s'|\mathrm{V}'=0$. Pero, como $g$ es \emph{cerrado}, $\mathrm{V}'$ contiene un conjunto
de la forma $g^{-1}(\mathrm{V})$, donde $\mathrm{V}$ es un entorno abierto de $y$ en $\mathrm{Y}$; ahora bien, $\mathrm{Y}$ es un preesquema local
y $y$ es su punto cerrado, de modo que $\mathrm{V}=\mathrm{Y}$ y, en consecuencia, $\mathrm{V}'=\mathrm{Y}'$,
$s'=0$ y, como $\mathrm{A}\to\Gamma(\mathrm{Y}',\sh{O}_{\mathrm{Y}'})$ es inyectivo por hipótesis, $s=0$.
\end{proof}

El siguiente caso particular de \sref{IV.11.5.3} será útil en el cap.~V:

\begin{corollary}[11.5.4]
\label{IV.11.5.4}
Sea $f=(\psi,\theta):\mathrm{X}\to\mathrm{Y}$ un morfismo propio de presentación finita,
y sea $p:\mathrm{X}\times_\mathrm{Y}\mathrm{X}\to\mathrm{X}$ la primera proyección. Supongamos que $\theta:\sh{O}_\mathrm{Y}\to f_*(\sh{O}_\mathrm{X})$ es inyectivo.
Para que $f$ sea plano, es necesario y suficiente que $p$ sea plano.
\end{corollary}

\begin{proposition}[11.5.5]
\label{IV.11.5.5}
Sean $\mathrm{A}$ un anillo, $\mathrm{Y}=\operatorname{Spec}(\mathrm{A})$, $f:\mathrm{X}\to\mathrm{Y}$ un morfismo
localmente de presentación finita, $\sh{F}$ un cuasi-coherente $\sh{O}_\mathrm{X}$-módulo de presentación finita, $x$ un punto
de $\mathrm{X}$. Sea $\mathrm{A}\to\mathrm{A}'$ un homomorfismo inyectivo que hace de $\mathrm{A}'$ un álgebra integral sobre $\mathrm{A}$. Pongamos
$\mathrm{Y}'=\operatorname{Spec}(\mathrm{A}')$, $\mathrm{X}'=\mathrm{X}\times_\mathrm{Y}\mathrm{Y}'$, $f'=f\times 1$, $\sh{F}'=\sh{F}\otimes_\mathrm{Y}\mathrm{Y}'$. Entonces, si $\sh{F}'$ es $f'$- plano en cada punto
de $\mathrm{X}'$ cuya proyección en $\mathrm{X}$ es $x$, $\sh{F}$ es $f$-plano en el punto $x$.
\end{proposition}

\begin{proof}
Supongamos primero que $\mathrm{A}'$ es un $\mathrm{A}$-álgebra \emph{finita} de \emph{presentación finita}; entonces, el
morfismo $\mathrm{Y}'\to\mathrm{Y}$ es propio (\textbf{II},~$6.1.11$) y de presentación finita, por lo que se satisfacen las hipótesis
de \sref{IV.11.5.3}, de donde resulta la conclusión. En el caso general, la proposición se deduce
de este caso particular, del hecho de que $\mathrm{A}'$ es el límite inductivo de sus sub-$\mathrm{A}$-álgebras \emph{finitas}
$\mathrm{A}'_\lambda$ y de los dos lemas siguientes:
\end{proof}

\begin{lemma}[11.5.5.1]
\label{IV.11.5.5.1}
Toda $\mathrm{A}$-álgebra finita $\mathrm{A}'$ es el límite inductivo de $\mathrm{A}$-álgebras $\mathrm{A}'_\lambda$
que son finitas y de presentación finita.
\end{lemma}

\begin{proof}
Se razona como en $(1.9.3.1)$. En efecto, $\mathrm{A}'=\mathrm{B}/\mathfrak{J}$, donde $\mathrm{B}$ es una $\mathrm{A}$-álgebra finita
que es libre como $\mathrm{A}$-módulo, y $\mathfrak{J}$ es un ideal de $\mathrm{B}$ $(1.4.7.1)$. Ahora bien, $\mathfrak{J}$ es el
límite inductivo de ideales $\mathfrak{J}_\lambda\subset\mathfrak{J}$ de $\mathrm{B}$ que son de tipo finito (y, por tanto, \emph{a fortiori}, son $\mathrm{A}$-módulos \emph{de tipo finito}),
por lo tanto,
por la exactitud del funtor $\varinjlim$, $\mathrm{A}'$ es el límite inductivo de las $\mathrm{A}$-álgebras $\mathrm{B}/\mathfrak{J}_\lambda$; ahora bien, $\mathrm{B}/\mathfrak{J}_\lambda$
es, por definición, un $\mathrm{A}$-módulo de presentación finita y, por tanto, también $(1.4.7)$ una $\mathrm{A}$-álgebra finita
de presentación finita.
\end{proof}

\oldpage[IV-3]{154}

Para aplicar este lema a la situación de \sref{IV.11.5.5}, observemos además
que si el homomorfismo $\mathrm{A}\to\mathrm{A}'$ es inyector, entonces
\emph{a fortiori}, también lo es $\mathrm{A}\to\mathrm{A}'_\lambda$, para todo
$\lambda$.

\begin{lemma}[11.5.5.2]
\label{IV.11.5.5.2}
Sean $\mathrm{A}$ un anillo, $\mathrm{A}'$ una $\mathrm{A}$-álgebra y $(\mathrm{A}'_\lambda)$ un sistema inductivo
de $\mathrm{A}$-álgebras tal que $\mathrm{A}'=\varinjlim\mathrm{A}'_\lambda$; pongamos $\mathrm{Y}=\operatorname{Spec}(\mathrm{A})$, $\mathrm{Y}'=\operatorname{Spec}(\mathrm{A}')$ y $\mathrm{Y}'_\lambda=\operatorname{Spec}(\mathrm{A}'_\lambda)$.
Sea $f:\mathrm{X}\to\mathrm{Y}$ un morfismo de presentación finita, y sea $\sh{F}$ un $\sh{O}_\mathrm{X}$-módulo cuasicoherente de
presentación finita; pongamos $\mathrm{X}'=\mathrm{X}\times_\mathrm{Y}\mathrm{Y}'$, $f'=f_{|\mathrm{Y}'}$, $\sh{F}'=\sh{F}\otimes_\mathrm{Y}\mathrm{Y}'$, $\mathrm{X}'_\lambda=\mathrm{X}\times_\mathrm{Y}\mathrm{Y}'_\lambda$, $f'_\lambda=f_{|\mathrm{Y}'_\lambda}$,
$\sh{F}'_\lambda=\sh{F}\otimes_\mathrm{Y}\mathrm{Y}'_\lambda$. Sea $x$ un punto de $\mathrm{X}$ tal que $\sh{F}'$ sea $f'$-plano en todo punto $x'\in\mathrm{X}'$ situado sobre $x$;
entonces existe $\lambda$ tal que $\sh{F}'_\lambda$ es $f'_\lambda$-plano en todo punto $x'\in\mathrm{X}'_\lambda$ situado sobre $x$.
\end{lemma}

\begin{proof}
Sea $\mathrm{U}'$ el conjunto de los $x'\in\mathrm{X}'$ tales que $\sh{F}'$ es $f'$-plano en el punto $x'$; se sabe por \sref{IV.11.3.1}
que $\mathrm{U}'$ es abierto en $\mathrm{X}'$, pues $f'$ es de presentación finita $(1.6.2)$; análogamente, el
conjunto $\mathrm{U}'_\lambda$ de puntos de $\mathrm{X}'_\lambda$ en los que $\sh{F}'_\lambda$ es $f'_\lambda$-plano es abierto en $\mathrm{X}'_\lambda$, y se sabe además por
\sref{IV.11.2.6} que $\mathrm{U}'$ es la unión de los $v_\lambda^{-1}(\mathrm{U}'_\lambda)$, donde $v_\lambda:\mathrm{X}'\to\mathrm{X}'_\lambda$ es la proyección canónica.
Consideremos el esquema $\mathrm{T}=\operatorname{Spec}(k(x))$; pongamos $\mathrm{T}'=\mathrm{T}\times_\mathrm{Y}\mathrm{Y}'$ y $\mathrm{T}'_\lambda=\mathrm{T}\times_\mathrm{Y}\mathrm{Y}'_\lambda$, y sean
$w_\lambda:\mathrm{T}'\to\mathrm{T}'_\lambda$, $g':\mathrm{T}'\to\mathrm{X}'$ y $g'_\lambda:\mathrm{T}'_\lambda\to\mathrm{X}'_\lambda$ las proyecciones canónicas. Pongamos $\mathrm{V}'_\lambda=g_\lambda'^{-1}(\mathrm{U}'_\lambda)$,
$\mathrm{V}'=g'^{-1}(\mathrm{U}')=\bigcup w_\lambda^{-1}(\mathrm{V}'_\lambda)$. Por hipótesis (y teniendo en cuenta (\textbf{I},~$3.6.1$)), $\mathrm{V}'=\mathrm{T}'$;
como $\mathrm{T}'$ es cuasicompacto, existe $\lambda$ tal que $w_\lambda^{-1}(\mathrm{V}'_\lambda)=\mathrm{T}'$. Se deduce entonces de $(8.3.3)$,
aplicado a los subconjuntos cerrados cuasicompactos $\mathrm{T}'_\lambda-\mathrm{V}'_\lambda$ de $\mathrm{T}'_\lambda$, que existe $\mu\geq\lambda$ tal
que $\mathrm{T}'_\mu=\mathrm{V}'_\mu$; esto significa que $\sh{F}'_\mu$ es $f'_\mu$-plano en cada punto de $\mathrm{X}'_\mu$ cuya proyección
en $\mathrm{X}$ es $x$.
\end{proof}

\subsection{Descenso de la planitud por morfismos arbitrarios: el caso de un preesquema base unirramificado}
\label{subsection:IV.11.6}

\begin{theorem}[11.6.1]
\label{IV.11.6.1}
Sea $\mathrm{A}$ un anillo local integral geométricamente unirramificado
$(\mathbf{0},~23.2.1)$, $\mathrm{Y}=\operatorname{Spec}(\mathrm{A})$, $f:\mathrm{X}\to\mathrm{Y}$ un morfismo localmente de presentación finita y $\sh{F}$ un
$\sh{O}_\mathrm{X}$-módulo cuasicoherente de presentación finita. Sea $\mathrm{A}'$ un anillo local y $u:\mathrm{A}\to\mathrm{A}'$ un
homomorfismo local inyectivo; pongamos $\mathrm{Y}'=\operatorname{Spec}(\mathrm{A}')$, $\mathrm{X}'=\mathrm{X}\times_\mathrm{Y}\mathrm{Y}'$, $f'=f_{|\mathrm{Y}'}$ y $\sh{F}'=\sh{F}\otimes_\mathrm{Y}\mathrm{Y}'$. Sea $x$
un punto de $\mathrm{X}$ cuya proyección $f(x)=y$ es el punto cerrado de $\mathrm{Y}$, y sea $x'$ un punto de $\mathrm{X}'$ cuyas proyecciones
en $\mathrm{X}$ y $\mathrm{Y}'$ son, respectivamente, $x$ y el punto cerrado $y'$ de $\mathrm{Y}'$. Entonces, si $\sh{F}'$ es $f'$-plano en el punto $x'$,
$\sh{F}$ es $f$-plano en el punto $x$.
\end{theorem}

\begin{proof}
Podemos limitarnos al caso en que $f$ es de presentación finita, puesto que la cuestión es local
en $\mathrm{X}$. Procederemos en varias etapas.

\medskip
\noindent\textbf{I)} \emph{Reducción al caso en que $\mathrm{A}$ y $\mathrm{A}'$ son anillos locales íntegramente cerrados.}

\medskip
Como $u$ es inyectivo y $\mathrm{A}$ es integral, existe un ideal primo $\mathfrak{p}'$ de $\mathrm{A}'$ tal que
$u^{-1}(\mathfrak{p}')=(0)$ $(\mathbf{0}_\mathrm{I},~1.5.8)$; por tanto, el homomorfismo compuesto $\mathrm{A}\to\mathrm{A}'\to\mathrm{A}''=\mathrm{A}'/\mathfrak{p}'$ es inyectivo
y local, y, si $\mathrm{Y}''=\operatorname{Spec}(\mathrm{A}'')$, $\mathrm{X}''=\mathrm{X}'\otimes_{\mathrm{A}'}\mathrm{A}''=\mathrm{X}\times_\mathrm{Y}\mathrm{Y}''$ y $\sh{F}''=\sh{F}'\otimes_{\mathrm{A}'}\mathrm{A}''$, entonces $\sh{F}''$ es $\mathrm{Y}''$-plano
en los puntos de $\mathrm{X}''$ situados sobre $x'$ $(2.1.4)$; sustituyendo, si es necesario, $\mathrm{A}'$ por $\mathrm{A}''$ y teniendo
en cuenta (\textbf{I},~$3.4.7$), podemos suponer en primer lugar que $\mathrm{A}'$ es integral. Si $\mathrm{K}'$ es el cuerpo de
fracciones de $\mathrm{A}'$, existe un anillo de valoración $\mathrm{B}'$ de $\mathrm{K}'$ que domina a $\mathrm{A}'$; como el homomorfismo compuesto
$\mathrm{A}\xrightarrow{u}\mathrm{A}'\to\mathrm{B}'$ es inyectivo y local, el mismo razonamiento que en
\oldpage[IV-3]{155}
el argumento precedente permite sustituir $\mathrm{A}'$ por $\mathrm{B}'$; podemos, pues, suponer que el anillo local $\mathrm{A}'$ está íntegramente cerrado,
siendo $\mathrm{A}$ un subanillo local de $\mathrm{A}'$ dominado por $\mathrm{A}'$. Sea $\mathrm{A}_1$ el cierre integral de $\mathrm{A}$; es claro que
$\mathrm{A}\subset\mathrm{A}_1\subset\mathrm{A}'$, y, por hipótesis, $\mathrm{A}_1$ es un anillo \emph{local};
si $\mathfrak{m}$, $\mathfrak{m}_1$ y $\mathfrak{m}'$ son los ideales máximos de $\mathrm{A}$, $\mathrm{A}_1$ y $\mathrm{A}'$, se tiene $\mathfrak{m}\subset\mathfrak{m}_1\subset\mathfrak{m}'$; en efecto, $\mathfrak{m}_1$ es
el \emph{único} ideal primo de $\mathrm{A}_1$ situado sobre $\mathfrak{m}$, puesto que $\mathrm{A}_1$ es un anillo local (Bourbaki,
\emph{Alg.\ comm.}, cap.~V, \S~2, n\textsuperscript{o}~1, prop.~1); como $\mathfrak{m}'\cap\mathrm{A}=\mathfrak{m}$, se tiene $\mathfrak{m}'\cap\mathrm{A}_1\cap\mathrm{A}=\mathfrak{m}$,
y, por consiguiente, $\mathfrak{m}'\cap\mathrm{A}_1=\mathfrak{m}_1$. Pongamos $\mathrm{Y}_1=\operatorname{Spec}(\mathrm{A}_1)$, $\mathrm{X}_1=\mathrm{X}\times_\mathrm{Y}\mathrm{Y}_1$, $f_1=f_{(\mathrm{Y}_1)}$ y $\sh{F}_1=\sh{F}\otimes_\mathrm{Y}\mathrm{Y}_1$, y
sea $x_1$ la proyección de $x'$ en $\mathrm{X}_1$; denotemos además por $y_1$ el único punto cerrado
de $\mathrm{Y}_1$, de modo que $y_1=f_1(x_1)$. Por hipótesis, el morfismo $\operatorname{Spec}(k(y_1))\to\operatorname{Spec}(k(y))$
es \emph{radicial}, de donde se deduce, por la transitividad de las fibras (\textbf{I},~$3.6.4$) y (\textbf{I},~$3.5.7$), que
el morfismo $f_1^{-1}(y_1)\to f^{-1}(y)$ es radicial y, en particular, que $x_1$ es el \emph{único} punto de $\mathrm{X}_1$
cuyas proyecciones en $\mathrm{X}$ y $\mathrm{Y}_1$ son, respectivamente, $x$ y $y_1$; además, como se ha visto, $y_1$
es el \emph{único} punto de $\mathrm{Y}_1$ cuya proyección en $\mathrm{Y}$ es $y$, por lo que $x_1$ es el \emph{único} punto de $\mathrm{X}_1$
cuya proyección en $\mathrm{X}$ es $x$. Si demostramos que $\sh{F}_1$ es $f_1$-plano en el punto $x_1$, podremos
aplicar \sref{IV.11.5.5}, de donde se obtiene la conclusión. Quedamos así reducidos al caso en que $\mathrm{A}$
es también \emph{íntegramente cerrado}.

\medskip
\noindent\textbf{II)} \emph{Reducción al caso en que $\mathrm{A}$ y $\mathrm{A}'$ son anillos locales íntegramente cerrados de $\mathbf{Z}$-álgebras de tipo
finito.}

\medskip
Podemos considerar $\mathrm{A}'$ como el límite inductivo filtrante de sus sub-$\mathbf{Z}$-álgebras de tipo finito
$\mathrm{B}'_\lambda$ que son dominios de integridad; además, como $\mathrm{A}'$ está íntegramente cerrado y el cierre
integral de una $\mathbf{Z}$-álgebra de tipo finito es también una $\mathbf{Z}$-álgebra de tipo finito $(7.8.3)$, se
ve que $\mathrm{A}'$ es el límite inductivo de sus sub-$\mathbf{Z}$-álgebras de tipo finito $\mathrm{B}'_\lambda$ \emph{íntegramente cerradas};
si $\mathfrak{m}'_\lambda=\mathfrak{m}'\cap\mathrm{B}'_\lambda$, $\mathrm{A}'$ es también el límite inductivo de los subanillos locales $(\mathrm{B}'_\lambda)_{\mathfrak{m}'_\lambda}=\mathrm{A}'_\lambda$
dominados por $\mathrm{A}'$ $(5.13.3)$. Para todo $\lambda$, $\mathrm{B}'_\lambda\cap\mathrm{A}$ es asimismo el límite inductivo de sus
sub-$\mathbf{Z}$-álgebras de tipo finito $\mathrm{B}_{\alpha\lambda}$, y, por tanto, $\mathrm{A}=\varinjlim_{\alpha,\lambda}\mathrm{B}_{\alpha\lambda}$; como antes, podemos sustituir $\mathrm{B}_{\alpha\lambda}$
en esta fórmula por su cierre integral (que, por hipótesis, está contenido en $\mathrm{B}'_\lambda\cap\mathrm{A}$), y después por el anillo local
$\mathrm{A}_{\alpha\lambda}=(\mathrm{B}_{\alpha\lambda})_{\mathfrak{m}_{\alpha\lambda}}$, donde
$\mathfrak{m}_{\alpha\lambda}=\mathfrak{m}\cap\mathrm{B}_{\alpha\lambda}=\mathfrak{m}'_\lambda\cap\mathrm{B}_{\alpha\lambda}$, de modo que $\mathrm{A}_{\alpha\lambda}$ está dominado por $\mathrm{A}'_\lambda$ y por $\mathrm{A}$.
Pongamos $\mathrm{Y}_{\alpha\lambda}=\operatorname{Spec}(\mathrm{A}_{\alpha\lambda})$; de $(8.9.1)$ se deduce que existen un par suficientemente grande
$(\alpha,\lambda)$, un morfismo $f_{\alpha\lambda}:\mathrm{X}_{\alpha\lambda}\to\mathrm{Y}_{\alpha\lambda}$ de tipo finito y un
$\sh{O}_{\mathrm{X}_{\alpha\lambda}}$-módulo coherente $\sh{F}_{\alpha\lambda}$ tales que
$\mathrm{X}=\mathrm{X}_{\alpha\lambda}\times_{\mathrm{Y}_{\alpha\lambda}}\mathrm{Y}$,
$f=f_{\alpha\lambda}\times 1_\mathrm{Y}$ y $\sh{F}=\sh{F}_{\alpha\lambda}\otimes_{\mathrm{Y}_{\alpha\lambda}}\mathrm{Y}$; si $x_{\alpha\lambda}$
es la proyección de $x$ en $\mathrm{X}_{\alpha\lambda}$, bastará demostrar que $\sh{F}_{\alpha\lambda}$ es $f_{\alpha\lambda}$-plano en el punto $x_{\alpha\lambda}$.
Como $\mathrm{A}'$ es el límite inductivo de los $\mathrm{A}'_\mu$ para $\mu\geq\lambda$, $\mathrm{X}'$ es el límite proyectivo de los
$\mathrm{X}_{\alpha\lambda}\times_{\mathrm{Y}_{\alpha\lambda}}\mathrm{Y}'_\mu=\mathrm{X}'_\mu$, y se tiene asimismo
$\sh{F}'=\sh{F}'_{\alpha\mu}\otimes_{\mathrm{Y}'_\mu}\mathrm{Y}'$, donde
$\sh{F}'_{\alpha\mu}=\sh{F}_{\alpha\lambda}\otimes_{\mathrm{Y}_{\alpha\lambda}}\mathrm{Y}'_\mu$. Aplicando
\sref{IV.11.2.6}, se ve que podemos tomar $\mu$ suficientemente grande para que $\sh{F}'_{\alpha\mu}$ sea $\mathrm{Y}'_\mu$-plano en el
punto $x'_\mu$, proyección de $x'$ en $\mathrm{X}'_\mu$; además, por la construcción de los $\mathrm{A}'_\mu$, la proyección
de $x'_\mu$ en $\mathrm{Y}'_\mu$ es el punto cerrado de $\mathrm{Y}'_\mu$.

\medskip
\noindent\textbf{III)} \emph{Reducción al caso en que el cuerpo residual $k'$ de $\mathrm{A}'$ es una extensión finita y radicial del
cuerpo residual $k$ de $\mathrm{A}$.}

\medskip
Podemos comenzar repitiendo el razonamiento de la parte I) de la
demostración de \sref{IV.11.4.11}, teniendo en cuenta que $\mathbf{Z}$ es un
\oldpage[IV-3]{156}
anillo de Jacobson; quedamos así reducidos al caso en que $k'$ es una extensión finita de $k$, como supondremos en lo sucesivo.
Sea $k''$ la mayor extensión separable de $k$ contenida en $k'$, y sea $k_1$ una extensión de Galois finita de $k$ que contiene a $k''$, de modo que $k''\otimes_k k_1$ es un producto directo de cuerpos isomorfos a $k_1$.
Como $k'$ es una extensión radicial de $k''$, $k'\otimes_k k_1$ es, por tanto, un producto directo de extensiones radiciales de $k_1$.

Existe un anillo local $\mathrm{A}_1$ que es una $\mathrm{A}$-álgebra y un $\mathrm{A}$-módulo libre de tipo finito, tal que $\mathrm{A}_1/\mathfrak{m}\mathrm{A}_1$ es $k$-isomorfo a $k_1$ $(\mathbf{0}_{\mathrm{III}},~10.3.1.2)$.
Más precisamente, escribamos $k_1=k[\mathrm{T}]/(r)$, donde $r$ es un polinomio mónico irreducible separable de grado $n$, y sea $\mathrm{R}\in\mathrm{A}[\mathrm{T}]$ un polinomio mónico cuya imagen canónica es $r$; podemos tomar $\mathrm{A}_1=\mathrm{A}[\mathrm{T}]/(\mathrm{R})$.
Si $\mathrm{K}$ es el cuerpo de fracciones de $\mathrm{A}$, es claro que $\mathrm{R}$ es un polinomio irreducible separable de $\mathrm{K}[\mathrm{T}]$; se deduce, en primer lugar, que $\mathrm{A}_1$ es un anillo integral cuyo cuerpo de fracciones $\mathrm{K}_1=\mathrm{K}[\mathrm{T}]/(\mathrm{R})$ es una extensión separable de $\mathrm{K}$.
Además, si $t$ es la imagen canónica de $\mathrm{T}$ en $\mathrm{A}_1$, los $t^j$ $(0\leq j<n)$ forman una base del $\mathrm{A}$-módulo $\mathrm{A}_1$, y sus imágenes en $k_1$ forman una base sobre $k$; por tanto, $d=\det\bigl(\operatorname{Tr}_{\mathrm{A}_1/\mathrm{A}}(t^{i+j})\bigr)$ es un elemento de $\mathrm{A}$ cuya clase en $k$ es no nula y que, por consiguiente, es invertible.
El mismo razonamiento que en \sref{IV.6.12.4.1},~I) demuestra entonces que el morfismo $\operatorname{Spec}(\mathrm{A}_1)\to\operatorname{Spec}(\mathrm{A})$ es plano y tiene fibras regulares; de \sref{IV.6.5.4},~(ii), se deduce que $\mathrm{A}_1$ es \emph{íntegramente cerrado}.

Consideremos ahora el anillo $\mathrm{A}'_1=\mathrm{A}'\otimes_\mathrm{A}\mathrm{A}_1$.
Es un $\mathrm{A}'$-módulo libre de tipo finito y, por tanto, un anillo semilocal (Bourbaki, \emph{Alg.\ comm.}, cap.~IV, \S~2, n\textsuperscript{o}~5, cor.~3 de prop.~9); además, todos los ideales máximos de esta $\mathrm{A}'$-álgebra finita están sobre el ideal máximo $\mathfrak{m}'$ de $\mathrm{A}'$ y, \emph{a fortiori}, contienen a $\mathfrak{m}\mathrm{A}'_1$.
Pero
\[
\mathrm{A}'_1/\mathfrak{m}\mathrm{A}'_1=(\mathrm{A}'/\mathfrak{m}\mathrm{A}')\otimes_k k_1,
\]
y, como $k_1$ es una extensión finita separable de $k$, el radical de $\mathrm{A}'_1/\mathfrak{m}\mathrm{A}'_1$ es igual a $(\mathfrak{m}'/\mathfrak{m}\mathrm{A}')\otimes_k k_1$ (Bourbaki, \emph{Alg.}, cap.~VIII, \S~7, n\textsuperscript{o}~2, cor.~2 de prop.~3).
Si $\mathfrak{n}_i$ $(1\leq i\leq r)$ son los ideales máximos de $\mathrm{A}'_1$, los cuerpos $\mathrm{A}'_1/\mathfrak{n}_i$ son, por tanto, los cuerpos que aparecen como factores del álgebra $k'\otimes_k k_1$; en otras palabras, son \emph{extensiones radiciales finitas de $k_1$}.

Como $\mathrm{A}\to\mathrm{A}'$ es inyectivo y $\mathrm{A}_1$ es un $\mathrm{A}$-módulo plano, el homomorfismo $\mathrm{A}_1\to\mathrm{A}'_1$ es inyectivo.
El homomorfismo canónico
\[
\mathrm{A}'_1\longrightarrow\prod_{i=1}^{r}(\mathrm{A}'_1)_{\mathfrak{n}_i}
\]
es también inyectivo (Bourbaki, \emph{Alg.\ comm.}, cap.~II, \S~3, n\textsuperscript{o}~3, th.~1), y, por consiguiente, también lo es el homomorfismo compuesto de $\mathrm{A}_1$ en este producto.
Ahora bien, $\mathrm{A}_1$ es integral y solo hay un número finito de núcleos de los homomorfismos $\mathrm{A}_1\to(\mathrm{A}'_1)_{\mathfrak{n}_i}$; como su intersección es cero, uno de ellos es ya cero.
En otras palabras, para algún $i$ existe un anillo $\mathrm{B}_1=(\mathrm{A}'_1)_{\mathfrak{n}_i}$ tal que el homomorfismo $\mathrm{A}_1\to\mathrm{B}_1$ es \emph{inyectivo} y \emph{local}.

Pongamos $\mathrm{Y}'_1=\operatorname{Spec}(\mathrm{B}_1)$, $\mathrm{X}'_1=\mathrm{X}'\times_{\mathrm{Y}'}\mathrm{Y}'_1$ y $\sh{F}'_1=\sh{F}'\otimes_{\mathrm{Y}'}\mathrm{Y}'_1$.
El módulo $\sh{F}'_1$ es $\mathrm{Y}'_1$-plano en todos los puntos de $\mathrm{X}'_1$ situados sobre $x'$; además, el ideal máximo de $\mathrm{B}_1$ es el único situado sobre $\mathfrak{m}'$, por lo que todos estos puntos se proyectan en el punto cerrado $y'_1$ de $\mathrm{Y}'_1$.
Sea $x'_1$ uno de estos puntos.
Por otra parte, pongamos $\mathrm{Y}_1=\operatorname{Spec}(\mathrm{A}_1)$, $\mathrm{X}_1=\mathrm{X}\times_\mathrm{Y}\mathrm{Y}_1$ y $\sh{F}_1=\sh{F}\otimes_\mathrm{Y}\mathrm{Y}_1$; si $x_1$ es la proyección de $x'_1$ en $\mathrm{X}_1$, sus proyecciones en $\mathrm{X}$ y $\mathrm{Y}_1$ son, respectivamente, $x$ y el punto cerrado $y_1$.
Si demostramos que $(\sh{F}_1)_{x_1}$ es un $\sh{O}_{y_1}$-módulo plano, se deduce que $\sh{F}_x$ es un $\sh{O}_y$-módulo plano: en efecto, $\sh{O}_{y_1}$ es un $\sh{O}_y$-módulo plano, de modo que $(\sh{F}_1)_{x_1}$ es un $\sh{O}_y$-módulo plano $(\mathbf{0}_\mathrm{I},~6.2.1)$.

\oldpage[IV-3]{157}
Pero $(\sh{F}_1)_{x_1}=\sh{F}_x\otimes_{\sh{O}_x}\sh{O}_{x_1}$, y $\sh{O}_{x_1}$ es un $\sh{O}_x$-módulo fielmente plano; por \sref{IV.2.2.11},~(iii), $\sh{F}_x$ es, por tanto, un $\sh{O}_y$-módulo plano.
Como $\mathrm{X}'_1=\mathrm{X}_1\times_{\mathrm{Y}_1}\mathrm{Y}'_1$ y $\sh{F}'_1=\sh{F}_1\otimes_{\mathrm{Y}_1}\mathrm{Y}'_1$, hemos quedado efectivamente reducidos a la situación de \sref{IV.11.6.1}, sustituyendo $\mathrm{A}$ y $\mathrm{A}'$, respectivamente, por $\mathrm{A}_1$ y $\mathrm{B}_1$.

\medskip
\noindent\textbf{IV)} \emph{Fin de la demostración.} --- Finalmente, quedamos reducidos a demostrar \sref{IV.11.6.1} bajo las siguientes hipótesis adicionales:
\begin{enumerate}[label=(\roman*)]
  \item $\mathrm{A}$ y $\mathrm{A}'$ son anillos locales de $\mathbf{Z}$-álgebras de tipo finito (y, por tanto, anillos \emph{excelentes} \sref{IV.7.8.3});
  \item $\mathrm{A}$ está íntegramente cerrado;
  \item el cuerpo residual $k'$ de $\mathrm{A}'$ es una extensión radicial finita del cuerpo residual $k$ de $\mathrm{A}$.
\end{enumerate}
Se sabe entonces (\sref{IV.7.8.3} y \sref{IV.2.3.8}) que, bajo las condiciones~(i) y~(ii), la topología $\mathfrak{m}$-ádica de $\mathrm{A}$ es inducida por la topología $\mathfrak{m}'$-ádica de $\mathrm{A}'$.
Por tanto, el homomorfismo entre los completados $\widehat{u}:\widehat{\mathrm{A}}\to\widehat{\mathrm{A}'}$ es inyectivo.
Por otra parte, como el morfismo $\operatorname{Spec}(k')\to\operatorname{Spec}(k)$ es radicial, lo mismo ocurre con $f'^{-1}(y')\to f^{-1}(y)$ \sref[I]{I.3.5.7}; existe, pues, un único punto $x'$ de $\mathrm{X}'$ cuyas proyecciones en $\mathrm{X}$ y $\mathrm{Y}'$ son, respectivamente, $x$ y $y'$.
Podemos, por tanto, aplicar \sref{IV.11.5.2}.
\end{proof}

\begin{corollary}[11.6.2]
\label{IV.11.6.2}
Sea $\mathrm{A}$ un anillo local integral unirramificado, $\mathrm{Y}=\operatorname{Spec}(\mathrm{A})$, $f:\mathrm{X}\to\mathrm{Y}$ un morfismo localmente de presentación finita y $\sh{F}$ un $\sh{O}_\mathrm{X}$-módulo cuasicoherente de presentación finita.
Sea $\mathrm{A}'$ un anillo local y $\mathrm{A}\to\mathrm{A}'$ un homomorfismo local inyectivo; pongamos $\mathrm{Y}'=\operatorname{Spec}(\mathrm{A}')$, $\mathrm{X}'=\mathrm{X}\times_\mathrm{Y}\mathrm{Y}'$, $f'=f_{(\mathrm{Y}')}$ y $\sh{F}'=\sh{F}\otimes_\mathrm{Y}\mathrm{Y}'$.
Sea $x$ un punto de $\mathrm{X}$ cuya proyección $f(x)=y$ es el punto cerrado de $\mathrm{Y}$, y supongamos que $\sh{F}'$ es $f'$-plano en todos los puntos $x'$ de $\mathrm{X}'$ cuyas proyecciones en $\mathrm{X}$ y $\mathrm{Y}'$ son, respectivamente, $x$ y el punto cerrado $y'$ de $\mathrm{Y}'$.
Entonces $\sh{F}$ es $f$-plano en $x$.
\end{corollary}

\begin{proof}
Podemos repetir la parte~I) de la demostración de \sref{IV.11.6.1}, que muestra (con la misma notación) que, si $\sh{F}_1$ es $f_1$-plano en \emph{todos} los puntos $x_1$ de $\mathrm{X}_1$ cuyas proyecciones en $\mathrm{X}$ y $\mathrm{Y}_1$ son, respectivamente, $x$ y $y_1$, entonces $\sh{F}$ es $f$-plano en $x$.
Quedamos así reducidos al caso en que $\mathrm{A}$ está íntegramente cerrado y, por tanto, es geométricamente unirramificado; la conclusión se sigue de \sref{IV.11.6.1}.
\end{proof}

\subsection{Contraejemplos}
\label{subsection:IV.11.7}

\begin{env}[11.7.1]
\label{IV.11.7.1}
Consideremos primero el caso en que $\mathrm{A}$ es un anillo local \emph{artiniano} y en que se satisfacen las hipótesis de \sref{IV.11.4.11}, \emph{salvo} la condición~(ii) relativa al cuerpo residual $k$ de $\mathrm{A}$. Veremos que, en tal caso, puede fallar la conclusión de \sref{IV.11.4.11}.
Sea $k$ un cuerpo que admite una extensión de Galois $k'$ de grado $[k':k]>1$, y sea $\Gamma$ su grupo de Galois.
Sea $\mathrm{A}$ una $k$-álgebra con una base de tres elementos $1,a,b$ y tabla de multiplicación $a^2=ab=ba=b^2=0$, de modo que $\mathrm{A}$ es un anillo local artiniano cuyo ideal máximo $\mathfrak{m}=ka+kb$ tiene cuadrado cero.
Sea $\mathrm{A}'=\mathrm{A}\otimes_k k'$, que es una $k'$-álgebra con base $1,a,b$ y un anillo local artiniano cuyo ideal máximo de cuadrado cero es $\mathfrak{m}'=k'a+k'b=\mathfrak{m}\mathrm{A}'$; el anillo $\mathrm{A}$ se identifica canónicamente con un subanillo de $\mathrm{A}'$.
Sea $\mathfrak{J}$ el subespacio vectorial sobre $k'$ de $\mathfrak{m}'$ generado por $a+\gamma b$, donde $\gamma\in k'$ no pertenece a $k$; es claro que $\mathfrak{J}$ es un ideal de $\mathrm{A}'$.
Pongamos $\mathrm{B}=\mathrm{A}'/\mathfrak{J}$; este es un anillo artiniano que no es plano como $\mathrm{A}$-módulo. De lo contrario (Bourbaki, \emph{Alg.\ comm.}, cap.~II, \S~3, n\textsuperscript{o}~2, cor.~2 de prop.~5), $\mathrm{B}$ sería un $\mathrm{A}$-módulo libre. Como $\mathrm{A}'$ es asimismo un $\mathrm{A}$-módulo libre y el homomorfismo canónico $\mathrm{A}'/\mathfrak{m}\mathrm{A}'\to\mathrm{B}/\mathfrak{m}\mathrm{B}$ es biyectivo, el homomorfismo canónico $\mathrm{A}'\to\mathrm{B}=\mathrm{A}'/\mathfrak{J}$ también sería
\oldpage[IV-3]{158}
biyectivo (loc.~cit., n\textsuperscript{o}~2, cor.~de prop.~6), lo cual es absurdo. En otras palabras, si $\mathrm{Y}=\Spec(\mathrm{A})$ y $\mathrm{X}=\Spec(\mathrm{B})$, entonces $\sh{O}_\mathrm{X}$ no es $\mathrm{Y}$-plano en el único punto $x$ de $\mathrm{X}$.
Pongamos ahora $\mathrm{A}_1=\mathrm{B}$ y $\mathrm{Y}_1=\Spec(\mathrm{A}_1)$, y consideremos el homomorfismo canónico $\mathrm{A}\to\mathrm{A}_1$, que es local e inyectivo porque $\mathfrak{J}\cap\mathrm{A}=0$. Si $\mathrm{B}_1=\mathrm{B}\otimes_\mathrm{A}\mathrm{A}_1=\mathrm{B}\otimes_\mathrm{A}\mathrm{B}$, veremos que existe un punto de $\mathrm{X}_1=\mathrm{X}\times_\mathrm{Y}\mathrm{Y}_1=\Spec(\mathrm{B}_1)$ en el que $\sh{O}_{\mathrm{X}_1}$ es $\mathrm{Y}_1$-plano.
Para ello, observemos que
\[
\mathrm{B}\otimes_\mathrm{A}\mathrm{B}
=\frac{\mathrm{A}'\otimes_\mathrm{A}\mathrm{A}'}
{\operatorname{Im}(\mathfrak{J}\otimes_\mathrm{A}\mathrm{A}')+\operatorname{Im}(\mathrm{A}'\otimes_\mathrm{A}\mathfrak{J})}.
\]
La estructura de $\mathrm{C}'=\mathrm{A}'\otimes_\mathrm{A}\mathrm{A}'$ se obtiene fácilmente. Consideremos la $\mathrm{A}$-álgebra producto $\mathrm{C}=\prod_{\sigma\in\Gamma}\mathrm{A}'_\sigma$, donde todos los $\mathrm{A}'_\sigma$ son iguales a $\mathrm{A}'$, y la aplicación canónica $\varphi:\mathrm{A}'\otimes_\mathrm{A}\mathrm{A}'\to\mathrm{C}$ definida por $\varphi(x\otimes y)=(\sigma(x)y)_{\sigma\in\Gamma}$, donde $\Gamma$ actúa canónicamente sobre $\mathrm{A}'$. Al pasar a los cocientes se obtiene un homomorfismo $\mathrm{C}'/\mathfrak{m}\mathrm{C}'\to\mathrm{C}/\mathfrak{m}\mathrm{C}$ que no es sino el homomorfismo canónico $k'\otimes_k k'\to\prod_\sigma k'_\sigma$, con $k'_\sigma=k'$ para todo $\sigma\in\Gamma$. Este último homomorfismo es biyectivo (Bourbaki, \emph{Alg.}, cap.~VIII, \S~8, prop.~4), y, por tanto, también lo es $\varphi$, pues $\mathrm{C}'$ y $\mathrm{C}$ son $\mathrm{A}$-módulos libres (Bourbaki, \emph{Alg.\ comm.}, cap.~II, \S~3, n\textsuperscript{o}~2, cor.~de prop.~6).
Se deduce que $\mathrm{B}\otimes_\mathrm{A}\mathrm{B}$ es una $\mathrm{A}$-álgebra semilocal que es producto directo de las $\mathrm{A}$-álgebras locales $\mathrm{A}'/(\sigma(\mathfrak{J})+\mathfrak{J})$. El factor correspondiente al elemento neutro de $\Gamma$ es isomorfo a $\mathrm{A}'/\mathfrak{J}$ y es, por tanto, un $\mathrm{A}_1$-módulo plano. Sin embargo, como $\gamma\notin k$, existe al menos un $\sigma\in\Gamma$ tal que $\sigma(\mathfrak{J})\ne\mathfrak{J}$; entonces $\sigma(\mathfrak{J})+\mathfrak{J}=\mathfrak{m}'$, y $\mathrm{A}'/\mathfrak{m}'$ no es un $(\mathrm{A}'/\mathfrak{J})$-módulo plano, pues es el cociente de $\mathrm{A}'/\mathfrak{J}$ por un ideal no nulo.
\end{env}

\begin{env}[11.7.2]
\label{IV.11.7.2}
Mostraremos ahora que el resultado de \sref{IV.11.5.4} deja de ser válido cuando ya no se supone que $f$ es propio (y, \emph{a fortiori}, \sref{IV.11.5.3} deja de ser exacto cuando ya no se supone que $g$ es propio).
Sea $k$ un cuerpo, sea $\mathrm{A}_0=k[\mathrm{S},\mathrm{T}]$ y sea $\mathrm{A}=\mathrm{A}_0/\mathrm{A}_0\mathrm{S}\mathrm{T}$. Así, $\mathrm{Y}=\Spec(\mathrm{A})$ es la curva reducible formada por los dos ``ejes coordenados'' del plano afín $\mathrm{V}_k^2=\Spec(\mathrm{A}_0)$.
El anillo $\mathrm{A}$ tiene dos ideales primos mínimos $\mathfrak{p}_1=\mathrm{A}_0\mathrm{S}/\mathrm{A}_0\mathrm{S}\mathrm{T}$ y $\mathfrak{p}_2=\mathrm{A}_0\mathrm{T}/\mathrm{A}_0\mathrm{S}\mathrm{T}$. Como $\mathrm{A}$ es reducido, se inyecta canónicamente en $\mathrm{B}=\mathrm{B}_1\oplus\mathrm{B}_2$, donde $\mathrm{B}_1=\mathrm{A}/\mathfrak{p}_1$ y $\mathrm{B}_2=\mathrm{A}/\mathfrak{p}_2$. Además, $\mathrm{B}_1$ se identifica con $k[\mathrm{T}]$ y $\mathrm{B}_2$ con $k[\mathrm{S}]$; estos son anillos integrales íntegramente cerrados y, por consiguiente, $\mathrm{Z}=\Spec(\mathrm{B})$ es precisamente la normalización de $\mathrm{Y}$ con respecto a $\mathrm{R}(\mathrm{Y})$ \sref[II]{II.6.3.8}, suma de los dos esquemas $\mathrm{Z}_1=\Spec(\mathrm{B}_1)$ y $\mathrm{Z}_2=\Spec(\mathrm{B}_2)$.
Denotemos por $y$ el ``punto doble'' de $\mathrm{Y}$ correspondiente al ideal máximo $\mathfrak{p}_1+\mathfrak{p}_2=\mathfrak{m}$ de $\mathrm{A}$, y por $z_1,z_2$ los puntos de $\mathrm{Z}$ situados sobre $y$, correspondientes respectivamente a los ideales máximos $\mathfrak{m}_1=(\mathrm{T})$ y $\mathfrak{m}_2=(\mathrm{S})$ de $\mathrm{B}_1$ y $\mathrm{B}_2$.
Sea $\mathrm{X}$ el subpreesquema de $\mathrm{Z}$ inducido en el complementario de $z_2$. Así, $\mathrm{X}=\Spec(\mathrm{B}_1\oplus(\mathrm{B}_2)_\mathrm{S})$. Es inmediato que el homomorfismo $\mathrm{A}\to\mathrm{A}'=\mathrm{B}_1\oplus(\mathrm{B}_2)_\mathrm{S}$ es inyectivo, pero el morfismo correspondiente $f:\mathrm{X}\to\mathrm{Y}$ no es cerrado, pues $f(\mathrm{Z}_2-\{z_2\})$ no es cerrado en $\mathrm{Y}$ aunque $\mathrm{Z}_2-\{z_2\}$ sea cerrado en $\mathrm{X}$; \emph{a fortiori}, $f$ no es propio.
Veamos ahora que $f$ no es plano en $z_1$. Basta demostrar \sref[0]{0.6.6.2} que $\sh{O}_{z_1}$ no es un $\sh{O}_y$-módulo fielmente plano, y para ello basta observar que el homomorfismo canónico $\sh{O}_y\to\sh{O}_{z_1}$ no es inyectivo. Esto es inmediato, pues $\sh{O}_{z_1}$ es un anillo integral, mientras que $\sh{O}_y$ tiene divisores de cero.
Sin embargo, la primera proyección $p:\mathrm{X}\times_\mathrm{Y}\mathrm{X}\to\mathrm{X}$ es un isomorfismo. En efecto,
\[
\mathrm{B}_1\otimes_\mathrm{A}\mathrm{B}_1
=(\mathrm{A}/\mathfrak{p}_1)\otimes_\mathrm{A}(\mathrm{A}/\mathfrak{p}_1)
=\mathrm{A}/\mathfrak{p}_1,
\]
mientras que $(\mathrm{B}_2)_\mathrm{S}\otimes_\mathrm{A}(\mathrm{B}_2)_\mathrm{S}=(\mathrm{B}_2\otimes_\mathrm{A}\mathrm{B}_2)_\mathrm{S}=(\mathrm{B}_2)_\mathrm{S}$ por la misma razón; finalmente, $\mathrm{B}_1\otimes_\mathrm{A}(\mathrm{B}_2)_\mathrm{S}=0$, pues la imagen canónica de $\mathrm{S}$ en $\mathrm{B}_1$ es cero.
\end{env}

\begin{env}[11.7.3]
\label{IV.11.7.3}
El ejemplo precedente puede generalizarse. Sobre un cuerpo $k$, consideremos una curva algebraica reducida $\mathrm{Y}$ que admite un único ``punto doble ordinario'' $y$ (noción que se definirá más adelante en general), junto con su normalización $\mathrm{Z}$, de modo que $g:\mathrm{Z}\to\mathrm{Y}$ sea finito, la restricción de $g$ a $\mathrm{Z}-g^{-1}(y)$ sea un isomorfismo sobre $\mathrm{Y}-\{y\}$ y $g^{-1}(y)$ conste de dos puntos ``simples'' $z_1,z_2$. Además, el preesquema $g^{-1}(y)$ es la suma de los dos preesquemas $\Spec(k(z_1))$ y $\Spec(k(z_2))$, cada uno de ellos canónicamente isomorfo a $\Spec(k(y))$.
Sea $\mathrm{X}$ el subpreesquema de $\mathrm{Z}$ inducido en el abierto $\mathrm{Z}-\{z_2\}$. El morfismo $f:\mathrm{X}\to\mathrm{Y}$ obtenido restringiendo $g$ no es propio; de lo contrario, \sref[II]{II.5.4.3} implicaría lo mismo para la inyección canónica $j:\mathrm{X}\to\mathrm{Z}$, que no es cerrada.
El morfismo $f$ es radicial: para todo $y'\in\mathrm{Y}$, la fibra $f^{-1}(y')$ consta de un único punto $x'$, y $k(y')\to k(x')$ es un isomorfismo. Se deduce primero que la diagonal $\Delta_f:\mathrm{X}\to\mathrm{X}\times_\mathrm{Y}\mathrm{X}$ es biyectiva \sref{IV.1.7.7.1}. Por otra parte, como $f$ es no ramificado \sref{IV.17.4.2}[(d')], $\Delta_f$ es una inmersión abierta \sref{IV.17.4.2}[(b')]. Por consiguiente, $\Delta_f$ es un isomorfismo y la primera proyección $p:\mathrm{X}\times_\mathrm{Y}\mathrm{X}\to\mathrm{X}$ es su inversa.
Sin embargo, $f$ no es plano en $z_1$. De lo contrario, $\sh{O}_{z_1}$ sería un $\sh{O}_y$-módulo fielmente plano \sref[0]{0.6.6.2}; como $\sh{O}_y$ contiene dos ideales primos mínimos distintos $\mathfrak{p}_1,\mathfrak{p}_2$, correspondientes a las dos ramas de $\mathrm{Y}$ en $y$, existirían entonces dos ideales primos distintos de $\sh{O}_{z_1}$ cuyas imágenes inversas por $u:\sh{O}_y\to\sh{O}_{z_1}$ serían $\mathfrak{p}_1$ y $\mathfrak{p}_2$ \sref[0]{0.6.5.1}. Esto es absurdo, pues $\sh{O}_{z_1}$ solo tiene dos ideales primos distintos, $0$ y su ideal máximo $\mathfrak{m}_1$, mientras que $u^{-1}(\mathfrak{m}_1)$ es el ideal máximo $\mathfrak{m}$ de $\sh{O}_y$.
\end{env}

\begin{env}[11.7.4]
\label{IV.11.7.4}
Obsérvese que, en el ejemplo precedente, el homomorfismo $u$ es \emph{inyectivo} cuando $\mathrm{Y}$ es \emph{irreducible} (se puede tomar, por ejemplo, $\mathrm{Y}=\operatorname{Spec}(k[\mathrm{S},\mathrm{T}]/(\mathrm{S}(\mathrm{S}^2+\mathrm{T}^2)-(\mathrm{S}^2-\mathrm{T}^2)))$, una ``cúbica con un punto doble''). En efecto, tras sustituir $\mathrm{Y}$ por un entorno afín de $y$, podemos suponer que $\mathrm{Y}=\operatorname{Spec}(\mathrm{A})$, donde $\mathrm{A}$ es integral, y, por tanto, que $\mathrm{Z}=\operatorname{Spec}(\mathrm{B})$, donde $\mathrm{B}$ es el cierre integral de $\mathrm{A}$. Como $\mathrm{B}$, $\sh{O}_y$ y $\sh{O}_{z_1}$ se identifican entonces con subanillos del
\oldpage[IV-3]{159}
cuerpo de fracciones de $\mathrm{A}$, $u$ es evidentemente inyectivo.
Por otra parte, el homomorfismo $\widehat{u}:\widehat{\sh{O}}_y\to\widehat{\sh{O}}_{z_1}$ no es inyectivo: $\widehat{\sh{O}}_{z_1}$ es un anillo local integral, pues $z_1$ es un punto simple, mientras que $\widehat{\sh{O}}_y$ tiene dos ideales primos mínimos distintos correspondientes a las dos ramas de $\mathrm{Y}$ y, por tanto, tiene divisores de cero.
Esto proporciona un ejemplo que muestra que, en el enunciado de \sref{IV.11.5.2}, la hipótesis de que $\widehat{u}$ es inyectivo no puede sustituirse por la hipótesis de que el propio $u$ es inyectivo, ni siquiera cuando $\mathrm{A}'$ es un anillo local.
En efecto, con la notación precedente, basta tomar $\mathrm{A}=\sh{O}_y$, $\mathrm{A}'=\sh{O}_{z_1}$, $\mathrm{Y}=\operatorname{Spec}(\mathrm{A})$ y $\mathrm{X}=\operatorname{Spec}(\mathrm{A}')$; el razonamiento de \sref{IV.11.7.3} sigue demostrando que la primera proyección $p:\mathrm{X}\times_\mathrm{Y}\mathrm{X}\to\mathrm{X}$ es un isomorfismo, aunque $f$ no es plano en $z_1$.
\end{env}

\begin{env}[11.7.5]
\label{IV.11.7.5}
Los ejemplos de \sref{IV.11.7.2} y \sref{IV.11.7.3} explican la restricción a los anillos locales \emph{unirramificados} en \sref{IV.11.6.1} y \sref{IV.11.6.2}.
Veremos ahora que, en \sref{IV.11.6.1}, no se puede debilitar la hipótesis sobre $\mathrm{A}$ suponiendo únicamente que $\mathrm{A}$ es unirramificado.
Consideremos el anillo local integral completo $\mathrm{A}=\mathbf{R}[[\mathrm{U},\mathrm{V}]]/(\mathrm{U}^2+\mathrm{V}^2)$, que es unirramificado, pero no geométricamente unirramificado \sref{IV.6.15.11}.
Se sabe (\emph{loc.\ cit.}) que, si $u$ y $v$ son las imágenes de $\mathrm{U}$ y $\mathrm{V}$ en $\mathrm{A}$, el cierre integral de $\mathrm{A}$ es $\bar{\mathrm{A}}=\mathrm{A}[t]$, donde $t=v/u$ y $t^2=-1$, de modo que $\bar{\mathrm{A}}$ es isomorfo a $\mathbf{C}[[\mathrm{U}]]$.
Pongamos $\mathrm{Y}=\operatorname{Spec}(\mathrm{A})$ y $\mathrm{X}=\operatorname{Spec}(\bar{\mathrm{A}})$, la normalización de $\mathrm{Y}$ \sref[II]{II.6.3.8}, y sean $y$ y $x$ los puntos cerrados de $\mathrm{Y}$ y $\mathrm{X}$, respectivamente.
Mostraremos que, para una $\mathrm{A}$-álgebra local adecuada $\mathrm{A}'$, si se pone $\mathrm{Y}'=\operatorname{Spec}(\mathrm{A}')$ y $\mathrm{X}'=\mathrm{X}\times_\mathrm{Y}\mathrm{Y}'$, y si $y'$ denota el punto cerrado de $\mathrm{Y}'$, entonces $\sh{O}_{\mathrm{X}'}$ es $\mathrm{Y}'$-plano en un punto de $\mathrm{X}'$ cuyas proyecciones en $\mathrm{X}$ y $\mathrm{Y}'$ son, respectivamente, $x$ y $y'$, pero no es $\mathrm{Y}'$-plano en todos los puntos que tienen estas proyecciones.
De \sref{IV.2.1.4} se seguirá entonces que $\sh{O}_\mathrm{X}$ no es $\mathrm{Y}$-plano en $x$ (lo cual, por lo demás, es trivial \emph{a priori}, pues $\bar{\mathrm{A}}$ no es un $\mathrm{A}$-módulo libre).

Sea $\mathrm{B}=\mathrm{A}\otimes_\mathbf{R}\mathbf{C}$, que es isomorfo a $\mathbf{C}[[\mathrm{U},\mathrm{V}]]/((\mathrm{U}+i\mathrm{V})(\mathrm{U}-i\mathrm{V}))$.
El anillo $\mathrm{B}$ tiene dos ideales primos mínimos $\mathfrak{p}'$ y $\mathfrak{p}''$, generados, respectivamente, por $u+iv$ y $u-iv$, y $\mathfrak{n}=\mathfrak{p}'+\mathfrak{p}''$ es el ideal máximo del anillo local completo $\mathrm{B}$.
Sea $\bar{\mathrm{B}}=\bar{\mathrm{A}}\otimes_\mathbf{R}\mathbf{C}$; es el producto directo de dos álgebras isomorfas a $\mathbf{C}[[\mathrm{U}]]$, generado por los idempotentes $e'=(1\otimes 1+t\otimes i)/2$ y $e''=(1\otimes 1-t\otimes i)/2$.
Como el homomorfismo $\mathrm{A}\to\bar{\mathrm{A}}$ es inyectivo, también lo es $\mathrm{B}\to\bar{\mathrm{B}}$; las imágenes de $u+iv$ y $u-iv$ por esta inyección son, respectivamente, $2ue'$ y $2ue''$.
Se deduce inmediatamente que $\bar{\mathrm{B}}$ se identifica canónicamente con $(\mathrm{B}/\mathfrak{p}')\oplus(\mathrm{B}/\mathfrak{p}'')$.
Tomemos ahora como $\mathrm{A}'$ el anillo local $\mathrm{B}/\mathfrak{p}'$; entonces $\bar{\mathrm{A}}\otimes_\mathrm{A}\mathrm{A}'$ se identifica con $\bar{\mathrm{B}}\otimes_\mathrm{B}\mathrm{A}'$.
Pero $(\mathrm{B}/\mathfrak{p}')\otimes_\mathrm{B}(\mathrm{B}/\mathfrak{p}')=\mathrm{B}/\mathfrak{p}'$ y $(\mathrm{B}/\mathfrak{p}')\otimes_\mathrm{B}(\mathrm{B}/\mathfrak{p}'')=\mathrm{B}/\mathfrak{n}$, de modo que $\bar{\mathrm{A}}\otimes_\mathrm{A}\mathrm{A}'$ es isomorfo a $\mathrm{A}'\oplus(\mathrm{B}/\mathfrak{n})$.
Esto demuestra nuestra afirmación, pues $\mathrm{B}/\mathfrak{n}=\mathrm{A}'/(\mathfrak{n}/\mathfrak{p}')$ no es un $\mathrm{A}'$-módulo plano; de lo contrario, sería un $\mathrm{A}'$-módulo libre (Bourbaki, \emph{Alg.\ comm.}, cap.~II, \S~3, no.~2, cor.~2 de prop.~5), lo cual es absurdo.
\end{env}

\subsection{Un criterio valuativo de planitud}
\label{subsection:IV.11.8}

\begin{theorem}[11.8.1]
\label{IV.11.8.1}
Sean $f:\mathrm{X}\to\mathrm{Y}$ un morfismo localmente de presentación finita, $\mathscr{F}$ un $\sh{O}_\mathrm{X}$-módulo cuasicoherente de presentación finita, $x$ un punto de $\mathrm{X}$ y $y=f(x)$.
Supongamos que el anillo local $\sh{O}_y$ es íntegro (resp.\ reducido y noetheriano).
Para que $\mathscr{F}$ sea $f$-plano en el punto $x$, es necesario y suficiente que, para todo anillo de valoración (resp.\ todo anillo de valoración discreta) $\mathrm{A}'$ y todo homomorfismo local $\sh{O}_y\to\mathrm{A}'$, se satisfaga la condición siguiente: poniendo $\mathrm{Y}'=\operatorname{Spec}(\mathrm{A}')$, $\mathrm{X}'=\mathrm{X}\times_\mathrm{Y}\mathrm{Y}'$, $f'=f_{(\mathrm{Y}')}$, el $\sh{O}_{\mathrm{X}'}$-módulo $\mathscr{F}'=\mathscr{F}\otimes_\mathrm{Y}\mathrm{Y}'$ es $f'$-plano en todos los puntos $x'$ de $\mathrm{X}'$ cuyas proyecciones en $\mathrm{X}$ y $\mathrm{Y}'$ son, respectivamente, $x$ y el punto cerrado $y'$ de $\mathrm{Y}'$.
\end{theorem}

\begin{proof}
La condición es evidentemente necesaria \sref{IV.2.1.4}; queda por demostrar que es suficiente.
Podemos limitarnos evidentemente (\sref[I]{I.3.6.5} y \sref[I]{I.2.4.4}) al caso en que $\mathrm{Y}=\operatorname{Spec}(\mathrm{A})$ es el espectro de un anillo \emph{local} $\mathrm{A}$ y $y$ es el punto cerrado de $\mathrm{Y}$.
\medskip

(i) \emph{Caso en que $\mathrm{A}$ es íntegro.} --- Sean $\mathrm{K}$ el cuerpo de fracciones de $\mathrm{A}$ y $\mathrm{A}_1$ el cierre integral de $\mathrm{A}$; si ponemos $\mathrm{Y}_1=\operatorname{Spec}(\mathrm{A}_1)$, $\mathrm{X}_1=\mathrm{X}\times_\mathrm{Y}\mathrm{Y}_1$, $f_1=f_{(\mathrm{Y}_1)}$, basta, en virtud de \sref{IV.11.5.5}, demostrar que $\mathscr{F}_1=\mathscr{F}\otimes_\mathrm{Y}\mathrm{Y}_1$ es $f_1$-plano en \emph{todos} los puntos $x_1$ de $\mathrm{X}_1$ cuya proyección es el punto $x$ de $\mathrm{X}$.
Ahora bien, si $f_1(x_1)=y_1$, entonces $\mathfrak{i}_{y_1}=\mathfrak{m}_1$, donde $\mathfrak{m}_1$ es un ideal primo de $\mathrm{A}_1$ cuya contracción a $\mathrm{A}$ es $\mathfrak{m}=\mathfrak{i}_y$ (y $\mathfrak{m}_1$ es, además, necesariamente un ideal maximal). Sea $\mathrm{A}'$ un anillo de valoración de $\mathrm{K}$ que domina $\sh{O}_{y_1}=(\mathrm{A}_1)_{\mathfrak{m}_1}$; como el homomorfismo $\mathrm{A}\to\sh{O}_{y_1}$ es local, también lo es $\mathrm{A}\to\mathrm{A}'$.
Existe entonces al menos un
\oldpage[IV-3]{160}
punto $x'\in\mathrm{X}'$ cuyas proyecciones en $\mathrm{X}_1$ y $\mathrm{Y}'$ son, respectivamente, $x_1$ e $y'$ \sref[I]{I.3.4.7}; como la hipótesis afirma que $\mathscr{F}'$ es $f'$-plano en el punto $x'$ y $\sh{O}_{y_1}$ es íntegramente cerrado, podemos aplicar \sref{IV.11.6.1}, y deducimos que $\mathscr{F}_1$ es $f_1$-plano en el punto $x_1$, de donde resulta el teorema en este caso.
\medskip

(ii) \emph{Caso en que $\mathrm{A}$ es reducido y noetheriano.} --- Sean $\mathfrak{p}_i$ ($1\leq i\leq m$) los ideales primos minimales de $\mathrm{A}$; como $\mathrm{A}$ es reducido, se identifica canónicamente con un subanillo del producto de los $\mathrm{A}_i=\mathrm{A}/\mathfrak{p}_i$, que son anillos locales noetherianos; poniendo $\mathrm{Y}_i=\operatorname{Spec}(\mathrm{A}_i)$, $\mathrm{X}_i=\mathrm{X}\times_\mathrm{Y}\mathrm{Y}_i$, $f_i=f_{(\mathrm{Y}_i)}$, se deduce de \sref{IV.11.5.2} que basta demostrar, para cada $i$, que $\mathscr{F}_i=\mathscr{F}\otimes_\mathrm{Y}\mathrm{Y}_i$ es $f_i$-plano en todo punto $x_i$ de $\mathrm{X}_i$ cuyas proyecciones en $\mathrm{X}$ y $\mathrm{Y}_i$ son, respectivamente, $x$ y el punto cerrado $y_i$ de $\mathrm{Y}_i$.
Ahora bien, como $\mathrm{A}_i$ es un anillo local noetheriano íntegro, existe en su cuerpo de fracciones $\mathrm{K}_i$ un subanillo $\mathrm{A}''_i$ que contiene a $\mathrm{A}_i$, que es una $\mathrm{A}_i$-álgebra finita (y, por tanto, un anillo semilocal noetheriano) y cuyos anillos locales son \emph{geométricamente unirramificados} (\sref{IV.6.15.5} y \sref[0]{0.23.2.5}).
Como los ideales maximales de $\mathrm{A}''_i$ están necesariamente sobre el ideal maximal de $\mathrm{A}_i$, se deduce además de \sref[I]{I.3.4.7} y \sref{IV.11.5.2} que basta, poniendo $\mathrm{Y}''_i=\operatorname{Spec}(\mathrm{A}''_i)$, $\mathrm{X}''_i=\mathrm{X}\times_\mathrm{Y}\mathrm{Y}''_i$, $f''_i=f_{(\mathrm{Y}''_i)}$, demostrar que $\mathscr{F}''_i=\mathscr{F}\otimes_\mathrm{Y}\mathrm{Y}''_i$ es $f''_i$-plano en todo punto $x''_i$ de $\mathrm{X}''_i$ cuyas proyecciones en $\mathrm{X}$ y $\mathrm{Y}''_i$ son, respectivamente, $x$ y el punto cerrado $y''_i$ de $\mathrm{Y}''_i$.
Sea ahora $\mathrm{A}'$ un anillo de valoración discreta de $\mathrm{K}_i$ que domina $\mathrm{A}''_i$, y sea $x'$ un punto de $\mathrm{X}'$ cuyas proyecciones en $\mathrm{X}''_i$ y $\mathrm{Y}'$ son, respectivamente, $x''_i$ e $y'$ \sref[I]{I.3.4.7}; como $\sh{O}_{y''_i}$ es geométricamente unirramificado, podemos aplicar de nuevo \sref{IV.11.6.1}, y deducimos que $\mathscr{F}''_i$ es efectivamente $f''_i$-plano en el punto $x''_i$.
\end{proof}

\begin{remarks}[11.8.2]
\label{IV.11.8.2}
\medskip\noindent
\begin{enumerate}[label=(\roman*)]
  \item En el enunciado de \sref{IV.11.8.1}, podemos limitarnos a suponer que la condición sobre $\mathscr{F}'$ se satisface para los anillos de valoración \emph{completos} $\mathrm{A}'$ cuyo cuerpo residual sea \emph{algebraicamente cerrado}.
  En efecto, se sabe que todo anillo de valoración $\mathrm{A}'$ está dominado por un anillo $\mathrm{A}''$ de esta clase \sref[II]{II.7.1.2}, y que, si $\mathrm{A}'$ es un anillo de valoración discreta, lo mismo ocurre con $\mathrm{A}''$ ($\mathbf{0}_{\mathrm{III}}$,~10.3.1).
  \item La demostración de \sref{IV.11.8.1} se simplifica cuando se supone no solo que $\mathrm{A}$ es íntegro y noetheriano, sino también que su completado $\hat{\mathrm{A}}$ es \emph{íntegro}.
  Sustituyendo $\mathrm{X}$ por $\mathrm{X}\otimes_\mathrm{A}\hat{\mathrm{A}}$ y razonando como en la demostración de \sref{IV.11.5.3}, podemos reducirnos, en efecto, a demostrar \sref{IV.11.8.1} cuando $\mathrm{A}=\sh{O}_y$ es íntegro, noetheriano y \emph{completo}.
  Ahora bien, se sabe por \sref[II]{II.7.1.7} que tal anillo $\mathrm{A}$ está dominado por un anillo de valoración discreta completo; la conclusión resulta, por tanto, directamente de \sref{IV.11.5.2}.
\end{enumerate}
\end{remarks}

\subsection{Familias separantes y universalmente separantes de homomorfismos de haces de módulos}
\label{subsection:IV.11.9}

\begin{env}[11.9.1]
\label{IV.11.9.1}
Sea $\mathrm{X}$ un preesquema, sea $(f_\lambda)_{\lambda\in\mathrm{L}}$ una familia de morfismos $f_\lambda:\mathrm{Z}_\lambda\to\mathrm{X}$ y sea $\mathscr{F}$ un $\sh{O}_\mathrm{X}$-módulo cuasicoherente; para cada $\lambda\in\mathrm{L}$, supongamos dados un $\sh{O}_{\mathrm{Z}_\lambda}$-módulo cuasicoherente $\sh{G}_\lambda$ y un homomorfismo
\[
  u_\lambda:\mathscr{F}\to(f_\lambda)_*(\sh{G}_\lambda).
\label{IV.11.9.1.1}
\]
\end{env}

\oldpage[IV-3]{161}
Decimos que la familia $(u_\lambda)$ (o la familia correspondiente de los $u_\lambda^{\#}:(f_\lambda)^*(\mathscr{F})\to\sh{G}_\lambda$) es \emph{separante} si la intersección de los núcleos de los $u_\lambda$ es \emph{cero}.
En otras palabras, esto significa que, para todo abierto $\mathrm{U}$ de $\mathrm{X}$ y toda sección $t$ de $\mathscr{F}$ sobre $\mathrm{U}$ tales que, para \emph{cada} $\lambda$, la sección $u_\lambda(t)$ (que, por definición, es una sección de $\sh{G}_\lambda$ sobre $f_\lambda^{-1}(\mathrm{U})$) sea cero, la propia $t$ es cero.

\begin{env}[11.9.2]
\label{IV.11.9.2}
Con las notaciones de \sref{IV.11.9.1}, sea $\mathrm{M}$ un segundo conjunto de índices; para cada $\lambda\in\mathrm{L}$, sea $(g_{\lambda\mu})_{\mu\in\mathrm{M}}$ una familia de morfismos $g_{\lambda\mu}:\mathrm{Z}'_{\lambda\mu}\to\mathrm{Z}_\lambda$; para cada par $(\lambda,\mu)$, supongamos dados un $\sh{O}_{\mathrm{Z}'_{\lambda\mu}}$-módulo cuasicoherente $\sh{H}_{\lambda\mu}$ y un homomorfismo $v_{\lambda\mu}:\sh{G}_\lambda\to(g_{\lambda\mu})_*(\sh{H}_{\lambda\mu})$; pongamos $h_{\lambda\mu}=f_\lambda\circ g_{\lambda\mu}:\mathrm{Z}'_{\lambda\mu}\to\mathrm{X}$ y consideremos el homomorfismo compuesto
\[
  w_{\lambda\mu}:\mathscr{F}\xrightarrow{u_\lambda}(f_\lambda)_*(\sh{G}_\lambda)\xrightarrow{(f_\lambda)_*(v_{\lambda\mu})}(h_{\lambda\mu})_*(\sh{H}_{\lambda\mu}).
\]
Supongamos que, para cada $\lambda\in\mathrm{L}$, la familia $(v_{\lambda\mu})_{\mu\in\mathrm{M}}$ es separante: entonces lo mismo ocurre con la familia $(f_\lambda)_*(v_{\lambda\mu})$ ($\mu\in\mathrm{M}$), como se ve inmediatamente.
Se concluye que, para que la familia $(u_\lambda)$ sea separante, es necesario y suficiente que lo sea la familia $(w_{\lambda\mu})_{(\lambda,\mu)\in\mathrm{L}\times\mathrm{M}}$.
\end{env}

\begin{env}[11.9.3]
\label{IV.11.9.3}
Para verificar que la familia $(u_\lambda)$ es separante (con las notaciones de \sref{IV.11.9.1}), podemos reducirnos evidentemente primero al caso en que $\mathrm{X}$ es \emph{afín}, pues la propiedad es local sobre $\mathrm{X}$.
Podemos suponer además que $\mathrm{Z}_\lambda=\mathrm{X}$ \emph{para todo} $\lambda\in\mathrm{L}$.
En efecto, sea $\mathrm{M}$ un conjunto de índices, suma de una familia $(\mathrm{M}_\lambda)_{\lambda\in\mathrm{L}}$, y, para cada $\lambda\in\mathrm{L}$, sea $(\mathrm{Y}_{\lambda\mu})_{\mu\in\mathrm{M}_\lambda}$ un recubrimiento abierto afín de $\mathrm{Z}_\lambda$; sea $j_{\lambda\mu}:\mathrm{Y}_{\lambda\mu}\to\mathrm{Z}_\lambda$ la inmersión canónica y pongamos $\sh{H}_{\lambda\mu}=j_{\lambda\mu}^*(\sh{G}_\lambda)=\sh{G}_\lambda|\mathrm{Y}_{\lambda\mu}$.
Si consideramos el homomorfismo canónico $v_{\lambda\mu}:\sh{G}_\lambda\to(j_{\lambda\mu})_*(j_{\lambda\mu}^*(\sh{G}_\lambda))=(j_{\lambda\mu})_*(\sh{H}_{\lambda\mu})$ relativo a $j_{\lambda\mu}$ \sref[0]{0.4.4.3.2}, es inmediato que, para cada $\lambda\in\mathrm{L}$, la familia $(v_{\lambda\mu})_{\mu\in\mathrm{M}_\lambda}$ es separante.
En virtud de \sref{IV.11.9.2}, nos reducimos, por tanto, a demostrar que es separante la familia de homomorfismos compuestos $((f_\lambda)_*(v_{\lambda\mu}))\circ u_\lambda$, donde los $f_\lambda$ son \emph{afines}.
Pero entonces los $(f_\lambda)_*(\sh{G}_\lambda)$ son $\sh{O}_\mathrm{X}$-módulos cuasicoherentes \sref[I]{I.1.6.2} y, en virtud de la definición, podemos sustituir los $\mathrm{Z}_\lambda$ por $\mathrm{X}$ y los $\sh{G}_\lambda$ por $(f_\lambda)_*(\sh{G}_\lambda)$, de donde resulta nuestra afirmación.

Observemos además que, si $\mathrm{L}$ es \emph{finito} y los $f_\lambda$ son \emph{cuasicompactos}, podemos suponer, en la reducción anterior, que los $\mathrm{M}_\lambda$ son también \emph{finitos}; en este caso nos reducimos así a verificar que una familia \emph{finita} de homomorfismos de $\mathscr{F}$ en $\sh{O}_\mathrm{X}$-módulos cuasicoherentes es separante.
\end{env}

\begin{env}[11.9.4]
\label{IV.11.9.4}
Consideremos, pues, el caso en que $\mathrm{Z}_\lambda=\mathrm{X}$ para todo $\lambda$ y en que $\mathrm{X}=\operatorname{Spec}(\mathrm{A})$ es afín; entonces $\mathscr{F}=\widetilde{\mathrm{M}}$ y $\sh{G}_\lambda=\widetilde{\mathrm{N}}_\lambda$, donde $\mathrm{M}$ y los $\mathrm{N}_\lambda$ son $\mathrm{A}$-módulos, y $u_\lambda=\widetilde{\varphi}_\lambda$, donde los $\varphi_\lambda:\mathrm{M}\to\mathrm{N}_\lambda$ son homomorfismos de $\mathrm{A}$-módulos.
Decir que la familia $(u_\lambda)$ es separante significa entonces que, para todo $s\in\mathrm{A}$, la intersección de los núcleos de los $(\varphi_\lambda)_s:\mathrm{M}_s\to(\mathrm{N}_\lambda)_s$ es cero.
Se dice también entonces que la familia $(\varphi_\lambda)$ es \emph{separante}.
Observemos que, si $\mathrm{L}$ es \emph{finito}, esto equivale a decir que la intersección de los núcleos de los $\varphi_\lambda$ es cero, pues entonces
\[
  \bigcap_{\lambda\in\mathrm{L}}\operatorname{Ker}((\varphi_\lambda)_s)
  =\left(\bigcap_{\lambda\in\mathrm{L}}\operatorname{Ker}(\varphi_\lambda)\right)_s
  \quad \sref[0]{0.1.3.2}.
\]
Pero esta relación deja de ser exacta en general
\oldpage[IV-3]{162}
cuando $\mathrm{L}$ es \emph{infinito}, y el hecho de que la intersección de los núcleos de los $\varphi_\lambda$ sea cero \emph{no implica, por tanto},
en general, que la familia $(\varphi_\lambda)$ sea separante.
Por ejemplo, supongamos que $\mathrm{A}$ es un anillo de valoración discreta de ideal maximal $\mathfrak{m}$, y consideremos la familia de homomorfismos $\varphi_k:\mathrm{A}\to\mathrm{A}/\mathfrak{m}^k$; sus núcleos son los ideales $\mathfrak{m}^k$, cuya intersección es cero. Sin embargo, esta familia no es separante, pues las fibras de todos los $\mathfrak{m}^k$ en el punto genérico $x$ de $\mathrm{X}=\operatorname{Spec}(\mathrm{A})$ (que es abierto en $\mathrm{X}$) son iguales a $\sh{O}_x=k(x)$, el cuerpo de fracciones de $\mathrm{A}$, y su intersección no es, por tanto, cero.
\end{env}

\begin{env}[11.9.5]
\label{IV.11.9.5}
En lo que sigue nos ocuparemos principalmente del problema del \emph{cambio de base} para las familias separantes.
Con las notaciones de \sref{IV.11.9.1}, consideremos un morfismo $g:\mathrm{X}'\to\mathrm{X}$ y pongamos $\mathscr{F}'=\mathscr{F}\otimes_{\sh{O}_\mathrm{X}}\sh{O}_{\mathrm{X}'}=g^*(\mathscr{F})$ y, para cada $\lambda$, $\mathrm{Z}'_\lambda=\mathrm{Z}_\lambda\times_\mathrm{X}\mathrm{X}'$, $f'_\lambda=f_\lambda\times 1$, $\sh{G}'_\lambda=\sh{G}_\lambda\otimes_{\sh{O}_{\mathrm{Z}_\lambda}}\sh{O}_{\mathrm{Z}'_\lambda}=(g'_\lambda)^*(\sh{G}_\lambda)$, donde $g'_\lambda:\mathrm{Z}'_\lambda\to\mathrm{Z}_\lambda$ es la proyección.
Para cada $\lambda$, denotamos por $u'_\lambda:\mathscr{F}'\to(f'_\lambda)_*(\sh{G}'_\lambda)$ el homomorfismo que se obtiene de la manera siguiente: sea
\[
  \psi:(f_\lambda)_*(\sh{G}_\lambda)\to(f_\lambda)_*((g'_\lambda)_*((g'_\lambda)^*(\sh{G}_\lambda)))=g_*((f'_\lambda)_*(\sh{G}'_\lambda))
\]
el homomorfismo $(f_\lambda)_*(\rho_{\sh{G}_\lambda})$, donde $\rho_{\sh{G}_\lambda}$ es el homomorfismo canónico \sref[0]{0.4.4.3.2} correspondiente a $g'_\lambda$.
Entonces $u'_\lambda$ se define como el compuesto
\[
  g^*(\mathscr{F})\xrightarrow{g^*(u_\lambda)}g^*((f_\lambda)_*(\sh{G}_\lambda))\xrightarrow{g^*(\psi)}g^*g_*((f'_\lambda)_*(\sh{G}'_\lambda))\xrightarrow{\sigma}(f'_\lambda)_*(\sh{G}'_\lambda)
\]
donde $\sigma=\sigma_{(f'_\lambda)_*(\sh{G}'_\lambda)}$ es el homomorfismo canónico \sref[0]{0.4.4.3.3} correspondiente a $g$.
Diremos, para abreviar, que $u'_\lambda$ se \emph{deduce de $u_\lambda$ por el cambio de base} $g$.
Cuando $f_\lambda:\mathrm{Z}_\lambda\to\mathrm{X}$ es un morfismo \emph{afín}, $(f'_\lambda)_*(\sh{G}'_\lambda)=g^*((f_\lambda)_*(\sh{G}_\lambda))$ y, en este caso, se tiene simplemente $u'_\lambda=g^*(u_\lambda)$.

También podemos interpretar $u'_\lambda$ de la manera siguiente: basta conocer el valor de $u'_\lambda(t')$, donde $t'$ es una sección de $\mathscr{F}'$ sobre un abierto $\mathrm{U}'$ de $\mathrm{X}'$ del tipo particular siguiente: $t'$ es la restricción a $\mathrm{U}'$ de la imagen canónica, por $\Gamma(\mathrm{U},\mathscr{F})\to\Gamma(g^{-1}(\mathrm{U}),\mathscr{F}')$, de una sección $t$ de $\mathscr{F}$ sobre un abierto $\mathrm{U}$ de $\mathrm{X}$ que contiene a $g(\mathrm{U}')$ (estas secciones generan el $\sh{O}_{\mathrm{X}'}$-módulo $\mathscr{F}'$ \sref[0]{0.3.7.1}).
Consideremos la sección $u_\lambda(t)$ de $\sh{G}_\lambda$ sobre $f_\lambda^{-1}(\mathrm{U})$ y su imagen canónica $t''$ por $\Gamma(f_\lambda^{-1}(\mathrm{U}),\sh{G}_\lambda)\to\Gamma(g_\lambda'^{-1}(f_\lambda^{-1}(\mathrm{U})),\sh{G}'_\lambda)$; entonces $u'_\lambda(t')$ es la restricción de $t''$ a $f_\lambda'^{-1}(\mathrm{U}')$.
Consideremos, en particular, el caso en que $\mathrm{Z}_\lambda$ es el subpreesquema inducido sobre un subconjunto abierto de $\mathrm{X}$, en que $\sh{G}_\lambda=j_\lambda^*(\mathscr{F})$, siendo $j_\lambda:\mathrm{Z}_\lambda\to\mathrm{X}$ la inmersión canónica, y en que $u_\lambda$ es el homomorfismo canónico $\rho_\mathscr{F}:\mathscr{F}\to(j_\lambda)_*(j_\lambda^*(\mathscr{F}))=(j_\lambda)_*(\sh{G}_\lambda)$ \sref[0]{0.4.4.3.2} correspondiente a $j_\lambda$.
Entonces $\mathrm{Z}'_\lambda$ está inducido sobre un subconjunto abierto de $\mathrm{X}'$, y la interpretación anterior muestra que $u'_\lambda$ no es sino el homomorfismo canónico $\rho_{\mathscr{F}'}:\mathscr{F}'\to(j'_\lambda)_*((j'_\lambda)^*(\mathscr{F}'))=(j'_\lambda)_*(\sh{G}'_\lambda)$ correspondiente a la inmersión canónica $j'_\lambda:\mathrm{Z}'_\lambda\to\mathrm{X}'$.
\end{env}

\begin{env}[11.9.6]
\label{IV.11.9.6}
En las condiciones de \sref{IV.11.9.5}, supongamos que $\mathrm{X}$ y $\mathrm{X}'$ son ambos afines y que se quiere demostrar que toda sección $t'$ de $\mathscr{F}'$ sobre $\mathrm{X}'$ cuyas imágenes por todos los $u'_\lambda$ sean cero es ella misma cero.
Entonces también podemos limitarnos al caso en que $\mathrm{Z}_\lambda=\mathrm{X}$ para todo $\lambda\in\mathrm{L}$.
En efecto, por \sref{IV.11.9.3} y \sref{IV.11.9.5}, si ponemos $\mathrm{Y}'_{\lambda\mu}=g_\lambda'^{-1}(\mathrm{Y}_{\lambda\mu})$, el homomorfismo $v'_{\lambda\mu}$ deducido de $v_{\lambda\mu}$ por el cambio de base $g$ no es sino el homomorfismo canónico
\oldpage[IV-3]{163}
$\rho_{\sh{G}'_\lambda}:\sh{G}'_\lambda\to(j'_{\lambda\mu})_*(j_{\lambda\mu}'^*(\sh{G}'_\lambda))$ correspondiente a la inmersión canónica $j'_{\lambda\mu}:\mathrm{Y}'_{\lambda\mu}\to\mathrm{Z}'_\lambda$, como vimos en \sref{IV.11.9.5}.
La afirmación se deduce entonces de los razonamientos de \sref{IV.11.9.2} y \sref{IV.11.9.3}, sustituyendo $\mathrm{Y}_{\lambda\mu}$ y $v_{\lambda\mu}$ por $\mathrm{Y}'_{\lambda\mu}$ y $v'_{\lambda\mu}$.

Se tiene una reducción análoga cuando se quiere demostrar que la familia $(u'_\lambda)$ es \emph{separante} (siendo $\mathrm{X}$ y $\mathrm{X}'$ afines): esto se deduce de nuevo de \sref{IV.11.9.2} y \sref{IV.11.9.3}.
\end{env}

\begin{proposition}[11.9.7]
\label{IV.11.9.7}
Con las notaciones de \sref{IV.11.9.1} y \sref{IV.11.9.5}, supongamos que $\mathrm{X}=\operatorname{Spec}(\mathrm{A})$ y $\mathrm{X}'=\operatorname{Spec}(\mathrm{A}')$ son afines y, además, que $\mathrm{A}'$ es un $\mathrm{A}$-módulo proyectivo.
Entonces, si la familia $(u_\lambda)$ es separante, toda sección $t'$ de $\mathscr{F}'$ sobre $\mathrm{X}'$ cuyas imágenes por todos los $u'_\lambda$ sean cero es ella misma cero.
\end{proposition}

\begin{proof}
Hemos visto en \sref{IV.11.9.6} que podemos limitarnos al caso en que todos los $\mathrm{Z}_\lambda$ son \emph{iguales a} $\mathrm{X}$.
La proposición es entonces una consecuencia del siguiente lema:
\end{proof}

\begin{lemma}[11.9.7.1]
\label{IV.11.9.7.1}
Sea $\mathrm{A}$ un anillo, sea $(\mathrm{M}_\lambda)_{\lambda\in\mathrm{L}}$ una familia de $\mathrm{A}$-módulos, sea $\mathrm{M}$ un $\mathrm{A}$-módulo y, para cada $\lambda$, sea $u_\lambda:\mathrm{M}\to\mathrm{M}_\lambda$ un homomorfismo.
Supongamos que la intersección de los núcleos de los $u_\lambda$ se reduce a cero.
Entonces, para todo $\mathrm{A}$-módulo proyectivo $\mathrm{N}$, la intersección de los núcleos de los homomorfismos $u_\lambda\otimes 1:\mathrm{M}\otimes_\mathrm{A}\mathrm{N}\to\mathrm{M}_\lambda\otimes_\mathrm{A}\mathrm{N}$ se reduce a cero.
\end{lemma}

\begin{proof}
En efecto, $\mathrm{N}$ es un factor directo de un $\mathrm{A}$-módulo libre $\mathrm{L}$, y basta evidentemente demostrar que la intersección de los núcleos de los homomorfismos $v_\lambda:\mathrm{M}\otimes_\mathrm{A}\mathrm{L}\to\mathrm{M}_\lambda\otimes_\mathrm{A}\mathrm{L}$ se reduce a cero, pues $u_\lambda\otimes 1$ es la restricción de $v_\lambda$.
Pero la afirmación resulta entonces trivialmente de la hipótesis.
\end{proof}

\begin{remarks}[11.9.8]
\label{IV.11.9.8}
No sabemos si, bajo las hipótesis de la proposición \sref{IV.11.9.7}, la familia $(u'_\lambda)$ es \emph{separante}: en efecto, por \sref{IV.11.9.4} habría que demostrar que toda sección $t'$ de $\mathscr{F}'$ sobre un subconjunto abierto $\mathrm{D}(h')\subset\mathrm{X}'$ (donde $h'\in\mathrm{A}'$) tal que todos los $u'_\lambda(t')$ sean cero es ella misma cero.
Ahora bien, no podemos aplicar la proposición \sref{IV.11.9.7} a $\mathrm{D}(h')=\operatorname{Spec}(\mathrm{A}'_{h'})$, pues el hecho de que $\mathrm{A}'$ sea un $\mathrm{A}$-módulo proyectivo (o incluso libre) no implica que $\mathrm{A}'_{h'}$ sea un $\mathrm{A}$-módulo proyectivo.
Por ejemplo, podemos tomar como $\mathrm{A}$ un anillo de valoración discreta, como $\mathrm{A}'$ un anillo de valoración discreta que sea un $\mathrm{A}$-módulo libre de rango $2$, y como $\mathrm{A}'_{h'}$ el cuerpo de fracciones de $\mathrm{A}'$.
\end{remarks}

Se tiene, no obstante, el resultado siguiente:

\begin{corollary}[11.9.9]
\label{IV.11.9.9}
Sea $\mathrm{X}$ un preesquema artiniano y sea $g:\mathrm{X}'\to\mathrm{X}$ un morfismo plano (observemos que estas dos condiciones se satisfacen si $\mathrm{X}$ es el espectro de un \emph{cuerpo} y $g$ es un morfismo arbitrario).
Entonces, con las notaciones de \sref{IV.11.9.1} y \sref{IV.11.9.5}, si la familia $(u_\lambda)$ es separante, lo mismo ocurre con la familia $(u'_\lambda)$.
\end{corollary}

\begin{proof}
Podemos limitarnos evidentemente al caso en que $\mathrm{X}=\operatorname{Spec}(\mathrm{A})$ es el espectro de un anillo local artiniano $\mathrm{A}$ \sref[I]{I.6.2.2}; observemos entonces que, para todo abierto afín $\mathrm{U}'=\operatorname{Spec}(\mathrm{A}')$ de $\mathrm{X}'$, $\mathrm{A}'$ es un $\mathrm{A}$-módulo plano y, por tanto, \emph{proyectivo} ($\mathbf{0}_{\mathrm{III}}$,~10.1.3).
Basta, por tanto, aplicar \sref{IV.11.9.7} a todo abierto afín de $\mathrm{X}'$ para obtener el corolario.
\end{proof}

\begin{theorem}[11.9.10]
\label{IV.11.9.10}
Sea $\mathrm{X}$ un preesquema, sea $(u_\lambda)$ una familia de homomorfismos \sref{IV.11.9.1.1}, sea $g:\mathrm{X}'\to\mathrm{X}$ un morfismo y sea $(u'_\lambda)$ la familia de homomorfismos deducida de $(u_\lambda)$ por el cambio de base $g$ \sref{IV.11.9.5}.
\begin{enumerate}[label=(\roman*)]
  \item Si $g$ es un morfismo fielmente plano y la familia $(u'_\lambda)$ es separante, entonces la familia $(u_\lambda)$ es separante.
  \item Supongamos que $g$ es un morfismo plano y que, además, se satisface una de las dos condiciones siguientes:
  \begin{enumerate}[label=\emph{\alph*)}]
    \item $\mathrm{L}$ es finito y los $f_\lambda$ son cuasicompactos.
    \item El morfismo $g$ es localmente de presentación finita.
  \end{enumerate}
  Entonces, si la familia $(u_\lambda)$ es separante, lo mismo ocurre con la familia $(u'_\lambda)$.
\end{enumerate}
\end{theorem}

\begin{proof}
\oldpage[IV-3]{164}
\begin{enumerate}[label=(\roman*)]
  \item En virtud de \sref{IV.2.2.8}, basta demostrar que, si una sección $t$ de $\mathscr{F}$ sobre un abierto $\mathrm{U}$ de $\mathrm{X}$ pertenece al núcleo de cada uno de los $u_\lambda$, su imagen $t'$ por el homomorfismo canónico $\Gamma(\rho):\Gamma(\mathrm{U},\mathscr{F})\to\Gamma(g^{-1}(\mathrm{U}),g^*(\mathscr{F}))=\Gamma(\mathrm{U},g_*(g^*(\mathscr{F})))$ es cero.
  Ahora bien, las imágenes de $t'$ por los $g^*(u_\lambda)$ son las imágenes de $u_\lambda(t)$ por el homomorfismo $\Gamma(\mathrm{U},(f_\lambda)_*(\sh{G}_\lambda))\to\Gamma(g^{-1}(\mathrm{U}),g^*((f_\lambda)_*(\sh{G}_\lambda)))$; por tanto, son cero y, \emph{a fortiori}, $u'_\lambda(t')=0$ para todo $\lambda$. Por hipótesis, se sigue que $t'=0$, lo que demuestra~(i).

  \item Como la cuestión es local sobre $\mathrm{X}$ y $\mathrm{X}'$, podemos restringirnos al caso en que $\mathrm{X}=\operatorname{Spec}(\mathrm{A})$ y $\mathrm{X}'=\operatorname{Spec}(\mathrm{A}')$ son afines y $\mathrm{A}'$ es un $\mathrm{A}$-módulo plano; basta probar que toda sección $z'$ de $\mathscr{F}'$ sobre $\mathrm{X}'$ cuyas imágenes por todos los $u'_\lambda$ sean cero es ella misma cero.
  Hemos visto en \sref{IV.11.9.6} que también podemos restringirnos al caso en que $\mathrm{Z}_\lambda=\mathrm{X}$ para todo $\lambda\in\mathrm{L}$.
  Distinguimos ahora los dos casos.
  \medskip

  \emph{a)} Si $\mathrm{L}$ es finito y los $f_\lambda$ son cuasicompactos, hemos visto en \sref{IV.11.9.3} que podemos reducirnos además al caso en que $\mathrm{Z}_\lambda=\mathrm{X}$ para todo $\lambda\in\mathrm{L}$ y en que, asimismo, $\mathrm{L}$ es finito.
  Entonces, decir que la intersección de los núcleos de los $u_\lambda:\mathscr{F}\to\sh{G}_\lambda$ es cero equivale a decir que el homomorfismo $u=(u_\lambda):\mathscr{F}\to\sh{G}=\bigoplus_\lambda\sh{G}_\lambda$ es inyectivo.
  Como $u'=g^*(u)$ es inyectivo porque $g$ es plano, la proposición queda demostrada en este caso.

  \emph{b)} Sean $\mathscr{F}=\widetilde{\mathrm{M}}$ y $\sh{G}_\lambda=\widetilde{\mathrm{N}}_\lambda$, y pongamos $\mathrm{M}'=\mathrm{M}\otimes_\mathrm{A}\mathrm{A}'$, $\mathrm{N}'_\lambda=\mathrm{N}_\lambda\otimes_\mathrm{A}\mathrm{A}'$, de modo que $\mathscr{F}'=\widetilde{\mathrm{M}'}$ y $\sh{G}'_\lambda=\widetilde{\mathrm{N}'_\lambda}$; por abuso de notación, seguiremos denotando por $u_\lambda$ el homomorfismo $\mathrm{M}\to\mathrm{N}_\lambda$, y por $u'_\lambda$ el homomorfismo $u_\lambda\otimes 1:\mathrm{M}'\to\mathrm{N}'_\lambda$.
  Sea $z'\in\mathrm{M}'$ un elemento tal que $u'_\lambda(z')=0$ para todo $\lambda$; queda por demostrar que $z'=0$.
  Ahora bien, la hipótesis de que $g$ es plano y de presentación finita implica, por \sref{IV.11.3.15}, que existe una sucesión finita $(s_i)_{1\leq i\leq n}$ de elementos de $\mathrm{A}$ tal que, si se pone $\mathfrak{J}_k=s_1\mathrm{A}+\ldots+s_k\mathrm{A}$ y $\mathrm{A}_i=\mathrm{A}_{s_i}/\mathfrak{J}_{i-1}\mathrm{A}_{s_i}$, el anillo $\mathrm{A}'_i=\mathrm{A}'\otimes_\mathrm{A}\mathrm{A}_i$ es \emph{libre} como $\mathrm{A}_i$-módulo para $1\leq i\leq n$, y $\mathfrak{J}_n=\mathrm{A}$.
  La proposición quedará demostrada si probamos la siguiente afirmación para $1\leq i\leq n$:
  \begin{enumerate}[label=$(*_i)$]
    \item \emph{Existe un entero $m_i>0$ tal que $s_j^{m_i}z'=0$ para $j\leq i$.}
  \end{enumerate}
  En efecto, poniendo $k=m_n$ y observando que los $s_i^k$ ($1\leq i\leq n$) generan el ideal unidad de $\mathrm{A}$, la afirmación $(*_n)$ mostrará que $z'$, por ser combinación lineal de los $s_i^k z'$, es cero.

  Demostremos $(*_i)$ por inducción sobre $i$; la afirmación es trivialmente cierta para $i=0$.
  Supongamos $i>0$ y sea $m$ un múltiplo común de los $m_j$ para $j\leq i-1$.
  Observemos (para $1\leq h\leq n$) que, si $\mathfrak{J}'_h$ es el ideal generado por los $s_j^m$ ($1\leq j\leq h$), entonces $\mathfrak{J}_h/\mathfrak{J}'_h$ es nilpotente.
  Sustituir los $s_j$ por $s_j^m$ para $1\leq j\leq n$ equivale, por tanto, a sustituir, para $1\leq h\leq n$, $\mathrm{A}_h$ por
  $\mathrm{B}_h=\mathrm{A}_{s_h}/\mathfrak{J}'_{h-1}\mathrm{A}_{s_h}$, para que
  $\mathrm{A}_h=\mathrm{B}_h/((\mathfrak{J}_{h-1}/\mathfrak{J}'_{h-1})\mathrm{B}_h)$; como $\mathrm{A}'$ es un $\mathrm{A}$-módulo plano, se deduce de ($\mathbf{0}_{\mathrm{III}}$,~10.1.2) que $\mathrm{B}'_h=\mathrm{A}'\otimes_\mathrm{A}\mathrm{B}_h$ vuelve a ser un $\mathrm{B}_h$-módulo libre.
  Por consiguiente, se pueden
  \oldpage[IV-3]{165}
  sustituir todos los $s_j$ ($1\leq j\leq n$) por $s_j^m$ sin alterar las propiedades de $\mathfrak{J}_h$ y $\mathrm{A}'_h$, y suponer que $m=1$.
  Entonces $z'$, al estar aniquilado por $\mathfrak{J}_{i-1}$, se identifica con un elemento de $\operatorname{Hom}_{\mathrm{A}'}(\mathrm{A}'/\mathfrak{J}_{i-1}\mathrm{A}',\mathrm{M}')$; y, como $\mathfrak{J}_{i-1}\mathrm{A}'$ es un ideal de tipo finito de $\mathrm{A}'$, este módulo de homomorfismos se identifica a su vez con $\operatorname{Hom}_\mathrm{A}(\mathrm{A}/\mathfrak{J}_{i-1},\mathrm{M})\otimes_\mathrm{A}\mathrm{A}'$ ($\mathbf{0}_{\mathrm{I}}$,~6.2.2).
  Sea $v_\lambda=\operatorname{Hom}(1,u_\lambda):\operatorname{Hom}_\mathrm{A}(\mathrm{A}/\mathfrak{J}_{i-1},\mathrm{M})\to\operatorname{Hom}_\mathrm{A}(\mathrm{A}/\mathfrak{J}_{i-1},\mathrm{N}_\lambda)$ el homomorfismo deducido de $u_\lambda$; la familia $(v_\lambda)$ también es \emph{separante}.
  En efecto, para todo $t\in\mathrm{A}$,
  $(\operatorname{Hom}_\mathrm{A}(\mathrm{A}/\mathfrak{J}_{i-1},\mathrm{M}))_t
  =\operatorname{Hom}_{\mathrm{A}_t}(\mathrm{A}_t/(\mathfrak{J}_{i-1})_t,\mathrm{M}_t)$,
  y lo mismo vale sustituyendo $\mathrm{M}$ por $\mathrm{N}_\lambda$, pues el ideal $\mathfrak{J}_{i-1}$ es de tipo finito ($\mathbf{0}_{\mathrm{I}}$,~1.3.5); como, por hipótesis, la intersección de los núcleos de los $(u_\lambda)_t$ es cero, lo mismo ocurre con la intersección de los núcleos de los $(v_\lambda)_t=\operatorname{Hom}(1,(u_\lambda)_t)$, de donde nuestra afirmación por \sref{IV.11.9.4}.
  Sustituyendo $\mathrm{A}$ por $\mathrm{A}/\mathfrak{J}_{i-1}$, $\mathrm{M}$ por $\operatorname{Hom}_\mathrm{A}(\mathrm{A}/\mathfrak{J}_{i-1},\mathrm{M})$, $\mathrm{N}_\lambda$ por $\operatorname{Hom}_\mathrm{A}(\mathrm{A}/\mathfrak{J}_{i-1},\mathrm{N}_\lambda)$ (que son $(\mathrm{A}/\mathfrak{J}_{i-1})$-módulos), $u_\lambda$ por $v_\lambda$ y, finalmente, $\mathrm{A}'$ por $\mathrm{A}'/\mathfrak{J}_{i-1}\mathrm{A}'$, vemos que podemos reducirnos al caso en que, en la situación inicial, el elemento $s=s_i\in\mathrm{A}$ es tal que $\mathrm{A}'_s$ es un $\mathrm{A}_s$-módulo libre.
  Ahora bien, la familia $(u_\lambda)_s:\mathrm{M}_s\to(\mathrm{N}_\lambda)_s$ es separante por hipótesis; como $(u'_\lambda)_s(z'/1)=0$, se deduce de \sref{IV.11.9.7} que $z'/1=0$ en $\mathrm{M}'_s$, lo que significa que existe un entero $r$ tal que $s^r z'=0$ en $\mathrm{M}'$. Esto concluye por inducción la demostración de $(*_i)$.
\end{enumerate}
\end{proof}

\begin{remark}[11.9.11]
\label{IV.11.9.11}
Limitémonos, por simplicidad, al caso en que $\mathrm{Z}_\lambda=\mathrm{X}$ para todo $\lambda$.
Conviene observar que, si la familia de homomorfismos $u_\lambda:\mathscr{F}\to\sh{G}_\lambda$ es separante, \emph{no se sigue necesariamente} que, para todo $x\in\mathrm{X}$, la intersección de los núcleos de los homomorfismos $(u_\lambda)_x:\mathscr{F}_x\to(\sh{G}_\lambda)_x$ se reduzca a $0$.
Por ejemplo, sea $\mathrm{X}$ un preesquema localmente noetheriano de Jacobson, de dimensión $\geq 1$, y sea $\mathscr{F}$ un $\sh{O}_\mathrm{X}$-módulo coherente; para todo punto cerrado $x$ de $\mathrm{X}$ y todo entero $n\geq 0$, $\mathscr{F}/\mathfrak{m}_x^n\mathscr{F}$ es un $\sh{O}_\mathrm{X}$-módulo coherente cuyo soporte está contenido en $\{x\}$.
La familia de homomorfismos canónicos $\mathscr{F}\to\mathscr{F}/\mathfrak{m}_x^n\mathscr{F}$ (donde $x$ recorre el conjunto $\mathrm{X}_0$ de los puntos cerrados de $\mathrm{X}$ y $n$ el conjunto de los enteros $\geq 0$) es \emph{separante}: en efecto, si $t$ es una sección de $\mathscr{F}$ sobre un abierto $\mathrm{U}$ de $\mathrm{X}$ cuyas imágenes en todos los $\Gamma(\mathrm{U},\mathscr{F}/\mathfrak{m}_x^n\mathscr{F})$ son nulas, se tiene inmediatamente $t_x=0$ para todo punto cerrado $x\in\mathrm{U}$; como el conjunto de los puntos cerrados contenidos en $\mathrm{U}$ es muy denso en $\mathrm{U}$, esto implica que $t=0$ \sref{IV.10.2.1}.
Sin embargo, si se toma $\mathscr{F}=\sh{O}_\mathrm{X}$ y si $y\in\mathrm{X}$ es un punto \emph{no cerrado}, se tiene $(\sh{O}_\mathrm{X}/\mathfrak{m}_x^n\sh{O}_\mathrm{X})_y=0$ para todo punto cerrado $x$ de $\mathrm{X}$, pero $\sh{O}_{\mathrm{X},y}\neq 0$, y la intersección de los núcleos de los homomorfismos $\sh{O}_{\mathrm{X},y}\to(\sh{O}_\mathrm{X}/\mathfrak{m}_x^n\sh{O}_\mathrm{X})_y$ es igual a $\sh{O}_{\mathrm{X},y}$.
\end{remark}

\begin{lemma}[11.9.12]
\label{IV.11.9.12}
Sean las notaciones las de \sref{IV.11.9.1} y \sref{IV.11.9.5}, y supongamos que la familia $(u_\lambda)$ es separante; supongamos además que $\mathrm{X}$ es un $\mathrm{S}$-preesquema, donde $\mathrm{S}=\operatorname{Spec}(\mathrm{A})$ es un esquema afín, y que $\mathrm{X}'=\mathrm{X}\times_\mathrm{S}\mathrm{S}'$, donde $\mathrm{S}'=\operatorname{Spec}(\mathrm{A}')$, siendo $\mathrm{A}'$ una $\mathrm{A}$-álgebra; supongamos, por último, que los $\sh{G}_\lambda$ son planos sobre $\mathrm{S}$.
Sea $(\mathrm{A}'_\alpha)_{\alpha\in\mathrm{I}}$ la familia filtrante de las sub-$\mathrm{A}$-álgebras de tipo finito de $\mathrm{A}'$ y, para todo $\alpha\in\mathrm{I}$, sean $\mathrm{S}''_\alpha=\operatorname{Spec}(\mathrm{A}'_\alpha)$, $\mathrm{X}''_\alpha=\mathrm{X}\times_\mathrm{S}\mathrm{S}''_\alpha$, $\mathrm{Z}''_{\alpha\lambda}=\mathrm{Z}_\lambda\times_\mathrm{X}\mathrm{X}''_\alpha$, y sean $f''_{\alpha\lambda}:\mathrm{Z}''_{\alpha\lambda}\to\mathrm{X}''_\alpha$, $\mathscr{F}''_\alpha$, $\sh{G}''_{\alpha\lambda}$ y $u''_{\alpha\lambda}$ los morfismos, módulos y homomorfismos de módulos deducidos de $f_\lambda$, $\mathscr{F}$, $\sh{G}_\lambda$ y $u_\lambda$ mediante el cambio de base $\mathrm{X}''_\alpha\to\mathrm{X}$.
Entonces, si para todo $\alpha\in\mathrm{I}$ la familia $(u''_{\alpha\lambda})_{\lambda\in\mathrm{L}}$ es separante, también lo es $(u'_\lambda)_{\lambda\in\mathrm{L}}$.
\end{lemma}

\begin{proof}
Basta probar que, si $t'$ es una sección de $\mathscr{F}'$ sobre un abierto afín $\mathrm{U}'$ de $\mathrm{X}'$ tal que $u'_\lambda(t')=0$ para todo $\lambda\in\mathrm{L}$, entonces $t'=0$.
Si $h_\alpha:\mathrm{X}'\to\mathrm{X}''_\alpha$ es la proyección canónica,
\oldpage[IV-3]{166}
se deduce de \sref{IV.8.2.11} que existen un índice $\alpha$ y un abierto cuasicompacto $\mathrm{U}''_\alpha\subset\mathrm{X}''_\alpha$ tales que $\mathrm{U}'=h_\alpha^{-1}(\mathrm{U}''_\alpha)$; además, en virtud de \sref{IV.8.5.2},~(i), se puede suponer que $t'$ es la imagen canónica de una sección $t''_\alpha$ de $\mathscr{F}''_\alpha$ sobre $\mathrm{U}''_\alpha$.
Reemplazando, si es necesario, $\mathrm{S}$, $\mathrm{X}$, $\mathrm{Z}_\lambda$, $f_\lambda$, $\mathscr{F}$, $\sh{G}_\lambda$ y $u_\lambda$ por $\mathrm{S}''_\alpha$, $\mathrm{U}''_\alpha$, $(f''_{\alpha\lambda})^{-1}(\mathrm{U}''_\alpha)$ y las restricciones correspondientes de $\mathscr{F}''_\alpha$, $\sh{G}''_{\alpha\lambda}$ y $u''_{\alpha\lambda}$, podemos suponer, por tanto, que $\mathrm{U}'=\mathrm{X}'$, que $t'$ es la imagen canónica de una sección $t$ de $\mathscr{F}$ sobre $\mathrm{X}$ y que el homomorfismo $\mathrm{A}\to\mathrm{A}'$ es inyectivo; de manera equivalente, si $p:\mathrm{S}'\to\mathrm{S}$ es el morfismo correspondiente, el homomorfismo $\sh{O}_\mathrm{S}\to p_*(\sh{O}_{\mathrm{S}'})$ que interviene en la definición de $p$ es \emph{inyectivo}.
De la planitud de $\sh{G}_\lambda$ sobre $\mathrm{S}$ se deduce inmediatamente (reduciendo, por ejemplo, al caso en que $\mathrm{Z}_\lambda$ es afín sobre $\mathrm{S}$ \sref[I]{I.1.6.3} y \sref[I]{I.1.6.5}) que el homomorfismo canónico $\rho:\sh{G}_\lambda\to(g'_\lambda)_*(\sh{G}'_\lambda)$ es también \emph{inyectivo}.
Pero el homomorfismo compuesto
\[
  \Gamma(\mathrm{X},\mathscr{F})\xrightarrow{\Gamma(u_\lambda)}\Gamma(\mathrm{Z}_\lambda,\sh{G}_\lambda)\xrightarrow{\Gamma(\rho)}\Gamma(\mathrm{Z}'_\lambda,\sh{G}'_\lambda)
\]
es igual, por definición, al homomorfismo compuesto
\[
  \Gamma(\mathrm{X},\mathscr{F})\xrightarrow{\Gamma(\rho)}\Gamma(\mathrm{X}',\mathscr{F}')\xrightarrow{\Gamma(u'_\lambda)}\Gamma(\mathrm{Z}'_\lambda,\sh{G}'_\lambda)
\]
Por consiguiente, la imagen de $t$ mediante estos homomorfismos compuestos es $u'_\lambda(t')=0$; en virtud de la inyectividad del homomorfismo $\Gamma(\mathrm{Z}_\lambda,\sh{G}_\lambda)\xrightarrow{\Gamma(\rho)}\Gamma(\mathrm{Z}'_\lambda,\sh{G}'_\lambda)$, se concluye que $u_\lambda(t)=0$ para todo $\lambda\in\mathrm{L}$, de donde $t=0$ por hipótesis y, finalmente, $t'=0$.
\end{proof}

\begin{proposition}[11.9.13]
\label{IV.11.9.13}
Sean las notaciones las de \sref{IV.11.9.1} y \sref{IV.11.9.5}, y supongamos que $\mathrm{X}$ es un preesquema sobre un cuerpo $k$ y que, poniendo $\mathrm{S}=\operatorname{Spec}(k)$, se tiene $\mathrm{X}'=\mathrm{X}\times_\mathrm{S}\mathrm{S}'$, donde $\mathrm{S}'$ es un $k$-preesquema arbitrario.
Entonces, si la familia $(u_\lambda)_{\lambda\in\mathrm{L}}$ es separante, también lo es $(u'_\lambda)$.
\end{proposition}

\begin{proof}
Podemos limitarnos al caso en que $\mathrm{S}'=\operatorname{Spec}(\mathrm{A}')$ es afín.
Si $\mathrm{A}'$ es una $k$-álgebra de tipo finito, el morfismo $g:\mathrm{X}'\to\mathrm{X}$ es plano y de presentación finita; estamos, por tanto, bajo las hipótesis de \sref{IV.11.9.10},~(ii),~$b$).
En el caso general, se considera $\mathrm{A}'$ como límite inductivo de sus sub-$k$-álgebras $\mathrm{A}'_\alpha$ de tipo finito y se aplica a cada $\mathrm{A}'_\alpha$ el resultado de \sref{IV.11.9.10},~(ii),~$b$); se concluye entonces mediante el lema \sref{IV.11.9.12}, pues los $\sh{G}_\lambda$ son planos sobre $\mathrm{S}$.
\end{proof}

\begin{env}[11.9.14]
\label{IV.11.9.14}
Conservemos las notaciones de \sref{IV.11.9.1} y \sref{IV.11.9.5}, y supongamos que $\mathrm{X}$ es un $\mathrm{S}$-preesquema.
Si, para todo cambio de base $g:\mathrm{X}\times_\mathrm{S}\mathrm{S}'\to\mathrm{X}$, donde $\mathrm{S}'\to\mathrm{S}$ es un morfismo arbitrario, la familia correspondiente $(u'_\lambda)$ es separante, se dice que la familia $(u_\lambda)$ es \emph{universalmente separante con respecto a} $\mathrm{S}$.
Cuando la familia $(u_\lambda)$ se reduce a un solo elemento $u$, se dice asimismo que $u$ es \emph{universalmente inyectivo} con respecto a $\mathrm{S}$.
Está claro entonces que, para todo morfismo $h:\mathrm{S}'\to\mathrm{S}$, la familia correspondiente $(u'_\lambda)$ es universalmente separante con respecto a $\mathrm{S}'$; recíprocamente, si $h$ es \emph{fielmente plano} y si $(u'_\lambda)$ es universalmente separante con respecto a $\mathrm{S}'$, entonces $(u_\lambda)$ es universalmente separante con respecto a $\mathrm{S}$, como se deduce inmediatamente de \sref{IV.11.9.10},~(i) y del hecho de que, para todo morfismo $\mathrm{S}''\to\mathrm{S}$, el morfismo correspondiente $\mathrm{S}''\times_\mathrm{S}\mathrm{S}'\to\mathrm{S}''$ es fielmente plano.
\end{env}

\begin{proposition}[11.9.15]
\label{IV.11.9.15}
Sean las notaciones las de \sref{IV.11.9.1}, y supongamos que $\mathrm{X}$ es un $\mathrm{S}$-preesquema, siendo los $\sh{G}_\lambda$ planos sobre $\mathrm{S}$.
Sea $\mathrm{S}_0$ un subpreesquema cerrado de $\mathrm{S}$ definido por un ideal cuasicoherente nilpotente $\sh{J}$ de $\sh{O}_\mathrm{S}$, tal que los $(\sh{O}_\mathrm{S}/\sh{J})$-módulos $\sh{J}^k/\sh{J}^{k+1}$ sean todos localmente libres.
\oldpage[IV-3]{167}
Sea $(u_{\lambda 0})$ la familia de homomorfismos obtenida de $(u_\lambda)$ por el cambio de base $\mathrm{X}\leftarrow\mathrm{X}_0=\mathrm{X}\times_\mathrm{S}\mathrm{S}_0$.
Entonces, si la familia $(u_{\lambda 0})$ es separante (resp.\ universalmente separante con respecto a $\mathrm{S}_0$), la familia $(u_\lambda)$ es separante (resp.\ universalmente separante con respecto a $\mathrm{S}$).
\end{proposition}

\begin{proof}
Observemos que, si $\mathrm{S}'\to\mathrm{S}$ es un cambio de base arbitrario y $\mathrm{S}'_0=\mathrm{S}'\times_\mathrm{S}\mathrm{S}_0$, entonces $\mathrm{S}'_0$ es un subpreesquema cerrado de $\mathrm{S}'$ definido por un ideal cuasicoherente nilpotente $\sh{J}'$ de $\sh{O}_{\mathrm{S}'}$, tal que $\sh{J}'^k/\sh{J}'^{k+1}$ es un $(\sh{O}_{\mathrm{S}'}/\sh{J}')$-módulo localmente libre para todo $k$ \sref{IV.2.1.8},~(i); como, además, los $\sh{G}'_\lambda=\sh{G}_\lambda\otimes_{\sh{O}_\mathrm{S}}\sh{O}_{\mathrm{S}'}$ son planos sobre $\mathrm{S}'$, se ve que la afirmación relativa a las familias universalmente separantes es consecuencia de la relativa a las familias separantes.
Para probar esta última, se puede reducir, mediante \sref{IV.11.9.3}, al caso en que $\mathrm{S}=\operatorname{Spec}(\mathrm{A})$ y $\mathrm{X}=\operatorname{Spec}(\mathrm{B})$ son afines, $\mathrm{Z}_\lambda=\mathrm{X}$ para todo $\lambda$, $\mathscr{F}=\widetilde{\mathrm{M}}$ y $\sh{G}_\lambda=\widetilde{\mathrm{N}}_\lambda$, donde $\mathrm{M}$ y los $\mathrm{N}_\lambda$ son $\mathrm{B}$-módulos, siendo los $\mathrm{N}_\lambda$ planos como $\mathrm{A}$-módulos.
Además, como la cuestión es local sobre $\mathrm{X}$ y $\mathrm{S}$, basta probar que, si $t\in\mathrm{M}$ es tal que $u_\lambda(t)=0$ para todo $\lambda$, entonces $t=0$.
Se tiene $\sh{J}=\widetilde{\mathfrak{J}}$, donde $\mathfrak{J}$ es un ideal nilpotente de $\mathrm{A}$ tal que los $\mathfrak{J}^k/\mathfrak{J}^{k+1}$ son $(\mathrm{A}/\mathfrak{J})$-módulos libres y, por hipótesis, los homomorfismos $u_{\lambda 0}:\mathrm{M}/\mathfrak{J}\mathrm{M}\to\mathrm{N}_\lambda/\mathfrak{J}\mathrm{N}_\lambda$ forman una familia separante.
Supongamos que $\mathfrak{J}^{n+1}=0$ ($n$ un entero $\geq 0$) y razonemos por inducción sobre $n$, siendo la afirmación trivial para $n=0$.
Si $\bar{t}$ es la clase de $t$ en $\mathrm{M}/\mathfrak{J}^n\mathrm{M}$, la clase de $u_\lambda(t)$ en $\mathrm{N}_\lambda/\mathfrak{J}^n\mathrm{N}_\lambda$ es nula para todo $\lambda$; de la hipótesis de inducción se deduce que $\bar{t}=0$, es decir, que $t\in\mathfrak{J}^n\mathrm{M}$.
Ahora bien, $\mathfrak{J}^n=\mathfrak{J}^n/\mathfrak{J}^{n+1}$ es un $(\mathrm{A}/\mathfrak{J})$-módulo libre; si $(e_\alpha)$ es una base de este módulo, se puede escribir, por tanto, $t=\sum_\alpha e_\alpha t_\alpha^0$, con $t_\alpha^0\in\mathrm{M}/\mathfrak{J}\mathrm{M}$, siendo nulos todos salvo para un número finito de índices.
Por otra parte, como $\mathrm{N}_\lambda$ es un $\mathrm{A}$-módulo plano, $\mathfrak{J}^n\otimes_{(\mathrm{A}/\mathfrak{J})}(\mathrm{N}_\lambda/\mathfrak{J}\mathrm{N}_\lambda)$ se identifica con $\mathfrak{J}^n\mathrm{N}_\lambda$ y se puede escribir, por tanto, $u_\lambda(t)=\sum_\alpha e_\alpha u_{\lambda 0}(t_\alpha^0)=\sum_\alpha e_\alpha\otimes u_{\lambda 0}(t_\alpha^0)$.
Como, por hipótesis, $u_\lambda(t)=0$, se deduce que $u_{\lambda 0}(t_\alpha^0)=0$ para todo $\alpha$ y todo $\lambda$; de donde $t_\alpha^0=0$ para todo $\alpha$ y, por tanto, $t=0$.
\end{proof}

\begin{theorem}[11.9.16]
\label{IV.11.9.16}
Sean las notaciones las de \sref{IV.11.9.1}, y supongamos que $\mathrm{X}$ es un $\mathrm{S}$-preesquema localmente noetheriano, que $\mathscr{F}$ es un $\sh{O}_\mathrm{X}$-módulo coherente y que los $\sh{G}_\lambda$ son planos sobre $\mathrm{S}$.
Para todo $s\in\mathrm{S}$, sea $((u_\lambda)_s)_{\lambda\in\mathrm{L}}$ la familia obtenida de $(u_\lambda)$ por el cambio de base $\mathrm{X}_s=\mathrm{X}\times_\mathrm{S}\operatorname{Spec}(k(s))\to\mathrm{X}$.
Para que la familia $(u_\lambda)$ sea universalmente separante con respecto a $\mathrm{S}$, es necesario y suficiente que, para todo $s\in\mathrm{S}$, la familia $((u_\lambda)_s)$ sea separante.
\end{theorem}

\begin{proof}
La necesidad de la condición se deduce inmediatamente de las definiciones.
Recíprocamente, supongamos satisfecha la condición y probemos primero que la familia $(u_\lambda)$ es separante.
Mediante \sref{IV.11.9.3}, podemos reducirnos al caso en que $\mathrm{S}=\operatorname{Spec}(\mathrm{A})$ y $\mathrm{X}=\operatorname{Spec}(\mathrm{B})$ son afines, $\mathrm{Z}_\lambda=\mathrm{X}$ para todo $\lambda$, $\mathscr{F}=\widetilde{\mathrm{M}}$ y $\sh{G}_\lambda=\widetilde{\mathrm{N}}_\lambda$, donde $\mathrm{B}$ es noetheriano, $\mathrm{M}$ es un $\mathrm{B}$-módulo de tipo finito y los $\mathrm{N}_\lambda$ son $\mathrm{A}$-módulos planos, y limitarnos a probar que, si $t\in\mathrm{M}$ es tal que $u_\lambda(t)=0$ para todo $\lambda$, entonces $t=0$.
Para demostrar que $t=0$, basta probar que la imagen $t_\mathfrak{p}$ de $t$ en $\mathrm{M}_\mathfrak{p}$ es nula para todo ideal maximal $\mathfrak{p}$ de $\mathrm{B}$ (Bourbaki, \emph{Alg.\ comm.}, cap.~II, \S~3, n.º~3, cor.~1 del teor.~1).
Podemos, pues, limitarnos a probar que la intersección
\oldpage[IV-3]{168}
de los núcleos de los $(u_\lambda)_\mathfrak{p}:\mathrm{M}_\mathfrak{p}\to(\mathrm{N}_\lambda)_\mathfrak{p}$ deducidos de $(u_\lambda)$ por el cambio de base $\operatorname{Spec}(\mathrm{B}_\mathfrak{p})\to\operatorname{Spec}(\mathrm{B})$ se reduce a $0$.
Dicho de otro modo, nos reducimos al caso en que $\mathrm{B}$ es un anillo local noetheriano; considerando además el ideal primo de $\mathrm{A}$ que es la imagen inversa del ideal maximal de $\mathrm{B}$, podemos suponer también que $\mathrm{A}$ es un anillo local con ideal maximal $\mathfrak{m}$.
Entonces, como $\mathfrak{m}\mathrm{B}$ está contenido en el ideal maximal de $\mathrm{B}$ y $\mathrm{M}$ es un $\mathrm{B}$-módulo de tipo finito, la intersección de los $\mathfrak{m}^{n}\mathrm{M}$ se reduce a $0$ \sref[0]{0.7.3.5}; por tanto, basta probar que, para todo $n$, la imagen de $t$ en $\mathrm{M}/\mathfrak{m}^{n+1}\mathrm{M}$ es nula.
Basta, pues, demostrar que la familia deducida de $(u_\lambda)$ por el cambio de base $\operatorname{Spec}(\mathrm{B}/\mathfrak{m}^{n+1}\mathrm{B})\to\operatorname{Spec}(\mathrm{B})$ es separante, lo que significa que podemos limitarnos al caso en que $\mathrm{A}$ es un anillo local cuyo ideal maximal $\mathfrak{m}$ es nilpotente.
Pero entonces los $\mathfrak{m}^k/\mathfrak{m}^{k+1}$ son $(\mathrm{A}/\mathfrak{m})$-módulos libres y, en virtud de la hipótesis sobre los $(u_\lambda)_s$, estamos precisamente en las condiciones de aplicación de \sref{IV.11.9.15}, de donde se sigue la conclusión anunciada.

Sea ahora $h:\mathrm{S}'\to\mathrm{S}$ un morfismo de cambio de base, y sea $(u'_\lambda)$ la familia obtenida de $(u_\lambda)$ por el cambio de base $h':\mathrm{X}\times_\mathrm{S}\mathrm{S}'\to\mathrm{X}$; probemos que $(u'_\lambda)$ también es separante.
Supongamos primero que $h$ es \emph{localmente de tipo finito}; entonces lo mismo ocurre con $h'$ y, por consiguiente, $\mathrm{X}'=\mathrm{X}\times_\mathrm{S}\mathrm{S}'$ es localmente noetheriano. Además, si $s'\in\mathrm{S}'$ está sobre $s\in\mathrm{S}$, se deduce de \sref{IV.11.9.13}, aplicado a $\mathrm{X}_s$ y a $\mathrm{X}_{s'}=\mathrm{X}_s\otimes_{k(s)}k(s')$, que la familia $((u'_\lambda)_{s'})$ es separante para todo $s'\in\mathrm{S}'$; podemos concluir, por tanto, de la primera parte de la demostración que $(u'_\lambda)$ es separante en este caso.

Por último, si $h$ es arbitrario, podemos limitarnos evidentemente al caso en que $\mathrm{S}=\operatorname{Spec}(\mathrm{A})$ y $\mathrm{S}'=\operatorname{Spec}(\mathrm{A}')$ son afines, y considerar $\mathrm{A}'$ como límite inductivo de sus sub-$\mathrm{A}$-álgebras de tipo finito.
Como los $\sh{G}_\lambda$ son planos sobre $\mathrm{S}$, basta aplicar lo anterior y el lema \sref{IV.11.9.12} para concluir la demostración.
\end{proof}

\begin{proposition}[11.9.17]
\label{IV.11.9.17}
Sea $f:\mathrm{X}\to\mathrm{S}$ un morfismo localmente de presentación finita, sea $\mathscr{F}$ un $\sh{O}_\mathrm{X}$-módulo cuasicoherente de presentación finita y $f$-plano, sea $\mathrm{U}$ un abierto de $\mathrm{X}$, sea $j:\mathrm{U}\to\mathrm{X}$ la inyección canónica y sea $u:\mathscr{F}\to j_*(j^*(\mathscr{F}))$ el homomorfismo canónico \sref[0]{0.4.4.3.2}.
Para todo $s\in\mathrm{S}$, sean $\mathrm{X}_s$ la fibra $\mathrm{X}\times_\mathrm{S}\operatorname{Spec}(k(s))$, $\mathrm{U}_s$ el abierto $\mathrm{U}\cap\mathrm{X}_s$ de $\mathrm{X}_s$, $\mathscr{F}_s=\mathscr{F}\otimes_{\sh{O}_\mathrm{S}}k(s)$, $j_s:\mathrm{U}_s\to\mathrm{X}_s$ la inyección canónica y $u_s:\mathscr{F}_s\to(j_s)_*((j_s)^*(\mathscr{F}_s))$ el homomorfismo canónico.
Para que $u$ sea universalmente inyectivo con respecto a $\mathrm{S}$, es necesario y suficiente que $u_s$ sea inyectivo para todo $s\in\mathrm{S}$.
\end{proposition}

\begin{proof}
Solo queda probar la suficiencia de la condición.
Cuando $\mathrm{X}$ es localmente noetheriano, la proposición es un corolario inmediato de \sref{IV.11.9.16}.
Nos reduciremos ahora a este caso en dos etapas, limitándonos, como podemos hacer, al caso en que $\mathrm{S}=\operatorname{Spec}(\mathrm{A})$ y $\mathrm{X}$ son afines.
\medskip

\emph{A) Caso en que $\mathrm{U}$ es cuasicompacto.} --- Utilizaremos el lema siguiente:

\begin{lemma}[11.9.17.1]
\label{IV.11.9.17.1}
Bajo las hipótesis generales de \sref{IV.11.9.17}, y suponiendo además que $\mathrm{S}$ y $\mathrm{X}$ son afines y que $\mathrm{U}$ es cuasicompacto, el conjunto $\mathrm{E}$ de los $s\in\mathrm{S}$ tales que $u_s$ es inyectivo es constructible.
\end{lemma}

\begin{proof}
En efecto, las fibras $\mathrm{X}_s$ son preesquemas localmente noetherianos; de donde $\mathrm{E}$ también puede, en virtud de \sref{IV.5.10.2}, definirse como el conjunto de los $s\in\mathrm{S}$ tales que $\operatorname{Ass}(\mathscr{F}_s)\subset\mathrm{U}_s$.
Observemos además que $\mathrm{U}$, por ser cuasicompacto en un esquema afín, es constructible. Así, la verificación de la condición \sref{IV.9.2.1},~(i) se deduce inmediatamente de \sref{IV.4.2.7}; por otra parte, la
\oldpage[IV-3]{169}
verificación de \sref{IV.9.2.1},~(ii) se efectúa fácilmente utilizando el estudio de los ciclos primos asociados en un entorno de la fibra genérica \sref{IV.9.8.3}, así como \sref{IV.9.5.2} y \sref{IV.9.5.3}.
\end{proof}

Establecido este lema, podemos suponer, en virtud de \sref{IV.8.9.1} y \sref{IV.8.2.11}, que existen un subanillo noetheriano $\mathrm{A}_0$ de $\mathrm{A}$, un preesquema $\mathrm{X}_0$ de tipo finito sobre $\mathrm{S}_0=\operatorname{Spec}(\mathrm{A}_0)$, un abierto $\mathrm{U}_0$ de $\mathrm{X}_0$ y un $\sh{O}_{\mathrm{X}_0}$-módulo coherente $\mathscr{F}_0$, tales que $\mathrm{X}=\mathrm{X}_0\times_{\mathrm{S}_0}\mathrm{S}$ y, si $p_0:\mathrm{X}\to\mathrm{X}_0$ es la proyección canónica, $\mathrm{U}=p_0^{-1}(\mathrm{U}_0)$ y $\mathscr{F}=p_0^*(\mathscr{F}_0)$.
Sea $(\mathrm{A}_\alpha)$ la familia filtrante de los subanillos de $\mathrm{A}$ que son $\mathrm{A}_0$-álgebras de tipo finito, y pongamos $\mathrm{S}_\alpha=\operatorname{Spec}(\mathrm{A}_\alpha)$, $\mathrm{X}_\alpha=\mathrm{X}_0\times_{\mathrm{S}_0}\mathrm{S}_\alpha$, $\mathrm{U}_\alpha=\mathrm{U}_0\times_{\mathrm{S}_0}\mathrm{S}_\alpha$ y $\mathscr{F}_\alpha=\mathscr{F}_0\otimes_{\mathrm{S}_0}\mathrm{S}_\alpha$; sea $u_\alpha$ el homomorfismo canónico relativo a $\mathscr{F}_\alpha$ y $\mathrm{U}_\alpha$, definido como en \sref{IV.11.9.17}.
Para todo $s\in\mathrm{S}$, la hipótesis de que $u_s$ es inyectivo implica que $(u_\alpha)_{s_\alpha}$ es inyectivo \sref{IV.11.9.10},~(i), donde $s_\alpha=q_\alpha(s)$ es la imagen de $s$ por el morfismo $q_\alpha:\mathrm{S}\to\mathrm{S}_\alpha$.
Si $\mathrm{E}_\alpha$ es el conjunto de los $s_\alpha\in\mathrm{S}_\alpha$ tales que $(u_\alpha)_{s_\alpha}$ es inyectivo, se tiene, por tanto, $\mathrm{S}=q_\alpha^{-1}(\mathrm{E}_\alpha)$, y los $\mathrm{E}_\alpha$ forman un sistema proyectivo de conjuntos.
Pero el lema \sref{IV.11.9.17.1}, aplicado a $u_\alpha$, muestra que $\mathrm{E}_\alpha$ es constructible en $\mathrm{S}_\alpha$; se deduce entonces de \sref{IV.8.3.4} que existe un índice $\alpha$ tal que $\mathrm{E}_\alpha=\mathrm{S}_\alpha$.
Pero entonces, como $\mathrm{X}_\alpha$ es noetheriano, $u_\alpha$ es universalmente inyectivo en virtud de \sref{IV.11.9.16}; de donde $u$ también es universalmente inyectivo.
\medskip

\emph{B) Caso general.} --- El abierto $\mathrm{U}$ es una unión filtrante creciente de abiertos cuasicompactos $\mathrm{U}_\lambda$; si $j_\lambda:\mathrm{U}_\lambda\to\mathrm{X}$ es la inyección canónica y $u_\lambda:\mathscr{F}\to(j_\lambda)_*(j_\lambda^*(\mathscr{F}))$ el homomorfismo canónico correspondiente, se deduce del lema \sref{IV.11.9.17.1}, aplicado a $u_\lambda$, que el conjunto $\mathrm{E}_\lambda$ de los $s\in\mathrm{S}$ tales que $(u_\lambda)_s$ es inyectivo es constructible.
Por otra parte, para todo $s\in\mathrm{S}$, decir que $u_s$ (resp.\ $(u_\lambda)_s$) es inyectivo significa que $\operatorname{Ass}(\mathscr{F}_s)\subset\mathrm{U}\cap\mathrm{X}_s$ (resp.\ $\operatorname{Ass}(\mathscr{F}_s)\subset\mathrm{U}_\lambda\cap\mathrm{X}_s$) \sref{IV.5.10.2}.
Como $\operatorname{Ass}(\mathscr{F}_s)$ es finito \sref{IV.3.1.6}, decir que $u_s$ es inyectivo significa, por tanto, que existe $\lambda$ tal que $s\in\mathrm{E}_\lambda$.
La hipótesis significa, pues, que $\mathrm{S}=\bigcup_\lambda\mathrm{E}_\lambda$.
En virtud de \sref{IV.1.9.9}, existe un índice $\lambda$ tal que $\mathrm{S}=\mathrm{E}_\lambda$; de la primera parte del razonamiento se deduce entonces que $u_\lambda$ es universalmente inyectivo.
Se sigue de la factorización de $u_\lambda$:
\[
  \mathscr{F}\xrightarrow{u}\,j_*(\mathscr{F}|\mathrm{U})\to(j_\lambda)_*(\mathscr{F}|\mathrm{U}_\lambda)
\]
que $u$ también es universalmente inyectivo.
\end{proof}

\begin{remarks}[11.9.18]
\label{IV.11.9.18}
Sea $\mathrm{X}$ un preesquema, sean $\mathscr{F}$ y $\sh{G}$ dos $\sh{O}_\mathrm{X}$-módulos cuasicoherentes de presentación finita, suponiéndose además que $\sh{G}$ es \emph{localmente libre}, y sea $u:\mathscr{F}\to\sh{G}$ un homomorfismo.
Las condiciones siguientes son equivalentes:
\begin{enumerate}[label=\emph{\alph*)}]
  \item para todo morfismo $g:\mathrm{X}'\to\mathrm{X}$, el homomorfismo $g^*(u):g^*(\mathscr{F})\to g^*(\sh{G})$ es inyectivo;
  \item para todo $x\in\mathrm{X}$, el homomorfismo $u\otimes 1_{k(x)}:\mathscr{F}\otimes_{\sh{O}_{\mathrm{X}}}k(x)\to\sh{G}\otimes_{\sh{O}_\mathrm{X}}k(x)$ es inyectivo;
  \item para todo $x\in\mathrm{X}$, existe un entorno $\mathrm{U}$ de $x$ tal que $u|\mathrm{U}:\mathscr{F}|\mathrm{U}\to\sh{G}|\mathrm{U}$ es un isomorfismo de $\mathscr{F}|\mathrm{U}$ sobre un factor directo de $\sh{G}|\mathrm{U}$.
\end{enumerate}

Está claro que $c)$ implica $a)$ y que $a)$ implica $b)$.
El hecho de que $b)$ implica $c)$ se deduce de \sref[0]{0.19.1.12}, \sref[0]{0.5.2.5} y \sref[0]{0.5.5.5}.

Cuando se cumplen las condiciones equivalentes anteriores, se dice que $u$ es \emph{universalmente inyectivo}.
\end{remarks}

\oldpage[IV-3]{170}
\subsection{Familias esquemáticamente dominantes de morfismos y familias esquemáticamente densas de subpreesquemas}
\label{subsection:IV.11.10}

\begin{proposition}[11.10.1]
\label{IV.11.10.1}
Sea $\mathrm{X}$ un preesquema, $(Z_\lambda)_{\lambda\in\mathrm{L}}$ una familia de preesquemas y, para cada $\lambda\in\mathrm{L}$, sea $f_\lambda=(\psi_\lambda,\theta_\lambda):\mathrm{Z}_\lambda\to\mathrm{X}$ un morfismo.
Las condiciones siguientes son equivalentes:
\begin{enumerate}[label=\emph{\alph*)}]
  \item La familia de homomorfismos $\theta_\lambda:\sh{O}_\mathrm{X}\to(f_\lambda)_*(\sh{O}_{\mathrm{Z}_\lambda})$ es separante (en otras palabras \sref{IV.11.9.1}, la intersección de los núcleos de los $\theta_\lambda$ es \emph{cero}).
  \item Para todo abierto $\mathrm{U}$ de $\mathrm{X}$, toda sección $t$ de $\sh{O}_\mathrm{X}$ sobre $\mathrm{U}$ cuyas imágenes por todos los homomorfismos canónicos
\end{enumerate}
\[
  (\theta_\lambda)_\mathrm{U}:\Gamma(\mathrm{U},\sh{O}_\mathrm{X})\to\Gamma(f_\lambda^{-1}(\mathrm{U}),\sh{O}_{\mathrm{Z}_\lambda})
\label{IV.11.10.1.1}
\]
son cero es ella misma cero.
\begin{enumerate}[label=\emph{\alph*)}]
\setcounter{enumi}{2}
  \item Para todo abierto $\mathrm{U}$ de $\mathrm{X}$ y todo subpreesquema cerrado $\mathrm{Y}$ de $\mathrm{U}$ tal que, para cada $\lambda\in\mathrm{L}$, existe una factorización
\end{enumerate}
\[
  f_\lambda^{-1}(\mathrm{U})\xrightarrow{g_\lambda}\mathrm{Y}\xrightarrow{j}\mathrm{U}
\label{IV.11.10.1.2}
\]
de la restricción $f_\lambda^{-1}(\mathrm{U})\to\mathrm{U}$ de $f_\lambda$ (donde $j$ es la inmersión canónica), se tiene $\mathrm{Y}=\mathrm{U}$.

\emph{Si además $\mathrm{X}$ es un $\mathrm{S}$-preesquema, estas condiciones también son equivalentes a las siguientes:}
\begin{enumerate}[label=\emph{\alph*)}]
\setcounter{enumi}{3}
  \item Para todo morfismo separado $g:\mathrm{X}'\to\mathrm{S}$ y todo par $u$, $v$ de $\mathrm{S}$-morfismos de un abierto $\mathrm{U}$ de $\mathrm{X}$ a $\mathrm{X}'$, tales que, para cada $\lambda$, los compuestos de $u$ y $v$ con la restricción de $f_\lambda$ a $f_\lambda^{-1}(\mathrm{U})$ sean iguales, se tiene $u=v$.
\end{enumerate}
\end{proposition}

\begin{proof}
La equivalencia de $a)$ y $b)$ se deduce de las definiciones.
Para ver que $b)$ implica $c)$, basta considerar el ideal cuasicoherente $\sh{J}$ de $\sh{O}_\mathrm{U}$ que define $\mathrm{Y}$ y observar que, para todo abierto $\mathrm{V}\subset\mathrm{U}$, la hipótesis implica que el homomorfismo $(\theta_\lambda)_\mathrm{V}$ se factoriza como
\[
  \Gamma(\mathrm{V},\sh{O}_\mathrm{U})\to\Gamma(\mathrm{V},\sh{O}_\mathrm{U}/\sh{J})\to\Gamma(f_\lambda^{-1}(\mathrm{V}),\sh{O}_{\mathrm{Z}_\lambda}).
\]
Se concluye que toda sección $t$ de $\sh{J}$ sobre $\mathrm{V}$ tiene imagen $0$ en todos los $\Gamma(f_\lambda^{-1}(\mathrm{V}),\sh{O}_{\mathrm{Z}_\lambda})$, de donde, en virtud de $b)$, $\sh{J}=0$ y $\mathrm{Y}=\mathrm{U}$.
Recíprocamente, si se satisface $c)$, para demostrar $b)$ basta aplicar $c)$ al subpreesquema cerrado $\mathrm{Y}$ de $\mathrm{U}$ definido por el ideal $t\sh{O}_\mathrm{U}$: la hipótesis de que las imágenes de $t$ por los $(\theta_\lambda)_\mathrm{U}$ son todas cero implica que se tienen las factorizaciones \sref{IV.11.10.1.2} para cada $\lambda$ \sref[I]{I.4.1.9}.
Para demostrar que $c)$ implica $d)$, basta aplicar $c)$ al subpreesquema cerrado $\mathrm{Y}$ de $\mathrm{U}$ que es la imagen inversa, por $(u,v)_\mathrm{S}$, de la diagonal de $\mathrm{X}'$ en $\mathrm{X}'\times_\mathrm{S}\mathrm{X}'$, y utilizar \sref[I]{I.4.4.1}.
Recíprocamente, se deduce $b)$ de $d)$ considerando el $\mathrm{S}$-esquema $\mathrm{X}'=\mathrm{S}[\mathrm{T}]=\mathrm{S}\otimes_\mathbf{Z}\mathbf{Z}[\mathrm{T}]$ ($\mathrm{T}$ indeterminada) y observando que las secciones de $\sh{O}_\mathrm{U}$ sobre $\mathrm{U}$ corresponden biyectivamente a los $\mathrm{S}$-morfismos $\mathrm{U}\to\mathrm{S}[\mathrm{T}]$ \sref[I]{I.3.3.15}; decir que dos secciones de $\sh{O}_\mathrm{U}$ sobre $\mathrm{U}$ tienen las mismas imágenes por todos los $(\theta_\lambda)_\mathrm{U}$ equivale a decir que los compuestos de los dos morfismos correspondientes con todos los morfismos $f_\lambda^{-1}(\mathrm{U})\to\mathrm{U}$ son iguales.
\end{proof}

\oldpage[IV-3]{171}
\begin{definition}[11.10.2]
\label{IV.11.10.2}
Cuando se satisfacen las condiciones equivalentes de \sref{IV.11.10.1}, se dice que la familia $(f_\lambda)$ es \emph{esquemáticamente dominante}.
Cuando los $f_\lambda$ son las inmersiones canónicas en $\mathrm{X}$ de una familia $(\mathrm{Z}_\lambda)$ de subpreesquemas de $\mathrm{X}$, se dice también que la familia $(\mathrm{Z}_\lambda)$ es \emph{esquemáticamente densa}.
\end{definition}

\begin{remarks}[11.10.3]
\label{IV.11.10.3}
\medskip\noindent
\begin{enumerate}[label=(\roman*)]
  \item La noción de familia esquemáticamente dominante sobre $\mathrm{X}$, como resulta, por ejemplo, de una condición de la forma $b)$ de \sref{IV.11.10.1}, es \emph{local} sobre $\mathrm{X}$: si $(\mathrm{W}_\alpha)$ es un recubrimiento abierto de $\mathrm{X}$, la familia $(f_\lambda)$ es esquemáticamente dominante si y solo si lo es cada una de las familias formadas por los morfismos $f_\lambda^{-1}(\mathrm{W}_\alpha)\to\mathrm{W}_\alpha$, restricciones de los $f_\lambda$.
  \item Si $\mathrm{Z}$ es la suma de los preesquemas $\mathrm{Z}_\lambda$ y $f:\mathrm{Z}\to\mathrm{X}$ el morfismo que coincide con $f_\lambda$ sobre cada $\mathrm{Z}_\lambda$, equivale a decir que la familia $(f_\lambda)$ es esquemáticamente dominante o que lo es la familia reducida al único elemento $f$.
  \item Sea $\mathrm{M}$ un segundo conjunto de índices y, para cada $\lambda\in\mathrm{L}$, sea $(g_{\lambda\mu})_{\mu\in\mathrm{M}}$ una familia de morfismos $g_{\lambda\mu}:\mathrm{Z}'_{\lambda\mu}\to\mathrm{Z}_\lambda$; si, para cada $\lambda\in\mathrm{L}$, la familia $(g_{\lambda\mu})_{\mu\in\mathrm{M}}$ es esquemáticamente dominante, entonces equivale a decir que la familia $(f_\lambda)_{\lambda\in\mathrm{L}}$ es esquemáticamente dominante o que lo es la familia $(f_\lambda\circ g_{\lambda\mu})_{(\lambda,\mu)\in\mathrm{L}\times\mathrm{M}}$ \sref{IV.11.9.2}.
  \item Sea $f:\mathrm{Z}\to\mathrm{X}$ un morfismo tal que $f_*(\sh{O}_\mathrm{Z})$ sea un $\sh{O}_\mathrm{X}$-módulo cuasicoherente (por ejemplo, un morfismo cuasicompacto y cuasiseparado \sref{IV.1.7.4}).
  Entonces, decir que $f$ es esquemáticamente dominante significa que la \emph{imagen cerrada} de $\mathrm{Z}$ \emph{por} $f$ \sref[I]{I.9.5.3} \emph{es idéntica a} $\mathrm{X}$.
\end{enumerate}
\end{remarks}

\begin{proposition}[11.10.4]
\label{IV.11.10.4}
Si la familia de morfismos $f_\lambda:\mathrm{Z}_\lambda\to\mathrm{X}$ es esquemáticamente dominante, la unión de los $f_\lambda(\mathrm{Z}_\lambda)$ es densa en $\mathrm{X}$.
Por el contrario, si esta unión es densa en $\mathrm{X}$ y si además $\mathrm{X}$ es \emph{reducido}, la familia $(f_\lambda)$ es esquemáticamente dominante.
\end{proposition}

\begin{proof}
La primera afirmación se deduce inmediatamente de \sref{IV.11.10.1},~$b$)).
Por otra parte, si $\mathrm{X}$ es reducido y la unión de los $f_\lambda(\mathrm{Z}_\lambda)$ es densa en $\mathrm{X}$, entonces, si se tienen las factorizaciones \sref{IV.11.10.1.2} para cada $\lambda\in\mathrm{L}$, se tienen también las factorizaciones $(f_\lambda^{-1}(\mathrm{U}))_\mathrm{red}\xrightarrow{(g_\lambda)_\mathrm{red}}\mathrm{Y}_\mathrm{red}\xrightarrow{j_\mathrm{red}}\mathrm{U}$, y la hipótesis implica que el espacio subyacente de $\mathrm{Y}$ es idéntico a $\mathrm{U}$; de donde $\mathrm{Y}=\mathrm{Y}_\mathrm{red}=\mathrm{U}$, pues $\mathrm{U}$ es reducido.
\end{proof}

Los resultados de {\hyperlink{ega4.c04.subsec.11-9}{\textsection~11.9}} sobre las familias separantes se traducen en resultados sobre las familias esquemáticamente dominantes:

\begin{theorem}[11.10.5]
\label{IV.11.10.5}
Sea $(f_\lambda)$ una familia de morfismos $f_\lambda:\mathrm{Z}_\lambda\to\mathrm{X}$, sea $g:\mathrm{X}'\to\mathrm{X}$ un morfismo, y pongamos $\mathrm{Z}'_\lambda=\mathrm{Z}_\lambda\times_\mathrm{X}\mathrm{X}'$, $f'_\lambda=(f_\lambda)_{(\mathrm{X}')}:\mathrm{Z}'_\lambda\to\mathrm{X}'$.
\begin{enumerate}[label=(\roman*)]
  \item Si $g$ es fielmente plano y la familia $(f'_\lambda)$ es esquemáticamente dominante, entonces la familia $(f_\lambda)$ es esquemáticamente dominante.
  \item Supongamos que $g$ es un morfismo plano y que, además, se satisface una de las dos condiciones siguientes:
  \begin{enumerate}[label=\emph{\alph*)}]
    \item $\mathrm{L}$ es finito y los $f_\lambda$ son cuasicompactos.
    \item El morfismo $g$ es localmente de presentación finita.
  \end{enumerate}
  Entonces, si la familia $(f_\lambda)$ es esquemáticamente dominante, lo mismo ocurre con la familia $(f'_\lambda)$.
\end{enumerate}
\end{theorem}

\oldpage[IV-3]{172}
\begin{proposition}[11.10.6]
\label{IV.11.10.6}
Sean las notaciones las de \sref{IV.11.10.5}, y supongamos que $\mathrm{X}$ es un preesquema sobre un cuerpo $k$ y que, poniendo $\mathrm{S}=\operatorname{Spec}(k)$, se tiene $\mathrm{X}'=\mathrm{X}\times_\mathrm{S}\mathrm{S}'$, donde $\mathrm{S}'$ es un $k$-preesquema arbitrario.
Entonces, si la familia $(f_\lambda)$ es esquemáticamente dominante, lo mismo ocurre con $(f'_\lambda)$.
\end{proposition}

\begin{corollary}[11.10.7]
\label{IV.11.10.7}
Sea $\mathrm{X}$ un preesquema sobre un cuerpo $k$, sea $(Z_\lambda)_{\lambda\in\mathrm{L}}$ una familia de $k$-preesquemas geométricamente reducidos sobre $k$ y, para cada $\lambda\in\mathrm{L}$, sea $f_\lambda:\mathrm{Z}_\lambda\to\mathrm{X}$ un $k$-morfismo.
Sea $\mathrm{Y}$ el subpreesquema reducido de $\mathrm{X}$ cuyo espacio subyacente es la adherencia de $\bigcup_{\lambda\in\mathrm{L}}f_\lambda(\mathrm{Z}_\lambda)$.
Sea $k'$ una extensión de $k$, y sean $\mathrm{X}'$, $\mathrm{Z}'_\lambda$, $f'_\lambda:\mathrm{Z}'_\lambda\to\mathrm{X}'$ los preesquemas y morfismos deducidos de $\mathrm{X}$, $\mathrm{Z}_\lambda$ y $f_\lambda$ por extensión de la base a $k'$.
Entonces, si $\mathrm{Y}'$ es el subpreesquema reducido de $\mathrm{X}'$ cuyo espacio subyacente es la adherencia de $\bigcup_{\lambda\in\mathrm{L}}f'_\lambda(\mathrm{Z}'_\lambda)$, se tiene $\mathrm{Y}'=\mathrm{Y}\otimes_k k'$.
En particular, $\mathrm{Y}$ es geométricamente reducido sobre $k$.
\end{corollary}

\begin{proof}
Como los $\mathrm{Z}_\lambda$ son reducidos, los morfismos $f_\lambda$ se factorizan como $\mathrm{Z}_\lambda\xrightarrow{g_\lambda}\mathrm{Y}\xrightarrow{j}\mathrm{X}$, donde $j$ es la inmersión canónica \sref[I]{I.5.2.2}.
Se sigue entonces de \sref{IV.11.10.4} que $(g_\lambda)$ es una familia esquemáticamente dominante.
Sea $\mathrm{Y}'_1=\mathrm{Y}\otimes_k k'$, un subpreesquema cerrado de $\mathrm{X}'$, y sea $g'_\lambda$ el morfismo deducido de $g_\lambda$ por extensión de la base a $k'$, de modo que $f'_\lambda$ se factoriza como $\mathrm{Z}'_\lambda\xrightarrow{g'_\lambda}\mathrm{Y}'_1\xrightarrow{j'}\mathrm{X}'$, donde $j'$ es la inmersión canónica.
Se deduce de \sref{IV.11.10.6} que la familia $(g'_\lambda)$ es esquemáticamente dominante.
Pero, por hipótesis, los $\mathrm{Z}'_\lambda$ son \emph{reducidos}, de donde, por \sref[I]{I.5.2.2}, los $g'_\lambda$ se factorizan como $\mathrm{Z}'_\lambda\xrightarrow{h'_\lambda}\mathrm{Y}'\to\mathrm{Y}'_1$; se concluye, por tanto, de \sref{IV.11.10.2} que $\mathrm{Y}'=\mathrm{Y}'_1$ y $h'_\lambda=g'_\lambda$, lo que demuestra el corolario.
\end{proof}

\begin{definition}[11.10.8]
\label{IV.11.10.8}
Sean las notaciones las de \sref{IV.11.10.5}, y supongamos que $\mathrm{X}$ es un $\mathrm{S}$-preesquema.
Se dice que la familia $(f_\lambda)$ es \emph{universalmente esquemáticamente dominante} (relativamente a $\mathrm{S}$) si, para todo cambio de base $\mathrm{S}'\to\mathrm{S}$, la familia $(f'_\lambda)$ correspondiente al cambio de base $\mathrm{X}'=\mathrm{X}\times_\mathrm{S}\mathrm{S}'\to\mathrm{X}$ es esquemáticamente dominante.
\end{definition}

Cuando $\mathrm{S}$ es el espectro de un cuerpo, \sref{IV.11.10.6} muestra que toda familia esquemáticamente dominante es universalmente esquemáticamente dominante (relativamente a $\mathrm{S}$).

Cuando los $f_\lambda$ son inmersiones canónicas $\mathrm{Z}_\lambda\to\mathrm{X}$, se dice también que la familia de subpreesquemas $\mathrm{Z}_\lambda$ es \emph{universalmente esquemáticamente densa en} $\mathrm{X}$ (relativamente a $\mathrm{S}$), en vez de decir que la familia $(f_\lambda)$ es universalmente esquemáticamente dominante (relativamente a $\mathrm{S}$).

\begin{theorem}[11.10.9]
\label{IV.11.10.9}
Sean las notaciones las de \sref{IV.11.10.5}, y supongamos que $\mathrm{X}$ es un $\mathrm{S}$-preesquema localmente noetheriano y que todos los $\mathrm{Z}_\lambda$ son $\mathrm{S}$-planos.
Para cada $s\in\mathrm{S}$, pongamos $\mathrm{X}_s=\mathrm{X}\times_\mathrm{S}\operatorname{Spec}(k(s))$, $(\mathrm{Z}_\lambda)_s=\mathrm{Z}_\lambda\times_\mathrm{S}\operatorname{Spec}(k(s))$, $(f_\lambda)_s=f_\lambda\times 1:(\mathrm{Z}_\lambda)_s\to\mathrm{X}_s$.
Para que la familia $(f_\lambda)$ sea universalmente esquemáticamente dominante relativamente a $\mathrm{S}$, es necesario y suficiente que, para todo $s\in\mathrm{S}$, la familia $((f_\lambda)_s)$ sea esquemáticamente dominante.
\end{theorem}

\begin{proposition}[11.10.10]
\label{IV.11.10.10}
Sea $f:\mathrm{X}\to\mathrm{S}$ un morfismo plano localmente de presentación finita, y sea $\mathrm{U}$ un abierto de $\mathrm{X}$.
Para que $\mathrm{U}$ sea universalmente esquemáticamente denso en $\mathrm{X}$ relativamente a $\mathrm{S}$, es necesario y suficiente que, para todo $s\in\mathrm{S}$, $\mathrm{U}_s=\mathrm{U}\cap\mathrm{X}_s$ sea esquemáticamente denso en $\mathrm{X}_s$ (o, lo que equivale a lo mismo, que se tenga $\operatorname{Ass}(\sh{O}_{\mathrm{X}_s})\subset\mathrm{U}_s$).
\end{proposition}

\oldpage[IV-3]{173}

\section{Estudio de las fibras de los morfismos planos de presentación finita}
\label{section:IV.12}

\setcounter{subsection}{-1}
\subsection{Introducción}
\label{subsection:IV.12.0}

\emph{En todo este parágrafo utilizaremos la notación general de} \sref{IV.9.4.1}.

\begin{env}[12.0.1]
\label{IV.12.0.1}
Dado un morfismo \emph{localmente de presentación finita} $f:X\to Y$, hemos visto \sref{IV.9.9} que, para ciertas propiedades \emph{locales} $\mathbf{P}$ de preesquemas sobre cuerpos o de Módulos sobre estos preesquemas, el conjunto $\mathrm{E}$ de los $x\in\mathrm{X}$ tales que la propiedad $\mathbf{P}$ se cumple, para la \emph{fibra} $\mathrm{X}_{f(x)}$, en el punto $x$ de esta fibra, es \emph{localmente construible en} $\mathrm{X}$.
Nos proponemos ahora demostrar que, para la mayoría de estas propiedades, si se supone además que el morfismo $f$ es \emph{plano}, entonces el conjunto $\mathrm{E}$ es de hecho \emph{abierto en} $\mathrm{X}$.
Análogamente (\sref{IV.9.2} a \sref{IV.9.8}), hemos demostrado que, si $f$ es \emph{de presentación finita} y esta vez $\mathbf{P}$ designa ciertas propiedades \emph{globales} de preesquemas sobre cuerpos o de Módulos sobre estos preesquemas, el conjunto $\mathrm{F}$ de los $y\in\mathrm{Y}$ tales que la propiedad $\mathbf{P}$ se cumple para la fibra $\mathrm{X}_y$ es \emph{localmente construible en} $\mathrm{Y}$.
Demostraremos que, si se supone además que el morfismo $f$ es \emph{propio} y \emph{plano}, $\mathrm{F}$ es de hecho \emph{abierto en} $\mathrm{Y}$.
\end{env}

\begin{env}[12.0.2]
\label{IV.12.0.2}
El método general de demostración de las propiedades consideradas comprende tres etapas.
Primero se reduce al caso en que $\mathrm{Y}$ es \emph{afín} y $\mathrm{X}$ es de presentación finita; después:

\medskip
\begin{enumerate}[label=\emph{\Alph*)}]
  \item Con ayuda de \sref{IV.8.9.1} (y eventualmente de otros resultados de \textsection\,8) y de \sref{IV.11.2.6}, se reduce al caso en que $\mathrm{X}$ y $\mathrm{Y}$ son \emph{noetherianos}.

  \item Se aplican los resultados de \textsection\,9 recordados en \sref{IV.12.0.1} para demostrar que $\mathrm{E}$ (resp.\ $\mathrm{F}$) es \emph{construible}.

  \item Para ver que $\mathrm{E}$ es abierto basta, en virtud de \sref[0]{0.9.2.5}, demostrar que, si $x\in\mathrm{E}$, toda \emph{generización} $x'$ de $x$ pertenece también a $\mathrm{E}$.
Mediante \sref[II]{II.7.1.7} se ve que, puesto que $\mathrm{Y}$ es \emph{noetheriano}, existe el espectro de un \emph{anillo de valoración discreta} $\mathrm{Y}_1$ y un morfismo $h:\mathrm{Y}_1\to\mathrm{X}$ tales que, si $y_1$ (resp.\ $y_1'$) es el punto cerrado (resp.\ el punto genérico) de $\mathrm{Y}_1$, se tiene $h(y_1)=x$, $h(y_1')=x'$.
Se efectúa entonces el cambio de base $g=f\circ h:\mathrm{Y}_1\to\mathrm{Y}$; teniendo en cuenta los resultados de \textsection\textsection\,4 y 6 sobre los preesquemas localmente noetherianos sobre cuerpos y los cambios del cuerpo base, se reduce a demostrar la afirmación considerada para un punto $x_1$ de $\mathrm{X}_1=\mathrm{X}\times_{\mathrm{Y}}\mathrm{Y}_1$ situado sobre $y_1$ y para una generización $x_1'$ de $x_1$ situada sobre $y_1'$.
Como $\mathrm{Y}_1=\Spec(\mathrm{A})$, donde $\mathrm{A}$ es un anillo de valoración discreta, de uniformizante $t$, hay que demostrar finalmente, para un $\mathrm{A}$-módulo $\mathrm{M}$, una propiedad de $\mathrm{M}$ sabiendo que $\mathrm{M}/t\mathrm{M}$ posee la misma propiedad y que $t$ es $\mathrm{M}$-\emph{regular} (lo cual resulta de la hipótesis de planitud); para ello se utilizan los resultados de \sref{IV.3.4} y de \sref{IV.5.12}.
Se procede de manera análoga para el conjunto $\mathrm{F}\subset\mathrm{Y}$.
\end{enumerate}
\end{env}

\oldpage[IV-3]{174}

\subsection{Propiedades locales de las fibras de un morfismo plano localmente de presentación finita}
\label{subsection:IV.12.1}

\begin{theorem}[12.1.1]
\label{IV.12.1.1}
Sea $f:X\to Y$ un morfismo localmente de presentación finita, $\sh{F}$ un $\sh{O}_X$-Módulo $f$-plano y de presentación finita, $\Phi$ una parte finita de $\mathbf{Z}\cup\{\pm\infty\}$ y $k$ un entero.
Los siguientes subconjuntos de $X$ son abiertos:
\begin{enumerate}[label=(\roman*)]
  \item El conjunto de los $x\in X$ tales que las dimensiones de los ciclos primos asociados a $\sh{F}_{f(x)}$ y que contienen a $x$ pertenecen a $\Phi$.
  \item El conjunto de los $x\in X$ tales que todos los ciclos primos asociados a $\sh{F}_{f(x)}$ y que contienen a $x$ tienen dimensión $k$.
  \item El conjunto de los $x\in X$ tales que $x$ no pertenece a ningún ciclo primo inmerso asociado de $\sh{F}_{f(x)}$.
  \item El conjunto de los $x\in X$ tales que $\sh{F}_{f(x)}$ es equidimensional en el punto $x$ y posee la propiedad $(\mathrm{S}_k)$ en el punto $x$ \sref{IV.5.7.2}.
  \item El conjunto de los $x\in X$ tales que $\operatorname{coprof}((\sh{F}_{f(x)})_x)\leqslant k$ \sref[0]{0.16.4.9}.
  \item El conjunto de los $x\in X$ tales que $\sh{F}_{f(x)}$ es un $\sh{O}_{X_{f(x)}}$-Módulo de Cohen--Macaulay en el punto $x$ \sref{IV.5.7.1}.
  \item El conjunto de los $x\in X$ tales que $\sh{F}_{f(x)}$ es geométricamente reducido en el punto $x$ \sref{IV.4.6.22}.
  \item El conjunto de los $x\in X$ tales que $\sh{F}_{f(x)}$ es geométricamente puntualmente íntegro en el punto $x$ \sref{IV.4.6.22}.
\end{enumerate}

Como las cuestiones son locales sobre $X$, se reduce primero al caso en que $X=\Spec(B)$ e $Y=\Spec(A)$ son afines y $f$ es un morfismo de presentación finita.
Existe entonces un subanillo noetheriano $A_0$ de $A$, un $A_0$-preesquema de tipo finito $X_0$ y un $\sh{O}_{X_0}$-Módulo coherente $\sh{F}_0$ tales que $\sh{F}_0\otimes_{A_0}A$ es isomorfo a $\sh{F}$ \sref{IV.8.9.1}; además, en virtud de \sref{IV.11.2.6}, se puede suponer que $\sh{F}_0$ es $Y_0$-plano (con $Y_0=\Spec(A_0)$).
Si $h:X\to X_0$ es la proyección canónica, el conjunto $E$ de los puntos $x\in X$ en los que se satisface una de las propiedades (i) a (viii) es igual a $h^{-1}(E_0)$, donde $E_0$ es el conjunto de los $x_0\in X_0$ en los que se satisface la propiedad correspondiente para $\sh{F}_0$: esto resulta, para las propiedades (i) a (iii), de \sref{IV.4.2.7}; para las propiedades (iv) a (vi), de \sref{IV.6.7.1}; para las propiedades (vii) y (viii), de \sref{IV.4.7.11}.
\end{theorem}

Se reduce así al caso en que $A$ y $B$ son noetherianos, $f$ es de tipo finito y $\sh{F}=\widetilde{M}$, donde $M$ es un $B$-módulo de tipo finito.
En virtud de \sref{IV.9.9.2}, se sabe que el conjunto $E$ es \emph{construible} para todas las propiedades consideradas, y queda en cada caso la etapa~C) de \sref{IV.12.0.2}, en la que hay que demostrar que $E$ es \emph{estable por generización}.

\begin{env}[12.1.1.1]
\label{IV.12.1.1.1}
Comencemos por el caso~(iii), que es el más sencillo.
Sea $x'\neq x$ una generización de $x$ en $X$.
Pongamos $y=f(x)$, $y'=f(x')$; existe el espectro de un anillo de valoración discreta $Y_1$ y un morfismo $u:Y_1\to X$ tales que, si $y_1$ e $y_1'$ son respectivamente el punto cerrado y el punto genérico de $Y_1$, se tiene $u(y_1)=x$ y $u(y_1')=x'$ \sref[II]{II.7.1.7}; pongamos $g=f\circ u:Y_1\to Y$; si $X_1=X\times_Y Y_1$, existe una $Y_1$-sección $u_1:Y_1\to X_1$ tal que $u=g_1\circ u_1$, donde $g_1:X_1\to X$ es la proyección canónica.
Si $x_1=u_1(y_1)$, $x_1'=u_1(y_1')$, se tiene $g_1(x_1)=x$ y $g_1(x_1')=x'$, y $x_1'$ es una generización de $x_1$.
Aplicando de nuevo \sref{IV.4.2.7}, se ve
\oldpage[IV-3]{175}
que se reduce al caso en que $Y=\Spec(A)$ es el espectro de un anillo de valoración discreta, siendo $y$ el punto cerrado e $y'$ el punto genérico de $Y$.
Si $t$ es un uniformizante de $A$, la hipótesis de que $\sh{F}$ es $f$-plano implica que $t$ es $\sh{F}$-regular \sref[0]{0.6.3.4}.
Por hipótesis, ninguno de los ciclos primos inmersos asociados de $\sh{F}/t\sh{F}$ contiene a $x$.
Entonces, de \sref{IV.3.4.4} se sigue que $x$ no pertenece a ningún ciclo primo inmerso asociado de $\sh{F}$, y lo mismo ocurre con toda generización de $x$.
En particular, $x'$ no pertenece a ningún ciclo primo inmerso asociado de $\sh{F}$, ni \emph{a fortiori} a ninguno de los asociados a $\sh{F}_{y'}$ (pues $X_{y'}$ es un subpreesquema que induce un subconjunto abierto de $X$).
\end{env}

\begin{env}[12.1.1.2]
\label{IV.12.1.1.2}
Consideremos ahora los casos~(iv) y~(v) ((vi) no es sino un caso particular de~(v)).
Se procede como antes (utilizando \sref{IV.6.7.1} en lugar de \sref{IV.4.2.7}) y se reduce al caso en que $Y$ es el espectro de un anillo de valoración discreta~$A$.

El anillo $C=\sh{O}_{X,x}$ es entonces el localizado de una $A$-álgebra de tipo finito, luego es catenario \sref{IV.5.6.4}, y por hipótesis el $C$-módulo $M_x/tM_x$ satisface $(\mathrm{S}_k)$ y es equidimensional (resp.\ se tiene $\operatorname{coprof}(M_x/tM_x)\leqslant k$); se deduce de \sref{IV.5.12.2} (resp.\ de \sref[0]{0.16.4.10}) que $M_x$ satisface $(\mathrm{S}_k)$ y es equidimensional (resp.\ que $\operatorname{coprof}(M_x)\leqslant k$), pues $t$ es $M_x$-regular y pertenece al ideal maximal de $C$.
De aquí resultan las conclusiones, puesto que $x'$ es una generización de $x$ y $(\sh{F}_{y'})_{x'}=\sh{F}_{x'}$.
\end{env}

\begin{env}[12.1.1.3]
\label{IV.12.1.1.3}
Pasemos a los casos~(vii) y~(viii).
Evidentemente se puede sustituir $Y$ por el subpreesquema íntegro cuyo espacio subyacente es $\overline{\{y'\}}$, y $X$ por $f^{-1}(\overline{\{y'\}})$, sin cambiar las fibras en $y$ e $y'$; se puede suponer, pues, que $A$ es íntegro, con cuerpo de fracciones $K=k(y')$.
Se sabe (\sref{IV.4.5.11} y \sref{IV.4.7.8}) que existe una extensión finita $K'$ de $K$ tal que, para el $(\sh{O}_X\otimes_A K')$-Módulo $\sh{F}'=\sh{F}\otimes_A K'$, los ciclos primos asociados son geométricamente irreducibles y las longitudes geométricas de $\sh{F}'$ en los puntos maximales $z_i'$ de su soporte son respectivamente iguales a las longitudes de $\sh{F}'_{z_i'}$ en estos puntos; por tanto, teniendo en cuenta \sref{IV.4.2.7}, equivale decir que $\sh{F}$ es geométricamente reducido (resp.\ geométricamente puntualmente íntegro) en un punto $x'$ de $X_{y'}=X\otimes_A K$, o decir que $\sh{F}'$ es \emph{reducido} (resp.\ \emph{íntegro}) en los puntos de $X\otimes_A K'$ situados sobre $x'$.
Sea $A'$ una $A$-álgebra finita cuyo cuerpo de fracciones es $K'$; de \sref{IV.4.7.11} se sigue que, en todo punto de $X\otimes_A A'$ situado sobre $x$, $\sh{F}\otimes_A A'$ posee la propiedad~(vii) (resp.~(viii)).
Sustituyendo $A$ por $A'$ y $X$ por $X\otimes_A A'$, se ve que se puede reducir a demostrar que $\sh{F}$ es \emph{reducido} (resp.\ \emph{íntegro}) en toda generización $x'$ de $x$.
Se procede entonces como en \sref{IV.12.1.1.1}, utilizando esta vez \sref{IV.4.7.11}, en el caso en que $A$ es un anillo de valoración discreta, $y$ es el punto cerrado, $y'$ el punto genérico de $Y$, y el uniformizante $t$ de $A$ es un elemento $M$-\emph{regular}.
Como por hipótesis $\sh{F}/t\sh{F}$ es reducido (resp.\ íntegro) en el punto $x$, de \sref{IV.3.4.6} (resp.\ \sref{IV.3.4.5}) se sigue que $\sh{F}$ es reducido (resp.\ íntegro) en el punto $x$, y por tanto también en un entorno de $x$, en particular en el punto $x'$, lo que termina la demostración en los casos~(vii) y~(viii).
\end{env}

\begin{env}[12.1.1.4]
\label{IV.12.1.1.4}
Queda examinar los casos~(i) y~(ii).
Sustituyendo $Y$ por el subpreesquema íntegro cuyo espacio subyacente es $\overline{\{y'\}}$, se puede reducir al caso en que $Y$ es irreducible e $y'$ es su punto genérico. Sea $z'$ un punto genérico de un ciclo primo asociado de $\sh{F}_{y'}$
\oldpage[IV-3]{176}
que contiene a $x'$, y sea $Z'$ la adherencia de $z'$ en $X$, de modo que $x\in Z'$.
Para tratar el caso~(i), bastará demostrar la

\begin{proposition}[12.1.1.5]
\label{IV.12.1.1.5}
Sea $Y$ un preesquema localmente noetheriano, $f:X\to Y$ un morfismo localmente de tipo finito, $\sh{F}$ un $\sh{O}_X$-Módulo coherente y $f$-plano; para todo $y\in Y$, pongamos $\sh{F}_y=\sh{F}\otimes_{\sh{O}_Y}k(y)$.
Sea $z'$ un punto de $X$, $y'=f(z')$, y supongamos que $z'\in\operatorname{Ass}(\sh{F}_{y'})$.
Sea $Z'$ \textnormal{(}resp.\ $Y'$\textnormal{)} el subpreesquema reducido de $X$ \textnormal{(}resp.\ de $Y$\textnormal{)} cuyo espacio subyacente es la adherencia de $\{z'\}$ en $X$ \textnormal{(}resp.\ de $\{y'\}$ en $Y$\textnormal{)}.
Entonces, para todo $y\in f(Z')$, las dimensiones de todas las componentes irreducibles de $Z'_y$ son iguales a $\dim(Z'_{y'})$ \textnormal{(}dimensión de la adherencia de $z'$ en $X_{y'}$\textnormal{)} \textnormal{(}lo que expresaremos más precisamente después \sref{IV.13.2.2} diciendo que la restricción $Z'\to Y'$ de $f$ es un morfismo equidimensional\textnormal{)}; además, en todo punto maximal $z$ de $Z'_y$, se tiene $z\in\operatorname{Ass}(\sh{F}_y)$.
\end{proposition}

Para ello, reduzcamos ahora al caso en que $Y$ es el espectro de un anillo de valoración discreta.
Tomemos el espectro de un anillo de valoración discreta $Y_1$ y un morfismo $h:Y_1\to X$ tales que $h(y_1)=z$, $h(y_1')=z'$, siendo $y_1$ e $y_1'$ respectivamente el punto cerrado y el punto genérico de $Y_1$ \sref[II]{II.7.1.7}.
Pongamos $g=f\circ h:Y_1\to Y$, $X_1=X\times_Y Y_1$, y sean $f_1:X_1\to Y_1$, $g_1:X_1\to X$ las proyecciones canónicas; existe una $Y_1$-sección $h_1:Y_1\to X_1$ tal que $h=g_1\circ h_1$ \sref[I]{I.3.3.14}.
Si $Z_1'=Z'_{y'}\otimes_{k(y')}k(y_1')$, se sabe \sref{IV.4.2.6} que las componentes irreducibles de $Z_1'$ dominan $Z'_{y'}$; sea $z_1'$ el punto genérico de una de estas componentes que contiene a $h_1(y_1')$, de modo que $f_1(z_1')=y_1'$ y $g_1(z_1')=z'$.
Como $h_1(y_1)$ es una especialización de $h_1(y_1')$, es \emph{a fortiori} una especialización de $z_1'$; sea $z_1$ un punto genérico de una de las componentes irreducibles de $\overline{\{z_1'\}}\cap(X_1)_{y_1}$ que contienen a $h_1(y_1)$.
Entonces $g_1(z_1)$ es una especialización de $g_1(z_1')=z'$ en $X$, y $z=h(y_1)=g_1(h_1(y_1))$ es una especialización de $g_1(z_1)$; pero, puesto que $f(g_1(z_1))=y$ y $z$ es un punto genérico de $Z'_y$, se tiene necesariamente $z=g_1(z_1)$.

Supongamos ahora demostrado que, si se pone $\sh{F}_1=\sh{F}\otimes_{\sh{O}_Y}\sh{O}_{Y_1}$, entonces $z_1\in\operatorname{Ass}((\sh{F}_1)_{y_1})$; de \sref{IV.4.2.7} se seguirá que $z\in\operatorname{Ass}(\sh{F}_y)$; además, las dimensiones de $\overline{\{z_1\}}\cap(X_1)_{y_1}$ y de $\overline{\{z\}}\cap X_y$ son iguales, y asimismo las dimensiones de $\overline{\{z_1'\}}\cap(X_1)_{y_1'}$ y de $\overline{\{z'\}}\cap X_{y'}$ son iguales \sref{IV.4.2.7}; se ha reducido, pues, como se anunció, al caso en que $Y$ es el espectro de un anillo de valoración discreta $A$, siendo $y$ e $y'$ respectivamente su punto cerrado y su punto genérico.

Como $z'\in\operatorname{Ass}(\sh{F}_{y'})$, se tiene \emph{a fortiori} $z'\in\operatorname{Ass}(\sh{F})$.
Si $t$ es un uniformizante de $A$, $t$ es $\sh{F}$-regular por planitud, por lo que \sref{IV.3.4.3} da $z\in\operatorname{Ass}(\sh{F}/t\sh{F})=\operatorname{Ass}(\sh{F}_y)$.
En cuanto a la afirmación relativa a las dimensiones, se deduce de \sref{IV.7.1.13}, aplicado a un subpreesquema cerrado de $X$ cuyo espacio subyacente es $Z'$.

Abordemos finalmente el caso~(ii) de \sref{IV.12.1.1}.
Con la misma notación, hay que demostrar que, si todos los ciclos primos asociados a $\sh{F}_y$ y que contienen a $x$ tienen dimensión $k$, lo mismo ocurre con todos los ciclos primos asociados a $\sh{F}_{y'}$ y que contienen a $x'$.
Ahora bien, acabamos de ver que \emph{todo} ciclo primo asociado a $\sh{F}_{y'}$ y que contiene a $x'$ tiene la misma dimensión que \emph{uno} de los ciclos primos asociados a $\sh{F}_y$ y que contienen a $x$; esto demuestra~(ii).
\end{env}

\begin{remarks}[12.1.2]
\label{IV.12.1.2}
\begin{enumerate}[label=(\roman*)]
  \item En las condiciones de \sref{IV.12.1.1.5}, con $\sh{F}=\sh{O}_X$ (de modo que $f$ es \emph{plano}), no se puede afirmar en general que la restricción $Z'\to Y'$ de $f$ sea un morfismo plano, ni siquiera un morfismo \emph{abierto}.
Esto es lo que muestra el ejemplo \sref{IV.6.5.5},
\oldpage[IV-3]{177}
(ii), en el que se toma por $Z'$ una de las dos componentes irreducibles de $X=\Spec(B)$.
Es inmediato que este morfismo restricción no es abierto en los puntos de $Z'$ situados sobre el «punto doble» de $Y'=Y$.

  \item Bajo las hipótesis de \sref{IV.12.1.1.5}, con $\sh{F}=\sh{O}_X$, no se puede debilitar la condición «$f$ es plano» sustituyéndola por «$f$ es universalmente abierto» (cf.\ \sref{IV.2.4.6}), como veremos más adelante con un ejemplo \sref{IV.14.4.10},~(i)).
\end{enumerate}
\end{remarks}

\begin{corollary}[12.1.3]
\label{IV.12.1.3}
Bajo las hipótesis de \sref{IV.12.1.1}, la función $x\mapsto\operatorname{coprof}((\sh{F}_{f(x)})_x)$ es semicontinua superiormente en $X$, y la función $x\mapsto\operatorname{prof}_x^*(\sh{F}_{f(x)})$ \sref{IV.10.8.1} es semicontinua inferiormente en $X$.
\end{corollary}

La primera afirmación no es sino una formulación equivalente de \sref{IV.12.1.1},~(v).
Para demostrar la segunda, se puede primero, teniendo en cuenta \sref{IV.10.8.7} y \sref{IV.10.8.8}, reducir como en \sref{IV.12.1.1} al caso en que $Y$ es el espectro de un anillo de valoración discreta $A$, de punto cerrado $y$ y punto genérico $y'$, siendo $X=\Spec(B)$ afín y $\sh{F}=\widetilde{M}$, donde $M$ es un $B$-módulo de tipo finito.
Como $\sh{F}$ es $f$-plano, toda componente irreducible de $\operatorname{Supp}(\sh{F})$ que contiene a $x\in X_y$ domina $Y$ \sref{IV.2.3.4}; en otras palabras, su punto genérico $z'$ pertenece a $X_{y'}$. Las componentes irreducibles $Z$ de $\operatorname{Supp}(\sh{F})$ que contienen una generización $x'$ de $x$ perteneciente a $X_{y'}$ son, por tanto, exactamente las que contienen a $x$; además, por \sref{IV.12.1.1.5}, se tiene $\dim(Z_{y'})=\dim(Z_y)$, de donde, si $T=\operatorname{Supp}(\sh{F})$, $\dim_x(T_y)=\dim_{x'}(T_{y'})$.
Además, para un uniformizante $t$ de $A$, se ha visto en \sref{IV.12.1.1},~(v) que $\operatorname{coprof}(M_x)=\operatorname{coprof}(M_x/tM_x)$; como, por otra parte, $\operatorname{coprof}(M_{x'})\leqslant\operatorname{coprof}(M_x)$ \sref{IV.6.11.5}, se ve que $\operatorname{coprof}((\sh{F}_{y'})_{x'})\leqslant\operatorname{coprof}((\sh{F}_y)_x)$; la relación $\operatorname{prof}_x^*(\sh{F}_y)\leqslant\operatorname{prof}_{x'}^*(\sh{F}_{y'})$ resulta entonces de \sref{IV.10.8.7}.

\begin{remark}[12.1.4]
\label{IV.12.1.4}
También se deduce de \sref{IV.12.1.1},~(i) que la función $x\mapsto\dim_x(X_{f(x)})$ es semicontinua superiormente en $X$, pero veremos más adelante \sref{IV.13.1.3} que esta propiedad es cierta incluso sin suponer $f$ plano.
\end{remark}

\begin{corollary}[12.1.5]
\label{IV.12.1.5}
Sea $Y$ un preesquema localmente noetheriano, $f:X\to Y$ un morfismo localmente de tipo finito, $\sh{F}$ un $\sh{O}_X$-Módulo coherente $f$-plano, $\sh{G}$ un $\sh{O}_Y$-Módulo coherente, $x$ un punto de $X$ e $y=f(x)$.
Supongamos que $\sh{G}$ posee la propiedad $(\mathrm{S}_k)$ en el punto $y$, y que $\sh{F}_y$ posee la propiedad $(\mathrm{S}_k)$ en el punto $x$ y es equidimensional en el punto $x$.
Entonces:
\begin{enumerate}[label=(\roman*)]
  \item $\sh{F}\otimes_{\sh{O}_Y}\sh{G}$ posee la propiedad $(\mathrm{S}_k)$ en el punto $x$.
  \item Si además $Y$ es localmente inmersible en un esquema regular, o es un preesquema excelente \sref{IV.7.8.5}, existe un entorno $V$ de $x$ en $X$ tal que $\sh{F}\otimes_{\sh{O}_Y}\sh{G}$ posee la propiedad $(\mathrm{S}_k)$ en ese entorno.
\end{enumerate}
\end{corollary}

\begin{proof}
En efecto, por \sref{IV.12.1.1},~(iv), existe un entorno abierto $V$ de $x$ tal que, para todo $x'\in V$, $\sh{F}_{f(x')}$ satisface $(\mathrm{S}_k)$ en el punto $x'$.
Por otra parte, $\sh{G}$ satisface $(\mathrm{S}_k)$ en todos los puntos de $\Spec(\sh{O}_y)$ \sref{IV.5.7.2}.
Sustituyendo $Y$ por $Y'=\Spec(\sh{O}_y)$ y $X$ por $V\times_Y Y'$, se reduce (teniendo en cuenta \sref[I]{I.3.6.5}) al caso en que $\sh{G}$ satisface $(\mathrm{S}_k)$ en todo $Y$ y en que, para todo $y\in f(X)$, $\sh{F}_y$ posee la propiedad $(\mathrm{S}_k)$ en todos los puntos de $f^{-1}(y)$; como $\sh{F}$ es $f$-plano, se sabe entonces por \sref{IV.6.4.1} que $\sh{F}\otimes_{\sh{O}_Y}\sh{G}$ posee la propiedad $(\mathrm{S}_k)$ en $X$, lo que demuestra~(i).
Para demostrar~(ii), basta observar que, bajo las hipótesis formuladas, $\sh{G}$ posee la
\oldpage[IV-3]{178}
propiedad $(\mathrm{S}_k)$ en un entorno $U$ de $y$ en $Y$ \sref{IV.6.11.2} y \sref{IV.7.8.3},~(iv)); se concluye entonces de la misma manera, sustituyendo esta vez $Y$ por $U$ y $X$ por $V\times_Y U$.
\end{proof}

\begin{theorem}[12.1.6]
\label{IV.12.1.6}
Sea $f:X\to Y$ un morfismo plano localmente de presentación finita, y sea $k$ un entero.
Los siguientes subconjuntos de $X$ son abiertos:
\begin{enumerate}[label=(\roman*)]
  \item El conjunto de los puntos $x\in X$ tales que $X_{f(x)}$ posee la propiedad $(\mathrm{S}_k)$ en el punto $x$.
  \item El conjunto de los $x\in X$ tales que $X_{f(x)}$ satisface geométricamente la propiedad $(\mathrm{R}_k)$ en el punto $x$, es equidimensional en el punto $x$ y, además, $x$ no pertenece a ningún ciclo primo inmerso asociado de $X_{f(x)}$.
  \item El conjunto de los $x\in X$ tales que $X_{f(x)}$ es geométricamente regular (es decir, liso) en el punto $x$.
  \item El conjunto de los $x\in X$ tales que $X_{f(x)}$ es geométricamente normal en el punto $x$.
\end{enumerate}

Las etapas~A) y~B) de \sref{IV.12.0.2} se efectúan aquí como en \sref{IV.12.1.1}; para la etapa~A), se utiliza \sref{IV.6.7.8}, y también \sref{IV.4.2.7} para~(ii); para la etapa~B), se utilizan \sref{IV.9.9.2} y \sref{IV.9.9.4}.
Queda examinar la etapa~C) en cada caso.
\end{theorem}

\begin{proof}
\begin{enumerate}[label=(\roman*)]
  \item[(i)] Como en \sref{IV.12.1.1.2}, se reduce (utilizando \sref{IV.6.7.8}) al caso en que $Y$ es el espectro de un anillo de valoración discreta $A$, $X=\Spec(B)$, donde $B$ es una $A$-álgebra de tipo finito, y $x,x'$ son dos puntos de $X$ tales que $f(x)=y$ es el punto cerrado e $y'=f(x')$ el punto genérico de $Y$, siendo además $x'$ una generización de $x$.
Como se quiere demostrar que $X$ satisface $(\mathrm{S}_k)$ en el punto $x'$, se puede sustituir $X$ por $\Spec(\sh{O}_{X,x})$; es decir, se puede suponer que el anillo $B$ es local, que el homomorfismo $A\to B$ es local y que $B$ es una $A$-álgebra esencialmente de tipo finito \sref{IV.1.3.8}.
Entonces, por hipótesis, si $t$ es un uniformizante de $A$, $t$ es un elemento $B$-regular por planitud, y $B/tB$ satisface $(\mathrm{S}_k)$.
Pero, como $A$ es un anillo universalmente catenario \sref{IV.5.6.4}, $B$ es catenario; por tanto, de \sref{IV.5.12.4} se sigue que $B$ satisface $(\mathrm{S}_k)$, lo que termina la demostración en el caso~(i).

  \item[(iii)] Aquí la etapa~C) es innecesaria; en efecto, como $f$ es plano, se sabe por \sref{IV.6.8.7} que, cuando $Y$ es localmente noetheriano y $f$ localmente de tipo finito, el conjunto de los $x\in X$ tales que $X_{f(x)}$ es geométricamente regular en el punto $x$ es abierto en $X$.

  \item[(ii)] Razonando como en \sref{IV.12.1.1.3}, para demostrar que la propiedad considerada es estable por generización, se puede primero, considerando una extensión finita arbitraria $K'$ de $k(y')$ y teniendo en cuenta la definición \sref{IV.6.7.6} de la propiedad geométrica $(\mathrm{R}_k)$, así como la invariancia por cambio del cuerpo base de las otras dos propiedades que aparecen en~(ii), sustituir en~(ii) la propiedad geométrica $(\mathrm{R}_k)$ por la propiedad $(\mathrm{R}_k)$.
Procediendo entonces como en~(i), se reduce al caso en que $Y$ es el espectro de un anillo de valoración discreta $A$; como $A$ es catenario, basta aplicar \sref{IV.5.12.5} para concluir.

  \item[(iv)] El conjunto de los puntos $x\in X$ tales que $X_{f(x)}$ es geométricamente normal en el punto $x$ está contenido en el de los $x\in X$ tales que $X_{f(x)}$ es geométricamente puntualmente íntegro y satisface $(\mathrm{S}_2)$ en el punto $x$, y este último conjunto es abierto en $X$ en virtud de \sref{IV.12.1.1},~(viii).
Se puede, pues, suponer ya que, para todo $x\in X$, $X_{f(x)}$ es geométricamente puntualmente íntegro y satisface $(\mathrm{S}_2)$ en el punto $x$, y \emph{a fortiori} es equidimensional.
Por otra parte, según el criterio de Serre \sref{IV.5.8.6}, decir que $X_{f(x)}$ es geométricamente normal en el punto $x$ significa que $X_{f(x)}$ satisface $(\mathrm{S}_2)$ y geométricamente la propiedad $(\mathrm{R}_1)$ en $x$;
\oldpage[IV-3]{179}
pero, en virtud de~(ii) y de las observaciones precedentes, este conjunto es la intersección de dos subconjuntos abiertos de $X$, luego es abierto en $X$.
\end{enumerate}
\end{proof}

\begin{corollary}[12.1.7]
\label{IV.12.1.7}
Sea $f:X\to Y$ un morfismo localmente de presentación finita.
Entonces, el conjunto de los puntos $x\in X$ en los que $f$ posee una de las propiedades siguientes \sref{IV.6.8.1}:
\begin{enumerate}[label=(\roman*)]
  \item satisface la propiedad $(\mathrm{S}_k)$;
  \item tiene coprofundidad $\leqslant n$;
  \item es de Cohen--Macaulay;
  \item es regular (o liso, lo que equivale a lo mismo);
  \item es normal;
  \item es reducido;
\end{enumerate}
es abierto en $X$.
\end{corollary}

\begin{proof}
En efecto, de \sref{IV.11.3.1} se sigue que el conjunto de los $x\in X$ en los que $f$ es \emph{plano} es \emph{abierto}.
Se puede, pues, restringir al caso en que $f$ es plano, y entonces el corolario resulta de \sref{IV.12.1.1},~(iv),~(vi) y~(vii), y de \sref{IV.12.1.6},~(i),~(ii) y~(iv).
\end{proof}

\begin{remarks}[12.1.8]
\label{IV.12.1.8}
\begin{enumerate}[label=(\roman*)]
  \item En \sref{IV.12.1.6},~(ii), no se puede suprimir la hipótesis de que $x$ no pertenece a ningún ciclo primo inmerso asociado de $X_{f(x)}$.
Lo muestra el ejemplo \sref{IV.5.12.3}: se toma $X=\Spec(A)$, siendo $Y$ el espectro del anillo local $A_0$ de $k[T]$ correspondiente al ideal primo $(T)$, y el morfismo $X\to Y$ correspondiente al homomorfismo inyectivo $A_0\to A$ obtenido, por localización y paso al cociente, de la inyección $k[T]\to k[T,U,V]$.
Este morfismo es plano porque $A$ es un $A_0$-módulo sin torsión \sref[0]{0.6.3.4}.
Aquí la fibra $X_y$ en el punto cerrado $y$ de $Y$, igual a $\Spec(A/tA)$, es irreducible, de dimensión~1 y satisface geométricamente la condición $(\mathrm{R}_0)$, pues el anillo local en su punto genérico es un cuerpo.
En cambio, en el punto genérico $y'$ de $Y$, la fibra $X_{y'}$ tiene dos componentes irreducibles de dimensiones~0 y~1, y la de dimensión~0 no es reducida, de modo que $X_{y'}$ no satisface la condición $(\mathrm{R}_0)$.

  \item En \sref{IV.12.1.6},~(ii), tampoco se puede suprimir la hipótesis de que $X_{f(x)}$ es equidimensional en el punto $x$.
Se ve mediante el ejemplo \sref{IV.5.12.6}, con $X=\Spec(A)$ e $Y=\Spec(A_0)$, donde $A_0$ se define como en~(i); el morfismo $X\to Y$ se obtiene de nuevo de la inyección $k[T]\to k[T,U,V,W]$ por localización y paso al cociente, y es plano porque $A$ es un $A_0$-módulo sin torsión \sref[0]{0.6.3.4}.
La fibra $X_y$ en el punto cerrado $y$ de $Y$ es reducida (por tanto satisface $(\mathrm{S}_1)$) y satisface $(\mathrm{R}_1)$, pero tiene dos componentes irreducibles de dimensiones~2 y~1, mientras que la fibra $X_{y'}$ en el punto genérico $y'$ de $Y$ no satisface $(\mathrm{R}_1)$.
\end{enumerate}
\end{remarks}


\subsection{Propiedades locales y globales de las fibras de un morfismo propio y plano de presentación finita}
\label{subsection:IV.12.2}

\oldpage[IV-3]{179}

\begin{theorem}[12.2.1]
\label{IV.12.2.1}
Sea $f:X\to Y$ un morfismo \emph{propio} de presentación finita, $\sh{F}$ un $f$-plano $\sh{O}_X$-módulo de presentación finita, $\Phi$ una parte finita de $\mathbf{Z}\cup\{\pm\infty\}$ y $k$ un entero.
Los siguientes subconjuntos de $Y$ son abiertos:
\begin{enumerate}[label=(\roman*)]
  \item El conjunto de los $y\in Y$ tales que el conjunto de las dimensiones de los ciclos primos asociados a $\sh{F}_y$ está contenido en $\Phi$.
  \item El conjunto de los $y\in Y$ tales que el conjunto de las dimensiones de las componentes irreducibles de $\operatorname{Supp}(\sh{F}_y)$ contiene a $\Phi$.
  \item El conjunto de los $y\in Y$ tales que todos los ciclos primos asociados a $\sh{F}_y$ tienen la misma dimensión, igual a $k$.
  \item El conjunto de los $y\in Y$ tales que $\sh{F}_y$ no tiene ningún ciclo primo inmerso asociado y el conjunto de las dimensiones de las componentes irreducibles de $\operatorname{Supp}(\sh{F}_y)$ es igual a $\Phi$.
\oldpage[IV-3]{180}
  \item El conjunto de los $y\in Y$ tales que $\sh{F}_y$ posee la propiedad $(\mathrm{S}_k)$ y es equidimensional en todo punto de $X_y$.
  \item El conjunto de los $y\in Y$ tales que $\operatorname{coprof}(\sh{F}_y)\leqslant k$.
  \item El conjunto de los $y\in Y$ tales que $\sh{F}_y$ es un $\sh{O}_{X_y}$-Módulo de Cohen--Macaulay.
  \item El conjunto de los $y\in Y$ tales que $\sh{F}_y$ es geométricamente reducido.
  \item El conjunto de los $y\in Y$ tales que $\sh{F}_y$ es geométricamente puntualmente íntegro en todo punto de $X_y$.
  \item El conjunto de los $y\in Y$ tales que $\sh{F}_y$ es geométricamente íntegro.
  \item El conjunto de los $y\in Y$ tales que $\sh{F}_y$ no tiene ningún ciclo primo inmerso asociado y la suma $l(y)$ de las multiplicidades totales \sref{IV.4.7.12} de $\sh{F}_y$ en los puntos maximales de $\operatorname{Supp}(\sh{F}_y)$ es $\leqslant k$.
\end{enumerate}
\end{theorem}

\begin{proof}
Con excepción de~(ii), (iii), (iv), (x) y~(xi), las propiedades $\mathbf{P}(y)$ consideradas son de la forma: «para todo $x\in X_y$, $\sh{F}_y$ posee la propiedad $\mathbf{Q}(x)$», donde $\mathbf{Q}(x)$ es una de las propiedades~(i) a~(viii) de \sref{IV.12.1.1}.
De \sref{IV.12.1.1} se sigue que el conjunto $U$ de los $x\in X$ tales que $\mathbf{Q}(x)$ es verdadera es abierto, y que el conjunto $V$ de los $y\in Y$ tales que $\mathbf{P}(y)$ es verdadera es precisamente $Y-f(X-U)$; en todos estos casos, el teorema queda, por tanto, demostrado bajo la sola hipótesis de que el morfismo $f$ sea \emph{cerrado}.
Para~(iii), se aplica~(i) con $\Phi$ reducido a un solo elemento.
Queda examinar por separado los casos~(ii), (iv) y~(xi) ((x) se deduce inmediatamente de~(xi) y \sref{IV.4.7.13}),\footnote{La fuente impresa cita (4.7.14), pero la {\EGAReadLink{IV.4.7}{Sección 4.7}} termina con la {\EGAReadLink{IV.4.7.13}{Proposición 4.7.13}} antes de que comience la {\EGAReadLink{IV.4.8}{Sección 4.8}}; la {\EGAReadLink{IV.4.7.13}{Proposición 4.7.13}} es el criterio pertinente de integralidad geométrica.} aplicando siempre el método descrito en \sref{IV.12.0.2}.
\end{proof}

\begin{env}[12.2.1.1]
\label{IV.12.2.1.1}
Caso~(ii): Las etapas~A) y~B) de la demostración se desarrollan como al comienzo de \sref{IV.12.1.1}; para la etapa~A), se utilizan \sref{IV.8.9.1}, \sref{IV.11.2.6}, así como \sref{IV.4.2.7}; para la etapa~B), se utiliza \sref{IV.9.5.5} aplicado a $\operatorname{Supp}(\sh{F})$.
Queda demostrar que, si (suponiendo $Y$ noetheriano) el conjunto de las dimensiones de las componentes irreducibles de $\operatorname{Supp}(\sh{F}_y)$ contiene a $\Phi$, ocurre lo mismo con el conjunto de las dimensiones de las componentes irreducibles de $\operatorname{Supp}(\sh{F}_{y'})$, para toda \emph{generización} $y'$ de $y$ en $Y$.
Sea $Y_1$ el espectro de un anillo de valoración discreta tal que, si $y_1$ e $y_1'$ son el punto cerrado y el punto genérico de $Y_1$, existe un morfismo $h:Y_1\to Y$ tal que $h(y_1)=y$, $h(y_1')=y'$ \sref[II]{II.7.1.7}.
Aplicando \sref{IV.4.2.7}, se ve que se puede sustituir $Y$ por $Y_1$ y $X$ por $X_1=X\times_Y Y_1$; en otras palabras, se puede restringir al caso en que $Y$ es el espectro de un anillo de valoración discreta, $y$ su punto cerrado e $y'$ su punto genérico.

Utilizando (\textbf{Err}\textsubscript{\textbf{III}},~30), se puede restringir al caso en que $\operatorname{Supp}(\sh{F})=X$, de modo que $f$ es cuasiplano \sref{IV.2.3.3}; las componentes irreducibles $Z_i$ de $X$ dominan entonces $Y$ \sref{IV.2.3.4}, es decir, sus puntos genéricos pertenecen a $X_{y'}$, y los $Z_i\cap X_{y'}$ son las componentes irreducibles de $X_{y'}$.
Pero toda componente irreducible de $X_y$ está contenida en uno de los $Z_i$, luego es una componente irreducible de $(Z_i)_y$; ahora bien, se sabe por \sref{IV.7.1.13} que las dimensiones de todas las componentes irreducibles de $(Z_i)_y$ son iguales a $\dim((Z_i)_{y'})$, lo que termina la demostración de~(ii).
\end{env}

\begin{env}[12.2.1.2]
\label{IV.12.2.1.2}
Caso~(iv): El mismo razonamiento que al comienzo de \sref{IV.12.2.1.1}, utilizando \sref{IV.12.1.1},~(iii), muestra que se puede suponer ya (sustituyendo $Y$ por un subconjunto abierto de $Y$) que, para todo $y\in Y$, $\sh{F}_y$ no tiene ningún ciclo primo inmerso asociado.
La afirmación del caso~(iv) es entonces consecuencia inmediata de las afirmaciones de los casos~(i) y~(ii).
\end{env}

\begin{env}[12.2.1.3]
\label{IV.12.2.1.3}
\oldpage[IV-3]{181}
Caso~(xi): Para la etapa~A) del razonamiento, se utilizan \sref{IV.8.9.1}, \sref{IV.11.2.6}, {\hyperlink{ega4.c04.thm.8-10-5}{(8.10.5,~(xii))}} (para conservar la hipótesis de que $f$ es propio), así como \sref{IV.4.2.7}, \sref{IV.4.5.6} y \sref{IV.4.7.9}.
Para la etapa~B), se ve, como en el caso~(iv), que se puede suponer que, para todo $y\in Y$, $\sh{F}_y$ no tiene ningún ciclo primo inmerso asociado, y se aplica \sref{IV.9.8.8}, que demuestra que la función $z\mapsto l(z)$ es construible.
Queda, por tanto, ver (suponiendo $Y$ noetheriano y $f$ propio) que, para toda generización $y'$ de un punto $y\in Y$, $l(y')\leqslant l(y)$.
Razonando como en \sref{IV.12.1.1.3}, se ve (con ayuda de \sref{IV.4.7.8} y \sref{IV.4.5.11}) que se puede suponer que las componentes irreducibles de $\operatorname{Supp}(\sh{F}_{y'})$ son geométricamente irreducibles y que $l(y')$ es la \emph{suma de las longitudes} de $(\sh{F}_{y'})_{z_j'}$ en los puntos maximales $z_j'$ de $\operatorname{Supp}(\sh{F}_{y'})$.
Aplicando de nuevo \sref{IV.4.2.7}, \sref{IV.4.5.6} y \sref{IV.4.7.9}, se reduce, como en \sref{IV.12.2.1.1}, al caso en que $Y$ es el espectro de un anillo de valoración discreta, $y$ su punto cerrado e $y'$ su punto genérico.
El hecho de que $l(y')\leqslant l(y)$ será entonces consecuencia del

\begin{lemma}[12.2.1.4]
\label{IV.12.2.1.4}
Sea $Y$ el espectro de un anillo de valoración discreta, $y$ su punto cerrado, $y'$ su punto genérico, $f:X\to Y$ un morfismo propio, $\sh{F}$ un $\sh{O}_X$-Módulo coherente y $f$-plano, y sean $z_i$ \textnormal{(}resp.\ $z_j'$\textnormal{)} los puntos maximales de $\operatorname{Supp}(\sh{F}_y)$ \textnormal{(}resp.\ de $\operatorname{Supp}(\sh{F}_{y'})$\textnormal{)}.
Supongamos que $\sh{F}_y$ no tiene ningún ciclo primo inmerso asociado.
Entonces se tiene
\[
  \sum_j \operatorname{long}((\sh{F}_{y'})_{z_j'}) \leqslant \sum_i \operatorname{long}((\sh{F}_y)_{z_i}).
\]
\end{lemma}

Se tiene $Y=\Spec(A)$, donde $A$ es un anillo de valoración discreta, cuyo uniformizante denotaremos por $t$, de modo que $\sh{F}_y=\sh{F}/t\sh{F}$.
Como $\sh{F}$ es $f$-plano, los $z_j'$ son también los puntos maximales de $\operatorname{Supp}(\sh{F})$ \sref{IV.2.3.4}; para cada $z_i$, denotemos por $z_{ik}'$ aquellos $z_j'$ que son generizaciones de $z_i$; de \sref{IV.3.4.1.1} se sigue que
\[
  \operatorname{long}((\sh{F}_y)_{z_i}) \geqslant \sum_k \operatorname{long}((\sh{F}_{y'})_{z_{ik}'})
\]
de donde, sumando,
\[
  \sum_i \operatorname{long}((\sh{F}_y)_{z_i}) \geqslant \sum_i \sum_k \operatorname{long}((\sh{F}_{y'})_{z_{ik}'}).
\]

El lema quedará, por tanto, demostrado si establecemos que, para todo $z_j'$, existe al menos un índice $i$ tal que $z_j'$ es uno de los $z_{ik}'$.
Ahora bien, como $f$ es propio (luego cerrado) y $f(z_j')=y'$, existe $x\in X$ que es una especialización de $z_j'$ y tal que $f(x)=y$; en otras palabras, $\overline{\{z_j'\}}\cap X_y$ no es vacío.
Como $t$ es $\sh{F}$-regular por planitud, se deduce de \sref{IV.3.4.3} que existe al menos un punto de $\operatorname{Ass}(\sh{F}_y)$ del cual $z_j'$ es una generización; pero, como los puntos de $\operatorname{Ass}(\sh{F}_y)$ son por hipótesis los $z_i$, esto termina la demostración del lema.
\end{env}

\begin{corollary}[12.2.2]
\label{IV.12.2.2}
Bajo las hipótesis de \sref{IV.12.2.1}:
\begin{enumerate}[label=(\roman*)]
  \item La función $y\mapsto\dim(\operatorname{Supp}(\sh{F}_y))$ es continua (y por tanto localmente constante) en $Y$.
  \item La función $y\mapsto\operatorname{coprof}(\sh{F}_y)$ \sref{IV.5.7.1} es semicontinua superiormente en $Y$.
  \item La función $y\mapsto l(y)$ \sref{IV.12.2.1},~(xi) es semicontinua superiormente en $Y$ cuando los $\sh{F}_y$ no tienen ningún ciclo primo inmerso asociado.
\oldpage[IV-3]{182}
  \item La función $y\mapsto\operatorname{prof}^*(\sh{F}_y)$ \sref{IV.10.8.1} es semicontinua inferiormente en $Y$.
\end{enumerate}
\end{corollary}

\begin{proof}
\begin{enumerate}[label=(\roman*)]
  \item De \sref{IV.12.2.1},~(i), aplicado a $\Phi$ igual al menor intervalo de $\mathbf{N}$ que contiene las dimensiones de los ciclos primos asociados a $\sh{F}_y$, se sigue que $z\mapsto\dim(\operatorname{Supp}(\sh{F}_z))$ es semicontinua superiormente en el punto $y$; por otra parte, de \sref{IV.12.2.1},~(ii), aplicado a $\Phi$ igual al conjunto de las dimensiones de las componentes irreducibles de $\operatorname{Supp}(\sh{F}_y)$, se sigue que esta misma función es semicontinua inferiormente en el punto $y$, de donde la conclusión.
Las afirmaciones~(ii) y~(iii) no son sino otras formulaciones de \sref{IV.12.2.1},~(vi) y~(xi).
Finalmente, por \sref{IV.12.1.3}, el conjunto de los $x\in X$ tales que $\operatorname{prof}_{f(x)}^*(\sh{F}_{f(x)})\leqslant k$ es abierto, y la conclusión de~(iv) resulta por el mismo razonamiento que al comienzo de \sref{IV.12.2.1}.
\end{enumerate}
\end{proof}

\begin{remarks}[12.2.3]
\label{IV.12.2.3}
\begin{enumerate}[label=(\roman*)]
  \item Obsérvese que, para \sref{IV.12.2.1},~(ii), se puede prescindir de la hipótesis de que $f$ es \emph{propio}.

  \item En \sref{IV.12.2.1},~(ii), no se puede sustituir la condición de que el conjunto $E(y)$ de las dimensiones de las componentes irreducibles de $\operatorname{Supp}(\sh{F}_y)$ contenga a $\Phi$ por la condición de que $E(y)$ sea \emph{igual a} $\Phi$.
En efecto, sea $k$ un cuerpo, y consideremos el espacio proyectivo $P=\mathbf{P}_k^3=\operatorname{Proj}(C)$ de dimensión~$3$, donde $C=k[t_0,t_1,t_2,t_3]$ (siendo los $t_i$ indeterminadas). En $P$, sea $X_0$ el subpreesquema cerrado que es la «unión de la recta $X_1$ definida por $t_1=t_3=0$ y el plano $X_2$ definido por $t_2-t_3=0$». En términos geométricos, esto significa que $X_0=\operatorname{Proj}(C/\mathfrak{a})$, donde el ideal graduado $\mathfrak{a}=\mathfrak{b}\mathfrak{c}$, con $\mathfrak{b}=Ct_1+Ct_3$ y $\mathfrak{c}=C(t_2-t_3)$. Consideremos el morfismo $f_0:X_0\to X_1$ que, geométricamente, es la «proyección de $X_0$ sobre $X_1$ desde la recta del infinito $D$ definida por $t_0=t_2=0$». Algebraicamente, $f_0$ corresponde al homomorfismo de álgebras graduadas
\[
  k[t_0,t_2]\longrightarrow k[t_0,t_1,t_2,t_3]/\mathfrak{a}
\]
obtenido, por paso al cociente, de la inyección canónica $k[t_0,t_2]\to k[t_0,t_1,t_2,t_3]$.
Está claro que $f_0$ es un morfismo proyectivo, luego \emph{propio}; lo mismo ocurre, por tanto, con su restricción $f:X\to Y$, donde $Y=D_+(t_0)$ y $X=f_0^{-1}(Y)$. Para ver que $f$ es \emph{plano}, basta observar que $k[t_2]$ es un dominio de ideales principales y que el anillo
\[
  C'=k[t_1,t_2,t_3]/(t_2-t_3)((t_1)+(t_3))
\]
es un $k[t_2]$-módulo sin torsión \sref[0]{0.6.3.4}.
Para todo $y\in Y$ distinto del punto $y_0$ definido por $t_2=0$, $f^{-1}(y)$ tiene entonces dos componentes irreducibles, de dimensiones~0 y~1, mientras que $f^{-1}(y_0)$ solo tiene una componente irreducible, de dimensión~1.

  \item En \sref{IV.12.2.1},~(i), tampoco se puede sustituir la condición de que el conjunto $E'(y)\supset E(y)$ de las dimensiones de los ciclos primos asociados a $\sh{F}_y$ esté contenido en $\Phi$ por la condición de que sea \emph{igual a} $\Phi$.
En efecto, sea $k$ un cuerpo, sea $A$ el anillo de polinomios $k[t]$, y sea $B$ el subanillo $k[t^4,t^3u,tu^3,u^4]$ del anillo de polinomios $k[t,u]$ (donde $t,u$ son indeterminadas). El anillo $B$ no es íntegramente cerrado: el elemento $t^2u^2$ pertenece a su clausura integral, pero no a $B$. Si $\mathfrak{m}$ es el ideal maximal de $B$ generado por $t^4,t^3u,tu^3,u^4$, entonces $\mathfrak{m}$ es un ideal primo inmerso asociado de $B/Bt^4$.\footnote{La fuente impresa dice $A/At^4$; la construcción y el cálculo de la fibra que sigue requieren $B/Bt^4$.}
Pongamos $Y=\Spec(A)$, $X=\Spec(B)$, y sea $f:X\to Y$ el morfismo correspondiente al homomorfismo $A\to B$ que envía $t$ en $t^4$.
Como este homomorfismo hace de $B$ un $A$-módulo sin torsión, $f$ es plano \sref[0]{0.6.3.4}.
Si $y$ es el punto de $Y$ correspondiente al ideal maximal $tA$, el preesquema $f^{-1}(y)$ es irreducible y de dimensión~1, pero tiene un ciclo primo inmerso asociado de dimensión~0. En cambio, si $y'$ es el punto genérico de $Y$, entonces $f^{-1}(y')$ no tiene ningún ciclo primo inmerso asociado, pues es el espectro de una $k(y')$-álgebra íntegra de dimensión~1. En este ejemplo, $f$ es un morfismo afín no propio; se obtiene inmediatamente un ejemplo análogo en el que $f$ es propio y plano considerando $X$ como un subconjunto abierto denso en todas partes de un esquema proyectivo sobre $Y$ \sref[II]{II.5.3.4} y \sref[II]{II.5.3.2}, o procediendo directamente como en la Observación~(ii).

Las Observaciones~(ii) y~(iii) muestran la necesidad de incluir en \sref{IV.12.2.1},~(iv) la condición de que $\sh{F}_y$ no tenga ningún ciclo primo inmerso asociado.

  \item En \sref{IV.12.2.1},~(xi), no se puede suprimir la hipótesis de que $f$ sea propio.
En efecto, sea $k$ un cuerpo, sea $Y$ la «recta afín», el espectro del anillo de polinomios $k[t]$ ($t$ una indeterminada), y sea $X$ la unión disjunta de $Y$ y del abierto complementario del punto cerrado $y$ de $Y$ definido por $t=0$. Está claro que el morfismo $f:X\to Y$, que es igual a las inyecciones canónicas en las dos componentes de $X$, es plano y que, tomando $\sh{F}=\sh{O}_X$, el Módulo $\sh{F}_z$ no tiene ningún ciclo primo inmerso asociado para ningún $z\in Y$. Sin embargo, se tiene $l(y)=1$ y $l(z)=2$ para todo $z\neq y$.

  \item En \sref{IV.12.2.1},~(xi), tampoco se puede suprimir la hipótesis de que $\sh{F}_y$ no tenga ningún ciclo primo inmerso asociado.
Lo muestra el ejemplo de la Observación~(ii): se ve inmediatamente que, con $\sh{F}=\sh{O}_X$, $l(y_0)=1$ y $l(y)=2$ para $y\neq y_0$ en $Y$.

  \item No sabemos si, en \sref{IV.12.2.1},~(xi), se puede sustituir la hipótesis de que $\sh{F}$ es $f$-plano por la hipótesis de que $\operatorname{Supp}(\sh{F})=X$ y $f$ es universalmente abierto.
Retomando la demostración de \sref{IV.12.2.1.4}, se ve (utilizando también \sref{IV.14.3.6}) que habría que resolver la cuestión siguiente: sea $A$ un anillo local noetheriano, $M$ un $A$-módulo de tipo finito,
\oldpage[IV-3]{183}
$t$ un elemento del ideal maximal $\mathfrak{m}$ de $A$ que no está contenido en ningún ideal primo minimal de $A$; si existe uno de estos ideales primos minimales $\mathfrak{p}$ tal que $\dim(A/\mathfrak{p})=1$, ¿es cierto que $\mathfrak{m}$ es un ideal primo asociado de $M/tM$ (sin suponer que $t$ es $M$-regular)?

Obsérvese, sin embargo, que en \sref{IV.12.2.1},~(xi) no se puede suprimir sin más la hipótesis de que $\sh{F}$ es $f$-plano. Tomemos aquí el $P$ de la Observación~(ii), y sea en $P$ el subpreesquema cerrado $X_0$ que es la unión de las tres rectas $X_1,X_2,X_3$ definidas respectivamente por $t_1=t_3=0$, $t_1-t_0=t_3=0$ y $t_2=t_3=0$. Definamos como antes la proyección $f_0$ de $X_0$ sobre $X_1$, y su restricción $f:X\to Y$, donde $Y=D_+(t_0)$ es la recta afín y $X=f_0^{-1}(Y)$. El morfismo $f$ es propio pero no plano y, tomando $\sh{F}=\sh{O}_X$, el Módulo $\sh{F}_y$ no tiene ningún ciclo primo inmerso asociado para ningún $y\in Y$. Pero se tiene $l(y_0)=1$ y $l(y)=2$ para $y\neq y_0$ en $Y$.
\end{enumerate}
\end{remarks}

\begin{theorem}[12.2.4]
\label{IV.12.2.4}
Sea $f:X\to Y$ un morfismo propio y plano de presentación finita, y sea $k$ un entero $\geqslant 1$.
Los siguientes subconjuntos de $Y$ son abiertos:
\begin{enumerate}[label=(\roman*)]
  \item El conjunto de los $y\in Y$ tales que $X_y$ posee la propiedad $(\mathrm{S}_k)$.
  \item El conjunto de los $y\in Y$ tales que $X_y$ satisface geométricamente la propiedad $(\mathrm{R}_k)$, es equidimensional en todo punto y no tiene ningún ciclo primo inmerso asociado.
  \item El conjunto de los $y\in Y$ tales que $X_y$ es geométricamente regular (es decir, liso sobre $k(y)$).
  \item El conjunto de los $y\in Y$ tales que $X_y$ es geométricamente normal.
  \item El conjunto de los $y\in Y$ tales que $X_y$ es geométricamente reducido.
  \item El conjunto de los $y\in Y$ tales que $X_y$ es geométricamente reducido y el número de componentes conexas geométricas de $X_y$ es igual a $k$.
  \item El conjunto de los $y\in Y$ tales que $X_y$ es geométricamente puntualmente íntegro \sref{IV.4.6.9}.
  \item El conjunto de los $y\in Y$ tales que $X_y$ es geométricamente íntegro.
  \item El conjunto de los $y\in Y$ tales que $X_y$ no tiene ningún ciclo primo inmerso asociado y la multiplicidad total \sref{IV.4.7.4} de $X_y$ sobre $k(y)$ es $\leqslant k$.
\end{enumerate}
\end{theorem}

\begin{proof}
Salvo~(vi), estas afirmaciones son casos particulares de afirmaciones de \sref{IV.12.2.1}, o se deducen de afirmaciones de \sref{IV.12.1.6} como al comienzo de la demostración de \sref{IV.12.2.1}.
Para~(vi), se reduce como siempre (teniendo en cuenta la invariancia del número geométrico de componentes conexas por cambio de base \sref{IV.4.5.6}) al caso en que $Y$ es noetheriano.
El conjunto $U$ de los $y\in Y$ tales que $X_y$ es geométricamente reducido es abierto en $Y$ en virtud de~(v); de \sref[III]{III.7.8.7} y~\sref[III]{III.7.8.6} se sigue entonces que, para todo $y\in U$, existe un entorno $V\subset U$ de $y$ y un entero $m$ tales que, para \emph{todo} $z\in V$, $\Gamma(X_z,\sh{O}_{X_z})$ es isomorfo a $(k(z))^m$; pero, en virtud de \sref[III]{III.4.3.4}, $m$ es entonces el número geométrico de componentes conexas de $X_z$, de donde la conclusión.
Daremos otra demostración de~(vi) en \sref{IV.15.5.9}.
\end{proof}


\subsection{Propiedades cohomológicas locales de las fibras de un morfismo plano localmente de presentación finita}
\label{subsection:IV.12.3}

\begin{lemma}[12.3.1]
\label{IV.12.3.1}
Sea $A$ un anillo, $B$ una $A$-álgebra de presentación finita que es un $A$-módulo plano, y $M$ un $B$-módulo de presentación finita que es un $A$-módulo plano.
Entonces las condiciones siguientes son equivalentes:
\begin{enumerate}[label=\emph{\alph*)}]
  \item $M$ es un $B$-módulo proyectivo.
  \item $M$ es un $B$-módulo plano.
  \item Para todo $y\in\Spec(A)$, $M\otimes_A k(y)$ es un $(B\otimes_A k(y))$-módulo proyectivo.
\end{enumerate}
\end{lemma}

\begin{proof}
\oldpage[IV-3]{184}

La equivalencia de \emph{a)} y \emph{b)} se deduce de Bourbaki, \emph{Alg.\ comm.}, cap.~II, \textsection\,5, n\textsuperscript{o}\,2, cor.~2 del teo.~1.
Como \emph{a)} implica trivialmente \emph{c)}, queda demostrar que \emph{c)} implica \emph{b)}, lo cual resulta del criterio de planitud por las fibras \sref{IV.11.3.10}, aplicado con $g=h$, $f=1_X$.
\end{proof}

\begin{proposition}[12.3.2]
\label{IV.12.3.2}
Sea $A$ un anillo, $B$ una $A$-álgebra de presentación finita que es un $A$-módulo plano, y $M$ un $B$-módulo de presentación finita que es un $A$-módulo plano.
Entonces:
\begin{enumerate}[label=(\roman*)]
  \item Existe una resolución por la izquierda de $M$ mediante $B$-módulos libres de tipo finito.
  \item Se tiene
  \begin{equation}
  \tag{12.3.2.1}
  \label{IV.12.3.2.1}
    \operatorname{dim.proj}_B(M)
    = \sup_{y\in\Spec(A)} \operatorname{dim.proj}_{B\otimes_A k(y)}(M\otimes_A k(y))
    = \operatorname{Tor.dim}_B(M).
  \end{equation}
  Aquí $\operatorname{Tor.dim}_B(M)$ es el menor entero $i$ tal que $\operatorname{Tor}_j^B(M,N)=0$ para todo $j>i$ y todo $B$-módulo $N$ (y es $+\infty$ si no existe tal entero $i$).
\end{enumerate}
\end{proposition}

\begin{proof}
\begin{enumerate}[label=(\roman*)]
  \item Por \sref{IV.8.9.1}, existen un subanillo noetheriano $A_0$ de $A$, una $A_0$-álgebra de tipo finito $B_0$ y un $B_0$-módulo de tipo finito $M_0$ tales que $B\simeq B_0\otimes_{A_0}A$ y $M\simeq M_0\otimes_{A_0}A$.
  Por \sref{IV.11.2.7}, podemos suponer además que $B_0$ y $M_0$ son $A_0$-módulos planos.
  Sea $(L_0)_\bullet$ una resolución por la izquierda de $M_0$ mediante $B_0$-módulos libres de tipo finito.
  Como $M_0$ y los $(L_0)_i$ son $A_0$-módulos planos, $(L_0)_\bullet\otimes_{A_0}A$ es una resolución por la izquierda de $M$ mediante $B$-módulos libres de tipo finito \sref{IV.2.1.10}.

  \item Por \sref[0]{0.17.2.2},~(ii),
  $\operatorname{Tor.dim}_B(M)\leqslant\operatorname{dim.proj}_B(M)$; y la definición de dimensión proyectiva da inmediatamente
  \[
    \operatorname{dim.proj}_{B\otimes_A k(y)}(M\otimes_A k(y))
    \leqslant \operatorname{dim.proj}_B(M)
  \]
  para todo $y\in\Spec(A)$.
  Para demostrar las desigualdades inversas, consideremos una resolución por la izquierda $(L_i)$ de $M$ mediante $B$-módulos libres de tipo finito, y supongamos que $\operatorname{Tor.dim}_B(M)=n$ (resp.\ $\operatorname{dim.proj}_{B\otimes_A k(y)}(M\otimes_A k(y))\leqslant n$ para todo $y\in\Spec(A)$).
  Entonces
  $R=\operatorname{Im}(L_n\to L_{n-1})=\operatorname{Ker}(L_{n-1}\to L_{n-2})$
  es un $B$-módulo de tipo finito y un $A$-módulo plano, por las hipótesis sobre $M$ y $B$ y por \sref{IV.2.1.10}.
  Además,
  \[
    \operatorname{Tor}_{n+1}^B(M,N)=\operatorname{Tor}_1^B(R,N)
  \]
  para todo $B$-módulo $N$ (M, V, 7).
  La hipótesis $\operatorname{Tor.dim}_B(M)=n$ implica, por tanto, $\operatorname{Tor}_1^B(R,N)=0$ para todo $B$-módulo $N$; es decir, $R$ es un $B$-módulo plano, luego proyectivo por \sref{IV.12.3.1}.
  Así, $\operatorname{dim.proj}_B(M)\leqslant n$.
  Por otra parte, la hipótesis $\operatorname{dim.proj}_{B\otimes_A k(y)}(M\otimes_A k(y))\leqslant n$ para todo $y\in\Spec(A)$ implica, al tensorizar por $k(y)$, que en cada una de las sucesiones siguientes (que son exactas por la planitud sobre $A$ de $M$, de los $L_i$ y de $R$ \sref{IV.2.1.10})
  \[
    0\to R\otimes_A k(y)\to L_{n-1}\otimes_A k(y)\to\cdots\to L_0\otimes_A k(y)\to M\otimes_A k(y)\to 0
  \]
  el $(B\otimes_A k(y))$-módulo $R\otimes_A k(y)$ es proyectivo, para todo $y\in\Spec(A)$.
  De nuevo se sigue de \sref{IV.12.3.1} que $R$ es un $B$-módulo proyectivo y, por tanto, que $\operatorname{dim.proj}_B(M)\leqslant n$.
\end{enumerate}
\end{proof}

\begin{proposition}[12.3.3]
\label{IV.12.3.3}
Sea $f:X\to S$ un morfismo localmente de presentación finita, y sea $\sh{L}_\bullet$ un complejo de $\sh{O}_X$-módulos cuasicoherentes de presentación finita; para todo $s\in S$, denotemos por $(\sh{L}_\bullet)_s$ el complejo $\sh{L}_\bullet\otimes_{\sh{O}_S}k(s)$ de $\sh{O}_{X_s}$-módulos de tipo finito.
Supongamos que $\sh{L}_{n-1}$ es $f$-plano.
Entonces el conjunto $U$ de los $x\in X$ tales que $(\sh{H}_n((\sh{L}_\bullet)_{f(x)}))_x=0$ es abierto en $X$.
Si, además, $\sh{L}_n$ es $f$-plano, entonces $\sh{H}_n(\sh{L}_\bullet)|_U=0$.
\end{proposition}

\begin{proof}
Evidentemente basta considerar el caso $n=1$, en el que el complejo $\sh{L}_\bullet$ se reduce a
$0\to\sh{L}_2\xrightarrow{u}\sh{L}_1\xrightarrow{v}\sh{L}_0\to 0$.
Podemos reducir primero al caso en que $S=\Spec(A)$ y $X$
\oldpage[IV-3]{185}
son afines, y después al caso en que $S$ es \emph{noetheriano}. En efecto, por \sref{IV.8.9.1} y \sref{IV.8.5.2}, existen un preesquema noetheriano $S'=\Spec(A')$, donde $A'$ es un subanillo de $A$, un morfismo de tipo finito $f':X'\to S'$ y tres $\sh{O}_{X'}$-módulos coherentes $\sh{L}'_i$ ($0\leqslant i\leqslant 2$), tales que
$X=X'\times_{S'}S$, $f=f'_{(S)}$ y $\sh{L}_i=\sh{L}'_i\otimes_{S'}S$, junto con dos homomorfismos $u'$, $v'$ tales que $u=u'\otimes 1$, $v=v'\otimes 1$ y $v'\circ u'=0$.
Podemos suponer además que $\sh{L}'_0$ es $f'$-plano \sref{IV.11.2.7}, y que $\sh{L}'_1$ es $f'$-plano cuando se supone que $\sh{L}_1$ es $f$-plano.
Por planitud fiel, la hipótesis
$(\sh{H}_n((\sh{L}_\bullet)_{f(x)}))_x=0$
es equivalente a
$(\sh{H}_n((\sh{L}'_\bullet)_{f'(x')}))_{x'}=0$,
donde $x'$ es la proyección de $x$ en $X'$.
Así, si $U'$ es el conjunto de los $x'\in X'$ tales que
$(\sh{H}_n((\sh{L}'_\bullet)_{f'(x')}))_{x'}=0$,
entonces $U$ es la imagen inversa de $U'$ por la proyección $X\to X'$; esto reduce la primera afirmación al caso en que $S$ es noetheriano.

Para la segunda afirmación, obsérvese también que $A$ es el límite inductivo filtrante de sus subanillos $A_\lambda$ que son $A'$-álgebras de tipo finito.
Pongamos $S_\lambda=\Spec(A_\lambda)$,
$X_\lambda=X'\times_{S'}S_\lambda$, y
$\sh{L}^{(\lambda)}_i=\sh{L}'_i\otimes_{S'}S_\lambda$; sean
$u_\lambda=u'\otimes 1:\sh{L}_2^{(\lambda)}\to\sh{L}_1^{(\lambda)}$ y
$v_\lambda=v'\otimes 1:\sh{L}_1^{(\lambda)}\to\sh{L}_0^{(\lambda)}$, de modo que
$\sh{H}_1(\sh{L}_\bullet^{(\lambda)})=\operatorname{Ker}(v_\lambda)/\operatorname{Im}(u_\lambda)$.
Como los límites inductivos filtrantes son exactos en la categoría de grupos conmutativos,
\[
  \operatorname{Ker}(v)=\varinjlim_\lambda\operatorname{Ker}(v_\lambda),\qquad
  \operatorname{Im}(u)=\varinjlim_\lambda\operatorname{Im}(u_\lambda),\qquad
  \sh{H}_1(\sh{L}_\bullet)=\varinjlim_\lambda\sh{H}_1(\sh{L}_\bullet^{(\lambda)}).
\]
Si se ha supuesto que $\sh{L}_1$ es $f$-plano y, tras reducir al caso $X=U$, se ha demostrado que
$\sh{H}_1(\sh{L}_\bullet^{(\lambda)})=0$ para todo $\lambda$, se sigue que
$\sh{H}_1(\sh{L}_\bullet)=0$.

\medskip
I) Supongamos en lo sucesivo que $S$ es noetheriano.
Sabemos (sin ninguna hipótesis de planitud sobre $\sh{L}_0$) que el conjunto $U$ es \emph{construible} en $X$ \sref{IV.9.9.6}.
Por \EGAReadLink{0.9.2.5}{($\mathbf{0}_{\mathrm{III}}$,~9.2.5)}, queda demostrar que toda \emph{generización} $x'$ en $X$ de un punto $x\in U$ pertenece también a $U$.
El método expuesto en \sref{IV.12.0.2} se aplica sin modificación (teniendo en cuenta que, para un preesquema $Z$ sobre un cuerpo $k$ y una extensión $k'$ de $k$, la proyección $Z\otimes_k k'\to Z$ es fielmente plana).
Podemos suponer, por tanto, que $S$ es el espectro de un anillo de valoración discreta $A$, de uniformizante $t$, que $x$ está situado sobre el punto cerrado de $S$ y que $x'$ está situado sobre el punto genérico de $S$.
La hipótesis de que $\sh{L}_0$ es $f$-plano implica que $t$ es $\sh{L}_0$-regular \EGAReadLink{0.6.3.4}{($\mathbf{0}_{\mathrm{I}}$,~6.3.4)}.
Se reduce así a demostrar el lema siguiente, en el que $B$, $M$, $N$ y $P$ serán sustituidos respectivamente por $\sh{O}_x$, $(\sh{L}_2)_x$, $(\sh{L}_1)_x$ y $(\sh{L}_0)_x$.

\begin{lemma}[12.3.3.1]
\label{IV.12.3.3.1}
Sea $B$ un anillo local noetheriano, sean $M$, $N$ y $P$ tres $B$-módulos, con $N$ de tipo finito, y sean
$M\xrightarrow{u}N\xrightarrow{v}P$ dos homomorfismos tales que $v\circ u=0$.
Sea $t$ un elemento del ideal maximal de $B$ que es $P$-regular.
Si la sucesión
\begin{equation}
\tag{12.3.3.2}
\label{IV.12.3.3.2}
M/tM \xrightarrow{u\otimes 1_{B/tB}} N/tN \xrightarrow{v\otimes 1_{B/tB}} P/tP
\end{equation}
es exacta, entonces la sucesión $M\xrightarrow{u}N\xrightarrow{v}P$ es exacta.
\end{lemma}

Obsérvese primero que, si se sustituye $M$ por su imagen $Q$ en $N$ y $u$ por la inyección $j:Q\to N$, la imagen de $j\otimes 1$ es la misma que la de $u\otimes 1$; por tanto, sin cambiar ni la hipótesis ni la conclusión, podemos suponer que $u$ es inyectivo.
Por otra parte, si $R$ es la imagen de $v$ y $\rho:N\to R$ la sobreyección canónica, entonces $t$ es evidentemente $R$-regular, $\rho\circ u=0$, y el núcleo de $\rho\otimes 1$ está contenido en el de $v\otimes 1$.
Por tanto, es igual a él si \sref{IV.12.3.3.2} es exacta; como además $\operatorname{Ker}(\rho)=\operatorname{Ker}(v)$,
\oldpage[IV-3]{186}
vemos que, para demostrar el lema, también podemos suponer que $v$ es sobreyectivo.
El lema es entonces consecuencia del siguiente.

\begin{lemma}[12.3.3.3]
\label{IV.12.3.3.3}
Sea $B$ un anillo, sean $M$, $N$ y $P$ tres $B$-módulos, y sean
$u:M\to N$, $v:N\to P$ dos homomorfismos tales que $u$ es inyectivo, $v$ es sobreyectivo y $v\circ u=0$.
Sea $t$ un elemento de $B$ que es $P$-regular.
Entonces se tiene, salvo isomorfismo canónico,
\begin{equation}
\tag{12.3.3.4}
\label{IV.12.3.3.4}
\operatorname{Ker}(v\otimes 1)/\operatorname{Im}(u\otimes 1) = (\operatorname{Ker}(v)/\operatorname{Im}(u))\otimes_B(B/tB)
\end{equation}
\end{lemma}

En efecto, la hipótesis de que la sucesión \sref{IV.12.3.3.2} es exacta implica entonces que $(\operatorname{Ker}\,v/\operatorname{Im}\,u)\otimes_B(B/tB)=0$ y, como $t$ pertenece al ideal maximal de $B$ y $\operatorname{Ker}(v)$ es un $B$-módulo de tipo finito (pues $B$ se supone noetheriano en \sref{IV.12.3.3.1}), el lema de Nakayama demuestra que $\operatorname{Ker}(v)=\operatorname{Im}(u)$.

Queda, pues, demostrar \sref{IV.12.3.3.3}.
Pongamos $u'=u\otimes 1$, $v'=v\otimes 1$, $Z=\operatorname{Ker}(v)$, $H=Z/\operatorname{Im}(u)$, de modo que las sucesiones
\begin{gather*}
  0\to M\to Z\to H\to 0\\
  0\to Z\to N\xrightarrow{v} P\to 0
\end{gather*}
dan, al tensorizar por $B/tB$ y utilizar el lema \sref{IV.3.4.1.4} y el hecho de que $t$ es $P$-regular, las sucesiones exactas
\begin{gather*}
  M/tM \xrightarrow{w'} Z/tZ \to H/tH \to 0\\
  0\to Z/tZ \to N/tN \xrightarrow{v'} P/tP \to 0
\end{gather*}
de donde $\operatorname{Ker}(v')=Z/tZ$, y como $\operatorname{Im}(w')=\operatorname{Im}(u')$, se obtiene \sref{IV.12.3.3.4}.

\medskip
II) Supongamos ahora además que $\sh{L}_1$ es $f$-plano; sustituyendo asimismo $X$ por $U$, se puede suponer que $X=U$.
Se trata de ver que, para todo $x\in X$, $(\sh{H}_1(\sh{L}_\bullet))_x=0$.
Pongamos $s=f(x)$, y sea $\mathfrak{n}$ el ideal $\mathfrak{m}_s\sh{O}_{X,x}$, que está contenido en el ideal maximal de $\sh{O}_{X,x}$.
Como $(\sh{H}_1(\sh{L}_\bullet))_x$ es un $\sh{O}_{X,x}$-módulo de tipo finito ($X$ es localmente noetheriano), es separado para la topología $\mathfrak{n}$-ádica \EGAReadLink{0.7.3.5}{($\mathbf{0}_{\mathrm{I}}$,~7.3.5)} y basta, por tanto, demostrar que su completado separado para esta topología es $0$.
Ahora bien, en virtud de \sref[III]{III.7.4.7.2}, este completado separado es igual a
$\varprojlim_k\sh{H}_1(\sh{L}_\bullet\otimes_{\sh{O}_S}(\sh{O}_{S,s}/\mathfrak{m}_s^{k+1}))$; bastará, por tanto, demostrar que cada término de este sistema proyectivo es cero.
Esto es cierto por hipótesis para $k=0$; razonamos, pues, por inducción sobre $k$.
La conclusión se deducirá del lema más general siguiente, que aplicamos con $A=\sh{O}_{S,s}/\mathfrak{m}_s^{k+1}$ y $\mathfrak{J}$ igual al ideal maximal $\mathfrak{m}_s/\mathfrak{m}_s^{k+1}$ de $A$.

\begin{lemma}[12.3.3.5]
\label{IV.12.3.3.5}
Sea $A$ un anillo, sea $\mathfrak{J}$ un ideal nilpotente de $A$, y sea
$L_\bullet:P\xrightarrow{u}Q\xrightarrow{v}R$ un complejo de $A$-módulos.
Pongamos $A_0=A/\mathfrak{J}$ y
$L_\bullet^{(0)}=L_\bullet\otimes_A A_0:P_0\xrightarrow{u_0}Q_0\xrightarrow{v_0}R_0$.
Supongamos que $\operatorname{gr}_{\mathfrak{J}}^{\bullet}(A)$ es un $A_0$-módulo plano, y que $Q$ y $R$ son $A$-módulos planos.
Entonces la relación $\mathrm{H}_1(L_\bullet^{(0)})=0$ implica $\mathrm{H}_1(L_\bullet)=0$.
\end{lemma}

Pongamos
$L_\bullet^{(k)}=L_\bullet\otimes_A(A/\mathfrak{J}^{k+1}):P_k\xrightarrow{u_k}Q_k\xrightarrow{v_k}R_k$.
Como existe un entero $k\geqslant 1$ tal que $L_\bullet^{(k)}=L_\bullet$, bastará demostrar, por inducción sobre $k$, que la sucesión
$P_k\xrightarrow{u_k}Q_k\xrightarrow{v_k}R_k$
\oldpage[IV-3]{187}
es exacta (esta afirmación resulta de la hipótesis para $k=0$).
Sea entonces $x\in Q_k$ un elemento tal que $v_k(x)=0$.
Como, por la hipótesis de inducción, la sucesión
$P_{k-1}\xrightarrow{u_{k-1}}Q_{k-1}\xrightarrow{v_{k-1}}R_{k-1}$ es exacta, la imagen canónica de $x$ en
$Q_{k-1}=Q_k/(\mathfrak{J}^kQ/\mathfrak{J}^{k+1}Q)$ pertenece a $\operatorname{Im}(u_{k-1})$.
Existe, por tanto, $z\in P_k$ tal que
$x=u_k(z)+x'$, con $x'\in\mathfrak{J}^kQ/\mathfrak{J}^{k+1}Q$.
Como $v_k\circ u_k=0$, la relación $v_k(x)=0$ implica $v_k(x')=0$; además, se tiene evidentemente
$v_k(x')\in\mathfrak{J}^kR/\mathfrak{J}^{k+1}R$.
Todo se reduce, por tanto, a demostrar que la sucesión
\[
  \mathfrak{J}^kP/\mathfrak{J}^{k+1}P
  \xrightarrow{u'}
  \mathfrak{J}^kQ/\mathfrak{J}^{k+1}Q
  \xrightarrow{v'}
  \mathfrak{J}^kR/\mathfrak{J}^{k+1}R
\]
donde $u'$ y $v'$ son inducidos por $u$ y $v$ por restricción y paso a los cocientes, es exacta.
Ahora bien, por hipótesis, $\mathfrak{J}^k/\mathfrak{J}^{k+1}$ es un $(A/\mathfrak{J})$-módulo plano, por lo que la sucesión
\[
  (P/\mathfrak{J}P)\otimes_{A/\mathfrak{J}}(\mathfrak{J}^k/\mathfrak{J}^{k+1})
  \xrightarrow{u''}
  (Q/\mathfrak{J}Q)\otimes_{A/\mathfrak{J}}(\mathfrak{J}^k/\mathfrak{J}^{k+1})
  \xrightarrow{v''}
  (R/\mathfrak{J}R)\otimes_{A/\mathfrak{J}}(\mathfrak{J}^k/\mathfrak{J}^{k+1})
\]
es exacta.
Pero
$(Q/\mathfrak{J}Q)\otimes_{A/\mathfrak{J}}(\mathfrak{J}^k/\mathfrak{J}^{k+1})
=Q\otimes_A(\mathfrak{J}^k/\mathfrak{J}^{k+1})$
se identifica con $\mathfrak{J}^kQ/\mathfrak{J}^{k+1}Q$ por la planitud de $Q$ sobre $A$; análogamente,
$(R/\mathfrak{J}R)\otimes_{A/\mathfrak{J}}(\mathfrak{J}^k/\mathfrak{J}^{k+1})$
se identifica con $\mathfrak{J}^kR/\mathfrak{J}^{k+1}R$.
Finalmente, la imagen de $u'$ se identifica con la de $u''$, lo que completa la demostración.
\end{proof}

\begin{corollary}[12.3.4]
\label{IV.12.3.4}
Sea $f:X\to S$ un morfismo plano y localmente de presentación finita, y sean $\sh{F}$ y $\sh{G}$ dos $\sh{O}_X$-módulos de presentación finita que son $f$-planos.
Sea $n$ un entero $\geq 0$, y sea $U$ el conjunto de los $x\in X$ tales que $(\mathscr{E}\!xt^n_{\sh{O}_{X_{f(x)}}}(\sh{F}_{f(x)},\sh{G}_{f(x)}))_x=0$ (resp.\ $(\mathscr{T}\!or_n^{\sh{O}_{X_{f(x)}}}(\sh{F}_{f(x)},\sh{G}_{f(x)}))_x=0$).
Entonces $U$ es abierto, y se tiene
$\mathscr{E}\!xt^n_{\sh{O}_X}(\sh{F},\sh{G})|_U=0$
(resp.\ $\mathscr{T}\!or_n^{\sh{O}_X}(\sh{F},\sh{G})|_U=0$).
\end{corollary}

\begin{proof}
Podemos evidentemente restringirnos al caso en que $S$ y $X$ son afines.
Entonces \sref{IV.12.3.2} proporciona una resolución por la izquierda $\sh{L}_\bullet=(\sh{L}_i)$ de $\sh{F}$ formada por $\sh{O}_X$-módulos \emph{libres de tipo finito}.
Si ponemos $\sh{M}^\bullet=\mathscr{H}\!om(\sh{L}_\bullet,\sh{G})$ (resp.\ $\sh{N}_\bullet=\sh{L}_\bullet\otimes_{\sh{O}_X}\sh{G}$), resulta que cada $\sh{M}^i$ (resp.\ $\sh{N}_i$) es isomorfo a un $\sh{O}_X$-módulo de la forma $\sh{G}^m$; por tanto, los $\sh{M}^i$ y los $\sh{N}_i$ son $\sh{O}_X$-módulos de presentación finita y son $f$-planos.
Además, para todo cambio de base $S'\to S$, si ponemos
$X'=X\times_S S'$, $\sh{F}'=\sh{F}\otimes_S S'$, $\sh{G}'=\sh{G}\otimes_S S'$, y
$\sh{L}'_\bullet=\sh{L}_\bullet\otimes_S S'$, entonces $\sh{L}'_\bullet$ sigue siendo una resolución por la izquierda de $\sh{F}'$ mediante $\sh{O}_{X'}$-módulos libres de tipo finito \sref{IV.2.1.10}; además,
$\sh{M}'^\bullet=\sh{M}^\bullet\otimes_S S'$ es igual a $\mathscr{H}\!om(\sh{L}'_\bullet,\sh{G}')$
(resp.\ $\sh{N}'_\bullet=\sh{N}_\bullet\otimes_S S'$ es igual a $\sh{L}'_\bullet\otimes_{\sh{O}_{X'}}\sh{G}'$), por lo que precede.
En particular, para todo $s\in S$, se tiene
$\mathscr{E}\!xt^n_{\sh{O}_{X_s}}(\sh{F}_s,\sh{G}_s)=\sh{H}^n((\sh{M}^\bullet)_s)$
(resp.\ $\mathscr{T}\!or_n^{\sh{O}_{X_s}}(\sh{F}_s,\sh{G}_s)=\sh{H}_n((\sh{N}_\bullet)_s)$).
Aplicando \sref{IV.12.3.3} a los complejos $f$-planos $\sh{M}^\bullet$ y $\sh{N}_\bullet$, se deduce inmediatamente el corolario.
\end{proof}

\section{Morfismos equidimensionales}
\label{section:IV.13}

\oldpage[IV-3]{187}

Esta sección está dedicada al estudio de la variación de la \emph{dimensión} de las fibras de un morfismo localmente de tipo finito $f:X\to Y$ (que ya intervino a propósito de la «fórmula de dimensión» en~\sref{IV.5.5} y~\sref{IV.5.6}).
Demostramos primero \sref{IV.13.1.3} que la función $x\mapsto\dim_x f^{-1}(f(x))$ es siempre \emph{semicontinua superiormente} en $X$ (teorema de semicontinuidad de Chevalley).
Estudiamos después más particularmente los morfismos, llamados «equidimensionales», para los cuales esta función es \emph{localmente constante}.
\oldpage[IV-3]{188}
Desafortunadamente, la noción de morfismo equidimensional no es estable por cambio de base; por ello, en muchas cuestiones es más conveniente trabajar con la noción de morfismo universalmente abierto, cuyo estudio constituye el objeto de \textsection\textsection~14 y~15.

\subsection{Teorema de semicontinuidad de Chevalley}
\label{subsection:IV.13.1}

\begin{lemma}[13.1.1]
\label{IV.13.1.1}
Sea $Y$ un preesquema irreducible localmente noetheriano, $X$ un preesquema irreducible y $f:X\to Y$ un morfismo dominante de tipo finito.
Sea $\xi$ (resp.\ $\eta$) el punto genérico de $X$ (resp.\ de $Y$), y sea $e=\dim(f^{-1}(\eta))$.
Entonces, para todo $x\in X$, todas las componentes irreducibles de $f^{-1}(f(x))$ son de dimensión $\geq e$.
\end{lemma}

\begin{proof}
La afirmación es inmediata cuando, para todo $y\in f(X)$, $\sh{O}_y$ es un anillo \emph{universalmente catenario}: en efecto, si $x$ es el punto genérico de una componente irreducible $Z$ de $f^{-1}(y)$, de \sref{IV.5.6.5}, combinado con \sref{IV.5.2.1}, se sigue que $e+\dim(\sh{O}_y)=\dim(Z)+\dim(\sh{O}_x)$; pero, por \sref[0]{0.16.3.9}, se tiene $\dim(\sh{O}_x)\leq\dim(\sh{O}_y)$, de donde la conclusión en este caso.

Reduzcamos ahora el caso general a este caso particular.
La cuestión es evidentemente local sobre $Y$ y, teniendo en cuenta \sref{IV.4.1.1.3}, también es local sobre $X$; podemos, por tanto, restringirnos al caso en que $Y=\Spec(A)$ y $X=\Spec(B)$ son afines e irreducibles, siendo $B$ una $A$-álgebra de tipo finito.
Además, por \sref{IV.1.5.4}, podemos suponer que $X$ e $Y$ son reducidos, de modo que $A$ y $B$ son dominios íntegros.
Como $f$ es dominante, $A$ es entonces un \emph{subanillo} de $B$.
Consideremos $A$ como límite inductivo de sus sub-$\bb{Z}$-álgebras de tipo finito; de \sref{IV.8.9.1} se sigue entonces que existe una tal subálgebra $A_0$ y una $A_0$-álgebra de tipo finito $B_0$ tales que $B=B_0\otimes_{A_0}A$.
Pongamos $Y_0=\Spec(A_0)$, $X_0=\Spec(B_0)$, y sea $f_0:X_0\to Y_0$ el morfismo correspondiente al homomorfismo $A_0\to B_0$, de modo que $f=f_0\times 1_Y$.
No es evidente \emph{a priori} que el preesquema $X_0$ sea íntegro, pero veremos que podemos reducirnos a este caso.
Sea $\eta_0$ el punto genérico de $Y_0$, de modo que, si $g$ denota el morfismo $Y\to Y_0$, se tiene $g(\eta)=\eta_0$; por transitividad de las fibras \sref[I]{I.3.6.4}, se tiene $f^{-1}(\eta)=f_0^{-1}(\eta_0)\otimes_{k(\eta_0)}k(\eta)$ y, como $f^{-1}(\eta)$ es irreducible por hipótesis, también lo es $f_0^{-1}(\eta_0)$ \sref{IV.4.4.1}.
Nuestra afirmación resultará entonces del lema siguiente:

\begin{lemma}[13.1.2]
\label{IV.13.1.2}
Sean $Y_0$, $Y$ dos preesquemas íntegros de puntos genéricos $\eta_0$, $\eta$, sea $g:Y\to Y_0$ un morfismo dominante, y sea $f_0:X_0\to Y_0$ un morfismo dominante tal que $f_0^{-1}(\eta_0)$ es irreducible.
Sea $X_0'$ la única componente irreducible de $X_0$ que corta a $f_0^{-1}(\eta_0)$ \EGAReadLink{0.2.1.8}{($\mathbf{0}_{\mathrm{I}}$,~2.1.8)}, y denotemos también por $X_0'$ el subpreesquema cerrado reducido de $X_0$ cuyo espacio subyacente es $X_0'$.
Supongamos que el preesquema $X=X_0\times_{Y_0}Y$ es íntegro; entonces $X$ es isomorfo a $X_0'\times_{Y_0}Y$.
\end{lemma}

En efecto, si $j_0:X_0'\hookrightarrow X_0$ es la inyección canónica, que es una inmersión cerrada, $j=j_0\times 1_Y:X_0'\times_{Y_0}Y\to X_0\times_{Y_0}Y=X$ es una inmersión cerrada.
Por otra parte, $X_0'$ contiene la fibra $f_0^{-1}(\eta_0)$, de modo que $X_0'\times_{Y_0}Y$ contiene a $f^{-1}(\eta)$; obsérvese también que $f^{-1}(\eta)$ no es vacío \sref[I]{I.3.4.7}, por lo que contiene el punto genérico de $X$; en consecuencia, la imagen de $j$ es necesariamente todo $X$.
Pero, como $X$ es íntegro, el único subpreesquema cerrado de $X$ cuyo espacio subyacente es $X$ es el propio $X$, de modo que $j$ es un isomorfismo.

Establecido este lema, podemos suponer, por tanto, que $X_0$ es íntegro; para todo $y_0\in Y_0$, $\sh{O}_{y_0}$ es una $\bb{Z}$-álgebra esencialmente de tipo finito, luego un anillo universalmente
\oldpage[IV-3]{189}
catenario \sref{IV.5.6.4}; en consecuencia, para todo $y_0\in f_0(X_0)$, todas las componentes irreducibles de $f_0^{-1}(y_0)$\footnote{La fuente impresa dice aquí $f_0^{-1}(y)$; el punto cuantificado y el argumento requieren $f_0^{-1}(y_0)$.} tienen dimensión $\geq e$, pues sabemos que $\dim(f_0^{-1}(\eta_0))=e$ por transitividad de las fibras \sref{IV.4.1.4}.
Para todo $y\in f(X)$, se tiene entonces $y_0=g(y)\in f_0(X_0)$ y, por transitividad de las fibras y \sref{IV.4.2.7}, esto completa la demostración.
\end{proof}

\begin{theorem}[13.1.3]
\label{IV.13.1.3}
\textnormal{(Chevalley).} ---
Sea $f:X\to Y$ un morfismo localmente de tipo finito.
Para todo entero $n$, el conjunto $\mathrm{F}_n(X)$ de los $x\in X$ tales que $\dim_x(f^{-1}(f(x)))\geq n$ es cerrado; en otras palabras, la función $x\mapsto\dim_x(f^{-1}(f(x)))$ es semicontinua superiormente en $X$.
\end{theorem}

\begin{proof}
I) Supongamos primero que $f$ es localmente de presentación finita.
La cuestión es evidentemente local sobre $X$ y sobre $Y$, y podemos suponer, por tanto, que $Y=\Spec(A)$, $X=\Spec(B)$ son afines, siendo $B$ una $A$-álgebra de presentación finita.
Sabemos entonces por \sref{IV.8.9.1} que existen un subanillo noetheriano $A_0$ de $A$ y una $A_0$-álgebra de tipo finito $B_0$ tales que, poniendo $Y_0=\Spec(A_0)$ y $X_0=\Spec(B_0)$, se tiene
\[
X=X_0\times_{Y_0}Y,\qquad f=f_0\times 1_Y,
\]
donde $f_0:X_0\to Y_0$ es el morfismo correspondiente al homomorfismo $A_0\to B_0$.
Sea $g:Y\to Y_0$ el morfismo correspondiente a la inyección canónica $A_0\to A$, sea $y$ un punto de $Y$, y pongamos $y_0=g(y)$; sabemos que $f^{-1}(y)=f_0^{-1}(y_0)\otimes_{k(y_0)}k(y)$, y de \sref{IV.4.2.7} se sigue que, si $p$ es la proyección canónica $f^{-1}(y)\to f_0^{-1}(y_0)$, las componentes irreducibles de $f^{-1}(y)$ son las componentes irreducibles de los espacios $p^{-1}(Z_0)$, donde $Z_0$ recorre el conjunto de las componentes irreducibles de $f_0^{-1}(y_0)$, y cada componente irreducible de $p^{-1}(Z_0)$ domina $Z_0$ y tiene dimensión $\dim(Z_0)$.
Teniendo en cuenta \sref[0]{0.14.1.5}, vemos, por tanto, que si $g':X\to X_0$ es la proyección canónica, $x$ un punto de $X$ y $x_0=g'(x)$, se tiene $\dim_x(f^{-1}(f(x)))=\dim_{x_0}(f_0^{-1}(f_0(x_0)))$, de donde $\mathrm{F}_n(X)={g'}^{-1}(\mathrm{F}_n(X_0))$, y se reduce a demostrar el teorema cuando $Y$ es \emph{noetheriano}, lo que supondremos en adelante.

Podemos suponer evidentemente que $X$ e $Y$ son reducidos \sref{IV.1.5.4}.
Considerando el conjunto de las partes cerradas $Y'$ de $Y$ tales que el teorema es cierto para el subpreesquema cerrado de $Y$ cuyo espacio subyacente es $Y'$ y para $X'=f^{-1}(Y')$, podemos razonar por recurrencia noetheriana \EGAReadLink{0.2.2.2}{($\mathbf{0}_{\mathrm{I}}$,~2.2.2)} y suponer que el teorema es cierto para toda parte cerrada $Y'\neq Y$ de $Y$.
Si $X_i$ ($1\leq i\leq m$) son los subpreesquemas cerrados reducidos de $X$ cuyos espacios subyacentes son las componentes irreducibles de $X$, se tiene $\mathrm{F}_n(X)=\bigcup_i\mathrm{F}_n(X_i)$ en virtud de \sref[0]{0.14.1.5}, y podemos, por tanto, restringirnos a demostrar el teorema para cada $X_i$; en otras palabras, podemos suponer que $X$ es \emph{irreducible}.
Si $Z$ es el subpreesquema cerrado de $Y$ cuyo espacio subyacente es $\overline{f(X)}$, $f$ se factoriza como $X\xrightarrow{g}Z\xrightarrow{j}Y$, donde $j$ es la inyección canónica \sref[I]{I.5.2.2} y $g$ es de tipo finito \sref{IV.1.5.4}; basta, por tanto, demostrar el teorema para $Z$ y $g$. Por la hipótesis de recurrencia, se reduce así al caso $Z=Y$, es decir, al caso en que $Y$ es irreducible y $f$ es dominante.
Sea entonces $\eta$ el punto genérico de $Y$ y pongamos $e=\dim(f^{-1}(\eta))$; de \sref{IV.13.1.1} se sigue que, para $n\leq e$, se tiene $\mathrm{F}_n(X)=X$, y en consecuencia podemos restringirnos al caso en que $n>e$.
Pero entonces, por \sref{IV.9.5.6}, existe un entorno abierto $U$ de $\eta$ en $Y$ tal que $\mathrm{F}_n(X)\subset f^{-1}(Y-U)$; como $Y-U\neq Y$, la hipótesis de recurrencia implica que $\mathrm{F}_n(X)$ es cerrado.

II) Pasemos ahora al caso general, escribiendo todavía $S=\Spec(A)$ para el objeto meta y suponiendo $X=\Spec(B)$ afín, con $B$ una $A$-álgebra de tipo finito, luego de la forma
\oldpage[IV-3]{190}
$B=A[T_1,\ldots,T_n]/\mathfrak{J}$.
Sea $(\mathfrak{J}_\lambda)$ la familia de los ideales finitamente generados de $A[T_1,\ldots,T_n]$ contenidos en $\mathfrak{J}$, de modo que $\mathfrak{J}$ es la unión filtrante de los $\mathfrak{J}_\lambda$; si se consideran $X$ y los $X_\lambda=\Spec(A[T_1,\ldots,T_n]/\mathfrak{J}_\lambda)$ como subpreesquemas cerrados de $Z=\Spec(A[T_1,\ldots,T_n])$, se tiene, para los espacios subyacentes, $X=\bigcap_\lambda X_\lambda$.
Si $f_\lambda:X_\lambda\to S$ es el morfismo estructural, se deduce que $f^{-1}(s)=\bigcap_\lambda f_\lambda^{-1}(s)$ para todo $s\in S$; y como los conjuntos $f_\lambda^{-1}(s)$ son cerrados en el espacio noetheriano $Z_s$, existe un índice $\lambda$ (que depende de $s$) tal que $f^{-1}(s)=f_\lambda^{-1}(s)$.
Para todo $x\in X$, pongamos $d(x)=\dim_x(f^{-1}(f(x)))$ y $d_\lambda(x)=\dim_x(f_\lambda^{-1}(f_\lambda(x)))$. Lo que precede muestra que $d(x)=\inf_\lambda d_\lambda(x)$; como las funciones $d_\lambda$ son semicontinuas superiormente por la primera parte de la demostración, también lo es $d$.
\end{proof}

\begin{corollary}[13.1.4]
\label{IV.13.1.4}
Bajo las hipótesis de \sref{IV.13.1.3}, el conjunto de los $x\in X$ tales que $x$ está aislado en $f^{-1}(f(x))$ es abierto en $X$.
\end{corollary}

\begin{proof}
En efecto, es el complementario de $\mathrm{F}_1(X)$ \sref[0]{0.14.1.10}.
\end{proof}

Observemos que recuperamos así, bajo hipótesis más generales, la consecuencia \sref[III]{III.4.4.10} del «Teorema principal» de Zariski.

\begin{corollary}[13.1.5]
\label{IV.13.1.5}
Bajo las hipótesis de \sref{IV.13.1.3}, supongamos además que $f$ es un morfismo cerrado.
Entonces, para todo entero $n$, el conjunto de los $y\in Y$ tales que $\dim(f^{-1}(y))\geq n$ es cerrado; en otras palabras, la aplicación $y\mapsto\dim(f^{-1}(y))$ es semicontinua superiormente; en particular, si $y$ es una especialización de $y'$, se tiene $\dim(f^{-1}(y))\geq\dim(f^{-1}(y'))$.
\end{corollary}

\begin{proof}
En efecto, decir que $\dim(f^{-1}(y))\geq n$ significa que $y\in f(\mathrm{F}_n(X))$ \sref[0]{0.14.1.6}.
\end{proof}

\begin{corollary}[13.1.6]
\label{IV.13.1.6}
Sean $X$, $Y$ dos preesquemas irreducibles, y sea $f:X\to Y$ un morfismo dominante localmente de tipo finito.
Sea $\xi$ (resp.\ $\eta$) el punto genérico de $X$ (resp.\ de $Y$), y sea $e=\dim(f^{-1}(\eta))$.
Entonces, para todo $x\in X$, se tiene $\dim_x(f^{-1}(f(x)))\geq e$.
\end{corollary}

\begin{proof}
En efecto, el conjunto de los $x$ tales que $\dim_x(f^{-1}(f(x)))<e$ es abierto por \sref{IV.13.1.3} y, como no puede contener a $\xi$, es vacío.
\end{proof}

Observemos finalmente el siguiente resultado sencillo:

\begin{proposition}[13.1.7]
\label{IV.13.1.7}
Sea $Y$ un preesquema cuasicompacto y $f:X\to Y$ un morfismo de tipo finito.
Existe un entero $n$ tal que, para todo $y\in Y$, se tiene $\dim(f^{-1}(y))\leq n$.
\end{proposition}

\begin{proof}
Como existe un recubrimiento abierto afín finito $(V_i)$ de $Y$ tal que cada $f^{-1}(V_i)$ es una unión finita de abiertos afines, se reduce inmediatamente al caso en que $Y=\Spec(A)$ y $X=\Spec(B)$ son afines, siendo $B$ una $A$-álgebra de tipo finito.
Si $B$ admite un sistema de $n$ generadores, entonces, para todo $y\in Y$, $B\otimes_A k(y)$ es una $k(y)$-álgebra que admite $n$ generadores, luego $\dim(f^{-1}(y))\leq n$ en virtud de \sref{IV.4.1.1}.
\end{proof}

\subsection{Morfismos equidimensionales: caso de los morfismos dominantes de preesquemas irreducibles}
\label{subsection:IV.13.2}

\begin{env}[13.2.1]
\label{IV.13.2.1}
Sea $Y$ un preesquema irreducible, $X$ un preesquema irreducible, $f:X\to Y$ un morfismo \emph{dominante} localmente de tipo finito; sea $\eta$ el punto genérico de $Y$.
Sabemos por \sref{IV.13.1.6} que, para todo $x\in X$, se tiene $\dim_x(f^{-1}(f(x)))\geq\dim(f^{-1}(\eta))$.
\end{env}

\oldpage[IV-3]{191}
\begin{definition}[13.2.2]
\label{IV.13.2.2}
Bajo las hipótesis de \sref{IV.13.2.1}, se dice que $f$ es \emph{equidimensional en el punto} $x$ (o que $X$ es \emph{equidimensional sobre} $Y$ \emph{en el punto} $x$) si
\[
  \dim_x(f^{-1}(f(x)))=\dim(f^{-1}(\eta)).
\]
Se dice que $f$ es \emph{equidimensional} (o que $X$ es \emph{equidimensional sobre} $Y$) si $f$ es equidimensional en todo punto $x\in X$.
\end{definition}

Del teorema de Chevalley \sref{IV.13.1.3} se sigue que el conjunto de los $x\in X$ en los que $f$ es equidimensional es \emph{abierto} y no vacío.
Además, si $f$ es equidimensional en el punto $x$, todas las componentes irreducibles de $f^{-1}(f(x))$ que contienen a $x$ tienen la \emph{misma dimensión}, pues cada una de ellas tiene dimensión $\geq\dim(f^{-1}(\eta))$ \sref{IV.13.1.6} y $\leq\dim_x(f^{-1}(f(x)))$ en virtud de \sref[0]{0.14.1.5}.

\begin{proposition}[13.2.3]
\label{IV.13.2.3}
Sea $Y$ un preesquema irreducible localmente noetheriano, $X$ un preesquema irreducible, $f:X\to Y$ un morfismo dominante localmente de tipo finito; sea $\eta$ el punto genérico de $Y$, $x$ un punto de $X$, $y=f(x)$, y supongamos que se tiene
\[
\label{IV.13.2.3.1}
\tag{13.2.3.1}
  \dim(\sh{O}_x)=\dim(\sh{O}_y)+\dim(\sh{O}_x\otimes_{\sh{O}_y}k(y)).
\]
Entonces $f$ es equidimensional en el punto $x$.
El recíproco es cierto si se cumple la igualdad en \sref{IV.5.6.5.2}, en particular si $\sh{O}_y$ es universalmente catenario.
\end{proposition}

\begin{proof}
Esto resulta inmediatamente de \sref{IV.5.6.5.2} y de la desigualdad $\delta(x)\geq 0$ \sref{IV.13.1.1}.
\end{proof}

\begin{env}[13.2.4]
\label{IV.13.2.4}
Más generalmente, sea ahora $Y$ un preesquema localmente noetheriano, $f:X\to Y$ un morfismo de tipo finito, $x$ un punto de $X$ y $X'$ una parte cerrada irreducible de $X$ que contiene a $x$; pongamos $y=f(x)$ e $Y'=\overline{f(X')}$, y sea $\eta'$ el punto genérico de $Y'$.
Denotemos también por $X'$ e $Y'$ los subpreesquemas cerrados reducidos de $X$ e $Y$ cuyos espacios subyacentes son respectivamente $X'$ e $Y'$; la restricción $X'\to Y$ de $f$ se factoriza como $X'\xrightarrow{f'}Y'\xrightarrow{j}Y$, donde $j$ es la inyección canónica \sref[I]{I.5.2.2}, y $f'$ es de tipo finito \sref{IV.1.5.4}.
Pongamos, para abreviar,
\[
\label{IV.13.2.4.1}
\tag{13.2.4.1}
  A=\sh{O}_{Y,y},\quad B=\sh{O}_{X,x},\quad A'=\sh{O}_{Y',y},\quad B'=\sh{O}_{X',x}.
\]
La fórmula \sref{IV.5.6.5.2}, aplicada al morfismo dominante $f'$ y a los preesquemas irreducibles $X'$, $Y'$, da
\[
\label{IV.13.2.4.2}
\tag{13.2.4.2}
  \dim(B')\leq\dim(A')+\dim(B'\otimes_{A'}k(y))-(\dim_x(f'^{-1}(y))-\dim(f'^{-1}(\eta')))
\]
Por otra parte, el anillo local $A'$ (resp.\ $B'$, $B'\otimes_{A'}k(y)$) es un cociente de $A$ (resp.\ de $B$, $B\otimes_A k(y)$), luego por \sref[0]{0.16.1.2.1} se tiene
\[
\tag{13.2.4.3}
\label{IV.13.2.4.3}
  \dim(A')\leq\dim(A),\quad \dim(B')\leq\dim(B),\quad \dim(B'\otimes_{A'}k(y))\leq\dim(B\otimes_A k(y)).
\]
Se deduce, por tanto, primero de \sref{IV.13.2.4.3} y \sref[0]{0.16.3.9}
\[
\label{IV.13.2.4.4}
\tag{13.2.4.4}
  \dim(B')\leq\dim(A')+\dim(B'\otimes_{A'}k(y))\leq\dim(A)+\dim(B\otimes_A k(y)).
\]
Además, en virtud de \sref[0]{0.16.3.9}, se tienen también las desigualdades
\[
\label{IV.13.2.4.5}
\tag{13.2.4.5}
  \dim(B')\leq\dim(B)\leq\dim(A)+\dim(B\otimes_A k(y)).
\]
\end{env}

\oldpage[IV-3]{192}

Comparando estas desigualdades se obtiene:

\begin{lemma}[13.2.5]
\label{IV.13.2.5}
Con la notación de \sref{IV.13.2.4}, las condiciones siguientes son equivalentes:
\begin{enumerate}[label=\emph{\alph*)}]
  \item $\dim(B')=\dim(A)+\dim(B\otimes_A k(y))$.
  \item Se cumplen simultáneamente las relaciones siguientes:
  \begin{enumerate}[label=(\roman*)]
    \item $\dim(A')=\dim(A)$.
    \item $\dim(B'\otimes_{A'}k(y))=\dim(B\otimes_A k(y))$, en otras palabras,
    \[
      \dim_x(X'\cap f^{-1}(y))=\dim_x(f^{-1}(y)).
    \]
    \item $\dim_x(f'^{-1}(y))=\dim(f'^{-1}(\eta'))$, es decir, $f':X'\to Y'$ es equidimensional \sref{IV.13.2.2} en el punto $x$.
    \item Se tiene la igualdad
    \[
      \dim(B')=\dim(A')+\dim(B'\otimes_{A'}k(y))-(\dim_x(f'^{-1}(y))-\dim(f'^{-1}(\eta')))
    \]
  \end{enumerate}
  (relación que se cumple siempre cuando $A=\sh{O}_{Y,y}$ es un anillo universalmente catenario, en virtud de \sref{IV.5.6.5}).
  \item Se cumplen simultáneamente las relaciones siguientes:
  \begin{enumerate}[label=(\roman*)]
    \item $\dim(B')=\dim(B)$.
    \item $\dim(B)=\dim(A)+\dim(B\otimes_A k(y))$.
  \end{enumerate}
\end{enumerate}
\end{lemma}

\begin{env}[13.2.6]
\label{IV.13.2.6}
Recordemos ahora que las componentes irreducibles $X_i$ de $X$ que contienen a $x$ son en número finito y que, por \sref{IV.5.1.2.1} y \sref[0]{0.14.2.1.1},
\[
\label{IV.13.2.6.1}
\tag{13.2.6.1}
  \dim(\sh{O}_{X,x})=\sup_i\dim(\sh{O}_{X_i,x}).
\]
La equivalencia de las condiciones $b)$ y $c)$ de \sref{IV.13.2.5} implica, por tanto, teniendo en cuenta \sref[0]{0.16.3.9} y \sref[0]{0.14.2.1}:
\end{env}

\begin{proposition}[13.2.7]
\label{IV.13.2.7}
Sea $Y$ un preesquema localmente noetheriano, $f:X\to Y$ un morfismo de tipo finito, $x$ un punto de $X$, $f(x)=y$; se tiene
\[
\tag{13.2.7.1}
\label{IV.13.2.7.1}
  \dim(\sh{O}_{X,x})\leq\dim(\sh{O}_{Y,y})+\dim(\sh{O}_{X,x}\otimes_{\sh{O}_{Y,y}}k(y)).
\]
Para que los dos miembros de \sref{IV.13.2.7.1} sean iguales, es necesario y suficiente que exista una parte cerrada irreducible $X'$ de $X$ que contenga a $x$ y satisfaga simultáneamente las condiciones siguientes:
\begin{enumerate}[label=(\roman*)]
  \item Si $Y'=\overline{f(X')}$, se tiene $\dim(\sh{O}_{Y',y})=\dim(\sh{O}_{Y,y})$.
  \item $\dim_x(X'\cap f^{-1}(y))=\dim_x(f^{-1}(y))$ (en otras palabras, $X'$ contiene una de las componentes irreducibles de $f^{-1}(y)$ de dimensión maximal entre las que contienen a $x$). Esto equivale a decir que $\dim(\sh{O}_{X',x}\otimes_{\sh{O}_{Y',y}}k(y))=\dim(\sh{O}_{X,x}\otimes_{\sh{O}_{Y,y}}k(y))$.
  \item $X'$ es equidimensional sobre $Y'$ en el punto $x$; en otras palabras, se tiene
  \[
    \dim_x(X'\cap f^{-1}(y))=\dim(X'\cap f^{-1}(\eta')),
  \]
  donde $\eta'$ es el punto genérico de $Y'$ (y, en consecuencia, todas las componentes irreducibles de $X'\cap f^{-1}(y)$ que contienen a $x$ tienen la misma dimensión \sref{IV.13.2.2}).
\oldpage[IV-3]{193}
  \item Se tiene la igualdad
  \[
    \dim(\sh{O}_{X',x})=\dim(\sh{O}_{Y',y})+\dim(\sh{O}_{X',x}\otimes_{\sh{O}_{Y',y}}k(y))
  \]
  (condición siempre implicada por (iii) cuando $\sh{O}_{Y,y}$ es un anillo universalmente catenario).
\end{enumerate}
Además, $X'$ es entonces una componente irreducible de $X$, e $Y'$ una componente irreducible de $Y$.
\end{proposition}

En el enunciado, los anillos locales $\sh{O}_{X',x}$ y $\sh{O}_{Y',y}$ son relativos a los subpreesquemas cerrados reducidos de $X$, $Y$ cuyos espacios subyacentes son respectivamente $X'$ e $Y'$.

Además, esto demuestra también:

\begin{corollary}[13.2.8]
\label{IV.13.2.8}
Si los dos miembros de \sref{IV.13.2.7.1} son iguales, las partes cerradas irreducibles $X'$ de $X$ que contienen a $x$ y satisfacen las condiciones (i) a (iv) de \sref{IV.13.2.7} son exactamente aquellas para las que $\dim(\sh{O}_{X',x})=\dim(\sh{O}_{X,x})$ (lo que implica necesariamente que $X'$ es una componente irreducible de $X$).
\end{corollary}

La Proposición \sref{IV.13.2.7} implica:

\begin{corollary}[13.2.9]
\label{IV.13.2.9}
Sea $Y$ un preesquema localmente noetheriano, $f:X\to Y$ un morfismo de tipo finito, $x$ un punto de $X$, $f(x)=y$.
Supongamos que $\sh{O}_{Y,y}$ es un anillo universalmente catenario.
Entonces las condiciones siguientes son equivalentes:
\begin{enumerate}[label=\emph{\alph*)}]
  \item Los dos miembros de \sref{IV.13.2.7.1} son iguales.
  \item Existe una componente irreducible $Z$ de $f^{-1}(y)$ que contiene a $x$, de dimensión $\dim_x(f^{-1}(y))$, tal que, para todo $x'$ de un entorno de $x$ en $Z$, se tiene
  \[
  \label{IV.13.2.9.1}
\tag{13.2.9.1}
    \dim(\sh{O}_{X,x'})=\dim(\sh{O}_{Y,y})+\dim(\sh{O}_{X,x'}\otimes_{\sh{O}_{Y,y}}k(y)).
  \]
  \item Existe una componente irreducible $Z$ de $f^{-1}(y)$ que contiene a $x$, de dimensión $\dim_x(f^{-1}(y))$, tal que, para el punto genérico $z$ de $Z$, se tiene
  \[
  \label{IV.13.2.9.2}
\tag{13.2.9.2}
    \dim(\sh{O}_{X,z})=\dim(\sh{O}_{Y,y}).
  \]
\end{enumerate}
\end{corollary}

Demostremos que $a)$ implica $b)$.
Sea $\dim_x(f^{-1}(y))=n$; en virtud de \sref{IV.13.2.7}, existe una componente irreducible $X'$ de $X$ que satisface las condiciones (i) a (iv) de \sref{IV.13.2.7}; sea $Z$ una componente irreducible de dimensión $n$ de $f^{-1}(y)$, que contiene a $x$ y está contenida en $X'$.
Como $f^{-1}(y)$ es localmente noetheriano, existe un entorno abierto $U$ de $x$ en $Z$ tal que $U$ no corta ninguna componente irreducible de $f^{-1}(y)$ distinta de las que contienen a $x$, luego \sref{IV.4.1.1.3} $\dim_{x'}(f^{-1}(y))=n$ para todo $x'\in U$; entonces está claro que las condiciones (i) a (iii) de \sref{IV.13.2.7} se satisfacen cuando se sustituye $x$ por cualquier punto $x'\in U$, y lo mismo ocurre con la condición (iv), pues $\sh{O}_{Y,y}$ es universalmente catenario; de donde la conclusión por \sref{IV.13.2.7}.
La condición $b)$ implica trivialmente $c)$ por \sref{IV.5.1.2}.
Finalmente, si se satisface $c)$ y si $X''$ es una componente irreducible de $X$ que contiene a $Z$ y tal que las condiciones (i) a (iv) de \sref{IV.13.2.7} se satisfacen al sustituir $X'$ por $X''$ y $x$ por $z$, está claro que estas condiciones también se satisfacen cuando se sustituye $X'$ por $X''$ y $x$ por $x$, pues $\sh{O}_{Y,y}$ es universalmente catenario; por tanto, $c)$ implica $a)$.

\begin{proposition}[13.2.10]
\label{IV.13.2.10}
Sea $Y$ un preesquema irreducible localmente noetheriano, $\eta$ su punto genérico, $f:X\to Y$ un morfismo de tipo finito, $y$ un punto de $f(X)$.
Sean $Z_i$ las
\oldpage[IV-3]{194}
componentes irreducibles de $f^{-1}(y)$, $z_i$ el punto genérico de $Z_i$, y consideremos las condiciones siguientes:
\begin{enumerate}[label=\emph{\alph*)}]
  \item Para todo $x\in f^{-1}(y)$, se tiene la relación
  \[
  \tag{13.2.10.1}
\label{IV.13.2.10.1}
    \dim(\sh{O}_{X,x})=\dim(\sh{O}_{Y,y})+\dim(\sh{O}_{X,x}\otimes_{\sh{O}_{Y,y}}k(y)).
  \]
  \item Para todo $i$, se tiene
  \[
  \label{IV.13.2.10.2}
\tag{13.2.10.2}
    \dim(\sh{O}_{X,z_i})=\dim(\sh{O}_{Y,y}).
  \]
  \item Para todo $i$, existe una componente irreducible $X_i$ de $X$ que contiene a $Z_i$ y tal que $\dim(X_i\cap f^{-1}(\eta))=\dim(Z_i)$ (en otras palabras, el \emph{subpreesquema cerrado reducido} $X_i$ de $X$ es equidimensional sobre $Y$ en el punto $z_i$).
\end{enumerate}
Entonces $a)$ implica $b)$ y $b)$ implica $c)$; además, si $\sh{O}_{Y,y}$ es universalmente catenario, las tres condiciones $a)$, $b)$, $c)$ son equivalentes.
\end{proposition}

Como el anillo $\sh{O}_{X,z_i}\otimes_{\sh{O}_{Y,y}}k(y)$ es de dimensión $0$ \sref{IV.5.1.2}, $a)$ implica evidentemente $b)$; $b)$ implica $c)$ en virtud de \sref{IV.13.2.7}, aplicado en el punto $z_i$.
Recíprocamente, supongamos que $\sh{O}_{Y,y}$ es universalmente catenario; como la condición $c)$ implica que $f(X_i)$ es denso en $Y$, de $c)$ se sigue que las condiciones (i) a (iv) de \sref{IV.13.2.7} se satisfacen cuando se sustituye $X'$ por $X_i$ y $x$ por $z_i$, luego $c)$ implica $b)$; finalmente, $b)$ implica $a)$ por \sref{IV.13.2.9}.

\begin{corollary}[13.2.11]
\label{IV.13.2.11}
Con la notación de \sref{IV.13.2.10}, supongamos que $\sh{O}_{Y,y}$ es universalmente catenario.
Para todo $y'\in Y$, sea $E(y')$ el conjunto de las dimensiones de las componentes irreducibles de $f^{-1}(y')$ y pongamos $d(y')=\dim(f^{-1}(y'))=\sup(E(y'))$.
Entonces, si se satisfacen las condiciones equivalentes $a)$, $b)$, $c)$ de \sref{IV.13.2.10}, se tiene $E(y)\subset E(\eta)$, de donde $d(y)\leq d(\eta)$.
\end{corollary}

En efecto, con la notación de \sref{IV.13.2.10}, $X_i\cap f^{-1}(\eta)$ no es vacío y, en consecuencia, es una componente irreducible de $f^{-1}(\eta)$ \sref[0]{0.2.1.8}.

\begin{remarks}[13.2.12]
\label{IV.13.2.12}
\begin{enumerate}[label=(\roman*)]
  \item Recordemos por \sref{IV.6.1.2} que la relación \sref{IV.13.2.10.1} se satisface siempre cuando el morfismo $f$ es \emph{plano} en el punto $x$.
  \item Con la notación de \sref{IV.13.2.10}, supongamos que $\sh{O}_{Y,y}$ es universalmente catenario y, además, que el morfismo $f$ es \emph{propio}; entonces, si se satisfacen las condiciones equivalentes $a)$, $b)$, $c)$ de \sref{IV.13.2.10}, se tiene $d(y)=d(\eta)$, pues de \sref{IV.13.1.5} se sigue que $d(y)\geq d(\eta)$.
  \item El morfismo de \sref{IV.12.2.3},~(ii) es \emph{propio} y \emph{plano}, y todos los anillos locales de $Y$ son universalmente catenarios; además, las dos componentes irreducibles $X_1$, $X_2$ de $X$ son \emph{equidimensionales} sobre $Y$ en todo punto; pero $E(y)$ tiene \emph{dos} elementos para todo $y\neq y_0$, mientras que $E(y_0)$ se reduce a un \emph{solo} elemento, luego $E(y)$ \emph{no es constante} en $Y$.
\end{enumerate}
\end{remarks}

\subsection{Morfismos equidimensionales: caso general}
\label{subsection:IV.13.3}

\begin{proposition}[13.3.1]
\label{IV.13.3.1}
Sea $Y$ un preesquema, $f:X\to Y$ un morfismo localmente de tipo finito, $x$ un punto de $X$, $y=f(x)$.
Denotemos por $Y_\alpha$ las componentes irreducibles de $Y$ que contienen a $y$.
Entonces las condiciones siguientes son equivalentes:
\begin{enumerate}[label=\emph{\alph*)}]
\oldpage[IV-3]{195}
  \item Existen un entero $e$ y un entorno abierto $U$ de $x$ tales que la imagen por $f$ de toda componente irreducible de $U$ es densa en algún $Y_\alpha$, y tales que, para todo $x'\in U$, el espacio $U\cap f^{-1}(f(x'))$ es equidimensional y de dimensión $e$.
  \item[\emph{a$'$)}] Existen un entero $e$ y un entorno abierto $U$ de $x$ tales que la imagen por $f$ de toda componente irreducible de $U$ es densa en algún $Y_\alpha$, y tales que, si se denota por $y_\alpha$ el punto genérico de $Y_\alpha$, toda componente irreducible de los espacios $U\cap f^{-1}(y)$, $U\cap f^{-1}(y_\alpha)$ es de dimensión $e$.
  \item[\emph{a$''$)}] Existen un entero $e$ y un entorno abierto $U$ de $x$ tales que, para cada una de las componentes irreducibles $U_\lambda$ de $U$, $f(U_\lambda)$ es denso en algún $Y_\alpha$ y, para todo $x'\in U_\lambda$, las componentes irreducibles de $U_\lambda\cap f^{-1}(f(x'))$ son todas de dimensión $e$.
  \item Existen un entero $e$, un entorno abierto $U$ de $x$ y un $Y$-morfismo cuasifinito $g:U\to Y\otimes_{\bb{Z}}\bb{Z}[T_1,\ldots,T_e]$ (preesquema que denotaremos también $Y[T_1,\ldots,T_e]$, para abreviar) tales que la imagen por $g$ de toda componente irreducible de $U$ es densa en una componente irreducible de $Y[T_1,\ldots,T_e]$.
\end{enumerate}
\end{proposition}

\begin{proof}
Es inmediato que $a'')$ implica $a)$, pues, para todo $x\in U$, cada componente irreducible de $U\cap f^{-1}(f(x))$ es una componente irreducible de uno de los $U_\lambda\cap f^{-1}(f(x))$, de donde la conclusión por \sref[0]{0.14.1.4}.
La condición $a)$ implica trivialmente $a')$.
Demostremos que $a')$ implica $a'')$; podemos restringirnos al caso en que $f$ es de tipo finito y $X$ e $Y$ son reducidos \sref{IV.1.5.4}.
Sea $U_\lambda$ una componente irreducible de $U$, y supongamos que $f(U_\lambda)$ es denso en $Y_\alpha$; entonces la restricción de $f$ a $U_\lambda$ se factoriza como $U_\lambda\xrightarrow{f_\alpha}Y_\alpha\longrightarrow Y$, donde $f_\alpha$ es de tipo finito y dominante \sref[I]{I.5.2.2}.
Sea $x_\lambda$ el punto genérico de $U_\lambda$; en virtud de \sref[0]{0.2.1.8}, $U_\lambda\cap f_\alpha^{-1}(y_\alpha)$ es la única componente irreducible de $U\cap f^{-1}(y_\alpha)$ que contiene a $x_\lambda$ y, por hipótesis, es de dimensión $e$, igual a las dimensiones de todas las componentes irreducibles de $U_\lambda\cap f^{-1}(y)$, en virtud de la hipótesis y de \sref{IV.13.1.1}.
Pero, por el teorema de Chevalley \sref{IV.13.1.3} y \sref{IV.13.1.1}, el conjunto $V_\lambda$ de los $x'\in U_\lambda$ tales que $\dim_{x'}(f_\lambda^{-1}(f_\lambda(x')))=e$ es abierto y contiene a $x$, y basta tomar la unión de los $V_\lambda$ como conjunto abierto que satisface las condiciones de $a'')$.

Demostremos ahora que $a)$ implica $b)$; podemos restringirnos al caso en que $Y=\Spec(A)$ y $X=\Spec(B)$ son afines y $U=X$.
Demostremos primero el lema siguiente:

\begin{lemma}[13.3.1.1]
\label{IV.13.3.1.1}
Sean $Y=\Spec(A)$, $X=\Spec(B)$ dos esquemas afines, $f:X\to Y$ un morfismo de tipo finito, $y$ un punto de $f(X)$, y pongamos $e=\dim(f^{-1}(y))$.
Entonces existe un $Y$-morfismo $g:X\to Y[T_1,\ldots,T_e]$ tal que (si se pone $X_y=X\times_Y\Spec(k(y))$) el morfismo $g_y:X_y\to\Spec(k(y)[T_1,\ldots,T_e])$ es finito.
Además, para todo morfismo $g$ de este tipo, $g_y$ es necesariamente sobreyectivo, y existe un entorno abierto $U$ de $f^{-1}(y)$ en $X$ tal que $g|U$ es cuasifinito.
\end{lemma}

\begin{proof}
Sea $\mathfrak{p}=\mathfrak{j}_y$; el anillo $B\otimes_A k(\mathfrak{p})$ es una $k(\mathfrak{p})$-álgebra de tipo finito, luego, por el lema de normalización (Bourbaki, \emph{Alg.\ comm.}, cap.~V, \textsection~3, n$^{\text{o}}$~1, teo.~1), existen en $B\otimes_A k(\mathfrak{p})$ una sucesión finita $(t_i)_{1\leq i\leq r}$ de elementos algebraicamente independientes sobre $k(\mathfrak{p})$ y tales que, si se pone $C'=k(\mathfrak{p})[t_1,\ldots,t_r]$, $B\otimes_A k(\mathfrak{p})$ es una $C'$-álgebra \emph{finita}; se tiene, por tanto, $\dim(B\otimes_A k(\mathfrak{p}))=\dim(C')$ \sref[0]{0.16.1.5} y, como $\dim(C')=r$ \sref{IV.5.2.1}, se tiene $r=e$.
Como $B\otimes_A k(\mathfrak{p})=(B/\mathfrak{p}B)\otimes_{A/\mathfrak{p}}k(\mathfrak{p})$, se puede, multiplicando los $t_i$ por un elemento $\neq0$
\oldpage[IV-3]{196}
conveniente de $A/\mathfrak{p}$, suponer que cada $t_i$ es la imagen canónica en $B\otimes_A k(\mathfrak{p})$ de un elemento $s_i\in B$.
Sea entonces $u:A[T_1,\ldots,T_e]\to B$ el homomorfismo tal que $u(T_i)=s_i$ para todo $i$, y sea $g:X\to Y[T_1,\ldots,T_e]$ el morfismo correspondiente.
Está claro que, por la elección de los $s_i$, $g_y$ es un morfismo finito.
Para todo morfismo $g:X\to Y[T_1,\ldots,T_e]$ tal que $g_y$ es finito, de \sref{IV.5.4.2} y \sref{IV.4.1.2.1} se sigue que $g_y$ es necesariamente sobreyectivo.
Por otra parte, por \sref{IV.13.1.4}, el conjunto $U$ de los $x\in X$ que están aislados en su fibra $g^{-1}(g(x))$ es abierto en $X$ y contiene a $f^{-1}(y)$, y por definición la restricción $g|U$ es un morfismo cuasifinito.
\end{proof}

Establecido este lema, para demostrar que $a)$ implica $b)$ queda mostrar que, si $x_\lambda$ es un punto maximal de $U$, entonces $g(x_\lambda)=z_\lambda$ es un punto maximal de $Z=Y[T_1,\ldots,T_e]$.
Por la hipótesis $a)$, podemos (disminuyendo $U$ si es necesario) suponer que $f(x_\lambda)$ es uno de los puntos genéricos $y_\alpha$ de las componentes irreducibles de $Y$ que contienen a $y$. Si $p:Z\to Y$ es el morfismo estructural, entonces $p(z_\lambda)=y_\alpha$, y $g$ induce, por tanto, un $k(y_\alpha)$-morfismo cuasifinito
\[
h:U\cap f^{-1}(y_\alpha)\longrightarrow p^{-1}(y_\alpha).
\]
Ahora bien, $p^{-1}(y_\alpha)=\Spec(k(y_\alpha)[T_1,\ldots,T_e])$ es íntegro y de dimensión $e$. Si $z_\lambda$ no fuera un punto maximal de $Z$, tampoco sería un punto maximal de $p^{-1}(y_\alpha)$ \sref[0]{0.2.1.8}, y su adherencia $Z'_\lambda$ en $p^{-1}(y_\alpha)$ tendría, por tanto, dimensión $<e$ \sref{IV.4.1.2.1}. Pero, como $h$ es cuasifinito, de \sref{IV.4.1.2} y de la hipótesis $a)$ se sigue que $\dim(Z'_\lambda)\geq e$ (la restricción de $h$ a $\overline{\{x_\lambda\}}$ se factoriza como
\[
\overline{\{x_\lambda\}}\longrightarrow Z'_\lambda\longrightarrow p^{-1}(y_\alpha)
\]
por \sref[I]{I.5.2.2}). Esto es una contradicción y demuestra que $a)$ implica $b)$.

Demostremos finalmente que $b)$ implica $a)$.
Obsérvese que el morfismo estructural $p:Z=Y[T_1,\ldots,T_e]\to Y$ es fielmente plano; por tanto, según \sref{IV.2.3.4}, las imágenes por $p$ de los puntos maximales de $Z$ son los puntos maximales de $Y$. Esto demuestra ya que los $f(x_\lambda)$ son puntos genéricos de los $Y_\alpha$.
Además, si $f(x_\lambda)=y_\alpha$, el $k(y_\alpha)$-morfismo $h:U\cap f^{-1}(y_\alpha)\to p^{-1}(y_\alpha)$ deducido de $g$ es dominante y cuasifinito por hipótesis; por tanto, según \sref{IV.4.1.2}, se tiene $\dim(U_\lambda\cap f^{-1}(y_\alpha))=e$ (siendo $U_\lambda$ la componente irreducible de $U$ de punto genérico $x_\lambda$).
Análogamente, para todo $x'\in U_\lambda$, el morfismo $U_\lambda\cap f^{-1}(f(x'))\to p^{-1}(f(x'))$ deducido de $g$ es un $k(f(x'))$-morfismo cuasifinito, luego por \sref{IV.4.1.2} se tiene $\dim(f^{-1}(f(x'))\cap U_\lambda)\leq e$; pero, por \sref{IV.13.1.6}, todas las componentes irreducibles de $U_\lambda\cap f^{-1}(f(x'))$ son de dimensión $\geq e$; se ve, por tanto, que estas componentes son exactamente de dimensión $e$, luego $b)$ implica $a)$.
\end{proof}

\begin{definition}[13.3.2]
\label{IV.13.3.2}
Sea $Y$ un preesquema, $f:X\to Y$ un morfismo localmente de tipo finito, $x$ un punto de $X$.
Se dice que $f$ es \emph{equidimensional en el punto} $x$ (o que $X$ es \emph{equidimensional sobre} $Y$ \emph{en el punto} $x$) si se satisfacen las condiciones equivalentes de \sref{IV.13.3.1}.
Se dice que $f$ es \emph{equidimensional} (o que $X$ es \emph{equidimensional sobre} $Y$) si $f$ es equidimensional en todo punto $x\in X$.
\end{definition}

Está claro que, para que $f$ sea equidimensional en un punto $x$, es necesario y suficiente que $f_{\mathrm{red}}$ lo sea.
Además, las condiciones de \sref{IV.13.3.1} muestran que el conjunto de los puntos en los que $f$ es equidimensional es \emph{abierto} en $X$.

Observemos que, cuando $X$ e $Y$ son \emph{irreducibles}, decir que $f$ es equidimen-
\oldpage[IV-3]{197}
sional en el punto $x$ significa, en virtud de \sref{IV.13.3.1},~$a'')$, que $f$ es dominante y que $\dim_x(f^{-1}(f(x)))=\dim(f^{-1}(\eta))$ (donde $\eta$ es el punto genérico de $Y$); la definición \sref{IV.13.3.2} coincide así, en este caso, con la definición \sref{IV.13.2.2}.

\begin{proposition}[13.3.3]
\label{IV.13.3.3}
Sea $Y$ un preesquema localmente noetheriano, $f:X\to Y$ un morfismo localmente de tipo finito, $x$ un punto de $X$ y $X_j$ las componentes irreducibles de $X$ (en número finito) que contienen a $x$.
Para que $f$ sea equidimensional en el punto $x$, es necesario y suficiente que, para todo $j$, $Y_j=\overline{f(X_j)}$ sea una componente irreducible de $Y$ y que, denotando por $X_j$ e $Y_j$ los subpreesquemas cerrados reducidos de $X$ e $Y$ de espacios subyacentes $X_j$ e $Y_j$, por $f_j:X_j\to Y_j$ el morfismo deducido de $f$ \sref[I]{I.5.2.2}, y por $y_j$ el punto genérico de $Y_j$, $f_j$ sea equidimensional en el punto $x$ y todos los números $\dim(f_j^{-1}(y_j))$ sean iguales.
\end{proposition}

\begin{proof}
Esto resulta inmediatamente de \sref{IV.13.3.1},~$a')$).
\end{proof}

\begin{corollary}[13.3.4]
\label{IV.13.3.4}
Con la notación de \sref{IV.13.3.3}, pongamos $y=f(x)$.
Si $\sh{O}_{Y,y}$ es un anillo universalmente catenario y $f$ es equidimensional en el punto $x$, se tiene
\[
\tag{13.3.4.1}
\label{IV.13.3.4.1}
  \dim(\sh{O}_{X_j,x})=\dim(\sh{O}_{Y_j,y})+e-\deg.\mathrm{tr}_{k(y)}k(x)
\]
donde $e$ es el valor común de los números $\dim(f_j^{-1}(y_j))$.
\end{corollary}

\begin{proof}
Como cada uno de los $\sh{O}_{Y_j,y}$, cociente de $\sh{O}_{Y,y}$, es un anillo universalmente catenario \sref{IV.5.6.1}, la igualdad es consecuencia de \sref{IV.5.6.5}.
\end{proof}

\begin{corollary}[13.3.5]
\label{IV.13.3.5}
Con la notación de \sref{IV.13.3.4}, supongamos que $f$ es equidimensional en el punto $x$ y que $\sh{O}_{Y,y}$ es un anillo universalmente catenario.
Si el anillo $\sh{O}_{Y,y}$ es equidimensional, también lo es $\sh{O}_{X,x}$; el recíproco es cierto si la imagen por $f$ de la unión de los $X_j$ es densa en un entorno de $y$.
\end{corollary}

\begin{proof}
En efecto, decir que $\sh{O}_{Y,y}$ es equidimensional significa que los números $\dim(\sh{O}_{Y'_i,y})$ son iguales para todas las componentes irreducibles $Y'_i$ de $Y$ que contienen a $y$, como resulta de \sref{IV.5.1.1.5}, aplicado al esquema local $\Spec(\sh{O}_{Y,y})$; como se tiene la relación \sref{IV.13.3.4.1}, del mismo razonamiento se sigue que $\sh{O}_{X,x}$ es entonces equidimensional.
Recíprocamente, si $\sh{O}_{X,x}$ es equidimensional, todos los números $\dim(\sh{O}_{Y_j,y})$ son iguales por \sref{IV.13.3.4.1}; se deduce que $\sh{O}_{Y,y}$ es equidimensional si los $Y_j$ son todas las componentes irreducibles de $Y$ que contienen a $y$, lo cual resulta de la hipótesis suplementaria.
\end{proof}

\begin{proposition}[13.3.6]
\label{IV.13.3.6}
Sea $Y$ un preesquema localmente noetheriano, $f:X\to Y$ un morfismo localmente de tipo finito, $x$ un punto de $X$, $y=f(x)$.
Supongamos que el anillo $\sh{O}_y$ es equidimensional. Entonces, si $\sh{O}_x$ es equidimensional y se tiene la igualdad
\[
\tag{13.3.6.1}
\label{IV.13.3.6.1}
  \dim(\sh{O}_x)=\dim(\sh{O}_y)+\dim(\sh{O}_x\otimes_{\sh{O}_y}k(y))
\]
(cf.\ \sref{IV.13.2.7.1}), $f$ es equidimensional en el punto $x$, y el recíproco es cierto si $\sh{O}_y$ es un anillo universalmente catenario.
\end{proposition}

\begin{proof}
Conservando la notación de \sref{IV.13.3.3}, de \sref{IV.13.2.8} y de la hipótesis de que $\sh{O}_x$ es equidimensional se sigue que cada $Y_j$ es una componente irreducible de $Y$ y que cada $f_j$ es equidimensional en el punto $x$; además, por \sref{IV.13.2.8}, se tiene (teniendo en cuenta \sref{IV.5.6.5})
\[
  \dim(\sh{O}_{X_j,x})=\dim(\sh{O}_{Y_j,y})+\dim_x(X_j\cap f^{-1}(y))-\deg.\mathrm{tr}_{k(y)}k(x).
\]

\oldpage[IV-3]{198}

Ahora bien, como $\sh{O}_y$ se supone equidimensional, esta igualdad puede escribirse
\[
\tag{13.3.6.2}
\label{IV.13.3.6.2}
  \dim_x(X_j\cap f^{-1}(y))=\dim(\sh{O}_{X,x})-\dim(\sh{O}_{Y,y})+\deg.\mathrm{tr}_{k(y)}k(x).
\]
El primer miembro de \sref{IV.13.3.6.2} es, por tanto, independiente de $j$; pero, como $f_j$ es equidimensional en el punto $x$, se tiene $\dim_x(X_j\cap f^{-1}(y))=\dim(f_j^{-1}(y_j))$, luego el criterio de \sref{IV.13.3.3} muestra que $f$ es equidimensional en el punto $x$.

Recíprocamente, supongamos que $\sh{O}_{Y,y}$ es un anillo universalmente catenario y que $f$ es equidimensional en el punto $x$; entonces, por \sref{IV.13.3.5}, $\sh{O}_{X,x}$ es equidimensional, y de \sref{IV.13.3.4} se sigue que se tiene la relación
\[
  \dim(\sh{O}_{X,x})=\dim(\sh{O}_{Y,y})+e-\deg.\mathrm{tr}_{k(y)}k(x)
\]
donde $e$ es el valor común de los números $\dim(f_j^{-1}(y_j))=\dim_x(X_j\cap f^{-1}(y))$; pero, por definición, $e$ es también igual a $\dim_x(f^{-1}(y))$, de donde la relación \sref{IV.13.3.6.1}, teniendo en cuenta \sref{IV.5.6.5.2}.
\end{proof}

\begin{proposition}[13.3.7]
\label{IV.13.3.7}
Sea $Y$ un preesquema, $f:X\to Y$ un morfismo localmente de tipo finito y equidimensional, y $Z$ una parte cerrada de $X$.
Entonces la función
\[
\label{IV.13.3.7.1}
\tag{13.3.7.1}
  x\mapsto\operatorname{codim}_x(Z\cap f^{-1}(f(x)),f^{-1}(f(x)))
\]
es semicontinua inferiormente en $X$.
\end{proposition}

\begin{proof}
Como todas las componentes irreducibles de $f^{-1}(f(x))$ que contienen a $x$ tienen la misma dimensión por hipótesis \sref{IV.13.3.1}, se tiene, en virtud de \sref{IV.5.2.1},
\[
  \operatorname{codim}_x(Z\cap f^{-1}(f(x)),f^{-1}(f(x)))=\dim_x(f^{-1}(f(x)))-\dim_x(Z\cap f^{-1}(f(x))).
\]
Pero, por hipótesis, el primer término del segundo miembro es una función continua de $x$ \sref{IV.13.3.1}, y el segundo es una función semicontinua superiormente de $x$ por \sref{IV.13.1.3}\footnote{La fuente impresa cita aquí (13.3.3); la semicontinuidad superior invocada es el teorema de Chevalley (13.1.3).}; de donde la conclusión.
\end{proof}

\begin{proposition}[13.3.8]
\label{IV.13.3.8}
Sea $Y$ un preesquema, $f:X\to Y$ un morfismo localmente de tipo finito, $x$ un punto de $X$.
Sea $Y'$ un preesquema, $g:Y'\to Y$ un morfismo plano, $X'=X\times_Y Y'$, $f'=f_{(Y')}:X'\to Y'$; si $f$ es equidimensional en el punto $x$, entonces $f'$ es equidimensional en todo punto $x'\in X'$ situado sobre $x$.
\end{proposition}

\begin{proof}
Como la cuestión es local sobre $X$ e $Y$, podemos restringirnos al caso en que toda componente irreducible de $X$ (resp.\ de $Y$) contiene a $x$ (resp.\ a $y$); como la imagen por $f$ de toda componente irreducible de $X$ es densa en una componente irreducible de $Y$ \sref{IV.13.3.1}, se sabe por \sref{IV.2.3.5} que la imagen por $f'$ de toda componente irreducible de $X'$ es densa en una componente irreducible de $Y'$.
Por otra parte, por transitividad de las fibras \sref[I]{I.3.6.4}, de \sref{IV.4.2.8} y \sref{IV.13.3.1} se sigue que, para todo $z'\in X'$, el conjunto de las dimensiones de las componentes irreducibles de ${f'}^{-1}(f'(z'))$ es el mismo que el conjunto de las dimensiones de las componentes irreducibles de $f^{-1}(f(z))$, donde $z$ es la proyección de $z'$ sobre $X$; de donde la conclusión en virtud de \sref{IV.13.3.1},~$a)$).
\end{proof}

\begin{remarks}[13.3.9]
\label{IV.13.3.9}
La propiedad de equidimensionalidad de un morfismo $f:X\to Y$ \emph{no} es estable por un cambio de base arbitrario $g:Y'\to Y$, ni siquiera cuando $g$ es la inyección canónica de una componente irreducible $Y'$ de $Y$ en $Y$ (siendo $Y'$
\oldpage[IV-3]{199}
considerado como subpreesquema cerrado reducido de $Y$).
Por ejemplo, sea $k$ un cuerpo, sea $A_0=k[S,T]$ el anillo de polinomios en dos indeterminadas, y pongamos $\mathfrak{a}=\mathfrak{p}_1\mathfrak{p}_2$, donde $\mathfrak{p}_1$ y $\mathfrak{p}_2$ son los ideales primos $A_0S$ y $A_0T$ de $A_0$; sea $A=A_0/\mathfrak{a}$ e $Y=\Spec(A)$, que tiene dos componentes irreducibles $Y_1=\Spec(A/\mathfrak{p}_1)$ e $Y_2=\Spec(A/\mathfrak{p}_2)$; tomemos $X=Y_1$, con $f:X\to Y$ la inyección canónica, que es evidentemente un morfismo equidimensional.
Por otra parte, tomemos $Y'=Y_2$, con $g:Y'\to Y$ la inyección canónica; entonces $X'=X\times_Y Y'=\Spec\bigl((A/\mathfrak{p}_1)\otimes_A(A/\mathfrak{p}_2)\bigr)=\Spec\bigl(A/(\mathfrak{p}_1+\mathfrak{p}_2)\bigr)$, y $\mathfrak{p}_1+\mathfrak{p}_2$ es un ideal maximal de $A$, de modo que $X'$ se reduce a un punto; el morfismo $f':X'\to Y'$ \emph{no es dominante}, pues la imagen por $f'$ del único punto de $X'$ es un punto cerrado de $Y'$; en consecuencia, $f'$ no es equidimensional.

También se puede dar un contraejemplo en el que $X$ e $Y$ son \emph{íntegros}, $f$ es \emph{finito} y \emph{birracional} (y \emph{a fortiori} equidimensional por \sref{IV.13.3.1},~$b)$), y $g:Y'\to Y$ es \emph{finito} y \emph{dominante}.
Sean $A$ y $\overline{A}$ los anillos locales definidos en \sref{IV.11.7.5}, y tomemos $Y=\Spec(A)$ y $X=\Spec(\overline{A})$; por otra parte, con la notación de \sref{IV.11.7.5}, tomemos $Y'=\Spec(\overline{B})$; entonces $X'=\Spec(\overline{A}\otimes_A\overline{B})=\Spec(\overline{B}\otimes_B\overline{B})$; pero se verifica inmediatamente que $\overline{B}\otimes_B\overline{B}$ es el producto directo de los anillos $B/\mathfrak{p}'$, $B/\mathfrak{p}''$ y de dos anillos isomorfos a $B/\mathfrak{n}$, cuyos espectros se reducen, por tanto, a un punto; como las proyecciones de estos puntos son puntos cerrados de $Y'$, se ve aquí de nuevo que $f'$ no lleva una componente irreducible de $X'$ sobre una parte densa de una componente irreducible de $Y'$, y por tanto $f'$ no es equidimensional.

En este ejemplo, el anillo $A$ no es geométricamente unirramificado; veremos en el parágrafo siguiente {\hyperlink{ega4.c04.cor.14-4-6}{(14.4.6)}} que tales fenómenos no pueden producirse cuando los puntos de $Y$ son \emph{geométricamente unirramificados}.
La falta de estabilidad de la noción de morfismo equidimensional restringe mucho su interés, en favor de la noción de morfismo universalmente abierto, que se estudiará en detalle en el parágrafo siguiente.
\end{remarks}

\section{Morfismos universalmente abiertos}
\label{section:IV.14}

Las secciones~\textsection\textsection~14 y~15 están dedicadas al estudio de la noción de morfismo \emph{universalmente abierto} \sref{IV.2.4.2}.
Ya hemos visto \sref{IV.2.4.6} que un morfismo \emph{plano} localmente de presentación finita es universalmente abierto, siendo falsa la recíproca.
En~\textsection~14 examinamos primero las propiedades de las dimensiones de las fibras de un morfismo universalmente abierto $f:X\to Y$; cuando $X$ e $Y$ son localmente noetherianos, $f$ se comporta a este respecto \sref{IV.14.2.1} como un morfismo plano (cf.~\sref{IV.6.1.2}), y es en particular \emph{equidimensional} cuando es localmente de tipo finito y dominante y $X$ es irreducible \sref{IV.14.2.2}.
Recíprocamente, un morfismo equidimensional localmente de tipo finito $f:X\to Y$ es universalmente abierto cuando se supone además que $Y$ es \emph{geométricamente unirramificado}, y en particular cuando $Y$ es \emph{normal} (criterio de Chevalley, \sref{IV.14.4.4}).
Mostramos también que los morfismos universalmente abiertos localmente de tipo finito $f:X\to Y$ (cuando $X$ e $Y$ son
\oldpage[IV-3]{200}
localmente noetherianos e $Y$ es irreducible) admiten ``suficientes'' \emph{cuasisecciones}, es decir, en un entorno de un punto cerrado $x$ de una fibra $f^{-1}(y)$, existe un subpreesquema cerrado $X'$ de $X$ que contiene a $x$ y tal que la restricción $X'\to Y$ de $f$ es un morfismo \emph{cuasifinito} (por tanto con fibras \emph{discretas}) y dominante \sref{IV.14.5.3}.

En~\textsection~15 estudiamos diversas propiedades de las fibras de los morfismos universalmente abiertos $f:X\to Y$, localmente de tipo finito, en particular cuando $X$ e $Y$ son localmente noetherianos.
Obtenemos así, en particular, un criterio para que un punto $x\in X$ pertenezca a una \emph{única} componente irreducible de $X$, en términos de propiedades de la \emph{fibra} $f^{-1}(f(x))$ de $x$: basta que $Y$ sea geométricamente unirramificado (por ejemplo normal) en el punto $f(x)$ y que $f^{-1}(f(x))$ sea geométricamente puntualmente íntegra en el punto $x$ \sref{IV.15.3.3}; si además $Y$ es localmente íntegro en el punto $f(x)$, entonces $f$ es plano en el punto $x$ y $X$ es localmente íntegro en el punto $x$.
Estudiamos también la variación del \emph{número de componentes conexas geométricas} de una fibra $f^{-1}(y)$; por ejemplo, si $f$ es universalmente abierto y \emph{propio}, y las fibras de $f$ son geométricamente reducidas, este número es \emph{localmente constante} \sref{IV.15.5.7}.
Finalmente, cuando $f$ admite una \emph{sección} $g$ (lo que ocurrirá cuando $X$ sea un \emph{esquema en grupos sobre $Y$}) y, para todo $y\in Y$, se denota por $X_y^0$ la componente conexa de la fibra $X_y=f^{-1}(y)$ en el punto $g(y)$ (la ``componente neutra'' en el caso de grupos), se estudia la unión $X^0$ de los $X_y^0$ para $y\in Y$, y se demuestra \sref{IV.15.6.4} que si $f$ es universalmente abierto y las fibras $X_y$ son geométricamente reducidas, entonces $X^0$ es un conjunto \emph{abierto} de $X$.

\subsection{Morfismos abiertos}
\label{subsection:IV.14.1}

\begin{env}[14.1.1]
\label{IV.14.1.1}
Recordemos \sref{IV.1.10.2} que se dice que una aplicación continua $\psi:X\to Y$ es \emph{abierta en un punto} $x\in X$ si la imagen por $\psi$ de \emph{todo} entorno de $x$ en $X$ es un entorno de $\psi(x)$ en $Y$.

Observemos que esto \emph{no} implica que exista un sistema fundamental de entornos de $x$ cuyas imágenes sean \emph{abiertas} en $Y$.
\end{env}

\begin{proposition}[14.1.2]
\label{IV.14.1.2}
Sean $X$, $Y$ dos espacios topológicos, $\psi:X\to Y$ una aplicación continua, $x$ un punto de $X$, $y=\psi(x)$.
\begin{enumerate}[label=(\roman*)]
  \item Si $\psi$ es abierta en $x$, entonces para todo subconjunto $Y'$ de $Y$ que contenga a $\psi(x)$ la restricción $\psi^{-1}(Y')\to Y'$ de $\psi$ es abierta en el punto $x$.
  \item Supongamos que $Y$ es la unión de una familia localmente finita de subconjuntos cerrados $(Y_i)$ y que para todo $i$ tal que $\psi(x)\in Y_i$, la restricción $\psi^{-1}(Y_i)\to Y_i$ de $\psi$ es abierta en el punto $x$; entonces $\psi$ es abierta en el punto $x$.
  \item Sea $\gamma:X'\to X$ una aplicación continua, $x'$ un punto de $X'$; si la aplicación compuesta $\psi\circ\gamma:X'\to Y$ es abierta en el punto $x'$, $\psi$ es abierta en el punto $\gamma(x')$.
\end{enumerate}
\end{proposition}

\begin{proof}
Si $U$ es un entorno de $x$, se tiene $\psi(U\cap\psi^{-1}(Y'))=\psi(U)\cap Y'$; de aquí se sigue inmediatamente (i).
Para demostrar (ii), observemos que existe un entorno $W$ de $\psi(x)$ en $Y$ que solo corta un número finito de los subconjuntos cerrados $Y_i$ que \emph{contienen} a $\psi(x)$; por tanto $W$ es la unión de los $W\cap Y_i$ para estos índices.
Ahora bien, si $U$ es un entorno de $x$ tal que $\psi(U)\cap Y_i$ es un entorno de $\psi(x)$ en $Y_i$, existe un entorno $V\subset W$ de $\psi(x)$ en $Y$ tal que, para todo $i$ tal que $\psi(x)\in Y_i$,
\oldpage[IV-3]{201}
se tiene $V\cap Y_i\subset\psi(U)\cap Y_i$ y, como la unión de los $V\cap Y_i$ para estos índices es $V$, se tiene $V\subset\psi(U)$; por tanto $\psi(U)$ es un entorno de $\psi(x)$ en $Y$.
La afirmación (iii) es trivial.
\end{proof}

\begin{remarks}[14.1.3]
\label{IV.14.1.3}
\begin{enumerate}[label=(\roman*)]
  \item El conjunto $Z$ de los puntos $x\in X$ en los que un morfismo es abierto \emph{no es necesariamente abierto}.
  Por ejemplo, sea $K$ un cuerpo, $A$ el anillo de polinomios $K[S,T]$, $V$ el plano afín $\Spec(A)$, $X$ el subpreesquema cerrado de $V$ que es la ``unión'' de la recta $X_1$ definida por $T=0$ y de la recta $X_2$ definida por $S=0$, es decir, $\Spec(A/\mathfrak{a})$, donde $\mathfrak{a}=AST$; sea $Y=X_2=\Spec(A/AS)$ y sea $f$ la proyección correspondiente a la inyección canónica $K[T]\to A/\mathfrak{a}$; entonces se tiene $Z=X_2$, que no es abierto en $X$.

  \item Sean $X$, $Y$ dos preesquemas noetherianos irreducibles con puntos genéricos $\xi$, $\eta$ respectivamente, $f:X\to Y$ un morfismo \emph{dominante localmente de tipo finito}; entonces $f$ \emph{es abierto en el punto $\xi$}.
  En efecto, podemos evidentemente reducirnos al caso en que $X$ e $Y$ son reducidos (por tanto íntegros) \sref{IV.1.5.4} y afines; en virtud del teorema de planitud genérica \sref{IV.6.9.1}, existe un abierto no vacío $V$ de $Y$ tal que la restricción $f^{-1}(V)\to V$ de $f$ es un morfismo \emph{plano}, y de \sref{IV.2.4.6} concluimos que esta restricción es un morfismo abierto.
  Sin embargo, existen morfismos dominantes de tipo finito $f:X\to Y$ (donde $X$ e $Y$ son \emph{irreducibles}) que no son abiertos (véase por ejemplo \sref[II]{II.8.1}); el conjunto de puntos en los que tal morfismo es abierto no es necesariamente \emph{cerrado} en $X$.

  \item No sabemos si, cuando $X$ e $Y$ son localmente noetherianos y $f:X\to Y$ es un morfismo localmente de tipo finito, el conjunto de puntos de $X$ en los que $f$ es abierto es \emph{localmente construible}.
\end{enumerate}
\end{remarks}

\begin{proposition}[14.1.4]
\label{IV.14.1.4}
Sean $X$, $Y$ dos espacios topológicos, $\psi:X\to Y$ una aplicación continua.
Para todo $y\in Y$, el conjunto de los $x\in\psi^{-1}(y)$ en los que $\psi$ es abierta es un subconjunto cerrado de $\psi^{-1}(y)$.
\end{proposition}

\begin{proof}
En efecto, supongamos que $\psi$ no es abierta en un punto $x\in\psi^{-1}(y)$; existe entonces un entorno $V$ de $x$ en $X$ tal que $\psi(V)$ no es un entorno de $y$; se sigue que para todo $x'\in V\cap\psi^{-1}(y)$, $\psi$ no es abierta en el punto $x'$.
\end{proof}

\begin{remark}[14.1.5]
\label{IV.14.1.5}
Incluso si $X$ e $Y$ son localmente noetherianos y $f$ es un morfismo \emph{finito}, puede ocurrir que $f$ sea abierto en todo punto de una fibra $f^{-1}(y)$, sin que exista un \emph{entorno} de $f^{-1}(y)$ en el que $f$ sea abierto en todos sus puntos.
Tomemos por ejemplo un cuerpo $K$, el anillo de polinomios $A$ igual a $K[\mathrm{T}_1,\mathrm{T}_2,\mathrm{T}_3]$, el espacio afín $V=\Spec(A)$ de dimensión $3$ sobre $K$, y $X=\Spec(A/\mathfrak{a})$, donde $\mathfrak{a}=\mathfrak{b}\mathfrak{c}$, con $\mathfrak{b}=(\mathrm{T}_3)$ y $\mathfrak{c}=(\mathrm{T}_1)+(\mathrm{T}_2-\mathrm{T}_3)$ en $A$, de modo que $X$ es la unión del plano $X_1=\Spec(A/\mathfrak{b})$ (``plano de ecuación $\mathrm{T}_3=0$'') y de la recta $X_2=\Spec(A/\mathfrak{c})$ (``recta de ecuaciones $\mathrm{T}_1=0$, $\mathrm{T}_2=\mathrm{T}_3$~''), que son sus componentes irreducibles.
Tomemos $Y=X_1$ y sea $f:X\to Y$ la proyección correspondiente a la inyección canónica $K[\mathrm{T}_1,\mathrm{T}_2]\to A/\mathfrak{a}$; si $y$ es el punto común a $X_1$ y $X_2$, $f^{-1}(y)$ se reduce a $y$ y $f$ es abierto en este punto, pero no es abierto en ningún punto de $X_2$ situado en un entorno de $y$ y distinto de $y$.
\end{remark}

\begin{env}[14.1.6]
\label{IV.14.1.6}
En lo que sigue, desempeñará un papel esencial el criterio \sref{IV.1.10.3} que, mediante el ``levantamiento de las generalizaciones'', caracteriza los morfismos \emph{localmente de presentación finita} $f:X\to Y$ que son abiertos en un punto $x$: para toda generalización $y'$ de $y=f(x)$, existe $x'\in X$, una generalización de $x$, tal que $y'=f(x')$.
\end{env}

\oldpage[IV-3]{202}
\subsection{Morfismos abiertos y fórmula de la dimensión}
\label{subsection:IV.14.2}

\begin{theorem}[14.2.1]
\label{IV.14.2.1}
Sean $X$, $Y$ dos preesquemas localmente noetherianos, $f:X\to Y$ un morfismo, $x$ un punto de $X$, $y=f(x)$.
Supongamos que $f$ es abierto en los puntos genéricos de las componentes irreducibles de $f^{-1}(y)$ que contienen a $x$.
Entonces se tiene la relación
\begin{equation}
\tag{14.2.1.1}
\label{IV.14.2.1.1}
  \dim(\sh{O}_x)=\dim(\sh{O}_y)+\dim(\sh{O}_x\otimes_{\sh{O}_y}k(y)).
\end{equation}
\end{theorem}

\begin{proof}
Sea $y'$ una generalización cualquiera de $y$ distinta de $y$, y consideremos el subpreesquema cerrado reducido $Y'$ de $Y$ cuyo espacio subyacente es $\overline{\{y'\}}$; ninguna componente irreducible $X'$ de $f^{-1}(Y')$ que contenga a $x$ puede estar contenida en $f^{-1}(y)$: en efecto, si $x'$ es el punto genérico de $X'$, se tendría $f(x')=y\neq y'$, y como $x'$ es su única generalización en $f^{-1}(Y')$, esto contradiría la hipótesis de que $f$ es abierto en $x'$, en virtud del hecho de que $a)$ implique $c)$ en \sref{IV.1.10.3}.
Podemos aplicar, por tanto, \sref{IV.6.1.2} al homomorfismo local de anillos locales noetherianos $\sh{O}_y\to\sh{O}_x$, de donde resulta la conclusión.
\end{proof}

\begin{corollary}[14.2.2]
\label{IV.14.2.2}
Sea $Y$ un preesquema localmente noetheriano, $f:X\to Y$ un morfismo localmente de tipo finito, $x$ un punto de $X$, $y=f(x)$.
Supongamos que $f$ es abierto en los puntos genéricos de las componentes irreducibles de $f^{-1}(y)$ que contienen a $x$.
Entonces $f$ es equidimensional en el punto $x$ en cada uno de los dos casos siguientes:
\begin{enumerate}[label=(\roman*)]
  \item $X$ es irreducible y $f$ es dominante.
  \item Los anillos $\sh{O}_x$ y $\sh{O}_y$ son equidimensionales.
\end{enumerate}
\end{corollary}

\begin{proof}
Para (i), esto se sigue de \sref{IV.14.2.1} y de \sref{IV.13.2.3}.
Para (ii), esto se sigue de \sref{IV.14.2.1} y de \sref{IV.13.3.6}.
\end{proof}

\begin{proposition}[14.2.3]
\label{IV.14.2.3}
Sea $Y$ un preesquema irreducible localmente noetheriano, $\eta$ su punto genérico, $f:X\to Y$ un morfismo localmente de tipo finito, $y$ un punto de $f(X)$, $Z$ una componente irreducible de $f^{-1}(y)$ tal que $f$ es abierto en el punto genérico de $Z$.
Entonces $Z$ está contenida en una componente irreducible $X'$ de $X$ dominante sobre $Y$ y tal que $\dim(Z)=\dim(X'\cap f^{-1}(y))=\dim(X'\cap f^{-1}(\eta))$.
\end{proposition}

\begin{proof}
Esto se sigue, en efecto, de \sref{IV.14.2.1} aplicado en el punto genérico de $Z$, y de \sref{IV.13.2.7}.
\end{proof}

\begin{corollary}[14.2.4]
\label{IV.14.2.4}
Con las notaciones de \sref{IV.14.2.3}, supongamos que $f$ es abierto en todos los puntos de $f^{-1}(y)$.
Para todo $z\in Y$, sea $\mathrm{E}(z)$ el conjunto de las dimensiones de las componentes irreducibles de $f^{-1}(z)$ y $d(z)=\sup\mathrm{E}(z)=\dim(f^{-1}(z))$.
Entonces $\mathrm{E}(y)\subset\mathrm{E}(\eta)$, de donde $d(y)\leq d(\eta)$.
\end{corollary}

\begin{proof}
Esto se sigue inmediatamente de \sref{IV.14.2.3}.
\end{proof}

\begin{corollary}[14.2.5]
\label{IV.14.2.5}
Sea $Y$ un preesquema localmente noetheriano, $f:X\to Y$ un morfismo propio, $y\in Y$ un punto tal que $f$ es abierto en todos los puntos de $f^{-1}(y)$.
Entonces la función $z\mapsto\dim(f^{-1}(z))$ es constante en un entorno de $y$.
\end{corollary}

\begin{proof}
Sean $Y_i$ ($1\leq i\leq n$) los subpreesquemas cerrados reducidos de $Y$ cuyos espacios subyacentes son las componentes irreducibles de $Y$ que contienen a $y$; si $f_i:f^{-1}(Y_i)\to Y_i$ es la restricción de $f$, sabemos que cada uno de los $f_i$ es propio \sref[II]{II.5.4.5} y abierto en todos los puntos de $f^{-1}(y)$ \sref{IV.14.1.2}.
Como la unión de los $Y_i$ es un entorno de $y$ en $Y$,
\oldpage[IV-3]{203}
vemos que podemos reducirnos a demostrar el corolario cuando $Y$ es irreducible; sea $\eta$ su punto genérico.
Se sigue de \sref{IV.14.2.4} que $d(y)\leq d(\eta)$; por otra parte, como $f$ es propio, deducimos de \sref{IV.13.1.5} que $z\mapsto\dim(f^{-1}(z))$ es semicontinua superiormente; en particular $d(y)\geq d(\eta)$, luego $d(y)=d(\eta)$.
Además, existe un entorno $V$ de $y$ tal que $d(z)\leq d(y)$ para todo $z\in V$; pero como todo $z$ es una especialización de $\eta$, se tiene por otra parte $d(z)\geq d(\eta)$, de donde finalmente $d(z)=d(y)$ para $z\in V$.
\end{proof}

\begin{corollary}[14.2.6]
\label{IV.14.2.6}
Sea $Y$ un preesquema localmente noetheriano, $f:X\to Y$ un morfismo abierto de tipo finito, $y$ un punto de $Y$.
Sean $y_i$ ($1\leq i\leq n$) los puntos genéricos de las componentes irreducibles de $Y$ que contienen a $y$.
Entonces, con las notaciones de \sref{IV.14.2.4}:
\begin{enumerate}[label=(\roman*)]
  \item Si para $1\leq i\leq n$ se tiene $\mathrm{E}(y)=\mathrm{E}(y_i)$ (resp.\ $d(y)=d(y_i)$), la función $\mathrm{E}$ (resp.\ $d$) es constante en un entorno de $y$.
  \item Existe un entorno abierto $U$ de $f^{-1}(y)$ tal que la función $z\mapsto\dim(U\cap f^{-1}(z))$ es constante en un entorno de $y$.
\end{enumerate}
\end{corollary}

\begin{proof}
\begin{enumerate}[label=(\roman*)]
  \item El mismo razonamiento que en \sref{IV.14.2.5} muestra que podemos reducirnos a considerar el caso en que $Y$ es irreducible, con punto genérico $\eta$.
  Si $y'$ es una generalización cualquiera de $y$, se tiene entonces (aplicando \sref{IV.14.2.4} a la restricción $f^{-1}(Y')\to Y'$ de $f$, donde $Y'=\overline{\{y'\}}$, y usando \sref{IV.14.1.2}) $\mathrm{E}(y)\subset\mathrm{E}(y')\subset\mathrm{E}(\eta)$ y $d(y)\leq d(y')\leq d(\eta)$; esto muestra que la relación $\mathrm{E}(y)=\mathrm{E}(\eta)$ (resp.\ $d(y)=d(\eta)$) implica $\mathrm{E}(y)=\mathrm{E}(y')$ (resp.\ $d(y)=d(y')$) para toda generalización $y'$ de $y$.
  Ahora bien, en virtud de \sref{IV.9.5.5}, las funciones $\mathrm{E}$ y $d$ son localmente construibles, luego se sigue de $(\mathbf{0}_{\mathrm{III}},~9.2.5)$ aplicado al conjunto de los $z$ tales que $\mathrm{E}(z)=\mathrm{E}(y)$ (resp.\ $d(z)=d(y)$) que este conjunto es un entorno de $y$.

  \item Para cada uno de los $y_i$, $f^{-1}(y_i)$ es un espacio noetheriano puesto que $f$ es de tipo finito; sean $x_{ij}$ ($1\leq j\leq m_i$) los puntos genéricos de sus componentes irreducibles, y pongamos $X_{ij}'=\overline{\{x_{ij}\}}$; mostremos que el complementario $U$ en $X$ de la unión de los $X_{ij}'$ que no cortan a $f^{-1}(y)$ (que evidentemente es un entorno abierto de $f^{-1}(y)$) satisface lo requerido.
  En efecto, la restricción $f|U:U\to Y$ de $f$ es un morfismo abierto de tipo finito \sref[I]{I.6.3.5}; para todo par $(i,j)$ tal que $x_{ij}\in U$, $x_{ij}$ es un punto maximal de $U\cap f^{-1}(y_i)$, y se sigue de \sref{IV.13.1.1}, aplicado a la restricción $X_{ij}'\to Y_i=\overline{\{y_i\}}$ de $f$ (teniendo en cuenta \sref[I]{I.5.2.2} y \sref{IV.1.5.4}), que todas las componentes irreducibles de $X_{ij}'\cap f^{-1}(y)$ tienen dimensión $\geq\dim(X_{ij}'\cap f^{-1}(y_i))$.
  Teniendo en cuenta \sref{IV.14.2.3} aplicado a la restricción $f^{-1}(Y_i)\to Y_i$ de $f$, que es un morfismo abierto \sref{IV.14.1.2}, $f^{-1}(y)$ es, para cada $i$, la unión de las componentes irreducibles de $X_{ij}'\cap f^{-1}(y)$ ($1\leq j\leq m_i$), luego $\dim(f^{-1}(y))\geq\dim(U\cap f^{-1}(y_i))$; pero como $f|U$ es abierto, también se tiene $\dim(f^{-1}(y))\leq\dim(f^{-1}(y_i)\cap U)$ en virtud de \sref{IV.14.2.4}; vemos, por tanto, que se puede aplicar (i) a $f|U$, de donde resulta la conclusión.
\end{enumerate}
\end{proof}

\begin{remark}[14.2.7]
\label{IV.14.2.7}
El ejemplo \sref{IV.13.2.12},~(iii)) muestra que, bajo las hipótesis de \sref{IV.14.2.5} o \sref{IV.14.2.6},~(ii)), no se puede reemplazar en la conclusión la función $z\mapsto\dim(f^{-1}(z))$ por la función $z\mapsto\mathrm{E}(z)$; en efecto, en este ejemplo, $f$ es propio y plano (por tanto universalmente abierto \sref{IV.2.4.6}) y todo entorno del único punto de $f^{-1}(y_0)$ contiene los puntos genéricos de las componentes irreducibles de $f^{-1}(\eta)$.
\end{remark}

\oldpage[IV-3]{204}

\subsection{Morfismos universalmente abiertos}
\label{subsection:IV.14.3}
\label{IV.14.3}

\begin{env}[14.3.1]
\label{IV.14.3.1}
Recordemos \sref{IV.2.4.2} que decir que un morfismo $f:X\to Y$ es universalmente abierto significa que para todo morfismo $g:Y'\to Y$, $f_{(Y')}:X_{(Y')}\to Y'$ es abierto; basta además que así sea cuando $Y'=Y[\mathrm{T}_1,\ldots,\mathrm{T}_e]$, para todo entero $e$ \sref{IV.8.10.2} (y si $Y$ es localmente noetheriano, basta por tanto que así sea para todo $Y'$ localmente noetheriano).
\end{env}

\begin{proposition}[14.3.2]
\label{IV.14.3.2}
Sea $f:X\to Y$ un morfismo de preesquemas.
Si $f$ es universalmente abierto, entonces, para todo morfismo $g:Y'\to Y$, donde $Y'$ es irreducible, la imagen por $f'=f_{(Y')}:X_{(Y')}\to Y'$ de toda componente irreducible de $X_{(Y')}$ es densa en $Y'$.
Recíprocamente, si esta condición se satisface para todo $Y'$ irreducible y todo morfismo de tipo finito $g:Y'\to Y$, y si además $f$ es localmente de presentación finita, entonces $f$ es universalmente abierto.
\end{proposition}

\begin{proof}
En efecto, se sigue de \sref{IV.1.10.4} que si $f$ es universalmente abierto, satisface la condición del enunciado.
Recíprocamente, supongamos que $f$ es localmente de presentación finita, y mostremos que para todo entero $e$, si se pone $Y''=Y[\mathrm{T}_1,\ldots,\mathrm{T}_e]$, $f_{(Y'')}$ es un morfismo abierto.
En efecto, sea $Y'$ un subpreesquema cerrado de $Y''$ cuyo espacio subyacente es un subconjunto cerrado irreducible de $Y''$; el morfismo compuesto $Y'\to Y''\to Y$ es de tipo finito, luego toda componente irreducible de $X_{(Y')}$ domina a $Y'$ por hipótesis; deducimos, por tanto, de \sref{IV.1.10.4} que $f_{(Y'')}$ es un morfismo abierto.
\end{proof}

Esta proposición muestra que la definición \sref[III]{III.4.3.9} coincide, en el caso considerado, con la definición general de los morfismos universalmente abiertos dada en \sref{IV.2.4.2}.

La noción de morfismo abierto en un punto \sref{IV.14.1.1} da lugar del mismo modo a lo siguiente:

\begin{definition}[14.3.3]
\label{IV.14.3.3}
Sea $f:X\to Y$ un morfismo de preesquemas, $x$ un punto de $X$.
Se dice que $f$ es \emph{universalmente abierto en el punto $x$} si, para todo morfismo $g:Y'\to Y$, poniendo $X'=X\times_Y Y'$, el morfismo $f'=f_{(Y')}:X'\to Y'$ es abierto en todo punto $x'$ de $X'$ cuya proyección en $X$ es $x$.
\end{definition}

\begin{remarks}[14.3.3.1]
\label{IV.14.3.3.1}
\begin{enumerate}[label=(\roman*)]
  \item El razonamiento de \sref{IV.8.10.1} muestra (con las mismas notaciones) que si $x$ es un punto de $X$ y $x_\lambda$ es su proyección en $X_\lambda$, entonces, si $f_\mu$ es abierto en el punto $x_\mu$ para todo $\mu\geq\lambda$, $f$ es abierto en el punto $x$; basta limitarse a los abiertos $U$ de $X$ que contienen a $x$ y observar que la hipótesis implica que $f_\mu(U_\mu)$ es un entorno de $f_\mu(x_\mu)$, luego $f(v_\mu^{-1}(U_\mu))$ es un entorno de $f(x)$.
  Deducimos que el enunciado \sref{IV.8.10.2} sigue siendo válido cuando se reemplaza ``universalmente abierto'' por ``universalmente abierto en el punto $x$'', y ``morfismo abierto'' por ``morfismo abierto en todo punto $x_n$ de $X_n$ cuya proyección en $X$ es $x$'': basta limitarse en la demostración a los abiertos $V$ que contienen a $x_n$.

  \item El resultado de \sref{IV.14.1.4} sigue siendo válido para un morfismo $f:X\to Y$, reemplazando ``abierto'' por ``universalmente abierto''.
  En efecto, supongamos que $f$ no es universalmente abierto en un punto $x\in f^{-1}(y)$; existe entonces un morfismo $Y'\to Y$ y un punto $x'\in X'$ que se proyecta sobre $x$, tal que $f'$ no es abierto en el punto $x'$.
\end{enumerate}
\end{remarks}

\oldpage[IV-3]{205}

Ahora bien, si $y'=f'(x')$, la proyección $f'^{-1}(y')\to f^{-1}(y)$ es un morfismo abierto \sref{IV.2.4.10}, y existe, en virtud de \sref{IV.14.1.4}, un entorno $V'$ de $x'$ en $f'^{-1}(y')$ en el que $f'$ no es abierto; por tanto, $f$ no es universalmente abierto en ningún punto de la imagen de $V'$ en $f^{-1}(y)$, que es un entorno de $x$ en $f^{-1}(y)$.

\begin{proposition}[14.3.4]
\label{IV.14.3.4}
\begin{enumerate}[label=(\roman*)]
  \item Sean $f:X\to Y$, $g:Y\to Z$ dos morfismos, $x$ un punto de $X$, $y=f(x)$.
  Si $f$ es universalmente abierto en el punto $x$ y $g$ es universalmente abierto en el punto $y$, entonces $g\circ f$ es universalmente abierto en el punto $x$.
  Recíprocamente, si $g\circ f$ es universalmente abierto en el punto $x$, $g$ es universalmente abierto en el punto $y$.

  \item Si $f:X\to Y$ es un $S$-morfismo universalmente abierto en el punto $x\in X$, entonces, para todo cambio de base $S'\to S$, $f_{(S')}:X_{(S')}\to Y_{(S')}$ es universalmente abierto en todo punto de $X_{(S')}$ situado sobre $x$.

  \item Para que $f:X\to Y$ sea universalmente abierto en el punto $x\in X$, es necesario y suficiente que $f_{\mathrm{red}}$ lo sea.

  \item Sea $f:X\to Y$ un morfismo localmente de presentación finita, $x$ un punto de $X$, $y=f(x)$; pongamos $Y_1=\Spec(\sh{O}_y)$, $X_1=X\times_Y Y_1$, $f_1=f_{(Y_1)}$.
  Para que $f$ sea universalmente abierto en el punto $x$, es necesario y suficiente que $f_1$ lo sea (recordemos \sref[I]{I.3.6.5} que $X_1$ se identifica canónicamente con un subespacio de $X$).
\end{enumerate}
\end{proposition}

\begin{proof}
En efecto, (ii) es una consecuencia evidente de la definición \sref{IV.14.3.3}; de la definición se sigue también que para demostrar la afirmación (i), basta hacerlo suprimiendo en todas partes la palabra ``universalmente'', y esto se sigue entonces de \sref{IV.14.1.2}.
La afirmación (iii) se sigue de (ii) y del hecho de que el morfismo canónico $X_{\mathrm{red}}\to X$ es sobreyectivo.
Finalmente, la condición (iv) es trivialmente necesaria.
Por otra parte, si se satisface y si $g:Y'\to Y$ es un morfismo cualquiera, pongamos $X'=X_{(Y')}$, $f'=f_{(Y')}$, y sea $x'$ un punto de $X'$ situado sobre $x$; entonces $f'$ es localmente de presentación finita, y para ver que es abierto en el punto $x'$, basta aplicar el criterio {\EGAReadLink{IV.1.10.3}{(1.10.3, $b)$)}}.
Sea $y'=f'(x')$, $Y_1'=\Spec(\sh{O}_{y'})$, $X_1'=X'\times_{Y'}Y_1'$; como $g(y')=y$, el morfismo compuesto $Y_1'\to Y'\xrightarrow{g}Y$ se factoriza como $Y_1'\to Y_1\to Y$ \sref[I]{I.2.4.4}, luego $X_1'=X_1\times_{Y_1}Y_1'$ y $f_1'=(f_1)_{(Y_1')}$; la conclusión se sigue entonces inmediatamente de la hipótesis de que $f_1'$ es abierto en el punto $x'$ y de \sref{IV.1.10.3}.
\end{proof}

\begin{proposition}[14.3.5]
\label{IV.14.3.5}
Sean $X$, $Y$ dos preesquemas, $f:X\to Y$ un morfismo, $x$ un punto de $X$.
Sea $(Y_i)_{1\leq i\leq n}$ una familia localmente finita de subpreesquemas cerrados de $Y$ tal que el espacio $Y$ es la unión de los $Y_i$, y supongamos que para todo $i$ tal que $f(x)\in Y_i$, la restricción $f^{-1}(Y_i)\to Y_i$ de $f$ es universalmente abierta en el punto $x$; entonces $f$ es universalmente abierto en el punto $x$.
\end{proposition}

\begin{proof}
Teniendo en cuenta la definición, esto se sigue de la proposición análoga {\hyperlink{ega4.c04.prop.14-1-2}{(14.1.2, (ii))}} para los morfismos abiertos en un punto.
\end{proof}

\begin{proposition}[14.3.6]
\label{IV.14.3.6}
Sea $Y$ un preesquema localmente noetheriano, $f:X\to Y$ un morfismo localmente de tipo finito.
Para que $f$ sea universalmente abierto en un punto $x\in X$, es necesario y suficiente que se satisfaga la condición siguiente: para todo morfismo $g:Y'\to Y$, donde $Y'=\Spec(A)$ es el espectro de un anillo de valoración discreta, tal que la imagen $g(y')$ del punto cerrado $y'$ de $Y'$ sea igual a $y=f(x)$, y para todo punto $x'\in X'=X\times_Y Y'$ cuyas proyecciones sobre $X$ e $Y'$ sean respectivamente $x$ e $y'$, existe una generalización $z'$ de $x'$ en $X'$ cuya proyección en $Y'$ es el punto genérico de $Y'$ (en otras palabras, \emph{existe una componente irreducible de $X'$ que contiene a $x'$ y domina a $Y'$}).

\oldpage[IV-3]{206}

Además, en la condición precedente, podemos limitarnos al caso en que $A$ es completo, su cuerpo residual es algebraicamente cerrado y $x'$ es racional sobre $k(y')$.
\end{proposition}

\begin{proof}
Si $Y'$ es como en el enunciado, la necesidad de la condición se sigue del hecho de que $f'$ debe ser abierto en el punto $x'$, y del criterio \sref{IV.1.10.3}.
Para ver que la condición es suficiente, consideremos un morfismo de tipo finito $g:Y''\to Y$\footnote{La fuente impresa tiene aquí $g:Y''\to U$, aunque $U$ no está definido; el producto fibrado posterior es sobre $Y$.}, y sean $X''=X\times_Y Y''$, $f''=f_{(Y'')}:X''\to Y''$, y $x''$ un punto de $X''$ situado sobre $x$.
Pongamos $y''=f''(x'')$, y sea $t$ una generalización de $y''$ en $Y''$, distinta de $y''$.
Como $Y''$ es localmente noetheriano, se sigue de \sref[II]{II.7.1.9} y $(\mathbf{0}_{\mathrm{III}},~10.3.1)$ que existe un esquema $Y'=\Spec(A)$, donde $A$ es un anillo de valoración discreta completo cuyo cuerpo residual es una clausura algebraica de $k(x'')$, y un morfismo $g:Y'\to Y''$ tal que, si $s$ e $y'$ son respectivamente el punto genérico y el punto cerrado, se tiene $g(s)=t$ y $g(y')=y''$.
Existe entonces un punto $x'$ de $X'=X\times_Y Y'=X''\times_{Y''}Y'$ cuyas proyecciones en $X''$ e $Y'$ son respectivamente $x''$ e $y'$, y que es racional sobre $k(y')$ \sref[I]{I.3.4.9}.
La hipótesis implica que existe una generalización $z'$ de $x'$ en $X'$ cuya proyección en $Y'$ es $s$; si $z''$ es la proyección de $z'$ en $X''$, $z''$ es una generalización de $x''$ y su proyección en $Y''$ es $t$; concluimos, por tanto, de \sref{IV.1.10.3} que $f''$ es abierto en el punto $x''$, luego $f$ es universalmente abierto en el punto $x$ {\hyperlink{ega4.c04.subsec.14-3-3-1}{(14.3.3.1, (i))}}.
\end{proof}

\begin{corollary}[14.3.7]
\label{IV.14.3.7}
Con las notaciones de \sref{IV.14.3.6}:
\begin{enumerate}[label=(\roman*)]
  \item Dado un punto $y\in Y$, para que $f$ sea universalmente abierto en todos los puntos de $f^{-1}(y)$, es necesario y suficiente que para todo morfismo $g:Y'\to Y$, donde $Y'$ es el espectro de un anillo de valoración discreta y donde la imagen $g(y')$ del punto cerrado $y'$ de $Y'$ es $y$, toda componente irreducible de $X'=X\times_Y Y'$ domine a $Y'$.
  \item Para que $f$ sea universalmente abierto, es necesario y suficiente que para todo morfismo $g:Y'\to Y$, donde $Y'$ es el espectro de un anillo de valoración discreta, toda componente irreducible de $X'=X\times_Y Y'$ domine a $Y'$.
\end{enumerate}
\end{corollary}

\begin{proof}
Está claro que basta demostrar (i); la necesidad de (i) se sigue de \sref{IV.14.3.2} y su suficiencia de \sref{IV.14.3.6}.
\end{proof}

\begin{proposition}[14.3.8]
\label{IV.14.3.8}
Sea $Y$ un preesquema localmente noetheriano, irreducible, regular y de dimensión $1$ (por ejemplo, el espectro de un anillo de Dedekind), $f:X\to Y$ un morfismo localmente de tipo finito, $y$ un punto de $Y$.
Las condiciones siguientes son equivalentes:
\begin{enumerate}[label=\emph{\alph*)}]
  \item $f_{\mathrm{red}}$ es plano en todo punto de $f^{-1}(y)$.
  \item $f$ es universalmente abierto en un entorno de $f^{-1}(y)$.
  \item $f$ es abierto en un entorno de $f^{-1}(y)$.
  \item Toda componente irreducible de $X$ que corta a $f^{-1}(y)$ domina a $Y$.
\end{enumerate}
\end{proposition}

\begin{proof}
Como $f_{\mathrm{red}}$ es localmente de tipo finito \sref{IV.1.3.4}, $a)$ implica que $f_{\mathrm{red}}$ es plano en un entorno de $f^{-1}(y)$ \sref{IV.11.1.1}, y basta aplicar \sref{IV.2.4.6} en tal entorno para ver que $a)$ implica $b)$.
La implicación $b)\Rightarrow c)$ es trivial, y la implicación $c)\Rightarrow d)$ se sigue de \sref{IV.1.10.4} aplicado a un entorno de $f^{-1}(y)$.
Queda por ver que $d)$ implica $a)$.
Podemos evidentemente, en virtud de \sref{IV.1.3.4}, reducirnos al caso en que $X$ es reducido.
Como además la cuestión es local en $X$ y en $Y$, podemos suponer que $Y=\Spec(A)$ y $X=\Spec(B)$
\oldpage[IV-3]{207}
afines; si $X_i$ ($1\leq i\leq n$) son los subpreesquemas cerrados íntegros de $X$ cuyos espacios subyacentes son las componentes irreducibles de $X$, entonces, para todo $x\in X$, $\sh{O}_{X,x}$ es un subanillo del producto directo de los $\sh{O}_{X_i,x}$; si $y=f(x)$, basta mostrar que cada $\sh{O}_{X_i,x}$ es un $\sh{O}_y$-módulo \emph{sin torsión}, pues entonces lo mismo será cierto para $\sh{O}_{X,x}$; como por hipótesis $\sh{O}_y$ es un anillo local regular de dimensión $1$, es decir, un anillo de valoración discreta \sref[II]{II.7.1.6}, se sigue entonces de $(\mathbf{0}_1,~6.3.4)$ que $\sh{O}_{X,x}$ es un $\sh{O}_y$-módulo \emph{plano}.
Pero si $X_i=\Spec(B_i)$, donde $B_i$ es un anillo íntegro, la hipótesis implica que el homomorfismo $A\to B_i$ es \emph{inyectivo} \sref[I]{I.1.2.7}; luego $B_i$ es un $A$-módulo sin torsión, y \emph{a fortiori} $\sh{O}_{X_i,x}$ es un $\sh{O}_y$-módulo sin torsión.
\end{proof}

\begin{remarks}[14.3.9]
\label{IV.14.3.9}
\begin{enumerate}[label=(\roman*)]
  \item En el enunciado de \sref{IV.14.3.8}, no se puede prescindir de la hipótesis de que $Y$ sea \emph{regular}.
  Con las notaciones de \sref{IV.11.7.5}, tomemos en efecto $Y=\Spec(A)$, $X=\Spec(\overline{A})$, de modo que $f:X\to Y$ es un morfismo \emph{finito sobreyectivo}; como $A$ es un anillo local íntegro de dimensión $1$, al igual que $\overline{A}$, se sigue inmediatamente de \sref{IV.1.10.4} que el morfismo $f$ es \emph{abierto}.
  Sin embargo, $f$ no es universalmente abierto (ni \emph{a fortiori} plano) \sref{IV.11.7.5}.
  Se podría dar un ejemplo análogo tomando como $Y$ el esquema local en un ``punto doble ordinario'' de una curva algebraica, y como $X$ la normalización de $Y$.

  \item El ejemplo de {\hyperlink{ega4.c04.item.14-1-3}{(14.1.3, (i))}} muestra que el conjunto de puntos de $X$ en los que un morfismo es universalmente abierto no es necesariamente abierto, puesto que el morfismo $f$ de ese ejemplo es universalmente abierto en todo punto en el que es abierto.
  El ejemplo $f:X\to Y$ considerado anteriormente en (i) muestra asimismo que el conjunto de puntos en los que un morfismo es universalmente abierto no es necesariamente cerrado, pues es inmediato que $f$ es universalmente abierto en todos los puntos de $X$ salvo uno (de hecho, en ellos es incluso un isomorfismo local).
  Sería interesante saber si el conjunto de puntos en los que un morfismo es universalmente abierto es localmente construible.
\end{enumerate}
\end{remarks}

Las dos proposiciones siguientes nos fueron señaladas por M.\ Artin:

\begin{lemma}[14.3.10]
\label{IV.14.3.10}
Sea $A$ un anillo de valoración (no necesariamente discreta), $k$ su cuerpo residual, $K$ su cuerpo de fracciones, y pongamos $S=\Spec(A)$.
Sea $X$ un preesquema irreducible, $f:X\to S$ un morfismo dominante de tipo finito; sea $X_0=X\otimes_A k$ (resp.\ $X_1=X\otimes_A K$) la fibra de $f$ en el punto cerrado (resp.\ en el punto genérico) de $S$.
Entonces, si $X_0\neq\varnothing$, se tiene $\dim(X_0)=\dim(X_1)$.
\end{lemma}

\begin{proof}
Podemos reducirnos al caso en que $X$ es afín, reemplazando si es necesario $X$ por un abierto afín que contenga un punto genérico de una componente irreducible de $X_0$ de dimensión máxima.
Sea $n=\dim(X_0)\geq 0$; usando \sref{IV.4.1.1.3}, se sigue de \sref{IV.13.3.1.1} que existe un entorno $U$ de $X_0$ en $X$ y un $S$-morfismo cuasifinito $g:U\to S[\mathrm{T}_1,\ldots,\mathrm{T}_n]=Z$ tal que la restricción $g_0:U_0=U\cap X_0\to Z_0=\Spec(k[\mathrm{T}_1,\ldots,\mathrm{T}_n])$ es finita y \emph{sobreyectiva}.
Por el cambio de base $\Spec(K)\to S$ y la restricción al subconjunto abierto $U_1=U\cap X_1$ de $X_1$, deducimos de $g$ un morfismo cuasifinito $g_1:U_1\to Z_1=\Spec(K[\mathrm{T}_1,\ldots,\mathrm{T}_n])$.
Como $U_1$ es denso en $X_1$ \sref{IV.4.1.1.3}, se tiene $\dim(U_1)=\dim(X_1)$; si podemos demostrar que el morfismo $g_1$ es \emph{dominante}, la proposición quedará establecida, pues entonces $\dim(U_1)\leq\dim(Z_1)=n$ y, por otra parte, $\dim(X_1)\geq n$ por \sref{IV.4.1.2}.
Supongamos lo contrario; existiría entonces un polinomio $\mathrm{F}_1\neq 0$ en $K[\mathrm{T}_1,\ldots,\mathrm{T}_n]$ tal que $g_1(U_1)\subset\mathrm{V}(\mathrm{F}_1)$.
Si $\omega$ es una valoración sobre $K$ asociada a $A$, y si $(a_\alpha)$ es la familia de coeficientes de $\mathrm{F}_1$, después de multiplicar $\mathrm{F}_1$ por un elemento no nulo de $K$, podemos suponer que $\inf_\alpha(\omega(a_\alpha))=0$; en otras palabras, $\mathrm{F}_1$ procede de un polinomio $\mathrm{F}\in A[\mathrm{T}_1,\ldots,\mathrm{T}_n]$ (con el que se identifica) tal que la imagen $\mathrm{F}_0$ de $\mathrm{F}$ en $k[\mathrm{T}_1,\ldots,\mathrm{T}_n]$ es \emph{no nula}.
Consideremos entonces el subconjunto cerrado $\mathrm{V}(\mathrm{F})$ de $Z$; se tiene $\mathrm{V}(\mathrm{F})\cap Z_1=\mathrm{V}(\mathrm{F}_1)$, y como $U_1$
\oldpage[IV-3]{208}
es denso en $U$ (pues contiene el punto genérico de $X$), $g(U)\subset\mathrm{V}(\mathrm{F})$; en particular, $g_0(U_0)\subset\mathrm{V}(\mathrm{F})\cap Z_0=\mathrm{V}(\mathrm{F}_0)$; pero como $\mathrm{F}_0\neq 0$, $\mathrm{V}(\mathrm{F}_0)$ es un subconjunto cerrado de $Z_0$ distinto de $Z_0$, y se llega a una contradicción.
\end{proof}

\begin{proposition}[14.3.11]
\label{IV.14.3.11}
Sea $f:X\to S$ un morfismo de tipo finito, $(g_\lambda)_{\lambda\in L}$ una familia de morfismos universalmente abiertos $g_\lambda:Y_\lambda\to S$, y para todo $\lambda\in L$, sea $u_\lambda:Y_\lambda\to X$ un $S$-morfismo.
Para todo $s\in S$, pongamos $X_s=X\times_S\Spec(k(s))$, $(Y_\lambda)_s=Y_\lambda\times_S\Spec(k(s))$, y sea $(u_\lambda)_s:(Y_\lambda)_s\to X_s$ el morfismo $u_\lambda\times 1$ inducido a partir de $u_\lambda$ por cambio de base.
Sea $Z(s)$ la adherencia en $X_s$ de la unión de los conjuntos $(u_\lambda)_s((Y_\lambda)_s)$, y pongamos $d(s)=\dim(Z(s))$.
Entonces, para toda generalización $s'$ de un punto $s\in S$, se tiene $d(s)\leq d(s')$.
\end{proposition}

\begin{proof}
Sabemos por \sref[II]{II.7.1.4} que existe un anillo de valoración $A'$ y un morfismo $h:S'=\Spec(A')\to S$ tal que, si $a$ (resp.\ $b$) es el punto cerrado (resp.\ el punto genérico) de $S'$, se tiene $h(a)=s$, $h(b)=s'$.
Además, la proyección $p:X_s\otimes_{k(s)}k(a)\to X_s$ es sobreyectiva y abierta \sref{IV.2.4.10}, y hace así de $X_s$ un espacio cociente de $X_s\otimes_{k(s)}k(a)$ por una relación de equivalencia abierta; por consiguiente, para todo subconjunto $M$ de $X_s$, $p^{-1}(\overline{M})=\overline{p^{-1}(M)}$ (Bourbaki, \emph{Topología general}, cap.~I, 4.ª ed., \S5, n.º~3, prop.~7).
Se razona del mismo modo para $X_{s'}$; teniendo en cuenta \sref[I]{I.3.4.8}, \sref{IV.4.2.7}, y el hecho de que los $g_\lambda$ son universalmente abiertos, vemos que basta demostrar la proposición después del cambio de base $S'\to S$.
Supongamos, por tanto, que $S'=S$, que $s$ es el punto cerrado y que $s'$ es el punto genérico de $S$.
La hipótesis de que $g_\lambda$ es abierto implica que toda componente irreducible de $Y_\lambda$ domina a $S$ \sref{IV.1.10.4}, luego su punto genérico es un punto maximal de $(Y_\lambda)_{s'}$; si $Z$ denota la adherencia en $X$ de $Z(s')$, se tiene por tanto $u_\lambda(Y_\lambda)\subset Z$, y por consiguiente $(u_\lambda)_s((Y_\lambda)_s)\subset Z_s=Z\cap X_s$; en otras palabras, $Z(s)\subset Z_s$, de donde $\dim(Z(s))\leq\dim(Z_s)$.
Pero aplicando \sref{IV.14.3.10} a un subpreesquema reducido de $X$ cuyo espacio subyacente sea una componente irreducible de $Z$, se obtiene $\dim(Z_s)\leq\dim(Z_{s'})$ (la desigualdad proviene del hecho de que puede haber componentes irreducibles de $Z$ que no corten a $X_s$).
Por otra parte, como $Z(s')$ es cerrado en $X_{s'}$ por definición, se tiene $Z_{s'}=Z(s')$, lo que termina la demostración.
\end{proof}

\begin{remark}[14.3.12]
\label{IV.14.3.12}
El caso considerado por M.\ Artin era aquel en que $Y_\lambda=S$ para todo $\lambda$, en otras palabras, el caso en que $(u_\lambda)$ es una familia de $S$-\emph{secciones} de $X$.
Otro caso útil es aquel en que la familia $(u_\lambda)$ se reduce a un único elemento; siempre se puede reducir a este caso considerando el preesquema $Y$ que es la suma de los $Y_\lambda$ y los morfismos $g:Y\to S$ y $u:Y\to X$ cuyas restricciones a cada $Y_\lambda$ son respectivamente $g_\lambda$ y $u_\lambda$.
\end{remark}

\begin{proposition}[14.3.13]
\label{IV.14.3.13}
Sea $f:X\to Y$ un morfismo localmente de tipo finito, $y$ un punto de $Y$, $x$ un punto maximal de la fibra $X_y=X\times_Y\Spec(k(y))$.
Consideremos las condiciones siguientes:
\begin{enumerate}[label=\emph{\alph*)}]
  \item $f$ es universalmente abierto en el punto $x$ (o incluso en todo punto de la componente irreducible de $X_y$ de punto genérico $x$ {\hyperlink{ega4.c04.subsec.14-3-3-1}{(14.3.3.1, (ii))}}).
  \item Para toda componente irreducible $Y_0$ de $Y$ que contenga a $y$, existe una componente irreducible $Z$ de $X$ que contiene a $x$, dominante sobre $Y_0$ y tal que $\dim_x(X_y)=\dim_x(Z\cap X_y)=\dim(Z\cap X_\eta)$, donde $\eta$ es el punto genérico de $Y_0$ (lo que implica que $Z$ es \emph{equidimensional} sobre $Y_0$ en el punto $x$ \sref{IV.13.2.2}).
\end{enumerate}

\oldpage[IV-3]{209}

\begin{enumerate}[label=\emph{\alph*$'$)}]
  \item[b$'$)] Para todo entorno abierto $U$ de $x$ en $X$ y toda generalización $y'$ de $y$, se tiene $\dim(U\cap X_{y'})\geq\dim_x(U\cap X_y)$.
\end{enumerate}

Entonces se tienen las implicaciones $a)\Rightarrow b)\Leftrightarrow b')$.
\end{proposition}

\begin{proof}
Para mostrar que $b)$ implica $b'$), basta observar que $y'$ pertenece a una componente irreducible $Y_0$ de $Y$ que contiene a $y$ y cuyo punto genérico es $\eta$; tomando $Z$ como en $b)$ y observando que el punto genérico de $Z$ (que es también el de $Z\cap X_\eta$ $(\mathbf{0}_1,~2.1.8)$) está contenido en $U$, se tiene $\dim(U\cap X_{y'})\geq\dim(U\cap Z\cap X_{y'})$, y, en virtud de \sref{IV.13.1.6}, $\dim(U\cap Z\cap X_{y'})\geq\dim(Z\cap X_\eta)$; de donde resulta la afirmación, ya que $\dim_x(U\cap X_y)=\dim_x(X_y)$ \sref{IV.4.1.1.3}.

Para demostrar que $b')$ implica $b)$, se puede reemplazar primero $X$ por $f^{-1}(Y_0)$, y suponer por tanto $Y_0=Y$; podemos reducirnos al caso en que $X$ e $Y$ son afines, puesto que $\dim_x(U\cap X_y)=\dim_x(X_y)$ \sref{IV.4.1.1.3}.
Las componentes irreducibles $W_i$ del preesquema noetheriano $X_\eta$ son entonces finitas en número, y el complementario de la unión de las $\overline{W}_i$ (adherencias \emph{en} $X$) que no contienen a $x$ es un entorno abierto $V$ de $x$.
Reemplazando $X$ por $V$, podemos suponer por tanto que $x\in\overline{W}_i$ para todo $i$ (la hipótesis $b'$) implica $\dim(V\cap X_\eta)\geq 0$, luego $x$ pertenece a $\overline{W}_i$ para al menos un $i$).
Además, las $\overline{W}_i$ son exactamente las componentes irreducibles de $X$ que dominan a $Y$ $(\mathbf{0}_1,~2.1.8)$; si se tuviera, para cada una de estas componentes $Z$, $\dim(Z\cap X_\eta)<\dim_x(X_y)$, se concluiría $\dim(X_\eta)<\dim_x(X_y)$, en contradicción con la hipótesis $b'$).
La relación $\dim_x(Z\cap X_y)=\dim(Z\cap X_\eta)$ se sigue entonces de \sref{IV.13.1.6}.

Queda por mostrar que $a)$ implica $b'$).
Teniendo en cuenta \sref[II]{II.7.1.4} y la invariancia de las hipótesis y de la conclusión por cambio de base (en virtud de \sref{IV.4.2.7}), podemos reducirnos al caso en que $Y$ es el espectro de un anillo de valoración, con punto cerrado $y$ y punto genérico $y'$, y en que $U=X$.
La hipótesis de que $f$ es abierto en el punto $x$ implica que existe una componente irreducible $Z$ de $X$ que contiene a $x$ y domina a $Y$ \sref{IV.1.10.3}.
Aplicando \sref{IV.14.3.10} a un entorno de $x$ en $Z$, concluimos que $\dim(Z\cap X_{y'})=\dim_x(Z\cap X_y)$; pero como $x$ es maximal en $X_y$, $Z\cap X_y$ contiene la componente irreducible de $X_y$ de punto genérico $x$, luego $\dim_x(Z\cap X_y)=\dim_x(X_y)$; por otra parte, se tiene $\dim(X_{y'})\geq\dim(Z\cap X_{y'})$, lo que termina la demostración de $b'$).
\end{proof}

\begin{remark}[14.3.14]
\label{IV.14.3.14}
No sabemos si, en \sref{IV.14.3.13}, la conclusión sigue siendo válida cuando se reemplaza la hipótesis $a)$ por la hipótesis más débil de que $f$ es \emph{abierto} en el punto $x$.
Sin embargo, se podría mostrar fácilmente que bastaría tratar el caso en que $Y$ es el espectro de un anillo local íntegro, cuyo punto genérico es aislado, y en que $X$ es un subpreesquema cerrado del fibrado vectorial $Y[\mathrm{T}]$.
\end{remark}
\subsection{Criterio de Chevalley para los morfismos universalmente abiertos}
\label{subsection:IV.14.4}

\begin{theorem}[14.4.1]
\label{IV.14.4.1}
Sea $f:X\to Y$ un morfismo localmente de tipo finito, $y$ un punto de $Y$, $x$ un punto maximal de la fibra $X_y=X\times_Y\Spec(k(y))$.
Supongamos que $y$ es geométricamente unirramificado.
Entonces las condiciones siguientes son equivalentes:
\begin{enumerate}[label=\emph{\alph*)}]
  \item $f$ es universalmente abierto en $x$ (o, equivalentemente, en todo punto de la componente irreducible de $X_y$ cuyo punto genérico es $x$ {\hyperlink{ega4.c04.subsec.14-3-3-1}{(14.3.3.1, (iii))}}).
\oldpage[IV-3]{210}
  \item Si $Y_0$ es la única componente irreducible de $Y$ que contiene a $y$, y $\eta$ es su punto genérico, existe una componente irreducible $Z$ de $X$ que contiene a $x$, domina a $Y_0$ y satisface $\dim_x(Z\cap X_y)=\dim(Z\cap X_\eta)$ (lo que significa que $Z$ es equidimensional sobre $Y_0$ en el punto $x$ \sref{IV.13.2.2}).
  \item[\emph{b$'$)}] Para todo entorno abierto $U$ de $x$ en $X$ y toda generalización $y'$ de $y$, se tiene $\dim(U\cap X_{y'})\geq\dim_x(U\cap X_y)$.
\end{enumerate}
Si además $Y$ es localmente noetheriano, estas condiciones son también equivalentes a la siguiente:
\begin{enumerate}[label=\emph{\alph*)}]
  \item[\emph{c)}] $f$ es abierto en $x$.
\end{enumerate}
\end{theorem}

Observemos primero que, como $(\sh{O}_y)_{\mathrm{red}}$ es íntegro, $y$ pertenece a una sola componente irreducible $Y_0$ de $Y$.
El hecho de que \emph{b)} y \emph{b$'$)} sean equivalentes y que \emph{a)} implique \emph{b$'$)} resulta de \sref{IV.14.3.13}; por otra parte, si $Y$ es localmente noetheriano, hemos visto en \sref{IV.14.2.3} que \emph{c)} implica \emph{b)}.
Queda por demostrar que, cuando $y$ es geométricamente unirramificado, \emph{b)} implica \emph{a)}.

\begin{lemma}[14.4.1.1]
\label{IV.14.4.1.1}
Sea $Y$ un preesquema, $Y'=Y[\mathrm{T}_1,\ldots,\mathrm{T}_n]$, $y'$ un punto de $Y'$ y $y$ su imagen en $Y$.
Si $Y$ es geométricamente unirramificado en el punto $y$, entonces $Y'$ es geométricamente unirramificado en el punto $y'$.
\end{lemma}

En efecto, para todo $y\in Y$, $\Spec(\kappa(y)[\mathrm{T}_1,\ldots,\mathrm{T}_n])$ es geométricamente regular sobre $\kappa(y)$ \sref[0]{0.17.3.7}, el morfismo estructural $Y'\to Y$ es liso \sref{IV.6.8.1}, y basta aplicar \sref{IV.11.3.14}.

\begin{lemma}[14.4.1.2]
\label{IV.14.4.1.2}
Sea $A$ un anillo local íntegro y unirramificado, $B$ un anillo íntegro que contiene a $A$ y es entero sobre él, y $\mathfrak{n}$ un ideal primo de $B$ situado sobre el ideal maximal $\mathfrak{m}$ de $A$.
Entonces el morfismo $\Spec(B_\mathfrak{n})\to\Spec(A)$ es sobreyectivo; dicho de otro modo, para todo ideal primo $\mathfrak{p}$ de $A$ existe un ideal primo $\mathfrak{q}$ de $B$ tal que $\mathfrak{q}\subset\mathfrak{n}$ y $\mathfrak{q}\cap A=\mathfrak{p}$.
\end{lemma}

Sea $K$ (resp.\ $L$) el cuerpo de fracciones de $A$ (resp.\ $B$), $A'$ el cierre entero de $A$, y $B'$ el subanillo de $L$ generado por $A'$ y $B$, de modo que tenemos un diagrama conmutativo de inyecciones canónicas
\[
  \xymatrix{
    B \ar[r] & B' \\
    A \ar[u] \ar[r] & A' \ar[u]
  }
\]
Como $B'$ es entero sobre $B$, existe un ideal primo $\mathfrak{n}'$ de $B'$ tal que $\mathfrak{n}'\cap B=\mathfrak{n}$ (Bourbaki, \emph{Alg.\ comm.}, cap.~V, \S~2, n\textsuperscript{o}~1, teo.~1), y (por la misma razón) $\Spec(A')\to\Spec(A)$ es sobreyectivo.
Por otra parte, como $A$ es unirramificado, $A'$ es un anillo \emph{local}, de modo que $\mathfrak{n}'\cap A'$, que está sobre el ideal maximal $\mathfrak{m}$ de $A$, es necesariamente igual al único ideal maximal $\mathfrak{m}'$ de $A'$.
En virtud del segundo teorema de Cohen--Seidenberg (\emph{loc.\ cit.}, \S~2, n\textsuperscript{o}~4, teo.~3), el morfismo $\Spec(B'_{\mathfrak{n}'})\to\Spec(A')$ es sobreyectivo, y lo mismo vale para el compuesto $\Spec(B'_{\mathfrak{n}'})\to\Spec(A')\to\Spec(A)$; pero este morfismo es también el compuesto $\Spec(B'_{\mathfrak{n}'})\to\Spec(B_\mathfrak{n})\to\Spec(A)$, por lo que el morfismo $\Spec(B_\mathfrak{n})\to\Spec(A)$ es sobreyectivo.

Establecidos estos lemas, volvamos a la demostración de la implicación $b)\Rightarrow a)$ en \sref{IV.14.4.1}.
En virtud de {\hyperlink{ega4.c04.subsec.14-3-3-1}{(14.3.3.1, (i))}}, basta demostrar que, para todo entero $n\geq 0$ y todo punto $x'$ de $X'=X[\mathrm{T}_1,\ldots,\mathrm{T}_n]$ situado sobre $x$, el morfismo
\oldpage[IV-3]{211}
$f':X'\to Y'=Y[\mathrm{T}_1,\ldots,\mathrm{T}_n]$, deducido de $f$ por cambio de base, es abierto en el punto $x'$.
Teniendo en cuenta el lema \sref{IV.14.4.1.1}, \sref{IV.2.3.4} y \sref{IV.4.2.7}, quedamos reducidos a demostrar que $f$ es \emph{abierto} en el punto $x$; además, por {\hyperlink{ega4.c04.prop.14-1-2}{(14.1.2, (iii))}}, basta evidentemente demostrar que la restricción de $f$ a un subpreesquema cerrado de $X$ cuyo espacio subyacente sea $Z$ es abierta en el punto $x$, de modo que podemos limitarnos al caso en que $X=Z$ es irreducible.
Reemplazando $X$ por un entorno abierto $V$ de $x$ tal que $V\cap X_y$ sea irreducible, podemos suponer, en virtud de \sref{IV.13.3.1}, que el morfismo $f$ se factoriza como $X\xrightarrow{g}Y''=Y[\mathrm{T}_1,\ldots,\mathrm{T}_m]\to Y$, donde $g$ es cuasifinito, dominante y localmente de tipo finito.
Como el morfismo estructural $Y''\to Y$ es abierto \sref{IV.2.4.6}, quedamos reducidos a demostrar que $g:X\to Y''$ es abierto en el punto $x$.
Por \sref{IV.14.4.1.1}, $g(x)$ es un punto geométricamente unirramificado de $Y''$.
Quedamos así reducidos a demostrar el lema siguiente:

\begin{lemma}[14.4.1.3]
\label{IV.14.4.1.3}
Sean $X$, $Y$ dos preesquemas irreducibles, y $f:X\to Y$ un morfismo localmente cuasifinito y dominante.
Si $x\in X$ es tal que $y=f(x)$ es un punto unirramificado de $Y$, entonces $f$ es abierto en el punto $x$.
\end{lemma}

Es suficiente demostrar que $f(\Spec(\sh{O}_{X,x}))=\Spec(\sh{O}_{Y,y})$ \sref{IV.1.10.3}.
Por tanto, mediante el cambio de base $\Spec(\sh{O}_{Y,y})\to Y$, podemos limitarnos al caso $Y=\Spec(A)$, donde $A$ es un anillo local y $y$ es el punto cerrado de $Y$ (teniendo en cuenta \sref[I]{I.3.6.5} y \sref[$0_{\mathrm{I}}$]{0.2.1.8}); esto prueba que $X\times_Y\Spec(\sh{O}_{Y,y})$ es irreducible; reemplazando $f$ por $f_{\mathrm{red}}$, podemos suponer que $X$ e $Y$ son reducidos, y por tanto íntegros.
Reemplazando finalmente $X$ e $Y$ por entornos afines de $x$ y $y$, respectivamente, podemos suponer, por \sref{IV.8.12.9}, que el morfismo $f$ se factoriza como $X\xrightarrow{j}X_1\xrightarrow{g}Y$, donde $j$ es una \emph{inmersión abierta} y $g$ un morfismo \emph{finito} (evidentemente dominante); como $X$ y $X_1$ son afines, $j$ es afín, luego separado y cuasicompacto, y por consiguiente $j$ se factoriza como $X\xrightarrow{u}X_2\xrightarrow{h}X_1$, donde $X_2$ es la imagen cerrada de $X$ por $j$, $h$ la inyección canónica y $u$ una inmersión abierta \sref[I]{I.9.5.3}.
En otros términos, podemos suponer que $X_1$ es íntegro, o incluso que $X_1=\Spec(B)$, donde $B$ es una $A$-álgebra íntegra y finita que contiene a $A$, puesto que $g$ es dominante. Si $\mathfrak{n}$ es el ideal primo de $B$ correspondiente al punto $x$, la hipótesis de que $A$ es unirramificado implica, por \sref{IV.14.4.1.2}, que el morfismo $\Spec(B_\mathfrak{n})\to\Spec(A)$ es sobreyectivo, es decir, que $\Spec(\sh{O}_{X,x})\to\Spec(\sh{O}_{Y,y})$ es sobreyectivo.
Q.E.D.

\begin{corollary}[14.4.2]
\label{IV.14.4.2}
Sea $f:X\to Y$ un morfismo localmente de tipo finito, $y$ un punto geométricamente unirramificado de $Y$, y $\eta$ el punto genérico de la única componente irreducible $Y_0$ de $Y$ que contiene a $y$.
Las condiciones siguientes son equivalentes:
\begin{enumerate}[label=\emph{\alph*)}]
  \item $f$ es universalmente abierto en todos los puntos de $X_y$ (o, lo que equivale a lo mismo, en los puntos maximales de $X_y$ {\hyperlink{ega4.c04.subsec.14-3-3-1}{(14.3.3.1, (ii))}}).
  \item Para todo $x\in X_y$, existe una componente irreducible $Z$ de $X$ que contiene a $x$ y es equidimensional sobre $Y$ en el punto $x$ \sref{IV.13.2.2}.
  \item[\emph{b$'$)}] Para todo $x\in X_y$ y todo entorno abierto $U$ de $x$ en $X$, se tiene
    \[
      \dim(U\cap X_{\eta})\geq\dim_x(U\cap X_y).
    \]
  \item[\emph{b$''$)}] Para todo abierto $U$ de $X$, se tiene $\dim(U\cap X_{\eta})\geq\dim(U\cap X_y)$.
\end{enumerate}
\end{corollary}
\oldpage[IV-3]{212}
Si además $Y$ es localmente noetheriano, estas condiciones son también equivalentes a la siguiente:
\begin{enumerate}[label=\emph{\alph*)}]
  \item[\emph{c)}] $f$ es abierto en todos los puntos de $X_y$ (o, lo que equivale a lo mismo, en los puntos maximales de $X_y$ {\hyperlink{ega4.c04.subsec.14-3-3-1}{(14.3.3.1, (ii))}}).
\end{enumerate}
La equivalencia de \emph{a)} y \emph{c)} cuando $Y$ es localmente noetheriano resulta de \sref{IV.14.4.1}; las condiciones \emph{b)} o \emph{b$'$)}, aplicadas a los puntos \emph{maximales} de $X_y$, implican \emph{a)} también en virtud de {\hyperlink{ega4.c04.subsec.14-3-3-1}{(14.3.3.1, (ii))}}; finalmente, \emph{b$'$)} y \emph{b$''$)} son equivalentes, ya que
\[
  \dim(U\cap X_y) = \sup_{x\in X_y}(\dim_x(U\cap X_y)).
\]

Queda por ver que la condición \emph{a)} implica \emph{b)} y \emph{b$'$)} en \emph{todo} punto $x\in X_y$.
Sea $d=\dim_x(X_y)$, y sea $x'$ el punto genérico de una componente irreducible de $X_y$ que contiene a $x$ y tiene dimensión $d$.
En virtud de \emph{a)} y de \sref{IV.14.4.1}, existe una componente irreducible $Z$ de $X$ que contiene a $x'$ y es equidimensional sobre $Y$ en el punto $x'$, de modo que $\dim_{x'}(Z\cap X_y)=\dim(Z\cap X_\eta)$.
Pero, por construcción, $\dim_{x'}(Z\cap X_y)=\dim_x(X_y)=d$, y $\dim_x(Z\cap X_y)\leq\dim(Z\cap X_y)=d$; teniendo en cuenta \sref{IV.13.1.6}, esto prueba que $Z$ es equidimensional sobre $Y$ en el punto $x$; por tanto, \emph{a)} implica \emph{b)}.
Además, tenemos $\dim(X_\eta)\geq\dim(Z\cap X_\eta)=d=\dim_x(X_y)$.
Reemplazando $X$ por un entorno abierto $U$ de $x$, vemos así que \emph{a)} implica \emph{b$'$)}.
Q.E.D.

\begin{corollary}[14.4.3]
\label{IV.14.4.3}
Sea $Y$ un preesquema geométricamente unirramificado, y $f:X\to Y$ un morfismo localmente de tipo finito.
Las condiciones siguientes son equivalentes:
\begin{enumerate}[label=\emph{\alph*)}]
  \item $f$ es universalmente abierto.
  \item Para todo abierto $U$ de $X$, todo $y\in Y$ y toda generalización $y'$ de $y$, se tiene
    \[
      \dim(U\cap X_{y'})\geq\dim(U\cap X_y).
    \]
\end{enumerate}
Si además $Y$ es localmente noetheriano, estas condiciones son también equivalentes a la siguiente:
\begin{enumerate}[label=\emph{\alph*)}]
  \item[\emph{c)}] $f$ es abierto.
\end{enumerate}
\end{corollary}

\begin{corollary}[14.4.4]
\label{IV.14.4.4}
(Criterio de Chevalley).
Sea $f:X\to Y$ un morfismo localmente de tipo finito.
\begin{enumerate}[label=(\roman*)]
  \item Si $f$ es equidimensional en un punto $x\in X$ \sref{IV.13.3.2} y si $y=f(x)$ es un punto geométricamente unirramificado de $Y$, entonces $f$ es universalmente abierto en el punto $x$.
  \item Si $Y$ es geométricamente unirramificado, $f$ es universalmente abierto en todos los puntos de $X$ donde $f$ es equidimensional, y el conjunto de estos puntos es abierto en $X$.
 En particular, si $f$ es equidimensional, es universalmente abierto.
\end{enumerate}
\end{corollary}

La afirmación (ii) es una consecuencia trivial de (i), pues ya sabemos que el conjunto de puntos donde $f$ es equidimensional es abierto \sref{IV.13.3.2}.
En cuanto a la afirmación (i), resulta del hecho de que la hipótesis implica que la condición \emph{b)} de \sref{IV.14.4.1} se satisface en el punto genérico de una componente irreducible de $X_y$ que contiene a $x$, teniendo en cuenta \sref{IV.13.3.1}; basta, por tanto, aplicar \sref{IV.14.4.1}.

\begin{remark}[14.4.5]
\label{IV.14.4.5}
Se puede demostrar que, si $Y$ es localmente noetheriano y todas las generalizaciones de $y=f(x)$ son puntos geométricamente unirramificados de $Y$ (cf.\ \sref{IV.6.15.2}), entonces, si $f$ es equidimensional en el punto $x$, es universalmente abierto \emph{en un entorno de $x$}.
\end{remark}

\begin{corollary}[14.4.6]
\label{IV.14.4.6}
\oldpage[IV-3]{213}
Sea $Y$ un preesquema localmente noetheriano, y $f:X\to Y$ un morfismo localmente de tipo finito.
Sea $y$ un punto geométricamente unirramificado de $Y$, y supongamos además que, para todo $x\in f^{-1}(y)$, el anillo $\sh{O}_{X,x}$ es equidimensional.
Entonces las condiciones siguientes son equivalentes:
\begin{enumerate}[label=\emph{\alph*)}]
  \item $f$ es equidimensional \sref{IV.13.3.2} en todos los puntos de $f^{-1}(y)$.
  \item $f$ es abierto en todos los puntos de $f^{-1}(y)$.
  \item $f$ es universalmente abierto en todos los puntos de $f^{-1}(y)$.
\end{enumerate}
\end{corollary}

En efecto, \emph{c)} implica trivialmente \emph{b)}, y \emph{a)} implica \emph{c)} en virtud de \sref{IV.14.4.4}; finalmente, debido a las hipótesis sobre $\sh{O}_y$ y $\sh{O}_x$, \emph{b)} implica \emph{a)} por \sref{IV.14.2.2}.
Más generalmente:

\begin{proposition}[14.4.7]
\label{IV.14.4.7}
Sea $Y$ un preesquema localmente noetheriano, $f:X\to Y$ un morfismo localmente de tipo finito, y $x$ un punto de $X$, tal que $y=f(x)$ es geométricamente unirramificado.
Las condiciones siguientes son equivalentes:
\begin{enumerate}[label=\emph{\alph*)}]
  \item $f$ es equidimensional \sref{IV.13.3.2} en el punto $x$.
  \item El anillo $\sh{O}_x$ es equidimensional y $f$ es abierto en los puntos genéricos de las componentes irreducibles de $f^{-1}(y)$ que contienen a $x$ (y, por tanto, también en todo punto de una componente de ese tipo).
  \item El anillo $\sh{O}_x$ es equidimensional y $f$ es universalmente abierto en los puntos genéricos de las componentes irreducibles de $f^{-1}(y)$ que contienen a $x$ (y, por tanto, también en todo punto de una componente de ese tipo).
\end{enumerate}
Además, cuando se satisfacen estas condiciones, para todo subpreesquema cerrado reducido $X_i$ de $X$ cuyo espacio subyacente sea una componente irreducible de $X$ que contiene a $x$, la restricción $f_i:X_i\to Y$ de $f$ es un morfismo equidimensional en el punto $x$, y universalmente abierto en todos los puntos de las componentes irreducibles de $f^{-1}(y)\cap X_i$ que contienen a $x$.
\end{proposition}

La condición \emph{a)} implica que $f$ es equidimensional en todos los puntos genéricos de las componentes irreducibles de $f^{-1}(y)$ que contienen a $x$ \sref{IV.13.3.1}, y, por consiguiente, en virtud de \sref{IV.14.4.4}, universalmente abierto en esos puntos; el mismo argumento aplicado a cada $f_i$ (teniendo en cuenta \sref{IV.13.3.3}) demuestra la última afirmación de la proposición, a la vista de {\hyperlink{ega4.c04.subsec.14-3-3-1}{(14.3.3.1, (iii))}}.
Además, por \sref{IV.14.2.1}, tenemos las relaciones
\begin{align}
  \dim(\sh{O}_{X_i,x}) &= \dim(\sh{O}_y) + \dim(\sh{O}_{X_i,x}\otimes_{\sh{O}_y}\kappa(y))
  \tag{14.4.7.1}\label{IV.14.4.7.1}\\
  \dim(\sh{O}_{X,x}) &= \dim(\sh{O}_y) + \dim(\sh{O}_{X,x}\otimes_{\sh{O}_y}\kappa(y))
  \tag{14.4.7.2}\label{IV.14.4.7.2}
\end{align}
Y como $f$ es equidimensional en el punto $x$, se deduce de \sref{IV.13.3.1} que $\dim_x(X_i\cap f^{-1}(y))=\dim_x(f^{-1}(y))$, y por tanto, mediante \sref{IV.5.2.3}, $\dim(\sh{O}_{X_i,x}\otimes_{\sh{O}_y}\kappa(y))=\dim(\sh{O}_{X,x}\otimes_{\sh{O}_y}\kappa(y))$.
Concluimos que $\dim(\sh{O}_{X_i,x})=\dim(\sh{O}_{X,x})$ para todo $i$; dicho de otro modo, $\sh{O}_{X,x}$ es equidimensional, y esto termina de demostrar que \emph{a)} implica \emph{c)}.
Es claro que \emph{c)} implica \emph{b)}; finalmente, \emph{b)} implica la relación \eqref{IV.14.4.7.2} en virtud de \sref{IV.14.2.1}; se deduce entonces de \sref{IV.13.3.6} que \emph{b)} implica \emph{a)}.

\begin{proposition}[14.4.8]
\label{IV.14.4.8}
Sea $Y$ un preesquema noetheriano, $f:X\to Y$ un morfismo localmente de tipo finito, $y$ un punto de $Y$, y $x$ un punto maximal de $X_y=f^{-1}(y)$.
Las condiciones siguientes son equivalentes:
\begin{enumerate}[label=\emph{\alph*)}]
  \item El morfismo $f$ es universalmente abierto en $x$; dicho de otro modo, para todo cambio de base $g:Y'\to Y$ se satisface la propiedad $\mathbf{P}(Y')$: para todo punto $x'$ de $X'=X\times_Y Y'$ situado sobre $x$, el morfismo $f'=f_{(Y')}:X'\to Y'$ es abierto en $x'$.
\oldpage[IV-3]{214}
  \item[\emph{a$'$)}] La propiedad $\mathbf{P}(Y')$ se satisface para todo morfismo finito $g:Y'\to Y$.
  \item[\emph{a$''$)}] La propiedad $\mathbf{P}(Y'')$ se satisface para la normalización $Y''$ de $Y_{\mathrm{red}}$ \sref[II]{II.6.3.8}.
  \item Para todo punto $x''$ de $X''=X\times_Y Y''$ situado sobre $x$, existe una componente irreducible $Z''$ de $X''$ que contiene a $x''$ y es equidimensional sobre $Y''$ en el punto $x''$.
\end{enumerate}
\end{proposition}

Es trivial que \emph{a)} implica \emph{a$'$)}.
Para demostrar que \emph{a$'$)} implica \emph{a$''$)}, observemos que podemos escribir $Y''=\Spec(\sh{B})$, donde $\sh{B}$ es una $\sh{O}_Y$-álgebra cuasicoherente entera sobre $\sh{O}_Y$; como $Y$ es noetheriano, $\sh{B}$ es el límite inductivo de sus sub-$\sh{O}_Y$-álgebras $\sh{B}_\lambda$, que son cuasicoherentes y de tipo finito \sref[I]{I.9.6.6}; pero entonces las $\sh{B}_\lambda$ son $\sh{O}_Y$-álgebras \emph{finitas} \sref[II]{II.6.1.2}; por tanto podemos escribir $Y''=\varprojlim Y'_\lambda$, donde $Y'_\lambda=\Spec(\sh{B}_\lambda)$, y en consecuencia $X''=X\times_Y Y''=\varprojlim X'_\lambda$, con $X'_\lambda=X\times_Y Y'_\lambda$.
En virtud de \emph{a$'$)}, los morfismos $f'_\lambda:X'_\lambda\to Y'_\lambda$ son abiertos en todos los puntos de $X'_\lambda$ situados sobre $x$; concluimos que $f''$ es abierto en todos los puntos de $X''$ situados sobre $x$, por \sref{IV.8.10.1} y {\hyperlink{ega4.c04.subsec.14-3-3-1}{(14.3.3.1, (i))}}.

Como el preesquema $Y''$ es \emph{normal} por definición, el hecho de que \emph{b)} implique \emph{a$''$)} se deduce de \sref{IV.14.4.4}, aplicado a la componente irreducible equidimensional del enunciado y a la restricción de $f''$ a dicha componente.
Queda por demostrar que \emph{a$''$)} implica \emph{a)} y \emph{b)}.
Teniendo en cuenta \sref{IV.1.10.3}, podemos limitarnos al caso $Y=\Spec(\sh{O}_y)$, observando que el morfismo canónico $\Spec(\sh{O}_y)\to Y$ es universalmente bicontinuo \sref[I]{I.3.6.5}, mientras que $Y''\times_Y\Spec(\sh{O}_y)$ es la normalización de $(\Spec(\sh{O}_y))_{\mathrm{red}}$, como se deduce de la conmutatividad de las operaciones de cierre entero y localización (Bourbaki, \emph{Alg.\ comm.}, cap.~V, \S~1, n\textsuperscript{o}~5, prop.~16).
Supongamos, pues, que $Y=\Spec(A)$, donde $A$ es un anillo local noetheriano, y $y$ es el punto cerrado de $Y$; sabemos por \sref[0]{0.23.2.5} que existe una factorización
\[
  Y'' \xrightarrow{u} Y_1 \xrightarrow{v} Y
\]
del morfismo estructural, tal que $v$ es un morfismo \emph{finito} y sobreyectivo, mientras que $u$ es entero, \emph{radicial} y dominante, luego sobreyectivo \sref[II]{II.6.1.10} y, por consiguiente, un \emph{homeomorfismo} universal \sref{IV.2.4.5}.
Si ponemos $X_1=X\times_Y Y_1$, la proyección $X''\to X_1$ es, por tanto, un homeomorfismo, y la hipótesis \emph{a$''$)} implica que $f_1=f_{(Y_1)}:X_1\to Y_1$ es abierto en todos los puntos de $X_1$ situados sobre $x$.
Además, esto muestra que, para demostrar la propiedad \emph{b)}, basta demostrar la misma propiedad reemplazando $Y''$, $X''$ y $x''$ por $Y_1$, $X_1$ y un punto $x_1$ de $X_1$ situado sobre $x$.
Pero $Y_1$ es noetheriano y, además, geométricamente unirramificado, puesto que $Y''$ es normal y $u$ es radicial \sref{IV.6.15.1}; la propiedad \emph{b)} resulta, por tanto, de \sref{IV.14.4.1}.
Queda por demostrar que $f$ es universalmente abierto en el punto $x$, lo que resultará del lema siguiente:

\begin{lemma}[14.4.8.1]
\label{IV.14.4.8.1}
Sea $v:Y_1\to Y$ un morfismo cerrado (resp.\ universalmente cerrado) y sobreyectivo.
Para que un morfismo $f:X\to Y$ sea abierto (resp.\ universalmente abierto) en un punto $x\in X$, basta que $f_1=f_{(Y_1)}:X_1=X\times_Y Y_1\to Y_1$ sea abierto (resp.\ universalmente abierto) en todos los puntos de $X_1$ situados sobre $x$.
\end{lemma}

La segunda afirmación resulta trivialmente de la primera y del hecho de que, para todo cambio de base $Y'\to Y$, el morfismo $v_{(Y')}:Y_1\times_Y Y'\to Y'$ sigue siendo sobreyectivo y universalmente cerrado.
Para demostrar la primera afirmación, consideremos
\oldpage[IV-3]{215}
un entorno abierto $U$ de $x$ en $X$; como $v$ es cerrado y sobreyectivo, para que $f(U)$ sea un entorno de $y=f(x)$ es necesario y suficiente que $v^{-1}(f(U))$ sea un entorno de $v^{-1}(y)$.
Pero si $p:X_1\to X$ es la proyección canónica, tenemos $v^{-1}(f(U))=f_1(p^{-1}(U))$ \sref[I]{I.3.4.8}, y la hipótesis implica que $f_1(p^{-1}(U))$ es un entorno de $f_1(p^{-1}(x))=v^{-1}(y)$ \sref[I]{I.3.4.8}.

\begin{corollary}[14.4.9]
\label{IV.14.4.9}
Sea $Y$ un preesquema noetheriano, y $f:X\to Y$ un morfismo localmente de tipo finito.
Las condiciones siguientes son equivalentes:
\begin{enumerate}[label=\emph{\alph*)}]
  \item $f$ es universalmente abierto; dicho de otro modo, para todo cambio de base $Y'\to Y$, el morfismo $f_{(Y')}:X\times_Y Y'\to Y'$ es abierto.
  \item[\emph{a$'$)}] Para todo morfismo finito $Y_1\to Y$, $f_{(Y_1)}$ es abierto.
  \item[\emph{a$''$)}] Si $Y''$ es la normalización de $Y_{\mathrm{red}}$, $f_{(Y'')}$ es abierto.
  \item Para todo punto $x''$ de $X''=X\times_Y Y''$, existe una componente irreducible $Z''$ de $X''$ que contiene a $x''$ y es equidimensional sobre $Y''$ en el punto $x''$ (cf.\ {\hyperlink{ega4.c04.item.14-4-10}{(14.4.10, (ii))}}).
\end{enumerate}
\end{corollary}

Esto resulta inmediatamente de \sref{IV.14.4.8} y \sref{IV.14.1.4}.

\begin{remarks}[14.4.10]
\label{IV.14.4.10}
(i) La equivalencia de las condiciones \emph{a)} y \emph{b)} en \sref{IV.14.4.8} (resp.\ \sref{IV.14.4.9}) es también válida para todo preesquema $Y$ y todo morfismo $f$ localmente de tipo finito.
En efecto, \emph{a)} implica \emph{b)} en virtud de \sref{IV.14.4.1}; recíprocamente, \emph{b)} implica que $f''$ es universalmente abierto en los puntos de $X''$ situados sobre $x$ (resp.\ en todo punto de $X''$) en virtud de \sref{IV.14.4.1}, y concluimos la propiedad \emph{a)} aplicando el lema \sref{IV.14.4.8.1} al morfismo entero y sobreyectivo $Y''\to Y$.

Puede que, en \sref{IV.14.4.1}, la hipótesis suplementaria de que $Y$ es noetheriano sea superflua para la equivalencia de \emph{a)} y \emph{c)} (cf.\ \sref{IV.14.3.14}).
En tal caso, las hipótesis de noetherianidad también serían superfluas en \sref{IV.14.4.2}, \sref{IV.14.4.3}, \sref{IV.14.4.8} y \sref{IV.14.4.9}.

(ii) Se pueden dar ejemplos de morfismos $f:X\to Y$ con las propiedades siguientes: $Y$ es noetheriano, regular y de dimensión~2, $f$ es \emph{universalmente abierto} y de tipo finito, $X$ tiene dos componentes irreducibles $X_1$, $X_2$, pero la restricción $X_1\to Y$ de $f$ a una de ellas \emph{no es un morfismo abierto}.
El principio de la construcción se basa en el método general de ``pegado'', que se presentará en el cap.~V y que aquí solo podemos esbozar.
Comencemos con un punto cerrado $y$ de $Y$, y consideremos el $Y$-esquema $Y_1$ obtenido haciendo la explosión de $y$ \sref[II]{II.8.1.3}. Si $f_1:Y_1\to Y$ es el morfismo estructural, sabemos que la restricción de $f_1$ a $f_1^{-1}(Y-\{y\})$ es un \emph{isomorfismo} sobre $Y-\{y\}$ (\emph{loc.\ cit.}), mientras que la fibra $f_1^{-1}(y)$ es isomorfa a $\operatorname{Proj}(S)$, donde $S=\bigoplus_{k=0}^{\infty}\mathfrak{m}_y^k/\mathfrak{m}_y^{k+1}$ \sref[II]{II.3.5.3}, es decir, aquí $\mathbf{P}^1_{\kappa(y)}$; se deduce de \sref{IV.14.4.1} que $f_1$ no es abierto en el punto genérico de $f_1^{-1}(y)$.
Por otra parte, pongamos $Y_2=\mathbf{P}^1_Y$, y sea $f_2:Y_2\to Y$ el morfismo estructural. Se deduce de \sref[II]{II.8.4.4} que $f_2$ es plano y, por tanto, universalmente abierto \sref{IV.2.4.6}; además, por \sref[II]{II.3.5.3}, $f_2^{-1}(y)$ es isomorfo a $\mathbf{P}^1_{\kappa(y)}$. Basta entonces ``pegar'' $Y_1$ y $Y_2$ a lo largo de las fibras isomorfas $f_1^{-1}(y)$ y $f_2^{-1}(y)$; esto da un morfismo $f:X\to Y$ cuyas componentes irreducibles $X_1$, $X_2$ se identifican canónicamente con $Y_1$, $Y_2$, respectivamente, y cuyas restricciones a estas componentes son $f_1$, $f_2$.

Recordemos, sin embargo, \sref{IV.12.1.1.5} que, si $Y$ es localmente noetheriano, $f:X\to Y$ es de tipo finito y \emph{plano}, toda componente irreducible de $X$ es equidimensional sobre $Y$ (y, en consecuencia, la restricción de $f$ a una componente de ese tipo es \emph{universalmente abierta} si todos los puntos de $Y$ son \emph{geométricamente unirramificados}).

(iii) Recordemos {\hyperlink{ega4.c04.item.12-1-2}{(12.1.2, (i))}} que existen morfismos $f:X\to Y$ con las propiedades siguientes: $Y$ es noetheriano (no geométricamente unirramificado), $f$ es finito y plano (e incluso \'etale \sref{IV.17.6.3}), pero la restricción de $f$ a una componente irreducible de $X$ \emph{no es un morfismo abierto} (aunque $f$ mismo sí lo es por \sref{IV.2.4.6}).

(iv) El criterio de Chevalley \sref{IV.14.4.4} explica la importancia de la noción de morfismo universalmente abierto.
Esta noción permite, en efecto, sustituir en numerosos resultados, más o menos clásicos, una hipótesis de \emph{normalidad} por la hipótesis de que cierto morfismo es universalmente abierto; el enunciado más general
\oldpage[IV-3]{216}
así obtenido se aplicará en particular a los morfismos \emph{planos} \sref{IV.2.4.6}, cuya importancia en geometría algebraica va en aumento.
Puede considerarse que los enunciados que contienen la hipótesis de que un morfismo es universalmente abierto son generalizaciones comunes de enunciados que contienen una hipótesis de normalidad y de enunciados que contienen una hipótesis de planitud.
\end{remarks}


\subsection{Morfismos universalmente abiertos y cuasisecciones}
\label{subsection:IV.14.5}
\label{IV.14.5}

\begin{lemma}[14.5.1]
\label{IV.14.5.1}
Sea $Y$ un preesquema irreducible y localmente noetheriano, $f:X\to Y$ un morfismo localmente de tipo finito, y $x$ un punto de $X$, tales que $f$ es equidimensional \sref{IV.13.3.2} en el punto $x$; pongamos $y=f(x)$, $e=\dim_x(f^{-1}(y))$.
Sea $X'$ un subespacio cerrado irreducible de $X$ que contiene a $x$, y sea $n$ un entero tal que $\codim(X',X)\leq n$ y $\dim_x(X'\cap f^{-1}(y))\leq e-n$.
Entonces necesariamente $\codim(X',X)=n$, $\dim_x(X'\cap f^{-1}(y))=e-n$, y la restricción $X'\to Y$ de $f$ es un morfismo equidimensional en $x$ (y, a fortiori, dominante).
\end{lemma}

Como la cuestión es local en $X$, podemos suponer que $f$ es de tipo finito y \emph{equidimensional}, de modo que para \emph{todo} $z\in Y$, \emph{todas} las componentes irreducibles de $f^{-1}(z)$ tienen dimensión $e$ {\hyperlink{ega4.c04.prop.13-3-1}{(13.3.1, $a''$))}}.
Sea $x'$ el punto genérico de $X'$, $y'=f(x')$, $Y'=\overline{\{y'\}}=\overline{f(X')}$ y pongamos $Z=f^{-1}(Y')$.
En virtud de \sref[0]{0.14.2.2}, tenemos
\begin{equation}
  \codim(X',Z)\leq n
  \tag{14.5.1.1}\label{IV.14.5.1.1}
\end{equation}
 y, si los dos miembros son iguales, necesariamente $\codim(X',X)=n$ y $Z$ contiene una componente irreducible de $X$, lo que implica por \sref{IV.13.3.1} que $f(X')$ es denso en $Y$ y, por tanto, $Z=X$; así, la igualdad $\codim(X',Z)=n$ equivale a la conjunción de la igualdad $\codim(X',X)=n$ y la relación $Z=X$.

Por otra parte, razonando en los preesquemas reducidos de $Y$ y $X$ que tienen respectivamente a $Y'$ y $Z$ como espacios subyacentes, deducimos de \sref{IV.5.1.2} y de \sref[I]{I.3.6.5} que
\begin{equation}
  \codim(X'\cap f^{-1}(y'),f^{-1}(y'))=\codim(X',Z).
  \tag{14.5.1.2}\label{IV.14.5.1.2}
\end{equation}
Por hipótesis, $f^{-1}(y')$ es biequidimensional \sref{IV.5.2.1} y de dimensión $e$, luego, por \sref[0]{0.14.3.5}, en virtud de \eqref{IV.14.5.1.2} y \eqref{IV.14.5.1.1},
\begin{equation}
  e'=\dim(X'\cap f^{-1}(y'))=e-\codim(X',Z)\geq e-n
  \tag{14.5.1.3}\label{IV.14.5.1.3}
\end{equation}
la igualdad se da si y solo si $\codim(X',X)=n$ y $f(X')$ es denso en $Y$.
Finalmente, por \sref{IV.13.1.1}, tenemos $\dim_x(X'\cap f^{-1}(y))\geq e'$, de donde, por \eqref{IV.14.5.1.3},
\begin{equation}
  \dim_x(X'\cap f^{-1}(y))\geq e'\geq e-n.
  \tag{14.5.1.4}\label{IV.14.5.1.4}
\end{equation}
Pero, por hipótesis, también tenemos $\dim_x(X'\cap f^{-1}(y))\leq e-n$, de donde se siguen las conclusiones de la proposición.

\begin{corollary}[14.5.2]
\label{IV.14.5.2}
Bajo las hipótesis de \sref{IV.14.5.1} sobre $f$, $X$, $Y$ y $x$, supongamos además que $x$ no es un punto maximal de $f^{-1}(y)$.
Entonces existen un entorno abierto afín $U$ de $x$ en $X$ y una sección $g\in\Gamma(U,\sh{O}_X)$ tales que el conjunto $X'$ de los $x'\in U$ que satisfacen $g(x')=0$ contiene a $x$ y no contiene ningún punto maximal de $f^{-1}(y)$.
Para todo $g$ con estas propiedades, $X'$ es equidimensional sobre $Y$ en el punto $x$, y tenemos
\begin{equation}
  \dim_x(X'\cap f^{-1}(y))=e-1 \qquad\text{y}\qquad \codim(X',X)=1.
  \tag{14.5.2.1}\label{IV.14.5.2.1}
\end{equation}
\end{corollary}

\oldpage[IV-3]{217}
Podemos limitarnos al caso en que $X=U$ es un entorno abierto afín de $x$ tal que todas las componentes irreducibles de $f^{-1}(y)$ contienen a $x$.
Estas componentes corresponden a los ideales primos mínimos de $\sh{O}_x/\mathfrak{m}_y\sh{O}_x$, y por hipótesis estos ideales son distintos de $\mathfrak{m}_x/\mathfrak{m}_y\sh{O}_x$ (Bourbaki, \emph{Alg.\ comm.}, cap.~II, \S~1, n\textsuperscript{o}~1, prop.~2); para obtener un $g\in\Gamma(U,\sh{O}_X)$ que satisfaga las condiciones del enunciado, basta tomar $g\in\mathfrak{j}_x$ tal que la imagen de $g$ en $\mathfrak{m}_x$ no pertenezca a ninguno de los ideales primos precedentes.
Además, como $\codim(X',X)\leq 1$ \sref{IV.5.1.8}, y como $X'$ no contiene ninguna de las componentes irreducibles de $f^{-1}(y)$ de dimensión $e$, y puesto que estas son todas las componentes irreducibles de $f^{-1}(y)$, tenemos, por \sref[0]{0.14.2.2.2}, $\dim_x(X'\cap f^{-1}(y))\leq e-1$.
Basta entonces aplicar \sref{IV.14.5.1}.

\begin{proposition}[14.5.3]
\label{IV.14.5.3}
Sea $Y$ un preesquema irreducible y localmente noetheriano, $f:X\to Y$ un morfismo localmente de tipo finito, y $x$ un punto de $X$.
Supongamos que $f$ es equidimensional en el punto $x$ y que $x$ es cerrado en $f^{-1}(f(x))$.
Entonces existe un subconjunto irreducible $X'$ de $X$, localmente cerrado en $X$, que contiene a $x$, tal que la restricción $X'\to Y$ de $f$ (donde $X'$ está dotado de la estructura de subesquema reducido inducida por $X$) es un morfismo \emph{cuasifinito y dominante}.
\end{proposition}

En efecto, con la notación de \sref{IV.14.5.2}, la hipótesis de que $x$ es cerrado en $f^{-1}(y)$ implica que $x$ no es un punto maximal de $X'\cap f^{-1}(y)$ mientras que $e-1\geq 1$.
Basta, por tanto, aplicar \sref{IV.14.5.2} por inducción descendente sobre $e=\dim_x(f^{-1}(f(x)))$ hasta llegar a $e=1$; la aplicación de \sref{IV.14.5.2} en este último caso proporciona un subespacio $X'$ tal que $X'\cap f^{-1}(f(x))$ es noetheriano y de dimensión $0$, luego finito y discreto; como $X'$ es equidimensional sobre $Y$, $X'\cap f^{-1}(f(x'))$ tiene dimensión $0$ para todo $x'\in X'$, lo que implica que la restricción $X'\to Y$ de $f$ es un morfismo cuasifinito \sref[II]{II.6.2.2}.

\begin{corollary}[14.5.4]
\label{IV.14.5.4}
Bajo las hipótesis de \sref{IV.14.5.3}, supongamos además que $Y=\Spec(A)$, donde $A$ es un anillo local noetheriano, íntegro y completo, y que $y=f(x)$ es el único punto cerrado de $Y$.
Entonces existe un anillo local íntegro $A'$, que contiene a $A$ y es una $A$-álgebra finita, tal que, si ponemos $Y'=\Spec(A')$ y $X'=X\times_Y Y'$, existe una $Y'$-sección $h:Y'\to X'$ cuyo morfismo compuesto $Y'\xrightarrow{h}X'\to X$ es una inmersión cuya imagen contiene a $x$.
\end{corollary}

Reemplazando si es necesario $X$ por un subpreesquema cerrado irreducible reducido de $X$, podemos, en virtud de \sref{IV.14.5.3}, limitarnos al caso en que el morfismo $f$ ya es cuasifinito y dominante.
Usando \sref[II]{II.6.2.5}, deducimos que $\sh{O}_x=A'$ es un anillo íntegro que es una $A$-álgebra finita, y que $X$ es una suma disjunta del subpreesquema cerrado $Y'=\Spec(A')$ y de un subpreesquema $Y''$; el esquema $Y'$ satisface lo pedido, pues el morfismo compuesto $Y'\to X'\to X$ no es otro que el morfismo canónico $\Spec(\sh{O}_x)\to X$.

\begin{remark}[14.5.5]
\label{IV.14.5.5}
Si, en el enunciado de \sref{IV.14.5.4}, no se exige que el morfismo $Y'\to X$ sea una inmersión, se puede suponer además que $A'$ es \emph{íntegramente cerrado}: en efecto, basta reemplazar $A'$ por su \emph{cierre entero} $A_1$, pues sabemos por \sref[0]{0.23.1.5} que $A_1$ es un $A$-módulo de tipo finito.
\end{remark}

\begin{proposition}[14.5.6]
\label{IV.14.5.6}
Sea $Y$ un preesquema localmente noetheriano, irreducible, regular y de dimensión $1$, $f:X\to Y$ un morfismo localmente de tipo finito, y $y$ un punto de $Y$.
Las condiciones siguientes son equivalentes:
\begin{enumerate}[label=\emph{\alph*)}]
  \item $f_{\mathrm{red}}$ es plano en todo punto de $f^{-1}(y)$.
\oldpage[IV-3]{218}
  \item $f$ es universalmente abierto en un entorno de $f^{-1}(y)$.
  \item $f$ es abierto en un entorno de $f^{-1}(y)$.
  \item Toda componente irreducible de $X$ que corta a $f^{-1}(y)$ domina a $Y$.
  \item Para todo punto $x\in f^{-1}(y)$ cerrado en $f^{-1}(y)$, y toda componente irreducible $X_i$ de $X$ que contiene a $x$, existe un subconjunto irreducible $X'$ de $X$, localmente cerrado en $X$, que contiene a $x$, está contenido en $X_i$, y tal que la restricción $X'\to Y$ de $f$ es un morfismo cuasifinito y dominante.
\end{enumerate}
\end{proposition}

La equivalencia de \emph{a)}, \emph{b)}, \emph{c)} y \emph{d)} ya se ha demostrado en \sref{IV.14.3.8}.
Es claro que \emph{e)} implica \emph{d)}, pues para toda componente irreducible $X_i$ de $X$ que corta a $f^{-1}(y)$ existe en $X_i\cap f^{-1}(y)$ un punto cerrado \sref{IV.5.1.11}.
Finalmente, para demostrar que \emph{c)} implica \emph{e)}, podemos limitarnos al caso en que $f$ es un morfismo de tipo finito; consideremos un punto cerrado $x$ de $f^{-1}(y)$ y sea $X_i$ una componente irreducible de $X$ que contiene a $x$.
Como la restricción $f_i:X_i\to Y$ de $f$ a $X_i$ es un morfismo dominante, de la equivalencia de \emph{c)} y \emph{d)} para $f_i$ se deduce que este morfismo es abierto en los puntos genéricos de $X_i\cap f^{-1}(y)$. Se deduce de \sref{IV.14.2.2} que $f_i$ es equidimensional en el punto $x$, y concluimos con ayuda de \sref{IV.14.5.3}.

\begin{remark}[14.5.7]
\label{IV.14.5.7}
Si, en el enunciado de \sref{IV.14.5.5}, suponemos que $Y=\Spec(A)$, donde $A$ es un anillo de valoración discreta completo, y que $y$ es el punto cerrado de $Y$, podemos suponer además que $X'=\Spec(A')$, donde $A'$ es un anillo de valoración discreta que es una $A$-álgebra finita, como se deduce de la demostración de \sref{IV.14.5.4} y del hecho de que un anillo local íntegro, regular y de dimensión $1$ es un anillo de valoración discreta \sref[II]{II.7.1.6}.
\end{remark}

\begin{proposition}[14.5.8]
\label{IV.14.5.8}
Sea $Y$ un preesquema localmente noetheriano, $f:X\to Y$ un morfismo localmente de tipo finito, y $y$ un punto de $Y$.
Para todo $Y$-preesquema $Y'$, pongamos $X(Y')=\Hom_Y(Y',X)$.
Sea $x$ un punto de $f^{-1}(y)$; para que $f$ sea universalmente abierto en $x$, es necesario y suficiente que se satisfaga la condición siguiente:

Para todo anillo de valoración discreta completo $A$ con cuerpo residual algebraicamente cerrado, todo morfismo $g:Y'=\Spec(A)\to Y$ tal que la imagen por $g$ del punto cerrado $y'$ de $Y'$ sea igual a $y$, y todo elemento $u_0\in X(\Spec(\kappa(y')))$ tal que $u_0(y')=x$, existen un anillo de valoración discreta $B$, un homomorfismo local $A\to B$ que hace de $B$ una $A$-álgebra finita y, si ponemos $Z'=\Spec(B)$, un elemento $u\in X(Z')$ tal que, si $z'$ es el punto cerrado de $Z'$, el diagrama
\[
  \xymatrix{
    \Spec(\kappa(z')) \ar[r] \ar[d]^{\ell} & Z' \ar[d]^{u} \\
    \Spec(\kappa(y')) \ar[r]_-{u_0} & X
  }
\]
es conmutativo.
\end{proposition}

\oldpage[IV-3]{219}
Observemos que, si $A$ y $B$ satisfacen las condiciones del enunciado, $B$ es un anillo de valoración discreta completo (Bourbaki, \emph{Alg.\ comm.}, cap.~III, \S~3, n\textsuperscript{o}~3, prop.~7 y cap.~IV, \S~2, n\textsuperscript{o}~2, cor.~3 de la prop.~9) cuyo cuerpo residual es isomorfo al de $A$ y, por tanto, algebraicamente cerrado. Además, como $u(z')=x$, existe en $X''=X\times_Y Z'$ un punto $x''$ cuyas proyecciones en $X$ y $Z'$ son, respectivamente, $x$ y $z'$, y que es racional sobre $\kappa(z')$; asimismo, como existe una $Z'$-sección $v$ de $X''$ tal que $v(z')=x''$, la imagen por $v$ del punto genérico $s$ de $Z'$ es una generalización $t$ de $x''$ cuya proyección en $Z'$ es $s$. Aplicando \sref{IV.14.3.6}, vemos que la condición del enunciado es \emph{suficiente}.
Demostremos ahora que es \emph{necesaria}.
Pongamos $X'=X\times_Y Y'$, $f'=f_{(Y')}:X'\to Y'$.
Por hipótesis existe un punto $x'\in X'$ situado sobre $x$ y $y'$ y racional sobre $\kappa(y')$ \sref[I]{I.3.3.14}, luego cerrado en $f'^{-1}(y')$.
En virtud de \sref{IV.14.3.6}, existe una componente irreducible $T'$ de $X'$ que contiene a $x'$ y domina a $Y'$; por tanto, por \sref{IV.14.3.8} y \sref{IV.14.3.13}, la restricción $T'\to Y'$ de $f'$ es equidimensional en $x'$.
Se deduce de \sref{IV.14.5.5} que existe una $A$-álgebra finita $B$ que es un dominio local íntegramente cerrado que domina a $A$ (y, por tanto, es un anillo de valoración discreta) y, poniendo $Z'=\Spec(B)$ y $X''=X\times_Y Z'=X'\times_{Y'}Z'$, una $Z'$-sección $v:Z'\to X''$ tal que la imagen del punto cerrado $z'$ de $Z'$ por el morfismo compuesto $Z'\xrightarrow{v}X''\to X'$ es $x'$; el morfismo compuesto $u:Z'\xrightarrow{v}X''\to X$ responde, por tanto, a la cuestión.

\begin{proposition}[14.5.9]
\label{IV.14.5.9}
Sea $Y$ un preesquema irreducible localmente noetheriano, $f:X\to Y$ un morfismo localmente de tipo finito, y sea $y$ un punto geométricamente unirramificado de $Y$.
Las condiciones siguientes son equivalentes:
\begin{enumerate}[label=\emph{\alph*)}]
  \item $f$ es universalmente abierto en todo punto de $f^{-1}(y)$.
  \item $f$ es abierto en todo punto de $f^{-1}(y)$.
  \item Para toda componente irreducible $Z$ de $f^{-1}(y)$, de punto genérico $z$, existe una componente irreducible $X_i$ de $X$ que contiene a $z$ y es equidimensional sobre $Y$ en $z$.
  \item Para todo punto cerrado $x$ de $f^{-1}(y)$, existe un subconjunto irreducible $X'$ de $X$, localmente cerrado en $X$, que contiene a $x$, y tal que la restricción $X'\to Y$ de $f$ es un morfismo cuasifinito dominante.
\end{enumerate}
\end{proposition}

La equivalencia de \emph{a)}, \emph{b)} y \emph{c)} ya se ha demostrado en \sref{IV.14.4.2}.
Para demostrar que \emph{a)} implica \emph{d)}, observemos que, por \sref{IV.14.4.2}, \emph{a)} implica que existe una componente irreducible $Z$ de $X$ que contiene a $x$ y es equidimensional sobre $Y$ en el punto $x$; la existencia de $X'$ resulta entonces de \sref{IV.14.5.3}, aplicado a la restricción $Z\to Y$ de $f$.
Recíprocamente, supongamos que se satisface \emph{d)}; por el criterio de Chevalley \sref{IV.14.4.4}, la restricción $X'\to Y$ de $f$ es universalmente abierta en $x$ y, \emph{a fortiori}, $f$ es abierto en $x$. Por tanto, $f$ es abierto en todo punto cerrado de $X_y=f^{-1}(y)$.
Pero $X_y$ es un $\kappa(y)$-preesquema localmente de tipo finito y, por tanto, un preesquema de Jacobson; el conjunto de puntos cerrados de $X_y$ es, pues, \emph{denso} en $X_y$ \sref{IV.10.3.1}, y se sigue de \sref{IV.14.1.4} que $f$ es abierto en \emph{todos} los puntos de $X_y$, lo que termina de demostrar que \emph{d)} implica \emph{b)}.

El resultado siguiente fue puesto de manifiesto por D.\ Mumford:

\begin{proposition}[14.5.10]
\label{IV.14.5.10}
\oldpage[IV-3]{220}
Sea $Y$ un preesquema noetheriano y $f:X\to Y$ un morfismo universalmente abierto, sobreyectivo y localmente de tipo finito.
Existe entonces un morfismo finito sobreyectivo $g:Y'\to Y$ tal que, si ponemos $X'=X\times_Y Y'$ y $f'=f_{(Y')}:X'\to Y'$, todo punto $y'\in Y'$ admite un entorno abierto $U'$ tal que existe una $U'$-sección de $f'^{-1}(U')$.
\end{proposition}

Demostraremos la proposición en varias etapas.

\emph{I) Reducción al caso en que $Y$ es íntegro.}
Si hemos demostrado la proposición para cada uno de los subpreesquemas reducidos $Y_i$ cuyo espacio subyacente es una componente irreducible de $Y$, y para las imágenes inversas $f^{-1}(Y_i)$, es claro que el preesquema $Y'$ que es la \emph{suma} de los $Y'_i$ correspondientes satisface lo requerido.
Podemos suponer, por tanto, que $Y$ es \emph{íntegro}, y nos limitaremos a este caso en lo que sigue.
Además, en la conclusión podremos tomar también $Y'$ íntegro (tras reemplazarlo, si es necesario, por una componente irreducible adecuada).

\emph{II) Carácter local sobre $Y$.}
Mostraremos que, si $Y$ puede recubrirse por un número finito de abiertos $U_j$ tales que, para todo $j$, la conclusión de la proposición se cumple para el morfismo $f^{-1}(U_j)\to U_j$, restricción de $f$, entonces se cumple también para $f$.
En efecto, podemos suponer evidentemente que los $U_j$ son afines, de modo que $U_j=\Spec(A_j)$, donde $A_j$ es un anillo noetheriano íntegro cuyo cuerpo de fracciones $K=R(Y)$ es el cuerpo de funciones racionales sobre $Y$.
Para todo $j$, existe por hipótesis una $A_j$-álgebra finita e íntegra $A'_j$ tal que el homomorfismo $A_j\to A'_j$ es inyectivo \sref[I]{I.1.2.7} y que el morfismo correspondiente $g_j:U'_j=\Spec(A'_j)\to\Spec(A_j)=U_j$ satisface las condiciones del enunciado (para $U_j$ y $f^{-1}(U_j)$).
Sea $K'$ una extensión finita de $K$ que contiene los cuerpos de fracciones de todas las $A'_j$ (que son extensiones finitas de $K$).
Consideremos la normalización $Y''$ de $Y$ en $K'$ \sref[II]{II.6.3.8}; es de la forma $\Spec(\sh{B})$, donde $\sh{B}$ es una $\sh{O}_Y$-álgebra entera cuasicoherente, cierre entero de $\sh{O}_Y$ en $K'$ \sref[II]{II.6.3.4}.
Estas definiciones muestran que, para todo $j$,
\[
 A'_j=\Gamma(U'_j,\sh{O}_{U'_j})
     =\Gamma\bigl(U_j,(g_j)_*(\sh{O}_{U'_j})\bigr)
\]
se identifica con una sub-$A_j$-álgebra finita de $\Gamma(U_j,\sh{B})$; dicho de otro modo,
\[
 (g_j)_*(\sh{O}_{U'_j})=\sh{A}'_j
\]
es una $\sh{O}_{U_j}$-álgebra coherente, subálgebra de $\sh{B}|U_j$.
Existe, por consiguiente, un sub-$\sh{O}_Y$-módulo coherente $\sh{C}'_j$ de $\sh{B}$ tal que $\sh{C}'_j|U_j=\sh{A}'_j$ \sref[I]{I.9.4.7}.
Si ponemos $\sh{C}'=\sum_j\sh{C}'_j$, la sub-$\sh{O}_Y$-álgebra $\sh{B}'$ de $\sh{B}$ generada por $\sh{C}'$ es coherente, pues $\sh{B}$ es una $\sh{O}_Y$-álgebra entera.
Mostremos ahora que $Y'=\Spec(\sh{B}')$ satisface lo requerido.
En efecto, es claro que el morfismo $g:Y'\to Y$ es finito y sobreyectivo y que $Y'$ es íntegro; además, para todo $j$, $\Gamma(U_j,\sh{B}')$ es una $A'_j$-álgebra finita.
Dicho de otro modo, la restricción $g^{-1}(U_j)\to U_j$ de $g$ se factoriza en la forma
\[
 g^{-1}(U_j)\longrightarrow U'_j\longrightarrow U_j,
\]
y, como la existencia local de secciones es estable por cambio de base, esto demuestra nuestra afirmación, pues todo $y\in Y$ pertenece a algún $U_j$.

\emph{III) Reducción al caso en que $Y$ es íntegro, local y geométricamente unirramificado.}
Supongamos primero que la proposición se ha demostrado cuando $Y$ es íntegro y \emph{local} (con $Y'$ íntegro), y mostremos que es válida cuando $Y$ es afín, íntegro y noetheriano.
En efecto, en virtud de la reducción II), basta demostrar que, para todo punto $y\in Y$, la proposición se satisface para un entorno abierto afín $V$ de $y$ en $Y$.
Sea $Y=\Spec(A)$, $\mathfrak{p}=\mathfrak{j}_y$, sea $Y_1=\Spec(A_\mathfrak{p})$, y pongamos $X_1=X\times_Y Y_1$.
Por hipótesis, existe un morfismo finito sobreyectivo $g_1:Y'_1\to Y_1$, donde $Y'_1=\Spec(B_1)$ y $B_1$ es una $A_\mathfrak{p}$-álgebra finita e íntegra, por tanto un anillo semilocal, tal que $g_1$ satisface las condiciones del enunciado para $Y_1$ y $X_1$.
Si $y'_j$ ($1\leq j\leq r$) son los puntos cerrados de $Y'_1$, existe entonces un recubrimiento de $Y'_1$ por abiertos $U'_j$ tales que $y'_j\in U'_j$ y existe una $U'_j$-sección $h'_j$ de $X\times_Y U'_j$ ($1\leq j\leq r$).
El $A_\mathfrak{p}$-módulo $B_1$ admite un sistema finito de generadores de la forma $z_i/s$ (con $s\in A-\mathfrak{p}$,
\oldpage[IV-3]{221}
$z_i$ enteros sobre $A$), y puede suponerse (multiplicando $s$, si es necesario, por un elemento de $A$) que estos generadores son elementos del cuerpo de fracciones de $B_1$ enteros sobre $A_s$, de modo que, si $V$ es el abierto afín $D(s)=\Spec(A_s)\subset Y$, $Y'_1$ se identifica con $Y_1\times_V V'$, donde $V'$ es el espectro de la $A_s$-álgebra finita generada por los $z_i/s$; $g:V'\to V$ es, por tanto, un morfismo finito sobreyectivo y $g_1=g_{(Y_1)}$.
Además, aplicando el método de \sref{IV.8.1.2},~$a$), puede suponerse que cada $U'_j$ es la imagen inversa por $p:Y'_1\to V'$ de un abierto $W'_j$ de $V'$, que los $W'_j$ recubren $V'$ \sref{IV.8.3.11}, y que cada sección $h'_j$ es de la forma $(v'_j)_{(Y_1)}$, donde $v'_j$ es una $W'_j$-sección de $X\times_Y W'_j$ \sref{IV.8.8.2},~(i).
Estamos, pues, efectivamente reducidos a demostrar la proposición cuando $Y=\Spec(A)$, donde $A$ es un anillo noetheriano local íntegro.
Sabemos entonces que existe una $A$-álgebra finita e íntegra $B$, con el mismo cuerpo de fracciones que $A$, tal que $A\subset B$ y $\Spec(B)$ es geométricamente unirramificado \sref[0]{0.23.2.5} y \sref{IV.6.15.5}; como el morfismo $\Spec(B)\to\Spec(A)$ es sobreyectivo, puede reemplazarse evidentemente $Y$ por $\Spec(B)$ y $X$ por $X\otimes_A B$ para demostrar la proposición.
Razonando como al comienzo de la reducción III), podemos suponer, por tanto, que $A$ es local e íntegro y que $Y=\Spec(A)$ es geométricamente unirramificado.

\emph{IV) Reducción al caso en que $X$ es íntegro y afín, y $f$ es cuasifinito, sobreyectivo, birracional y universalmente abierto.}
Supongamos, pues, $Y=\Spec(A)$ íntegro, local y geométricamente unirramificado.
Existe entonces un subpreesquema irreducible $X_0$ de $X$ tal que la restricción $f_0:X_0\to Y$ de $f$ es un morfismo cuasifinito y \emph{dominante}, y $f_0(X_0)$ contiene al punto cerrado $y$ de $Y$ \sref{IV.14.5.9}; como $Y$ es geométricamente unirramificado, se sigue de \sref{IV.14.4.1} que $f_0$ sigue siendo \emph{universalmente abierto}.
Además, podemos suponer $X_0$ reducido, por tanto íntegro, y vemos que, reemplazando $X$ por $X_0$, podemos suponer $X$ íntegro, $f$ \emph{cuasifinito y dominante}, y que $f(X)$ contiene a $y$; como $f$ es \emph{abierto} y todo abierto de $Y$ que contiene a $y$ es igual a $Y$, $f$ es sobreyectivo.

Sean $\xi$, $\eta$ los puntos genéricos de $X$ e $Y$, respectivamente; $\kappa(\xi)$ es, por tanto, una extensión finita de $K=\kappa(\eta)$.
Por \sref{IV.4.6.8}, existe una extensión \emph{finita} $K'$ de $K$ tal que $(\Spec(\kappa(\xi)\otimes_K K'))_{\mathrm{red}}$ es geométricamente reducido sobre $K'$\footnote{La fuente impresa escribe aquí $K$; el cambio de base y la conclusión siguiente exigen $K'$.} y sus componentes irreducibles son geométricamente irreducibles; se sigue de \sref{IV.4.5.9} y \sref{IV.4.6.1} que los cuerpos residuales de $\Spec(\kappa(\xi)\otimes_K K')$ son extensiones finitas de $K'$ primarias y separables, y son, por tanto, iguales a $K'$.
Aplicando una vez más \sref[0]{0.23.2.5} y \sref{IV.6.15.5}, existe una $A$-álgebra finita e íntegra $B$, que tiene a $K'$ por cuerpo de fracciones, contiene a $A$ y es tal que $\Spec(B)$ es geométricamente unirramificado; como el morfismo $\Spec(B)\to\Spec(A)$ es sobreyectivo, podemos reemplazar de nuevo $Y$ por $\Spec(B)$ y $X$ por $(X\otimes_A B)_{\mathrm{red}}$ para demostrar la proposición; la reducción III) permite entonces suponer una vez más que $A$ es local e íntegro y que $Y=\Spec(A)$ es geométricamente unirramificado, de punto cerrado $y$.
Además, $f$ sigue siendo un morfismo cuasifinito, \emph{sobreyectivo} y \emph{universalmente abierto}; cada una de las componentes irreducibles $X_i$ de $X$ ($1\leq i\leq m$) domina a $Y$ \sref{IV.1.10.4} y es \emph{birracional} sobre $Y$.
Sea, pues, $x$ un punto de $f^{-1}(y)$, y sea $X_i$ una componente irreducible de $X$ que contiene a $x$; denotemos también por $X_i$ el subpreesquema reducido (por tanto, íntegro) de $X$ cuyo espacio subyacente es $X_i$, y sea $f_i:X_i\to Y$ la restricción de $f$.
Aplicando una vez más \sref{IV.14.4.1} al morfismo cuasifinito y dominante $f_i$, vemos que $f_i$ es universalmente abierto; si $U$ es un abierto afín de $X_i$ que contiene a $x$, $f_i(U)$ es, por tanto,
\oldpage[IV-3]{222}
abierto en $Y$ y contiene a $y$, luego es igual a $Y$.
Podemos, pues, reemplazar $X$ por $U$ para demostrar la proposición.

\emph{V) Fin de la demostración.}
Suponemos, pues, $Y$ local, íntegro y geométricamente unirramificado, $X$ íntegro y afín, y $f$ cuasifinito, sobreyectivo, birracional y universalmente abierto; se sigue que $f$ es automáticamente \emph{separado}.
En virtud del teorema principal \sref{IV.8.12.6}, existe una factorización $X\xrightarrow{j}Z\xrightarrow{u}Y$, donde $j$ es una inmersión abierta y $u$ un morfismo \emph{finito}.
Reemplazando además $Z$ por la imagen cerrada de $j$ \sref[I]{I.9.5.10}, podemos suponer que $Z$ es \emph{íntegro}; como $u$ es finito y birracional e $Y$ geométricamente unirramificado, se sigue de \sref[III]{III.4.3.5} y \sref[III]{III.4.3.4} que $u$ (y, por consiguiente, $f$) es un morfismo \emph{radicial}; como $f$ es universalmente abierto y sobreyectivo, es, por tanto, un \emph{homeomorfismo} universal.
En consecuencia, por \sref{IV.2.4.5},~(ii), $f$ es un morfismo \emph{finito}.
Pero entonces se satisfacen las condiciones del enunciado con $Y'$ \emph{íntegro} tomando simplemente $Y'=X$ y $g=f$, siendo la $Y'$-sección el morfismo diagonal $\Delta_f$.

Este resultado puede utilizarse para desarrollar criterios de ``descenso'' para las propiedades de los morfismos universalmente abiertos y sobreyectivos.
Señalemos en particular el criterio siguiente, debido a D.\ Mumford:

\begin{corollary}[14.5.11]
\label{IV.14.5.11}
Sea $Y$ un preesquema localmente noetheriano, $f:X\to Y$ un morfismo localmente de tipo finito y $g:Y'\to Y$ un morfismo localmente de tipo finito, universalmente abierto y sobreyectivo; pongamos $X'=X\times_Y Y'$, $f'=f_{(Y')}:X'\to Y'$.
Para que $f$ sea afín, es necesario y suficiente que $f'$ lo sea.
\end{corollary}

Basta demostrar que la condición es suficiente.
Como la cuestión es local sobre $Y$, podemos suponer que $Y$ es afín, por tanto noetheriano.
En virtud de \sref{IV.14.5.10}, existe un morfismo finito sobreyectivo $h:Y_1\to Y$ tal que, si ponemos $Y'_1=Y'\times_Y Y_1$ y $g'=g_{(Y_1)}:Y'_1\to Y_1$, todo punto de $Y_1$ admite un entorno abierto $U_1$ tal que existe una $U_1$-sección de $g'^{-1}(U_1)$.
Si ponemos $X_1=X\times_Y Y_1$, la proyección canónica $p:X_1\to X$ es un morfismo finito sobreyectivo; por consiguiente, si demostramos que $X_1$ es un esquema afín, se seguirá ante todo que $X$ es cuasicompacto, luego noetheriano, y después que $X$ es afín en virtud del teorema de Chevalley \sref[II]{II.6.7.1}.
Basta, pues, demostrar que el morfismo $f_1=f_{(Y_1)}:X_1\to Y_1$ es afín.
Ahora bien, si ponemos
\[
  X'_1=X_1\times_{Y_1}Y'_1=X'\times_{Y'}Y'_1 \qquad\text{y}\qquad f'_1=(f_1)_{(Y'_1)}=f'_{(Y'_1)}:X'_1\to Y'_1,
\]
$f'_1$ es afín en virtud de la hipótesis.
Estamos, pues, reducidos a demostrar el corolario reemplazando $Y$, $X$, $f$ e $Y'$ por $Y_1$, $X_1$, $f_1$ e $Y'_1$; dicho de otro modo, basta demostrar que $f$ es afín cuando, en el enunciado de \sref{IV.14.5.11}, se supone además que todo punto de $Y$ admite un entorno abierto $U$ tal que existe una $U$-sección de $g^{-1}(U)$.
Como la cuestión es local sobre $Y$, podemos suponer incluso que existe una $Y$-sección $s$ de $Y'$.
Disponemos ahora del lema elemental siguiente (válido en toda categoría que admita productos fibrados):

\begin{lemma}[14.5.11.1]
\label{IV.14.5.11.1}
Sean $f:X\to S$, $g:Y\to S$ dos morfismos, y $p_1:X\times_S Y\to X$, $p_2:X\times_S Y\to Y$ las proyecciones canónicas.
Si $s:S\to Y$ es una sección de $g$, entonces $s'=(1,s\circ f)_S$ es una sección de $p_1$, y $X$, provisto de los morfismos $f:X\to S$ y $s':X\to X\times_S Y$, se identifica con el producto de los $Y$-preesquemas $S$ y $X\times_S Y$ para los morfismos $s:S\to Y$ y $p_2:X\times_S Y\to Y$.
\end{lemma}

\begin{proof}
Este es un caso particular de \sref[I]{I.3.3.11}, donde se reemplaza el diagrama por
\[
\xymatrix{
  X \ar[r]^-{s'} \ar[d]_f & X\times_S Y \ar[r]^-{p_1} \ar[d]^{p_2} & X \ar[d]^f \\
  S \ar[r]_-{s} & Y \ar[r]_-{g} & S
}
\]
\end{proof}

Aplicando este lema con $S$ e $Y$ reemplazados, respectivamente, por $Y$ e $Y'$, se ve que puede escribirse $f=(f')_{(Y)}$ para el cambio de base $s:Y\to Y'$; por tanto, $f$ es afín, puesto que $f'$ lo es.
Q.E.D.

Una variante de este criterio es la siguiente:

\begin{corollary}[14.5.12]
\label{IV.14.5.12}
\oldpage[IV-3]{223}
Con las notaciones e hipótesis generales de \sref{IV.14.5.11}, sea $\sh{L}$ un $\sh{O}_X$-módulo invertible.
Supongamos que existe un subconjunto cerrado $Z$ de $X$, propio sobre $Y$ \sref[II]{II.5.4.10}, y tal que el preesquema inducido por $X$ en el abierto $X-Z$ es normal.
Para que $\sh{L}$ sea amplio con respecto a $f$, es necesario y suficiente que $\sh{L}'=\sh{L}\otimes_{\sh{O}_Y}\sh{O}_{Y'}$ sea amplio con respecto a $f'$.
\end{corollary}

\begin{proof}
Conservando las notaciones de la demostración de \sref{IV.14.5.11}, pongamos $\sh{L}_1=\sh{L}\otimes_{\sh{O}_Y}\sh{O}_{Y_1}$.
Por \sref[III]{III.2.6.2}, bastará demostrar que $\sh{L}_1$ es amplio con respecto a $f\circ p$; pero $f\circ p=h\circ f_1$, y, a la vista de \sref[II]{II.4.6.13},~(v), bastará demostrar que $\sh{L}_1$ es amplio con respecto a $f_1$.
Como la cuestión es local sobre $Y_1$, podemos suponer de nuevo que $g'$ admite una sección $s$.
Si $\sh{L}'_1=\sh{L}'\otimes_{\sh{O}_{Y'}}\sh{O}_{Y'_1}$, puede escribirse entonces, en virtud del lema \sref{IV.14.5.11.1}, $\sh{L}_1=\sh{L}'_1\otimes_{\sh{O}_{Y'_1}}\sh{O}_{Y_1}$ para el cambio de base $s:Y_1\to Y'_1$; la conclusión resulta, por tanto, de dos aplicaciones de \sref[II]{II.4.6.13},~(iii).
\end{proof}

\section{Estudio de las fibras de un morfismo universalmente abierto}
\label{section:IV.15}

\subsection{Multiplicidades de las fibras de un morfismo universalmente abierto}
\label{subsection:IV.15.1}

\begin{proposition}[15.1.1]
\label{IV.15.1.1}
\oldpage[IV-3]{223}
Sea $\mathrm{Y}$ un preesquema irreducible localmente noetheriano, de punto genérico $\eta$, sea $f:\mathrm{X}\to\mathrm{Y}$ un morfismo localmente de tipo finito, sea $y$ un punto de $\mathrm{Y}$ y sea $\sh{F}$ un $\sh{O}_\mathrm{X}$-módulo coherente;
ponemos $\sh{F}_y=\sh{F}\otimes_{\sh{O}_\mathrm{Y}}\mathbf{k}(y)$ (haz de módulos sobre la fibra $f^{-1}(y)$).
Sea $z$ el punto genérico de una componente irreducible de $\operatorname{Supp}(\sh{F}_y)$, y sean $\mathrm{X}_i$ ($1\leqslant i\leqslant n$) aquellos subpreesquemas cerrados reducidos de $\mathrm{X}$ cuyo espacio subyacente es una componente irreducible de $\operatorname{Supp}(\sh{F})$ que contiene a $z$, y tales que la restricción $\mathrm{X}_i\to\mathrm{Y}$ de $f$ es un morfismo universalmente abierto en el punto $z$; sea $z_i$ el punto genérico de $\mathrm{X}_i$ ($1\leqslant i\leqslant n$).
Si $\lambda_x(\sh{F}_y)$ (resp.~$\lambda_t(\sh{F}_\eta)$) denota la longitud geométrica de $\sh{F}_y$ en un punto $x\in f^{-1}(y)$ (resp.~la de $\sh{F}_\eta$ en un punto $t\in f^{-1}(\eta)$) \sref{IV.4.7.5}, entonces se tiene
\[
\tag{15.1.1.1}
\label{IV.15.1.1.1}
  \lambda_z(\sh{F}_y)\geqslant\sum_{i=1}^{n}\lambda_{z_i}(\sh{F}_\eta).
\]
Si además se supone que el anillo $\sh{O}_y$ es regular, se tiene
\[
\tag{15.1.1.2}
\label{IV.15.1.1.2}
  \operatorname{length}((\sh{F}_y)_z)\geqslant\sum_{i=1}^{n}\operatorname{length}((\sh{F}_\eta)_{z_i}).
\]
\end{proposition}

\begin{proof}
\label{IV.15.1.1-proof}
\oldpage[IV-3]{224}
Las fibras $f^{-1}(y')$ de $f$ (para todo $y'\in\mathrm{Y}$) no cambian cuando se sustituye $f$ por $f_\mathrm{red}$, de modo que se puede suponer que $\mathrm{X}$ e $\mathrm{Y}$ son reducidos \sref{IV.2.4.3},~(vi), y por tanto que $\mathrm{Y}$ es íntegro; pondremos $\mathrm{K}=\mathbf{k}(\eta)$.
Además, se puede sustituir $f$ por la restricción $\operatorname{Supp}(\sh{F})\to\mathrm{Y}$ de $f$ al subpreesquema cerrado reducido de $\mathrm{X}$ cuyo espacio subyacente es $\operatorname{Supp}(\sh{F})$, y limitarse así al caso en que $\mathrm{X}=\operatorname{Supp}(\sh{F})$.
Por último, si $y=\eta$, hay una sola componente irreducible $\mathrm{X}_i$ de $\mathrm{X}$ que contiene a $z$ \sref[0]{0.2.1.8}, y se tiene $z_i=z$, lo que hace triviales \sref{IV.15.1.1.1} y \sref{IV.15.1.1.2} (sin hipótesis sobre $f$).
Podemos, pues, limitarnos al caso $y\neq\eta$; se deduce entonces de la hipótesis y de \sref{IV.1.10.3} que los $z_i$ pertenecen a $f^{-1}(\eta)$, de modo que los segundos miembros de \sref{IV.15.1.1.1} y \sref{IV.15.1.1.2} están bien definidos.
Sea $\mathrm{Z}_i$ el subpreesquema cerrado reducido de $f^{-1}(\eta)$ cuyo espacio subyacente es $\mathrm{X}_i\cap f^{-1}(\eta)$; sabemos por \sref{IV.4.6.6} que existe una extensión radicial finita $\mathrm{K}'$ de $\mathrm{K}$ tal que, para todo $i$, el $\mathrm{K}'$-preesquema $(\mathrm{Z}_i\otimes_\mathrm{K}\mathrm{K}')_\mathrm{red}$ es geométricamente reducido sobre $\mathrm{K}'$.
Además, el morfismo de proyección $f^{-1}(\eta)\otimes_\mathrm{K}\mathrm{K}'\to f^{-1}(\eta)$ es finito, dominante y radicial \sref[I]{I.3.5.7}, y por tanto es un homeomorfismo \sref{IV.2.4.5}.
Si $z_i'$ denota el único punto de $f^{-1}(\eta)\otimes_\mathrm{K}\mathrm{K}'$ situado sobre $z_i$, se deduce de \sref{IV.4.7.9} y \sref{IV.4.7.5} que se tiene
\[
\tag{15.1.1.3}
\label{IV.15.1.1.3}
  \lambda_{z_i}(\sh{F}_\eta)=\lambda_{z_i'}(\sh{F}_\eta\otimes_\mathrm{K}\mathrm{K}')=\operatorname{length}((\sh{F}_\eta\otimes_\mathrm{K}\mathrm{K}')_{z_i'}).
\]

Sea $\mathrm{V}$ un anillo de valoración discreta, de cuerpo de fracciones $\mathrm{K}'$, que domina a $\sh{O}_y$ \sref[II]{II.7.1.7}; pongamos $\mathrm{Y}'=\operatorname{Spec}(\mathrm{V})$, $\mathrm{X}'=\mathrm{X}\times_\mathrm{Y}\mathrm{Y}'$, $\sh{F}'=\sh{F}\otimes_\mathrm{Y}\mathrm{Y}'$, $f'=f_{(\mathrm{Y}')}:\mathrm{X}'\to\mathrm{Y}'$; si $g:\mathrm{Y}'\to\mathrm{Y}$ es el morfismo estructural, $\eta'$ el punto genérico de $\mathrm{Y}'$ e $y'$ su punto cerrado, se tiene $g(y')=y$, $g(\eta')=\eta$; como el morfismo $\operatorname{Spec}(\mathbf{k}(y'))\to\operatorname{Spec}(\mathbf{k}(y))$ es fielmente plano, lo mismo ocurre con la proyección $g':f'^{-1}(y')\to f^{-1}(y)$ \sref{IV.2.2.13}, y por tanto \sref{IV.2.3.4} existe un punto genérico $z'$ de una componente irreducible de $f'^{-1}(y')$ tal que $g'(z')=z$.
Por construcción, se tiene $\mathbf{k}(\eta')=\mathrm{K}'$, y por tanto $f'^{-1}(\eta')=f^{-1}(\eta)\otimes_\mathrm{K}\mathrm{K}'$; por otra parte, si se pone $\mathrm{X}_i'=\mathrm{X}_i\times_\mathrm{Y}\mathrm{Y}'$, la restricción $\mathrm{X}_i'\to\mathrm{Y}'$ de $f'$ es un morfismo universalmente abierto en el punto $z'$, y por consiguiente \sref{IV.1.10.3} existe un punto $z_i''\in f'^{-1}(\eta')$ que es una generalización de $z'$; evidentemente, se puede suponer que $z_i''$ es el punto genérico de una componente irreducible de $\mathrm{X}_i'\cap f'^{-1}(\eta')$; como la proyección de $\mathrm{X}_i'\cap f'^{-1}(\eta')$ en $f^{-1}(\eta)$ es $\mathrm{Z}_i$, se tiene $z_i''=z_i'$.
En virtud de \sref{IV.4.7.9} y \sref{IV.4.7.5}, se tiene
\[
\tag{15.1.1.4}
\label{IV.15.1.1.4}
  \lambda_z(\sh{F}_y)=\lambda_{z'}(\sh{F}'_{y'})\geqslant\operatorname{length}((\sh{F}'_{y'})_{z'})
\]
y de esta desigualdad y de \sref{IV.15.1.1.3} se deduce que, para establecer \sref{IV.15.1.1.1}, basta probar la desigualdad
\[
\tag{15.1.1.5}
\label{IV.15.1.1.5}
  \operatorname{length}((\sh{F}_{y'}')_{z'})\geqslant\sum_{i=1}^{n}\operatorname{length}((\sh{F}_{\eta'}')_{z_i'}).
\]

Consideremos ahora la segunda desigualdad \sref{IV.15.1.1.2}; utilizaremos el lema siguiente:

\begin{lemma}[15.1.1.6]
\label{IV.15.1.1.6}
Sea $\mathrm{A}$ un anillo local regular que no es un cuerpo, sea $\mathrm{K}$ su cuerpo de fracciones y sea $k$ su cuerpo residual.
Existe un anillo de valoración discreta $\mathrm{V}$ que domina a $\mathrm{A}$, tiene a $\mathrm{K}$ como cuerpo de fracciones y cuyo cuerpo residual es una extensión puramente trascendente de $k$.
\end{lemma}

\oldpage[IV-3]{225}
Pongamos $\mathrm{S}=\operatorname{Spec}(\mathrm{A})$ y consideremos el $\mathrm{S}$-esquema $\mathrm{P}$ obtenido haciendo explotar el punto cerrado $s$ de $\mathrm{S}$ \sref[II]{II.8.1.3}; si $\mathfrak{m}$ es el ideal maximal de $\mathrm{A}$, se tiene por definición $\mathrm{P}=\operatorname{Proj}(\mathrm{B})$, donde $\mathrm{B}$ es la $\mathrm{A}$-álgebra graduada $\bigoplus_{n\geqslant 0}\mathfrak{m}^n$.
Si $h:\mathrm{P}\to\mathrm{S}$ es el morfismo estructural, la fibra $h^{-1}(s)$ es isomorfa a $\operatorname{Proj}(\mathrm{B}\otimes_\mathrm{A} k)$ \sref[II]{II.2.8.10}, y por definición $\mathrm{B}\otimes_\mathrm{A} k$ es la $k$-álgebra graduada $\bigoplus_{n\geqslant 0}\mathfrak{m}^n/\mathfrak{m}^{n+1}$; pero, como $\mathrm{A}$ es regular, esta álgebra es isomorfa a un álgebra de polinomios $k[t_1,\ldots,t_e]$ ($t_i$ indeterminadas) \sref[0]{0.17.1.1}, y se deduce que $h^{-1}(s)$ es isomorfa a $\mathbf{P}_k^{e-1}$.
Sea $\zeta$ el punto genérico de $h^{-1}(s)$; si $t_i'$ ($1\leqslant i\leqslant e$) son los elementos de $\mathfrak{m}$ cuyas clases $\bmod\mathfrak{m}^2$ son los $t_i$, $\mathrm{B}_{(t_e')}$ es el anillo de un entorno abierto afín de $\zeta$ en $\mathrm{P}$, y se tiene $\mathbf{k}(\zeta)=k(t_1,\ldots,t_{e-1})$; por otra parte, como $\mathrm{A}$ es íntegramente cerrado, se verifica fácilmente que $\mathrm{B}$ también lo es (Bourbaki, \emph{Alg.\ comm.}, cap.~V, \S\,1, n.º~8, cor.~1 de la prop.~21), de modo que $\mathrm{B}_{(t_e')}$ también es íntegramente cerrado y, por tanto, lo mismo ocurre con el anillo local $\sh{O}_{\mathrm{P},\zeta}$.
Por último, como $\mathrm{A}=\sh{O}_s$ es regular, y por tanto un anillo universalmente catenario \sref{IV.5.6.4}, se tiene, en virtud de \sref{IV.5.6.5}, $\dim(\sh{O}_{\mathrm{P},\zeta})=\dim(\sh{O}_s)-(e-1)$, puesto que $h$ es un morfismo birracional, $\dim(h^{-1}(s))=e-1$, y el anillo local de la fibra $h^{-1}(s)$ en su punto genérico $\zeta$ es de dimensión $0$.
Pero $\dim(\sh{O}_s)=e$, luego $\dim(\sh{O}_{\mathrm{P},\zeta})=1$, y $\sh{O}_{\mathrm{P},\zeta}=\mathrm{V}$ es un anillo de valoración discreta, noetheriano, de dimensión $1$ e íntegramente cerrado \sref[II]{II.7.1.6}; cumple claramente lo pedido.

Establecido este lema, se puede retomar el razonamiento para la desigualdad \sref{IV.15.1.1.1} con $\mathrm{Y}'=\operatorname{Spec}(\mathrm{V})$; esta vez $\mathbf{k}(y')$ es una extensión \emph{separable} de $\mathbf{k}(y)$, y por \sref{IV.4.7.9} se tiene $\operatorname{length}((\sh{F}_y)_z)=\operatorname{length}((\sh{F}_{y'}')_{z'})$; como además $f'^{-1}(\eta')=f^{-1}(\eta)$ y $\sh{F}_{\eta'}'=\sh{F}_\eta$, se queda de nuevo reducido a probar la desigualdad \sref{IV.15.1.1.5}.
Ahora bien, si $t$ es un uniformizante de $\sh{O}_{y'}$, $f'^{-1}(y')$ es el subpreesquema cerrado de $\mathrm{X}'$ definido por el ideal $t\sh{O}_{\mathrm{X}'}$; por otra parte, se tiene $\sh{F}_{y'}'=\sh{F}'/t\sh{F}'$; se verifica asimismo inmediatamente \sref[I]{I.9.1.12} que $(\sh{F}_\eta')_{z_i'}=\sh{F}_{z_i'}'$; la desigualdad por probar \sref{IV.15.1.1.5} no es, pues, sino un caso particular de \sref{IV.3.4.1.1}.
\end{proof}

\begin{corollary}[15.1.2]
\label{IV.15.1.2}
Siendo las hipótesis las de \sref{IV.15.1.1}, supongamos además que $\lambda_z(\sh{F}_y)=1$ (resp.~que $\sh{O}_y$ es regular y $\operatorname{length}((\sh{F}_y)_z)=1$).
Entonces existe a lo sumo una componente irreducible $\mathrm{X}_i$ de $\operatorname{Supp}(\sh{F})$ que contiene a $z$ y tal que la restricción $\mathrm{X}_i\to\mathrm{Y}$ de $f$ es universalmente abierta en el punto $z$; además, si $z_i$ es el punto genérico de esta componente, entonces $\lambda_{z_i}(\sh{F}_\eta)=1$ (resp.~$\operatorname{length}((\sh{F}_\eta)_{z_i})=1$).
\end{corollary}

Esto se deduce inmediatamente de \sref{IV.15.1.1}, observando que, por definición de $\mathrm{X}_i$, se tiene necesariamente $(\sh{F}_\eta)_{z_i}\neq 0$ \sref[I]{I.9.1.13}.

\begin{corollary}[15.1.3]
\label{IV.15.1.3}
Sea $\mathrm{Y}$ un preesquema irreducible localmente noetheriano de punto genérico $\eta$, sea $f:\mathrm{X}\to\mathrm{Y}$ un morfismo localmente de tipo finito, sea $\sh{F}$ un $\sh{O}_\mathrm{X}$-módulo coherente de soporte $\mathrm{X}$, sea $x$ un punto de $\mathrm{X}$ y sea $y=f(x)$.
Supongamos que: $1^\circ$~$x$ pertenece a una única componente irreducible $\mathrm{Z}$ de $f^{-1}(y)$; $2^\circ$~$\sh{F}_y=\sh{F}\otimes_{\sh{O}_\mathrm{Y}}\mathbf{k}(y)$ es geométricamente reducido sobre $\mathbf{k}(y)$, es decir, si $z$ es el punto genérico de $\mathrm{Z}$, $\lambda_z(\sh{F}_y)=1$ (resp.~$\sh{O}_y$ es regular y $\operatorname{length}((\sh{F}_y)_z)=1$).
Entonces existe a lo sumo una componente irreducible $\mathrm{X}'$ de $\mathrm{X}$ que contiene a $x$, tal que $\dim_x(\mathrm{X}'\cap f^{-1}(y))=\dim_x(f^{-1}(y))$ y tal que la restricción $\mathrm{X}'\to\mathrm{Y}$ de $f$ es universalmente abierta en los puntos genéricos de las componentes irreducibles de $\mathrm{X}'\cap f^{-1}(y)$ que contienen a $x$.
Además, si existe tal componente $\mathrm{X}'$ y si $z'$ es su punto genérico, entonces $\lambda_{z'}(\sh{F}_\eta)=1$ (resp.~$\operatorname{length}((\sh{F}_\eta)_{z'})=1$).
\end{corollary}

\oldpage[IV-3]{226}
En efecto, una componente irreducible $\mathrm{X}'$ de $\mathrm{X}$ que contiene a $x$ y tal que $\dim_x(\mathrm{X}'\cap f^{-1}(y))=\dim_x(f^{-1}(y))$ contiene necesariamente a $\mathrm{Z}$, puesto que una componente irreducible de $\mathrm{X}'\cap f^{-1}(y)$ que contiene a $x$ debe tener la misma dimensión que la única componente irreducible $\mathrm{Z}$ de $f^{-1}(y)$ que contiene a $x$.
Se puede aplicar entonces \sref{IV.15.1.2} a $z$.

\begin{remarks}[15.1.4]
\label{IV.15.1.4}
\begin{enumerate}[label=(\roman*)]
  \item En el enunciado de \sref{IV.15.1.1}, no se puede sustituir ``universalmente abierto'' por ``equidimensional'', como muestra el ejemplo \sref{IV.14.4.10},~(ii), donde se toma $\sh{F}=\sh{O}_\mathrm{X}$; las fibras de $f$ son entonces esquemas artinianos reducidos y, por tanto (con las notaciones introducidas \emph{loc.\ cit.}), se tiene $\lambda_{x'}(\sh{F}_y)=1$, pero hay dos componentes irreducibles $\mathrm{X}_1$, $\mathrm{X}_2$ de $\mathrm{X}$ que contienen a $x'$, por lo que el miembro derecho de \sref{IV.15.1.1.1} es igual a $2$, aunque la restricción de $f$ a cada una de las componentes $\mathrm{X}_1$, $\mathrm{X}_2$ es un morfismo equidimensional en todo punto.
Este mismo ejemplo muestra también que en \sref{IV.15.1.2} y \sref{IV.15.1.3} no se puede sustituir la hipótesis ``universalmente abierto'' por ``equidimensional''.
  \item Es verosímil que, para la validez de la desigualdad \sref{IV.15.1.1.2}, no se pueda suprimir la hipótesis de que $\sh{O}_y$ sea \emph{regular}.
\end{enumerate}
\end{remarks}

\subsection{Planitud de los morfismos universalmente abiertos de fibras geométricamente reducidas}
\label{subsection:IV.15.2}

\begin{env}[15.2.1]
\label{IV.15.2.1}
Vimos en \sref{IV.2.4.6} que un morfismo \emph{plano} y localmente de presentación finita es \emph{universalmente abierto}.
Se tiene un recíproco parcial de este resultado bajo hipótesis adicionales:
\end{env}

\begin{theorem}[15.2.2]
\label{IV.15.2.2}
Sea $\mathrm{Y}$ un preesquema localmente noetheriano, sea $f:\mathrm{X}\to\mathrm{Y}$ un morfismo localmente de tipo finito, sea $\sh{F}$ un $\sh{O}_\mathrm{X}$-módulo coherente, sea $x$ un punto de $\mathrm{X}$ y sea $y=f(x)$.
Supongamos satisfechas las condiciones siguientes:
\begin{enumerate}[label=(\roman*)]
  \item $\operatorname{Supp}(\sh{F})=\mathrm{X}$.
  \item $f$ es universalmente abierto en los puntos genéricos de las componentes irreducibles de $f^{-1}(y)$ que contienen a $x$.
  \item $\sh{F}_y=\sh{F}\otimes_{\sh{O}_\mathrm{Y}}\mathbf{k}(y)$ es geométricamente reducido en el punto $x$ \sref{IV.4.6.22}.
  \item El anillo $\sh{O}_y$ es reducido.
\end{enumerate}
Entonces $\sh{F}$ es $f$-plano en el punto $x$.
\end{theorem}

\begin{proof}
\label{IV.15.2.2-proof}
Como $\sh{O}_y$ es reducido y noetheriano, aplicaremos el criterio valuativo de planitud \sref{IV.11.8.1}.
Consideremos, pues, un homomorfismo local $\sh{O}_y\to\mathrm{A}'$, donde $\mathrm{A}'$ es un anillo de valoración discreta; pongamos $\mathrm{Y}'=\operatorname{Spec}(\mathrm{A}')$, $\mathrm{X}'=\mathrm{X}\times_\mathrm{Y}\mathrm{Y}'$, $f'=f_{(\mathrm{Y}')}:\mathrm{X}'\to\mathrm{Y}'$, y sea $x'$ un punto de $\mathrm{X}'$ cuyas proyecciones en $\mathrm{X}$ y $\mathrm{Y}'$ sean respectivamente $x$ y el punto cerrado $y'$ de $\mathrm{Y}'$; si se pone $\sh{F}'=\sh{F}\otimes_\mathrm{Y}\mathrm{Y}'$, entonces $\operatorname{Supp}(\sh{F}')=\mathrm{X}'$ por \sref[I]{I.9.1.13}; todo punto genérico $z'$ de una componente irreducible de $f'^{-1}(y')$ que contiene a $x'$ se proyecta en $f^{-1}(y)$ sobre un punto genérico $z$ de una componente irreducible de $f^{-1}(y)$ que contiene a $x$, en virtud de \sref{IV.4.2.6}; por tanto, $f'$ es universalmente abierto en el punto $z'$.
Finalmente, se deduce de \sref{IV.4.7.11} que $\sh{F}_{y'}'$ es geométricamente reducido en el punto $x'$.
Así se ve que quedamos reducidos a probar el siguiente

\oldpage[IV-3]{227}
\begin{lemma}[15.2.2.1]
\label{IV.15.2.2.1}
Sea $\mathrm{Y}=\operatorname{Spec}(\mathrm{A})$ el espectro de un anillo de valoración discreta, sea $y$ su punto cerrado, sea $f:\mathrm{X}\to\mathrm{Y}$ un morfismo localmente de tipo finito, sea $x$ un punto de $f^{-1}(y)$ y sea $\sh{F}$ un $\sh{O}_\mathrm{X}$-módulo coherente.
Supongamos satisfechas las condiciones siguientes:
\begin{enumerate}[label=(\roman*)]
  \item $\operatorname{Supp}(\sh{F})=\mathrm{X}$.
  \item $f$ es abierto en los puntos genéricos de las componentes irreducibles de $f^{-1}(y)$ que contienen a $x$.
  \item $\sh{F}_y=\sh{F}\otimes_{\sh{O}_y}\mathbf{k}(y)$ es reducido en el punto $x$ \sref{IV.3.2.2}.
\end{enumerate}
Entonces $\sh{F}$ es $f$-plano en el punto $x$.
\end{lemma}

Denotemos por $t$ un uniformizante de $\mathrm{A}$, pongamos $\mathrm{B}=\sh{O}_x$ y $\sh{F}_x=\mathrm{M}$; la condición~(i) implica que $\operatorname{Supp}(\mathrm{M})=\operatorname{Spec}(\mathrm{B})$ \sref[I]{I.9.1.13}.
La condición~(ii) implica, en virtud de \sref{IV.14.3.2}, que toda componente irreducible de $\mathrm{X}$ que contiene a $x$ domina a $\mathrm{Y}$; dicho de otro modo, la imagen inversa en $\mathrm{A}$ de todo ideal primo minimal $\mathfrak{p}_i$ de $\mathrm{B}$ es $0$, lo que significa también que $t\notin\mathfrak{p}_i$ para todo $i$. Por último, (iii) significa que $\mathrm{M}/t\mathrm{M}$ es un $(\mathrm{B}/t\mathrm{B})$-módulo \emph{reducido} \sref{IV.3.2.2}.
Queda por mostrar que, en estas condiciones, $\mathrm{M}$ es un $\mathrm{A}$-módulo \emph{sin torsión} \sref[0]{0.6.3.4}, y para ello basta mostrar que $t$ es un elemento $\mathrm{M}$-\emph{regular}; pero esto se deduce de \sref{IV.3.4.6.1}.
\end{proof}

\begin{corollary}[15.2.3]
\label{IV.15.2.3}
Sea $\mathrm{Y}$ un preesquema localmente noetheriano, sea $f:\mathrm{X}\to\mathrm{Y}$ un morfismo localmente de tipo finito, sea $x$ un punto de $\mathrm{X}$ y sea $y=f(x)$.
Supongamos satisfechas las condiciones siguientes:
\begin{enumerate}[label=(\roman*)]
  \item $f$ es universalmente abierto en los puntos genéricos de las componentes irreducibles de $f^{-1}(y)$ que contienen a $x$.
  \item $f^{-1}(y)$ es geométricamente reducido sobre $\mathbf{k}(y)$ en el punto $x$ \sref{IV.4.6.9}.
  \item El anillo $\sh{O}_y$ es reducido.
\end{enumerate}
En estas condiciones, $f$ es plano en el punto $x$.
Basta aplicar \sref{IV.15.2.2} con $\sh{F}=\sh{O}_\mathrm{X}$ \sref{IV.4.6.22}.
\end{corollary}

\begin{remarks}[15.2.4]
\label{IV.15.2.4}
\begin{enumerate}[label=(\roman*)]
  \item Las tres primeras condiciones de \sref{IV.15.2.2} no cambian cuando se sustituye $f$ por $f_\mathrm{red}$; pero, si $\mathfrak{N}$ es el nilradical de un anillo local noetheriano $\mathrm{A}$, $\mathrm{A}/\mathfrak{N}$ es plano sobre $\mathrm{A}/\mathfrak{N}$, pero no sobre el propio $\mathrm{A}$ cuando $\mathfrak{N}\neq 0$ (Bourbaki, \emph{Alg.\ comm.}, cap.~II, \S\,3, n.º~2, cor.~2 de la prop.~5).
Se ve así que en \sref{IV.15.2.2} no se puede suprimir la condición~(iv).
  \item Se deduce de \sref{IV.2.4.6} que, si se cumple la conclusión de \sref{IV.15.2.2}, así como la hipótesis~(i), $f$ es universalmente abierto en un entorno de $x$, y en particular en los puntos genéricos de las componentes irreducibles de $f^{-1}(y)$ que contienen a $x$.
Aunque se satisfagan las condiciones~(i), (iii) y (iv), no se puede sustituir (ii) por la hipótesis más débil de que $f$ sea universalmente abierto \emph{únicamente en el punto $x$}.
Esto es lo que muestra el ejemplo \sref{IV.14.1.3},~(i), donde $f$ es evidentemente universalmente abierto en todo punto de $\mathrm{X}_2$, y por tanto en particular en el punto $x$, intersección de $\mathrm{X}_1$ y $\mathrm{X}_2$; las condiciones~(ii) y (iii) de \sref{IV.15.2.3} también se satisfacen en este ejemplo.
Sin embargo, $\sh{O}_x$ no es un $\sh{O}_y$-módulo sin torsión: $\sh{O}_y$ se identifica con un subanillo de $\sh{O}_x$ que contiene divisores de cero en $\sh{O}_x$; por tanto, $f$ no es plano en el punto $x$.
  \item En \sref{IV.15.2.3}, la conclusión no es necesariamente válida cuando se sustituye la hipótesis~(ii) por la hipótesis más débil de que $f^{-1}(y)$, considerado como $\mathbf{k}(y)$-preesquema, \emph{no tenga ciclos primos asociados inmersos}.
Se obtiene un ejemplo tomando como $\mathrm{Y}$ una curva con una cúspide, por ejemplo $\mathrm{Y}=\operatorname{Spec}(\mathrm{K}[\mathrm{S},\mathrm{T}]/(\mathrm{S}^3-\mathrm{T}^2))$, siendo $\mathrm{K}$ un cuerpo algebraicamente cerrado, y como $\mathrm{X}$ su \emph{normalización} \sref[II]{II.6.3.8}.
Si $s$ y $t$ son las clases de $\mathrm{S}$ y $\mathrm{T}$ en $\mathrm{A}=\mathrm{K}[\mathrm{S},\mathrm{T}]/(\mathrm{S}^3-\mathrm{T}^2)$, se verifica inmediatamente que $\mathrm{X}=\operatorname{Spec}(\mathrm{B})$, donde $\mathrm{B}=\mathrm{K}[s,t,u]$, siendo $u$ el elemento $t/s$ del cuerpo de fracciones de $\mathrm{A}$; como $s=u^2$, $t=u^3$, se tiene también $\mathrm{B}=\mathrm{K}[u]$, isomorfo a un anillo de polinomios en una indeterminada sobre $\mathrm{K}$.
El único punto $x$ de $\mathrm{X}$ situado sobre el punto $y$ correspondiente al ideal maximal $(s)+(t)$ corresponde al ideal maximal $(u)$; pero, como la clase $\bar{u}$ de $u$ en $\sh{O}_x/\mathfrak{m}_y\sh{O}_x$ satisface $\bar{u}^2=0$, $f^{-1}(y)$ no es un $\mathbf{k}(y)$-preesquema reducido.
Aquí $f$ es un morfismo finito, sobreyectivo y radicial, y por tanto un homeomorfismo universal \sref{IV.2.4.5}, pero no es plano en el punto $x$, pues si $\sh{O}_x$ fuese un $\sh{O}_y$-módulo plano, sería un $\sh{O}_y$-módulo
\oldpage[IV-3]{228}
libre de rango $1$ generado por el elemento $1$ de $\sh{O}_x$ (Bourbaki, \emph{Alg.\ comm.}, cap.~II, \S\,3, n.º~2, prop.~5), lo que es absurdo.
  \item Mostremos finalmente que en \sref{IV.15.2.2} o \sref{IV.15.2.3} no se puede sustituir ``geométricamente reducido'' por ``reducido''.
Definiremos en efecto dos anillos $\mathrm{A}$, $\mathrm{A}'$ con las propiedades siguientes:
\begin{enumerate}[label=\arabic*)]
  \item $\mathrm{A}$ es un anillo local noetheriano de ideal maximal $\mathfrak{m}$, íntegro, completo, de dimensión $1$ y geométricamente unirramificado.
  \item $\mathrm{A}'$ es la clausura integral de $\mathrm{A}$, una $\mathrm{A}$-álgebra finita cuyo ideal maximal es $\mathfrak{m}\mathrm{A}'$, y su cuerpo residual $k'$ es una extensión radicial finita no trivial del cuerpo residual $k=\mathrm{A}/\mathfrak{m}$ de $\mathrm{A}$.
\end{enumerate}
Se puede tomar entonces $\mathrm{Y}=\operatorname{Spec}(\mathrm{A})$, $\mathrm{X}=\operatorname{Spec}(\mathrm{A}')$, $\sh{F}=\sh{O}_\mathrm{X}$, siendo $y$ y $x$ los puntos cerrados de $\mathrm{Y}$ y $\mathrm{X}$ respectivamente; las hipótesis~(i) y (iv) de \sref{IV.15.2.2} se satisfacen trivialmente; como $\mathrm{A}$ es de dimensión $1$, $f:\mathrm{X}\to\mathrm{Y}$ es radicial, finito y sobreyectivo, y por tanto un homeomorfismo universal \sref{IV.2.4.5}, lo que prueba la hipótesis~(i) de \sref{IV.15.2.3}.
Finalmente, está claro que $f^{-1}(y)$ es reducido.
Sin embargo, $f$ no es plano en el punto $x$, pues si $\mathrm{A}'$ fuese un $\mathrm{A}$-módulo plano, sería un $\mathrm{A}$-módulo libre de rango $1$ (puesto que tiene el mismo cuerpo de fracciones que $\mathrm{A}$) generado por el elemento $1$ de $\mathrm{A}$ (Bourbaki, \emph{Alg.\ comm.}, cap.~II, \S\,3, n.º~2, prop.~5), lo que es absurdo.

Para construir los anillos $\mathrm{A}$ y $\mathrm{A}'$, partamos de un cuerpo imperfecto $k$ de característica $p>0$; sea $\mathrm{K}_0=k((\mathrm{T}))$ el cuerpo de series formales de Laurent sobre $k$, y sea $\mathrm{A}_0$ su anillo de valoración para la valoración discreta dada por el orden de una serie formal de Laurent; el cuerpo residual de $\mathrm{A}_0$ es, pues, $k$.
Sea $\alpha$ un elemento de $k$ que no es una potencia $p$-ésima, y tomemos $k'=k(\alpha^{1/p})$; $\mathrm{K}=k'((\mathrm{T}))$ es entonces una extensión radicial finita de $\mathrm{K}_0$, y $\mathrm{A}'=\mathrm{A}_0\otimes_k k'$ es el anillo de valoración discreta de $\mathrm{K}$.
Si $\mathfrak{m}'$ es el ideal maximal de $\mathrm{A}'$, las condiciones~$1)$ y $2)$ se satisfacen tomando $\mathrm{A}=\mathrm{A}_0+\mathfrak{m}'$: en efecto, como $\mathrm{A}_0$ es noetheriano y $\mathrm{A}'$ es un $\mathrm{A}_0$-módulo de tipo finito, $\mathrm{A}$ es una $\mathrm{A}_0$-álgebra finita, y por tanto un anillo noetheriano; como $\mathrm{A}'$ es finito sobre $\mathrm{A}$, $\mathfrak{m}'$ es el único ideal maximal de $\mathrm{A}$, que es así un anillo local completo, evidentemente de dimensión $1$ (por ser finito sobre $\mathrm{A}_0$) y geométricamente unirramificado, puesto que $\mathrm{A}'$ es íntegramente cerrado y tiene el mismo cuerpo de fracciones $\mathrm{K}$ que $\mathrm{A}$.
  \item La demostración de \sref{IV.15.2.2} muestra, sin embargo, que se puede debilitar la condición~(iii) sustituyendo ``geométricamente reducido'' por ``reducido'' cuando se puede aplicar el criterio valuativo de planitud \sref{IV.11.8.1} utilizando únicamente anillos de valoración discreta $\mathrm{A}'$ cuyo cuerpo residual $k'$ sea \emph{separable} sobre $\mathbf{k}(y)$; en efecto, si $x'$ y $z'$ tienen el mismo significado que en \sref{IV.15.2.2}, las multiplicidades radiciales $\mu_i(\mathbf{k}(z')|k')$ y $\mu_i(\mathbf{k}(z)|\mathbf{k}(y))$ son iguales, en virtud de \sref{IV.4.7.3}, y por tanto se deduce de \sref{IV.4.7.9} que $\operatorname{length}((\sh{F}_y)_z)$ y $\operatorname{length}((\sh{F}_{y'}')_{z'})$ son iguales, y quedamos de nuevo reducidos al caso en que $\mathrm{Y}$ es el espectro de un anillo de valoración discreta e $y$ su punto cerrado.
Por ejemplo, se obtiene este resultado más fuerte cuando $\sh{O}_y$ es \emph{unirramificado} y está dominado por un anillo de valoración discreta cuyo cuerpo residual es una extensión separable de $\mathbf{k}(y)$, en virtud de \sref{IV.11.6.2}.\footnote{El texto impreso cita (11.6.4), que no aparece en {\hyperlink{ega4.c04.subsec.11-6}{\S11.6}}; el resultado aplicable de descenso unirramificado es el Corolario~{\hyperlink{ega4.c04.cor.11-6-2}{(11.6.2)}}.}
Este es el caso, por ejemplo, cuando $\sh{O}_y$ es un anillo \emph{regular}, por \sref{IV.15.1.1.6}.
Pero, como ha señalado Hironaka, existen anillos locales noetherianos \emph{íntegramente cerrados}, de dimensión $2$ (provenientes de esquemas locales algebraicos sobre cuerpos imperfectos), que no satisfacen la condición precedente.
\end{enumerate}
\end{remarks}

\subsection{Aplicaciones: criterios de reducción e irreducibilidad}
\label{subsection:IV.15.3}

\begin{proposition}[15.3.1]
\label{IV.15.3.1}
Sea $\mathrm{Y}$ un preesquema localmente noetheriano, sea $f:\mathrm{X}\to\mathrm{Y}$ un morfismo localmente de tipo finito, sea $\sh{F}$ un $\sh{O}_\mathrm{X}$-módulo coherente, sea $x$ un punto de $\mathrm{X}$ y sea $y=f(x)$.
Supongamos satisfechas las condiciones~\emph{(i)}, \emph{(ii)} y \emph{(iii)} de \sref{IV.15.2.2}.
Entonces:
\begin{enumerate}[label=(\roman*)]
  \item Existe un entorno $\mathrm{U}$ de $x$ en $\mathrm{X}$ tal que $f|\mathrm{U}$ es universalmente abierto y tal que, para todo $x'\in\mathrm{U}$, $\sh{F}_{f(x')}$ es geométricamente reducido sobre $\mathbf{k}(f(x'))$ en el punto $x'$.
  \item Si además $\mathrm{Y}$ es reducido en el punto $y$, existe un entorno $\mathrm{U}_1\subset\mathrm{U}$ de $x$ tal que $\sh{F}|\mathrm{U}_1$ es $f$-plano y reducido.
\end{enumerate}
\end{proposition}

\begin{proof}
\label{IV.15.3.1-proof}
Teniendo en cuenta \sref[I]{I.5.1.8}, \sref{IV.4.2.7} y \sref{IV.4.7.11}, ni las hipótesis ni la conclusión~(i) del enunciado cambian cuando se sustituye $\mathrm{Y}$ por $\mathrm{Y}_\mathrm{red}$ y $\mathrm{X}$ por $\mathrm{X}\times_\mathrm{Y}\mathrm{Y}_\mathrm{red}$, de modo que se puede suponer que $\mathrm{Y}$ es \emph{reducido}.
Se deduce entonces de \sref{IV.15.2.2} que $\sh{F}$ es $f$-plano en el punto $x$, y por tanto, en virtud de \sref{IV.11.1.1}, en un entorno de $x$; \emph{a~fortiori}, $f$ es universalmente abierto en este entorno \sref{IV.2.4.6}.
El hecho de que $\sh{F}_{f(x')}$ sea entonces geométricamente reducido en el punto $x'$
\oldpage[IV-3]{229}
para todos los puntos de un entorno de $x$ se deduce de \sref{IV.12.1.1},~(vii).
La afirmación~(ii) se deduce de esto y de \sref{IV.3.3.4}.
\end{proof}

\begin{proposition}[15.3.2]
\label{IV.15.3.2}
Siendo las notaciones las de \sref{IV.15.3.1}, supongamos satisfechas las condiciones~\emph{(i)}, \emph{(ii)}, \emph{(iii)} de \sref{IV.15.2.2}, y supongamos además que $f^{-1}(y)$ es equidimensional en el punto $x$.
Entonces $f$ es equidimensional en el punto $x$.
\end{proposition}

\begin{proof}
\label{IV.15.3.2-proof}
Se tiene, para todo $y'\in\mathrm{Y}$, $\operatorname{Supp}(\sh{F}_{y'})=f^{-1}(y')$.
Por hipótesis, $\sh{F}_y$ es reducido (y satisface $(S_1)$) y es equidimensional en el punto $x$.
En virtud de \sref{IV.12.1.1},~(ii), se puede suponer, pues, que $\sh{F}_{f(x')}$ es equidimensional y de dimensión constante $e=\dim_x(f^{-1}(y))$ en el punto $x'$ para todo $x'\in\mathrm{U}$, lo que significa que $f^{-1}(f(x'))$ es equidimensional y de dimensión $e$ en el punto $x'$.
Como además $\sh{F}|\mathrm{U}$ es $f$-plano, se deduce de \sref{IV.2.3.4} y \sref{IV.13.3.1},~$a)$ que $f$ es equidimensional en el punto $x$.
\end{proof}

\begin{corollary}[15.3.3]
\label{IV.15.3.3}
Siendo las notaciones las de \sref{IV.15.3.1}, supongamos satisfechas las condiciones siguientes:
\begin{enumerate}[label=(\roman*)]
  \item $\operatorname{Supp}(\sh{F})=\mathrm{X}$.
  \item $f$ es universalmente abierto en los puntos genéricos de las componentes irreducibles de $f^{-1}(y)$ que contienen a $x$.
  \item $\sh{F}_y$ es geométricamente puntualmente íntegro sobre $\mathbf{k}(y)$ en el punto $x$ \sref{IV.4.6.22}.
  \item $\mathrm{Y}$ es geométricamente unirramificado en el punto $y$.
\end{enumerate}
Entonces $x$ pertenece a una única componente irreducible de $\mathrm{X}$.
Si además $\mathrm{Y}$ es íntegro en el punto $y$ y $\sh{F}=\sh{O}_\mathrm{X}$, entonces $\mathrm{X}$ es íntegro en el punto $x$ y $f$ es plano en el punto $x$.
\end{corollary}

\begin{proof}
\label{IV.15.3.3-proof}
Observemos primero que (i) y (iii) implican que $x$ pertenece a una única componente irreducible $\mathrm{Z}$ de $f^{-1}(y)$.
Se deduce entonces de \sref{IV.14.4.7} y \sref{IV.15.3.2} que, para toda componente irreducible $\mathrm{X}_i$ de $\mathrm{X}$ que contiene a $x$ (y por tanto a $\mathrm{Z}$), la restricción $f_i$ de $f$ a $\mathrm{X}_i$ es equidimensional en el punto $x$ y universalmente abierta en el punto genérico de $\mathrm{Z}$.
La primera afirmación del enunciado se deduce entonces de \sref{IV.15.1.3}, y la segunda de \sref{IV.15.3.1}.
\end{proof}

\begin{remarks}[15.3.4]
\label{IV.15.3.4}
\begin{enumerate}[label=(\roman*)]
  \item El ejemplo \sref{IV.14.4.10},~(ii) muestra que, en \sref{IV.15.3.3}, no se puede suprimir la hipótesis de que $\mathrm{Y}$ sea geométricamente unirramificado en el punto $y$.
  \item La observación \sref{IV.15.2.4},~(v) muestra que la conclusión de \sref{IV.15.3.3} sigue siendo cierta bajo las hipótesis siguientes: $\sh{O}_y$ es regular, $\operatorname{Supp}(\sh{F})=\mathrm{X}$ y $\sh{F}_y$ es íntegro en el punto $x$.
Se desconoce si, en este enunciado, la hipótesis de que $\sh{O}_y$ sea regular puede sustituirse por la hipótesis de que sea geométricamente unirramificado, o incluso íntegro e íntegramente cerrado.
\end{enumerate}
\end{remarks}

\subsection{Complementos sobre los morfismos de Cohen--Macaulay}
\label{subsection:IV.15.4}

\begin{proposition}[15.4.1]
\label{IV.15.4.1}
Sea $\mathrm{Y}$ un preesquema localmente noetheriano, sea $f:\mathrm{X}\to\mathrm{Y}$ un morfismo localmente de tipo finito, sea $\sh{F}$ un $\sh{O}_\mathrm{X}$-módulo coherente tal que $\operatorname{Supp}(\sh{F})=\mathrm{X}$, sea $x$ un punto de $\mathrm{X}$ y sea $y=f(x)$.
Ponemos, para todo $z\in\mathrm{X}$, $\mathrm{X}_{f(z)}=f^{-1}(f(z))$, $\sh{F}_{f(z)}=\sh{F}\otimes_{\sh{O}_{f(z)}}\mathbf{k}(f(z))$.
Supongamos que $\sh{F}$ es $f$-plano en el punto $x$ y que $\sh{F}_y$ es un $\sh{O}_{\mathrm{X}_y}$-módulo de Cohen--Macaulay en el punto $x$.
Entonces $f$ es
\oldpage[IV-3]{230}
equidimensional en el punto $x$.
En efecto, por \sref{IV.11.1.1} y \sref{IV.12.1.1},~(vi), existe un entorno $\mathrm{U}$ de $x$ tal que $\sh{F}|\mathrm{U}$ es $f$-plano y, para todo $x'\in\mathrm{U}$, $\sh{F}_{f(x')}$ es un $\sh{O}_{\mathrm{X}_{f(x')}}$-módulo de Cohen--Macaulay en el punto $x'$; esta última propiedad implica que $\sh{F}_{f(x')}$ es equidimensional en el punto $x'$ y que $x'$ no pertenece a ningún ciclo primo asociado inmerso de $\sh{F}_{f(x')}$ \sref[0]{0.16.5.4}.
Además, en virtud de \sref{IV.12.1.1},~(ii), se puede suponer que las dimensiones de las componentes irreducibles de $\operatorname{Supp}(\sh{F}_{f(x')})|(\mathrm{U}\cap\mathrm{X}_{f(x')})$ tienen un valor (común) independiente de $x'$.
Como $\operatorname{Supp}(\sh{F}_{f(x')})=\mathrm{X}_{f(x')}$ por hipótesis, se deduce de \sref{IV.2.3.4} y \sref{IV.13.3.1},~$a)$ que $f$ es equidimensional en el punto $x$.
\end{proposition}

\begin{proposition}[15.4.2]
\label{IV.15.4.2}
Sea $\mathrm{Y}$ un preesquema localmente noetheriano, sea $f:\mathrm{X}\to\mathrm{Y}$ un morfismo localmente de tipo finito, sea $\sh{F}$ un $\sh{O}_\mathrm{X}$-módulo coherente tal que $\operatorname{Supp}(\sh{F})=\mathrm{X}$, sea $x$ un punto de $\mathrm{X}$ y sea $y=f(x)$.
Supongamos que $\sh{O}_y$ es regular y que $\sh{F}$ es un $\sh{O}_\mathrm{X}$-módulo de Cohen--Macaulay en el punto $x$.
Entonces (con las notaciones de \sref{IV.15.4.1}) las condiciones siguientes son equivalentes:
\begin{enumerate}[label=\emph{\alph*})]
  \item $\sh{F}$ es $f$-plano en el punto $x$ y $\sh{F}_y$ es un $\sh{O}_{\mathrm{X}_y}$-módulo de Cohen--Macaulay en el punto $x$.
  \item $\sh{F}$ es $f$-plano en el punto $x$.
  \item $f$ es universalmente abierto en un entorno de $x$.
  \item $f$ es abierto en los puntos genéricos de las componentes irreducibles de $\mathrm{X}_y$ que contienen a $x$.
  \item $\dim(\sh{F}_x)=\dim(\sh{O}_y)+\dim((\sh{F}_y)_x)$.
  \item[\emph{e}$'$)] $\dim(\sh{O}_{\mathrm{X},x})=\dim(\sh{O}_y)+\dim(\sh{O}_{\mathrm{X},x}\otimes_{\sh{O}_y}\mathbf{k}(y))$.
\end{enumerate}
Como $\operatorname{Supp}(\sh{F})=\mathrm{X}$, las condiciones $e)$ y $e')$ son equivalentes por definición \sref{IV.5.1.12}.
La condición $a)$ implica trivialmente $b)$; $b)$ implica que $\sh{F}$ es $f$-plano en los puntos de un entorno de $x$ \sref{IV.11.1.1}, luego $b)$ implica $c)$ \sref{IV.2.4.6}; $c)$ implica trivialmente $d)$, y $d)$ implica $e')$ por \sref{IV.14.2.1}.
Finalmente, como $\sh{O}_y$ es regular y $\sh{F}_x$ es un $\sh{O}_{\mathrm{X},x}$-módulo de Cohen--Macaulay, se deduce de \sref{IV.6.1.4} que $e)$ implica $a)$.
\end{proposition}

\begin{proposition}[15.4.3]
\label{IV.15.4.3}
Sea $f:\mathrm{X}\to\mathrm{Y}$ un morfismo localmente de presentación finita, sea $\sh{F}$ un $\sh{O}_\mathrm{X}$-módulo de presentación finita tal que $\operatorname{Supp}(\sh{F})=\mathrm{X}$, sea $x$ un punto de $\mathrm{X}$ y sea $y=f(x)$.
Entonces (con las notaciones de \sref{IV.15.4.1}) las condiciones siguientes son equivalentes:
\begin{enumerate}[label=\emph{\alph*})]
  \item $\sh{F}$ es $f$-plano en el punto $x$ y $\sh{F}_y$ es un $\sh{O}_{\mathrm{X}_y}$-módulo de Cohen--Macaulay en el punto $x$.
  \item Existe un entorno abierto $\mathrm{U}$ de $x$ en $\mathrm{X}$ y un $\mathrm{Y}$-morfismo cuasifinito $g:\mathrm{U}\to\mathrm{Y}'=\mathrm{Y}[\mathrm{T}_1,\ldots,\mathrm{T}_n]$, donde los $\mathrm{T}_i$ son indeterminadas, tal que $\sh{F}|\mathrm{U}$ es $g$-plano en el punto $x$.
\end{enumerate}
\end{proposition}

\begin{proof}
\label{IV.15.4.3-proof}
Mostremos primero que $b)$ implica $a)$.
Pongamos $h:\mathrm{Y}'\to\mathrm{Y}$; el $\mathbf{k}(y)$-preesquema $\mathrm{Y}_y'=h^{-1}(y)$ es regular \sref[0]{0.17.3.7}; sea $\mathrm{A}$ el anillo local de este preesquema en el punto $y'=g(x)$, sea $k$ su cuerpo residual y sea $\mathrm{B}$ el anillo local de $\mathrm{X}_y$ en el punto $x$; pongamos además $(\sh{F}_y)_x=\mathrm{M}$, que es un $\mathrm{B}$-módulo de tipo finito.
La hipótesis de que el morfismo $g$ es cuasifinito implica lo mismo para $g_y:\mathrm{U}_y\to\mathrm{Y}_y'$ \sref[II]{II.6.2.4}, luego $\mathrm{M}\otimes_\mathrm{A} k$ es un $(\mathrm{B}\otimes_\mathrm{A} k)$-módulo de longitud finita \sref[II]{II.6.2.2}; como $\mathrm{M}$ es un $\mathrm{A}$-módulo plano por hipótesis y $\mathrm{A}$ es un anillo de Cohen--Macaulay, $\mathrm{M}$ es un $\mathrm{B}$-módulo de Cohen--Macaulay \sref{IV.6.3.3}.
El hecho de que $\sh{F}$ sea $f$-plano en el punto $x$ se deduce de que es $g$-plano y de que el morfismo $h$ es plano \sref{IV.2.1.6}.

Mostremos ahora que $a)$ implica $b)$.
Como la cuestión es local sobre $\mathrm{X}$ y $\mathrm{Y}$, se puede suponer que $\mathrm{X}$ e $\mathrm{Y}$ son afines; utilizando \sref{IV.8.9.1} y \sref{IV.11.2.7}, se reduce al caso en que $\mathrm{X}$ e $\mathrm{Y}$ son noetherianos, teniendo en cuenta \sref{IV.6.7.1}.
\oldpage[IV-3]{231}
La hipótesis implica que $f$ es equidimensional en el punto $x$ \sref{IV.15.4.1}, de donde existe un morfismo cuasifinito $g:\mathrm{U}\to\mathrm{Y}'$ tal que toda componente irreducible de $\mathrm{U}$ domina a $\mathrm{Y}'$ \sref{IV.13.3.1},~$b)$ y tal que $g_y:\mathrm{U}_y\to\mathrm{Y}'_y$ es finito y sobreyectivo \sref{IV.13.3.1.1}.
Por otra parte, para probar que $\sh{F}|\mathrm{U}$ es $g$-plano en el punto $x$, basta, puesto que $h:\mathrm{Y}'\to\mathrm{Y}$ es plano, mostrar mediante el criterio de planitud por fibras \sref{IV.11.3.10} que $\sh{F}_y$ es $g_y$-plano en el punto $x$.
Pero, por hipótesis, $\sh{F}_y$ es un $\sh{O}_{\mathrm{X}_y}$-módulo de Cohen--Macaulay en el punto $x$, y $\mathrm{Y}'_y$ es un preesquema regular; además, como $g_y$ es finito y sobreyectivo, satisface la condición $e')$ de \sref{IV.15.4.2} por \sref{IV.5.6.10}.
Para concluir, basta entonces aplicar \sref{IV.15.4.2} a $g_y$ y $\sh{F}_y$.
\end{proof}
\subsection{Rango separable de las fibras de un morfismo cuasifinito y universalmente abierto. Aplicación a las componentes conexas geométricas de las fibras de un morfismo propio}
\label{subsection:IV.15.5}

\begin{proposition}[15.5.1]
\label{IV.15.5.1}
Sea $Y$ un preesquema localmente noetheriano y sea $f:X\to Y$ un morfismo separado y cuasifinito.
Para todo $z\in Y$, sea $n(z)$ el número geométrico de puntos de $f^{-1}(z)$ (o \emph{rango separable} de $f^{-1}(z)$ sobre $\kappa(z)$; cf.\ \sref[I]{I.6.4.8}).
Se tienen las propiedades siguientes:
\begin{enumerate}[label=(\roman*)]
  \item La función $z\mapsto n(z)$ es semicontinua inferiormente en todo punto $y\in Y$ tal que $f$ sea universalmente abierto en los puntos de $f^{-1}(y)$.
  \item Sea $y$ un punto de $Y$ tal que $f$ sea universalmente abierto en los puntos de $f^{-1}(y)$; si la función $z\mapsto n(z)$ es constante en un entorno de $y$, entonces $f$ es propio en el punto $y$ \sref{IV.15.7.1}.
  \item Supongamos que $f$ es universalmente abierto y que todo punto de $Y$ es geométricamente unirramificado; entonces, si la función $z\mapsto n(z)$ es localmente constante en $Y$, las componentes irreducibles de $X$ son dos a dos disjuntas.
\end{enumerate}
\end{proposition}

\begin{proof}
(i)
Observemos que, como $f$ es cuasifinito, el número $n(z)$ es el número geométrico de componentes conexas de $f^{-1}(z)$; sabemos que la función $z\mapsto n(z)$ es localmente construible \sref{IV.9.7.9}; teniendo en cuenta $(\mathbf{0}_{\mathrm{III}},~9.3.4)$, quedamos reducidos a probar el

\begin{corollary}[15.5.2]
\label{IV.15.5.2}
Bajo las hipótesis generales de \sref{IV.15.5.1}, sea $y$ un punto de $Y$ tal que $f$ sea universalmente abierto en los puntos de $f^{-1}(y)$.
Entonces, si $y'$ es una generalización de $y$, se tiene
\end{corollary}

\begin{equation}
\tag{15.5.2.1}
\label{IV.15.5.2.1}
  n(y) \leq n(y').
\end{equation}

Para todo cambio de base $g:Z\to Y$, $f_{(\mathrm{Z})}:X_{(\mathrm{Z})}\to Z$ es separado, cuasifinito y universalmente abierto en los puntos de $X_{(\mathrm{Z})}$ que se proyectan en $f^{-1}(y)$; además, si $z$, $z'$ son dos puntos de $Z$ tales que $g(z)=y$, $g(z')=y'$ y $z'$ es una generalización de $z$, se tiene $n(z)=n(y)$ y $n(z')=n(y')$ (\sref[I]{I.6.4.12}); basta, pues, probar \sref{IV.15.5.2.1} para un $g$ conveniente.
Evidentemente, se puede suponer $y\neq y'$.
Utilizaremos el lema siguiente:

\begin{lemma}[15.5.2.2]
\label{IV.15.5.2.2}
Bajo las hipótesis generales de \sref{IV.15.5.1}, si $y$ es un punto de $Y$ e $y'$ es una generalización de $y$ distinta de $y$, existe un esquema $Z$ que es el espectro de un anillo de valoración discreta, de punto cerrado $z$ y punto genérico $z'$, y un morfismo $g:Z\to Y$ tales que:
$1^\circ$ $g(z)=y$, $g(z')=y'$;
$2^\circ$ si se pone $f'=f_{(Z)}$, $n(y)$ es el número de puntos de $f'^{-1}(z)$ y $n(y')$ es el número de puntos de $f'^{-1}(z')$.
\end{lemma}

\oldpage[IV-3]{232}
En efecto, existe un anillo de valoración discreta $A_1$ y un morfismo $g_1$ de $Z_1=\Spec(A_1)$ en $Y$ tales que, si $z_1$ y $z_1'$ son el punto cerrado y el punto genérico de $Z_1$, se tiene $g_1(z_1)=y$ y $g_1(z_1')=y'$ (\sref[II]{II.7.1.9}); se puede suponer ya, por tanto, que $Y$ es el espectro de un anillo de valoración discreta $A$, que $y$ es su punto cerrado e $y'$ su punto genérico.
Para todo $x\in f^{-1}(y)$, $\kappa(x)$ es por hipótesis una extensión finita de $\kappa(y)$ (\sref[II]{II.6.2.2}); se deduce de \sref{IV.4.5.11} que existe una extensión algebraica finita $k$ de $\kappa(y)$ tal que el \emph{número de puntos} de $f^{-1}(y)\otimes_{\kappa(y)}k$ es igual a $n(y)$.
Como existe un anillo de valoración discreta $A_2$ que domina a $A$, cuyo cuerpo residual es finito sobre $\kappa(y)$ y contiene a $k$ (\sref[II]{II.7.1.2}), se puede suponer, en segundo lugar, que $n(y)$ es el número de puntos de $f^{-1}(y)$.
Finalmente, el mismo razonamiento muestra que existe una extensión \emph{finita} $K$ de $\kappa(y')$ tal que el número de puntos de $f^{-1}(y')\otimes_{\kappa(y')}K$ es igual a $n(y')$; si $w$ es una valoración de $\kappa(y')$ asociada a $A$, existe una valoración discreta de $K$ que prolonga a $w$, y el anillo $A_3$ de esta valoración domina a $A$, tiene a $K$ como cuerpo de fracciones y satisface las condiciones del lema.

En adelante se puede suponer, pues, que $Y$ es el espectro de un anillo de valoración discreta, que $y$ es su punto cerrado, $y'$ su punto genérico, y que $n(y)$ y $n(y')$ son los números de puntos de $f^{-1}(y)$ y $f^{-1}(y')$, de modo que para todo $x\in f^{-1}(y)$ (resp.\ todo $x'\in f^{-1}(y')$) se tiene $\kappa(x)=\kappa(y)$ (resp.\ $\kappa(x')=\kappa(y')$).

Denotemos entonces por $X_i$ los subpreesquemas reducidos de $X$ cuyos espacios subyacentes son las componentes irreducibles de $X$ ($1\leq i\leq m$).
En virtud de \sref{IV.1.10.4}, la restricción $X_i\to Y$ de $f$ a cada $X_i$ es un morfismo \emph{dominante} y, como $f^{-1}(y')$ es discreto, cada intersección $X_i\cap f^{-1}(y')$ consta únicamente del punto genérico de $X_i$; por tanto (denotando por $n_i(z)$ el número geométrico de puntos de $X_i\cap f^{-1}(z)$ para todo $z\in Y$), $n_i(y')=1$ para todo $i$, y $n(y')=\sum_{i=1}^{m} n_i(y')=m$.
Además, $f^{-1}(y)$ es la unión de los $X_i\cap f^{-1}(y)$, de donde
\begin{equation}
\tag{15.5.2.3}
\label{IV.15.5.2.3}
  n(y) \leq \sum_{i=1}^{m} n_i(y),
\end{equation}
y, para probar \sref{IV.15.5.2.1}, basta establecer que $n_i(y)\leq 1$ para todo $i$.
Dicho de otro modo, se puede suponer que $X$ es \emph{íntegro} y que el morfismo $f$ es birracional.
Pero entonces, para todo $x\in f^{-1}(y)$, se tiene $\sh{O}_y\subset\sh{O}_x\subset K$, donde $K=\kappa(y')$ es el cuerpo de fracciones de $\sh{O}_y$.
Como $\sh{O}_y$ es un anillo de valoración y $\sh{O}_x$ domina a $\sh{O}_y$, se tiene necesariamente $\sh{O}_y=\sh{O}_x$ (\sref[II]{II.7.1.1}).
Existe entonces una $Y$-sección $h:Y\to X$ tal que $h(y)=x$ (\sref[I]{I.2.4.4}); como $Y$ es reducido y $X$ es separado sobre $Y$, tal sección es \emph{única} (\sref[I]{I.7.2.2}), de modo que el conjunto $f^{-1}(y)$ contiene un solo punto, lo que termina la demostración de (i).

(ii)
Como al comienzo de (i), quedamos reducidos a mostrar que, si los dos miembros de \sref{IV.15.5.2.1} son iguales para \emph{toda} generalización $y'$ de $y$, $f$ es propio en el punto $y$.
Gracias al criterio de propiedad local \sref{IV.15.7.5} y a las observaciones iniciales, se puede suponer además que $Y=\Spec(A)$ es el espectro de un anillo de valoración discreta $A$, que $y$ es su punto cerrado e $y'$ su punto genérico, y queda por mostrar que $f$ es propio.

Utilicemos \sref{IV.15.5.2.2} y observemos que, si $A'$ es un anillo de valoración discreta que domina a $A$, $A'$ es un $A$-módulo sin torsión, de modo que el morfismo $\Spec(A')\to\Spec(A)$ es fielmente plano y cuasicompacto; por tanto, equivale a decir que $f$ es propio o que el morfismo deducido de $f$ por el cambio de base $\Spec(A')\to\Spec(A)$ es
\oldpage[IV-3]{233}
propio \sref{IV.2.7.1},~(vii); se puede suponer, por tanto, que $n(y)$ y $n(y')$ son los números de puntos de las fibras $f^{-1}(y)$ y $f'^{-1}(y')$, respectivamente.
Con las notaciones de (i), se tiene entonces, por \sref{IV.15.5.2.3} y \sref{IV.15.5.2.1},
\begin{equation}
\tag{15.5.2.4}
\label{IV.15.5.2.4}
  n(y) \leq \sum_i n_i(y) \leq \sum_i n_i(y') = n(y')
\end{equation}
y, para que los miembros extremos sean iguales, es necesario que $n_i(y)=n_i(y')=1$ para todo $i$.
Como vimos en (i), si $n_i(y)=1$ para todo $i$, $X_i\cap f^{-1}(y)$ no es vacío y la restricción $X_i\to Y$ de $f$ es un isomorfismo; como los $X_i$ son subpreesquemas cerrados de $X$, esto implica que $f$ es propio (\sref[II]{II.5.4.5}).

(iii)
Podemos limitarnos al caso en que $Y$ es afín y reducido y, como $Y$ es por hipótesis localmente íntegro (\sref[I]{I.5.1.4}), se puede suponer que $Y$ es íntegro y que la función $y\mapsto n(y)$ es constante en $Y$.
Sean de nuevo $X_i$ ($1\leq i\leq m$) las componentes irreducibles de $X$.
Se deduce de \sref{IV.1.10.4} que, si $y'$ es el punto genérico de $Y$, cada uno de los $X_i\cap f^{-1}(y')$ consta de un solo punto; como la restricción $f_i:X_i\to Y$ de $f$ es un morfismo cuasifinito y dominante y todo punto de $Y$ es geométricamente unirramificado, se deduce del criterio de Chevalley \sref{IV.14.4.4} que $f_i$ es \emph{universalmente abierto}.
Así, para todo punto $y$ de $Y$, se tiene $n_i(y)\leq n_i(y')$; por otra parte, como los $X_i\cap f^{-1}(y')$\footnote{La fuente impresa tiene aquí $X_i\cap f^{-1}(y)$; la igualdad siguiente $\sum_i n_i(y')=n(y')$ y el argumento de la fibra genérica exigen $y'$.} son dos a dos disjuntos, $\sum_i n_i(y')=n(y')$; se siguen teniendo, por tanto, las relaciones \sref{IV.15.5.2.4}.
Pero, por hipótesis, los miembros extremos son iguales, de modo que $n(y)=\sum_i n_i(y)$, lo que implica que los $X_i\cap f^{-1}(y)$ son dos a dos disjuntos, y esto termina la demostración.
\end{proof}

Combinando esta proposición con el teorema de conexidad de Zariski y sus consecuencias (\sref[III]{III.4.3}), se obtienen los resultados siguientes, que completan (\sref[III]{III.4.3.7} y \sref[III]{III.4.3.10}):

\begin{proposition}[15.5.3]
\label{IV.15.5.3}
Sea $Y$ un preesquema irreducible localmente noetheriano, sea $f:X\to Y$ un morfismo propio; para todo $y\in Y$, sea $n(y)$ el número geométrico de componentes conexas \sref{IV.4.5.2} de $f^{-1}(y)$.
Supongamos que la restricción de $f$ a toda componente irreducible de $X$ es dominante sobre $Y$.
Entonces, si $y_0$ es un punto geométricamente unirramificado de $Y$, la función $y\mapsto n(y)$ es semicontinua inferiormente en el punto $y_0$.
\end{proposition}

\begin{proof}
Consideremos la factorización de Stein $f=g\circ f'$ (\sref[III]{III.4.3.3}), donde $g:Y'\to Y$ es un morfismo finito y $f':X\to Y'$\footnote{La fuente impresa tiene aquí $f':X\to Y$; la factorización de Stein enunciada y el resto de la demostración exigen como destino $Y'$.} es un morfismo sobreyectivo cuyas fibras son geométricamente conexas (\sref[III]{III.4.3.4}).
Como $f'$ es sobreyectivo, la imagen inversa por $f'$ de toda componente irreducible de $Y'$ contiene al menos una componente irreducible de $X$; por tanto, cada componente irreducible de $Y'$ domina a $Y$.
Se puede aplicar entonces el criterio de Chevalley \sref{IV.14.4.4} a las restricciones de $g$ a cada una de estas componentes irreducibles, y se concluye que $g$ es \emph{universalmente abierto en todo punto} de $g^{-1}(y_0)$.
La conclusión se deduce, pues, de \sref{IV.15.5.2} y del hecho de que el número de componentes geométricamente conexas de $f^{-1}(y)$ es igual al número geométrico de puntos de $g^{-1}(y)$ (\sref[III]{III.4.3.4}).
\end{proof}

\begin{corollary}[15.5.4]
\label{IV.15.5.4}
Sea $Y$ un preesquema localmente noetheriano, sea $f:X\to Y$ un morfismo propio; sea $y$ un punto de $Y$ tal que $f$ sea universalmente abierto en los puntos de $f^{-1}(y)$;
\oldpage[IV-3]{234}
entonces, con las notaciones de \sref{IV.15.5.3}, la función $z\mapsto n(z)$ es semicontinua inferiormente en el punto $y$.
\end{corollary}

\begin{proof}
Con las notaciones de \sref{IV.15.5.3}, observemos que en la factorización de Stein $f=g\circ f'$, $f'$ es sobreyectivo; como $f$ es universalmente abierto en los puntos de $f^{-1}(y)$, $g$ es universalmente abierto en los puntos de $g^{-1}(y)$ (\sref{IV.14.3.4},~(i)); basta entonces aplicar a $g$ la proposición \sref{IV.15.5.1},~(i), teniendo en cuenta (\sref[III]{III.4.3.4}).
\end{proof}

\begin{remarks}[15.5.5]
\label{IV.15.5.5}
(i)
Aunque $Y$ sea íntegro y normal, $X$ íntegro y $f$ finito y sobreyectivo (y por tanto universalmente abierto en virtud del criterio de Chevalley \sref{IV.14.4.4}), la función $y\mapsto n(y)$ no es necesariamente localmente constante.
Se obtiene un ejemplo tomando $Y=\Spec(k[\mathrm{T}])$, donde $k$ es algebraicamente cerrado, y $X=\Spec(A)$, donde $A=k[\mathrm{S},\mathrm{T}]/(\mathrm{S}^2+\mathrm{T}^2-1)$; en los puntos $y'$, $y''$ de $Y$ correspondientes a los ideales maximales $(\mathrm{T}-1)$ y $(\mathrm{T}+1)$, se tiene $n(y')=n(y'')=1$, pero $n(y)=2$ en todos los demás puntos de $Y$.
Más adelante daremos una condición adicional que asegura que $y\mapsto n(y)$ sea localmente constante \sref{IV.15.5.7}.

(ii)
El ejemplo \sref{IV.14.4.10} (ii) muestra que en \sref{IV.15.5.1},~(iii), no se puede suprimir la hipótesis de que los puntos de $Y$ sean geométricamente unirramificados.

(iii)
El ejemplo \sref{IV.11.7.5} muestra que la conclusión de \sref{IV.15.5.1},~(i) deja de ser válida si solo se supone que el morfismo $f$ es \emph{abierto} (incluso si, como ocurre en el ejemplo citado, $f$ es finito y sobreyectivo).

(iv)
Finalmente, en \sref{IV.15.5.1},~(iv), no se puede prescindir de la hipótesis de que el morfismo $f$ sea \emph{separado}, como muestra el ejemplo en que $X$ es la recta afín en la que se ha ``duplicado'' un punto $x$ (\sref[I]{I.5.5.11}) e $Y$ es la recta afín: aquí $f$ es un isomorfismo local (y por tanto plano) y es cuasifinito, pero $z\mapsto n(z)$ no es semicontinua inferiormente.
\end{remarks}

\begin{lemma}[15.5.6]
\label{IV.15.5.6}
Sea $Y$ el espectro de un anillo de valoración discreta, sea $a$ su punto cerrado y $b$ su punto genérico.
Sea $f:X\to Y$ un morfismo abierto localmente de tipo finito.
Supongamos que $X$ es conexo y que la fibra $f^{-1}(a)$ es un preesquema reducido.
Entonces $f^{-1}(b)$ es conexo.
\end{lemma}

\begin{proof}
Utilizaremos el siguiente lema puramente topológico (caso particular de un resultado más general del cap.\ III, 3.ª Parte):

\begin{lemma}[15.5.6.1]
\label{IV.15.5.6.1}
Sea $X$ un espacio topológico localmente noetheriano en el que todo subconjunto cerrado irreducible admite un punto genérico.
Sea $Z$ un subconjunto cerrado de $X$ de interior vacío que posee la propiedad siguiente: para todo $x\in Z$, si se denota por $\mathrm{T}_x$ el conjunto de generalizaciones de $x$ en $X$, entonces $\mathrm{T}_x-\{x\}$ es conexo.
En estas condiciones, si $X$ es conexo, $X-Z$ es conexo.
\end{lemma}

\begin{proof}
Demos una demostración independiente.
Razonando por contradicción, supongamos que el abierto $U=X-Z$, que es denso en $X$ puesto que $Z$ tiene interior vacío, es la unión de dos abiertos no vacíos disjuntos $U'$ y $U''$, y denotemos por $X'$ y $X''$ las adherencias en $X$ de $U'$ y $U''$.
Como $X=\overline{U}=X'\cup X''$ y $X$ es conexo, se tiene $X'\cap X''\neq\varnothing$, y está claro que $X'\cap X''\subset Z$.
Sea $x$ un punto genérico de una componente irreducible de $X'\cap X''$.
Como $X$ es localmente noetheriano, existe un entorno abierto $V$ de $x$ en $X$ tal que $V\cap U'$ y $V\cap U''$ solo tienen un número finito de componentes irreducibles; como $x$ es adherente a $U'$ y a $U''$, pertenece necesariamente a la adherencia de una componente irreducible $Z'$ de $V\cap U'$ y a la adherencia de una componente irreducible $Z''$ de $V\cap U''$.
Pero $\overline{Z'}$ (resp.\ $\overline{Z''}$) es cerrado e irreducible en $X$, y por tanto admite un punto genérico $z'$ (resp.\ $z''$), y se tiene necesariamente $z'\in U'$ y $z''\in U''$.
Esto prueba que las intersecciones de $\mathrm{T}_x-\{x\}$ con $U'$ y $U''$ no son vacías.
Pero, por definición, $(X'\cap X'')\cap(\mathrm{T}_x-\{x\})=\varnothing$; por tanto, las intersecciones de $X'$ y de $X''$ con $\mathrm{T}_x-\{x\}$ son dos cerrados no vacíos disjuntos de $\mathrm{T}_x-\{x\}$ cuya unión es $\mathrm{T}_x-\{x\}$; esto contradice la hipótesis de que $\mathrm{T}_x-\{x\}$ sea conexo.
\end{proof}
\end{proof}

\oldpage[IV-3]{235}
\begin{proof}[Demostración de {\rm {\hyperlink{ega4.c04.lem.15-5-6}{(15.5.6)}}}]
Para probar \sref{IV.15.5.6}, aplicamos el lema \sref{IV.15.5.6.1} a $X$ y a su subconjunto cerrado de interior vacío $Z=f^{-1}(a)$, que tiene interior vacío porque $f$ es abierto.
Observemos que, por hipótesis y teniendo en cuenta \sref{IV.15.2.2.1}, esto implica que $f$ es \emph{plano} en los puntos de $f^{-1}(a)$.
Para todo punto $x$ de esta índole, $\sh{O}_x$ es, pues, un $\sh{O}_a$-módulo sin torsión ($\mathbf{0}_{\mathrm{I}},~6.3.4$), y en particular, si $t$ es un uniformizante de $\sh{O}_a$, $t$ es un elemento $\sh{O}_x$-regular; además, $\sh{O}_x/t\sh{O}_x$ es reducido por hipótesis. Si es de profundidad $0$, es por tanto un cuerpo ($\mathbf{0},~16.4.7$), y como $t\sh{O}_x$ es entonces el ideal maximal de $\sh{O}_x$, $\sh{O}_x$ es un anillo de valoración discreta ($\mathbf{0},~17.1.4$); por consiguiente, $\Spec(\sh{O}_x)-\{x\}$ consta de un solo punto y es, \emph{a fortiori}, conexo.
Si, por el contrario, $\operatorname{depth}(\sh{O}_x/t\sh{O}_x)\geq 1$, se tiene $\operatorname{depth}(\sh{O}_x)\geq 2$, puesto que $t$ es $\sh{O}_x$-regular ($\mathbf{0},~16.4.6$); se deduce entonces del teorema de Hartshorne \sref{IV.5.10.7} que $\Spec(\sh{O}_x)-\{x\}$ es conexo, lo que termina la demostración.
\end{proof}

\begin{proposition}[15.5.7]
\label{IV.15.5.7}
Sea $Y$ un preesquema localmente noetheriano y sea $f:X\to Y$ un morfismo propio.
Para todo $z\in Y$, sea $n(z)$ el número geométrico de componentes conexas de $f^{-1}(z)$.
Sea $y$ un punto de $Y$ tal que $f$ sea universalmente abierto en los puntos de $f^{-1}(y)$ y que $f^{-1}(y)$ sea geométricamente reducido sobre $\kappa(y)$ \sref{IV.4.6.2}.
Entonces la función $z\mapsto n(z)$ es constante en un entorno de $y$.
\end{proposition}

\begin{proof}
Ya sabemos por \sref{IV.15.5.4} que, en las condiciones del enunciado, $z\mapsto n(z)$ es semicontinua inferiormente en el punto $y$; queda por ver que también es semicontinua superiormente y, por el mismo razonamiento que al comienzo de la demostración de \sref{IV.15.5.1},~(i), basta mostrar que, si $y'$ es una generalización de $y$, se tiene
\begin{equation}
\tag{15.5.7.1}
\label{IV.15.5.7.1}
  n(y') \leq n(y).
\end{equation}

Utilizando el argumento del comienzo de la demostración de \sref{IV.15.5.2} y el Lema \sref{IV.15.5.2.2}, aplicados al morfismo finito $g$ de la factorización de Stein $f=g\circ f'$ (recordando que el número geométrico de puntos de $g^{-1}(y)$ es igual al número de componentes conexas de $f^{-1}(y)$ (\sref[III]{III.4.3.4})), se reduce al caso en que $n(y)$ y $n(y')$ son respectivamente los números de componentes conexas de $f^{-1}(y)$ y $f^{-1}(y')$, y en que $Y$ es el espectro de un anillo de valoración discreta, $y$ es su punto cerrado e $y'$ su punto genérico.
Además, se puede sustituir $X$ por cualquiera de sus componentes conexas, pues estas son abiertas y cerradas en $X$ y propias sobre $Y$.
Supongamos, pues, que $X$ es conexo. La hipótesis de que $f$ es abierto en los puntos de $f^{-1}(y)$ implica que, si $f^{-1}(y)\neq\varnothing$, entonces $f^{-1}(y')\neq\varnothing$ (1.10.3); la hipótesis de que $f$ es cerrado implica que, si $f^{-1}(y')\neq\varnothing$, entonces $f^{-1}(y)\neq\varnothing$ (\sref[II]{II.7.2.1}).
Por tanto, si $X\neq\varnothing$ (lo que evidentemente se puede suponer, pues la proposición es trivialmente cierta en el caso contrario), $f^{-1}(y)$ y $f^{-1}(y')$ son ambos no vacíos.
Además, la hipótesis implica que el preesquema $f^{-1}(y)$ es \emph{reducido} \sref{IV.4.6.1}; como $f$ es abierto en los puntos de $f^{-1}(y)$, el lema \sref{IV.15.5.6} implica que $f^{-1}(y')$ es conexo, es decir, $n(y')=1$.
Como $f^{-1}(y)$ no es vacío, esto prueba \sref{IV.15.5.7.1}.
\end{proof}

\begin{remarks}[15.5.8]
\label{IV.15.5.8}
La relación \sref{IV.15.5.7.1} no es necesariamente válida si no se supone que $f$ sea propio, incluso si satisface las demás hipótesis de \sref{IV.15.5.7}.
Con las notaciones de \sref{IV.15.5.5},~(i), basta considerar la restricción de $f$ a $X-\{x'\}$, donde $x'$ es uno de los dos puntos de $X$ situados sobre un punto $y$ distinto de $y'$ e $y''$; se tiene entonces $n(y)=1$, mientras que, si $\eta$ es el punto genérico de $Y$, $n(\eta)=2$.
\end{remarks}

\oldpage[IV-3]{236}

\begin{proposition}[15.5.9]
\label{IV.15.5.9}
Sea $f:X\to Y$ un morfismo plano y localmente de presentación finita; para todo $y\in Y$, sea $n(y)$ el número geométrico de componentes conexas de $f^{-1}(y)$.
\begin{enumerate}[label=(\roman*)]
  \item Si $f$ es separado y cuasifinito, la función $y\mapsto n(y)$ (igual entonces al número geométrico de puntos de $f^{-1}(y)$) es semicontinua inferiormente sobre $Y$; si es constante en un entorno de $y_0$, $f$ es propio en $y_0$.
  \item Si $f$ es propio, la función $y\mapsto n(y)$ es semicontinua inferiormente sobre $Y$; si además $f^{-1}(y_0)$ es geométricamente reducido sobre $\kappa(y_0)$, $n(y)$ es constante en un entorno de $y_0$.
\end{enumerate}
\end{proposition}

\begin{proof}
En todos los casos, las cuestiones son locales sobre $Y$, de modo que se puede suponer que $Y$ es afín y $f$ de presentación finita; utilizando \sref{IV.8.9.1}, \sref{IV.8.10.5} y \sref{IV.11.2.7}, se reduce al caso en que $Y$ es noetheriano (utilizando además \sref{IV.8.2.11} para la segunda afirmación de (i)).
Como entonces $f$ es además universalmente abierto \sref{IV.2.4.6}, basta aplicar \sref{IV.15.5.1}, \sref{IV.15.5.4} y \sref{IV.15.5.7} para concluir.
\end{proof}

\begin{remarks}[15.5.10]
\label{IV.15.5.10}
Observemos que así se ha obtenido otra demostración de \sref{IV.12.2.4},~(vi).
Recíprocamente, se puede deducir \sref{IV.15.5.7} de \sref{IV.12.2.4},~(vi): en efecto, se puede limitar al caso en que $Y$ es reducido (sustituyendo $f$ por $f_\mathrm{red}$), y se deduce entonces de las hipótesis de \sref{IV.15.5.7} y de \sref{IV.15.2.3} que $f$ es plano en un entorno abierto de $f^{-1}(y)$; como además $f$ es propio, se puede suponer que este entorno es de la forma $f^{-1}(U)$, donde $U$ es un entorno abierto de $y$ en $Y$.
Se concluye entonces por \sref{IV.12.2.4},~(vi).
La demostración de \sref{IV.15.5.7} dada anteriormente tiene la ventaja de poner de manifiesto el resultado \sref{IV.15.5.6}, que posee interés independiente.
\end{remarks}


\subsection{Componentes conexas de las fibras a lo largo de una sección}
\label{subsection:IV.15.6}

\begin{proposition}[15.6.1]
\label{IV.15.6.1}
Sea $Y$ un preesquema localmente noetheriano, sea $f:X\to Y$ un morfismo localmente de tipo finito y sea $g:Y\to X$ una $Y$-sección de $X$ (\sref[I]{I.2.5.5}).
Para todo $y\in Y$, denotemos por $X_y^0$ la componente conexa de $f^{-1}(y)$ que contiene al punto $g(y)$.
Sea $y$ un punto de $Y$ y sea $y'$ una generalización de $y$ en $Y$.
Entonces:
\begin{enumerate}[label=(\roman*)]
  \item Si existe una componente irreducible $Z$ de $X_{y'}^0$, de dimensión igual a $\dim(X_{y'}^0)$ y que contiene a $g(y')$ (como ocurre si $X_{y'}^0$ es irreducible), entonces
\begin{equation}
\tag{15.6.1.1}
\label{IV.15.6.1.1}
  \dim(X_y^0) \geq \dim(X_{y'}^0).
\end{equation}
  \item Si además $X_y^0$ es irreducible y si los dos miembros de \sref{IV.15.6.1.1} son iguales, entonces $X_y^0\subset\overline{X_{y'}^0}$ (adherencia en $X$).
\end{enumerate}
\end{proposition}

\begin{proof}
(i)
Sea $\overline{Z}$ la adherencia de $Z$ en $X$; como $y$ es una especialización de $y'$ y $g$ es continua, $g(y)\in\overline{Z}$.
Como $Z=\overline{Z}\cap f^{-1}(y')$, se deduce del teorema de semicontinuidad de Chevalley (\sref{IV.13.1.3}), aplicado a la restricción de $f$ al subpreesquema cerrado reducido de $X$ cuyo espacio subyacente es $\overline{Z}$, que $\dim_{g(y)}(\overline{Z}\cap f^{-1}(y))\geq\dim(Z)$; pero las componentes irreducibles de $\overline{Z}\cap f^{-1}(y)$ que contienen a $g(y)$ están evidentemente contenidas en $X_y^0$, de donde \emph{a fortiori} la desigualdad \sref{IV.15.6.1.1}.

(ii)
El razonamiento de (i) muestra además que, si los dos miembros de \sref{IV.15.6.1.1} son iguales, se tiene necesariamente $\dim_{g(y)}(\overline{Z}\cap f^{-1}(y))=\dim_{g(y)}(X_y^0)$; si $X_y^0$ es irreducible, $\overline{Z}\cap f^{-1}(y)=X_y^0$, de modo que $X_y^0$ está contenido en $\overline{Z}$ y, \emph{a fortiori}, en $\overline{X_{y'}^0}$.
\end{proof}

\oldpage[IV-3]{237}
\begin{corollary}[15.6.2]
\label{IV.15.6.2}
Con las notaciones de \sref{IV.15.6.1}, supongamos además que, para todo $y\in Y$, $X_y^0$ es irreducible.
Entonces la función $z\mapsto\dim(X_z^0)$ es semicontinua superiormente sobre $Y$.
Si además esta función es constante en un entorno de un punto $y\in Y$, entonces $X_y^0\subset\overline{X_{y'}^0}$ para toda generalización $y'$ de $y$.
\end{corollary}

\begin{proof}
Sea $y$ un punto de $Y$ y sea $U$ un entorno abierto afín de $g(y)$ en $X$; existe entonces un entorno abierto $V$ de $y$ en $Y$ tal que $g(V)\subset U$, y para todo $z\in V$, $U\cap X_z^0$ es denso en $X_z^0$, luego $\dim(X_z^0)=\dim(U\cap X_z^0)$ (\sref{IV.4.1.1.3}).
Sustituyendo $X$ por $U$, se puede suponer que $f$ es un morfismo de tipo finito.
Se sabe entonces por \sref{IV.9.7.10} que la función $z\mapsto\dim(X_z^0)$ es localmente construible, y las afirmaciones del corolario se deducen de \sref{IV.15.6.1} y de $(\mathbf{0}_{\mathrm{III}},~9.3.4)$.
\end{proof}

\begin{proposition}[15.6.3]
\label{IV.15.6.3}
Con las notaciones de \sref{IV.15.6.1}, supongamos que, para todo $y\in Y$, $X_y^0$ es geométricamente irreducible \sref{IV.4.5.2}.
Entonces, para todo $y\in Y$ tal que la función $z\mapsto\dim(X_z^0)$ sea constante en un entorno de $y$, $f$ es universalmente abierto en los puntos de $X_y^0$.
\end{proposition}

\begin{proof}
Apliquemos el criterio \sref{IV.14.3.7}.
Sea, pues, $h:Y'\to Y$ un morfismo, siendo $Y'$ el espectro de un anillo de valoración discreta cuyos puntos cerrado y genérico denotamos por $t$ y $t'$, respectivamente, y supongamos que $h(t)=y$, de modo que $y'=h(t')$ es una generalización de $y$.
Pongamos $X'=X\times_Y Y'$, $f'=f_{(Y')}:X'\to Y'$, $g'=g_{(Y')}:Y'\to X'$.
Con las mismas notaciones que en \sref{IV.15.6.1}, la hipótesis sobre $X_y^0$ implica que $X_t^{\prime 0}=X_y^0\otimes_{\kappa(y)}\kappa(t)$ y $X_{t'}^{\prime 0}=X_{y'}^0\otimes_{\kappa(y')}\kappa(t')$, y que $X_t^{\prime 0}$ y $X_{t'}^{\prime 0}$ son irreducibles \sref{IV.4.4.1}; como $\dim(X_y^0)=\dim(X_t^{\prime 0})$ y $\dim(X_{y'}^0)=\dim(X_{t'}^{\prime 0})$ \sref{IV.4.1.4}, se ve que basta probar la proposición para $f'$ y $X_t^{\prime 0}$. Dicho de otro modo, se puede limitar al caso en que $Y$ es el espectro de un anillo de valoración discreta, $y$ su punto cerrado e $y'$ su punto genérico; si $x'$ es el punto genérico de $X_{y'}^0$, se deduce de \sref{IV.15.6.2} que todo punto de $X_y^0$ es una especialización de $x'$, de donde la conclusión.
\end{proof}

\begin{proposition}[15.6.4]
\label{IV.15.6.4}
Bajo las hipótesis generales de \sref{IV.15.6.1}, sea $X^0$ la unión de los $X_y^0$ cuando $y$ recorre $Y$.
Si $y\in Y$ es tal que $X_y^0$ es geométricamente reducido sobre $\kappa(y)$ \sref{IV.4.6.2} y si $f$ es universalmente abierto en todo punto de $X_y^0$, entonces $X^0$ es un entorno de $X_y^0$ en $X$.
En particular, si $f$ es universalmente abierto y si, para todo $y\in Y$, $X_y^0$ es geométricamente reducido sobre $\kappa(y)$, $X^0$ es abierto en $X$.
\end{proposition}

\begin{proof}
Mostremos primero que se puede reducir al caso en que $f$ es un morfismo de \emph{tipo finito}.
Sea $x$ un punto de $X_y^0$; se deduce de \sref{IV.5.10.8.1} (con $\mathfrak{F}=\{\varnothing\}$) que existe una sucesión finita $(Z_i)_{0\leq i\leq n}$ de componentes irreducibles de $f^{-1}(y)$ tal que $g(y)\in Z_0$, $x\in Z_n$ y $Z_i\cap Z_{i+1}\neq\varnothing$ para $0\leq i\leq n-1$; como los $Z_i$ son irreducibles, existe una sucesión finita $(U_i)$ ($1\leq i\leq n$) de abiertos afines de $f^{-1}(y)$ tal que $g(y)\in U_1$, $x\in U_n$, y $U_i$ es un entorno abierto afín de un punto $x_i$ de $Z_i\cap Z_{i-1}$ para $1\leq i\leq n$.
Para cada $i$ existe un abierto cuasicompacto $W_i$ de $X$ tal que $U_i=f^{-1}(y)\cap W_i$; sea $W$ el abierto cuasicompacto de $X$ que es la unión de los $W_i$.
Como $g$ es continua, existe un entorno abierto afín $V$ de $y$ en $Y$ tal que $g(V)\subset W$; si $X'=W\cap f^{-1}(V)$, la restricción de $f$ a $X'$ es de tipo finito. Además, para todo $z\in V$, $X_z^{\prime 0}$ es la componente conexa de $X'\cap f^{-1}(z)$ que contiene a $g(z)$, y se tiene $X_z^{\prime 0}\subset X_z^0$, mientras que $x$ pertenece a $X_y^{\prime 0}$.
En efecto, cada $U_i$ contiene los dos conjuntos irreducibles $U_i\cap Z_{i-1}$
\oldpage[IV-3]{238}
y $U_i\cap Z_i$, que se cortan, y los conjuntos irreducibles $U_i\cap Z_{i-1}$ y $U_{i-1}\cap Z_{i-1}$ se cortan para $2\leq i\leq n$; por tanto, la unión de los $U_i\cap Z_{i-1}$ y los $U_i\cap Z_i$, para $1\leq i\leq n$, es conexa y contiene a $x$ y a $g(y)$.
Como $f|_{X'}$ es universalmente abierto en los puntos de $X_y^{\prime 0}$, esto termina la demostración de nuestra afirmación: si la proposición está probada cuando $f$ es de tipo finito, la unión $X^{\prime 0}$ de los $X_z^{\prime 0}$, para $z\in V$, es un entorno de $x$, y por tanto lo mismo ocurre con $X^0$.

Supongamos, pues, que $f$ es de tipo finito.
Entonces $X^0$ es localmente construible \sref{IV.9.7.12}; basta, por tanto, mostrar que $X^0$ contiene toda generalización $x'$ de $x$ ($\mathbf{0}_{\mathrm{III}},~9.2.5$).
Pongamos $y'=f(x')$; existe un anillo de valoración discreta $A$ y un morfismo $u$ de $Y'=\Spec(A)$ en $X$ tales que, si $z$ es el punto cerrado y $z'$ el punto genérico de $Y'$, se tiene $u(z)=x$ y $u(z')=x'$ (\sref[II]{II.7.1.9}); por tanto, $h=f\circ u$, y así $h(z)=y$ y $h(z')=y'$.
Pongamos $X'=X\times_Y Y'$, $f'=f_{(Y')}:X'\to Y'$, que es un morfismo de tipo finito universalmente abierto en los puntos situados sobre $X_y^0$,\footnote{La fuente impresa afirma que $f'$ es universalmente abierto en los puntos de $f'^{-1}(z)$; la hipótesis solo da apertura universal sobre la componente conexa distinguida $X_y^0$.} y $g'=g_{(Y')}:Y'\to X'$, que es una $Y'$-sección de $X'$.
Si $p_y:f'^{-1}(z)\to f^{-1}(y)$ (resp.\ $p_{y'}:f'^{-1}(z')\to f^{-1}(y')$) es la proyección canónica, $p_y^{-1}(X_y^0)$ (resp.\ $p_{y'}^{-1}(X_{y'}^0)$) es la componente conexa de $f'^{-1}(z)$ (resp.\ $f'^{-1}(z')$) que contiene a $g'(z)$ (resp.\ $g'(z')$), puesto que las componentes conexas $X_t^0$ son geométricamente conexas \sref{IV.4.5.13} y basta aplicar \sref{IV.4.4.1}.
Además, el morfismo $u$ corresponde a una $Y'$-sección $v$ de $X'$ tal que $p_y(v(z))=x$ y $p_{y'}(v(z'))=x'$, de modo que $v(z')$ es una generalización de $v(z)$; bastará, pues, probar que $v(z')$ pertenece a $p_{y'}^{-1}(X_{y'}^0)$.\footnote{La fuente impresa escribe $v(y)$ y $v(y')$ en esta frase; los puntos cerrado y genérico de $Y'$ se denotaron por $z$ y $z'$.}
Finalmente, está claro que $p_y^{-1}(X_y^0)$ es geométricamente reducido sobre $\kappa(z)$.

Dicho de otro modo, quedamos reducidos al caso en que $Y$ es el espectro de un anillo de valoración discreta, $y$ su punto cerrado e $y'$ su punto genérico.
Como $f^{-1}(y)-X_y^0$ es entonces cerrado en $X$, sustituyendo $X$ por el abierto $X-(f^{-1}(y)-X_y^0)$ se puede suponer que $X_y^0=f^{-1}(y)$; de manera análoga, se puede sustituir $X$ por el abierto que es la componente conexa de $X$ que contiene a $g(Y)$, componente conexa que contiene evidentemente a $X^0$; dicho de otro modo, se puede suponer que $X$ es \emph{conexo}.
La proposición quedará entonces establecida si se muestra que $X^0=X$, o equivalentemente que $X_{y'}$ es conexo; pero $f^{-1}(y')$ es geométricamente reducido sobre $\kappa(y')$ y, \emph{a fortiori}, \emph{reducido}; además, $f$ es abierto en los puntos de $f^{-1}(y)$, de modo que se puede aplicar el lema \sref{IV.15.5.6}, de donde la conclusión.
\end{proof}

\begin{corollary}[15.6.5]
\label{IV.15.6.5}
Sea $f:X\to Y$ un morfismo plano de presentación finita y sea $g$ una $Y$-sección de $X$; para todo $y\in Y$, sea $X_y^0$ la componente conexa de $f^{-1}(y)$ que contiene a $g(y)$, y sea $X^0$ la unión de los $X_y^0$ cuando $y$ recorre $Y$.
Entonces, para todo $y\in Y$ tal que $X_y^0$ sea geométricamente reducido sobre $\kappa(y)$, $X^0$ es un entorno de $X_y^0$ en $X$.
En particular, si $f$ es reducido \sref{IV.6.8.1}, $X^0$ es abierto en $X$.
\end{corollary}

\begin{proof}
Como la cuestión es local sobre $Y$, se puede suponer $Y=\Spec(A)$ afín.
Existen, pues, un subanillo noetheriano $A_0$ y un morfismo plano de tipo finito $f_0:X_0\to Y_0=\Spec(A_0)$ tales que $X=X_0\times_{Y_0}Y$ y $f=(f_0)_{(Y)}$ (\sref{IV.8.9.1} y \sref{IV.11.2.7}); además, se puede suponer que existe una $Y_0$-sección $g_0$ de $X_0$ tal que $g=(g_0)_{(Y)}$ \sref{IV.8.8.2}.
Para todo $y\in Y$, sea $y_0$ su proyección en $Y_0$; se deduce de \sref{IV.4.5.13} que $(X_0)_{y_0}^0$ es geométricamente conexo, de modo que, si $p:f^{-1}(y)\to f_0^{-1}(y_0)$ es la proyección canónica, se tiene $X_y^0=p^{-1}((X_0)_{y_0}^0)$ \sref{IV.4.4.1}; además, como $f$ es plano, la hipótesis de que $X_y^0$ es geométricamente reducido sobre $\kappa(y)$ implica que $(X_0)_{y_0}^0$ es geométricamente reducido sobre $\kappa(y_0)$ \sref{IV.4.6.10}.
Quedamos así reducidos al caso en que
\oldpage[IV-3]{239}
$Y$ es además \emph{noetheriano}; como $f$ es universalmente abierto (\sref{IV.11.1.1}), basta entonces aplicar \sref{IV.15.6.4}.
\end{proof}

\begin{proposition}[15.6.6]
\label{IV.15.6.6}
Sea $f:X\to Y$ un morfismo tal que $Y$ sea localmente noetheriano y $f$ localmente de tipo finito, o tal que $f$ sea de presentación finita.
Sea $g$ una $Y$-sección de $X$; para todo $y\in Y$, sea $X_y^0$ la componente conexa de $f^{-1}(y)$ que contiene a $g(y)$; finalmente, sea $X^0$ la unión de los $X_y^0$ cuando $y$ recorre $Y$.
Supongamos que, para todo $y\in Y$, $X_y^0$ es geométricamente irreducible.
Entonces:
\begin{enumerate}[label=(\roman*)]
  \item La función $y\mapsto\dim(X_y^0)$ es semicontinua superiormente en $Y$.
  \item Sea $x_0\in X^0$, $y_0=f(x_0)$.
Si la función $y\mapsto\dim(X_y^0)$ es constante en un entorno de $y_0$, existe un entorno abierto $U$ de $y_0$ tal que $f$ es universalmente abierto en todos los puntos de $X^0\cap f^{-1}(U)$.
Si además $Y$ es reducido y $f^{-1}(y_0)$ es geométricamente reducido sobre $\kappa(y_0)$ en el punto $x_0$, $f$ es plano en el punto $x_0$.
  \item Recíprocamente, supongamos que $f$ es universalmente abierto en el punto genérico de $X_{y_0}^0$ y que, además, se satisface una de las condiciones siguientes:
  \begin{enumerate}
    \item[$\alpha$)] La fibra $f^{-1}(y_0)$ es geométricamente reducida sobre $\kappa(y_0)$ en el punto $g(y_0)$ y el anillo $\sh{O}_{Y,y_0}$ es noetheriano.
    \item[$\beta$)] Para todo punto genérico $y'$ de una componente irreducible de $Y$ que contiene a $y_0$, se tiene $\dim(f^{-1}(y'))=\dim(X_{y'}^0)$.
  \end{enumerate}
Entonces la función $y\mapsto\dim(X_y^0)$ es constante en un entorno de $y_0$.
  \item El conjunto $W$ de puntos $x\in X^0$ tales que la función $y\mapsto\dim(X_y^0)$ es constante en un entorno de $f(x)$ y que $X_{f(x)}^0$ es geométricamente reducido sobre $\kappa(f(x))$, es abierto en $X$.
\end{enumerate}
\end{proposition}

\begin{proof}
Como las cuestiones son locales sobre $Y$, se puede suponer que $Y=\Spec(A)$ es afín.

Cuando $f$ es de presentación finita, existe un subanillo noetheriano $A_1$ de $A$ y un morfismo de tipo finito $f_1:X_1\to Y_1=\Spec(A_1)$ tales que $X=X_1\times_{Y_1}Y$ y $f=(f_1)_{(Y)}$ \sref{IV.8.9.1}; además, se puede suponer que existe una $Y_1$-sección $g_1$ de $X_1$ tal que $g=(g_1)_{(Y)}$ \sref{IV.8.8.2}.
Para todo $y\in Y$, sea $y_1$ su proyección en $Y_1$; se deduce de \sref{IV.4.5.13} que $(X_1)_{y_1}^0$ es geométricamente conexo, de modo que, si $p:f^{-1}(y)\to f_1^{-1}(y_1)$ es la proyección canónica, $X_y^0=p^{-1}((X_1)_{y_1}^0)$ \sref{IV.4.4.1}.
Observemos además que, en virtud de \sref{IV.9.7.11} y \sref{IV.9.7.12}, se puede aplicar \sref{IV.9.3.3} a la propiedad de ser geométricamente irreducible; se puede suponer, por tanto, que $A_1$ se ha elegido de modo que $(X_1)_{y_1}^0$ sea geométricamente irreducible para todo $y_1\in Y_1$.
Se tiene $\dim(X_y^0)=\dim((X_1)_{y_1}^0)$ \sref{IV.4.1.4}, y, aplicando una vez más \sref{IV.9.3.3} a la propiedad de tener una dimensión prescrita (y utilizando \sref{IV.9.5.5} y \sref{IV.9.7.12}), se puede suponer (restringiendo $Y$, si es necesario, a un entorno de $y_0$) que, si $\dim(X_y^0)$ es constante en $Y$, entonces $\dim((X_1)_{y_1}^0)$ es constante en $Y_1$.
Si $Y$ es reducido, también lo es $Y_1$ puesto que $A_1\subset A$; además, si $f^{-1}(y_0)$ es geométricamente reducido sobre $\kappa(y_0)$ en el punto $x_0$, entonces $f_1^{-1}(y_{01})$ es geométricamente reducido en el punto $x_{01}$ (siendo $x_{01}$ e $y_{01}$ las proyecciones respectivas de $x_0$ e $y_0$ en $X_1$ e $Y_1$) \sref{IV.4.6.10}.

Estas observaciones muestran que, para probar (i) y (ii), se puede limitar al caso en que $Y$ es noetheriano y $f$ localmente de tipo finito.
Pero, en este caso, la afirmación (i) se deduce de \sref{IV.15.6.2}, y la primera afirmación de (ii), de \sref{IV.15.6.3}.
Además, si $Y$ es reducido y $f^{-1}(y_0)$ es geométricamente reducido sobre $\kappa(y_0)$ en el punto $x_0$ (suponiendo aún $Y$ noetheriano), entonces, como $X_{y_0}^0$ es la única componente irreducible de $f^{-1}(y_0)$ que contiene a $x_0$,
\footnote{La fuente impresa tiene aquí $X_y^0$, $f^{-1}(y)$ y $x$; la afirmación~(ii) concierne a los puntos fijados $x_0$ e $y_0=f(x_0)$.} se
\oldpage[IV-3]{240}
deduce de \sref{IV.15.6.3} y \sref{IV.15.2.3} que, si $\dim(X_{y_0}^0)$ es constante en un entorno de $y_0$, $f$ es plano en el punto $x_0$.

Para probar (iii), observemos primero que el conjunto $X^0$ es localmente construible en $X$ \sref{IV.9.7.12}, de modo que, por \sref{IV.9.5.5}, el conjunto $E$ de los $y\in Y$ tales que $\dim(X_y^0)=\dim(X_{y_0}^0)$ es localmente construible en $Y$.
Apliquemos entonces \sref{IV.1.10.1} a $E$: teniendo en cuenta (i), basta probar que $\dim(X_{y'}^0)=\dim(X_{y_0}^0)$ para el punto genérico $y'$ de una componente irreducible de $Y$ que contiene a $y_0$; esto permite ya suponer que $Y=\Spec(\sh{O}_{Y,y_0})$.

Si estamos en el caso $\alpha$), entonces, en virtud de \sref[II]{II.7.1.9}, y utilizando \sref{IV.4.5.13} y \sref{IV.4.4.1} como en \sref{IV.15.6.4}, se reduce inmediatamente al caso en que $Y$ es el espectro de un anillo de valoración discreta, $y_0$ su punto cerrado e $y'$ su punto genérico.
Pero entonces las hipótesis implican, en virtud de \sref{IV.15.2.2.1}, que $f$ es \emph{plano} en el punto $g(y_0)$, y también en un entorno $V$ de este punto en $X$ (\sref{IV.11.1.1}).
Para probar nuestra afirmación, se puede sustituir $X$ por $V$, pues $V\cap X_{y'}^0$ no es vacío en virtud de \sref{IV.2.3.4}, y por tanto es un abierto por todas partes denso en $X_{y'}^0$ y tiene la misma dimensión \sref{IV.4.1.1.3}; además, como $X_{y'}^0$ es también la componente conexa de $f^{-1}(y')$ que contiene a $g(y')$, se puede suponer que $V$ está elegido de modo que no encuentre ninguna otra componente irreducible de $f^{-1}(y_0)$ ni de $f^{-1}(y')$; dicho de otro modo, se puede suponer que $X^0=X$; pero entonces la conclusión se deduce de \sref{IV.12.1.1},~(i), puesto que por hipótesis $X_{y_0}^0$ es íntegro.

Supongamos ahora que estamos en el caso $\beta$).
Apliquemos esta vez (\sref[II]{II.7.1.4}) y (\sref[II]{II.7.1.9}) como en el caso $\alpha$): quedamos entonces reducidos al caso en que $Y$ es el espectro de un anillo de valoración (no necesariamente discreta), $y_0$ su punto cerrado e $y'$ su punto genérico.
En virtud de \sref{IV.14.3.13}, existe una componente irreducible $Z$ de $X$ que contiene a $X_{y_0}^0$ y domina a $Y$, y tal que $\dim(Z\cap f^{-1}(y'))=\dim(X_{y_0}^0)$; pero la hipótesis $\beta$) y el hecho de que $\dim(Z\cap f^{-1}(y'))\leq\dim(f^{-1}(y'))$ muestran que $\dim(X_{y_0}^0)\leq\dim(X_{y'}^0)$, de donde $\dim(X_{y_0}^0)=\dim(X_{y'}^0)$ en virtud de (i).

Queda por probar (iv).
Observemos primero que los conjuntos considerados no cambian cuando se sustituye $f$ por $f_{(Y_\mathrm{red})}:X\times_Y Y_\mathrm{red}\to Y_\mathrm{red}$, pues la proyección $X\times_Y Y_\mathrm{red}\to X$ es un homeomorfismo.
Dicho de otro modo, se puede suponer que $Y$ es reducido, y se deduce entonces de (ii) que $f$ es \emph{plano} en los puntos de $W$.
Consideremos un punto $x_0$ de $W$ y probemos que $W$ es un entorno de $x_0$; procediendo como al comienzo de la demostración y utilizando \sref{IV.11.2.7}, se puede suponer también que $Y$ es \emph{noetheriano}; se deduce entonces de (ii) y de \sref{IV.15.6.4} que $X^0$ es un entorno de $x_0$ en $X$, y de \sref{IV.12.1.1},~(vii), que $W$ también es un entorno de $x_0$ en $X$.
\end{proof}

\begin{corollary}[15.6.7]
\label{IV.15.6.7}
Supongamos que $f$ satisface las condiciones preliminares de \sref{IV.15.6.6}, y supongamos además que, para todo $y\in Y$, $X_y^0$ es geométricamente íntegro sobre $\kappa(y)$.
Entonces las condiciones siguientes son equivalentes:
\begin{enumerate}[label=\emph{\alph*)}]
  \item La función $y\mapsto\dim(X_y^0)$ es localmente constante en $Y$.
  \item El morfismo $f_{(Y_\mathrm{red})}:X\times_Y Y_\mathrm{red}\to Y_\mathrm{red}$ deducido de $f$ por cambio de base es plano en los puntos de $X^0$.
\end{enumerate}
\oldpage[IV-3]{241}
Además, estas condiciones implican que $X^0$ es abierto en $X$.
Finalmente, cuando $Y$ es localmente noetheriano, $a)$ y $b)$ son también equivalentes a
\begin{enumerate}
  \item[\rm b$'$)] $f$ es universalmente abierto en los puntos de $X^0$.
\end{enumerate}
\end{corollary}

\begin{proof}
El hecho de que $a)$ implique $b)$ se deduce de \sref{IV.15.6.6},~(ii), junto con el hecho de que $X^0$ es entonces abierto en $X$.
Si se satisface $b)$, se puede limitar al caso en que $Y$ es reducido y $f$ plano; entonces se reduce, como al comienzo de \sref{IV.15.6.6}, y utilizando además \sref{IV.11.2.7}, al caso en que $Y$ es noetheriano, caso en el que la conclusión se deduce de \sref{IV.15.6.6},~(iii), caso $\alpha$)).
La equivalencia de $b)$ y $b')$ ya ha sido probada cuando $Y$ es localmente noetheriano \sref{IV.15.2.2.1}, teniendo en cuenta que el morfismo $X\times_Y Y_\mathrm{red}\to X$ es un homeomorfismo universal.
\end{proof}

\begin{proposition}[15.6.8]
\label{IV.15.6.8}
Sea $f:X\to Y$ un morfismo separado de presentación finita y sea $g$ una $Y$-sección de $X$; para todo $y\in Y$, sea $X_y^0$ la componente conexa de $f^{-1}(y)$ que contiene a $g(y)$, y sea $X^0$ la unión de los $X_y^0$ cuando $y$ recorre $Y$.
Entonces, si $y\in Y$ es tal que $X_y^0$ es propio sobre $\kappa(y)$, $X^0$ es propio sobre $Y$ en el punto $y$ (es decir, en el sentido de \sref{IV.15.7.1}, existe un entorno abierto $U$ de $y$ en $Y$ tal que $X^0\cap f^{-1}(U)$ es cerrado en $f^{-1}(U)$ y propio sobre $U$).
\end{proposition}

\begin{proof}
Procediendo como al comienzo de \sref{IV.15.6.6} (donde se utiliza \sref{IV.8.10.5},~(v), y se sustituye \sref{IV.9.7.11} por \sref{IV.9.6.7}), se reduce al caso en que $Y$ es noetheriano y $f$ de tipo finito. Se sabe entonces por \sref{IV.9.7.12} que $X^0$ es localmente construible en $X$; además, los $X_z^0$ son geométricamente conexos (para todo $z\in Y$) en virtud de \sref{IV.4.5.13}; finalmente, como $g(Y)\subset X^0$, $X^0$ es universalmente submersivo sobre $Y$ \sref{IV.15.7.8}.
Basta entonces aplicar a $X^0$ el criterio \sref{IV.15.7.9}.
\end{proof}

\begin{corollary}[15.6.9]
\label{IV.15.6.9}
Sea $f:X\to Y$ un morfismo separado tal que $Y$ sea localmente noetheriano y $f$ localmente de tipo finito, o tal que $f$ sea de presentación finita.
Sea $g$ una $Y$-sección de $X$, y para todo $y\in Y$, sea $X_y^0$ la componente conexa de $f^{-1}(y)$ que contiene a $g(y)$.
Sea $y$ un punto de $Y$ tal que $X_y^0$ sea propio sobre $\kappa(y)$.
Entonces, para toda generalización $y'$ de $y$, $\dim(X_y^0)\geq\dim(X_{y'}^0)$.
\end{corollary}

\begin{proof}
Supongamos primero que $f$ es de presentación finita; entonces, sustituyendo $Y$ por un entorno de $y$ si es necesario, se puede suponer que $f|_{X^0}$ es \emph{propio} \sref{IV.15.6.8}, y basta aplicar \sref{IV.13.1.5} a esta restricción.
Si $Y$ es localmente noetheriano y $f$ localmente de tipo finito, $X$ es localmente noetheriano; además, la hipótesis de que $X_y^0$ es propio sobre $\kappa(y)$ implica que $X_y^0$ es noetheriano. Existe, pues, un abierto noetheriano $V\subset X$ que contiene a $X_y^0$; por tanto, existe un entorno abierto $W$ de $y$ en $Y$ tal que $g(W)\subset V$.
Además, si $z$ es un punto maximal de $X_{y'}^0$, se puede suponer que $z\in V$. La restricción de $f$ a $V\cap f^{-1}(W)$ es entonces un morfismo de tipo finito, y $g|_W$ es una $W$-sección; aplicando a este morfismo el resultado ya obtenido, se ve que, si $Z$ es la componente irreducible de $X_{y'}^0$ de punto genérico $z$, entonces $\dim(Z)\leq\dim(X_y^0)$. Como $z$ es un punto maximal arbitrario de $X_{y'}^0$, se tiene en efecto $\dim(X_{y'}^0)\leq\dim(X_y^0)$.
\end{proof}

\begin{remarks}[15.6.10]
\label{IV.15.6.10}
(i)
La hipótesis suplementaria sobre la existencia de una componente irreducible de $X_{y'}^0$ que contiene a $g(y')$ y de dimensión igual a $\dim(X_{y'}^0)$, formulada en \sref{IV.15.6.1}, no es superflua, como muestra el ejemplo siguiente.
Sea $A$ un anillo de valoración discreta, sea $\pi$ un uniformizante de $A$, sea $K$ el cuerpo de fracciones de $A$ y sea $k$ su cuerpo residual. Pongamos $Y=\Spec(A)$ y denotemos por $y$ e $y'$ el punto cerrado y el punto genérico de $Y$.
Consideremos los dos $Y$-esquemas: $X_1=\Spec(A[t])$, $X_2=\Spec(K[u,v])$, donde $t$, $u$, $v$ son indeterminadas (observemos que la proyección de $X_2$ sobre $Y$ es, por tanto, $\{y'\}$).
En el anillo $A[t]$, el ideal principal $(\pi t-1)$ es \emph{maximal}, y corresponde, pues, a un punto cerrado $x_1$ de $X_1$, que se proyecta sobre el \emph{punto genérico} $y'$ de $Y$ y es tal que $\kappa(x_1)=K$.
\oldpage[IV-3]{242}
Consideremos también el punto cerrado $x_2$ de $X_2$ correspondiente al ideal maximal $\left(u-\frac{1}{\pi}\right)+(v)$ de $K[u,v]$; se tiene asimismo $\kappa(x_2)=K$.
Como se verá en el cap.\ V, se pueden entonces \emph{repegar} $X_1$ y $X_2$ a lo largo de dos subpreesquemas cerrados $Z_1$ y $Z_2$ de $X_1$ y $X_2$ cuyos espacios subyacentes son, respectivamente, $\{x_1\}$ y $\{x_2\}$, por medio del único $Y$-isomorfismo de $Z_1$ sobre $Z_2$.
Denotemos por $X$ el $Y$-preesquema así obtenido y por $x$ el punto de $X$ que es la imagen de $x_1$ y $x_2$.
La ``sección nula'' $g$ de $X_1$ (\sref[II]{II.1.7.9}) es entonces también una $Y$-sección de $X$; $f^{-1}(y)$ es igual a $(X_1)_y$, mientras que $f^{-1}(y')$ tiene dos componentes irreducibles, isomorfas respectivamente a $(X_1)_{y'}^0$ y $(X_2)_{y'}^0$, que tienen como único punto común a $x$ y cuyas dimensiones respectivas son~1 y~2. Véase, sin embargo, la Proposición~\sref{IV.15.6.9}.

(ii)
Consideremos el morfismo $f$ definido en \sref{IV.12.2.3},~(ii), que es propio y plano; además, la restricción de $f$ a $X_1\cap f_0^{-1}(Y)$ es un isomorfismo, de modo que el morfismo inverso $g:Y\to X$ es una $Y$-sección de $X$.
Se tiene entonces $\dim(X_{y_0})=1$, mientras que $\dim(X_y^0)=0$ para $y\neq y_0$, aunque $X_y^0$ es geométricamente irreducible para todo $y\in Y$ (pero $X_{y_0}^0$ no es reducido); se ve, por tanto, que en \sref{IV.15.6.6},~(iii), no se pueden omitir las hipótesis $\alpha$) y $\beta$).
Además, $X^0$ no es un entorno de $X_{y_0}^0$ en $X$, lo que muestra que en \sref{IV.15.6.4} no se puede prescindir de la hipótesis de que $X_{y_0}^0$ sea reducido.\footnote{La fuente impresa tiene aquí $X^y$; la frase precedente y la Proposición~\sref{IV.15.6.4} exigen $X_{y_0}^0$.}

(iii)
En el cap.\ VI aplicaremos los resultados precedentes a los $Y$-\emph{preesquemas en grupos} localmente de tipo finito sobre un preesquema localmente noetheriano $Y$.
Si $G$ es un preesquema de esta índole, existe la $Y$-sección canónica $g$, la ``sección unidad'', y se muestra que $G_y^0$ es siempre geométricamente irreducible y $\dim(G_y^0)=\dim(G_y)$, poniendo $G_y=f^{-1}(y)$.
Se deduce entonces de \sref{IV.15.6.2} y \sref{IV.15.6.3} que la función $z\mapsto\dim(G_z)$ es semicontinua superiormente y que, si es constante en un entorno de un punto $y$, $f$ es universalmente abierto en los puntos de $G_y^0$.
Si además $G_y$ es geométricamente reducido sobre $\kappa(y)$ en uno de sus puntos, se muestra que $G_y$ es \emph{liso} \sref{IV.6.8.1} sobre $\kappa(y)$ en todos sus puntos; se deducirá entonces de \sref{IV.15.6.6} que, si además $\sh{O}_y$ es reducido, también $f$ es liso en todos los puntos de $G_y^0$ (aunque no necesariamente en todos los puntos de $G_y$).
Estos resultados se aplicarán, por ejemplo, en la teoría de los esquemas de Picard, donde utilizaremos un teorema general que asegura que, bajo ciertas condiciones, la función $z\mapsto\dim(G_z)$ es localmente constante.
Observemos que, teniendo en cuenta aplicaciones de esta índole, los enunciados como \sref{IV.15.6.1} se formulan para morfismos \emph{localmente} de tipo finito, y no solo para morfismos de tipo finito.

Observemos también que, en el caso de un $Y$-preesquema en grupos $G$, la hipótesis $\beta$) de \sref{IV.15.6.6} siempre se satisface.
\end{remarks}


\subsection{Apéndice: criterios valuativos de propiedad local}
\label{subsection:IV.15.7}

Esta sección aporta complementos al criterio valuativo de propiedad probado en (\sref[II]{II.7.3.10}); es independiente del resto del \S~15.

\begin{definition}[15.7.1]
\label{IV.15.7.1}
Sea $f:X\to Y$ un morfismo de preesquemas y sea $y$ un punto de $Y$.
Se dice que $f$ es \emph{propio en el punto $y$} si existe un entorno abierto $U$ de $y$ en $Y$ tal que la restricción $f^{-1}(U)\to U$ de $f$ es un morfismo propio.
Se dice que una parte $Z$ de $X$ es \emph{propia sobre $Y$ en el punto $y$} si existe un entorno abierto $U$ de $y$ tal que $Z\cap f^{-1}(U)$ es cerrado en $f^{-1}(U)$ y propio sobre $U$ (\sref[II]{II.5.4.10}) (lo que equivale a decir que, para todo subpreesquema cerrado de $X$ cuyo espacio subyacente es $Z$, la restricción de $f$ a $Z$ es propia en el punto $y$).
\end{definition}

Sea $g:Y'\to Y$ un morfismo cualquiera; pongamos $X'=X\times_Y Y'$, $f'=f_{(Y')}$ y, si $p:X'\to X$ es la proyección canónica, $Z'=p^{-1}(Z)$.
Entonces, si $Z$ es propio sobre $Y$ en el punto $y$, $Z'$ es propio sobre $Y'$ en todo punto $y'$ situado sobre $y$ (\sref[II]{II.5.4.2}).

\begin{proposition}[15.7.2]
\label{IV.15.7.2}
Sea $Y$ un preesquema íntegro localmente noetheriano, sea $X$ un preesquema íntegro, sea $f:X\to Y$ un morfismo separado, dominante y de tipo finito, y sea $y$ un punto de $Y$.
Las condiciones siguientes son equivalentes:
\begin{enumerate}[label=\emph{\alph*)}]
  \item $f$ es propio en el punto $y$.
  \item Si $Y_1=\Spec(\sh{O}_y)$, $X_1=X\times_Y Y_1$, el morfismo $f_1=f_{(Y_1)}:X_1\to Y_1$ es propio.
  \item Todo anillo de valoración discreta que tiene a $R(X)$ como cuerpo de fracciones y domina a $\sh{O}_y$\footnote{La fuente impresa escribe aquí únicamente $\mathcal O$; el anillo local exigido por el criterio es $\mathcal O_y$.} domina un anillo local $\sh{O}_x$ de $X$ (en cuyo caso se tiene necesariamente $f(x)=y$).
\end{enumerate}
\end{proposition}

\oldpage[IV-3]{243}
\begin{proof}
El hecho de que $a)$ implique $b)$ se deduce de (\sref[II]{II.5.4.2}), y la implicación $b)\Rightarrow c)$ se deduce de (\sref[II]{II.7.3.10}).
Queda, pues, por mostrar que $c)$ implica $a)$.
La cuestión es local sobre $Y$, de modo que se puede suponer que $Y$ es afín, y por tanto noetheriano. Por el lema de Chow (\sref[II]{II.5.6.1}), existen un preesquema íntegro $X'$, un morfismo proyectivo $p:P\to Y$, una inmersión abierta dominante $j:X'\to P$ y un morfismo proyectivo birracional (y por tanto sobreyectivo) $g:X'\to X$ tales que el diagrama
\[
  \xymatrix{
    P \ar[d]_-{p} & X' \ar[l]_-{j} \ar[d]^-{g} \\
    Y & X \ar[l]^-{f}
  }
\]
es conmutativo.
Como $X'$ es íntegro y $j$ dominante, $P$ es irreducible y, sustituyendo $P$ por $P_\mathrm{red}$, se puede suponer que $P$ es íntegro (\sref[I]{I.5.2.2} y \sref[II]{II.5.5.5},~(vi)).
Basta probar que existe un entorno abierto $U$ de $y$ tal que $p^{-1}(U)\subset j(X')$, pues entonces la restricción de $j$ a
\[
  j^{-1}\bigl(p^{-1}(U)\bigr)=g^{-1}\bigl(f^{-1}(U)\bigr)=V
\]
será un isomorfismo sobre $p^{-1}(U)$. Por tanto, la restricción $V\to U$ de $f\circ g$ será propia y, como la restricción $V\to f^{-1}(U)$ de $g$ es sobreyectiva, la restricción $f^{-1}(U)\to U$ de $f$ será propia (\sref[II]{II.5.4.3}).
Como $T=P-j(X')$ es cerrado en $P$, $p(T)$ es cerrado en $Y$, y basta mostrar que $y\notin p(T)$, o incluso que todo $z'\in P$ tal que $p(z')=y$ \emph{pertenece} a $j(X')$.
Ahora bien, el cuerpo de funciones racionales $R(P)$ es igual a $R(X')$, y por tanto a $R(X)$ por construcción; como se puede limitar al caso en que $z'$ no es el punto genérico de $P$, existe un anillo de valoración discreta $A$ que tiene a $R(X)$ como cuerpo de fracciones y domina a $\sh{O}_{z'}$ (\sref[II]{II.7.1.7}); como $p(z')=y$, $A$ domina a $\sh{O}_y$, y por hipótesis existe, pues, $x\in X$ tal que $f(x)=y$ y $A$ domina a $\sh{O}_x$.
Como $g$ es propio, existe $x'\in X'$ tal que $g(x')=x$ y $A$ domina a $\sh{O}_{x'}$ (\sref[II]{II.7.3.10}); por tanto, $z'$ y $j(x')$ son dos puntos de $P$ cuyos anillos locales $\sh{O}_{z'}$ y $\sh{O}_{j(x')}=\sh{O}_{x'}$ están relacionados por (\sref[I]{I.8.1.4}); como $P$ es un esquema, esto implica $z'=j(x')$ (\sref[I]{I.8.2.2}), lo que termina la demostración.
\end{proof}

\begin{remark}[15.7.3]
\label{IV.15.7.3}
En \sref{IV.15.7.2}, se puede omitir la hipótesis de que $Y$ sea localmente noetheriano sustituyendo, en la condición $c)$, los anillos de valoración discreta por anillos de valoración arbitrarios; la demostración queda entonces sin cambios, teniendo en cuenta (\sref[II]{II.7.3.10}).
\end{remark}

\begin{corollary}[15.7.4]
\label{IV.15.7.4}
Sea $Y$ un preesquema noetheriano, sea $f:X\to Y$ un morfismo separado de tipo finito y sea $Z$ una parte de $X$ tal que los puntos maximales de su adherencia $\overline{Z}$ pertenezcan a $Z$ (lo que ocurre cuando $Z$ es finito o cuando $Z$ es construible ($\mathbf{0}_{\mathrm{III}},~9.2.2$)).
Sea $y$ un punto de $Y$.
Las condiciones siguientes son equivalentes:
\begin{enumerate}[label=\emph{\alph*)}]
  \item El conjunto $\overline{Z}$ es propio sobre $Y$ en el punto $y$.
  \item Si $Y_1=\Spec(\sh{O}_y)$, $X_1=X\times_Y Y_1$, y si $g_1:X_1\to X$ es la proyección canónica y $Z_1=g_1^{-1}(Z)$, la adherencia $\overline{Z_1}$ de $Z_1$ en $X_1$ es propia sobre $Y_1$.
  \item Para todo esquema $Y'=\Spec(A')$, donde $A'$ es un anillo de valoración discreta, y todo morfismo $g:Y'\to Y$ tal que la imagen $g(y')$ del punto cerrado $y'$ de $Y'$ sea $y$, se tiene la propiedad siguiente: poniendo $X'=X\times_Y Y'$, $f'=f_{(Y')}:X'\to Y'$, $g'=g_{(\mathrm{X})}:X'\to X$, $Z'=g'^{-1}(Z)$, y denotando por $\eta'$
\oldpage[IV-3]{244}
el punto genérico de $Y'$, para todo punto $x'\in Z'\cap f'^{-1}(\eta')$ racional sobre $\kappa(\eta')$, existe una $Y'$-sección de $X'$ que contiene a $x'$.
\end{enumerate}
\end{corollary}

\begin{proof}
El hecho de que $a)$ implique $b)$ se deduce de que $\overline{Z_1}\subset g_1^{-1}(\overline{Z})$ y de que $g_1^{-1}(\overline{Z})$ es propio sobre $Y_1$ (\sref{IV.15.7.1}).
El hecho de que $b)$ implique $c)$ se deduce de (\sref[II]{II.7.3.3}).
Queda, pues, por ver que $c)$ implica $a)$.
Está claro que basta probar que toda componente irreducible de $\overline{Z}$ es propia sobre $Y$ en el punto $y$, y la hipótesis permite, por tanto, limitarse al caso en que $Z$ se reduce a un solo punto $z$.
También se puede, teniendo en cuenta (\sref[I]{I.5.2.2}) y (\sref[II]{II.5.4.6}), suponer que $X$ e $Y$ son reducidos; entonces, por la definición de una parte propia sobre $Y$ (\sref[II]{II.5.4.10}), supongamos que $\overline{Z}=X=\overline{\{z\}}$ y, aplicando una vez más (\sref[I]{I.5.2.2}), que $f$ es dominante. Así, $X$ e $Y$ son íntegros y noetherianos, y queda por probar que $f$ es propio en el punto $y$ de $Y$. Para ello, mostremos que $f$ satisface la condición $c)$ de (\sref{IV.15.7.2}).
Sea $A'$ un anillo de valoración discreta que tiene a $R(X)=\kappa(z)$ como cuerpo de fracciones y domina a $\sh{O}_y$; con la notación de $c)$, se tiene $g(y')=y$, y $g(\eta')$ es el punto genérico $\eta=f(z)$ de $Y$.
Como por definición $\kappa(\eta')=\kappa(z)$, existe un $\kappa(\eta)$-morfismo $\Spec(\kappa(\eta'))\to\Spec(\kappa(z))$ que transforma $\eta'$ en $z$, y por tanto una $\kappa(\eta')$-sección de $f'^{-1}(\eta')$ (\sref[I]{I.3.3.14}); si $z'$ es la imagen de $\eta'$ por esta sección, $z'$ es, pues, racional sobre $\kappa(\eta')$ y $g'(z')=z$.
Se concluye de la hipótesis $c)$ que existe una $Y'$-sección $h'$ de $X'$ tal que $h'(\eta')=z'$; sea $h=g'\circ h':Y'\to X$ el $Y$-morfismo correspondiente, y sea $x=h(y')$. Está claro que $f(x)=y$ y que $A'=\sh{O}_{y'}$ domina a $\sh{O}_x$, lo que termina la demostración.
\end{proof}

\begin{corollary}[15.7.5]
\label{IV.15.7.5}
Sea $Y$ un preesquema localmente noetheriano, sea $f:X\to Y$ un morfismo separado de tipo finito y sea $y$ un punto de $Y$.
Las condiciones siguientes son equivalentes:
\begin{enumerate}[label=\emph{\alph*)}]
  \item $f$ es propio en el punto $y$.
  \item Si $Y_1=\Spec(\sh{O}_y)$ y $X_1=X\times_Y Y_1$, el morfismo $f_1=f_{(Y_1)}:X_1\to Y_1$ es propio.
  \item Para todo esquema $Y'=\Spec(A')$, donde $A'$ es un anillo de valoración discreta, y todo morfismo $g:Y'\to Y$ tal que la imagen $g(y')$ del punto cerrado $y'$ de $Y'$ sea $y$, se tiene la propiedad siguiente: poniendo $X'=X\times_Y Y'$, $f'=f_{(Y')}:X'\to Y'$, y denotando por $\eta'$ el punto genérico de $Y'$, para todo $x'\in f'^{-1}(\eta')$ racional sobre $\kappa(\eta')$, existe una $Y'$-sección de $X'$ que contiene a $x'$.
\end{enumerate}
\end{corollary}

\begin{corollary}[15.7.6]
\label{IV.15.7.6}
Sea $Y$ un preesquema localmente noetheriano y sea $f:X\to Y$ un morfismo de tipo finito.
Las condiciones siguientes son equivalentes:
\begin{enumerate}[label=\emph{\alph*)}]
  \item Existe un morfismo propio $g:X\to Z$ y una inmersión $h:Z\to Y$ tales que $f=h\circ g$.
  \item El subespacio $f(X)$ es localmente cerrado en $Y$; si $Z'$ es el subpreesquema imagen cerrada de $Y$ determinado por $X$ mediante $f$ (\sref[I]{I.9.5.3} y \sref[I]{I.9.5.1}), y $Z''$ es el preesquema inducido por $Z'$ sobre el subconjunto abierto $f(X)$ de $Z'$, de modo que $f$ se factoriza de manera única como $f=j\circ g$ (donde $j:Z''\to Y$ es la inyección canónica), entonces $g:X\to Z''$ es propio.
  \item Para todo $y\in f(X)$, $f$ es propio en el punto $y$.
  \item Para todo esquema $Y'=\Spec(A')$, donde $A'$ es un anillo de valoración discreta, y todo morfismo $g:Y'\to Y$ tal que $g(Y')\subset f(X)$, toda $Y'$-sección racional (\sref[I]{I.7.1.2}) de $X'=X\times_Y Y'$ se prolonga de manera única a una $Y'$-sección de $X'$.
\end{enumerate}
\end{corollary}

\begin{proof}
Está claro que $b)$ implica $a)$. Para probar que $a)$ implica $c)$, observemos primero que $a)$ implica que $f(X)$ es localmente cerrado en $Y$. Para todo $y\in f(X)$ existe, por tanto, un entorno abierto $U$
\oldpage[IV-3]{245}
de $y$ en $Y$ tal que $f(X)\cap U$ es cerrado en $U$; la restricción $h^{-1}(U)\to U$ de $h$ es entonces una inmersión cerrada, y la restricción
\[
  f^{-1}(U)=g^{-1}\bigl(h^{-1}(U)\bigr)\longrightarrow h^{-1}(U)
\]
de $g$ es propia. Por consiguiente, la restricción $f^{-1}(U)\to U$ de $f$ es propia (\sref[II]{II.5.4.2}).

Para probar que $c)$ implica $b)$, observemos primero que, para todo $y\in f(X)$, existe por $c)$ un entorno abierto $U$ de $y$ en $Y$ tal que $f(X)\cap U$ es cerrado en $U$. Así, $f(X)$ es localmente cerrado, y por tanto abierto en su adherencia. Como esta adherencia es el espacio subyacente de $Z'$, y como $f$ se factoriza como $f=j_0\circ g_0$, donde $j_0:Z'\to Y$ es la inyección canónica y $g_0:X\to Z'$, la igualdad $g_0(X)=Z''$ (que es abierto en $Z'$) da la factorización del enunciado de $b)$. El hecho de que $g$ sea propio se deduce entonces del carácter local de esta propiedad (sobre $Z''$) y de \sref[II]{II.5.4.3}.

Está claro que $c)$ implica $d)$ por \sref{IV.15.7.5}; la unicidad de la prolongación de una $Y'$-sección racional se deduce de la hipótesis de que $f$ es separado (\sref[I]{I.7.2.2}).
Finalmente, probemos que $d)$ implica $c)$. En primer lugar, se deduce de $d)$, teniendo en cuenta \sref[II]{II.7.2.3} y \sref[II]{II.7.2.4}, que $f$ es separado; se puede aplicar entonces el criterio \sref{IV.15.7.5},~$c$), que muestra que $f$ es propio en todo punto de $f(X)$.
\end{proof}

\begin{remarks}[15.7.7]
\label{IV.15.7.7}
(i)
En \sref{IV.15.7.4},~$c$), \sref{IV.15.7.5},~$c$), y \sref{IV.15.7.6},~$d$), se puede limitar al caso en que el anillo de valoración discreta $A'$ sea \emph{completo} y tenga un cuerpo residual algebraicamente cerrado.
En efecto, en la demostración de \sref{IV.15.7.4}, consideremos un anillo de valoración discreta completo $A''$ que domine a $A'$ y tenga un cuerpo residual algebraicamente cerrado ($\mathbf{0}_{\mathrm{III}},~10.3.1$). La condición $c)$ se satisface entonces para $Y''=\Spec(A'')$, $X''=X\times_Y Y''$ y $f''=f_{(Y'')}$; pero $X''=X'\times_{Y'}Y''$, y se deduce de la observación \sref[II]{II.7.3.9},~(i), que $f'$ es propio. Por consiguiente, $Y'$, $X'$ y $f'$ también satisfacen la condición $c)$ (\sref[II]{II.7.3.8}).

(ii)
Teniendo en cuenta (\sref[II]{II.7.3.8}), la condición $c)$ de \sref{IV.15.7.5} puede sustituirse por la condición de que $f'$ sea \emph{propio}.
\end{remarks}

\begin{env}[15.7.8]
\label{IV.15.7.8}
Se dice que un morfismo de preesquemas $f:X\to Y$ es \emph{submersivo} si es sobreyectivo y si la topología de $Y$ es igual al cociente de la topología de $X$ por la relación de equivalencia definida por $f$.
Se dice que $f$ es \emph{universalmente submersivo} si, para todo cambio de base $Y'\to Y$, el morfismo $f'=f_{(Y')}:X\times_Y Y'\to Y'$ es submersivo.
Todo morfismo sobreyectivo que sea \emph{universalmente abierto}, o \emph{universalmente cerrado}, o \emph{fielmente plano y cuasicompacto} es universalmente submersivo \sref{IV.2.3.12}.
Dado un morfismo $f:X\to Y$, se dice que una parte $Z$ de $X$ es \emph{submersiva sobre $f(Z)$} si la topología del subespacio $f(Z)$ de $Y$ es igual al cociente de la topología del subespacio $Z$ de $X$ por la relación de equivalencia definida por $f|Z$.
Se dice que $Z$ es \emph{universalmente submersiva sobre $f(Z)$} si, para todo cambio de base $g:Y'\to Y$, la imagen inversa $Z'=p^{-1}(Z)$ de $Z$ por la proyección $p:X\times_Y Y'\to X$ es submersiva sobre $f'(Z')=g^{-1}(f(Z))$.
Observemos que, si el morfismo $f$ admite una $Y$-sección $h:Y\to X$, todo conjunto $Z$ que contiene a $h(Y)$ es universalmente submersivo sobre $Y$.
\end{env}

\begin{proposition}[15.7.9]
\label{IV.15.7.9}
Sea $Y$ un preesquema localmente noetheriano, sea $y$ un punto de $Y$ y sea $f:X\to Y$ un morfismo separado de tipo finito.
Sea $Z$ una parte localmente construible de $X$ que posee las propiedades siguientes:
\begin{enumerate}[label=(\roman*)]
  \item Para toda generalización $y'$ de $y$, $Z_{y'}=Z\cap f^{-1}(y')$ es una componente conexa de $f^{-1}(y')$ y es geométricamente conexa sobre $\kappa(y')$ \sref{IV.4.5.2}.
\oldpage[IV-3]{246}
  \item $Z_y$ es propio sobre $\Spec(\kappa(y))$.
  \item $Z$ es universalmente submersivo sobre $f(Z)$ \sref{IV.15.7.8}.
\end{enumerate}
Entonces $Z$ es propio sobre $Y$ en el punto $y$ (dicho de otro modo, en el sentido de \sref{IV.15.7.1}, existe un entorno abierto $U$ de $y$ tal que $Z\cap f^{-1}(U)$ es cerrado en $f^{-1}(U)$ y propio sobre $U$).
\end{proposition}

\begin{proof}
Utilizando el método de (\sref{IV.8.1.2},~$a$)) y (\sref{IV.8.10.5},~(xii)), se reduce a probar que, cuando $Y=\Spec(A)$ es el espectro de un anillo noetheriano \emph{local}, las hipótesis (i), (ii) y (iii) implican que $Z$ es \emph{propio} sobre $Y$.

Observemos primero que, si $A'$ es un anillo local noetheriano, $\varphi:A\to A'$ un homomorfismo local, $Y'=\Spec(A')$ y $g:Y'\to Y$ el morfismo correspondiente a $\varphi$, la imagen inversa $Z'$ de $Z$ por la proyección canónica $p:X\times_Y Y'\to X$ posee, con respecto a $Y'$, las propiedades (i), (ii) y (iii) del enunciado, teniendo en cuenta \sref{IV.1.8.2}. Utilizando \sref{IV.2.7.1},~(vii), equivale, por tanto, a probar que $Z'$ es propio sobre $Y'$, siempre que $A'$ sea un $A$-módulo fielmente plano.
Utilizaremos esta observación tomando $A'=\hat{A}$ ($\mathbf{0}_{\mathrm{I}},~7.3.5$); dicho de otro modo, nos limitaremos en adelante al caso en que $A$ es \emph{completo}.
Como $Z_y$ es una componente conexa de $f^{-1}(y)$, propia sobre $\kappa(y)$, se deduce de \sref[III]{III.5.5.2} que $X$ es la \emph{suma} de dos subpreesquemas $X_0$, $X_1$ inducidos sobre partes abiertas y cerradas de $X$, tales que $X_0$ es propio sobre $Y$ y $X_0\cap f^{-1}(y)=Z_y$.
Pongamos $Z_0=X_0\cap Z$, $Z_1=X_1\cap Z$, de modo que estos conjuntos forman una partición de $Z$ en dos partes abiertas y cerradas.
Como los $Z_{y'}$ son conexos, se tiene necesariamente $Z_{y'}\subset Z_0$ o $Z_{y'}\subset Z_1$, y por tanto $f(Z_0)\cap f(Z_1)=\varnothing$.
Pero, como $Z_0$ y $Z_1$ están saturados para la relación de equivalencia definida por $f|Z$, se deduce de (iii) que $f(Z_0)$ y $f(Z_1)$ son ambos abiertos en $f(Z)$.
Ahora bien, $f(Z_0)$ contiene a $y$, y todo entorno de $y$ en $Y$ es igual a $Y$, de modo que $f(Z_0)=f(Z)$, y por tanto $f(Z_1)=\varnothing$; así, $Z_1=\varnothing$ y $Z\subset X_0$.

Se puede suponer, pues, en adelante que además $f$ es \emph{propio}, y queda por probar que $Z$ es cerrado en $X$.
Dicho de otro modo, bastará probar que una parte \emph{construible} $Z$ de un preesquema $X$ \emph{propio} sobre $Y$, que solo satisface las condiciones (i) y (iii), es \emph{cerrada} en $X$.
En virtud de ($\mathbf{0}_{\mathrm{III}},~9.2.5$), basta, por tanto, probar que para todo $x'\in Z$ y toda especialización $x$ de $x'$ en $X$, se tiene $x\in Z$.
Sea $A_1$ un anillo de valoración discreta completo y sea $u$ un morfismo de $Y_1=\Spec(A_1)$ en $X$ tal que, si $y_1$ e $y_1'$ son el punto cerrado y el punto genérico de $Y_1$, entonces $u(y_1)=x$ y $u(y_1')=x'$ (\sref[II]{II.7.1.7}); pongamos $g=f\circ u:Y_1\to Y$ y $X_1=X\times_Y Y_1$. Existe una $Y_1$-sección $u_1:Y_1\to X_1$ de $f_1=f_{(Y_1)}$ tal que $u=p\circ u_1$, donde $p:X_1\to X$ es la proyección canónica.
Además, $x_1=u_1(y_1)$, $x_1'=u_1(y_1')$, y entonces $p(x_1)=x$ y $p(x_1')=x'$, y $x_1$ es una especialización de $x_1'$.
Asimismo, $x_1'\in Z_1=p^{-1}(Z)$. Como hemos observado que las condiciones (i) y (iii) son estables bajo \emph{todo} cambio de base, al igual que la propiedad de ser construible, quedamos reducidos a la situación siguiente: $Y$ es el espectro de un anillo de valoración discreta completo, $X$ es propio sobre $Y$, $f(x)=y$ es el punto cerrado de $Y$, $y'=f(x')$ es su punto genérico, existe una $Y$-sección $h$ de $X$ tal que $h(y)=x$ y $h(y')=x'$, y finalmente $x'\in Z$; debemos probar que $x\in Z$.
Ahora bien, $Z_y$ es una componente conexa de $f^{-1}(y)$, \emph{propia} sobre $\Spec(\kappa(y))$, puesto que $X$ es propio sobre $Y$.
Aplicando una vez más \sref[III]{III.5.5.2}, se ve, como antes, que existe una parte abierta y cerrada $X'$ de $X$ tal que $X'\cap f^{-1}(y)=Z_y$; como $Z_y$ es conexo,
\oldpage[IV-3]{247}
se puede suponer que $X'$ es conexo (sustituyendo $X'$, si es necesario, por la componente conexa que contiene a $Z_y$).
Pero entonces el mismo razonamiento que al comienzo prueba que necesariamente $Z\subset X'$; como $h(Y)$ es conexo y contiene a $x'$, está necesariamente contenido en $X'$, luego $x=h(y)\in X'\cap f^{-1}(y)=Z_y$, lo que termina la demostración.
\end{proof}

\begin{corollary}[15.7.10]
\label{IV.15.7.10}
Sea $Y$ un preesquema localmente noetheriano y sea $f:X\to Y$ un morfismo separado de tipo finito, universalmente submersivo \sref{IV.15.7.8}, tal que todas las fibras $X_y=f^{-1}(y)$ sean geométricamente conexas.
Entonces, si $y\in Y$ es tal que $X_y$ es propio sobre $\kappa(y)$, $f$ es propio en el punto $y$.
\end{corollary}

\begin{proof}
Teniendo en cuenta \sref{IV.15.7.5}, se reduce al caso en que $Y=\Spec(\sh{O}_y)$, puesto que las hipótesis son estables bajo cambio de base.
Pero, en este caso, basta aplicar \sref{IV.15.7.9} a $Z=X$.
\end{proof}

\begin{corollary}[15.7.11]
\label{IV.15.7.11}
Sea $f:X\to Y$ un morfismo separado de presentación finita; supongamos que existe un morfismo sobreyectivo de presentación finita $g:Y'\to Y$, que es además propio o plano, y tal que, si se pone $X'=X\times_Y Y'$, existe una $Y'$-sección de $X'$ (o incluso un $Y$-morfismo $Y'\to X$ (\sref[I]{I.3.3.14})).
Supongamos finalmente que las fibras $X_y=f^{-1}(y)$ son geométricamente conexas.
Entonces el conjunto $U$ de los $y\in Y$ tales que $X_y$ es propio sobre $\kappa(y)$ es abierto en $Y$, y la restricción $f^{-1}(U)\to U$ de $f$ es propia.
\end{corollary}

\begin{proof}
Como la cuestión es local sobre $Y$, se puede suponer $Y=\Spec(A)$ afín.
Aplicando \sref{IV.8.9.1} a $f$ y a $g$, se ve que existe un subanillo noetheriano $A_0$ de $A$ y dos morfismos de tipo finito $f_0:X_0\to Y_0=\Spec(A_0)$, $g_0:Y_0'\to Y_0$ tales que $X=X_0\times_{Y_0}Y$, $Y'=Y_0'\times_{Y_0}Y$, $f=(f_0)_{(Y)}$ y $g=(g_0)_{(Y)}$; además, utilizando \sref{IV.8.10.5},~(v),~(vi) y (xii), y \sref{IV.11.2.7}, se puede suponer que $f_0$ es separado y que $g_0$ es sobreyectivo y propio (resp.\ plano) si $g$ es propio (resp.\ plano); utilizando \sref{IV.8.8.2}, se puede suponer además que existe una $Y_0'$-sección de $X_0$.

Por otra parte, utilizando \sref{IV.9.7.7}, se puede aplicar \sref{IV.9.3.3} a la propiedad de ser geométricamente conexo, y por tanto suponer que $A_0$ está elegido de modo que los $f_0^{-1}(y_0)$ sean geométricamente conexos para todo $y_0\in Y_0$.
Además, utilizando \sref{IV.2.7.1},~(vii), equivale a decir que $f^{-1}(y)$ es propio sobre $\kappa(y)$ o que $f_0^{-1}(y_0)$ es propio sobre $\kappa(y_0)$, donde $y_0$ es la proyección de $y$ en $Y_0$.
Estas observaciones muestran que quedamos reducidos a probar el corolario cuando $Y$ es \emph{noetheriano}.
Probemos ahora el

\begin{lemma}[15.7.11.1]
\label{IV.15.7.11.1}
Sea $f:X\to Y$ un morfismo y sea $g:Y'\to Y$ un morfismo sobreyectivo que es además propio o plano y cuasicompacto.
Pongamos $X'=X\times_Y Y'$, $f'=f_{(Y')}:X'\to Y'$.
Entonces, si $f'$ es submersivo (resp.\ universalmente submersivo), $f$ también lo es.
\end{lemma}

\begin{proof}
Se puede limitar al caso en que $f'$ es submersivo.
En efecto, supongámoslo probado en este caso y mostremos que, si $f'$ es universalmente submersivo, entonces $f$ también lo es.
Sea $Y_1\to Y$ un morfismo cualquiera y pongamos $Y_1'=Y'\times_Y Y_1$; si $X_1'=X'\times_{Y'} Y_1'$, se tiene $f_1'=(f')_{(Y_1')}:X_1'\to Y_1'$; $f_1'$ es submersivo por hipótesis; además, el morfismo $g_1:Y_1'\to Y_1$ es sobreyectivo y propio (resp.\ plano y cuasicompacto) si $g$ lo es.
Por tanto, si se pone $X_1=X\times_Y Y_1$, $f_1=f_{(Y_1)}$, se deduce de la hipótesis y del hecho de que $X_1'=X_1\times_{Y_1} Y_1'$ que $f_1$ es submersivo, y por tanto $f$ universalmente submersivo.
Supongamos, pues, únicamente que $f'$ es submersivo.
Está claro, ante todo, que $f$ es sobreyectivo.
Sea $E$ una parte de $Y$ tal que $f^{-1}(E)$ sea cerrado en $X$; entonces, si $p:X'\to X$ es la proyección canónica,
\oldpage[IV-3]{248}
$p^{-1}(f^{-1}(E))=f'^{-1}(g^{-1}(E))$ es cerrado en $X'$, luego, como $f'$ es submersivo, $g^{-1}(E)$ es cerrado en $Y'$; pero, como $g$ también es universalmente submersivo \sref{IV.15.7.8}, $E$ es cerrado en $Y$, lo que prueba el lema.
\end{proof}

Como por hipótesis existe una $Y'$-sección de $X'$, $f'$ es universalmente submersivo \sref{IV.15.7.8}, y por tanto $f$ también lo es por el lema \sref{IV.15.7.11.1}.
Basta entonces aplicar \sref{IV.15.7.10}, observando que el conjunto de los $y\in Y$ tales que $f$ es propio en $y$ es abierto en $Y$ por definición.
\end{proof}


\clearpage
\oldpage[IV-3]{249}
\phantomsection
\addcontentsline{toc}{section}{Bibliografía histórica (continuación)}
\section*{Bibliografía histórica (continuación)}

\begin{description}[style=multiline,leftmargin=3.5em,labelwidth=2.8em]
  \item[{[39]}] J.-P. Serre, \emph{Corps locaux}, París (Hermann), 1962.
  \item[{[40]}] J.-P. Serre, \emph{Groupes proalgébriques},
  \emph{Publ. Inst. Hautes Études Sci.}, n.º~7 (1960).
\end{description}

\clearpage
\oldpage[IV-3]{250}
\phantomsection
\addcontentsline{toc}{section}{Índice de notaciones}
\section*{Índice de notaciones}

\begingroup
\small
\setlength{\parindent}{0pt}
\setlength{\parskip}{2pt}

$\mathfrak{L}(X)$, $\mathfrak{O}c(X)$, $\mathfrak{G}c(X)$,
$\mathfrak{F}gc(X)$ ($X$ preesquema): 8.3.9.

$Q(\mathcal F)$ ($\mathcal F$ $\mathcal O_X$-Módulo): 8.5.10.

$\operatorname{Spr}(Y)$, $\operatorname{Spro}(Y)$,
$\operatorname{Sprf}(Y)$ ($Y$ preesquema): 8.6.1.

$\operatorname{Pro}(C)$ ($C$ categoría): 8.13.3.

$X_S$, $\mathcal F_S$, $u_S$, $s_S$, $Z_S$, $h_S$
($X$ $S$-preesquema, $s\in S$, $\mathcal F$ $\mathcal O_X$-Módulo,
$u$ homomorfismo de $\mathcal O_X$-Módulos, $s$ sección de un
$\mathcal O_X$-Módulo, $Z$ parte de $X$, $h:X\to Y$ $S$-morfismo): 9.4.1.

$Qc(X)$, $Sqc(X)$, $\mathfrak O(X)$, $\mathfrak G(X)$,
$\mathfrak U(X)$, $\mathfrak C(X)$ ($X$ espacio topológico): 10.1.1.

$\operatorname{prof}^{*}_{x}(\mathcal F)$,
$\operatorname{prof}_{*}(\mathcal F)$,
$\operatorname{prof}^{*}(\mathcal F)$
($\mathcal F$ $\mathcal O_X$-Módulo): 10.8.1.

$S(X)$ ($X$ preesquema de Jacobson): 10.9.1.

$S(f)$ ($f$ morfismo de preesquemas de Jacobson): 10.9.2.

$\operatorname{Spm}(A)$, $D_m(a)$ ($A$ anillo de Jacobson, $a\in A$): 10.9.3.

$\operatorname{Tor.dim}_{B}(M)$ ($M$ $B$-módulo): 12.3.2.
\endgroup

\clearpage
\oldpage[IV-3]{251}
\phantomsection
\addcontentsline{toc}{section}{Índice terminológico}
\section*{Índice terminológico}

\begingroup
\small
\setlength{\parindent}{0pt}
\setlength{\parskip}{2pt}
\newcommand{\EGAIVthreeTerm}[2]{\textbf{#1}: #2.\par}

\EGAIVthreeTerm{Criterio de planitud por las fibras}{11.3.10}
\EGAIVthreeTerm{Espacio algebraico sobre $k$, espacio $k$-algebraico,
espacio prealgebraico sobre $k$, espacio $k$-prealgebraico}{10.10.2}
\EGAIVthreeTerm{Espacio de Jacobson}{10.3.1}
\EGAIVthreeTerm{Espacio $k$-prealgebraico reducido}{10.10.4}
\EGAIVthreeTerm{Familia de morfismos esquemáticamente dominante}{11.10.2}
\EGAIVthreeTerm{Familia de morfismos universalmente esquemáticamente dominante}{11.10.7}
\EGAIVthreeTerm{Familia de subpreesquemas esquemáticamente densa}{11.10.2}
\EGAIVthreeTerm{Familia separante de homomorfismos de haces de módulos}{11.9.1}
\EGAIVthreeTerm{Familia separante de homomorfismos de módulos}{11.9.4}
\EGAIVthreeTerm{Familia universalmente separante de homomorfismos de haces de módulos}{11.9.14}
\EGAIVthreeTerm{Geométricamente isomorfos (preesquemas sobre $k$)}{9.1.4}
\EGAIVthreeTerm{Homomorfismo de Módulos universalmente inyectivo}{11.9.14 y 11.9.18}
\EGAIVthreeTerm{Lema de Chow para los morfismos de presentación finita}{8.10.5.1}
\EGAIVthreeTerm{Teorema principal de Zariski para los morfismos de presentación finita}{8.12.6}
\EGAIVthreeTerm{Morfismo equidimensional en un punto, morfismo equidimensional}{13.2.2 y 13.3.2}
\EGAIVthreeTerm{Morfismo propio en un punto}{15.7.1}
\EGAIVthreeTerm{Morfismo pseudofinito}{8.12.3}
\EGAIVthreeTerm{Morfismo submersivo, morfismo universalmente submersivo}{15.7.8}
\EGAIVthreeTerm{Morfismo universalmente abierto en un punto}{14.3.3}
\EGAIVthreeTerm{Abierto ultraafín}{10.9.5}
\EGAIVthreeTerm{Parte finita saturada de un preesquema algebraico sobre $k$}{9.8.9}
\EGAIVthreeTerm{Parte propia en un punto (de un $S$-preesquema)}{15.7.1}
\EGAIVthreeTerm{Parte cuasiconstruible, parte localmente cuasiconstruible de un espacio topológico}{10.1.1}
\EGAIVthreeTerm{Parte muy densa de un espacio topológico}{10.1.3}
\EGAIVthreeTerm{Polinomio geométricamente irreducible}{9.7.4}
\EGAIVthreeTerm{Preesquema de Jacobson}{10.4.1}
\EGAIVthreeTerm{Preesquema equidimensional sobre otro en un punto, equidimensional sobre otro}{13.2.2 y 13.3.2}
\EGAIVthreeTerm{Preesquema esencialmente afín sobre otro}{8.13.4}
\EGAIVthreeTerm{Principio de extensión finita}{9.1.1}
\EGAIVthreeTerm{Profundidad rectificada}{10.8.1}
\EGAIVthreeTerm{Proobjeto, morfismo de proobjetos}{8.13.3}
\EGAIVthreeTerm{Proobjeto esencialmente afín}{8.13.4, 8.14.1}
\EGAIVthreeTerm{Propiedad construible, propiedad ind-construible}{9.2.1 y 9.2.2, (i) y (ii)}
\EGAIVthreeTerm{Cuasihomeomorfismo de espacios topológicos}{10.2.2}
\EGAIVthreeTerm{Cuasiisomorfismo de espacios anillados}{10.2.8}
\EGAIVthreeTerm{Saturado de una parte finita de un preesquema algebraico sobre $k$}{9.8.9}
\EGAIVthreeTerm{Subconjunto de un $S$-preesquema submersivo, universalmente submersivo sobre su imagen}{15.7.8}
\EGAIVthreeTerm{Esqueleto primario de un Módulo, esqueleto virtual}{9.8.9}
\EGAIVthreeTerm{Espectro maximal de un anillo de Jacobson}{10.9.3}
\EGAIVthreeTerm{Tipo primario de un Módulo, tipo primario geométrico de un Módulo}{9.8.9}
\EGAIVthreeTerm{Ultrapreesquema, morfismo de ultrapreesquemas}{10.9.5}
\EGAIVthreeTerm{$R$-ultrapreesquema}{10.10.2}
\endgroup

\clearpage
\oldpage[IV-3]{253}
\phantomsection
\addcontentsline{toc}{section}{Tabla de materias original}
\section*{Tabla de materias original}

\begingroup
\small
\newcommand{\EGAIVthreeContents}[3]{\textbf{#1}\quad #2\dotfill #3\par}

\EGAIVthreeContents{Capítulo IV}{Estudio local de los esquemas y de los morfismos de esquemas (continuación)}{5}
\EGAIVthreeContents{§8}{Límites proyectivos de preesquemas}{5}
\EGAIVthreeContents{8.1}{Introducción}{5}
\EGAIVthreeContents{8.2}{Límites proyectivos de preesquemas}{7}
\EGAIVthreeContents{8.3}{Partes construibles en un límite proyectivo de preesquemas}{12}
\EGAIVthreeContents{8.4}{Criterios de irreducibilidad y conexión para los límites proyectivos de preesquemas}{17}
\EGAIVthreeContents{8.5}{Módulos de presentación finita sobre un límite proyectivo de preesquemas}{19}
\EGAIVthreeContents{8.6}{Subpreesquemas de presentación finita de un límite proyectivo de preesquemas}{25}
\EGAIVthreeContents{8.7}{Criterios para que un límite proyectivo de preesquemas sea reducido (resp. íntegro)}{27}
\EGAIVthreeContents{8.8}{Preesquemas de presentación finita sobre un límite proyectivo de preesquemas}{28}
\EGAIVthreeContents{8.9}{Primeras aplicaciones a la eliminación de hipótesis noetherianas}{34}
\EGAIVthreeContents{8.10}{Propiedades de permanencia de los morfismos por paso al límite proyectivo}{36}
\EGAIVthreeContents{8.11}{Aplicación a los morfismos cuasifinitos}{41}
\EGAIVthreeContents{8.12}{Nueva demostración y generalización del teorema principal de Zariski}{43}
\EGAIVthreeContents{8.13}{Traducción en términos de proobjetos}{49}
\EGAIVthreeContents{8.14}{Caracterización de un preesquema localmente de presentación finita sobre otro mediante el funtor que representa}{52}

\EGAIVthreeContents{§9}{Propiedades construibles}{54}
\EGAIVthreeContents{9.1}{El principio de extensión finita}{54}
\EGAIVthreeContents{9.2}{Propiedades construibles e ind-construibles}{56}
\EGAIVthreeContents{9.3}{Propiedades construibles de morfismos de preesquemas algebraicos}{60}

\oldpage[IV-3]{254}
\EGAIVthreeContents{9.4}{Construibilidad de ciertas propiedades de los Módulos}{62}
\EGAIVthreeContents{9.5}{Construibilidad de propiedades topológicas}{67}
\EGAIVthreeContents{9.6}{Construibilidad de ciertas propiedades de los morfismos}{71}
\EGAIVthreeContents{9.7}{Construibilidad de la separabilidad, la irreducibilidad geométrica y la conexión geométrica}{76}
\EGAIVthreeContents{9.8}{Descomposición primaria en un entorno de una fibra genérica}{83}
\EGAIVthreeContents{9.9}{Construibilidad de las propiedades locales de las fibras}{88}

\EGAIVthreeContents{§10}{Preesquemas de Jacobson}{95}
\EGAIVthreeContents{10.1}{Partes muy densas de un espacio topológico}{95}
\EGAIVthreeContents{10.2}{Cuasihomeomorfismos}{97}
\EGAIVthreeContents{10.3}{Espacios de Jacobson}{101}
\EGAIVthreeContents{10.4}{Preesquemas de Jacobson y anillos de Jacobson}{101}
\EGAIVthreeContents{10.5}{Preesquemas de Jacobson noetherianos}{104}
\EGAIVthreeContents{10.6}{Dimensión en los preesquemas de Jacobson}{107}
\EGAIVthreeContents{10.7}{Ejemplos y contraejemplos}{109}
\EGAIVthreeContents{10.8}{Profundidad rectificada}{110}
\EGAIVthreeContents{10.9}{Espectros maximales y ultrapreesquemas}{112}
\EGAIVthreeContents{10.10}{Espacios algebraicos de Serre}{114}

\EGAIVthreeContents{§11}{Propiedades topológicas de los morfismos planos de presentación finita. Criterios de planitud}{116}
\EGAIVthreeContents{11.1}{Lugares de planitud (caso noetheriano)}{117}
\EGAIVthreeContents{11.2}{Planitud de un límite proyectivo de preesquemas}{119}
\EGAIVthreeContents{11.3}{Aplicación a la eliminación de hipótesis noetherianas}{132}
\EGAIVthreeContents{11.4}{Descenso de la planitud por morfismos cualesquiera: caso de un preesquema base artiniano}{143}
\EGAIVthreeContents{11.5}{Descenso de la planitud por morfismos cualesquiera: caso general}{150}
\EGAIVthreeContents{11.6}{Descenso de la planitud por morfismos cualesquiera: caso de un preesquema base unirrama}{154}
\EGAIVthreeContents{11.7}{Contraejemplos}{157}
\EGAIVthreeContents{11.8}{Un criterio valuativo de planitud}{159}
\EGAIVthreeContents{11.9}{Familias separantes y universalmente separantes de homomorfismos de haces de módulos}{160}
\EGAIVthreeContents{11.10}{Familias esquemáticamente dominantes de morfismos y familias esquemáticamente densas de subpreesquemas}{170}

\EGAIVthreeContents{§12}{Estudio de las fibras de los morfismos planos de presentación finita}{173}
\EGAIVthreeContents{12.0}{Introducción}{173}
\EGAIVthreeContents{12.1}{Propiedades locales de las fibras de un morfismo plano localmente de presentación finita}{174}

\oldpage[IV-3]{255}
\EGAIVthreeContents{12.2}{Propiedades locales y globales de las fibras de un morfismo propio, plano y de presentación finita}{179}
\EGAIVthreeContents{12.3}{Propiedades cohomológicas locales de las fibras de un morfismo plano y localmente de presentación finita}{183}

\EGAIVthreeContents{§13}{Morfismos equidimensionales}{187}
\EGAIVthreeContents{13.1}{Teorema de semicontinuidad de Chevalley}{188}
\EGAIVthreeContents{13.2}{Morfismos equidimensionales: caso de morfismos dominantes de preesquemas irreducibles}{190}
\EGAIVthreeContents{13.3}{Morfismos equidimensionales: caso general}{194}

\EGAIVthreeContents{§14}{Morfismos universalmente abiertos}{199}
\EGAIVthreeContents{14.1}{Morfismos abiertos}{200}
\EGAIVthreeContents{14.2}{Morfismos abiertos y fórmula de las dimensiones}{202}
\EGAIVthreeContents{14.3}{Morfismos universalmente abiertos}{204}
\EGAIVthreeContents{14.4}{Criterio de Chevalley para los morfismos universalmente abiertos}{209}
\EGAIVthreeContents{14.5}{Morfismos universalmente abiertos y cuasisecciones}{216}

\EGAIVthreeContents{§15}{Estudio de las fibras de un morfismo universalmente abierto}{223}
\EGAIVthreeContents{15.1}{Multiplicidades de las fibras de un morfismo universalmente abierto}{223}
\EGAIVthreeContents{15.2}{Planitud de los morfismos universalmente abiertos con fibras geométricamente reducidas}{226}
\EGAIVthreeContents{15.3}{Aplicación: criterios de reducción e irreducibilidad}{228}
\EGAIVthreeContents{15.4}{Complementos sobre los morfismos de Cohen--Macaulay}{229}
\EGAIVthreeContents{15.5}{Rango separable de las fibras de un morfismo cuasifinito y universalmente abierto. Aplicación a las componentes conexas geométricas de las fibras de un morfismo propio}{231}
\EGAIVthreeContents{15.6}{Componentes conexas de las fibras a lo largo de una sección}{236}
\EGAIVthreeContents{15.7}{Apéndice: criterios valuativos de propiedad local}{242}
\EGAIVthreeContents{}{Bibliografía (continuación)}{249}
\EGAIVthreeContents{}{Índice de notaciones}{250}
\EGAIVthreeContents{}{Índice terminológico}{251}

\vspace{2em}
\begin{flushright}
Manuscrito recibido el 15 de noviembre de 1964.
\end{flushright}
\endgroup

\clearpage
\phantomsection
\addcontentsline{toc}{section}{Bibliografía consolidada de la edición}
\section*{Bibliografía consolidada de la edición}
\nocite{*}
\providecommand{\bysame}{\leavevmode\hbox to3em{\hrulefill}\thinspace}
\providecommand{\MR}{\relax\ifhmode\unskip\space\fi MR }
% \MRhref is called by the amsart/book/proc definition of \MR.
\providecommand{\MRhref}[2]{%
  \href{http://www.ams.org/mathscinet-getitem?mr=#1}{#2}
}
\providecommand{\href}[2]{#2}
\begin{thebibliography}{Sam53b}

\bibitem[AG64]{IV-43}
M.~Artin and A.~Grothendieck, \emph{Cohomologie \'etale des sch\'emas},
  1963--1964.

\bibitem[Cara]{IV-37}
H.~Cartan, \emph{S\'eminaire de l'\'ecole normale sup\'erieure, 13th year
  (1960--1961)}.

\bibitem[Carb]{III-29}
\bysame, \emph{S\'eminaire de l'\'ecole normale sup\'erieure, 13th year
  (1960--61), expos\'e no 11}.

\bibitem[CC]{I-1}
H.~Cartan and C.~Chevalley, \emph{S\'eminaire de l'\'ecole normale
  sup\'erieure, 8th year (1955--56), g\'eom\'etrie alg\'ebrique}.

\bibitem[CE56]{I-2}
H.~Cartan and S.~Eilenberg, \emph{Homological {A}lgebra}, Princeton Math.
  Series (Princeton University Press), 1956.

\bibitem[Che51]{II-23}
C.~Chevalley, \emph{Introduction to the theory of algebraic functions of one
  variable}, Math. Surveys, vol.~6, Amer. Math. Soc., New York, 1951.

\bibitem[CI58]{I-3}
W.~L. Chow and J.~Igusa, \emph{Cohomology theory of varieties over rings},
  Proc. Nat. Acad. Sci. U.S.A. \textbf{XLIV} (1958), 1244--1248.

\bibitem[eAG64]{IV-42}
M.~Demazure et~A.~Grothendieck, \emph{Sch\'emas en groupes}, 1963--1964.

\bibitem[Gab61]{III-27}
P.~Gabriel, \emph{Des cat\'egories ab\'eliennes}, Ph.D. thesis, Paris, 1961.

\bibitem[God58]{I-4}
R.~Godement, \emph{Th\'eorie des faisceaux}, Actual. Scient. et Ind., no. 1252,
  Paris (Hermann), 1958.

\bibitem[Gra60]{I-5}
H.~Grauert, \emph{Ein {T}heorem der analytischen {G}arbentheorie und die
  {M}odulr\"aume komplexer {S}trukturen}, Publ. Math. Inst. Hautes \'Etudes
  Scient., no. 5, 1960.

\bibitem[Gro]{I-8}
A.~Grothendieck, \emph{G\'eom\'etrie formelle et g\'eom\'etrie alg\'ebrique},
  S\'eminaire Bourbaki, 11th year (1958--59), expos\'e 182.

\bibitem[Gro57]{I-6}
\bysame, \emph{Sur quelques points d'alg\`ebre homologique}, T\^ohoku Math.
  Journ. \textbf{IX} (1957), 119--221.

\bibitem[Gro58]{I-7}
\bysame, \emph{Cohomology theory of abstract algebraic varieties}, Proc.
  Intern. Congress of Math., Edinburgh (1958), 103--118.

\bibitem[Gro62]{IV-41}
\bysame, \emph{Cohomologie locale des faisceaux coh\'erents, et th\'eor\`emes
  de lefschetz locaux et globaux}, 1962.

\bibitem[Gro71]{SGA1}
A.~Grothendieck (ed.), \emph{Rev\^etements \'etales et groupe fondamental
  ({SGA} 1)}, Lecture Notes in Mathematics, vol. 224, Springer-Verlag, 1971.

\bibitem[Hir]{IV-34}
H.~Hironaka, \emph{Projective morphisms}, Announced as forthcoming in the
  printed EGA bibliography.

\bibitem[Jaf60]{II-24}
P.~Jaffard, \emph{Les syst{\`e}mes d'id{\'e}aux}, Paris (Dunod), 1960.

\bibitem[K{\"a}h58]{I-22}
E.~K{\"a}hler, \emph{Geometria {A}rithmetica}, Ann. di Mat. (4) \textbf{XLV}
  (1958), 1--368.

\bibitem[Nag55]{II-25}
M.~Nagata, \emph{On the derived normal rings of noetherian integral domains},
  Mem. Coll. Sci. Kyoto \textbf{A, XXIX} (1955), 293--303.

\bibitem[Nag57]{IV-31}
\bysame, \emph{A {J}acobian criterion of simple points}, Illinois J. Math.
  \textbf{1} (1957), 427--432.

\bibitem[Nag58a]{I-9}
\bysame, \emph{A general theory of algebraic geometry over {D}edekind domains},
  Amer. Math. Journ. \textbf{I: LXXVIII, II: LXXX} (1956, 1958), 78--116,
  382--420.

\bibitem[Nag58b]{II-26}
\bysame, \emph{Existence theorems for non projective complete algebraic
  varieties}, Ill. J. Math \textbf{II} (1958), 490--498.

\bibitem[Nag62]{IV-30}
\bysame, \emph{Local rings}, Interscience Tracts, vol.~13, Interscience, New
  York, 1962.

\bibitem[N{\'e}r64]{IV-32}
A.~N{\'e}ron, \emph{Mod\`eles minimaux des vari\'et\'es ab\'eliennes sur les
  corps locaux et globaux}, Publications Math\'ematiques de l'IH\'ES
  \textbf{21} (1964).

\bibitem[Nor53]{I-10}
D.~G. Northcott, \emph{Ideal theory}, Cambridge Univ. Press, 1953.

\bibitem[Roo61]{III-28}
J.~E. Roos, \emph{Sur les foncteurs d\'eriv\'es de $\varprojlim$.
  applications}, C. R. Acad. Sci. Paris \textbf{CCLII} (1961), 3702--3704.

\bibitem[Sam53a]{I-12}
P.~Samuel, \emph{Alg\`ebre locale}, M\'em. Sci. Math., no. 123, Paris, 1953.

\bibitem[Sam53b]{I-11}
\bysame, \emph{Commutative algebra (notes by {D}. {H}erzig)}, Cornell Univ.,
  1953.

\bibitem[Ser56]{I-16}
J.-P. Serre, \emph{G\'eom\'etrie alg\'ebrique et g\'eom\'etrie analytique},
  Ann. Inst. Fourier \textbf{VI} (1955--56), 1--42.

\bibitem[Ser55a]{I-14}
\bysame, \emph{Faisceaux alg\'ebriques coh\'erents}, Ann. of Math. \textbf{LXI}
  (1955), 197--278.

\bibitem[Ser55b]{I-17}
\bysame, \emph{Sur la dimension homologique des anneaux et des modules
  noeth\'eriens}, Proc. Intern. Symp. on Alg. Number theory, Tokyo--Nikko
  (1955), 176--189.

\bibitem[Ser57]{I-15}
\bysame, \emph{Sur la cohomologie des vari\'et\'es alg\'ebriques}, Journ. de
  Math. (9) \textbf{XXXVI} (1957), 1--16.

\bibitem[Ses62]{IV-44}
C.~S. Seshadri, \emph{Quotient space by an abelian variety}, Mathematische
  Annalen \textbf{152} (1962), 185--194.

\bibitem[SZ60]{I-13}
P.~Samuel and O.~Zariski, \emph{Commutative algebra, 2 vol.}, New York (Van
  Nostrand), 1958--60.

\bibitem[Wei46]{I-18}
A.~Weil, \emph{Foundations of algebraic geometry}, Amer. Math. Soc. Coll.
  Publ., no. 29, 1946.

\bibitem[Wei49]{I-19}
\bysame, \emph{Numbers of solutions of equations in finite fields}, Bull. Amer.
  Math. Soc. \textbf{LV} (1949), 497--508.

\bibitem[Zar47]{I-21}
O.~Zariski, \emph{A new proof of {H}ilbert's {N}ullstellensatz}, Bull. Amer.
  Math. Soc. \textbf{LIII} (1947), 362--368.

\bibitem[Zar51]{I-20}
\bysame, \emph{Theory and applications of holomorphic functions on algebraic
  varieties over arbitrary ground fields}, Mem. Amer. Math. Soc., no. 5, 1951.

\end{thebibliography}

\bibliographystyle{amsalpha}

\end{document}
